% Complete English EGA IV1 editable text.
%
% Standalone release driver. External destinations owned by EGA 0 sections
% 1--13 and EGA IV Parts Two--Four remain deliberately non-linking here.
\PassOptionsToPackage{hyperfootnotes=false,draft}{hyperref}
% BEGIN NATIVE INPUT preamble
\documentclass[oneside]{amsart}

\usepackage[all]{xy}
\usepackage{tikz-cd}
\usepackage[T1]{fontenc}
\usepackage{xstring}
\usepackage{xparse}
\usepackage{xcolor}
\definecolor{brightmaroon}{rgb}{0.76, 0.13, 0.28}
\usepackage[linktocpage=true,colorlinks=true,hyperindex,hypertexnames=false,citecolor=blue,linkcolor=brightmaroon]{hyperref}
\usepackage[left=1.25in,right=1.25in,top=0.75in,bottom=0.75in]{geometry}
%\usepackage[charter,greekfamily=didot]{mathdesign}
%\usepackage{Baskervaldx}
\usepackage{amssymb}
\usepackage{mathrsfs}
\usepackage{mathpazo}
\linespread{1.05}

\usepackage[nobottomtitles]{titlesec}
\usepackage{marginnote}
\usepackage[shortlabels]{enumitem}
\usepackage{longtable}
\usepackage{multicol}
\usepackage{aurical}
\usepackage{microtype}

\newtheoremstyle{ega-env-style}%
  {}{}{\rmfamily}{}{\bfseries}{.}{ }{\thmnote{(#3)}}%

\newtheoremstyle{ega-thm-env-style}%
  {}{}{\itshape}{}{\bfseries}{. --- }{ }{\thmname{#1}\thmnote{ (#3)}}%

\newtheoremstyle{ega-defn-env-style}%
  {}{}{\rmfamily}{}{\bfseries}{. --- }{ }{\thmname{#1}\thmnote{ (#3)}}%

\theoremstyle{ega-env-style}
\newtheorem*{env}{---}

\theoremstyle{ega-thm-env-style}
\newtheorem*{theorem}{Theorem}
\newtheorem*{proposition}{Proposition}
\newtheorem*{lemma}{Lemma}
\newtheorem*{corollary}{Corollary}

\theoremstyle{ega-defn-env-style}
\newtheorem*{definition}{Definition}
\newtheorem*{example}{Example}
\newtheorem*{examples}{Examples}
\newtheorem*{remark}{Remark}
\newtheorem*{remarks}{Remarks}
\newtheorem*{scholium}{Scholium}
\newtheorem*{notation}{Notation}

% indent subsections, see https://tex.stackexchange.com/questions/177290/.
% also make section titles bigger.
% also add § to \thesection, https://tex.stackexchange.com/questions/119667/ and https://tex.stackexchange.com/questions/308737/.
\makeatletter
\def\l@subsection{\@tocline{2}{0pt}{2.5pc}{2.2pc}{}}
\def\section{\@startsection{section}{1}%
  \z@{.7\linespacing\@plus\linespacing}{.5\linespacing}%
  {\normalfont\bfseries\Large\scshape\centering}}
\renewcommand{\@seccntformat}[1]{%
  \ifnum\pdfstrcmp{#1}{section}=0\textsection\fi%
  \csname the#1\endcsname.~}
\makeatother

%\allowdisplaybreaks[1]
%\binoppenalty=9999
%\relpenalty=9999

% for Chapter 0, Chapter I, etc.
% credit for ZeroRoman https://tex.stackexchange.com/questions/211414/
% added into scripts/make_book.py
%\newcommand{\ZeroRoman}[1]{\ifcase\value{#1}\relax 0\else\Roman{#1}\fi}
%\renewcommand{\thechapter}{\ZeroRoman{chapter}}

\def\mathcal{\mathscr}
% BEGIN NATIVE INPUT preamble-base
\def\sh{\mathcal}                   % sheaf font
\def\bb{\mathbf}                    % bold font
\def\cat{\mathtt}                   % category font
\def\leq{\leqslant}                 % <=
\def\geq{\geqslant}                 % >=
\def\setmin{-}                      % set minus
\def\rad{\mathfrak{r}}              % radical
\def\nilrad{\mathfrak{N}}           % nilradical
\def\emp{\varnothing}               % empty set
\def\vphi{\phi}                     % for switching \phi and \varphi, change if needed
\def\HH{\mathrm{H}}                 % cohomology H
\def\CHH{\check{\HH}}               % Čech cohomology H
\def\RR{\mathrm{R}}                 % right derived R
\def\LL{\mathrm{L}}                 % left derived L
\def\dual#1{{#1}^\vee}              % dual
\def\kres{k}                        % residue field k
\def\C{\cat{C}}                     % category C
\def\op{^\cat{op}}                  % opposite category
\def\Set{\cat{Set}}                 % category of sets
\def\CHom{\cat{Hom}}                % functor category
\def\supertilde{{\,\widetilde{\,}\,}}   % use \supertilde instead of ^\sim
\def\GL{\bb{GL}}
\def\bbGamma{\boldsymbol\Gamma}
\def\red{\mathrm{red}}
\def\rg{\operatorname{rg}}
\def\gr{\operatorname{gr}}
\def\Gr{\operatorname{Gr}}
\def\Hom{\operatorname{Hom}}
\def\Proj{\operatorname{Proj}}
\def\Tor{\operatorname{Tor}}
\def\Ext{\operatorname{Ext}}
\def\Supp{\operatorname{Supp}}
\def\Ker{\operatorname{Ker}}
\def\Im{\operatorname{Im}}
\def\Coker{\operatorname{Coker}}
\def\Spec{\operatorname{Spec}}
\def\Spm{\operatorname{Spm}}
\def\Spf{\operatorname{Spf}}
\def\grad{\operatorname{grad}}
\def\dimc{\operatorname{dimc}}
\def\codim{\operatorname{codim}}
\def\id{\operatorname{id}}
\def\Der{\operatorname{Der}}
\def\Diff{\operatorname{Diff}}
\def\Hyp{\operatorname{Hyp}}
\def\Tr{\operatorname{Tr}}

\renewcommand{\to}{\mathchoice{\longrightarrow}{\rightarrow}{\rightarrow}{\rightarrow}}
\newcommand{\from}{\mathchoice{\longleftarrow}{\leftarrow}{\leftarrow}{\leftarrow}}
\let\mapstoo\mapsto
\renewcommand{\mapsto}{\mathchoice{\longmapsto}{\mapstoo}{\mapstoo}{\mapstoo}}
\def\isoto{\simeq} % isomorphism
\def\simto{\xrightarrow{\sim}}

\renewcommand{\hat}[1]{\widehat{#1}}
\renewcommand{\tilde}[1]{\widetilde{#1}}
% END NATIVE INPUT preamble-base

\def\shHom{\sh{H}\textup{\kern-2.2pt{\Fontauri\slshape om}}}   % sheaf Hom
\def\shProj{\sh{P}\textup{\kern-2.2pt{\Fontauri\slshape roj}}} % sheaf Proj
\def\shExt{\sh{E}\textup{\kern-2.2pt{\Fontauri\slshape xt}}}   % sheaf Ext
\def\shTor{\sh{T}\textup{\kern-2.2pt{\Fontauri\slshape or}}}   % sheaf Tor
\def\shGr{\sh{G}\textup{\kern-2.2pt{\Fontauri\slshape r}}}     % sheaf Gr
\def\shDer{\sh{D}\,\textup{\kern-2.2pt{\Fontauri\slshape er}}} % sheaf Der
\def\shDiff{\sh{D}\,\textup{\kern-2.2pt{\Fontauri\slshape if{}f}}\,} % sheaf Diff
\def\shHomcont{\sh{H}\textup{\kern-2.2pt{\Fontauri\slshape om.\,cont}}}   % sheaf Hom.cont
\def\shAut{\sh{A}\textup{\kern-2.2pt{\Fontauri\slshape ut}}}   % sheaf Aut

% if unsure of a translation
%\newcommand{\unsure}[2][]{\hl{#2}\marginpar{#1}}
%\newcommand{\completelyunsure}{\unsure{[\ldots]}}
\def\unsure#1{#1 {\color{red}(?)}}
\def\completelyunsure{{\color{red}(???)}}

% use to mark where original page starts
\newcommand{\oldpage}[2][]{{\marginnote{\normalfont{\textbf{#1}~|~#2}}}\ignorespaces}
\def\sectionbreak{\begin{center}***\end{center}}

% Give every source label its own zero-size PDF destination.  The inherited
% unnumbered theorem environments otherwise leave many labels pointing to the
% preceding subsection or display anchor rather than to their exact statement.
\let\EGAOriginalLabel\label
\renewcommand{\label}[1]{\phantomsection\EGAOriginalLabel{#1}}

% for referencing environments.
% use as \sref{chapter-number.x.y.z}, with optional args
% for volume and indices, e.g. \sref[volume]{chapter-number.x.y.z}[i].
\makeatletter
\newcommand{\EGAUnavailableHyperref}[2]{%
  \begingroup
  \let\@latex@warning\@gobble
  \let\@latex@warning@no@line\@gobble
  \let\G@refundefinedtrue\relax
  \hyperref[#1]{#2}%
  \endgroup
}
\NewDocumentCommand{\sref}{o m o}{%
  \@ifundefined{r@#2}%
    {\IfNoValueTF{#1}%
      {\IfNoValueTF{#3}%
        {\EGAUnavailableHyperref{#2}{\normalfont{(\StrBehind{#2}{.})}}}%
        {\EGAUnavailableHyperref{#2}{\normalfont{(\StrBehind{#2}{.},~{#3})}}}}%
      {\IfNoValueTF{#3}%
        {\EGAUnavailableHyperref{#2}{\normalfont{(\textbf{#1},~\StrBehind{#2}{.})}}}%
        {\EGAUnavailableHyperref{#2}{\normalfont{(\textbf{#1},~\StrBehind{#2}{.},~{#3})}}}}}%
    {\IfNoValueTF{#1}%
      {\IfNoValueTF{#3}%
        {\hyperref[#2]{\normalfont{(\StrBehind{#2}{.})}}}%
        {\hyperref[#2]{\normalfont{(\StrBehind{#2}{.},~{#3})}}}}%
      {\IfNoValueTF{#3}%
        {\hyperref[#2]{\normalfont{(\textbf{#1},~\StrBehind{#2}{.})}}}%
        {\hyperref[#2]{\normalfont{(\textbf{#1},~\StrBehind{#2}{.},~{#3})}}}}}%
}
\makeatother

% for marking changes following errata
% use as \erratum[volume]{correction} to say where the erratum is given
\newcommand{\erratum}[2][]{{#2}\marginpar{\footnotesize\textbf{Err}\textsubscript{#1}}}

% for referencing equations, use as \eref{eq:chapter-number.x.y.z}.
\newcommand{\eref}[1]{\hyperref[#1]{\normalfont{(\StrBehind{#1}{.})}}}
% END NATIVE INPUT preamble
% BEGIN NATIVE INPUT EGA4_reference_v2_registry
% Generated stable semantic destinations for the EGA IV reader.
% Visible text and numbering are unchanged.
\makeatletter
\newcommand{\sgaStableTarget}[1]{%
  \Hy@raisedlink{\hyper@anchorstart{#1}\hyper@anchorend}}
\expandafter\def\csname sga@target@IV.1.1.1\endcsname{ega4.c04.def.1-1-1}
\expandafter\def\csname sga@target@IV.1.1.2\endcsname{ega4.c04.prop.1-1-2}
\expandafter\def\csname sga@target@IV.1.1.3\endcsname{ega4.c04.prop.1-1-3}
\expandafter\def\csname sga@target@IV.1.1.4\endcsname{ega4.c04.cor.1-1-4}
\expandafter\def\csname sga@target@IV.1.1.5\endcsname{ega4.c04.prop.1-1-5}
\expandafter\def\csname sga@target@IV.1.1.6\endcsname{ega4.c04.prop.1-1-6}
\expandafter\def\csname sga@target@IV.1.10.1\endcsname{ega4.c04.thm.1-10-1}
\expandafter\def\csname sga@target@IV.1.10.2\endcsname{ega4.c04.cor.1-10-2}
\expandafter\def\csname sga@target@IV.1.10.3\endcsname{ega4.c04.prop.1-10-3}
\expandafter\def\csname sga@target@IV.1.10.4\endcsname{ega4.c04.cor.1-10-4}
\expandafter\def\csname sga@target@IV.1.2.1\endcsname{ega4.c04.def.1-2-1}
\expandafter\def\csname sga@target@IV.1.2.2\endcsname{ega4.c04.prop.1-2-2}
\expandafter\def\csname sga@target@IV.1.2.3\endcsname{ega4.c04.cor.1-2-3}
\expandafter\def\csname sga@target@IV.1.2.4\endcsname{ega4.c04.prop.1-2-4}
\expandafter\def\csname sga@target@IV.1.2.4.1\endcsname{ega4.c04.eq.1-2-4-1}
\expandafter\def\csname sga@target@IV.1.2.5\endcsname{ega4.c04.prop.1-2-5}
\expandafter\def\csname sga@target@IV.1.2.6\endcsname{ega4.c04.prop.1-2-6}
\expandafter\def\csname sga@target@IV.1.2.7\endcsname{ega4.c04.prop.1-2-7}
\expandafter\def\csname sga@target@IV.1.2.8\endcsname{ega4.c04.cor.1-2-8}
\expandafter\def\csname sga@target@IV.1.2.9\endcsname{ega4.c04.prop.1-2-9}
\expandafter\def\csname sga@target@IV.1.3.1\endcsname{ega4.c04.statement.1-3-1}
\expandafter\def\csname sga@target@IV.1.3.10\endcsname{ega4.c04.statement.1-3-10}
\expandafter\def\csname sga@target@IV.1.3.11\endcsname{ega4.c04.prop.1-3-11}
\expandafter\def\csname sga@target@IV.1.3.2\endcsname{ega4.c04.def.1-3-2}
\expandafter\def\csname sga@target@IV.1.3.3\endcsname{ega4.c04.prop.1-3-3}
\expandafter\def\csname sga@target@IV.1.3.4\endcsname{ega4.c04.prop.1-3-4}
\expandafter\def\csname sga@target@IV.1.3.5\endcsname{ega4.c04.cor.1-3-5}
\expandafter\def\csname sga@target@IV.1.3.6\endcsname{ega4.c04.prop.1-3-6}
\expandafter\def\csname sga@target@IV.1.3.7\endcsname{ega4.c04.prop.1-3-7}
\expandafter\def\csname sga@target@IV.1.3.8\endcsname{ega4.c04.statement.1-3-8}
\expandafter\def\csname sga@target@IV.1.3.9\endcsname{ega4.c04.prop.1-3-9}
\expandafter\def\csname sga@target@IV.1.4.1\endcsname{ega4.c04.statement.1-4-1}
\expandafter\def\csname sga@target@IV.1.4.2\endcsname{ega4.c04.def.1-4-2}
\expandafter\def\csname sga@target@IV.1.4.3\endcsname{ega4.c04.prop.1-4-3}
\expandafter\def\csname sga@target@IV.1.4.3.1\endcsname{ega4.c04.cor.1-4-3-1}
\expandafter\def\csname sga@target@IV.1.4.4\endcsname{ega4.c04.prop.1-4-4}
\expandafter\def\csname sga@target@IV.1.4.4.1\endcsname{ega4.c04.lem.1-4-4-1}
\expandafter\def\csname sga@target@IV.1.4.5\endcsname{ega4.c04.cor.1-4-5}
\expandafter\def\csname sga@target@IV.1.4.6\endcsname{ega4.c04.prop.1-4-6}
\expandafter\def\csname sga@target@IV.1.4.7\endcsname{ega4.c04.prop.1-4-7}
\expandafter\def\csname sga@target@IV.1.4.7.1\endcsname{ega4.c04.lem.1-4-7-1}
\expandafter\def\csname sga@target@IV.1.5.1\endcsname{ega4.c04.prop.1-5-1}
\expandafter\def\csname sga@target@IV.1.5.2\endcsname{ega4.c04.def.1-5-2}
\expandafter\def\csname sga@target@IV.1.5.3\endcsname{ega4.c04.prop.1-5-3}
\expandafter\def\csname sga@target@IV.1.5.4\endcsname{ega4.c04.prop.1-5-4}
\expandafter\def\csname sga@target@IV.1.5.5\endcsname{ega4.c04.cor.1-5-5}
\expandafter\def\csname sga@target@IV.1.5.6\endcsname{ega4.c04.cor.1-5-6}
\expandafter\def\csname sga@target@IV.1.5.7\endcsname{ega4.c04.prop.1-5-7}
\expandafter\def\csname sga@target@IV.1.6.1\endcsname{ega4.c04.def.1-6-1}
\expandafter\def\csname sga@target@IV.1.6.2\endcsname{ega4.c04.prop.1-6-2}
\expandafter\def\csname sga@target@IV.1.6.3\endcsname{ega4.c04.cor.1-6-3}
\expandafter\def\csname sga@target@IV.1.6.4\endcsname{ega4.c04.rem.1-6-4}
\expandafter\def\csname sga@target@IV.1.6.5\endcsname{ega4.c04.prop.1-6-5}
\expandafter\def\csname sga@target@IV.1.7.1\endcsname{ega4.c04.statement.1-7-1}
\expandafter\def\csname sga@target@IV.1.7.10\endcsname{ega4.c04.statement.1-7-10}
\expandafter\def\csname sga@target@IV.1.7.11\endcsname{ega4.c04.statement.1-7-11}
\expandafter\def\csname sga@target@IV.1.7.12\endcsname{ega4.c04.statement.1-7-12}
\expandafter\def\csname sga@target@IV.1.7.13\endcsname{ega4.c04.statement.1-7-13}
\expandafter\def\csname sga@target@IV.1.7.14\endcsname{ega4.c04.statement.1-7-14}
\expandafter\def\csname sga@target@IV.1.7.15\endcsname{ega4.c04.statement.1-7-15}
\expandafter\def\csname sga@target@IV.1.7.16\endcsname{ega4.c04.statement.1-7-16}
\expandafter\def\csname sga@target@IV.1.7.17\endcsname{ega4.c04.statement.1-7-17}
\expandafter\def\csname sga@target@IV.1.7.18\endcsname{ega4.c04.statement.1-7-18}
\expandafter\def\csname sga@target@IV.1.7.19\endcsname{ega4.c04.statement.1-7-19}
\expandafter\def\csname sga@target@IV.1.7.2\endcsname{ega4.c04.statement.1-7-2}
\expandafter\def\csname sga@target@IV.1.7.20\endcsname{ega4.c04.statement.1-7-20}
\expandafter\def\csname sga@target@IV.1.7.21\endcsname{ega4.c04.statement.1-7-21}
\expandafter\def\csname sga@target@IV.1.7.22\endcsname{ega4.c04.statement.1-7-22}
\expandafter\def\csname sga@target@IV.1.7.3\endcsname{ega4.c04.statement.1-7-3}
\expandafter\def\csname sga@target@IV.1.7.4\endcsname{ega4.c04.statement.1-7-4}
\expandafter\def\csname sga@target@IV.1.7.5\endcsname{ega4.c04.statement.1-7-5}
\expandafter\def\csname sga@target@IV.1.7.6\endcsname{ega4.c04.statement.1-7-6}
\expandafter\def\csname sga@target@IV.1.7.7\endcsname{ega4.c04.statement.1-7-7}
\expandafter\def\csname sga@target@IV.1.7.8\endcsname{ega4.c04.statement.1-7-8}
\expandafter\def\csname sga@target@IV.1.7.9\endcsname{ega4.c04.statement.1-7-9}
\expandafter\def\csname sga@target@IV.1.8.1\endcsname{ega4.c04.statement.1-8-1}
\expandafter\def\csname sga@target@IV.1.8.2\endcsname{ega4.c04.prop.1-8-2}
\expandafter\def\csname sga@target@IV.1.8.3\endcsname{ega4.c04.lem.1-8-3}
\expandafter\def\csname sga@target@IV.1.8.4\endcsname{ega4.c04.thm.1-8-4}
\expandafter\def\csname sga@target@IV.1.8.4.1\endcsname{ega4.c04.lem.1-8-4-1}
\expandafter\def\csname sga@target@IV.1.8.4.2\endcsname{ega4.c04.lem.1-8-4-2}
\expandafter\def\csname sga@target@IV.1.8.5\endcsname{ega4.c04.cor.1-8-5}
\expandafter\def\csname sga@target@IV.1.8.5.1\endcsname{ega4.c04.lem.1-8-5-1}
\expandafter\def\csname sga@target@IV.1.8.6\endcsname{ega4.c04.cor.1-8-6}
\expandafter\def\csname sga@target@IV.1.8.7\endcsname{ega4.c04.prop.1-8-7}
\expandafter\def\csname sga@target@IV.1.8.7.1\endcsname{ega4.c04.lem.1-8-7-1}
\expandafter\def\csname sga@target@IV.1.9.1\endcsname{ega4.c04.lem.1-9-1}
\expandafter\def\csname sga@target@IV.1.9.10\endcsname{ega4.c04.prop.1-9-10}
\expandafter\def\csname sga@target@IV.1.9.11\endcsname{ega4.c04.prop.1-9-11}
\expandafter\def\csname sga@target@IV.1.9.12\endcsname{ega4.c04.cor.1-9-12}
\expandafter\def\csname sga@target@IV.1.9.13\endcsname{ega4.c04.statement.1-9-13}
\expandafter\def\csname sga@target@IV.1.9.14\endcsname{ega4.c04.prop.1-9-14}
\expandafter\def\csname sga@target@IV.1.9.15\endcsname{ega4.c04.prop.1-9-15}
\expandafter\def\csname sga@target@IV.1.9.16\endcsname{ega4.c04.statement.1-9-16}
\expandafter\def\csname sga@target@IV.1.9.2\endcsname{ega4.c04.prop.1-9-2}
\expandafter\def\csname sga@target@IV.1.9.2.1\endcsname{ega4.c04.statement.1-9-2-1}
\expandafter\def\csname sga@target@IV.1.9.2.i\endcsname{ega4.c04.statement.1-9-2-i}
\expandafter\def\csname sga@target@IV.1.9.2.ii\endcsname{ega4.c04.statement.1-9-2-ii}
\expandafter\def\csname sga@target@IV.1.9.3\endcsname{ega4.c04.prop.1-9-3}
\expandafter\def\csname sga@target@IV.1.9.3.1\endcsname{ega4.c04.lem.1-9-3-1}
\expandafter\def\csname sga@target@IV.1.9.3.2\endcsname{ega4.c04.lem.1-9-3-2}
\expandafter\def\csname sga@target@IV.1.9.4\endcsname{ega4.c04.def.1-9-4}
\expandafter\def\csname sga@target@IV.1.9.5\endcsname{ega4.c04.prop.1-9-5}
\expandafter\def\csname sga@target@IV.1.9.6\endcsname{ega4.c04.statement.1-9-6}
\expandafter\def\csname sga@target@IV.1.9.7\endcsname{ega4.c04.thm.1-9-7}
\expandafter\def\csname sga@target@IV.1.9.8\endcsname{ega4.c04.cor.1-9-8}
\expandafter\def\csname sga@target@IV.1.9.9\endcsname{ega4.c04.cor.1-9-9}
\expandafter\def\csname sga@target@IV.10.1.1\endcsname{ega4.c04.statement.10-1-1}
\expandafter\def\csname sga@target@IV.10.1.2\endcsname{ega4.c04.prop.10-1-2}
\expandafter\def\csname sga@target@IV.10.1.3\endcsname{ega4.c04.def.10-1-3}
\expandafter\def\csname sga@target@IV.10.1.4\endcsname{ega4.c04.cor.10-1-4}
\expandafter\def\csname sga@target@IV.10.10.1\endcsname{ega4.c04.statement.10-10-1}
\expandafter\def\csname sga@target@IV.10.10.2\endcsname{ega4.c04.statement.10-10-2}
\expandafter\def\csname sga@target@IV.10.10.3\endcsname{ega4.c04.statement.10-10-3}
\expandafter\def\csname sga@target@IV.10.10.4\endcsname{ega4.c04.statement.10-10-4}
\expandafter\def\csname sga@target@IV.10.10.5\endcsname{ega4.c04.statement.10-10-5}
\expandafter\def\csname sga@target@IV.10.2.1\endcsname{ega4.c04.prop.10-2-1}
\expandafter\def\csname sga@target@IV.10.2.2\endcsname{ega4.c04.def.10-2-2}
\expandafter\def\csname sga@target@IV.10.2.3\endcsname{ega4.c04.cor.10-2-3}
\expandafter\def\csname sga@target@IV.10.2.4\endcsname{ega4.c04.cor.10-2-4}
\expandafter\def\csname sga@target@IV.10.2.5\endcsname{ega4.c04.cor.10-2-5}
\expandafter\def\csname sga@target@IV.10.2.6\endcsname{ega4.c04.cor.10-2-6}
\expandafter\def\csname sga@target@IV.10.2.7\endcsname{ega4.c04.prop.10-2-7}
\expandafter\def\csname sga@target@IV.10.2.8\endcsname{ega4.c04.item.10-2-8}
\expandafter\def\csname sga@target@IV.10.2.8.1\endcsname{ega4.c04.eq.10-2-8-1}
\expandafter\def\csname sga@target@IV.10.3.1\endcsname{ega4.c04.def.10-3-1}
\expandafter\def\csname sga@target@IV.10.3.2\endcsname{ega4.c04.prop.10-3-2}
\expandafter\def\csname sga@target@IV.10.3.3\endcsname{ega4.c04.prop.10-3-3}
\expandafter\def\csname sga@target@IV.10.4.1\endcsname{ega4.c04.def.10-4-1}
\expandafter\def\csname sga@target@IV.10.4.10\endcsname{ega4.c04.statement.10-4-10}
\expandafter\def\csname sga@target@IV.10.4.11\endcsname{ega4.c04.statement.10-4-11}
\expandafter\def\csname sga@target@IV.10.4.11.1\endcsname{ega4.c04.lem.10-4-11-1}
\expandafter\def\csname sga@target@IV.10.4.12\endcsname{ega4.c04.cor.10-4-12}
\expandafter\def\csname sga@target@IV.10.4.12.1\endcsname{ega4.c04.sec.10-4-12-1}
\expandafter\def\csname sga@target@IV.10.4.12.1.errata\endcsname{ega4.c04.eq.10-4-12-1-errata}
\expandafter\def\csname sga@target@IV.10.4.13\endcsname{ega4.c04.cor.10-4-13}
\expandafter\def\csname sga@target@IV.10.4.14\endcsname{ega4.c04.rem.10-4-14}
\expandafter\def\csname sga@target@IV.10.4.2\endcsname{ega4.c04.prop.10-4-2}
\expandafter\def\csname sga@target@IV.10.4.3\endcsname{ega4.c04.statement.10-4-3}
\expandafter\def\csname sga@target@IV.10.4.4\endcsname{ega4.c04.prop.10-4-4}
\expandafter\def\csname sga@target@IV.10.4.5\endcsname{ega4.c04.prop.10-4-5}
\expandafter\def\csname sga@target@IV.10.4.5.1\endcsname{ega4.c04.lem.10-4-5-1}
\expandafter\def\csname sga@target@IV.10.4.6\endcsname{ega4.c04.cor.10-4-6}
\expandafter\def\csname sga@target@IV.10.4.7\endcsname{ega4.c04.cor.10-4-7}
\expandafter\def\csname sga@target@IV.10.4.8\endcsname{ega4.c04.cor.10-4-8}
\expandafter\def\csname sga@target@IV.10.4.9\endcsname{ega4.c04.statement.10-4-9}
\expandafter\def\csname sga@target@IV.10.5.1\endcsname{ega4.c04.prop.10-5-1}
\expandafter\def\csname sga@target@IV.10.5.10\endcsname{ega4.c04.prop.10-5-10}
\expandafter\def\csname sga@target@IV.10.5.2\endcsname{ega4.c04.cor.10-5-2}
\expandafter\def\csname sga@target@IV.10.5.3\endcsname{ega4.c04.cor.10-5-3}
\expandafter\def\csname sga@target@IV.10.5.4\endcsname{ega4.c04.rem.10-5-4}
\expandafter\def\csname sga@target@IV.10.5.5\endcsname{ega4.c04.prop.10-5-5}
\expandafter\def\csname sga@target@IV.10.5.6\endcsname{ega4.c04.rem.10-5-6}
\expandafter\def\csname sga@target@IV.10.5.7\endcsname{ega4.c04.prop.10-5-7}
\expandafter\def\csname sga@target@IV.10.5.8\endcsname{ega4.c04.cor.10-5-8}
\expandafter\def\csname sga@target@IV.10.5.9\endcsname{ega4.c04.cor.10-5-9}
\expandafter\def\csname sga@target@IV.10.6.1\endcsname{ega4.c04.prop.10-6-1}
\expandafter\def\csname sga@target@IV.10.6.1.1\endcsname{ega4.c04.eq.10-6-1-1}
\expandafter\def\csname sga@target@IV.10.6.1.2\endcsname{ega4.c04.eq.10-6-1-2}
\expandafter\def\csname sga@target@IV.10.6.1.3\endcsname{ega4.c04.eq.10-6-1-3}
\expandafter\def\csname sga@target@IV.10.6.1.4\endcsname{ega4.c04.eq.10-6-1-4}
\expandafter\def\csname sga@target@IV.10.6.2\endcsname{ega4.c04.cor.10-6-2}
\expandafter\def\csname sga@target@IV.10.6.3\endcsname{ega4.c04.prop.10-6-3}
\expandafter\def\csname sga@target@IV.10.6.3.1\endcsname{ega4.c04.eq.10-6-3-1}
\expandafter\def\csname sga@target@IV.10.6.3.2\endcsname{ega4.c04.eq.10-6-3-2}
\expandafter\def\csname sga@target@IV.10.6.3.3\endcsname{ega4.c04.eq.10-6-3-3}
\expandafter\def\csname sga@target@IV.10.6.4\endcsname{ega4.c04.cor.10-6-4}
\expandafter\def\csname sga@target@IV.10.6.4.1\endcsname{ega4.c04.eq.10-6-4-1}
\expandafter\def\csname sga@target@IV.10.7.1\endcsname{ega4.c04.statement.10-7-1}
\expandafter\def\csname sga@target@IV.10.7.2\endcsname{ega4.c04.statement.10-7-2}
\expandafter\def\csname sga@target@IV.10.7.3\endcsname{ega4.c04.statement.10-7-3}
\expandafter\def\csname sga@target@IV.10.8.1\endcsname{ega4.c04.def.10-8-1}
\expandafter\def\csname sga@target@IV.10.8.1.1\endcsname{ega4.c04.eq.10-8-1-1}
\expandafter\def\csname sga@target@IV.10.8.1.2\endcsname{ega4.c04.eq.10-8-1-2}
\expandafter\def\csname sga@target@IV.10.8.2\endcsname{ega4.c04.statement.10-8-2}
\expandafter\def\csname sga@target@IV.10.8.3\endcsname{ega4.c04.lem.10-8-3}
\expandafter\def\csname sga@target@IV.10.8.3.1\endcsname{ega4.c04.eq.10-8-3-1}
\expandafter\def\csname sga@target@IV.10.8.4\endcsname{ega4.c04.cor.10-8-4}
\expandafter\def\csname sga@target@IV.10.8.5\endcsname{ega4.c04.prop.10-8-5}
\expandafter\def\csname sga@target@IV.10.8.6\endcsname{ega4.c04.cor.10-8-6}
\expandafter\def\csname sga@target@IV.10.8.7\endcsname{ega4.c04.prop.10-8-7}
\expandafter\def\csname sga@target@IV.10.8.7.1\endcsname{ega4.c04.eq.10-8-7-1}
\expandafter\def\csname sga@target@IV.10.8.8\endcsname{ega4.c04.cor.10-8-8}
\expandafter\def\csname sga@target@IV.10.9.1\endcsname{ega4.c04.statement.10-9-1}
\expandafter\def\csname sga@target@IV.10.9.2\endcsname{ega4.c04.statement.10-9-2}
\expandafter\def\csname sga@target@IV.10.9.3\endcsname{ega4.c04.statement.10-9-3}
\expandafter\def\csname sga@target@IV.10.9.4\endcsname{ega4.c04.prop.10-9-4}
\expandafter\def\csname sga@target@IV.10.9.5\endcsname{ega4.c04.statement.10-9-5}
\expandafter\def\csname sga@target@IV.10.9.6\endcsname{ega4.c04.prop.10-9-6}
\expandafter\def\csname sga@target@IV.11.1.1\endcsname{ega4.c04.thm.11-1-1}
\expandafter\def\csname sga@target@IV.11.1.1.1\endcsname{ega4.c04.sec.11-1-1-1}
\expandafter\def\csname sga@target@IV.11.1.1.2\endcsname{ega4.c04.lem.11-1-1-2}
\expandafter\def\csname sga@target@IV.11.1.2\endcsname{ega4.c04.cor.11-1-2}
\expandafter\def\csname sga@target@IV.11.1.3\endcsname{ega4.c04.cor.11-1-3}
\expandafter\def\csname sga@target@IV.11.1.4\endcsname{ega4.c04.cor.11-1-4}
\expandafter\def\csname sga@target@IV.11.1.5\endcsname{ega4.c04.cor.11-1-5}
\expandafter\def\csname sga@target@IV.11.1.6\endcsname{ega4.c04.cor.11-1-6}
\expandafter\def\csname sga@target@IV.11.10.1\endcsname{ega4.c04.prop.11-10-1}
\expandafter\def\csname sga@target@IV.11.10.1.1\endcsname{ega4.c04.eq.11-10-1-1}
\expandafter\def\csname sga@target@IV.11.10.1.2\endcsname{ega4.c04.eq.11-10-1-2}
\expandafter\def\csname sga@target@IV.11.10.10\endcsname{ega4.c04.prop.11-10-10}
\expandafter\def\csname sga@target@IV.11.10.2\endcsname{ega4.c04.def.11-10-2}
\expandafter\def\csname sga@target@IV.11.10.3\endcsname{ega4.c04.item.11-10-3}
\expandafter\def\csname sga@target@IV.11.10.4\endcsname{ega4.c04.prop.11-10-4}
\expandafter\def\csname sga@target@IV.11.10.5\endcsname{ega4.c04.thm.11-10-5}
\expandafter\def\csname sga@target@IV.11.10.6\endcsname{ega4.c04.prop.11-10-6}
\expandafter\def\csname sga@target@IV.11.10.7\endcsname{ega4.c04.cor.11-10-7}
\expandafter\def\csname sga@target@IV.11.10.8\endcsname{ega4.c04.def.11-10-8}
\expandafter\def\csname sga@target@IV.11.10.9\endcsname{ega4.c04.thm.11-10-9}
\expandafter\def\csname sga@target@IV.11.2.1\endcsname{ega4.c04.statement.11-2-1}
\expandafter\def\csname sga@target@IV.11.2.1.1\endcsname{ega4.c04.eq.11-2-1-1}
\expandafter\def\csname sga@target@IV.11.2.1.2\endcsname{ega4.c04.eq.11-2-1-2}
\expandafter\def\csname sga@target@IV.11.2.1.3\endcsname{ega4.c04.eq.11-2-1-3}
\expandafter\def\csname sga@target@IV.11.2.1.4\endcsname{ega4.c04.eq.11-2-1-4}
\expandafter\def\csname sga@target@IV.11.2.1.5\endcsname{ega4.c04.eq.11-2-1-5}
\expandafter\def\csname sga@target@IV.11.2.10\endcsname{ega4.c04.rem.11-2-10}
\expandafter\def\csname sga@target@IV.11.2.2\endcsname{ega4.c04.statement.11-2-2}
\expandafter\def\csname sga@target@IV.11.2.2.1\endcsname{ega4.c04.eq.11-2-2-1}
\expandafter\def\csname sga@target@IV.11.2.2.2\endcsname{ega4.c04.lem.11-2-2-2}
\expandafter\def\csname sga@target@IV.11.2.3\endcsname{ega4.c04.statement.11-2-3}
\expandafter\def\csname sga@target@IV.11.2.3.1\endcsname{ega4.c04.eq.11-2-3-1}
\expandafter\def\csname sga@target@IV.11.2.3.2\endcsname{ega4.c04.eq.11-2-3-2}
\expandafter\def\csname sga@target@IV.11.2.3.3\endcsname{ega4.c04.eq.11-2-3-3}
\expandafter\def\csname sga@target@IV.11.2.4\endcsname{ega4.c04.lem.11-2-4}
\expandafter\def\csname sga@target@IV.11.2.5\endcsname{ega4.c04.lem.11-2-5}
\expandafter\def\csname sga@target@IV.11.2.6\endcsname{ega4.c04.thm.11-2-6}
\expandafter\def\csname sga@target@IV.11.2.6.1\endcsname{ega4.c04.cor.11-2-6-1}
\expandafter\def\csname sga@target@IV.11.2.7\endcsname{ega4.c04.cor.11-2-7}
\expandafter\def\csname sga@target@IV.11.2.8\endcsname{ega4.c04.prop.11-2-8}
\expandafter\def\csname sga@target@IV.11.2.9\endcsname{ega4.c04.thm.11-2-9}
\expandafter\def\csname sga@target@IV.11.2.9.1\endcsname{ega4.c04.statement.11-2-9-1}
\expandafter\def\csname sga@target@IV.11.2.9.2\endcsname{ega4.c04.lem.11-2-9-2}
\expandafter\def\csname sga@target@IV.11.2.9.3\endcsname{ega4.c04.lem.11-2-9-3}
\expandafter\def\csname sga@target@IV.11.2.9.4\endcsname{ega4.c04.cor.11-2-9-4}
\expandafter\def\csname sga@target@IV.11.2.9.5\endcsname{ega4.c04.lem.11-2-9-5}
\expandafter\def\csname sga@target@IV.11.2.9.6\endcsname{ega4.c04.lem.11-2-9-6}
\expandafter\def\csname sga@target@IV.11.2.9.7\endcsname{ega4.c04.eq.11-2-9-7}
\expandafter\def\csname sga@target@IV.11.2.9.8\endcsname{ega4.c04.eq.11-2-9-8}
\expandafter\def\csname sga@target@IV.11.2.9.9\endcsname{ega4.c04.eq.11-2-9-9}
\expandafter\def\csname sga@target@IV.11.3\endcsname{ega4.c04.subsec.11-3.key-iv-11-3}
\expandafter\def\csname sga@target@IV.11.3.1\endcsname{ega4.c04.thm.11-3-1}
\expandafter\def\csname sga@target@IV.11.3.10\endcsname{ega4.c04.thm.11-3-10}
\expandafter\def\csname sga@target@IV.11.3.10.1\endcsname{ega4.c04.lem.11-3-10-1}
\expandafter\def\csname sga@target@IV.11.3.10.2\endcsname{ega4.c04.lem.11-3-10-2}
\expandafter\def\csname sga@target@IV.11.3.10.3\endcsname{ega4.c04.eq.11-3-10-3}
\expandafter\def\csname sga@target@IV.11.3.10.4\endcsname{ega4.c04.sec.11-3-10-4}
\expandafter\def\csname sga@target@IV.11.3.11\endcsname{ega4.c04.cor.11-3-11}
\expandafter\def\csname sga@target@IV.11.3.12\endcsname{ega4.c04.rem.11-3-12}
\expandafter\def\csname sga@target@IV.11.3.13\endcsname{ega4.c04.prop.11-3-13}
\expandafter\def\csname sga@target@IV.11.3.14\endcsname{ega4.c04.cor.11-3-14}
\expandafter\def\csname sga@target@IV.11.3.15\endcsname{ega4.c04.prop.11-3-15}
\expandafter\def\csname sga@target@IV.11.3.16\endcsname{ega4.c04.prop.11-3-16}
\expandafter\def\csname sga@target@IV.11.3.17\endcsname{ega4.c04.cor.11-3-17}
\expandafter\def\csname sga@target@IV.11.3.2\endcsname{ega4.c04.cor.11-3-2}
\expandafter\def\csname sga@target@IV.11.3.3\endcsname{ega4.c04.cor.11-3-3}
\expandafter\def\csname sga@target@IV.11.3.4\endcsname{ega4.c04.prop.11-3-4}
\expandafter\def\csname sga@target@IV.11.3.5\endcsname{ega4.c04.cor.11-3-5}
\expandafter\def\csname sga@target@IV.11.3.6\endcsname{ega4.c04.prop.11-3-6}
\expandafter\def\csname sga@target@IV.11.3.7\endcsname{ega4.c04.prop.11-3-7}
\expandafter\def\csname sga@target@IV.11.3.8\endcsname{ega4.c04.thm.11-3-8}
\expandafter\def\csname sga@target@IV.11.3.9\endcsname{ega4.c04.prop.11-3-9}
\expandafter\def\csname sga@target@IV.11.3.9.1\endcsname{ega4.c04.lem.11-3-9-1}
\expandafter\def\csname sga@target@IV.11.4.1\endcsname{ega4.c04.thm.11-4-1}
\expandafter\def\csname sga@target@IV.11.4.1.1\endcsname{ega4.c04.lem.11-4-1-1}
\expandafter\def\csname sga@target@IV.11.4.1.2\endcsname{ega4.c04.lem.11-4-1-2}
\expandafter\def\csname sga@target@IV.11.4.10\endcsname{ega4.c04.cor.11-4-10}
\expandafter\def\csname sga@target@IV.11.4.11\endcsname{ega4.c04.prop.11-4-11}
\expandafter\def\csname sga@target@IV.11.4.12\endcsname{ega4.c04.prop.11-4-12}
\expandafter\def\csname sga@target@IV.11.4.13\endcsname{ega4.c04.cor.11-4-13}
\expandafter\def\csname sga@target@IV.11.4.2\endcsname{ega4.c04.rem.11-4-2}
\expandafter\def\csname sga@target@IV.11.4.3\endcsname{ega4.c04.cor.11-4-3}
\expandafter\def\csname sga@target@IV.11.4.4\endcsname{ega4.c04.cor.11-4-4}
\expandafter\def\csname sga@target@IV.11.4.5\endcsname{ega4.c04.prop.11-4-5}
\expandafter\def\csname sga@target@IV.11.4.6\endcsname{ega4.c04.cor.11-4-6}
\expandafter\def\csname sga@target@IV.11.4.7\endcsname{ega4.c04.prop.11-4-7}
\expandafter\def\csname sga@target@IV.11.4.8\endcsname{ega4.c04.rem.11-4-8}
\expandafter\def\csname sga@target@IV.11.4.9\endcsname{ega4.c04.cor.11-4-9}
\expandafter\def\csname sga@target@IV.11.5.1\endcsname{ega4.c04.thm.11-5-1}
\expandafter\def\csname sga@target@IV.11.5.2\endcsname{ega4.c04.cor.11-5-2}
\expandafter\def\csname sga@target@IV.11.5.3\endcsname{ega4.c04.prop.11-5-3}
\expandafter\def\csname sga@target@IV.11.5.4\endcsname{ega4.c04.cor.11-5-4}
\expandafter\def\csname sga@target@IV.11.5.5\endcsname{ega4.c04.prop.11-5-5}
\expandafter\def\csname sga@target@IV.11.5.5.1\endcsname{ega4.c04.lem.11-5-5-1}
\expandafter\def\csname sga@target@IV.11.5.5.2\endcsname{ega4.c04.lem.11-5-5-2}
\expandafter\def\csname sga@target@IV.11.6.1\endcsname{ega4.c04.thm.11-6-1}
\expandafter\def\csname sga@target@IV.11.6.2\endcsname{ega4.c04.cor.11-6-2}
\expandafter\def\csname sga@target@IV.11.7.1\endcsname{ega4.c04.statement.11-7-1}
\expandafter\def\csname sga@target@IV.11.7.2\endcsname{ega4.c04.statement.11-7-2}
\expandafter\def\csname sga@target@IV.11.7.3\endcsname{ega4.c04.statement.11-7-3}
\expandafter\def\csname sga@target@IV.11.7.4\endcsname{ega4.c04.statement.11-7-4}
\expandafter\def\csname sga@target@IV.11.7.5\endcsname{ega4.c04.statement.11-7-5}
\expandafter\def\csname sga@target@IV.11.8.1\endcsname{ega4.c04.thm.11-8-1}
\expandafter\def\csname sga@target@IV.11.8.2\endcsname{ega4.c04.subsec.11-8-2}
\expandafter\def\csname sga@target@IV.11.9.1\endcsname{ega4.c04.statement.11-9-1}
\expandafter\def\csname sga@target@IV.11.9.1.1\endcsname{ega4.c04.statement.11-9-1-1}
\expandafter\def\csname sga@target@IV.11.9.10\endcsname{ega4.c04.thm.11-9-10}
\expandafter\def\csname sga@target@IV.11.9.11\endcsname{ega4.c04.rem.11-9-11}
\expandafter\def\csname sga@target@IV.11.9.12\endcsname{ega4.c04.lem.11-9-12}
\expandafter\def\csname sga@target@IV.11.9.13\endcsname{ega4.c04.prop.11-9-13}
\expandafter\def\csname sga@target@IV.11.9.14\endcsname{ega4.c04.statement.11-9-14}
\expandafter\def\csname sga@target@IV.11.9.15\endcsname{ega4.c04.prop.11-9-15}
\expandafter\def\csname sga@target@IV.11.9.16\endcsname{ega4.c04.thm.11-9-16}
\expandafter\def\csname sga@target@IV.11.9.17\endcsname{ega4.c04.prop.11-9-17}
\expandafter\def\csname sga@target@IV.11.9.17.1\endcsname{ega4.c04.lem.11-9-17-1}
\expandafter\def\csname sga@target@IV.11.9.18\endcsname{ega4.c04.item.11-9-18}
\expandafter\def\csname sga@target@IV.11.9.2\endcsname{ega4.c04.statement.11-9-2}
\expandafter\def\csname sga@target@IV.11.9.3\endcsname{ega4.c04.statement.11-9-3}
\expandafter\def\csname sga@target@IV.11.9.4\endcsname{ega4.c04.statement.11-9-4}
\expandafter\def\csname sga@target@IV.11.9.5\endcsname{ega4.c04.statement.11-9-5}
\expandafter\def\csname sga@target@IV.11.9.6\endcsname{ega4.c04.statement.11-9-6}
\expandafter\def\csname sga@target@IV.11.9.7\endcsname{ega4.c04.prop.11-9-7}
\expandafter\def\csname sga@target@IV.11.9.7.1\endcsname{ega4.c04.lem.11-9-7-1}
\expandafter\def\csname sga@target@IV.11.9.8\endcsname{ega4.c04.subsec.11-9-8}
\expandafter\def\csname sga@target@IV.11.9.9\endcsname{ega4.c04.cor.11-9-9}
\expandafter\def\csname sga@target@IV.12.0.1\endcsname{ega4.c04.statement.12-0-1}
\expandafter\def\csname sga@target@IV.12.0.2\endcsname{ega4.c04.statement.12-0-2}
\expandafter\def\csname sga@target@IV.12.1.1\endcsname{ega4.c04.thm.12-1-1}
\expandafter\def\csname sga@target@IV.12.1.1.1\endcsname{ega4.c04.statement.12-1-1-1}
\expandafter\def\csname sga@target@IV.12.1.1.2\endcsname{ega4.c04.statement.12-1-1-2}
\expandafter\def\csname sga@target@IV.12.1.1.3\endcsname{ega4.c04.statement.12-1-1-3}
\expandafter\def\csname sga@target@IV.12.1.1.4\endcsname{ega4.c04.statement.12-1-1-4}
\expandafter\def\csname sga@target@IV.12.1.1.5\endcsname{ega4.c04.prop.12-1-1-5}
\expandafter\def\csname sga@target@IV.12.1.2\endcsname{ega4.c04.item.12-1-2}
\expandafter\def\csname sga@target@IV.12.1.3\endcsname{ega4.c04.cor.12-1-3}
\expandafter\def\csname sga@target@IV.12.1.4\endcsname{ega4.c04.rem.12-1-4}
\expandafter\def\csname sga@target@IV.12.1.5\endcsname{ega4.c04.cor.12-1-5}
\expandafter\def\csname sga@target@IV.12.1.6\endcsname{ega4.c04.thm.12-1-6}
\expandafter\def\csname sga@target@IV.12.1.7\endcsname{ega4.c04.cor.12-1-7}
\expandafter\def\csname sga@target@IV.12.1.8\endcsname{ega4.c04.item.12-1-8}
\expandafter\def\csname sga@target@IV.12.2.1\endcsname{ega4.c04.thm.12-2-1}
\expandafter\def\csname sga@target@IV.12.2.1.1\endcsname{ega4.c04.statement.12-2-1-1}
\expandafter\def\csname sga@target@IV.12.2.1.2\endcsname{ega4.c04.statement.12-2-1-2}
\expandafter\def\csname sga@target@IV.12.2.1.3\endcsname{ega4.c04.statement.12-2-1-3}
\expandafter\def\csname sga@target@IV.12.2.1.4\endcsname{ega4.c04.lem.12-2-1-4}
\expandafter\def\csname sga@target@IV.12.2.2\endcsname{ega4.c04.cor.12-2-2}
\expandafter\def\csname sga@target@IV.12.2.3\endcsname{ega4.c04.item.12-2-3}
\expandafter\def\csname sga@target@IV.12.2.4\endcsname{ega4.c04.thm.12-2-4}
\expandafter\def\csname sga@target@IV.12.3.1\endcsname{ega4.c04.lem.12-3-1}
\expandafter\def\csname sga@target@IV.12.3.2\endcsname{ega4.c04.prop.12-3-2}
\expandafter\def\csname sga@target@IV.12.3.2.1\endcsname{ega4.c04.eq.12-3-2-1}
\expandafter\def\csname sga@target@IV.12.3.3\endcsname{ega4.c04.prop.12-3-3}
\expandafter\def\csname sga@target@IV.12.3.3.1\endcsname{ega4.c04.lem.12-3-3-1}
\expandafter\def\csname sga@target@IV.12.3.3.2\endcsname{ega4.c04.eq.12-3-3-2}
\expandafter\def\csname sga@target@IV.12.3.3.3\endcsname{ega4.c04.lem.12-3-3-3}
\expandafter\def\csname sga@target@IV.12.3.3.4\endcsname{ega4.c04.eq.12-3-3-4}
\expandafter\def\csname sga@target@IV.12.3.3.5\endcsname{ega4.c04.lem.12-3-3-5}
\expandafter\def\csname sga@target@IV.12.3.4\endcsname{ega4.c04.cor.12-3-4}
\expandafter\def\csname sga@target@IV.13.1.1\endcsname{ega4.c04.lem.13-1-1}
\expandafter\def\csname sga@target@IV.13.1.2\endcsname{ega4.c04.lem.13-1-2}
\expandafter\def\csname sga@target@IV.13.1.3\endcsname{ega4.c04.thm.13-1-3}
\expandafter\def\csname sga@target@IV.13.1.4\endcsname{ega4.c04.cor.13-1-4}
\expandafter\def\csname sga@target@IV.13.1.5\endcsname{ega4.c04.cor.13-1-5}
\expandafter\def\csname sga@target@IV.13.1.6\endcsname{ega4.c04.cor.13-1-6}
\expandafter\def\csname sga@target@IV.13.1.7\endcsname{ega4.c04.prop.13-1-7}
\expandafter\def\csname sga@target@IV.13.2.1\endcsname{ega4.c04.statement.13-2-1}
\expandafter\def\csname sga@target@IV.13.2.10\endcsname{ega4.c04.prop.13-2-10}
\expandafter\def\csname sga@target@IV.13.2.10.1\endcsname{ega4.c04.eq.13-2-10-1}
\expandafter\def\csname sga@target@IV.13.2.10.2\endcsname{ega4.c04.eq.13-2-10-2}
\expandafter\def\csname sga@target@IV.13.2.11\endcsname{ega4.c04.cor.13-2-11}
\expandafter\def\csname sga@target@IV.13.2.12\endcsname{ega4.c04.item.13-2-12}
\expandafter\def\csname sga@target@IV.13.2.2\endcsname{ega4.c04.def.13-2-2}
\expandafter\def\csname sga@target@IV.13.2.3\endcsname{ega4.c04.prop.13-2-3}
\expandafter\def\csname sga@target@IV.13.2.3.1\endcsname{ega4.c04.prop.13-2-3-1}
\expandafter\def\csname sga@target@IV.13.2.4\endcsname{ega4.c04.statement.13-2-4}
\expandafter\def\csname sga@target@IV.13.2.4.1\endcsname{ega4.c04.statement.13-2-4-1}
\expandafter\def\csname sga@target@IV.13.2.4.2\endcsname{ega4.c04.eq.13-2-4-2}
\expandafter\def\csname sga@target@IV.13.2.4.3\endcsname{ega4.c04.eq.13-2-4-3}
\expandafter\def\csname sga@target@IV.13.2.4.4\endcsname{ega4.c04.eq.13-2-4-4}
\expandafter\def\csname sga@target@IV.13.2.4.5\endcsname{ega4.c04.eq.13-2-4-5}
\expandafter\def\csname sga@target@IV.13.2.5\endcsname{ega4.c04.lem.13-2-5}
\expandafter\def\csname sga@target@IV.13.2.6\endcsname{ega4.c04.statement.13-2-6}
\expandafter\def\csname sga@target@IV.13.2.6.1\endcsname{ega4.c04.statement.13-2-6-1}
\expandafter\def\csname sga@target@IV.13.2.7\endcsname{ega4.c04.prop.13-2-7}
\expandafter\def\csname sga@target@IV.13.2.7.1\endcsname{ega4.c04.prop.13-2-7-1}
\expandafter\def\csname sga@target@IV.13.2.8\endcsname{ega4.c04.cor.13-2-8}
\expandafter\def\csname sga@target@IV.13.2.9\endcsname{ega4.c04.cor.13-2-9}
\expandafter\def\csname sga@target@IV.13.2.9.1\endcsname{ega4.c04.eq.13-2-9-1}
\expandafter\def\csname sga@target@IV.13.2.9.2\endcsname{ega4.c04.eq.13-2-9-2}
\expandafter\def\csname sga@target@IV.13.3.1\endcsname{ega4.c04.prop.13-3-1}
\expandafter\def\csname sga@target@IV.13.3.1.1\endcsname{ega4.c04.lem.13-3-1-1}
\expandafter\def\csname sga@target@IV.13.3.2\endcsname{ega4.c04.def.13-3-2}
\expandafter\def\csname sga@target@IV.13.3.3\endcsname{ega4.c04.prop.13-3-3}
\expandafter\def\csname sga@target@IV.13.3.4\endcsname{ega4.c04.cor.13-3-4}
\expandafter\def\csname sga@target@IV.13.3.4.1\endcsname{ega4.c04.cor.13-3-4-1}
\expandafter\def\csname sga@target@IV.13.3.5\endcsname{ega4.c04.cor.13-3-5}
\expandafter\def\csname sga@target@IV.13.3.6\endcsname{ega4.c04.prop.13-3-6}
\expandafter\def\csname sga@target@IV.13.3.6.1\endcsname{ega4.c04.prop.13-3-6-1}
\expandafter\def\csname sga@target@IV.13.3.6.2\endcsname{ega4.c04.eq.13-3-6-2}
\expandafter\def\csname sga@target@IV.13.3.7\endcsname{ega4.c04.prop.13-3-7}
\expandafter\def\csname sga@target@IV.13.3.7.1\endcsname{ega4.c04.prop.13-3-7-1}
\expandafter\def\csname sga@target@IV.13.3.8\endcsname{ega4.c04.prop.13-3-8}
\expandafter\def\csname sga@target@IV.13.3.9\endcsname{ega4.c04.statement.13-3-9}
\expandafter\def\csname sga@target@IV.14.1.1\endcsname{ega4.c04.statement.14-1-1}
\expandafter\def\csname sga@target@IV.14.1.2\endcsname{ega4.c04.prop.14-1-2}
\expandafter\def\csname sga@target@IV.14.1.3\endcsname{ega4.c04.item.14-1-3}
\expandafter\def\csname sga@target@IV.14.1.4\endcsname{ega4.c04.prop.14-1-4}
\expandafter\def\csname sga@target@IV.14.1.5\endcsname{ega4.c04.rem.14-1-5}
\expandafter\def\csname sga@target@IV.14.1.6\endcsname{ega4.c04.statement.14-1-6}
\expandafter\def\csname sga@target@IV.14.2.1\endcsname{ega4.c04.thm.14-2-1}
\expandafter\def\csname sga@target@IV.14.2.1.1\endcsname{ega4.c04.eq.14-2-1-1}
\expandafter\def\csname sga@target@IV.14.2.2\endcsname{ega4.c04.cor.14-2-2}
\expandafter\def\csname sga@target@IV.14.2.3\endcsname{ega4.c04.prop.14-2-3}
\expandafter\def\csname sga@target@IV.14.2.4\endcsname{ega4.c04.cor.14-2-4}
\expandafter\def\csname sga@target@IV.14.2.5\endcsname{ega4.c04.cor.14-2-5}
\expandafter\def\csname sga@target@IV.14.2.6\endcsname{ega4.c04.cor.14-2-6}
\expandafter\def\csname sga@target@IV.14.2.7\endcsname{ega4.c04.rem.14-2-7}
\expandafter\def\csname sga@target@IV.14.3\endcsname{ega4.c04.subsec.14-3.key-iv-14-3}
\expandafter\def\csname sga@target@IV.14.3.1\endcsname{ega4.c04.statement.14-3-1}
\expandafter\def\csname sga@target@IV.14.3.10\endcsname{ega4.c04.lem.14-3-10}
\expandafter\def\csname sga@target@IV.14.3.11\endcsname{ega4.c04.prop.14-3-11}
\expandafter\def\csname sga@target@IV.14.3.12\endcsname{ega4.c04.rem.14-3-12}
\expandafter\def\csname sga@target@IV.14.3.13\endcsname{ega4.c04.prop.14-3-13}
\expandafter\def\csname sga@target@IV.14.3.14\endcsname{ega4.c04.rem.14-3-14}
\expandafter\def\csname sga@target@IV.14.3.2\endcsname{ega4.c04.prop.14-3-2}
\expandafter\def\csname sga@target@IV.14.3.3\endcsname{ega4.c04.def.14-3-3}
\expandafter\def\csname sga@target@IV.14.3.3.1\endcsname{ega4.c04.subsec.14-3-3-1}
\expandafter\def\csname sga@target@IV.14.3.4\endcsname{ega4.c04.prop.14-3-4}
\expandafter\def\csname sga@target@IV.14.3.5\endcsname{ega4.c04.prop.14-3-5}
\expandafter\def\csname sga@target@IV.14.3.6\endcsname{ega4.c04.prop.14-3-6}
\expandafter\def\csname sga@target@IV.14.3.7\endcsname{ega4.c04.cor.14-3-7}
\expandafter\def\csname sga@target@IV.14.3.8\endcsname{ega4.c04.prop.14-3-8}
\expandafter\def\csname sga@target@IV.14.3.9\endcsname{ega4.c04.item.14-3-9}
\expandafter\def\csname sga@target@IV.14.4.1\endcsname{ega4.c04.thm.14-4-1}
\expandafter\def\csname sga@target@IV.14.4.1.1\endcsname{ega4.c04.lem.14-4-1-1}
\expandafter\def\csname sga@target@IV.14.4.1.2\endcsname{ega4.c04.lem.14-4-1-2}
\expandafter\def\csname sga@target@IV.14.4.1.3\endcsname{ega4.c04.lem.14-4-1-3}
\expandafter\def\csname sga@target@IV.14.4.10\endcsname{ega4.c04.item.14-4-10}
\expandafter\def\csname sga@target@IV.14.4.2\endcsname{ega4.c04.cor.14-4-2}
\expandafter\def\csname sga@target@IV.14.4.3\endcsname{ega4.c04.cor.14-4-3}
\expandafter\def\csname sga@target@IV.14.4.4\endcsname{ega4.c04.cor.14-4-4}
\expandafter\def\csname sga@target@IV.14.4.5\endcsname{ega4.c04.rem.14-4-5}
\expandafter\def\csname sga@target@IV.14.4.6\endcsname{ega4.c04.cor.14-4-6}
\expandafter\def\csname sga@target@IV.14.4.7\endcsname{ega4.c04.prop.14-4-7}
\expandafter\def\csname sga@target@IV.14.4.7.1\endcsname{ega4.c04.eq.14-4-7-1}
\expandafter\def\csname sga@target@IV.14.4.7.2\endcsname{ega4.c04.eq.14-4-7-2}
\expandafter\def\csname sga@target@IV.14.4.8\endcsname{ega4.c04.prop.14-4-8}
\expandafter\def\csname sga@target@IV.14.4.8.1\endcsname{ega4.c04.lem.14-4-8-1}
\expandafter\def\csname sga@target@IV.14.4.9\endcsname{ega4.c04.cor.14-4-9}
\expandafter\def\csname sga@target@IV.14.5\endcsname{ega4.c04.subsec.14-5.key-iv-14-5}
\expandafter\def\csname sga@target@IV.14.5.1\endcsname{ega4.c04.lem.14-5-1}
\expandafter\def\csname sga@target@IV.14.5.1.1\endcsname{ega4.c04.eq.14-5-1-1}
\expandafter\def\csname sga@target@IV.14.5.1.2\endcsname{ega4.c04.eq.14-5-1-2}
\expandafter\def\csname sga@target@IV.14.5.1.3\endcsname{ega4.c04.eq.14-5-1-3}
\expandafter\def\csname sga@target@IV.14.5.1.4\endcsname{ega4.c04.eq.14-5-1-4}
\expandafter\def\csname sga@target@IV.14.5.10\endcsname{ega4.c04.prop.14-5-10}
\expandafter\def\csname sga@target@IV.14.5.11\endcsname{ega4.c04.cor.14-5-11}
\expandafter\def\csname sga@target@IV.14.5.11.1\endcsname{ega4.c04.lem.14-5-11-1}
\expandafter\def\csname sga@target@IV.14.5.12\endcsname{ega4.c04.cor.14-5-12}
\expandafter\def\csname sga@target@IV.14.5.2\endcsname{ega4.c04.cor.14-5-2}
\expandafter\def\csname sga@target@IV.14.5.2.1\endcsname{ega4.c04.eq.14-5-2-1}
\expandafter\def\csname sga@target@IV.14.5.3\endcsname{ega4.c04.prop.14-5-3}
\expandafter\def\csname sga@target@IV.14.5.4\endcsname{ega4.c04.cor.14-5-4}
\expandafter\def\csname sga@target@IV.14.5.5\endcsname{ega4.c04.rem.14-5-5}
\expandafter\def\csname sga@target@IV.14.5.6\endcsname{ega4.c04.prop.14-5-6}
\expandafter\def\csname sga@target@IV.14.5.7\endcsname{ega4.c04.rem.14-5-7}
\expandafter\def\csname sga@target@IV.14.5.8\endcsname{ega4.c04.prop.14-5-8}
\expandafter\def\csname sga@target@IV.14.5.9\endcsname{ega4.c04.prop.14-5-9}
\expandafter\def\csname sga@target@IV.15.1.1\endcsname{ega4.c04.prop.15-1-1}
\expandafter\def\csname sga@target@IV.15.1.1-proof\endcsname{ega4.c04.statement.15-1-1-proof}
\expandafter\def\csname sga@target@IV.15.1.1.1\endcsname{ega4.c04.eq.15-1-1-1}
\expandafter\def\csname sga@target@IV.15.1.1.2\endcsname{ega4.c04.eq.15-1-1-2}
\expandafter\def\csname sga@target@IV.15.1.1.3\endcsname{ega4.c04.eq.15-1-1-3}
\expandafter\def\csname sga@target@IV.15.1.1.4\endcsname{ega4.c04.eq.15-1-1-4}
\expandafter\def\csname sga@target@IV.15.1.1.5\endcsname{ega4.c04.eq.15-1-1-5}
\expandafter\def\csname sga@target@IV.15.1.1.6\endcsname{ega4.c04.lem.15-1-1-6}
\expandafter\def\csname sga@target@IV.15.1.2\endcsname{ega4.c04.cor.15-1-2}
\expandafter\def\csname sga@target@IV.15.1.3\endcsname{ega4.c04.cor.15-1-3}
\expandafter\def\csname sga@target@IV.15.1.4\endcsname{ega4.c04.statement.15-1-4}
\expandafter\def\csname sga@target@IV.15.2.1\endcsname{ega4.c04.statement.15-2-1}
\expandafter\def\csname sga@target@IV.15.2.2\endcsname{ega4.c04.thm.15-2-2}
\expandafter\def\csname sga@target@IV.15.2.2-proof\endcsname{ega4.c04.item.15-2-2-proof}
\expandafter\def\csname sga@target@IV.15.2.2.1\endcsname{ega4.c04.lem.15-2-2-1}
\expandafter\def\csname sga@target@IV.15.2.3\endcsname{ega4.c04.cor.15-2-3}
\expandafter\def\csname sga@target@IV.15.2.4\endcsname{ega4.c04.item.15-2-4}
\expandafter\def\csname sga@target@IV.15.3.1\endcsname{ega4.c04.prop.15-3-1}
\expandafter\def\csname sga@target@IV.15.3.1-proof\endcsname{ega4.c04.item.15-3-1-proof}
\expandafter\def\csname sga@target@IV.15.3.2\endcsname{ega4.c04.prop.15-3-2}
\expandafter\def\csname sga@target@IV.15.3.2-proof\endcsname{ega4.c04.item.15-3-2-proof}
\expandafter\def\csname sga@target@IV.15.3.3\endcsname{ega4.c04.cor.15-3-3}
\expandafter\def\csname sga@target@IV.15.3.3-proof\endcsname{ega4.c04.item.15-3-3-proof}
\expandafter\def\csname sga@target@IV.15.3.4\endcsname{ega4.c04.item.15-3-4}
\expandafter\def\csname sga@target@IV.15.4.1\endcsname{ega4.c04.prop.15-4-1}
\expandafter\def\csname sga@target@IV.15.4.2\endcsname{ega4.c04.prop.15-4-2}
\expandafter\def\csname sga@target@IV.15.4.3\endcsname{ega4.c04.prop.15-4-3}
\expandafter\def\csname sga@target@IV.15.4.3-proof\endcsname{ega4.c04.item.15-4-3-proof}
\expandafter\def\csname sga@target@IV.15.5.1\endcsname{ega4.c04.prop.15-5-1}
\expandafter\def\csname sga@target@IV.15.5.10\endcsname{ega4.c04.item.15-5-10}
\expandafter\def\csname sga@target@IV.15.5.2\endcsname{ega4.c04.cor.15-5-2}
\expandafter\def\csname sga@target@IV.15.5.2.1\endcsname{ega4.c04.eq.15-5-2-1}
\expandafter\def\csname sga@target@IV.15.5.2.2\endcsname{ega4.c04.lem.15-5-2-2}
\expandafter\def\csname sga@target@IV.15.5.2.3\endcsname{ega4.c04.eq.15-5-2-3}
\expandafter\def\csname sga@target@IV.15.5.2.4\endcsname{ega4.c04.eq.15-5-2-4}
\expandafter\def\csname sga@target@IV.15.5.3\endcsname{ega4.c04.prop.15-5-3}
\expandafter\def\csname sga@target@IV.15.5.4\endcsname{ega4.c04.cor.15-5-4}
\expandafter\def\csname sga@target@IV.15.5.5\endcsname{ega4.c04.statement.15-5-5}
\expandafter\def\csname sga@target@IV.15.5.6\endcsname{ega4.c04.lem.15-5-6}
\expandafter\def\csname sga@target@IV.15.5.6.1\endcsname{ega4.c04.lem.15-5-6-1}
\expandafter\def\csname sga@target@IV.15.5.7\endcsname{ega4.c04.prop.15-5-7}
\expandafter\def\csname sga@target@IV.15.5.7.1\endcsname{ega4.c04.eq.15-5-7-1}
\expandafter\def\csname sga@target@IV.15.5.8\endcsname{ega4.c04.statement.15-5-8}
\expandafter\def\csname sga@target@IV.15.5.9\endcsname{ega4.c04.prop.15-5-9}
\expandafter\def\csname sga@target@IV.15.6.1\endcsname{ega4.c04.prop.15-6-1}
\expandafter\def\csname sga@target@IV.15.6.1.1\endcsname{ega4.c04.eq.15-6-1-1}
\expandafter\def\csname sga@target@IV.15.6.10\endcsname{ega4.c04.item.15-6-10}
\expandafter\def\csname sga@target@IV.15.6.2\endcsname{ega4.c04.cor.15-6-2}
\expandafter\def\csname sga@target@IV.15.6.3\endcsname{ega4.c04.prop.15-6-3}
\expandafter\def\csname sga@target@IV.15.6.4\endcsname{ega4.c04.prop.15-6-4}
\expandafter\def\csname sga@target@IV.15.6.5\endcsname{ega4.c04.cor.15-6-5}
\expandafter\def\csname sga@target@IV.15.6.6\endcsname{ega4.c04.prop.15-6-6}
\expandafter\def\csname sga@target@IV.15.6.7\endcsname{ega4.c04.cor.15-6-7}
\expandafter\def\csname sga@target@IV.15.6.8\endcsname{ega4.c04.prop.15-6-8}
\expandafter\def\csname sga@target@IV.15.6.9\endcsname{ega4.c04.cor.15-6-9}
\expandafter\def\csname sga@target@IV.15.7.1\endcsname{ega4.c04.def.15-7-1}
\expandafter\def\csname sga@target@IV.15.7.10\endcsname{ega4.c04.cor.15-7-10}
\expandafter\def\csname sga@target@IV.15.7.11\endcsname{ega4.c04.cor.15-7-11}
\expandafter\def\csname sga@target@IV.15.7.11.1\endcsname{ega4.c04.lem.15-7-11-1}
\expandafter\def\csname sga@target@IV.15.7.2\endcsname{ega4.c04.prop.15-7-2}
\expandafter\def\csname sga@target@IV.15.7.3\endcsname{ega4.c04.rem.15-7-3}
\expandafter\def\csname sga@target@IV.15.7.4\endcsname{ega4.c04.cor.15-7-4}
\expandafter\def\csname sga@target@IV.15.7.5\endcsname{ega4.c04.cor.15-7-5}
\expandafter\def\csname sga@target@IV.15.7.6\endcsname{ega4.c04.cor.15-7-6}
\expandafter\def\csname sga@target@IV.15.7.7\endcsname{ega4.c04.item.15-7-7}
\expandafter\def\csname sga@target@IV.15.7.8\endcsname{ega4.c04.statement.15-7-8}
\expandafter\def\csname sga@target@IV.15.7.9\endcsname{ega4.c04.prop.15-7-9}
\expandafter\def\csname sga@target@IV.16.1\endcsname{ega4.c04.subsec.16-1}
\expandafter\def\csname sga@target@IV.16.1.1\endcsname{ega4.c04.statement.16-1-1}
\expandafter\def\csname sga@target@IV.16.1.10\endcsname{ega4.c04.cor.16-1-10}
\expandafter\def\csname sga@target@IV.16.1.11\endcsname{ega4.c04.rem.16-1-11}
\expandafter\def\csname sga@target@IV.16.1.2\endcsname{ega4.c04.def.16-1-2}
\expandafter\def\csname sga@target@IV.16.1.2.1\endcsname{ega4.c04.def.16-1-2-1}
\expandafter\def\csname sga@target@IV.16.1.2.2\endcsname{ega4.c04.eq.16-1-2-2}
\expandafter\def\csname sga@target@IV.16.1.3\endcsname{ega4.c04.statement.16-1-3}
\expandafter\def\csname sga@target@IV.16.1.4\endcsname{ega4.c04.statement.16-1-4}
\expandafter\def\csname sga@target@IV.16.1.5\endcsname{ega4.c04.prop.16-1-5}
\expandafter\def\csname sga@target@IV.16.1.6\endcsname{ega4.c04.cor.16-1-6}
\expandafter\def\csname sga@target@IV.16.1.7\endcsname{ega4.c04.cor.16-1-7}
\expandafter\def\csname sga@target@IV.16.1.8\endcsname{ega4.c04.prop.16-1-8}
\expandafter\def\csname sga@target@IV.16.1.9\endcsname{ega4.c04.prop.16-1-9}
\expandafter\def\csname sga@target@IV.16.10\endcsname{ega4.c04.subsec.16-10}
\expandafter\def\csname sga@target@IV.16.10.1\endcsname{ega4.c04.def.16-10-1}
\expandafter\def\csname sga@target@IV.16.10.2\endcsname{ega4.c04.prop.16-10-2}
\expandafter\def\csname sga@target@IV.16.10.3\endcsname{ega4.c04.prop.16-10-3}
\expandafter\def\csname sga@target@IV.16.10.4\endcsname{ega4.c04.prop.16-10-4}
\expandafter\def\csname sga@target@IV.16.10.5\endcsname{ega4.c04.prop.16-10-5}
\expandafter\def\csname sga@target@IV.16.10.6\endcsname{ega4.c04.rem.16-10-6}
\expandafter\def\csname sga@target@IV.16.11\endcsname{ega4.c04.subsec.16-11}
\expandafter\def\csname sga@target@IV.16.11.1\endcsname{ega4.c04.statement.16-11-1}
\expandafter\def\csname sga@target@IV.16.11.1.1\endcsname{ega4.c04.statement.16-11-1-1}
\expandafter\def\csname sga@target@IV.16.11.1.2\endcsname{ega4.c04.statement.16-11-1-2}
\expandafter\def\csname sga@target@IV.16.11.1.3\endcsname{ega4.c04.eq.16-11-1-3}
\expandafter\def\csname sga@target@IV.16.11.1.4\endcsname{ega4.c04.eq.16-11-1-4}
\expandafter\def\csname sga@target@IV.16.11.1.5\endcsname{ega4.c04.eq.16-11-1-5}
\expandafter\def\csname sga@target@IV.16.11.1.6\endcsname{ega4.c04.eq.16-11-1-6}
\expandafter\def\csname sga@target@IV.16.11.1.7\endcsname{ega4.c04.eq.16-11-1-7}
\expandafter\def\csname sga@target@IV.16.11.2\endcsname{ega4.c04.thm.16-11-2}
\expandafter\def\csname sga@target@IV.16.11.2.1\endcsname{ega4.c04.eq.16-11-2-1}
\expandafter\def\csname sga@target@IV.16.11.2.2\endcsname{ega4.c04.eq.16-11-2-2}
\expandafter\def\csname sga@target@IV.16.11.2.3\endcsname{ega4.c04.eq.16-11-2-3}
\expandafter\def\csname sga@target@IV.16.11.3\endcsname{ega4.c04.statement.16-11-3}
\expandafter\def\csname sga@target@IV.16.12\endcsname{ega4.c04.subsec.16-12}
\expandafter\def\csname sga@target@IV.16.12.1\endcsname{ega4.c04.statement.16-12-1}
\expandafter\def\csname sga@target@IV.16.12.2\endcsname{ega4.c04.thm.16-12-2}
\expandafter\def\csname sga@target@IV.16.12.3\endcsname{ega4.c04.statement.16-12-3}
\expandafter\def\csname sga@target@IV.16.2\endcsname{ega4.c04.subsec.16-2}
\expandafter\def\csname sga@target@IV.16.2.1\endcsname{ega4.c04.statement.16-2-1}
\expandafter\def\csname sga@target@IV.16.2.1.1\endcsname{ega4.c04.statement.16-2-1-1}
\expandafter\def\csname sga@target@IV.16.2.1.2\endcsname{ega4.c04.eq.16-2-1-2}
\expandafter\def\csname sga@target@IV.16.2.1.3\endcsname{ega4.c04.eq.16-2-1-3}
\expandafter\def\csname sga@target@IV.16.2.1.4\endcsname{ega4.c04.eq.16-2-1-4}
\expandafter\def\csname sga@target@IV.16.2.2\endcsname{ega4.c04.prop.16-2-2}
\expandafter\def\csname sga@target@IV.16.2.3\endcsname{ega4.c04.cor.16-2-3}
\expandafter\def\csname sga@target@IV.16.2.3.1\endcsname{ega4.c04.eq.16-2-3-1}
\expandafter\def\csname sga@target@IV.16.2.3.2\endcsname{ega4.c04.eq.16-2-3-2}
\expandafter\def\csname sga@target@IV.16.2.4\endcsname{ega4.c04.cor.16-2-4}
\expandafter\def\csname sga@target@IV.16.2.5\endcsname{ega4.c04.statement.16-2-5}
\expandafter\def\csname sga@target@IV.16.2.6\endcsname{ega4.c04.statement.16-2-6}
\expandafter\def\csname sga@target@IV.16.2.6.1\endcsname{ega4.c04.statement.16-2-6-1}
\expandafter\def\csname sga@target@IV.16.2.6.2\endcsname{ega4.c04.eq.16-2-6-2}
\expandafter\def\csname sga@target@IV.16.2.7\endcsname{ega4.c04.prop.16-2-7}
\expandafter\def\csname sga@target@IV.16.2.7.1\endcsname{ega4.c04.prop.16-2-7-1}
\expandafter\def\csname sga@target@IV.16.3\endcsname{ega4.c04.subsec.16-3}
\expandafter\def\csname sga@target@IV.16.3.1\endcsname{ega4.c04.def.16-3-1}
\expandafter\def\csname sga@target@IV.16.3.1.1\endcsname{ega4.c04.eq.16-3-1-1}
\expandafter\def\csname sga@target@IV.16.3.2\endcsname{ega4.c04.statement.16-3-2}
\expandafter\def\csname sga@target@IV.16.3.3\endcsname{ega4.c04.statement.16-3-3}
\expandafter\def\csname sga@target@IV.16.3.3.1\endcsname{ega4.c04.statement.16-3-3-1}
\expandafter\def\csname sga@target@IV.16.3.4\endcsname{ega4.c04.prop.16-3-4}
\expandafter\def\csname sga@target@IV.16.3.5\endcsname{ega4.c04.statement.16-3-5}
\expandafter\def\csname sga@target@IV.16.3.6\endcsname{ega4.c04.def.16-3-6}
\expandafter\def\csname sga@target@IV.16.3.6.1\endcsname{ega4.c04.eq.16-3-6-1}
\expandafter\def\csname sga@target@IV.16.3.7\endcsname{ega4.c04.statement.16-3-7}
\expandafter\def\csname sga@target@IV.16.3.7.1\endcsname{ega4.c04.eq.16-3-7-1}
\expandafter\def\csname sga@target@IV.16.3.8\endcsname{ega4.c04.prop.16-3-8}
\expandafter\def\csname sga@target@IV.16.3.9\endcsname{ega4.c04.prop.16-3-9}
\expandafter\def\csname sga@target@IV.16.4\endcsname{ega4.c04.subsec.16-4}
\expandafter\def\csname sga@target@IV.16.4.1\endcsname{ega4.c04.statement.16-4-1}
\expandafter\def\csname sga@target@IV.16.4.1.1\endcsname{ega4.c04.statement.16-4-1-1}
\expandafter\def\csname sga@target@IV.16.4.1.2\endcsname{ega4.c04.statement.16-4-1-2}
\expandafter\def\csname sga@target@IV.16.4.1.3\endcsname{ega4.c04.eq.16-4-1-3}
\expandafter\def\csname sga@target@IV.16.4.1.4\endcsname{ega4.c04.eq.16-4-1-4}
\expandafter\def\csname sga@target@IV.16.4.1.5\endcsname{ega4.c04.eq.16-4-1-5}
\expandafter\def\csname sga@target@IV.16.4.10\endcsname{ega4.c04.cor.16-4-10}
\expandafter\def\csname sga@target@IV.16.4.11\endcsname{ega4.c04.prop.16-4-11}
\expandafter\def\csname sga@target@IV.16.4.11.1\endcsname{ega4.c04.prop.16-4-11-1}
\expandafter\def\csname sga@target@IV.16.4.11.2\endcsname{ega4.c04.prop.16-4-11-2}
\expandafter\def\csname sga@target@IV.16.4.11.3\endcsname{ega4.c04.eq.16-4-11-3}
\expandafter\def\csname sga@target@IV.16.4.11.4\endcsname{ega4.c04.eq.16-4-11-4}
\expandafter\def\csname sga@target@IV.16.4.12\endcsname{ega4.c04.cor.16-4-12}
\expandafter\def\csname sga@target@IV.16.4.13\endcsname{ega4.c04.cor.16-4-13}
\expandafter\def\csname sga@target@IV.16.4.14\endcsname{ega4.c04.prop.16-4-14}
\expandafter\def\csname sga@target@IV.16.4.14.1\endcsname{ega4.c04.prop.16-4-14-1}
\expandafter\def\csname sga@target@IV.16.4.15\endcsname{ega4.c04.cor.16-4-15}
\expandafter\def\csname sga@target@IV.16.4.15.1\endcsname{ega4.c04.cor.16-4-15-1}
\expandafter\def\csname sga@target@IV.16.4.15.2\endcsname{ega4.c04.eq.16-4-15-2}
\expandafter\def\csname sga@target@IV.16.4.16\endcsname{ega4.c04.cor.16-4-16}
\expandafter\def\csname sga@target@IV.16.4.16.1\endcsname{ega4.c04.cor.16-4-16-1}
\expandafter\def\csname sga@target@IV.16.4.16.2\endcsname{ega4.c04.eq.16-4-16-2}
\expandafter\def\csname sga@target@IV.16.4.17\endcsname{ega4.c04.cor.16-4-17}
\expandafter\def\csname sga@target@IV.16.4.18\endcsname{ega4.c04.statement.16-4-18}
\expandafter\def\csname sga@target@IV.16.4.18.1\endcsname{ega4.c04.statement.16-4-18-1}
\expandafter\def\csname sga@target@IV.16.4.18.2\endcsname{ega4.c04.statement.16-4-18-2}
\expandafter\def\csname sga@target@IV.16.4.18.3\endcsname{ega4.c04.eq.16-4-18-3}
\expandafter\def\csname sga@target@IV.16.4.19\endcsname{ega4.c04.cor.16-4-19}
\expandafter\def\csname sga@target@IV.16.4.19.1\endcsname{ega4.c04.cor.16-4-19-1}
\expandafter\def\csname sga@target@IV.16.4.2\endcsname{ega4.c04.statement.16-4-2}
\expandafter\def\csname sga@target@IV.16.4.20\endcsname{ega4.c04.prop.16-4-20}
\expandafter\def\csname sga@target@IV.16.4.21\endcsname{ega4.c04.cor.16-4-21}
\expandafter\def\csname sga@target@IV.16.4.21.1\endcsname{ega4.c04.cor.16-4-21-1}
\expandafter\def\csname sga@target@IV.16.4.22\endcsname{ega4.c04.cor.16-4-22}
\expandafter\def\csname sga@target@IV.16.4.23\endcsname{ega4.c04.prop.16-4-23}
\expandafter\def\csname sga@target@IV.16.4.23.1\endcsname{ega4.c04.prop.16-4-23-1}
\expandafter\def\csname sga@target@IV.16.4.24\endcsname{ega4.c04.statement.16-4-24}
\expandafter\def\csname sga@target@IV.16.4.3\endcsname{ega4.c04.statement.16-4-3}
\expandafter\def\csname sga@target@IV.16.4.3.1\endcsname{ega4.c04.statement.16-4-3-1}
\expandafter\def\csname sga@target@IV.16.4.3.2\endcsname{ega4.c04.statement.16-4-3-2}
\expandafter\def\csname sga@target@IV.16.4.3.3\endcsname{ega4.c04.eq.16-4-3-3}
\expandafter\def\csname sga@target@IV.16.4.3.4\endcsname{ega4.c04.eq.16-4-3-4}
\expandafter\def\csname sga@target@IV.16.4.3.5\endcsname{ega4.c04.eq.16-4-3-5}
\expandafter\def\csname sga@target@IV.16.4.3.6\endcsname{ega4.c04.eq.16-4-3-6}
\expandafter\def\csname sga@target@IV.16.4.3.7\endcsname{ega4.c04.eq.16-4-3-7}
\expandafter\def\csname sga@target@IV.16.4.4\endcsname{ega4.c04.statement.16-4-4}
\expandafter\def\csname sga@target@IV.16.4.5\endcsname{ega4.c04.prop.16-4-5}
\expandafter\def\csname sga@target@IV.16.4.6\endcsname{ega4.c04.cor.16-4-6}
\expandafter\def\csname sga@target@IV.16.4.7\endcsname{ega4.c04.statement.16-4-7}
\expandafter\def\csname sga@target@IV.16.4.7.1\endcsname{ega4.c04.statement.16-4-7-1}
\expandafter\def\csname sga@target@IV.16.4.7.2\endcsname{ega4.c04.eq.16-4-7-2}
\expandafter\def\csname sga@target@IV.16.4.7.3\endcsname{ega4.c04.eq.16-4-7-3}
\expandafter\def\csname sga@target@IV.16.4.8\endcsname{ega4.c04.prop.16-4-8}
\expandafter\def\csname sga@target@IV.16.4.8.1\endcsname{ega4.c04.prop.16-4-8-1}
\expandafter\def\csname sga@target@IV.16.4.8.2\endcsname{ega4.c04.eq.16-4-8-2}
\expandafter\def\csname sga@target@IV.16.4.8.3\endcsname{ega4.c04.lem.16-4-8-3}
\expandafter\def\csname sga@target@IV.16.4.8.4\endcsname{ega4.c04.eq.16-4-8-4}
\expandafter\def\csname sga@target@IV.16.4.9\endcsname{ega4.c04.cor.16-4-9}
\expandafter\def\csname sga@target@IV.16.4.9.1\endcsname{ega4.c04.cor.16-4-9-1}
\expandafter\def\csname sga@target@IV.16.5\endcsname{ega4.c04.subsec.16-5}
\expandafter\def\csname sga@target@IV.16.5.1\endcsname{ega4.c04.statement.16-5-1}
\expandafter\def\csname sga@target@IV.16.5.1.1\endcsname{ega4.c04.eq.16-5-1-1}
\expandafter\def\csname sga@target@IV.16.5.1.2\endcsname{ega4.c04.eq.16-5-1-2}
\expandafter\def\csname sga@target@IV.16.5.10\endcsname{ega4.c04.statement.16-5-10}
\expandafter\def\csname sga@target@IV.16.5.10.1\endcsname{ega4.c04.statement.16-5-10-1}
\expandafter\def\csname sga@target@IV.16.5.10.2\endcsname{ega4.c04.statement.16-5-10-2}
\expandafter\def\csname sga@target@IV.16.5.11\endcsname{ega4.c04.prop.16-5-11}
\expandafter\def\csname sga@target@IV.16.5.12\endcsname{ega4.c04.statement.16-5-12}
\expandafter\def\csname sga@target@IV.16.5.12.1\endcsname{ega4.c04.statement.16-5-12-1}
\expandafter\def\csname sga@target@IV.16.5.12.2\endcsname{ega4.c04.eq.16-5-12-2}
\expandafter\def\csname sga@target@IV.16.5.13\endcsname{ega4.c04.statement.16-5-13}
\expandafter\def\csname sga@target@IV.16.5.13.1\endcsname{ega4.c04.statement.16-5-13-1}
\expandafter\def\csname sga@target@IV.16.5.13.2\endcsname{ega4.c04.eq.16-5-13-2}
\expandafter\def\csname sga@target@IV.16.5.13.3\endcsname{ega4.c04.eq.16-5-13-3}
\expandafter\def\csname sga@target@IV.16.5.13.4\endcsname{ega4.c04.eq.16-5-13-4}
\expandafter\def\csname sga@target@IV.16.5.13.5\endcsname{ega4.c04.eq.16-5-13-5}
\expandafter\def\csname sga@target@IV.16.5.14\endcsname{ega4.c04.statement.16-5-14}
\expandafter\def\csname sga@target@IV.16.5.14.1\endcsname{ega4.c04.statement.16-5-14-1}
\expandafter\def\csname sga@target@IV.16.5.14.2\endcsname{ega4.c04.eq.16-5-14-2}
\expandafter\def\csname sga@target@IV.16.5.14.3\endcsname{ega4.c04.eq.16-5-14-3}
\expandafter\def\csname sga@target@IV.16.5.15\endcsname{ega4.c04.statement.16-5-15}
\expandafter\def\csname sga@target@IV.16.5.16\endcsname{ega4.c04.prop.16-5-16}
\expandafter\def\csname sga@target@IV.16.5.17\endcsname{ega4.c04.prop.16-5-17}
\expandafter\def\csname sga@target@IV.16.5.18\endcsname{ega4.c04.cor.16-5-18}
\expandafter\def\csname sga@target@IV.16.5.2\endcsname{ega4.c04.prop.16-5-2}
\expandafter\def\csname sga@target@IV.16.5.3\endcsname{ega4.c04.prop.16-5-3}
\expandafter\def\csname sga@target@IV.16.5.3.1\endcsname{ega4.c04.eq.16-5-3-1}
\expandafter\def\csname sga@target@IV.16.5.4\endcsname{ega4.c04.statement.16-5-4}
\expandafter\def\csname sga@target@IV.16.5.5\endcsname{ega4.c04.cor.16-5-5}
\expandafter\def\csname sga@target@IV.16.5.5.1\endcsname{ega4.c04.cor.16-5-5-1}
\expandafter\def\csname sga@target@IV.16.5.6\endcsname{ega4.c04.cor.16-5-6}
\expandafter\def\csname sga@target@IV.16.5.7\endcsname{ega4.c04.statement.16-5-7}
\expandafter\def\csname sga@target@IV.16.5.7.1\endcsname{ega4.c04.statement.16-5-7-1}
\expandafter\def\csname sga@target@IV.16.5.8\endcsname{ega4.c04.statement.16-5-8}
\expandafter\def\csname sga@target@IV.16.5.8.1\endcsname{ega4.c04.eq.16-5-8-1}
\expandafter\def\csname sga@target@IV.16.5.8.2\endcsname{ega4.c04.eq.16-5-8-2}
\expandafter\def\csname sga@target@IV.16.5.9\endcsname{ega4.c04.statement.16-5-9}
\expandafter\def\csname sga@target@IV.16.5.9.1\endcsname{ega4.c04.statement.16-5-9-1}
\expandafter\def\csname sga@target@IV.16.6\endcsname{ega4.c04.subsec.16-6}
\expandafter\def\csname sga@target@IV.16.6.1\endcsname{ega4.c04.statement.16-6-1}
\expandafter\def\csname sga@target@IV.16.6.1.1\endcsname{ega4.c04.statement.16-6-1-1}
\expandafter\def\csname sga@target@IV.16.6.1.2\endcsname{ega4.c04.statement.16-6-1-2}
\expandafter\def\csname sga@target@IV.16.6.2\endcsname{ega4.c04.thm.16-6-2}
\expandafter\def\csname sga@target@IV.16.6.2.1\endcsname{ega4.c04.eq.16-6-2-1}
\expandafter\def\csname sga@target@IV.16.6.2.2\endcsname{ega4.c04.eq.16-6-2-2}
\expandafter\def\csname sga@target@IV.16.6.2.3\endcsname{ega4.c04.eq.16-6-2-3}
\expandafter\def\csname sga@target@IV.16.6.2.4\endcsname{ega4.c04.eq.16-6-2-4}
\expandafter\def\csname sga@target@IV.16.6.2.5\endcsname{ega4.c04.eq.16-6-2-5}
\expandafter\def\csname sga@target@IV.16.6.3\endcsname{ega4.c04.def.16-6-3}
\expandafter\def\csname sga@target@IV.16.6.4\endcsname{ega4.c04.prop.16-6-4}
\expandafter\def\csname sga@target@IV.16.6.4.1\endcsname{ega4.c04.prop.16-6-4-1}
\expandafter\def\csname sga@target@IV.16.6.5\endcsname{ega4.c04.statement.16-6-5}
\expandafter\def\csname sga@target@IV.16.6.5.1\endcsname{ega4.c04.statement.16-6-5-1}
\expandafter\def\csname sga@target@IV.16.7\endcsname{ega4.c04.subsec.16-7}
\expandafter\def\csname sga@target@IV.16.7.1\endcsname{ega4.c04.statement.16-7-1}
\expandafter\def\csname sga@target@IV.16.7.1.1\endcsname{ega4.c04.eq.16-7-1-1}
\expandafter\def\csname sga@target@IV.16.7.1.2\endcsname{ega4.c04.eq.16-7-1-2}
\expandafter\def\csname sga@target@IV.16.7.10\endcsname{ega4.c04.rem.16-7-10}
\expandafter\def\csname sga@target@IV.16.7.2\endcsname{ega4.c04.statement.16-7-2}
\expandafter\def\csname sga@target@IV.16.7.2.1\endcsname{ega4.c04.statement.16-7-2-1}
\expandafter\def\csname sga@target@IV.16.7.3\endcsname{ega4.c04.prop.16-7-3}
\expandafter\def\csname sga@target@IV.16.7.4\endcsname{ega4.c04.statement.16-7-4}
\expandafter\def\csname sga@target@IV.16.7.4.1\endcsname{ega4.c04.statement.16-7-4-1}
\expandafter\def\csname sga@target@IV.16.7.5\endcsname{ega4.c04.statement.16-7-5}
\expandafter\def\csname sga@target@IV.16.7.5.1\endcsname{ega4.c04.statement.16-7-5-1}
\expandafter\def\csname sga@target@IV.16.7.5.2\endcsname{ega4.c04.statement.16-7-5-2}
\expandafter\def\csname sga@target@IV.16.7.5.3\endcsname{ega4.c04.statement.16-7-5-3}
\expandafter\def\csname sga@target@IV.16.7.5.4\endcsname{ega4.c04.statement.16-7-5-4}
\expandafter\def\csname sga@target@IV.16.7.6\endcsname{ega4.c04.prop.16-7-6}
\expandafter\def\csname sga@target@IV.16.7.7\endcsname{ega4.c04.statement.16-7-7}
\expandafter\def\csname sga@target@IV.16.7.7.1\endcsname{ega4.c04.eq.16-7-7-1}
\expandafter\def\csname sga@target@IV.16.7.7.2\endcsname{ega4.c04.eq.16-7-7-2}
\expandafter\def\csname sga@target@IV.16.7.8\endcsname{ega4.c04.statement.16-7-8}
\expandafter\def\csname sga@target@IV.16.7.8.1\endcsname{ega4.c04.statement.16-7-8-1}
\expandafter\def\csname sga@target@IV.16.7.8.2\endcsname{ega4.c04.eq.16-7-8-2}
\expandafter\def\csname sga@target@IV.16.7.9\endcsname{ega4.c04.statement.16-7-9}
\expandafter\def\csname sga@target@IV.16.7.9.1\endcsname{ega4.c04.statement.16-7-9-1}
\expandafter\def\csname sga@target@IV.16.8\endcsname{ega4.c04.subsec.16-8}
\expandafter\def\csname sga@target@IV.16.8.1\endcsname{ega4.c04.def.16-8-1}
\expandafter\def\csname sga@target@IV.16.8.10\endcsname{ega4.c04.cor.16-8-10}
\expandafter\def\csname sga@target@IV.16.8.11\endcsname{ega4.c04.statement.16-8-11}
\expandafter\def\csname sga@target@IV.16.8.2\endcsname{ega4.c04.statement.16-8-2}
\expandafter\def\csname sga@target@IV.16.8.3\endcsname{ega4.c04.statement.16-8-3}
\expandafter\def\csname sga@target@IV.16.8.3.1\endcsname{ega4.c04.eq.16-8-3-1}
\expandafter\def\csname sga@target@IV.16.8.4\endcsname{ega4.c04.prop.16-8-4}
\expandafter\def\csname sga@target@IV.16.8.4.1\endcsname{ega4.c04.prop.16-8-4-1}
\expandafter\def\csname sga@target@IV.16.8.5\endcsname{ega4.c04.statement.16-8-5}
\expandafter\def\csname sga@target@IV.16.8.5.1\endcsname{ega4.c04.eq.16-8-5-1}
\expandafter\def\csname sga@target@IV.16.8.5.2\endcsname{ega4.c04.eq.16-8-5-2}
\expandafter\def\csname sga@target@IV.16.8.6\endcsname{ega4.c04.prop.16-8-6}
\expandafter\def\csname sga@target@IV.16.8.7\endcsname{ega4.c04.statement.16-8-7}
\expandafter\def\csname sga@target@IV.16.8.7.1\endcsname{ega4.c04.eq.16-8-7-1}
\expandafter\def\csname sga@target@IV.16.8.8\endcsname{ega4.c04.prop.16-8-8}
\expandafter\def\csname sga@target@IV.16.8.8.1\endcsname{ega4.c04.eq.16-8-8-1}
\expandafter\def\csname sga@target@IV.16.8.8.2\endcsname{ega4.c04.eq.16-8-8-2}
\expandafter\def\csname sga@target@IV.16.8.9\endcsname{ega4.c04.prop.16-8-9}
\expandafter\def\csname sga@target@IV.16.8.9.1\endcsname{ega4.c04.lem.16-8-9-1}
\expandafter\def\csname sga@target@IV.16.8.9.2\endcsname{ega4.c04.lem.16-8-9-2}
\expandafter\def\csname sga@target@IV.16.8.9.3\endcsname{ega4.c04.lem.16-8-9-3}
\expandafter\def\csname sga@target@IV.16.8.9.4\endcsname{ega4.c04.eq.16-8-9-4}
\expandafter\def\csname sga@target@IV.16.9\endcsname{ega4.c04.subsec.16-9}
\expandafter\def\csname sga@target@IV.16.9.1\endcsname{ega4.c04.def.16-9-1}
\expandafter\def\csname sga@target@IV.16.9.10\endcsname{ega4.c04.prop.16-9-10}
\expandafter\def\csname sga@target@IV.16.9.11\endcsname{ega4.c04.cor.16-9-11}
\expandafter\def\csname sga@target@IV.16.9.12\endcsname{ega4.c04.statement.16-9-12}
\expandafter\def\csname sga@target@IV.16.9.13\endcsname{ega4.c04.prop.16-9-13}
\expandafter\def\csname sga@target@IV.16.9.13.1\endcsname{ega4.c04.prop.16-9-13-1}
\expandafter\def\csname sga@target@IV.16.9.13.2\endcsname{ega4.c04.lem.16-9-13-2}
\expandafter\def\csname sga@target@IV.16.9.13.3\endcsname{ega4.c04.lem.16-9-13-3}
\expandafter\def\csname sga@target@IV.16.9.13.4\endcsname{ega4.c04.eq.16-9-13-4}
\expandafter\def\csname sga@target@IV.16.9.13.5\endcsname{ega4.c04.cor.16-9-13-5}
\expandafter\def\csname sga@target@IV.16.9.13.6\endcsname{ega4.c04.cor.16-9-13-6}
\expandafter\def\csname sga@target@IV.16.9.13.7\endcsname{ega4.c04.eq.16-9-13-7}
\expandafter\def\csname sga@target@IV.16.9.2\endcsname{ega4.c04.def.16-9-2}
\expandafter\def\csname sga@target@IV.16.9.3\endcsname{ega4.c04.prop.16-9-3}
\expandafter\def\csname sga@target@IV.16.9.4\endcsname{ega4.c04.cor.16-9-4}
\expandafter\def\csname sga@target@IV.16.9.4.1\endcsname{ega4.c04.eq.16-9-4-1}
\expandafter\def\csname sga@target@IV.16.9.5\endcsname{ega4.c04.cor.16-9-5}
\expandafter\def\csname sga@target@IV.16.9.6\endcsname{ega4.c04.statement.16-9-6}
\expandafter\def\csname sga@target@IV.16.9.7\endcsname{ega4.c04.statement.16-9-7}
\expandafter\def\csname sga@target@IV.16.9.8\endcsname{ega4.c04.prop.16-9-8}
\expandafter\def\csname sga@target@IV.16.9.9\endcsname{ega4.c04.cor.16-9-9}
\expandafter\def\csname sga@target@IV.17.1.1\endcsname{ega4.c04.def.17-1-1}
\expandafter\def\csname sga@target@IV.17.1.1.1\endcsname{ega4.c04.def.17-1-1-1}
\expandafter\def\csname sga@target@IV.17.1.2\endcsname{ega4.c04.rem.17-1-2}
\expandafter\def\csname sga@target@IV.17.1.2.1\endcsname{ega4.c04.eq.17-1-2-1}
\expandafter\def\csname sga@target@IV.17.1.3\endcsname{ega4.c04.prop.17-1-3}
\expandafter\def\csname sga@target@IV.17.1.4\endcsname{ega4.c04.prop.17-1-4}
\expandafter\def\csname sga@target@IV.17.1.5\endcsname{ega4.c04.cor.17-1-5}
\expandafter\def\csname sga@target@IV.17.1.6\endcsname{ega4.c04.prop.17-1-6}
\expandafter\def\csname sga@target@IV.17.10.1\endcsname{ega4.c04.def.17-10-1}
\expandafter\def\csname sga@target@IV.17.10.2\endcsname{ega4.c04.prop.17-10-2}
\expandafter\def\csname sga@target@IV.17.10.3\endcsname{ega4.c04.cor.17-10-3}
\expandafter\def\csname sga@target@IV.17.10.3.1\endcsname{ega4.c04.cor.17-10-3-1}
\expandafter\def\csname sga@target@IV.17.10.4\endcsname{ega4.c04.cor.17-10-4}
\expandafter\def\csname sga@target@IV.17.10.4.1\endcsname{ega4.c04.eq.17-10-4-1}
\expandafter\def\csname sga@target@IV.17.11.1\endcsname{ega4.c04.thm.17-11-1}
\expandafter\def\csname sga@target@IV.17.11.1.1\endcsname{ega4.c04.eq.17-11-1-1}
\expandafter\def\csname sga@target@IV.17.11.1.2\endcsname{ega4.c04.eq.17-11-1-2}
\expandafter\def\csname sga@target@IV.17.11.2\endcsname{ega4.c04.cor.17-11-2}
\expandafter\def\csname sga@target@IV.17.11.3\endcsname{ega4.c04.cor.17-11-3}
\expandafter\def\csname sga@target@IV.17.11.4\endcsname{ega4.c04.cor.17-11-4}
\expandafter\def\csname sga@target@IV.17.11.5\endcsname{ega4.c04.prop.17-11-5}
\expandafter\def\csname sga@target@IV.17.12.1\endcsname{ega4.c04.thm.17-12-1}
\expandafter\def\csname sga@target@IV.17.12.2\endcsname{ega4.c04.cor.17-12-2}
\expandafter\def\csname sga@target@IV.17.12.3\endcsname{ega4.c04.cor.17-12-3}
\expandafter\def\csname sga@target@IV.17.12.4\endcsname{ega4.c04.prop.17-12-4}
\expandafter\def\csname sga@target@IV.17.12.5\endcsname{ega4.c04.prop.17-12-5}
\expandafter\def\csname sga@target@IV.17.12.6\endcsname{ega4.c04.cor.17-12-6}
\expandafter\def\csname sga@target@IV.17.12.7\endcsname{ega4.c04.ex.17-12-7}
\expandafter\def\csname sga@target@IV.17.12.8\endcsname{ega4.c04.rem.17-12-8}
\expandafter\def\csname sga@target@IV.17.13\endcsname{ega4.c04.subsec.17-13.key-iv-17-13}
\expandafter\def\csname sga@target@IV.17.13.1\endcsname{ega4.c04.statement.17-13-1}
\expandafter\def\csname sga@target@IV.17.13.1.1\endcsname{ega4.c04.statement.17-13-1-1}
\expandafter\def\csname sga@target@IV.17.13.1.2\endcsname{ega4.c04.statement.17-13-1-2}
\expandafter\def\csname sga@target@IV.17.13.1.3\endcsname{ega4.c04.eq.17-13-1-3}
\expandafter\def\csname sga@target@IV.17.13.10\endcsname{ega4.c04.cor.17-13-10}
\expandafter\def\csname sga@target@IV.17.13.11\endcsname{ega4.c04.statement.17-13-11}
\expandafter\def\csname sga@target@IV.17.13.11.1\endcsname{ega4.c04.statement.17-13-11-1}
\expandafter\def\csname sga@target@IV.17.13.12\endcsname{ega4.c04.cor.17-13-12}
\expandafter\def\csname sga@target@IV.17.13.12.1\endcsname{ega4.c04.cor.17-13-12-1}
\expandafter\def\csname sga@target@IV.17.13.12.2\endcsname{ega4.c04.eq.17-13-12-2}
\expandafter\def\csname sga@target@IV.17.13.2\endcsname{ega4.c04.prop.17-13-2}
\expandafter\def\csname sga@target@IV.17.13.2.1\endcsname{ega4.c04.eq.17-13-2-1}
\expandafter\def\csname sga@target@IV.17.13.2.2\endcsname{ega4.c04.eq.17-13-2-2}
\expandafter\def\csname sga@target@IV.17.13.3\endcsname{ega4.c04.def.17-13-3}
\expandafter\def\csname sga@target@IV.17.13.4\endcsname{ega4.c04.statement.17-13-4}
\expandafter\def\csname sga@target@IV.17.13.5\endcsname{ega4.c04.statement.17-13-5}
\expandafter\def\csname sga@target@IV.17.13.5.1\endcsname{ega4.c04.statement.17-13-5-1}
\expandafter\def\csname sga@target@IV.17.13.5.2\endcsname{ega4.c04.eq.17-13-5-2}
\expandafter\def\csname sga@target@IV.17.13.5.3\endcsname{ega4.c04.eq.17-13-5-3}
\expandafter\def\csname sga@target@IV.17.13.5.4\endcsname{ega4.c04.eq.17-13-5-4}
\expandafter\def\csname sga@target@IV.17.13.6\endcsname{ega4.c04.cor.17-13-6}
\expandafter\def\csname sga@target@IV.17.13.6.1\endcsname{ega4.c04.eq.17-13-6-1}
\expandafter\def\csname sga@target@IV.17.13.6.2\endcsname{ega4.c04.eq.17-13-6-2}
\expandafter\def\csname sga@target@IV.17.13.7\endcsname{ega4.c04.statement.17-13-7}
\expandafter\def\csname sga@target@IV.17.13.7.1\endcsname{ega4.c04.eq.17-13-7-1}
\expandafter\def\csname sga@target@IV.17.13.8\endcsname{ega4.c04.prop.17-13-8}
\expandafter\def\csname sga@target@IV.17.13.8.1\endcsname{ega4.c04.eq.17-13-8-1}
\expandafter\def\csname sga@target@IV.17.13.9\endcsname{ega4.c04.statement.17-13-9}
\expandafter\def\csname sga@target@IV.17.13.9.1\endcsname{ega4.c04.eq.17-13-9-1}
\expandafter\def\csname sga@target@IV.17.13.9.2\endcsname{ega4.c04.eq.17-13-9-2}
\expandafter\def\csname sga@target@IV.17.14.1\endcsname{ega4.c04.prop.17-14-1}
\expandafter\def\csname sga@target@IV.17.14.2\endcsname{ega4.c04.prop.17-14-2}
\expandafter\def\csname sga@target@IV.17.14.3\endcsname{ega4.c04.rem.17-14-3}
\expandafter\def\csname sga@target@IV.17.15.1\endcsname{ega4.c04.prop.17-15-1}
\expandafter\def\csname sga@target@IV.17.15.10\endcsname{ega4.c04.statement.17-15-10}
\expandafter\def\csname sga@target@IV.17.15.11\endcsname{ega4.c04.prop.17-15-11}
\expandafter\def\csname sga@target@IV.17.15.12\endcsname{ega4.c04.prop.17-15-12}
\expandafter\def\csname sga@target@IV.17.15.13\endcsname{ega4.c04.cor.17-15-13}
\expandafter\def\csname sga@target@IV.17.15.14\endcsname{ega4.c04.prop.17-15-14}
\expandafter\def\csname sga@target@IV.17.15.14.1\endcsname{ega4.c04.lem.17-15-14-1}
\expandafter\def\csname sga@target@IV.17.15.14.2\endcsname{ega4.c04.lem.17-15-14-2}
\expandafter\def\csname sga@target@IV.17.15.15\endcsname{ega4.c04.prop.17-15-15}
\expandafter\def\csname sga@target@IV.17.15.2\endcsname{ega4.c04.cor.17-15-2}
\expandafter\def\csname sga@target@IV.17.15.3\endcsname{ega4.c04.prop.17-15-3}
\expandafter\def\csname sga@target@IV.17.15.4\endcsname{ega4.c04.cor.17-15-4}
\expandafter\def\csname sga@target@IV.17.15.5\endcsname{ega4.c04.prop.17-15-5}
\expandafter\def\csname sga@target@IV.17.15.6\endcsname{ega4.c04.cor.17-15-6}
\expandafter\def\csname sga@target@IV.17.15.7\endcsname{ega4.c04.cor.17-15-7}
\expandafter\def\csname sga@target@IV.17.15.8\endcsname{ega4.c04.prop.17-15-8}
\expandafter\def\csname sga@target@IV.17.15.8.1\endcsname{ega4.c04.eq.17-15-8-1}
\expandafter\def\csname sga@target@IV.17.15.8.2\endcsname{ega4.c04.eq.17-15-8-2}
\expandafter\def\csname sga@target@IV.17.15.9\endcsname{ega4.c04.cor.17-15-9}
\expandafter\def\csname sga@target@IV.17.16\endcsname{ega4.c04.subsec.17-16.key-iv-17-16}
\expandafter\def\csname sga@target@IV.17.16.1\endcsname{ega4.c04.prop.17-16-1}
\expandafter\def\csname sga@target@IV.17.16.2\endcsname{ega4.c04.cor.17-16-2}
\expandafter\def\csname sga@target@IV.17.16.3\endcsname{ega4.c04.cor.17-16-3}
\expandafter\def\csname sga@target@IV.17.16.4\endcsname{ega4.c04.prop.17-16-4}
\expandafter\def\csname sga@target@IV.17.16.5\endcsname{ega4.c04.cor.17-16-5}
\expandafter\def\csname sga@target@IV.17.16.6\endcsname{ega4.c04.prop.17-16-6}
\expandafter\def\csname sga@target@IV.17.2.1\endcsname{ega4.c04.prop.17-2-1}
\expandafter\def\csname sga@target@IV.17.2.2\endcsname{ega4.c04.cor.17-2-2}
\expandafter\def\csname sga@target@IV.17.2.3\endcsname{ega4.c04.prop.17-2-3}
\expandafter\def\csname sga@target@IV.17.2.3.1\endcsname{ega4.c04.eq.17-2-3-1}
\expandafter\def\csname sga@target@IV.17.2.4\endcsname{ega4.c04.cor.17-2-4}
\expandafter\def\csname sga@target@IV.17.2.5\endcsname{ega4.c04.prop.17-2-5}
\expandafter\def\csname sga@target@IV.17.2.5.1\endcsname{ega4.c04.prop.17-2-5-1}
\expandafter\def\csname sga@target@IV.17.2.6\endcsname{ega4.c04.prop.17-2-6}
\expandafter\def\csname sga@target@IV.17.3.1\endcsname{ega4.c04.def.17-3-1}
\expandafter\def\csname sga@target@IV.17.3.2\endcsname{ega4.c04.subsec.17-3-2}
\expandafter\def\csname sga@target@IV.17.3.3\endcsname{ega4.c04.prop.17-3-3}
\expandafter\def\csname sga@target@IV.17.3.4\endcsname{ega4.c04.prop.17-3-4}
\expandafter\def\csname sga@target@IV.17.3.5\endcsname{ega4.c04.cor.17-3-5}
\expandafter\def\csname sga@target@IV.17.3.6\endcsname{ega4.c04.prop.17-3-6}
\expandafter\def\csname sga@target@IV.17.3.7\endcsname{ega4.c04.def.17-3-7}
\expandafter\def\csname sga@target@IV.17.3.8\endcsname{ega4.c04.prop.17-3-8}
\expandafter\def\csname sga@target@IV.17.3.9\endcsname{ega4.c04.cor.17-3-9}
\expandafter\def\csname sga@target@IV.17.4.1\endcsname{ega4.c04.thm.17-4-1}
\expandafter\def\csname sga@target@IV.17.4.1.1\endcsname{ega4.c04.eq.17-4-1-1}
\expandafter\def\csname sga@target@IV.17.4.1.2\endcsname{ega4.c04.rem.17-4-1-2}
\expandafter\def\csname sga@target@IV.17.4.10\endcsname{ega4.c04.rem.17-4-10}
\expandafter\def\csname sga@target@IV.17.4.2\endcsname{ega4.c04.cor.17-4-2}
\expandafter\def\csname sga@target@IV.17.4.3\endcsname{ega4.c04.cor.17-4-3}
\expandafter\def\csname sga@target@IV.17.4.4\endcsname{ega4.c04.prop.17-4-4}
\expandafter\def\csname sga@target@IV.17.4.5\endcsname{ega4.c04.statement.17-4-5}
\expandafter\def\csname sga@target@IV.17.4.6\endcsname{ega4.c04.prop.17-4-6}
\expandafter\def\csname sga@target@IV.17.4.7\endcsname{ega4.c04.cor.17-4-7}
\expandafter\def\csname sga@target@IV.17.4.8\endcsname{ega4.c04.cor.17-4-8}
\expandafter\def\csname sga@target@IV.17.4.9\endcsname{ega4.c04.prop.17-4-9}
\expandafter\def\csname sga@target@IV.17.5.1\endcsname{ega4.c04.thm.17-5-1}
\expandafter\def\csname sga@target@IV.17.5.1.1\endcsname{ega4.c04.eq.17-5-1-1}
\expandafter\def\csname sga@target@IV.17.5.1.2\endcsname{ega4.c04.eq.17-5-1-2}
\expandafter\def\csname sga@target@IV.17.5.2\endcsname{ega4.c04.cor.17-5-2}
\expandafter\def\csname sga@target@IV.17.5.3\endcsname{ega4.c04.prop.17-5-3}
\expandafter\def\csname sga@target@IV.17.5.4\endcsname{ega4.c04.rem.17-5-4}
\expandafter\def\csname sga@target@IV.17.5.5\endcsname{ega4.c04.prop.17-5-5}
\expandafter\def\csname sga@target@IV.17.5.6\endcsname{ega4.c04.cor.17-5-6}
\expandafter\def\csname sga@target@IV.17.5.7\endcsname{ega4.c04.prop.17-5-7}
\expandafter\def\csname sga@target@IV.17.5.8\endcsname{ega4.c04.prop.17-5-8}
\expandafter\def\csname sga@target@IV.17.6.1\endcsname{ega4.c04.thm.17-6-1}
\expandafter\def\csname sga@target@IV.17.6.2\endcsname{ega4.c04.cor.17-6-2}
\expandafter\def\csname sga@target@IV.17.6.3\endcsname{ega4.c04.prop.17-6-3}
\expandafter\def\csname sga@target@IV.17.6.4\endcsname{ega4.c04.prop.17-6-4}
\expandafter\def\csname sga@target@IV.17.7\endcsname{ega4.c04.subsec.17-7.key-iv-17-7}
\expandafter\def\csname sga@target@IV.17.7.1\endcsname{ega4.c04.prop.17-7-1}
\expandafter\def\csname sga@target@IV.17.7.10\endcsname{ega4.c04.prop.17-7-10}
\expandafter\def\csname sga@target@IV.17.7.11\endcsname{ega4.c04.prop.17-7-11}
\expandafter\def\csname sga@target@IV.17.7.11.1\endcsname{ega4.c04.lem.17-7-11-1}
\expandafter\def\csname sga@target@IV.17.7.2\endcsname{ega4.c04.cor.17-7-2}
\expandafter\def\csname sga@target@IV.17.7.3\endcsname{ega4.c04.cor.17-7-3}
\expandafter\def\csname sga@target@IV.17.7.4\endcsname{ega4.c04.prop.17-7-4}
\expandafter\def\csname sga@target@IV.17.7.5\endcsname{ega4.c04.lem.17-7-5}
\expandafter\def\csname sga@target@IV.17.7.6\endcsname{ega4.c04.rem.17-7-6}
\expandafter\def\csname sga@target@IV.17.7.7\endcsname{ega4.c04.prop.17-7-7}
\expandafter\def\csname sga@target@IV.17.7.8\endcsname{ega4.c04.prop.17-7-8}
\expandafter\def\csname sga@target@IV.17.7.9\endcsname{ega4.c04.cor.17-7-9}
\expandafter\def\csname sga@target@IV.17.8\endcsname{ega4.c04.subsec.17-8.key-iv-17-8}
\expandafter\def\csname sga@target@IV.17.8.1\endcsname{ega4.c04.prop.17-8-1}
\expandafter\def\csname sga@target@IV.17.8.2\endcsname{ega4.c04.prop.17-8-2}
\expandafter\def\csname sga@target@IV.17.8.3\endcsname{ega4.c04.rem.17-8-3}
\expandafter\def\csname sga@target@IV.17.9\endcsname{ega4.c04.subsec.17-9.key-iv-17-9}
\expandafter\def\csname sga@target@IV.17.9.1\endcsname{ega4.c04.thm.17-9-1}
\expandafter\def\csname sga@target@IV.17.9.2\endcsname{ega4.c04.cor.17-9-2}
\expandafter\def\csname sga@target@IV.17.9.3\endcsname{ega4.c04.cor.17-9-3}
\expandafter\def\csname sga@target@IV.17.9.4\endcsname{ega4.c04.cor.17-9-4}
\expandafter\def\csname sga@target@IV.17.9.5\endcsname{ega4.c04.cor.17-9-5}
\expandafter\def\csname sga@target@IV.17.9.6\endcsname{ega4.c04.prop.17-9-6}
\expandafter\def\csname sga@target@IV.18\endcsname{ega4.c04.sec.18.key-iv-18}
\expandafter\def\csname sga@target@IV.18.1.1\endcsname{ega4.c04.prop.18-1-1}
\expandafter\def\csname sga@target@IV.18.1.1.1\endcsname{ega4.c04.eq.18-1-1-1}
\expandafter\def\csname sga@target@IV.18.1.2\endcsname{ega4.c04.thm.18-1-2}
\expandafter\def\csname sga@target@IV.18.1.3\endcsname{ega4.c04.cor.18-1-3}
\expandafter\def\csname sga@target@IV.18.10.1\endcsname{ega4.c04.thm.18-10-1}
\expandafter\def\csname sga@target@IV.18.10.10\endcsname{ega4.c04.statement.18-10-10}
\expandafter\def\csname sga@target@IV.18.10.11\endcsname{ega4.c04.rem.18-10-11}
\expandafter\def\csname sga@target@IV.18.10.12\endcsname{ega4.c04.cor.18-10-12}
\expandafter\def\csname sga@target@IV.18.10.13\endcsname{ega4.c04.prop.18-10-13}
\expandafter\def\csname sga@target@IV.18.10.14\endcsname{ega4.c04.cor.18-10-14}
\expandafter\def\csname sga@target@IV.18.10.15\endcsname{ega4.c04.statement.18-10-15}
\expandafter\def\csname sga@target@IV.18.10.16\endcsname{ega4.c04.thm.18-10-16}
\expandafter\def\csname sga@target@IV.18.10.17\endcsname{ega4.c04.rem.18-10-17}
\expandafter\def\csname sga@target@IV.18.10.17.1\endcsname{ega4.c04.statement.18-10-17-1}
\expandafter\def\csname sga@target@IV.18.10.17.2\endcsname{ega4.c04.statement.18-10-17-2}
\expandafter\def\csname sga@target@IV.18.10.18\endcsname{ega4.c04.lem.18-10-18}
\expandafter\def\csname sga@target@IV.18.10.19\endcsname{ega4.c04.prop.18-10-19}
\expandafter\def\csname sga@target@IV.18.10.2\endcsname{ega4.c04.rem.18-10-2}
\expandafter\def\csname sga@target@IV.18.10.20\endcsname{ega4.c04.rem.18-10-20}
\expandafter\def\csname sga@target@IV.18.10.21\endcsname{ega4.c04.prop.18-10-21}
\expandafter\def\csname sga@target@IV.18.10.3\endcsname{ega4.c04.cor.18-10-3}
\expandafter\def\csname sga@target@IV.18.10.3.1\endcsname{ega4.c04.rem.18-10-3-1}
\expandafter\def\csname sga@target@IV.18.10.4\endcsname{ega4.c04.cor.18-10-4}
\expandafter\def\csname sga@target@IV.18.10.5\endcsname{ega4.c04.cor.18-10-5}
\expandafter\def\csname sga@target@IV.18.10.6\endcsname{ega4.c04.lem.18-10-6}
\expandafter\def\csname sga@target@IV.18.10.7\endcsname{ega4.c04.prop.18-10-7}
\expandafter\def\csname sga@target@IV.18.10.8\endcsname{ega4.c04.prop.18-10-8}
\expandafter\def\csname sga@target@IV.18.10.9\endcsname{ega4.c04.cor.18-10-9}
\expandafter\def\csname sga@target@IV.18.11.1\endcsname{ega4.c04.lem.18-11-1}
\expandafter\def\csname sga@target@IV.18.11.10\endcsname{ega4.c04.thm.18-11-10}
\expandafter\def\csname sga@target@IV.18.11.11\endcsname{ega4.c04.rem.18-11-11}
\expandafter\def\csname sga@target@IV.18.11.12\endcsname{ega4.c04.cor.18-11-12}
\expandafter\def\csname sga@target@IV.18.11.12.1\endcsname{ega4.c04.lem.18-11-12-1}
\expandafter\def\csname sga@target@IV.18.11.2\endcsname{ega4.c04.lem.18-11-2}
\expandafter\def\csname sga@target@IV.18.11.2.1\endcsname{ega4.c04.eq.18-11-2-1}
\expandafter\def\csname sga@target@IV.18.11.2.2\endcsname{ega4.c04.eq.18-11-2-2}
\expandafter\def\csname sga@target@IV.18.11.3\endcsname{ega4.c04.lem.18-11-3}
\expandafter\def\csname sga@target@IV.18.11.3.1\endcsname{ega4.c04.eq.18-11-3-1}
\expandafter\def\csname sga@target@IV.18.11.3.2\endcsname{ega4.c04.eq.18-11-3-2}
\expandafter\def\csname sga@target@IV.18.11.3.3\endcsname{ega4.c04.lem.18-11-3-3}
\expandafter\def\csname sga@target@IV.18.11.3.4\endcsname{ega4.c04.lem.18-11-3-4}
\expandafter\def\csname sga@target@IV.18.11.3.5\endcsname{ega4.c04.lem.18-11-3-5}
\expandafter\def\csname sga@target@IV.18.11.3.6\endcsname{ega4.c04.eq.18-11-3-6}
\expandafter\def\csname sga@target@IV.18.11.4\endcsname{ega4.c04.lem.18-11-4}
\expandafter\def\csname sga@target@IV.18.11.5\endcsname{ega4.c04.prop.18-11-5}
\expandafter\def\csname sga@target@IV.18.11.5.1\endcsname{ega4.c04.lem.18-11-5-1}
\expandafter\def\csname sga@target@IV.18.11.5.2\endcsname{ega4.c04.lem.18-11-5-2}
\expandafter\def\csname sga@target@IV.18.11.6\endcsname{ega4.c04.rem.18-11-6}
\expandafter\def\csname sga@target@IV.18.11.7\endcsname{ega4.c04.cor.18-11-7}
\expandafter\def\csname sga@target@IV.18.11.8\endcsname{ega4.c04.cor.18-11-8}
\expandafter\def\csname sga@target@IV.18.11.9\endcsname{ega4.c04.cor.18-11-9}
\expandafter\def\csname sga@target@IV.18.12.1\endcsname{ega4.c04.thm.18-12-1}
\expandafter\def\csname sga@target@IV.18.12.10\endcsname{ega4.c04.cor.18-12-10}
\expandafter\def\csname sga@target@IV.18.12.11\endcsname{ega4.c04.cor.18-12-11}
\expandafter\def\csname sga@target@IV.18.12.12\endcsname{ega4.c04.prop.18-12-12}
\expandafter\def\csname sga@target@IV.18.12.13\endcsname{ega4.c04.cor.18-12-13}
\expandafter\def\csname sga@target@IV.18.12.14\endcsname{ega4.c04.rem.18-12-14}
\expandafter\def\csname sga@target@IV.18.12.15\endcsname{ega4.c04.prop.18-12-15}
\expandafter\def\csname sga@target@IV.18.12.16\endcsname{ega4.c04.prop.18-12-16}
\expandafter\def\csname sga@target@IV.18.12.17\endcsname{ega4.c04.cor.18-12-17}
\expandafter\def\csname sga@target@IV.18.12.2\endcsname{ega4.c04.rem.18-12-2}
\expandafter\def\csname sga@target@IV.18.12.3\endcsname{ega4.c04.cor.18-12-3}
\expandafter\def\csname sga@target@IV.18.12.4\endcsname{ega4.c04.cor.18-12-4}
\expandafter\def\csname sga@target@IV.18.12.5\endcsname{ega4.c04.rem.18-12-5}
\expandafter\def\csname sga@target@IV.18.12.6\endcsname{ega4.c04.cor.18-12-6}
\expandafter\def\csname sga@target@IV.18.12.7\endcsname{ega4.c04.prop.18-12-7}
\expandafter\def\csname sga@target@IV.18.12.7.1\endcsname{ega4.c04.lem.18-12-7-1}
\expandafter\def\csname sga@target@IV.18.12.8\endcsname{ega4.c04.prop.18-12-8}
\expandafter\def\csname sga@target@IV.18.12.9\endcsname{ega4.c04.rem.18-12-9}
\expandafter\def\csname sga@target@IV.18.2.1\endcsname{ega4.c04.statement.18-2-1}
\expandafter\def\csname sga@target@IV.18.2.1.1\endcsname{ega4.c04.statement.18-2-1-1}
\expandafter\def\csname sga@target@IV.18.2.1.2\endcsname{ega4.c04.eq.18-2-1-2}
\expandafter\def\csname sga@target@IV.18.2.10\endcsname{ega4.c04.statement.18-2-10}
\expandafter\def\csname sga@target@IV.18.2.2\endcsname{ega4.c04.statement.18-2-2}
\expandafter\def\csname sga@target@IV.18.2.2.1\endcsname{ega4.c04.eq.18-2-2-1}
\expandafter\def\csname sga@target@IV.18.2.2.2\endcsname{ega4.c04.eq.18-2-2-2}
\expandafter\def\csname sga@target@IV.18.2.2.3\endcsname{ega4.c04.eq.18-2-2-3}
\expandafter\def\csname sga@target@IV.18.2.2.4\endcsname{ega4.c04.eq.18-2-2-4}
\expandafter\def\csname sga@target@IV.18.2.3\endcsname{ega4.c04.prop.18-2-3}
\expandafter\def\csname sga@target@IV.18.2.4\endcsname{ega4.c04.cor.18-2-4}
\expandafter\def\csname sga@target@IV.18.2.5\endcsname{ega4.c04.cor.18-2-5}
\expandafter\def\csname sga@target@IV.18.2.6\endcsname{ega4.c04.cor.18-2-6}
\expandafter\def\csname sga@target@IV.18.2.7\endcsname{ega4.c04.item.18-2-7}
\expandafter\def\csname sga@target@IV.18.2.7.1\endcsname{ega4.c04.eq.18-2-7-1}
\expandafter\def\csname sga@target@IV.18.2.7.2\endcsname{ega4.c04.eq.18-2-7-2}
\expandafter\def\csname sga@target@IV.18.2.8\endcsname{ega4.c04.prop.18-2-8}
\expandafter\def\csname sga@target@IV.18.2.9\endcsname{ega4.c04.cor.18-2-9}
\expandafter\def\csname sga@target@IV.18.3.1\endcsname{ega4.c04.prop.18-3-1}
\expandafter\def\csname sga@target@IV.18.3.2\endcsname{ega4.c04.prop.18-3-2}
\expandafter\def\csname sga@target@IV.18.3.2.1\endcsname{ega4.c04.lem.18-3-2-1}
\expandafter\def\csname sga@target@IV.18.3.2.2\endcsname{ega4.c04.eq.18-3-2-2}
\expandafter\def\csname sga@target@IV.18.3.3\endcsname{ega4.c04.statement.18-3-3}
\expandafter\def\csname sga@target@IV.18.3.4\endcsname{ega4.c04.thm.18-3-4}
\expandafter\def\csname sga@target@IV.18.4.1\endcsname{ega4.c04.lem.18-4-1}
\expandafter\def\csname sga@target@IV.18.4.10\endcsname{ega4.c04.prop.18-4-10}
\expandafter\def\csname sga@target@IV.18.4.11\endcsname{ega4.c04.statement.18-4-11}
\expandafter\def\csname sga@target@IV.18.4.12\endcsname{ega4.c04.cor.18-4-12}
\expandafter\def\csname sga@target@IV.18.4.13\endcsname{ega4.c04.cor.18-4-13}
\expandafter\def\csname sga@target@IV.18.4.14\endcsname{ega4.c04.cor.18-4-14}
\expandafter\def\csname sga@target@IV.18.4.2\endcsname{ega4.c04.prop.18-4-2}
\expandafter\def\csname sga@target@IV.18.4.3\endcsname{ega4.c04.cor.18-4-3}
\expandafter\def\csname sga@target@IV.18.4.4\endcsname{ega4.c04.lem.18-4-4}
\expandafter\def\csname sga@target@IV.18.4.5\endcsname{ega4.c04.prop.18-4-5}
\expandafter\def\csname sga@target@IV.18.4.6\endcsname{ega4.c04.thm.18-4-6}
\expandafter\def\csname sga@target@IV.18.4.6.1\endcsname{ega4.c04.lem.18-4-6-1}
\expandafter\def\csname sga@target@IV.18.4.7\endcsname{ega4.c04.cor.18-4-7}
\expandafter\def\csname sga@target@IV.18.4.8\endcsname{ega4.c04.cor.18-4-8}
\expandafter\def\csname sga@target@IV.18.4.9\endcsname{ega4.c04.prop.18-4-9}
\expandafter\def\csname sga@target@IV.18.5.1\endcsname{ega4.c04.statement.18-5-1}
\expandafter\def\csname sga@target@IV.18.5.1.1\endcsname{ega4.c04.eq.18-5-1-1}
\expandafter\def\csname sga@target@IV.18.5.10\endcsname{ega4.c04.prop.18-5-10}
\expandafter\def\csname sga@target@IV.18.5.11\endcsname{ega4.c04.thm.18-5-11}
\expandafter\def\csname sga@target@IV.18.5.12\endcsname{ega4.c04.cor.18-5-12}
\expandafter\def\csname sga@target@IV.18.5.13\endcsname{ega4.c04.rem.18-5-13}
\expandafter\def\csname sga@target@IV.18.5.13.1\endcsname{ega4.c04.lem.18-5-13-1}
\expandafter\def\csname sga@target@IV.18.5.14\endcsname{ega4.c04.prop.18-5-14}
\expandafter\def\csname sga@target@IV.18.5.15\endcsname{ega4.c04.prop.18-5-15}
\expandafter\def\csname sga@target@IV.18.5.16\endcsname{ega4.c04.item.18-5-16}
\expandafter\def\csname sga@target@IV.18.5.17\endcsname{ega4.c04.thm.18-5-17}
\expandafter\def\csname sga@target@IV.18.5.18\endcsname{ega4.c04.rem.18-5-18}
\expandafter\def\csname sga@target@IV.18.5.19\endcsname{ega4.c04.prop.18-5-19}
\expandafter\def\csname sga@target@IV.18.5.2\endcsname{ega4.c04.statement.18-5-2}
\expandafter\def\csname sga@target@IV.18.5.3\endcsname{ega4.c04.lem.18-5-3}
\expandafter\def\csname sga@target@IV.18.5.3.1\endcsname{ega4.c04.lem.18-5-3-1}
\expandafter\def\csname sga@target@IV.18.5.4\endcsname{ega4.c04.prop.18-5-4}
\expandafter\def\csname sga@target@IV.18.5.4.1\endcsname{ega4.c04.eq.18-5-4-1}
\expandafter\def\csname sga@target@IV.18.5.4.2\endcsname{ega4.c04.eq.18-5-4-2}
\expandafter\def\csname sga@target@IV.18.5.4.3\endcsname{ega4.c04.lem.18-5-4-3}
\expandafter\def\csname sga@target@IV.18.5.4.4\endcsname{ega4.c04.rem.18-5-4-4}
\expandafter\def\csname sga@target@IV.18.5.5\endcsname{ega4.c04.def.18-5-5}
\expandafter\def\csname sga@target@IV.18.5.6\endcsname{ega4.c04.prop.18-5-6}
\expandafter\def\csname sga@target@IV.18.5.7\endcsname{ega4.c04.rem.18-5-7}
\expandafter\def\csname sga@target@IV.18.5.8\endcsname{ega4.c04.def.18-5-8}
\expandafter\def\csname sga@target@IV.18.5.9\endcsname{ega4.c04.prop.18-5-9}
\expandafter\def\csname sga@target@IV.18.6.1\endcsname{ega4.c04.statement.18-6-1}
\expandafter\def\csname sga@target@IV.18.6.10\endcsname{ega4.c04.cor.18-6-10}
\expandafter\def\csname sga@target@IV.18.6.11\endcsname{ega4.c04.cor.18-6-11}
\expandafter\def\csname sga@target@IV.18.6.12\endcsname{ega4.c04.prop.18-6-12}
\expandafter\def\csname sga@target@IV.18.6.13\endcsname{ega4.c04.cor.18-6-13}
\expandafter\def\csname sga@target@IV.18.6.14\endcsname{ega4.c04.prop.18-6-14}
\expandafter\def\csname sga@target@IV.18.6.15\endcsname{ega4.c04.cor.18-6-15}
\expandafter\def\csname sga@target@IV.18.6.2\endcsname{ega4.c04.lem.18-6-2}
\expandafter\def\csname sga@target@IV.18.6.3\endcsname{ega4.c04.lem.18-6-3}
\expandafter\def\csname sga@target@IV.18.6.4\endcsname{ega4.c04.statement.18-6-4}
\expandafter\def\csname sga@target@IV.18.6.5\endcsname{ega4.c04.def.18-6-5}
\expandafter\def\csname sga@target@IV.18.6.6\endcsname{ega4.c04.thm.18-6-6}
\expandafter\def\csname sga@target@IV.18.6.7\endcsname{ega4.c04.statement.18-6-7}
\expandafter\def\csname sga@target@IV.18.6.8\endcsname{ega4.c04.prop.18-6-8}
\expandafter\def\csname sga@target@IV.18.6.9\endcsname{ega4.c04.thm.18-6-9}
\expandafter\def\csname sga@target@IV.18.6.9.1\endcsname{ega4.c04.lem.18-6-9-1}
\expandafter\def\csname sga@target@IV.18.6.9.2\endcsname{ega4.c04.lem.18-6-9-2}
\expandafter\def\csname sga@target@IV.18.7.1\endcsname{ega4.c04.statement.18-7-1}
\expandafter\def\csname sga@target@IV.18.7.2\endcsname{ega4.c04.prop.18-7-2}
\expandafter\def\csname sga@target@IV.18.7.3\endcsname{ega4.c04.cor.18-7-3}
\expandafter\def\csname sga@target@IV.18.7.4\endcsname{ega4.c04.cor.18-7-4}
\expandafter\def\csname sga@target@IV.18.7.5\endcsname{ega4.c04.prop.18-7-5}
\expandafter\def\csname sga@target@IV.18.7.5.1\endcsname{ega4.c04.lem.18-7-5-1}
\expandafter\def\csname sga@target@IV.18.7.6\endcsname{ega4.c04.cor.18-7-6}
\expandafter\def\csname sga@target@IV.18.7.7\endcsname{ega4.c04.rem.18-7-7}
\expandafter\def\csname sga@target@IV.18.8.1\endcsname{ega4.c04.prop.18-8-1}
\expandafter\def\csname sga@target@IV.18.8.10\endcsname{ega4.c04.prop.18-8-10}
\expandafter\def\csname sga@target@IV.18.8.11\endcsname{ega4.c04.rem.18-8-11}
\expandafter\def\csname sga@target@IV.18.8.12\endcsname{ega4.c04.prop.18-8-12}
\expandafter\def\csname sga@target@IV.18.8.13\endcsname{ega4.c04.cor.18-8-13}
\expandafter\def\csname sga@target@IV.18.8.14\endcsname{ega4.c04.cor.18-8-14}
\expandafter\def\csname sga@target@IV.18.8.15\endcsname{ega4.c04.prop.18-8-15}
\expandafter\def\csname sga@target@IV.18.8.16\endcsname{ega4.c04.cor.18-8-16}
\expandafter\def\csname sga@target@IV.18.8.17\endcsname{ega4.c04.prop.18-8-17}
\expandafter\def\csname sga@target@IV.18.8.18\endcsname{ega4.c04.prop.18-8-18}
\expandafter\def\csname sga@target@IV.18.8.2\endcsname{ega4.c04.def.18-8-2}
\expandafter\def\csname sga@target@IV.18.8.3\endcsname{ega4.c04.rem.18-8-3}
\expandafter\def\csname sga@target@IV.18.8.4\endcsname{ega4.c04.lem.18-8-4}
\expandafter\def\csname sga@target@IV.18.8.5\endcsname{ega4.c04.lem.18-8-5}
\expandafter\def\csname sga@target@IV.18.8.6\endcsname{ega4.c04.statement.18-8-6}
\expandafter\def\csname sga@target@IV.18.8.6.1\endcsname{ega4.c04.eq.18-8-6-1}
\expandafter\def\csname sga@target@IV.18.8.6.2\endcsname{ega4.c04.eq.18-8-6-2}
\expandafter\def\csname sga@target@IV.18.8.7\endcsname{ega4.c04.def.18-8-7}
\expandafter\def\csname sga@target@IV.18.8.8\endcsname{ega4.c04.prop.18-8-8}
\expandafter\def\csname sga@target@IV.18.8.8.1\endcsname{ega4.c04.statement.18-8-8-1}
\expandafter\def\csname sga@target@IV.18.8.9\endcsname{ega4.c04.statement.18-8-9}
\expandafter\def\csname sga@target@IV.18.9.1\endcsname{ega4.c04.thm.18-9-1}
\expandafter\def\csname sga@target@IV.18.9.10\endcsname{ega4.c04.rem.18-9-10}
\expandafter\def\csname sga@target@IV.18.9.11\endcsname{ega4.c04.cor.18-9-11}
\expandafter\def\csname sga@target@IV.18.9.2\endcsname{ega4.c04.cor.18-9-2}
\expandafter\def\csname sga@target@IV.18.9.3\endcsname{ega4.c04.cor.18-9-3}
\expandafter\def\csname sga@target@IV.18.9.4\endcsname{ega4.c04.cor.18-9-4}
\expandafter\def\csname sga@target@IV.18.9.5\endcsname{ega4.c04.cor.18-9-5}
\expandafter\def\csname sga@target@IV.18.9.5.1\endcsname{ega4.c04.eq.18-9-5-1}
\expandafter\def\csname sga@target@IV.18.9.6\endcsname{ega4.c04.statement.18-9-6}
\expandafter\def\csname sga@target@IV.18.9.7\endcsname{ega4.c04.thm.18-9-7}
\expandafter\def\csname sga@target@IV.18.9.7.1\endcsname{ega4.c04.rem.18-9-7-1}
\expandafter\def\csname sga@target@IV.18.9.7.2\endcsname{ega4.c04.statement.18-9-7-2}
\expandafter\def\csname sga@target@IV.18.9.7.3\endcsname{ega4.c04.lem.18-9-7-3}
\expandafter\def\csname sga@target@IV.18.9.7.4\endcsname{ega4.c04.eq.18-9-7-4}
\expandafter\def\csname sga@target@IV.18.9.7.5\endcsname{ega4.c04.lem.18-9-7-5}
\expandafter\def\csname sga@target@IV.18.9.7.6\endcsname{ega4.c04.rem.18-9-7-6}
\expandafter\def\csname sga@target@IV.18.9.7.7\endcsname{ega4.c04.lem.18-9-7-7}
\expandafter\def\csname sga@target@IV.18.9.7.8\endcsname{ega4.c04.rem.18-9-7-8}
\expandafter\def\csname sga@target@IV.18.9.8\endcsname{ega4.c04.cor.18-9-8}
\expandafter\def\csname sga@target@IV.18.9.9\endcsname{ega4.c04.cor.18-9-9}
\expandafter\def\csname sga@target@IV.19.1.1\endcsname{ega4.c04.prop.19-1-1}
\expandafter\def\csname sga@target@IV.19.1.2\endcsname{ega4.c04.cor.19-1-2}
\expandafter\def\csname sga@target@IV.19.1.3\endcsname{ega4.c04.def.19-1-3}
\expandafter\def\csname sga@target@IV.19.1.4\endcsname{ega4.c04.prop.19-1-4}
\expandafter\def\csname sga@target@IV.19.1.4.1\endcsname{ega4.c04.eq.19-1-4-1}
\expandafter\def\csname sga@target@IV.19.1.5\endcsname{ega4.c04.prop.19-1-5}
\expandafter\def\csname sga@target@IV.19.1.5.1\endcsname{ega4.c04.item.19-1-5-1}
\expandafter\def\csname sga@target@IV.19.1.6\endcsname{ega4.c04.cor.19-1-6}
\expandafter\def\csname sga@target@IV.19.1.6.1\endcsname{ega4.c04.lem.19-1-6-1}
\expandafter\def\csname sga@target@IV.19.1.7\endcsname{ega4.c04.item.19-1-7}
\expandafter\def\csname sga@target@IV.19.1.8\endcsname{ega4.c04.statement.19-1-8}
\expandafter\def\csname sga@target@IV.19.2.1\endcsname{ega4.c04.def.19-2-1}
\expandafter\def\csname sga@target@IV.19.2.2\endcsname{ega4.c04.def.19-2-2}
\expandafter\def\csname sga@target@IV.19.2.3\endcsname{ega4.c04.prop.19-2-3}
\expandafter\def\csname sga@target@IV.19.2.4\endcsname{ega4.c04.prop.19-2-4}
\expandafter\def\csname sga@target@IV.19.2.5\endcsname{ega4.c04.item.19-2-5}
\expandafter\def\csname sga@target@IV.19.2.6\endcsname{ega4.c04.cor.19-2-6}
\expandafter\def\csname sga@target@IV.19.2.7\endcsname{ega4.c04.prop.19-2-7}
\expandafter\def\csname sga@target@IV.19.2.8\endcsname{ega4.c04.prop.19-2-8}
\expandafter\def\csname sga@target@IV.19.2.9\endcsname{ega4.c04.prop.19-2-9}
\expandafter\def\csname sga@target@IV.19.3.1\endcsname{ega4.c04.def.19-3-1}
\expandafter\def\csname sga@target@IV.19.3.10\endcsname{ega4.c04.cor.19-3-10}
\expandafter\def\csname sga@target@IV.19.3.2\endcsname{ega4.c04.prop.19-3-2}
\expandafter\def\csname sga@target@IV.19.3.2.1\endcsname{ega4.c04.lem.19-3-2-1}
\expandafter\def\csname sga@target@IV.19.3.3\endcsname{ega4.c04.cor.19-3-3}
\expandafter\def\csname sga@target@IV.19.3.4\endcsname{ega4.c04.cor.19-3-4}
\expandafter\def\csname sga@target@IV.19.3.5\endcsname{ega4.c04.cor.19-3-5}
\expandafter\def\csname sga@target@IV.19.3.6\endcsname{ega4.c04.def.19-3-6}
\expandafter\def\csname sga@target@IV.19.3.7\endcsname{ega4.c04.prop.19-3-7}
\expandafter\def\csname sga@target@IV.19.3.8\endcsname{ega4.c04.cor.19-3-8}
\expandafter\def\csname sga@target@IV.19.3.9\endcsname{ega4.c04.prop.19-3-9}
\expandafter\def\csname sga@target@IV.19.4.1\endcsname{ega4.c04.statement.19-4-1}
\expandafter\def\csname sga@target@IV.19.4.1.1\endcsname{ega4.c04.subsec.19-4-1-1}
\expandafter\def\csname sga@target@IV.19.4.10\endcsname{ega4.c04.rem.19-4-10}
\expandafter\def\csname sga@target@IV.19.4.11\endcsname{ega4.c04.prop.19-4-11}
\expandafter\def\csname sga@target@IV.19.4.11.1\endcsname{ega4.c04.eq.19-4-11-1}
\expandafter\def\csname sga@target@IV.19.4.2\endcsname{ega4.c04.lem.19-4-2}
\expandafter\def\csname sga@target@IV.19.4.3\endcsname{ega4.c04.cor.19-4-3}
\expandafter\def\csname sga@target@IV.19.4.4\endcsname{ega4.c04.prop.19-4-4}
\expandafter\def\csname sga@target@IV.19.4.5\endcsname{ega4.c04.cor.19-4-5}
\expandafter\def\csname sga@target@IV.19.4.6\endcsname{ega4.c04.prop.19-4-6}
\expandafter\def\csname sga@target@IV.19.4.7\endcsname{ega4.c04.cor.19-4-7}
\expandafter\def\csname sga@target@IV.19.4.8\endcsname{ega4.c04.prop.19-4-8}
\expandafter\def\csname sga@target@IV.19.4.9\endcsname{ega4.c04.cor.19-4-9}
\expandafter\def\csname sga@target@IV.19.5.1\endcsname{ega4.c04.thm.19-5-1}
\expandafter\def\csname sga@target@IV.19.5.1.1\endcsname{ega4.c04.eq.19-5-1-1}
\expandafter\def\csname sga@target@IV.19.5.2\endcsname{ega4.c04.cor.19-5-2}
\expandafter\def\csname sga@target@IV.19.5.3\endcsname{ega4.c04.prop.19-5-3}
\expandafter\def\csname sga@target@IV.19.5.4\endcsname{ega4.c04.statement.19-5-4}
\expandafter\def\csname sga@target@IV.19.5.5\endcsname{ega4.c04.thm.19-5-5}
\expandafter\def\csname sga@target@IV.19.5.5.1\endcsname{ega4.c04.statement.19-5-5-1}
\expandafter\def\csname sga@target@IV.19.5.5.2\endcsname{ega4.c04.statement.19-5-5-2}
\expandafter\def\csname sga@target@IV.19.5.5.3\endcsname{ega4.c04.statement.19-5-5-3}
\expandafter\def\csname sga@target@IV.19.5.5.4\endcsname{ega4.c04.statement.19-5-5-4}
\expandafter\def\csname sga@target@IV.19.5.5.5\endcsname{ega4.c04.statement.19-5-5-5}
\expandafter\def\csname sga@target@IV.19.5.6\endcsname{ega4.c04.statement.19-5-6}
\expandafter\def\csname sga@target@IV.19.6.1\endcsname{ega4.c04.statement.19-6-1}
\expandafter\def\csname sga@target@IV.19.6.1.1\endcsname{ega4.c04.eq.19-6-1-1}
\expandafter\def\csname sga@target@IV.19.6.2\endcsname{ega4.c04.statement.19-6-2}
\expandafter\def\csname sga@target@IV.19.6.2.1\endcsname{ega4.c04.statement.19-6-2-1}
\expandafter\def\csname sga@target@IV.19.6.3\endcsname{ega4.c04.prop.19-6-3}
\expandafter\def\csname sga@target@IV.19.6.4\endcsname{ega4.c04.cor.19-6-4}
\expandafter\def\csname sga@target@IV.19.7\endcsname{ega4.c04.subsec.19-7.key-iv-19-7}
\expandafter\def\csname sga@target@IV.19.7.1\endcsname{ega4.c04.thm.19-7-1}
\expandafter\def\csname sga@target@IV.19.7.1.1\endcsname{ega4.c04.statement.19-7-1-1}
\expandafter\def\csname sga@target@IV.19.7.1.10\endcsname{ega4.c04.statement.19-7-1-10}
\expandafter\def\csname sga@target@IV.19.7.1.11\endcsname{ega4.c04.statement.19-7-1-11}
\expandafter\def\csname sga@target@IV.19.7.1.12\endcsname{ega4.c04.statement.19-7-1-12}
\expandafter\def\csname sga@target@IV.19.7.1.13\endcsname{ega4.c04.statement.19-7-1-13}
\expandafter\def\csname sga@target@IV.19.7.1.14\endcsname{ega4.c04.statement.19-7-1-14}
\expandafter\def\csname sga@target@IV.19.7.1.15\endcsname{ega4.c04.statement.19-7-1-15}
\expandafter\def\csname sga@target@IV.19.7.1.16\endcsname{ega4.c04.statement.19-7-1-16}
\expandafter\def\csname sga@target@IV.19.7.1.17\endcsname{ega4.c04.statement.19-7-1-17}
\expandafter\def\csname sga@target@IV.19.7.1.18\endcsname{ega4.c04.statement.19-7-1-18}
\expandafter\def\csname sga@target@IV.19.7.1.19\endcsname{ega4.c04.statement.19-7-1-19}
\expandafter\def\csname sga@target@IV.19.7.1.2\endcsname{ega4.c04.statement.19-7-1-2}
\expandafter\def\csname sga@target@IV.19.7.1.20\endcsname{ega4.c04.statement.19-7-1-20}
\expandafter\def\csname sga@target@IV.19.7.1.21\endcsname{ega4.c04.statement.19-7-1-21}
\expandafter\def\csname sga@target@IV.19.7.1.22\endcsname{ega4.c04.statement.19-7-1-22}
\expandafter\def\csname sga@target@IV.19.7.1.23\endcsname{ega4.c04.statement.19-7-1-23}
\expandafter\def\csname sga@target@IV.19.7.1.3\endcsname{ega4.c04.statement.19-7-1-3}
\expandafter\def\csname sga@target@IV.19.7.1.5\endcsname{ega4.c04.statement.19-7-1-5}
\expandafter\def\csname sga@target@IV.19.7.1.6\endcsname{ega4.c04.statement.19-7-1-6}
\expandafter\def\csname sga@target@IV.19.7.1.7\endcsname{ega4.c04.lem.19-7-1-7}
\expandafter\def\csname sga@target@IV.19.7.1.8\endcsname{ega4.c04.statement.19-7-1-8}
\expandafter\def\csname sga@target@IV.19.7.1.9\endcsname{ega4.c04.statement.19-7-1-9}
\expandafter\def\csname sga@target@IV.19.7.2\endcsname{ega4.c04.statement.19-7-2}
\expandafter\def\csname sga@target@IV.19.7.3\endcsname{ega4.c04.cor.19-7-3}
\expandafter\def\csname sga@target@IV.19.7.3.1\endcsname{ega4.c04.eq.19-7-3-1}
\expandafter\def\csname sga@target@IV.19.7.4\endcsname{ega4.c04.cor.19-7-4}
\expandafter\def\csname sga@target@IV.19.7.4.1\endcsname{ega4.c04.eq.19-7-4-1}
\expandafter\def\csname sga@target@IV.19.7.5\endcsname{ega4.c04.cor.19-7-5}
\expandafter\def\csname sga@target@IV.19.8.1\endcsname{ega4.c04.prop.19-8-1}
\expandafter\def\csname sga@target@IV.19.8.2\endcsname{ega4.c04.prop.19-8-2}
\expandafter\def\csname sga@target@IV.19.8.3\endcsname{ega4.c04.rem.19-8-3}
\expandafter\def\csname sga@target@IV.19.9.1\endcsname{ega4.c04.statement.19-9-1}
\expandafter\def\csname sga@target@IV.19.9.1.1\endcsname{ega4.c04.eq.19-9-1-1}
\expandafter\def\csname sga@target@IV.19.9.1.2\endcsname{ega4.c04.eq.19-9-1-2}
\expandafter\def\csname sga@target@IV.19.9.1.3\endcsname{ega4.c04.eq.19-9-1-3}
\expandafter\def\csname sga@target@IV.19.9.2\endcsname{ega4.c04.lem.19-9-2}
\expandafter\def\csname sga@target@IV.19.9.3\endcsname{ega4.c04.prop.19-9-3}
\expandafter\def\csname sga@target@IV.19.9.4\endcsname{ega4.c04.cor.19-9-4}
\expandafter\def\csname sga@target@IV.19.9.5\endcsname{ega4.c04.prop.19-9-5}
\expandafter\def\csname sga@target@IV.19.9.6\endcsname{ega4.c04.prop.19-9-6}
\expandafter\def\csname sga@target@IV.19.9.7\endcsname{ega4.c04.cor.19-9-7}
\expandafter\def\csname sga@target@IV.19.9.8\endcsname{ega4.c04.prop.19-9-8}
\expandafter\def\csname sga@target@IV.19.9.9\endcsname{ega4.c04.rem.19-9-9}
\expandafter\def\csname sga@target@IV.2.1.1\endcsname{ega4.c04.statement.2-1-1}
\expandafter\def\csname sga@target@IV.2.1.10\endcsname{ega4.c04.cor.2-1-10}
\expandafter\def\csname sga@target@IV.2.1.11\endcsname{ega4.c04.prop.2-1-11}
\expandafter\def\csname sga@target@IV.2.1.12\endcsname{ega4.c04.prop.2-1-12}
\expandafter\def\csname sga@target@IV.2.1.13\endcsname{ega4.c04.prop.2-1-13}
\expandafter\def\csname sga@target@IV.2.1.14\endcsname{ega4.c04.prop.2-1-14}
\expandafter\def\csname sga@target@IV.2.1.2\endcsname{ega4.c04.prop.2-1-2}
\expandafter\def\csname sga@target@IV.2.1.3\endcsname{ega4.c04.prop.2-1-3}
\expandafter\def\csname sga@target@IV.2.1.4\endcsname{ega4.c04.prop.2-1-4}
\expandafter\def\csname sga@target@IV.2.1.5\endcsname{ega4.c04.prop.2-1-5}
\expandafter\def\csname sga@target@IV.2.1.6\endcsname{ega4.c04.cor.2-1-6}
\expandafter\def\csname sga@target@IV.2.1.7\endcsname{ega4.c04.cor.2-1-7}
\expandafter\def\csname sga@target@IV.2.1.8\endcsname{ega4.c04.prop.2-1-8}
\expandafter\def\csname sga@target@IV.2.1.9\endcsname{ega4.c04.cor.2-1-9}
\expandafter\def\csname sga@target@IV.2.2.1\endcsname{ega4.c04.prop.2-2-1}
\expandafter\def\csname sga@target@IV.2.2.10\endcsname{ega4.c04.prop.2-2-10}
\expandafter\def\csname sga@target@IV.2.2.10.1\endcsname{ega4.c04.eq.2-2-10-1}
\expandafter\def\csname sga@target@IV.2.2.11\endcsname{ega4.c04.cor.2-2-11}
\expandafter\def\csname sga@target@IV.2.2.12\endcsname{ega4.c04.cor.2-2-12}
\expandafter\def\csname sga@target@IV.2.2.13\endcsname{ega4.c04.cor.2-2-13}
\expandafter\def\csname sga@target@IV.2.2.14\endcsname{ega4.c04.prop.2-2-14}
\expandafter\def\csname sga@target@IV.2.2.15\endcsname{ega4.c04.prop.2-2-15}
\expandafter\def\csname sga@target@IV.2.2.16\endcsname{ega4.c04.cor.2-2-16}
\expandafter\def\csname sga@target@IV.2.2.17\endcsname{ega4.c04.prop.2-2-17}
\expandafter\def\csname sga@target@IV.2.2.2\endcsname{ega4.c04.prop.2-2-2}
\expandafter\def\csname sga@target@IV.2.2.2.1\endcsname{ega4.c04.eq.2-2-2-1}
\expandafter\def\csname sga@target@IV.2.2.3\endcsname{ega4.c04.cor.2-2-3}
\expandafter\def\csname sga@target@IV.2.2.4\endcsname{ega4.c04.def.2-2-4}
\expandafter\def\csname sga@target@IV.2.2.5\endcsname{ega4.c04.cor.2-2-5}
\expandafter\def\csname sga@target@IV.2.2.6\endcsname{ega4.c04.cor.2-2-6}
\expandafter\def\csname sga@target@IV.2.2.7\endcsname{ega4.c04.prop.2-2-7}
\expandafter\def\csname sga@target@IV.2.2.8\endcsname{ega4.c04.cor.2-2-8}
\expandafter\def\csname sga@target@IV.2.2.8.1\endcsname{ega4.c04.eq.2-2-8-1}
\expandafter\def\csname sga@target@IV.2.2.9\endcsname{ega4.c04.rem.2-2-9}
\expandafter\def\csname sga@target@IV.2.3.1\endcsname{ega4.c04.lem.2-3-1}
\expandafter\def\csname sga@target@IV.2.3.1.1\endcsname{ega4.c04.eq.2-3-1-1}
\expandafter\def\csname sga@target@IV.2.3.10\endcsname{ega4.c04.thm.2-3-10}
\expandafter\def\csname sga@target@IV.2.3.11\endcsname{ega4.c04.cor.2-3-11}
\expandafter\def\csname sga@target@IV.2.3.12\endcsname{ega4.c04.cor.2-3-12}
\expandafter\def\csname sga@target@IV.2.3.13\endcsname{ega4.c04.cor.2-3-13}
\expandafter\def\csname sga@target@IV.2.3.14\endcsname{ega4.c04.cor.2-3-14}
\expandafter\def\csname sga@target@IV.2.3.15\endcsname{ega4.c04.rem.2-3-15}
\expandafter\def\csname sga@target@IV.2.3.2\endcsname{ega4.c04.prop.2-3-2}
\expandafter\def\csname sga@target@IV.2.3.3\endcsname{ega4.c04.statement.2-3-3}
\expandafter\def\csname sga@target@IV.2.3.4\endcsname{ega4.c04.prop.2-3-4}
\expandafter\def\csname sga@target@IV.2.3.5\endcsname{ega4.c04.cor.2-3-5}
\expandafter\def\csname sga@target@IV.2.3.6\endcsname{ega4.c04.prop.2-3-6}
\expandafter\def\csname sga@target@IV.2.3.7\endcsname{ega4.c04.prop.2-3-7}
\expandafter\def\csname sga@target@IV.2.3.8\endcsname{ega4.c04.cor.2-3-8}
\expandafter\def\csname sga@target@IV.2.3.9\endcsname{ega4.c04.rem.2-3-9}
\expandafter\def\csname sga@target@IV.2.4.1\endcsname{ega4.c04.statement.2-4-1}
\expandafter\def\csname sga@target@IV.2.4.10\endcsname{ega4.c04.cor.2-4-10}
\expandafter\def\csname sga@target@IV.2.4.11\endcsname{ega4.c04.rem.2-4-11}
\expandafter\def\csname sga@target@IV.2.4.2\endcsname{ega4.c04.def.2-4-2}
\expandafter\def\csname sga@target@IV.2.4.3\endcsname{ega4.c04.prop.2-4-3}
\expandafter\def\csname sga@target@IV.2.4.4\endcsname{ega4.c04.prop.2-4-4}
\expandafter\def\csname sga@target@IV.2.4.5\endcsname{ega4.c04.prop.2-4-5}
\expandafter\def\csname sga@target@IV.2.4.6\endcsname{ega4.c04.thm.2-4-6}
\expandafter\def\csname sga@target@IV.2.4.7\endcsname{ega4.c04.cor.2-4-7}
\expandafter\def\csname sga@target@IV.2.4.8\endcsname{ega4.c04.rem.2-4-8}
\expandafter\def\csname sga@target@IV.2.4.9\endcsname{ega4.c04.prop.2-4-9}
\expandafter\def\csname sga@target@IV.2.5.1\endcsname{ega4.c04.prop.2-5-1}
\expandafter\def\csname sga@target@IV.2.5.2\endcsname{ega4.c04.prop.2-5-2}
\expandafter\def\csname sga@target@IV.2.5.3\endcsname{ega4.c04.lem.2-5-3}
\expandafter\def\csname sga@target@IV.2.5.4\endcsname{ega4.c04.rem.2-5-4}
\expandafter\def\csname sga@target@IV.2.5.4.1\endcsname{ega4.c04.lem.2-5-4-1}
\expandafter\def\csname sga@target@IV.2.5.5\endcsname{ega4.c04.prop.2-5-5}
\expandafter\def\csname sga@target@IV.2.5.5.1\endcsname{ega4.c04.statement.2-5-5-1}
\expandafter\def\csname sga@target@IV.2.5.5.2\endcsname{ega4.c04.lem.2-5-5-2}
\expandafter\def\csname sga@target@IV.2.5.5.3\endcsname{ega4.c04.statement.2-5-5-3}
\expandafter\def\csname sga@target@IV.2.5.5.4\endcsname{ega4.c04.rem.2-5-5-4}
\expandafter\def\csname sga@target@IV.2.5.6\endcsname{ega4.c04.lem.2-5-6}
\expandafter\def\csname sga@target@IV.2.5.7\endcsname{ega4.c04.cor.2-5-7}
\expandafter\def\csname sga@target@IV.2.5.8\endcsname{ega4.c04.prop.2-5-8}
\expandafter\def\csname sga@target@IV.2.5.8.1\endcsname{ega4.c04.statement.2-5-8-1}
\expandafter\def\csname sga@target@IV.2.5.8.2\endcsname{ega4.c04.statement.2-5-8-2}
\expandafter\def\csname sga@target@IV.2.5.8.3\endcsname{ega4.c04.statement.2-5-8-3}
\expandafter\def\csname sga@target@IV.2.5.8.4\endcsname{ega4.c04.statement.2-5-8-4}
\expandafter\def\csname sga@target@IV.2.5.8.5\endcsname{ega4.c04.statement.2-5-8-5}
\expandafter\def\csname sga@target@IV.2.5.8.6\endcsname{ega4.c04.statement.2-5-8-6}
\expandafter\def\csname sga@target@IV.2.5.8.7\endcsname{ega4.c04.statement.2-5-8-7}
\expandafter\def\csname sga@target@IV.2.5.8.8\endcsname{ega4.c04.statement.2-5-8-8}
\expandafter\def\csname sga@target@IV.2.5.8.9\endcsname{ega4.c04.statement.2-5-8-9}
\expandafter\def\csname sga@target@IV.2.6.1\endcsname{ega4.c04.prop.2-6-1}
\expandafter\def\csname sga@target@IV.2.6.2\endcsname{ega4.c04.prop.2-6-2}
\expandafter\def\csname sga@target@IV.2.6.3\endcsname{ega4.c04.rem.2-6-3}
\expandafter\def\csname sga@target@IV.2.6.4\endcsname{ega4.c04.cor.2-6-4}
\expandafter\def\csname sga@target@IV.2.7.1\endcsname{ega4.c04.prop.2-7-1}
\expandafter\def\csname sga@target@IV.2.7.1.1\endcsname{ega4.c04.lem.2-7-1-1}
\expandafter\def\csname sga@target@IV.2.7.1.viii\endcsname{ega4.c04.item.2-7-1-viii}
\expandafter\def\csname sga@target@IV.2.7.2\endcsname{ega4.c04.cor.2-7-2}
\expandafter\def\csname sga@target@IV.2.7.3\endcsname{ega4.c04.sec.2-7-3}
\expandafter\def\csname sga@target@IV.2.8.1\endcsname{ega4.c04.prop.2-8-1}
\expandafter\def\csname sga@target@IV.2.8.1.1\endcsname{ega4.c04.lem.2-8-1-1}
\expandafter\def\csname sga@target@IV.2.8.2\endcsname{ega4.c04.cor.2-8-2}
\expandafter\def\csname sga@target@IV.2.8.3\endcsname{ega4.c04.statement.2-8-3}
\expandafter\def\csname sga@target@IV.2.8.4\endcsname{ega4.c04.prop.2-8-4}
\expandafter\def\csname sga@target@IV.2.8.4.1\endcsname{ega4.c04.statement.2-8-4-1}
\expandafter\def\csname sga@target@IV.2.8.5\endcsname{ega4.c04.prop.2-8-5}
\expandafter\def\csname sga@target@IV.2.8.6\endcsname{ega4.c04.cor.2-8-6}
\expandafter\def\csname sga@target@IV.2.8.6.1\endcsname{ega4.c04.statement.2-8-6-1}
\expandafter\def\csname sga@target@IV.20\endcsname{ega4.c04.sec.20}
\expandafter\def\csname sga@target@IV.20.0\endcsname{ega4.c04.subsec.20-0}
\expandafter\def\csname sga@target@IV.20.1\endcsname{ega4.c04.subsec.20-1}
\expandafter\def\csname sga@target@IV.20.1.1\endcsname{ega4.c04.statement.20-1-1}
\expandafter\def\csname sga@target@IV.20.1.1.1\endcsname{ega4.c04.eq.20-1-1-1}
\expandafter\def\csname sga@target@IV.20.1.10\endcsname{ega4.c04.statement.20-1-10}
\expandafter\def\csname sga@target@IV.20.1.11\endcsname{ega4.c04.statement.20-1-11}
\expandafter\def\csname sga@target@IV.20.1.11.1\endcsname{ega4.c04.eq.20-1-11-1}
\expandafter\def\csname sga@target@IV.20.1.11.2\endcsname{ega4.c04.eq.20-1-11-2}
\expandafter\def\csname sga@target@IV.20.1.12\endcsname{ega4.c04.prop.20-1-12}
\expandafter\def\csname sga@target@IV.20.1.13\endcsname{ega4.c04.rem.20-1-13}
\expandafter\def\csname sga@target@IV.20.1.2\endcsname{ega4.c04.statement.20-1-2}
\expandafter\def\csname sga@target@IV.20.1.2.1\endcsname{ega4.c04.statement.20-1-2-1}
\expandafter\def\csname sga@target@IV.20.1.2.2\endcsname{ega4.c04.statement.20-1-2-2}
\expandafter\def\csname sga@target@IV.20.1.3\endcsname{ega4.c04.statement.20-1-3}
\expandafter\def\csname sga@target@IV.20.1.3.1\endcsname{ega4.c04.statement.20-1-3-1}
\expandafter\def\csname sga@target@IV.20.1.4\endcsname{ega4.c04.statement.20-1-4}
\expandafter\def\csname sga@target@IV.20.1.4.1\endcsname{ega4.c04.statement.20-1-4-1}
\expandafter\def\csname sga@target@IV.20.1.5\endcsname{ega4.c04.statement.20-1-5}
\expandafter\def\csname sga@target@IV.20.1.5.1\endcsname{ega4.c04.statement.20-1-5-1}
\expandafter\def\csname sga@target@IV.20.1.6\endcsname{ega4.c04.prop.20-1-6}
\expandafter\def\csname sga@target@IV.20.1.7\endcsname{ega4.c04.statement.20-1-7}
\expandafter\def\csname sga@target@IV.20.1.8\endcsname{ega4.c04.statement.20-1-8}
\expandafter\def\csname sga@target@IV.20.1.9\endcsname{ega4.c04.statement.20-1-9}
\expandafter\def\csname sga@target@IV.20.2\endcsname{ega4.c04.subsec.20-2}
\expandafter\def\csname sga@target@IV.20.2.1\endcsname{ega4.c04.statement.20-2-1}
\expandafter\def\csname sga@target@IV.20.2.10\endcsname{ega4.c04.statement.20-2-10}
\expandafter\def\csname sga@target@IV.20.2.10.1\endcsname{ega4.c04.eq.20-2-10-1}
\expandafter\def\csname sga@target@IV.20.2.10.2\endcsname{ega4.c04.eq.20-2-10-2}
\expandafter\def\csname sga@target@IV.20.2.11\endcsname{ega4.c04.prop.20-2-11}
\expandafter\def\csname sga@target@IV.20.2.11.1\endcsname{ega4.c04.eq.20-2-11-1}
\expandafter\def\csname sga@target@IV.20.2.12\endcsname{ega4.c04.cor.20-2-12}
\expandafter\def\csname sga@target@IV.20.2.13\endcsname{ega4.c04.statement.20-2-13}
\expandafter\def\csname sga@target@IV.20.2.14\endcsname{ega4.c04.statement.20-2-14}
\expandafter\def\csname sga@target@IV.20.2.15\endcsname{ega4.c04.prop.20-2-15}
\expandafter\def\csname sga@target@IV.20.2.16\endcsname{ega4.c04.rem.20-2-16}
\expandafter\def\csname sga@target@IV.20.2.2\endcsname{ega4.c04.statement.20-2-2}
\expandafter\def\csname sga@target@IV.20.2.3\endcsname{ega4.c04.statement.20-2-3}
\expandafter\def\csname sga@target@IV.20.2.3.1\endcsname{ega4.c04.statement.20-2-3-1}
\expandafter\def\csname sga@target@IV.20.2.4\endcsname{ega4.c04.prop.20-2-4}
\expandafter\def\csname sga@target@IV.20.2.5\endcsname{ega4.c04.cor.20-2-5}
\expandafter\def\csname sga@target@IV.20.2.6\endcsname{ega4.c04.cor.20-2-6}
\expandafter\def\csname sga@target@IV.20.2.7\endcsname{ega4.c04.rem.20-2-7}
\expandafter\def\csname sga@target@IV.20.2.8\endcsname{ega4.c04.statement.20-2-8}
\expandafter\def\csname sga@target@IV.20.2.9\endcsname{ega4.c04.lem.20-2-9}
\expandafter\def\csname sga@target@IV.20.3\endcsname{ega4.c04.subsec.20-3}
\expandafter\def\csname sga@target@IV.20.3.1\endcsname{ega4.c04.statement.20-3-1}
\expandafter\def\csname sga@target@IV.20.3.10\endcsname{ega4.c04.lem.20-3-10}
\expandafter\def\csname sga@target@IV.20.3.11\endcsname{ega4.c04.prop.20-3-11}
\expandafter\def\csname sga@target@IV.20.3.11.1\endcsname{ega4.c04.eq.20-3-11-1}
\expandafter\def\csname sga@target@IV.20.3.12\endcsname{ega4.c04.cor.20-3-12}
\expandafter\def\csname sga@target@IV.20.3.13\endcsname{ega4.c04.rem.20-3-13}
\expandafter\def\csname sga@target@IV.20.3.2\endcsname{ega4.c04.statement.20-3-2}
\expandafter\def\csname sga@target@IV.20.3.2.1\endcsname{ega4.c04.eq.20-3-2-1}
\expandafter\def\csname sga@target@IV.20.3.3\endcsname{ega4.c04.statement.20-3-3}
\expandafter\def\csname sga@target@IV.20.3.4\endcsname{ega4.c04.statement.20-3-4}
\expandafter\def\csname sga@target@IV.20.3.5\endcsname{ega4.c04.prop.20-3-5}
\expandafter\def\csname sga@target@IV.20.3.6\endcsname{ega4.c04.statement.20-3-6}
\expandafter\def\csname sga@target@IV.20.3.6.1\endcsname{ega4.c04.eq.20-3-6-1}
\expandafter\def\csname sga@target@IV.20.3.7\endcsname{ega4.c04.prop.20-3-7}
\expandafter\def\csname sga@target@IV.20.3.8\endcsname{ega4.c04.prop.20-3-8}
\expandafter\def\csname sga@target@IV.20.3.8.1\endcsname{ega4.c04.eq.20-3-8-1}
\expandafter\def\csname sga@target@IV.20.3.8.2\endcsname{ega4.c04.lem.20-3-8-2}
\expandafter\def\csname sga@target@IV.20.3.8.3\endcsname{ega4.c04.lem.20-3-8-3}
\expandafter\def\csname sga@target@IV.20.3.8.4\endcsname{ega4.c04.lem.20-3-8-4}
\expandafter\def\csname sga@target@IV.20.3.8.5\endcsname{ega4.c04.sec.20-3-8-5}
\expandafter\def\csname sga@target@IV.20.3.8.6\endcsname{ega4.c04.sec.20-3-8-6}
\expandafter\def\csname sga@target@IV.20.3.8.7\endcsname{ega4.c04.sec.20-3-8-7}
\expandafter\def\csname sga@target@IV.20.3.8.8\endcsname{ega4.c04.lem.20-3-8-8}
\expandafter\def\csname sga@target@IV.20.3.9\endcsname{ega4.c04.cor.20-3-9}
\expandafter\def\csname sga@target@IV.20.4\endcsname{ega4.c04.subsec.20-4}
\expandafter\def\csname sga@target@IV.20.4.1\endcsname{ega4.c04.statement.20-4-1}
\expandafter\def\csname sga@target@IV.20.4.10\endcsname{ega4.c04.cor.20-4-10}
\expandafter\def\csname sga@target@IV.20.4.11\endcsname{ega4.c04.prop.20-4-11}
\expandafter\def\csname sga@target@IV.20.4.12\endcsname{ega4.c04.cor.20-4-12}
\expandafter\def\csname sga@target@IV.20.4.2\endcsname{ega4.c04.statement.20-4-2}
\expandafter\def\csname sga@target@IV.20.4.3\endcsname{ega4.c04.statement.20-4-3}
\expandafter\def\csname sga@target@IV.20.4.4\endcsname{ega4.c04.prop.20-4-4}
\expandafter\def\csname sga@target@IV.20.4.5\endcsname{ega4.c04.prop.20-4-5}
\expandafter\def\csname sga@target@IV.20.4.6\endcsname{ega4.c04.prop.20-4-6}
\expandafter\def\csname sga@target@IV.20.4.7\endcsname{ega4.c04.cor.20-4-7}
\expandafter\def\csname sga@target@IV.20.4.8\endcsname{ega4.c04.rem.20-4-8}
\expandafter\def\csname sga@target@IV.20.4.9\endcsname{ega4.c04.cor.20-4-9}
\expandafter\def\csname sga@target@IV.20.5\endcsname{ega4.c04.subsec.20-5}
\expandafter\def\csname sga@target@IV.20.5.1\endcsname{ega4.c04.statement.20-5-1}
\expandafter\def\csname sga@target@IV.20.5.2\endcsname{ega4.c04.statement.20-5-2}
\expandafter\def\csname sga@target@IV.20.5.2.1\endcsname{ega4.c04.statement.20-5-2-1}
\expandafter\def\csname sga@target@IV.20.5.3\endcsname{ega4.c04.prop.20-5-3}
\expandafter\def\csname sga@target@IV.20.5.4\endcsname{ega4.c04.statement.20-5-4}
\expandafter\def\csname sga@target@IV.20.5.5\endcsname{ega4.c04.statement.20-5-5}
\expandafter\def\csname sga@target@IV.20.5.6\endcsname{ega4.c04.statement.20-5-6}
\expandafter\def\csname sga@target@IV.20.5.7\endcsname{ega4.c04.statement.20-5-7}
\expandafter\def\csname sga@target@IV.20.5.8\endcsname{ega4.c04.statement.20-5-8}
\expandafter\def\csname sga@target@IV.20.6\endcsname{ega4.c04.subsec.20-6}
\expandafter\def\csname sga@target@IV.20.6.1\endcsname{ega4.c04.statement.20-6-1}
\expandafter\def\csname sga@target@IV.20.6.1.1\endcsname{ega4.c04.eq.20-6-1-1}
\expandafter\def\csname sga@target@IV.20.6.2\endcsname{ega4.c04.statement.20-6-2}
\expandafter\def\csname sga@target@IV.20.6.2.1\endcsname{ega4.c04.statement.20-6-2-1}
\expandafter\def\csname sga@target@IV.20.6.2.2\endcsname{ega4.c04.eq.20-6-2-2}
\expandafter\def\csname sga@target@IV.20.6.3\endcsname{ega4.c04.prop.20-6-3}
\expandafter\def\csname sga@target@IV.20.6.4\endcsname{ega4.c04.prop.20-6-4}
\expandafter\def\csname sga@target@IV.20.6.5\endcsname{ega4.c04.statement.20-6-5}
\expandafter\def\csname sga@target@IV.20.6.5.1\endcsname{ega4.c04.eq.20-6-5-1}
\expandafter\def\csname sga@target@IV.20.6.6\endcsname{ega4.c04.prop.20-6-6}
\expandafter\def\csname sga@target@IV.20.6.7\endcsname{ega4.c04.statement.20-6-7}
\expandafter\def\csname sga@target@IV.20.6.8\endcsname{ega4.c04.prop.20-6-8}
\expandafter\def\csname sga@target@IV.20.6.8.1\endcsname{ega4.c04.eq.20-6-8-1}
\expandafter\def\csname sga@target@IV.20.6.9\endcsname{ega4.c04.statement.20-6-9}
\expandafter\def\csname sga@target@IV.20.6.9.1\endcsname{ega4.c04.statement.20-6-9-1}
\expandafter\def\csname sga@target@IV.21.1\endcsname{ega4.c04.subsec.21-1.key-iv-21-1}
\expandafter\def\csname sga@target@IV.21.1.1\endcsname{ega4.c04.statement.21-1-1}
\expandafter\def\csname sga@target@IV.21.1.10\endcsname{ega4.c04.prop.21-1-10}
\expandafter\def\csname sga@target@IV.21.1.2\endcsname{ega4.c04.statement.21-1-2}
\expandafter\def\csname sga@target@IV.21.1.3\endcsname{ega4.c04.statement.21-1-3}
\expandafter\def\csname sga@target@IV.21.1.3.1\endcsname{ega4.c04.statement.21-1-3-1}
\expandafter\def\csname sga@target@IV.21.1.3.2\endcsname{ega4.c04.statement.21-1-3-2}
\expandafter\def\csname sga@target@IV.21.1.3.3\endcsname{ega4.c04.statement.21-1-3-3}
\expandafter\def\csname sga@target@IV.21.1.3.4\endcsname{ega4.c04.statement.21-1-3-4}
\expandafter\def\csname sga@target@IV.21.1.4\endcsname{ega4.c04.statement.21-1-4}
\expandafter\def\csname sga@target@IV.21.1.4.1\endcsname{ega4.c04.eq.21-1-4-1}
\expandafter\def\csname sga@target@IV.21.1.4.2\endcsname{ega4.c04.eq.21-1-4-2}
\expandafter\def\csname sga@target@IV.21.1.4.3\endcsname{ega4.c04.eq.21-1-4-3}
\expandafter\def\csname sga@target@IV.21.1.5\endcsname{ega4.c04.statement.21-1-5}
\expandafter\def\csname sga@target@IV.21.1.5.1\endcsname{ega4.c04.statement.21-1-5-1}
\expandafter\def\csname sga@target@IV.21.1.5.2\endcsname{ega4.c04.eq.21-1-5-2}
\expandafter\def\csname sga@target@IV.21.1.6\endcsname{ega4.c04.def.21-1-6}
\expandafter\def\csname sga@target@IV.21.1.6.1\endcsname{ega4.c04.eq.21-1-6-1}
\expandafter\def\csname sga@target@IV.21.1.6.2\endcsname{ega4.c04.eq.21-1-6-2}
\expandafter\def\csname sga@target@IV.21.1.6.3\endcsname{ega4.c04.eq.21-1-6-3}
\expandafter\def\csname sga@target@IV.21.1.6.4\endcsname{ega4.c04.eq.21-1-6-4}
\expandafter\def\csname sga@target@IV.21.1.7\endcsname{ega4.c04.statement.21-1-7}
\expandafter\def\csname sga@target@IV.21.1.7.1\endcsname{ega4.c04.statement.21-1-7-1}
\expandafter\def\csname sga@target@IV.21.1.8\endcsname{ega4.c04.prop.21-1-8}
\expandafter\def\csname sga@target@IV.21.1.9\endcsname{ega4.c04.cor.21-1-9}
\expandafter\def\csname sga@target@IV.21.10.1\endcsname{ega4.c04.statement.21-10-1}
\expandafter\def\csname sga@target@IV.21.10.1.1\endcsname{ega4.c04.eq.21-10-1-1}
\expandafter\def\csname sga@target@IV.21.10.1.2\endcsname{ega4.c04.eq.21-10-1-2}
\expandafter\def\csname sga@target@IV.21.10.10\endcsname{ega4.c04.statement.21-10-10}
\expandafter\def\csname sga@target@IV.21.10.11\endcsname{ega4.c04.statement.21-10-11}
\expandafter\def\csname sga@target@IV.21.10.11.1\endcsname{ega4.c04.eq.21-10-11-1}
\expandafter\def\csname sga@target@IV.21.10.12\endcsname{ega4.c04.rem.21-10-12}
\expandafter\def\csname sga@target@IV.21.10.13\endcsname{ega4.c04.lem.21-10-13}
\expandafter\def\csname sga@target@IV.21.10.13.1\endcsname{ega4.c04.prop.21-10-13-1}
\expandafter\def\csname sga@target@IV.21.10.13bis\endcsname{ega4.c04.prop.21-10-13bis}
\expandafter\def\csname sga@target@IV.21.10.14\endcsname{ega4.c04.statement.21-10-14}
\expandafter\def\csname sga@target@IV.21.10.14.1\endcsname{ega4.c04.eq.21-10-14-1}
\expandafter\def\csname sga@target@IV.21.10.14.2\endcsname{ega4.c04.eq.21-10-14-2}
\expandafter\def\csname sga@target@IV.21.10.15\endcsname{ega4.c04.prop.21-10-15}
\expandafter\def\csname sga@target@IV.21.10.16\endcsname{ega4.c04.prop.21-10-16}
\expandafter\def\csname sga@target@IV.21.10.16.1\endcsname{ega4.c04.prop.21-10-16-1}
\expandafter\def\csname sga@target@IV.21.10.16.2\endcsname{ega4.c04.eq.21-10-16-2}
\expandafter\def\csname sga@target@IV.21.10.17\endcsname{ega4.c04.prop.21-10-17}
\expandafter\def\csname sga@target@IV.21.10.17.1\endcsname{ega4.c04.prop.21-10-17-1}
\expandafter\def\csname sga@target@IV.21.10.17.10\endcsname{ega4.c04.eq.21-10-17-10}
\expandafter\def\csname sga@target@IV.21.10.17.2\endcsname{ega4.c04.eq.21-10-17-2}
\expandafter\def\csname sga@target@IV.21.10.17.3\endcsname{ega4.c04.lem.21-10-17-3}
\expandafter\def\csname sga@target@IV.21.10.17.4\endcsname{ega4.c04.lem.21-10-17-4}
\expandafter\def\csname sga@target@IV.21.10.17.5\endcsname{ega4.c04.eq.21-10-17-5}
\expandafter\def\csname sga@target@IV.21.10.17.6\endcsname{ega4.c04.eq.21-10-17-6}
\expandafter\def\csname sga@target@IV.21.10.17.7\endcsname{ega4.c04.lem.21-10-17-7}
\expandafter\def\csname sga@target@IV.21.10.17.8\endcsname{ega4.c04.lem.21-10-17-8}
\expandafter\def\csname sga@target@IV.21.10.17.9\endcsname{ega4.c04.eq.21-10-17-9}
\expandafter\def\csname sga@target@IV.21.10.18\endcsname{ega4.c04.prop.21-10-18}
\expandafter\def\csname sga@target@IV.21.10.18.1\endcsname{ega4.c04.prop.21-10-18-1}
\expandafter\def\csname sga@target@IV.21.10.19\endcsname{ega4.c04.cor.21-10-19}
\expandafter\def\csname sga@target@IV.21.10.19.1\endcsname{ega4.c04.cor.21-10-19-1}
\expandafter\def\csname sga@target@IV.21.10.2\endcsname{ega4.c04.statement.21-10-2}
\expandafter\def\csname sga@target@IV.21.10.3\endcsname{ega4.c04.statement.21-10-3}
\expandafter\def\csname sga@target@IV.21.10.3.1\endcsname{ega4.c04.statement.21-10-3-1}
\expandafter\def\csname sga@target@IV.21.10.4\endcsname{ega4.c04.prop.21-10-4}
\expandafter\def\csname sga@target@IV.21.10.4.1\endcsname{ega4.c04.prop.21-10-4-1}
\expandafter\def\csname sga@target@IV.21.10.5\endcsname{ega4.c04.statement.21-10-5}
\expandafter\def\csname sga@target@IV.21.10.6\endcsname{ega4.c04.prop.21-10-6}
\expandafter\def\csname sga@target@IV.21.10.7\endcsname{ega4.c04.rem.21-10-7}
\expandafter\def\csname sga@target@IV.21.10.8\endcsname{ega4.c04.prop.21-10-8}
\expandafter\def\csname sga@target@IV.21.10.9\endcsname{ega4.c04.statement.21-10-9}
\expandafter\def\csname sga@target@IV.21.11\endcsname{ega4.c04.subsec.21-11.key-iv-21-11}
\expandafter\def\csname sga@target@IV.21.11.1\endcsname{ega4.c04.thm.21-11-1}
\expandafter\def\csname sga@target@IV.21.11.1.1\endcsname{ega4.c04.eq.21-11-1-1}
\expandafter\def\csname sga@target@IV.21.11.2\endcsname{ega4.c04.lem.21-11-2}
\expandafter\def\csname sga@target@IV.21.12\endcsname{ega4.c04.subsec.21-12.key-iv-21-12}
\expandafter\def\csname sga@target@IV.21.12.1\endcsname{ega4.c04.statement.21-12-1}
\expandafter\def\csname sga@target@IV.21.12.1.1\endcsname{ega4.c04.eq.21-12-1-1}
\expandafter\def\csname sga@target@IV.21.12.10\endcsname{ega4.c04.prop.21-12-10}
\expandafter\def\csname sga@target@IV.21.12.12\endcsname{ega4.c04.thm.21-12-12}
\expandafter\def\csname sga@target@IV.21.12.13\endcsname{ega4.c04.cor.21-12-13}
\expandafter\def\csname sga@target@IV.21.12.14\endcsname{ega4.c04.item.21-12-14}
\expandafter\def\csname sga@target@IV.21.12.15\endcsname{ega4.c04.prop.21-12-15}
\expandafter\def\csname sga@target@IV.21.12.16\endcsname{ega4.c04.cor.21-12-16}
\expandafter\def\csname sga@target@IV.21.12.17\endcsname{ega4.c04.item.21-12-17}
\expandafter\def\csname sga@target@IV.21.12.2\endcsname{ega4.c04.prop.21-12-2}
\expandafter\def\csname sga@target@IV.21.12.2.1\endcsname{ega4.c04.prop.21-12-2-1}
\expandafter\def\csname sga@target@IV.21.12.3\endcsname{ega4.c04.cor.21-12-3}
\expandafter\def\csname sga@target@IV.21.12.4\endcsname{ega4.c04.eq.21-12-4}
\expandafter\def\csname sga@target@IV.21.12.5\endcsname{ega4.c04.statement.21-12-5}
\expandafter\def\csname sga@target@IV.21.12.5.1\endcsname{ega4.c04.eq.21-12-5-1}
\expandafter\def\csname sga@target@IV.21.12.5.2\endcsname{ega4.c04.eq.21-12-5-2}
\expandafter\def\csname sga@target@IV.21.12.5.3\endcsname{ega4.c04.eq.21-12-5-3}
\expandafter\def\csname sga@target@IV.21.12.6\endcsname{ega4.c04.thm.21-12-6}
\expandafter\def\csname sga@target@IV.21.12.7\endcsname{ega4.c04.cor.21-12-7}
\expandafter\def\csname sga@target@IV.21.12.8\endcsname{ega4.c04.statement.21-12-8}
\expandafter\def\csname sga@target@IV.21.12.9\endcsname{ega4.c04.statement.21-12-9}
\expandafter\def\csname sga@target@IV.21.13\endcsname{ega4.c04.subsec.21-13.key-iv-21-13}
\expandafter\def\csname sga@target@IV.21.13.1\endcsname{ega4.c04.def.21-13-1}
\expandafter\def\csname sga@target@IV.21.13.10\endcsname{ega4.c04.prop.21-13-10}
\expandafter\def\csname sga@target@IV.21.13.11\endcsname{ega4.c04.cor.21-13-11}
\expandafter\def\csname sga@target@IV.21.13.12\endcsname{ega4.c04.prop.21-13-12}
\expandafter\def\csname sga@target@IV.21.13.12.1\endcsname{ega4.c04.statement.21-13-12-1}
\expandafter\def\csname sga@target@IV.21.13.13\endcsname{ega4.c04.rem.21-13-13}
\expandafter\def\csname sga@target@IV.21.13.14\endcsname{ega4.c04.prop.21-13-14}
\expandafter\def\csname sga@target@IV.21.13.15\endcsname{ega4.c04.cor.21-13-15}
\expandafter\def\csname sga@target@IV.21.13.16\endcsname{ega4.c04.rem.21-13-16}
\expandafter\def\csname sga@target@IV.21.13.2\endcsname{ega4.c04.lem.21-13-2}
\expandafter\def\csname sga@target@IV.21.13.2.1\endcsname{ega4.c04.eq.21-13-2-1}
\expandafter\def\csname sga@target@IV.21.13.3\endcsname{ega4.c04.cor.21-13-3}
\expandafter\def\csname sga@target@IV.21.13.4\endcsname{ega4.c04.lem.21-13-4}
\expandafter\def\csname sga@target@IV.21.13.5\endcsname{ega4.c04.prop.21-13-5}
\expandafter\def\csname sga@target@IV.21.13.6\endcsname{ega4.c04.cor.21-13-6}
\expandafter\def\csname sga@target@IV.21.13.7\endcsname{ega4.c04.def.21-13-7}
\expandafter\def\csname sga@target@IV.21.13.8\endcsname{ega4.c04.prop.21-13-8}
\expandafter\def\csname sga@target@IV.21.13.9\endcsname{ega4.c04.statement.21-13-9}
\expandafter\def\csname sga@target@IV.21.13.9.1\endcsname{ega4.c04.eq.21-13-9-1}
\expandafter\def\csname sga@target@IV.21.14\endcsname{ega4.c04.subsec.21-14.key-iv-21-14}
\expandafter\def\csname sga@target@IV.21.14.1\endcsname{ega4.c04.thm.21-14-1}
\expandafter\def\csname sga@target@IV.21.14.2\endcsname{ega4.c04.cor.21-14-2}
\expandafter\def\csname sga@target@IV.21.14.3\endcsname{ega4.c04.prop.21-14-3}
\expandafter\def\csname sga@target@IV.21.14.4\endcsname{ega4.c04.statement.21-14-4}
\expandafter\def\csname sga@target@IV.21.14.4bis\endcsname{ega4.c04.statement.21-14-4bis}
\expandafter\def\csname sga@target@IV.21.15\endcsname{ega4.c04.subsec.21-15.key-iv-21-15}
\expandafter\def\csname sga@target@IV.21.15.1\endcsname{ega4.c04.statement.21-15-1}
\expandafter\def\csname sga@target@IV.21.15.2\endcsname{ega4.c04.def.21-15-2}
\expandafter\def\csname sga@target@IV.21.15.3\endcsname{ega4.c04.statement.21-15-3}
\expandafter\def\csname sga@target@IV.21.15.3.1\endcsname{ega4.c04.eq.21-15-3-1}
\expandafter\def\csname sga@target@IV.21.15.3.2\endcsname{ega4.c04.eq.21-15-3-2}
\expandafter\def\csname sga@target@IV.21.15.3bis\endcsname{ega4.c04.statement.21-15-3bis}
\expandafter\def\csname sga@target@IV.21.15.4\endcsname{ega4.c04.prop.21-15-4}
\expandafter\def\csname sga@target@IV.21.15.5\endcsname{ega4.c04.prop.21-15-5}
\expandafter\def\csname sga@target@IV.21.15.6\endcsname{ega4.c04.prop.21-15-6}
\expandafter\def\csname sga@target@IV.21.15.7\endcsname{ega4.c04.statement.21-15-7}
\expandafter\def\csname sga@target@IV.21.15.8\endcsname{ega4.c04.statement.21-15-8}
\expandafter\def\csname sga@target@IV.21.15.9\endcsname{ega4.c04.statement.21-15-9}
\expandafter\def\csname sga@target@IV.21.15.9.1\endcsname{ega4.c04.eq.21-15-9-1}
\expandafter\def\csname sga@target@IV.21.2.1\endcsname{ega4.c04.statement.21-2-1}
\expandafter\def\csname sga@target@IV.21.2.10\endcsname{ega4.c04.statement.21-2-10}
\expandafter\def\csname sga@target@IV.21.2.10.1\endcsname{ega4.c04.eq.21-2-10-1}
\expandafter\def\csname sga@target@IV.21.2.11\endcsname{ega4.c04.prop.21-2-11}
\expandafter\def\csname sga@target@IV.21.2.11.1\endcsname{ega4.c04.prop.21-2-11-1}
\expandafter\def\csname sga@target@IV.21.2.12\endcsname{ega4.c04.statement.21-2-12}
\expandafter\def\csname sga@target@IV.21.2.2\endcsname{ega4.c04.prop.21-2-2}
\expandafter\def\csname sga@target@IV.21.2.3\endcsname{ega4.c04.cor.21-2-3}
\expandafter\def\csname sga@target@IV.21.2.4\endcsname{ega4.c04.statement.21-2-4}
\expandafter\def\csname sga@target@IV.21.2.5\endcsname{ega4.c04.statement.21-2-5}
\expandafter\def\csname sga@target@IV.21.2.5.1\endcsname{ega4.c04.statement.21-2-5-1}
\expandafter\def\csname sga@target@IV.21.2.5.2\endcsname{ega4.c04.eq.21-2-5-2}
\expandafter\def\csname sga@target@IV.21.2.5.3\endcsname{ega4.c04.eq.21-2-5-3}
\expandafter\def\csname sga@target@IV.21.2.5.4\endcsname{ega4.c04.eq.21-2-5-4}
\expandafter\def\csname sga@target@IV.21.2.6\endcsname{ega4.c04.prop.21-2-6}
\expandafter\def\csname sga@target@IV.21.2.6.1\endcsname{ega4.c04.eq.21-2-6-1}
\expandafter\def\csname sga@target@IV.21.2.7\endcsname{ega4.c04.statement.21-2-7}
\expandafter\def\csname sga@target@IV.21.2.7.1\endcsname{ega4.c04.statement.21-2-7-1}
\expandafter\def\csname sga@target@IV.21.2.7.2\endcsname{ega4.c04.eq.21-2-7-2}
\expandafter\def\csname sga@target@IV.21.2.8\endcsname{ega4.c04.statement.21-2-8}
\expandafter\def\csname sga@target@IV.21.2.8.1\endcsname{ega4.c04.statement.21-2-8-1}
\expandafter\def\csname sga@target@IV.21.2.8.2\endcsname{ega4.c04.statement.21-2-8-2}
\expandafter\def\csname sga@target@IV.21.2.8.3\endcsname{ega4.c04.statement.21-2-8-3}
\expandafter\def\csname sga@target@IV.21.2.9\endcsname{ega4.c04.statement.21-2-9}
\expandafter\def\csname sga@target@IV.21.2.9.1\endcsname{ega4.c04.statement.21-2-9-1}
\expandafter\def\csname sga@target@IV.21.2.9.2\endcsname{ega4.c04.eq.21-2-9-2}
\expandafter\def\csname sga@target@IV.21.2.9.3\endcsname{ega4.c04.eq.21-2-9-3}
\expandafter\def\csname sga@target@IV.21.2.9.4\endcsname{ega4.c04.eq.21-2-9-4}
\expandafter\def\csname sga@target@IV.21.3.1\endcsname{ega4.c04.statement.21-3-1}
\expandafter\def\csname sga@target@IV.21.3.2\endcsname{ega4.c04.statement.21-3-2}
\expandafter\def\csname sga@target@IV.21.3.2.1\endcsname{ega4.c04.statement.21-3-2-1}
\expandafter\def\csname sga@target@IV.21.3.2.2\endcsname{ega4.c04.eq.21-3-2-2}
\expandafter\def\csname sga@target@IV.21.3.2.3\endcsname{ega4.c04.eq.21-3-2-3}
\expandafter\def\csname sga@target@IV.21.3.2.4\endcsname{ega4.c04.eq.21-3-2-4}
\expandafter\def\csname sga@target@IV.21.3.3\endcsname{ega4.c04.prop.21-3-3}
\expandafter\def\csname sga@target@IV.21.3.3.1\endcsname{ega4.c04.eq.21-3-3-1}
\expandafter\def\csname sga@target@IV.21.3.4\endcsname{ega4.c04.prop.21-3-4}
\expandafter\def\csname sga@target@IV.21.3.5\endcsname{ega4.c04.cor.21-3-5}
\expandafter\def\csname sga@target@IV.21.3.6\endcsname{ega4.c04.rem.21-3-6}
\expandafter\def\csname sga@target@IV.21.3.6.1\endcsname{ega4.c04.eq.21-3-6-1}
\expandafter\def\csname sga@target@IV.21.4\endcsname{ega4.c04.subsec.21-4.key-iv-21-4}
\expandafter\def\csname sga@target@IV.21.4.1\endcsname{ega4.c04.statement.21-4-1}
\expandafter\def\csname sga@target@IV.21.4.10\endcsname{ega4.c04.cor.21-4-10}
\expandafter\def\csname sga@target@IV.21.4.11\endcsname{ega4.c04.cor.21-4-11}
\expandafter\def\csname sga@target@IV.21.4.12\endcsname{ega4.c04.rem.21-4-12}
\expandafter\def\csname sga@target@IV.21.4.13\endcsname{ega4.c04.cor.21-4-13}
\expandafter\def\csname sga@target@IV.21.4.2\endcsname{ega4.c04.def.21-4-2}
\expandafter\def\csname sga@target@IV.21.4.2.1\endcsname{ega4.c04.eq.21-4-2-1}
\expandafter\def\csname sga@target@IV.21.4.3\endcsname{ega4.c04.statement.21-4-3}
\expandafter\def\csname sga@target@IV.21.4.4\endcsname{ega4.c04.statement.21-4-4}
\expandafter\def\csname sga@target@IV.21.4.5\endcsname{ega4.c04.prop.21-4-5}
\expandafter\def\csname sga@target@IV.21.4.6\endcsname{ega4.c04.cor.21-4-6}
\expandafter\def\csname sga@target@IV.21.4.6.1\endcsname{ega4.c04.eq.21-4-6-1}
\expandafter\def\csname sga@target@IV.21.4.7\endcsname{ega4.c04.statement.21-4-7}
\expandafter\def\csname sga@target@IV.21.4.8\endcsname{ega4.c04.prop.21-4-8}
\expandafter\def\csname sga@target@IV.21.4.9\endcsname{ega4.c04.prop.21-4-9}
\expandafter\def\csname sga@target@IV.21.5\endcsname{ega4.c04.subsec.21-5.key-iv-21-5}
\expandafter\def\csname sga@target@IV.21.5.1\endcsname{ega4.c04.statement.21-5-1}
\expandafter\def\csname sga@target@IV.21.5.2\endcsname{ega4.c04.lem.21-5-2}
\expandafter\def\csname sga@target@IV.21.5.3\endcsname{ega4.c04.statement.21-5-3}
\expandafter\def\csname sga@target@IV.21.5.3.1\endcsname{ega4.c04.eq.21-5-3-1}
\expandafter\def\csname sga@target@IV.21.5.3.2\endcsname{ega4.c04.eq.21-5-3-2}
\expandafter\def\csname sga@target@IV.21.5.4\endcsname{ega4.c04.statement.21-5-4}
\expandafter\def\csname sga@target@IV.21.5.5\endcsname{ega4.c04.def.21-5-5}
\expandafter\def\csname sga@target@IV.21.5.5.1\endcsname{ega4.c04.eq.21-5-5-1}
\expandafter\def\csname sga@target@IV.21.5.5.2\endcsname{ega4.c04.eq.21-5-5-2}
\expandafter\def\csname sga@target@IV.21.5.5.3\endcsname{ega4.c04.eq.21-5-5-3}
\expandafter\def\csname sga@target@IV.21.5.6\endcsname{ega4.c04.prop.21-5-6}
\expandafter\def\csname sga@target@IV.21.5.6.1\endcsname{ega4.c04.prop.21-5-6-1}
\expandafter\def\csname sga@target@IV.21.5.7\endcsname{ega4.c04.prop.21-5-7}
\expandafter\def\csname sga@target@IV.21.5.7.1\endcsname{ega4.c04.prop.21-5-7-1}
\expandafter\def\csname sga@target@IV.21.5.7.2\endcsname{ega4.c04.lem.21-5-7-2}
\expandafter\def\csname sga@target@IV.21.5.7.3\endcsname{ega4.c04.lem.21-5-7-3}
\expandafter\def\csname sga@target@IV.21.5.7.4\endcsname{ega4.c04.eq.21-5-7-4}
\expandafter\def\csname sga@target@IV.21.5.8\endcsname{ega4.c04.prop.21-5-8}
\expandafter\def\csname sga@target@IV.21.5.8.1\endcsname{ega4.c04.eq.21-5-8-1}
\expandafter\def\csname sga@target@IV.21.6\endcsname{ega4.c04.subsec.21-6.key-iv-21-6}
\expandafter\def\csname sga@target@IV.21.6.1\endcsname{ega4.c04.statement.21-6-1}
\expandafter\def\csname sga@target@IV.21.6.10\endcsname{ega4.c04.cor.21-6-10}
\expandafter\def\csname sga@target@IV.21.6.10.1\endcsname{ega4.c04.eq.21-6-10-1}
\expandafter\def\csname sga@target@IV.21.6.11\endcsname{ega4.c04.cor.21-6-11}
\expandafter\def\csname sga@target@IV.21.6.12\endcsname{ega4.c04.prop.21-6-12}
\expandafter\def\csname sga@target@IV.21.6.12.1\endcsname{ega4.c04.eq.21-6-12-1}
\expandafter\def\csname sga@target@IV.21.6.12.2\endcsname{ega4.c04.eq.21-6-12-2}
\expandafter\def\csname sga@target@IV.21.6.13\endcsname{ega4.c04.cor.21-6-13}
\expandafter\def\csname sga@target@IV.21.6.14\endcsname{ega4.c04.cor.21-6-14}
\expandafter\def\csname sga@target@IV.21.6.2\endcsname{ega4.c04.statement.21-6-2}
\expandafter\def\csname sga@target@IV.21.6.3\endcsname{ega4.c04.statement.21-6-3}
\expandafter\def\csname sga@target@IV.21.6.4\endcsname{ega4.c04.statement.21-6-4}
\expandafter\def\csname sga@target@IV.21.6.4.1\endcsname{ega4.c04.statement.21-6-4-1}
\expandafter\def\csname sga@target@IV.21.6.4.2\endcsname{ega4.c04.eq.21-6-4-2}
\expandafter\def\csname sga@target@IV.21.6.5\endcsname{ega4.c04.statement.21-6-5}
\expandafter\def\csname sga@target@IV.21.6.5.1\endcsname{ega4.c04.eq.21-6-5-1}
\expandafter\def\csname sga@target@IV.21.6.6\endcsname{ega4.c04.prop.21-6-6}
\expandafter\def\csname sga@target@IV.21.6.6.2\endcsname{ega4.c04.statement.21-6-6-2}
\expandafter\def\csname sga@target@IV.21.6.7\endcsname{ega4.c04.statement.21-6-7}
\expandafter\def\csname sga@target@IV.21.6.7.1\endcsname{ega4.c04.eq.21-6-7-1}
\expandafter\def\csname sga@target@IV.21.6.7.2\endcsname{ega4.c04.eq.21-6-7-2}
\expandafter\def\csname sga@target@IV.21.6.7.3\endcsname{ega4.c04.eq.21-6-7-3}
\expandafter\def\csname sga@target@IV.21.6.7.4\endcsname{ega4.c04.eq.21-6-7-4}
\expandafter\def\csname sga@target@IV.21.6.7.5\endcsname{ega4.c04.eq.21-6-7-5}
\expandafter\def\csname sga@target@IV.21.6.8\endcsname{ega4.c04.statement.21-6-8}
\expandafter\def\csname sga@target@IV.21.6.9\endcsname{ega4.c04.thm.21-6-9}
\expandafter\def\csname sga@target@IV.21.6.9.1\endcsname{ega4.c04.eq.21-6-9-1}
\expandafter\def\csname sga@target@IV.21.6.9.2\endcsname{ega4.c04.statement.21-6-9-2}
\expandafter\def\csname sga@target@IV.21.7\endcsname{ega4.c04.subsec.21-7.key-iv-21-7}
\expandafter\def\csname sga@target@IV.21.7.1\endcsname{ega4.c04.statement.21-7-1}
\expandafter\def\csname sga@target@IV.21.7.2\endcsname{ega4.c04.prop.21-7-2}
\expandafter\def\csname sga@target@IV.21.7.2.1\endcsname{ega4.c04.eq.21-7-2-1}
\expandafter\def\csname sga@target@IV.21.7.3\endcsname{ega4.c04.cor.21-7-3}
\expandafter\def\csname sga@target@IV.21.7.3.1\endcsname{ega4.c04.statement.21-7-3-1}
\expandafter\def\csname sga@target@IV.21.7.4\endcsname{ega4.c04.prop.21-7-4}
\expandafter\def\csname sga@target@IV.21.7.5\endcsname{ega4.c04.rem.21-7-5}
\expandafter\def\csname sga@target@IV.21.7.6\endcsname{ega4.c04.cor.21-7-6}
\expandafter\def\csname sga@target@IV.21.7.7\endcsname{ega4.c04.rem.21-7-7}
\expandafter\def\csname sga@target@IV.21.8.1\endcsname{ega4.c04.lem.21-8-1}
\expandafter\def\csname sga@target@IV.21.8.2\endcsname{ega4.c04.cor.21-8-2}
\expandafter\def\csname sga@target@IV.21.8.3\endcsname{ega4.c04.prop.21-8-3}
\expandafter\def\csname sga@target@IV.21.8.3.1\endcsname{ega4.c04.eq.21-8-3-1}
\expandafter\def\csname sga@target@IV.21.8.4\endcsname{ega4.c04.cor.21-8-4}
\expandafter\def\csname sga@target@IV.21.8.5\endcsname{ega4.c04.prop.21-8-5}
\expandafter\def\csname sga@target@IV.21.8.5.1\endcsname{ega4.c04.eq.21-8-5-1}
\expandafter\def\csname sga@target@IV.21.8.5.2\endcsname{ega4.c04.eq.21-8-5-2}
\expandafter\def\csname sga@target@IV.21.8.5.3\endcsname{ega4.c04.eq.21-8-5-3}
\expandafter\def\csname sga@target@IV.21.8.6\endcsname{ega4.c04.rem.21-8-6}
\expandafter\def\csname sga@target@IV.21.9\endcsname{ega4.c04.subsec.21-9.key-iv-21-9}
\expandafter\def\csname sga@target@IV.21.9.1\endcsname{ega4.c04.statement.21-9-1}
\expandafter\def\csname sga@target@IV.21.9.1.1\endcsname{ega4.c04.eq.21-9-1-1}
\expandafter\def\csname sga@target@IV.21.9.1.2\endcsname{ega4.c04.eq.21-9-1-2}
\expandafter\def\csname sga@target@IV.21.9.10\endcsname{ega4.c04.statement.21-9-10}
\expandafter\def\csname sga@target@IV.21.9.11\endcsname{ega4.c04.prop.21-9-11}
\expandafter\def\csname sga@target@IV.21.9.12\endcsname{ega4.c04.cor.21-9-12}
\expandafter\def\csname sga@target@IV.21.9.13\endcsname{ega4.c04.rem.21-9-13}
\expandafter\def\csname sga@target@IV.21.9.2\endcsname{ega4.c04.lem.21-9-2}
\expandafter\def\csname sga@target@IV.21.9.3\endcsname{ega4.c04.rem.21-9-3}
\expandafter\def\csname sga@target@IV.21.9.4\endcsname{ega4.c04.prop.21-9-4}
\expandafter\def\csname sga@target@IV.21.9.4.1\endcsname{ega4.c04.prop.21-9-4-1}
\expandafter\def\csname sga@target@IV.21.9.5\endcsname{ega4.c04.cor.21-9-5}
\expandafter\def\csname sga@target@IV.21.9.6\endcsname{ega4.c04.cor.21-9-6}
\expandafter\def\csname sga@target@IV.21.9.7\endcsname{ega4.c04.cor.21-9-7}
\expandafter\def\csname sga@target@IV.21.9.8\endcsname{ega4.c04.statement.21-9-8}
\expandafter\def\csname sga@target@IV.21.9.9\endcsname{ega4.c04.prop.21-9-9}
\expandafter\def\csname sga@target@IV.3.1.1\endcsname{ega4.c04.def.3-1-1}
\expandafter\def\csname sga@target@IV.3.1.1.1\endcsname{ega4.c04.eq.3-1-1-1}
\expandafter\def\csname sga@target@IV.3.1.10\endcsname{ega4.c04.prop.3-1-10}
\expandafter\def\csname sga@target@IV.3.1.10.1\endcsname{ega4.c04.lem.3-1-10-1}
\expandafter\def\csname sga@target@IV.3.1.10.2\endcsname{ega4.c04.rem.3-1-10-2}
\expandafter\def\csname sga@target@IV.3.1.10.3\endcsname{ega4.c04.eq.3-1-10-3}
\expandafter\def\csname sga@target@IV.3.1.10.4\endcsname{ega4.c04.eq.3-1-10-4}
\expandafter\def\csname sga@target@IV.3.1.10.5\endcsname{ega4.c04.rem.3-1-10-5}
\expandafter\def\csname sga@target@IV.3.1.11\endcsname{ega4.c04.cor.3-1-11}
\expandafter\def\csname sga@target@IV.3.1.12\endcsname{ega4.c04.rem.3-1-12}
\expandafter\def\csname sga@target@IV.3.1.13\endcsname{ega4.c04.prop.3-1-13}
\expandafter\def\csname sga@target@IV.3.1.2\endcsname{ega4.c04.prop.3-1-2}
\expandafter\def\csname sga@target@IV.3.1.3\endcsname{ega4.c04.prop.3-1-3}
\expandafter\def\csname sga@target@IV.3.1.4\endcsname{ega4.c04.cor.3-1-4}
\expandafter\def\csname sga@target@IV.3.1.5\endcsname{ega4.c04.cor.3-1-5}
\expandafter\def\csname sga@target@IV.3.1.6\endcsname{ega4.c04.prop.3-1-6}
\expandafter\def\csname sga@target@IV.3.1.7\endcsname{ega4.c04.prop.3-1-7}
\expandafter\def\csname sga@target@IV.3.1.8\endcsname{ega4.c04.prop.3-1-8}
\expandafter\def\csname sga@target@IV.3.1.9\endcsname{ega4.c04.cor.3-1-9}
\expandafter\def\csname sga@target@IV.3.2.1\endcsname{ega4.c04.prop.3-2-1}
\expandafter\def\csname sga@target@IV.3.2.2\endcsname{ega4.c04.def.3-2-2}
\expandafter\def\csname sga@target@IV.3.2.3\endcsname{ega4.c04.prop.3-2-3}
\expandafter\def\csname sga@target@IV.3.2.4\endcsname{ega4.c04.def.3-2-4}
\expandafter\def\csname sga@target@IV.3.2.5\endcsname{ega4.c04.def.3-2-5}
\expandafter\def\csname sga@target@IV.3.2.6\endcsname{ega4.c04.prop.3-2-6}
\expandafter\def\csname sga@target@IV.3.2.7\endcsname{ega4.c04.cor.3-2-7}
\expandafter\def\csname sga@target@IV.3.2.8\endcsname{ega4.c04.cor.3-2-8}
\expandafter\def\csname sga@target@IV.3.2.9\endcsname{ega4.c04.rem.3-2-9}
\expandafter\def\csname sga@target@IV.3.3.1\endcsname{ega4.c04.prop.3-3-1}
\expandafter\def\csname sga@target@IV.3.3.1.1\endcsname{ega4.c04.prop.3-3-1-1}
\expandafter\def\csname sga@target@IV.3.3.2\endcsname{ega4.c04.cor.3-3-2}
\expandafter\def\csname sga@target@IV.3.3.3\endcsname{ega4.c04.cor.3-3-3}
\expandafter\def\csname sga@target@IV.3.3.4\endcsname{ega4.c04.prop.3-3-4}
\expandafter\def\csname sga@target@IV.3.3.5\endcsname{ega4.c04.cor.3-3-5}
\expandafter\def\csname sga@target@IV.3.3.6\endcsname{ega4.c04.prop.3-3-6}
\expandafter\def\csname sga@target@IV.3.3.6.1\endcsname{ega4.c04.eq.3-3-6-1}
\expandafter\def\csname sga@target@IV.3.3.6.2\endcsname{ega4.c04.eq.3-3-6-2}
\expandafter\def\csname sga@target@IV.3.3.6.3\endcsname{ega4.c04.eq.3-3-6-3}
\expandafter\def\csname sga@target@IV.3.3.6.4\endcsname{ega4.c04.eq.3-3-6-4}
\expandafter\def\csname sga@target@IV.3.3.7\endcsname{ega4.c04.cor.3-3-7}
\expandafter\def\csname sga@target@IV.3.3.7.1\endcsname{ega4.c04.cor.3-3-7-1}
\expandafter\def\csname sga@target@IV.3.3.8\endcsname{ega4.c04.sec.3-3-8}
\expandafter\def\csname sga@target@IV.3.3.9\endcsname{ega4.c04.cor.3-3-9}
\expandafter\def\csname sga@target@IV.3.3.9.1\endcsname{ega4.c04.cor.3-3-9-1}
\expandafter\def\csname sga@target@IV.3.4\endcsname{ega4.c04.subsec.3-4.key-iv-3-4}
\expandafter\def\csname sga@target@IV.3.4.1\endcsname{ega4.c04.prop.3-4-1}
\expandafter\def\csname sga@target@IV.3.4.1.1\endcsname{ega4.c04.eq.3-4-1-1}
\expandafter\def\csname sga@target@IV.3.4.1.2\endcsname{ega4.c04.lem.3-4-1-2}
\expandafter\def\csname sga@target@IV.3.4.1.3\endcsname{ega4.c04.lem.3-4-1-3}
\expandafter\def\csname sga@target@IV.3.4.1.4\endcsname{ega4.c04.lem.3-4-1-4}
\expandafter\def\csname sga@target@IV.3.4.1.5\endcsname{ega4.c04.eq.3-4-1-5}
\expandafter\def\csname sga@target@IV.3.4.2\endcsname{ega4.c04.cor.3-4-2}
\expandafter\def\csname sga@target@IV.3.4.3\endcsname{ega4.c04.prop.3-4-3}
\expandafter\def\csname sga@target@IV.3.4.4\endcsname{ega4.c04.prop.3-4-4}
\expandafter\def\csname sga@target@IV.3.4.5\endcsname{ega4.c04.prop.3-4-5}
\expandafter\def\csname sga@target@IV.3.4.6\endcsname{ega4.c04.prop.3-4-6}
\expandafter\def\csname sga@target@IV.3.4.6.1\endcsname{ega4.c04.lem.3-4-6-1}
\expandafter\def\csname sga@target@IV.3.4.7\endcsname{ega4.c04.cor.3-4-7}
\expandafter\def\csname sga@target@IV.3.4.8\endcsname{ega4.c04.rem.3-4-8}
\expandafter\def\csname sga@target@IV.3.4.9\endcsname{ega4.c04.prop.3-4-9}
\expandafter\def\csname sga@target@IV.4.1.1\endcsname{ega4.c04.def.4-1-1}
\expandafter\def\csname sga@target@IV.4.1.1.1\endcsname{ega4.c04.eq.4-1-1-1}
\expandafter\def\csname sga@target@IV.4.1.1.2\endcsname{ega4.c04.eq.4-1-1-2}
\expandafter\def\csname sga@target@IV.4.1.1.3\endcsname{ega4.c04.eq.4-1-1-3}
\expandafter\def\csname sga@target@IV.4.1.2\endcsname{ega4.c04.thm.4-1-2}
\expandafter\def\csname sga@target@IV.4.1.2.1\endcsname{ega4.c04.cor.4-1-2-1}
\expandafter\def\csname sga@target@IV.4.1.3\endcsname{ega4.c04.rem.4-1-3}
\expandafter\def\csname sga@target@IV.4.1.4\endcsname{ega4.c04.cor.4-1-4}
\expandafter\def\csname sga@target@IV.4.1.5\endcsname{ega4.c04.cor.4-1-5}
\expandafter\def\csname sga@target@IV.4.2.1\endcsname{ega4.c04.prop.4-2-1}
\expandafter\def\csname sga@target@IV.4.2.1.1\endcsname{ega4.c04.eq.4-2-1-1}
\expandafter\def\csname sga@target@IV.4.2.1.2\endcsname{ega4.c04.eq.4-2-1-2}
\expandafter\def\csname sga@target@IV.4.2.1.3\endcsname{ega4.c04.eq.4-2-1-3}
\expandafter\def\csname sga@target@IV.4.2.1.4\endcsname{ega4.c04.rem.4-2-1-4}
\expandafter\def\csname sga@target@IV.4.2.1.5\endcsname{ega4.c04.eq.4-2-1-5}
\expandafter\def\csname sga@target@IV.4.2.10\endcsname{ega4.c04.cor.4-2-10}
\expandafter\def\csname sga@target@IV.4.2.2\endcsname{ega4.c04.cor.4-2-2}
\expandafter\def\csname sga@target@IV.4.2.3\endcsname{ega4.c04.cor.4-2-3}
\expandafter\def\csname sga@target@IV.4.2.4\endcsname{ega4.c04.prop.4-2-4}
\expandafter\def\csname sga@target@IV.4.2.4.i\endcsname{ega4.c04.sec.4-2-4-i}
\expandafter\def\csname sga@target@IV.4.2.4.ii\endcsname{ega4.c04.sec.4-2-4-ii}
\expandafter\def\csname sga@target@IV.4.2.4.iii\endcsname{ega4.c04.sec.4-2-4-iii}
\expandafter\def\csname sga@target@IV.4.2.5\endcsname{ega4.c04.prop.4-2-5}
\expandafter\def\csname sga@target@IV.4.2.6\endcsname{ega4.c04.cor.4-2-6}
\expandafter\def\csname sga@target@IV.4.2.7\endcsname{ega4.c04.prop.4-2-7}
\expandafter\def\csname sga@target@IV.4.2.7.i\endcsname{ega4.c04.sec.4-2-7-i}
\expandafter\def\csname sga@target@IV.4.2.7.ii\endcsname{ega4.c04.sec.4-2-7-ii}
\expandafter\def\csname sga@target@IV.4.2.7.iii\endcsname{ega4.c04.sec.4-2-7-iii}
\expandafter\def\csname sga@target@IV.4.2.7.iv\endcsname{ega4.c04.sec.4-2-7-iv}
\expandafter\def\csname sga@target@IV.4.2.8\endcsname{ega4.c04.cor.4-2-8}
\expandafter\def\csname sga@target@IV.4.2.9\endcsname{ega4.c04.prop.4-2-9}
\expandafter\def\csname sga@target@IV.4.3.1\endcsname{ega4.c04.statement.4-3-1}
\expandafter\def\csname sga@target@IV.4.3.2\endcsname{ega4.c04.prop.4-3-2}
\expandafter\def\csname sga@target@IV.4.3.3\endcsname{ega4.c04.cor.4-3-3}
\expandafter\def\csname sga@target@IV.4.3.4\endcsname{ega4.c04.cor.4-3-4}
\expandafter\def\csname sga@target@IV.4.3.5\endcsname{ega4.c04.prop.4-3-5}
\expandafter\def\csname sga@target@IV.4.3.6\endcsname{ega4.c04.cor.4-3-6}
\expandafter\def\csname sga@target@IV.4.3.7\endcsname{ega4.c04.cor.4-3-7}
\expandafter\def\csname sga@target@IV.4.3.8\endcsname{ega4.c04.cor.4-3-8}
\expandafter\def\csname sga@target@IV.4.3.9\endcsname{ega4.c04.prop.4-3-9}
\expandafter\def\csname sga@target@IV.4.4.1\endcsname{ega4.c04.statement.4-4-1}
\expandafter\def\csname sga@target@IV.4.4.2\endcsname{ega4.c04.lem.4-4-2}
\expandafter\def\csname sga@target@IV.4.4.2.i\endcsname{ega4.c04.sec.4-4-2-i}
\expandafter\def\csname sga@target@IV.4.4.2.ii\endcsname{ega4.c04.sec.4-4-2-ii}
\expandafter\def\csname sga@target@IV.4.4.3\endcsname{ega4.c04.rem.4-4-3}
\expandafter\def\csname sga@target@IV.4.4.4\endcsname{ega4.c04.thm.4-4-4}
\expandafter\def\csname sga@target@IV.4.4.5\endcsname{ega4.c04.cor.4-4-5}
\expandafter\def\csname sga@target@IV.4.4.6\endcsname{ega4.c04.cor.4-4-6}
\expandafter\def\csname sga@target@IV.4.5\endcsname{ega4.c04.subsec.4-5.key-iv-4-5}
\expandafter\def\csname sga@target@IV.4.5.1\endcsname{ega4.c04.prop.4-5-1}
\expandafter\def\csname sga@target@IV.4.5.10\endcsname{ega4.c04.cor.4-5-10}
\expandafter\def\csname sga@target@IV.4.5.11\endcsname{ega4.c04.cor.4-5-11}
\expandafter\def\csname sga@target@IV.4.5.12\endcsname{ega4.c04.sec.4-5-12}
\expandafter\def\csname sga@target@IV.4.5.12.i\endcsname{ega4.c04.sec.4-5-12-i}
\expandafter\def\csname sga@target@IV.4.5.12.ii\endcsname{ega4.c04.sec.4-5-12-ii}
\expandafter\def\csname sga@target@IV.4.5.13\endcsname{ega4.c04.prop.4-5-13}
\expandafter\def\csname sga@target@IV.4.5.13.1\endcsname{ega4.c04.cor.4-5-13-1}
\expandafter\def\csname sga@target@IV.4.5.13.1.i\endcsname{ega4.c04.sec.4-5-13-1-i}
\expandafter\def\csname sga@target@IV.4.5.13.1.ii\endcsname{ega4.c04.sec.4-5-13-1-ii}
\expandafter\def\csname sga@target@IV.4.5.14\endcsname{ega4.c04.cor.4-5-14}
\expandafter\def\csname sga@target@IV.4.5.15\endcsname{ega4.c04.prop.4-5-15}
\expandafter\def\csname sga@target@IV.4.5.15.i\endcsname{ega4.c04.sec.4-5-15-i}
\expandafter\def\csname sga@target@IV.4.5.15.ii\endcsname{ega4.c04.sec.4-5-15-ii}
\expandafter\def\csname sga@target@IV.4.5.16\endcsname{ega4.c04.cor.4-5-16}
\expandafter\def\csname sga@target@IV.4.5.16.i\endcsname{ega4.c04.sec.4-5-16-i}
\expandafter\def\csname sga@target@IV.4.5.16.ii\endcsname{ega4.c04.sec.4-5-16-ii}
\expandafter\def\csname sga@target@IV.4.5.17\endcsname{ega4.c04.cor.4-5-17}
\expandafter\def\csname sga@target@IV.4.5.18\endcsname{ega4.c04.rem.4-5-18}
\expandafter\def\csname sga@target@IV.4.5.19\endcsname{ega4.c04.prop.4-5-19}
\expandafter\def\csname sga@target@IV.4.5.19.1\endcsname{ega4.c04.lem.4-5-19-1}
\expandafter\def\csname sga@target@IV.4.5.19.2\endcsname{ega4.c04.eq.4-5-19-2}
\expandafter\def\csname sga@target@IV.4.5.19.3\endcsname{ega4.c04.cor.4-5-19-3}
\expandafter\def\csname sga@target@IV.4.5.2\endcsname{ega4.c04.def.4-5-2}
\expandafter\def\csname sga@target@IV.4.5.2.a\endcsname{ega4.c04.sec.4-5-2-a}
\expandafter\def\csname sga@target@IV.4.5.2.b\endcsname{ega4.c04.sec.4-5-2-b}
\expandafter\def\csname sga@target@IV.4.5.2.c\endcsname{ega4.c04.sec.4-5-2-c}
\expandafter\def\csname sga@target@IV.4.5.20\endcsname{ega4.c04.rem.4-5-20}
\expandafter\def\csname sga@target@IV.4.5.21\endcsname{ega4.c04.prop.4-5-21}
\expandafter\def\csname sga@target@IV.4.5.3\endcsname{ega4.c04.statement.4-5-3}
\expandafter\def\csname sga@target@IV.4.5.4\endcsname{ega4.c04.prop.4-5-4}
\expandafter\def\csname sga@target@IV.4.5.5\endcsname{ega4.c04.def.4-5-5}
\expandafter\def\csname sga@target@IV.4.5.6\endcsname{ega4.c04.prop.4-5-6}
\expandafter\def\csname sga@target@IV.4.5.6.i\endcsname{ega4.c04.sec.4-5-6-i}
\expandafter\def\csname sga@target@IV.4.5.6.ii\endcsname{ega4.c04.sec.4-5-6-ii}
\expandafter\def\csname sga@target@IV.4.5.7\endcsname{ega4.c04.prop.4-5-7}
\expandafter\def\csname sga@target@IV.4.5.8\endcsname{ega4.c04.cor.4-5-8}
\expandafter\def\csname sga@target@IV.4.5.8.i\endcsname{ega4.c04.sec.4-5-8-i}
\expandafter\def\csname sga@target@IV.4.5.8.ii\endcsname{ega4.c04.sec.4-5-8-ii}
\expandafter\def\csname sga@target@IV.4.5.9\endcsname{ega4.c04.prop.4-5-9}
\expandafter\def\csname sga@target@IV.4.5.9.a\endcsname{ega4.c04.sec.4-5-9-a}
\expandafter\def\csname sga@target@IV.4.5.9.b\endcsname{ega4.c04.sec.4-5-9-b}
\expandafter\def\csname sga@target@IV.4.5.9.c\endcsname{ega4.c04.sec.4-5-9-c}
\expandafter\def\csname sga@target@IV.4.6\endcsname{ega4.c04.subsec.4-6.key-iv-4-6}
\expandafter\def\csname sga@target@IV.4.6.1\endcsname{ega4.c04.prop.4-6-1}
\expandafter\def\csname sga@target@IV.4.6.1.a\endcsname{ega4.c04.sec.4-6-1-a}
\expandafter\def\csname sga@target@IV.4.6.1.b\endcsname{ega4.c04.sec.4-6-1-b}
\expandafter\def\csname sga@target@IV.4.6.1.c\endcsname{ega4.c04.sec.4-6-1-c}
\expandafter\def\csname sga@target@IV.4.6.1.d\endcsname{ega4.c04.sec.4-6-1-d}
\expandafter\def\csname sga@target@IV.4.6.1.e\endcsname{ega4.c04.sec.4-6-1-e}
\expandafter\def\csname sga@target@IV.4.6.10\endcsname{ega4.c04.prop.4-6-10}
\expandafter\def\csname sga@target@IV.4.6.11\endcsname{ega4.c04.cor.4-6-11}
\expandafter\def\csname sga@target@IV.4.6.12\endcsname{ega4.c04.prop.4-6-12}
\expandafter\def\csname sga@target@IV.4.6.12.a\endcsname{ega4.c04.sec.4-6-12-a}
\expandafter\def\csname sga@target@IV.4.6.12.b\endcsname{ega4.c04.sec.4-6-12-b}
\expandafter\def\csname sga@target@IV.4.6.12.c\endcsname{ega4.c04.sec.4-6-12-c}
\expandafter\def\csname sga@target@IV.4.6.12.d\endcsname{ega4.c04.sec.4-6-12-d}
\expandafter\def\csname sga@target@IV.4.6.12.e\endcsname{ega4.c04.sec.4-6-12-e}
\expandafter\def\csname sga@target@IV.4.6.13\endcsname{ega4.c04.cor.4-6-13}
\expandafter\def\csname sga@target@IV.4.6.13.f\endcsname{ega4.c04.sec.4-6-13-f}
\expandafter\def\csname sga@target@IV.4.6.14\endcsname{ega4.c04.statement.4-6-14}
\expandafter\def\csname sga@target@IV.4.6.15\endcsname{ega4.c04.prop.4-6-15}
\expandafter\def\csname sga@target@IV.4.6.15.i\endcsname{ega4.c04.sec.4-6-15-i}
\expandafter\def\csname sga@target@IV.4.6.15.ii\endcsname{ega4.c04.sec.4-6-15-ii}
\expandafter\def\csname sga@target@IV.4.6.16\endcsname{ega4.c04.prop.4-6-16}
\expandafter\def\csname sga@target@IV.4.6.16.a\endcsname{ega4.c04.sec.4-6-16-a}
\expandafter\def\csname sga@target@IV.4.6.16.b\endcsname{ega4.c04.sec.4-6-16-b}
\expandafter\def\csname sga@target@IV.4.6.16.c\endcsname{ega4.c04.sec.4-6-16-c}
\expandafter\def\csname sga@target@IV.4.6.16.d\endcsname{ega4.c04.sec.4-6-16-d}
\expandafter\def\csname sga@target@IV.4.6.17\endcsname{ega4.c04.def.4-6-17}
\expandafter\def\csname sga@target@IV.4.6.18\endcsname{ega4.c04.prop.4-6-18}
\expandafter\def\csname sga@target@IV.4.6.18.a\endcsname{ega4.c04.sec.4-6-18-a}
\expandafter\def\csname sga@target@IV.4.6.18.b\endcsname{ega4.c04.sec.4-6-18-b}
\expandafter\def\csname sga@target@IV.4.6.18.c\endcsname{ega4.c04.sec.4-6-18-c}
\expandafter\def\csname sga@target@IV.4.6.18.d\endcsname{ega4.c04.sec.4-6-18-d}
\expandafter\def\csname sga@target@IV.4.6.19\endcsname{ega4.c04.def.4-6-19}
\expandafter\def\csname sga@target@IV.4.6.2\endcsname{ega4.c04.def.4-6-2}
\expandafter\def\csname sga@target@IV.4.6.20\endcsname{ega4.c04.prop.4-6-20}
\expandafter\def\csname sga@target@IV.4.6.21\endcsname{ega4.c04.prop.4-6-21}
\expandafter\def\csname sga@target@IV.4.6.21.i\endcsname{ega4.c04.sec.4-6-21-i}
\expandafter\def\csname sga@target@IV.4.6.21.ii\endcsname{ega4.c04.sec.4-6-21-ii}
\expandafter\def\csname sga@target@IV.4.6.22\endcsname{ega4.c04.def.4-6-22}
\expandafter\def\csname sga@target@IV.4.6.3\endcsname{ega4.c04.cor.4-6-3}
\expandafter\def\csname sga@target@IV.4.6.4\endcsname{ega4.c04.cor.4-6-4}
\expandafter\def\csname sga@target@IV.4.6.5\endcsname{ega4.c04.prop.4-6-5}
\expandafter\def\csname sga@target@IV.4.6.5.i\endcsname{ega4.c04.sec.4-6-5-i}
\expandafter\def\csname sga@target@IV.4.6.5.ii\endcsname{ega4.c04.sec.4-6-5-ii}
\expandafter\def\csname sga@target@IV.4.6.6\endcsname{ega4.c04.prop.4-6-6}
\expandafter\def\csname sga@target@IV.4.6.7\endcsname{ega4.c04.cor.4-6-7}
\expandafter\def\csname sga@target@IV.4.6.8\endcsname{ega4.c04.cor.4-6-8}
\expandafter\def\csname sga@target@IV.4.6.9\endcsname{ega4.c04.def.4-6-9}
\expandafter\def\csname sga@target@IV.4.7.1\endcsname{ega4.c04.lem.4-7-1}
\expandafter\def\csname sga@target@IV.4.7.1.1\endcsname{ega4.c04.eq.4-7-1-1}
\expandafter\def\csname sga@target@IV.4.7.10\endcsname{ega4.c04.prop.4-7-10}
\expandafter\def\csname sga@target@IV.4.7.11\endcsname{ega4.c04.prop.4-7-11}
\expandafter\def\csname sga@target@IV.4.7.12\endcsname{ega4.c04.def.4-7-12}
\expandafter\def\csname sga@target@IV.4.7.13\endcsname{ega4.c04.prop.4-7-13}
\expandafter\def\csname sga@target@IV.4.7.2\endcsname{ega4.c04.cor.4-7-2}
\expandafter\def\csname sga@target@IV.4.7.2.1\endcsname{ega4.c04.eq.4-7-2-1}
\expandafter\def\csname sga@target@IV.4.7.3\endcsname{ega4.c04.prop.4-7-3}
\expandafter\def\csname sga@target@IV.4.7.3.1\endcsname{ega4.c04.eq.4-7-3-1}
\expandafter\def\csname sga@target@IV.4.7.3.i\endcsname{ega4.c04.sec.4-7-3-i}
\expandafter\def\csname sga@target@IV.4.7.3.ii\endcsname{ega4.c04.sec.4-7-3-ii}
\expandafter\def\csname sga@target@IV.4.7.3.iii\endcsname{ega4.c04.sec.4-7-3-iii}
\expandafter\def\csname sga@target@IV.4.7.4\endcsname{ega4.c04.def.4-7-4}
\expandafter\def\csname sga@target@IV.4.7.5\endcsname{ega4.c04.def.4-7-5}
\expandafter\def\csname sga@target@IV.4.7.6\endcsname{ega4.c04.prop.4-7-6}
\expandafter\def\csname sga@target@IV.4.7.6.1\endcsname{ega4.c04.eq.4-7-6-1}
\expandafter\def\csname sga@target@IV.4.7.7\endcsname{ega4.c04.statement.4-7-7}
\expandafter\def\csname sga@target@IV.4.7.8\endcsname{ega4.c04.prop.4-7-8}
\expandafter\def\csname sga@target@IV.4.7.9\endcsname{ega4.c04.cor.4-7-9}
\expandafter\def\csname sga@target@IV.4.8.1\endcsname{ega4.c04.statement.4-8-1}
\expandafter\def\csname sga@target@IV.4.8.10\endcsname{ega4.c04.cor.4-8-10}
\expandafter\def\csname sga@target@IV.4.8.11\endcsname{ega4.c04.cor.4-8-11}
\expandafter\def\csname sga@target@IV.4.8.12\endcsname{ega4.c04.cor.4-8-12}
\expandafter\def\csname sga@target@IV.4.8.13\endcsname{ega4.c04.prop.4-8-13}
\expandafter\def\csname sga@target@IV.4.8.14\endcsname{ega4.c04.prop.4-8-14}
\expandafter\def\csname sga@target@IV.4.8.15\endcsname{ega4.c04.sec.4-8-15}
\expandafter\def\csname sga@target@IV.4.8.15.i\endcsname{ega4.c04.sec.4-8-15-i}
\expandafter\def\csname sga@target@IV.4.8.15.ii\endcsname{ega4.c04.sec.4-8-15-ii}
\expandafter\def\csname sga@target@IV.4.8.2\endcsname{ega4.c04.statement.4-8-2}
\expandafter\def\csname sga@target@IV.4.8.2.1\endcsname{ega4.c04.eq.4-8-2-1}
\expandafter\def\csname sga@target@IV.4.8.2.2\endcsname{ega4.c04.eq.4-8-2-2}
\expandafter\def\csname sga@target@IV.4.8.2.3\endcsname{ega4.c04.eq.4-8-2-3}
\expandafter\def\csname sga@target@IV.4.8.2.4\endcsname{ega4.c04.eq.4-8-2-4}
\expandafter\def\csname sga@target@IV.4.8.2.5\endcsname{ega4.c04.eq.4-8-2-5}
\expandafter\def\csname sga@target@IV.4.8.2.6\endcsname{ega4.c04.eq.4-8-2-6}
\expandafter\def\csname sga@target@IV.4.8.3\endcsname{ega4.c04.prop.4-8-3}
\expandafter\def\csname sga@target@IV.4.8.4\endcsname{ega4.c04.def.4-8-4}
\expandafter\def\csname sga@target@IV.4.8.5\endcsname{ega4.c04.prop.4-8-5}
\expandafter\def\csname sga@target@IV.4.8.6\endcsname{ega4.c04.lem.4-8-6}
\expandafter\def\csname sga@target@IV.4.8.6.a\endcsname{ega4.c04.sec.4-8-6-a}
\expandafter\def\csname sga@target@IV.4.8.6.b\endcsname{ega4.c04.sec.4-8-6-b}
\expandafter\def\csname sga@target@IV.4.8.6.c\endcsname{ega4.c04.sec.4-8-6-c}
\expandafter\def\csname sga@target@IV.4.8.7\endcsname{ega4.c04.cor.4-8-7}
\expandafter\def\csname sga@target@IV.4.8.8\endcsname{ega4.c04.lem.4-8-8}
\expandafter\def\csname sga@target@IV.4.8.8.i\endcsname{ega4.c04.sec.4-8-8-i}
\expandafter\def\csname sga@target@IV.4.8.8.ii\endcsname{ega4.c04.sec.4-8-8-ii}
\expandafter\def\csname sga@target@IV.4.8.8.iii\endcsname{ega4.c04.sec.4-8-8-iii}
\expandafter\def\csname sga@target@IV.4.8.9\endcsname{ega4.c04.prop.4-8-9}
\expandafter\def\csname sga@target@IV.4.9.1\endcsname{ega4.c04.prop.4-9-1}
\expandafter\def\csname sga@target@IV.4.9.1.a\endcsname{ega4.c04.sec.4-9-1-a}
\expandafter\def\csname sga@target@IV.4.9.1.b\endcsname{ega4.c04.sec.4-9-1-b}
\expandafter\def\csname sga@target@IV.4.9.1.bprime\endcsname{ega4.c04.sec.4-9-1-bprime}
\expandafter\def\csname sga@target@IV.4.9.1.c\endcsname{ega4.c04.sec.4-9-1-c}
\expandafter\def\csname sga@target@IV.4.9.2\endcsname{ega4.c04.cor.4-9-2}
\expandafter\def\csname sga@target@IV.4.9.3\endcsname{ega4.c04.rem.4-9-3}
\expandafter\def\csname sga@target@IV.4.9.4\endcsname{ega4.c04.prop.4-9-4}
\expandafter\def\csname sga@target@IV.4.9.5\endcsname{ega4.c04.cor.4-9-5}
\expandafter\def\csname sga@target@IV.4.9.5.a\endcsname{ega4.c04.sec.4-9-5-a}
\expandafter\def\csname sga@target@IV.4.9.5.b\endcsname{ega4.c04.sec.4-9-5-b}
\expandafter\def\csname sga@target@IV.4.9.6\endcsname{ega4.c04.prop.4-9-6}
\expandafter\def\csname sga@target@IV.4.9.6.a\endcsname{ega4.c04.sec.4-9-6-a}
\expandafter\def\csname sga@target@IV.4.9.6.aprime\endcsname{ega4.c04.sec.4-9-6-aprime}
\expandafter\def\csname sga@target@IV.4.9.6.b\endcsname{ega4.c04.sec.4-9-6-b}
\expandafter\def\csname sga@target@IV.4.9.6.bprime\endcsname{ega4.c04.sec.4-9-6-bprime}
\expandafter\def\csname sga@target@IV.4.9.7\endcsname{ega4.c04.cor.4-9-7}
\expandafter\def\csname sga@target@IV.5.1.1\endcsname{ega4.c04.statement.5-1-1}
\expandafter\def\csname sga@target@IV.5.1.1.1\endcsname{ega4.c04.eq.5-1-1-1}
\expandafter\def\csname sga@target@IV.5.1.1.2\endcsname{ega4.c04.eq.5-1-1-2}
\expandafter\def\csname sga@target@IV.5.1.1.3\endcsname{ega4.c04.eq.5-1-1-3}
\expandafter\def\csname sga@target@IV.5.1.1.4\endcsname{ega4.c04.sec.5-1-1-4}
\expandafter\def\csname sga@target@IV.5.1.1.5\endcsname{ega4.c04.sec.5-1-1-5}
\expandafter\def\csname sga@target@IV.5.1.1.6\endcsname{ega4.c04.sec.5-1-1-6}
\expandafter\def\csname sga@target@IV.5.1.10\endcsname{ega4.c04.prop.5-1-10}
\expandafter\def\csname sga@target@IV.5.1.10.1\endcsname{ega4.c04.lem.5-1-10-1}
\expandafter\def\csname sga@target@IV.5.1.10.i\endcsname{ega4.c04.sec.5-1-10-i}
\expandafter\def\csname sga@target@IV.5.1.10.ii\endcsname{ega4.c04.sec.5-1-10-ii}
\expandafter\def\csname sga@target@IV.5.1.11\endcsname{ega4.c04.cor.5-1-11}
\expandafter\def\csname sga@target@IV.5.1.12\endcsname{ega4.c04.statement.5-1-12}
\expandafter\def\csname sga@target@IV.5.1.12.1\endcsname{ega4.c04.eq.5-1-12-1}
\expandafter\def\csname sga@target@IV.5.1.12.2\endcsname{ega4.c04.eq.5-1-12-2}
\expandafter\def\csname sga@target@IV.5.1.12.3\endcsname{ega4.c04.eq.5-1-12-3}
\expandafter\def\csname sga@target@IV.5.1.13\endcsname{ega4.c04.prop.5-1-13}
\expandafter\def\csname sga@target@IV.5.1.13.1\endcsname{ega4.c04.eq.5-1-13-1}
\expandafter\def\csname sga@target@IV.5.1.2\endcsname{ega4.c04.prop.5-1-2}
\expandafter\def\csname sga@target@IV.5.1.2.1\endcsname{ega4.c04.eq.5-1-2-1}
\expandafter\def\csname sga@target@IV.5.1.3\endcsname{ega4.c04.cor.5-1-3}
\expandafter\def\csname sga@target@IV.5.1.3.1\endcsname{ega4.c04.eq.5-1-3-1}
\expandafter\def\csname sga@target@IV.5.1.3.2\endcsname{ega4.c04.eq.5-1-3-2}
\expandafter\def\csname sga@target@IV.5.1.4\endcsname{ega4.c04.prop.5-1-4}
\expandafter\def\csname sga@target@IV.5.1.4.1\endcsname{ega4.c04.eq.5-1-4-1}
\expandafter\def\csname sga@target@IV.5.1.4.2\endcsname{ega4.c04.eq.5-1-4-2}
\expandafter\def\csname sga@target@IV.5.1.5\endcsname{ega4.c04.cor.5-1-5}
\expandafter\def\csname sga@target@IV.5.1.6\endcsname{ega4.c04.statement.5-1-6}
\expandafter\def\csname sga@target@IV.5.1.7\endcsname{ega4.c04.prop.5-1-7}
\expandafter\def\csname sga@target@IV.5.1.8\endcsname{ega4.c04.cor.5-1-8}
\expandafter\def\csname sga@target@IV.5.1.9\endcsname{ega4.c04.prop.5-1-9}
\expandafter\def\csname sga@target@IV.5.1.9.1\endcsname{ega4.c04.statement.5-1-9-1}
\expandafter\def\csname sga@target@IV.5.10.1\endcsname{ega4.c04.statement.5-10-1}
\expandafter\def\csname sga@target@IV.5.10.1.1\endcsname{ega4.c04.eq.5-10-1-1}
\expandafter\def\csname sga@target@IV.5.10.10\endcsname{ega4.c04.prop.5-10-10}
\expandafter\def\csname sga@target@IV.5.10.11\endcsname{ega4.c04.statement.5-10-11}
\expandafter\def\csname sga@target@IV.5.10.11.1\endcsname{ega4.c04.eq.5-10-11-1}
\expandafter\def\csname sga@target@IV.5.10.12\endcsname{ega4.c04.prop.5-10-12}
\expandafter\def\csname sga@target@IV.5.10.13\endcsname{ega4.c04.statement.5-10-13}
\expandafter\def\csname sga@target@IV.5.10.14\endcsname{ega4.c04.prop.5-10-14}
\expandafter\def\csname sga@target@IV.5.10.15\endcsname{ega4.c04.cor.5-10-15}
\expandafter\def\csname sga@target@IV.5.10.16\endcsname{ega4.c04.prop.5-10-16}
\expandafter\def\csname sga@target@IV.5.10.17\endcsname{ega4.c04.prop.5-10-17}
\expandafter\def\csname sga@target@IV.5.10.2\endcsname{ega4.c04.prop.5-10-2}
\expandafter\def\csname sga@target@IV.5.10.3\endcsname{ega4.c04.cor.5-10-3}
\expandafter\def\csname sga@target@IV.5.10.4\endcsname{ega4.c04.cor.5-10-4}
\expandafter\def\csname sga@target@IV.5.10.5\endcsname{ega4.c04.thm.5-10-5}
\expandafter\def\csname sga@target@IV.5.10.6\endcsname{ega4.c04.cor.5-10-6}
\expandafter\def\csname sga@target@IV.5.10.7\endcsname{ega4.c04.thm.5-10-7}
\expandafter\def\csname sga@target@IV.5.10.8\endcsname{ega4.c04.cor.5-10-8}
\expandafter\def\csname sga@target@IV.5.10.8.1\endcsname{ega4.c04.lem.5-10-8-1}
\expandafter\def\csname sga@target@IV.5.10.9\endcsname{ega4.c04.cor.5-10-9}
\expandafter\def\csname sga@target@IV.5.11.1\endcsname{ega4.c04.prop.5-11-1}
\expandafter\def\csname sga@target@IV.5.11.1.1\endcsname{ega4.c04.lem.5-11-1-1}
\expandafter\def\csname sga@target@IV.5.11.2\endcsname{ega4.c04.cor.5-11-2}
\expandafter\def\csname sga@target@IV.5.11.3\endcsname{ega4.c04.rem.5-11-3}
\expandafter\def\csname sga@target@IV.5.11.4\endcsname{ega4.c04.cor.5-11-4}
\expandafter\def\csname sga@target@IV.5.11.5\endcsname{ega4.c04.cor.5-11-5}
\expandafter\def\csname sga@target@IV.5.11.6\endcsname{ega4.c04.cor.5-11-6}
\expandafter\def\csname sga@target@IV.5.11.7\endcsname{ega4.c04.item.5-11-7}
\expandafter\def\csname sga@target@IV.5.12\endcsname{ega4.c04.subsec.5-12.key-iv-5-12}
\expandafter\def\csname sga@target@IV.5.12.1\endcsname{ega4.c04.prop.5-12-1}
\expandafter\def\csname sga@target@IV.5.12.1.1\endcsname{ega4.c04.eq.5-12-1-1}
\expandafter\def\csname sga@target@IV.5.12.10\endcsname{ega4.c04.cor.5-12-10}
\expandafter\def\csname sga@target@IV.5.12.2\endcsname{ega4.c04.prop.5-12-2}
\expandafter\def\csname sga@target@IV.5.12.3\endcsname{ega4.c04.rem.5-12-3}
\expandafter\def\csname sga@target@IV.5.12.4\endcsname{ega4.c04.cor.5-12-4}
\expandafter\def\csname sga@target@IV.5.12.5\endcsname{ega4.c04.prop.5-12-5}
\expandafter\def\csname sga@target@IV.5.12.6\endcsname{ega4.c04.rem.5-12-6}
\expandafter\def\csname sga@target@IV.5.12.7\endcsname{ega4.c04.cor.5-12-7}
\expandafter\def\csname sga@target@IV.5.12.8\endcsname{ega4.c04.prop.5-12-8}
\expandafter\def\csname sga@target@IV.5.12.9\endcsname{ega4.c04.rem.5-12-9}
\expandafter\def\csname sga@target@IV.5.13.1\endcsname{ega4.c04.prop.5-13-1}
\expandafter\def\csname sga@target@IV.5.13.2\endcsname{ega4.c04.cor.5-13-2}
\expandafter\def\csname sga@target@IV.5.13.3\endcsname{ega4.c04.prop.5-13-3}
\expandafter\def\csname sga@target@IV.5.13.4\endcsname{ega4.c04.cor.5-13-4}
\expandafter\def\csname sga@target@IV.5.13.5\endcsname{ega4.c04.statement.5-13-5}
\expandafter\def\csname sga@target@IV.5.13.6\endcsname{ega4.c04.cor.5-13-6}
\expandafter\def\csname sga@target@IV.5.13.7\endcsname{ega4.c04.prop.5-13-7}
\expandafter\def\csname sga@target@IV.5.13.7.1\endcsname{ega4.c04.lem.5-13-7-1}
\expandafter\def\csname sga@target@IV.5.13.8\endcsname{ega4.c04.cor.5-13-8}
\expandafter\def\csname sga@target@IV.5.2.1\endcsname{ega4.c04.prop.5-2-1}
\expandafter\def\csname sga@target@IV.5.2.1.1\endcsname{ega4.c04.eq.5-2-1-1}
\expandafter\def\csname sga@target@IV.5.2.2\endcsname{ega4.c04.cor.5-2-2}
\expandafter\def\csname sga@target@IV.5.2.3\endcsname{ega4.c04.cor.5-2-3}
\expandafter\def\csname sga@target@IV.5.2.3.1\endcsname{ega4.c04.eq.5-2-3-1}
\expandafter\def\csname sga@target@IV.5.2.4\endcsname{ega4.c04.cor.5-2-4}
\expandafter\def\csname sga@target@IV.5.2.5\endcsname{ega4.c04.sec.5-2-5}
\expandafter\def\csname sga@target@IV.5.2.5.i\endcsname{ega4.c04.sec.5-2-5-i}
\expandafter\def\csname sga@target@IV.5.2.5.ii\endcsname{ega4.c04.sec.5-2-5-ii}
\expandafter\def\csname sga@target@IV.5.3.1\endcsname{ega4.c04.prop.5-3-1}
\expandafter\def\csname sga@target@IV.5.3.1.1\endcsname{ega4.c04.eq.5-3-1-1}
\expandafter\def\csname sga@target@IV.5.3.1.2\endcsname{ega4.c04.eq.5-3-1-2}
\expandafter\def\csname sga@target@IV.5.3.1.3\endcsname{ega4.c04.lem.5-3-1-3}
\expandafter\def\csname sga@target@IV.5.4.1\endcsname{ega4.c04.prop.5-4-1}
\expandafter\def\csname sga@target@IV.5.4.1.i\endcsname{ega4.c04.sec.5-4-1-i}
\expandafter\def\csname sga@target@IV.5.4.1.ii\endcsname{ega4.c04.sec.5-4-1-ii}
\expandafter\def\csname sga@target@IV.5.4.2\endcsname{ega4.c04.cor.5-4-2}
\expandafter\def\csname sga@target@IV.5.4.3\endcsname{ega4.c04.sec.5-4-3}
\expandafter\def\csname sga@target@IV.5.4.3.1\endcsname{ega4.c04.eq.5-4-3-1}
\expandafter\def\csname sga@target@IV.5.4.3.i\endcsname{ega4.c04.sec.5-4-3-i}
\expandafter\def\csname sga@target@IV.5.4.3.ii\endcsname{ega4.c04.sec.5-4-3-ii}
\expandafter\def\csname sga@target@IV.5.5\endcsname{ega4.c04.sec.5-5}
\expandafter\def\csname sga@target@IV.5.5.1\endcsname{ega4.c04.statement.5-5-1}
\expandafter\def\csname sga@target@IV.5.5.1.1\endcsname{ega4.c04.eq.5-5-1-1}
\expandafter\def\csname sga@target@IV.5.5.2\endcsname{ega4.c04.prop.5-5-2}
\expandafter\def\csname sga@target@IV.5.5.2.1\endcsname{ega4.c04.eq.5-5-2-1}
\expandafter\def\csname sga@target@IV.5.5.2.2\endcsname{ega4.c04.eq.5-5-2-2}
\expandafter\def\csname sga@target@IV.5.5.3\endcsname{ega4.c04.prop.5-5-3}
\expandafter\def\csname sga@target@IV.5.5.3.1\endcsname{ega4.c04.eq.5-5-3-1}
\expandafter\def\csname sga@target@IV.5.5.4\endcsname{ega4.c04.cor.5-5-4}
\expandafter\def\csname sga@target@IV.5.5.4.1\endcsname{ega4.c04.eq.5-5-4-1}
\expandafter\def\csname sga@target@IV.5.5.5\endcsname{ega4.c04.cor.5-5-5}
\expandafter\def\csname sga@target@IV.5.5.6\endcsname{ega4.c04.cor.5-5-6}
\expandafter\def\csname sga@target@IV.5.5.6.1\endcsname{ega4.c04.eq.5-5-6-1}
\expandafter\def\csname sga@target@IV.5.5.7\endcsname{ega4.c04.lem.5-5-7}
\expandafter\def\csname sga@target@IV.5.5.7.1\endcsname{ega4.c04.eq.5-5-7-1}
\expandafter\def\csname sga@target@IV.5.5.7.2\endcsname{ega4.c04.statement.5-5-7-2}
\expandafter\def\csname sga@target@IV.5.5.7.3\endcsname{ega4.c04.eq.5-5-7-3}
\expandafter\def\csname sga@target@IV.5.5.7.i\endcsname{ega4.c04.sec.5-5-7-i}
\expandafter\def\csname sga@target@IV.5.5.7.ii\endcsname{ega4.c04.sec.5-5-7-ii}
\expandafter\def\csname sga@target@IV.5.5.8\endcsname{ega4.c04.thm.5-5-8}
\expandafter\def\csname sga@target@IV.5.5.8.1\endcsname{ega4.c04.eq.5-5-8-1}
\expandafter\def\csname sga@target@IV.5.5.8.2\endcsname{ega4.c04.eq.5-5-8-2}
\expandafter\def\csname sga@target@IV.5.5.8.3\endcsname{ega4.c04.eq.5-5-8-3}
\expandafter\def\csname sga@target@IV.5.6\endcsname{ega4.c04.sec.5-6}
\expandafter\def\csname sga@target@IV.5.6.1\endcsname{ega4.c04.prop.5-6-1}
\expandafter\def\csname sga@target@IV.5.6.1.1\endcsname{ega4.c04.eq.5-6-1-1}
\expandafter\def\csname sga@target@IV.5.6.1.2\endcsname{ega4.c04.eq.5-6-1-2}
\expandafter\def\csname sga@target@IV.5.6.1.3\endcsname{ega4.c04.statement.5-6-1-3}
\expandafter\def\csname sga@target@IV.5.6.1.4\endcsname{ega4.c04.statement.5-6-1-4}
\expandafter\def\csname sga@target@IV.5.6.1.5\endcsname{ega4.c04.statement.5-6-1-5}
\expandafter\def\csname sga@target@IV.5.6.1.a\endcsname{ega4.c04.sec.5-6-1-a}
\expandafter\def\csname sga@target@IV.5.6.1.b\endcsname{ega4.c04.sec.5-6-1-b}
\expandafter\def\csname sga@target@IV.5.6.1.c\endcsname{ega4.c04.sec.5-6-1-c}
\expandafter\def\csname sga@target@IV.5.6.10\endcsname{ega4.c04.prop.5-6-10}
\expandafter\def\csname sga@target@IV.5.6.11\endcsname{ega4.c04.ex.5-6-11}
\expandafter\def\csname sga@target@IV.5.6.11.1\endcsname{ega4.c04.lem.5-6-11-1}
\expandafter\def\csname sga@target@IV.5.6.11.1.i\endcsname{ega4.c04.sec.5-6-11-1-i}
\expandafter\def\csname sga@target@IV.5.6.11.1.ii\endcsname{ega4.c04.sec.5-6-11-1-ii}
\expandafter\def\csname sga@target@IV.5.6.2\endcsname{ega4.c04.def.5-6-2}
\expandafter\def\csname sga@target@IV.5.6.3\endcsname{ega4.c04.sec.5-6-3}
\expandafter\def\csname sga@target@IV.5.6.3.i\endcsname{ega4.c04.sec.5-6-3-i}
\expandafter\def\csname sga@target@IV.5.6.3.ii\endcsname{ega4.c04.sec.5-6-3-ii}
\expandafter\def\csname sga@target@IV.5.6.3.iii\endcsname{ega4.c04.sec.5-6-3-iii}
\expandafter\def\csname sga@target@IV.5.6.3.iv\endcsname{ega4.c04.sec.5-6-3-iv}
\expandafter\def\csname sga@target@IV.5.6.4\endcsname{ega4.c04.prop.5-6-4}
\expandafter\def\csname sga@target@IV.5.6.5\endcsname{ega4.c04.prop.5-6-5}
\expandafter\def\csname sga@target@IV.5.6.5.1\endcsname{ega4.c04.eq.5-6-5-1}
\expandafter\def\csname sga@target@IV.5.6.5.2\endcsname{ega4.c04.eq.5-6-5-2}
\expandafter\def\csname sga@target@IV.5.6.5.3\endcsname{ega4.c04.eq.5-6-5-3}
\expandafter\def\csname sga@target@IV.5.6.6\endcsname{ega4.c04.cor.5-6-6}
\expandafter\def\csname sga@target@IV.5.6.6.1\endcsname{ega4.c04.eq.5-6-6-1}
\expandafter\def\csname sga@target@IV.5.6.6.2\endcsname{ega4.c04.eq.5-6-6-2}
\expandafter\def\csname sga@target@IV.5.6.7\endcsname{ega4.c04.cor.5-6-7}
\expandafter\def\csname sga@target@IV.5.6.7.1\endcsname{ega4.c04.eq.5-6-7-1}
\expandafter\def\csname sga@target@IV.5.6.8\endcsname{ega4.c04.cor.5-6-8}
\expandafter\def\csname sga@target@IV.5.6.8.1\endcsname{ega4.c04.eq.5-6-8-1}
\expandafter\def\csname sga@target@IV.5.6.8.2\endcsname{ega4.c04.eq.5-6-8-2}
\expandafter\def\csname sga@target@IV.5.6.9\endcsname{ega4.c04.sec.5-6-9}
\expandafter\def\csname sga@target@IV.5.6.9.i\endcsname{ega4.c04.sec.5-6-9-i}
\expandafter\def\csname sga@target@IV.5.6.9.ii\endcsname{ega4.c04.sec.5-6-9-ii}
\expandafter\def\csname sga@target@IV.5.7.1\endcsname{ega4.c04.def.5-7-1}
\expandafter\def\csname sga@target@IV.5.7.1.1\endcsname{ega4.c04.eq.5-7-1-1}
\expandafter\def\csname sga@target@IV.5.7.10\endcsname{ega4.c04.cor.5-7-10}
\expandafter\def\csname sga@target@IV.5.7.11\endcsname{ega4.c04.cor.5-7-11}
\expandafter\def\csname sga@target@IV.5.7.11.i\endcsname{ega4.c04.sec.5-7-11-i}
\expandafter\def\csname sga@target@IV.5.7.11.ii\endcsname{ega4.c04.sec.5-7-11-ii}
\expandafter\def\csname sga@target@IV.5.7.12\endcsname{ega4.c04.statement.5-7-12}
\expandafter\def\csname sga@target@IV.5.7.13\endcsname{ega4.c04.cor.5-7-13}
\expandafter\def\csname sga@target@IV.5.7.13.1\endcsname{ega4.c04.eq.5-7-13-1}
\expandafter\def\csname sga@target@IV.5.7.2\endcsname{ega4.c04.def.5-7-2}
\expandafter\def\csname sga@target@IV.5.7.2.1\endcsname{ega4.c04.eq.5-7-2-1}
\expandafter\def\csname sga@target@IV.5.7.2.2\endcsname{ega4.c04.eq.5-7-2-2}
\expandafter\def\csname sga@target@IV.5.7.3\endcsname{ega4.c04.sec.5-7-3}
\expandafter\def\csname sga@target@IV.5.7.3.i\endcsname{ega4.c04.sec.5-7-3-i}
\expandafter\def\csname sga@target@IV.5.7.3.ii\endcsname{ega4.c04.sec.5-7-3-ii}
\expandafter\def\csname sga@target@IV.5.7.3.iii\endcsname{ega4.c04.sec.5-7-3-iii}
\expandafter\def\csname sga@target@IV.5.7.3.iv\endcsname{ega4.c04.sec.5-7-3-iv}
\expandafter\def\csname sga@target@IV.5.7.3.v\endcsname{ega4.c04.sec.5-7-3-v}
\expandafter\def\csname sga@target@IV.5.7.3.vi\endcsname{ega4.c04.sec.5-7-3-vi}
\expandafter\def\csname sga@target@IV.5.7.4\endcsname{ega4.c04.prop.5-7-4}
\expandafter\def\csname sga@target@IV.5.7.4.1\endcsname{ega4.c04.eq.5-7-4-1}
\expandafter\def\csname sga@target@IV.5.7.4.2\endcsname{ega4.c04.eq.5-7-4-2}
\expandafter\def\csname sga@target@IV.5.7.4.i\endcsname{ega4.c04.sec.5-7-4-i}
\expandafter\def\csname sga@target@IV.5.7.4.ii\endcsname{ega4.c04.sec.5-7-4-ii}
\expandafter\def\csname sga@target@IV.5.7.5\endcsname{ega4.c04.prop.5-7-5}
\expandafter\def\csname sga@target@IV.5.7.6\endcsname{ega4.c04.cor.5-7-6}
\expandafter\def\csname sga@target@IV.5.7.7\endcsname{ega4.c04.cor.5-7-7}
\expandafter\def\csname sga@target@IV.5.7.8\endcsname{ega4.c04.sec.5-7-8}
\expandafter\def\csname sga@target@IV.5.7.9\endcsname{ega4.c04.prop.5-7-9}
\expandafter\def\csname sga@target@IV.5.7.9.1\endcsname{ega4.c04.eq.5-7-9-1}
\expandafter\def\csname sga@target@IV.5.7.9.2\endcsname{ega4.c04.eq.5-7-9-2}
\expandafter\def\csname sga@target@IV.5.8.1\endcsname{ega4.c04.statement.5-8-1}
\expandafter\def\csname sga@target@IV.5.8.2\endcsname{ega4.c04.def.5-8-2}
\expandafter\def\csname sga@target@IV.5.8.3\endcsname{ega4.c04.prop.5-8-3}
\expandafter\def\csname sga@target@IV.5.8.4\endcsname{ega4.c04.statement.5-8-4}
\expandafter\def\csname sga@target@IV.5.8.5\endcsname{ega4.c04.prop.5-8-5}
\expandafter\def\csname sga@target@IV.5.8.6\endcsname{ega4.c04.thm.5-8-6}
\expandafter\def\csname sga@target@IV.5.8.6.1\endcsname{ega4.c04.eq.5-8-6-1}
\expandafter\def\csname sga@target@IV.5.8.6.i\endcsname{ega4.c04.item.5-8-6-i}
\expandafter\def\csname sga@target@IV.5.8.6.ii\endcsname{ega4.c04.item.5-8-6-ii}
\expandafter\def\csname sga@target@IV.5.9.1\endcsname{ega4.c04.statement.5-9-1}
\expandafter\def\csname sga@target@IV.5.9.1.1\endcsname{ega4.c04.eq.5-9-1-1}
\expandafter\def\csname sga@target@IV.5.9.1.2\endcsname{ega4.c04.eq.5-9-1-2}
\expandafter\def\csname sga@target@IV.5.9.1.3\endcsname{ega4.c04.eq.5-9-1-3}
\expandafter\def\csname sga@target@IV.5.9.10\endcsname{ega4.c04.cor.5-9-10}
\expandafter\def\csname sga@target@IV.5.9.10.i\endcsname{ega4.c04.item.5-9-10-i}
\expandafter\def\csname sga@target@IV.5.9.10.ii\endcsname{ega4.c04.item.5-9-10-ii}
\expandafter\def\csname sga@target@IV.5.9.11\endcsname{ega4.c04.prop.5-9-11}
\expandafter\def\csname sga@target@IV.5.9.12\endcsname{ega4.c04.item.5-9-12}
\expandafter\def\csname sga@target@IV.5.9.12.i\endcsname{ega4.c04.item.5-9-12-i}
\expandafter\def\csname sga@target@IV.5.9.12.ii\endcsname{ega4.c04.item.5-9-12-ii}
\expandafter\def\csname sga@target@IV.5.9.12.iii\endcsname{ega4.c04.item.5-9-12-iii}
\expandafter\def\csname sga@target@IV.5.9.2\endcsname{ega4.c04.prop.5-9-2}
\expandafter\def\csname sga@target@IV.5.9.2.i\endcsname{ega4.c04.item.5-9-2-i}
\expandafter\def\csname sga@target@IV.5.9.2.ii\endcsname{ega4.c04.item.5-9-2-ii}
\expandafter\def\csname sga@target@IV.5.9.3\endcsname{ega4.c04.rem.5-9-3}
\expandafter\def\csname sga@target@IV.5.9.3.1\endcsname{ega4.c04.eq.5-9-3-1}
\expandafter\def\csname sga@target@IV.5.9.4\endcsname{ega4.c04.prop.5-9-4}
\expandafter\def\csname sga@target@IV.5.9.4.1\endcsname{ega4.c04.eq.5-9-4-1}
\expandafter\def\csname sga@target@IV.5.9.5\endcsname{ega4.c04.cor.5-9-5}
\expandafter\def\csname sga@target@IV.5.9.6\endcsname{ega4.c04.cor.5-9-6}
\expandafter\def\csname sga@target@IV.5.9.6.1\endcsname{ega4.c04.eq.5-9-6-1}
\expandafter\def\csname sga@target@IV.5.9.7\endcsname{ega4.c04.statement.5-9-7}
\expandafter\def\csname sga@target@IV.5.9.7.1\endcsname{ega4.c04.eq.5-9-7-1}
\expandafter\def\csname sga@target@IV.5.9.8\endcsname{ega4.c04.prop.5-9-8}
\expandafter\def\csname sga@target@IV.5.9.8.1\endcsname{ega4.c04.eq.5-9-8-1}
\expandafter\def\csname sga@target@IV.5.9.8.a\endcsname{ega4.c04.item.5-9-8-a}
\expandafter\def\csname sga@target@IV.5.9.8.a-prime\endcsname{ega4.c04.item.5-9-8-a-prime}
\expandafter\def\csname sga@target@IV.5.9.8.b\endcsname{ega4.c04.item.5-9-8-b}
\expandafter\def\csname sga@target@IV.5.9.8.b-prime\endcsname{ega4.c04.item.5-9-8-b-prime}
\expandafter\def\csname sga@target@IV.5.9.9\endcsname{ega4.c04.def.5-9-9}
\expandafter\def\csname sga@target@IV.6.1.1\endcsname{ega4.c04.prop.6-1-1}
\expandafter\def\csname sga@target@IV.6.1.1.1\endcsname{ega4.c04.eq.6-1-1-1}
\expandafter\def\csname sga@target@IV.6.1.2\endcsname{ega4.c04.cor.6-1-2}
\expandafter\def\csname sga@target@IV.6.1.2.1\endcsname{ega4.c04.eq.6-1-2-1}
\expandafter\def\csname sga@target@IV.6.1.3\endcsname{ega4.c04.cor.6-1-3}
\expandafter\def\csname sga@target@IV.6.1.4\endcsname{ega4.c04.cor.6-1-4}
\expandafter\def\csname sga@target@IV.6.1.4.1\endcsname{ega4.c04.eq.6-1-4-1}
\expandafter\def\csname sga@target@IV.6.1.5\endcsname{ega4.c04.prop.6-1-5}
\expandafter\def\csname sga@target@IV.6.10.1\endcsname{ega4.c04.statement.6-10-1}
\expandafter\def\csname sga@target@IV.6.10.2\endcsname{ega4.c04.prop.6-10-2}
\expandafter\def\csname sga@target@IV.6.10.3\endcsname{ega4.c04.prop.6-10-3}
\expandafter\def\csname sga@target@IV.6.10.3.1\endcsname{ega4.c04.eq.6-10-3-1}
\expandafter\def\csname sga@target@IV.6.10.3.2\endcsname{ega4.c04.eq.6-10-3-2}
\expandafter\def\csname sga@target@IV.6.10.3.3\endcsname{ega4.c04.eq.6-10-3-3}
\expandafter\def\csname sga@target@IV.6.10.4\endcsname{ega4.c04.cor.6-10-4}
\expandafter\def\csname sga@target@IV.6.10.5\endcsname{ega4.c04.prop.6-10-5}
\expandafter\def\csname sga@target@IV.6.10.5.1\endcsname{ega4.c04.eq.6-10-5-1}
\expandafter\def\csname sga@target@IV.6.10.5.2\endcsname{ega4.c04.eq.6-10-5-2}
\expandafter\def\csname sga@target@IV.6.10.5.3\endcsname{ega4.c04.eq.6-10-5-3}
\expandafter\def\csname sga@target@IV.6.10.5.4\endcsname{ega4.c04.eq.6-10-5-4}
\expandafter\def\csname sga@target@IV.6.10.6\endcsname{ega4.c04.prop.6-10-6}
\expandafter\def\csname sga@target@IV.6.10.6.1\endcsname{ega4.c04.eq.6-10-6-1}
\expandafter\def\csname sga@target@IV.6.10.6.2\endcsname{ega4.c04.eq.6-10-6-2}
\expandafter\def\csname sga@target@IV.6.10.6.3\endcsname{ega4.c04.eq.6-10-6-3}
\expandafter\def\csname sga@target@IV.6.11\endcsname{ega4.c04.subsec.6-11.key-iv-6-11}
\expandafter\def\csname sga@target@IV.6.11.1\endcsname{ega4.c04.lem.6-11-1}
\expandafter\def\csname sga@target@IV.6.11.2\endcsname{ega4.c04.prop.6-11-2}
\expandafter\def\csname sga@target@IV.6.11.2.1\endcsname{ega4.c04.eq.6-11-2-1}
\expandafter\def\csname sga@target@IV.6.11.2.2\endcsname{ega4.c04.eq.6-11-2-2}
\expandafter\def\csname sga@target@IV.6.11.3\endcsname{ega4.c04.cor.6-11-3}
\expandafter\def\csname sga@target@IV.6.11.4\endcsname{ega4.c04.statement.6-11-4}
\expandafter\def\csname sga@target@IV.6.11.5\endcsname{ega4.c04.prop.6-11-5}
\expandafter\def\csname sga@target@IV.6.11.5.1\endcsname{ega4.c04.eq.6-11-5-1}
\expandafter\def\csname sga@target@IV.6.11.6\endcsname{ega4.c04.prop.6-11-6}
\expandafter\def\csname sga@target@IV.6.11.6.1\endcsname{ega4.c04.lem.6-11-6-1}
\expandafter\def\csname sga@target@IV.6.11.7\endcsname{ega4.c04.cor.6-11-7}
\expandafter\def\csname sga@target@IV.6.11.8\endcsname{ega4.c04.prop.6-11-8}
\expandafter\def\csname sga@target@IV.6.11.8.1\endcsname{ega4.c04.eq.6-11-8-1}
\expandafter\def\csname sga@target@IV.6.11.9\endcsname{ega4.c04.statement.6-11-9}
\expandafter\def\csname sga@target@IV.6.12\endcsname{ega4.c04.subsec.6-12.key-iv-6-12}
\expandafter\def\csname sga@target@IV.6.12.1\endcsname{ega4.c04.statement.6-12-1}
\expandafter\def\csname sga@target@IV.6.12.10\endcsname{ega4.c04.prop.6-12-10}
\expandafter\def\csname sga@target@IV.6.12.2\endcsname{ega4.c04.prop.6-12-2}
\expandafter\def\csname sga@target@IV.6.12.3\endcsname{ega4.c04.cor.6-12-3}
\expandafter\def\csname sga@target@IV.6.12.4\endcsname{ega4.c04.thm.6-12-4}
\expandafter\def\csname sga@target@IV.6.12.4.1\endcsname{ega4.c04.lem.6-12-4-1}
\expandafter\def\csname sga@target@IV.6.12.5\endcsname{ega4.c04.cor.6-12-5}
\expandafter\def\csname sga@target@IV.6.12.6\endcsname{ega4.c04.cor.6-12-6}
\expandafter\def\csname sga@target@IV.6.12.7\endcsname{ega4.c04.thm.6-12-7}
\expandafter\def\csname sga@target@IV.6.12.8\endcsname{ega4.c04.cor.6-12-8}
\expandafter\def\csname sga@target@IV.6.12.9\endcsname{ega4.c04.prop.6-12-9}
\expandafter\def\csname sga@target@IV.6.13\endcsname{ega4.c04.subsec.6-13.key-iv-6-13}
\expandafter\def\csname sga@target@IV.6.13.1\endcsname{ega4.c04.statement.6-13-1}
\expandafter\def\csname sga@target@IV.6.13.2\endcsname{ega4.c04.prop.6-13-2}
\expandafter\def\csname sga@target@IV.6.13.3\endcsname{ega4.c04.cor.6-13-3}
\expandafter\def\csname sga@target@IV.6.13.4\endcsname{ega4.c04.prop.6-13-4}
\expandafter\def\csname sga@target@IV.6.13.5\endcsname{ega4.c04.cor.6-13-5}
\expandafter\def\csname sga@target@IV.6.13.6\endcsname{ega4.c04.prop.6-13-6}
\expandafter\def\csname sga@target@IV.6.13.7\endcsname{ega4.c04.prop.6-13-7}
\expandafter\def\csname sga@target@IV.6.13.7.1\endcsname{ega4.c04.lem.6-13-7-1}
\expandafter\def\csname sga@target@IV.6.14.1\endcsname{ega4.c04.prop.6-14-1}
\expandafter\def\csname sga@target@IV.6.14.1.1\endcsname{ega4.c04.lem.6-14-1-1}
\expandafter\def\csname sga@target@IV.6.14.1.2\endcsname{ega4.c04.lem.6-14-1-2}
\expandafter\def\csname sga@target@IV.6.14.1.2.env\endcsname{ega4.c04.statement.6-14-1-2-env}
\expandafter\def\csname sga@target@IV.6.14.1.3\endcsname{ega4.c04.eq.6-14-1-3}
\expandafter\def\csname sga@target@IV.6.14.1.4\endcsname{ega4.c04.lem.6-14-1-4}
\expandafter\def\csname sga@target@IV.6.14.2\endcsname{ega4.c04.cor.6-14-2}
\expandafter\def\csname sga@target@IV.6.14.3\endcsname{ega4.c04.cor.6-14-3}
\expandafter\def\csname sga@target@IV.6.14.4\endcsname{ega4.c04.prop.6-14-4}
\expandafter\def\csname sga@target@IV.6.14.4.1\endcsname{ega4.c04.lem.6-14-4-1}
\expandafter\def\csname sga@target@IV.6.14.5\endcsname{ega4.c04.cor.6-14-5}
\expandafter\def\csname sga@target@IV.6.15.1\endcsname{ega4.c04.statement.6-15-1}
\expandafter\def\csname sga@target@IV.6.15.10\endcsname{ega4.c04.prop.6-15-10}
\expandafter\def\csname sga@target@IV.6.15.11\endcsname{ega4.c04.item.6-15-11}
\expandafter\def\csname sga@target@IV.6.15.2\endcsname{ega4.c04.statement.6-15-2}
\expandafter\def\csname sga@target@IV.6.15.3\endcsname{ega4.c04.statement.6-15-3}
\expandafter\def\csname sga@target@IV.6.15.3.1\endcsname{ega4.c04.lem.6-15-3-1}
\expandafter\def\csname sga@target@IV.6.15.4\endcsname{ega4.c04.statement.6-15-4}
\expandafter\def\csname sga@target@IV.6.15.4.1\endcsname{ega4.c04.lem.6-15-4-1}
\expandafter\def\csname sga@target@IV.6.15.5\endcsname{ega4.c04.prop.6-15-5}
\expandafter\def\csname sga@target@IV.6.15.6\endcsname{ega4.c04.prop.6-15-6}
\expandafter\def\csname sga@target@IV.6.15.7\endcsname{ega4.c04.prop.6-15-7}
\expandafter\def\csname sga@target@IV.6.15.8\endcsname{ega4.c04.lem.6-15-8}
\expandafter\def\csname sga@target@IV.6.15.9\endcsname{ega4.c04.cor.6-15-9}
\expandafter\def\csname sga@target@IV.6.2.1\endcsname{ega4.c04.prop.6-2-1}
\expandafter\def\csname sga@target@IV.6.2.1.1\endcsname{ega4.c04.eq.6-2-1-1}
\expandafter\def\csname sga@target@IV.6.2.1.2\endcsname{ega4.c04.eq.6-2-1-2}
\expandafter\def\csname sga@target@IV.6.2.2\endcsname{ega4.c04.cor.6-2-2}
\expandafter\def\csname sga@target@IV.6.2.3\endcsname{ega4.c04.prop.6-2-3}
\expandafter\def\csname sga@target@IV.6.2.3.1\endcsname{ega4.c04.eq.6-2-3-1}
\expandafter\def\csname sga@target@IV.6.2.4\endcsname{ega4.c04.prop.6-2-4}
\expandafter\def\csname sga@target@IV.6.2.4.1\endcsname{ega4.c04.eq.6-2-4-1}
\expandafter\def\csname sga@target@IV.6.3.1\endcsname{ega4.c04.prop.6-3-1}
\expandafter\def\csname sga@target@IV.6.3.1.1\endcsname{ega4.c04.eq.6-3-1-1}
\expandafter\def\csname sga@target@IV.6.3.2\endcsname{ega4.c04.cor.6-3-2}
\expandafter\def\csname sga@target@IV.6.3.2.1\endcsname{ega4.c04.eq.6-3-2-1}
\expandafter\def\csname sga@target@IV.6.3.3\endcsname{ega4.c04.cor.6-3-3}
\expandafter\def\csname sga@target@IV.6.3.4\endcsname{ega4.c04.cor.6-3-4}
\expandafter\def\csname sga@target@IV.6.3.4.1\endcsname{ega4.c04.eq.6-3-4-1}
\expandafter\def\csname sga@target@IV.6.3.4.2\endcsname{ega4.c04.eq.6-3-4-2}
\expandafter\def\csname sga@target@IV.6.3.5\endcsname{ega4.c04.cor.6-3-5}
\expandafter\def\csname sga@target@IV.6.3.6\endcsname{ega4.c04.cor.6-3-6}
\expandafter\def\csname sga@target@IV.6.3.7\endcsname{ega4.c04.cor.6-3-7}
\expandafter\def\csname sga@target@IV.6.3.8\endcsname{ega4.c04.prop.6-3-8}
\expandafter\def\csname sga@target@IV.6.4.1\endcsname{ega4.c04.prop.6-4-1}
\expandafter\def\csname sga@target@IV.6.4.2\endcsname{ega4.c04.cor.6-4-2}
\expandafter\def\csname sga@target@IV.6.4.3\endcsname{ega4.c04.rem.6-4-3}
\expandafter\def\csname sga@target@IV.6.5.1\endcsname{ega4.c04.prop.6-5-1}
\expandafter\def\csname sga@target@IV.6.5.2\endcsname{ega4.c04.cor.6-5-2}
\expandafter\def\csname sga@target@IV.6.5.3\endcsname{ega4.c04.prop.6-5-3}
\expandafter\def\csname sga@target@IV.6.5.4\endcsname{ega4.c04.cor.6-5-4}
\expandafter\def\csname sga@target@IV.6.5.5\endcsname{ega4.c04.item.6-5-5}
\expandafter\def\csname sga@target@IV.6.6.1\endcsname{ega4.c04.prop.6-6-1}
\expandafter\def\csname sga@target@IV.6.6.2\endcsname{ega4.c04.rem.6-6-2}
\expandafter\def\csname sga@target@IV.6.7.1\endcsname{ega4.c04.prop.6-7-1}
\expandafter\def\csname sga@target@IV.6.7.1.1\endcsname{ega4.c04.lem.6-7-1-1}
\expandafter\def\csname sga@target@IV.6.7.2\endcsname{ega4.c04.cor.6-7-2}
\expandafter\def\csname sga@target@IV.6.7.3\endcsname{ega4.c04.cor.6-7-3}
\expandafter\def\csname sga@target@IV.6.7.4\endcsname{ega4.c04.prop.6-7-4}
\expandafter\def\csname sga@target@IV.6.7.4.1\endcsname{ega4.c04.lem.6-7-4-1}
\expandafter\def\csname sga@target@IV.6.7.5\endcsname{ega4.c04.item.6-7-5}
\expandafter\def\csname sga@target@IV.6.7.6\endcsname{ega4.c04.def.6-7-6}
\expandafter\def\csname sga@target@IV.6.7.7\endcsname{ega4.c04.prop.6-7-7}
\expandafter\def\csname sga@target@IV.6.7.8\endcsname{ega4.c04.cor.6-7-8}
\expandafter\def\csname sga@target@IV.6.8.1\endcsname{ega4.c04.def.6-8-1}
\expandafter\def\csname sga@target@IV.6.8.2\endcsname{ega4.c04.prop.6-8-2}
\expandafter\def\csname sga@target@IV.6.8.3\endcsname{ega4.c04.prop.6-8-3}
\expandafter\def\csname sga@target@IV.6.8.4\endcsname{ega4.c04.item.6-8-4}
\expandafter\def\csname sga@target@IV.6.8.5\endcsname{ega4.c04.prop.6-8-5}
\expandafter\def\csname sga@target@IV.6.8.6\endcsname{ega4.c04.thm.6-8-6}
\expandafter\def\csname sga@target@IV.6.8.7\endcsname{ega4.c04.cor.6-8-7}
\expandafter\def\csname sga@target@IV.6.8.8\endcsname{ega4.c04.rem.6-8-8}
\expandafter\def\csname sga@target@IV.6.9.1\endcsname{ega4.c04.thm.6-9-1}
\expandafter\def\csname sga@target@IV.6.9.2\endcsname{ega4.c04.lem.6-9-2}
\expandafter\def\csname sga@target@IV.6.9.3\endcsname{ega4.c04.cor.6-9-3}
\expandafter\def\csname sga@target@IV.7.1.1\endcsname{ega4.c04.def.7-1-1}
\expandafter\def\csname sga@target@IV.7.1.10\endcsname{ega4.c04.prop.7-1-10}
\expandafter\def\csname sga@target@IV.7.1.11\endcsname{ega4.c04.prop.7-1-11}
\expandafter\def\csname sga@target@IV.7.1.12\endcsname{ega4.c04.item.7-1-12}
\expandafter\def\csname sga@target@IV.7.1.13\endcsname{ega4.c04.cor.7-1-13}
\expandafter\def\csname sga@target@IV.7.1.2\endcsname{ega4.c04.prop.7-1-2}
\expandafter\def\csname sga@target@IV.7.1.3\endcsname{ega4.c04.prop.7-1-3}
\expandafter\def\csname sga@target@IV.7.1.3.1\endcsname{ega4.c04.eq.7-1-3-1}
\expandafter\def\csname sga@target@IV.7.1.3.2\endcsname{ega4.c04.eq.7-1-3-2}
\expandafter\def\csname sga@target@IV.7.1.3.3\endcsname{ega4.c04.eq.7-1-3-3}
\expandafter\def\csname sga@target@IV.7.1.3.4\endcsname{ega4.c04.eq.7-1-3-4}
\expandafter\def\csname sga@target@IV.7.1.3.5\endcsname{ega4.c04.eq.7-1-3-5}
\expandafter\def\csname sga@target@IV.7.1.3.6\endcsname{ega4.c04.eq.7-1-3-6}
\expandafter\def\csname sga@target@IV.7.1.3.7\endcsname{ega4.c04.eq.7-1-3-7}
\expandafter\def\csname sga@target@IV.7.1.4\endcsname{ega4.c04.cor.7-1-4}
\expandafter\def\csname sga@target@IV.7.1.5\endcsname{ega4.c04.cor.7-1-5}
\expandafter\def\csname sga@target@IV.7.1.6\endcsname{ega4.c04.rem.7-1-6}
\expandafter\def\csname sga@target@IV.7.1.7\endcsname{ega4.c04.lem.7-1-7}
\expandafter\def\csname sga@target@IV.7.1.8\endcsname{ega4.c04.thm.7-1-8}
\expandafter\def\csname sga@target@IV.7.1.9\endcsname{ega4.c04.def.7-1-9}
\expandafter\def\csname sga@target@IV.7.2.1\endcsname{ega4.c04.statement.7-2-1}
\expandafter\def\csname sga@target@IV.7.2.1.1\endcsname{ega4.c04.statement.7-2-1-1}
\expandafter\def\csname sga@target@IV.7.2.10\endcsname{ega4.c04.statement.7-2-10}
\expandafter\def\csname sga@target@IV.7.2.2\endcsname{ega4.c04.prop.7-2-2}
\expandafter\def\csname sga@target@IV.7.2.3\endcsname{ega4.c04.prop.7-2-3}
\expandafter\def\csname sga@target@IV.7.2.4\endcsname{ega4.c04.prop.7-2-4}
\expandafter\def\csname sga@target@IV.7.2.5\endcsname{ega4.c04.thm.7-2-5}
\expandafter\def\csname sga@target@IV.7.2.5.1\endcsname{ega4.c04.eq.7-2-5-1}
\expandafter\def\csname sga@target@IV.7.2.6\endcsname{ega4.c04.def.7-2-6}
\expandafter\def\csname sga@target@IV.7.2.7\endcsname{ega4.c04.cor.7-2-7}
\expandafter\def\csname sga@target@IV.7.2.8\endcsname{ega4.c04.cor.7-2-8}
\expandafter\def\csname sga@target@IV.7.2.9\endcsname{ega4.c04.cor.7-2-9}
\expandafter\def\csname sga@target@IV.7.3.1\endcsname{ega4.c04.statement.7-3-1}
\expandafter\def\csname sga@target@IV.7.3.10\endcsname{ega4.c04.statement.7-3-10}
\expandafter\def\csname sga@target@IV.7.3.11\endcsname{ega4.c04.rem.7-3-11}
\expandafter\def\csname sga@target@IV.7.3.12\endcsname{ega4.c04.prop.7-3-12}
\expandafter\def\csname sga@target@IV.7.3.12.1\endcsname{ega4.c04.eq.7-3-12-1}
\expandafter\def\csname sga@target@IV.7.3.13\endcsname{ega4.c04.statement.7-3-13}
\expandafter\def\csname sga@target@IV.7.3.14\endcsname{ega4.c04.prop.7-3-14}
\expandafter\def\csname sga@target@IV.7.3.15\endcsname{ega4.c04.prop.7-3-15}
\expandafter\def\csname sga@target@IV.7.3.16\endcsname{ega4.c04.prop.7-3-16}
\expandafter\def\csname sga@target@IV.7.3.16.1\endcsname{ega4.c04.eq.7-3-16-1}
\expandafter\def\csname sga@target@IV.7.3.16.2\endcsname{ega4.c04.lem.7-3-16-2}
\expandafter\def\csname sga@target@IV.7.3.17\endcsname{ega4.c04.prop.7-3-17}
\expandafter\def\csname sga@target@IV.7.3.18\endcsname{ega4.c04.prop.7-3-18}
\expandafter\def\csname sga@target@IV.7.3.19\endcsname{ega4.c04.statement.7-3-19}
\expandafter\def\csname sga@target@IV.7.3.2\endcsname{ega4.c04.lem.7-3-2}
\expandafter\def\csname sga@target@IV.7.3.3\endcsname{ega4.c04.cor.7-3-3}
\expandafter\def\csname sga@target@IV.7.3.4\endcsname{ega4.c04.statement.7-3-4}
\expandafter\def\csname sga@target@IV.7.3.5\endcsname{ega4.c04.item.7-3-5}
\expandafter\def\csname sga@target@IV.7.3.6\endcsname{ega4.c04.statement.7-3-6}
\expandafter\def\csname sga@target@IV.7.3.7\endcsname{ega4.c04.lem.7-3-7}
\expandafter\def\csname sga@target@IV.7.3.8\endcsname{ega4.c04.statement.7-3-8}
\expandafter\def\csname sga@target@IV.7.3.9\endcsname{ega4.c04.item.7-3-9}
\expandafter\def\csname sga@target@IV.7.4.1\endcsname{ega4.c04.statement.7-4-1}
\expandafter\def\csname sga@target@IV.7.4.2\endcsname{ega4.c04.lem.7-4-2}
\expandafter\def\csname sga@target@IV.7.4.3\endcsname{ega4.c04.cor.7-4-3}
\expandafter\def\csname sga@target@IV.7.4.4\endcsname{ega4.c04.thm.7-4-4}
\expandafter\def\csname sga@target@IV.7.4.4.1\endcsname{ega4.c04.lem.7-4-4-1}
\expandafter\def\csname sga@target@IV.7.4.4.2\endcsname{ega4.c04.lem.7-4-4-2}
\expandafter\def\csname sga@target@IV.7.4.5\endcsname{ega4.c04.cor.7-4-5}
\expandafter\def\csname sga@target@IV.7.4.6\endcsname{ega4.c04.prop.7-4-6}
\expandafter\def\csname sga@target@IV.7.4.7\endcsname{ega4.c04.cor.7-4-7}
\expandafter\def\csname sga@target@IV.7.4.8\endcsname{ega4.c04.rem.7-4-8}
\expandafter\def\csname sga@target@IV.7.5.0\endcsname{ega4.c04.statement.7-5-0}
\expandafter\def\csname sga@target@IV.7.5.1\endcsname{ega4.c04.thm.7-5-1}
\expandafter\def\csname sga@target@IV.7.5.1.1\endcsname{ega4.c04.lem.7-5-1-1}
\expandafter\def\csname sga@target@IV.7.5.2\endcsname{ega4.c04.cor.7-5-2}
\expandafter\def\csname sga@target@IV.7.5.3\endcsname{ega4.c04.item.7-5-3}
\expandafter\def\csname sga@target@IV.7.5.4\endcsname{ega4.c04.item.7-5-4}
\expandafter\def\csname sga@target@IV.7.5.5\endcsname{ega4.c04.lem.7-5-5}
\expandafter\def\csname sga@target@IV.7.5.6\endcsname{ega4.c04.prop.7-5-6}
\expandafter\def\csname sga@target@IV.7.5.7\endcsname{ega4.c04.cor.7-5-7}
\expandafter\def\csname sga@target@IV.7.5.7.1\endcsname{ega4.c04.lem.7-5-7-1}
\expandafter\def\csname sga@target@IV.7.5.8\endcsname{ega4.c04.rem.7-5-8}
\expandafter\def\csname sga@target@IV.7.6.1\endcsname{ega4.c04.prop.7-6-1}
\expandafter\def\csname sga@target@IV.7.6.2\endcsname{ega4.c04.cor.7-6-2}
\expandafter\def\csname sga@target@IV.7.6.3\endcsname{ega4.c04.cor.7-6-3}
\expandafter\def\csname sga@target@IV.7.6.4\endcsname{ega4.c04.thm.7-6-4}
\expandafter\def\csname sga@target@IV.7.6.4.1\endcsname{ega4.c04.lem.7-6-4-1}
\expandafter\def\csname sga@target@IV.7.6.5\endcsname{ega4.c04.cor.7-6-5}
\expandafter\def\csname sga@target@IV.7.6.6\endcsname{ega4.c04.cor.7-6-6}
\expandafter\def\csname sga@target@IV.7.6.7\endcsname{ega4.c04.item.7-6-7}
\expandafter\def\csname sga@target@IV.7.7.1\endcsname{ega4.c04.statement.7-7-1}
\expandafter\def\csname sga@target@IV.7.7.2\endcsname{ega4.c04.thm.7-7-2}
\expandafter\def\csname sga@target@IV.7.7.3\endcsname{ega4.c04.cor.7-7-3}
\expandafter\def\csname sga@target@IV.7.7.4\endcsname{ega4.c04.cor.7-7-4}
\expandafter\def\csname sga@target@IV.7.8.1\endcsname{ega4.c04.statement.7-8-1}
\expandafter\def\csname sga@target@IV.7.8.2\endcsname{ega4.c04.def.7-8-2}
\expandafter\def\csname sga@target@IV.7.8.3\endcsname{ega4.c04.statement.7-8-3}
\expandafter\def\csname sga@target@IV.7.8.3.1\endcsname{ega4.c04.statement.7-8-3-1}
\expandafter\def\csname sga@target@IV.7.8.4\endcsname{ega4.c04.item.7-8-4}
\expandafter\def\csname sga@target@IV.7.8.4.1\endcsname{ega4.c04.statement.7-8-4-1}
\expandafter\def\csname sga@target@IV.7.8.5\endcsname{ega4.c04.statement.7-8-5}
\expandafter\def\csname sga@target@IV.7.8.6\endcsname{ega4.c04.prop.7-8-6}
\expandafter\def\csname sga@target@IV.7.9.1\endcsname{ega4.c04.statement.7-9-1}
\expandafter\def\csname sga@target@IV.7.9.10\endcsname{ega4.c04.item.7-9-10}
\expandafter\def\csname sga@target@IV.7.9.2\endcsname{ega4.c04.statement.7-9-2}
\expandafter\def\csname sga@target@IV.7.9.3\endcsname{ega4.c04.prop.7-9-3}
\expandafter\def\csname sga@target@IV.7.9.3.1\endcsname{ega4.c04.lem.7-9-3-1}
\expandafter\def\csname sga@target@IV.7.9.4\endcsname{ega4.c04.cor.7-9-4}
\expandafter\def\csname sga@target@IV.7.9.5\endcsname{ega4.c04.prop.7-9-5}
\expandafter\def\csname sga@target@IV.7.9.6\endcsname{ega4.c04.rem.7-9-6}
\expandafter\def\csname sga@target@IV.7.9.7\endcsname{ega4.c04.statement.7-9-7}
\expandafter\def\csname sga@target@IV.7.9.8\endcsname{ega4.c04.prop.7-9-8}
\expandafter\def\csname sga@target@IV.7.9.8.1\endcsname{ega4.c04.statement.7-9-8-1}
\expandafter\def\csname sga@target@IV.7.9.9\endcsname{ega4.c04.cor.7-9-9}
\expandafter\def\csname sga@target@IV.8.1.1\endcsname{ega4.c04.statement.8-1-1}
\expandafter\def\csname sga@target@IV.8.1.2\endcsname{ega4.c04.statement.8-1-2}
\expandafter\def\csname sga@target@IV.8.1.3\endcsname{ega4.c04.statement.8-1-3}
\expandafter\def\csname sga@target@IV.8.10.1\endcsname{ega4.c04.prop.8-10-1}
\expandafter\def\csname sga@target@IV.8.10.2\endcsname{ega4.c04.cor.8-10-2}
\expandafter\def\csname sga@target@IV.8.10.3\endcsname{ega4.c04.prop.8-10-3}
\expandafter\def\csname sga@target@IV.8.10.3.1\endcsname{ega4.c04.rem.8-10-3-1}
\expandafter\def\csname sga@target@IV.8.10.4\endcsname{ega4.c04.prop.8-10-4}
\expandafter\def\csname sga@target@IV.8.10.5\endcsname{ega4.c04.thm.8-10-5}
\expandafter\def\csname sga@target@IV.8.10.5.1\endcsname{ega4.c04.prop.8-10-5-1}
\expandafter\def\csname sga@target@IV.8.10.5.2\endcsname{ega4.c04.lem.8-10-5-2}
\expandafter\def\csname sga@target@IV.8.11\endcsname{ega4.c04.subsec.8-11.key-iv-8-11}
\expandafter\def\csname sga@target@IV.8.11.1\endcsname{ega4.c04.thm.8-11-1}
\expandafter\def\csname sga@target@IV.8.11.2\endcsname{ega4.c04.thm.8-11-2}
\expandafter\def\csname sga@target@IV.8.11.3\endcsname{ega4.c04.rem.8-11-3}
\expandafter\def\csname sga@target@IV.8.11.4\endcsname{ega4.c04.statement.8-11-4}
\expandafter\def\csname sga@target@IV.8.11.5\endcsname{ega4.c04.prop.8-11-5}
\expandafter\def\csname sga@target@IV.8.11.5.1\endcsname{ega4.c04.rem.8-11-5-1}
\expandafter\def\csname sga@target@IV.8.11.6\endcsname{ega4.c04.prop.8-11-6}
\expandafter\def\csname sga@target@IV.8.12\endcsname{ega4.c04.subsec.8-12.key-iv-8-12}
\expandafter\def\csname sga@target@IV.8.12.1\endcsname{ega4.c04.lem.8-12-1}
\expandafter\def\csname sga@target@IV.8.12.10\endcsname{ega4.c04.cor.8-12-10}
\expandafter\def\csname sga@target@IV.8.12.10.1\endcsname{ega4.c04.lem.8-12-10-1}
\expandafter\def\csname sga@target@IV.8.12.11\endcsname{ega4.c04.cor.8-12-11}
\expandafter\def\csname sga@target@IV.8.12.12\endcsname{ega4.c04.cor.8-12-12}
\expandafter\def\csname sga@target@IV.8.12.2\endcsname{ega4.c04.lem.8-12-2}
\expandafter\def\csname sga@target@IV.8.12.3\endcsname{ega4.c04.prop.8-12-3}
\expandafter\def\csname sga@target@IV.8.12.4\endcsname{ega4.c04.cor.8-12-4}
\expandafter\def\csname sga@target@IV.8.12.5\endcsname{ega4.c04.statement.8-12-5}
\expandafter\def\csname sga@target@IV.8.12.6\endcsname{ega4.c04.thm.8-12-6}
\expandafter\def\csname sga@target@IV.8.12.6.1\endcsname{ega4.c04.eq.8-12-6-1}
\expandafter\def\csname sga@target@IV.8.12.7\endcsname{ega4.c04.rem.8-12-7}
\expandafter\def\csname sga@target@IV.8.12.8\endcsname{ega4.c04.cor.8-12-8}
\expandafter\def\csname sga@target@IV.8.12.9\endcsname{ega4.c04.cor.8-12-9}
\expandafter\def\csname sga@target@IV.8.13.1\endcsname{ega4.c04.prop.8-13-1}
\expandafter\def\csname sga@target@IV.8.13.1.1\endcsname{ega4.c04.eq.8-13-1-1}
\expandafter\def\csname sga@target@IV.8.13.2\endcsname{ega4.c04.cor.8-13-2}
\expandafter\def\csname sga@target@IV.8.13.2.1\endcsname{ega4.c04.eq.8-13-2-1}
\expandafter\def\csname sga@target@IV.8.13.3\endcsname{ega4.c04.statement.8-13-3}
\expandafter\def\csname sga@target@IV.8.13.4\endcsname{ega4.c04.statement.8-13-4}
\expandafter\def\csname sga@target@IV.8.13.4.1\endcsname{ega4.c04.eq.8-13-4-1}
\expandafter\def\csname sga@target@IV.8.13.4.2\endcsname{ega4.c04.eq.8-13-4-2}
\expandafter\def\csname sga@target@IV.8.13.5\endcsname{ega4.c04.prop.8-13-5}
\expandafter\def\csname sga@target@IV.8.13.6\endcsname{ega4.c04.rem.8-13-6}
\expandafter\def\csname sga@target@IV.8.14.1\endcsname{ega4.c04.statement.8-14-1}
\expandafter\def\csname sga@target@IV.8.14.2\endcsname{ega4.c04.prop.8-14-2}
\expandafter\def\csname sga@target@IV.8.14.2.1\endcsname{ega4.c04.eq.8-14-2-1}
\expandafter\def\csname sga@target@IV.8.14.2.2\endcsname{ega4.c04.cor.8-14-2-2}
\expandafter\def\csname sga@target@IV.8.14.2.3\endcsname{ega4.c04.eq.8-14-2-3}
\expandafter\def\csname sga@target@IV.8.14.2.4\endcsname{ega4.c04.eq.8-14-2-4}
\expandafter\def\csname sga@target@IV.8.14.3\endcsname{ega4.c04.rem.8-14-3}
\expandafter\def\csname sga@target@IV.8.2.1\endcsname{ega4.c04.statement.8-2-1}
\expandafter\def\csname sga@target@IV.8.2.1.1\endcsname{ega4.c04.eq.8-2-1-1}
\expandafter\def\csname sga@target@IV.8.2.10\endcsname{ega4.c04.cor.8-2-10}
\expandafter\def\csname sga@target@IV.8.2.11\endcsname{ega4.c04.cor.8-2-11}
\expandafter\def\csname sga@target@IV.8.2.12\endcsname{ega4.c04.cor.8-2-12}
\expandafter\def\csname sga@target@IV.8.2.12.1\endcsname{ega4.c04.eq.8-2-12-1}
\expandafter\def\csname sga@target@IV.8.2.13\endcsname{ega4.c04.prop.8-2-13}
\expandafter\def\csname sga@target@IV.8.2.14\endcsname{ega4.c04.rem.8-2-14}
\expandafter\def\csname sga@target@IV.8.2.2\endcsname{ega4.c04.statement.8-2-2}
\expandafter\def\csname sga@target@IV.8.2.3\endcsname{ega4.c04.prop.8-2-3}
\expandafter\def\csname sga@target@IV.8.2.3.1\endcsname{ega4.c04.eq.8-2-3-1}
\expandafter\def\csname sga@target@IV.8.2.4\endcsname{ega4.c04.lem.8-2-4}
\expandafter\def\csname sga@target@IV.8.2.5\endcsname{ega4.c04.prop.8-2-5}
\expandafter\def\csname sga@target@IV.8.2.6\endcsname{ega4.c04.lem.8-2-6}
\expandafter\def\csname sga@target@IV.8.2.7\endcsname{ega4.c04.rem.8-2-7}
\expandafter\def\csname sga@target@IV.8.2.8\endcsname{ega4.c04.statement.8-2-8}
\expandafter\def\csname sga@target@IV.8.2.9\endcsname{ega4.c04.prop.8-2-9}
\expandafter\def\csname sga@target@IV.8.3.1\endcsname{ega4.c04.statement.8-3-1}
\expandafter\def\csname sga@target@IV.8.3.10\endcsname{ega4.c04.statement.8-3-10}
\expandafter\def\csname sga@target@IV.8.3.10.1\endcsname{ega4.c04.eq.8-3-10-1}
\expandafter\def\csname sga@target@IV.8.3.11\endcsname{ega4.c04.thm.8-3-11}
\expandafter\def\csname sga@target@IV.8.3.12\endcsname{ega4.c04.cor.8-3-12}
\expandafter\def\csname sga@target@IV.8.3.13\endcsname{ega4.c04.prop.8-3-13}
\expandafter\def\csname sga@target@IV.8.3.2\endcsname{ega4.c04.thm.8-3-2}
\expandafter\def\csname sga@target@IV.8.3.2.1\endcsname{ega4.c04.eq.8-3-2-1}
\expandafter\def\csname sga@target@IV.8.3.3\endcsname{ega4.c04.cor.8-3-3}
\expandafter\def\csname sga@target@IV.8.3.4\endcsname{ega4.c04.cor.8-3-4}
\expandafter\def\csname sga@target@IV.8.3.5\endcsname{ega4.c04.cor.8-3-5}
\expandafter\def\csname sga@target@IV.8.3.6\endcsname{ega4.c04.cor.8-3-6}
\expandafter\def\csname sga@target@IV.8.3.7\endcsname{ega4.c04.cor.8-3-7}
\expandafter\def\csname sga@target@IV.8.3.7.1\endcsname{ega4.c04.eq.8-3-7-1}
\expandafter\def\csname sga@target@IV.8.3.8\endcsname{ega4.c04.prop.8-3-8}
\expandafter\def\csname sga@target@IV.8.3.9\endcsname{ega4.c04.statement.8-3-9}
\expandafter\def\csname sga@target@IV.8.3.9.1\endcsname{ega4.c04.eq.8-3-9-1}
\expandafter\def\csname sga@target@IV.8.3.9.2\endcsname{ega4.c04.eq.8-3-9-2}
\expandafter\def\csname sga@target@IV.8.3.9.3\endcsname{ega4.c04.eq.8-3-9-3}
\expandafter\def\csname sga@target@IV.8.3.9.4\endcsname{ega4.c04.eq.8-3-9-4}
\expandafter\def\csname sga@target@IV.8.3.9.5\endcsname{ega4.c04.eq.8-3-9-5}
\expandafter\def\csname sga@target@IV.8.4.1\endcsname{ega4.c04.prop.8-4-1}
\expandafter\def\csname sga@target@IV.8.4.2\endcsname{ega4.c04.prop.8-4-2}
\expandafter\def\csname sga@target@IV.8.4.2.1\endcsname{ega4.c04.lem.8-4-2-1}
\expandafter\def\csname sga@target@IV.8.4.3\endcsname{ega4.c04.cor.8-4-3}
\expandafter\def\csname sga@target@IV.8.4.4\endcsname{ega4.c04.prop.8-4-4}
\expandafter\def\csname sga@target@IV.8.4.5\endcsname{ega4.c04.cor.8-4-5}
\expandafter\def\csname sga@target@IV.8.4.6\endcsname{ega4.c04.rem.8-4-6}
\expandafter\def\csname sga@target@IV.8.5.1\endcsname{ega4.c04.statement.8-5-1}
\expandafter\def\csname sga@target@IV.8.5.1.1\endcsname{ega4.c04.eq.8-5-1-1}
\expandafter\def\csname sga@target@IV.8.5.1.2\endcsname{ega4.c04.eq.8-5-1-2}
\expandafter\def\csname sga@target@IV.8.5.1.3\endcsname{ega4.c04.eq.8-5-1-3}
\expandafter\def\csname sga@target@IV.8.5.1.4\endcsname{ega4.c04.eq.8-5-1-4}
\expandafter\def\csname sga@target@IV.8.5.10\endcsname{ega4.c04.statement.8-5-10}
\expandafter\def\csname sga@target@IV.8.5.10.1\endcsname{ega4.c04.eq.8-5-10-1}
\expandafter\def\csname sga@target@IV.8.5.10.2\endcsname{ega4.c04.eq.8-5-10-2}
\expandafter\def\csname sga@target@IV.8.5.11\endcsname{ega4.c04.prop.8-5-11}
\expandafter\def\csname sga@target@IV.8.5.2\endcsname{ega4.c04.thm.8-5-2}
\expandafter\def\csname sga@target@IV.8.5.2.1\endcsname{ega4.c04.lem.8-5-2-1}
\expandafter\def\csname sga@target@IV.8.5.2.2\endcsname{ega4.c04.eq.8-5-2-2}
\expandafter\def\csname sga@target@IV.8.5.2.3\endcsname{ega4.c04.eq.8-5-2-3}
\expandafter\def\csname sga@target@IV.8.5.2.4\endcsname{ega4.c04.cor.8-5-2-4}
\expandafter\def\csname sga@target@IV.8.5.2.5\endcsname{ega4.c04.cor.8-5-2-5}
\expandafter\def\csname sga@target@IV.8.5.3\endcsname{ega4.c04.statement.8-5-3}
\expandafter\def\csname sga@target@IV.8.5.4\endcsname{ega4.c04.cor.8-5-4}
\expandafter\def\csname sga@target@IV.8.5.5\endcsname{ega4.c04.prop.8-5-5}
\expandafter\def\csname sga@target@IV.8.5.6\endcsname{ega4.c04.prop.8-5-6}
\expandafter\def\csname sga@target@IV.8.5.7\endcsname{ega4.c04.cor.8-5-7}
\expandafter\def\csname sga@target@IV.8.5.8\endcsname{ega4.c04.cor.8-5-8}
\expandafter\def\csname sga@target@IV.8.5.9\endcsname{ega4.c04.lem.8-5-9}
\expandafter\def\csname sga@target@IV.8.6.1\endcsname{ega4.c04.statement.8-6-1}
\expandafter\def\csname sga@target@IV.8.6.2\endcsname{ega4.c04.statement.8-6-2}
\expandafter\def\csname sga@target@IV.8.6.2.1\endcsname{ega4.c04.eq.8-6-2-1}
\expandafter\def\csname sga@target@IV.8.6.2.2\endcsname{ega4.c04.eq.8-6-2-2}
\expandafter\def\csname sga@target@IV.8.6.2.3\endcsname{ega4.c04.eq.8-6-2-3}
\expandafter\def\csname sga@target@IV.8.6.3\endcsname{ega4.c04.prop.8-6-3}
\expandafter\def\csname sga@target@IV.8.6.3.1\endcsname{ega4.c04.cor.8-6-3-1}
\expandafter\def\csname sga@target@IV.8.6.4\endcsname{ega4.c04.cor.8-6-4}
\expandafter\def\csname sga@target@IV.8.7.1\endcsname{ega4.c04.prop.8-7-1}
\expandafter\def\csname sga@target@IV.8.7.2\endcsname{ega4.c04.prop.8-7-2}
\expandafter\def\csname sga@target@IV.8.7.3\endcsname{ega4.c04.cor.8-7-3}
\expandafter\def\csname sga@target@IV.8.7.4\endcsname{ega4.c04.rem.8-7-4}
\expandafter\def\csname sga@target@IV.8.8.1\endcsname{ega4.c04.statement.8-8-1}
\expandafter\def\csname sga@target@IV.8.8.1.1\endcsname{ega4.c04.eq.8-8-1-1}
\expandafter\def\csname sga@target@IV.8.8.2\endcsname{ega4.c04.thm.8-8-2}
\expandafter\def\csname sga@target@IV.8.8.2.1\endcsname{ega4.c04.lem.8-8-2-1}
\expandafter\def\csname sga@target@IV.8.8.2.2\endcsname{ega4.c04.lem.8-8-2-2}
\expandafter\def\csname sga@target@IV.8.8.2.3\endcsname{ega4.c04.lem.8-8-2-3}
\expandafter\def\csname sga@target@IV.8.8.2.4\endcsname{ega4.c04.cor.8-8-2-4}
\expandafter\def\csname sga@target@IV.8.8.2.5\endcsname{ega4.c04.cor.8-8-2-5}
\expandafter\def\csname sga@target@IV.8.8.3\endcsname{ega4.c04.statement.8-8-3}
\expandafter\def\csname sga@target@IV.8.9.1\endcsname{ega4.c04.prop.8-9-1}
\expandafter\def\csname sga@target@IV.8.9.2\endcsname{ega4.c04.statement.8-9-2}
\expandafter\def\csname sga@target@IV.8.9.3\endcsname{ega4.c04.prop.8-9-3}
\expandafter\def\csname sga@target@IV.8.9.4\endcsname{ega4.c04.prop.8-9-4}
\expandafter\def\csname sga@target@IV.8.9.4.1\endcsname{ega4.c04.lem.8-9-4-1}
\expandafter\def\csname sga@target@IV.8.9.5\endcsname{ega4.c04.cor.8-9-5}
\expandafter\def\csname sga@target@IV.9.1.1\endcsname{ega4.c04.prop.9-1-1}
\expandafter\def\csname sga@target@IV.9.1.2\endcsname{ega4.c04.cor.9-1-2}
\expandafter\def\csname sga@target@IV.9.1.3\endcsname{ega4.c04.rem.9-1-3}
\expandafter\def\csname sga@target@IV.9.1.4\endcsname{ega4.c04.prop.9-1-4}
\expandafter\def\csname sga@target@IV.9.1.5\endcsname{ega4.c04.cor.9-1-5}
\expandafter\def\csname sga@target@IV.9.1.6\endcsname{ega4.c04.prop.9-1-6}
\expandafter\def\csname sga@target@IV.9.2\endcsname{ega4.c04.subsec.9-2.key-iv-9-2}
\expandafter\def\csname sga@target@IV.9.2.1\endcsname{ega4.c04.def.9-2-1}
\expandafter\def\csname sga@target@IV.9.2.2\endcsname{ega4.c04.item.9-2-2}
\expandafter\def\csname sga@target@IV.9.2.3\endcsname{ega4.c04.prop.9-2-3}
\expandafter\def\csname sga@target@IV.9.2.4\endcsname{ega4.c04.cor.9-2-4}
\expandafter\def\csname sga@target@IV.9.2.5\endcsname{ega4.c04.rem.9-2-5}
\expandafter\def\csname sga@target@IV.9.2.6\endcsname{ega4.c04.prop.9-2-6}
\expandafter\def\csname sga@target@IV.9.2.6.1\endcsname{ega4.c04.prop.9-2-6-1}
\expandafter\def\csname sga@target@IV.9.2.6.2\endcsname{ega4.c04.cor.9-2-6-2}
\expandafter\def\csname sga@target@IV.9.3.1\endcsname{ega4.c04.prop.9-3-1}
\expandafter\def\csname sga@target@IV.9.3.2\endcsname{ega4.c04.prop.9-3-2}
\expandafter\def\csname sga@target@IV.9.3.3\endcsname{ega4.c04.prop.9-3-3}
\expandafter\def\csname sga@target@IV.9.3.4\endcsname{ega4.c04.rem.9-3-4}
\expandafter\def\csname sga@target@IV.9.3.5\endcsname{ega4.c04.prop.9-3-5}
\expandafter\def\csname sga@target@IV.9.3.6\endcsname{ega4.c04.statement.9-3-6}
\expandafter\def\csname sga@target@IV.9.4.1\endcsname{ega4.c04.statement.9-4-1}
\expandafter\def\csname sga@target@IV.9.4.2\endcsname{ega4.c04.prop.9-4-2}
\expandafter\def\csname sga@target@IV.9.4.3\endcsname{ega4.c04.cor.9-4-3}
\expandafter\def\csname sga@target@IV.9.4.3.1\endcsname{ega4.c04.eq.9-4-3-1}
\expandafter\def\csname sga@target@IV.9.4.4\endcsname{ega4.c04.prop.9-4-4}
\expandafter\def\csname sga@target@IV.9.4.5\endcsname{ega4.c04.cor.9-4-5}
\expandafter\def\csname sga@target@IV.9.4.6\endcsname{ega4.c04.cor.9-4-6}
\expandafter\def\csname sga@target@IV.9.4.7\endcsname{ega4.c04.prop.9-4-7}
\expandafter\def\csname sga@target@IV.9.4.7.1\endcsname{ega4.c04.lem.9-4-7-1}
\expandafter\def\csname sga@target@IV.9.4.7.2\endcsname{ega4.c04.rem.9-4-7-2}
\expandafter\def\csname sga@target@IV.9.4.8\endcsname{ega4.c04.prop.9-4-8}
\expandafter\def\csname sga@target@IV.9.4.9\endcsname{ega4.c04.rem.9-4-9}
\expandafter\def\csname sga@target@IV.9.5.1\endcsname{ega4.c04.prop.9-5-1}
\expandafter\def\csname sga@target@IV.9.5.2\endcsname{ega4.c04.cor.9-5-2}
\expandafter\def\csname sga@target@IV.9.5.3\endcsname{ega4.c04.prop.9-5-3}
\expandafter\def\csname sga@target@IV.9.5.3.case2\endcsname{ega4.c04.item.9-5-3-case2}
\expandafter\def\csname sga@target@IV.9.5.4\endcsname{ega4.c04.cor.9-5-4}
\expandafter\def\csname sga@target@IV.9.5.5\endcsname{ega4.c04.prop.9-5-5}
\expandafter\def\csname sga@target@IV.9.5.6\endcsname{ega4.c04.cor.9-5-6}
\expandafter\def\csname sga@target@IV.9.5.7\endcsname{ega4.c04.rem.9-5-7}
\expandafter\def\csname sga@target@IV.9.6.1\endcsname{ega4.c04.prop.9-6-1}
\expandafter\def\csname sga@target@IV.9.6.2\endcsname{ega4.c04.prop.9-6-2}
\expandafter\def\csname sga@target@IV.9.6.3\endcsname{ega4.c04.prop.9-6-3}
\expandafter\def\csname sga@target@IV.9.6.3.1\endcsname{ega4.c04.eq.9-6-3-1}
\expandafter\def\csname sga@target@IV.9.6.4\endcsname{ega4.c04.cor.9-6-4}
\expandafter\def\csname sga@target@IV.9.6.5\endcsname{ega4.c04.cor.9-6-5}
\expandafter\def\csname sga@target@IV.9.6.6\endcsname{ega4.c04.sec.9-6-6}
\expandafter\def\csname sga@target@IV.9.6.7\endcsname{ega4.c04.sec.9-6-7}
\expandafter\def\csname sga@target@IV.9.7\endcsname{ega4.c04.subsec.9-7.key-iv-9-7}
\expandafter\def\csname sga@target@IV.9.7.1\endcsname{ega4.c04.lem.9-7-1}
\expandafter\def\csname sga@target@IV.9.7.10\endcsname{ega4.c04.cor.9-7-10}
\expandafter\def\csname sga@target@IV.9.7.11\endcsname{ega4.c04.rem.9-7-11}
\expandafter\def\csname sga@target@IV.9.7.12\endcsname{ega4.c04.prop.9-7-12}
\expandafter\def\csname sga@target@IV.9.7.2\endcsname{ega4.c04.lem.9-7-2}
\expandafter\def\csname sga@target@IV.9.7.3\endcsname{ega4.c04.subsec.9-7-3}
\expandafter\def\csname sga@target@IV.9.7.4\endcsname{ega4.c04.lem.9-7-4}
\expandafter\def\csname sga@target@IV.9.7.5\endcsname{ega4.c04.lem.9-7-5}
\expandafter\def\csname sga@target@IV.9.7.5.1\endcsname{ega4.c04.lem.9-7-5-1}
\expandafter\def\csname sga@target@IV.9.7.6\endcsname{ega4.c04.prop.9-7-6}
\expandafter\def\csname sga@target@IV.9.7.7\endcsname{ega4.c04.thm.9-7-7}
\expandafter\def\csname sga@target@IV.9.7.8\endcsname{ega4.c04.prop.9-7-8}
\expandafter\def\csname sga@target@IV.9.7.9\endcsname{ega4.c04.cor.9-7-9}
\expandafter\def\csname sga@target@IV.9.8\endcsname{ega4.c04.subsec.9-8.key-iv-9-8}
\expandafter\def\csname sga@target@IV.9.8.1\endcsname{ega4.c04.prop.9-8-1}
\expandafter\def\csname sga@target@IV.9.8.1-proof\endcsname{ega4.c04.subsec.9-8-1-proof}
\expandafter\def\csname sga@target@IV.9.8.2\endcsname{ega4.c04.subsec.9-8-2}
\expandafter\def\csname sga@target@IV.9.8.3\endcsname{ega4.c04.prop.9-8-3}
\expandafter\def\csname sga@target@IV.9.8.3-proof\endcsname{ega4.c04.subsec.9-8-3-proof}
\expandafter\def\csname sga@target@IV.9.8.4\endcsname{ega4.c04.cor.9-8-4}
\expandafter\def\csname sga@target@IV.9.8.4-proof\endcsname{ega4.c04.item.9-8-4-proof}
\expandafter\def\csname sga@target@IV.9.8.5\endcsname{ega4.c04.prop.9-8-5}
\expandafter\def\csname sga@target@IV.9.8.5-proof\endcsname{ega4.c04.item.9-8-5-proof}
\expandafter\def\csname sga@target@IV.9.8.6\endcsname{ega4.c04.prop.9-8-6}
\expandafter\def\csname sga@target@IV.9.8.6-proof\endcsname{ega4.c04.item.9-8-6-proof}
\expandafter\def\csname sga@target@IV.9.8.7\endcsname{ega4.c04.cor.9-8-7}
\expandafter\def\csname sga@target@IV.9.8.7-proof\endcsname{ega4.c04.item.9-8-7-proof}
\expandafter\def\csname sga@target@IV.9.8.8\endcsname{ega4.c04.prop.9-8-8}
\expandafter\def\csname sga@target@IV.9.8.8-proof\endcsname{ega4.c04.item.9-8-8-proof}
\expandafter\def\csname sga@target@IV.9.8.9\endcsname{ega4.c04.rem.9-8-9}
\expandafter\def\csname sga@target@IV.9.8.9.1\endcsname{ega4.c04.item.9-8-9-1}
\expandafter\def\csname sga@target@IV.9.9\endcsname{ega4.c04.subsec.9-9.key-iv-9-9}
\expandafter\def\csname sga@target@IV.9.9.1\endcsname{ega4.c04.prop.9-9-1}
\expandafter\def\csname sga@target@IV.9.9.1-proof\endcsname{ega4.c04.item.9-9-1-proof}
\expandafter\def\csname sga@target@IV.9.9.1.1\endcsname{ega4.c04.eq.9-9-1-1}
\expandafter\def\csname sga@target@IV.9.9.2\endcsname{ega4.c04.prop.9-9-2}
\expandafter\def\csname sga@target@IV.9.9.2-proof\endcsname{ega4.c04.item.9-9-2-proof}
\expandafter\def\csname sga@target@IV.9.9.2.1\endcsname{ega4.c04.eq.9-9-2-1}
\expandafter\def\csname sga@target@IV.9.9.2.2\endcsname{ega4.c04.eq.9-9-2-2}
\expandafter\def\csname sga@target@IV.9.9.3\endcsname{ega4.c04.cor.9-9-3}
\expandafter\def\csname sga@target@IV.9.9.3-proof\endcsname{ega4.c04.statement.9-9-3-proof}
\expandafter\def\csname sga@target@IV.9.9.4\endcsname{ega4.c04.prop.9-9-4}
\expandafter\def\csname sga@target@IV.9.9.4-proof\endcsname{ega4.c04.item.9-9-4-proof}
\expandafter\def\csname sga@target@IV.9.9.5\endcsname{ega4.c04.cor.9-9-5}
\expandafter\def\csname sga@target@IV.9.9.5-proof\endcsname{ega4.c04.item.9-9-5-proof}
\expandafter\def\csname sga@target@IV.9.9.6\endcsname{ega4.c04.prop.9-9-6}
\expandafter\def\csname sga@target@IV.9.9.6-proof\endcsname{ega4.c04.item.9-9-6-proof}
\expandafter\def\csname sga@target@section:IV.1\endcsname{ega4.c04.sec.1}
\expandafter\def\csname sga@target@section:IV.10\endcsname{ega4.c04.sec.10}
\expandafter\def\csname sga@target@section:IV.11\endcsname{ega4.c04.sec.11}
\expandafter\def\csname sga@target@section:IV.12\endcsname{ega4.c04.sec.12}
\expandafter\def\csname sga@target@section:IV.13\endcsname{ega4.c04.sec.13}
\expandafter\def\csname sga@target@section:IV.14\endcsname{ega4.c04.sec.14}
\expandafter\def\csname sga@target@section:IV.15\endcsname{ega4.c04.sec.15}
\expandafter\def\csname sga@target@section:IV.16\endcsname{ega4.c04.sec.16}
\expandafter\def\csname sga@target@section:IV.17\endcsname{ega4.c04.sec.17}
\expandafter\def\csname sga@target@section:IV.18\endcsname{ega4.c04.sec.18}
\expandafter\def\csname sga@target@section:IV.19\endcsname{ega4.c04.sec.19}
\expandafter\def\csname sga@target@section:IV.2\endcsname{ega4.c04.sec.2}
\expandafter\def\csname sga@target@section:IV.21\endcsname{ega4.c04.sec.21}
\expandafter\def\csname sga@target@section:IV.3\endcsname{ega4.c04.sec.3}
\expandafter\def\csname sga@target@section:IV.4\endcsname{ega4.c04.sec.4}
\expandafter\def\csname sga@target@section:IV.5\endcsname{ega4.c04.sec.5}
\expandafter\def\csname sga@target@section:IV.6\endcsname{ega4.c04.sec.6}
\expandafter\def\csname sga@target@section:IV.7\endcsname{ega4.c04.sec.7}
\expandafter\def\csname sga@target@section:IV.8\endcsname{ega4.c04.sec.8}
\expandafter\def\csname sga@target@section:IV.9\endcsname{ega4.c04.sec.9}
\expandafter\def\csname sga@target@subsection:IV.1.1\endcsname{ega4.c04.subsec.1-1}
\expandafter\def\csname sga@target@subsection:IV.1.10\endcsname{ega4.c04.subsec.1-10}
\expandafter\def\csname sga@target@subsection:IV.1.2\endcsname{ega4.c04.subsec.1-2}
\expandafter\def\csname sga@target@subsection:IV.1.3\endcsname{ega4.c04.subsec.1-3}
\expandafter\def\csname sga@target@subsection:IV.1.4\endcsname{ega4.c04.subsec.1-4}
\expandafter\def\csname sga@target@subsection:IV.1.5\endcsname{ega4.c04.subsec.1-5}
\expandafter\def\csname sga@target@subsection:IV.1.6\endcsname{ega4.c04.subsec.1-6}
\expandafter\def\csname sga@target@subsection:IV.1.7\endcsname{ega4.c04.subsec.1-7}
\expandafter\def\csname sga@target@subsection:IV.1.8\endcsname{ega4.c04.subsec.1-8}
\expandafter\def\csname sga@target@subsection:IV.1.9\endcsname{ega4.c04.subsec.1-9}
\expandafter\def\csname sga@target@subsection:IV.10.1\endcsname{ega4.c04.subsec.10-1}
\expandafter\def\csname sga@target@subsection:IV.10.10\endcsname{ega4.c04.subsec.10-10}
\expandafter\def\csname sga@target@subsection:IV.10.2\endcsname{ega4.c04.subsec.10-2}
\expandafter\def\csname sga@target@subsection:IV.10.3\endcsname{ega4.c04.subsec.10-3}
\expandafter\def\csname sga@target@subsection:IV.10.4\endcsname{ega4.c04.subsec.10-4}
\expandafter\def\csname sga@target@subsection:IV.10.5\endcsname{ega4.c04.subsec.10-5}
\expandafter\def\csname sga@target@subsection:IV.10.6\endcsname{ega4.c04.subsec.10-6}
\expandafter\def\csname sga@target@subsection:IV.10.7\endcsname{ega4.c04.subsec.10-7}
\expandafter\def\csname sga@target@subsection:IV.10.8\endcsname{ega4.c04.subsec.10-8}
\expandafter\def\csname sga@target@subsection:IV.10.9\endcsname{ega4.c04.subsec.10-9}
\expandafter\def\csname sga@target@subsection:IV.11.1\endcsname{ega4.c04.subsec.11-1}
\expandafter\def\csname sga@target@subsection:IV.11.10\endcsname{ega4.c04.subsec.11-10}
\expandafter\def\csname sga@target@subsection:IV.11.2\endcsname{ega4.c04.subsec.11-2}
\expandafter\def\csname sga@target@subsection:IV.11.3\endcsname{ega4.c04.subsec.11-3}
\expandafter\def\csname sga@target@subsection:IV.11.4\endcsname{ega4.c04.subsec.11-4}
\expandafter\def\csname sga@target@subsection:IV.11.5\endcsname{ega4.c04.subsec.11-5}
\expandafter\def\csname sga@target@subsection:IV.11.6\endcsname{ega4.c04.subsec.11-6}
\expandafter\def\csname sga@target@subsection:IV.11.7\endcsname{ega4.c04.subsec.11-7}
\expandafter\def\csname sga@target@subsection:IV.11.8\endcsname{ega4.c04.subsec.11-8}
\expandafter\def\csname sga@target@subsection:IV.11.9\endcsname{ega4.c04.subsec.11-9}
\expandafter\def\csname sga@target@subsection:IV.12.0\endcsname{ega4.c04.subsec.12-0}
\expandafter\def\csname sga@target@subsection:IV.12.1\endcsname{ega4.c04.subsec.12-1}
\expandafter\def\csname sga@target@subsection:IV.12.2\endcsname{ega4.c04.subsec.12-2}
\expandafter\def\csname sga@target@subsection:IV.12.3\endcsname{ega4.c04.subsec.12-3}
\expandafter\def\csname sga@target@subsection:IV.13.1\endcsname{ega4.c04.subsec.13-1}
\expandafter\def\csname sga@target@subsection:IV.13.2\endcsname{ega4.c04.subsec.13-2}
\expandafter\def\csname sga@target@subsection:IV.13.3\endcsname{ega4.c04.subsec.13-3}
\expandafter\def\csname sga@target@subsection:IV.14.1\endcsname{ega4.c04.subsec.14-1}
\expandafter\def\csname sga@target@subsection:IV.14.2\endcsname{ega4.c04.subsec.14-2}
\expandafter\def\csname sga@target@subsection:IV.14.3\endcsname{ega4.c04.subsec.14-3}
\expandafter\def\csname sga@target@subsection:IV.14.4\endcsname{ega4.c04.subsec.14-4}
\expandafter\def\csname sga@target@subsection:IV.14.5\endcsname{ega4.c04.subsec.14-5}
\expandafter\def\csname sga@target@subsection:IV.15.1\endcsname{ega4.c04.subsec.15-1}
\expandafter\def\csname sga@target@subsection:IV.15.2\endcsname{ega4.c04.subsec.15-2}
\expandafter\def\csname sga@target@subsection:IV.15.3\endcsname{ega4.c04.subsec.15-3}
\expandafter\def\csname sga@target@subsection:IV.15.4\endcsname{ega4.c04.subsec.15-4}
\expandafter\def\csname sga@target@subsection:IV.15.5\endcsname{ega4.c04.subsec.15-5}
\expandafter\def\csname sga@target@subsection:IV.15.6\endcsname{ega4.c04.subsec.15-6}
\expandafter\def\csname sga@target@subsection:IV.15.7\endcsname{ega4.c04.subsec.15-7}
\expandafter\def\csname sga@target@subsection:IV.17.1\endcsname{ega4.c04.subsec.17-1}
\expandafter\def\csname sga@target@subsection:IV.17.10\endcsname{ega4.c04.subsec.17-10}
\expandafter\def\csname sga@target@subsection:IV.17.11\endcsname{ega4.c04.subsec.17-11}
\expandafter\def\csname sga@target@subsection:IV.17.12\endcsname{ega4.c04.subsec.17-12}
\expandafter\def\csname sga@target@subsection:IV.17.13\endcsname{ega4.c04.subsec.17-13}
\expandafter\def\csname sga@target@subsection:IV.17.14\endcsname{ega4.c04.subsec.17-14}
\expandafter\def\csname sga@target@subsection:IV.17.15\endcsname{ega4.c04.subsec.17-15}
\expandafter\def\csname sga@target@subsection:IV.17.16\endcsname{ega4.c04.subsec.17-16}
\expandafter\def\csname sga@target@subsection:IV.17.2\endcsname{ega4.c04.subsec.17-2}
\expandafter\def\csname sga@target@subsection:IV.17.3\endcsname{ega4.c04.subsec.17-3}
\expandafter\def\csname sga@target@subsection:IV.17.4\endcsname{ega4.c04.subsec.17-4}
\expandafter\def\csname sga@target@subsection:IV.17.5\endcsname{ega4.c04.subsec.17-5}
\expandafter\def\csname sga@target@subsection:IV.17.6\endcsname{ega4.c04.subsec.17-6}
\expandafter\def\csname sga@target@subsection:IV.17.7\endcsname{ega4.c04.subsec.17-7}
\expandafter\def\csname sga@target@subsection:IV.17.8\endcsname{ega4.c04.subsec.17-8}
\expandafter\def\csname sga@target@subsection:IV.17.9\endcsname{ega4.c04.subsec.17-9}
\expandafter\def\csname sga@target@subsection:IV.18.1\endcsname{ega4.c04.subsec.18-1}
\expandafter\def\csname sga@target@subsection:IV.18.10\endcsname{ega4.c04.subsec.18-10}
\expandafter\def\csname sga@target@subsection:IV.18.11\endcsname{ega4.c04.subsec.18-11}
\expandafter\def\csname sga@target@subsection:IV.18.12\endcsname{ega4.c04.subsec.18-12}
\expandafter\def\csname sga@target@subsection:IV.18.2\endcsname{ega4.c04.subsec.18-2}
\expandafter\def\csname sga@target@subsection:IV.18.3\endcsname{ega4.c04.subsec.18-3}
\expandafter\def\csname sga@target@subsection:IV.18.4\endcsname{ega4.c04.subsec.18-4}
\expandafter\def\csname sga@target@subsection:IV.18.5\endcsname{ega4.c04.subsec.18-5}
\expandafter\def\csname sga@target@subsection:IV.18.6\endcsname{ega4.c04.subsec.18-6}
\expandafter\def\csname sga@target@subsection:IV.18.7\endcsname{ega4.c04.subsec.18-7}
\expandafter\def\csname sga@target@subsection:IV.18.8\endcsname{ega4.c04.subsec.18-8}
\expandafter\def\csname sga@target@subsection:IV.18.9\endcsname{ega4.c04.subsec.18-9}
\expandafter\def\csname sga@target@subsection:IV.19.1\endcsname{ega4.c04.subsec.19-1}
\expandafter\def\csname sga@target@subsection:IV.19.2\endcsname{ega4.c04.subsec.19-2}
\expandafter\def\csname sga@target@subsection:IV.19.3\endcsname{ega4.c04.subsec.19-3}
\expandafter\def\csname sga@target@subsection:IV.19.4\endcsname{ega4.c04.subsec.19-4}
\expandafter\def\csname sga@target@subsection:IV.19.5\endcsname{ega4.c04.subsec.19-5}
\expandafter\def\csname sga@target@subsection:IV.19.6\endcsname{ega4.c04.subsec.19-6}
\expandafter\def\csname sga@target@subsection:IV.19.7\endcsname{ega4.c04.subsec.19-7}
\expandafter\def\csname sga@target@subsection:IV.19.8\endcsname{ega4.c04.subsec.19-8}
\expandafter\def\csname sga@target@subsection:IV.19.9\endcsname{ega4.c04.subsec.19-9}
\expandafter\def\csname sga@target@subsection:IV.2.1\endcsname{ega4.c04.subsec.2-1}
\expandafter\def\csname sga@target@subsection:IV.2.2\endcsname{ega4.c04.subsec.2-2}
\expandafter\def\csname sga@target@subsection:IV.2.3\endcsname{ega4.c04.subsec.2-3}
\expandafter\def\csname sga@target@subsection:IV.2.4\endcsname{ega4.c04.subsec.2-4}
\expandafter\def\csname sga@target@subsection:IV.2.5\endcsname{ega4.c04.subsec.2-5}
\expandafter\def\csname sga@target@subsection:IV.2.6\endcsname{ega4.c04.subsec.2-6}
\expandafter\def\csname sga@target@subsection:IV.2.7\endcsname{ega4.c04.subsec.2-7}
\expandafter\def\csname sga@target@subsection:IV.2.8\endcsname{ega4.c04.subsec.2-8}
\expandafter\def\csname sga@target@subsection:IV.21.1\endcsname{ega4.c04.subsec.21-1}
\expandafter\def\csname sga@target@subsection:IV.21.10\endcsname{ega4.c04.subsec.21-10}
\expandafter\def\csname sga@target@subsection:IV.21.11\endcsname{ega4.c04.subsec.21-11}
\expandafter\def\csname sga@target@subsection:IV.21.12\endcsname{ega4.c04.subsec.21-12}
\expandafter\def\csname sga@target@subsection:IV.21.13\endcsname{ega4.c04.subsec.21-13}
\expandafter\def\csname sga@target@subsection:IV.21.14\endcsname{ega4.c04.subsec.21-14}
\expandafter\def\csname sga@target@subsection:IV.21.15\endcsname{ega4.c04.subsec.21-15}
\expandafter\def\csname sga@target@subsection:IV.21.2\endcsname{ega4.c04.subsec.21-2}
\expandafter\def\csname sga@target@subsection:IV.21.3\endcsname{ega4.c04.subsec.21-3}
\expandafter\def\csname sga@target@subsection:IV.21.4\endcsname{ega4.c04.subsec.21-4}
\expandafter\def\csname sga@target@subsection:IV.21.5\endcsname{ega4.c04.subsec.21-5}
\expandafter\def\csname sga@target@subsection:IV.21.6\endcsname{ega4.c04.subsec.21-6}
\expandafter\def\csname sga@target@subsection:IV.21.7\endcsname{ega4.c04.subsec.21-7}
\expandafter\def\csname sga@target@subsection:IV.21.8\endcsname{ega4.c04.subsec.21-8}
\expandafter\def\csname sga@target@subsection:IV.21.9\endcsname{ega4.c04.subsec.21-9}
\expandafter\def\csname sga@target@subsection:IV.3.1\endcsname{ega4.c04.subsec.3-1}
\expandafter\def\csname sga@target@subsection:IV.3.2\endcsname{ega4.c04.subsec.3-2}
\expandafter\def\csname sga@target@subsection:IV.3.3\endcsname{ega4.c04.subsec.3-3}
\expandafter\def\csname sga@target@subsection:IV.3.4\endcsname{ega4.c04.subsec.3-4}
\expandafter\def\csname sga@target@subsection:IV.4.1\endcsname{ega4.c04.subsec.4-1}
\expandafter\def\csname sga@target@subsection:IV.4.2\endcsname{ega4.c04.subsec.4-2}
\expandafter\def\csname sga@target@subsection:IV.4.3\endcsname{ega4.c04.subsec.4-3}
\expandafter\def\csname sga@target@subsection:IV.4.4\endcsname{ega4.c04.subsec.4-4}
\expandafter\def\csname sga@target@subsection:IV.4.5\endcsname{ega4.c04.subsec.4-5}
\expandafter\def\csname sga@target@subsection:IV.4.6\endcsname{ega4.c04.subsec.4-6}
\expandafter\def\csname sga@target@subsection:IV.4.7\endcsname{ega4.c04.subsec.4-7}
\expandafter\def\csname sga@target@subsection:IV.4.8\endcsname{ega4.c04.subsec.4-8}
\expandafter\def\csname sga@target@subsection:IV.4.9\endcsname{ega4.c04.subsec.4-9}
\expandafter\def\csname sga@target@subsection:IV.5.1\endcsname{ega4.c04.subsec.5-1}
\expandafter\def\csname sga@target@subsection:IV.5.10\endcsname{ega4.c04.subsec.5-10}
\expandafter\def\csname sga@target@subsection:IV.5.11\endcsname{ega4.c04.subsec.5-11}
\expandafter\def\csname sga@target@subsection:IV.5.12\endcsname{ega4.c04.subsec.5-12}
\expandafter\def\csname sga@target@subsection:IV.5.13\endcsname{ega4.c04.subsec.5-13}
\expandafter\def\csname sga@target@subsection:IV.5.2\endcsname{ega4.c04.subsec.5-2}
\expandafter\def\csname sga@target@subsection:IV.5.3\endcsname{ega4.c04.subsec.5-3}
\expandafter\def\csname sga@target@subsection:IV.5.4\endcsname{ega4.c04.subsec.5-4}
\expandafter\def\csname sga@target@subsection:IV.5.5\endcsname{ega4.c04.subsec.5-5}
\expandafter\def\csname sga@target@subsection:IV.5.6\endcsname{ega4.c04.subsec.5-6}
\expandafter\def\csname sga@target@subsection:IV.5.7\endcsname{ega4.c04.subsec.5-7}
\expandafter\def\csname sga@target@subsection:IV.5.8\endcsname{ega4.c04.subsec.5-8}
\expandafter\def\csname sga@target@subsection:IV.5.9\endcsname{ega4.c04.subsec.5-9}
\expandafter\def\csname sga@target@subsection:IV.6.1\endcsname{ega4.c04.subsec.6-1}
\expandafter\def\csname sga@target@subsection:IV.6.10\endcsname{ega4.c04.subsec.6-10}
\expandafter\def\csname sga@target@subsection:IV.6.11\endcsname{ega4.c04.subsec.6-11}
\expandafter\def\csname sga@target@subsection:IV.6.12\endcsname{ega4.c04.subsec.6-12}
\expandafter\def\csname sga@target@subsection:IV.6.13\endcsname{ega4.c04.subsec.6-13}
\expandafter\def\csname sga@target@subsection:IV.6.14\endcsname{ega4.c04.subsec.6-14}
\expandafter\def\csname sga@target@subsection:IV.6.15\endcsname{ega4.c04.subsec.6-15}
\expandafter\def\csname sga@target@subsection:IV.6.2\endcsname{ega4.c04.subsec.6-2}
\expandafter\def\csname sga@target@subsection:IV.6.3\endcsname{ega4.c04.subsec.6-3}
\expandafter\def\csname sga@target@subsection:IV.6.4\endcsname{ega4.c04.subsec.6-4}
\expandafter\def\csname sga@target@subsection:IV.6.5\endcsname{ega4.c04.subsec.6-5}
\expandafter\def\csname sga@target@subsection:IV.6.6\endcsname{ega4.c04.subsec.6-6}
\expandafter\def\csname sga@target@subsection:IV.6.7\endcsname{ega4.c04.subsec.6-7}
\expandafter\def\csname sga@target@subsection:IV.6.8\endcsname{ega4.c04.subsec.6-8}
\expandafter\def\csname sga@target@subsection:IV.6.9\endcsname{ega4.c04.subsec.6-9}
\expandafter\def\csname sga@target@subsection:IV.7.1\endcsname{ega4.c04.subsec.7-1}
\expandafter\def\csname sga@target@subsection:IV.7.2\endcsname{ega4.c04.subsec.7-2}
\expandafter\def\csname sga@target@subsection:IV.7.3\endcsname{ega4.c04.subsec.7-3}
\expandafter\def\csname sga@target@subsection:IV.7.4\endcsname{ega4.c04.subsec.7-4}
\expandafter\def\csname sga@target@subsection:IV.7.5\endcsname{ega4.c04.subsec.7-5}
\expandafter\def\csname sga@target@subsection:IV.7.6\endcsname{ega4.c04.subsec.7-6}
\expandafter\def\csname sga@target@subsection:IV.7.7\endcsname{ega4.c04.subsec.7-7}
\expandafter\def\csname sga@target@subsection:IV.7.8\endcsname{ega4.c04.subsec.7-8}
\expandafter\def\csname sga@target@subsection:IV.7.9\endcsname{ega4.c04.subsec.7-9}
\expandafter\def\csname sga@target@subsection:IV.8.1\endcsname{ega4.c04.subsec.8-1}
\expandafter\def\csname sga@target@subsection:IV.8.10\endcsname{ega4.c04.subsec.8-10}
\expandafter\def\csname sga@target@subsection:IV.8.11\endcsname{ega4.c04.subsec.8-11}
\expandafter\def\csname sga@target@subsection:IV.8.12\endcsname{ega4.c04.subsec.8-12}
\expandafter\def\csname sga@target@subsection:IV.8.13\endcsname{ega4.c04.subsec.8-13}
\expandafter\def\csname sga@target@subsection:IV.8.14\endcsname{ega4.c04.subsec.8-14}
\expandafter\def\csname sga@target@subsection:IV.8.2\endcsname{ega4.c04.subsec.8-2}
\expandafter\def\csname sga@target@subsection:IV.8.3\endcsname{ega4.c04.subsec.8-3}
\expandafter\def\csname sga@target@subsection:IV.8.4\endcsname{ega4.c04.subsec.8-4}
\expandafter\def\csname sga@target@subsection:IV.8.5\endcsname{ega4.c04.subsec.8-5}
\expandafter\def\csname sga@target@subsection:IV.8.6\endcsname{ega4.c04.subsec.8-6}
\expandafter\def\csname sga@target@subsection:IV.8.7\endcsname{ega4.c04.subsec.8-7}
\expandafter\def\csname sga@target@subsection:IV.8.8\endcsname{ega4.c04.subsec.8-8}
\expandafter\def\csname sga@target@subsection:IV.8.9\endcsname{ega4.c04.subsec.8-9}
\expandafter\def\csname sga@target@subsection:IV.9.1\endcsname{ega4.c04.subsec.9-1}
\expandafter\def\csname sga@target@subsection:IV.9.2\endcsname{ega4.c04.subsec.9-2}
\expandafter\def\csname sga@target@subsection:IV.9.3\endcsname{ega4.c04.subsec.9-3}
\expandafter\def\csname sga@target@subsection:IV.9.4\endcsname{ega4.c04.subsec.9-4}
\expandafter\def\csname sga@target@subsection:IV.9.5\endcsname{ega4.c04.subsec.9-5}
\expandafter\def\csname sga@target@subsection:IV.9.6\endcsname{ega4.c04.subsec.9-6}
\expandafter\def\csname sga@target@subsection:IV.9.7\endcsname{ega4.c04.subsec.9-7}
\expandafter\def\csname sga@target@subsection:IV.9.8\endcsname{ega4.c04.subsec.9-8}
\expandafter\def\csname sga@target@subsection:IV.9.9\endcsname{ega4.c04.subsec.9-9}
\expandafter\def\csname sga@bib@I-4\endcsname{ega4.reader.bib.i-4}
\expandafter\def\csname sga@bib@IV-32\endcsname{ega4.reader.bib.iv-32}
\expandafter\def\csname sga@bib@IV-34\endcsname{ega4.reader.bib.iv-34}
\expandafter\def\csname sga@bib@IV-37\endcsname{ega4.reader.bib.iv-37}
\expandafter\def\csname sga@bib@IV-41\endcsname{ega4.reader.bib.iv-41}
\expandafter\def\csname sga@bib@IV-42\endcsname{ega4.reader.bib.iv-42}
\expandafter\def\csname sga@bib@IV-43\endcsname{ega4.reader.bib.iv-43}
\expandafter\def\csname sga@bib@IV-44\endcsname{ega4.reader.bib.iv-44}
\expandafter\def\csname sga@bib@SGA1\endcsname{ega4.reader.bib.sga1}
\let\sgaOriginalLabel\label
\renewcommand{\label}[1]{%
  \sgaOriginalLabel{#1}%
  \ifcsname sga@target@#1\endcsname
    \expandafter\sgaStableTarget\expandafter{\csname sga@target@#1\endcsname}%
  \fi
}
\let\sgaOriginalBibitem\bibitem
\newcommand{\sgaBibTarget}[1]{%
  \ifcsname sga@bib@#1\endcsname
    \expandafter\sgaStableTarget\expandafter{\csname sga@bib@#1\endcsname}%
  \fi
}
\newcommand{\sgaBibitemPlain}[1]{%
  \sgaOriginalBibitem{#1}\sgaBibTarget{#1}}
\newcommand{\sgaBibitemOptional}[2][]{%
  \sgaOriginalBibitem[#1]{#2}\sgaBibTarget{#2}}
\renewcommand{\bibitem}{%
  \@ifnextchar[\sgaBibitemOptional\sgaBibitemPlain}
\makeatother
% END NATIVE INPUT EGA4_reference_v2_registry
% BEGIN NATIVE INPUT EGA4_reference_v2_links
% Generated stable-link dispatch for the EGA IV reader.
% Stable destinations are emitted by EGA4_reference_v2_registry.tex.
\makeatletter
\expandafter\def\csname sga@target@IV.1.1.1\endcsname{ega4.c04.def.1-1-1}
\expandafter\def\csname sga@target@IV.1.1.2\endcsname{ega4.c04.prop.1-1-2}
\expandafter\def\csname sga@target@IV.1.1.3\endcsname{ega4.c04.prop.1-1-3}
\expandafter\def\csname sga@target@IV.1.1.4\endcsname{ega4.c04.cor.1-1-4}
\expandafter\def\csname sga@target@IV.1.1.5\endcsname{ega4.c04.prop.1-1-5}
\expandafter\def\csname sga@target@IV.1.1.6\endcsname{ega4.c04.prop.1-1-6}
\expandafter\def\csname sga@target@IV.1.10.1\endcsname{ega4.c04.thm.1-10-1}
\expandafter\def\csname sga@target@IV.1.10.2\endcsname{ega4.c04.cor.1-10-2}
\expandafter\def\csname sga@target@IV.1.10.3\endcsname{ega4.c04.prop.1-10-3}
\expandafter\def\csname sga@target@IV.1.10.4\endcsname{ega4.c04.cor.1-10-4}
\expandafter\def\csname sga@target@IV.1.2.1\endcsname{ega4.c04.def.1-2-1}
\expandafter\def\csname sga@target@IV.1.2.2\endcsname{ega4.c04.prop.1-2-2}
\expandafter\def\csname sga@target@IV.1.2.3\endcsname{ega4.c04.cor.1-2-3}
\expandafter\def\csname sga@target@IV.1.2.4\endcsname{ega4.c04.prop.1-2-4}
\expandafter\def\csname sga@target@IV.1.2.4.1\endcsname{ega4.c04.eq.1-2-4-1}
\expandafter\def\csname sga@target@IV.1.2.5\endcsname{ega4.c04.prop.1-2-5}
\expandafter\def\csname sga@target@IV.1.2.6\endcsname{ega4.c04.prop.1-2-6}
\expandafter\def\csname sga@target@IV.1.2.7\endcsname{ega4.c04.prop.1-2-7}
\expandafter\def\csname sga@target@IV.1.2.8\endcsname{ega4.c04.cor.1-2-8}
\expandafter\def\csname sga@target@IV.1.2.9\endcsname{ega4.c04.prop.1-2-9}
\expandafter\def\csname sga@target@IV.1.3.1\endcsname{ega4.c04.statement.1-3-1}
\expandafter\def\csname sga@target@IV.1.3.10\endcsname{ega4.c04.statement.1-3-10}
\expandafter\def\csname sga@target@IV.1.3.11\endcsname{ega4.c04.prop.1-3-11}
\expandafter\def\csname sga@target@IV.1.3.2\endcsname{ega4.c04.def.1-3-2}
\expandafter\def\csname sga@target@IV.1.3.3\endcsname{ega4.c04.prop.1-3-3}
\expandafter\def\csname sga@target@IV.1.3.4\endcsname{ega4.c04.prop.1-3-4}
\expandafter\def\csname sga@target@IV.1.3.5\endcsname{ega4.c04.cor.1-3-5}
\expandafter\def\csname sga@target@IV.1.3.6\endcsname{ega4.c04.prop.1-3-6}
\expandafter\def\csname sga@target@IV.1.3.7\endcsname{ega4.c04.prop.1-3-7}
\expandafter\def\csname sga@target@IV.1.3.8\endcsname{ega4.c04.statement.1-3-8}
\expandafter\def\csname sga@target@IV.1.3.9\endcsname{ega4.c04.prop.1-3-9}
\expandafter\def\csname sga@target@IV.1.4.1\endcsname{ega4.c04.statement.1-4-1}
\expandafter\def\csname sga@target@IV.1.4.2\endcsname{ega4.c04.def.1-4-2}
\expandafter\def\csname sga@target@IV.1.4.3\endcsname{ega4.c04.prop.1-4-3}
\expandafter\def\csname sga@target@IV.1.4.3.1\endcsname{ega4.c04.cor.1-4-3-1}
\expandafter\def\csname sga@target@IV.1.4.4\endcsname{ega4.c04.prop.1-4-4}
\expandafter\def\csname sga@target@IV.1.4.4.1\endcsname{ega4.c04.lem.1-4-4-1}
\expandafter\def\csname sga@target@IV.1.4.5\endcsname{ega4.c04.cor.1-4-5}
\expandafter\def\csname sga@target@IV.1.4.6\endcsname{ega4.c04.prop.1-4-6}
\expandafter\def\csname sga@target@IV.1.4.7\endcsname{ega4.c04.prop.1-4-7}
\expandafter\def\csname sga@target@IV.1.4.7.1\endcsname{ega4.c04.lem.1-4-7-1}
\expandafter\def\csname sga@target@IV.1.5.1\endcsname{ega4.c04.prop.1-5-1}
\expandafter\def\csname sga@target@IV.1.5.2\endcsname{ega4.c04.def.1-5-2}
\expandafter\def\csname sga@target@IV.1.5.3\endcsname{ega4.c04.prop.1-5-3}
\expandafter\def\csname sga@target@IV.1.5.4\endcsname{ega4.c04.prop.1-5-4}
\expandafter\def\csname sga@target@IV.1.5.5\endcsname{ega4.c04.cor.1-5-5}
\expandafter\def\csname sga@target@IV.1.5.6\endcsname{ega4.c04.cor.1-5-6}
\expandafter\def\csname sga@target@IV.1.5.7\endcsname{ega4.c04.prop.1-5-7}
\expandafter\def\csname sga@target@IV.1.6.1\endcsname{ega4.c04.def.1-6-1}
\expandafter\def\csname sga@target@IV.1.6.2\endcsname{ega4.c04.prop.1-6-2}
\expandafter\def\csname sga@target@IV.1.6.3\endcsname{ega4.c04.cor.1-6-3}
\expandafter\def\csname sga@target@IV.1.6.4\endcsname{ega4.c04.rem.1-6-4}
\expandafter\def\csname sga@target@IV.1.6.5\endcsname{ega4.c04.prop.1-6-5}
\expandafter\def\csname sga@target@IV.1.7.1\endcsname{ega4.c04.statement.1-7-1}
\expandafter\def\csname sga@target@IV.1.7.10\endcsname{ega4.c04.statement.1-7-10}
\expandafter\def\csname sga@target@IV.1.7.11\endcsname{ega4.c04.statement.1-7-11}
\expandafter\def\csname sga@target@IV.1.7.12\endcsname{ega4.c04.statement.1-7-12}
\expandafter\def\csname sga@target@IV.1.7.13\endcsname{ega4.c04.statement.1-7-13}
\expandafter\def\csname sga@target@IV.1.7.14\endcsname{ega4.c04.statement.1-7-14}
\expandafter\def\csname sga@target@IV.1.7.15\endcsname{ega4.c04.statement.1-7-15}
\expandafter\def\csname sga@target@IV.1.7.16\endcsname{ega4.c04.statement.1-7-16}
\expandafter\def\csname sga@target@IV.1.7.17\endcsname{ega4.c04.statement.1-7-17}
\expandafter\def\csname sga@target@IV.1.7.18\endcsname{ega4.c04.statement.1-7-18}
\expandafter\def\csname sga@target@IV.1.7.19\endcsname{ega4.c04.statement.1-7-19}
\expandafter\def\csname sga@target@IV.1.7.2\endcsname{ega4.c04.statement.1-7-2}
\expandafter\def\csname sga@target@IV.1.7.20\endcsname{ega4.c04.statement.1-7-20}
\expandafter\def\csname sga@target@IV.1.7.21\endcsname{ega4.c04.statement.1-7-21}
\expandafter\def\csname sga@target@IV.1.7.22\endcsname{ega4.c04.statement.1-7-22}
\expandafter\def\csname sga@target@IV.1.7.3\endcsname{ega4.c04.statement.1-7-3}
\expandafter\def\csname sga@target@IV.1.7.4\endcsname{ega4.c04.statement.1-7-4}
\expandafter\def\csname sga@target@IV.1.7.5\endcsname{ega4.c04.statement.1-7-5}
\expandafter\def\csname sga@target@IV.1.7.6\endcsname{ega4.c04.statement.1-7-6}
\expandafter\def\csname sga@target@IV.1.7.7\endcsname{ega4.c04.statement.1-7-7}
\expandafter\def\csname sga@target@IV.1.7.8\endcsname{ega4.c04.statement.1-7-8}
\expandafter\def\csname sga@target@IV.1.7.9\endcsname{ega4.c04.statement.1-7-9}
\expandafter\def\csname sga@target@IV.1.8.1\endcsname{ega4.c04.statement.1-8-1}
\expandafter\def\csname sga@target@IV.1.8.2\endcsname{ega4.c04.prop.1-8-2}
\expandafter\def\csname sga@target@IV.1.8.3\endcsname{ega4.c04.lem.1-8-3}
\expandafter\def\csname sga@target@IV.1.8.4\endcsname{ega4.c04.thm.1-8-4}
\expandafter\def\csname sga@target@IV.1.8.4.1\endcsname{ega4.c04.lem.1-8-4-1}
\expandafter\def\csname sga@target@IV.1.8.4.2\endcsname{ega4.c04.lem.1-8-4-2}
\expandafter\def\csname sga@target@IV.1.8.5\endcsname{ega4.c04.cor.1-8-5}
\expandafter\def\csname sga@target@IV.1.8.5.1\endcsname{ega4.c04.lem.1-8-5-1}
\expandafter\def\csname sga@target@IV.1.8.6\endcsname{ega4.c04.cor.1-8-6}
\expandafter\def\csname sga@target@IV.1.8.7\endcsname{ega4.c04.prop.1-8-7}
\expandafter\def\csname sga@target@IV.1.8.7.1\endcsname{ega4.c04.lem.1-8-7-1}
\expandafter\def\csname sga@target@IV.1.9.1\endcsname{ega4.c04.lem.1-9-1}
\expandafter\def\csname sga@target@IV.1.9.10\endcsname{ega4.c04.prop.1-9-10}
\expandafter\def\csname sga@target@IV.1.9.11\endcsname{ega4.c04.prop.1-9-11}
\expandafter\def\csname sga@target@IV.1.9.12\endcsname{ega4.c04.cor.1-9-12}
\expandafter\def\csname sga@target@IV.1.9.13\endcsname{ega4.c04.statement.1-9-13}
\expandafter\def\csname sga@target@IV.1.9.14\endcsname{ega4.c04.prop.1-9-14}
\expandafter\def\csname sga@target@IV.1.9.15\endcsname{ega4.c04.prop.1-9-15}
\expandafter\def\csname sga@target@IV.1.9.16\endcsname{ega4.c04.statement.1-9-16}
\expandafter\def\csname sga@target@IV.1.9.2\endcsname{ega4.c04.prop.1-9-2}
\expandafter\def\csname sga@target@IV.1.9.2.1\endcsname{ega4.c04.statement.1-9-2-1}
\expandafter\def\csname sga@target@IV.1.9.2.i\endcsname{ega4.c04.statement.1-9-2-i}
\expandafter\def\csname sga@target@IV.1.9.2.ii\endcsname{ega4.c04.statement.1-9-2-ii}
\expandafter\def\csname sga@target@IV.1.9.3\endcsname{ega4.c04.prop.1-9-3}
\expandafter\def\csname sga@target@IV.1.9.3.1\endcsname{ega4.c04.lem.1-9-3-1}
\expandafter\def\csname sga@target@IV.1.9.3.2\endcsname{ega4.c04.lem.1-9-3-2}
\expandafter\def\csname sga@target@IV.1.9.4\endcsname{ega4.c04.def.1-9-4}
\expandafter\def\csname sga@target@IV.1.9.5\endcsname{ega4.c04.prop.1-9-5}
\expandafter\def\csname sga@target@IV.1.9.6\endcsname{ega4.c04.statement.1-9-6}
\expandafter\def\csname sga@target@IV.1.9.7\endcsname{ega4.c04.thm.1-9-7}
\expandafter\def\csname sga@target@IV.1.9.8\endcsname{ega4.c04.cor.1-9-8}
\expandafter\def\csname sga@target@IV.1.9.9\endcsname{ega4.c04.cor.1-9-9}
\expandafter\def\csname sga@target@IV.10.1.1\endcsname{ega4.c04.statement.10-1-1}
\expandafter\def\csname sga@target@IV.10.1.2\endcsname{ega4.c04.prop.10-1-2}
\expandafter\def\csname sga@target@IV.10.1.3\endcsname{ega4.c04.def.10-1-3}
\expandafter\def\csname sga@target@IV.10.1.4\endcsname{ega4.c04.cor.10-1-4}
\expandafter\def\csname sga@target@IV.10.10.1\endcsname{ega4.c04.statement.10-10-1}
\expandafter\def\csname sga@target@IV.10.10.2\endcsname{ega4.c04.statement.10-10-2}
\expandafter\def\csname sga@target@IV.10.10.3\endcsname{ega4.c04.statement.10-10-3}
\expandafter\def\csname sga@target@IV.10.10.4\endcsname{ega4.c04.statement.10-10-4}
\expandafter\def\csname sga@target@IV.10.10.5\endcsname{ega4.c04.statement.10-10-5}
\expandafter\def\csname sga@target@IV.10.2.1\endcsname{ega4.c04.prop.10-2-1}
\expandafter\def\csname sga@target@IV.10.2.2\endcsname{ega4.c04.def.10-2-2}
\expandafter\def\csname sga@target@IV.10.2.3\endcsname{ega4.c04.cor.10-2-3}
\expandafter\def\csname sga@target@IV.10.2.4\endcsname{ega4.c04.cor.10-2-4}
\expandafter\def\csname sga@target@IV.10.2.5\endcsname{ega4.c04.cor.10-2-5}
\expandafter\def\csname sga@target@IV.10.2.6\endcsname{ega4.c04.cor.10-2-6}
\expandafter\def\csname sga@target@IV.10.2.7\endcsname{ega4.c04.prop.10-2-7}
\expandafter\def\csname sga@target@IV.10.2.8\endcsname{ega4.c04.item.10-2-8}
\expandafter\def\csname sga@target@IV.10.2.8.1\endcsname{ega4.c04.eq.10-2-8-1}
\expandafter\def\csname sga@target@IV.10.3.1\endcsname{ega4.c04.def.10-3-1}
\expandafter\def\csname sga@target@IV.10.3.2\endcsname{ega4.c04.prop.10-3-2}
\expandafter\def\csname sga@target@IV.10.3.3\endcsname{ega4.c04.prop.10-3-3}
\expandafter\def\csname sga@target@IV.10.4.1\endcsname{ega4.c04.def.10-4-1}
\expandafter\def\csname sga@target@IV.10.4.10\endcsname{ega4.c04.statement.10-4-10}
\expandafter\def\csname sga@target@IV.10.4.11\endcsname{ega4.c04.statement.10-4-11}
\expandafter\def\csname sga@target@IV.10.4.11.1\endcsname{ega4.c04.lem.10-4-11-1}
\expandafter\def\csname sga@target@IV.10.4.12\endcsname{ega4.c04.cor.10-4-12}
\expandafter\def\csname sga@target@IV.10.4.12.1\endcsname{ega4.c04.sec.10-4-12-1}
\expandafter\def\csname sga@target@IV.10.4.12.1.errata\endcsname{ega4.c04.eq.10-4-12-1-errata}
\expandafter\def\csname sga@target@IV.10.4.13\endcsname{ega4.c04.cor.10-4-13}
\expandafter\def\csname sga@target@IV.10.4.14\endcsname{ega4.c04.rem.10-4-14}
\expandafter\def\csname sga@target@IV.10.4.2\endcsname{ega4.c04.prop.10-4-2}
\expandafter\def\csname sga@target@IV.10.4.3\endcsname{ega4.c04.statement.10-4-3}
\expandafter\def\csname sga@target@IV.10.4.4\endcsname{ega4.c04.prop.10-4-4}
\expandafter\def\csname sga@target@IV.10.4.5\endcsname{ega4.c04.prop.10-4-5}
\expandafter\def\csname sga@target@IV.10.4.5.1\endcsname{ega4.c04.lem.10-4-5-1}
\expandafter\def\csname sga@target@IV.10.4.6\endcsname{ega4.c04.cor.10-4-6}
\expandafter\def\csname sga@target@IV.10.4.7\endcsname{ega4.c04.cor.10-4-7}
\expandafter\def\csname sga@target@IV.10.4.8\endcsname{ega4.c04.cor.10-4-8}
\expandafter\def\csname sga@target@IV.10.4.9\endcsname{ega4.c04.statement.10-4-9}
\expandafter\def\csname sga@target@IV.10.5.1\endcsname{ega4.c04.prop.10-5-1}
\expandafter\def\csname sga@target@IV.10.5.10\endcsname{ega4.c04.prop.10-5-10}
\expandafter\def\csname sga@target@IV.10.5.2\endcsname{ega4.c04.cor.10-5-2}
\expandafter\def\csname sga@target@IV.10.5.3\endcsname{ega4.c04.cor.10-5-3}
\expandafter\def\csname sga@target@IV.10.5.4\endcsname{ega4.c04.rem.10-5-4}
\expandafter\def\csname sga@target@IV.10.5.5\endcsname{ega4.c04.prop.10-5-5}
\expandafter\def\csname sga@target@IV.10.5.6\endcsname{ega4.c04.rem.10-5-6}
\expandafter\def\csname sga@target@IV.10.5.7\endcsname{ega4.c04.prop.10-5-7}
\expandafter\def\csname sga@target@IV.10.5.8\endcsname{ega4.c04.cor.10-5-8}
\expandafter\def\csname sga@target@IV.10.5.9\endcsname{ega4.c04.cor.10-5-9}
\expandafter\def\csname sga@target@IV.10.6.1\endcsname{ega4.c04.prop.10-6-1}
\expandafter\def\csname sga@target@IV.10.6.1.1\endcsname{ega4.c04.eq.10-6-1-1}
\expandafter\def\csname sga@target@IV.10.6.1.2\endcsname{ega4.c04.eq.10-6-1-2}
\expandafter\def\csname sga@target@IV.10.6.1.3\endcsname{ega4.c04.eq.10-6-1-3}
\expandafter\def\csname sga@target@IV.10.6.1.4\endcsname{ega4.c04.eq.10-6-1-4}
\expandafter\def\csname sga@target@IV.10.6.2\endcsname{ega4.c04.cor.10-6-2}
\expandafter\def\csname sga@target@IV.10.6.3\endcsname{ega4.c04.prop.10-6-3}
\expandafter\def\csname sga@target@IV.10.6.3.1\endcsname{ega4.c04.eq.10-6-3-1}
\expandafter\def\csname sga@target@IV.10.6.3.2\endcsname{ega4.c04.eq.10-6-3-2}
\expandafter\def\csname sga@target@IV.10.6.3.3\endcsname{ega4.c04.eq.10-6-3-3}
\expandafter\def\csname sga@target@IV.10.6.4\endcsname{ega4.c04.cor.10-6-4}
\expandafter\def\csname sga@target@IV.10.6.4.1\endcsname{ega4.c04.eq.10-6-4-1}
\expandafter\def\csname sga@target@IV.10.7.1\endcsname{ega4.c04.statement.10-7-1}
\expandafter\def\csname sga@target@IV.10.7.2\endcsname{ega4.c04.statement.10-7-2}
\expandafter\def\csname sga@target@IV.10.7.3\endcsname{ega4.c04.statement.10-7-3}
\expandafter\def\csname sga@target@IV.10.8.1\endcsname{ega4.c04.def.10-8-1}
\expandafter\def\csname sga@target@IV.10.8.1.1\endcsname{ega4.c04.eq.10-8-1-1}
\expandafter\def\csname sga@target@IV.10.8.1.2\endcsname{ega4.c04.eq.10-8-1-2}
\expandafter\def\csname sga@target@IV.10.8.2\endcsname{ega4.c04.statement.10-8-2}
\expandafter\def\csname sga@target@IV.10.8.3\endcsname{ega4.c04.lem.10-8-3}
\expandafter\def\csname sga@target@IV.10.8.3.1\endcsname{ega4.c04.eq.10-8-3-1}
\expandafter\def\csname sga@target@IV.10.8.4\endcsname{ega4.c04.cor.10-8-4}
\expandafter\def\csname sga@target@IV.10.8.5\endcsname{ega4.c04.prop.10-8-5}
\expandafter\def\csname sga@target@IV.10.8.6\endcsname{ega4.c04.cor.10-8-6}
\expandafter\def\csname sga@target@IV.10.8.7\endcsname{ega4.c04.prop.10-8-7}
\expandafter\def\csname sga@target@IV.10.8.7.1\endcsname{ega4.c04.eq.10-8-7-1}
\expandafter\def\csname sga@target@IV.10.8.8\endcsname{ega4.c04.cor.10-8-8}
\expandafter\def\csname sga@target@IV.10.9.1\endcsname{ega4.c04.statement.10-9-1}
\expandafter\def\csname sga@target@IV.10.9.2\endcsname{ega4.c04.statement.10-9-2}
\expandafter\def\csname sga@target@IV.10.9.3\endcsname{ega4.c04.statement.10-9-3}
\expandafter\def\csname sga@target@IV.10.9.4\endcsname{ega4.c04.prop.10-9-4}
\expandafter\def\csname sga@target@IV.10.9.5\endcsname{ega4.c04.statement.10-9-5}
\expandafter\def\csname sga@target@IV.10.9.6\endcsname{ega4.c04.prop.10-9-6}
\expandafter\def\csname sga@target@IV.11.1.1\endcsname{ega4.c04.thm.11-1-1}
\expandafter\def\csname sga@target@IV.11.1.1.1\endcsname{ega4.c04.sec.11-1-1-1}
\expandafter\def\csname sga@target@IV.11.1.1.2\endcsname{ega4.c04.lem.11-1-1-2}
\expandafter\def\csname sga@target@IV.11.1.2\endcsname{ega4.c04.cor.11-1-2}
\expandafter\def\csname sga@target@IV.11.1.3\endcsname{ega4.c04.cor.11-1-3}
\expandafter\def\csname sga@target@IV.11.1.4\endcsname{ega4.c04.cor.11-1-4}
\expandafter\def\csname sga@target@IV.11.1.5\endcsname{ega4.c04.cor.11-1-5}
\expandafter\def\csname sga@target@IV.11.1.6\endcsname{ega4.c04.cor.11-1-6}
\expandafter\def\csname sga@target@IV.11.10.1\endcsname{ega4.c04.prop.11-10-1}
\expandafter\def\csname sga@target@IV.11.10.1.1\endcsname{ega4.c04.eq.11-10-1-1}
\expandafter\def\csname sga@target@IV.11.10.1.2\endcsname{ega4.c04.eq.11-10-1-2}
\expandafter\def\csname sga@target@IV.11.10.10\endcsname{ega4.c04.prop.11-10-10}
\expandafter\def\csname sga@target@IV.11.10.2\endcsname{ega4.c04.def.11-10-2}
\expandafter\def\csname sga@target@IV.11.10.3\endcsname{ega4.c04.item.11-10-3}
\expandafter\def\csname sga@target@IV.11.10.4\endcsname{ega4.c04.prop.11-10-4}
\expandafter\def\csname sga@target@IV.11.10.5\endcsname{ega4.c04.thm.11-10-5}
\expandafter\def\csname sga@target@IV.11.10.6\endcsname{ega4.c04.prop.11-10-6}
\expandafter\def\csname sga@target@IV.11.10.7\endcsname{ega4.c04.cor.11-10-7}
\expandafter\def\csname sga@target@IV.11.10.8\endcsname{ega4.c04.def.11-10-8}
\expandafter\def\csname sga@target@IV.11.10.9\endcsname{ega4.c04.thm.11-10-9}
\expandafter\def\csname sga@target@IV.11.2.1\endcsname{ega4.c04.statement.11-2-1}
\expandafter\def\csname sga@target@IV.11.2.1.1\endcsname{ega4.c04.eq.11-2-1-1}
\expandafter\def\csname sga@target@IV.11.2.1.2\endcsname{ega4.c04.eq.11-2-1-2}
\expandafter\def\csname sga@target@IV.11.2.1.3\endcsname{ega4.c04.eq.11-2-1-3}
\expandafter\def\csname sga@target@IV.11.2.1.4\endcsname{ega4.c04.eq.11-2-1-4}
\expandafter\def\csname sga@target@IV.11.2.1.5\endcsname{ega4.c04.eq.11-2-1-5}
\expandafter\def\csname sga@target@IV.11.2.10\endcsname{ega4.c04.rem.11-2-10}
\expandafter\def\csname sga@target@IV.11.2.2\endcsname{ega4.c04.statement.11-2-2}
\expandafter\def\csname sga@target@IV.11.2.2.1\endcsname{ega4.c04.eq.11-2-2-1}
\expandafter\def\csname sga@target@IV.11.2.2.2\endcsname{ega4.c04.lem.11-2-2-2}
\expandafter\def\csname sga@target@IV.11.2.3\endcsname{ega4.c04.statement.11-2-3}
\expandafter\def\csname sga@target@IV.11.2.3.1\endcsname{ega4.c04.eq.11-2-3-1}
\expandafter\def\csname sga@target@IV.11.2.3.2\endcsname{ega4.c04.eq.11-2-3-2}
\expandafter\def\csname sga@target@IV.11.2.3.3\endcsname{ega4.c04.eq.11-2-3-3}
\expandafter\def\csname sga@target@IV.11.2.4\endcsname{ega4.c04.lem.11-2-4}
\expandafter\def\csname sga@target@IV.11.2.5\endcsname{ega4.c04.lem.11-2-5}
\expandafter\def\csname sga@target@IV.11.2.6\endcsname{ega4.c04.thm.11-2-6}
\expandafter\def\csname sga@target@IV.11.2.6.1\endcsname{ega4.c04.cor.11-2-6-1}
\expandafter\def\csname sga@target@IV.11.2.7\endcsname{ega4.c04.cor.11-2-7}
\expandafter\def\csname sga@target@IV.11.2.8\endcsname{ega4.c04.prop.11-2-8}
\expandafter\def\csname sga@target@IV.11.2.9\endcsname{ega4.c04.thm.11-2-9}
\expandafter\def\csname sga@target@IV.11.2.9.1\endcsname{ega4.c04.statement.11-2-9-1}
\expandafter\def\csname sga@target@IV.11.2.9.2\endcsname{ega4.c04.lem.11-2-9-2}
\expandafter\def\csname sga@target@IV.11.2.9.3\endcsname{ega4.c04.lem.11-2-9-3}
\expandafter\def\csname sga@target@IV.11.2.9.4\endcsname{ega4.c04.cor.11-2-9-4}
\expandafter\def\csname sga@target@IV.11.2.9.5\endcsname{ega4.c04.lem.11-2-9-5}
\expandafter\def\csname sga@target@IV.11.2.9.6\endcsname{ega4.c04.lem.11-2-9-6}
\expandafter\def\csname sga@target@IV.11.2.9.7\endcsname{ega4.c04.eq.11-2-9-7}
\expandafter\def\csname sga@target@IV.11.2.9.8\endcsname{ega4.c04.eq.11-2-9-8}
\expandafter\def\csname sga@target@IV.11.2.9.9\endcsname{ega4.c04.eq.11-2-9-9}
\expandafter\def\csname sga@target@IV.11.3\endcsname{ega4.c04.subsec.11-3.key-iv-11-3}
\expandafter\def\csname sga@target@IV.11.3.1\endcsname{ega4.c04.thm.11-3-1}
\expandafter\def\csname sga@target@IV.11.3.10\endcsname{ega4.c04.thm.11-3-10}
\expandafter\def\csname sga@target@IV.11.3.10.1\endcsname{ega4.c04.lem.11-3-10-1}
\expandafter\def\csname sga@target@IV.11.3.10.2\endcsname{ega4.c04.lem.11-3-10-2}
\expandafter\def\csname sga@target@IV.11.3.10.3\endcsname{ega4.c04.eq.11-3-10-3}
\expandafter\def\csname sga@target@IV.11.3.10.4\endcsname{ega4.c04.sec.11-3-10-4}
\expandafter\def\csname sga@target@IV.11.3.11\endcsname{ega4.c04.cor.11-3-11}
\expandafter\def\csname sga@target@IV.11.3.12\endcsname{ega4.c04.rem.11-3-12}
\expandafter\def\csname sga@target@IV.11.3.13\endcsname{ega4.c04.prop.11-3-13}
\expandafter\def\csname sga@target@IV.11.3.14\endcsname{ega4.c04.cor.11-3-14}
\expandafter\def\csname sga@target@IV.11.3.15\endcsname{ega4.c04.prop.11-3-15}
\expandafter\def\csname sga@target@IV.11.3.16\endcsname{ega4.c04.prop.11-3-16}
\expandafter\def\csname sga@target@IV.11.3.17\endcsname{ega4.c04.cor.11-3-17}
\expandafter\def\csname sga@target@IV.11.3.2\endcsname{ega4.c04.cor.11-3-2}
\expandafter\def\csname sga@target@IV.11.3.3\endcsname{ega4.c04.cor.11-3-3}
\expandafter\def\csname sga@target@IV.11.3.4\endcsname{ega4.c04.prop.11-3-4}
\expandafter\def\csname sga@target@IV.11.3.5\endcsname{ega4.c04.cor.11-3-5}
\expandafter\def\csname sga@target@IV.11.3.6\endcsname{ega4.c04.prop.11-3-6}
\expandafter\def\csname sga@target@IV.11.3.7\endcsname{ega4.c04.prop.11-3-7}
\expandafter\def\csname sga@target@IV.11.3.8\endcsname{ega4.c04.thm.11-3-8}
\expandafter\def\csname sga@target@IV.11.3.9\endcsname{ega4.c04.prop.11-3-9}
\expandafter\def\csname sga@target@IV.11.3.9.1\endcsname{ega4.c04.lem.11-3-9-1}
\expandafter\def\csname sga@target@IV.11.4.1\endcsname{ega4.c04.thm.11-4-1}
\expandafter\def\csname sga@target@IV.11.4.1.1\endcsname{ega4.c04.lem.11-4-1-1}
\expandafter\def\csname sga@target@IV.11.4.1.2\endcsname{ega4.c04.lem.11-4-1-2}
\expandafter\def\csname sga@target@IV.11.4.10\endcsname{ega4.c04.cor.11-4-10}
\expandafter\def\csname sga@target@IV.11.4.11\endcsname{ega4.c04.prop.11-4-11}
\expandafter\def\csname sga@target@IV.11.4.12\endcsname{ega4.c04.prop.11-4-12}
\expandafter\def\csname sga@target@IV.11.4.13\endcsname{ega4.c04.cor.11-4-13}
\expandafter\def\csname sga@target@IV.11.4.2\endcsname{ega4.c04.rem.11-4-2}
\expandafter\def\csname sga@target@IV.11.4.3\endcsname{ega4.c04.cor.11-4-3}
\expandafter\def\csname sga@target@IV.11.4.4\endcsname{ega4.c04.cor.11-4-4}
\expandafter\def\csname sga@target@IV.11.4.5\endcsname{ega4.c04.prop.11-4-5}
\expandafter\def\csname sga@target@IV.11.4.6\endcsname{ega4.c04.cor.11-4-6}
\expandafter\def\csname sga@target@IV.11.4.7\endcsname{ega4.c04.prop.11-4-7}
\expandafter\def\csname sga@target@IV.11.4.8\endcsname{ega4.c04.rem.11-4-8}
\expandafter\def\csname sga@target@IV.11.4.9\endcsname{ega4.c04.cor.11-4-9}
\expandafter\def\csname sga@target@IV.11.5.1\endcsname{ega4.c04.thm.11-5-1}
\expandafter\def\csname sga@target@IV.11.5.2\endcsname{ega4.c04.cor.11-5-2}
\expandafter\def\csname sga@target@IV.11.5.3\endcsname{ega4.c04.prop.11-5-3}
\expandafter\def\csname sga@target@IV.11.5.4\endcsname{ega4.c04.cor.11-5-4}
\expandafter\def\csname sga@target@IV.11.5.5\endcsname{ega4.c04.prop.11-5-5}
\expandafter\def\csname sga@target@IV.11.5.5.1\endcsname{ega4.c04.lem.11-5-5-1}
\expandafter\def\csname sga@target@IV.11.5.5.2\endcsname{ega4.c04.lem.11-5-5-2}
\expandafter\def\csname sga@target@IV.11.6.1\endcsname{ega4.c04.thm.11-6-1}
\expandafter\def\csname sga@target@IV.11.6.2\endcsname{ega4.c04.cor.11-6-2}
\expandafter\def\csname sga@target@IV.11.7.1\endcsname{ega4.c04.statement.11-7-1}
\expandafter\def\csname sga@target@IV.11.7.2\endcsname{ega4.c04.statement.11-7-2}
\expandafter\def\csname sga@target@IV.11.7.3\endcsname{ega4.c04.statement.11-7-3}
\expandafter\def\csname sga@target@IV.11.7.4\endcsname{ega4.c04.statement.11-7-4}
\expandafter\def\csname sga@target@IV.11.7.5\endcsname{ega4.c04.statement.11-7-5}
\expandafter\def\csname sga@target@IV.11.8.1\endcsname{ega4.c04.thm.11-8-1}
\expandafter\def\csname sga@target@IV.11.8.2\endcsname{ega4.c04.subsec.11-8-2}
\expandafter\def\csname sga@target@IV.11.9.1\endcsname{ega4.c04.statement.11-9-1}
\expandafter\def\csname sga@target@IV.11.9.1.1\endcsname{ega4.c04.statement.11-9-1-1}
\expandafter\def\csname sga@target@IV.11.9.10\endcsname{ega4.c04.thm.11-9-10}
\expandafter\def\csname sga@target@IV.11.9.11\endcsname{ega4.c04.rem.11-9-11}
\expandafter\def\csname sga@target@IV.11.9.12\endcsname{ega4.c04.lem.11-9-12}
\expandafter\def\csname sga@target@IV.11.9.13\endcsname{ega4.c04.prop.11-9-13}
\expandafter\def\csname sga@target@IV.11.9.14\endcsname{ega4.c04.statement.11-9-14}
\expandafter\def\csname sga@target@IV.11.9.15\endcsname{ega4.c04.prop.11-9-15}
\expandafter\def\csname sga@target@IV.11.9.16\endcsname{ega4.c04.thm.11-9-16}
\expandafter\def\csname sga@target@IV.11.9.17\endcsname{ega4.c04.prop.11-9-17}
\expandafter\def\csname sga@target@IV.11.9.17.1\endcsname{ega4.c04.lem.11-9-17-1}
\expandafter\def\csname sga@target@IV.11.9.18\endcsname{ega4.c04.item.11-9-18}
\expandafter\def\csname sga@target@IV.11.9.2\endcsname{ega4.c04.statement.11-9-2}
\expandafter\def\csname sga@target@IV.11.9.3\endcsname{ega4.c04.statement.11-9-3}
\expandafter\def\csname sga@target@IV.11.9.4\endcsname{ega4.c04.statement.11-9-4}
\expandafter\def\csname sga@target@IV.11.9.5\endcsname{ega4.c04.statement.11-9-5}
\expandafter\def\csname sga@target@IV.11.9.6\endcsname{ega4.c04.statement.11-9-6}
\expandafter\def\csname sga@target@IV.11.9.7\endcsname{ega4.c04.prop.11-9-7}
\expandafter\def\csname sga@target@IV.11.9.7.1\endcsname{ega4.c04.lem.11-9-7-1}
\expandafter\def\csname sga@target@IV.11.9.8\endcsname{ega4.c04.subsec.11-9-8}
\expandafter\def\csname sga@target@IV.11.9.9\endcsname{ega4.c04.cor.11-9-9}
\expandafter\def\csname sga@target@IV.12.0.1\endcsname{ega4.c04.statement.12-0-1}
\expandafter\def\csname sga@target@IV.12.0.2\endcsname{ega4.c04.statement.12-0-2}
\expandafter\def\csname sga@target@IV.12.1.1\endcsname{ega4.c04.thm.12-1-1}
\expandafter\def\csname sga@target@IV.12.1.1.1\endcsname{ega4.c04.statement.12-1-1-1}
\expandafter\def\csname sga@target@IV.12.1.1.2\endcsname{ega4.c04.statement.12-1-1-2}
\expandafter\def\csname sga@target@IV.12.1.1.3\endcsname{ega4.c04.statement.12-1-1-3}
\expandafter\def\csname sga@target@IV.12.1.1.4\endcsname{ega4.c04.statement.12-1-1-4}
\expandafter\def\csname sga@target@IV.12.1.1.5\endcsname{ega4.c04.prop.12-1-1-5}
\expandafter\def\csname sga@target@IV.12.1.2\endcsname{ega4.c04.item.12-1-2}
\expandafter\def\csname sga@target@IV.12.1.3\endcsname{ega4.c04.cor.12-1-3}
\expandafter\def\csname sga@target@IV.12.1.4\endcsname{ega4.c04.rem.12-1-4}
\expandafter\def\csname sga@target@IV.12.1.5\endcsname{ega4.c04.cor.12-1-5}
\expandafter\def\csname sga@target@IV.12.1.6\endcsname{ega4.c04.thm.12-1-6}
\expandafter\def\csname sga@target@IV.12.1.7\endcsname{ega4.c04.cor.12-1-7}
\expandafter\def\csname sga@target@IV.12.1.8\endcsname{ega4.c04.item.12-1-8}
\expandafter\def\csname sga@target@IV.12.2.1\endcsname{ega4.c04.thm.12-2-1}
\expandafter\def\csname sga@target@IV.12.2.1.1\endcsname{ega4.c04.statement.12-2-1-1}
\expandafter\def\csname sga@target@IV.12.2.1.2\endcsname{ega4.c04.statement.12-2-1-2}
\expandafter\def\csname sga@target@IV.12.2.1.3\endcsname{ega4.c04.statement.12-2-1-3}
\expandafter\def\csname sga@target@IV.12.2.1.4\endcsname{ega4.c04.lem.12-2-1-4}
\expandafter\def\csname sga@target@IV.12.2.2\endcsname{ega4.c04.cor.12-2-2}
\expandafter\def\csname sga@target@IV.12.2.3\endcsname{ega4.c04.item.12-2-3}
\expandafter\def\csname sga@target@IV.12.2.4\endcsname{ega4.c04.thm.12-2-4}
\expandafter\def\csname sga@target@IV.12.3.1\endcsname{ega4.c04.lem.12-3-1}
\expandafter\def\csname sga@target@IV.12.3.2\endcsname{ega4.c04.prop.12-3-2}
\expandafter\def\csname sga@target@IV.12.3.2.1\endcsname{ega4.c04.eq.12-3-2-1}
\expandafter\def\csname sga@target@IV.12.3.3\endcsname{ega4.c04.prop.12-3-3}
\expandafter\def\csname sga@target@IV.12.3.3.1\endcsname{ega4.c04.lem.12-3-3-1}
\expandafter\def\csname sga@target@IV.12.3.3.2\endcsname{ega4.c04.eq.12-3-3-2}
\expandafter\def\csname sga@target@IV.12.3.3.3\endcsname{ega4.c04.lem.12-3-3-3}
\expandafter\def\csname sga@target@IV.12.3.3.4\endcsname{ega4.c04.eq.12-3-3-4}
\expandafter\def\csname sga@target@IV.12.3.3.5\endcsname{ega4.c04.lem.12-3-3-5}
\expandafter\def\csname sga@target@IV.12.3.4\endcsname{ega4.c04.cor.12-3-4}
\expandafter\def\csname sga@target@IV.13.1.1\endcsname{ega4.c04.lem.13-1-1}
\expandafter\def\csname sga@target@IV.13.1.2\endcsname{ega4.c04.lem.13-1-2}
\expandafter\def\csname sga@target@IV.13.1.3\endcsname{ega4.c04.thm.13-1-3}
\expandafter\def\csname sga@target@IV.13.1.4\endcsname{ega4.c04.cor.13-1-4}
\expandafter\def\csname sga@target@IV.13.1.5\endcsname{ega4.c04.cor.13-1-5}
\expandafter\def\csname sga@target@IV.13.1.6\endcsname{ega4.c04.cor.13-1-6}
\expandafter\def\csname sga@target@IV.13.1.7\endcsname{ega4.c04.prop.13-1-7}
\expandafter\def\csname sga@target@IV.13.2.1\endcsname{ega4.c04.statement.13-2-1}
\expandafter\def\csname sga@target@IV.13.2.10\endcsname{ega4.c04.prop.13-2-10}
\expandafter\def\csname sga@target@IV.13.2.10.1\endcsname{ega4.c04.eq.13-2-10-1}
\expandafter\def\csname sga@target@IV.13.2.10.2\endcsname{ega4.c04.eq.13-2-10-2}
\expandafter\def\csname sga@target@IV.13.2.11\endcsname{ega4.c04.cor.13-2-11}
\expandafter\def\csname sga@target@IV.13.2.12\endcsname{ega4.c04.item.13-2-12}
\expandafter\def\csname sga@target@IV.13.2.2\endcsname{ega4.c04.def.13-2-2}
\expandafter\def\csname sga@target@IV.13.2.3\endcsname{ega4.c04.prop.13-2-3}
\expandafter\def\csname sga@target@IV.13.2.3.1\endcsname{ega4.c04.prop.13-2-3-1}
\expandafter\def\csname sga@target@IV.13.2.4\endcsname{ega4.c04.statement.13-2-4}
\expandafter\def\csname sga@target@IV.13.2.4.1\endcsname{ega4.c04.statement.13-2-4-1}
\expandafter\def\csname sga@target@IV.13.2.4.2\endcsname{ega4.c04.eq.13-2-4-2}
\expandafter\def\csname sga@target@IV.13.2.4.3\endcsname{ega4.c04.eq.13-2-4-3}
\expandafter\def\csname sga@target@IV.13.2.4.4\endcsname{ega4.c04.eq.13-2-4-4}
\expandafter\def\csname sga@target@IV.13.2.4.5\endcsname{ega4.c04.eq.13-2-4-5}
\expandafter\def\csname sga@target@IV.13.2.5\endcsname{ega4.c04.lem.13-2-5}
\expandafter\def\csname sga@target@IV.13.2.6\endcsname{ega4.c04.statement.13-2-6}
\expandafter\def\csname sga@target@IV.13.2.6.1\endcsname{ega4.c04.statement.13-2-6-1}
\expandafter\def\csname sga@target@IV.13.2.7\endcsname{ega4.c04.prop.13-2-7}
\expandafter\def\csname sga@target@IV.13.2.7.1\endcsname{ega4.c04.prop.13-2-7-1}
\expandafter\def\csname sga@target@IV.13.2.8\endcsname{ega4.c04.cor.13-2-8}
\expandafter\def\csname sga@target@IV.13.2.9\endcsname{ega4.c04.cor.13-2-9}
\expandafter\def\csname sga@target@IV.13.2.9.1\endcsname{ega4.c04.eq.13-2-9-1}
\expandafter\def\csname sga@target@IV.13.2.9.2\endcsname{ega4.c04.eq.13-2-9-2}
\expandafter\def\csname sga@target@IV.13.3.1\endcsname{ega4.c04.prop.13-3-1}
\expandafter\def\csname sga@target@IV.13.3.1.1\endcsname{ega4.c04.lem.13-3-1-1}
\expandafter\def\csname sga@target@IV.13.3.2\endcsname{ega4.c04.def.13-3-2}
\expandafter\def\csname sga@target@IV.13.3.3\endcsname{ega4.c04.prop.13-3-3}
\expandafter\def\csname sga@target@IV.13.3.4\endcsname{ega4.c04.cor.13-3-4}
\expandafter\def\csname sga@target@IV.13.3.4.1\endcsname{ega4.c04.cor.13-3-4-1}
\expandafter\def\csname sga@target@IV.13.3.5\endcsname{ega4.c04.cor.13-3-5}
\expandafter\def\csname sga@target@IV.13.3.6\endcsname{ega4.c04.prop.13-3-6}
\expandafter\def\csname sga@target@IV.13.3.6.1\endcsname{ega4.c04.prop.13-3-6-1}
\expandafter\def\csname sga@target@IV.13.3.6.2\endcsname{ega4.c04.eq.13-3-6-2}
\expandafter\def\csname sga@target@IV.13.3.7\endcsname{ega4.c04.prop.13-3-7}
\expandafter\def\csname sga@target@IV.13.3.7.1\endcsname{ega4.c04.prop.13-3-7-1}
\expandafter\def\csname sga@target@IV.13.3.8\endcsname{ega4.c04.prop.13-3-8}
\expandafter\def\csname sga@target@IV.13.3.9\endcsname{ega4.c04.statement.13-3-9}
\expandafter\def\csname sga@target@IV.14.1.1\endcsname{ega4.c04.statement.14-1-1}
\expandafter\def\csname sga@target@IV.14.1.2\endcsname{ega4.c04.prop.14-1-2}
\expandafter\def\csname sga@target@IV.14.1.3\endcsname{ega4.c04.item.14-1-3}
\expandafter\def\csname sga@target@IV.14.1.4\endcsname{ega4.c04.prop.14-1-4}
\expandafter\def\csname sga@target@IV.14.1.5\endcsname{ega4.c04.rem.14-1-5}
\expandafter\def\csname sga@target@IV.14.1.6\endcsname{ega4.c04.statement.14-1-6}
\expandafter\def\csname sga@target@IV.14.2.1\endcsname{ega4.c04.thm.14-2-1}
\expandafter\def\csname sga@target@IV.14.2.1.1\endcsname{ega4.c04.eq.14-2-1-1}
\expandafter\def\csname sga@target@IV.14.2.2\endcsname{ega4.c04.cor.14-2-2}
\expandafter\def\csname sga@target@IV.14.2.3\endcsname{ega4.c04.prop.14-2-3}
\expandafter\def\csname sga@target@IV.14.2.4\endcsname{ega4.c04.cor.14-2-4}
\expandafter\def\csname sga@target@IV.14.2.5\endcsname{ega4.c04.cor.14-2-5}
\expandafter\def\csname sga@target@IV.14.2.6\endcsname{ega4.c04.cor.14-2-6}
\expandafter\def\csname sga@target@IV.14.2.7\endcsname{ega4.c04.rem.14-2-7}
\expandafter\def\csname sga@target@IV.14.3\endcsname{ega4.c04.subsec.14-3.key-iv-14-3}
\expandafter\def\csname sga@target@IV.14.3.1\endcsname{ega4.c04.statement.14-3-1}
\expandafter\def\csname sga@target@IV.14.3.10\endcsname{ega4.c04.lem.14-3-10}
\expandafter\def\csname sga@target@IV.14.3.11\endcsname{ega4.c04.prop.14-3-11}
\expandafter\def\csname sga@target@IV.14.3.12\endcsname{ega4.c04.rem.14-3-12}
\expandafter\def\csname sga@target@IV.14.3.13\endcsname{ega4.c04.prop.14-3-13}
\expandafter\def\csname sga@target@IV.14.3.14\endcsname{ega4.c04.rem.14-3-14}
\expandafter\def\csname sga@target@IV.14.3.2\endcsname{ega4.c04.prop.14-3-2}
\expandafter\def\csname sga@target@IV.14.3.3\endcsname{ega4.c04.def.14-3-3}
\expandafter\def\csname sga@target@IV.14.3.3.1\endcsname{ega4.c04.subsec.14-3-3-1}
\expandafter\def\csname sga@target@IV.14.3.4\endcsname{ega4.c04.prop.14-3-4}
\expandafter\def\csname sga@target@IV.14.3.5\endcsname{ega4.c04.prop.14-3-5}
\expandafter\def\csname sga@target@IV.14.3.6\endcsname{ega4.c04.prop.14-3-6}
\expandafter\def\csname sga@target@IV.14.3.7\endcsname{ega4.c04.cor.14-3-7}
\expandafter\def\csname sga@target@IV.14.3.8\endcsname{ega4.c04.prop.14-3-8}
\expandafter\def\csname sga@target@IV.14.3.9\endcsname{ega4.c04.item.14-3-9}
\expandafter\def\csname sga@target@IV.14.4.1\endcsname{ega4.c04.thm.14-4-1}
\expandafter\def\csname sga@target@IV.14.4.1.1\endcsname{ega4.c04.lem.14-4-1-1}
\expandafter\def\csname sga@target@IV.14.4.1.2\endcsname{ega4.c04.lem.14-4-1-2}
\expandafter\def\csname sga@target@IV.14.4.1.3\endcsname{ega4.c04.lem.14-4-1-3}
\expandafter\def\csname sga@target@IV.14.4.10\endcsname{ega4.c04.item.14-4-10}
\expandafter\def\csname sga@target@IV.14.4.2\endcsname{ega4.c04.cor.14-4-2}
\expandafter\def\csname sga@target@IV.14.4.3\endcsname{ega4.c04.cor.14-4-3}
\expandafter\def\csname sga@target@IV.14.4.4\endcsname{ega4.c04.cor.14-4-4}
\expandafter\def\csname sga@target@IV.14.4.5\endcsname{ega4.c04.rem.14-4-5}
\expandafter\def\csname sga@target@IV.14.4.6\endcsname{ega4.c04.cor.14-4-6}
\expandafter\def\csname sga@target@IV.14.4.7\endcsname{ega4.c04.prop.14-4-7}
\expandafter\def\csname sga@target@IV.14.4.7.1\endcsname{ega4.c04.eq.14-4-7-1}
\expandafter\def\csname sga@target@IV.14.4.7.2\endcsname{ega4.c04.eq.14-4-7-2}
\expandafter\def\csname sga@target@IV.14.4.8\endcsname{ega4.c04.prop.14-4-8}
\expandafter\def\csname sga@target@IV.14.4.8.1\endcsname{ega4.c04.lem.14-4-8-1}
\expandafter\def\csname sga@target@IV.14.4.9\endcsname{ega4.c04.cor.14-4-9}
\expandafter\def\csname sga@target@IV.14.5\endcsname{ega4.c04.subsec.14-5.key-iv-14-5}
\expandafter\def\csname sga@target@IV.14.5.1\endcsname{ega4.c04.lem.14-5-1}
\expandafter\def\csname sga@target@IV.14.5.1.1\endcsname{ega4.c04.eq.14-5-1-1}
\expandafter\def\csname sga@target@IV.14.5.1.2\endcsname{ega4.c04.eq.14-5-1-2}
\expandafter\def\csname sga@target@IV.14.5.1.3\endcsname{ega4.c04.eq.14-5-1-3}
\expandafter\def\csname sga@target@IV.14.5.1.4\endcsname{ega4.c04.eq.14-5-1-4}
\expandafter\def\csname sga@target@IV.14.5.10\endcsname{ega4.c04.prop.14-5-10}
\expandafter\def\csname sga@target@IV.14.5.11\endcsname{ega4.c04.cor.14-5-11}
\expandafter\def\csname sga@target@IV.14.5.11.1\endcsname{ega4.c04.lem.14-5-11-1}
\expandafter\def\csname sga@target@IV.14.5.12\endcsname{ega4.c04.cor.14-5-12}
\expandafter\def\csname sga@target@IV.14.5.2\endcsname{ega4.c04.cor.14-5-2}
\expandafter\def\csname sga@target@IV.14.5.2.1\endcsname{ega4.c04.eq.14-5-2-1}
\expandafter\def\csname sga@target@IV.14.5.3\endcsname{ega4.c04.prop.14-5-3}
\expandafter\def\csname sga@target@IV.14.5.4\endcsname{ega4.c04.cor.14-5-4}
\expandafter\def\csname sga@target@IV.14.5.5\endcsname{ega4.c04.rem.14-5-5}
\expandafter\def\csname sga@target@IV.14.5.6\endcsname{ega4.c04.prop.14-5-6}
\expandafter\def\csname sga@target@IV.14.5.7\endcsname{ega4.c04.rem.14-5-7}
\expandafter\def\csname sga@target@IV.14.5.8\endcsname{ega4.c04.prop.14-5-8}
\expandafter\def\csname sga@target@IV.14.5.9\endcsname{ega4.c04.prop.14-5-9}
\expandafter\def\csname sga@target@IV.15.1.1\endcsname{ega4.c04.prop.15-1-1}
\expandafter\def\csname sga@target@IV.15.1.1-proof\endcsname{ega4.c04.statement.15-1-1-proof}
\expandafter\def\csname sga@target@IV.15.1.1.1\endcsname{ega4.c04.eq.15-1-1-1}
\expandafter\def\csname sga@target@IV.15.1.1.2\endcsname{ega4.c04.eq.15-1-1-2}
\expandafter\def\csname sga@target@IV.15.1.1.3\endcsname{ega4.c04.eq.15-1-1-3}
\expandafter\def\csname sga@target@IV.15.1.1.4\endcsname{ega4.c04.eq.15-1-1-4}
\expandafter\def\csname sga@target@IV.15.1.1.5\endcsname{ega4.c04.eq.15-1-1-5}
\expandafter\def\csname sga@target@IV.15.1.1.6\endcsname{ega4.c04.lem.15-1-1-6}
\expandafter\def\csname sga@target@IV.15.1.2\endcsname{ega4.c04.cor.15-1-2}
\expandafter\def\csname sga@target@IV.15.1.3\endcsname{ega4.c04.cor.15-1-3}
\expandafter\def\csname sga@target@IV.15.1.4\endcsname{ega4.c04.statement.15-1-4}
\expandafter\def\csname sga@target@IV.15.2.1\endcsname{ega4.c04.statement.15-2-1}
\expandafter\def\csname sga@target@IV.15.2.2\endcsname{ega4.c04.thm.15-2-2}
\expandafter\def\csname sga@target@IV.15.2.2-proof\endcsname{ega4.c04.item.15-2-2-proof}
\expandafter\def\csname sga@target@IV.15.2.2.1\endcsname{ega4.c04.lem.15-2-2-1}
\expandafter\def\csname sga@target@IV.15.2.3\endcsname{ega4.c04.cor.15-2-3}
\expandafter\def\csname sga@target@IV.15.2.4\endcsname{ega4.c04.item.15-2-4}
\expandafter\def\csname sga@target@IV.15.3.1\endcsname{ega4.c04.prop.15-3-1}
\expandafter\def\csname sga@target@IV.15.3.1-proof\endcsname{ega4.c04.item.15-3-1-proof}
\expandafter\def\csname sga@target@IV.15.3.2\endcsname{ega4.c04.prop.15-3-2}
\expandafter\def\csname sga@target@IV.15.3.2-proof\endcsname{ega4.c04.item.15-3-2-proof}
\expandafter\def\csname sga@target@IV.15.3.3\endcsname{ega4.c04.cor.15-3-3}
\expandafter\def\csname sga@target@IV.15.3.3-proof\endcsname{ega4.c04.item.15-3-3-proof}
\expandafter\def\csname sga@target@IV.15.3.4\endcsname{ega4.c04.item.15-3-4}
\expandafter\def\csname sga@target@IV.15.4.1\endcsname{ega4.c04.prop.15-4-1}
\expandafter\def\csname sga@target@IV.15.4.2\endcsname{ega4.c04.prop.15-4-2}
\expandafter\def\csname sga@target@IV.15.4.3\endcsname{ega4.c04.prop.15-4-3}
\expandafter\def\csname sga@target@IV.15.4.3-proof\endcsname{ega4.c04.item.15-4-3-proof}
\expandafter\def\csname sga@target@IV.15.5.1\endcsname{ega4.c04.prop.15-5-1}
\expandafter\def\csname sga@target@IV.15.5.10\endcsname{ega4.c04.item.15-5-10}
\expandafter\def\csname sga@target@IV.15.5.2\endcsname{ega4.c04.cor.15-5-2}
\expandafter\def\csname sga@target@IV.15.5.2.1\endcsname{ega4.c04.eq.15-5-2-1}
\expandafter\def\csname sga@target@IV.15.5.2.2\endcsname{ega4.c04.lem.15-5-2-2}
\expandafter\def\csname sga@target@IV.15.5.2.3\endcsname{ega4.c04.eq.15-5-2-3}
\expandafter\def\csname sga@target@IV.15.5.2.4\endcsname{ega4.c04.eq.15-5-2-4}
\expandafter\def\csname sga@target@IV.15.5.3\endcsname{ega4.c04.prop.15-5-3}
\expandafter\def\csname sga@target@IV.15.5.4\endcsname{ega4.c04.cor.15-5-4}
\expandafter\def\csname sga@target@IV.15.5.5\endcsname{ega4.c04.statement.15-5-5}
\expandafter\def\csname sga@target@IV.15.5.6\endcsname{ega4.c04.lem.15-5-6}
\expandafter\def\csname sga@target@IV.15.5.6.1\endcsname{ega4.c04.lem.15-5-6-1}
\expandafter\def\csname sga@target@IV.15.5.7\endcsname{ega4.c04.prop.15-5-7}
\expandafter\def\csname sga@target@IV.15.5.7.1\endcsname{ega4.c04.eq.15-5-7-1}
\expandafter\def\csname sga@target@IV.15.5.8\endcsname{ega4.c04.statement.15-5-8}
\expandafter\def\csname sga@target@IV.15.5.9\endcsname{ega4.c04.prop.15-5-9}
\expandafter\def\csname sga@target@IV.15.6.1\endcsname{ega4.c04.prop.15-6-1}
\expandafter\def\csname sga@target@IV.15.6.1.1\endcsname{ega4.c04.eq.15-6-1-1}
\expandafter\def\csname sga@target@IV.15.6.10\endcsname{ega4.c04.item.15-6-10}
\expandafter\def\csname sga@target@IV.15.6.2\endcsname{ega4.c04.cor.15-6-2}
\expandafter\def\csname sga@target@IV.15.6.3\endcsname{ega4.c04.prop.15-6-3}
\expandafter\def\csname sga@target@IV.15.6.4\endcsname{ega4.c04.prop.15-6-4}
\expandafter\def\csname sga@target@IV.15.6.5\endcsname{ega4.c04.cor.15-6-5}
\expandafter\def\csname sga@target@IV.15.6.6\endcsname{ega4.c04.prop.15-6-6}
\expandafter\def\csname sga@target@IV.15.6.7\endcsname{ega4.c04.cor.15-6-7}
\expandafter\def\csname sga@target@IV.15.6.8\endcsname{ega4.c04.prop.15-6-8}
\expandafter\def\csname sga@target@IV.15.6.9\endcsname{ega4.c04.cor.15-6-9}
\expandafter\def\csname sga@target@IV.15.7.1\endcsname{ega4.c04.def.15-7-1}
\expandafter\def\csname sga@target@IV.15.7.10\endcsname{ega4.c04.cor.15-7-10}
\expandafter\def\csname sga@target@IV.15.7.11\endcsname{ega4.c04.cor.15-7-11}
\expandafter\def\csname sga@target@IV.15.7.11.1\endcsname{ega4.c04.lem.15-7-11-1}
\expandafter\def\csname sga@target@IV.15.7.2\endcsname{ega4.c04.prop.15-7-2}
\expandafter\def\csname sga@target@IV.15.7.3\endcsname{ega4.c04.rem.15-7-3}
\expandafter\def\csname sga@target@IV.15.7.4\endcsname{ega4.c04.cor.15-7-4}
\expandafter\def\csname sga@target@IV.15.7.5\endcsname{ega4.c04.cor.15-7-5}
\expandafter\def\csname sga@target@IV.15.7.6\endcsname{ega4.c04.cor.15-7-6}
\expandafter\def\csname sga@target@IV.15.7.7\endcsname{ega4.c04.item.15-7-7}
\expandafter\def\csname sga@target@IV.15.7.8\endcsname{ega4.c04.statement.15-7-8}
\expandafter\def\csname sga@target@IV.15.7.9\endcsname{ega4.c04.prop.15-7-9}
\expandafter\def\csname sga@target@IV.16.1\endcsname{ega4.c04.subsec.16-1}
\expandafter\def\csname sga@target@IV.16.1.1\endcsname{ega4.c04.statement.16-1-1}
\expandafter\def\csname sga@target@IV.16.1.10\endcsname{ega4.c04.cor.16-1-10}
\expandafter\def\csname sga@target@IV.16.1.11\endcsname{ega4.c04.rem.16-1-11}
\expandafter\def\csname sga@target@IV.16.1.2\endcsname{ega4.c04.def.16-1-2}
\expandafter\def\csname sga@target@IV.16.1.2.1\endcsname{ega4.c04.def.16-1-2-1}
\expandafter\def\csname sga@target@IV.16.1.2.2\endcsname{ega4.c04.eq.16-1-2-2}
\expandafter\def\csname sga@target@IV.16.1.3\endcsname{ega4.c04.statement.16-1-3}
\expandafter\def\csname sga@target@IV.16.1.4\endcsname{ega4.c04.statement.16-1-4}
\expandafter\def\csname sga@target@IV.16.1.5\endcsname{ega4.c04.prop.16-1-5}
\expandafter\def\csname sga@target@IV.16.1.6\endcsname{ega4.c04.cor.16-1-6}
\expandafter\def\csname sga@target@IV.16.1.7\endcsname{ega4.c04.cor.16-1-7}
\expandafter\def\csname sga@target@IV.16.1.8\endcsname{ega4.c04.prop.16-1-8}
\expandafter\def\csname sga@target@IV.16.1.9\endcsname{ega4.c04.prop.16-1-9}
\expandafter\def\csname sga@target@IV.16.10\endcsname{ega4.c04.subsec.16-10}
\expandafter\def\csname sga@target@IV.16.10.1\endcsname{ega4.c04.def.16-10-1}
\expandafter\def\csname sga@target@IV.16.10.2\endcsname{ega4.c04.prop.16-10-2}
\expandafter\def\csname sga@target@IV.16.10.3\endcsname{ega4.c04.prop.16-10-3}
\expandafter\def\csname sga@target@IV.16.10.4\endcsname{ega4.c04.prop.16-10-4}
\expandafter\def\csname sga@target@IV.16.10.5\endcsname{ega4.c04.prop.16-10-5}
\expandafter\def\csname sga@target@IV.16.10.6\endcsname{ega4.c04.rem.16-10-6}
\expandafter\def\csname sga@target@IV.16.11\endcsname{ega4.c04.subsec.16-11}
\expandafter\def\csname sga@target@IV.16.11.1\endcsname{ega4.c04.statement.16-11-1}
\expandafter\def\csname sga@target@IV.16.11.1.1\endcsname{ega4.c04.statement.16-11-1-1}
\expandafter\def\csname sga@target@IV.16.11.1.2\endcsname{ega4.c04.statement.16-11-1-2}
\expandafter\def\csname sga@target@IV.16.11.1.3\endcsname{ega4.c04.eq.16-11-1-3}
\expandafter\def\csname sga@target@IV.16.11.1.4\endcsname{ega4.c04.eq.16-11-1-4}
\expandafter\def\csname sga@target@IV.16.11.1.5\endcsname{ega4.c04.eq.16-11-1-5}
\expandafter\def\csname sga@target@IV.16.11.1.6\endcsname{ega4.c04.eq.16-11-1-6}
\expandafter\def\csname sga@target@IV.16.11.1.7\endcsname{ega4.c04.eq.16-11-1-7}
\expandafter\def\csname sga@target@IV.16.11.2\endcsname{ega4.c04.thm.16-11-2}
\expandafter\def\csname sga@target@IV.16.11.2.1\endcsname{ega4.c04.eq.16-11-2-1}
\expandafter\def\csname sga@target@IV.16.11.2.2\endcsname{ega4.c04.eq.16-11-2-2}
\expandafter\def\csname sga@target@IV.16.11.2.3\endcsname{ega4.c04.eq.16-11-2-3}
\expandafter\def\csname sga@target@IV.16.11.3\endcsname{ega4.c04.statement.16-11-3}
\expandafter\def\csname sga@target@IV.16.12\endcsname{ega4.c04.subsec.16-12}
\expandafter\def\csname sga@target@IV.16.12.1\endcsname{ega4.c04.statement.16-12-1}
\expandafter\def\csname sga@target@IV.16.12.2\endcsname{ega4.c04.thm.16-12-2}
\expandafter\def\csname sga@target@IV.16.12.3\endcsname{ega4.c04.statement.16-12-3}
\expandafter\def\csname sga@target@IV.16.2\endcsname{ega4.c04.subsec.16-2}
\expandafter\def\csname sga@target@IV.16.2.1\endcsname{ega4.c04.statement.16-2-1}
\expandafter\def\csname sga@target@IV.16.2.1.1\endcsname{ega4.c04.statement.16-2-1-1}
\expandafter\def\csname sga@target@IV.16.2.1.2\endcsname{ega4.c04.eq.16-2-1-2}
\expandafter\def\csname sga@target@IV.16.2.1.3\endcsname{ega4.c04.eq.16-2-1-3}
\expandafter\def\csname sga@target@IV.16.2.1.4\endcsname{ega4.c04.eq.16-2-1-4}
\expandafter\def\csname sga@target@IV.16.2.2\endcsname{ega4.c04.prop.16-2-2}
\expandafter\def\csname sga@target@IV.16.2.3\endcsname{ega4.c04.cor.16-2-3}
\expandafter\def\csname sga@target@IV.16.2.3.1\endcsname{ega4.c04.eq.16-2-3-1}
\expandafter\def\csname sga@target@IV.16.2.3.2\endcsname{ega4.c04.eq.16-2-3-2}
\expandafter\def\csname sga@target@IV.16.2.4\endcsname{ega4.c04.cor.16-2-4}
\expandafter\def\csname sga@target@IV.16.2.5\endcsname{ega4.c04.statement.16-2-5}
\expandafter\def\csname sga@target@IV.16.2.6\endcsname{ega4.c04.statement.16-2-6}
\expandafter\def\csname sga@target@IV.16.2.6.1\endcsname{ega4.c04.statement.16-2-6-1}
\expandafter\def\csname sga@target@IV.16.2.6.2\endcsname{ega4.c04.eq.16-2-6-2}
\expandafter\def\csname sga@target@IV.16.2.7\endcsname{ega4.c04.prop.16-2-7}
\expandafter\def\csname sga@target@IV.16.2.7.1\endcsname{ega4.c04.prop.16-2-7-1}
\expandafter\def\csname sga@target@IV.16.3\endcsname{ega4.c04.subsec.16-3}
\expandafter\def\csname sga@target@IV.16.3.1\endcsname{ega4.c04.def.16-3-1}
\expandafter\def\csname sga@target@IV.16.3.1.1\endcsname{ega4.c04.eq.16-3-1-1}
\expandafter\def\csname sga@target@IV.16.3.2\endcsname{ega4.c04.statement.16-3-2}
\expandafter\def\csname sga@target@IV.16.3.3\endcsname{ega4.c04.statement.16-3-3}
\expandafter\def\csname sga@target@IV.16.3.3.1\endcsname{ega4.c04.statement.16-3-3-1}
\expandafter\def\csname sga@target@IV.16.3.4\endcsname{ega4.c04.prop.16-3-4}
\expandafter\def\csname sga@target@IV.16.3.5\endcsname{ega4.c04.statement.16-3-5}
\expandafter\def\csname sga@target@IV.16.3.6\endcsname{ega4.c04.def.16-3-6}
\expandafter\def\csname sga@target@IV.16.3.6.1\endcsname{ega4.c04.eq.16-3-6-1}
\expandafter\def\csname sga@target@IV.16.3.7\endcsname{ega4.c04.statement.16-3-7}
\expandafter\def\csname sga@target@IV.16.3.7.1\endcsname{ega4.c04.eq.16-3-7-1}
\expandafter\def\csname sga@target@IV.16.3.8\endcsname{ega4.c04.prop.16-3-8}
\expandafter\def\csname sga@target@IV.16.3.9\endcsname{ega4.c04.prop.16-3-9}
\expandafter\def\csname sga@target@IV.16.4\endcsname{ega4.c04.subsec.16-4}
\expandafter\def\csname sga@target@IV.16.4.1\endcsname{ega4.c04.statement.16-4-1}
\expandafter\def\csname sga@target@IV.16.4.1.1\endcsname{ega4.c04.statement.16-4-1-1}
\expandafter\def\csname sga@target@IV.16.4.1.2\endcsname{ega4.c04.statement.16-4-1-2}
\expandafter\def\csname sga@target@IV.16.4.1.3\endcsname{ega4.c04.eq.16-4-1-3}
\expandafter\def\csname sga@target@IV.16.4.1.4\endcsname{ega4.c04.eq.16-4-1-4}
\expandafter\def\csname sga@target@IV.16.4.1.5\endcsname{ega4.c04.eq.16-4-1-5}
\expandafter\def\csname sga@target@IV.16.4.10\endcsname{ega4.c04.cor.16-4-10}
\expandafter\def\csname sga@target@IV.16.4.11\endcsname{ega4.c04.prop.16-4-11}
\expandafter\def\csname sga@target@IV.16.4.11.1\endcsname{ega4.c04.prop.16-4-11-1}
\expandafter\def\csname sga@target@IV.16.4.11.2\endcsname{ega4.c04.prop.16-4-11-2}
\expandafter\def\csname sga@target@IV.16.4.11.3\endcsname{ega4.c04.eq.16-4-11-3}
\expandafter\def\csname sga@target@IV.16.4.11.4\endcsname{ega4.c04.eq.16-4-11-4}
\expandafter\def\csname sga@target@IV.16.4.12\endcsname{ega4.c04.cor.16-4-12}
\expandafter\def\csname sga@target@IV.16.4.13\endcsname{ega4.c04.cor.16-4-13}
\expandafter\def\csname sga@target@IV.16.4.14\endcsname{ega4.c04.prop.16-4-14}
\expandafter\def\csname sga@target@IV.16.4.14.1\endcsname{ega4.c04.prop.16-4-14-1}
\expandafter\def\csname sga@target@IV.16.4.15\endcsname{ega4.c04.cor.16-4-15}
\expandafter\def\csname sga@target@IV.16.4.15.1\endcsname{ega4.c04.cor.16-4-15-1}
\expandafter\def\csname sga@target@IV.16.4.15.2\endcsname{ega4.c04.eq.16-4-15-2}
\expandafter\def\csname sga@target@IV.16.4.16\endcsname{ega4.c04.cor.16-4-16}
\expandafter\def\csname sga@target@IV.16.4.16.1\endcsname{ega4.c04.cor.16-4-16-1}
\expandafter\def\csname sga@target@IV.16.4.16.2\endcsname{ega4.c04.eq.16-4-16-2}
\expandafter\def\csname sga@target@IV.16.4.17\endcsname{ega4.c04.cor.16-4-17}
\expandafter\def\csname sga@target@IV.16.4.18\endcsname{ega4.c04.statement.16-4-18}
\expandafter\def\csname sga@target@IV.16.4.18.1\endcsname{ega4.c04.statement.16-4-18-1}
\expandafter\def\csname sga@target@IV.16.4.18.2\endcsname{ega4.c04.statement.16-4-18-2}
\expandafter\def\csname sga@target@IV.16.4.18.3\endcsname{ega4.c04.eq.16-4-18-3}
\expandafter\def\csname sga@target@IV.16.4.19\endcsname{ega4.c04.cor.16-4-19}
\expandafter\def\csname sga@target@IV.16.4.19.1\endcsname{ega4.c04.cor.16-4-19-1}
\expandafter\def\csname sga@target@IV.16.4.2\endcsname{ega4.c04.statement.16-4-2}
\expandafter\def\csname sga@target@IV.16.4.20\endcsname{ega4.c04.prop.16-4-20}
\expandafter\def\csname sga@target@IV.16.4.21\endcsname{ega4.c04.cor.16-4-21}
\expandafter\def\csname sga@target@IV.16.4.21.1\endcsname{ega4.c04.cor.16-4-21-1}
\expandafter\def\csname sga@target@IV.16.4.22\endcsname{ega4.c04.cor.16-4-22}
\expandafter\def\csname sga@target@IV.16.4.23\endcsname{ega4.c04.prop.16-4-23}
\expandafter\def\csname sga@target@IV.16.4.23.1\endcsname{ega4.c04.prop.16-4-23-1}
\expandafter\def\csname sga@target@IV.16.4.24\endcsname{ega4.c04.statement.16-4-24}
\expandafter\def\csname sga@target@IV.16.4.3\endcsname{ega4.c04.statement.16-4-3}
\expandafter\def\csname sga@target@IV.16.4.3.1\endcsname{ega4.c04.statement.16-4-3-1}
\expandafter\def\csname sga@target@IV.16.4.3.2\endcsname{ega4.c04.statement.16-4-3-2}
\expandafter\def\csname sga@target@IV.16.4.3.3\endcsname{ega4.c04.eq.16-4-3-3}
\expandafter\def\csname sga@target@IV.16.4.3.4\endcsname{ega4.c04.eq.16-4-3-4}
\expandafter\def\csname sga@target@IV.16.4.3.5\endcsname{ega4.c04.eq.16-4-3-5}
\expandafter\def\csname sga@target@IV.16.4.3.6\endcsname{ega4.c04.eq.16-4-3-6}
\expandafter\def\csname sga@target@IV.16.4.3.7\endcsname{ega4.c04.eq.16-4-3-7}
\expandafter\def\csname sga@target@IV.16.4.4\endcsname{ega4.c04.statement.16-4-4}
\expandafter\def\csname sga@target@IV.16.4.5\endcsname{ega4.c04.prop.16-4-5}
\expandafter\def\csname sga@target@IV.16.4.6\endcsname{ega4.c04.cor.16-4-6}
\expandafter\def\csname sga@target@IV.16.4.7\endcsname{ega4.c04.statement.16-4-7}
\expandafter\def\csname sga@target@IV.16.4.7.1\endcsname{ega4.c04.statement.16-4-7-1}
\expandafter\def\csname sga@target@IV.16.4.7.2\endcsname{ega4.c04.eq.16-4-7-2}
\expandafter\def\csname sga@target@IV.16.4.7.3\endcsname{ega4.c04.eq.16-4-7-3}
\expandafter\def\csname sga@target@IV.16.4.8\endcsname{ega4.c04.prop.16-4-8}
\expandafter\def\csname sga@target@IV.16.4.8.1\endcsname{ega4.c04.prop.16-4-8-1}
\expandafter\def\csname sga@target@IV.16.4.8.2\endcsname{ega4.c04.eq.16-4-8-2}
\expandafter\def\csname sga@target@IV.16.4.8.3\endcsname{ega4.c04.lem.16-4-8-3}
\expandafter\def\csname sga@target@IV.16.4.8.4\endcsname{ega4.c04.eq.16-4-8-4}
\expandafter\def\csname sga@target@IV.16.4.9\endcsname{ega4.c04.cor.16-4-9}
\expandafter\def\csname sga@target@IV.16.4.9.1\endcsname{ega4.c04.cor.16-4-9-1}
\expandafter\def\csname sga@target@IV.16.5\endcsname{ega4.c04.subsec.16-5}
\expandafter\def\csname sga@target@IV.16.5.1\endcsname{ega4.c04.statement.16-5-1}
\expandafter\def\csname sga@target@IV.16.5.1.1\endcsname{ega4.c04.eq.16-5-1-1}
\expandafter\def\csname sga@target@IV.16.5.1.2\endcsname{ega4.c04.eq.16-5-1-2}
\expandafter\def\csname sga@target@IV.16.5.10\endcsname{ega4.c04.statement.16-5-10}
\expandafter\def\csname sga@target@IV.16.5.10.1\endcsname{ega4.c04.statement.16-5-10-1}
\expandafter\def\csname sga@target@IV.16.5.10.2\endcsname{ega4.c04.statement.16-5-10-2}
\expandafter\def\csname sga@target@IV.16.5.11\endcsname{ega4.c04.prop.16-5-11}
\expandafter\def\csname sga@target@IV.16.5.12\endcsname{ega4.c04.statement.16-5-12}
\expandafter\def\csname sga@target@IV.16.5.12.1\endcsname{ega4.c04.statement.16-5-12-1}
\expandafter\def\csname sga@target@IV.16.5.12.2\endcsname{ega4.c04.eq.16-5-12-2}
\expandafter\def\csname sga@target@IV.16.5.13\endcsname{ega4.c04.statement.16-5-13}
\expandafter\def\csname sga@target@IV.16.5.13.1\endcsname{ega4.c04.statement.16-5-13-1}
\expandafter\def\csname sga@target@IV.16.5.13.2\endcsname{ega4.c04.eq.16-5-13-2}
\expandafter\def\csname sga@target@IV.16.5.13.3\endcsname{ega4.c04.eq.16-5-13-3}
\expandafter\def\csname sga@target@IV.16.5.13.4\endcsname{ega4.c04.eq.16-5-13-4}
\expandafter\def\csname sga@target@IV.16.5.13.5\endcsname{ega4.c04.eq.16-5-13-5}
\expandafter\def\csname sga@target@IV.16.5.14\endcsname{ega4.c04.statement.16-5-14}
\expandafter\def\csname sga@target@IV.16.5.14.1\endcsname{ega4.c04.statement.16-5-14-1}
\expandafter\def\csname sga@target@IV.16.5.14.2\endcsname{ega4.c04.eq.16-5-14-2}
\expandafter\def\csname sga@target@IV.16.5.14.3\endcsname{ega4.c04.eq.16-5-14-3}
\expandafter\def\csname sga@target@IV.16.5.15\endcsname{ega4.c04.statement.16-5-15}
\expandafter\def\csname sga@target@IV.16.5.16\endcsname{ega4.c04.prop.16-5-16}
\expandafter\def\csname sga@target@IV.16.5.17\endcsname{ega4.c04.prop.16-5-17}
\expandafter\def\csname sga@target@IV.16.5.18\endcsname{ega4.c04.cor.16-5-18}
\expandafter\def\csname sga@target@IV.16.5.2\endcsname{ega4.c04.prop.16-5-2}
\expandafter\def\csname sga@target@IV.16.5.3\endcsname{ega4.c04.prop.16-5-3}
\expandafter\def\csname sga@target@IV.16.5.3.1\endcsname{ega4.c04.eq.16-5-3-1}
\expandafter\def\csname sga@target@IV.16.5.4\endcsname{ega4.c04.statement.16-5-4}
\expandafter\def\csname sga@target@IV.16.5.5\endcsname{ega4.c04.cor.16-5-5}
\expandafter\def\csname sga@target@IV.16.5.5.1\endcsname{ega4.c04.cor.16-5-5-1}
\expandafter\def\csname sga@target@IV.16.5.6\endcsname{ega4.c04.cor.16-5-6}
\expandafter\def\csname sga@target@IV.16.5.7\endcsname{ega4.c04.statement.16-5-7}
\expandafter\def\csname sga@target@IV.16.5.7.1\endcsname{ega4.c04.statement.16-5-7-1}
\expandafter\def\csname sga@target@IV.16.5.8\endcsname{ega4.c04.statement.16-5-8}
\expandafter\def\csname sga@target@IV.16.5.8.1\endcsname{ega4.c04.eq.16-5-8-1}
\expandafter\def\csname sga@target@IV.16.5.8.2\endcsname{ega4.c04.eq.16-5-8-2}
\expandafter\def\csname sga@target@IV.16.5.9\endcsname{ega4.c04.statement.16-5-9}
\expandafter\def\csname sga@target@IV.16.5.9.1\endcsname{ega4.c04.statement.16-5-9-1}
\expandafter\def\csname sga@target@IV.16.6\endcsname{ega4.c04.subsec.16-6}
\expandafter\def\csname sga@target@IV.16.6.1\endcsname{ega4.c04.statement.16-6-1}
\expandafter\def\csname sga@target@IV.16.6.1.1\endcsname{ega4.c04.statement.16-6-1-1}
\expandafter\def\csname sga@target@IV.16.6.1.2\endcsname{ega4.c04.statement.16-6-1-2}
\expandafter\def\csname sga@target@IV.16.6.2\endcsname{ega4.c04.thm.16-6-2}
\expandafter\def\csname sga@target@IV.16.6.2.1\endcsname{ega4.c04.eq.16-6-2-1}
\expandafter\def\csname sga@target@IV.16.6.2.2\endcsname{ega4.c04.eq.16-6-2-2}
\expandafter\def\csname sga@target@IV.16.6.2.3\endcsname{ega4.c04.eq.16-6-2-3}
\expandafter\def\csname sga@target@IV.16.6.2.4\endcsname{ega4.c04.eq.16-6-2-4}
\expandafter\def\csname sga@target@IV.16.6.2.5\endcsname{ega4.c04.eq.16-6-2-5}
\expandafter\def\csname sga@target@IV.16.6.3\endcsname{ega4.c04.def.16-6-3}
\expandafter\def\csname sga@target@IV.16.6.4\endcsname{ega4.c04.prop.16-6-4}
\expandafter\def\csname sga@target@IV.16.6.4.1\endcsname{ega4.c04.prop.16-6-4-1}
\expandafter\def\csname sga@target@IV.16.6.5\endcsname{ega4.c04.statement.16-6-5}
\expandafter\def\csname sga@target@IV.16.6.5.1\endcsname{ega4.c04.statement.16-6-5-1}
\expandafter\def\csname sga@target@IV.16.7\endcsname{ega4.c04.subsec.16-7}
\expandafter\def\csname sga@target@IV.16.7.1\endcsname{ega4.c04.statement.16-7-1}
\expandafter\def\csname sga@target@IV.16.7.1.1\endcsname{ega4.c04.eq.16-7-1-1}
\expandafter\def\csname sga@target@IV.16.7.1.2\endcsname{ega4.c04.eq.16-7-1-2}
\expandafter\def\csname sga@target@IV.16.7.10\endcsname{ega4.c04.rem.16-7-10}
\expandafter\def\csname sga@target@IV.16.7.2\endcsname{ega4.c04.statement.16-7-2}
\expandafter\def\csname sga@target@IV.16.7.2.1\endcsname{ega4.c04.statement.16-7-2-1}
\expandafter\def\csname sga@target@IV.16.7.3\endcsname{ega4.c04.prop.16-7-3}
\expandafter\def\csname sga@target@IV.16.7.4\endcsname{ega4.c04.statement.16-7-4}
\expandafter\def\csname sga@target@IV.16.7.4.1\endcsname{ega4.c04.statement.16-7-4-1}
\expandafter\def\csname sga@target@IV.16.7.5\endcsname{ega4.c04.statement.16-7-5}
\expandafter\def\csname sga@target@IV.16.7.5.1\endcsname{ega4.c04.statement.16-7-5-1}
\expandafter\def\csname sga@target@IV.16.7.5.2\endcsname{ega4.c04.statement.16-7-5-2}
\expandafter\def\csname sga@target@IV.16.7.5.3\endcsname{ega4.c04.statement.16-7-5-3}
\expandafter\def\csname sga@target@IV.16.7.5.4\endcsname{ega4.c04.statement.16-7-5-4}
\expandafter\def\csname sga@target@IV.16.7.6\endcsname{ega4.c04.prop.16-7-6}
\expandafter\def\csname sga@target@IV.16.7.7\endcsname{ega4.c04.statement.16-7-7}
\expandafter\def\csname sga@target@IV.16.7.7.1\endcsname{ega4.c04.eq.16-7-7-1}
\expandafter\def\csname sga@target@IV.16.7.7.2\endcsname{ega4.c04.eq.16-7-7-2}
\expandafter\def\csname sga@target@IV.16.7.8\endcsname{ega4.c04.statement.16-7-8}
\expandafter\def\csname sga@target@IV.16.7.8.1\endcsname{ega4.c04.statement.16-7-8-1}
\expandafter\def\csname sga@target@IV.16.7.8.2\endcsname{ega4.c04.eq.16-7-8-2}
\expandafter\def\csname sga@target@IV.16.7.9\endcsname{ega4.c04.statement.16-7-9}
\expandafter\def\csname sga@target@IV.16.7.9.1\endcsname{ega4.c04.statement.16-7-9-1}
\expandafter\def\csname sga@target@IV.16.8\endcsname{ega4.c04.subsec.16-8}
\expandafter\def\csname sga@target@IV.16.8.1\endcsname{ega4.c04.def.16-8-1}
\expandafter\def\csname sga@target@IV.16.8.10\endcsname{ega4.c04.cor.16-8-10}
\expandafter\def\csname sga@target@IV.16.8.11\endcsname{ega4.c04.statement.16-8-11}
\expandafter\def\csname sga@target@IV.16.8.2\endcsname{ega4.c04.statement.16-8-2}
\expandafter\def\csname sga@target@IV.16.8.3\endcsname{ega4.c04.statement.16-8-3}
\expandafter\def\csname sga@target@IV.16.8.3.1\endcsname{ega4.c04.eq.16-8-3-1}
\expandafter\def\csname sga@target@IV.16.8.4\endcsname{ega4.c04.prop.16-8-4}
\expandafter\def\csname sga@target@IV.16.8.4.1\endcsname{ega4.c04.prop.16-8-4-1}
\expandafter\def\csname sga@target@IV.16.8.5\endcsname{ega4.c04.statement.16-8-5}
\expandafter\def\csname sga@target@IV.16.8.5.1\endcsname{ega4.c04.eq.16-8-5-1}
\expandafter\def\csname sga@target@IV.16.8.5.2\endcsname{ega4.c04.eq.16-8-5-2}
\expandafter\def\csname sga@target@IV.16.8.6\endcsname{ega4.c04.prop.16-8-6}
\expandafter\def\csname sga@target@IV.16.8.7\endcsname{ega4.c04.statement.16-8-7}
\expandafter\def\csname sga@target@IV.16.8.7.1\endcsname{ega4.c04.eq.16-8-7-1}
\expandafter\def\csname sga@target@IV.16.8.8\endcsname{ega4.c04.prop.16-8-8}
\expandafter\def\csname sga@target@IV.16.8.8.1\endcsname{ega4.c04.eq.16-8-8-1}
\expandafter\def\csname sga@target@IV.16.8.8.2\endcsname{ega4.c04.eq.16-8-8-2}
\expandafter\def\csname sga@target@IV.16.8.9\endcsname{ega4.c04.prop.16-8-9}
\expandafter\def\csname sga@target@IV.16.8.9.1\endcsname{ega4.c04.lem.16-8-9-1}
\expandafter\def\csname sga@target@IV.16.8.9.2\endcsname{ega4.c04.lem.16-8-9-2}
\expandafter\def\csname sga@target@IV.16.8.9.3\endcsname{ega4.c04.lem.16-8-9-3}
\expandafter\def\csname sga@target@IV.16.8.9.4\endcsname{ega4.c04.eq.16-8-9-4}
\expandafter\def\csname sga@target@IV.16.9\endcsname{ega4.c04.subsec.16-9}
\expandafter\def\csname sga@target@IV.16.9.1\endcsname{ega4.c04.def.16-9-1}
\expandafter\def\csname sga@target@IV.16.9.10\endcsname{ega4.c04.prop.16-9-10}
\expandafter\def\csname sga@target@IV.16.9.11\endcsname{ega4.c04.cor.16-9-11}
\expandafter\def\csname sga@target@IV.16.9.12\endcsname{ega4.c04.statement.16-9-12}
\expandafter\def\csname sga@target@IV.16.9.13\endcsname{ega4.c04.prop.16-9-13}
\expandafter\def\csname sga@target@IV.16.9.13.1\endcsname{ega4.c04.prop.16-9-13-1}
\expandafter\def\csname sga@target@IV.16.9.13.2\endcsname{ega4.c04.lem.16-9-13-2}
\expandafter\def\csname sga@target@IV.16.9.13.3\endcsname{ega4.c04.lem.16-9-13-3}
\expandafter\def\csname sga@target@IV.16.9.13.4\endcsname{ega4.c04.eq.16-9-13-4}
\expandafter\def\csname sga@target@IV.16.9.13.5\endcsname{ega4.c04.cor.16-9-13-5}
\expandafter\def\csname sga@target@IV.16.9.13.6\endcsname{ega4.c04.cor.16-9-13-6}
\expandafter\def\csname sga@target@IV.16.9.13.7\endcsname{ega4.c04.eq.16-9-13-7}
\expandafter\def\csname sga@target@IV.16.9.2\endcsname{ega4.c04.def.16-9-2}
\expandafter\def\csname sga@target@IV.16.9.3\endcsname{ega4.c04.prop.16-9-3}
\expandafter\def\csname sga@target@IV.16.9.4\endcsname{ega4.c04.cor.16-9-4}
\expandafter\def\csname sga@target@IV.16.9.4.1\endcsname{ega4.c04.eq.16-9-4-1}
\expandafter\def\csname sga@target@IV.16.9.5\endcsname{ega4.c04.cor.16-9-5}
\expandafter\def\csname sga@target@IV.16.9.6\endcsname{ega4.c04.statement.16-9-6}
\expandafter\def\csname sga@target@IV.16.9.7\endcsname{ega4.c04.statement.16-9-7}
\expandafter\def\csname sga@target@IV.16.9.8\endcsname{ega4.c04.prop.16-9-8}
\expandafter\def\csname sga@target@IV.16.9.9\endcsname{ega4.c04.cor.16-9-9}
\expandafter\def\csname sga@target@IV.17.1.1\endcsname{ega4.c04.def.17-1-1}
\expandafter\def\csname sga@target@IV.17.1.1.1\endcsname{ega4.c04.def.17-1-1-1}
\expandafter\def\csname sga@target@IV.17.1.2\endcsname{ega4.c04.rem.17-1-2}
\expandafter\def\csname sga@target@IV.17.1.2.1\endcsname{ega4.c04.eq.17-1-2-1}
\expandafter\def\csname sga@target@IV.17.1.3\endcsname{ega4.c04.prop.17-1-3}
\expandafter\def\csname sga@target@IV.17.1.4\endcsname{ega4.c04.prop.17-1-4}
\expandafter\def\csname sga@target@IV.17.1.5\endcsname{ega4.c04.cor.17-1-5}
\expandafter\def\csname sga@target@IV.17.1.6\endcsname{ega4.c04.prop.17-1-6}
\expandafter\def\csname sga@target@IV.17.10.1\endcsname{ega4.c04.def.17-10-1}
\expandafter\def\csname sga@target@IV.17.10.2\endcsname{ega4.c04.prop.17-10-2}
\expandafter\def\csname sga@target@IV.17.10.3\endcsname{ega4.c04.cor.17-10-3}
\expandafter\def\csname sga@target@IV.17.10.3.1\endcsname{ega4.c04.cor.17-10-3-1}
\expandafter\def\csname sga@target@IV.17.10.4\endcsname{ega4.c04.cor.17-10-4}
\expandafter\def\csname sga@target@IV.17.10.4.1\endcsname{ega4.c04.eq.17-10-4-1}
\expandafter\def\csname sga@target@IV.17.11.1\endcsname{ega4.c04.thm.17-11-1}
\expandafter\def\csname sga@target@IV.17.11.1.1\endcsname{ega4.c04.eq.17-11-1-1}
\expandafter\def\csname sga@target@IV.17.11.1.2\endcsname{ega4.c04.eq.17-11-1-2}
\expandafter\def\csname sga@target@IV.17.11.2\endcsname{ega4.c04.cor.17-11-2}
\expandafter\def\csname sga@target@IV.17.11.3\endcsname{ega4.c04.cor.17-11-3}
\expandafter\def\csname sga@target@IV.17.11.4\endcsname{ega4.c04.cor.17-11-4}
\expandafter\def\csname sga@target@IV.17.11.5\endcsname{ega4.c04.prop.17-11-5}
\expandafter\def\csname sga@target@IV.17.12.1\endcsname{ega4.c04.thm.17-12-1}
\expandafter\def\csname sga@target@IV.17.12.2\endcsname{ega4.c04.cor.17-12-2}
\expandafter\def\csname sga@target@IV.17.12.3\endcsname{ega4.c04.cor.17-12-3}
\expandafter\def\csname sga@target@IV.17.12.4\endcsname{ega4.c04.prop.17-12-4}
\expandafter\def\csname sga@target@IV.17.12.5\endcsname{ega4.c04.prop.17-12-5}
\expandafter\def\csname sga@target@IV.17.12.6\endcsname{ega4.c04.cor.17-12-6}
\expandafter\def\csname sga@target@IV.17.12.7\endcsname{ega4.c04.ex.17-12-7}
\expandafter\def\csname sga@target@IV.17.12.8\endcsname{ega4.c04.rem.17-12-8}
\expandafter\def\csname sga@target@IV.17.13\endcsname{ega4.c04.subsec.17-13.key-iv-17-13}
\expandafter\def\csname sga@target@IV.17.13.1\endcsname{ega4.c04.statement.17-13-1}
\expandafter\def\csname sga@target@IV.17.13.1.1\endcsname{ega4.c04.statement.17-13-1-1}
\expandafter\def\csname sga@target@IV.17.13.1.2\endcsname{ega4.c04.statement.17-13-1-2}
\expandafter\def\csname sga@target@IV.17.13.1.3\endcsname{ega4.c04.eq.17-13-1-3}
\expandafter\def\csname sga@target@IV.17.13.10\endcsname{ega4.c04.cor.17-13-10}
\expandafter\def\csname sga@target@IV.17.13.11\endcsname{ega4.c04.statement.17-13-11}
\expandafter\def\csname sga@target@IV.17.13.11.1\endcsname{ega4.c04.statement.17-13-11-1}
\expandafter\def\csname sga@target@IV.17.13.12\endcsname{ega4.c04.cor.17-13-12}
\expandafter\def\csname sga@target@IV.17.13.12.1\endcsname{ega4.c04.cor.17-13-12-1}
\expandafter\def\csname sga@target@IV.17.13.12.2\endcsname{ega4.c04.eq.17-13-12-2}
\expandafter\def\csname sga@target@IV.17.13.2\endcsname{ega4.c04.prop.17-13-2}
\expandafter\def\csname sga@target@IV.17.13.2.1\endcsname{ega4.c04.eq.17-13-2-1}
\expandafter\def\csname sga@target@IV.17.13.2.2\endcsname{ega4.c04.eq.17-13-2-2}
\expandafter\def\csname sga@target@IV.17.13.3\endcsname{ega4.c04.def.17-13-3}
\expandafter\def\csname sga@target@IV.17.13.4\endcsname{ega4.c04.statement.17-13-4}
\expandafter\def\csname sga@target@IV.17.13.5\endcsname{ega4.c04.statement.17-13-5}
\expandafter\def\csname sga@target@IV.17.13.5.1\endcsname{ega4.c04.statement.17-13-5-1}
\expandafter\def\csname sga@target@IV.17.13.5.2\endcsname{ega4.c04.eq.17-13-5-2}
\expandafter\def\csname sga@target@IV.17.13.5.3\endcsname{ega4.c04.eq.17-13-5-3}
\expandafter\def\csname sga@target@IV.17.13.5.4\endcsname{ega4.c04.eq.17-13-5-4}
\expandafter\def\csname sga@target@IV.17.13.6\endcsname{ega4.c04.cor.17-13-6}
\expandafter\def\csname sga@target@IV.17.13.6.1\endcsname{ega4.c04.eq.17-13-6-1}
\expandafter\def\csname sga@target@IV.17.13.6.2\endcsname{ega4.c04.eq.17-13-6-2}
\expandafter\def\csname sga@target@IV.17.13.7\endcsname{ega4.c04.statement.17-13-7}
\expandafter\def\csname sga@target@IV.17.13.7.1\endcsname{ega4.c04.eq.17-13-7-1}
\expandafter\def\csname sga@target@IV.17.13.8\endcsname{ega4.c04.prop.17-13-8}
\expandafter\def\csname sga@target@IV.17.13.8.1\endcsname{ega4.c04.eq.17-13-8-1}
\expandafter\def\csname sga@target@IV.17.13.9\endcsname{ega4.c04.statement.17-13-9}
\expandafter\def\csname sga@target@IV.17.13.9.1\endcsname{ega4.c04.eq.17-13-9-1}
\expandafter\def\csname sga@target@IV.17.13.9.2\endcsname{ega4.c04.eq.17-13-9-2}
\expandafter\def\csname sga@target@IV.17.14.1\endcsname{ega4.c04.prop.17-14-1}
\expandafter\def\csname sga@target@IV.17.14.2\endcsname{ega4.c04.prop.17-14-2}
\expandafter\def\csname sga@target@IV.17.14.3\endcsname{ega4.c04.rem.17-14-3}
\expandafter\def\csname sga@target@IV.17.15.1\endcsname{ega4.c04.prop.17-15-1}
\expandafter\def\csname sga@target@IV.17.15.10\endcsname{ega4.c04.statement.17-15-10}
\expandafter\def\csname sga@target@IV.17.15.11\endcsname{ega4.c04.prop.17-15-11}
\expandafter\def\csname sga@target@IV.17.15.12\endcsname{ega4.c04.prop.17-15-12}
\expandafter\def\csname sga@target@IV.17.15.13\endcsname{ega4.c04.cor.17-15-13}
\expandafter\def\csname sga@target@IV.17.15.14\endcsname{ega4.c04.prop.17-15-14}
\expandafter\def\csname sga@target@IV.17.15.14.1\endcsname{ega4.c04.lem.17-15-14-1}
\expandafter\def\csname sga@target@IV.17.15.14.2\endcsname{ega4.c04.lem.17-15-14-2}
\expandafter\def\csname sga@target@IV.17.15.15\endcsname{ega4.c04.prop.17-15-15}
\expandafter\def\csname sga@target@IV.17.15.2\endcsname{ega4.c04.cor.17-15-2}
\expandafter\def\csname sga@target@IV.17.15.3\endcsname{ega4.c04.prop.17-15-3}
\expandafter\def\csname sga@target@IV.17.15.4\endcsname{ega4.c04.cor.17-15-4}
\expandafter\def\csname sga@target@IV.17.15.5\endcsname{ega4.c04.prop.17-15-5}
\expandafter\def\csname sga@target@IV.17.15.6\endcsname{ega4.c04.cor.17-15-6}
\expandafter\def\csname sga@target@IV.17.15.7\endcsname{ega4.c04.cor.17-15-7}
\expandafter\def\csname sga@target@IV.17.15.8\endcsname{ega4.c04.prop.17-15-8}
\expandafter\def\csname sga@target@IV.17.15.8.1\endcsname{ega4.c04.eq.17-15-8-1}
\expandafter\def\csname sga@target@IV.17.15.8.2\endcsname{ega4.c04.eq.17-15-8-2}
\expandafter\def\csname sga@target@IV.17.15.9\endcsname{ega4.c04.cor.17-15-9}
\expandafter\def\csname sga@target@IV.17.16\endcsname{ega4.c04.subsec.17-16.key-iv-17-16}
\expandafter\def\csname sga@target@IV.17.16.1\endcsname{ega4.c04.prop.17-16-1}
\expandafter\def\csname sga@target@IV.17.16.2\endcsname{ega4.c04.cor.17-16-2}
\expandafter\def\csname sga@target@IV.17.16.3\endcsname{ega4.c04.cor.17-16-3}
\expandafter\def\csname sga@target@IV.17.16.4\endcsname{ega4.c04.prop.17-16-4}
\expandafter\def\csname sga@target@IV.17.16.5\endcsname{ega4.c04.cor.17-16-5}
\expandafter\def\csname sga@target@IV.17.16.6\endcsname{ega4.c04.prop.17-16-6}
\expandafter\def\csname sga@target@IV.17.2.1\endcsname{ega4.c04.prop.17-2-1}
\expandafter\def\csname sga@target@IV.17.2.2\endcsname{ega4.c04.cor.17-2-2}
\expandafter\def\csname sga@target@IV.17.2.3\endcsname{ega4.c04.prop.17-2-3}
\expandafter\def\csname sga@target@IV.17.2.3.1\endcsname{ega4.c04.eq.17-2-3-1}
\expandafter\def\csname sga@target@IV.17.2.4\endcsname{ega4.c04.cor.17-2-4}
\expandafter\def\csname sga@target@IV.17.2.5\endcsname{ega4.c04.prop.17-2-5}
\expandafter\def\csname sga@target@IV.17.2.5.1\endcsname{ega4.c04.prop.17-2-5-1}
\expandafter\def\csname sga@target@IV.17.2.6\endcsname{ega4.c04.prop.17-2-6}
\expandafter\def\csname sga@target@IV.17.3.1\endcsname{ega4.c04.def.17-3-1}
\expandafter\def\csname sga@target@IV.17.3.2\endcsname{ega4.c04.subsec.17-3-2}
\expandafter\def\csname sga@target@IV.17.3.3\endcsname{ega4.c04.prop.17-3-3}
\expandafter\def\csname sga@target@IV.17.3.4\endcsname{ega4.c04.prop.17-3-4}
\expandafter\def\csname sga@target@IV.17.3.5\endcsname{ega4.c04.cor.17-3-5}
\expandafter\def\csname sga@target@IV.17.3.6\endcsname{ega4.c04.prop.17-3-6}
\expandafter\def\csname sga@target@IV.17.3.7\endcsname{ega4.c04.def.17-3-7}
\expandafter\def\csname sga@target@IV.17.3.8\endcsname{ega4.c04.prop.17-3-8}
\expandafter\def\csname sga@target@IV.17.3.9\endcsname{ega4.c04.cor.17-3-9}
\expandafter\def\csname sga@target@IV.17.4.1\endcsname{ega4.c04.thm.17-4-1}
\expandafter\def\csname sga@target@IV.17.4.1.1\endcsname{ega4.c04.eq.17-4-1-1}
\expandafter\def\csname sga@target@IV.17.4.1.2\endcsname{ega4.c04.rem.17-4-1-2}
\expandafter\def\csname sga@target@IV.17.4.10\endcsname{ega4.c04.rem.17-4-10}
\expandafter\def\csname sga@target@IV.17.4.2\endcsname{ega4.c04.cor.17-4-2}
\expandafter\def\csname sga@target@IV.17.4.3\endcsname{ega4.c04.cor.17-4-3}
\expandafter\def\csname sga@target@IV.17.4.4\endcsname{ega4.c04.prop.17-4-4}
\expandafter\def\csname sga@target@IV.17.4.5\endcsname{ega4.c04.statement.17-4-5}
\expandafter\def\csname sga@target@IV.17.4.6\endcsname{ega4.c04.prop.17-4-6}
\expandafter\def\csname sga@target@IV.17.4.7\endcsname{ega4.c04.cor.17-4-7}
\expandafter\def\csname sga@target@IV.17.4.8\endcsname{ega4.c04.cor.17-4-8}
\expandafter\def\csname sga@target@IV.17.4.9\endcsname{ega4.c04.prop.17-4-9}
\expandafter\def\csname sga@target@IV.17.5.1\endcsname{ega4.c04.thm.17-5-1}
\expandafter\def\csname sga@target@IV.17.5.1.1\endcsname{ega4.c04.eq.17-5-1-1}
\expandafter\def\csname sga@target@IV.17.5.1.2\endcsname{ega4.c04.eq.17-5-1-2}
\expandafter\def\csname sga@target@IV.17.5.2\endcsname{ega4.c04.cor.17-5-2}
\expandafter\def\csname sga@target@IV.17.5.3\endcsname{ega4.c04.prop.17-5-3}
\expandafter\def\csname sga@target@IV.17.5.4\endcsname{ega4.c04.rem.17-5-4}
\expandafter\def\csname sga@target@IV.17.5.5\endcsname{ega4.c04.prop.17-5-5}
\expandafter\def\csname sga@target@IV.17.5.6\endcsname{ega4.c04.cor.17-5-6}
\expandafter\def\csname sga@target@IV.17.5.7\endcsname{ega4.c04.prop.17-5-7}
\expandafter\def\csname sga@target@IV.17.5.8\endcsname{ega4.c04.prop.17-5-8}
\expandafter\def\csname sga@target@IV.17.6.1\endcsname{ega4.c04.thm.17-6-1}
\expandafter\def\csname sga@target@IV.17.6.2\endcsname{ega4.c04.cor.17-6-2}
\expandafter\def\csname sga@target@IV.17.6.3\endcsname{ega4.c04.prop.17-6-3}
\expandafter\def\csname sga@target@IV.17.6.4\endcsname{ega4.c04.prop.17-6-4}
\expandafter\def\csname sga@target@IV.17.7\endcsname{ega4.c04.subsec.17-7.key-iv-17-7}
\expandafter\def\csname sga@target@IV.17.7.1\endcsname{ega4.c04.prop.17-7-1}
\expandafter\def\csname sga@target@IV.17.7.10\endcsname{ega4.c04.prop.17-7-10}
\expandafter\def\csname sga@target@IV.17.7.11\endcsname{ega4.c04.prop.17-7-11}
\expandafter\def\csname sga@target@IV.17.7.11.1\endcsname{ega4.c04.lem.17-7-11-1}
\expandafter\def\csname sga@target@IV.17.7.2\endcsname{ega4.c04.cor.17-7-2}
\expandafter\def\csname sga@target@IV.17.7.3\endcsname{ega4.c04.cor.17-7-3}
\expandafter\def\csname sga@target@IV.17.7.4\endcsname{ega4.c04.prop.17-7-4}
\expandafter\def\csname sga@target@IV.17.7.5\endcsname{ega4.c04.lem.17-7-5}
\expandafter\def\csname sga@target@IV.17.7.6\endcsname{ega4.c04.rem.17-7-6}
\expandafter\def\csname sga@target@IV.17.7.7\endcsname{ega4.c04.prop.17-7-7}
\expandafter\def\csname sga@target@IV.17.7.8\endcsname{ega4.c04.prop.17-7-8}
\expandafter\def\csname sga@target@IV.17.7.9\endcsname{ega4.c04.cor.17-7-9}
\expandafter\def\csname sga@target@IV.17.8\endcsname{ega4.c04.subsec.17-8.key-iv-17-8}
\expandafter\def\csname sga@target@IV.17.8.1\endcsname{ega4.c04.prop.17-8-1}
\expandafter\def\csname sga@target@IV.17.8.2\endcsname{ega4.c04.prop.17-8-2}
\expandafter\def\csname sga@target@IV.17.8.3\endcsname{ega4.c04.rem.17-8-3}
\expandafter\def\csname sga@target@IV.17.9\endcsname{ega4.c04.subsec.17-9.key-iv-17-9}
\expandafter\def\csname sga@target@IV.17.9.1\endcsname{ega4.c04.thm.17-9-1}
\expandafter\def\csname sga@target@IV.17.9.2\endcsname{ega4.c04.cor.17-9-2}
\expandafter\def\csname sga@target@IV.17.9.3\endcsname{ega4.c04.cor.17-9-3}
\expandafter\def\csname sga@target@IV.17.9.4\endcsname{ega4.c04.cor.17-9-4}
\expandafter\def\csname sga@target@IV.17.9.5\endcsname{ega4.c04.cor.17-9-5}
\expandafter\def\csname sga@target@IV.17.9.6\endcsname{ega4.c04.prop.17-9-6}
\expandafter\def\csname sga@target@IV.18\endcsname{ega4.c04.sec.18.key-iv-18}
\expandafter\def\csname sga@target@IV.18.1.1\endcsname{ega4.c04.prop.18-1-1}
\expandafter\def\csname sga@target@IV.18.1.1.1\endcsname{ega4.c04.eq.18-1-1-1}
\expandafter\def\csname sga@target@IV.18.1.2\endcsname{ega4.c04.thm.18-1-2}
\expandafter\def\csname sga@target@IV.18.1.3\endcsname{ega4.c04.cor.18-1-3}
\expandafter\def\csname sga@target@IV.18.10.1\endcsname{ega4.c04.thm.18-10-1}
\expandafter\def\csname sga@target@IV.18.10.10\endcsname{ega4.c04.statement.18-10-10}
\expandafter\def\csname sga@target@IV.18.10.11\endcsname{ega4.c04.rem.18-10-11}
\expandafter\def\csname sga@target@IV.18.10.12\endcsname{ega4.c04.cor.18-10-12}
\expandafter\def\csname sga@target@IV.18.10.13\endcsname{ega4.c04.prop.18-10-13}
\expandafter\def\csname sga@target@IV.18.10.14\endcsname{ega4.c04.cor.18-10-14}
\expandafter\def\csname sga@target@IV.18.10.15\endcsname{ega4.c04.statement.18-10-15}
\expandafter\def\csname sga@target@IV.18.10.16\endcsname{ega4.c04.thm.18-10-16}
\expandafter\def\csname sga@target@IV.18.10.17\endcsname{ega4.c04.rem.18-10-17}
\expandafter\def\csname sga@target@IV.18.10.17.1\endcsname{ega4.c04.statement.18-10-17-1}
\expandafter\def\csname sga@target@IV.18.10.17.2\endcsname{ega4.c04.statement.18-10-17-2}
\expandafter\def\csname sga@target@IV.18.10.18\endcsname{ega4.c04.lem.18-10-18}
\expandafter\def\csname sga@target@IV.18.10.19\endcsname{ega4.c04.prop.18-10-19}
\expandafter\def\csname sga@target@IV.18.10.2\endcsname{ega4.c04.rem.18-10-2}
\expandafter\def\csname sga@target@IV.18.10.20\endcsname{ega4.c04.rem.18-10-20}
\expandafter\def\csname sga@target@IV.18.10.21\endcsname{ega4.c04.prop.18-10-21}
\expandafter\def\csname sga@target@IV.18.10.3\endcsname{ega4.c04.cor.18-10-3}
\expandafter\def\csname sga@target@IV.18.10.3.1\endcsname{ega4.c04.rem.18-10-3-1}
\expandafter\def\csname sga@target@IV.18.10.4\endcsname{ega4.c04.cor.18-10-4}
\expandafter\def\csname sga@target@IV.18.10.5\endcsname{ega4.c04.cor.18-10-5}
\expandafter\def\csname sga@target@IV.18.10.6\endcsname{ega4.c04.lem.18-10-6}
\expandafter\def\csname sga@target@IV.18.10.7\endcsname{ega4.c04.prop.18-10-7}
\expandafter\def\csname sga@target@IV.18.10.8\endcsname{ega4.c04.prop.18-10-8}
\expandafter\def\csname sga@target@IV.18.10.9\endcsname{ega4.c04.cor.18-10-9}
\expandafter\def\csname sga@target@IV.18.11.1\endcsname{ega4.c04.lem.18-11-1}
\expandafter\def\csname sga@target@IV.18.11.10\endcsname{ega4.c04.thm.18-11-10}
\expandafter\def\csname sga@target@IV.18.11.11\endcsname{ega4.c04.rem.18-11-11}
\expandafter\def\csname sga@target@IV.18.11.12\endcsname{ega4.c04.cor.18-11-12}
\expandafter\def\csname sga@target@IV.18.11.12.1\endcsname{ega4.c04.lem.18-11-12-1}
\expandafter\def\csname sga@target@IV.18.11.2\endcsname{ega4.c04.lem.18-11-2}
\expandafter\def\csname sga@target@IV.18.11.2.1\endcsname{ega4.c04.eq.18-11-2-1}
\expandafter\def\csname sga@target@IV.18.11.2.2\endcsname{ega4.c04.eq.18-11-2-2}
\expandafter\def\csname sga@target@IV.18.11.3\endcsname{ega4.c04.lem.18-11-3}
\expandafter\def\csname sga@target@IV.18.11.3.1\endcsname{ega4.c04.eq.18-11-3-1}
\expandafter\def\csname sga@target@IV.18.11.3.2\endcsname{ega4.c04.eq.18-11-3-2}
\expandafter\def\csname sga@target@IV.18.11.3.3\endcsname{ega4.c04.lem.18-11-3-3}
\expandafter\def\csname sga@target@IV.18.11.3.4\endcsname{ega4.c04.lem.18-11-3-4}
\expandafter\def\csname sga@target@IV.18.11.3.5\endcsname{ega4.c04.lem.18-11-3-5}
\expandafter\def\csname sga@target@IV.18.11.3.6\endcsname{ega4.c04.eq.18-11-3-6}
\expandafter\def\csname sga@target@IV.18.11.4\endcsname{ega4.c04.lem.18-11-4}
\expandafter\def\csname sga@target@IV.18.11.5\endcsname{ega4.c04.prop.18-11-5}
\expandafter\def\csname sga@target@IV.18.11.5.1\endcsname{ega4.c04.lem.18-11-5-1}
\expandafter\def\csname sga@target@IV.18.11.5.2\endcsname{ega4.c04.lem.18-11-5-2}
\expandafter\def\csname sga@target@IV.18.11.6\endcsname{ega4.c04.rem.18-11-6}
\expandafter\def\csname sga@target@IV.18.11.7\endcsname{ega4.c04.cor.18-11-7}
\expandafter\def\csname sga@target@IV.18.11.8\endcsname{ega4.c04.cor.18-11-8}
\expandafter\def\csname sga@target@IV.18.11.9\endcsname{ega4.c04.cor.18-11-9}
\expandafter\def\csname sga@target@IV.18.12.1\endcsname{ega4.c04.thm.18-12-1}
\expandafter\def\csname sga@target@IV.18.12.10\endcsname{ega4.c04.cor.18-12-10}
\expandafter\def\csname sga@target@IV.18.12.11\endcsname{ega4.c04.cor.18-12-11}
\expandafter\def\csname sga@target@IV.18.12.12\endcsname{ega4.c04.prop.18-12-12}
\expandafter\def\csname sga@target@IV.18.12.13\endcsname{ega4.c04.cor.18-12-13}
\expandafter\def\csname sga@target@IV.18.12.14\endcsname{ega4.c04.rem.18-12-14}
\expandafter\def\csname sga@target@IV.18.12.15\endcsname{ega4.c04.prop.18-12-15}
\expandafter\def\csname sga@target@IV.18.12.16\endcsname{ega4.c04.prop.18-12-16}
\expandafter\def\csname sga@target@IV.18.12.17\endcsname{ega4.c04.cor.18-12-17}
\expandafter\def\csname sga@target@IV.18.12.2\endcsname{ega4.c04.rem.18-12-2}
\expandafter\def\csname sga@target@IV.18.12.3\endcsname{ega4.c04.cor.18-12-3}
\expandafter\def\csname sga@target@IV.18.12.4\endcsname{ega4.c04.cor.18-12-4}
\expandafter\def\csname sga@target@IV.18.12.5\endcsname{ega4.c04.rem.18-12-5}
\expandafter\def\csname sga@target@IV.18.12.6\endcsname{ega4.c04.cor.18-12-6}
\expandafter\def\csname sga@target@IV.18.12.7\endcsname{ega4.c04.prop.18-12-7}
\expandafter\def\csname sga@target@IV.18.12.7.1\endcsname{ega4.c04.lem.18-12-7-1}
\expandafter\def\csname sga@target@IV.18.12.8\endcsname{ega4.c04.prop.18-12-8}
\expandafter\def\csname sga@target@IV.18.12.9\endcsname{ega4.c04.rem.18-12-9}
\expandafter\def\csname sga@target@IV.18.2.1\endcsname{ega4.c04.statement.18-2-1}
\expandafter\def\csname sga@target@IV.18.2.1.1\endcsname{ega4.c04.statement.18-2-1-1}
\expandafter\def\csname sga@target@IV.18.2.1.2\endcsname{ega4.c04.eq.18-2-1-2}
\expandafter\def\csname sga@target@IV.18.2.10\endcsname{ega4.c04.statement.18-2-10}
\expandafter\def\csname sga@target@IV.18.2.2\endcsname{ega4.c04.statement.18-2-2}
\expandafter\def\csname sga@target@IV.18.2.2.1\endcsname{ega4.c04.eq.18-2-2-1}
\expandafter\def\csname sga@target@IV.18.2.2.2\endcsname{ega4.c04.eq.18-2-2-2}
\expandafter\def\csname sga@target@IV.18.2.2.3\endcsname{ega4.c04.eq.18-2-2-3}
\expandafter\def\csname sga@target@IV.18.2.2.4\endcsname{ega4.c04.eq.18-2-2-4}
\expandafter\def\csname sga@target@IV.18.2.3\endcsname{ega4.c04.prop.18-2-3}
\expandafter\def\csname sga@target@IV.18.2.4\endcsname{ega4.c04.cor.18-2-4}
\expandafter\def\csname sga@target@IV.18.2.5\endcsname{ega4.c04.cor.18-2-5}
\expandafter\def\csname sga@target@IV.18.2.6\endcsname{ega4.c04.cor.18-2-6}
\expandafter\def\csname sga@target@IV.18.2.7\endcsname{ega4.c04.item.18-2-7}
\expandafter\def\csname sga@target@IV.18.2.7.1\endcsname{ega4.c04.eq.18-2-7-1}
\expandafter\def\csname sga@target@IV.18.2.7.2\endcsname{ega4.c04.eq.18-2-7-2}
\expandafter\def\csname sga@target@IV.18.2.8\endcsname{ega4.c04.prop.18-2-8}
\expandafter\def\csname sga@target@IV.18.2.9\endcsname{ega4.c04.cor.18-2-9}
\expandafter\def\csname sga@target@IV.18.3.1\endcsname{ega4.c04.prop.18-3-1}
\expandafter\def\csname sga@target@IV.18.3.2\endcsname{ega4.c04.prop.18-3-2}
\expandafter\def\csname sga@target@IV.18.3.2.1\endcsname{ega4.c04.lem.18-3-2-1}
\expandafter\def\csname sga@target@IV.18.3.2.2\endcsname{ega4.c04.eq.18-3-2-2}
\expandafter\def\csname sga@target@IV.18.3.3\endcsname{ega4.c04.statement.18-3-3}
\expandafter\def\csname sga@target@IV.18.3.4\endcsname{ega4.c04.thm.18-3-4}
\expandafter\def\csname sga@target@IV.18.4.1\endcsname{ega4.c04.lem.18-4-1}
\expandafter\def\csname sga@target@IV.18.4.10\endcsname{ega4.c04.prop.18-4-10}
\expandafter\def\csname sga@target@IV.18.4.11\endcsname{ega4.c04.statement.18-4-11}
\expandafter\def\csname sga@target@IV.18.4.12\endcsname{ega4.c04.cor.18-4-12}
\expandafter\def\csname sga@target@IV.18.4.13\endcsname{ega4.c04.cor.18-4-13}
\expandafter\def\csname sga@target@IV.18.4.14\endcsname{ega4.c04.cor.18-4-14}
\expandafter\def\csname sga@target@IV.18.4.2\endcsname{ega4.c04.prop.18-4-2}
\expandafter\def\csname sga@target@IV.18.4.3\endcsname{ega4.c04.cor.18-4-3}
\expandafter\def\csname sga@target@IV.18.4.4\endcsname{ega4.c04.lem.18-4-4}
\expandafter\def\csname sga@target@IV.18.4.5\endcsname{ega4.c04.prop.18-4-5}
\expandafter\def\csname sga@target@IV.18.4.6\endcsname{ega4.c04.thm.18-4-6}
\expandafter\def\csname sga@target@IV.18.4.6.1\endcsname{ega4.c04.lem.18-4-6-1}
\expandafter\def\csname sga@target@IV.18.4.7\endcsname{ega4.c04.cor.18-4-7}
\expandafter\def\csname sga@target@IV.18.4.8\endcsname{ega4.c04.cor.18-4-8}
\expandafter\def\csname sga@target@IV.18.4.9\endcsname{ega4.c04.prop.18-4-9}
\expandafter\def\csname sga@target@IV.18.5.1\endcsname{ega4.c04.statement.18-5-1}
\expandafter\def\csname sga@target@IV.18.5.1.1\endcsname{ega4.c04.eq.18-5-1-1}
\expandafter\def\csname sga@target@IV.18.5.10\endcsname{ega4.c04.prop.18-5-10}
\expandafter\def\csname sga@target@IV.18.5.11\endcsname{ega4.c04.thm.18-5-11}
\expandafter\def\csname sga@target@IV.18.5.12\endcsname{ega4.c04.cor.18-5-12}
\expandafter\def\csname sga@target@IV.18.5.13\endcsname{ega4.c04.rem.18-5-13}
\expandafter\def\csname sga@target@IV.18.5.13.1\endcsname{ega4.c04.lem.18-5-13-1}
\expandafter\def\csname sga@target@IV.18.5.14\endcsname{ega4.c04.prop.18-5-14}
\expandafter\def\csname sga@target@IV.18.5.15\endcsname{ega4.c04.prop.18-5-15}
\expandafter\def\csname sga@target@IV.18.5.16\endcsname{ega4.c04.item.18-5-16}
\expandafter\def\csname sga@target@IV.18.5.17\endcsname{ega4.c04.thm.18-5-17}
\expandafter\def\csname sga@target@IV.18.5.18\endcsname{ega4.c04.rem.18-5-18}
\expandafter\def\csname sga@target@IV.18.5.19\endcsname{ega4.c04.prop.18-5-19}
\expandafter\def\csname sga@target@IV.18.5.2\endcsname{ega4.c04.statement.18-5-2}
\expandafter\def\csname sga@target@IV.18.5.3\endcsname{ega4.c04.lem.18-5-3}
\expandafter\def\csname sga@target@IV.18.5.3.1\endcsname{ega4.c04.lem.18-5-3-1}
\expandafter\def\csname sga@target@IV.18.5.4\endcsname{ega4.c04.prop.18-5-4}
\expandafter\def\csname sga@target@IV.18.5.4.1\endcsname{ega4.c04.eq.18-5-4-1}
\expandafter\def\csname sga@target@IV.18.5.4.2\endcsname{ega4.c04.eq.18-5-4-2}
\expandafter\def\csname sga@target@IV.18.5.4.3\endcsname{ega4.c04.lem.18-5-4-3}
\expandafter\def\csname sga@target@IV.18.5.4.4\endcsname{ega4.c04.rem.18-5-4-4}
\expandafter\def\csname sga@target@IV.18.5.5\endcsname{ega4.c04.def.18-5-5}
\expandafter\def\csname sga@target@IV.18.5.6\endcsname{ega4.c04.prop.18-5-6}
\expandafter\def\csname sga@target@IV.18.5.7\endcsname{ega4.c04.rem.18-5-7}
\expandafter\def\csname sga@target@IV.18.5.8\endcsname{ega4.c04.def.18-5-8}
\expandafter\def\csname sga@target@IV.18.5.9\endcsname{ega4.c04.prop.18-5-9}
\expandafter\def\csname sga@target@IV.18.6.1\endcsname{ega4.c04.statement.18-6-1}
\expandafter\def\csname sga@target@IV.18.6.10\endcsname{ega4.c04.cor.18-6-10}
\expandafter\def\csname sga@target@IV.18.6.11\endcsname{ega4.c04.cor.18-6-11}
\expandafter\def\csname sga@target@IV.18.6.12\endcsname{ega4.c04.prop.18-6-12}
\expandafter\def\csname sga@target@IV.18.6.13\endcsname{ega4.c04.cor.18-6-13}
\expandafter\def\csname sga@target@IV.18.6.14\endcsname{ega4.c04.prop.18-6-14}
\expandafter\def\csname sga@target@IV.18.6.15\endcsname{ega4.c04.cor.18-6-15}
\expandafter\def\csname sga@target@IV.18.6.2\endcsname{ega4.c04.lem.18-6-2}
\expandafter\def\csname sga@target@IV.18.6.3\endcsname{ega4.c04.lem.18-6-3}
\expandafter\def\csname sga@target@IV.18.6.4\endcsname{ega4.c04.statement.18-6-4}
\expandafter\def\csname sga@target@IV.18.6.5\endcsname{ega4.c04.def.18-6-5}
\expandafter\def\csname sga@target@IV.18.6.6\endcsname{ega4.c04.thm.18-6-6}
\expandafter\def\csname sga@target@IV.18.6.7\endcsname{ega4.c04.statement.18-6-7}
\expandafter\def\csname sga@target@IV.18.6.8\endcsname{ega4.c04.prop.18-6-8}
\expandafter\def\csname sga@target@IV.18.6.9\endcsname{ega4.c04.thm.18-6-9}
\expandafter\def\csname sga@target@IV.18.6.9.1\endcsname{ega4.c04.lem.18-6-9-1}
\expandafter\def\csname sga@target@IV.18.6.9.2\endcsname{ega4.c04.lem.18-6-9-2}
\expandafter\def\csname sga@target@IV.18.7.1\endcsname{ega4.c04.statement.18-7-1}
\expandafter\def\csname sga@target@IV.18.7.2\endcsname{ega4.c04.prop.18-7-2}
\expandafter\def\csname sga@target@IV.18.7.3\endcsname{ega4.c04.cor.18-7-3}
\expandafter\def\csname sga@target@IV.18.7.4\endcsname{ega4.c04.cor.18-7-4}
\expandafter\def\csname sga@target@IV.18.7.5\endcsname{ega4.c04.prop.18-7-5}
\expandafter\def\csname sga@target@IV.18.7.5.1\endcsname{ega4.c04.lem.18-7-5-1}
\expandafter\def\csname sga@target@IV.18.7.6\endcsname{ega4.c04.cor.18-7-6}
\expandafter\def\csname sga@target@IV.18.7.7\endcsname{ega4.c04.rem.18-7-7}
\expandafter\def\csname sga@target@IV.18.8.1\endcsname{ega4.c04.prop.18-8-1}
\expandafter\def\csname sga@target@IV.18.8.10\endcsname{ega4.c04.prop.18-8-10}
\expandafter\def\csname sga@target@IV.18.8.11\endcsname{ega4.c04.rem.18-8-11}
\expandafter\def\csname sga@target@IV.18.8.12\endcsname{ega4.c04.prop.18-8-12}
\expandafter\def\csname sga@target@IV.18.8.13\endcsname{ega4.c04.cor.18-8-13}
\expandafter\def\csname sga@target@IV.18.8.14\endcsname{ega4.c04.cor.18-8-14}
\expandafter\def\csname sga@target@IV.18.8.15\endcsname{ega4.c04.prop.18-8-15}
\expandafter\def\csname sga@target@IV.18.8.16\endcsname{ega4.c04.cor.18-8-16}
\expandafter\def\csname sga@target@IV.18.8.17\endcsname{ega4.c04.prop.18-8-17}
\expandafter\def\csname sga@target@IV.18.8.18\endcsname{ega4.c04.prop.18-8-18}
\expandafter\def\csname sga@target@IV.18.8.2\endcsname{ega4.c04.def.18-8-2}
\expandafter\def\csname sga@target@IV.18.8.3\endcsname{ega4.c04.rem.18-8-3}
\expandafter\def\csname sga@target@IV.18.8.4\endcsname{ega4.c04.lem.18-8-4}
\expandafter\def\csname sga@target@IV.18.8.5\endcsname{ega4.c04.lem.18-8-5}
\expandafter\def\csname sga@target@IV.18.8.6\endcsname{ega4.c04.statement.18-8-6}
\expandafter\def\csname sga@target@IV.18.8.6.1\endcsname{ega4.c04.eq.18-8-6-1}
\expandafter\def\csname sga@target@IV.18.8.6.2\endcsname{ega4.c04.eq.18-8-6-2}
\expandafter\def\csname sga@target@IV.18.8.7\endcsname{ega4.c04.def.18-8-7}
\expandafter\def\csname sga@target@IV.18.8.8\endcsname{ega4.c04.prop.18-8-8}
\expandafter\def\csname sga@target@IV.18.8.8.1\endcsname{ega4.c04.statement.18-8-8-1}
\expandafter\def\csname sga@target@IV.18.8.9\endcsname{ega4.c04.statement.18-8-9}
\expandafter\def\csname sga@target@IV.18.9.1\endcsname{ega4.c04.thm.18-9-1}
\expandafter\def\csname sga@target@IV.18.9.10\endcsname{ega4.c04.rem.18-9-10}
\expandafter\def\csname sga@target@IV.18.9.11\endcsname{ega4.c04.cor.18-9-11}
\expandafter\def\csname sga@target@IV.18.9.2\endcsname{ega4.c04.cor.18-9-2}
\expandafter\def\csname sga@target@IV.18.9.3\endcsname{ega4.c04.cor.18-9-3}
\expandafter\def\csname sga@target@IV.18.9.4\endcsname{ega4.c04.cor.18-9-4}
\expandafter\def\csname sga@target@IV.18.9.5\endcsname{ega4.c04.cor.18-9-5}
\expandafter\def\csname sga@target@IV.18.9.5.1\endcsname{ega4.c04.eq.18-9-5-1}
\expandafter\def\csname sga@target@IV.18.9.6\endcsname{ega4.c04.statement.18-9-6}
\expandafter\def\csname sga@target@IV.18.9.7\endcsname{ega4.c04.thm.18-9-7}
\expandafter\def\csname sga@target@IV.18.9.7.1\endcsname{ega4.c04.rem.18-9-7-1}
\expandafter\def\csname sga@target@IV.18.9.7.2\endcsname{ega4.c04.statement.18-9-7-2}
\expandafter\def\csname sga@target@IV.18.9.7.3\endcsname{ega4.c04.lem.18-9-7-3}
\expandafter\def\csname sga@target@IV.18.9.7.4\endcsname{ega4.c04.eq.18-9-7-4}
\expandafter\def\csname sga@target@IV.18.9.7.5\endcsname{ega4.c04.lem.18-9-7-5}
\expandafter\def\csname sga@target@IV.18.9.7.6\endcsname{ega4.c04.rem.18-9-7-6}
\expandafter\def\csname sga@target@IV.18.9.7.7\endcsname{ega4.c04.lem.18-9-7-7}
\expandafter\def\csname sga@target@IV.18.9.7.8\endcsname{ega4.c04.rem.18-9-7-8}
\expandafter\def\csname sga@target@IV.18.9.8\endcsname{ega4.c04.cor.18-9-8}
\expandafter\def\csname sga@target@IV.18.9.9\endcsname{ega4.c04.cor.18-9-9}
\expandafter\def\csname sga@target@IV.19.1.1\endcsname{ega4.c04.prop.19-1-1}
\expandafter\def\csname sga@target@IV.19.1.2\endcsname{ega4.c04.cor.19-1-2}
\expandafter\def\csname sga@target@IV.19.1.3\endcsname{ega4.c04.def.19-1-3}
\expandafter\def\csname sga@target@IV.19.1.4\endcsname{ega4.c04.prop.19-1-4}
\expandafter\def\csname sga@target@IV.19.1.4.1\endcsname{ega4.c04.eq.19-1-4-1}
\expandafter\def\csname sga@target@IV.19.1.5\endcsname{ega4.c04.prop.19-1-5}
\expandafter\def\csname sga@target@IV.19.1.5.1\endcsname{ega4.c04.item.19-1-5-1}
\expandafter\def\csname sga@target@IV.19.1.6\endcsname{ega4.c04.cor.19-1-6}
\expandafter\def\csname sga@target@IV.19.1.6.1\endcsname{ega4.c04.lem.19-1-6-1}
\expandafter\def\csname sga@target@IV.19.1.7\endcsname{ega4.c04.item.19-1-7}
\expandafter\def\csname sga@target@IV.19.1.8\endcsname{ega4.c04.statement.19-1-8}
\expandafter\def\csname sga@target@IV.19.2.1\endcsname{ega4.c04.def.19-2-1}
\expandafter\def\csname sga@target@IV.19.2.2\endcsname{ega4.c04.def.19-2-2}
\expandafter\def\csname sga@target@IV.19.2.3\endcsname{ega4.c04.prop.19-2-3}
\expandafter\def\csname sga@target@IV.19.2.4\endcsname{ega4.c04.prop.19-2-4}
\expandafter\def\csname sga@target@IV.19.2.5\endcsname{ega4.c04.item.19-2-5}
\expandafter\def\csname sga@target@IV.19.2.6\endcsname{ega4.c04.cor.19-2-6}
\expandafter\def\csname sga@target@IV.19.2.7\endcsname{ega4.c04.prop.19-2-7}
\expandafter\def\csname sga@target@IV.19.2.8\endcsname{ega4.c04.prop.19-2-8}
\expandafter\def\csname sga@target@IV.19.2.9\endcsname{ega4.c04.prop.19-2-9}
\expandafter\def\csname sga@target@IV.19.3.1\endcsname{ega4.c04.def.19-3-1}
\expandafter\def\csname sga@target@IV.19.3.10\endcsname{ega4.c04.cor.19-3-10}
\expandafter\def\csname sga@target@IV.19.3.2\endcsname{ega4.c04.prop.19-3-2}
\expandafter\def\csname sga@target@IV.19.3.2.1\endcsname{ega4.c04.lem.19-3-2-1}
\expandafter\def\csname sga@target@IV.19.3.3\endcsname{ega4.c04.cor.19-3-3}
\expandafter\def\csname sga@target@IV.19.3.4\endcsname{ega4.c04.cor.19-3-4}
\expandafter\def\csname sga@target@IV.19.3.5\endcsname{ega4.c04.cor.19-3-5}
\expandafter\def\csname sga@target@IV.19.3.6\endcsname{ega4.c04.def.19-3-6}
\expandafter\def\csname sga@target@IV.19.3.7\endcsname{ega4.c04.prop.19-3-7}
\expandafter\def\csname sga@target@IV.19.3.8\endcsname{ega4.c04.cor.19-3-8}
\expandafter\def\csname sga@target@IV.19.3.9\endcsname{ega4.c04.prop.19-3-9}
\expandafter\def\csname sga@target@IV.19.4.1\endcsname{ega4.c04.statement.19-4-1}
\expandafter\def\csname sga@target@IV.19.4.1.1\endcsname{ega4.c04.subsec.19-4-1-1}
\expandafter\def\csname sga@target@IV.19.4.10\endcsname{ega4.c04.rem.19-4-10}
\expandafter\def\csname sga@target@IV.19.4.11\endcsname{ega4.c04.prop.19-4-11}
\expandafter\def\csname sga@target@IV.19.4.11.1\endcsname{ega4.c04.eq.19-4-11-1}
\expandafter\def\csname sga@target@IV.19.4.2\endcsname{ega4.c04.lem.19-4-2}
\expandafter\def\csname sga@target@IV.19.4.3\endcsname{ega4.c04.cor.19-4-3}
\expandafter\def\csname sga@target@IV.19.4.4\endcsname{ega4.c04.prop.19-4-4}
\expandafter\def\csname sga@target@IV.19.4.5\endcsname{ega4.c04.cor.19-4-5}
\expandafter\def\csname sga@target@IV.19.4.6\endcsname{ega4.c04.prop.19-4-6}
\expandafter\def\csname sga@target@IV.19.4.7\endcsname{ega4.c04.cor.19-4-7}
\expandafter\def\csname sga@target@IV.19.4.8\endcsname{ega4.c04.prop.19-4-8}
\expandafter\def\csname sga@target@IV.19.4.9\endcsname{ega4.c04.cor.19-4-9}
\expandafter\def\csname sga@target@IV.19.5.1\endcsname{ega4.c04.thm.19-5-1}
\expandafter\def\csname sga@target@IV.19.5.1.1\endcsname{ega4.c04.eq.19-5-1-1}
\expandafter\def\csname sga@target@IV.19.5.2\endcsname{ega4.c04.cor.19-5-2}
\expandafter\def\csname sga@target@IV.19.5.3\endcsname{ega4.c04.prop.19-5-3}
\expandafter\def\csname sga@target@IV.19.5.4\endcsname{ega4.c04.statement.19-5-4}
\expandafter\def\csname sga@target@IV.19.5.5\endcsname{ega4.c04.thm.19-5-5}
\expandafter\def\csname sga@target@IV.19.5.5.1\endcsname{ega4.c04.statement.19-5-5-1}
\expandafter\def\csname sga@target@IV.19.5.5.2\endcsname{ega4.c04.statement.19-5-5-2}
\expandafter\def\csname sga@target@IV.19.5.5.3\endcsname{ega4.c04.statement.19-5-5-3}
\expandafter\def\csname sga@target@IV.19.5.5.4\endcsname{ega4.c04.statement.19-5-5-4}
\expandafter\def\csname sga@target@IV.19.5.5.5\endcsname{ega4.c04.statement.19-5-5-5}
\expandafter\def\csname sga@target@IV.19.5.6\endcsname{ega4.c04.statement.19-5-6}
\expandafter\def\csname sga@target@IV.19.6.1\endcsname{ega4.c04.statement.19-6-1}
\expandafter\def\csname sga@target@IV.19.6.1.1\endcsname{ega4.c04.eq.19-6-1-1}
\expandafter\def\csname sga@target@IV.19.6.2\endcsname{ega4.c04.statement.19-6-2}
\expandafter\def\csname sga@target@IV.19.6.2.1\endcsname{ega4.c04.statement.19-6-2-1}
\expandafter\def\csname sga@target@IV.19.6.3\endcsname{ega4.c04.prop.19-6-3}
\expandafter\def\csname sga@target@IV.19.6.4\endcsname{ega4.c04.cor.19-6-4}
\expandafter\def\csname sga@target@IV.19.7\endcsname{ega4.c04.subsec.19-7.key-iv-19-7}
\expandafter\def\csname sga@target@IV.19.7.1\endcsname{ega4.c04.thm.19-7-1}
\expandafter\def\csname sga@target@IV.19.7.1.1\endcsname{ega4.c04.statement.19-7-1-1}
\expandafter\def\csname sga@target@IV.19.7.1.10\endcsname{ega4.c04.statement.19-7-1-10}
\expandafter\def\csname sga@target@IV.19.7.1.11\endcsname{ega4.c04.statement.19-7-1-11}
\expandafter\def\csname sga@target@IV.19.7.1.12\endcsname{ega4.c04.statement.19-7-1-12}
\expandafter\def\csname sga@target@IV.19.7.1.13\endcsname{ega4.c04.statement.19-7-1-13}
\expandafter\def\csname sga@target@IV.19.7.1.14\endcsname{ega4.c04.statement.19-7-1-14}
\expandafter\def\csname sga@target@IV.19.7.1.15\endcsname{ega4.c04.statement.19-7-1-15}
\expandafter\def\csname sga@target@IV.19.7.1.16\endcsname{ega4.c04.statement.19-7-1-16}
\expandafter\def\csname sga@target@IV.19.7.1.17\endcsname{ega4.c04.statement.19-7-1-17}
\expandafter\def\csname sga@target@IV.19.7.1.18\endcsname{ega4.c04.statement.19-7-1-18}
\expandafter\def\csname sga@target@IV.19.7.1.19\endcsname{ega4.c04.statement.19-7-1-19}
\expandafter\def\csname sga@target@IV.19.7.1.2\endcsname{ega4.c04.statement.19-7-1-2}
\expandafter\def\csname sga@target@IV.19.7.1.20\endcsname{ega4.c04.statement.19-7-1-20}
\expandafter\def\csname sga@target@IV.19.7.1.21\endcsname{ega4.c04.statement.19-7-1-21}
\expandafter\def\csname sga@target@IV.19.7.1.22\endcsname{ega4.c04.statement.19-7-1-22}
\expandafter\def\csname sga@target@IV.19.7.1.23\endcsname{ega4.c04.statement.19-7-1-23}
\expandafter\def\csname sga@target@IV.19.7.1.3\endcsname{ega4.c04.statement.19-7-1-3}
\expandafter\def\csname sga@target@IV.19.7.1.5\endcsname{ega4.c04.statement.19-7-1-5}
\expandafter\def\csname sga@target@IV.19.7.1.6\endcsname{ega4.c04.statement.19-7-1-6}
\expandafter\def\csname sga@target@IV.19.7.1.7\endcsname{ega4.c04.lem.19-7-1-7}
\expandafter\def\csname sga@target@IV.19.7.1.8\endcsname{ega4.c04.statement.19-7-1-8}
\expandafter\def\csname sga@target@IV.19.7.1.9\endcsname{ega4.c04.statement.19-7-1-9}
\expandafter\def\csname sga@target@IV.19.7.2\endcsname{ega4.c04.statement.19-7-2}
\expandafter\def\csname sga@target@IV.19.7.3\endcsname{ega4.c04.cor.19-7-3}
\expandafter\def\csname sga@target@IV.19.7.3.1\endcsname{ega4.c04.eq.19-7-3-1}
\expandafter\def\csname sga@target@IV.19.7.4\endcsname{ega4.c04.cor.19-7-4}
\expandafter\def\csname sga@target@IV.19.7.4.1\endcsname{ega4.c04.eq.19-7-4-1}
\expandafter\def\csname sga@target@IV.19.7.5\endcsname{ega4.c04.cor.19-7-5}
\expandafter\def\csname sga@target@IV.19.8.1\endcsname{ega4.c04.prop.19-8-1}
\expandafter\def\csname sga@target@IV.19.8.2\endcsname{ega4.c04.prop.19-8-2}
\expandafter\def\csname sga@target@IV.19.8.3\endcsname{ega4.c04.rem.19-8-3}
\expandafter\def\csname sga@target@IV.19.9.1\endcsname{ega4.c04.statement.19-9-1}
\expandafter\def\csname sga@target@IV.19.9.1.1\endcsname{ega4.c04.eq.19-9-1-1}
\expandafter\def\csname sga@target@IV.19.9.1.2\endcsname{ega4.c04.eq.19-9-1-2}
\expandafter\def\csname sga@target@IV.19.9.1.3\endcsname{ega4.c04.eq.19-9-1-3}
\expandafter\def\csname sga@target@IV.19.9.2\endcsname{ega4.c04.lem.19-9-2}
\expandafter\def\csname sga@target@IV.19.9.3\endcsname{ega4.c04.prop.19-9-3}
\expandafter\def\csname sga@target@IV.19.9.4\endcsname{ega4.c04.cor.19-9-4}
\expandafter\def\csname sga@target@IV.19.9.5\endcsname{ega4.c04.prop.19-9-5}
\expandafter\def\csname sga@target@IV.19.9.6\endcsname{ega4.c04.prop.19-9-6}
\expandafter\def\csname sga@target@IV.19.9.7\endcsname{ega4.c04.cor.19-9-7}
\expandafter\def\csname sga@target@IV.19.9.8\endcsname{ega4.c04.prop.19-9-8}
\expandafter\def\csname sga@target@IV.19.9.9\endcsname{ega4.c04.rem.19-9-9}
\expandafter\def\csname sga@target@IV.2.1.1\endcsname{ega4.c04.statement.2-1-1}
\expandafter\def\csname sga@target@IV.2.1.10\endcsname{ega4.c04.cor.2-1-10}
\expandafter\def\csname sga@target@IV.2.1.11\endcsname{ega4.c04.prop.2-1-11}
\expandafter\def\csname sga@target@IV.2.1.12\endcsname{ega4.c04.prop.2-1-12}
\expandafter\def\csname sga@target@IV.2.1.13\endcsname{ega4.c04.prop.2-1-13}
\expandafter\def\csname sga@target@IV.2.1.14\endcsname{ega4.c04.prop.2-1-14}
\expandafter\def\csname sga@target@IV.2.1.2\endcsname{ega4.c04.prop.2-1-2}
\expandafter\def\csname sga@target@IV.2.1.3\endcsname{ega4.c04.prop.2-1-3}
\expandafter\def\csname sga@target@IV.2.1.4\endcsname{ega4.c04.prop.2-1-4}
\expandafter\def\csname sga@target@IV.2.1.5\endcsname{ega4.c04.prop.2-1-5}
\expandafter\def\csname sga@target@IV.2.1.6\endcsname{ega4.c04.cor.2-1-6}
\expandafter\def\csname sga@target@IV.2.1.7\endcsname{ega4.c04.cor.2-1-7}
\expandafter\def\csname sga@target@IV.2.1.8\endcsname{ega4.c04.prop.2-1-8}
\expandafter\def\csname sga@target@IV.2.1.9\endcsname{ega4.c04.cor.2-1-9}
\expandafter\def\csname sga@target@IV.2.2.1\endcsname{ega4.c04.prop.2-2-1}
\expandafter\def\csname sga@target@IV.2.2.10\endcsname{ega4.c04.prop.2-2-10}
\expandafter\def\csname sga@target@IV.2.2.10.1\endcsname{ega4.c04.eq.2-2-10-1}
\expandafter\def\csname sga@target@IV.2.2.11\endcsname{ega4.c04.cor.2-2-11}
\expandafter\def\csname sga@target@IV.2.2.12\endcsname{ega4.c04.cor.2-2-12}
\expandafter\def\csname sga@target@IV.2.2.13\endcsname{ega4.c04.cor.2-2-13}
\expandafter\def\csname sga@target@IV.2.2.14\endcsname{ega4.c04.prop.2-2-14}
\expandafter\def\csname sga@target@IV.2.2.15\endcsname{ega4.c04.prop.2-2-15}
\expandafter\def\csname sga@target@IV.2.2.16\endcsname{ega4.c04.cor.2-2-16}
\expandafter\def\csname sga@target@IV.2.2.17\endcsname{ega4.c04.prop.2-2-17}
\expandafter\def\csname sga@target@IV.2.2.2\endcsname{ega4.c04.prop.2-2-2}
\expandafter\def\csname sga@target@IV.2.2.2.1\endcsname{ega4.c04.eq.2-2-2-1}
\expandafter\def\csname sga@target@IV.2.2.3\endcsname{ega4.c04.cor.2-2-3}
\expandafter\def\csname sga@target@IV.2.2.4\endcsname{ega4.c04.def.2-2-4}
\expandafter\def\csname sga@target@IV.2.2.5\endcsname{ega4.c04.cor.2-2-5}
\expandafter\def\csname sga@target@IV.2.2.6\endcsname{ega4.c04.cor.2-2-6}
\expandafter\def\csname sga@target@IV.2.2.7\endcsname{ega4.c04.prop.2-2-7}
\expandafter\def\csname sga@target@IV.2.2.8\endcsname{ega4.c04.cor.2-2-8}
\expandafter\def\csname sga@target@IV.2.2.8.1\endcsname{ega4.c04.eq.2-2-8-1}
\expandafter\def\csname sga@target@IV.2.2.9\endcsname{ega4.c04.rem.2-2-9}
\expandafter\def\csname sga@target@IV.2.3.1\endcsname{ega4.c04.lem.2-3-1}
\expandafter\def\csname sga@target@IV.2.3.1.1\endcsname{ega4.c04.eq.2-3-1-1}
\expandafter\def\csname sga@target@IV.2.3.10\endcsname{ega4.c04.thm.2-3-10}
\expandafter\def\csname sga@target@IV.2.3.11\endcsname{ega4.c04.cor.2-3-11}
\expandafter\def\csname sga@target@IV.2.3.12\endcsname{ega4.c04.cor.2-3-12}
\expandafter\def\csname sga@target@IV.2.3.13\endcsname{ega4.c04.cor.2-3-13}
\expandafter\def\csname sga@target@IV.2.3.14\endcsname{ega4.c04.cor.2-3-14}
\expandafter\def\csname sga@target@IV.2.3.15\endcsname{ega4.c04.rem.2-3-15}
\expandafter\def\csname sga@target@IV.2.3.2\endcsname{ega4.c04.prop.2-3-2}
\expandafter\def\csname sga@target@IV.2.3.3\endcsname{ega4.c04.statement.2-3-3}
\expandafter\def\csname sga@target@IV.2.3.4\endcsname{ega4.c04.prop.2-3-4}
\expandafter\def\csname sga@target@IV.2.3.5\endcsname{ega4.c04.cor.2-3-5}
\expandafter\def\csname sga@target@IV.2.3.6\endcsname{ega4.c04.prop.2-3-6}
\expandafter\def\csname sga@target@IV.2.3.7\endcsname{ega4.c04.prop.2-3-7}
\expandafter\def\csname sga@target@IV.2.3.8\endcsname{ega4.c04.cor.2-3-8}
\expandafter\def\csname sga@target@IV.2.3.9\endcsname{ega4.c04.rem.2-3-9}
\expandafter\def\csname sga@target@IV.2.4.1\endcsname{ega4.c04.statement.2-4-1}
\expandafter\def\csname sga@target@IV.2.4.10\endcsname{ega4.c04.cor.2-4-10}
\expandafter\def\csname sga@target@IV.2.4.11\endcsname{ega4.c04.rem.2-4-11}
\expandafter\def\csname sga@target@IV.2.4.2\endcsname{ega4.c04.def.2-4-2}
\expandafter\def\csname sga@target@IV.2.4.3\endcsname{ega4.c04.prop.2-4-3}
\expandafter\def\csname sga@target@IV.2.4.4\endcsname{ega4.c04.prop.2-4-4}
\expandafter\def\csname sga@target@IV.2.4.5\endcsname{ega4.c04.prop.2-4-5}
\expandafter\def\csname sga@target@IV.2.4.6\endcsname{ega4.c04.thm.2-4-6}
\expandafter\def\csname sga@target@IV.2.4.7\endcsname{ega4.c04.cor.2-4-7}
\expandafter\def\csname sga@target@IV.2.4.8\endcsname{ega4.c04.rem.2-4-8}
\expandafter\def\csname sga@target@IV.2.4.9\endcsname{ega4.c04.prop.2-4-9}
\expandafter\def\csname sga@target@IV.2.5.1\endcsname{ega4.c04.prop.2-5-1}
\expandafter\def\csname sga@target@IV.2.5.2\endcsname{ega4.c04.prop.2-5-2}
\expandafter\def\csname sga@target@IV.2.5.3\endcsname{ega4.c04.lem.2-5-3}
\expandafter\def\csname sga@target@IV.2.5.4\endcsname{ega4.c04.rem.2-5-4}
\expandafter\def\csname sga@target@IV.2.5.4.1\endcsname{ega4.c04.lem.2-5-4-1}
\expandafter\def\csname sga@target@IV.2.5.5\endcsname{ega4.c04.prop.2-5-5}
\expandafter\def\csname sga@target@IV.2.5.5.1\endcsname{ega4.c04.statement.2-5-5-1}
\expandafter\def\csname sga@target@IV.2.5.5.2\endcsname{ega4.c04.lem.2-5-5-2}
\expandafter\def\csname sga@target@IV.2.5.5.3\endcsname{ega4.c04.statement.2-5-5-3}
\expandafter\def\csname sga@target@IV.2.5.5.4\endcsname{ega4.c04.rem.2-5-5-4}
\expandafter\def\csname sga@target@IV.2.5.6\endcsname{ega4.c04.lem.2-5-6}
\expandafter\def\csname sga@target@IV.2.5.7\endcsname{ega4.c04.cor.2-5-7}
\expandafter\def\csname sga@target@IV.2.5.8\endcsname{ega4.c04.prop.2-5-8}
\expandafter\def\csname sga@target@IV.2.5.8.1\endcsname{ega4.c04.statement.2-5-8-1}
\expandafter\def\csname sga@target@IV.2.5.8.2\endcsname{ega4.c04.statement.2-5-8-2}
\expandafter\def\csname sga@target@IV.2.5.8.3\endcsname{ega4.c04.statement.2-5-8-3}
\expandafter\def\csname sga@target@IV.2.5.8.4\endcsname{ega4.c04.statement.2-5-8-4}
\expandafter\def\csname sga@target@IV.2.5.8.5\endcsname{ega4.c04.statement.2-5-8-5}
\expandafter\def\csname sga@target@IV.2.5.8.6\endcsname{ega4.c04.statement.2-5-8-6}
\expandafter\def\csname sga@target@IV.2.5.8.7\endcsname{ega4.c04.statement.2-5-8-7}
\expandafter\def\csname sga@target@IV.2.5.8.8\endcsname{ega4.c04.statement.2-5-8-8}
\expandafter\def\csname sga@target@IV.2.5.8.9\endcsname{ega4.c04.statement.2-5-8-9}
\expandafter\def\csname sga@target@IV.2.6.1\endcsname{ega4.c04.prop.2-6-1}
\expandafter\def\csname sga@target@IV.2.6.2\endcsname{ega4.c04.prop.2-6-2}
\expandafter\def\csname sga@target@IV.2.6.3\endcsname{ega4.c04.rem.2-6-3}
\expandafter\def\csname sga@target@IV.2.6.4\endcsname{ega4.c04.cor.2-6-4}
\expandafter\def\csname sga@target@IV.2.7.1\endcsname{ega4.c04.prop.2-7-1}
\expandafter\def\csname sga@target@IV.2.7.1.1\endcsname{ega4.c04.lem.2-7-1-1}
\expandafter\def\csname sga@target@IV.2.7.1.viii\endcsname{ega4.c04.item.2-7-1-viii}
\expandafter\def\csname sga@target@IV.2.7.2\endcsname{ega4.c04.cor.2-7-2}
\expandafter\def\csname sga@target@IV.2.7.3\endcsname{ega4.c04.sec.2-7-3}
\expandafter\def\csname sga@target@IV.2.8.1\endcsname{ega4.c04.prop.2-8-1}
\expandafter\def\csname sga@target@IV.2.8.1.1\endcsname{ega4.c04.lem.2-8-1-1}
\expandafter\def\csname sga@target@IV.2.8.2\endcsname{ega4.c04.cor.2-8-2}
\expandafter\def\csname sga@target@IV.2.8.3\endcsname{ega4.c04.statement.2-8-3}
\expandafter\def\csname sga@target@IV.2.8.4\endcsname{ega4.c04.prop.2-8-4}
\expandafter\def\csname sga@target@IV.2.8.4.1\endcsname{ega4.c04.statement.2-8-4-1}
\expandafter\def\csname sga@target@IV.2.8.5\endcsname{ega4.c04.prop.2-8-5}
\expandafter\def\csname sga@target@IV.2.8.6\endcsname{ega4.c04.cor.2-8-6}
\expandafter\def\csname sga@target@IV.2.8.6.1\endcsname{ega4.c04.statement.2-8-6-1}
\expandafter\def\csname sga@target@IV.20\endcsname{ega4.c04.sec.20}
\expandafter\def\csname sga@target@IV.20.0\endcsname{ega4.c04.subsec.20-0}
\expandafter\def\csname sga@target@IV.20.1\endcsname{ega4.c04.subsec.20-1}
\expandafter\def\csname sga@target@IV.20.1.1\endcsname{ega4.c04.statement.20-1-1}
\expandafter\def\csname sga@target@IV.20.1.1.1\endcsname{ega4.c04.eq.20-1-1-1}
\expandafter\def\csname sga@target@IV.20.1.10\endcsname{ega4.c04.statement.20-1-10}
\expandafter\def\csname sga@target@IV.20.1.11\endcsname{ega4.c04.statement.20-1-11}
\expandafter\def\csname sga@target@IV.20.1.11.1\endcsname{ega4.c04.eq.20-1-11-1}
\expandafter\def\csname sga@target@IV.20.1.11.2\endcsname{ega4.c04.eq.20-1-11-2}
\expandafter\def\csname sga@target@IV.20.1.12\endcsname{ega4.c04.prop.20-1-12}
\expandafter\def\csname sga@target@IV.20.1.13\endcsname{ega4.c04.rem.20-1-13}
\expandafter\def\csname sga@target@IV.20.1.2\endcsname{ega4.c04.statement.20-1-2}
\expandafter\def\csname sga@target@IV.20.1.2.1\endcsname{ega4.c04.statement.20-1-2-1}
\expandafter\def\csname sga@target@IV.20.1.2.2\endcsname{ega4.c04.statement.20-1-2-2}
\expandafter\def\csname sga@target@IV.20.1.3\endcsname{ega4.c04.statement.20-1-3}
\expandafter\def\csname sga@target@IV.20.1.3.1\endcsname{ega4.c04.statement.20-1-3-1}
\expandafter\def\csname sga@target@IV.20.1.4\endcsname{ega4.c04.statement.20-1-4}
\expandafter\def\csname sga@target@IV.20.1.4.1\endcsname{ega4.c04.statement.20-1-4-1}
\expandafter\def\csname sga@target@IV.20.1.5\endcsname{ega4.c04.statement.20-1-5}
\expandafter\def\csname sga@target@IV.20.1.5.1\endcsname{ega4.c04.statement.20-1-5-1}
\expandafter\def\csname sga@target@IV.20.1.6\endcsname{ega4.c04.prop.20-1-6}
\expandafter\def\csname sga@target@IV.20.1.7\endcsname{ega4.c04.statement.20-1-7}
\expandafter\def\csname sga@target@IV.20.1.8\endcsname{ega4.c04.statement.20-1-8}
\expandafter\def\csname sga@target@IV.20.1.9\endcsname{ega4.c04.statement.20-1-9}
\expandafter\def\csname sga@target@IV.20.2\endcsname{ega4.c04.subsec.20-2}
\expandafter\def\csname sga@target@IV.20.2.1\endcsname{ega4.c04.statement.20-2-1}
\expandafter\def\csname sga@target@IV.20.2.10\endcsname{ega4.c04.statement.20-2-10}
\expandafter\def\csname sga@target@IV.20.2.10.1\endcsname{ega4.c04.eq.20-2-10-1}
\expandafter\def\csname sga@target@IV.20.2.10.2\endcsname{ega4.c04.eq.20-2-10-2}
\expandafter\def\csname sga@target@IV.20.2.11\endcsname{ega4.c04.prop.20-2-11}
\expandafter\def\csname sga@target@IV.20.2.11.1\endcsname{ega4.c04.eq.20-2-11-1}
\expandafter\def\csname sga@target@IV.20.2.12\endcsname{ega4.c04.cor.20-2-12}
\expandafter\def\csname sga@target@IV.20.2.13\endcsname{ega4.c04.statement.20-2-13}
\expandafter\def\csname sga@target@IV.20.2.14\endcsname{ega4.c04.statement.20-2-14}
\expandafter\def\csname sga@target@IV.20.2.15\endcsname{ega4.c04.prop.20-2-15}
\expandafter\def\csname sga@target@IV.20.2.16\endcsname{ega4.c04.rem.20-2-16}
\expandafter\def\csname sga@target@IV.20.2.2\endcsname{ega4.c04.statement.20-2-2}
\expandafter\def\csname sga@target@IV.20.2.3\endcsname{ega4.c04.statement.20-2-3}
\expandafter\def\csname sga@target@IV.20.2.3.1\endcsname{ega4.c04.statement.20-2-3-1}
\expandafter\def\csname sga@target@IV.20.2.4\endcsname{ega4.c04.prop.20-2-4}
\expandafter\def\csname sga@target@IV.20.2.5\endcsname{ega4.c04.cor.20-2-5}
\expandafter\def\csname sga@target@IV.20.2.6\endcsname{ega4.c04.cor.20-2-6}
\expandafter\def\csname sga@target@IV.20.2.7\endcsname{ega4.c04.rem.20-2-7}
\expandafter\def\csname sga@target@IV.20.2.8\endcsname{ega4.c04.statement.20-2-8}
\expandafter\def\csname sga@target@IV.20.2.9\endcsname{ega4.c04.lem.20-2-9}
\expandafter\def\csname sga@target@IV.20.3\endcsname{ega4.c04.subsec.20-3}
\expandafter\def\csname sga@target@IV.20.3.1\endcsname{ega4.c04.statement.20-3-1}
\expandafter\def\csname sga@target@IV.20.3.10\endcsname{ega4.c04.lem.20-3-10}
\expandafter\def\csname sga@target@IV.20.3.11\endcsname{ega4.c04.prop.20-3-11}
\expandafter\def\csname sga@target@IV.20.3.11.1\endcsname{ega4.c04.eq.20-3-11-1}
\expandafter\def\csname sga@target@IV.20.3.12\endcsname{ega4.c04.cor.20-3-12}
\expandafter\def\csname sga@target@IV.20.3.13\endcsname{ega4.c04.rem.20-3-13}
\expandafter\def\csname sga@target@IV.20.3.2\endcsname{ega4.c04.statement.20-3-2}
\expandafter\def\csname sga@target@IV.20.3.2.1\endcsname{ega4.c04.eq.20-3-2-1}
\expandafter\def\csname sga@target@IV.20.3.3\endcsname{ega4.c04.statement.20-3-3}
\expandafter\def\csname sga@target@IV.20.3.4\endcsname{ega4.c04.statement.20-3-4}
\expandafter\def\csname sga@target@IV.20.3.5\endcsname{ega4.c04.prop.20-3-5}
\expandafter\def\csname sga@target@IV.20.3.6\endcsname{ega4.c04.statement.20-3-6}
\expandafter\def\csname sga@target@IV.20.3.6.1\endcsname{ega4.c04.eq.20-3-6-1}
\expandafter\def\csname sga@target@IV.20.3.7\endcsname{ega4.c04.prop.20-3-7}
\expandafter\def\csname sga@target@IV.20.3.8\endcsname{ega4.c04.prop.20-3-8}
\expandafter\def\csname sga@target@IV.20.3.8.1\endcsname{ega4.c04.eq.20-3-8-1}
\expandafter\def\csname sga@target@IV.20.3.8.2\endcsname{ega4.c04.lem.20-3-8-2}
\expandafter\def\csname sga@target@IV.20.3.8.3\endcsname{ega4.c04.lem.20-3-8-3}
\expandafter\def\csname sga@target@IV.20.3.8.4\endcsname{ega4.c04.lem.20-3-8-4}
\expandafter\def\csname sga@target@IV.20.3.8.5\endcsname{ega4.c04.sec.20-3-8-5}
\expandafter\def\csname sga@target@IV.20.3.8.6\endcsname{ega4.c04.sec.20-3-8-6}
\expandafter\def\csname sga@target@IV.20.3.8.7\endcsname{ega4.c04.sec.20-3-8-7}
\expandafter\def\csname sga@target@IV.20.3.8.8\endcsname{ega4.c04.lem.20-3-8-8}
\expandafter\def\csname sga@target@IV.20.3.9\endcsname{ega4.c04.cor.20-3-9}
\expandafter\def\csname sga@target@IV.20.4\endcsname{ega4.c04.subsec.20-4}
\expandafter\def\csname sga@target@IV.20.4.1\endcsname{ega4.c04.statement.20-4-1}
\expandafter\def\csname sga@target@IV.20.4.10\endcsname{ega4.c04.cor.20-4-10}
\expandafter\def\csname sga@target@IV.20.4.11\endcsname{ega4.c04.prop.20-4-11}
\expandafter\def\csname sga@target@IV.20.4.12\endcsname{ega4.c04.cor.20-4-12}
\expandafter\def\csname sga@target@IV.20.4.2\endcsname{ega4.c04.statement.20-4-2}
\expandafter\def\csname sga@target@IV.20.4.3\endcsname{ega4.c04.statement.20-4-3}
\expandafter\def\csname sga@target@IV.20.4.4\endcsname{ega4.c04.prop.20-4-4}
\expandafter\def\csname sga@target@IV.20.4.5\endcsname{ega4.c04.prop.20-4-5}
\expandafter\def\csname sga@target@IV.20.4.6\endcsname{ega4.c04.prop.20-4-6}
\expandafter\def\csname sga@target@IV.20.4.7\endcsname{ega4.c04.cor.20-4-7}
\expandafter\def\csname sga@target@IV.20.4.8\endcsname{ega4.c04.rem.20-4-8}
\expandafter\def\csname sga@target@IV.20.4.9\endcsname{ega4.c04.cor.20-4-9}
\expandafter\def\csname sga@target@IV.20.5\endcsname{ega4.c04.subsec.20-5}
\expandafter\def\csname sga@target@IV.20.5.1\endcsname{ega4.c04.statement.20-5-1}
\expandafter\def\csname sga@target@IV.20.5.2\endcsname{ega4.c04.statement.20-5-2}
\expandafter\def\csname sga@target@IV.20.5.2.1\endcsname{ega4.c04.statement.20-5-2-1}
\expandafter\def\csname sga@target@IV.20.5.3\endcsname{ega4.c04.prop.20-5-3}
\expandafter\def\csname sga@target@IV.20.5.4\endcsname{ega4.c04.statement.20-5-4}
\expandafter\def\csname sga@target@IV.20.5.5\endcsname{ega4.c04.statement.20-5-5}
\expandafter\def\csname sga@target@IV.20.5.6\endcsname{ega4.c04.statement.20-5-6}
\expandafter\def\csname sga@target@IV.20.5.7\endcsname{ega4.c04.statement.20-5-7}
\expandafter\def\csname sga@target@IV.20.5.8\endcsname{ega4.c04.statement.20-5-8}
\expandafter\def\csname sga@target@IV.20.6\endcsname{ega4.c04.subsec.20-6}
\expandafter\def\csname sga@target@IV.20.6.1\endcsname{ega4.c04.statement.20-6-1}
\expandafter\def\csname sga@target@IV.20.6.1.1\endcsname{ega4.c04.eq.20-6-1-1}
\expandafter\def\csname sga@target@IV.20.6.2\endcsname{ega4.c04.statement.20-6-2}
\expandafter\def\csname sga@target@IV.20.6.2.1\endcsname{ega4.c04.statement.20-6-2-1}
\expandafter\def\csname sga@target@IV.20.6.2.2\endcsname{ega4.c04.eq.20-6-2-2}
\expandafter\def\csname sga@target@IV.20.6.3\endcsname{ega4.c04.prop.20-6-3}
\expandafter\def\csname sga@target@IV.20.6.4\endcsname{ega4.c04.prop.20-6-4}
\expandafter\def\csname sga@target@IV.20.6.5\endcsname{ega4.c04.statement.20-6-5}
\expandafter\def\csname sga@target@IV.20.6.5.1\endcsname{ega4.c04.eq.20-6-5-1}
\expandafter\def\csname sga@target@IV.20.6.6\endcsname{ega4.c04.prop.20-6-6}
\expandafter\def\csname sga@target@IV.20.6.7\endcsname{ega4.c04.statement.20-6-7}
\expandafter\def\csname sga@target@IV.20.6.8\endcsname{ega4.c04.prop.20-6-8}
\expandafter\def\csname sga@target@IV.20.6.8.1\endcsname{ega4.c04.eq.20-6-8-1}
\expandafter\def\csname sga@target@IV.20.6.9\endcsname{ega4.c04.statement.20-6-9}
\expandafter\def\csname sga@target@IV.20.6.9.1\endcsname{ega4.c04.statement.20-6-9-1}
\expandafter\def\csname sga@target@IV.21.1\endcsname{ega4.c04.subsec.21-1.key-iv-21-1}
\expandafter\def\csname sga@target@IV.21.1.1\endcsname{ega4.c04.statement.21-1-1}
\expandafter\def\csname sga@target@IV.21.1.10\endcsname{ega4.c04.prop.21-1-10}
\expandafter\def\csname sga@target@IV.21.1.2\endcsname{ega4.c04.statement.21-1-2}
\expandafter\def\csname sga@target@IV.21.1.3\endcsname{ega4.c04.statement.21-1-3}
\expandafter\def\csname sga@target@IV.21.1.3.1\endcsname{ega4.c04.statement.21-1-3-1}
\expandafter\def\csname sga@target@IV.21.1.3.2\endcsname{ega4.c04.statement.21-1-3-2}
\expandafter\def\csname sga@target@IV.21.1.3.3\endcsname{ega4.c04.statement.21-1-3-3}
\expandafter\def\csname sga@target@IV.21.1.3.4\endcsname{ega4.c04.statement.21-1-3-4}
\expandafter\def\csname sga@target@IV.21.1.4\endcsname{ega4.c04.statement.21-1-4}
\expandafter\def\csname sga@target@IV.21.1.4.1\endcsname{ega4.c04.eq.21-1-4-1}
\expandafter\def\csname sga@target@IV.21.1.4.2\endcsname{ega4.c04.eq.21-1-4-2}
\expandafter\def\csname sga@target@IV.21.1.4.3\endcsname{ega4.c04.eq.21-1-4-3}
\expandafter\def\csname sga@target@IV.21.1.5\endcsname{ega4.c04.statement.21-1-5}
\expandafter\def\csname sga@target@IV.21.1.5.1\endcsname{ega4.c04.statement.21-1-5-1}
\expandafter\def\csname sga@target@IV.21.1.5.2\endcsname{ega4.c04.eq.21-1-5-2}
\expandafter\def\csname sga@target@IV.21.1.6\endcsname{ega4.c04.def.21-1-6}
\expandafter\def\csname sga@target@IV.21.1.6.1\endcsname{ega4.c04.eq.21-1-6-1}
\expandafter\def\csname sga@target@IV.21.1.6.2\endcsname{ega4.c04.eq.21-1-6-2}
\expandafter\def\csname sga@target@IV.21.1.6.3\endcsname{ega4.c04.eq.21-1-6-3}
\expandafter\def\csname sga@target@IV.21.1.6.4\endcsname{ega4.c04.eq.21-1-6-4}
\expandafter\def\csname sga@target@IV.21.1.7\endcsname{ega4.c04.statement.21-1-7}
\expandafter\def\csname sga@target@IV.21.1.7.1\endcsname{ega4.c04.statement.21-1-7-1}
\expandafter\def\csname sga@target@IV.21.1.8\endcsname{ega4.c04.prop.21-1-8}
\expandafter\def\csname sga@target@IV.21.1.9\endcsname{ega4.c04.cor.21-1-9}
\expandafter\def\csname sga@target@IV.21.10.1\endcsname{ega4.c04.statement.21-10-1}
\expandafter\def\csname sga@target@IV.21.10.1.1\endcsname{ega4.c04.eq.21-10-1-1}
\expandafter\def\csname sga@target@IV.21.10.1.2\endcsname{ega4.c04.eq.21-10-1-2}
\expandafter\def\csname sga@target@IV.21.10.10\endcsname{ega4.c04.statement.21-10-10}
\expandafter\def\csname sga@target@IV.21.10.11\endcsname{ega4.c04.statement.21-10-11}
\expandafter\def\csname sga@target@IV.21.10.11.1\endcsname{ega4.c04.eq.21-10-11-1}
\expandafter\def\csname sga@target@IV.21.10.12\endcsname{ega4.c04.rem.21-10-12}
\expandafter\def\csname sga@target@IV.21.10.13\endcsname{ega4.c04.lem.21-10-13}
\expandafter\def\csname sga@target@IV.21.10.13.1\endcsname{ega4.c04.prop.21-10-13-1}
\expandafter\def\csname sga@target@IV.21.10.13bis\endcsname{ega4.c04.prop.21-10-13bis}
\expandafter\def\csname sga@target@IV.21.10.14\endcsname{ega4.c04.statement.21-10-14}
\expandafter\def\csname sga@target@IV.21.10.14.1\endcsname{ega4.c04.eq.21-10-14-1}
\expandafter\def\csname sga@target@IV.21.10.14.2\endcsname{ega4.c04.eq.21-10-14-2}
\expandafter\def\csname sga@target@IV.21.10.15\endcsname{ega4.c04.prop.21-10-15}
\expandafter\def\csname sga@target@IV.21.10.16\endcsname{ega4.c04.prop.21-10-16}
\expandafter\def\csname sga@target@IV.21.10.16.1\endcsname{ega4.c04.prop.21-10-16-1}
\expandafter\def\csname sga@target@IV.21.10.16.2\endcsname{ega4.c04.eq.21-10-16-2}
\expandafter\def\csname sga@target@IV.21.10.17\endcsname{ega4.c04.prop.21-10-17}
\expandafter\def\csname sga@target@IV.21.10.17.1\endcsname{ega4.c04.prop.21-10-17-1}
\expandafter\def\csname sga@target@IV.21.10.17.10\endcsname{ega4.c04.eq.21-10-17-10}
\expandafter\def\csname sga@target@IV.21.10.17.2\endcsname{ega4.c04.eq.21-10-17-2}
\expandafter\def\csname sga@target@IV.21.10.17.3\endcsname{ega4.c04.lem.21-10-17-3}
\expandafter\def\csname sga@target@IV.21.10.17.4\endcsname{ega4.c04.lem.21-10-17-4}
\expandafter\def\csname sga@target@IV.21.10.17.5\endcsname{ega4.c04.eq.21-10-17-5}
\expandafter\def\csname sga@target@IV.21.10.17.6\endcsname{ega4.c04.eq.21-10-17-6}
\expandafter\def\csname sga@target@IV.21.10.17.7\endcsname{ega4.c04.lem.21-10-17-7}
\expandafter\def\csname sga@target@IV.21.10.17.8\endcsname{ega4.c04.lem.21-10-17-8}
\expandafter\def\csname sga@target@IV.21.10.17.9\endcsname{ega4.c04.eq.21-10-17-9}
\expandafter\def\csname sga@target@IV.21.10.18\endcsname{ega4.c04.prop.21-10-18}
\expandafter\def\csname sga@target@IV.21.10.18.1\endcsname{ega4.c04.prop.21-10-18-1}
\expandafter\def\csname sga@target@IV.21.10.19\endcsname{ega4.c04.cor.21-10-19}
\expandafter\def\csname sga@target@IV.21.10.19.1\endcsname{ega4.c04.cor.21-10-19-1}
\expandafter\def\csname sga@target@IV.21.10.2\endcsname{ega4.c04.statement.21-10-2}
\expandafter\def\csname sga@target@IV.21.10.3\endcsname{ega4.c04.statement.21-10-3}
\expandafter\def\csname sga@target@IV.21.10.3.1\endcsname{ega4.c04.statement.21-10-3-1}
\expandafter\def\csname sga@target@IV.21.10.4\endcsname{ega4.c04.prop.21-10-4}
\expandafter\def\csname sga@target@IV.21.10.4.1\endcsname{ega4.c04.prop.21-10-4-1}
\expandafter\def\csname sga@target@IV.21.10.5\endcsname{ega4.c04.statement.21-10-5}
\expandafter\def\csname sga@target@IV.21.10.6\endcsname{ega4.c04.prop.21-10-6}
\expandafter\def\csname sga@target@IV.21.10.7\endcsname{ega4.c04.rem.21-10-7}
\expandafter\def\csname sga@target@IV.21.10.8\endcsname{ega4.c04.prop.21-10-8}
\expandafter\def\csname sga@target@IV.21.10.9\endcsname{ega4.c04.statement.21-10-9}
\expandafter\def\csname sga@target@IV.21.11\endcsname{ega4.c04.subsec.21-11.key-iv-21-11}
\expandafter\def\csname sga@target@IV.21.11.1\endcsname{ega4.c04.thm.21-11-1}
\expandafter\def\csname sga@target@IV.21.11.1.1\endcsname{ega4.c04.eq.21-11-1-1}
\expandafter\def\csname sga@target@IV.21.11.2\endcsname{ega4.c04.lem.21-11-2}
\expandafter\def\csname sga@target@IV.21.12\endcsname{ega4.c04.subsec.21-12.key-iv-21-12}
\expandafter\def\csname sga@target@IV.21.12.1\endcsname{ega4.c04.statement.21-12-1}
\expandafter\def\csname sga@target@IV.21.12.1.1\endcsname{ega4.c04.eq.21-12-1-1}
\expandafter\def\csname sga@target@IV.21.12.10\endcsname{ega4.c04.prop.21-12-10}
\expandafter\def\csname sga@target@IV.21.12.12\endcsname{ega4.c04.thm.21-12-12}
\expandafter\def\csname sga@target@IV.21.12.13\endcsname{ega4.c04.cor.21-12-13}
\expandafter\def\csname sga@target@IV.21.12.14\endcsname{ega4.c04.item.21-12-14}
\expandafter\def\csname sga@target@IV.21.12.15\endcsname{ega4.c04.prop.21-12-15}
\expandafter\def\csname sga@target@IV.21.12.16\endcsname{ega4.c04.cor.21-12-16}
\expandafter\def\csname sga@target@IV.21.12.17\endcsname{ega4.c04.item.21-12-17}
\expandafter\def\csname sga@target@IV.21.12.2\endcsname{ega4.c04.prop.21-12-2}
\expandafter\def\csname sga@target@IV.21.12.2.1\endcsname{ega4.c04.prop.21-12-2-1}
\expandafter\def\csname sga@target@IV.21.12.3\endcsname{ega4.c04.cor.21-12-3}
\expandafter\def\csname sga@target@IV.21.12.4\endcsname{ega4.c04.eq.21-12-4}
\expandafter\def\csname sga@target@IV.21.12.5\endcsname{ega4.c04.statement.21-12-5}
\expandafter\def\csname sga@target@IV.21.12.5.1\endcsname{ega4.c04.eq.21-12-5-1}
\expandafter\def\csname sga@target@IV.21.12.5.2\endcsname{ega4.c04.eq.21-12-5-2}
\expandafter\def\csname sga@target@IV.21.12.5.3\endcsname{ega4.c04.eq.21-12-5-3}
\expandafter\def\csname sga@target@IV.21.12.6\endcsname{ega4.c04.thm.21-12-6}
\expandafter\def\csname sga@target@IV.21.12.7\endcsname{ega4.c04.cor.21-12-7}
\expandafter\def\csname sga@target@IV.21.12.8\endcsname{ega4.c04.statement.21-12-8}
\expandafter\def\csname sga@target@IV.21.12.9\endcsname{ega4.c04.statement.21-12-9}
\expandafter\def\csname sga@target@IV.21.13\endcsname{ega4.c04.subsec.21-13.key-iv-21-13}
\expandafter\def\csname sga@target@IV.21.13.1\endcsname{ega4.c04.def.21-13-1}
\expandafter\def\csname sga@target@IV.21.13.10\endcsname{ega4.c04.prop.21-13-10}
\expandafter\def\csname sga@target@IV.21.13.11\endcsname{ega4.c04.cor.21-13-11}
\expandafter\def\csname sga@target@IV.21.13.12\endcsname{ega4.c04.prop.21-13-12}
\expandafter\def\csname sga@target@IV.21.13.12.1\endcsname{ega4.c04.statement.21-13-12-1}
\expandafter\def\csname sga@target@IV.21.13.13\endcsname{ega4.c04.rem.21-13-13}
\expandafter\def\csname sga@target@IV.21.13.14\endcsname{ega4.c04.prop.21-13-14}
\expandafter\def\csname sga@target@IV.21.13.15\endcsname{ega4.c04.cor.21-13-15}
\expandafter\def\csname sga@target@IV.21.13.16\endcsname{ega4.c04.rem.21-13-16}
\expandafter\def\csname sga@target@IV.21.13.2\endcsname{ega4.c04.lem.21-13-2}
\expandafter\def\csname sga@target@IV.21.13.2.1\endcsname{ega4.c04.eq.21-13-2-1}
\expandafter\def\csname sga@target@IV.21.13.3\endcsname{ega4.c04.cor.21-13-3}
\expandafter\def\csname sga@target@IV.21.13.4\endcsname{ega4.c04.lem.21-13-4}
\expandafter\def\csname sga@target@IV.21.13.5\endcsname{ega4.c04.prop.21-13-5}
\expandafter\def\csname sga@target@IV.21.13.6\endcsname{ega4.c04.cor.21-13-6}
\expandafter\def\csname sga@target@IV.21.13.7\endcsname{ega4.c04.def.21-13-7}
\expandafter\def\csname sga@target@IV.21.13.8\endcsname{ega4.c04.prop.21-13-8}
\expandafter\def\csname sga@target@IV.21.13.9\endcsname{ega4.c04.statement.21-13-9}
\expandafter\def\csname sga@target@IV.21.13.9.1\endcsname{ega4.c04.eq.21-13-9-1}
\expandafter\def\csname sga@target@IV.21.14\endcsname{ega4.c04.subsec.21-14.key-iv-21-14}
\expandafter\def\csname sga@target@IV.21.14.1\endcsname{ega4.c04.thm.21-14-1}
\expandafter\def\csname sga@target@IV.21.14.2\endcsname{ega4.c04.cor.21-14-2}
\expandafter\def\csname sga@target@IV.21.14.3\endcsname{ega4.c04.prop.21-14-3}
\expandafter\def\csname sga@target@IV.21.14.4\endcsname{ega4.c04.statement.21-14-4}
\expandafter\def\csname sga@target@IV.21.14.4bis\endcsname{ega4.c04.statement.21-14-4bis}
\expandafter\def\csname sga@target@IV.21.15\endcsname{ega4.c04.subsec.21-15.key-iv-21-15}
\expandafter\def\csname sga@target@IV.21.15.1\endcsname{ega4.c04.statement.21-15-1}
\expandafter\def\csname sga@target@IV.21.15.2\endcsname{ega4.c04.def.21-15-2}
\expandafter\def\csname sga@target@IV.21.15.3\endcsname{ega4.c04.statement.21-15-3}
\expandafter\def\csname sga@target@IV.21.15.3.1\endcsname{ega4.c04.eq.21-15-3-1}
\expandafter\def\csname sga@target@IV.21.15.3.2\endcsname{ega4.c04.eq.21-15-3-2}
\expandafter\def\csname sga@target@IV.21.15.3bis\endcsname{ega4.c04.statement.21-15-3bis}
\expandafter\def\csname sga@target@IV.21.15.4\endcsname{ega4.c04.prop.21-15-4}
\expandafter\def\csname sga@target@IV.21.15.5\endcsname{ega4.c04.prop.21-15-5}
\expandafter\def\csname sga@target@IV.21.15.6\endcsname{ega4.c04.prop.21-15-6}
\expandafter\def\csname sga@target@IV.21.15.7\endcsname{ega4.c04.statement.21-15-7}
\expandafter\def\csname sga@target@IV.21.15.8\endcsname{ega4.c04.statement.21-15-8}
\expandafter\def\csname sga@target@IV.21.15.9\endcsname{ega4.c04.statement.21-15-9}
\expandafter\def\csname sga@target@IV.21.15.9.1\endcsname{ega4.c04.eq.21-15-9-1}
\expandafter\def\csname sga@target@IV.21.2.1\endcsname{ega4.c04.statement.21-2-1}
\expandafter\def\csname sga@target@IV.21.2.10\endcsname{ega4.c04.statement.21-2-10}
\expandafter\def\csname sga@target@IV.21.2.10.1\endcsname{ega4.c04.eq.21-2-10-1}
\expandafter\def\csname sga@target@IV.21.2.11\endcsname{ega4.c04.prop.21-2-11}
\expandafter\def\csname sga@target@IV.21.2.11.1\endcsname{ega4.c04.prop.21-2-11-1}
\expandafter\def\csname sga@target@IV.21.2.12\endcsname{ega4.c04.statement.21-2-12}
\expandafter\def\csname sga@target@IV.21.2.2\endcsname{ega4.c04.prop.21-2-2}
\expandafter\def\csname sga@target@IV.21.2.3\endcsname{ega4.c04.cor.21-2-3}
\expandafter\def\csname sga@target@IV.21.2.4\endcsname{ega4.c04.statement.21-2-4}
\expandafter\def\csname sga@target@IV.21.2.5\endcsname{ega4.c04.statement.21-2-5}
\expandafter\def\csname sga@target@IV.21.2.5.1\endcsname{ega4.c04.statement.21-2-5-1}
\expandafter\def\csname sga@target@IV.21.2.5.2\endcsname{ega4.c04.eq.21-2-5-2}
\expandafter\def\csname sga@target@IV.21.2.5.3\endcsname{ega4.c04.eq.21-2-5-3}
\expandafter\def\csname sga@target@IV.21.2.5.4\endcsname{ega4.c04.eq.21-2-5-4}
\expandafter\def\csname sga@target@IV.21.2.6\endcsname{ega4.c04.prop.21-2-6}
\expandafter\def\csname sga@target@IV.21.2.6.1\endcsname{ega4.c04.eq.21-2-6-1}
\expandafter\def\csname sga@target@IV.21.2.7\endcsname{ega4.c04.statement.21-2-7}
\expandafter\def\csname sga@target@IV.21.2.7.1\endcsname{ega4.c04.statement.21-2-7-1}
\expandafter\def\csname sga@target@IV.21.2.7.2\endcsname{ega4.c04.eq.21-2-7-2}
\expandafter\def\csname sga@target@IV.21.2.8\endcsname{ega4.c04.statement.21-2-8}
\expandafter\def\csname sga@target@IV.21.2.8.1\endcsname{ega4.c04.statement.21-2-8-1}
\expandafter\def\csname sga@target@IV.21.2.8.2\endcsname{ega4.c04.statement.21-2-8-2}
\expandafter\def\csname sga@target@IV.21.2.8.3\endcsname{ega4.c04.statement.21-2-8-3}
\expandafter\def\csname sga@target@IV.21.2.9\endcsname{ega4.c04.statement.21-2-9}
\expandafter\def\csname sga@target@IV.21.2.9.1\endcsname{ega4.c04.statement.21-2-9-1}
\expandafter\def\csname sga@target@IV.21.2.9.2\endcsname{ega4.c04.eq.21-2-9-2}
\expandafter\def\csname sga@target@IV.21.2.9.3\endcsname{ega4.c04.eq.21-2-9-3}
\expandafter\def\csname sga@target@IV.21.2.9.4\endcsname{ega4.c04.eq.21-2-9-4}
\expandafter\def\csname sga@target@IV.21.3.1\endcsname{ega4.c04.statement.21-3-1}
\expandafter\def\csname sga@target@IV.21.3.2\endcsname{ega4.c04.statement.21-3-2}
\expandafter\def\csname sga@target@IV.21.3.2.1\endcsname{ega4.c04.statement.21-3-2-1}
\expandafter\def\csname sga@target@IV.21.3.2.2\endcsname{ega4.c04.eq.21-3-2-2}
\expandafter\def\csname sga@target@IV.21.3.2.3\endcsname{ega4.c04.eq.21-3-2-3}
\expandafter\def\csname sga@target@IV.21.3.2.4\endcsname{ega4.c04.eq.21-3-2-4}
\expandafter\def\csname sga@target@IV.21.3.3\endcsname{ega4.c04.prop.21-3-3}
\expandafter\def\csname sga@target@IV.21.3.3.1\endcsname{ega4.c04.eq.21-3-3-1}
\expandafter\def\csname sga@target@IV.21.3.4\endcsname{ega4.c04.prop.21-3-4}
\expandafter\def\csname sga@target@IV.21.3.5\endcsname{ega4.c04.cor.21-3-5}
\expandafter\def\csname sga@target@IV.21.3.6\endcsname{ega4.c04.rem.21-3-6}
\expandafter\def\csname sga@target@IV.21.3.6.1\endcsname{ega4.c04.eq.21-3-6-1}
\expandafter\def\csname sga@target@IV.21.4\endcsname{ega4.c04.subsec.21-4.key-iv-21-4}
\expandafter\def\csname sga@target@IV.21.4.1\endcsname{ega4.c04.statement.21-4-1}
\expandafter\def\csname sga@target@IV.21.4.10\endcsname{ega4.c04.cor.21-4-10}
\expandafter\def\csname sga@target@IV.21.4.11\endcsname{ega4.c04.cor.21-4-11}
\expandafter\def\csname sga@target@IV.21.4.12\endcsname{ega4.c04.rem.21-4-12}
\expandafter\def\csname sga@target@IV.21.4.13\endcsname{ega4.c04.cor.21-4-13}
\expandafter\def\csname sga@target@IV.21.4.2\endcsname{ega4.c04.def.21-4-2}
\expandafter\def\csname sga@target@IV.21.4.2.1\endcsname{ega4.c04.eq.21-4-2-1}
\expandafter\def\csname sga@target@IV.21.4.3\endcsname{ega4.c04.statement.21-4-3}
\expandafter\def\csname sga@target@IV.21.4.4\endcsname{ega4.c04.statement.21-4-4}
\expandafter\def\csname sga@target@IV.21.4.5\endcsname{ega4.c04.prop.21-4-5}
\expandafter\def\csname sga@target@IV.21.4.6\endcsname{ega4.c04.cor.21-4-6}
\expandafter\def\csname sga@target@IV.21.4.6.1\endcsname{ega4.c04.eq.21-4-6-1}
\expandafter\def\csname sga@target@IV.21.4.7\endcsname{ega4.c04.statement.21-4-7}
\expandafter\def\csname sga@target@IV.21.4.8\endcsname{ega4.c04.prop.21-4-8}
\expandafter\def\csname sga@target@IV.21.4.9\endcsname{ega4.c04.prop.21-4-9}
\expandafter\def\csname sga@target@IV.21.5\endcsname{ega4.c04.subsec.21-5.key-iv-21-5}
\expandafter\def\csname sga@target@IV.21.5.1\endcsname{ega4.c04.statement.21-5-1}
\expandafter\def\csname sga@target@IV.21.5.2\endcsname{ega4.c04.lem.21-5-2}
\expandafter\def\csname sga@target@IV.21.5.3\endcsname{ega4.c04.statement.21-5-3}
\expandafter\def\csname sga@target@IV.21.5.3.1\endcsname{ega4.c04.eq.21-5-3-1}
\expandafter\def\csname sga@target@IV.21.5.3.2\endcsname{ega4.c04.eq.21-5-3-2}
\expandafter\def\csname sga@target@IV.21.5.4\endcsname{ega4.c04.statement.21-5-4}
\expandafter\def\csname sga@target@IV.21.5.5\endcsname{ega4.c04.def.21-5-5}
\expandafter\def\csname sga@target@IV.21.5.5.1\endcsname{ega4.c04.eq.21-5-5-1}
\expandafter\def\csname sga@target@IV.21.5.5.2\endcsname{ega4.c04.eq.21-5-5-2}
\expandafter\def\csname sga@target@IV.21.5.5.3\endcsname{ega4.c04.eq.21-5-5-3}
\expandafter\def\csname sga@target@IV.21.5.6\endcsname{ega4.c04.prop.21-5-6}
\expandafter\def\csname sga@target@IV.21.5.6.1\endcsname{ega4.c04.prop.21-5-6-1}
\expandafter\def\csname sga@target@IV.21.5.7\endcsname{ega4.c04.prop.21-5-7}
\expandafter\def\csname sga@target@IV.21.5.7.1\endcsname{ega4.c04.prop.21-5-7-1}
\expandafter\def\csname sga@target@IV.21.5.7.2\endcsname{ega4.c04.lem.21-5-7-2}
\expandafter\def\csname sga@target@IV.21.5.7.3\endcsname{ega4.c04.lem.21-5-7-3}
\expandafter\def\csname sga@target@IV.21.5.7.4\endcsname{ega4.c04.eq.21-5-7-4}
\expandafter\def\csname sga@target@IV.21.5.8\endcsname{ega4.c04.prop.21-5-8}
\expandafter\def\csname sga@target@IV.21.5.8.1\endcsname{ega4.c04.eq.21-5-8-1}
\expandafter\def\csname sga@target@IV.21.6\endcsname{ega4.c04.subsec.21-6.key-iv-21-6}
\expandafter\def\csname sga@target@IV.21.6.1\endcsname{ega4.c04.statement.21-6-1}
\expandafter\def\csname sga@target@IV.21.6.10\endcsname{ega4.c04.cor.21-6-10}
\expandafter\def\csname sga@target@IV.21.6.10.1\endcsname{ega4.c04.eq.21-6-10-1}
\expandafter\def\csname sga@target@IV.21.6.11\endcsname{ega4.c04.cor.21-6-11}
\expandafter\def\csname sga@target@IV.21.6.12\endcsname{ega4.c04.prop.21-6-12}
\expandafter\def\csname sga@target@IV.21.6.12.1\endcsname{ega4.c04.eq.21-6-12-1}
\expandafter\def\csname sga@target@IV.21.6.12.2\endcsname{ega4.c04.eq.21-6-12-2}
\expandafter\def\csname sga@target@IV.21.6.13\endcsname{ega4.c04.cor.21-6-13}
\expandafter\def\csname sga@target@IV.21.6.14\endcsname{ega4.c04.cor.21-6-14}
\expandafter\def\csname sga@target@IV.21.6.2\endcsname{ega4.c04.statement.21-6-2}
\expandafter\def\csname sga@target@IV.21.6.3\endcsname{ega4.c04.statement.21-6-3}
\expandafter\def\csname sga@target@IV.21.6.4\endcsname{ega4.c04.statement.21-6-4}
\expandafter\def\csname sga@target@IV.21.6.4.1\endcsname{ega4.c04.statement.21-6-4-1}
\expandafter\def\csname sga@target@IV.21.6.4.2\endcsname{ega4.c04.eq.21-6-4-2}
\expandafter\def\csname sga@target@IV.21.6.5\endcsname{ega4.c04.statement.21-6-5}
\expandafter\def\csname sga@target@IV.21.6.5.1\endcsname{ega4.c04.eq.21-6-5-1}
\expandafter\def\csname sga@target@IV.21.6.6\endcsname{ega4.c04.prop.21-6-6}
\expandafter\def\csname sga@target@IV.21.6.6.2\endcsname{ega4.c04.statement.21-6-6-2}
\expandafter\def\csname sga@target@IV.21.6.7\endcsname{ega4.c04.statement.21-6-7}
\expandafter\def\csname sga@target@IV.21.6.7.1\endcsname{ega4.c04.eq.21-6-7-1}
\expandafter\def\csname sga@target@IV.21.6.7.2\endcsname{ega4.c04.eq.21-6-7-2}
\expandafter\def\csname sga@target@IV.21.6.7.3\endcsname{ega4.c04.eq.21-6-7-3}
\expandafter\def\csname sga@target@IV.21.6.7.4\endcsname{ega4.c04.eq.21-6-7-4}
\expandafter\def\csname sga@target@IV.21.6.7.5\endcsname{ega4.c04.eq.21-6-7-5}
\expandafter\def\csname sga@target@IV.21.6.8\endcsname{ega4.c04.statement.21-6-8}
\expandafter\def\csname sga@target@IV.21.6.9\endcsname{ega4.c04.thm.21-6-9}
\expandafter\def\csname sga@target@IV.21.6.9.1\endcsname{ega4.c04.eq.21-6-9-1}
\expandafter\def\csname sga@target@IV.21.6.9.2\endcsname{ega4.c04.statement.21-6-9-2}
\expandafter\def\csname sga@target@IV.21.7\endcsname{ega4.c04.subsec.21-7.key-iv-21-7}
\expandafter\def\csname sga@target@IV.21.7.1\endcsname{ega4.c04.statement.21-7-1}
\expandafter\def\csname sga@target@IV.21.7.2\endcsname{ega4.c04.prop.21-7-2}
\expandafter\def\csname sga@target@IV.21.7.2.1\endcsname{ega4.c04.eq.21-7-2-1}
\expandafter\def\csname sga@target@IV.21.7.3\endcsname{ega4.c04.cor.21-7-3}
\expandafter\def\csname sga@target@IV.21.7.3.1\endcsname{ega4.c04.statement.21-7-3-1}
\expandafter\def\csname sga@target@IV.21.7.4\endcsname{ega4.c04.prop.21-7-4}
\expandafter\def\csname sga@target@IV.21.7.5\endcsname{ega4.c04.rem.21-7-5}
\expandafter\def\csname sga@target@IV.21.7.6\endcsname{ega4.c04.cor.21-7-6}
\expandafter\def\csname sga@target@IV.21.7.7\endcsname{ega4.c04.rem.21-7-7}
\expandafter\def\csname sga@target@IV.21.8.1\endcsname{ega4.c04.lem.21-8-1}
\expandafter\def\csname sga@target@IV.21.8.2\endcsname{ega4.c04.cor.21-8-2}
\expandafter\def\csname sga@target@IV.21.8.3\endcsname{ega4.c04.prop.21-8-3}
\expandafter\def\csname sga@target@IV.21.8.3.1\endcsname{ega4.c04.eq.21-8-3-1}
\expandafter\def\csname sga@target@IV.21.8.4\endcsname{ega4.c04.cor.21-8-4}
\expandafter\def\csname sga@target@IV.21.8.5\endcsname{ega4.c04.prop.21-8-5}
\expandafter\def\csname sga@target@IV.21.8.5.1\endcsname{ega4.c04.eq.21-8-5-1}
\expandafter\def\csname sga@target@IV.21.8.5.2\endcsname{ega4.c04.eq.21-8-5-2}
\expandafter\def\csname sga@target@IV.21.8.5.3\endcsname{ega4.c04.eq.21-8-5-3}
\expandafter\def\csname sga@target@IV.21.8.6\endcsname{ega4.c04.rem.21-8-6}
\expandafter\def\csname sga@target@IV.21.9\endcsname{ega4.c04.subsec.21-9.key-iv-21-9}
\expandafter\def\csname sga@target@IV.21.9.1\endcsname{ega4.c04.statement.21-9-1}
\expandafter\def\csname sga@target@IV.21.9.1.1\endcsname{ega4.c04.eq.21-9-1-1}
\expandafter\def\csname sga@target@IV.21.9.1.2\endcsname{ega4.c04.eq.21-9-1-2}
\expandafter\def\csname sga@target@IV.21.9.10\endcsname{ega4.c04.statement.21-9-10}
\expandafter\def\csname sga@target@IV.21.9.11\endcsname{ega4.c04.prop.21-9-11}
\expandafter\def\csname sga@target@IV.21.9.12\endcsname{ega4.c04.cor.21-9-12}
\expandafter\def\csname sga@target@IV.21.9.13\endcsname{ega4.c04.rem.21-9-13}
\expandafter\def\csname sga@target@IV.21.9.2\endcsname{ega4.c04.lem.21-9-2}
\expandafter\def\csname sga@target@IV.21.9.3\endcsname{ega4.c04.rem.21-9-3}
\expandafter\def\csname sga@target@IV.21.9.4\endcsname{ega4.c04.prop.21-9-4}
\expandafter\def\csname sga@target@IV.21.9.4.1\endcsname{ega4.c04.prop.21-9-4-1}
\expandafter\def\csname sga@target@IV.21.9.5\endcsname{ega4.c04.cor.21-9-5}
\expandafter\def\csname sga@target@IV.21.9.6\endcsname{ega4.c04.cor.21-9-6}
\expandafter\def\csname sga@target@IV.21.9.7\endcsname{ega4.c04.cor.21-9-7}
\expandafter\def\csname sga@target@IV.21.9.8\endcsname{ega4.c04.statement.21-9-8}
\expandafter\def\csname sga@target@IV.21.9.9\endcsname{ega4.c04.prop.21-9-9}
\expandafter\def\csname sga@target@IV.3.1.1\endcsname{ega4.c04.def.3-1-1}
\expandafter\def\csname sga@target@IV.3.1.1.1\endcsname{ega4.c04.eq.3-1-1-1}
\expandafter\def\csname sga@target@IV.3.1.10\endcsname{ega4.c04.prop.3-1-10}
\expandafter\def\csname sga@target@IV.3.1.10.1\endcsname{ega4.c04.lem.3-1-10-1}
\expandafter\def\csname sga@target@IV.3.1.10.2\endcsname{ega4.c04.rem.3-1-10-2}
\expandafter\def\csname sga@target@IV.3.1.10.3\endcsname{ega4.c04.eq.3-1-10-3}
\expandafter\def\csname sga@target@IV.3.1.10.4\endcsname{ega4.c04.eq.3-1-10-4}
\expandafter\def\csname sga@target@IV.3.1.10.5\endcsname{ega4.c04.rem.3-1-10-5}
\expandafter\def\csname sga@target@IV.3.1.11\endcsname{ega4.c04.cor.3-1-11}
\expandafter\def\csname sga@target@IV.3.1.12\endcsname{ega4.c04.rem.3-1-12}
\expandafter\def\csname sga@target@IV.3.1.13\endcsname{ega4.c04.prop.3-1-13}
\expandafter\def\csname sga@target@IV.3.1.2\endcsname{ega4.c04.prop.3-1-2}
\expandafter\def\csname sga@target@IV.3.1.3\endcsname{ega4.c04.prop.3-1-3}
\expandafter\def\csname sga@target@IV.3.1.4\endcsname{ega4.c04.cor.3-1-4}
\expandafter\def\csname sga@target@IV.3.1.5\endcsname{ega4.c04.cor.3-1-5}
\expandafter\def\csname sga@target@IV.3.1.6\endcsname{ega4.c04.prop.3-1-6}
\expandafter\def\csname sga@target@IV.3.1.7\endcsname{ega4.c04.prop.3-1-7}
\expandafter\def\csname sga@target@IV.3.1.8\endcsname{ega4.c04.prop.3-1-8}
\expandafter\def\csname sga@target@IV.3.1.9\endcsname{ega4.c04.cor.3-1-9}
\expandafter\def\csname sga@target@IV.3.2.1\endcsname{ega4.c04.prop.3-2-1}
\expandafter\def\csname sga@target@IV.3.2.2\endcsname{ega4.c04.def.3-2-2}
\expandafter\def\csname sga@target@IV.3.2.3\endcsname{ega4.c04.prop.3-2-3}
\expandafter\def\csname sga@target@IV.3.2.4\endcsname{ega4.c04.def.3-2-4}
\expandafter\def\csname sga@target@IV.3.2.5\endcsname{ega4.c04.def.3-2-5}
\expandafter\def\csname sga@target@IV.3.2.6\endcsname{ega4.c04.prop.3-2-6}
\expandafter\def\csname sga@target@IV.3.2.7\endcsname{ega4.c04.cor.3-2-7}
\expandafter\def\csname sga@target@IV.3.2.8\endcsname{ega4.c04.cor.3-2-8}
\expandafter\def\csname sga@target@IV.3.2.9\endcsname{ega4.c04.rem.3-2-9}
\expandafter\def\csname sga@target@IV.3.3.1\endcsname{ega4.c04.prop.3-3-1}
\expandafter\def\csname sga@target@IV.3.3.1.1\endcsname{ega4.c04.prop.3-3-1-1}
\expandafter\def\csname sga@target@IV.3.3.2\endcsname{ega4.c04.cor.3-3-2}
\expandafter\def\csname sga@target@IV.3.3.3\endcsname{ega4.c04.cor.3-3-3}
\expandafter\def\csname sga@target@IV.3.3.4\endcsname{ega4.c04.prop.3-3-4}
\expandafter\def\csname sga@target@IV.3.3.5\endcsname{ega4.c04.cor.3-3-5}
\expandafter\def\csname sga@target@IV.3.3.6\endcsname{ega4.c04.prop.3-3-6}
\expandafter\def\csname sga@target@IV.3.3.6.1\endcsname{ega4.c04.eq.3-3-6-1}
\expandafter\def\csname sga@target@IV.3.3.6.2\endcsname{ega4.c04.eq.3-3-6-2}
\expandafter\def\csname sga@target@IV.3.3.6.3\endcsname{ega4.c04.eq.3-3-6-3}
\expandafter\def\csname sga@target@IV.3.3.6.4\endcsname{ega4.c04.eq.3-3-6-4}
\expandafter\def\csname sga@target@IV.3.3.7\endcsname{ega4.c04.cor.3-3-7}
\expandafter\def\csname sga@target@IV.3.3.7.1\endcsname{ega4.c04.cor.3-3-7-1}
\expandafter\def\csname sga@target@IV.3.3.8\endcsname{ega4.c04.sec.3-3-8}
\expandafter\def\csname sga@target@IV.3.3.9\endcsname{ega4.c04.cor.3-3-9}
\expandafter\def\csname sga@target@IV.3.3.9.1\endcsname{ega4.c04.cor.3-3-9-1}
\expandafter\def\csname sga@target@IV.3.4\endcsname{ega4.c04.subsec.3-4.key-iv-3-4}
\expandafter\def\csname sga@target@IV.3.4.1\endcsname{ega4.c04.prop.3-4-1}
\expandafter\def\csname sga@target@IV.3.4.1.1\endcsname{ega4.c04.eq.3-4-1-1}
\expandafter\def\csname sga@target@IV.3.4.1.2\endcsname{ega4.c04.lem.3-4-1-2}
\expandafter\def\csname sga@target@IV.3.4.1.3\endcsname{ega4.c04.lem.3-4-1-3}
\expandafter\def\csname sga@target@IV.3.4.1.4\endcsname{ega4.c04.lem.3-4-1-4}
\expandafter\def\csname sga@target@IV.3.4.1.5\endcsname{ega4.c04.eq.3-4-1-5}
\expandafter\def\csname sga@target@IV.3.4.2\endcsname{ega4.c04.cor.3-4-2}
\expandafter\def\csname sga@target@IV.3.4.3\endcsname{ega4.c04.prop.3-4-3}
\expandafter\def\csname sga@target@IV.3.4.4\endcsname{ega4.c04.prop.3-4-4}
\expandafter\def\csname sga@target@IV.3.4.5\endcsname{ega4.c04.prop.3-4-5}
\expandafter\def\csname sga@target@IV.3.4.6\endcsname{ega4.c04.prop.3-4-6}
\expandafter\def\csname sga@target@IV.3.4.6.1\endcsname{ega4.c04.lem.3-4-6-1}
\expandafter\def\csname sga@target@IV.3.4.7\endcsname{ega4.c04.cor.3-4-7}
\expandafter\def\csname sga@target@IV.3.4.8\endcsname{ega4.c04.rem.3-4-8}
\expandafter\def\csname sga@target@IV.3.4.9\endcsname{ega4.c04.prop.3-4-9}
\expandafter\def\csname sga@target@IV.4.1.1\endcsname{ega4.c04.def.4-1-1}
\expandafter\def\csname sga@target@IV.4.1.1.1\endcsname{ega4.c04.eq.4-1-1-1}
\expandafter\def\csname sga@target@IV.4.1.1.2\endcsname{ega4.c04.eq.4-1-1-2}
\expandafter\def\csname sga@target@IV.4.1.1.3\endcsname{ega4.c04.eq.4-1-1-3}
\expandafter\def\csname sga@target@IV.4.1.2\endcsname{ega4.c04.thm.4-1-2}
\expandafter\def\csname sga@target@IV.4.1.2.1\endcsname{ega4.c04.cor.4-1-2-1}
\expandafter\def\csname sga@target@IV.4.1.3\endcsname{ega4.c04.rem.4-1-3}
\expandafter\def\csname sga@target@IV.4.1.4\endcsname{ega4.c04.cor.4-1-4}
\expandafter\def\csname sga@target@IV.4.1.5\endcsname{ega4.c04.cor.4-1-5}
\expandafter\def\csname sga@target@IV.4.2.1\endcsname{ega4.c04.prop.4-2-1}
\expandafter\def\csname sga@target@IV.4.2.1.1\endcsname{ega4.c04.eq.4-2-1-1}
\expandafter\def\csname sga@target@IV.4.2.1.2\endcsname{ega4.c04.eq.4-2-1-2}
\expandafter\def\csname sga@target@IV.4.2.1.3\endcsname{ega4.c04.eq.4-2-1-3}
\expandafter\def\csname sga@target@IV.4.2.1.4\endcsname{ega4.c04.rem.4-2-1-4}
\expandafter\def\csname sga@target@IV.4.2.1.5\endcsname{ega4.c04.eq.4-2-1-5}
\expandafter\def\csname sga@target@IV.4.2.10\endcsname{ega4.c04.cor.4-2-10}
\expandafter\def\csname sga@target@IV.4.2.2\endcsname{ega4.c04.cor.4-2-2}
\expandafter\def\csname sga@target@IV.4.2.3\endcsname{ega4.c04.cor.4-2-3}
\expandafter\def\csname sga@target@IV.4.2.4\endcsname{ega4.c04.prop.4-2-4}
\expandafter\def\csname sga@target@IV.4.2.4.i\endcsname{ega4.c04.sec.4-2-4-i}
\expandafter\def\csname sga@target@IV.4.2.4.ii\endcsname{ega4.c04.sec.4-2-4-ii}
\expandafter\def\csname sga@target@IV.4.2.4.iii\endcsname{ega4.c04.sec.4-2-4-iii}
\expandafter\def\csname sga@target@IV.4.2.5\endcsname{ega4.c04.prop.4-2-5}
\expandafter\def\csname sga@target@IV.4.2.6\endcsname{ega4.c04.cor.4-2-6}
\expandafter\def\csname sga@target@IV.4.2.7\endcsname{ega4.c04.prop.4-2-7}
\expandafter\def\csname sga@target@IV.4.2.7.i\endcsname{ega4.c04.sec.4-2-7-i}
\expandafter\def\csname sga@target@IV.4.2.7.ii\endcsname{ega4.c04.sec.4-2-7-ii}
\expandafter\def\csname sga@target@IV.4.2.7.iii\endcsname{ega4.c04.sec.4-2-7-iii}
\expandafter\def\csname sga@target@IV.4.2.7.iv\endcsname{ega4.c04.sec.4-2-7-iv}
\expandafter\def\csname sga@target@IV.4.2.8\endcsname{ega4.c04.cor.4-2-8}
\expandafter\def\csname sga@target@IV.4.2.9\endcsname{ega4.c04.prop.4-2-9}
\expandafter\def\csname sga@target@IV.4.3.1\endcsname{ega4.c04.statement.4-3-1}
\expandafter\def\csname sga@target@IV.4.3.2\endcsname{ega4.c04.prop.4-3-2}
\expandafter\def\csname sga@target@IV.4.3.3\endcsname{ega4.c04.cor.4-3-3}
\expandafter\def\csname sga@target@IV.4.3.4\endcsname{ega4.c04.cor.4-3-4}
\expandafter\def\csname sga@target@IV.4.3.5\endcsname{ega4.c04.prop.4-3-5}
\expandafter\def\csname sga@target@IV.4.3.6\endcsname{ega4.c04.cor.4-3-6}
\expandafter\def\csname sga@target@IV.4.3.7\endcsname{ega4.c04.cor.4-3-7}
\expandafter\def\csname sga@target@IV.4.3.8\endcsname{ega4.c04.cor.4-3-8}
\expandafter\def\csname sga@target@IV.4.3.9\endcsname{ega4.c04.prop.4-3-9}
\expandafter\def\csname sga@target@IV.4.4.1\endcsname{ega4.c04.statement.4-4-1}
\expandafter\def\csname sga@target@IV.4.4.2\endcsname{ega4.c04.lem.4-4-2}
\expandafter\def\csname sga@target@IV.4.4.2.i\endcsname{ega4.c04.sec.4-4-2-i}
\expandafter\def\csname sga@target@IV.4.4.2.ii\endcsname{ega4.c04.sec.4-4-2-ii}
\expandafter\def\csname sga@target@IV.4.4.3\endcsname{ega4.c04.rem.4-4-3}
\expandafter\def\csname sga@target@IV.4.4.4\endcsname{ega4.c04.thm.4-4-4}
\expandafter\def\csname sga@target@IV.4.4.5\endcsname{ega4.c04.cor.4-4-5}
\expandafter\def\csname sga@target@IV.4.4.6\endcsname{ega4.c04.cor.4-4-6}
\expandafter\def\csname sga@target@IV.4.5\endcsname{ega4.c04.subsec.4-5.key-iv-4-5}
\expandafter\def\csname sga@target@IV.4.5.1\endcsname{ega4.c04.prop.4-5-1}
\expandafter\def\csname sga@target@IV.4.5.10\endcsname{ega4.c04.cor.4-5-10}
\expandafter\def\csname sga@target@IV.4.5.11\endcsname{ega4.c04.cor.4-5-11}
\expandafter\def\csname sga@target@IV.4.5.12\endcsname{ega4.c04.sec.4-5-12}
\expandafter\def\csname sga@target@IV.4.5.12.i\endcsname{ega4.c04.sec.4-5-12-i}
\expandafter\def\csname sga@target@IV.4.5.12.ii\endcsname{ega4.c04.sec.4-5-12-ii}
\expandafter\def\csname sga@target@IV.4.5.13\endcsname{ega4.c04.prop.4-5-13}
\expandafter\def\csname sga@target@IV.4.5.13.1\endcsname{ega4.c04.cor.4-5-13-1}
\expandafter\def\csname sga@target@IV.4.5.13.1.i\endcsname{ega4.c04.sec.4-5-13-1-i}
\expandafter\def\csname sga@target@IV.4.5.13.1.ii\endcsname{ega4.c04.sec.4-5-13-1-ii}
\expandafter\def\csname sga@target@IV.4.5.14\endcsname{ega4.c04.cor.4-5-14}
\expandafter\def\csname sga@target@IV.4.5.15\endcsname{ega4.c04.prop.4-5-15}
\expandafter\def\csname sga@target@IV.4.5.15.i\endcsname{ega4.c04.sec.4-5-15-i}
\expandafter\def\csname sga@target@IV.4.5.15.ii\endcsname{ega4.c04.sec.4-5-15-ii}
\expandafter\def\csname sga@target@IV.4.5.16\endcsname{ega4.c04.cor.4-5-16}
\expandafter\def\csname sga@target@IV.4.5.16.i\endcsname{ega4.c04.sec.4-5-16-i}
\expandafter\def\csname sga@target@IV.4.5.16.ii\endcsname{ega4.c04.sec.4-5-16-ii}
\expandafter\def\csname sga@target@IV.4.5.17\endcsname{ega4.c04.cor.4-5-17}
\expandafter\def\csname sga@target@IV.4.5.18\endcsname{ega4.c04.rem.4-5-18}
\expandafter\def\csname sga@target@IV.4.5.19\endcsname{ega4.c04.prop.4-5-19}
\expandafter\def\csname sga@target@IV.4.5.19.1\endcsname{ega4.c04.lem.4-5-19-1}
\expandafter\def\csname sga@target@IV.4.5.19.2\endcsname{ega4.c04.eq.4-5-19-2}
\expandafter\def\csname sga@target@IV.4.5.19.3\endcsname{ega4.c04.cor.4-5-19-3}
\expandafter\def\csname sga@target@IV.4.5.2\endcsname{ega4.c04.def.4-5-2}
\expandafter\def\csname sga@target@IV.4.5.2.a\endcsname{ega4.c04.sec.4-5-2-a}
\expandafter\def\csname sga@target@IV.4.5.2.b\endcsname{ega4.c04.sec.4-5-2-b}
\expandafter\def\csname sga@target@IV.4.5.2.c\endcsname{ega4.c04.sec.4-5-2-c}
\expandafter\def\csname sga@target@IV.4.5.20\endcsname{ega4.c04.rem.4-5-20}
\expandafter\def\csname sga@target@IV.4.5.21\endcsname{ega4.c04.prop.4-5-21}
\expandafter\def\csname sga@target@IV.4.5.3\endcsname{ega4.c04.statement.4-5-3}
\expandafter\def\csname sga@target@IV.4.5.4\endcsname{ega4.c04.prop.4-5-4}
\expandafter\def\csname sga@target@IV.4.5.5\endcsname{ega4.c04.def.4-5-5}
\expandafter\def\csname sga@target@IV.4.5.6\endcsname{ega4.c04.prop.4-5-6}
\expandafter\def\csname sga@target@IV.4.5.6.i\endcsname{ega4.c04.sec.4-5-6-i}
\expandafter\def\csname sga@target@IV.4.5.6.ii\endcsname{ega4.c04.sec.4-5-6-ii}
\expandafter\def\csname sga@target@IV.4.5.7\endcsname{ega4.c04.prop.4-5-7}
\expandafter\def\csname sga@target@IV.4.5.8\endcsname{ega4.c04.cor.4-5-8}
\expandafter\def\csname sga@target@IV.4.5.8.i\endcsname{ega4.c04.sec.4-5-8-i}
\expandafter\def\csname sga@target@IV.4.5.8.ii\endcsname{ega4.c04.sec.4-5-8-ii}
\expandafter\def\csname sga@target@IV.4.5.9\endcsname{ega4.c04.prop.4-5-9}
\expandafter\def\csname sga@target@IV.4.5.9.a\endcsname{ega4.c04.sec.4-5-9-a}
\expandafter\def\csname sga@target@IV.4.5.9.b\endcsname{ega4.c04.sec.4-5-9-b}
\expandafter\def\csname sga@target@IV.4.5.9.c\endcsname{ega4.c04.sec.4-5-9-c}
\expandafter\def\csname sga@target@IV.4.6\endcsname{ega4.c04.subsec.4-6.key-iv-4-6}
\expandafter\def\csname sga@target@IV.4.6.1\endcsname{ega4.c04.prop.4-6-1}
\expandafter\def\csname sga@target@IV.4.6.1.a\endcsname{ega4.c04.sec.4-6-1-a}
\expandafter\def\csname sga@target@IV.4.6.1.b\endcsname{ega4.c04.sec.4-6-1-b}
\expandafter\def\csname sga@target@IV.4.6.1.c\endcsname{ega4.c04.sec.4-6-1-c}
\expandafter\def\csname sga@target@IV.4.6.1.d\endcsname{ega4.c04.sec.4-6-1-d}
\expandafter\def\csname sga@target@IV.4.6.1.e\endcsname{ega4.c04.sec.4-6-1-e}
\expandafter\def\csname sga@target@IV.4.6.10\endcsname{ega4.c04.prop.4-6-10}
\expandafter\def\csname sga@target@IV.4.6.11\endcsname{ega4.c04.cor.4-6-11}
\expandafter\def\csname sga@target@IV.4.6.12\endcsname{ega4.c04.prop.4-6-12}
\expandafter\def\csname sga@target@IV.4.6.12.a\endcsname{ega4.c04.sec.4-6-12-a}
\expandafter\def\csname sga@target@IV.4.6.12.b\endcsname{ega4.c04.sec.4-6-12-b}
\expandafter\def\csname sga@target@IV.4.6.12.c\endcsname{ega4.c04.sec.4-6-12-c}
\expandafter\def\csname sga@target@IV.4.6.12.d\endcsname{ega4.c04.sec.4-6-12-d}
\expandafter\def\csname sga@target@IV.4.6.12.e\endcsname{ega4.c04.sec.4-6-12-e}
\expandafter\def\csname sga@target@IV.4.6.13\endcsname{ega4.c04.cor.4-6-13}
\expandafter\def\csname sga@target@IV.4.6.13.f\endcsname{ega4.c04.sec.4-6-13-f}
\expandafter\def\csname sga@target@IV.4.6.14\endcsname{ega4.c04.statement.4-6-14}
\expandafter\def\csname sga@target@IV.4.6.15\endcsname{ega4.c04.prop.4-6-15}
\expandafter\def\csname sga@target@IV.4.6.15.i\endcsname{ega4.c04.sec.4-6-15-i}
\expandafter\def\csname sga@target@IV.4.6.15.ii\endcsname{ega4.c04.sec.4-6-15-ii}
\expandafter\def\csname sga@target@IV.4.6.16\endcsname{ega4.c04.prop.4-6-16}
\expandafter\def\csname sga@target@IV.4.6.16.a\endcsname{ega4.c04.sec.4-6-16-a}
\expandafter\def\csname sga@target@IV.4.6.16.b\endcsname{ega4.c04.sec.4-6-16-b}
\expandafter\def\csname sga@target@IV.4.6.16.c\endcsname{ega4.c04.sec.4-6-16-c}
\expandafter\def\csname sga@target@IV.4.6.16.d\endcsname{ega4.c04.sec.4-6-16-d}
\expandafter\def\csname sga@target@IV.4.6.17\endcsname{ega4.c04.def.4-6-17}
\expandafter\def\csname sga@target@IV.4.6.18\endcsname{ega4.c04.prop.4-6-18}
\expandafter\def\csname sga@target@IV.4.6.18.a\endcsname{ega4.c04.sec.4-6-18-a}
\expandafter\def\csname sga@target@IV.4.6.18.b\endcsname{ega4.c04.sec.4-6-18-b}
\expandafter\def\csname sga@target@IV.4.6.18.c\endcsname{ega4.c04.sec.4-6-18-c}
\expandafter\def\csname sga@target@IV.4.6.18.d\endcsname{ega4.c04.sec.4-6-18-d}
\expandafter\def\csname sga@target@IV.4.6.19\endcsname{ega4.c04.def.4-6-19}
\expandafter\def\csname sga@target@IV.4.6.2\endcsname{ega4.c04.def.4-6-2}
\expandafter\def\csname sga@target@IV.4.6.20\endcsname{ega4.c04.prop.4-6-20}
\expandafter\def\csname sga@target@IV.4.6.21\endcsname{ega4.c04.prop.4-6-21}
\expandafter\def\csname sga@target@IV.4.6.21.i\endcsname{ega4.c04.sec.4-6-21-i}
\expandafter\def\csname sga@target@IV.4.6.21.ii\endcsname{ega4.c04.sec.4-6-21-ii}
\expandafter\def\csname sga@target@IV.4.6.22\endcsname{ega4.c04.def.4-6-22}
\expandafter\def\csname sga@target@IV.4.6.3\endcsname{ega4.c04.cor.4-6-3}
\expandafter\def\csname sga@target@IV.4.6.4\endcsname{ega4.c04.cor.4-6-4}
\expandafter\def\csname sga@target@IV.4.6.5\endcsname{ega4.c04.prop.4-6-5}
\expandafter\def\csname sga@target@IV.4.6.5.i\endcsname{ega4.c04.sec.4-6-5-i}
\expandafter\def\csname sga@target@IV.4.6.5.ii\endcsname{ega4.c04.sec.4-6-5-ii}
\expandafter\def\csname sga@target@IV.4.6.6\endcsname{ega4.c04.prop.4-6-6}
\expandafter\def\csname sga@target@IV.4.6.7\endcsname{ega4.c04.cor.4-6-7}
\expandafter\def\csname sga@target@IV.4.6.8\endcsname{ega4.c04.cor.4-6-8}
\expandafter\def\csname sga@target@IV.4.6.9\endcsname{ega4.c04.def.4-6-9}
\expandafter\def\csname sga@target@IV.4.7.1\endcsname{ega4.c04.lem.4-7-1}
\expandafter\def\csname sga@target@IV.4.7.1.1\endcsname{ega4.c04.eq.4-7-1-1}
\expandafter\def\csname sga@target@IV.4.7.10\endcsname{ega4.c04.prop.4-7-10}
\expandafter\def\csname sga@target@IV.4.7.11\endcsname{ega4.c04.prop.4-7-11}
\expandafter\def\csname sga@target@IV.4.7.12\endcsname{ega4.c04.def.4-7-12}
\expandafter\def\csname sga@target@IV.4.7.13\endcsname{ega4.c04.prop.4-7-13}
\expandafter\def\csname sga@target@IV.4.7.2\endcsname{ega4.c04.cor.4-7-2}
\expandafter\def\csname sga@target@IV.4.7.2.1\endcsname{ega4.c04.eq.4-7-2-1}
\expandafter\def\csname sga@target@IV.4.7.3\endcsname{ega4.c04.prop.4-7-3}
\expandafter\def\csname sga@target@IV.4.7.3.1\endcsname{ega4.c04.eq.4-7-3-1}
\expandafter\def\csname sga@target@IV.4.7.3.i\endcsname{ega4.c04.sec.4-7-3-i}
\expandafter\def\csname sga@target@IV.4.7.3.ii\endcsname{ega4.c04.sec.4-7-3-ii}
\expandafter\def\csname sga@target@IV.4.7.3.iii\endcsname{ega4.c04.sec.4-7-3-iii}
\expandafter\def\csname sga@target@IV.4.7.4\endcsname{ega4.c04.def.4-7-4}
\expandafter\def\csname sga@target@IV.4.7.5\endcsname{ega4.c04.def.4-7-5}
\expandafter\def\csname sga@target@IV.4.7.6\endcsname{ega4.c04.prop.4-7-6}
\expandafter\def\csname sga@target@IV.4.7.6.1\endcsname{ega4.c04.eq.4-7-6-1}
\expandafter\def\csname sga@target@IV.4.7.7\endcsname{ega4.c04.statement.4-7-7}
\expandafter\def\csname sga@target@IV.4.7.8\endcsname{ega4.c04.prop.4-7-8}
\expandafter\def\csname sga@target@IV.4.7.9\endcsname{ega4.c04.cor.4-7-9}
\expandafter\def\csname sga@target@IV.4.8.1\endcsname{ega4.c04.statement.4-8-1}
\expandafter\def\csname sga@target@IV.4.8.10\endcsname{ega4.c04.cor.4-8-10}
\expandafter\def\csname sga@target@IV.4.8.11\endcsname{ega4.c04.cor.4-8-11}
\expandafter\def\csname sga@target@IV.4.8.12\endcsname{ega4.c04.cor.4-8-12}
\expandafter\def\csname sga@target@IV.4.8.13\endcsname{ega4.c04.prop.4-8-13}
\expandafter\def\csname sga@target@IV.4.8.14\endcsname{ega4.c04.prop.4-8-14}
\expandafter\def\csname sga@target@IV.4.8.15\endcsname{ega4.c04.sec.4-8-15}
\expandafter\def\csname sga@target@IV.4.8.15.i\endcsname{ega4.c04.sec.4-8-15-i}
\expandafter\def\csname sga@target@IV.4.8.15.ii\endcsname{ega4.c04.sec.4-8-15-ii}
\expandafter\def\csname sga@target@IV.4.8.2\endcsname{ega4.c04.statement.4-8-2}
\expandafter\def\csname sga@target@IV.4.8.2.1\endcsname{ega4.c04.eq.4-8-2-1}
\expandafter\def\csname sga@target@IV.4.8.2.2\endcsname{ega4.c04.eq.4-8-2-2}
\expandafter\def\csname sga@target@IV.4.8.2.3\endcsname{ega4.c04.eq.4-8-2-3}
\expandafter\def\csname sga@target@IV.4.8.2.4\endcsname{ega4.c04.eq.4-8-2-4}
\expandafter\def\csname sga@target@IV.4.8.2.5\endcsname{ega4.c04.eq.4-8-2-5}
\expandafter\def\csname sga@target@IV.4.8.2.6\endcsname{ega4.c04.eq.4-8-2-6}
\expandafter\def\csname sga@target@IV.4.8.3\endcsname{ega4.c04.prop.4-8-3}
\expandafter\def\csname sga@target@IV.4.8.4\endcsname{ega4.c04.def.4-8-4}
\expandafter\def\csname sga@target@IV.4.8.5\endcsname{ega4.c04.prop.4-8-5}
\expandafter\def\csname sga@target@IV.4.8.6\endcsname{ega4.c04.lem.4-8-6}
\expandafter\def\csname sga@target@IV.4.8.6.a\endcsname{ega4.c04.sec.4-8-6-a}
\expandafter\def\csname sga@target@IV.4.8.6.b\endcsname{ega4.c04.sec.4-8-6-b}
\expandafter\def\csname sga@target@IV.4.8.6.c\endcsname{ega4.c04.sec.4-8-6-c}
\expandafter\def\csname sga@target@IV.4.8.7\endcsname{ega4.c04.cor.4-8-7}
\expandafter\def\csname sga@target@IV.4.8.8\endcsname{ega4.c04.lem.4-8-8}
\expandafter\def\csname sga@target@IV.4.8.8.i\endcsname{ega4.c04.sec.4-8-8-i}
\expandafter\def\csname sga@target@IV.4.8.8.ii\endcsname{ega4.c04.sec.4-8-8-ii}
\expandafter\def\csname sga@target@IV.4.8.8.iii\endcsname{ega4.c04.sec.4-8-8-iii}
\expandafter\def\csname sga@target@IV.4.8.9\endcsname{ega4.c04.prop.4-8-9}
\expandafter\def\csname sga@target@IV.4.9.1\endcsname{ega4.c04.prop.4-9-1}
\expandafter\def\csname sga@target@IV.4.9.1.a\endcsname{ega4.c04.sec.4-9-1-a}
\expandafter\def\csname sga@target@IV.4.9.1.b\endcsname{ega4.c04.sec.4-9-1-b}
\expandafter\def\csname sga@target@IV.4.9.1.bprime\endcsname{ega4.c04.sec.4-9-1-bprime}
\expandafter\def\csname sga@target@IV.4.9.1.c\endcsname{ega4.c04.sec.4-9-1-c}
\expandafter\def\csname sga@target@IV.4.9.2\endcsname{ega4.c04.cor.4-9-2}
\expandafter\def\csname sga@target@IV.4.9.3\endcsname{ega4.c04.rem.4-9-3}
\expandafter\def\csname sga@target@IV.4.9.4\endcsname{ega4.c04.prop.4-9-4}
\expandafter\def\csname sga@target@IV.4.9.5\endcsname{ega4.c04.cor.4-9-5}
\expandafter\def\csname sga@target@IV.4.9.5.a\endcsname{ega4.c04.sec.4-9-5-a}
\expandafter\def\csname sga@target@IV.4.9.5.b\endcsname{ega4.c04.sec.4-9-5-b}
\expandafter\def\csname sga@target@IV.4.9.6\endcsname{ega4.c04.prop.4-9-6}
\expandafter\def\csname sga@target@IV.4.9.6.a\endcsname{ega4.c04.sec.4-9-6-a}
\expandafter\def\csname sga@target@IV.4.9.6.aprime\endcsname{ega4.c04.sec.4-9-6-aprime}
\expandafter\def\csname sga@target@IV.4.9.6.b\endcsname{ega4.c04.sec.4-9-6-b}
\expandafter\def\csname sga@target@IV.4.9.6.bprime\endcsname{ega4.c04.sec.4-9-6-bprime}
\expandafter\def\csname sga@target@IV.4.9.7\endcsname{ega4.c04.cor.4-9-7}
\expandafter\def\csname sga@target@IV.5.1.1\endcsname{ega4.c04.statement.5-1-1}
\expandafter\def\csname sga@target@IV.5.1.1.1\endcsname{ega4.c04.eq.5-1-1-1}
\expandafter\def\csname sga@target@IV.5.1.1.2\endcsname{ega4.c04.eq.5-1-1-2}
\expandafter\def\csname sga@target@IV.5.1.1.3\endcsname{ega4.c04.eq.5-1-1-3}
\expandafter\def\csname sga@target@IV.5.1.1.4\endcsname{ega4.c04.sec.5-1-1-4}
\expandafter\def\csname sga@target@IV.5.1.1.5\endcsname{ega4.c04.sec.5-1-1-5}
\expandafter\def\csname sga@target@IV.5.1.1.6\endcsname{ega4.c04.sec.5-1-1-6}
\expandafter\def\csname sga@target@IV.5.1.10\endcsname{ega4.c04.prop.5-1-10}
\expandafter\def\csname sga@target@IV.5.1.10.1\endcsname{ega4.c04.lem.5-1-10-1}
\expandafter\def\csname sga@target@IV.5.1.10.i\endcsname{ega4.c04.sec.5-1-10-i}
\expandafter\def\csname sga@target@IV.5.1.10.ii\endcsname{ega4.c04.sec.5-1-10-ii}
\expandafter\def\csname sga@target@IV.5.1.11\endcsname{ega4.c04.cor.5-1-11}
\expandafter\def\csname sga@target@IV.5.1.12\endcsname{ega4.c04.statement.5-1-12}
\expandafter\def\csname sga@target@IV.5.1.12.1\endcsname{ega4.c04.eq.5-1-12-1}
\expandafter\def\csname sga@target@IV.5.1.12.2\endcsname{ega4.c04.eq.5-1-12-2}
\expandafter\def\csname sga@target@IV.5.1.12.3\endcsname{ega4.c04.eq.5-1-12-3}
\expandafter\def\csname sga@target@IV.5.1.13\endcsname{ega4.c04.prop.5-1-13}
\expandafter\def\csname sga@target@IV.5.1.13.1\endcsname{ega4.c04.eq.5-1-13-1}
\expandafter\def\csname sga@target@IV.5.1.2\endcsname{ega4.c04.prop.5-1-2}
\expandafter\def\csname sga@target@IV.5.1.2.1\endcsname{ega4.c04.eq.5-1-2-1}
\expandafter\def\csname sga@target@IV.5.1.3\endcsname{ega4.c04.cor.5-1-3}
\expandafter\def\csname sga@target@IV.5.1.3.1\endcsname{ega4.c04.eq.5-1-3-1}
\expandafter\def\csname sga@target@IV.5.1.3.2\endcsname{ega4.c04.eq.5-1-3-2}
\expandafter\def\csname sga@target@IV.5.1.4\endcsname{ega4.c04.prop.5-1-4}
\expandafter\def\csname sga@target@IV.5.1.4.1\endcsname{ega4.c04.eq.5-1-4-1}
\expandafter\def\csname sga@target@IV.5.1.4.2\endcsname{ega4.c04.eq.5-1-4-2}
\expandafter\def\csname sga@target@IV.5.1.5\endcsname{ega4.c04.cor.5-1-5}
\expandafter\def\csname sga@target@IV.5.1.6\endcsname{ega4.c04.statement.5-1-6}
\expandafter\def\csname sga@target@IV.5.1.7\endcsname{ega4.c04.prop.5-1-7}
\expandafter\def\csname sga@target@IV.5.1.8\endcsname{ega4.c04.cor.5-1-8}
\expandafter\def\csname sga@target@IV.5.1.9\endcsname{ega4.c04.prop.5-1-9}
\expandafter\def\csname sga@target@IV.5.1.9.1\endcsname{ega4.c04.statement.5-1-9-1}
\expandafter\def\csname sga@target@IV.5.10.1\endcsname{ega4.c04.statement.5-10-1}
\expandafter\def\csname sga@target@IV.5.10.1.1\endcsname{ega4.c04.eq.5-10-1-1}
\expandafter\def\csname sga@target@IV.5.10.10\endcsname{ega4.c04.prop.5-10-10}
\expandafter\def\csname sga@target@IV.5.10.11\endcsname{ega4.c04.statement.5-10-11}
\expandafter\def\csname sga@target@IV.5.10.11.1\endcsname{ega4.c04.eq.5-10-11-1}
\expandafter\def\csname sga@target@IV.5.10.12\endcsname{ega4.c04.prop.5-10-12}
\expandafter\def\csname sga@target@IV.5.10.13\endcsname{ega4.c04.statement.5-10-13}
\expandafter\def\csname sga@target@IV.5.10.14\endcsname{ega4.c04.prop.5-10-14}
\expandafter\def\csname sga@target@IV.5.10.15\endcsname{ega4.c04.cor.5-10-15}
\expandafter\def\csname sga@target@IV.5.10.16\endcsname{ega4.c04.prop.5-10-16}
\expandafter\def\csname sga@target@IV.5.10.17\endcsname{ega4.c04.prop.5-10-17}
\expandafter\def\csname sga@target@IV.5.10.2\endcsname{ega4.c04.prop.5-10-2}
\expandafter\def\csname sga@target@IV.5.10.3\endcsname{ega4.c04.cor.5-10-3}
\expandafter\def\csname sga@target@IV.5.10.4\endcsname{ega4.c04.cor.5-10-4}
\expandafter\def\csname sga@target@IV.5.10.5\endcsname{ega4.c04.thm.5-10-5}
\expandafter\def\csname sga@target@IV.5.10.6\endcsname{ega4.c04.cor.5-10-6}
\expandafter\def\csname sga@target@IV.5.10.7\endcsname{ega4.c04.thm.5-10-7}
\expandafter\def\csname sga@target@IV.5.10.8\endcsname{ega4.c04.cor.5-10-8}
\expandafter\def\csname sga@target@IV.5.10.8.1\endcsname{ega4.c04.lem.5-10-8-1}
\expandafter\def\csname sga@target@IV.5.10.9\endcsname{ega4.c04.cor.5-10-9}
\expandafter\def\csname sga@target@IV.5.11.1\endcsname{ega4.c04.prop.5-11-1}
\expandafter\def\csname sga@target@IV.5.11.1.1\endcsname{ega4.c04.lem.5-11-1-1}
\expandafter\def\csname sga@target@IV.5.11.2\endcsname{ega4.c04.cor.5-11-2}
\expandafter\def\csname sga@target@IV.5.11.3\endcsname{ega4.c04.rem.5-11-3}
\expandafter\def\csname sga@target@IV.5.11.4\endcsname{ega4.c04.cor.5-11-4}
\expandafter\def\csname sga@target@IV.5.11.5\endcsname{ega4.c04.cor.5-11-5}
\expandafter\def\csname sga@target@IV.5.11.6\endcsname{ega4.c04.cor.5-11-6}
\expandafter\def\csname sga@target@IV.5.11.7\endcsname{ega4.c04.item.5-11-7}
\expandafter\def\csname sga@target@IV.5.12\endcsname{ega4.c04.subsec.5-12.key-iv-5-12}
\expandafter\def\csname sga@target@IV.5.12.1\endcsname{ega4.c04.prop.5-12-1}
\expandafter\def\csname sga@target@IV.5.12.1.1\endcsname{ega4.c04.eq.5-12-1-1}
\expandafter\def\csname sga@target@IV.5.12.10\endcsname{ega4.c04.cor.5-12-10}
\expandafter\def\csname sga@target@IV.5.12.2\endcsname{ega4.c04.prop.5-12-2}
\expandafter\def\csname sga@target@IV.5.12.3\endcsname{ega4.c04.rem.5-12-3}
\expandafter\def\csname sga@target@IV.5.12.4\endcsname{ega4.c04.cor.5-12-4}
\expandafter\def\csname sga@target@IV.5.12.5\endcsname{ega4.c04.prop.5-12-5}
\expandafter\def\csname sga@target@IV.5.12.6\endcsname{ega4.c04.rem.5-12-6}
\expandafter\def\csname sga@target@IV.5.12.7\endcsname{ega4.c04.cor.5-12-7}
\expandafter\def\csname sga@target@IV.5.12.8\endcsname{ega4.c04.prop.5-12-8}
\expandafter\def\csname sga@target@IV.5.12.9\endcsname{ega4.c04.rem.5-12-9}
\expandafter\def\csname sga@target@IV.5.13.1\endcsname{ega4.c04.prop.5-13-1}
\expandafter\def\csname sga@target@IV.5.13.2\endcsname{ega4.c04.cor.5-13-2}
\expandafter\def\csname sga@target@IV.5.13.3\endcsname{ega4.c04.prop.5-13-3}
\expandafter\def\csname sga@target@IV.5.13.4\endcsname{ega4.c04.cor.5-13-4}
\expandafter\def\csname sga@target@IV.5.13.5\endcsname{ega4.c04.statement.5-13-5}
\expandafter\def\csname sga@target@IV.5.13.6\endcsname{ega4.c04.cor.5-13-6}
\expandafter\def\csname sga@target@IV.5.13.7\endcsname{ega4.c04.prop.5-13-7}
\expandafter\def\csname sga@target@IV.5.13.7.1\endcsname{ega4.c04.lem.5-13-7-1}
\expandafter\def\csname sga@target@IV.5.13.8\endcsname{ega4.c04.cor.5-13-8}
\expandafter\def\csname sga@target@IV.5.2.1\endcsname{ega4.c04.prop.5-2-1}
\expandafter\def\csname sga@target@IV.5.2.1.1\endcsname{ega4.c04.eq.5-2-1-1}
\expandafter\def\csname sga@target@IV.5.2.2\endcsname{ega4.c04.cor.5-2-2}
\expandafter\def\csname sga@target@IV.5.2.3\endcsname{ega4.c04.cor.5-2-3}
\expandafter\def\csname sga@target@IV.5.2.3.1\endcsname{ega4.c04.eq.5-2-3-1}
\expandafter\def\csname sga@target@IV.5.2.4\endcsname{ega4.c04.cor.5-2-4}
\expandafter\def\csname sga@target@IV.5.2.5\endcsname{ega4.c04.sec.5-2-5}
\expandafter\def\csname sga@target@IV.5.2.5.i\endcsname{ega4.c04.sec.5-2-5-i}
\expandafter\def\csname sga@target@IV.5.2.5.ii\endcsname{ega4.c04.sec.5-2-5-ii}
\expandafter\def\csname sga@target@IV.5.3.1\endcsname{ega4.c04.prop.5-3-1}
\expandafter\def\csname sga@target@IV.5.3.1.1\endcsname{ega4.c04.eq.5-3-1-1}
\expandafter\def\csname sga@target@IV.5.3.1.2\endcsname{ega4.c04.eq.5-3-1-2}
\expandafter\def\csname sga@target@IV.5.3.1.3\endcsname{ega4.c04.lem.5-3-1-3}
\expandafter\def\csname sga@target@IV.5.4.1\endcsname{ega4.c04.prop.5-4-1}
\expandafter\def\csname sga@target@IV.5.4.1.i\endcsname{ega4.c04.sec.5-4-1-i}
\expandafter\def\csname sga@target@IV.5.4.1.ii\endcsname{ega4.c04.sec.5-4-1-ii}
\expandafter\def\csname sga@target@IV.5.4.2\endcsname{ega4.c04.cor.5-4-2}
\expandafter\def\csname sga@target@IV.5.4.3\endcsname{ega4.c04.sec.5-4-3}
\expandafter\def\csname sga@target@IV.5.4.3.1\endcsname{ega4.c04.eq.5-4-3-1}
\expandafter\def\csname sga@target@IV.5.4.3.i\endcsname{ega4.c04.sec.5-4-3-i}
\expandafter\def\csname sga@target@IV.5.4.3.ii\endcsname{ega4.c04.sec.5-4-3-ii}
\expandafter\def\csname sga@target@IV.5.5\endcsname{ega4.c04.sec.5-5}
\expandafter\def\csname sga@target@IV.5.5.1\endcsname{ega4.c04.statement.5-5-1}
\expandafter\def\csname sga@target@IV.5.5.1.1\endcsname{ega4.c04.eq.5-5-1-1}
\expandafter\def\csname sga@target@IV.5.5.2\endcsname{ega4.c04.prop.5-5-2}
\expandafter\def\csname sga@target@IV.5.5.2.1\endcsname{ega4.c04.eq.5-5-2-1}
\expandafter\def\csname sga@target@IV.5.5.2.2\endcsname{ega4.c04.eq.5-5-2-2}
\expandafter\def\csname sga@target@IV.5.5.3\endcsname{ega4.c04.prop.5-5-3}
\expandafter\def\csname sga@target@IV.5.5.3.1\endcsname{ega4.c04.eq.5-5-3-1}
\expandafter\def\csname sga@target@IV.5.5.4\endcsname{ega4.c04.cor.5-5-4}
\expandafter\def\csname sga@target@IV.5.5.4.1\endcsname{ega4.c04.eq.5-5-4-1}
\expandafter\def\csname sga@target@IV.5.5.5\endcsname{ega4.c04.cor.5-5-5}
\expandafter\def\csname sga@target@IV.5.5.6\endcsname{ega4.c04.cor.5-5-6}
\expandafter\def\csname sga@target@IV.5.5.6.1\endcsname{ega4.c04.eq.5-5-6-1}
\expandafter\def\csname sga@target@IV.5.5.7\endcsname{ega4.c04.lem.5-5-7}
\expandafter\def\csname sga@target@IV.5.5.7.1\endcsname{ega4.c04.eq.5-5-7-1}
\expandafter\def\csname sga@target@IV.5.5.7.2\endcsname{ega4.c04.statement.5-5-7-2}
\expandafter\def\csname sga@target@IV.5.5.7.3\endcsname{ega4.c04.eq.5-5-7-3}
\expandafter\def\csname sga@target@IV.5.5.7.i\endcsname{ega4.c04.sec.5-5-7-i}
\expandafter\def\csname sga@target@IV.5.5.7.ii\endcsname{ega4.c04.sec.5-5-7-ii}
\expandafter\def\csname sga@target@IV.5.5.8\endcsname{ega4.c04.thm.5-5-8}
\expandafter\def\csname sga@target@IV.5.5.8.1\endcsname{ega4.c04.eq.5-5-8-1}
\expandafter\def\csname sga@target@IV.5.5.8.2\endcsname{ega4.c04.eq.5-5-8-2}
\expandafter\def\csname sga@target@IV.5.5.8.3\endcsname{ega4.c04.eq.5-5-8-3}
\expandafter\def\csname sga@target@IV.5.6\endcsname{ega4.c04.sec.5-6}
\expandafter\def\csname sga@target@IV.5.6.1\endcsname{ega4.c04.prop.5-6-1}
\expandafter\def\csname sga@target@IV.5.6.1.1\endcsname{ega4.c04.eq.5-6-1-1}
\expandafter\def\csname sga@target@IV.5.6.1.2\endcsname{ega4.c04.eq.5-6-1-2}
\expandafter\def\csname sga@target@IV.5.6.1.3\endcsname{ega4.c04.statement.5-6-1-3}
\expandafter\def\csname sga@target@IV.5.6.1.4\endcsname{ega4.c04.statement.5-6-1-4}
\expandafter\def\csname sga@target@IV.5.6.1.5\endcsname{ega4.c04.statement.5-6-1-5}
\expandafter\def\csname sga@target@IV.5.6.1.a\endcsname{ega4.c04.sec.5-6-1-a}
\expandafter\def\csname sga@target@IV.5.6.1.b\endcsname{ega4.c04.sec.5-6-1-b}
\expandafter\def\csname sga@target@IV.5.6.1.c\endcsname{ega4.c04.sec.5-6-1-c}
\expandafter\def\csname sga@target@IV.5.6.10\endcsname{ega4.c04.prop.5-6-10}
\expandafter\def\csname sga@target@IV.5.6.11\endcsname{ega4.c04.ex.5-6-11}
\expandafter\def\csname sga@target@IV.5.6.11.1\endcsname{ega4.c04.lem.5-6-11-1}
\expandafter\def\csname sga@target@IV.5.6.11.1.i\endcsname{ega4.c04.sec.5-6-11-1-i}
\expandafter\def\csname sga@target@IV.5.6.11.1.ii\endcsname{ega4.c04.sec.5-6-11-1-ii}
\expandafter\def\csname sga@target@IV.5.6.2\endcsname{ega4.c04.def.5-6-2}
\expandafter\def\csname sga@target@IV.5.6.3\endcsname{ega4.c04.sec.5-6-3}
\expandafter\def\csname sga@target@IV.5.6.3.i\endcsname{ega4.c04.sec.5-6-3-i}
\expandafter\def\csname sga@target@IV.5.6.3.ii\endcsname{ega4.c04.sec.5-6-3-ii}
\expandafter\def\csname sga@target@IV.5.6.3.iii\endcsname{ega4.c04.sec.5-6-3-iii}
\expandafter\def\csname sga@target@IV.5.6.3.iv\endcsname{ega4.c04.sec.5-6-3-iv}
\expandafter\def\csname sga@target@IV.5.6.4\endcsname{ega4.c04.prop.5-6-4}
\expandafter\def\csname sga@target@IV.5.6.5\endcsname{ega4.c04.prop.5-6-5}
\expandafter\def\csname sga@target@IV.5.6.5.1\endcsname{ega4.c04.eq.5-6-5-1}
\expandafter\def\csname sga@target@IV.5.6.5.2\endcsname{ega4.c04.eq.5-6-5-2}
\expandafter\def\csname sga@target@IV.5.6.5.3\endcsname{ega4.c04.eq.5-6-5-3}
\expandafter\def\csname sga@target@IV.5.6.6\endcsname{ega4.c04.cor.5-6-6}
\expandafter\def\csname sga@target@IV.5.6.6.1\endcsname{ega4.c04.eq.5-6-6-1}
\expandafter\def\csname sga@target@IV.5.6.6.2\endcsname{ega4.c04.eq.5-6-6-2}
\expandafter\def\csname sga@target@IV.5.6.7\endcsname{ega4.c04.cor.5-6-7}
\expandafter\def\csname sga@target@IV.5.6.7.1\endcsname{ega4.c04.eq.5-6-7-1}
\expandafter\def\csname sga@target@IV.5.6.8\endcsname{ega4.c04.cor.5-6-8}
\expandafter\def\csname sga@target@IV.5.6.8.1\endcsname{ega4.c04.eq.5-6-8-1}
\expandafter\def\csname sga@target@IV.5.6.8.2\endcsname{ega4.c04.eq.5-6-8-2}
\expandafter\def\csname sga@target@IV.5.6.9\endcsname{ega4.c04.sec.5-6-9}
\expandafter\def\csname sga@target@IV.5.6.9.i\endcsname{ega4.c04.sec.5-6-9-i}
\expandafter\def\csname sga@target@IV.5.6.9.ii\endcsname{ega4.c04.sec.5-6-9-ii}
\expandafter\def\csname sga@target@IV.5.7.1\endcsname{ega4.c04.def.5-7-1}
\expandafter\def\csname sga@target@IV.5.7.1.1\endcsname{ega4.c04.eq.5-7-1-1}
\expandafter\def\csname sga@target@IV.5.7.10\endcsname{ega4.c04.cor.5-7-10}
\expandafter\def\csname sga@target@IV.5.7.11\endcsname{ega4.c04.cor.5-7-11}
\expandafter\def\csname sga@target@IV.5.7.11.i\endcsname{ega4.c04.sec.5-7-11-i}
\expandafter\def\csname sga@target@IV.5.7.11.ii\endcsname{ega4.c04.sec.5-7-11-ii}
\expandafter\def\csname sga@target@IV.5.7.12\endcsname{ega4.c04.statement.5-7-12}
\expandafter\def\csname sga@target@IV.5.7.13\endcsname{ega4.c04.cor.5-7-13}
\expandafter\def\csname sga@target@IV.5.7.13.1\endcsname{ega4.c04.eq.5-7-13-1}
\expandafter\def\csname sga@target@IV.5.7.2\endcsname{ega4.c04.def.5-7-2}
\expandafter\def\csname sga@target@IV.5.7.2.1\endcsname{ega4.c04.eq.5-7-2-1}
\expandafter\def\csname sga@target@IV.5.7.2.2\endcsname{ega4.c04.eq.5-7-2-2}
\expandafter\def\csname sga@target@IV.5.7.3\endcsname{ega4.c04.sec.5-7-3}
\expandafter\def\csname sga@target@IV.5.7.3.i\endcsname{ega4.c04.sec.5-7-3-i}
\expandafter\def\csname sga@target@IV.5.7.3.ii\endcsname{ega4.c04.sec.5-7-3-ii}
\expandafter\def\csname sga@target@IV.5.7.3.iii\endcsname{ega4.c04.sec.5-7-3-iii}
\expandafter\def\csname sga@target@IV.5.7.3.iv\endcsname{ega4.c04.sec.5-7-3-iv}
\expandafter\def\csname sga@target@IV.5.7.3.v\endcsname{ega4.c04.sec.5-7-3-v}
\expandafter\def\csname sga@target@IV.5.7.3.vi\endcsname{ega4.c04.sec.5-7-3-vi}
\expandafter\def\csname sga@target@IV.5.7.4\endcsname{ega4.c04.prop.5-7-4}
\expandafter\def\csname sga@target@IV.5.7.4.1\endcsname{ega4.c04.eq.5-7-4-1}
\expandafter\def\csname sga@target@IV.5.7.4.2\endcsname{ega4.c04.eq.5-7-4-2}
\expandafter\def\csname sga@target@IV.5.7.4.i\endcsname{ega4.c04.sec.5-7-4-i}
\expandafter\def\csname sga@target@IV.5.7.4.ii\endcsname{ega4.c04.sec.5-7-4-ii}
\expandafter\def\csname sga@target@IV.5.7.5\endcsname{ega4.c04.prop.5-7-5}
\expandafter\def\csname sga@target@IV.5.7.6\endcsname{ega4.c04.cor.5-7-6}
\expandafter\def\csname sga@target@IV.5.7.7\endcsname{ega4.c04.cor.5-7-7}
\expandafter\def\csname sga@target@IV.5.7.8\endcsname{ega4.c04.sec.5-7-8}
\expandafter\def\csname sga@target@IV.5.7.9\endcsname{ega4.c04.prop.5-7-9}
\expandafter\def\csname sga@target@IV.5.7.9.1\endcsname{ega4.c04.eq.5-7-9-1}
\expandafter\def\csname sga@target@IV.5.7.9.2\endcsname{ega4.c04.eq.5-7-9-2}
\expandafter\def\csname sga@target@IV.5.8.1\endcsname{ega4.c04.statement.5-8-1}
\expandafter\def\csname sga@target@IV.5.8.2\endcsname{ega4.c04.def.5-8-2}
\expandafter\def\csname sga@target@IV.5.8.3\endcsname{ega4.c04.prop.5-8-3}
\expandafter\def\csname sga@target@IV.5.8.4\endcsname{ega4.c04.statement.5-8-4}
\expandafter\def\csname sga@target@IV.5.8.5\endcsname{ega4.c04.prop.5-8-5}
\expandafter\def\csname sga@target@IV.5.8.6\endcsname{ega4.c04.thm.5-8-6}
\expandafter\def\csname sga@target@IV.5.8.6.1\endcsname{ega4.c04.eq.5-8-6-1}
\expandafter\def\csname sga@target@IV.5.8.6.i\endcsname{ega4.c04.item.5-8-6-i}
\expandafter\def\csname sga@target@IV.5.8.6.ii\endcsname{ega4.c04.item.5-8-6-ii}
\expandafter\def\csname sga@target@IV.5.9.1\endcsname{ega4.c04.statement.5-9-1}
\expandafter\def\csname sga@target@IV.5.9.1.1\endcsname{ega4.c04.eq.5-9-1-1}
\expandafter\def\csname sga@target@IV.5.9.1.2\endcsname{ega4.c04.eq.5-9-1-2}
\expandafter\def\csname sga@target@IV.5.9.1.3\endcsname{ega4.c04.eq.5-9-1-3}
\expandafter\def\csname sga@target@IV.5.9.10\endcsname{ega4.c04.cor.5-9-10}
\expandafter\def\csname sga@target@IV.5.9.10.i\endcsname{ega4.c04.item.5-9-10-i}
\expandafter\def\csname sga@target@IV.5.9.10.ii\endcsname{ega4.c04.item.5-9-10-ii}
\expandafter\def\csname sga@target@IV.5.9.11\endcsname{ega4.c04.prop.5-9-11}
\expandafter\def\csname sga@target@IV.5.9.12\endcsname{ega4.c04.item.5-9-12}
\expandafter\def\csname sga@target@IV.5.9.12.i\endcsname{ega4.c04.item.5-9-12-i}
\expandafter\def\csname sga@target@IV.5.9.12.ii\endcsname{ega4.c04.item.5-9-12-ii}
\expandafter\def\csname sga@target@IV.5.9.12.iii\endcsname{ega4.c04.item.5-9-12-iii}
\expandafter\def\csname sga@target@IV.5.9.2\endcsname{ega4.c04.prop.5-9-2}
\expandafter\def\csname sga@target@IV.5.9.2.i\endcsname{ega4.c04.item.5-9-2-i}
\expandafter\def\csname sga@target@IV.5.9.2.ii\endcsname{ega4.c04.item.5-9-2-ii}
\expandafter\def\csname sga@target@IV.5.9.3\endcsname{ega4.c04.rem.5-9-3}
\expandafter\def\csname sga@target@IV.5.9.3.1\endcsname{ega4.c04.eq.5-9-3-1}
\expandafter\def\csname sga@target@IV.5.9.4\endcsname{ega4.c04.prop.5-9-4}
\expandafter\def\csname sga@target@IV.5.9.4.1\endcsname{ega4.c04.eq.5-9-4-1}
\expandafter\def\csname sga@target@IV.5.9.5\endcsname{ega4.c04.cor.5-9-5}
\expandafter\def\csname sga@target@IV.5.9.6\endcsname{ega4.c04.cor.5-9-6}
\expandafter\def\csname sga@target@IV.5.9.6.1\endcsname{ega4.c04.eq.5-9-6-1}
\expandafter\def\csname sga@target@IV.5.9.7\endcsname{ega4.c04.statement.5-9-7}
\expandafter\def\csname sga@target@IV.5.9.7.1\endcsname{ega4.c04.eq.5-9-7-1}
\expandafter\def\csname sga@target@IV.5.9.8\endcsname{ega4.c04.prop.5-9-8}
\expandafter\def\csname sga@target@IV.5.9.8.1\endcsname{ega4.c04.eq.5-9-8-1}
\expandafter\def\csname sga@target@IV.5.9.8.a\endcsname{ega4.c04.item.5-9-8-a}
\expandafter\def\csname sga@target@IV.5.9.8.a-prime\endcsname{ega4.c04.item.5-9-8-a-prime}
\expandafter\def\csname sga@target@IV.5.9.8.b\endcsname{ega4.c04.item.5-9-8-b}
\expandafter\def\csname sga@target@IV.5.9.8.b-prime\endcsname{ega4.c04.item.5-9-8-b-prime}
\expandafter\def\csname sga@target@IV.5.9.9\endcsname{ega4.c04.def.5-9-9}
\expandafter\def\csname sga@target@IV.6.1.1\endcsname{ega4.c04.prop.6-1-1}
\expandafter\def\csname sga@target@IV.6.1.1.1\endcsname{ega4.c04.eq.6-1-1-1}
\expandafter\def\csname sga@target@IV.6.1.2\endcsname{ega4.c04.cor.6-1-2}
\expandafter\def\csname sga@target@IV.6.1.2.1\endcsname{ega4.c04.eq.6-1-2-1}
\expandafter\def\csname sga@target@IV.6.1.3\endcsname{ega4.c04.cor.6-1-3}
\expandafter\def\csname sga@target@IV.6.1.4\endcsname{ega4.c04.cor.6-1-4}
\expandafter\def\csname sga@target@IV.6.1.4.1\endcsname{ega4.c04.eq.6-1-4-1}
\expandafter\def\csname sga@target@IV.6.1.5\endcsname{ega4.c04.prop.6-1-5}
\expandafter\def\csname sga@target@IV.6.10.1\endcsname{ega4.c04.statement.6-10-1}
\expandafter\def\csname sga@target@IV.6.10.2\endcsname{ega4.c04.prop.6-10-2}
\expandafter\def\csname sga@target@IV.6.10.3\endcsname{ega4.c04.prop.6-10-3}
\expandafter\def\csname sga@target@IV.6.10.3.1\endcsname{ega4.c04.eq.6-10-3-1}
\expandafter\def\csname sga@target@IV.6.10.3.2\endcsname{ega4.c04.eq.6-10-3-2}
\expandafter\def\csname sga@target@IV.6.10.3.3\endcsname{ega4.c04.eq.6-10-3-3}
\expandafter\def\csname sga@target@IV.6.10.4\endcsname{ega4.c04.cor.6-10-4}
\expandafter\def\csname sga@target@IV.6.10.5\endcsname{ega4.c04.prop.6-10-5}
\expandafter\def\csname sga@target@IV.6.10.5.1\endcsname{ega4.c04.eq.6-10-5-1}
\expandafter\def\csname sga@target@IV.6.10.5.2\endcsname{ega4.c04.eq.6-10-5-2}
\expandafter\def\csname sga@target@IV.6.10.5.3\endcsname{ega4.c04.eq.6-10-5-3}
\expandafter\def\csname sga@target@IV.6.10.5.4\endcsname{ega4.c04.eq.6-10-5-4}
\expandafter\def\csname sga@target@IV.6.10.6\endcsname{ega4.c04.prop.6-10-6}
\expandafter\def\csname sga@target@IV.6.10.6.1\endcsname{ega4.c04.eq.6-10-6-1}
\expandafter\def\csname sga@target@IV.6.10.6.2\endcsname{ega4.c04.eq.6-10-6-2}
\expandafter\def\csname sga@target@IV.6.10.6.3\endcsname{ega4.c04.eq.6-10-6-3}
\expandafter\def\csname sga@target@IV.6.11\endcsname{ega4.c04.subsec.6-11.key-iv-6-11}
\expandafter\def\csname sga@target@IV.6.11.1\endcsname{ega4.c04.lem.6-11-1}
\expandafter\def\csname sga@target@IV.6.11.2\endcsname{ega4.c04.prop.6-11-2}
\expandafter\def\csname sga@target@IV.6.11.2.1\endcsname{ega4.c04.eq.6-11-2-1}
\expandafter\def\csname sga@target@IV.6.11.2.2\endcsname{ega4.c04.eq.6-11-2-2}
\expandafter\def\csname sga@target@IV.6.11.3\endcsname{ega4.c04.cor.6-11-3}
\expandafter\def\csname sga@target@IV.6.11.4\endcsname{ega4.c04.statement.6-11-4}
\expandafter\def\csname sga@target@IV.6.11.5\endcsname{ega4.c04.prop.6-11-5}
\expandafter\def\csname sga@target@IV.6.11.5.1\endcsname{ega4.c04.eq.6-11-5-1}
\expandafter\def\csname sga@target@IV.6.11.6\endcsname{ega4.c04.prop.6-11-6}
\expandafter\def\csname sga@target@IV.6.11.6.1\endcsname{ega4.c04.lem.6-11-6-1}
\expandafter\def\csname sga@target@IV.6.11.7\endcsname{ega4.c04.cor.6-11-7}
\expandafter\def\csname sga@target@IV.6.11.8\endcsname{ega4.c04.prop.6-11-8}
\expandafter\def\csname sga@target@IV.6.11.8.1\endcsname{ega4.c04.eq.6-11-8-1}
\expandafter\def\csname sga@target@IV.6.11.9\endcsname{ega4.c04.statement.6-11-9}
\expandafter\def\csname sga@target@IV.6.12\endcsname{ega4.c04.subsec.6-12.key-iv-6-12}
\expandafter\def\csname sga@target@IV.6.12.1\endcsname{ega4.c04.statement.6-12-1}
\expandafter\def\csname sga@target@IV.6.12.10\endcsname{ega4.c04.prop.6-12-10}
\expandafter\def\csname sga@target@IV.6.12.2\endcsname{ega4.c04.prop.6-12-2}
\expandafter\def\csname sga@target@IV.6.12.3\endcsname{ega4.c04.cor.6-12-3}
\expandafter\def\csname sga@target@IV.6.12.4\endcsname{ega4.c04.thm.6-12-4}
\expandafter\def\csname sga@target@IV.6.12.4.1\endcsname{ega4.c04.lem.6-12-4-1}
\expandafter\def\csname sga@target@IV.6.12.5\endcsname{ega4.c04.cor.6-12-5}
\expandafter\def\csname sga@target@IV.6.12.6\endcsname{ega4.c04.cor.6-12-6}
\expandafter\def\csname sga@target@IV.6.12.7\endcsname{ega4.c04.thm.6-12-7}
\expandafter\def\csname sga@target@IV.6.12.8\endcsname{ega4.c04.cor.6-12-8}
\expandafter\def\csname sga@target@IV.6.12.9\endcsname{ega4.c04.prop.6-12-9}
\expandafter\def\csname sga@target@IV.6.13\endcsname{ega4.c04.subsec.6-13.key-iv-6-13}
\expandafter\def\csname sga@target@IV.6.13.1\endcsname{ega4.c04.statement.6-13-1}
\expandafter\def\csname sga@target@IV.6.13.2\endcsname{ega4.c04.prop.6-13-2}
\expandafter\def\csname sga@target@IV.6.13.3\endcsname{ega4.c04.cor.6-13-3}
\expandafter\def\csname sga@target@IV.6.13.4\endcsname{ega4.c04.prop.6-13-4}
\expandafter\def\csname sga@target@IV.6.13.5\endcsname{ega4.c04.cor.6-13-5}
\expandafter\def\csname sga@target@IV.6.13.6\endcsname{ega4.c04.prop.6-13-6}
\expandafter\def\csname sga@target@IV.6.13.7\endcsname{ega4.c04.prop.6-13-7}
\expandafter\def\csname sga@target@IV.6.13.7.1\endcsname{ega4.c04.lem.6-13-7-1}
\expandafter\def\csname sga@target@IV.6.14.1\endcsname{ega4.c04.prop.6-14-1}
\expandafter\def\csname sga@target@IV.6.14.1.1\endcsname{ega4.c04.lem.6-14-1-1}
\expandafter\def\csname sga@target@IV.6.14.1.2\endcsname{ega4.c04.lem.6-14-1-2}
\expandafter\def\csname sga@target@IV.6.14.1.2.env\endcsname{ega4.c04.statement.6-14-1-2-env}
\expandafter\def\csname sga@target@IV.6.14.1.3\endcsname{ega4.c04.eq.6-14-1-3}
\expandafter\def\csname sga@target@IV.6.14.1.4\endcsname{ega4.c04.lem.6-14-1-4}
\expandafter\def\csname sga@target@IV.6.14.2\endcsname{ega4.c04.cor.6-14-2}
\expandafter\def\csname sga@target@IV.6.14.3\endcsname{ega4.c04.cor.6-14-3}
\expandafter\def\csname sga@target@IV.6.14.4\endcsname{ega4.c04.prop.6-14-4}
\expandafter\def\csname sga@target@IV.6.14.4.1\endcsname{ega4.c04.lem.6-14-4-1}
\expandafter\def\csname sga@target@IV.6.14.5\endcsname{ega4.c04.cor.6-14-5}
\expandafter\def\csname sga@target@IV.6.15.1\endcsname{ega4.c04.statement.6-15-1}
\expandafter\def\csname sga@target@IV.6.15.10\endcsname{ega4.c04.prop.6-15-10}
\expandafter\def\csname sga@target@IV.6.15.11\endcsname{ega4.c04.item.6-15-11}
\expandafter\def\csname sga@target@IV.6.15.2\endcsname{ega4.c04.statement.6-15-2}
\expandafter\def\csname sga@target@IV.6.15.3\endcsname{ega4.c04.statement.6-15-3}
\expandafter\def\csname sga@target@IV.6.15.3.1\endcsname{ega4.c04.lem.6-15-3-1}
\expandafter\def\csname sga@target@IV.6.15.4\endcsname{ega4.c04.statement.6-15-4}
\expandafter\def\csname sga@target@IV.6.15.4.1\endcsname{ega4.c04.lem.6-15-4-1}
\expandafter\def\csname sga@target@IV.6.15.5\endcsname{ega4.c04.prop.6-15-5}
\expandafter\def\csname sga@target@IV.6.15.6\endcsname{ega4.c04.prop.6-15-6}
\expandafter\def\csname sga@target@IV.6.15.7\endcsname{ega4.c04.prop.6-15-7}
\expandafter\def\csname sga@target@IV.6.15.8\endcsname{ega4.c04.lem.6-15-8}
\expandafter\def\csname sga@target@IV.6.15.9\endcsname{ega4.c04.cor.6-15-9}
\expandafter\def\csname sga@target@IV.6.2.1\endcsname{ega4.c04.prop.6-2-1}
\expandafter\def\csname sga@target@IV.6.2.1.1\endcsname{ega4.c04.eq.6-2-1-1}
\expandafter\def\csname sga@target@IV.6.2.1.2\endcsname{ega4.c04.eq.6-2-1-2}
\expandafter\def\csname sga@target@IV.6.2.2\endcsname{ega4.c04.cor.6-2-2}
\expandafter\def\csname sga@target@IV.6.2.3\endcsname{ega4.c04.prop.6-2-3}
\expandafter\def\csname sga@target@IV.6.2.3.1\endcsname{ega4.c04.eq.6-2-3-1}
\expandafter\def\csname sga@target@IV.6.2.4\endcsname{ega4.c04.prop.6-2-4}
\expandafter\def\csname sga@target@IV.6.2.4.1\endcsname{ega4.c04.eq.6-2-4-1}
\expandafter\def\csname sga@target@IV.6.3.1\endcsname{ega4.c04.prop.6-3-1}
\expandafter\def\csname sga@target@IV.6.3.1.1\endcsname{ega4.c04.eq.6-3-1-1}
\expandafter\def\csname sga@target@IV.6.3.2\endcsname{ega4.c04.cor.6-3-2}
\expandafter\def\csname sga@target@IV.6.3.2.1\endcsname{ega4.c04.eq.6-3-2-1}
\expandafter\def\csname sga@target@IV.6.3.3\endcsname{ega4.c04.cor.6-3-3}
\expandafter\def\csname sga@target@IV.6.3.4\endcsname{ega4.c04.cor.6-3-4}
\expandafter\def\csname sga@target@IV.6.3.4.1\endcsname{ega4.c04.eq.6-3-4-1}
\expandafter\def\csname sga@target@IV.6.3.4.2\endcsname{ega4.c04.eq.6-3-4-2}
\expandafter\def\csname sga@target@IV.6.3.5\endcsname{ega4.c04.cor.6-3-5}
\expandafter\def\csname sga@target@IV.6.3.6\endcsname{ega4.c04.cor.6-3-6}
\expandafter\def\csname sga@target@IV.6.3.7\endcsname{ega4.c04.cor.6-3-7}
\expandafter\def\csname sga@target@IV.6.3.8\endcsname{ega4.c04.prop.6-3-8}
\expandafter\def\csname sga@target@IV.6.4.1\endcsname{ega4.c04.prop.6-4-1}
\expandafter\def\csname sga@target@IV.6.4.2\endcsname{ega4.c04.cor.6-4-2}
\expandafter\def\csname sga@target@IV.6.4.3\endcsname{ega4.c04.rem.6-4-3}
\expandafter\def\csname sga@target@IV.6.5.1\endcsname{ega4.c04.prop.6-5-1}
\expandafter\def\csname sga@target@IV.6.5.2\endcsname{ega4.c04.cor.6-5-2}
\expandafter\def\csname sga@target@IV.6.5.3\endcsname{ega4.c04.prop.6-5-3}
\expandafter\def\csname sga@target@IV.6.5.4\endcsname{ega4.c04.cor.6-5-4}
\expandafter\def\csname sga@target@IV.6.5.5\endcsname{ega4.c04.item.6-5-5}
\expandafter\def\csname sga@target@IV.6.6.1\endcsname{ega4.c04.prop.6-6-1}
\expandafter\def\csname sga@target@IV.6.6.2\endcsname{ega4.c04.rem.6-6-2}
\expandafter\def\csname sga@target@IV.6.7.1\endcsname{ega4.c04.prop.6-7-1}
\expandafter\def\csname sga@target@IV.6.7.1.1\endcsname{ega4.c04.lem.6-7-1-1}
\expandafter\def\csname sga@target@IV.6.7.2\endcsname{ega4.c04.cor.6-7-2}
\expandafter\def\csname sga@target@IV.6.7.3\endcsname{ega4.c04.cor.6-7-3}
\expandafter\def\csname sga@target@IV.6.7.4\endcsname{ega4.c04.prop.6-7-4}
\expandafter\def\csname sga@target@IV.6.7.4.1\endcsname{ega4.c04.lem.6-7-4-1}
\expandafter\def\csname sga@target@IV.6.7.5\endcsname{ega4.c04.item.6-7-5}
\expandafter\def\csname sga@target@IV.6.7.6\endcsname{ega4.c04.def.6-7-6}
\expandafter\def\csname sga@target@IV.6.7.7\endcsname{ega4.c04.prop.6-7-7}
\expandafter\def\csname sga@target@IV.6.7.8\endcsname{ega4.c04.cor.6-7-8}
\expandafter\def\csname sga@target@IV.6.8.1\endcsname{ega4.c04.def.6-8-1}
\expandafter\def\csname sga@target@IV.6.8.2\endcsname{ega4.c04.prop.6-8-2}
\expandafter\def\csname sga@target@IV.6.8.3\endcsname{ega4.c04.prop.6-8-3}
\expandafter\def\csname sga@target@IV.6.8.4\endcsname{ega4.c04.item.6-8-4}
\expandafter\def\csname sga@target@IV.6.8.5\endcsname{ega4.c04.prop.6-8-5}
\expandafter\def\csname sga@target@IV.6.8.6\endcsname{ega4.c04.thm.6-8-6}
\expandafter\def\csname sga@target@IV.6.8.7\endcsname{ega4.c04.cor.6-8-7}
\expandafter\def\csname sga@target@IV.6.8.8\endcsname{ega4.c04.rem.6-8-8}
\expandafter\def\csname sga@target@IV.6.9.1\endcsname{ega4.c04.thm.6-9-1}
\expandafter\def\csname sga@target@IV.6.9.2\endcsname{ega4.c04.lem.6-9-2}
\expandafter\def\csname sga@target@IV.6.9.3\endcsname{ega4.c04.cor.6-9-3}
\expandafter\def\csname sga@target@IV.7.1.1\endcsname{ega4.c04.def.7-1-1}
\expandafter\def\csname sga@target@IV.7.1.10\endcsname{ega4.c04.prop.7-1-10}
\expandafter\def\csname sga@target@IV.7.1.11\endcsname{ega4.c04.prop.7-1-11}
\expandafter\def\csname sga@target@IV.7.1.12\endcsname{ega4.c04.item.7-1-12}
\expandafter\def\csname sga@target@IV.7.1.13\endcsname{ega4.c04.cor.7-1-13}
\expandafter\def\csname sga@target@IV.7.1.2\endcsname{ega4.c04.prop.7-1-2}
\expandafter\def\csname sga@target@IV.7.1.3\endcsname{ega4.c04.prop.7-1-3}
\expandafter\def\csname sga@target@IV.7.1.3.1\endcsname{ega4.c04.eq.7-1-3-1}
\expandafter\def\csname sga@target@IV.7.1.3.2\endcsname{ega4.c04.eq.7-1-3-2}
\expandafter\def\csname sga@target@IV.7.1.3.3\endcsname{ega4.c04.eq.7-1-3-3}
\expandafter\def\csname sga@target@IV.7.1.3.4\endcsname{ega4.c04.eq.7-1-3-4}
\expandafter\def\csname sga@target@IV.7.1.3.5\endcsname{ega4.c04.eq.7-1-3-5}
\expandafter\def\csname sga@target@IV.7.1.3.6\endcsname{ega4.c04.eq.7-1-3-6}
\expandafter\def\csname sga@target@IV.7.1.3.7\endcsname{ega4.c04.eq.7-1-3-7}
\expandafter\def\csname sga@target@IV.7.1.4\endcsname{ega4.c04.cor.7-1-4}
\expandafter\def\csname sga@target@IV.7.1.5\endcsname{ega4.c04.cor.7-1-5}
\expandafter\def\csname sga@target@IV.7.1.6\endcsname{ega4.c04.rem.7-1-6}
\expandafter\def\csname sga@target@IV.7.1.7\endcsname{ega4.c04.lem.7-1-7}
\expandafter\def\csname sga@target@IV.7.1.8\endcsname{ega4.c04.thm.7-1-8}
\expandafter\def\csname sga@target@IV.7.1.9\endcsname{ega4.c04.def.7-1-9}
\expandafter\def\csname sga@target@IV.7.2.1\endcsname{ega4.c04.statement.7-2-1}
\expandafter\def\csname sga@target@IV.7.2.1.1\endcsname{ega4.c04.statement.7-2-1-1}
\expandafter\def\csname sga@target@IV.7.2.10\endcsname{ega4.c04.statement.7-2-10}
\expandafter\def\csname sga@target@IV.7.2.2\endcsname{ega4.c04.prop.7-2-2}
\expandafter\def\csname sga@target@IV.7.2.3\endcsname{ega4.c04.prop.7-2-3}
\expandafter\def\csname sga@target@IV.7.2.4\endcsname{ega4.c04.prop.7-2-4}
\expandafter\def\csname sga@target@IV.7.2.5\endcsname{ega4.c04.thm.7-2-5}
\expandafter\def\csname sga@target@IV.7.2.5.1\endcsname{ega4.c04.eq.7-2-5-1}
\expandafter\def\csname sga@target@IV.7.2.6\endcsname{ega4.c04.def.7-2-6}
\expandafter\def\csname sga@target@IV.7.2.7\endcsname{ega4.c04.cor.7-2-7}
\expandafter\def\csname sga@target@IV.7.2.8\endcsname{ega4.c04.cor.7-2-8}
\expandafter\def\csname sga@target@IV.7.2.9\endcsname{ega4.c04.cor.7-2-9}
\expandafter\def\csname sga@target@IV.7.3.1\endcsname{ega4.c04.statement.7-3-1}
\expandafter\def\csname sga@target@IV.7.3.10\endcsname{ega4.c04.statement.7-3-10}
\expandafter\def\csname sga@target@IV.7.3.11\endcsname{ega4.c04.rem.7-3-11}
\expandafter\def\csname sga@target@IV.7.3.12\endcsname{ega4.c04.prop.7-3-12}
\expandafter\def\csname sga@target@IV.7.3.12.1\endcsname{ega4.c04.eq.7-3-12-1}
\expandafter\def\csname sga@target@IV.7.3.13\endcsname{ega4.c04.statement.7-3-13}
\expandafter\def\csname sga@target@IV.7.3.14\endcsname{ega4.c04.prop.7-3-14}
\expandafter\def\csname sga@target@IV.7.3.15\endcsname{ega4.c04.prop.7-3-15}
\expandafter\def\csname sga@target@IV.7.3.16\endcsname{ega4.c04.prop.7-3-16}
\expandafter\def\csname sga@target@IV.7.3.16.1\endcsname{ega4.c04.eq.7-3-16-1}
\expandafter\def\csname sga@target@IV.7.3.16.2\endcsname{ega4.c04.lem.7-3-16-2}
\expandafter\def\csname sga@target@IV.7.3.17\endcsname{ega4.c04.prop.7-3-17}
\expandafter\def\csname sga@target@IV.7.3.18\endcsname{ega4.c04.prop.7-3-18}
\expandafter\def\csname sga@target@IV.7.3.19\endcsname{ega4.c04.statement.7-3-19}
\expandafter\def\csname sga@target@IV.7.3.2\endcsname{ega4.c04.lem.7-3-2}
\expandafter\def\csname sga@target@IV.7.3.3\endcsname{ega4.c04.cor.7-3-3}
\expandafter\def\csname sga@target@IV.7.3.4\endcsname{ega4.c04.statement.7-3-4}
\expandafter\def\csname sga@target@IV.7.3.5\endcsname{ega4.c04.item.7-3-5}
\expandafter\def\csname sga@target@IV.7.3.6\endcsname{ega4.c04.statement.7-3-6}
\expandafter\def\csname sga@target@IV.7.3.7\endcsname{ega4.c04.lem.7-3-7}
\expandafter\def\csname sga@target@IV.7.3.8\endcsname{ega4.c04.statement.7-3-8}
\expandafter\def\csname sga@target@IV.7.3.9\endcsname{ega4.c04.item.7-3-9}
\expandafter\def\csname sga@target@IV.7.4.1\endcsname{ega4.c04.statement.7-4-1}
\expandafter\def\csname sga@target@IV.7.4.2\endcsname{ega4.c04.lem.7-4-2}
\expandafter\def\csname sga@target@IV.7.4.3\endcsname{ega4.c04.cor.7-4-3}
\expandafter\def\csname sga@target@IV.7.4.4\endcsname{ega4.c04.thm.7-4-4}
\expandafter\def\csname sga@target@IV.7.4.4.1\endcsname{ega4.c04.lem.7-4-4-1}
\expandafter\def\csname sga@target@IV.7.4.4.2\endcsname{ega4.c04.lem.7-4-4-2}
\expandafter\def\csname sga@target@IV.7.4.5\endcsname{ega4.c04.cor.7-4-5}
\expandafter\def\csname sga@target@IV.7.4.6\endcsname{ega4.c04.prop.7-4-6}
\expandafter\def\csname sga@target@IV.7.4.7\endcsname{ega4.c04.cor.7-4-7}
\expandafter\def\csname sga@target@IV.7.4.8\endcsname{ega4.c04.rem.7-4-8}
\expandafter\def\csname sga@target@IV.7.5.0\endcsname{ega4.c04.statement.7-5-0}
\expandafter\def\csname sga@target@IV.7.5.1\endcsname{ega4.c04.thm.7-5-1}
\expandafter\def\csname sga@target@IV.7.5.1.1\endcsname{ega4.c04.lem.7-5-1-1}
\expandafter\def\csname sga@target@IV.7.5.2\endcsname{ega4.c04.cor.7-5-2}
\expandafter\def\csname sga@target@IV.7.5.3\endcsname{ega4.c04.item.7-5-3}
\expandafter\def\csname sga@target@IV.7.5.4\endcsname{ega4.c04.item.7-5-4}
\expandafter\def\csname sga@target@IV.7.5.5\endcsname{ega4.c04.lem.7-5-5}
\expandafter\def\csname sga@target@IV.7.5.6\endcsname{ega4.c04.prop.7-5-6}
\expandafter\def\csname sga@target@IV.7.5.7\endcsname{ega4.c04.cor.7-5-7}
\expandafter\def\csname sga@target@IV.7.5.7.1\endcsname{ega4.c04.lem.7-5-7-1}
\expandafter\def\csname sga@target@IV.7.5.8\endcsname{ega4.c04.rem.7-5-8}
\expandafter\def\csname sga@target@IV.7.6.1\endcsname{ega4.c04.prop.7-6-1}
\expandafter\def\csname sga@target@IV.7.6.2\endcsname{ega4.c04.cor.7-6-2}
\expandafter\def\csname sga@target@IV.7.6.3\endcsname{ega4.c04.cor.7-6-3}
\expandafter\def\csname sga@target@IV.7.6.4\endcsname{ega4.c04.thm.7-6-4}
\expandafter\def\csname sga@target@IV.7.6.4.1\endcsname{ega4.c04.lem.7-6-4-1}
\expandafter\def\csname sga@target@IV.7.6.5\endcsname{ega4.c04.cor.7-6-5}
\expandafter\def\csname sga@target@IV.7.6.6\endcsname{ega4.c04.cor.7-6-6}
\expandafter\def\csname sga@target@IV.7.6.7\endcsname{ega4.c04.item.7-6-7}
\expandafter\def\csname sga@target@IV.7.7.1\endcsname{ega4.c04.statement.7-7-1}
\expandafter\def\csname sga@target@IV.7.7.2\endcsname{ega4.c04.thm.7-7-2}
\expandafter\def\csname sga@target@IV.7.7.3\endcsname{ega4.c04.cor.7-7-3}
\expandafter\def\csname sga@target@IV.7.7.4\endcsname{ega4.c04.cor.7-7-4}
\expandafter\def\csname sga@target@IV.7.8.1\endcsname{ega4.c04.statement.7-8-1}
\expandafter\def\csname sga@target@IV.7.8.2\endcsname{ega4.c04.def.7-8-2}
\expandafter\def\csname sga@target@IV.7.8.3\endcsname{ega4.c04.statement.7-8-3}
\expandafter\def\csname sga@target@IV.7.8.3.1\endcsname{ega4.c04.statement.7-8-3-1}
\expandafter\def\csname sga@target@IV.7.8.4\endcsname{ega4.c04.item.7-8-4}
\expandafter\def\csname sga@target@IV.7.8.4.1\endcsname{ega4.c04.statement.7-8-4-1}
\expandafter\def\csname sga@target@IV.7.8.5\endcsname{ega4.c04.statement.7-8-5}
\expandafter\def\csname sga@target@IV.7.8.6\endcsname{ega4.c04.prop.7-8-6}
\expandafter\def\csname sga@target@IV.7.9.1\endcsname{ega4.c04.statement.7-9-1}
\expandafter\def\csname sga@target@IV.7.9.10\endcsname{ega4.c04.item.7-9-10}
\expandafter\def\csname sga@target@IV.7.9.2\endcsname{ega4.c04.statement.7-9-2}
\expandafter\def\csname sga@target@IV.7.9.3\endcsname{ega4.c04.prop.7-9-3}
\expandafter\def\csname sga@target@IV.7.9.3.1\endcsname{ega4.c04.lem.7-9-3-1}
\expandafter\def\csname sga@target@IV.7.9.4\endcsname{ega4.c04.cor.7-9-4}
\expandafter\def\csname sga@target@IV.7.9.5\endcsname{ega4.c04.prop.7-9-5}
\expandafter\def\csname sga@target@IV.7.9.6\endcsname{ega4.c04.rem.7-9-6}
\expandafter\def\csname sga@target@IV.7.9.7\endcsname{ega4.c04.statement.7-9-7}
\expandafter\def\csname sga@target@IV.7.9.8\endcsname{ega4.c04.prop.7-9-8}
\expandafter\def\csname sga@target@IV.7.9.8.1\endcsname{ega4.c04.statement.7-9-8-1}
\expandafter\def\csname sga@target@IV.7.9.9\endcsname{ega4.c04.cor.7-9-9}
\expandafter\def\csname sga@target@IV.8.1.1\endcsname{ega4.c04.statement.8-1-1}
\expandafter\def\csname sga@target@IV.8.1.2\endcsname{ega4.c04.statement.8-1-2}
\expandafter\def\csname sga@target@IV.8.1.3\endcsname{ega4.c04.statement.8-1-3}
\expandafter\def\csname sga@target@IV.8.10.1\endcsname{ega4.c04.prop.8-10-1}
\expandafter\def\csname sga@target@IV.8.10.2\endcsname{ega4.c04.cor.8-10-2}
\expandafter\def\csname sga@target@IV.8.10.3\endcsname{ega4.c04.prop.8-10-3}
\expandafter\def\csname sga@target@IV.8.10.3.1\endcsname{ega4.c04.rem.8-10-3-1}
\expandafter\def\csname sga@target@IV.8.10.4\endcsname{ega4.c04.prop.8-10-4}
\expandafter\def\csname sga@target@IV.8.10.5\endcsname{ega4.c04.thm.8-10-5}
\expandafter\def\csname sga@target@IV.8.10.5.1\endcsname{ega4.c04.prop.8-10-5-1}
\expandafter\def\csname sga@target@IV.8.10.5.2\endcsname{ega4.c04.lem.8-10-5-2}
\expandafter\def\csname sga@target@IV.8.11\endcsname{ega4.c04.subsec.8-11.key-iv-8-11}
\expandafter\def\csname sga@target@IV.8.11.1\endcsname{ega4.c04.thm.8-11-1}
\expandafter\def\csname sga@target@IV.8.11.2\endcsname{ega4.c04.thm.8-11-2}
\expandafter\def\csname sga@target@IV.8.11.3\endcsname{ega4.c04.rem.8-11-3}
\expandafter\def\csname sga@target@IV.8.11.4\endcsname{ega4.c04.statement.8-11-4}
\expandafter\def\csname sga@target@IV.8.11.5\endcsname{ega4.c04.prop.8-11-5}
\expandafter\def\csname sga@target@IV.8.11.5.1\endcsname{ega4.c04.rem.8-11-5-1}
\expandafter\def\csname sga@target@IV.8.11.6\endcsname{ega4.c04.prop.8-11-6}
\expandafter\def\csname sga@target@IV.8.12\endcsname{ega4.c04.subsec.8-12.key-iv-8-12}
\expandafter\def\csname sga@target@IV.8.12.1\endcsname{ega4.c04.lem.8-12-1}
\expandafter\def\csname sga@target@IV.8.12.10\endcsname{ega4.c04.cor.8-12-10}
\expandafter\def\csname sga@target@IV.8.12.10.1\endcsname{ega4.c04.lem.8-12-10-1}
\expandafter\def\csname sga@target@IV.8.12.11\endcsname{ega4.c04.cor.8-12-11}
\expandafter\def\csname sga@target@IV.8.12.12\endcsname{ega4.c04.cor.8-12-12}
\expandafter\def\csname sga@target@IV.8.12.2\endcsname{ega4.c04.lem.8-12-2}
\expandafter\def\csname sga@target@IV.8.12.3\endcsname{ega4.c04.prop.8-12-3}
\expandafter\def\csname sga@target@IV.8.12.4\endcsname{ega4.c04.cor.8-12-4}
\expandafter\def\csname sga@target@IV.8.12.5\endcsname{ega4.c04.statement.8-12-5}
\expandafter\def\csname sga@target@IV.8.12.6\endcsname{ega4.c04.thm.8-12-6}
\expandafter\def\csname sga@target@IV.8.12.6.1\endcsname{ega4.c04.eq.8-12-6-1}
\expandafter\def\csname sga@target@IV.8.12.7\endcsname{ega4.c04.rem.8-12-7}
\expandafter\def\csname sga@target@IV.8.12.8\endcsname{ega4.c04.cor.8-12-8}
\expandafter\def\csname sga@target@IV.8.12.9\endcsname{ega4.c04.cor.8-12-9}
\expandafter\def\csname sga@target@IV.8.13.1\endcsname{ega4.c04.prop.8-13-1}
\expandafter\def\csname sga@target@IV.8.13.1.1\endcsname{ega4.c04.eq.8-13-1-1}
\expandafter\def\csname sga@target@IV.8.13.2\endcsname{ega4.c04.cor.8-13-2}
\expandafter\def\csname sga@target@IV.8.13.2.1\endcsname{ega4.c04.eq.8-13-2-1}
\expandafter\def\csname sga@target@IV.8.13.3\endcsname{ega4.c04.statement.8-13-3}
\expandafter\def\csname sga@target@IV.8.13.4\endcsname{ega4.c04.statement.8-13-4}
\expandafter\def\csname sga@target@IV.8.13.4.1\endcsname{ega4.c04.eq.8-13-4-1}
\expandafter\def\csname sga@target@IV.8.13.4.2\endcsname{ega4.c04.eq.8-13-4-2}
\expandafter\def\csname sga@target@IV.8.13.5\endcsname{ega4.c04.prop.8-13-5}
\expandafter\def\csname sga@target@IV.8.13.6\endcsname{ega4.c04.rem.8-13-6}
\expandafter\def\csname sga@target@IV.8.14.1\endcsname{ega4.c04.statement.8-14-1}
\expandafter\def\csname sga@target@IV.8.14.2\endcsname{ega4.c04.prop.8-14-2}
\expandafter\def\csname sga@target@IV.8.14.2.1\endcsname{ega4.c04.eq.8-14-2-1}
\expandafter\def\csname sga@target@IV.8.14.2.2\endcsname{ega4.c04.cor.8-14-2-2}
\expandafter\def\csname sga@target@IV.8.14.2.3\endcsname{ega4.c04.eq.8-14-2-3}
\expandafter\def\csname sga@target@IV.8.14.2.4\endcsname{ega4.c04.eq.8-14-2-4}
\expandafter\def\csname sga@target@IV.8.14.3\endcsname{ega4.c04.rem.8-14-3}
\expandafter\def\csname sga@target@IV.8.2.1\endcsname{ega4.c04.statement.8-2-1}
\expandafter\def\csname sga@target@IV.8.2.1.1\endcsname{ega4.c04.eq.8-2-1-1}
\expandafter\def\csname sga@target@IV.8.2.10\endcsname{ega4.c04.cor.8-2-10}
\expandafter\def\csname sga@target@IV.8.2.11\endcsname{ega4.c04.cor.8-2-11}
\expandafter\def\csname sga@target@IV.8.2.12\endcsname{ega4.c04.cor.8-2-12}
\expandafter\def\csname sga@target@IV.8.2.12.1\endcsname{ega4.c04.eq.8-2-12-1}
\expandafter\def\csname sga@target@IV.8.2.13\endcsname{ega4.c04.prop.8-2-13}
\expandafter\def\csname sga@target@IV.8.2.14\endcsname{ega4.c04.rem.8-2-14}
\expandafter\def\csname sga@target@IV.8.2.2\endcsname{ega4.c04.statement.8-2-2}
\expandafter\def\csname sga@target@IV.8.2.3\endcsname{ega4.c04.prop.8-2-3}
\expandafter\def\csname sga@target@IV.8.2.3.1\endcsname{ega4.c04.eq.8-2-3-1}
\expandafter\def\csname sga@target@IV.8.2.4\endcsname{ega4.c04.lem.8-2-4}
\expandafter\def\csname sga@target@IV.8.2.5\endcsname{ega4.c04.prop.8-2-5}
\expandafter\def\csname sga@target@IV.8.2.6\endcsname{ega4.c04.lem.8-2-6}
\expandafter\def\csname sga@target@IV.8.2.7\endcsname{ega4.c04.rem.8-2-7}
\expandafter\def\csname sga@target@IV.8.2.8\endcsname{ega4.c04.statement.8-2-8}
\expandafter\def\csname sga@target@IV.8.2.9\endcsname{ega4.c04.prop.8-2-9}
\expandafter\def\csname sga@target@IV.8.3.1\endcsname{ega4.c04.statement.8-3-1}
\expandafter\def\csname sga@target@IV.8.3.10\endcsname{ega4.c04.statement.8-3-10}
\expandafter\def\csname sga@target@IV.8.3.10.1\endcsname{ega4.c04.eq.8-3-10-1}
\expandafter\def\csname sga@target@IV.8.3.11\endcsname{ega4.c04.thm.8-3-11}
\expandafter\def\csname sga@target@IV.8.3.12\endcsname{ega4.c04.cor.8-3-12}
\expandafter\def\csname sga@target@IV.8.3.13\endcsname{ega4.c04.prop.8-3-13}
\expandafter\def\csname sga@target@IV.8.3.2\endcsname{ega4.c04.thm.8-3-2}
\expandafter\def\csname sga@target@IV.8.3.2.1\endcsname{ega4.c04.eq.8-3-2-1}
\expandafter\def\csname sga@target@IV.8.3.3\endcsname{ega4.c04.cor.8-3-3}
\expandafter\def\csname sga@target@IV.8.3.4\endcsname{ega4.c04.cor.8-3-4}
\expandafter\def\csname sga@target@IV.8.3.5\endcsname{ega4.c04.cor.8-3-5}
\expandafter\def\csname sga@target@IV.8.3.6\endcsname{ega4.c04.cor.8-3-6}
\expandafter\def\csname sga@target@IV.8.3.7\endcsname{ega4.c04.cor.8-3-7}
\expandafter\def\csname sga@target@IV.8.3.7.1\endcsname{ega4.c04.eq.8-3-7-1}
\expandafter\def\csname sga@target@IV.8.3.8\endcsname{ega4.c04.prop.8-3-8}
\expandafter\def\csname sga@target@IV.8.3.9\endcsname{ega4.c04.statement.8-3-9}
\expandafter\def\csname sga@target@IV.8.3.9.1\endcsname{ega4.c04.eq.8-3-9-1}
\expandafter\def\csname sga@target@IV.8.3.9.2\endcsname{ega4.c04.eq.8-3-9-2}
\expandafter\def\csname sga@target@IV.8.3.9.3\endcsname{ega4.c04.eq.8-3-9-3}
\expandafter\def\csname sga@target@IV.8.3.9.4\endcsname{ega4.c04.eq.8-3-9-4}
\expandafter\def\csname sga@target@IV.8.3.9.5\endcsname{ega4.c04.eq.8-3-9-5}
\expandafter\def\csname sga@target@IV.8.4.1\endcsname{ega4.c04.prop.8-4-1}
\expandafter\def\csname sga@target@IV.8.4.2\endcsname{ega4.c04.prop.8-4-2}
\expandafter\def\csname sga@target@IV.8.4.2.1\endcsname{ega4.c04.lem.8-4-2-1}
\expandafter\def\csname sga@target@IV.8.4.3\endcsname{ega4.c04.cor.8-4-3}
\expandafter\def\csname sga@target@IV.8.4.4\endcsname{ega4.c04.prop.8-4-4}
\expandafter\def\csname sga@target@IV.8.4.5\endcsname{ega4.c04.cor.8-4-5}
\expandafter\def\csname sga@target@IV.8.4.6\endcsname{ega4.c04.rem.8-4-6}
\expandafter\def\csname sga@target@IV.8.5.1\endcsname{ega4.c04.statement.8-5-1}
\expandafter\def\csname sga@target@IV.8.5.1.1\endcsname{ega4.c04.eq.8-5-1-1}
\expandafter\def\csname sga@target@IV.8.5.1.2\endcsname{ega4.c04.eq.8-5-1-2}
\expandafter\def\csname sga@target@IV.8.5.1.3\endcsname{ega4.c04.eq.8-5-1-3}
\expandafter\def\csname sga@target@IV.8.5.1.4\endcsname{ega4.c04.eq.8-5-1-4}
\expandafter\def\csname sga@target@IV.8.5.10\endcsname{ega4.c04.statement.8-5-10}
\expandafter\def\csname sga@target@IV.8.5.10.1\endcsname{ega4.c04.eq.8-5-10-1}
\expandafter\def\csname sga@target@IV.8.5.10.2\endcsname{ega4.c04.eq.8-5-10-2}
\expandafter\def\csname sga@target@IV.8.5.11\endcsname{ega4.c04.prop.8-5-11}
\expandafter\def\csname sga@target@IV.8.5.2\endcsname{ega4.c04.thm.8-5-2}
\expandafter\def\csname sga@target@IV.8.5.2.1\endcsname{ega4.c04.lem.8-5-2-1}
\expandafter\def\csname sga@target@IV.8.5.2.2\endcsname{ega4.c04.eq.8-5-2-2}
\expandafter\def\csname sga@target@IV.8.5.2.3\endcsname{ega4.c04.eq.8-5-2-3}
\expandafter\def\csname sga@target@IV.8.5.2.4\endcsname{ega4.c04.cor.8-5-2-4}
\expandafter\def\csname sga@target@IV.8.5.2.5\endcsname{ega4.c04.cor.8-5-2-5}
\expandafter\def\csname sga@target@IV.8.5.3\endcsname{ega4.c04.statement.8-5-3}
\expandafter\def\csname sga@target@IV.8.5.4\endcsname{ega4.c04.cor.8-5-4}
\expandafter\def\csname sga@target@IV.8.5.5\endcsname{ega4.c04.prop.8-5-5}
\expandafter\def\csname sga@target@IV.8.5.6\endcsname{ega4.c04.prop.8-5-6}
\expandafter\def\csname sga@target@IV.8.5.7\endcsname{ega4.c04.cor.8-5-7}
\expandafter\def\csname sga@target@IV.8.5.8\endcsname{ega4.c04.cor.8-5-8}
\expandafter\def\csname sga@target@IV.8.5.9\endcsname{ega4.c04.lem.8-5-9}
\expandafter\def\csname sga@target@IV.8.6.1\endcsname{ega4.c04.statement.8-6-1}
\expandafter\def\csname sga@target@IV.8.6.2\endcsname{ega4.c04.statement.8-6-2}
\expandafter\def\csname sga@target@IV.8.6.2.1\endcsname{ega4.c04.eq.8-6-2-1}
\expandafter\def\csname sga@target@IV.8.6.2.2\endcsname{ega4.c04.eq.8-6-2-2}
\expandafter\def\csname sga@target@IV.8.6.2.3\endcsname{ega4.c04.eq.8-6-2-3}
\expandafter\def\csname sga@target@IV.8.6.3\endcsname{ega4.c04.prop.8-6-3}
\expandafter\def\csname sga@target@IV.8.6.3.1\endcsname{ega4.c04.cor.8-6-3-1}
\expandafter\def\csname sga@target@IV.8.6.4\endcsname{ega4.c04.cor.8-6-4}
\expandafter\def\csname sga@target@IV.8.7.1\endcsname{ega4.c04.prop.8-7-1}
\expandafter\def\csname sga@target@IV.8.7.2\endcsname{ega4.c04.prop.8-7-2}
\expandafter\def\csname sga@target@IV.8.7.3\endcsname{ega4.c04.cor.8-7-3}
\expandafter\def\csname sga@target@IV.8.7.4\endcsname{ega4.c04.rem.8-7-4}
\expandafter\def\csname sga@target@IV.8.8.1\endcsname{ega4.c04.statement.8-8-1}
\expandafter\def\csname sga@target@IV.8.8.1.1\endcsname{ega4.c04.eq.8-8-1-1}
\expandafter\def\csname sga@target@IV.8.8.2\endcsname{ega4.c04.thm.8-8-2}
\expandafter\def\csname sga@target@IV.8.8.2.1\endcsname{ega4.c04.lem.8-8-2-1}
\expandafter\def\csname sga@target@IV.8.8.2.2\endcsname{ega4.c04.lem.8-8-2-2}
\expandafter\def\csname sga@target@IV.8.8.2.3\endcsname{ega4.c04.lem.8-8-2-3}
\expandafter\def\csname sga@target@IV.8.8.2.4\endcsname{ega4.c04.cor.8-8-2-4}
\expandafter\def\csname sga@target@IV.8.8.2.5\endcsname{ega4.c04.cor.8-8-2-5}
\expandafter\def\csname sga@target@IV.8.8.3\endcsname{ega4.c04.statement.8-8-3}
\expandafter\def\csname sga@target@IV.8.9.1\endcsname{ega4.c04.prop.8-9-1}
\expandafter\def\csname sga@target@IV.8.9.2\endcsname{ega4.c04.statement.8-9-2}
\expandafter\def\csname sga@target@IV.8.9.3\endcsname{ega4.c04.prop.8-9-3}
\expandafter\def\csname sga@target@IV.8.9.4\endcsname{ega4.c04.prop.8-9-4}
\expandafter\def\csname sga@target@IV.8.9.4.1\endcsname{ega4.c04.lem.8-9-4-1}
\expandafter\def\csname sga@target@IV.8.9.5\endcsname{ega4.c04.cor.8-9-5}
\expandafter\def\csname sga@target@IV.9.1.1\endcsname{ega4.c04.prop.9-1-1}
\expandafter\def\csname sga@target@IV.9.1.2\endcsname{ega4.c04.cor.9-1-2}
\expandafter\def\csname sga@target@IV.9.1.3\endcsname{ega4.c04.rem.9-1-3}
\expandafter\def\csname sga@target@IV.9.1.4\endcsname{ega4.c04.prop.9-1-4}
\expandafter\def\csname sga@target@IV.9.1.5\endcsname{ega4.c04.cor.9-1-5}
\expandafter\def\csname sga@target@IV.9.1.6\endcsname{ega4.c04.prop.9-1-6}
\expandafter\def\csname sga@target@IV.9.2\endcsname{ega4.c04.subsec.9-2.key-iv-9-2}
\expandafter\def\csname sga@target@IV.9.2.1\endcsname{ega4.c04.def.9-2-1}
\expandafter\def\csname sga@target@IV.9.2.2\endcsname{ega4.c04.item.9-2-2}
\expandafter\def\csname sga@target@IV.9.2.3\endcsname{ega4.c04.prop.9-2-3}
\expandafter\def\csname sga@target@IV.9.2.4\endcsname{ega4.c04.cor.9-2-4}
\expandafter\def\csname sga@target@IV.9.2.5\endcsname{ega4.c04.rem.9-2-5}
\expandafter\def\csname sga@target@IV.9.2.6\endcsname{ega4.c04.prop.9-2-6}
\expandafter\def\csname sga@target@IV.9.2.6.1\endcsname{ega4.c04.prop.9-2-6-1}
\expandafter\def\csname sga@target@IV.9.2.6.2\endcsname{ega4.c04.cor.9-2-6-2}
\expandafter\def\csname sga@target@IV.9.3.1\endcsname{ega4.c04.prop.9-3-1}
\expandafter\def\csname sga@target@IV.9.3.2\endcsname{ega4.c04.prop.9-3-2}
\expandafter\def\csname sga@target@IV.9.3.3\endcsname{ega4.c04.prop.9-3-3}
\expandafter\def\csname sga@target@IV.9.3.4\endcsname{ega4.c04.rem.9-3-4}
\expandafter\def\csname sga@target@IV.9.3.5\endcsname{ega4.c04.prop.9-3-5}
\expandafter\def\csname sga@target@IV.9.3.6\endcsname{ega4.c04.statement.9-3-6}
\expandafter\def\csname sga@target@IV.9.4.1\endcsname{ega4.c04.statement.9-4-1}
\expandafter\def\csname sga@target@IV.9.4.2\endcsname{ega4.c04.prop.9-4-2}
\expandafter\def\csname sga@target@IV.9.4.3\endcsname{ega4.c04.cor.9-4-3}
\expandafter\def\csname sga@target@IV.9.4.3.1\endcsname{ega4.c04.eq.9-4-3-1}
\expandafter\def\csname sga@target@IV.9.4.4\endcsname{ega4.c04.prop.9-4-4}
\expandafter\def\csname sga@target@IV.9.4.5\endcsname{ega4.c04.cor.9-4-5}
\expandafter\def\csname sga@target@IV.9.4.6\endcsname{ega4.c04.cor.9-4-6}
\expandafter\def\csname sga@target@IV.9.4.7\endcsname{ega4.c04.prop.9-4-7}
\expandafter\def\csname sga@target@IV.9.4.7.1\endcsname{ega4.c04.lem.9-4-7-1}
\expandafter\def\csname sga@target@IV.9.4.7.2\endcsname{ega4.c04.rem.9-4-7-2}
\expandafter\def\csname sga@target@IV.9.4.8\endcsname{ega4.c04.prop.9-4-8}
\expandafter\def\csname sga@target@IV.9.4.9\endcsname{ega4.c04.rem.9-4-9}
\expandafter\def\csname sga@target@IV.9.5.1\endcsname{ega4.c04.prop.9-5-1}
\expandafter\def\csname sga@target@IV.9.5.2\endcsname{ega4.c04.cor.9-5-2}
\expandafter\def\csname sga@target@IV.9.5.3\endcsname{ega4.c04.prop.9-5-3}
\expandafter\def\csname sga@target@IV.9.5.3.case2\endcsname{ega4.c04.item.9-5-3-case2}
\expandafter\def\csname sga@target@IV.9.5.4\endcsname{ega4.c04.cor.9-5-4}
\expandafter\def\csname sga@target@IV.9.5.5\endcsname{ega4.c04.prop.9-5-5}
\expandafter\def\csname sga@target@IV.9.5.6\endcsname{ega4.c04.cor.9-5-6}
\expandafter\def\csname sga@target@IV.9.5.7\endcsname{ega4.c04.rem.9-5-7}
\expandafter\def\csname sga@target@IV.9.6.1\endcsname{ega4.c04.prop.9-6-1}
\expandafter\def\csname sga@target@IV.9.6.2\endcsname{ega4.c04.prop.9-6-2}
\expandafter\def\csname sga@target@IV.9.6.3\endcsname{ega4.c04.prop.9-6-3}
\expandafter\def\csname sga@target@IV.9.6.3.1\endcsname{ega4.c04.eq.9-6-3-1}
\expandafter\def\csname sga@target@IV.9.6.4\endcsname{ega4.c04.cor.9-6-4}
\expandafter\def\csname sga@target@IV.9.6.5\endcsname{ega4.c04.cor.9-6-5}
\expandafter\def\csname sga@target@IV.9.6.6\endcsname{ega4.c04.sec.9-6-6}
\expandafter\def\csname sga@target@IV.9.6.7\endcsname{ega4.c04.sec.9-6-7}
\expandafter\def\csname sga@target@IV.9.7\endcsname{ega4.c04.subsec.9-7.key-iv-9-7}
\expandafter\def\csname sga@target@IV.9.7.1\endcsname{ega4.c04.lem.9-7-1}
\expandafter\def\csname sga@target@IV.9.7.10\endcsname{ega4.c04.cor.9-7-10}
\expandafter\def\csname sga@target@IV.9.7.11\endcsname{ega4.c04.rem.9-7-11}
\expandafter\def\csname sga@target@IV.9.7.12\endcsname{ega4.c04.prop.9-7-12}
\expandafter\def\csname sga@target@IV.9.7.2\endcsname{ega4.c04.lem.9-7-2}
\expandafter\def\csname sga@target@IV.9.7.3\endcsname{ega4.c04.subsec.9-7-3}
\expandafter\def\csname sga@target@IV.9.7.4\endcsname{ega4.c04.lem.9-7-4}
\expandafter\def\csname sga@target@IV.9.7.5\endcsname{ega4.c04.lem.9-7-5}
\expandafter\def\csname sga@target@IV.9.7.5.1\endcsname{ega4.c04.lem.9-7-5-1}
\expandafter\def\csname sga@target@IV.9.7.6\endcsname{ega4.c04.prop.9-7-6}
\expandafter\def\csname sga@target@IV.9.7.7\endcsname{ega4.c04.thm.9-7-7}
\expandafter\def\csname sga@target@IV.9.7.8\endcsname{ega4.c04.prop.9-7-8}
\expandafter\def\csname sga@target@IV.9.7.9\endcsname{ega4.c04.cor.9-7-9}
\expandafter\def\csname sga@target@IV.9.8\endcsname{ega4.c04.subsec.9-8.key-iv-9-8}
\expandafter\def\csname sga@target@IV.9.8.1\endcsname{ega4.c04.prop.9-8-1}
\expandafter\def\csname sga@target@IV.9.8.1-proof\endcsname{ega4.c04.subsec.9-8-1-proof}
\expandafter\def\csname sga@target@IV.9.8.2\endcsname{ega4.c04.subsec.9-8-2}
\expandafter\def\csname sga@target@IV.9.8.3\endcsname{ega4.c04.prop.9-8-3}
\expandafter\def\csname sga@target@IV.9.8.3-proof\endcsname{ega4.c04.subsec.9-8-3-proof}
\expandafter\def\csname sga@target@IV.9.8.4\endcsname{ega4.c04.cor.9-8-4}
\expandafter\def\csname sga@target@IV.9.8.4-proof\endcsname{ega4.c04.item.9-8-4-proof}
\expandafter\def\csname sga@target@IV.9.8.5\endcsname{ega4.c04.prop.9-8-5}
\expandafter\def\csname sga@target@IV.9.8.5-proof\endcsname{ega4.c04.item.9-8-5-proof}
\expandafter\def\csname sga@target@IV.9.8.6\endcsname{ega4.c04.prop.9-8-6}
\expandafter\def\csname sga@target@IV.9.8.6-proof\endcsname{ega4.c04.item.9-8-6-proof}
\expandafter\def\csname sga@target@IV.9.8.7\endcsname{ega4.c04.cor.9-8-7}
\expandafter\def\csname sga@target@IV.9.8.7-proof\endcsname{ega4.c04.item.9-8-7-proof}
\expandafter\def\csname sga@target@IV.9.8.8\endcsname{ega4.c04.prop.9-8-8}
\expandafter\def\csname sga@target@IV.9.8.8-proof\endcsname{ega4.c04.item.9-8-8-proof}
\expandafter\def\csname sga@target@IV.9.8.9\endcsname{ega4.c04.rem.9-8-9}
\expandafter\def\csname sga@target@IV.9.8.9.1\endcsname{ega4.c04.item.9-8-9-1}
\expandafter\def\csname sga@target@IV.9.9\endcsname{ega4.c04.subsec.9-9.key-iv-9-9}
\expandafter\def\csname sga@target@IV.9.9.1\endcsname{ega4.c04.prop.9-9-1}
\expandafter\def\csname sga@target@IV.9.9.1-proof\endcsname{ega4.c04.item.9-9-1-proof}
\expandafter\def\csname sga@target@IV.9.9.1.1\endcsname{ega4.c04.eq.9-9-1-1}
\expandafter\def\csname sga@target@IV.9.9.2\endcsname{ega4.c04.prop.9-9-2}
\expandafter\def\csname sga@target@IV.9.9.2-proof\endcsname{ega4.c04.item.9-9-2-proof}
\expandafter\def\csname sga@target@IV.9.9.2.1\endcsname{ega4.c04.eq.9-9-2-1}
\expandafter\def\csname sga@target@IV.9.9.2.2\endcsname{ega4.c04.eq.9-9-2-2}
\expandafter\def\csname sga@target@IV.9.9.3\endcsname{ega4.c04.cor.9-9-3}
\expandafter\def\csname sga@target@IV.9.9.3-proof\endcsname{ega4.c04.statement.9-9-3-proof}
\expandafter\def\csname sga@target@IV.9.9.4\endcsname{ega4.c04.prop.9-9-4}
\expandafter\def\csname sga@target@IV.9.9.4-proof\endcsname{ega4.c04.item.9-9-4-proof}
\expandafter\def\csname sga@target@IV.9.9.5\endcsname{ega4.c04.cor.9-9-5}
\expandafter\def\csname sga@target@IV.9.9.5-proof\endcsname{ega4.c04.item.9-9-5-proof}
\expandafter\def\csname sga@target@IV.9.9.6\endcsname{ega4.c04.prop.9-9-6}
\expandafter\def\csname sga@target@IV.9.9.6-proof\endcsname{ega4.c04.item.9-9-6-proof}
\expandafter\def\csname sga@target@section:IV.1\endcsname{ega4.c04.sec.1}
\expandafter\def\csname sga@target@section:IV.10\endcsname{ega4.c04.sec.10}
\expandafter\def\csname sga@target@section:IV.11\endcsname{ega4.c04.sec.11}
\expandafter\def\csname sga@target@section:IV.12\endcsname{ega4.c04.sec.12}
\expandafter\def\csname sga@target@section:IV.13\endcsname{ega4.c04.sec.13}
\expandafter\def\csname sga@target@section:IV.14\endcsname{ega4.c04.sec.14}
\expandafter\def\csname sga@target@section:IV.15\endcsname{ega4.c04.sec.15}
\expandafter\def\csname sga@target@section:IV.16\endcsname{ega4.c04.sec.16}
\expandafter\def\csname sga@target@section:IV.17\endcsname{ega4.c04.sec.17}
\expandafter\def\csname sga@target@section:IV.18\endcsname{ega4.c04.sec.18}
\expandafter\def\csname sga@target@section:IV.19\endcsname{ega4.c04.sec.19}
\expandafter\def\csname sga@target@section:IV.2\endcsname{ega4.c04.sec.2}
\expandafter\def\csname sga@target@section:IV.21\endcsname{ega4.c04.sec.21}
\expandafter\def\csname sga@target@section:IV.3\endcsname{ega4.c04.sec.3}
\expandafter\def\csname sga@target@section:IV.4\endcsname{ega4.c04.sec.4}
\expandafter\def\csname sga@target@section:IV.5\endcsname{ega4.c04.sec.5}
\expandafter\def\csname sga@target@section:IV.6\endcsname{ega4.c04.sec.6}
\expandafter\def\csname sga@target@section:IV.7\endcsname{ega4.c04.sec.7}
\expandafter\def\csname sga@target@section:IV.8\endcsname{ega4.c04.sec.8}
\expandafter\def\csname sga@target@section:IV.9\endcsname{ega4.c04.sec.9}
\expandafter\def\csname sga@target@subsection:IV.1.1\endcsname{ega4.c04.subsec.1-1}
\expandafter\def\csname sga@target@subsection:IV.1.10\endcsname{ega4.c04.subsec.1-10}
\expandafter\def\csname sga@target@subsection:IV.1.2\endcsname{ega4.c04.subsec.1-2}
\expandafter\def\csname sga@target@subsection:IV.1.3\endcsname{ega4.c04.subsec.1-3}
\expandafter\def\csname sga@target@subsection:IV.1.4\endcsname{ega4.c04.subsec.1-4}
\expandafter\def\csname sga@target@subsection:IV.1.5\endcsname{ega4.c04.subsec.1-5}
\expandafter\def\csname sga@target@subsection:IV.1.6\endcsname{ega4.c04.subsec.1-6}
\expandafter\def\csname sga@target@subsection:IV.1.7\endcsname{ega4.c04.subsec.1-7}
\expandafter\def\csname sga@target@subsection:IV.1.8\endcsname{ega4.c04.subsec.1-8}
\expandafter\def\csname sga@target@subsection:IV.1.9\endcsname{ega4.c04.subsec.1-9}
\expandafter\def\csname sga@target@subsection:IV.10.1\endcsname{ega4.c04.subsec.10-1}
\expandafter\def\csname sga@target@subsection:IV.10.10\endcsname{ega4.c04.subsec.10-10}
\expandafter\def\csname sga@target@subsection:IV.10.2\endcsname{ega4.c04.subsec.10-2}
\expandafter\def\csname sga@target@subsection:IV.10.3\endcsname{ega4.c04.subsec.10-3}
\expandafter\def\csname sga@target@subsection:IV.10.4\endcsname{ega4.c04.subsec.10-4}
\expandafter\def\csname sga@target@subsection:IV.10.5\endcsname{ega4.c04.subsec.10-5}
\expandafter\def\csname sga@target@subsection:IV.10.6\endcsname{ega4.c04.subsec.10-6}
\expandafter\def\csname sga@target@subsection:IV.10.7\endcsname{ega4.c04.subsec.10-7}
\expandafter\def\csname sga@target@subsection:IV.10.8\endcsname{ega4.c04.subsec.10-8}
\expandafter\def\csname sga@target@subsection:IV.10.9\endcsname{ega4.c04.subsec.10-9}
\expandafter\def\csname sga@target@subsection:IV.11.1\endcsname{ega4.c04.subsec.11-1}
\expandafter\def\csname sga@target@subsection:IV.11.10\endcsname{ega4.c04.subsec.11-10}
\expandafter\def\csname sga@target@subsection:IV.11.2\endcsname{ega4.c04.subsec.11-2}
\expandafter\def\csname sga@target@subsection:IV.11.3\endcsname{ega4.c04.subsec.11-3}
\expandafter\def\csname sga@target@subsection:IV.11.4\endcsname{ega4.c04.subsec.11-4}
\expandafter\def\csname sga@target@subsection:IV.11.5\endcsname{ega4.c04.subsec.11-5}
\expandafter\def\csname sga@target@subsection:IV.11.6\endcsname{ega4.c04.subsec.11-6}
\expandafter\def\csname sga@target@subsection:IV.11.7\endcsname{ega4.c04.subsec.11-7}
\expandafter\def\csname sga@target@subsection:IV.11.8\endcsname{ega4.c04.subsec.11-8}
\expandafter\def\csname sga@target@subsection:IV.11.9\endcsname{ega4.c04.subsec.11-9}
\expandafter\def\csname sga@target@subsection:IV.12.0\endcsname{ega4.c04.subsec.12-0}
\expandafter\def\csname sga@target@subsection:IV.12.1\endcsname{ega4.c04.subsec.12-1}
\expandafter\def\csname sga@target@subsection:IV.12.2\endcsname{ega4.c04.subsec.12-2}
\expandafter\def\csname sga@target@subsection:IV.12.3\endcsname{ega4.c04.subsec.12-3}
\expandafter\def\csname sga@target@subsection:IV.13.1\endcsname{ega4.c04.subsec.13-1}
\expandafter\def\csname sga@target@subsection:IV.13.2\endcsname{ega4.c04.subsec.13-2}
\expandafter\def\csname sga@target@subsection:IV.13.3\endcsname{ega4.c04.subsec.13-3}
\expandafter\def\csname sga@target@subsection:IV.14.1\endcsname{ega4.c04.subsec.14-1}
\expandafter\def\csname sga@target@subsection:IV.14.2\endcsname{ega4.c04.subsec.14-2}
\expandafter\def\csname sga@target@subsection:IV.14.3\endcsname{ega4.c04.subsec.14-3}
\expandafter\def\csname sga@target@subsection:IV.14.4\endcsname{ega4.c04.subsec.14-4}
\expandafter\def\csname sga@target@subsection:IV.14.5\endcsname{ega4.c04.subsec.14-5}
\expandafter\def\csname sga@target@subsection:IV.15.1\endcsname{ega4.c04.subsec.15-1}
\expandafter\def\csname sga@target@subsection:IV.15.2\endcsname{ega4.c04.subsec.15-2}
\expandafter\def\csname sga@target@subsection:IV.15.3\endcsname{ega4.c04.subsec.15-3}
\expandafter\def\csname sga@target@subsection:IV.15.4\endcsname{ega4.c04.subsec.15-4}
\expandafter\def\csname sga@target@subsection:IV.15.5\endcsname{ega4.c04.subsec.15-5}
\expandafter\def\csname sga@target@subsection:IV.15.6\endcsname{ega4.c04.subsec.15-6}
\expandafter\def\csname sga@target@subsection:IV.15.7\endcsname{ega4.c04.subsec.15-7}
\expandafter\def\csname sga@target@subsection:IV.17.1\endcsname{ega4.c04.subsec.17-1}
\expandafter\def\csname sga@target@subsection:IV.17.10\endcsname{ega4.c04.subsec.17-10}
\expandafter\def\csname sga@target@subsection:IV.17.11\endcsname{ega4.c04.subsec.17-11}
\expandafter\def\csname sga@target@subsection:IV.17.12\endcsname{ega4.c04.subsec.17-12}
\expandafter\def\csname sga@target@subsection:IV.17.13\endcsname{ega4.c04.subsec.17-13}
\expandafter\def\csname sga@target@subsection:IV.17.14\endcsname{ega4.c04.subsec.17-14}
\expandafter\def\csname sga@target@subsection:IV.17.15\endcsname{ega4.c04.subsec.17-15}
\expandafter\def\csname sga@target@subsection:IV.17.16\endcsname{ega4.c04.subsec.17-16}
\expandafter\def\csname sga@target@subsection:IV.17.2\endcsname{ega4.c04.subsec.17-2}
\expandafter\def\csname sga@target@subsection:IV.17.3\endcsname{ega4.c04.subsec.17-3}
\expandafter\def\csname sga@target@subsection:IV.17.4\endcsname{ega4.c04.subsec.17-4}
\expandafter\def\csname sga@target@subsection:IV.17.5\endcsname{ega4.c04.subsec.17-5}
\expandafter\def\csname sga@target@subsection:IV.17.6\endcsname{ega4.c04.subsec.17-6}
\expandafter\def\csname sga@target@subsection:IV.17.7\endcsname{ega4.c04.subsec.17-7}
\expandafter\def\csname sga@target@subsection:IV.17.8\endcsname{ega4.c04.subsec.17-8}
\expandafter\def\csname sga@target@subsection:IV.17.9\endcsname{ega4.c04.subsec.17-9}
\expandafter\def\csname sga@target@subsection:IV.18.1\endcsname{ega4.c04.subsec.18-1}
\expandafter\def\csname sga@target@subsection:IV.18.10\endcsname{ega4.c04.subsec.18-10}
\expandafter\def\csname sga@target@subsection:IV.18.11\endcsname{ega4.c04.subsec.18-11}
\expandafter\def\csname sga@target@subsection:IV.18.12\endcsname{ega4.c04.subsec.18-12}
\expandafter\def\csname sga@target@subsection:IV.18.2\endcsname{ega4.c04.subsec.18-2}
\expandafter\def\csname sga@target@subsection:IV.18.3\endcsname{ega4.c04.subsec.18-3}
\expandafter\def\csname sga@target@subsection:IV.18.4\endcsname{ega4.c04.subsec.18-4}
\expandafter\def\csname sga@target@subsection:IV.18.5\endcsname{ega4.c04.subsec.18-5}
\expandafter\def\csname sga@target@subsection:IV.18.6\endcsname{ega4.c04.subsec.18-6}
\expandafter\def\csname sga@target@subsection:IV.18.7\endcsname{ega4.c04.subsec.18-7}
\expandafter\def\csname sga@target@subsection:IV.18.8\endcsname{ega4.c04.subsec.18-8}
\expandafter\def\csname sga@target@subsection:IV.18.9\endcsname{ega4.c04.subsec.18-9}
\expandafter\def\csname sga@target@subsection:IV.19.1\endcsname{ega4.c04.subsec.19-1}
\expandafter\def\csname sga@target@subsection:IV.19.2\endcsname{ega4.c04.subsec.19-2}
\expandafter\def\csname sga@target@subsection:IV.19.3\endcsname{ega4.c04.subsec.19-3}
\expandafter\def\csname sga@target@subsection:IV.19.4\endcsname{ega4.c04.subsec.19-4}
\expandafter\def\csname sga@target@subsection:IV.19.5\endcsname{ega4.c04.subsec.19-5}
\expandafter\def\csname sga@target@subsection:IV.19.6\endcsname{ega4.c04.subsec.19-6}
\expandafter\def\csname sga@target@subsection:IV.19.7\endcsname{ega4.c04.subsec.19-7}
\expandafter\def\csname sga@target@subsection:IV.19.8\endcsname{ega4.c04.subsec.19-8}
\expandafter\def\csname sga@target@subsection:IV.19.9\endcsname{ega4.c04.subsec.19-9}
\expandafter\def\csname sga@target@subsection:IV.2.1\endcsname{ega4.c04.subsec.2-1}
\expandafter\def\csname sga@target@subsection:IV.2.2\endcsname{ega4.c04.subsec.2-2}
\expandafter\def\csname sga@target@subsection:IV.2.3\endcsname{ega4.c04.subsec.2-3}
\expandafter\def\csname sga@target@subsection:IV.2.4\endcsname{ega4.c04.subsec.2-4}
\expandafter\def\csname sga@target@subsection:IV.2.5\endcsname{ega4.c04.subsec.2-5}
\expandafter\def\csname sga@target@subsection:IV.2.6\endcsname{ega4.c04.subsec.2-6}
\expandafter\def\csname sga@target@subsection:IV.2.7\endcsname{ega4.c04.subsec.2-7}
\expandafter\def\csname sga@target@subsection:IV.2.8\endcsname{ega4.c04.subsec.2-8}
\expandafter\def\csname sga@target@subsection:IV.21.1\endcsname{ega4.c04.subsec.21-1}
\expandafter\def\csname sga@target@subsection:IV.21.10\endcsname{ega4.c04.subsec.21-10}
\expandafter\def\csname sga@target@subsection:IV.21.11\endcsname{ega4.c04.subsec.21-11}
\expandafter\def\csname sga@target@subsection:IV.21.12\endcsname{ega4.c04.subsec.21-12}
\expandafter\def\csname sga@target@subsection:IV.21.13\endcsname{ega4.c04.subsec.21-13}
\expandafter\def\csname sga@target@subsection:IV.21.14\endcsname{ega4.c04.subsec.21-14}
\expandafter\def\csname sga@target@subsection:IV.21.15\endcsname{ega4.c04.subsec.21-15}
\expandafter\def\csname sga@target@subsection:IV.21.2\endcsname{ega4.c04.subsec.21-2}
\expandafter\def\csname sga@target@subsection:IV.21.3\endcsname{ega4.c04.subsec.21-3}
\expandafter\def\csname sga@target@subsection:IV.21.4\endcsname{ega4.c04.subsec.21-4}
\expandafter\def\csname sga@target@subsection:IV.21.5\endcsname{ega4.c04.subsec.21-5}
\expandafter\def\csname sga@target@subsection:IV.21.6\endcsname{ega4.c04.subsec.21-6}
\expandafter\def\csname sga@target@subsection:IV.21.7\endcsname{ega4.c04.subsec.21-7}
\expandafter\def\csname sga@target@subsection:IV.21.8\endcsname{ega4.c04.subsec.21-8}
\expandafter\def\csname sga@target@subsection:IV.21.9\endcsname{ega4.c04.subsec.21-9}
\expandafter\def\csname sga@target@subsection:IV.3.1\endcsname{ega4.c04.subsec.3-1}
\expandafter\def\csname sga@target@subsection:IV.3.2\endcsname{ega4.c04.subsec.3-2}
\expandafter\def\csname sga@target@subsection:IV.3.3\endcsname{ega4.c04.subsec.3-3}
\expandafter\def\csname sga@target@subsection:IV.3.4\endcsname{ega4.c04.subsec.3-4}
\expandafter\def\csname sga@target@subsection:IV.4.1\endcsname{ega4.c04.subsec.4-1}
\expandafter\def\csname sga@target@subsection:IV.4.2\endcsname{ega4.c04.subsec.4-2}
\expandafter\def\csname sga@target@subsection:IV.4.3\endcsname{ega4.c04.subsec.4-3}
\expandafter\def\csname sga@target@subsection:IV.4.4\endcsname{ega4.c04.subsec.4-4}
\expandafter\def\csname sga@target@subsection:IV.4.5\endcsname{ega4.c04.subsec.4-5}
\expandafter\def\csname sga@target@subsection:IV.4.6\endcsname{ega4.c04.subsec.4-6}
\expandafter\def\csname sga@target@subsection:IV.4.7\endcsname{ega4.c04.subsec.4-7}
\expandafter\def\csname sga@target@subsection:IV.4.8\endcsname{ega4.c04.subsec.4-8}
\expandafter\def\csname sga@target@subsection:IV.4.9\endcsname{ega4.c04.subsec.4-9}
\expandafter\def\csname sga@target@subsection:IV.5.1\endcsname{ega4.c04.subsec.5-1}
\expandafter\def\csname sga@target@subsection:IV.5.10\endcsname{ega4.c04.subsec.5-10}
\expandafter\def\csname sga@target@subsection:IV.5.11\endcsname{ega4.c04.subsec.5-11}
\expandafter\def\csname sga@target@subsection:IV.5.12\endcsname{ega4.c04.subsec.5-12}
\expandafter\def\csname sga@target@subsection:IV.5.13\endcsname{ega4.c04.subsec.5-13}
\expandafter\def\csname sga@target@subsection:IV.5.2\endcsname{ega4.c04.subsec.5-2}
\expandafter\def\csname sga@target@subsection:IV.5.3\endcsname{ega4.c04.subsec.5-3}
\expandafter\def\csname sga@target@subsection:IV.5.4\endcsname{ega4.c04.subsec.5-4}
\expandafter\def\csname sga@target@subsection:IV.5.5\endcsname{ega4.c04.subsec.5-5}
\expandafter\def\csname sga@target@subsection:IV.5.6\endcsname{ega4.c04.subsec.5-6}
\expandafter\def\csname sga@target@subsection:IV.5.7\endcsname{ega4.c04.subsec.5-7}
\expandafter\def\csname sga@target@subsection:IV.5.8\endcsname{ega4.c04.subsec.5-8}
\expandafter\def\csname sga@target@subsection:IV.5.9\endcsname{ega4.c04.subsec.5-9}
\expandafter\def\csname sga@target@subsection:IV.6.1\endcsname{ega4.c04.subsec.6-1}
\expandafter\def\csname sga@target@subsection:IV.6.10\endcsname{ega4.c04.subsec.6-10}
\expandafter\def\csname sga@target@subsection:IV.6.11\endcsname{ega4.c04.subsec.6-11}
\expandafter\def\csname sga@target@subsection:IV.6.12\endcsname{ega4.c04.subsec.6-12}
\expandafter\def\csname sga@target@subsection:IV.6.13\endcsname{ega4.c04.subsec.6-13}
\expandafter\def\csname sga@target@subsection:IV.6.14\endcsname{ega4.c04.subsec.6-14}
\expandafter\def\csname sga@target@subsection:IV.6.15\endcsname{ega4.c04.subsec.6-15}
\expandafter\def\csname sga@target@subsection:IV.6.2\endcsname{ega4.c04.subsec.6-2}
\expandafter\def\csname sga@target@subsection:IV.6.3\endcsname{ega4.c04.subsec.6-3}
\expandafter\def\csname sga@target@subsection:IV.6.4\endcsname{ega4.c04.subsec.6-4}
\expandafter\def\csname sga@target@subsection:IV.6.5\endcsname{ega4.c04.subsec.6-5}
\expandafter\def\csname sga@target@subsection:IV.6.6\endcsname{ega4.c04.subsec.6-6}
\expandafter\def\csname sga@target@subsection:IV.6.7\endcsname{ega4.c04.subsec.6-7}
\expandafter\def\csname sga@target@subsection:IV.6.8\endcsname{ega4.c04.subsec.6-8}
\expandafter\def\csname sga@target@subsection:IV.6.9\endcsname{ega4.c04.subsec.6-9}
\expandafter\def\csname sga@target@subsection:IV.7.1\endcsname{ega4.c04.subsec.7-1}
\expandafter\def\csname sga@target@subsection:IV.7.2\endcsname{ega4.c04.subsec.7-2}
\expandafter\def\csname sga@target@subsection:IV.7.3\endcsname{ega4.c04.subsec.7-3}
\expandafter\def\csname sga@target@subsection:IV.7.4\endcsname{ega4.c04.subsec.7-4}
\expandafter\def\csname sga@target@subsection:IV.7.5\endcsname{ega4.c04.subsec.7-5}
\expandafter\def\csname sga@target@subsection:IV.7.6\endcsname{ega4.c04.subsec.7-6}
\expandafter\def\csname sga@target@subsection:IV.7.7\endcsname{ega4.c04.subsec.7-7}
\expandafter\def\csname sga@target@subsection:IV.7.8\endcsname{ega4.c04.subsec.7-8}
\expandafter\def\csname sga@target@subsection:IV.7.9\endcsname{ega4.c04.subsec.7-9}
\expandafter\def\csname sga@target@subsection:IV.8.1\endcsname{ega4.c04.subsec.8-1}
\expandafter\def\csname sga@target@subsection:IV.8.10\endcsname{ega4.c04.subsec.8-10}
\expandafter\def\csname sga@target@subsection:IV.8.11\endcsname{ega4.c04.subsec.8-11}
\expandafter\def\csname sga@target@subsection:IV.8.12\endcsname{ega4.c04.subsec.8-12}
\expandafter\def\csname sga@target@subsection:IV.8.13\endcsname{ega4.c04.subsec.8-13}
\expandafter\def\csname sga@target@subsection:IV.8.14\endcsname{ega4.c04.subsec.8-14}
\expandafter\def\csname sga@target@subsection:IV.8.2\endcsname{ega4.c04.subsec.8-2}
\expandafter\def\csname sga@target@subsection:IV.8.3\endcsname{ega4.c04.subsec.8-3}
\expandafter\def\csname sga@target@subsection:IV.8.4\endcsname{ega4.c04.subsec.8-4}
\expandafter\def\csname sga@target@subsection:IV.8.5\endcsname{ega4.c04.subsec.8-5}
\expandafter\def\csname sga@target@subsection:IV.8.6\endcsname{ega4.c04.subsec.8-6}
\expandafter\def\csname sga@target@subsection:IV.8.7\endcsname{ega4.c04.subsec.8-7}
\expandafter\def\csname sga@target@subsection:IV.8.8\endcsname{ega4.c04.subsec.8-8}
\expandafter\def\csname sga@target@subsection:IV.8.9\endcsname{ega4.c04.subsec.8-9}
\expandafter\def\csname sga@target@subsection:IV.9.1\endcsname{ega4.c04.subsec.9-1}
\expandafter\def\csname sga@target@subsection:IV.9.2\endcsname{ega4.c04.subsec.9-2}
\expandafter\def\csname sga@target@subsection:IV.9.3\endcsname{ega4.c04.subsec.9-3}
\expandafter\def\csname sga@target@subsection:IV.9.4\endcsname{ega4.c04.subsec.9-4}
\expandafter\def\csname sga@target@subsection:IV.9.5\endcsname{ega4.c04.subsec.9-5}
\expandafter\def\csname sga@target@subsection:IV.9.6\endcsname{ega4.c04.subsec.9-6}
\expandafter\def\csname sga@target@subsection:IV.9.7\endcsname{ega4.c04.subsec.9-7}
\expandafter\def\csname sga@target@subsection:IV.9.8\endcsname{ega4.c04.subsec.9-8}
\expandafter\def\csname sga@target@subsection:IV.9.9\endcsname{ega4.c04.subsec.9-9}
\let\sgaOriginalRef\ref
\let\sgaOriginalPageref\pageref
\renewcommand{\ref}[1]{%
  \ifcsname sga@target@#1\endcsname
    \hyperlink{\csname sga@target@#1\endcsname}{\sgaOriginalRef*{#1}}%
  \else
    \sgaOriginalRef{#1}%
  \fi
}
\renewcommand{\pageref}[1]{%
  \ifcsname sga@target@#1\endcsname
    \hyperlink{\csname sga@target@#1\endcsname}{\sgaOriginalPageref*{#1}}%
  \else
    \sgaOriginalPageref{#1}%
  \fi
}
\renewcommand{\eqref}[1]{%
  \ifcsname sga@target@#1\endcsname
    \hyperlink{\csname sga@target@#1\endcsname}{\textup{\tagform@{\sgaOriginalRef*{#1}}}}%
  \else
    \textup{\tagform@{\sgaOriginalRef{#1}}}%
  \fi
}
% Stable semantic routing for the reader's prose reference helpers.
\newcommand{\sgaStableOrOriginalLink}[2]{%
  \ifcsname sga@target@#1\endcsname
    \hyperlink{\csname sga@target@#1\endcsname}{#2}%
  \else
    \@ifundefined{r@#1}{#2}{\hyperref[#1]{#2}}%
  \fi
}
\RenewDocumentCommand{\sref}{o m o}{%
  \IfNoValueTF{#1}%
    {\IfNoValueTF{#3}%
      {\sgaStableOrOriginalLink{#2}{\normalfont{(\StrBehind{#2}{.})}}}%
      {\sgaStableOrOriginalLink{#2}{\normalfont{(\StrBehind{#2}{.},~{#3})}}}}%
    {\IfNoValueTF{#3}%
      {\sgaStableOrOriginalLink{#2}{\normalfont{(\textbf{#1},~\StrBehind{#2}{.})}}}%
      {\sgaStableOrOriginalLink{#2}{\normalfont{(\textbf{#1},~\StrBehind{#2}{.},~{#3})}}}}%
}
\renewcommand{\eref}[1]{%
  \sgaStableOrOriginalLink{#1}{\normalfont{(\StrBehind{#1}{.})}}}
\makeatother
% END NATIVE INPUT EGA4_reference_v2_links
\setlength{\emergencystretch}{3em}

% Register the three external labels referenced by IV.1 without inventing
% false landing pages in this standalone volume.
\makeatletter
\newcommand{\EGAIVOneKnownExternalLabel}[1]{%
  \expandafter\gdef\csname r@#1\endcsname{{}{0}{}{}{}}}
\EGAIVOneKnownExternalLabel{subsection:I.6.3}
\EGAIVOneKnownExternalLabel{subsection:I.6.6}
\EGAIVOneKnownExternalLabel{section:0.9}
\makeatother

\hypersetup{
  pdftitle={Elements of Algebraic Geometry IV, Part One -- English Edition},
  pdfauthor={Alexandre Grothendieck and Jean Dieudonne; English translation},
  pdfsubject={Standalone English reader of EGA IV, Part One}
}

\begin{document}

\title{Elements of Algebraic Geometry IV\\Part One}
\author{Alexandre Grothendieck and Jean Dieudonn\'e}
\date{}
\maketitle

\begin{center}
\large Standalone English edition
\end{center}

\tableofcontents

% Chapter 0 continuation printed in EGA IV, Part One.
\setcounter{section}{13}
% BEGIN NATIVE INPUT ega0/ega0-14
\section{Combinatorial dimension of a topological space}
\label{section:0.14}

\begin{gather*}
  \textbf{Summary}
\end{gather*}

\oldpage[IV-1]{5}
\begin{longtable}{ll}
  \textsection14. & Combinatorial dimension of a topological space.\\
  \textsection15. & $M$-regular sequences and $\sh{F}$-regular sequences.\\
  \textsection16. & Dimension and depth of Noetherian local rings.\\
  \textsection17. & Regular rings.\\
  \textsection18. & Supplement on extensions of algebras.\\
  \textsection19. & Formally smooth algebras and Cohen rings.\\
  \textsection20. & Derivations and differentials.\\
  \textsection21. & Differentials in rings of characteristic $p$.\\
  \textsection22. & Differential criteria for smoothness and regularity.\\
  \textsection23. & Japanese rings.
\end{longtable}
\bigskip

Almost all of the preceding sections have been focused on the exposition of ideas of commutative algebra that will be used throughout Chapter~IV.
Even though a large amount of these ideas already appear in multiple works~(\cite{I-1}, \cite{I-12}, \cite{I-13}, \cite{I-17}, \cite{IV-30}), we thought that it would be more practical for the reader to have a coherent, vaguely independent exposition.
Together with \textsection\textsection5,~6, and~7 of Chapter~IV (where we use the language of schemes), these sections constitute, in the middle of our treatise, a miniature special treatise, somewhat independent of Chapters~I to~III, and one that aims to present, in a coherent manner, the properties of rings that ``behave well'' with respect to operations such as completion, or integral closure, by systematically associating these properties to more general ideas.\footnote{The majority of properties which we discuss were discovered by Chevalley, Zariski, Nagata, and Serre. The method used here was first developed in the autumn of 1961, in a course taught at Harvard University by A.~Grothendieck.}

\subsection{Combinatorial dimension of a topological space}
\label{subsection:0.14.1}

\oldpage[0\textsubscript{IV-1}]{6}
\begin{env}[14.1.1]
\label{0.14.1.1}
Let $I$ be an ordered set; a \emph{chain} of elements of $I$ is, by definition, a strictly-increasing finite sequence $i_0<i_1<\cdots<i_n$ of elements of $I$ ($n\geq 0$); by definition, the \emph{length} of this chain is $n$.
If $X$ is a topological space, the set of its \emph{irreducible closed} subsets is ordered by inclusion, and so we have the notion of a \emph{chain} of irreducible closed subsets of $X$.
\end{env}

\begin{definition}[14.1.2]
\label{0.14.1.2}
Let $X$ be a topological space.
We define the combinatorial dimension of $X$ (or simply the dimension of $X$, if there is no risk of confusion), denoted by $\dimc(X)$ (or simply $\dim(X)$), to be the upper bound of lengths of chains of irreducible closed subsets of $X$.
For all $x\in X$, we define the combinatorial dimension of $X$ at $x$ (or simply the dimension of $X$ at $x$), denoted by $\dim_x(X)$, to be the number $\inf_U(\dim(U))$, where $U$ varies over the open neighbourhoods of $x$ in $X$.
\end{definition}

It follows from this definition that we have
\[
  \dim(\emp)=-\infty
\]
(the upper bound in $\overline{\bb{R}}$ of the empty set being $-\infty$).
If $(X_\alpha)$ is the family of irreducible components of $X$, then we have
\[
\label{0.14.1.2.1}
  \dim(X)=\sup_\alpha\dim(X_\alpha),
  \tag{14.1.2.1}
\]
because every chain of irreducible closed subsets of $X$ is, by definition, contained in some irreducible component of $X$, and, conversely, the irreducible components are closed in $X$, so every irreducible closed subset of an $X_\alpha$ is an irreducible closed subset of $X$.

\begin{definition}[14.1.3]
\label{0.14.1.3}
We say that a topological space $X$ is equidimensional if all its irreducible components have the same dimension (which is thus equal to $\dim(X)$, by \sref{0.14.1.2.1}).
\end{definition}

\begin{proposition}[14.1.4]
\label{0.14.1.4}
\medskip\noindent
\begin{enumerate}
  \item[{\rm(i)}] For every subset $Y$ of a topological space $X$, we have $\dim(Y)\leq\dim(X)$.
  \item[{\rm(ii)}] If a topological space $X$ is a finite union of closed subsets $X_i$, then we have $\dim(X)=\sup_i\dim(X_i)$.
\end{enumerate}
\end{proposition}

\begin{proof}
For every irreducible closed subset $Z$ of $Y$, the closure $\overline{Z}$ of $Z$ in $X$ is irreducible \sref[0\textsubscript{I}]{0.2.1.2}, and $\overline{Z}\cap Y=Z$, whence (i).
Now, if $X=\bigcup_{i=1}^nX_i$, where the $X_i$ are closed, then every irreducible closed subset of $X$ is contained in one of the $X_i$ \sref[0\textsubscript{I}]{0.2.1.1}, and so every chain of irreducible closed subsets of $X$ is contained in one of the $X_i$, whence (ii).
\end{proof}

From \sref{0.14.1.4}[i], we see that, for all $x\in X$, we can also write
\[
\label{0.14.1.4.1}
  \dim_x(X)=\lim_U\dim(U),
  \tag{14.1.4.1}
\]
where the limit is taken over the downward-directed set of open neighbourhoods of $x$ in $X$.

\oldpage[0\textsubscript{IV-1}]{7}
\begin{corollary}[14.1.5]
\label{0.14.1.5}
Let $X$ be a topological space, $x$ a point of $X$, $U$ a neighbourhood of $x$, and $Y_i$ ($1\leq i\leq n$) closed subsets of $U$ such that, for all $i$, $x\in Y_i$, and such that $U$ is the union of the $Y_i$.
Then we have
\[
\label{0.14.1.5.1}
  \dim_x(X)=\sup_i(\dim_x(Y_i)).
  \tag{14.1.5.1}
\]
\end{corollary}

\begin{proof}
It follows from \sref{0.14.1.4}[ii] that we have $\dim_x(X) = \inf_V(\sup_i(\dim(Y_i\cap V)))$, where $V$ ranges over the set of open neighbourhoods of $x$ that are contained in $U$;
similarly, we have $\dim_x(Y_i) = \inf_V(\dim(Y_i\cap V))$ for all $i$.
The corollary is thus evident if
\[
  \sup_i(\dim_x(Y_i))=+\infty;
\]
if this were not the case, then there would be an open neighbourhood $V_0\subset U$ of $x$ such that $\dim(Y_i\cap V)=\dim_x(Y_i)$ for $1\leq i\leq n$ and for all $V\subset V_0$, whence the conclusion.
\end{proof}

\begin{proposition}[14.1.6]
\label{0.14.1.6}
For every topological space $X$, we have $\dim(X)=\sup_{x\in X}\dim_x(X)$.
\end{proposition}

\begin{proof}
It follows from Definition~\sref{0.14.1.2} and Proposition~\sref{0.14.1.4} that $\dim_x(X)\leq\dim(X)$ for all $x\in X$.
Now, let $Z_0\subset Z_1\subset\ldots\subset Z_n$ be a chain of irreducible closed subsets of $X$, and let $x\in Z_0$;
for every open subset $U\subset X$ that contains $x$, $U\cap Z_i$ is irreducible \sref[0\textsubscript{I}]{0.2.1.6} and closed in $U$, and since we have $\overline{U\cap Z_i}=Z_i$ in $X$, the $U\cap Z_i$ are pairwise distinct;
thus $\dim(U)\geq n$, which finishes the proof.
\end{proof}

\begin{corollary}[14.1.7]
\label{0.14.1.7}
If $(X_\alpha)$ is an open, or closed and locally finite, cover of $X$, then $\dim(X)=\sup_\alpha(\dim(X_\alpha))$.
\end{corollary}

\begin{proof}
If $X_\alpha$ is a neighbourhood of $x\in X$, then $\dim_x(X)\leq\dim(X_\alpha)$, whence the claim for open covers.
On the other hand, if the $X_\alpha$ are closed, and $U$ is a neighbourhood of $x\in X$ which meets only finitely many of the $X_\alpha$, then
\[
  \dim_x(X)\leq\dim(U)=\sup_\alpha(\dim(U\cap X_\alpha))\leq\sup_\alpha(\dim(X_\alpha))
\]
by \sref{0.14.1.4}, whence the other claim.
\end{proof}

\begin{corollary}[14.1.8]
\label{0.14.1.8}
Let $X$ be a Noetherian Kolmogoroff space \sref[0\textsubscript{I}]{0.2.1.3}, and $F$ the set of closed points of $X$.
Then $\dim(X)=\sup_{x\in F}\dim_x(X)$.
\end{corollary}

\begin{proof}
With the notation from the proof of \sref{0.14.1.6}, it suffices to note that there exists a closed point in $Z_0$ \sref[0\textsubscript{I}]{0.2.1.3}.
\end{proof}

\begin{proposition}[14.1.9]
\label{0.14.1.9}
Let $X$ be a nonempty Noetherian Kolmogoroff space.
To have $\dim(X)=0$, it is necessary and sufficient for $X$ to be finite and discrete.
\end{proposition}

\begin{proof}
If a space $X$ is separated (and \emph{a fortiori} if $X$ is a discrete space), then all the irreducible closed subsets of $X$ are single points, and so $\dim(X)=0$.
Conversely, suppose that $X$ is a Noetherian Kolmogoroff space such that $\dim(X)=0$;
since every irreducible component of $X$ contains a closed point \sref[0\textsubscript{I}]{0.2.1.3}, it must be exactly this single point.
Since $X$ has only a finite number of irreducible components, it is thus finite and discrete.
\end{proof}

\begin{corollary}[14.1.10]
\label{0.14.1.10}
Let $X$ be a Noetherian Kolmogoroff space.
For a point $x\in X$ to be isolated, it is necessary and sufficient to have $\dim_x(X)=0$.
\end{corollary}

\begin{proof}
The condition is clearly necessary (without any hypotheses on $X$).
It is also sufficient,
\oldpage[0\textsubscript{IV-1}]{8}
because it implies that $\dim(U)=0$ for any open neighbourhood $U$ of $x$, and since $U$ is a Noetherian Kolmogoroff space, $U$ is finite and discrete.
\end{proof}

\begin{proposition}[14.1.11]
\label{0.14.1.11}
The function $x\mapsto\dim_x(X)$ is upper semi-continuous on $X$.
\end{proposition}

\begin{proof}
It is clear that this function is upper semi-continuous at every point where its value is $+\infty$.
So suppose that $\dim_x(X)=n<+\infty$;
then Equation~\sref{0.14.1.4.1} shows that there exists an open neighbourhood $U_0$ of $x$ such that $\dim(U)=n$ for every open neighbourhood $U\subset U_0$ of $x$.
So, for all $y\in U_0$ and every open neighbourhood $V\subset U_0$ of $y$, we have $\dim(V)\leq\dim(U_0)=n$ \sref{0.14.1.4};
we thus deduce from \sref{0.14.1.4.1} that $\dim_y(X)\leq n$.
\end{proof}

\begin{remark}[14.1.12]
\label{0.14.1.12}
If $X$ and $Y$ are topological spaces, and $f:X\to Y$ a continuous map, then it can be the case that $\dim(f(X))>\dim(X)$;
we obtain such an example by taking $X$ to be a discrete space with 2 points, $a$ and $b$, and $Y$ to be the set $\{a,b\}$ endowed with the topology for which\footnote{\emph{[Trans.] This is now often referred to as the \emph{Sierpi\'nski space}, or the \emph{connected two-point set}.}} the closed sets are $\emp$, $\{a\}$, and $\{a,b\}$;
if $f:X\to Y$ is the identity map, then $\dim(Y)=1$ and $\dim(X)=0$.
We note that $Y$ is the spectrum of a discrete valuation ring $A$, of which $a$ is the unique closed point, and $b$ the generic point;
if $K$ and $k$ are the field of fractions and the residue field of $A$ (respectively), then $X$ is the spectrum of the ring $k\times K$, and $f$ is the continuous map corresponding to the homomorphism $(\vphi,\psi):A\to k\times K$, where $\vphi:A\to k$ and $\psi:A\to K$ are the canonical homomorphism (cf.~\sref[IV]{IV.5.4.3}).
\end{remark}

\subsection{Codimension of a closed subset}
\label{subsection:0.14.2}

\begin{definition}[14.2.1]
\label{0.14.2.1}
Given an irreducible closed subset $Y$ of a topological space $X$, we define the combinatorial codimension (or simply codimension) of $Y$ in $X$, denoted by $\codim(Y,X)$, to be the upper bound of the lengths of chains of irreducible closed subsets of $X$ of which $Y$ is the smallest element.
If $Y$ is an arbitrary closed subset of $X$, then we define the codimension of $Y$ in $X$, again denoted by $\codim(Y,X)$, to be the lower bound of the codimensions in $X$ of the irreducible components of $Y$.
We say that $X$ is equicodimensional if all the minimal irreducible closed subsets of $X$ have the same codimension in $X$.
\end{definition}

It follows from this definition that $\codim(\emp,X)=+\infty$, since the lower bound of the empty set of $\overline{\bb{R}}$ is $+\infty$.
If $Y$ is closed in $X$, and if $(X_\alpha)$ (resp.~$(Y_\alpha)$) is the family of irreducible components of $X$ (resp.~$Y$), then every $Y_\beta$ is contained in some $X_\alpha$, and, more generally, every chain of irreducible closed subsets of $X$ of which $Y_\beta$ is the smallest element is formed of subsets of some $X_\alpha$; we thus have
\[
\label{0.14.2.1.1}
  \codim(Y,X)=\inf_\beta(\sup_\alpha(\codim(Y_\beta,X_\alpha))),
  \tag{14.2.1.1}
\]
where, for every $\beta$, $\alpha$ ranges over the set of indices such that $Y_\beta\subset X_\alpha$.

\begin{proposition}[14.2.2]
\label{0.14.2.2}
Let $X$ be a topological space.
\begin{enumerate}
  \item[{\rm(i)}] If $\Phi$ is the set of irreducible closed subsets of $X$, then
    \[
    \label{0.14.2.2.1}
      \dim(X)=\sup_{Y\in\Phi}(\codim(Y,X)).
      \tag{14.2.2.1}
    \]
\oldpage[0\textsubscript{IV-1}]{9}
  \item[{\rm(ii)}] For every nonempty closed subset $Y$ of $X$, we have
    \[
    \label{0.14.2.2.2}
      \dim(Y)+\codim(Y,X)\leq\dim(X).
      \tag{14.2.2.2}
    \]
  \item[{\rm(iii)}] If $Y$, $Z$, and $T$ are closed subsets of $X$ such that $Y\subset Z\subset T$, then
    \[
    \label{0.14.2.2.3}
      \codim(Y,Z)+\codim(Z,T)\leq\codim(Y,T).
       \tag{14.2.2.3}
    \]
  \item[{\rm(iv)}] For a closed subset $Y$ of $X$ to be such that $\codim(Y,X)=0$, it is necessary and sufficient for $Y$ to contain an irreducible component of $X$.
\end{enumerate}
\end{proposition}

\begin{proof}
Claims~(i) and (iv) are immediate consequences of Definition~\sref{0.14.2.1}.
To show (ii), we can restrict to the case where $Y$ is irreducible, and then the equation follows from Definitions~\sref{0.14.1.1} and \sref{0.14.2.1}.
Finally, to show (iii), we can, by Definition~\sref{0.14.2.1}, first restrict to the case where $Y$ is irreducible;
then $\codim(Y,Z)=\sup_\alpha(\codim(Y,Z_\alpha))$ for the irreducible components $Z_\alpha$ of $Z$ that contain $Y$;
it is clear that $\codim(Y,T)\geq\codim(Y,Z)$, so the inequality is true if $\codim(Y,Z)=+\infty$;
but if this were not the case, then there would exist some $\alpha$ such that $\codim(Y,Z)=\codim(Y,Z_\alpha)$, and by \sref{0.14.2.1}, we can restrict to the case where $Z$ itself is irreducible;
but then the inequality in \sref{0.14.2.2.3} is an evident consequence of Definition~\sref{0.14.2.1}.
\end{proof}

\begin{proposition}[14.2.3]
\label{0.14.2.3}
Let $X$ be a topological space, and $Y$ a closed subset of $X$.
For every open subset $U$ of $X$, we have
\[
\label{0.14.2.3.1}
  \codim(Y\cap U,U)\geq\codim(Y,X).
  \tag{14.2.3.1}
\]

Furthermore, for this inequality \sref{0.14.2.3.1} to be an equality, it is necessary and sufficient to have $\codim(Y,X)=\inf_\alpha(\codim(Y_\alpha,X))$, where $(Y_\alpha)$ is the family of irreducible components of $Y$ that meet $U$.
\end{proposition}

\begin{proof}
We know \sref[0\textsubscript{I}]{0.2.1.6} that $Z\mapsto\overline{Z}$ is a bijection from the set of irreducible closed subsets of $U$ to the set of irreducible closed subsets of $X$ that meet $U$, and, in particular, induces a correspondence between the irreducible components of $Y\cap U$ and the irreducible components of $Y$ that meet $U$;
if $Y_\alpha$ is one of the latter such components, then we have $\codim(Y_\alpha,X)=\codim(Y_\alpha\cap U,U)$, and the proposition then follows from Definition~\sref{0.14.2.1}.
\end{proof}

\begin{definition}[14.2.4]
\label{0.14.2.4}
Let $X$ be a topological space, $Y$ a closed subset of $X$, and $x$ a point of $X$.
We define the codimension of $Y$ in $X$ at the point $x$, denoted by $\codim_x(Y,X)$, to be the number $\sup_U(\codim(Y\cap U,U))$, where $U$ ranges over the set of open neighbourhoods of $x$ in $X$.
\end{definition}

By \sref{0.14.2.3}, we can also write
\[
\label{0.14.2.4.1}
  \codim_x(Y,X)=\lim_U(\codim(Y\cap U,U)),
  \tag{14.2.4.1}
\]
where the limit is taken over the downward-directed set of open neighbourhoods of $x$ in $X$.
We note that we have
\[
  \codim_x(Y,X)=+\infty\text{ if }x\in X\setmin Y.
\]

\oldpage[0\textsubscript{IV-1}]{10}
\begin{proposition}[14.2.5]
\label{0.14.2.5}
If $(Y_i)_{1\leq i\leq n}$ is a finite family of closed subsets of a topological space $X$, and $Y$ is the union of this family, then
\[
\label{0.14.2.5.1}
  \codim(Y,X)=\inf_i(\codim(Y_i,X)).
  \tag{14.2.5.1}
\]
\end{proposition}

\begin{proof}
Every irreducible component of one of the $Y_i$ is contained in an irreducible component of $Y$, and, conversely, every irreducible component of $Y$ is also an irreducible component of one of the $Y_i$ \sref[0\textsubscript{I}]{0.2.1.1};
the conclusion then follows from Definition~\sref{0.14.2.1} and the inequality in \sref{0.14.2.2.3}.
\end{proof}

\begin{corollary}[14.2.6]
\label{0.14.2.6}
Let $X$ be a topological space, and $Y$ a locally-Noetherian closed subspace of $X$.
\begin{enumerate}
  \item[{\rm(i)}] For all $x\in X$, there exists only a finite number of irreducible components $Y_i$ ($1\leq i\leq n$) of $Y$ that contain $x$, and we have $\codim_x(Y,X)=\inf_i(\codim(Y_i,X))$.
  \item[{\rm(ii)}] The function $x\mapsto\codim_x(Y,X)$ is lower semi-continuous on $X$.
\end{enumerate}
\end{corollary}

\begin{proof}
By hypothesis, there exists an open neighbourhood $U_0$ of $x$ in $X$ such that $Y\cap U_0$ is Noetherian, and so $U_0$ has only a finite number of irreducible components, which are the intersections of $U_0$ with the irreducible components of $Y$;
\emph{a fortiori} there are only a finite number of irreducible components $Y_i$ ($1\leq i\leq n$) of $Y$ that contain $x$, and we can, by replacing $U_0$ with an open neighbourhood $U\subset U_0$ of $x$ that doesn't meet any of the $Y_j$ that do not contain $x$, assume that the $Y_i\cap U$ are the irreducible components of $Y\cap U$;
for every open neighbourhood $V\subset U$ of $x$ in $X$, the $Y_i\cap V$ are thus the irreducible components of $Y\cap V$, and \sref{0.14.2.3} then shows that $\codim(Y_i,X)=\codim(Y_i\cap V,V)$, which proves (i).
Further, for every $x'\in U$, the irreducible components of $Y$ that contain $x'$ are certain $Y_i$, and so $\codim_{x'}(Y,X)\geq\codim_x(Y,X)$, which proves (ii).
\end{proof}

\subsection{The chain condition}
\label{subsection:0.14.3}

\begin{env}[14.3.1]
\label{0.14.3.1}
In a topological space $X$, we say that a chain $Z_0\subset Z_1\subset\cdots\subset Z_n$ of irreducible closed subsets is \emph{saturated} if there does not exist an irreducible closed subset $Z'$, distinct from each of the $Z_i$, such that $Z_k\subset Z'\subset Z_{k+1}$ for any $k$.
\end{env}

\begin{proposition}[14.3.2]
\label{0.14.3.2}
Let $X$ be a topological space such that, for any two irreducible closed subsets $Y$ and $Z$ of $X$ with $Y\subset Z$, we have $\codim(Y,Z)<+\infty$.
The following two conditions are equivalent.
\begin{enumerate}
  \item[{\rm(a)}] Any two saturated chains of closed irreducible subsets of $X$ that have the same first and last elements as one another have the same length.
  \item[{\rm(b)}] If $Y$, $Z$, and $T$ are irreducible closed subsets of $X$ such that $Y\subset Z\subset T$, then
    \[
    \label{0.14.3.2.1}
      \codim(Y,T)=\codim(Y,Z)+\codim(Z,T).
      \tag{14.3.2.1}
    \]
\end{enumerate}
\end{proposition}

\begin{proof}
It is immediate that (a) implies (b).
Conversely, suppose that (b) is satisfied, and we will show that if we have two saturated chains with the same first and last elements as one another, of lengths $m$ and $n\leq m$ (respectively), then $m=n$.
We proceed by induction on $n$, with the proposition being clear for $n=1$.
So suppose that $1<n<m$, and let $Z_0\subset Z_1\subset\cdots\subset Z_n$ be a
\oldpage[0\textsubscript{IV-1}]{11}
saturated chain such that there exists another saturated chain, with first element $Z_0$ and last element $Z_n$, of length $m$.
Since $\codim(Z_0,Z_n)\geq m>n$, and $\codim(Z_0,Z_1)=1$, it follows from \emph{(b)} that $\codim(Z_1,Z_n)=\codim(Z_0,Z_n)-1>n-1$, which contradicts our induction hypothesis.
\end{proof}

When the conditions of \sref{0.14.3.2} are satisfied, we say that $X$ satisfies the \emph{chain condition}, or that it is a \emph{catenary space}.
It is clear that every closed subspace of a catenary space is catenary.

\begin{proposition}[14.3.3]
\label{0.14.3.3}
Let $X$ be a Noetherian Kolmogoroff space of finite dimension.
The following conditions are equivalent.
\begin{enumerate}
  \item[{\rm(a)}] Any two maximal chains of irreducible closed subsets of $X$ have the same length.
  \item[{\rm(b)}] $X$ is catenary and equidimensional, and its irreducible components are equicodimensional.
  \item[{\rm(c)}] $X$ is equidimensional, and, for any irreducible closed subsets $Y$ and $Z$ of $X$ with $Y\subset Z$, we have
    \[
    \label{0.14.3.3.1}
      \dim(Z)=\dim(Y)+\codim(Y,Z).
      \tag{14.3.3.1}
    \]
  \item[{\rm(d)}] $X$ is equicodimensional, and, for any irreducible closed subsets $Y$ and $Z$ of $X$ with $Y\subset Z$, we have
    \[
    \label{0.14.3.3.2}
      \codim(Y,X)=\codim(Y,Z)+\codim(Z,X).
      \tag{14.3.3.2}
    \]
\end{enumerate}
\end{proposition}

\begin{proof}
The hypotheses on $X$ imply that the first and last elements of a maximal chain of irreducible closed subsets of $X$ are necessarily a closed point and an irreducible component of $X$ (respectively) \sref[0\textsubscript{I}]{0.2.1.3};
further, every saturated chain with first element $Y$ and last element $Z$ (thus $Y\subset Z$) is contained in a maximal chain whose elements differ from those of the given chain, or are contained in $Y$, or contain $Z$.
These remarks immediately establish the equivalence between (a) and (b), and also show that if (a) is satisfied, then we have, for every irreducible closed subset $Y$ of $X$,
\[
\label{0.14.3.3.3}
  \dim(Y)+\codim(Y,X)=\dim(X);
  \tag{14.3.3.3}
\]
from \sref{0.14.3.2.1}, we immediately deduce \sref{0.14.3.3.1} and \sref{0.14.3.3.2} from \sref{0.14.3.3.3}.
Conversely, \sref{0.14.3.3.1} implies \sref{0.14.3.2.1}, and so \sref{0.14.3.3.1} implies the chain condition, by \sref{0.14.3.2};
further, by applying \sref{0.14.3.3.1} to the case where $Y$ is a single closed point $x$ of $X$, and $Z$ is an irreducible component of $X$, we get that $\codim(\{x\},X)=\dim(Z)$;
we thus conclude that (c) implies (b).
Similarly, \sref{0.14.3.3.2} implies \sref{0.14.3.2.1}, and thus the chain condition;
further, with the same choice of $Y$ and $Z$ as above, \sref{0.14.3.3.2} again implies that $\codim(\{x\},X)=\dim(Z)$, and so (since every irreducible component of $X$ contains a closed point, by \sref[0\textsubscript{I}]{0.2.1.3}), (d) implies (b).
\end{proof}

We say that a Noetherian Kolmogoroff space is \emph{biequidimensional} if it is of \emph{finite dimension} and if it verifies any of the (equivalent) conditions of \sref{0.14.3.3}.

\begin{corollary}[14.3.4]
\label{0.14.3.4}
Let $X$ be a biequidimensional Noetherian Kolmogoroff space;
then, for every closed point $x$ of $X$, and every irreducible component $Z$ of $X$, we have
\[
\label{0.14.3.4.1}
  \dim(X)=\dim(Z)=\codim(\{x\},X)=\dim_x(X).
  \tag{14.3.4.1}
\]
\end{corollary}

\oldpage[0\textsubscript{IV-1}]{12}
\begin{proof}
The last equality follows from the fact that, if $Y_0=\{x\}\subset Y_1\subset\cdots\subset Y_m$ is a maximal chain of irreducible closed subsets of $X$, and $U$ an open neighbourhood of $x$, then the $U\cap Y_i$ are pairwise distinct irreducible closed subsets of $U$ (because $\overline{U\cap Y_i}=Y_i$), whence $\dim(U)=\dim(X)$, by \sref{0.14.1.4}.
\end{proof}

\begin{corollary}[14.3.5]
\label{0.14.3.5}
Let $X$ be a Noetherian Kolmogoroff space; if $X$ is biequidimensional, then so is every union of irreducible components of $X$, and every irreducible closed subset of $X$.
In addition, for every closed subset $Y$ of $X$, we have
\[
\label{0.14.3.5.1}
  \dim(Y)+\codim(Y,X)=\dim(X).
  \tag{14.3.5.1}
\]
\end{corollary}

\begin{proof}
Every chain of irreducible closed subsets of $X$ is contained in an irreducible component of $X$, and so the first claim follows immediately from \sref{0.14.3.3}.
Further, if $X'$ is an irreducible closed subset of $X$, then $X'$ trivially satisfies the conditions of \sref{0.14.3.3}[c], whence the second claim.

Finally, to show \sref{0.14.3.5.1}, note that we have seen, in the proof of \sref{0.14.3.3}, that this equation is true whenever $Y$ is irreducible;
if $Y_i$ ($1\leq i\leq m$) are the irreducible components of $Y$, then the $Y_i$ for which $\dim(Y_i)$ is the largest are also those for which $\codim(Y_i,X)$ is the smallest;
so \sref{0.14.3.5.1} follows from the definitions of $\dim(Y)$ and $\codim(Y,X)$.
\end{proof}

\begin{remark}[14.3.6]
\label{0.14.3.6}
The reader will note that the proof of \sref{0.14.3.2} applies to any ordered set, and the fact that we are working with the example of a set of irreducible closed subsets of a topological space is not used at all in the proof.
It is the same in the proof of \sref{0.14.3.3}, which holds, more generally, for any ordered set $E$ such that, for all $x\in E$, there exists some $z\leq x$ which is \emph{minimal} in $E$, and such that the length of chains of elements of $E$ is bounded.
\end{remark}
% END NATIVE INPUT ega0/ega0-14
% BEGIN NATIVE INPUT ega0/ega0-15
\section{\texorpdfstring{$M$}{M}-regular sequences and \texorpdfstring{$\sh{F}$}{F}-regular sequences}
\label{section:0.15}

\subsection{M-regular sequences and M-quasi-regular sequences}
\label{subsection:0.15.1}
\label{0.15.1}

\begin{env}[15.1.1]
\label{0.15.1.1}
\oldpage[0\textsubscript{IV-1}]{12}
Let $A$ be a ring, $N$ an $A$-module;
in what follows, we will write $N[T_1,\ldots,T_n]$, for short, for the graded $A$-module $N\otimes_A A[T_1,\ldots,T_n]$ ($T_i$ indeterminates), where $N$ is graded by the trivial graduation and $A[T_1,\ldots,T_n]$ by the total degree.

Let $M$ be an $A$-module, $f_i$ ($1\leq i\leq n$) elements of $A$, $\mathfrak{J}=\sum_{i=1}^n f_iA$ the ideal of $A$ generated by the $f_i$, and let us equip $M$ with the $\mathfrak{J}$-preadic filtration;
we then define a \emph{surjective graded homomorphism of graded $A$-modules}, of degree $0$,
\[
\label{0.15.1.1.1}
  \varphi : (\gr_0(M))[T_1,\ldots,T_n] \to \gr_\bullet(M)
\tag{15.1.1.1}
\]
\oldpage[0\textsubscript{IV-1}]{13}
in the following way:
if $\xi\in\gr_0(M)$ is the class mod~$\mathfrak{J}M$ of an element $x\in M$, to the pair formed by $\xi$ and a homogeneous polynomial $P[T_1,\ldots,T_n]$ of degree $k$, we associate the class mod~$\mathfrak{J}^{k+1}M$ of $P(f_1,\ldots,f_n)x$;
it is immediate that this class depends only on $\xi$ and on $P$, and if $S_k$ is the set of homogeneous polynomials of degree $k$ in $A[T_1,\ldots,T_n]$, we have thus defined an $A$-linear map
\[
  \varphi_k : \gr_0(M)\otimes_A S_k \to \gr_k(M),
\]
which is surjective by definition of $\mathfrak{J}^k M$.

We will say in what follows that an element $f\in A$ is \emph{not a zero divisor in an $A$-module} $M$ if the homothety $f_M:x\mapsto fx$ of $M$ is injective.
\end{env}

\begin{lemma}[15.1.2]
\label{0.15.1.2}
Let $M$ be an $A$-module, $f$ an element of $A$, and let us equip $M$ with the $(fA)$-preadic filtration.
If $f$ is not a zero divisor in $M$, the canonical homomorphism defined in \sref{0.15.1.1}:
\[
\label{0.15.1.2.1}
  \varphi : (M/fM)[T] \to \gr_\bullet(M)
\tag{15.1.2.1}
\]
is bijective.
The converse is true if $M$ is separated for the $(fA)$-preadic topology.
\end{lemma}

\begin{proof}
We have here $\gr_k(M)=f^kM/f^{k+1}M$;
to say that $\varphi$ is injective means that for all $x\in M$ and all $k$, the relation $f^kx\in f^{k+1}M$ implies $x\in fM$;
it is clear that this is indeed the case if $f$ is not a zero divisor in $M$.
Conversely, note that the homothety $f_M:x\mapsto fx$ defines, by passing to quotients, homomorphisms $\gr_k(M)\to\gr_{k+1}(M)$, hence a graded homomorphism $g$ of degree $1$ of $\gr_\bullet(M)$ into itself;
and to say that $\varphi$ is injective also means that $g$ is injective.
If we further suppose the filtration of $M$ separated, it follows that $f_M$ is injective
(Bourbaki, \emph{Alg.~comm.}, chap.~III, \S~2, no.~8, cor.~1 of th.~1).
\end{proof}

\begin{corollary}[15.1.3]
\label{0.15.1.3}
Let $A$ be a Noetherian ring, $M$ a finitely generated $A$-module.
For $f$ to be not a zero divisor in $M$, it is necessary and sufficient that $\varphi$ be bijective.
\end{corollary}

\begin{proof}
We need only show that the condition is sufficient;
to prove that $f_M$ is injective, it suffices to show that for every prime ideal $\mathfrak{p}$ of $A$, the homothety of ratio $f'=f/1\in A'=A_\mathfrak{p}$ in the $A_\mathfrak{p}$-module $M_\mathfrak{p}=M'$ is injective
(Bourbaki, \emph{Alg.~comm.}, chap.~II, \S~3, no.~4, th.~1).
Now $f'^kM'/f'^{k+1}M'=(f^kM/f^{k+1}M)\otimes_A A'$, and it is immediate that the canonical homomorphism $\varphi':(M'/f'M')[T]\to\gr_\bullet(M')$ is equal to $\varphi\otimes 1_{A'}$;
if $\varphi$ is injective, so is $\varphi'$ ($\mathbf{0_I}$, 1.3.2).
We are thus reduced to the case where $A$ is a Noetherian \emph{local} ring, whose maximal ideal we denote by $\mathfrak{m}$.
If $f\notin\mathfrak{m}$, then $f$ is invertible and there is nothing to prove.
If $f\in\mathfrak{m}$, we know that $M$ is separated for the $(fA)$-preadic topology ($\mathbf{0_I}$, 7.3.5), and it suffices therefore to apply \sref{0.15.1.2}.
\end{proof}

\begin{definition}[15.1.4]
\label{0.15.1.4}
Let $A$ be a ring, $M$ an $A$-module.
We say that an element $f\in A$ is \emph{$M$-regular} (resp. \emph{$M$-quasi-regular}) if $f$ is not a zero divisor in $M$ (resp. if the homomorphism \sref{0.15.1.2.1} is bijective).
\end{definition}

We have thus just shown that every $M$-regular element is $M$-quasi-regular, and that the converse is true if $A$ is Noetherian and $M$ a finitely generated $A$-module.
When $M$ is no longer assumed to be finitely generated, the converse can fail:
for example, if $A$ is a principal ideal domain that is not a field, $K$ its field of fractions, and $M$ the $A$-module $K/A$, we have $fM=M$ for all $f\neq 0$ in $A$, hence both sides of \sref{0.15.1.2.1} reduce to $0$, but $f$ is then a zero divisor in $M$.

\begin{env}[15.1.5]
\label{0.15.1.5}
\oldpage[0\textsubscript{IV-1}]{14}
Suppose now that $M$ is equipped with a \emph{filtration} $(M_k)$ such that $M_0=M$, and denote by $\gr_\bullet(M)=\bigoplus_{k\geq 0}M_k/M_{k+1}$ the associated graded $A$-module.
The notations being those of \sref{0.15.1.1}, set, for all $k\geq 0$,
\[
\label{0.15.1.5.1}
  M'_k = M_k + \mathfrak{J}M_{k-1} + \cdots + \mathfrak{J}^{k-1}M_1 + \mathfrak{J}^k M_0.
\tag{15.1.5.1}
\]
It is clear that $(M'_k)$ is a filtration on $M$.
When $M_k=\mathfrak{Q}^k M$, where $\mathfrak{Q}$ is an ideal of $A$, the filtration $(M'_k)$ is none other than the $(\mathfrak{J}+\mathfrak{Q})$-preadic filtration.
We will again write $(\gr_\bullet(M)\otimes_A(A/\mathfrak{J}))[T_1,\ldots,T_n]$ for the graded tensor product of graded $A$-modules ($\mathbf{II}$, 2.1.2)
$(\gr_\bullet(M)\otimes_A(A/\mathfrak{J}))\otimes_A A[T_1,\ldots,T_n]$.
We define a \emph{surjective graded homomorphism of graded $A$-modules}, of degree $0$,
\[
\label{0.15.1.5.2}
  \psi : (\gr_\bullet(M)\otimes_A(A/\mathfrak{J}))[T_1,\ldots,T_n] \to \gr'_\bullet(M)
\tag{15.1.5.2}
\]
in the following way:
if $\alpha\in A/\mathfrak{J}$ is the class of an element $a\in A$, $\xi\in M_k/M_{k+1}$ the class of an element $x\in M_k$, and $P[T_1,\ldots,T_n]$ a homogeneous polynomial of degree $h$, we associate to the triplet $(\alpha,\xi,P)$ the class mod~$M'_{k+h+1}$ of $P(f_1,\ldots,f_n)ax$;
one verifies immediately that this element depends only on $\xi$, $\alpha$, and $P$, and this defines the surjective graded homomorphism \sref{0.15.1.5.2}, by virtue of the definition \sref{0.15.1.5.1}.
One recovers the homomorphism \sref{0.15.1.1.1} by considering the filtration on $M$ such that $M_1=0$.
\end{env}

\begin{lemma}[15.1.6]
\label{0.15.1.6}
We keep the notations of \sref{0.15.1.5} and suppose $n=1$, setting $f=f_1$.
If $f$ is $\gr_\bullet(M)$-regular, the homomorphism \sref{0.15.1.5.2} is bijective.
\end{lemma}

\begin{proof}
Let $Q_k$ (resp. $Q'_k$) be the sub-$A$-module of elements of degree $k$ in the first (resp. second) member of \sref{0.15.1.5.2}.
We equip $Q_k$ with the filtration
\[
\label{0.15.1.6.1}
  (Q_k)_i = \sum_{j\leq k-i} (\gr_{k-j}(M)\otimes_A(A/fA))T^j
\tag{15.1.6.1}
\]
so that $(Q_k)_0=Q_k$ and $(Q_k)_{k+1}=0$, and for this filtration, one has
\[
  \gr_i(Q_k) = (\gr_i(M)\otimes_A(A/fA))T^{k-i}.
\]
We equip $Q'_k$ with the filtration formed by $\psi((Q_k)_i)=(Q'_k)_i$.
As these filtrations are \emph{finite}, it suffices to prove that the homomorphisms $\psi_{ki}:\gr_i(Q_k)\to\gr_i(Q'_k)$ induced by $\psi_k$ are injective
(Bourbaki, \emph{Alg.~comm.}, chap.~III, \S~3, no.~8, cor.~1 of th.~1).
Now $\gr_i(Q_k)=((M_i/M_{i+1})\otimes_A(A/fA))T^{k-i}$;
on the other hand, $(Q'_k)_{i+1}$ is the image of
\[
  M_k+fM_{k-1}+\cdots+f^{k-i-1}M_{i+1}
\]
\oldpage[0\textsubscript{IV-1}]{15}
in $M'_k/M'_{k+1}$.
We are thus led to write that, for $x\in M_i$, the relation
\[
\label{0.15.1.6.2}
  f^{k-i}x\in M_k+fM_{k-1}+\cdots+f^{k-i-1}M_{i+1}+M'_{k+1}
\tag{15.1.6.2}
\]
implies $x\in fM_i+M_{i+1}$.
Now, the second member of \sref{0.15.1.6.2} is contained in $M_{i+1}+M'_{k+1}$, and by virtue of \sref{0.15.1.5.1}, $M'_{k+1}$ is contained in $M_{i+1}+f^{k+1-i}M$.
The hypothesis on $f$ implies that $f$ is not a zero divisor in $M/M_h$ for all $h$
(Bourbaki, \emph{loc.~cit.}).
The relation $f^{k-i}x\in f^{k+1-i}M+M_{i+1}$ thus implies (reasoning by induction) $x\in fM_i+M_{i+1}$.
\end{proof}

\begin{definition}[15.1.7]
\label{0.15.1.7}
Let $(f_i)_{1\leq i\leq n}$ be a sequence of elements of $A$.
We say that it is \emph{$M$-regular} if, for $1\leq i\leq n$, $f_i$ is $(M/(\sum_{1\leq j\leq i-1}f_jM))$-regular.
We say that the sequence $(f_i)$ is \emph{$M$-quasi-regular} if the canonical homomorphism $\varphi$ defined in \sref{0.15.1.1.1} is bijective.

When $M=A$, we say simply ``regular sequence'' and ``quasi-regular sequence''.
\end{definition}

\begin{proposition}[15.1.8]
\label{0.15.1.8}
The notations being those of \sref{0.15.1.5}, suppose that the sequence $(f_i)_{1\leq i\leq n}$ is $\gr_\bullet(M)$-regular.
Then the canonical homomorphism \sref{0.15.1.5.2} is bijective.
\end{proposition}

\begin{proof}
The proposition is none other than \sref{0.15.1.6} for $n=1$;
let us reason by induction on $n$.
Set $\mathfrak{J}''=\sum_{i=1}^{n-1}f_iA$, so that $\mathfrak{J}=\mathfrak{J}''+f_nA$, and let
\[
  M''_k = M_k+\mathfrak{J}''M_{k-1}+\cdots+\mathfrak{J}''^{k-1}M_1+\mathfrak{J}''^kM_0;
\]
one can then write
\[
  M'_k = M''_k + f_nM''_{k-1}+\cdots+f_n^{k-1}M''_1+f_n^kM''_0.
\]
Let us write $\gr''_\bullet(M)$ for the graded $A$-module associated to $M$ filtered by the filtration $(M''_k)$.
One can write, up to canonical isomorphisms,
\[
  (\gr_\bullet(M)\otimes(A/\mathfrak{J}))\otimes_A A[T_1,\ldots,T_n]
  = (\gr_\bullet(M)\otimes(A/\mathfrak{J}''))[T_1,\ldots,T_{n-1}]\otimes(A/f_nA)\otimes_A A[T_n]
\]
and by \sref{0.15.1.5.2}, $\psi$ factorizes as
\[
\label{0.15.1.8.1}
  (\gr''_\bullet(M)\otimes_A(A/f_nA))[T_n] \to \gr'_\bullet(M)
\tag{15.1.8.1}
\]
and
\[
  ((\gr_\bullet(M)\otimes(A/\mathfrak{J}''))[T_1,\ldots,T_{n-1}])\otimes(A/f_nA)\otimes_A A[T_n]
  \xrightarrow{\psi'\otimes 1} \gr''_\bullet(M)\otimes(A/f_nA)\otimes_A A[T_n]
\]
where $\psi'$ is the map \sref{0.15.1.5.2} in which one replaces $n$ by $n-1$, $\mathfrak{J}$ by $\mathfrak{J}''$, and $\gr'_\bullet(M)$ by $\gr''_\bullet(M)$.
The induction hypothesis implies that $\psi'$ is bijective, and it remains to see, by applying \sref{0.15.1.6}, that it suffices to show that $f_n$ is $\gr''_\bullet(M)$-regular.
But $\gr''_\bullet(M)$ identifies with $(\gr_\bullet(M)\otimes(A/\mathfrak{J}''))[T_1,\ldots,T_{n-1}]$;
by hypothesis, the homothety of ratio $f_n$ in $\gr_\bullet(M)\otimes(A/\mathfrak{J}'')$ is injective;
since $A[T_1,\ldots,T_{n-1}]$ is a free $A$-module, it follows that $f_n$ is $\gr''_\bullet(M)$-regular.
\end{proof}

\begin{proposition}[15.1.9]
\label{0.15.1.9}
Let $M$ be an $A$-module, $(f_i)_{1\leq i\leq n}$ a sequence of elements of $A$.
If the sequence $(f_i)$ is $M$-regular, it is $M$-quasi-regular.
The converse is true when one supposes in addition that the $A$-modules $M$, $M/f_1M$, \ldots, $M/(\sum_{1\leq i\leq n-1}f_iM)$ are separated for the $\mathfrak{J}$-preadic topology.
\end{proposition}

\begin{proof}
The first assertion follows from \sref{0.15.1.8}, taking $M_1=0$, and then the homomorphism \sref{0.15.1.5.2} reduces to \sref{0.15.1.1.1}.
Conversely, suppose that the homomorphism $\varphi$ of \sref{0.15.1.1.1} is bijective.
Note that if we set
\[
  \mathfrak{J}''=\sum_{i=1}^{n-1}f_iA,
\]
\oldpage[0\textsubscript{IV-1}]{16}
and denote by $\gr''_\bullet(M)$ the graded $A$-module associated to $M$ filtered by the $\mathfrak{J}''$-preadic filtration, $\varphi$ factorizes as
\[
\label{0.15.1.9.1}
  (\gr''_\bullet(M)\otimes(A/f_nA))[T_n] \to \gr_\bullet(M)
\tag{15.1.9.1}
\]
and $\varphi'\otimes 1$, where $\varphi'$ is the canonical homomorphism
\[
  (\gr''_0(M))[T_1,\ldots,T_{n-1}] \to \gr''_\bullet(M).
\]
As these two homomorphisms are surjective, they are both bijective if $\varphi$ is supposed bijective, hence $\varphi'$ is necessarily injective, and thus bijective since it is surjective.
One can then reason by induction on $n$ (the case $n=1$ follows from \sref{0.15.1.2}), since on $M$, $M/f_1M$, \ldots, $M/(\sum_{j=1}^{n-2}f_jM)$ the $\mathfrak{J}''$-preadic topology is finer than the $\mathfrak{J}$-preadic topology, and is thus separated.
One sees therefore that the sequence $(f_i)_{1\leq i\leq n-1}$ is $M$-regular.
On the other hand, let us show that if $y\in M/\mathfrak{J}''M$ is such that $f_n^ky\in f_n^{k+1}(M/\mathfrak{J}''M)$, then $y\in f_n(M/\mathfrak{J}''M)$.
Indeed, if $x$ is a representative of $y$, let us show that $f_n^kx\in\mathfrak{J}^{k+1}M$;
our assertion will then follow from the fact that \sref{0.15.1.9.1} is injective.
The hypothesis implies $f_n^kx\in f_n^{k+1}M+\mathfrak{J}''M\subset\mathfrak{J}M$, whence, since $\varphi$ is injective, $f_n^{k-1}x\in\mathfrak{J}^kM$;
by induction, one thus obtains $x\in\mathfrak{J}M$, and finally $f_n^kx\in\mathfrak{J}^{k+1}M$.
We finish the reasoning as in \sref{0.15.1.2}:
on $M/\mathfrak{J}''M$, the $\mathfrak{J}$-preadic filtration is identical to the $(f_nA)$-preadic filtration, which is separated by hypothesis, and one concludes that the homothety of ratio $f_n$ in $M/\mathfrak{J}''M$ is injective, which finishes showing that the sequence $(f_i)_{1\leq i\leq n}$ is $M$-regular.
\end{proof}

\begin{corollary}[15.1.10]
\label{0.15.1.10}
With the notations of \sref{0.15.1.9}, suppose that for every finitely generated $A$-module $N$, $M\otimes_AN$ is separated for the $\mathfrak{J}$-preadic topology.
Then, for the sequence $(f_i)$ to be $M$-regular, it is necessary and sufficient that it be $M$-quasi-regular.
\end{corollary}

One will note that the hypothesis of \sref{0.15.1.10} implies that, if the sequence $(f_i)_{1\leq i\leq n}$ is $M$-regular, so is the sequence $(f_{\sigma(i)})$ for every permutation $\sigma$ of the interval $[1,n]$.

\begin{corollary}[15.1.11]
\label{0.15.1.11}
Let $A$ be a Noetherian ring, $M$ a finitely generated $A$-module, $(f_i)_{1\leq i\leq n}$ a finite sequence of elements contained in the radical of $A$.
Then, for the sequence $(f_i)$ to be $M$-regular, it is necessary and sufficient that it be $M$-quasi-regular.
\end{corollary}

\begin{proof}
This is a particular case of \sref{0.15.1.10}, since every finitely generated $A$-module is then separated for the $\mathfrak{J}$-preadic topology ($\mathbf{0_I}$, 7.3.5).
\end{proof}

\begin{remarks}[15.1.12]
\label{0.15.1.12}
\mbox{}
\begin{enumerate}
  \item[(i)] When $A$ is a Noetherian local ring, $M$ a finitely generated $A$-module, and when the $f_i$ are contained in the maximal ideal of $A$, the notion of $M$-regular sequence was introduced by J.-P.~Serre under the name of $M$-\emph{suite} \cite{I-17}.
  \item[(ii)] Even if $A$ is Noetherian and $M=A$, a sequence $(f,g)$ of two elements of $A$ that is quasi-regular is not necessarily regular.
  For example, the sequence $(0,1)$ is quasi-regular, since then $\mathfrak{J}=fA+gA=A$, hence $\gr_\bullet(A)=0$;
  but if $A\neq 0$ it is not regular, since $0$ is not $A$-regular.
  One will note however that in this case the sequence $(1,0)$ is regular, since $1$ is not a zero divisor in $A$, and since $A/A=0$, $0$ is not a zero divisor in this module (cf. \sref{0.15.1.1}).
  \item[(iii)]
\oldpage[0\textsubscript{IV-1}]{17}
  In \sref{0.15.1.6}, one cannot in general replace the hypothesis that $f$ is $\gr_\bullet(M)$-regular by the hypothesis that it is $\gr_\bullet(M)$-quasi-regular.
  To see this, let us take the example \sref{0.15.1.4} where $A$ is a principal ideal domain with field of fractions $K\neq A$, and let us take $M=K$, the filtration $(M_k)$ of $M$ being defined by $M_0=M=K$, $M_1=A$, $M_k=0$ for $k\geq 2$, whence $\gr_0(M)=K/A$, $\gr_1(M)=A$, $\gr_k(M)=0$ for $k\geq 2$.
  If $f\neq 0$ in $A$, one sees that $f$ is not $\gr_\bullet(M)$-regular, but since $f^k\gr_\bullet(M)/f^{k+1}\gr_\bullet(M)=f^kA/f^{k+1}A$, it is immediate that the homomorphism \sref{0.15.1.2.1} (where $M$ is replaced by $\gr_\bullet(M)$) is injective, hence $f$ is $\gr_\bullet(M)$-quasi-regular.
  But with the notations of \sref{0.15.1.6}, one then has $M'_k=K$ for all $k\geq 0$, hence $\gr'_\bullet(M)=0$, while the first member of \sref{0.15.1.5.2} is not reduced to $0$.
\end{enumerate}
\end{remarks}

\begin{proposition}[15.1.13]
\label{0.15.1.13}
Let $A$ be a ring, $M$, $N$ two $A$-modules, $(f_i)_{1\leq i\leq n}$ a sequence of elements of $A$.
If the sequence $(f_i)$ is $M$-regular and if $N$ is a flat $A$-module, then the sequence $(f_i)$ is $(M\otimes_AN)$-regular.
The converse is true if $N$ is a faithfully flat $A$-module.
\end{proposition}

\begin{proof}
Indeed (without any hypothesis on $N$), $(M\otimes_AN)/(\sum_{j=1}^{i-1}f_j(M\otimes_AN))$ is isomorphic to $(M/(\sum_{j=1}^{i-1}f_jM))\otimes_AN$;
it suffices then to apply ($\mathbf{0_I}$, 6.1.1) and ($\mathbf{0_I}$, 6.4.1) to the homothety of ratio $f_i$ in $M/(\sum_{j=1}^{i-1}f_jM)$.
\end{proof}

\begin{corollary}[15.1.14]
\label{0.15.1.14}
Let $A$ be a ring, $M$ an $A$-module, $(f_i)_{1\leq i\leq n}$ a finite sequence of elements of $A$.
Let $\rho:A\to B$ be a ring homomorphism, $M'=M_{(B)}=M\otimes_AB$, $f'_i=\rho(f_i)$ for $1\leq i\leq n$.
If $B$ is a flat $A$-module and if the sequence $(f_i)$ is $M$-regular, then the sequence $(f'_i)$ is $M'$-regular;
the converse is true if $B$ is a faithfully flat $A$-module.
\end{corollary}

\begin{proof}
Indeed, $M'/(\sum_{j=1}^{i-1}f'_jM')$ identifies with $(M/(\sum_{j=1}^{i-1}f_jM))\otimes_AB$, and the homothety of ratio $f'_i$ in the $B$-module $M'/(\sum_{j=1}^{i-1}f'_jM')$ is identified with the homothety of ratio $f_i$ in the $A$-module $(M/(\sum_{j=1}^{i-1}f_jM))\otimes_AB$;
the assertion thus follows from \sref{0.15.1.13}.
\end{proof}

\begin{proposition}[15.1.15]
\label{0.15.1.15}
Let $A$ be a ring, $B$ a commutative $A$-algebra, $(f_i)_{1\leq i\leq n}$ a finite sequence of elements of $B$, $M$ a $B$-module.
Suppose that the sequence $(f_i)$ is $M$-regular and that each of the $A$-modules $M/(\sum_{j=1}^{i-1}f_jM)$ ($1\leq i\leq n+1$) is flat.
Then, for every ring homomorphism $\rho:A\to A'$, setting $B'=B\otimes_AA'$, $M'=M\otimes_AA'$, and $f'_i=f_i\otimes 1$ ($1\leq i\leq n$), the sequence $(f'_i)$ is $M'$-regular and the $A'$-modules $M'/(\sum_{j=1}^{i-1}f'_jM')$ ($1\leq i\leq n+1$) are flat.
\end{proposition}

\begin{proof}
Consider the exact sequence (by hypothesis)
\[
  0 \to M \xrightarrow{f_1} M \to M/f_1M \to 0
\]
(the arrow $M\xrightarrow{f_1}M$ being the homothety of ratio $f_1$ in $M$).
The hypothesis that $M/f_1M$ is $A$-flat implies that the sequence
\[
  0 \to M\otimes_AA' \xrightarrow{f_1\otimes 1} M\otimes_AA' \to (M/f_1M)\otimes_AA' \to 0
\]
\oldpage[0\textsubscript{IV-1}]{18}
is exact ($\mathbf{0_I}$, 6.1.2).
This shows that the homothety of ratio $f'_1$ in $M'$ is injective.

Set then, for $1\leq i\leq n-1$, $M_i=M/(\sum_{j=1}^i f_jM)$, $M'_i=M'/(\sum_{j=1}^i f'_jM')$;
we have $M'_i=M_i\otimes_AA'$, $M_i/f_{i+1}M_i=M_{i+1}$, $M'_i/f'_{i+1}M'_i=M'_{i+1}$;
it suffices to replace $M$ and $f_1$ by $M_i$ and $f_{i+1}$ to conclude that the homothety of ratio $f'_{i+1}$ in $M'_i$ is injective, by reason of the flatness of $M_i/f_{i+1}M_i$.
The last assertion follows from ($\mathbf{0_I}$, 6.2.1).
\end{proof}

\begin{proposition}[15.1.16]
\label{0.15.1.16}
Let $A$, $B$ be two Noetherian local rings, $k$ the residue field of $A$, $\varphi:A\to B$ a local homomorphism, $M$ a finitely generated $B$-module.
Let $(f_i)_{1\leq i\leq n}$ be a sequence of elements of the maximal ideal of $B$, and for each $i$, let $g_i$ be the image of $f_i$ in $B\otimes_Ak$.
The following conditions are equivalent:
\begin{enumerate}
  \item[a)] The sequence $(f_i)$ is $M$-regular and the quotients $M_i=M/(\sum_{j=1}^i f_jM)$ are flat $A$-modules for $1\leq i\leq n$.
  \item[b)] The sequence $(f_i)$ is $M$-regular and the $A$-module $M_n$ is flat.
  \item[c)] The $A$-module $M$ is flat and the sequence $(g_i)_{1\leq i\leq n}$ is $(M\otimes_Ak)$-regular.
  \item[d)] The $A$-module $M$ is flat, and for every ring homomorphism $\rho:A\to A'$, if we set $B'=B\otimes_AA'$, $M'=M\otimes_AA'$, $f'_i=f_i\otimes 1$ ($1\leq i\leq n$), the sequence $(f'_i)$ is $M'$-regular.
\end{enumerate}
\end{proposition}

\begin{proof}
It is clear that a) implies b);
to show that b) implies d), it suffices, by virtue of \sref{0.15.1.15}, to show that b) implies that the $M_i$ are flat $A$-modules for $0\leq i\leq n-1$ (setting $M_0=M$).
As $M_{i+1}=M_i/f_{i+1}M_i$, we are reduced, by descending induction, to proving that if $h$ is an element of the maximal ideal of $B$ such that the homothety of ratio $h$ in $M$ is injective and $M/hM$ is a flat $A$-module, then $M$ is a flat $A$-module;
but this is none other than ($\mathbf{0_{III}}$, 10.2.7).
It is clear that c) is a particular case of d) obtained by taking $A'=k$.
It thus remains to prove that c) implies a).
As $M$ is a flat $A$-module and a finitely generated $B$-module, it follows from c) and from ($\mathbf{0_{III}}$, 10.2.4) that the homothety of ratio $f_1$ in $M$ is injective and that its cokernel $M/f_1M=M_1$ is a flat $A$-module.
Let us reason by induction on $i$ and suppose shown that $M_i$ is a flat $A$-module;
as the homothety of ratio $g_{i+1}$ in $M_i\otimes_Ak$ is injective by hypothesis, it follows again from ($\mathbf{0_{III}}$, 10.2.4) that the homothety of ratio $f_{i+1}$ in $M_i$ is injective and that $M_i/f_{i+1}M_i=M_{i+1}$ is a flat $A$-module, which finishes the proof.
One will note that the proof that c) implies a) does not use the fact that the $f_i$ belong to the maximal ideal of $B$.
\end{proof}

\begin{corollary}[15.1.17]
\label{0.15.1.17}
The hypotheses on $A$ and $B$ being those of \sref{0.15.1.16}, let $M$ be a finitely generated $B$-module, flat over $A$, $R$ a sub-$B$-module of $M$ such that $N=M/R$ is a flat $A$-module.
Let $(f_i)_{1\leq i\leq n}$ be a sequence of elements of the maximal ideal of $B$ having the following properties:
\begin{enumerate}
  \item[(i)] $f_iM\subset R$ for $1\leq i\leq n$.
  \item[(ii)] If $g_i$ is the image of $f_i$ in $B\otimes_Ak$, the sequence $(g_i)_{1\leq i\leq n}$ is $(M\otimes_Ak)$-regular.
  \item[(iii)] The sum of the sub-$(B\otimes_Ak)$-modules $g_i(M\otimes_Ak)$ of $M\otimes_Ak$ is the image of $R\otimes_Ak$ in $M\otimes_Ak$ (equal to the kernel of the canonical homomorphism $M\otimes_Ak\to N\otimes_Ak$).
\end{enumerate}
Then the sequence $(f_i)$ is $M$-regular and $R=\sum_{i=1}^n f_iM$.
\end{corollary}

\begin{proof}
\oldpage[0\textsubscript{IV-1}]{19}
As $N$ is a flat $A$-module, it follows from the exactness of the sequence $0\to R\to M\to N\to 0$ that the sequence $0\to R\otimes_Ak\to M\otimes_Ak\to N\otimes_Ak\to 0$ is exact ($\mathbf{0_I}$, 6.1.2);
we thus have $R\otimes_Ak=\sum g_i(M\otimes_Ak)$.
As $M$ is a finitely generated $B$-module and $B$ is Noetherian, $R$ is also a finitely generated $B$-module, whence $R=\sum f_iM$ by Nakayama's lemma, since the $f_i$ belong to the maximal ideal of $B$.
\end{proof}

\begin{lemma}[15.1.18]
\label{0.15.1.18}
Let $A$ be a ring, $M$ an $A$-module equipped with a separated and exhaustive filtration $(M_k)_{k\geq 0}$, and let $f\in A$ be a $\gr_\bullet(M)$-regular element.
Then:
\begin{enumerate}
  \item[(i)] $f$ is $M$-regular.
  \item[(ii)] For all $k$, one has $M_k\cap fM=fM_k$, and the canonical image $M'_k$ of $M_k$ in $M'=M/fM$ is thus canonically isomorphic to $M_k/fM_k$.
  \item[(iii)] The sequence
  \[
    0 \to M_{k+1}/fM_{k+1} \to M_k/fM_k \to \gr_k(M)/f\cdot\gr_k(M) \to 0
  \]
  is exact, and if one equips $M'$ with the quotient filtration $(M'_k)$, $\gr_\bullet(M')$ is thus canonically isomorphic to $\gr_\bullet(M)/f\cdot\gr_\bullet(M)$.
\end{enumerate}
\end{lemma}

\begin{proof}
Assertion (i) follows from the hypothesis on $f$ and from Bourbaki, \emph{Alg.~comm.}, chap.~III, \S~2, no.~8, cor.~1 of th.~1.
As $\gr_\bullet(M/M_k)$ is a direct factor of $\gr_\bullet(M)$, $f$ is $\gr_\bullet(M/M_k)$-regular for every $k$, hence $(M/M_k)$-regular by (i), which establishes (ii).
To prove (iii), it suffices to see that the map $M_{k+1}/fM_{k+1}\to M_k/fM_k$ is injective:
now, if $x\in M_{k+1}$ is such that $x\in fM_k$, it follows from (ii) that $x\in fM_{k+1}$, which finishes proving the lemma.
\end{proof}

\begin{proposition}[15.1.19]
\label{0.15.1.19}
Let $A$ be a ring, $M$ an $A$-module equipped with a filtration $(M_k)_{k\geq 0}$ such that $M_0=M$, $(f_i)_{1\leq i\leq n}$ a $\gr_\bullet(M)$-regular sequence of elements of $A$.
Suppose in addition that, for $0\leq i\leq n-1$, the quotient filtration of $(M_k)$ on $M/(\sum_{j\leq i}f_jM)$ is separated.
Then the sequence $(f_i)$ is $M$-regular, and if one sets $M'=M/(\sum_{i\leq n}f_iM)$ and equips $M'$ with the quotient filtration of $(M_k)$, $\gr_\bullet(M')$ is canonically isomorphic to $\gr_\bullet(M)/(\sum_{i\leq n}f_i\gr_\bullet(M))$.
\end{proposition}

\begin{proof}
For $n=1$, the proposition follows from the lemma \sref{0.15.1.18}.
Let us reason by induction on $n$, setting $N=M/(\sum_{i\leq n-1}f_iM)$ and denoting by $(N_k)$ the quotient filtration of $(M_k)$ on $N$, which is separated by hypothesis;
moreover, $\gr_\bullet(N)$ is isomorphic to
\[
  \gr_\bullet(M)/\Big(\sum_{i\leq n-1}f_i\gr_\bullet(M)\Big).
\]
By hypothesis, $f_n$ is thus $\gr_\bullet(N)$-regular, and it follows from the lemma \sref{0.15.1.18} applied to $N$ that $f_n$ is $N$-regular and $\gr_\bullet(M')$ is isomorphic to $\gr_\bullet(N)/f_n\gr_\bullet(N)$, which proves the proposition.
\end{proof}

The proposition \sref{0.15.1.19} will apply in particular for filtrations satisfying one or the other of the following hypotheses:
\begin{enumerate}
  \item[$1^\circ$] The filtration $(M_k)$ is finite and separated (since this implies $M_k=0$ for $k$ large enough);
  \item[$2^\circ$] $A$ is a Noetherian ring, $\mathfrak{J}$ an ideal contained in the radical of $A$, $M$ a finitely generated $A$-module, and $(M_k)$ the $\mathfrak{J}$-preadic filtration ($\mathbf{0_I}$, 7.3.5).
\end{enumerate}

\begin{corollary}[15.1.20]
\label{0.15.1.20}
\oldpage[0\textsubscript{IV-1}]{20}
Let $A$ be a ring, $M$ an $A$-module, $(f_i)_{1\leq i\leq n}$ an $M$-regular sequence of elements of $A$.
Then, for any integers $\nu_i\geq 1$, the sequence $(f_i^{\nu_i})$ is $M$-regular, and $M/(\sum_{i\leq n}f_i^{\nu_i}M)$ has a composition series whose quotients are isomorphic to $M/(\sum_{i\leq n}f_iM)$.
\end{corollary}

\begin{proof}
Let us reason by induction on $n$;
for $n=1$ the assertions are immediate since $f^jM/f^{j+1}M$ is isomorphic to $M/fM$, the homothety of ratio $f$ in $M$ being injective by hypothesis.
Set then $N=M/(\sum_{i\leq n-1}f_i^{\nu_i}M)$;
by the induction hypothesis, $N$ admits a finite filtration $(N_k)$ for which $\gr_k(N)$ are isomorphic to $M/(\sum_{i\leq n-1}f_iM)$;
by hypothesis $f_n$ is thus $\gr_\bullet(N)$-regular;
it follows by \sref{0.15.1.19} that $f_n$ is $N$-regular.
Applying the case $n=1$ proved at the start, one concludes that $f_n^{\nu_n}$ is $N$-regular and that $N/f_n^{\nu_n}N$ admits a composition series whose quotients are isomorphic to $N/f_nN$;
but for the filtration of $N/f_nN$ given by the quotient filtration of $(N_k)$, $\gr_k(N/f_nN)$ is isomorphic to $\gr_k(N)/f_n\gr_k(N)$ by \sref{0.15.1.19} for all $k$, hence to $M/(\sum_{i\leq n}f_iM)$, which proves the corollary.
\end{proof}

\begin{proposition}[15.1.21]
\label{0.15.1.21}
Let $A$, $B$ be two Noetherian local rings, $\varphi:A\to B$ a local homomorphism, $M$ a finitely generated $B$-module, $(f_i)_{1\leq i\leq n}$ an $A$-regular sequence of elements of the maximal ideal of $A$.
The following properties are equivalent:
\begin{enumerate}
  \item[a)] $M$ is a flat $A$-module.
  \item[b)] If $\mathfrak{J}$ is the ideal $\sum_{i\leq n}f_iA$, $M/\mathfrak{J}M$ is a flat $A/\mathfrak{J}$-module and the sequence $(f_i)$ is $M$-regular.
\end{enumerate}
\end{proposition}

\begin{proof}
The fact that a) implies b) follows from \sref{0.15.1.13} and from ($\mathbf{0_{III}}$, 10.2.2).
Conversely, if the sequence $(f_i)$ is $M$-regular, it is $M$-quasi-regular \sref{0.15.1.9}, hence the homomorphism \sref{0.15.1.1.1} is bijective.
But as the sequence $(f_i)_{1\leq i\leq n}$ is also supposed $A$-regular, it is $A$-quasi-regular \sref{0.15.1.9}, and the corresponding canonical homomorphism $(A/\mathfrak{J})[T_1,\ldots,T_n]\to\gr_\bullet(A)$ is bijective.
One concludes that the canonical morphism ($\mathbf{0_{III}}$, 10.1.1.2)
\[
  \gr_0(M)\otimes_{A/\mathfrak{J}}\gr_\bullet(A) \to \gr_\bullet(M)
\]
is bijective.
The conclusion then follows from ($\mathbf{0_{III}}$, 10.2.2), the $f_i$ being in the maximal ideal of $A$ and the homomorphism $\varphi$ being local.
\end{proof}

\subsection{\texorpdfstring{$\sh{F}$}{F}-regular sequences}
\label{subsection:0.15.2}

\begin{env}[15.2.1]
\label{0.15.2.1}
Let $X$ be a ringed space, $\sh{F}$ an $\sh{O}_X$-Module, $f_i$ ($1\leq i\leq n$) a sequence of sections of $\sh{O}_X$ over $X$, $\sh{J}$ the Ideal of $\sh{O}_X$ generated by the $f_i$ (image of $\sh{O}_X^n$ by the homomorphism $\sh{O}_X^n\to\sh{O}_X$ defined canonically by the $f_i$ ($\mathbf{0_I}$, 5.1.1)).
The sub-$\sh{O}_X$-Modules $\sh{J}^k\sh{F}$ define a \emph{filtration} on $\sh{F}$, and one sets further $\gr_k(\sh{F})=\sh{J}^k\sh{F}/\sh{J}^{k+1}\sh{F}$;
if $S_k$ is the set of homogeneous polynomials of total degree $k$ in the ring $\bb{Z}[T_1,\ldots,T_n]$, $S_k$ can be considered as a simple sheaf on $X$, and one defines as in \sref{0.15.1.1} a canonical homomorphism
\[
\label{0.15.2.1.1}
  \varphi_k : \gr_0(\sh{F})\otimes_{\bb{Z}} S_k \to \gr_k(\sh{F})
\tag{15.2.1.1}
\]
\oldpage[0\textsubscript{IV-1}]{21}
in the following way:
for every open set $U$ of $X$ (or of a base of the topology of $X$), one considers the homomorphism \sref{0.15.1.1.1}
\[
  \varphi_{k,U} : (\gr_0(\Gamma(U,\sh{F})))\otimes_{\bb{Z}}S_k \to \gr_k(\Gamma(U,\sh{F}))
\]
defined by the $n$ sections $f_i|U\in\Gamma(U,\sh{O}_X)$, the $\Gamma(U,\sh{O}_X)$-module $\Gamma(U,\sh{F})$ being filtered by the submodules $(\Gamma(U,\sh{J}))^k\Gamma(U,\sh{F})$.
The $\sh{O}_X$-Module $\gr_k(\sh{F})$ being associated to the presheaf $U\mapsto\gr_k(\Gamma(U,\sh{F}))$ for all $k$, the $\varphi_{k,U}$ define the homomorphism of $\sh{O}_X$-Modules \sref{0.15.2.1.1}.
By virtue of the exactness properties of inductive limits of modules and their commutation with tensor products, it follows from the above that for all $x\in X$, the homomorphism $(\varphi_k)_x$ of the stalks of the two members of \sref{0.15.2.1.1} at the point $x$ identifies with the homomorphism of \sref{0.15.1.1}
\[
\label{0.15.2.1.2}
  \gr_0(\sh{F}_x)\otimes_{\bb{Z}}S_k \to \gr_k(\sh{F}_x)
\tag{15.2.1.2}
\]
defined by the sequence $((f_i)_x)_{1\leq i\leq n}$ of elements of $\sh{O}_x$.
\end{env}

\begin{env}[15.2.2]
\label{0.15.2.2}
The notations being those of \sref{0.15.2.1}, we say that the sequence $(f_i)_{1\leq i\leq n}$ is \emph{$\sh{F}$-regular} if, denoting by $f_j\sh{F}$ the sub-$\sh{O}_X$-Module of $\sh{F}$ that is the image of the endomorphism of $\sh{F}$ defined by $f_j$, the endomorphism of $\sh{F}/(\sum_{j=1}^{i-1}f_j\sh{F})$ defined by $f_i$ is \emph{injective} for $1\leq i\leq n$;
we say that the sequence $(f_i)$ is \emph{$\sh{F}$-quasi-regular} if the homomorphisms \sref{0.15.2.1.1} are injective.
It follows from these definitions and the remarks of \sref{0.15.2.1} that for the sequence $(f_i)$ to be $\sh{F}$-regular (resp. $\sh{F}$-quasi-regular), it is necessary and sufficient that, for all $x\in X$, the sequence $((f_i)_x)_{1\leq i\leq n}$ be $\sh{F}_x$-regular (resp. $\sh{F}_x$-quasi-regular).
Every $\sh{F}$-regular sequence is thus $\sh{F}$-quasi-regular \sref{0.15.1.9};
the converse is true if $\sh{F}$ is an $\sh{O}_X$-Module \emph{of finite type}, if the $\sh{O}_x$ are \emph{Noetherian} rings for all $x\in X$, and if moreover the $(f_i)_x$ belong to the \emph{radical} of $\sh{O}_x$ for all $x\in X$ \sref{0.15.1.11}.

If $X=\Spec(A)$ is an affine scheme, and $\sh{F}=\widetilde{M}$, it follows from ($\mathbf{I}$, 1.3) that for the sequence $(f_i)$ to be $\sh{F}$-regular (resp. $\sh{F}$-quasi-regular), it is necessary and sufficient that it be $M$-regular (resp. $M$-quasi-regular).

When $\sh{F}=\sh{O}_X$, one omits the mention of $\sh{F}$ in the preceding definitions.
\end{env}

\begin{remark}[15.2.3]
\label{0.15.2.3}
When $X=\Spec(A)$ and $\mathfrak{J}=\sum_{i=1}^n f_iA$, the sequence $(f_j)$ is quasi-regular for \emph{every} $A$-module $M$ in every open set $U$ not meeting $V(\mathfrak{J})$, since for every point $x\in U$, $\mathfrak{J}_x=\sh{O}_x$ and the two members of \sref{0.15.2.1.2} are zero.
The notion of $\widetilde{M}$-regular sequence is thus of interest only in a neighbourhood of the points of $V(\mathfrak{J})$.
\end{remark}

\begin{proposition}[15.2.4]
\label{0.15.2.4}
Let $X$ be a ringed space, $\sh{F}$ a coherent $\sh{O}_X$-Module, $x$ a point of $X$, $(f_i)_{1\leq i\leq n}$ a sequence of elements of $\Gamma(X,\sh{O}_X)$.
If the sequence $((f_i)_x)_{1\leq i\leq n}$ is $\sh{F}_x$-regular, there exists an open neighbourhood $U$ of $x$ such that the sequence $(f_i|U)_{1\leq i\leq n}$ is $(\sh{F}|U)$-regular.
\end{proposition}

\begin{proof}
Indeed, for $1\leq i\leq n$, the $\sh{O}_X$-Module $\sh{F}/(\sum_{j=1}^{i-1}f_j\sh{F})$ is coherent ($\mathbf{0_I}$, 5.3.4), hence so is the kernel of the endomorphism of this $\sh{O}_X$-Module defined by $f_i$ ($\mathbf{0_I}$, 5.3.4);
as the stalk of this kernel at the point $x$ is zero, the same holds for its restriction to a neighbourhood of $x$ ($\mathbf{0_I}$, 5.2.2).
\end{proof}

\begin{proposition}[15.2.5]
\label{0.15.2.5}
\oldpage[0\textsubscript{IV-1}]{22}
Let $u=(\psi,\theta):X\to Y$ be a flat morphism of ringed spaces ($\mathbf{0_I}$, 6.7.1);
let $\sh{F}$ be an $\sh{O}_Y$-Module, $\sh{G}=u^*(\sh{F})$, $(f_i)_{1\leq i\leq n}$ a sequence of elements of $\Gamma(Y,\sh{O}_Y)$, $(g_i)_{1\leq i\leq n}$ the sequence of their images $g_i=\theta(f_i)\in\Gamma(Y,u_*(\sh{O}_X))=\Gamma(X,\sh{O}_X)$.
If the sequence $(f_i)$ is $\sh{F}$-regular, the sequence $(g_i)$ is $\sh{G}$-regular.
The converse is true if $u$ is a faithfully flat morphism ($\mathbf{0_I}$, 6.7.8).
\end{proposition}

\begin{proof}
Let $x$ be a point of $X$, $y=u(x)=\psi(x)$;
we have
\[
  \sh{G}_x=\psi_x^*(\sh{F}_y)\otimes_{\psi_x^*(\sh{O}_y)}\sh{O}_x,
  \qquad\text{and}\qquad
  (g_i)_x=\theta_x^\#(\psi_x^*((f_i)_y));
\]
as $\theta_x^\#:\psi_x^*(\sh{O}_y)\to\sh{O}_x$ makes $\sh{O}_x$ a flat $\psi_x^*(\sh{O}_y)$-module, and since the functor $\psi^*$ is exact, the first assertion follows from the corresponding assertion of \sref{0.15.1.14};
the same holds for the second, taking into account the fact that $u$ being faithfully flat means that $\psi$ is surjective and that $\theta_x^\#$ makes $\sh{O}_x$ a faithfully flat $\psi_x^*(\sh{O}_y)$-module.
\end{proof}
% END NATIVE INPUT ega0/ega0-15
% BEGIN NATIVE INPUT ega0/ega0-16
\section{Dimension and depth of Noetherian local rings}
\label{section:0.16}

\subsection{Dimension of a ring}
\label{subsection:0.16.1}
\label{0.16.1}

\begin{env}[16.1.1]
\label{0.16.1.1}
We call the \emph{dimension} (or \emph{Krull dimension}) of a ring $\mathrm{A}$ and denote by $\dim\mathrm{A}$ the (combinatorial) dimension of its spectrum $\Spec(\mathrm{A})$ \sref{0.14.2.1}; as the irreducible closed subsets of $\Spec(\mathrm{A})$ are the $\mathrm{V}(\mathfrak{p})$, where $\mathfrak{p}$ is a prime ideal of $\mathrm{A}$ (Bourbaki, \emph{Alg.~comm.}, chap.~II, \textsection~4, n\textsuperscript{o}~3, cor.~1 of prop.~14), it amounts to saying that $\dim(\mathrm{A})$ is \emph{the upper bound of the lengths of chains} \sref{0.14.1.1}
\[
  \mathfrak{p}_0\subset\mathfrak{p}_1\subset\ldots\subset\mathfrak{p}_n
\]
of \emph{prime ideals} of $\mathrm{A}$.
If $\mathrm{A}\neq 0$, we have $\dim\mathrm{A}\geq 0$.
A nonzero Artinian ring (in particular a field) is of dimension $0$ (Bourbaki, \emph{Alg.~comm.}, chap.~IV, \textsection~2, n\textsuperscript{o}~5, prop.~7); a Dedekind ring that is not a field (in particular $\mathbf{Z}$ or a polynomial ring $k[\mathrm{T}]$ over a field $k$) is of dimension 1 (Bourbaki, \emph{Alg.~comm.}, chap.~VII, \textsection~2).
A Noetherian ring is not necessarily of finite dimension [30].
If $(\mathfrak{p}_\alpha)$ is the family of minimal prime ideals of $\mathrm{A}$, we have \sref{0.14.1.2.1}
\begin{equation}
\label{0.16.1.1.1}
\tag{16.1.1.1}
  \dim(\mathrm{A})=\sup_\alpha(\dim(\mathrm{A}/\mathfrak{p}_\alpha)).
\end{equation}
\end{env}

\begin{env}[16.1.2]
\label{0.16.1.2}
For every ideal $\mathfrak{a}$ of $\mathrm{A}$, the prime ideals of $\mathrm{A}/\mathfrak{a}$ correspond bijectively to the prime ideals of $\mathrm{A}$ containing $\mathfrak{a}$, hence
\begin{equation}
\label{0.16.1.2.1}
\tag{16.1.2.1}
  \dim(\mathrm{A}/\mathfrak{a})\leq\dim(\mathrm{A}).
\end{equation}
If in addition $\mathfrak{a}$ \emph{is not contained in any minimal prime ideal} $\mathfrak{p}_\alpha$ of $\mathrm{A}$, every chain of prime ideals of $\mathrm{A}$ containing $\mathfrak{a}$ can be completed by one of the $\mathfrak{p}_\alpha$, hence
\begin{equation}
\label{0.16.1.2.2}
\tag{16.1.2.2}
  1+\dim(\mathrm{A}/\mathfrak{a})\leq\dim(\mathrm{A}).
\end{equation}
\end{env}

\begin{env}[16.1.3]
\label{0.16.1.3}
\oldpage[0\textsubscript{IV-1}]{23}
If $\mathrm{S}$ is a multiplicative subset of $\mathrm{A}$, the prime ideals of $\mathrm{S}^{-1}\mathrm{A}$ correspond bijectively to the prime ideals of $\mathrm{A}$ not meeting $\mathrm{S}$, hence
\begin{equation}
\label{0.16.1.3.1}
\tag{16.1.3.1}
  \dim(\mathrm{S}^{-1}\mathrm{A})\leq\dim(\mathrm{A}).
\end{equation}
For every prime ideal $\mathfrak{p}$ of $\mathrm{A}$, we have, by definition \sref{0.14.2.1}
\begin{equation}
\label{0.16.1.3.2}
\tag{16.1.3.2}
  \dim(\mathrm{A}_\mathfrak{p})=\codim(\mathrm{V}(\mathfrak{p}),\Spec(\mathrm{A})).
\end{equation}
This number is also called the \emph{height} of the prime ideal $\mathfrak{p}$ and denoted $\operatorname{ht}(\mathfrak{p})$.
More generally, for every ideal $\mathfrak{a}$ of $\mathrm{A}$, we call the \emph{height} of $\mathfrak{a}$ the number
\begin{equation}
\label{0.16.1.3.3}
\tag{16.1.3.3}
  \operatorname{ht}(\mathfrak{a})=\codim(\mathrm{V}(\mathfrak{a}),\Spec(\mathrm{A})).
\end{equation}
If $(\mathfrak{p}_\lambda)$ is the family of minimal prime ideals among those containing $\mathfrak{a}$, we have then \sref{0.14.2.1}
\begin{equation}
\label{0.16.1.3.4}
\tag{16.1.3.4}
  \operatorname{ht}(\mathfrak{a})=\inf_\lambda(\operatorname{ht}(\mathfrak{p}_\lambda)).
\end{equation}
We evidently have
\begin{equation}
\label{0.16.1.3.5}
\tag{16.1.3.5}
  \dim(\mathrm{A})=\sup_\mathfrak{m}(\dim(\mathrm{A}_\mathfrak{m}))
\end{equation}
where $\mathfrak{m}$ ranges over the set of maximal ideals of $\mathrm{A}$.
\end{env}

\begin{env}[16.1.4]
\label{0.16.1.4}
For every prime ideal $\mathfrak{p}$ of $\mathrm{A}$, we have \sref{0.14.2.2.2}
\begin{equation}
\label{0.16.1.4.1}
\tag{16.1.4.1}
  \dim(\mathrm{A}_\mathfrak{p})+\dim(\mathrm{A}/\mathfrak{p})\leq\dim\mathrm{A}.
\end{equation}
We say that $\mathrm{A}$ is \emph{catenary} (resp.~\emph{equidimensional}, \emph{equicodimensional}, \emph{biequidimensional}) when $\Spec(\mathrm{A})$ is catenary (resp.~equidimensional, equicodimensional, biequidimensional).
When $\mathrm{A}$ is Noetherian and biequidimensional, the two members of \sref{0.16.1.4.1} are equal for every prime ideal $\mathfrak{p}$ of $\mathrm{A}$ \sref{0.14.3.5}.

For every ideal $\mathfrak{a}$ of $\mathrm{A}$, the prime ideals of $\mathrm{A}/\mathfrak{a}$ correspond bijectively to the prime ideals of $\mathrm{A}$ containing $\mathfrak{a}$, hence, \emph{if $\mathrm{A}$ is catenary, so is $\mathrm{A}/\mathfrak{a}$}.
Similarly, for every multiplicative subset $\mathrm{S}$ of $\mathrm{A}$, the prime ideals of $\mathrm{S}^{-1}\mathrm{A}$ correspond bijectively to the prime ideals of $\mathrm{A}$ not meeting $\mathrm{S}$, hence \emph{if $\mathrm{A}$ is catenary, so is $\mathrm{S}^{-1}\mathrm{A}$}.

Furthermore, as the prime ideals of $\mathrm{A}$ contained in a prime ideal $\mathfrak{p}$ correspond bijectively to the prime ideals of $\mathrm{A}_\mathfrak{p}$, for $\mathrm{A}$ to be catenary it is necessary and sufficient that $\mathrm{A}_\mathfrak{p}$ be so for every $\mathfrak{p}$.
To say that $\mathrm{A}$ is catenary means (\sref{0.14.3.2} and \sref{0.16.1.3.2}) that for every pair of prime ideals $\mathfrak{p}$, $\mathfrak{q}$ of $\mathrm{A}$ such that $\mathfrak{q}\subset\mathfrak{p}$, we have
\begin{equation}
\label{0.16.1.4.2}
\tag{16.1.4.2}
  \dim(\mathrm{A}_\mathfrak{q})+\dim(\mathrm{A}_\mathfrak{p}/\mathfrak{q}\mathrm{A}_\mathfrak{p})=\dim(\mathrm{A}_\mathfrak{p}).
\end{equation}
If $\mathrm{A}$ is biequidimensional, it is the same for $\mathrm{A}_\mathfrak{p}$ for every prime ideal $\mathfrak{p}$ of $\mathrm{A}$.
\end{env}

\begin{proposition}[16.1.5]
\label{0.16.1.5}
Let $\rho:\mathrm{A}\to\mathrm{B}$ be a ring homomorphism making $\mathrm{B}$ an integral $\mathrm{A}$-algebra.
Then $\dim(\mathrm{B})\leq\dim(\mathrm{A})$; if in addition $\rho$ is injective, we have $\dim(\mathrm{B})=\dim(\mathrm{A})$.
\end{proposition}

\begin{proof}
\oldpage[0\textsubscript{IV-1}]{24}
If $\mathfrak{n}$ is the kernel of $\rho$, $\rho(\mathrm{A})$ is isomorphic to $\mathrm{A}/\mathfrak{n}$, hence $\dim\rho(\mathrm{A})\leq\dim(\mathrm{A})$ by \sref{0.16.1.2.1}, and since $\mathrm{B}$ is integral over $\rho(\mathrm{A})$, one can reduce to proving the second assertion.
If $\mathfrak{p}_0'\subset\mathfrak{p}_1'\subset\ldots\subset\mathfrak{p}_n'$ is a chain of prime ideals of $\mathrm{B}$, the $\mathfrak{p}_i'\cap\mathrm{A}$ are then distinct prime ideals in $\mathrm{A}$ (Bourbaki, \emph{Alg.~comm.}, chap.~V, \textsection~2, n\textsuperscript{o}~1, cor.~1 of prop.~1), hence they form a chain of prime ideals.
Conversely, for every chain $\mathfrak{p}_0\subset\mathfrak{p}_1\subset\ldots\subset\mathfrak{p}_m$ of prime ideals of $\mathrm{A}$, we know that for each $i$ there exists a prime ideal $\mathfrak{p}_i'$ of $\mathrm{B}$ such that $\mathfrak{p}_i'\cap\mathrm{A}=\mathfrak{p}_i$ and $\mathfrak{p}_i'\subset\mathfrak{p}_{i+1}'$ for $0\leq i\leq m-1$ (\emph{loc.~cit.}, cor.~2 of th.~1).
Hence the conclusion.
\end{proof}

\begin{proposition}[16.1.6]
\label{0.16.1.6}
Let $\mathrm{B}$ be an integral ring, $\mathrm{A}$ a local subring of $\mathrm{B}$, $\mathfrak{m}$ its maximal ideal.
Suppose that $\mathrm{A}$ is integrally closed and that $\mathrm{B}$ is integral over $\mathrm{A}$; then, for every maximal ideal $\mathfrak{n}$ of $\mathrm{B}$, we have $\dim(\mathrm{B}_\mathfrak{n})=\dim(\mathrm{A})$.
\end{proposition}

\begin{proof}
The first part of the reasoning of \sref{0.16.1.5} shows that $\dim(\mathrm{B}_\mathfrak{n})\leq\dim(\mathrm{A})$.
Conversely, consider a chain of prime ideals $\mathfrak{p}_0\subset\mathfrak{p}_1\subset\ldots\subset\mathfrak{p}_k=\mathfrak{m}$ of $\mathrm{A}$; we know that $\mathfrak{n}\cap\mathrm{A}=\mathfrak{m}$ (Bourbaki, \emph{Alg.~comm.}, chap.~V, \textsection~2, n\textsuperscript{o}~1, prop.~1) and by virtue of the hypotheses made, one can apply the second theorem of Cohen-Seidenberg (Bourbaki, \emph{Alg.~comm.}, chap.~V, \textsection~2, n\textsuperscript{o}~4, th.~3), which proves the existence of a chain of prime ideals $\mathfrak{p}_0'\subset\mathfrak{p}_1'\subset\ldots\subset\mathfrak{p}_k'=\mathfrak{n}$ of $\mathrm{B}$ such that $\mathfrak{p}_i'\cap\mathrm{A}=\mathfrak{p}_i$ for each $i$; from this one deduces immediately $\dim(\mathrm{B}_\mathfrak{n})\geq\dim(\mathrm{A})$, whence the proposition.
\end{proof}

\begin{env}[16.1.7]
\label{0.16.1.7}
Let $\mathrm{M}$ be an $\mathrm{A}$-module; we call the \emph{dimension} of $\mathrm{M}$ and denote by $\dim_\mathrm{A}(\mathrm{M})$ or $\dim(\mathrm{M})$ the dimension of the ring $\mathrm{A}/\operatorname{Ann}(\mathrm{M})$ (isomorphic to the ring of homotheties $\mathrm{A}_\mathrm{M}$ of $\mathrm{M}$).
We thus have $\dim(\mathrm{M}^k)=\dim(\mathrm{M})$ for every $k>0$.
If $\mathrm{M}$ is \emph{of finite type}, the prime ideals of $\mathrm{A}$ containing $\operatorname{Ann}(\mathrm{M})$ are exactly those belonging to $\Supp(\mathrm{M})$ $(\mathbf{0}_\mathrm{I},~1.7.4)$ and the latter is closed in $\Spec(\mathrm{A})$, hence we then have
\begin{equation}
\label{0.16.1.7.1}
\tag{16.1.7.1}
  \dim(\mathrm{M})=\dim(\Supp(\mathrm{M}))\leq\dim(\mathrm{A}).
\end{equation}
If $\mathrm{M}$, $\mathrm{N}$ are two $\mathrm{A}$-modules such that $\mathrm{N}\subset\mathrm{M}$, we have
\begin{equation}
\label{0.16.1.7.2}
\tag{16.1.7.2}
  \dim(\mathrm{N})\leq\dim(\mathrm{M}),\qquad\dim(\mathrm{M}/\mathrm{N})\leq\dim(\mathrm{M}).
\end{equation}
If $\mathrm{S}$ is a multiplicative subset of $\mathrm{A}$ and if $\mathrm{M}$ is of finite type, we have
\begin{equation}
\label{0.16.1.7.3}
\tag{16.1.7.3}
  \dim_{\mathrm{S}^{-1}\mathrm{A}}(\mathrm{S}^{-1}\mathrm{M})\leq\dim_\mathrm{A}(\mathrm{M}).
\end{equation}
Furthermore, every maximal chain of prime ideals containing $\operatorname{Ann}(\mathrm{M})$ terminates at a maximal ideal $\mathfrak{m}$, hence its length is equal to the length of the chain of corresponding prime ideals of $\mathrm{A}_\mathfrak{m}$ (which contain $(\operatorname{Ann}(\mathrm{M}))_\mathfrak{m}$); from this remark and from \sref{0.16.1.7.3}, one obtains
\begin{equation}
\label{0.16.1.7.4}
\tag{16.1.7.4}
  \dim_\mathrm{A}(\mathrm{M})=\sup_\mathfrak{m}(\dim_{\mathrm{A}_\mathfrak{m}}(\mathrm{M}_\mathfrak{m}))
\end{equation}
where $\mathfrak{m}$ ranges over the set of maximal ideals of $\mathrm{A}$.
We say that $\mathrm{M}$ is \emph{catenary} (resp.~\emph{equidimensional}, \emph{equicodimensional}, \emph{biequidimensional}) when $\mathrm{A}/\operatorname{Ann}(\mathrm{M})$ is.
\end{env}

\begin{proposition}[16.1.8]
\label{0.16.1.8}
Let $\mathrm{A}$ be a ring, $\mathrm{M}$ an $\mathrm{A}$-module of finite type, $\mathfrak{p}$ a prime ideal of $\mathrm{A}$.
We have
\begin{equation}
\label{0.16.1.8.1}
\tag{16.1.8.1}
  \dim_{\mathrm{A}_\mathfrak{p}}(\mathrm{M}_\mathfrak{p})+\dim_\mathrm{A}(\mathrm{M}/\mathfrak{p}\mathrm{M})\leq\dim_\mathrm{A}(\mathrm{M}).
\end{equation}
\end{proposition}

\begin{proof}
\oldpage[0\textsubscript{IV-1}]{25}
We have $\operatorname{Ann}(\mathrm{M}_\mathfrak{p})=(\operatorname{Ann}(\mathrm{M}))_\mathfrak{p}$, hence the prime ideals of $\mathrm{A}_\mathfrak{p}$ containing $\operatorname{Ann}(\mathrm{M}_\mathfrak{p})$ correspond bijectively to the prime ideals of $\mathrm{A}$ containing $\operatorname{Ann}(\mathrm{M})$ and contained in $\mathfrak{p}$.
On the other hand, if $\operatorname{Ann}(\mathrm{M})=\mathfrak{n}$, $\mathrm{V}(\operatorname{Ann}(\mathrm{M}/\mathfrak{p}\mathrm{M}))=\mathrm{V}(\mathfrak{p}+\mathfrak{n})$ $(\mathbf{0}_\mathrm{I},~1.7.5)$, hence the prime ideals of $\mathrm{A}$ containing $\operatorname{Ann}(\mathrm{M}/\mathfrak{p}\mathrm{M})$ contain $\mathfrak{p}+\mathfrak{n}$; whence the conclusion.
We note in addition that
\begin{equation}
\label{0.16.1.8.2}
\tag{16.1.8.2}
  \dim_{\mathrm{A}/\mathfrak{p}}(\mathrm{M}/\mathfrak{p}\mathrm{M})=\dim_\mathrm{A}(\mathrm{M}/\mathfrak{p}\mathrm{M}).
\end{equation}
\end{proof}

\begin{proposition}[16.1.9]
\label{0.16.1.9}
Let $\mathrm{A}$ be a Noetherian ring, $\rho:\mathrm{A}\to\mathrm{B}$ a ring homomorphism, $\mathrm{M}$ a $\mathrm{B}$-module such that $\mathrm{M}_{[\rho]}$ is an $\mathrm{A}$-module of finite type.
Then
\begin{equation}
\label{0.16.1.9.1}
\tag{16.1.9.1}
  \dim_\mathrm{A}(\mathrm{M}_{[\rho]})=\dim_\mathrm{B}(\mathrm{M}).
\end{equation}
\end{proposition}

\begin{proof}
Indeed, the ring $\mathrm{B}_\mathrm{M}$ of homotheties of the $\mathrm{B}$-module $\mathrm{M}$ is a subring of the $\mathrm{A}$-module $\mathrm{C}=\operatorname{End}_\mathrm{A}(\mathrm{M}_{[\rho]})$, and $\mathrm{C}$ is an $\mathrm{A}$-module of finite type (Bourbaki, \emph{Alg.~comm.}, chap.~III, \textsection~3, n\textsuperscript{o}~1, lemma~2), hence it is the same for $\mathrm{B}_\mathrm{M}$; moreover the ring $\mathrm{A}_\mathrm{M}$ of homotheties of the $\mathrm{A}$-module $\mathrm{M}_{[\rho]}$ is a subring of $\mathrm{B}_\mathrm{M}$, canonical image of $\mathrm{A}$ in $\mathrm{C}$, hence $\mathrm{B}_\mathrm{M}$ is finite over $\mathrm{A}_\mathrm{M}$; the conclusion then follows from \sref{0.16.1.5}.
\end{proof}

\begin{env}[16.1.10]
\label{0.16.1.10}
Let $\mathrm{A}$ be a Noetherian ring; if $\mathrm{M}$ is an $\mathrm{A}$-module of \emph{finite length}, we know (Bourbaki, \emph{Alg.~comm.}, chap.~IV, \textsection~2, n\textsuperscript{o}~5, prop.~7 and cor.~1 of prop.~7) that $\Supp(\mathrm{M})$ is a finite discrete set, hence $\dim(\mathrm{M})=0$ if $\mathrm{M}\neq 0$; conversely, if $\mathrm{M}$ is an $\mathrm{A}$-module of finite type such that $\dim(\mathrm{M})=0$, every point of $\Supp(\mathrm{M})$ is closed, in other words it is a maximal ideal of $\mathrm{A}$, hence (\emph{loc.~cit.}, prop.~7) $\mathrm{M}$ is of finite length.
\end{env}

\subsection[Dimension of a semi-local Noetherian ring]{Dimension of a semi-local Noetherian ring\protect\footnote{This section reproduces the lecture given by Serre in a course at the Collège de France in 1958.}}
\label{subsection:0.16.2}

\begin{env}[16.2.1]
\label{0.16.2.1}
Let $\mathrm{A}$ be a \emph{semi-local Noetherian} ring, $\mathfrak{r}$ its radical; recall that an \emph{ideal of definition} $\mathfrak{q}$ of $\mathrm{A}$ is an ideal such that $\mathfrak{q}\subset\mathfrak{r}$ and $\mathfrak{q}$ contains a power of $\mathfrak{r}$, or equivalently, that the only prime ideals containing $\mathfrak{q}$ are \emph{maximal}, or that $\mathrm{A}/\mathfrak{q}$ is an Artinian ring.
Let $\mathrm{M}$ be an $\mathrm{A}$-module of finite type; for every integer $n>0$, $\mathrm{M}/\mathfrak{q}^n\mathrm{M}$ is an $\mathrm{A}$-module of \emph{finite length}, equal to $\sum_{i=0}^{n-1}\operatorname{long}(\gr_i(\mathrm{M}))$, where $\gr_\bullet(\mathrm{M})$ is the $(\mathrm{A}/\mathfrak{q})$-graded module associated to $\mathrm{M}$ equipped with the $\mathfrak{q}$-preadic filtration.
Recall that there exists a \emph{polynomial} $\mathrm{Q}(n)$ with rational coefficients such that $\operatorname{long}(\gr_n(\mathrm{M}))=\mathrm{Q}(n)$ for $n$ sufficiently large $(\mathbf{III},~2.5.4)$; from this one deduces immediately that there exists a polynomial $\mathrm{P}_\mathfrak{q}(\mathrm{M},n)$ in $n$, with rational coefficients, such that $\operatorname{long}(\mathrm{M}/\mathfrak{q}^n\mathrm{M})=\mathrm{P}_\mathfrak{q}(\mathrm{M},n)$ \emph{for $n$ sufficiently large}; it is clear that this polynomial is unique and that its leading coefficient is $>0$ if $\mathrm{M}\neq 0$; we call it the \emph{Hilbert--Samuel polynomial} of $\mathrm{M}$ for the $\mathfrak{q}$-preadic filtration.
\end{env}

\oldpage[0\textsubscript{IV-1}]{26}

\begin{lemma}[16.2.2.1]
\label{0.16.2.2.1}
Let $(\mathrm{M}_n)$ be a $\mathfrak{q}$-good filtration of $\mathrm{M}$ (Bourbaki, \emph{Alg.~comm.}, chap.~III, \textsection~3, n\textsuperscript{o}~1); then $\operatorname{long}(\mathrm{M}/\mathrm{M}_n)$ is equal, for $n$ sufficiently large, to a polynomial in $n$ whose degree and leading coefficient are the same as those of $\mathrm{P}_\mathfrak{q}(\mathrm{M},n)$.
\end{lemma}

\begin{proof}
Indeed, there exists $n_0$ such that $\mathrm{M}_{n+1}=\mathfrak{q}\mathrm{M}_n$ for $n\geq n_0$, whence the inclusions $\mathfrak{q}^{n+n_0}\mathrm{M}\subset\mathrm{M}_{n+n_0}=\mathfrak{q}^n\mathrm{M}_{n_0}\subset\mathfrak{q}^n\mathrm{M}\subset\mathrm{M}_n$ for $n$ sufficiently large, which implies
\[
  \mathrm{P}_\mathfrak{q}(\mathrm{M},n+n_0)\geq\operatorname{long}(\mathrm{M}/\mathrm{M}_{n+n_0})\geq\mathrm{P}_\mathfrak{q}(\mathrm{M},n)\geq\operatorname{long}(\mathrm{M}/\mathrm{M}_n)
\]
hence the conclusion.
\end{proof}

If $\mathfrak{q}'$ is a second ideal of definition, one has $\mathfrak{q}'\supset\mathfrak{q}^m$ for some integer $m$, hence $\mathrm{P}_{\mathfrak{q}'}(\mathrm{M},n)\leq\mathrm{P}_\mathfrak{q}(\mathrm{M},mn)$; the \emph{degree} of $\mathrm{P}_\mathfrak{q}(\mathrm{M},n)$ therefore \emph{does not depend on the ideal of definition} $\mathfrak{q}$; we denote it $d(\mathrm{M})$.

\begin{lemma}[16.2.2.2]
\label{0.16.2.2.2}
Let $0\to\mathrm{M}'\to\mathrm{M}\to\mathrm{M}''\to 0$ be a short exact sequence of $\mathrm{A}$-modules of finite type.
Then the polynomial $\mathrm{P}_\mathfrak{q}(\mathrm{M},n)-\mathrm{P}_\mathfrak{q}(\mathrm{M}',n)-\mathrm{P}_\mathfrak{q}(\mathrm{M}'',n)$ has degree $\leq d(\mathrm{M}')-1\leq d(\mathrm{M})-1$.
\end{lemma}

\begin{proof}
Indeed, the filtration of the $\mathrm{M}_n'=\mathrm{M}'\cap\mathfrak{q}^n\mathrm{M}$ is $\mathfrak{q}$-good (Bourbaki, \emph{Alg.~comm.}, chap.~III, \textsection~3, n\textsuperscript{o}~1, prop.~1); as $\operatorname{long}(\mathrm{M}/\mathfrak{q}^n\mathrm{M})=\operatorname{long}(\mathrm{M}''/\mathfrak{q}^n\mathrm{M}'')+\operatorname{long}(\mathrm{M}'/\mathrm{M}_n')$, the lemma follows immediately from \sref{0.16.2.2.1} applied to $\mathrm{M}'$.
\end{proof}

\begin{theorem}[16.2.3]
\label{0.16.2.3}
\textup{(Krull-Chevalley-Samuel).}
Let $\mathrm{A}$ be a semi-local Noetherian ring, $\mathfrak{r}$ its radical, $\mathrm{M}\neq 0$ an $\mathrm{A}$-module of finite type, $d(\mathrm{M})$ the degree of the Hilbert--Samuel polynomial of $\mathrm{M}$ (for any ideal of definition of $\mathrm{A}$); let $s(\mathrm{M})$ denote the lower bound of the number of elements $x_1,\ldots,x_n$ of $\mathfrak{r}$ such that $\mathrm{M}/(x_1\mathrm{M}+\ldots+x_n\mathrm{M})$ is of finite length.
Then
\[
  d(\mathrm{M})=s(\mathrm{M})=\dim(\mathrm{M}).
\]
\end{theorem}

\begin{lemma}[16.2.3.1]
\label{0.16.2.3.1}
For every $x\in\mathfrak{r}$, let ${}_x\mathrm{M}$ be the submodule of $\mathrm{M}$ of elements annihilated by $x$.
Then:
\begin{enumerate}
  \item[\rm{(i)}] $s(\mathrm{M})\leq s(\mathrm{M}/x\mathrm{M})+1$.
  \item[\rm{(ii)}] Let $\mathfrak{p}_i$ be the prime ideals belonging to $\Supp(\mathrm{M})$ and such that
  \[
    \dim(\mathrm{A}/\mathfrak{p}_i)=\dim(\mathrm{M})\quad(1\leq i\leq m).
  \]
\end{enumerate}
If $x\notin\mathfrak{p}_i$ for every $i$, we have $\dim(\mathrm{M}/x\mathrm{M})\leq\dim(\mathrm{M})-1$.
\begin{enumerate}
  \item[\rm{(iii)}] If $\mathfrak{q}$ is an ideal of definition of $\mathrm{A}$, the polynomial
  \[
    \mathrm{P}_\mathfrak{q}({}_x\mathrm{M},n)-\mathrm{P}_\mathfrak{q}(\mathrm{M}/x\mathrm{M},n)
  \]
  has degree $\leq d(\mathrm{M})-1$.
\end{enumerate}
\end{lemma}

\begin{proof}
The assertion (i) is trivial, since if $(x_i)_{1\leq i\leq m}$ is such that for the module $\mathrm{N}=\mathrm{M}/x\mathrm{M}$, $\mathrm{N}/(x_1\mathrm{N}+\ldots+x_m\mathrm{N})$ is of finite length, this amounts to saying that $\mathrm{M}/(x\mathrm{M}+x_1\mathrm{M}+\ldots+x_m\mathrm{M})$ is isomorphic to this module.

For (ii), note that, if $\mathfrak{a}$ denotes the annihilator of $\mathrm{M}$, the prime ideals containing the annihilator of $\mathrm{M}/x\mathrm{M}$ are those containing $\mathfrak{a}+\mathrm{A}x$ $(\mathbf{0}_\mathrm{I},~1.7.5)$; none of the $\mathfrak{p}_i$ can therefore contain $\mathfrak{a}+\mathrm{A}x$ and by definition of $\dim(\mathrm{M})$ \sref{0.16.1.7} every chain of prime ideals containing $\mathfrak{a}+\mathrm{A}x$ has therefore length $\leq\dim(\mathrm{M})-1$.

\oldpage[0\textsubscript{IV-1}]{27}

Finally, to prove (iii), note that one has two short exact sequences
\[
  0\to{}_x\mathrm{M}\to\mathrm{M}\xrightarrow{x}x\mathrm{M}\to 0,\qquad 0\to x\mathrm{M}\to\mathrm{M}\to\mathrm{M}/x\mathrm{M}\to 0
\]
and the assertion follows immediately from the lemma \sref{0.16.2.2.2}.
\end{proof}

The proof of \sref{0.16.2.3} proceeds in three steps:

\begin{env}[16.2.3.2]
\label{0.16.2.3.2}
$\dim(\mathrm{M})\leq d(\mathrm{M})$.
\end{env}

We argue by induction on $d(\mathrm{M})$, the case $d(\mathrm{M})=0$ being trivial.
Suppose thus $d(\mathrm{M})\geq 1$, and let $\mathfrak{p}_0\in\Supp(\mathrm{M})$ such that $\dim(\mathrm{A}/\mathfrak{p}_0)=\dim(\mathrm{M})$.
This implies that $\mathfrak{p}_0$ is a \emph{minimal} element of $\Supp(\mathrm{M})$, hence is associated to $\mathrm{M}$ (Bourbaki, \emph{Alg.~comm.}, chap.~IV, \textsection~1, n\textsuperscript{o}~4, th.~2) and $\mathrm{M}$ therefore contains a submodule $\mathrm{N}$ isomorphic to $\mathrm{A}/\mathfrak{p}_0$ (\emph{loc.~cit.}, \textsection~1, n\textsuperscript{o}~1); as $d(\mathrm{M})\geq d(\mathrm{N})$ by \sref{0.16.2.2.2}, it suffices to prove that $\dim(\mathrm{N})\leq d(\mathrm{N})$.
Let then $\mathfrak{p}_0\subset\mathfrak{p}_1\subset\ldots\subset\mathfrak{p}_n$ be a chain of distinct prime ideals in $\mathrm{A}$, and let us show that $n\leq d(\mathrm{N})$.
This is evident if $n=0$.
Otherwise, there exists $x\in\mathfrak{p}_1\cap\mathfrak{r}$ such that $x\notin\mathfrak{p}_0$, for the relation $\mathfrak{p}_0\supset\mathfrak{p}_1\cap\mathfrak{r}$ would imply $\mathfrak{p}_0\supset\mathfrak{r}$ since $\mathfrak{p}_0$ does not contain $\mathfrak{p}_1$, hence $\mathfrak{p}_0$ would be maximal, which is absurd.
We have $\mathfrak{p}_i\in\Supp(\mathrm{N}/x\mathrm{N})$, for $i\geq 1$, hence $n-1\leq\dim(\mathrm{N}/x\mathrm{N})$ \sref{0.16.1.7}.
As ${}_x\mathrm{N}=0$ by the choice of $x$, we have $d(\mathrm{N}/x\mathrm{N})\leq d(\mathrm{N})-1$ by \sref{0.16.2.3.1}, (iii).
The induction hypothesis then implies $d(\mathrm{N}/x\mathrm{N})=\dim(\mathrm{N}/x\mathrm{N})\geq n-1$, hence $n-1\leq d(\mathrm{N})-1$.

\begin{env}[16.2.3.3]
\label{0.16.2.3.3}
$d(\mathrm{M})\leq s(\mathrm{M})$.
\end{env}

Let indeed $(x_i)_{1\leq i\leq n}$ be a family of elements of $\mathfrak{r}$ such that, if we set $\mathfrak{a}=x_1\mathrm{A}+\ldots+x_n\mathrm{A}$, $\mathrm{M}/\mathfrak{a}\mathrm{M}$ is of finite length.
The ideals associated to $\mathrm{M}/\mathfrak{a}\mathrm{M}$ are then maximal (Bourbaki, \emph{Alg.~comm.}, chap.~IV, \textsection~2, n\textsuperscript{o}~5, prop.~7) and since these are the only prime ideals containing $\mathfrak{q}=\mathfrak{a}+(\mathfrak{r}\cap\operatorname{Ann}(\mathrm{M}))$, $\mathrm{A}/\mathfrak{q}$ is Artinian (\emph{loc.~cit.}, prop.~9) and $\mathfrak{q}$ is an ideal of definition of $\mathrm{A}$.
But as $\mathfrak{q}^m\mathrm{M}/\mathfrak{q}^{m+1}\mathrm{M}=\mathfrak{a}^m\mathrm{M}/\mathfrak{a}^{m+1}\mathrm{M}$, and if $\mathrm{M}$ is generated by $r$ elements $z_j$, $\mathfrak{a}^m\mathrm{M}/\mathfrak{a}^{m+1}\mathrm{M}$ is an $(\mathrm{A}/\mathfrak{q})$-module generated by the canonical images of $x_1^{\nu_1}x_2^{\nu_2}\cdots x_n^{\nu_n}z_j$ where $\sum_i\nu_i=m$, hence by $r\binom{m+n-1}{n-1}$ elements; its length is therefore $\leq ar\binom{m+n-1}{n-1}$, where $a$ is a constant, and one deduces immediately that the degree of $\mathrm{P}_\mathfrak{q}(\mathrm{M},n)$ is $\leq n$; whence $d(\mathrm{M})\leq s(\mathrm{M})$.

\begin{env}[16.2.3.4]
\label{0.16.2.3.4}
$s(\mathrm{M})\leq\dim(\mathrm{M})$.
\end{env}

Let us note first that $n=\dim(\mathrm{M})$ is finite by \sref{0.16.2.3.2}.
We argue by induction on $n$; the proposition is evident for $n=0$, since $\mathrm{M}$ is then of finite length \sref{0.16.1.10}.
Suppose $n\geq 1$, and let $\mathfrak{p}_i$ $(1\leq i\leq m)$ be the prime ideals of $\Supp(\mathrm{M})$ such that $\dim(\mathrm{A}/\mathfrak{p}_i)=n$; the definition of the $\mathfrak{p}_i$ shows that they are minimal elements of $\Supp(\mathrm{M})$, hence finite in number (Bourbaki, \emph{Alg.~comm.}, chap.~IV, \textsection~1, n\textsuperscript{o}~4, th.~2).
As $n\geq 1$, the $\mathfrak{p}_i$ are not maximal, and there therefore exists $x\in\mathfrak{r}$ such that $x\notin\mathfrak{p}_i$ for every $i$ (Bourbaki, \emph{Alg.~comm.}, chap.~II, \textsection~1, n\textsuperscript{o}~1, prop.~2).
By virtue of \sref{0.16.2.3.1}, we have $s(\mathrm{M})\leq s(\mathrm{M}/x\mathrm{M})+1$ and $\dim(\mathrm{M})\geq\dim(\mathrm{M}/x\mathrm{M})+1$; the induction hypothesis gives the conclusion.

\begin{corollary}[16.2.4]
\label{0.16.2.4}
Under the hypotheses of \sref{0.16.2.3}, we have
\begin{equation}
\label{0.16.2.4.1}
\tag{16.2.4.1}
  \dim_{\widehat{\mathrm{A}}}(\widehat{\mathrm{M}})=\dim_\mathrm{A}(\mathrm{M}).
\end{equation}
\end{corollary}

\oldpage[0\textsubscript{IV-1}]{28}

\begin{proof}
Indeed, $\mathrm{M}/\mathfrak{r}^n\mathrm{M}$ and $\widehat{\mathrm{M}}/\mathfrak{r}^n\widehat{\mathrm{M}}$ are isomorphic, hence $d(\widehat{\mathrm{M}})=d(\mathrm{M})$.
\end{proof}

\begin{corollary}[16.2.5]
\label{0.16.2.5}
The dimension of a semi-local Noetherian ring $\mathrm{A}$ is finite and equal to the minimum number of elements of the radical of $\mathrm{A}$ generating an ideal of definition.
\end{corollary}

\begin{proof}
This is the equality $s(\mathrm{M})=\dim(\mathrm{M})$ for $\mathrm{M}=\mathrm{A}$.
\end{proof}

In particular:

\begin{corollary}[16.2.6]
\label{0.16.2.6}
Let $\mathrm{A}$ be a local Noetherian ring, $\mathfrak{m}$ its maximal ideal, $k=\mathrm{A}/\mathfrak{m}$ its residue field; we have
\begin{equation}
\label{0.16.2.6.1}
\tag{16.2.6.1}
  \dim(\mathrm{A})\leq\rg_k(\mathfrak{m}/\mathfrak{m}^2).
\end{equation}
\end{corollary}

\begin{proof}
Indeed, we know that if the $x_i$ $(1\leq i\leq n)$ are elements of $\mathfrak{m}$ whose classes mod.~$\mathfrak{m}^2$ form a basis of the $k$-vector space $\mathfrak{m}/\mathfrak{m}^2$, the $x_i$ generate $\mathfrak{m}$ (Bourbaki, \emph{Alg.~comm.}, chap.~II, \textsection~3, n\textsuperscript{o}~2, cor.~2 of prop.~4).
\end{proof}

\begin{proposition}[16.2.7]
\label{0.16.2.7}
Let $k$ be a field, $\mathrm{B}$ a graded $k$-algebra of finite type generated by its elements of degree $1$ and such that $\mathrm{B}_0=k$; let $\mathrm{E}$ be a graded $\mathrm{B}$-module of finite type, such that $(\mathbf{III},~2.5.4)$ for $n$ sufficiently large, $\rg_k(\mathrm{E}_n)=r(n)$ is of the form $\Phi(n)-\Phi(n-1)$, where $\Phi$ is a polynomial in $n$ of degree $d$.
Then there exist $d$ homogeneous elements of degree $\geq 1$ in $\mathrm{B}$, $t_1,\ldots,t_d$, forming a family $(t_i)_{1\leq i\leq d}$ such that $\mathrm{E}/(\sum_i t_i\mathrm{E})$ is of finite rank over $k$.
\end{proposition}

\begin{proof}
Let $\mathfrak{q}$ be the maximal ideal $\bigoplus_{n\geq 1}\mathrm{B}_n$ of $\mathrm{B}$, and, for every $n$, let $\mathrm{E}_n'$ be the sub-$\mathrm{B}$-module $\bigoplus_{h\geq n+1}\mathrm{E}_h$ of $\mathrm{E}$.
Every element of $\mathrm{B}-\mathfrak{q}$ determines an injective homothety in the Artinian $\mathrm{B}$-module $\mathrm{E}/\mathrm{E}_n'$, and this homothety is thus bijective (Bourbaki, \emph{Alg.}, chap.~VIII, \textsection~2, n\textsuperscript{o}~2, lemma~3), which implies that $(\mathrm{E}/\mathrm{E}_n')_\mathfrak{q}$ identifies with $\mathrm{E}/\mathrm{E}_n'$; as the residue field of $\mathrm{B}_\mathfrak{q}$ is $k$, one deduces from the hypothesis that $\operatorname{long}_{\mathrm{B}_\mathfrak{q}}((\mathrm{E}/\mathrm{E}_n')_\mathfrak{q})$ is a polynomial in $n$ of degree $d$. Moreover, $\mathrm{E}$ is a $\mathrm{B}$-module of finite type and $\mathrm{B}$ is generated over $k$ by its homogeneous elements of degree $1$, hence the filtration $(\mathrm{E}_n')$ is $\mathfrak{q}$-good (Bourbaki, \emph{Alg.~comm.}, chap.~III, \textsection~1, n\textsuperscript{o}~3, lemma~1), hence the filtration $((\mathrm{E}_n')_\mathfrak{q})$ on $\mathrm{E}_\mathfrak{q}$ is $\mathfrak{q}\mathrm{B}_\mathfrak{q}$-good.
Since $\mathrm{B}_\mathfrak{q}$ is a local Noetherian ring, one concludes from \sref{0.16.2.3} and \sref{0.16.2.2.1} that $\dim(\mathrm{E}_\mathfrak{q})=d$. We argue by induction on $d$, the assertion being evident if $d=0$.
Suppose $d\geq 1$, and observe that the minimal prime ideals $\mathfrak{p}_i$ $(1\leq i\leq r)$ in $\operatorname{Ass}(\mathrm{E})$ are \emph{graded} (Bourbaki, \emph{Alg.~comm.}, chap.~IV, \textsection~3, n\textsuperscript{o}~1, prop.~1), and that $\mathfrak{q}$ is distinct from the union of the $\mathfrak{p}_i$ since $d\geq 1$; there is therefore a \emph{homogeneous} element $t_1\in\mathfrak{q}$ not belonging to any of the $\mathfrak{p}_i$ (Bourbaki, \emph{Alg.~comm.}, chap.~III, \textsection~1, n\textsuperscript{o}~4, prop.~8).
As the prime ideals $(\mathfrak{p}_i)_\mathfrak{q}$ are the minimal elements of $\operatorname{Ass}(\mathrm{E}_\mathfrak{q})$ (Bourbaki, \emph{Alg.~comm.}, chap.~IV, \textsection~1, n\textsuperscript{o}~2, prop.~5), we have $\dim(\mathrm{E}_\mathfrak{q}/t_1\mathrm{E}_\mathfrak{q})=d-1$ \sref{0.16.3.4}.
It suffices then to apply the induction hypothesis to the graded $\mathrm{B}$-module $\mathrm{E}/t_1\mathrm{E}$ to obtain the conclusion.
\end{proof}

\subsection{Systems of parameters in a Noetherian local ring}
\label{subsection:0.16.3}

\begin{proposition}[16.3.1]
\label{0.16.3.1}
Let $\mathrm{A}$ be a Noetherian ring, $\mathfrak{p}$ a prime ideal of $\mathrm{A}$, and $n$ an integer.
The following conditions are equivalent:
\begin{enumerate}
  \item[a)] $\operatorname{ht}(\mathfrak{p})=\dim(\mathrm{A}_\mathfrak{p})\leq n$;
  \item[b)] There exist $n$ elements $x_i$ of $\mathrm{A}$ $(1\leq i\leq n)$ such that $\mathfrak{p}$ is minimal in the set of prime ideals containing the ideal $\mathfrak{a}$ generated by the $x_i$ (in other words, $\mathrm{V}(\mathfrak{p})$ is an irreducible component of $\mathrm{V}(\mathfrak{a})$).
\end{enumerate}
\end{proposition}

\begin{proof}
Condition b) implies that $\mathfrak{a}\mathrm{A}_\mathfrak{p}$ is an ideal of definition of $\mathrm{A}_\mathfrak{p}$, whence a) by virtue of \sref{0.16.2.5}.
Conversely, if a) holds, there exists an ideal of definition $\mathfrak{b}$ of $\mathrm{A}_\mathfrak{p}$ generated by $n$ elements $x_i/s$, with $s\in\mathrm{A}-\mathfrak{p}$.
By virtue of the one-to-one correspondence between prime ideals of $\mathrm{A}_\mathfrak{p}$ and prime ideals of $\mathrm{A}$ contained in $\mathfrak{p}$, the ideal $\mathfrak{a}$ of $\mathrm{A}$ generated by the $x_i$ satisfies b).
\end{proof}

\oldpage[0\textsubscript{IV-1}]{29}
For $n=1$, one obtains the special case:

\begin{corollary}[16.3.2]
\label{0.16.3.2}
\textup{(``Hauptidealsatz'').}
For a prime ideal in a Noetherian ring to be of height $\leq 1$, it is necessary and sufficient that it be minimal in the set of prime ideals containing a suitable principal ideal.
\end{corollary}

\begin{corollary}[16.3.3]
\label{0.16.3.3}
\textup{(Artin-Tate).}
Let $\mathrm{A}$ be an integral Noetherian ring.
The following conditions are equivalent:
\begin{enumerate}
  \item[a)] $\mathrm{A}$ is a semi-local ring of dimension $\leq 1$.
  \item[b)] There exists $f\neq 0$ in $\mathrm{A}$ such that $\mathrm{A}_f$ is a field.
\end{enumerate}
\end{corollary}

\begin{proof}
The condition $a)$ is equivalent to saying that $\mathrm{A}$ has only finitely many prime ideals $\mathfrak{m}_i\neq 0$, and that these ideals are all maximal (0 being a prime ideal).
As the product of the $\mathfrak{m}_i$ is not zero, since $\mathrm{A}$ is integral, an element $f\neq 0$ of this product belongs to every maximal ideal, hence $(0)$ is the only prime ideal of the integral ring $\mathrm{A}_f$, in other words $\mathrm{A}_f$ is a field.
(One notes that this part of the argument does not use the fact that $\mathrm{A}$ is Noetherian.)
Let us now prove that $b)$ implies $a)$; the hypothesis $b)$ means that every prime ideal $\mathfrak{p}\neq 0$ of $\mathrm{A}$ contains $f$.
Let $\mathfrak{p}_i$ $(1\leq i\leq n)$ be the minimal prime ideals among those containing $f$ (which are finite in number since $\mathrm{A}/f\mathrm{A}$ is Noetherian); it suffices to prove that the $\mathfrak{p}_i$ are \emph{maximal} ideals, for then every nonzero prime ideal necessarily contains one of the $\mathfrak{p}_i$ and is therefore equal to it.
Suppose then that there exists a maximal ideal $\mathfrak{m}$ containing one of the $\mathfrak{p}_i$ and distinct from the $\mathfrak{p}_i$; $\mathfrak{m}$ is therefore not contained in the union of the $\mathfrak{p}_i$ (Bourbaki, \emph{Alg.~comm.}, chap.~II, \textsection~1, n\textsuperscript{o}~1, prop.~2), in other words there exists $g\in\mathfrak{m}$ not belonging to any of the $\mathfrak{p}_i$.
Let $\mathfrak{q}$ be a minimal prime ideal among those containing $g$; it follows from the ``Hauptidealsatz'' \sref{0.16.3.2} that $\mathfrak{q}$ is of height 1, hence $\neq 0$; it therefore contains by assumption $f$, hence also one of the $\mathfrak{p}_i$, and as $g$ does not belong to any of the $\mathfrak{p}_i$, $\mathfrak{q}$ is distinct from the $\mathfrak{p}_i$; it would thus be of height $\geq 2$, which is a contradiction.
\end{proof}

\begin{proposition}[16.3.4]
\label{0.16.3.4}
Let $\mathrm{A}$ be a semi-local Noetherian ring, $\mathfrak{r}$ its radical, $\mathrm{M}$ an $\mathrm{A}$-module of finite type, $\mathfrak{p}_i$ the elements of $\Supp(\mathrm{M})$ such that $\dim(\mathrm{A}/\mathfrak{p}_i)=\dim(\mathrm{M})$.
Then the $\mathfrak{p}_i$ are minimal elements of $\operatorname{Ass}(\mathrm{M})$, and $\dim(\mathrm{M}/\mathfrak{p}_i\mathrm{M})=\dim(\mathrm{M})$.
For every $x\in\mathfrak{r}$, we have
\begin{equation}
\label{0.16.3.4.1}
\tag{16.3.4.1}
  \dim(\mathrm{M}/x\mathrm{M})\geq\dim(\mathrm{M})-1.
\end{equation}
Furthermore, for the two members of \sref{0.16.3.4.1} to be equal, it is necessary and sufficient that $x$ not belong to any of the $\mathfrak{p}_i$.
\end{proposition}

\begin{proof}
The first assertion is immediate, since the $\mathfrak{p}_i$ are by definition the minimal elements of $\Supp(\mathrm{M})$ \sref{0.16.1.6}, and they are also the minimal elements of $\operatorname{Ass}(\mathrm{M})$ (Bourbaki, \emph{Alg.~comm.}, chap.~IV, \textsection~1, n\textsuperscript{o}~4, th.~2).
Furthermore, $\Supp(\mathrm{M}/\mathfrak{p}_i\mathrm{M})$ is the set of prime ideals of $\mathrm{A}$ containing $\mathfrak{p}_i$ $(\mathbf{0}_\mathrm{I},~1.7.5)$, hence $\dim(\mathrm{M}/\mathfrak{p}_i\mathrm{M})=\dim(\mathrm{A}/\mathfrak{p}_i)=\dim(\mathrm{M})$.

Set $\mathrm{N}=\mathrm{M}/x\mathrm{M}$.
If $x_1,\ldots,x_n$ are elements of $\mathfrak{r}$ such that $\mathrm{M}/(x\mathrm{M}+x_1\mathrm{M}+\ldots+x_n\mathrm{M})$ is of finite length, we have $s(\mathrm{M})\leq n+1$, hence $\dim(\mathrm{M})\leq n+1$ by \sref{0.16.2.3}, and as one can choose them so that $n=s(\mathrm{N})=\dim(\mathrm{N})$, we get \sref{0.16.3.4.1}.

The fact that the two members of \sref{0.16.3.4.1} are equal when $x$ does not belong to any of the $\mathfrak{p}_i$ follows from \sref{0.16.3.4.1} and \sref{0.16.2.3.1}.
Conversely, if \oldpage[0\textsubscript{IV-1}]{30} $\dim(\mathrm{M}/x\mathrm{M})=\dim(\mathrm{M})-1$, none of the $\mathfrak{p}_i$ can belong to $\Supp(\mathrm{M}/x\mathrm{M})$, and the prime ideals in the support of the latter are those containing $\operatorname{Ann}(\mathrm{M})+\mathrm{A}x$; since the $\mathfrak{p}_i$ contain $\operatorname{Ann}(\mathrm{M})$, they cannot contain $x$.
\end{proof}

\begin{corollary}[16.3.5]
\label{0.16.3.5}
With the notations of \sref{0.16.3.4}, for every ideal $\mathfrak{a}\subset\mathfrak{p}_i$, we have $\dim(\mathrm{M}/\mathfrak{a}\mathrm{M})=\dim(\mathrm{M})$ and if one sets $\mathrm{N}=\mathrm{M}/\mathfrak{a}\mathrm{M}$, the relation $\dim(\mathrm{M}/x\mathrm{M})=\dim(\mathrm{M})-1$ implies $\dim(\mathrm{N}/x\mathrm{N})=\dim(\mathrm{N})-1$.
\end{corollary}

\begin{proof}
Indeed, $\Supp(\mathrm{N})$ is the set of prime ideals containing $\mathfrak{a}+\operatorname{Ann}(\mathrm{M})$, hence $\mathfrak{p}_i$ belongs to $\Supp(\mathrm{N})$ and $\dim(\mathrm{N})\leq\dim(\mathrm{M})=\dim(\mathrm{A}/\mathfrak{p}_i)$, hence $\dim(\mathrm{A}/\mathfrak{p}_i)=\dim(\mathrm{N})$; furthermore the prime ideals $\mathfrak{q}\in\Supp(\mathrm{N})$ such that $\dim(\mathrm{A}/\mathfrak{q})=\dim(\mathrm{N})$ are certain of the $\mathfrak{p}_j$; if $x$ does not belong to any of the $\mathfrak{p}_j$, we have then $\dim(\mathrm{N}/x\mathrm{N})=\dim(\mathrm{N})-1$ by virtue of \sref{0.16.3.4}.
\end{proof}

\begin{definition}[16.3.6]
\label{0.16.3.6}
Let $\mathrm{A}$ be a semi-local Noetherian ring, $\mathfrak{r}$ its radical, $\mathrm{M}$ an $\mathrm{A}$-module of finite type, and set $n=\dim(\mathrm{M})$.
A \emph{system of parameters} for $\mathrm{M}$ is any family $(x_i)_{1\leq i\leq n}$ of $n$ elements of $\mathfrak{r}$ such that $\mathrm{M}/(x_1\mathrm{M}+\ldots+x_n\mathrm{M})$ is of \emph{finite length}.
\end{definition}

We have seen \sref{0.16.2.3} that such systems always exist.
We have seen in the course of the proof of \sref{0.16.2.3.3} that if $\mathfrak{J}=\sum_{i=1}^n x_i\mathrm{A}$ is such that $\mathrm{M}/\mathfrak{J}\mathrm{M}$ is of finite length, $\mathfrak{q}=\mathfrak{J}+\operatorname{Ann}(\mathrm{M})$ is an \emph{ideal of definition} of $\mathrm{A}$, and conversely, if this is so, $\mathrm{M}/\mathfrak{J}\mathrm{M}$ is a module of finite type over the Artinian ring $\mathrm{A}/\mathfrak{q}$, hence is of finite length.
It amounts to the same to say that $(x_i)_{1\leq i\leq n}$ is a system of parameters \emph{for} $\mathrm{M}$ or \emph{for} $\mathrm{A}/\operatorname{Ann}(\mathrm{M})$ (or equivalently that their images in $\mathrm{A}/\operatorname{Ann}(\mathrm{M})$ form a system of parameters for this ring).
If $(x_i)_{1\leq i\leq n}$ is a system of parameters for $\mathrm{M}$, so is $(x_i^k)$ for every integer $k\geq 1$, since one can reduce, by the preceding, to the case $\mathrm{M}=\mathrm{A}$, and if $\mathfrak{J}$ is an ideal of definition of $\mathrm{A}$, $\sum_i x_i^k\mathrm{A}$ contains $\mathfrak{J}^{kn}$, hence is also an ideal of definition.

\begin{proposition}[16.3.7]
\label{0.16.3.7}
Let $\mathrm{A}$ be a semi-local Noetherian ring, $\mathfrak{r}$ its radical, $x_1,\ldots,x_k$ elements of $\mathfrak{r}$, $\mathrm{M}$ an $\mathrm{A}$-module of finite type.
Then
\begin{equation}
\label{0.16.3.7.1}
\tag{16.3.7.1}
  \dim(\mathrm{M}/(x_1\mathrm{M}+\ldots+x_k\mathrm{M}))\geq\dim(\mathrm{M})-k.
\end{equation}
Furthermore, for the two members of \sref{0.16.3.7.1} to be equal, it is necessary and sufficient that there exist elements $x_i$ $(k+1\leq i\leq n=\dim(\mathrm{M}))$ of $\mathfrak{r}$ such that $(x_i)_{1\leq i\leq n}$ is a system of parameters for $\mathrm{M}$.
\end{proposition}

\begin{proof}
The inequality \sref{0.16.3.7.1} follows from \sref{0.16.3.4.1} by induction on $k$.
If the two members of \sref{0.16.3.7.1} are equal, and if $x_{k+1},\ldots,x_n$ is a system of parameters for $\mathrm{M}/(x_1\mathrm{M}+\ldots+x_k\mathrm{M})$, it is clear that $\mathrm{M}/(x_1\mathrm{M}+\ldots+x_k\mathrm{M}+x_{k+1}\mathrm{M}+\ldots+x_n\mathrm{M})$ is of finite length, hence $(x_i)_{1\leq i\leq n}$ is a system of parameters for $\mathrm{M}$; conversely, if there exist $x_i$ $(k+1\leq i\leq n)$ having this property and if one sets $\mathrm{N}=\mathrm{M}/(x_1\mathrm{M}+\ldots+x_k\mathrm{M})$, it is clear that $\mathrm{N}/(x_{k+1}\mathrm{N}+\ldots+x_n\mathrm{N})$ is of finite length, hence $\dim(\mathrm{N})\leq n-k$, and the two members are equal by virtue of \sref{0.16.3.7.1}.
\end{proof}

\oldpage[0\textsubscript{IV-1}]{31}

\begin{proposition}[16.3.8]
\label{0.16.3.8}
Let $\mathrm{A}$ be a semi-local Noetherian ring, $\mathrm{X}=\Spec(\mathrm{A})$ its spectrum, $\mathfrak{a}$ an ideal of $\mathrm{A}$ distinct from $\mathrm{A}$ such that $\codim(\mathrm{V}(\mathfrak{a}),\mathrm{X})=r>0$.
Then there exist $r$ elements $x_1,\ldots,x_r$ of $\mathfrak{a}$, forming part of a system of parameters of $\mathrm{A}$, such that
\[
  \codim(\mathrm{V}(\mathfrak{a}),\mathrm{V}(\mathrm{A}x_1+\ldots+\mathrm{A}x_r))=0.
\]
\end{proposition}

\begin{proof}
Let $\mathfrak{p}_i$ $(1\leq i\leq m)$ be the minimal prime ideals of the ring $\mathrm{A}$, $\mathfrak{q}_j$ $(1\leq j\leq n)$ the minimal prime ideals among those containing $\mathfrak{a}$; we have \sref{0.14.2.1} $\codim(\mathrm{V}(\mathfrak{a}),\mathrm{X})=\inf_j(\dim(\mathrm{A}_{\mathfrak{q}_j}))$ \sref{0.16.1.3.2}.
The hypothesis $r>0$ implies that $\mathfrak{a}$ is not contained in any of the $\mathfrak{p}_i$; hence there exists $x\in\mathfrak{a}$ not belonging to any of the $\mathfrak{p}_i$ (Bourbaki, \emph{Alg.~comm.}, chap.~II, \textsection~1, n\textsuperscript{o}~1, prop.~2).
One has as above $\codim(\mathrm{V}(\mathfrak{a}),\mathrm{V}(\mathrm{A}x))=\inf_j(\dim((\mathrm{A}/\mathrm{A}x)_{\mathfrak{q}_j/\mathrm{A}x}))=\inf_j(\dim(\mathrm{A}_{\mathfrak{q}_j}/\mathrm{A}_{\mathfrak{q}_j}x))$; but as $x/1$ is not in any of the minimal prime ideals of $\mathrm{A}_{\mathfrak{q}_j}$, $\dim(\mathrm{A}_{\mathfrak{q}_j}/\mathrm{A}_{\mathfrak{q}_j}x)=\dim(\mathrm{A}_{\mathfrak{q}_j})-1$ \sref{0.16.3.4}, whence $\codim(\mathrm{V}(\mathfrak{a}),\mathrm{V}(\mathrm{A}x))=r-1$.
For the same reason, $\dim(\mathrm{A}/\mathrm{A}x)=\dim(\mathrm{A})-1$.
We argue by induction on $r$; if $r>1$, there exist $r-1$ elements $x_i'$ $(2\leq i\leq r)$ of $\mathrm{A}'=\mathrm{A}/\mathrm{A}x$, forming part of a system of parameters of $\mathrm{A}'$ and such that $\codim(\mathrm{V}(\mathfrak{a}/\mathrm{A}x),\mathrm{V}(\mathrm{A}'x_2'+\ldots+\mathrm{A}'x_r'))=0$.
Let $(x_i')_{r+1\leq i\leq s}$ be a family of elements of $\mathrm{A}'$ such that $(x_i')_{2\leq i\leq s}$ is a system of parameters of $\mathrm{A}'$; for every $i$ such that $2\leq i\leq s$, let $x_i$ be a representative of the class $x_i'$ in $\mathrm{A}$, with $x_i\in\mathfrak{a}$ for $2\leq i\leq r$. It is immediate that $\mathrm{A}/(\mathrm{A}x+\mathrm{A}x_2+\ldots+\mathrm{A}x_s)$ is of finite length, and as $\dim(\mathrm{A})=s$, we see that $x_1=x$ and the $x_i$ for $i\geq 2$ satisfy the conditions of the statement.
\end{proof}

\begin{proposition}[16.3.9]
\label{0.16.3.9}
Let $\mathrm{A}$, $\mathrm{B}$ be two local Noetherian rings, $\mathfrak{m}$ the maximal ideal of $\mathrm{A}$, $k$ its residue field, $\varphi:\mathrm{A}\to\mathrm{B}$ a local homomorphism.
\begin{enumerate}
  \item[\rm{(i)}] We have
  \begin{equation}
  \label{0.16.3.9.1}
  \tag{16.3.9.1}
    \dim(\mathrm{B})\leq\dim(\mathrm{A})+\dim(\mathrm{B}\otimes_\mathrm{A}k).
  \end{equation}
  \item[\rm{(ii)}] Let $\mathrm{M}\neq 0$ be an $\mathrm{A}$-module of finite type, $\mathrm{N}\neq 0$ a $\mathrm{B}$-module of finite type.
  We have
  \begin{equation}
  \label{0.16.3.9.2}
  \tag{16.3.9.2}
    \dim_\mathrm{B}(\mathrm{M}\otimes_\mathrm{A}\mathrm{N})\leq\dim_\mathrm{A}(\mathrm{M})+\dim_{\mathrm{B}\otimes_\mathrm{A}k}(\mathrm{N}\otimes_\mathrm{A}k).
  \end{equation}
\end{enumerate}
\end{proposition}

\begin{proof}
Let $m$ be the dimension of $\mathrm{A}$, $(s_i)_{1\leq i\leq m}$ a system of parameters of $\mathrm{A}$ \sref{0.16.3.6}, so that $\mathfrak{a}=\sum_i s_i\mathrm{A}$, $\mathrm{A}/\mathfrak{a}$ is an $\mathrm{A}$-module of finite length.
Then the Artinian ring $\mathrm{A}/\mathfrak{a}$ has nilpotent maximal ideal $\mathfrak{m}/\mathfrak{a}$; the ideal $\mathfrak{m}\mathrm{B}/\mathfrak{a}\mathrm{B}$ of $\mathrm{B}/\mathfrak{a}\mathrm{B}$, image of $(\mathfrak{m}/\mathfrak{a})\otimes_\mathrm{A}\mathrm{B}$, is therefore also nilpotent, and hence $\dim(\mathrm{B}/\mathfrak{a}\mathrm{B})=\dim(\mathrm{B}/\mathfrak{m}\mathrm{B})$ (the spectra of these two rings having the same underlying space).
On the other hand $\mathrm{B}/\mathfrak{a}\mathrm{B}=\mathrm{B}/(s_1\mathrm{B}+\ldots+s_m\mathrm{B})$, and the images of the $s_i$ in $\mathrm{B}$ belong to the maximal ideal of $\mathrm{B}$.
Hence by \sref{0.16.3.7}, we have
\[
  \dim(\mathrm{B})\leq m+\dim(\mathrm{B}\otimes_\mathrm{A}k)
\]
which is \sref{0.16.3.9.1}.

To prove \sref{0.16.3.9.2}, note that if $\mathfrak{r}$ (resp.~$\mathfrak{s}$) is the annihilator of $\mathrm{M}$ (resp.~$\mathrm{N}$) \oldpage[0\textsubscript{IV-1}]{32} in $\mathrm{A}$ (resp.~$\mathrm{B}$), we have $\dim_\mathrm{A}(\mathrm{M})=\dim(\mathrm{A}/\mathfrak{r})$, $\dim_\mathrm{B}(\mathrm{N})=\dim(\mathrm{B}/\mathfrak{s})$ by definition \sref{0.16.1.7}; on the other hand, $\Supp(\mathrm{M}\otimes_\mathrm{A}\mathrm{N})$ is a closed subset of $\Spec(\mathrm{B})$ that identifies $(\mathbf{I},~9.1.13)$ to $\Supp(\mathrm{M})\times_{\Spec(\mathrm{A})}\Supp(\mathrm{N})=\Spec((\mathrm{B}/\mathfrak{s})\otimes_\mathrm{A}(\mathrm{A}/\mathfrak{r}))=\Spec(\mathrm{B}/(\mathfrak{s}+\mathfrak{r}\mathrm{B}))$.
As $\mathrm{N}\otimes_\mathrm{A}k$ is a $\mathrm{B}$-module of finite type, we have $\dim_{\mathrm{B}\otimes_\mathrm{A}k}(\mathrm{N}\otimes_\mathrm{A}k)=\dim_\mathrm{B}(\mathrm{N}\otimes_\mathrm{A}k)$ \sref{0.16.1.9}, and $\Supp(\mathrm{N}\otimes_\mathrm{A}k)$ is equal to $\Spec(\mathrm{B}/(\mathfrak{s}+\mathfrak{m}\mathrm{B}))$ by the same reasoning as above.
Applying \sref{0.16.3.9.1} to $\mathrm{A}/\mathfrak{r}$ and to $\mathrm{B}/(\mathfrak{s}+\mathfrak{r}\mathrm{B})$, we get
\[
  \dim(\mathrm{B}/(\mathfrak{s}+\mathfrak{r}\mathrm{B}))\leq\dim(\mathrm{A}/\mathfrak{r})+\dim(\mathrm{B}/(\mathfrak{s}+\mathfrak{m}\mathrm{B}))
\]
which is nothing other than \sref{0.16.3.9.2}.
\end{proof}

\begin{corollary}[16.3.10]
\label{0.16.3.10}
Under the hypotheses of \sref{0.16.3.9}, if $\mathfrak{m}\mathrm{B}$ is an ideal of definition of $\mathrm{B}$, we have $\dim(\mathrm{B})\leq\dim(\mathrm{A})$; if in addition $\mathrm{A}$ is integral and $\dim(\mathrm{A})=\dim(\mathrm{B})$, $\varphi$ is injective.
\end{corollary}

\begin{proof}
The first assertion follows from \sref{0.16.3.9} since then $\dim(\mathrm{B}/\mathfrak{m}\mathrm{B})=0$.
One can moreover replace $\mathrm{A}$ by $\varphi(\mathrm{A})$, hence $\dim(\mathrm{B})\leq\dim(\varphi(\mathrm{A}))$; if $\mathfrak{a}=\Ker(\varphi)\neq 0$, and if $\mathrm{A}$ is integral, we have $\dim(\varphi(\mathrm{A}))=\dim(\mathrm{A}/\mathfrak{a})<\dim(\mathrm{A})$ \sref{0.16.1.2.2}, hence one can have $\dim(\mathrm{A})=\dim(\mathrm{B})$ only if $\mathfrak{a}=0$.
\end{proof}

\subsection[Depth and codepth]{Depth and
codepth\protect\footnote{In the third part of Chapter~III, we develop the
notion of depth from the cohomological point of view and in a more general
framework.}}
\label{subsection:0.16.4}

\begin{proposition}[16.4.1]
\label{0.16.4.1}
Let $\mathrm{A}$ be a local Noetherian ring, $\mathfrak{m}$ its maximal ideal, $\mathrm{M}$ an $\mathrm{A}$-module of finite type.
Then every $\mathrm{M}$-regular sequence $(x_i)_{1\leq i\leq r}$ \sref{0.15.1.7} formed of elements of $\mathfrak{m}$, is part of a system of parameters for $\mathrm{M}$.
\end{proposition}

\begin{proof}
We argue by induction on $r$; consider the $\mathrm{A}$-module
\[
  \mathrm{N}=\mathrm{M}/(x_1\mathrm{M}+\ldots+x_{r-1}\mathrm{M});
\]
by the induction hypothesis $x_1,\ldots,x_{r-1}$ form part of a system of parameters for $\mathrm{M}$, hence $\dim(\mathrm{N})=\dim(\mathrm{M})-(r-1)$ \sref{0.16.3.7}.
As by hypothesis the homothety of ratio $x_r$ in $\mathrm{N}$ is injective, $x_r$ does not belong to any of the prime ideals associated to $\mathrm{N}$, hence \sref{0.16.3.4} $\dim(\mathrm{N}/x_r\mathrm{N})=\dim(\mathrm{N})-1$, which can also be written
\[
  \dim(\mathrm{M}/(x_1\mathrm{M}+\ldots+x_r\mathrm{M}))=\dim(\mathrm{M})-r;
\]
the conclusion then follows from \sref{0.16.3.7}.
\end{proof}

For an $\mathrm{M}$-regular sequence formed of elements of $\mathfrak{m}$, $\mathrm{M}\neq 0$, there are at most $\dim(\mathrm{M})$ elements.

\begin{lemma}[16.4.2]
\label{0.16.4.2}
Let $\mathrm{A}$ be a local Noetherian ring, $\mathfrak{m}$ its maximal ideal, $k=\mathrm{A}/\mathfrak{m}$ its residue field, $\mathrm{M}$ an $\mathrm{A}$-module.
For an $\mathrm{M}$-regular sequence $(x_i)_{1\leq i\leq n}$ of elements of $\mathfrak{m}$ to be maximal, it is necessary and sufficient that $\Hom_\mathrm{A}(k,\mathrm{M}/(x_1\mathrm{M}+\ldots+x_n\mathrm{M}))\neq 0$.
\end{lemma}

\begin{proof}
To say that $(x_i)$ is maximal means that, for no $x\in\mathfrak{m}$, the homothety of ratio $x$ in $\mathrm{N}=\mathrm{M}/(x_1\mathrm{M}+\ldots+x_n\mathrm{M})$ is injective, or equivalently that $\mathfrak{m}$ is contained in the union of the prime ideals associated to $\mathrm{N}$ (Bourbaki, \emph{Alg.~comm.}, chap.~IV, \textsection~1, n\textsuperscript{o}~1, cor.~2 of prop.~2), hence equal to one of them, since it is a \oldpage[0\textsubscript{IV-1}]{33} maximal ideal (Bourbaki, \emph{Alg.~comm.}, chap.~II, \textsection~1, n\textsuperscript{o}~1, prop.~2); in other words, there is an element $z\in\mathrm{N}$ whose annihilator is $\mathfrak{m}$, and hence the submodule $\mathrm{A}z$ of $\mathrm{N}$ is isomorphic to $\mathrm{A}/\mathfrak{m}=k$; whence the lemma.
\end{proof}

\begin{lemma}[16.4.3]
\label{0.16.4.3}
Let $\mathrm{A}$ be a local Noetherian ring, $\mathfrak{m}$ its maximal ideal, $k=\mathrm{A}/\mathfrak{m}$ its residue field, $\mathrm{M}$ an $\mathrm{A}$-module, $(x_i)_{1\leq i\leq n}$ an $\mathrm{M}$-regular sequence of elements of $\mathfrak{m}$.
Then the $\mathrm{A}$-modules $\Hom_\mathrm{A}(k,\mathrm{M}/(x_1\mathrm{M}+\ldots+x_n\mathrm{M}))$ and $\Ext^n_\mathrm{A}(k,\mathrm{M})$ are isomorphic, and for $n\geq 1$, they are also isomorphic to $\Ext^{n-1}_\mathrm{A}(k,\mathrm{M}/x_1\mathrm{M})$.
\end{lemma}

\begin{proof}
We argue by induction on $n$, the proposition being evident for $n=0$.
Set $\mathrm{N}=\mathrm{M}/x_1\mathrm{M}$; the sequence $(x_i)_{2\leq i\leq n}$ is $\mathrm{N}$-regular, hence
\[
  \Hom_\mathrm{A}(k,\mathrm{M}/(x_1\mathrm{M}+\ldots+x_n\mathrm{M}))=\Hom_\mathrm{A}(k,\mathrm{N}/(x_2\mathrm{N}+\ldots+x_n\mathrm{N}))
\]
is isomorphic to $\Ext^{n-1}_\mathrm{A}(k,\mathrm{N})$ by the induction hypothesis.
Consider then the short exact sequence $0\to\mathrm{M}\xrightarrow{x_1}\mathrm{M}\to\mathrm{M}/x_1\mathrm{M}=\mathrm{N}\to 0$; one deduces the exact sequence of Ext
\begin{equation}
\label{0.16.4.3.1}
\tag{16.4.3.1}
  \Ext^{n-1}_\mathrm{A}(k,\mathrm{M})\to\Ext^{n-1}_\mathrm{A}(k,\mathrm{N})\to\Ext^n_\mathrm{A}(k,\mathrm{M})\xrightarrow{x_1}\Ext^n_\mathrm{A}(k,\mathrm{M}).
\end{equation}
By the induction hypothesis, $\Ext^{n-1}_\mathrm{A}(k,\mathrm{M})$ is isomorphic to $\Hom_\mathrm{A}(k,\mathrm{M}/(x_1\mathrm{M}+\ldots+x_{n-1}\mathrm{M}))$, which is zero since $(x_i)_{1\leq i\leq n-1}$ is not a maximal $\mathrm{M}$-regular sequence \sref{0.16.4.2}.
On the other hand, as $x_1\in\mathfrak{m}$, the homothety of ratio $x_1$ in the $\mathrm{A}$-module $\Hom_\mathrm{A}(k,\mathrm{T})$ is zero for every $\mathrm{A}$-module $\mathrm{T}$, hence it is the same for every $\mathrm{A}$-module $\Ext^n_\mathrm{A}(k,\mathrm{T})$; the assertions of the lemma then follow immediately from the exact sequence \sref{0.16.4.3.1}.
\end{proof}

\begin{corollary}[16.4.4]
\label{0.16.4.4}
Let $\mathrm{A}$ be a local Noetherian ring, $\mathfrak{m}$ its maximal ideal, $k=\mathrm{A}/\mathfrak{m}$ its residue field, $\mathrm{M}$ an $\mathrm{A}$-module of finite type $\neq 0$.
All maximal $\mathrm{M}$-regular sequences of elements of $\mathfrak{m}$ have the same number of elements, which is the smallest integer $n$ such that $\Ext^n_\mathrm{A}(k,\mathrm{M})\neq 0$.
\end{corollary}

\begin{definition}[16.4.5]
\label{0.16.4.5}
Let $\mathrm{A}$ be a local Noetherian ring, $\mathfrak{m}$ its maximal ideal, $\mathrm{M}$ an $\mathrm{A}$-module of finite type.
We call the upper bound of the numbers of elements in $\mathrm{M}$-regular sequences formed of elements of $\mathfrak{m}$ the \emph{depth} of $\mathrm{M}$, and denote it $\operatorname{prof}_\mathrm{A}\mathrm{M}$ or $\operatorname{prof}(\mathrm{M})$.
We thus have $\operatorname{prof}(\mathrm{M})=+\infty$ if and only if $\mathrm{M}=0$, and, by \sref{0.16.4.1},
\begin{equation}
\label{0.16.4.5.1}
\tag{16.4.5.1}
  \operatorname{prof}(\mathrm{M})\leq\dim(\mathrm{M})\quad\text{if}\quad\mathrm{M}\neq 0.
\end{equation}
We evidently have $\operatorname{prof}(\mathrm{M}^k)=\operatorname{prof}(\mathrm{M})$ for every $k>0$.
\end{definition}

\begin{proposition}[16.4.6]
\label{0.16.4.6}
Let $\mathrm{A}$ be a local Noetherian ring, $\mathfrak{m}$ its maximal ideal, $\mathrm{M}$ an $\mathrm{A}$-module of finite type.
\begin{enumerate}
  \item[\rm{(i)}]\label{0.16.4.6.i} For $\operatorname{prof}(\mathrm{M})=0$, it is necessary and sufficient that $\mathfrak{m}\in\operatorname{Ass}(\mathrm{M})$ (or equivalently that there exists a submodule of $\mathrm{M}$ isomorphic to the residue field $k=\mathrm{A}/\mathfrak{m}$ of $\mathrm{A}$).
  \item[\rm{(ii)}] For every $\mathrm{M}$-regular element $x$ belonging to the maximal ideal of $\mathrm{A}$, we have
  \begin{equation}
  \label{0.16.4.6.1}
  \tag{16.4.6.1}
    \operatorname{prof}(\mathrm{M}/x\mathrm{M})=\operatorname{prof}(\mathrm{M})-1.
  \end{equation}
  \item[\rm{(iii)}] If $\mathrm{M}\neq 0$, we have
  \begin{equation}
  \label{0.16.4.6.2}
  \tag{16.4.6.2}
    \operatorname{prof}(\mathrm{M})\leq\inf_{\mathfrak{p}\in\operatorname{Ass}(\mathrm{M})}\dim(\mathrm{A}/\mathfrak{p})\leq\dim(\mathrm{M}).
  \end{equation}
  \oldpage[0\textsubscript{IV-1}]{34}
  \item[\rm{(iv)}] We have $\operatorname{prof}_\mathrm{A}(\mathrm{M})=\operatorname{prof}_{\widehat{\mathrm{A}}}(\widehat{\mathrm{M}})$.
\end{enumerate}
\end{proposition}

\begin{proof}
(i) To say $\operatorname{prof}(\mathrm{M})=0$ means that there is no $x\in\mathfrak{m}$ such that the homothety of ratio $x$ in $\mathrm{M}$ is injective, that is, every $x\in\mathfrak{m}$ belongs to a prime ideal associated to $\mathrm{M}$ (Bourbaki, \emph{Alg.~comm.}, chap.~IV, \textsection~1, n\textsuperscript{o}~1, cor.~2 of prop.~2); but $\mathfrak{m}$ cannot be the union of prime ideals associated to $\mathrm{M}$ unless it is equal to one of them (Bourbaki, \emph{Alg.~comm.}, chap.~II, \textsection~1, n\textsuperscript{o}~1, prop.~2).

(ii) If $(x_i)_{1\leq i\leq m}$ is a $(\mathrm{M}/x\mathrm{M})$-regular sequence of elements of $\mathfrak{m}$, the sequence formed of $x$ and the $x_i$ is $\mathrm{M}$-regular, and it follows from the definition that it is maximal when the first sequence is; whence the conclusion.

(iii) We argue by induction on $n=\operatorname{prof}(\mathrm{M})$, the assertion being evident for $n=0$.
We use the following lemma:
\end{proof}

\begin{lemma}[16.4.6.3]
\label{0.16.4.6.3}
Let $\mathrm{A}$ be a Noetherian ring, $\mathrm{M}$ an $\mathrm{A}$-module of finite type, $t\in\mathrm{A}$ an $\mathrm{M}$-regular element.
Then, for every $\mathfrak{p}\in\operatorname{Ass}(\mathrm{M})$, every \emph{minimal} prime ideal among those containing $\mathfrak{p}+\mathrm{A}t$ belongs to $\operatorname{Ass}(\mathrm{M}/t\mathrm{M})$.
\end{lemma}

\begin{proof}[Proof of the lemma]
We know that there exists a short exact sequence
\[
  0\to\mathrm{M}'\to\mathrm{M}\to\mathrm{M}''\to 0
\]
such that $\operatorname{Ass}(\mathrm{M}')=\{\mathfrak{p}\}$, $\operatorname{Ass}(\mathrm{M}'')=\operatorname{Ass}(\mathrm{M})-\{\mathfrak{p}\}$ (Bourbaki, \emph{Alg.~comm.}, chap.~IV, \textsection~1, n\textsuperscript{o}~1, prop.~4); as $t$ is $\mathrm{M}$-regular, it is also $\mathrm{M}'$-regular and $\mathrm{M}''$-regular (\emph{loc.~cit.}, \textsection~1, n\textsuperscript{o}~1, cor.~2 of prop.~2).
One deduces that the sequence
\[
  0\to\mathrm{M}'/t\mathrm{M}'\to\mathrm{M}/t\mathrm{M}\to\mathrm{M}''/t\mathrm{M}''\to 0
\]
is exact \sref{0.15.1.18}.
Consequently, $\operatorname{Ass}(\mathrm{M}'/t\mathrm{M}')\subset\operatorname{Ass}(\mathrm{M}/t\mathrm{M})$.
But the prime ideals of $\mathrm{A}$ containing $\mathfrak{p}$ are the points of $\Supp(\mathrm{M}')$ (\emph{loc.~cit.}, \textsection~1, n\textsuperscript{o}~3, prop.~7), hence the prime ideals containing $\mathfrak{p}+\mathrm{A}t$ are the points of the support of $\mathrm{M}'/t\mathrm{M}'=\mathrm{M}'\otimes_\mathrm{A}(\mathrm{A}/t\mathrm{A})$ $(\mathbf{0}_\mathrm{I},~1.7.5)$; a minimal element of the set of these prime ideals belongs then to $\operatorname{Ass}(\mathrm{M}'/t\mathrm{M}')$ (Bourbaki, \emph{loc.~cit.}, \textsection~1, n\textsuperscript{o}~3, cor.~1 of prop.~7), and \emph{a fortiori} to $\operatorname{Ass}(\mathrm{M}/t\mathrm{M})$.
\end{proof}

\begin{proof}[Continuation of the proof of \sref{0.16.4.6}]
This lemma established, we return to the proof of (iii).
Let $x\in\mathfrak{m}$ be an $\mathrm{M}$-regular element, so that $\operatorname{prof}(\mathrm{M}/x\mathrm{M})=n-1$ by (ii); for every $\mathfrak{p}\in\operatorname{Ass}(\mathrm{M})$, it follows from the lemma that there is a prime ideal $\mathfrak{p}'$ associated to $\mathrm{M}/x\mathrm{M}$ and containing $\mathfrak{p}+x\mathrm{A}$; but as $x$ is $\mathrm{M}$-regular, we have $x\notin\mathfrak{p}$, whence $\mathfrak{p}'\neq\mathfrak{p}$ and hence $\dim(\mathrm{A}/\mathfrak{p}')\leq\dim(\mathrm{A}/\mathfrak{p})-1$.
The induction hypothesis shows that $n-1\leq\dim(\mathrm{A}/\mathfrak{p}')$; one concludes that $n\leq\dim(\mathrm{A}/\mathfrak{p})$ for every $\mathfrak{p}\in\operatorname{Ass}(\mathrm{M})$.

(iv) One can restrict to the case $\mathrm{M}\neq 0$.
Recall that $\widehat{\mathrm{M}}=\mathrm{M}\otimes_\mathrm{A}\widehat{\mathrm{A}}$ up to canonical isomorphism $(\mathbf{0}_\mathrm{I},~7.3.3)$; as $\widehat{\mathrm{A}}$ is a flat $\mathrm{A}$-module, there are canonical isomorphisms $\Ext^i_\mathrm{A}(k,\mathrm{M})\otimes_\mathrm{A}\widehat{\mathrm{A}}\xrightarrow{\sim}\Ext^i_{\widehat{\mathrm{A}}}(k,\mathrm{M}\otimes_\mathrm{A}\widehat{\mathrm{A}})$ since $k=k\otimes_\mathrm{A}\widehat{\mathrm{A}}$ up to canonical isomorphism $(\mathbf{0}_\mathrm{III},~12.3.5)$.
The assertion (iv) then follows from \sref{0.16.4.4} and the fact that $\widehat{\mathrm{A}}$ is a faithfully flat $\mathrm{A}$-module $(\mathbf{0}_\mathrm{I},~7.3.5)$.
\end{proof}

\begin{env}[16.4.6.4]
\label{0.16.4.6.4}
One notes that if $\mathfrak{p}$ is a prime ideal of $\mathrm{A}$, one can \emph{also} have $\operatorname{prof}_{\mathrm{A}_\mathfrak{p}}(\mathrm{M}_\mathfrak{p})<\operatorname{prof}_\mathrm{A}(\mathrm{M})$ and $\operatorname{prof}_{\mathrm{A}_\mathfrak{p}}(\mathrm{M}_\mathfrak{p})>\operatorname{prof}_\mathrm{A}(\mathrm{M})$; for an example of the first case,
\oldpage[0\textsubscript{IV-1}]{35}
it suffices to take $\mathrm{A}$ a local integral ring of dimension $\geq 1$; we have $\operatorname{prof}(\mathrm{A})\geq 1$ but $\operatorname{prof}(\mathrm{K})=0$ for the field of fractions $\mathrm{K}$ of $\mathrm{A}$.
For an example of the second case, consider a Noetherian integral local ring $\mathrm{A}_0$ of dimension $\geq 2$, with maximal ideal $\mathfrak{m}_0$ and residue field $k=\mathrm{A}_0/\mathfrak{m}_0$; let $\mathrm{A}=\mathrm{A}_0\oplus\mathfrak{J}$ be the trivial extension of $\mathrm{A}_0$ by the $\mathrm{A}_0$-module $k$ (so that the multiplication in $\mathrm{A}$ is given by $(x,z)(x',z')=(xx',xz'+x'z)$).
It is clear that $\mathfrak{J}$ is the nilradical of $\mathrm{A}$ and $\mathrm{A}/\mathfrak{J}$ is isomorphic to $\mathrm{A}_0$, so that every prime ideal $\mathfrak{p}$ of $\mathrm{A}$ is of the form $\mathfrak{p}_0\oplus\mathfrak{J}$, where $\mathfrak{p}_0$ is a prime ideal of $\mathrm{A}_0$; in particular $\mathfrak{m}=\mathfrak{m}_0\oplus\mathfrak{J}$ is the only maximal ideal of $\mathrm{A}$.
As every element of $\mathfrak{m}_0$ annihilates $\mathfrak{J}$, we have $\operatorname{prof}(\mathrm{A})=\operatorname{prof}(\mathrm{A}_\mathfrak{m})=0$; on the contrary, if $\mathfrak{p}_0\neq\mathfrak{m}_0$, the elements of $\mathfrak{J}$ are annihilated by those of $\mathfrak{m}_0-\mathfrak{p}_0$, hence $\mathrm{A}_\mathfrak{p}$ is canonically isomorphic to $(\mathrm{A}_0)_{\mathfrak{p}_0}$, and as $\mathrm{A}_0$ is integral, $\operatorname{prof}((\mathrm{A}_0)_{\mathfrak{p}_0})\geq 1$ if $\mathfrak{p}_0\neq 0$.
\end{env}

\begin{corollary}[16.4.7]
\label{0.16.4.7}
For a local Noetherian reduced ring $\mathrm{A}$ to be of depth $0$, it is necessary and sufficient that $\mathrm{A}$ be a field.
\end{corollary}

\begin{proof}
Indeed, as the intersection of the minimal prime ideals of $\mathrm{A}$ is $0$, $\mathrm{A}$ has no embedded associated prime ideals; by \sref{0.16.4.6}, (i), if $\operatorname{prof}(\mathrm{A})=0$, the maximal ideal $\mathfrak{m}$ of $\mathrm{A}$ is also a minimal prime ideal, hence it is the only prime ideal of $\mathrm{A}$, and as $\mathrm{A}$ is reduced, $\mathfrak{m}=0$.
\end{proof}

\begin{proposition}[16.4.8]
\label{0.16.4.8}
Let $\mathrm{A}$, $\mathrm{B}$ be two local Noetherian rings, $\rho:\mathrm{A}\to\mathrm{B}$ a local homomorphism, $\mathrm{M}$ a $\mathrm{B}$-module such that $\mathrm{M}_{[\rho]}$ is an $\mathrm{A}$-module of finite type.
Then
\begin{equation}
\label{0.16.4.8.1}
\tag{16.4.8.1}
  \operatorname{prof}_\mathrm{A}(\mathrm{M}_{[\rho]})=\operatorname{prof}_\mathrm{B}(\mathrm{M}).
\end{equation}
\end{proposition}

\begin{proof}
One can restrict to the case $\mathrm{M}\neq 0$; suppose $\operatorname{prof}_\mathrm{A}(\mathrm{M}_{[\rho]})=n$.
If $(x_i)_{1\leq i\leq n}$ is a $(\mathrm{M}_{[\rho]})$-regular maximal sequence formed of elements of the maximal ideal of $\mathrm{A}$, the $\rho(x_i)$ constitute an $\mathrm{M}$-regular sequence formed of elements of the maximal ideal of $\mathrm{B}$, and we have then, setting $\mathrm{N}=\mathrm{M}/(\rho(x_1)\mathrm{M}+\ldots+\rho(x_n)\mathrm{M})$, $\operatorname{prof}_\mathrm{B}(\mathrm{N})=\operatorname{prof}_\mathrm{B}(\mathrm{M})-n$ \sref{0.16.4.6}, (ii); similarly $\mathrm{N}_{[\rho]}=\mathrm{M}_{[\rho]}/(x_1\mathrm{M}_{[\rho]}+\ldots+x_n\mathrm{M}_{[\rho]})$ and $\operatorname{prof}_\mathrm{A}(\mathrm{N}_{[\rho]})=0$; one is thus reduced to proving the proposition when $n=0$.
Let $k$ be the residue field of $\mathrm{A}$; since $k$ is a quotient of $\mathrm{A}$, the $\mathrm{A}$-module $\mathrm{P}=\operatorname{Hom}_\mathrm{A}(k,\mathrm{M}_{[\rho]})$ is a submodule of $\operatorname{Hom}_\mathrm{A}(\mathrm{A},\mathrm{M}_{[\rho]})=\mathrm{M}_{[\rho]}$, and the hypothesis implies that $\mathrm{P}\neq 0$ \sref{0.16.4.6}, (i); furthermore, $\mathrm{P}$ is an $\mathrm{A}$-module of finite type (since $\mathrm{M}_{[\rho]}$ is an $\mathrm{A}$-module of finite type), hence a finite-dimensional $k$-vector space, annihilated by the maximal ideal of $\mathrm{A}$, and finally an $\mathrm{A}$-module of \emph{finite length}.
Moreover, $\mathrm{P}$ is also a submodule of the $\mathrm{B}$-module $\mathrm{M}$, and since it is of finite length as an $\mathrm{A}$-module, it is \emph{a fortiori} of finite length as a $\mathrm{B}$-module.
As it is $\neq 0$, it contains a simple sub-$\mathrm{B}$-module, that is to say one isomorphic to the residue field of $\mathrm{B}$; one concludes then by \sref{0.16.4.6}, (i) that $\operatorname{prof}_\mathrm{B}(\mathrm{M})=0$.
\end{proof}

\begin{definition}[16.4.9]
\label{0.16.4.9}
Let $\mathrm{A}$ be a local Noetherian ring, $\mathrm{M}$ an $\mathrm{A}$-module of finite type.
If $\mathrm{M}\neq 0$, we call the \emph{codepth} of $\mathrm{M}$ and denote $\operatorname{coprof}_\mathrm{A}(\mathrm{M})$ or $\operatorname{coprof}(\mathrm{M})$, the finite integer $\dim_\mathrm{A}(\mathrm{M})-\operatorname{prof}_\mathrm{A}(\mathrm{M})\geq 0$ \sref{0.16.4.5.1}.
If $\mathrm{M}=0$, we set $\operatorname{coprof}(\mathrm{M})=0$.
We evidently have $\operatorname{coprof}(\mathrm{M}^k)=\operatorname{coprof}(\mathrm{M})$ for every $k>0$.
\end{definition}

\begin{proposition}[16.4.10]
\label{0.16.4.10}
Let $\mathrm{A}$ be a local Noetherian ring, $\mathrm{M}$ an $\mathrm{A}$-module of finite type.
\oldpage[0\textsubscript{IV-1}]{36}
\begin{enumerate}
  \item[\rm{(i)}] For every $\mathrm{M}$-regular element $x$ belonging to the maximal ideal of $\mathrm{A}$, we have $\operatorname{coprof}(\mathrm{M}/x\mathrm{M})=\operatorname{coprof}(\mathrm{M})$.
  \item[\rm{(ii)}] We have $\operatorname{coprof}_\mathrm{A}(\mathrm{M})=\operatorname{coprof}_{\widehat{\mathrm{A}}}(\widehat{\mathrm{M}})$.
\end{enumerate}
\end{proposition}

\begin{proof}
Indeed, (i) follows from \sref{0.16.3.4} and \sref{0.16.4.6}, (ii), and (ii) from \sref{0.16.2.4} and \sref{0.16.4.6}, (iv).
\end{proof}

We shall show later $(\mathbf{IV},~6.11.5)$ that for every prime ideal $\mathfrak{p}$ of $\mathrm{A}$, we have $\operatorname{coprof}_{\mathrm{A}_\mathfrak{p}}(\mathrm{M}_\mathfrak{p})\leq\operatorname{coprof}_\mathrm{A}(\mathrm{M})$.

\begin{proposition}[16.4.11]
\label{0.16.4.11}
Under the hypotheses of \sref{0.16.4.8}, we have
\begin{equation}
\label{0.16.4.11.1}
\tag{16.4.11.1}
  \operatorname{coprof}_\mathrm{A}(\mathrm{M}_{[\rho]})=\operatorname{coprof}_\mathrm{B}(\mathrm{M}).
\end{equation}
\end{proposition}

\begin{proof}
This follows from \sref{0.16.1.9} and \sref{0.16.4.8}.
\end{proof}

\subsection{Cohen-Macaulay modules}
\label{subsection:0.16.5}

\begin{definition}[16.5.1]
\label{0.16.5.1}
Let $\mathrm{A}$ be a local Noetherian ring.
An $\mathrm{A}$-module of finite type $\mathrm{M}$ is called a \emph{Cohen-Macaulay module} if $\operatorname{coprof}(\mathrm{M})=0$, that is, if $\mathrm{M}=0$ or $\dim(\mathrm{M})=\operatorname{prof}(\mathrm{M})$.
We say that $\mathrm{A}$ is a \emph{Cohen-Macaulay ring} if the $\mathrm{A}$-module $\mathrm{A}$ is a Cohen-Macaulay module.
\end{definition}

An $\mathrm{A}$-module of finite length is a Cohen-Macaulay module since it is of dimension 0 \sref{0.16.1.10}.
A local reduced ring of dimension 1 is a Cohen-Macaulay ring, since then its maximal ideal $\mathfrak{m}$ cannot be associated to $\mathrm{A}$ (Bourbaki, \emph{Alg.~comm.}, chap.~IV, \textsection~2, n\textsuperscript{o}~5, prop.~10), hence $\operatorname{prof}(\mathrm{A})=1$ by virtue of \sref{0.16.4.6}, (i).
If $\mathrm{A}$ is an \emph{integral local integrally closed} ring of dimension $\geq 2$, we have $\operatorname{prof}(\mathrm{A})\geq 2$: indeed, if $x\neq 0$ is an element of $\mathfrak{m}$, we know (Bourbaki, \emph{Alg.~comm.}, chap.~VII, \textsection~1, n\textsuperscript{o}~4, prop.~8) that the prime ideals associated to $\mathrm{A}/x\mathrm{A}$ are of height 1, hence distinct from $\mathfrak{m}$ and their union is therefore distinct from $\mathfrak{m}$ (Bourbaki, \emph{Alg.~comm.}, chap.~II, \textsection~1, n\textsuperscript{o}~1, prop.~2); there thus exist $(x,y)$ sequences that are $\mathrm{A}$-regular and formed of elements of $\mathfrak{m}$, whence our assertion.
One concludes that an integral local integrally closed ring of dimension 2 is a Cohen-Macaulay ring.

\begin{proposition}[16.5.2]
\label{0.16.5.2}
Let $\mathrm{A}$ be a local Noetherian ring, $\mathrm{M}$ an $\mathrm{A}$-module of finite type.
For $\mathrm{M}$ to be a Cohen-Macaulay module, it is necessary and sufficient that $\widehat{\mathrm{M}}$ be a $\widehat{\mathrm{A}}$-module of Cohen-Macaulay.
\end{proposition}

\begin{proof}
This follows immediately from \sref{0.16.4.10}, (ii).
\end{proof}

\begin{proposition}[16.5.3]
\label{0.16.5.3}
Under the hypotheses of \sref{0.16.4.8}, for $\mathrm{M}$ to be a Cohen-Macaulay $\mathrm{B}$-module, it is necessary and sufficient that $\mathrm{M}_{[\rho]}$ be a Cohen-Macaulay $\mathrm{A}$-module.
\end{proposition}

\begin{proof}
This follows from \sref{0.16.4.11}.
\end{proof}

\begin{proposition}[16.5.4]
\label{0.16.5.4}
Let $\mathrm{A}$ be a local Noetherian ring, $\mathrm{M}$ an $\mathrm{A}$-module of finite type $\neq 0$.
If $\mathrm{M}$ is a Cohen-Macaulay module, we have $\dim(\mathrm{A}/\mathfrak{p})=\dim(\mathrm{M})$ for every prime ideal $\mathfrak{p}\in\operatorname{Ass}(\mathrm{M})$; in particular no prime ideal associated to $\mathrm{M}$ is embedded.
\end{proposition}

\begin{proof}
This follows immediately from the inequalities \sref{0.16.4.6.2}.
\end{proof}

One can also say that $\mathrm{M}$ (or $\Supp(\mathrm{M})$) is \emph{equidimensional} \sref{0.16.1.7}.

\oldpage[0\textsubscript{IV-1}]{37}

\begin{proposition}[16.5.5]
\label{0.16.5.5}
Let $\mathrm{A}$ be a local Noetherian ring, $\mathfrak{m}$ its maximal ideal, $\mathrm{M}$ an $\mathrm{A}$-module of finite type $\neq 0$, $x$ an element of $\mathfrak{m}$.
Suppose that $\mathrm{M}$ is a Cohen-Macaulay module; then the following conditions are equivalent:
\begin{enumerate}
  \item[a)] $x$ is $\mathrm{M}$-regular;
  \item[b)] $\dim(\mathrm{M}/x\mathrm{M})=\dim(\mathrm{M})-1$;
  \item[c)] $x$ does not belong to any \emph{minimal} element of $\Supp(\mathrm{M})$.
\end{enumerate}
Furthermore, $\mathrm{M}/x\mathrm{M}$ is then a Cohen-Macaulay module.
\end{proposition}

\begin{proof}
We have seen \sref{0.16.5.4} that all the $\mathfrak{p}$ in $\operatorname{Ass}(\mathrm{M})$ are minimal in $\Supp(\mathrm{M})$ and such that $\dim(\mathrm{A}/\mathfrak{p})=\dim(\mathrm{M})$; the equivalence of a) and c) then follows from Bourbaki, \emph{Alg.~comm.}, chap.~IV, \textsection~1, n\textsuperscript{o}~1, cor.~2 of prop.~2; the equivalence of b) and c) follows from \sref{0.16.3.4}.
The fact that $\mathrm{M}/x\mathrm{M}$ is a Cohen-Macaulay module follows finally from \sref{0.16.4.10}, (i).
\end{proof}

\begin{corollary}[16.5.6]
\label{0.16.5.6}
Under the hypotheses of \sref{0.16.5.5}, let $(x_i)_{1\leq i\leq r}$ be a sequence of elements of $\mathfrak{m}$.
The following conditions are equivalent:
\begin{enumerate}
  \item[a)] The sequence $(x_i)$ is $\mathrm{M}$-regular.
  \item[b)] $\dim(\mathrm{M}/(x_1\mathrm{M}+\ldots+x_r\mathrm{M}))=\dim(\mathrm{M})-r$.
  \item[c)] The $x_i$ form part of a system of parameters for $\mathrm{M}$.
\end{enumerate}
Furthermore, when these conditions are satisfied, $\mathrm{N}=\mathrm{M}/(x_1\mathrm{M}+\ldots+x_r\mathrm{M})$ is a Cohen-Macaulay module, hence for every prime ideal $\mathfrak{p}\in\operatorname{Ass}(\mathrm{N})$, we have $\dim(\mathrm{A}/\mathfrak{p})=\dim(\mathrm{M})-r$.
\end{corollary}

\begin{proof}
We already know that b) and c) are equivalent \sref{0.16.3.6} and that a) implies c) \sref{0.16.4.1} without any hypothesis on $\mathrm{M}$. To see that c) implies a), and to prove the last assertion of the statement, we argue by induction on $r$, the assertion being trivial for $r=0$.
Since $x_1$ forms part of a system of parameters, $\dim(\mathrm{M}/x_1\mathrm{M})=\dim(\mathrm{M})-1$ \sref{0.16.3.6}, hence $x_1$ is $\mathrm{M}$-regular and $\mathrm{M}/x_1\mathrm{M}$ is Cohen-Macaulay by \sref{0.16.5.5}.
Since the sequence $(x_i)_{2\leq i\leq r}$ forms part of a system of parameters for $\mathrm{P}=\mathrm{M}/x_1\mathrm{M}$, the induction hypothesis shows that $(x_i)_{2\leq i\leq r}$ is a $\mathrm{P}$-regular sequence and that $\mathrm{P}/(x_2\mathrm{P}+\ldots+x_r\mathrm{P})=\mathrm{N}$ is a Cohen-Macaulay module; furthermore $(x_i)_{1\leq i\leq r}$ is $\mathrm{M}$-regular.
\end{proof}

\begin{remark}[16.5.7]
\label{0.16.5.7}
It follows from \sref{0.16.5.6} that there is an identity between \emph{systems of parameters} for $\mathrm{M}$ and the maximal $\mathrm{M}$-regular sequences of elements of $\mathfrak{m}$ when $\mathrm{M}$ is a Cohen-Macaulay module.
Conversely, it follows from \sref{0.16.4.1} that if there is any $\mathrm{M}$-regular sequence that is a system of parameters for $\mathrm{M}$, then $\mathrm{M}$ is a Cohen-Macaulay module.
The last property stated in \sref{0.16.5.6} also \emph{characterizes} Cohen-Macaulay modules:
\end{remark}

\begin{proposition}[16.5.8]
\label{0.16.5.8}
Let $\mathrm{A}$ be a local Noetherian ring, $\mathrm{M}$ an $\mathrm{A}$-module of finite type.
Suppose that for every sequence $(x_i)_{1\leq i\leq r}$ of elements of $\mathrm{A}$ forming part of a system of parameters for $\mathrm{M}$, and for every prime ideal $\mathfrak{p}$ associated to $\mathrm{N}=\mathrm{M}/(x_1\mathrm{M}+\ldots+x_r\mathrm{M})$, we have $\dim(\mathrm{A}/\mathfrak{p})=\dim(\mathrm{N})$.
Then $\mathrm{M}$ is a Cohen-Macaulay module.
\end{proposition}

\begin{proof}
Let $(y_i)_{1\leq i\leq n}$ be a system of parameters for $\mathrm{M}$; it suffices to show that the sequence $(y_i)$ is $\mathrm{M}$-regular.
We argue by induction on $n=\dim(\mathrm{M})$, the proposition being trivial for $n=0$.
By hypothesis (applied with $r=0$) the prime ideals $\mathfrak{p}_j$ associated to $\mathrm{M}$ are all such that $\dim(\mathrm{A}/\mathfrak{p}_j)=\dim(\mathrm{M})$; as $y_1$ forms part of a system of parameters
\oldpage[0\textsubscript{IV-1}]{38}
for $\mathrm{M}$, it follows from \sref{0.16.3.7} and \sref{0.16.3.4} that $y_1$ does not belong to any of the $\mathfrak{p}_j$, hence is $\mathrm{M}$-regular.
If we set $\mathrm{P}=\mathrm{M}/y_1\mathrm{M}$, then $(y_i)_{2\leq i\leq n}$ is a system of parameters for $\mathrm{P}$, and it is immediate that $\mathrm{P}$ satisfies the hypothesis of the statement; applying the induction hypothesis, one sees that $(y_i)_{2\leq i\leq n}$ is a $\mathrm{P}$-regular sequence, which proves the proposition.
\end{proof}

\begin{proposition}[16.5.9]
\label{0.16.5.9}
Let $\mathrm{A}$ be a local Noetherian ring, $\mathrm{M}$ an $\mathrm{A}$-module of finite type.
Suppose that $\mathrm{M}$ is a Cohen-Macaulay module.
Let $\mathfrak{p}\in\Supp(\mathrm{M})$, and set $r=\dim(\mathrm{M})-\dim(\mathrm{M}/\mathfrak{p}\mathrm{M})$.
Then there exists an $\mathrm{M}$-regular sequence $(x_i)_{1\leq i\leq r}$, formed of elements of $\mathfrak{p}$; for every such sequence, we have
\[
  \dim(\mathrm{M}/(x_1\mathrm{M}+\ldots+x_r\mathrm{M}))=\dim(\mathrm{M}/\mathfrak{p}\mathrm{M})=\dim(\mathrm{A}/\mathfrak{p}),
\]
and $\mathfrak{p}$ is a \emph{minimal} element of $\operatorname{Ass}(\mathrm{M}/(x_1\mathrm{M}+\ldots+x_r\mathrm{M}))$.
\end{proposition}

\begin{proof}
Let us first prove the existence of the $x_i$; there is nothing to prove for $r=0$.
For $r>0$, we proceed by induction.
As for every prime ideal $\mathfrak{q}_i$ associated to $\mathrm{M}$ we have $\dim(\mathrm{A}/\mathfrak{q}_i)=\dim(\mathrm{M})$ \sref{0.16.5.4}, the hypothesis $\dim(\mathrm{M}/\mathfrak{p}\mathrm{M})<\dim(\mathrm{M})$ implies that $\mathfrak{p}$ cannot be contained in any of the $\mathfrak{q}_i$, the latter being the minimal prime ideals among those containing $\operatorname{Ann}(\mathrm{M})$ \sref{0.16.1.2.2} and \sref{0.16.1.7}; one concludes that $\mathfrak{p}$ is not contained in the union of the $\mathfrak{q}_i$ (Bourbaki, \emph{Alg.~comm.}, chap.~II, \textsection~1, n\textsuperscript{o}~1, prop.~2), and hence there exists $x_1\in\mathfrak{p}$ that is $\mathrm{M}$-regular (Bourbaki, \emph{Alg.~comm.}, chap.~IV, \textsection~1, n\textsuperscript{o}~1, cor.~2 of prop.~2).
It follows from \sref{0.16.5.5} that $\mathrm{N}=\mathrm{M}/x_1\mathrm{M}$ is a Cohen-Macaulay module, of dimension $\dim(\mathrm{M})-1$; as $x_1\in\mathfrak{p}$, moreover $\mathrm{N}/\mathfrak{p}\mathrm{N}=\mathrm{M}/\mathfrak{p}\mathrm{M}$, whence $\dim(\mathrm{N})-\dim(\mathrm{N}/\mathfrak{p}\mathrm{N})=r-1$; one can therefore apply the induction hypothesis to $\mathrm{N}$, and if $(x_i)_{2\leq i\leq r}$ is an $\mathrm{N}$-regular sequence formed of elements of $\mathfrak{p}$, it is clear that $(x_i)_{1\leq i\leq r}$ is an $\mathrm{M}$-regular sequence formed of elements of $\mathfrak{p}$.

Now if one sets $\mathrm{P}=\mathrm{M}/(x_1\mathrm{M}+\ldots+x_r\mathrm{M})$, we have $\mathrm{P}/\mathfrak{p}\mathrm{P}=\mathrm{M}/\mathfrak{p}\mathrm{M}$ since the $x_i$ are in $\mathfrak{p}$, and $\dim(\mathrm{P})=\dim(\mathrm{P}/\mathfrak{p}\mathrm{P})$ by \sref{0.16.5.6}.
But since $\mathrm{P}$ is a Cohen-Macaulay module, we also have $\dim(\mathrm{P})=\dim(\mathrm{A}/\mathfrak{p}')$ for every $\mathfrak{p}'\in\operatorname{Ass}(\mathrm{P})$ \sref{0.16.5.4}.
As $\mathfrak{p}\in\Supp(\mathrm{P})$ $(\mathbf{0}_\mathrm{I},~1.7.5)$, $\mathfrak{p}$ contains some $\mathfrak{p}'\in\operatorname{Ass}(\mathrm{P})$ (Bourbaki, \emph{Alg.~comm.}, chap.~IV, \textsection~1, n\textsuperscript{o}~4, th.~2), and as $\dim(\mathrm{P})=\dim(\mathrm{P}/\mathfrak{p}\mathrm{P})\leq\dim(\mathrm{A}/\mathfrak{p})\leq\dim(\mathrm{A}/\mathfrak{p}')=\dim(\mathrm{P})$, one necessarily has $\mathfrak{p}=\mathfrak{p}'$ \sref{0.16.1.2.2}, which completes the proof.
\end{proof}

\begin{corollary}[16.5.10]
\label{0.16.5.10}
Under the same hypotheses as in \sref{0.16.5.9}:
\begin{enumerate}
  \item[\rm{(i)}] $\mathrm{M}_\mathfrak{p}$ is an $\mathrm{A}_\mathfrak{p}$-module of Cohen-Macaulay.
  \item[\rm{(ii)}] We have
  \begin{equation}
  \label{0.16.5.10.1}
  \tag{16.5.10.1}
    \dim(\mathrm{M})=\dim_\mathrm{A}(\mathrm{M}/\mathfrak{p}\mathrm{M})+\dim_{\mathrm{A}_\mathfrak{p}}(\mathrm{M}_\mathfrak{p}).
  \end{equation}
\end{enumerate}
\end{corollary}

\begin{proof}
With the notations of \sref{0.16.5.9}, let $x_i'$ be the canonical images of the $x_i$ in $\mathrm{A}_\mathfrak{p}$ $(1\leq i\leq r)$; as the $x_i'$ belong to the maximal ideal of $\mathrm{A}_\mathfrak{p}$ and form a $\mathrm{M}_\mathfrak{p}$-regular sequence by flatness \sref{0.15.1.14}, we have
\[
  \operatorname{prof}_{\mathrm{A}_\mathfrak{p}}(\mathrm{M}_\mathfrak{p})\geq r=\dim(\mathrm{M})-\dim(\mathrm{M}/\mathfrak{p}\mathrm{M})\geq\dim_{\mathrm{A}_\mathfrak{p}}(\mathrm{M}_\mathfrak{p})
\]
(by virtue of \sref{0.16.1.8.1}).
But taking into account \sref{0.16.4.5.1}, the three terms of this inequality are necessarily equal, whence the corollary.
\end{proof}

\oldpage[0\textsubscript{IV-1}]{39}

\begin{corollary}[16.5.11]
\label{0.16.5.11}
Let $\mathrm{A}$ be a local Noetherian ring; suppose that there exists a Cohen-Macaulay $\mathrm{A}$-module of finite type $\mathrm{M}$ with support $\Spec(\mathrm{A})$.
Then, for every prime ideal $\mathfrak{p}$ of $\mathrm{A}$, we have
\begin{equation}
\label{0.16.5.11.1}
\tag{16.5.11.1}
  \dim(\mathrm{A})=\dim(\mathrm{A}/\mathfrak{p})+\dim(\mathrm{A}_\mathfrak{p}).
\end{equation}
\end{corollary}

\begin{proof}
Indeed, $\Supp(\mathrm{M}/\mathfrak{p}\mathrm{M})=\Spec(\mathrm{A}/\mathfrak{p})$ $(\mathbf{0}_\mathrm{I},~1.7.5)$ and $\Supp(\mathrm{M}_\mathfrak{p})=\Spec(\mathrm{A}_\mathfrak{p})$; the relation \sref{0.16.5.11.1} is thus a special case of \sref{0.16.5.10.1}.
\end{proof}

\begin{corollary}[16.5.12]
\label{0.16.5.12}
Every quotient ring of a local Noetherian ring $\mathrm{A}$ satisfying the conditions of \sref{0.16.5.11} (in particular every quotient ring of a local Cohen-Macaulay ring) is \emph{catenary}.
\end{corollary}

\begin{proof}
It suffices evidently to prove that $\mathrm{A}$ itself is catenary, that is to say for two prime ideals $\mathfrak{p}\subset\mathfrak{q}$ of $\mathrm{A}$, we have \sref{0.16.1.4.2}
\[
  \dim(\mathrm{A}_\mathfrak{q})=\dim(\mathrm{A}_\mathfrak{q}/\mathfrak{p}\mathrm{A}_\mathfrak{q})+\dim(\mathrm{A}_\mathfrak{p}).
\]
But, by virtue of \sref{0.16.5.10}, (i) $\mathrm{A}_\mathfrak{q}$ satisfies the same hypotheses as $\mathrm{A}$, and the preceding relation is then none other than \sref{0.16.5.11.1} applied to $\mathrm{A}_\mathfrak{q}$ and to the prime ideal $\mathfrak{p}\mathrm{A}_\mathfrak{q}$ of $\mathrm{A}_\mathfrak{q}$.
\end{proof}

\begin{env}[16.5.13]
\label{0.16.5.13}
Let $\mathrm{A}$ be a Noetherian ring, $\mathrm{M}$ an $\mathrm{A}$-module of finite type.
We say that $\mathrm{M}$ is a \emph{Cohen-Macaulay $\mathrm{A}$-module} if, for every prime ideal $\mathfrak{p}$ of $\mathrm{A}$, $\mathrm{M}_\mathfrak{p}$ is an $\mathrm{A}_\mathfrak{p}$-module of Cohen-Macaulay; by virtue of \sref{0.16.5.10} this definition coincides with \sref{0.16.5.1} when $\mathrm{A}$ is \emph{local}.
We say that $\mathrm{A}$ is a \emph{Cohen-Macaulay ring} if all the $\mathrm{A}_\mathfrak{p}$ are.
It is clear that if $\mathrm{M}$ (resp.~$\mathrm{A}$) is a Cohen-Macaulay $\mathrm{A}$-module (resp.~a Cohen-Macaulay ring), $\mathrm{S}^{-1}\mathrm{M}$ is a Cohen-Macaulay $\mathrm{S}^{-1}\mathrm{A}$-module (resp.~$\mathrm{S}^{-1}\mathrm{A}$ is a Cohen-Macaulay ring) for every multiplicative subset $\mathrm{S}$ of $\mathrm{A}$.
\end{env}
% END NATIVE INPUT ega0/ega0-16
% BEGIN NATIVE INPUT ega0/ega0-17
\section{Regular rings}
\label{section:0.17}

\subsection{Definition of regular rings}
\label{subsection:0.17.1}

\oldpage[0\textsubscript{IV-1}]{39}
\begin{proposition}[17.1.1]
\label{0.17.1.1}
Let $A$ be a Noetherian local ring of dimension $n$, $\mathfrak{m}$ its maximal ideal, $k=A/\mathfrak{m}$ its residue field.
The following conditions are equivalent:
\begin{enumerate}
  \item[{\rm a)}] The canonical surjective homomorphism of graded $k$-modules \sref{0.15.1.1}
  \[
    \varphi:\mathbf{S}^*_k(\mathfrak{m}/\mathfrak{m}^2)\to\operatorname{gr}^*_\mathfrak{m}(A)
  \]
  (where the second term is the graded $k$-module associated to $A$ equipped with the $\mathfrak{m}$-preadic filtration) is bijective.
  \item[{\rm b)}] We have $\operatorname{rg}_k(\mathfrak{m}/\mathfrak{m}^2)=n$.
  \item[{\rm c)}] The ideal $\mathfrak{m}$ admits a system of generators of $n$ elements.
  \item[{\rm d)}] The ideal $\mathfrak{m}$ admits a system of generators that is an $A$-regular sequence.
\end{enumerate}
\end{proposition}

\begin{proof}
By Nakayama's lemma, $\operatorname{rg}_k(\mathfrak{m}/\mathfrak{m}^2)$ is the smallest number of elements of a system of generators of $\mathfrak{m}$, so b) and c) are equivalent.
On the other hand, if $(x_i)_{1\leq i\leq r}$ is a system of generators of $\mathfrak{m}$ whose classes $\overline{x}_i$ mod.~$\mathfrak{m}^2$ form a basis of the $k$-vector space $\mathfrak{m}/\mathfrak{m}^2$, $\mathbf{S}^*(\mathfrak{m}/\mathfrak{m}^2)$ is isomorphic to the polynomial ring $k[\mathrm{T}_1,\ldots,\mathrm{T}_r]$;
\oldpage[0\textsubscript{IV-1}]{40}
taking into account that every $A$-module of finite type is separated for the $\mathfrak{m}$-preadic filtration \sref[0\textsubscript{I}]{0.7.3.5}, it follows from \sref{0.15.1.9} that conditions a) and d) are equivalent;
moreover, since every $A$-regular sequence of elements of $\mathfrak{m}$ has at most $n$ elements \sref{0.16.4.1}, we see that d) implies c).
It remains to prove that c) implies a).

Let us set, for brevity, $\mathrm{S}=\mathbf{S}^*(\mathfrak{m}/\mathfrak{m}^2)=k[\mathrm{T}_1,\ldots,\mathrm{T}_n]$, $\mathrm{G}=\operatorname{gr}^*_\mathfrak{m}(A)$, and consider the exact sequence $0\to\mathfrak{J}\to\mathrm{S}\xrightarrow{\varphi}\mathrm{G}\to 0$, where the kernel $\mathfrak{J}$ of $\varphi$ is thus a graded ideal of $\mathrm{S}$.
For every integer $s>0$, we thus have
\[
\label{0.17.1.1.1}
  \binom{s+n-1}{n-1}=\operatorname{long}(\mathrm{S}_s)=\operatorname{long}(\mathrm{G}_s)+\operatorname{long}(\mathfrak{J}_s).
  \tag{17.1.1.1}
\]

Suppose $\mathfrak{J}\neq 0$, so that there exists a homogeneous element $u\in\mathfrak{J}$ of degree $h>0$;
$\mathfrak{J}_s$ contains $u\mathrm{S}_{s-h}$ for $s\geq h$, and since $u\neq 0$ and $\mathrm{S}$ is integral, $u\mathrm{S}_{s-h}$ is isomorphic to $\mathrm{S}_{s-h}$ as a $k$-vector space;
we would thus have $\operatorname{long}(\mathfrak{J}_s)\geq\operatorname{long}(\mathrm{S}_{s-h})=\binom{s-h+n-1}{n-1}$, whence for $s\geq h$,
\[
\label{0.17.1.1.2}
  \operatorname{long}(\mathrm{G}_s)\leq\binom{s+n-1}{n-1}-\binom{s-h+n-1}{n-1}.
  \tag{17.1.1.2}
\]
But the second term of \sref{0.17.1.1.2} is a polynomial in $s$ of degree $\leq n-2$;
since $\mathrm{G}_s=\mathfrak{m}^s/\mathfrak{m}^{s+1}$, and for $s$ sufficiently large, $\operatorname{long}(\mathrm{G}_s)$ is a polynomial in $s$ of degree exactly equal to $n-1$ whose leading coefficient is $>0$ \sref{0.16.2.1}, which contradicts the inequality \sref{0.17.1.1.2}.
\end{proof}

\begin{definition}[17.1.2]
\label{0.17.1.2}
We say that a Noetherian local ring that satisfies the equivalent conditions of \sref{0.17.1.1} is \emph{regular}.
\end{definition}

\begin{corollary}[17.1.3]
\label{0.17.1.3}
A regular local ring is integral, integrally closed, and a Cohen-Macaulay ring.
\end{corollary}

\begin{proof}
The last assertion follows from \sref{0.17.1.1}[d)];
on the other hand, if $A$ is regular, $\operatorname{gr}^*_\mathfrak{m}(A)$ is integral and completely integrally closed, being isomorphic to $k[\mathrm{T}_1,\ldots,\mathrm{T}_n]$;
we conclude that $A$ also possesses these two properties (Bourbaki, \emph{Alg.~comm.}, chap.~III, \textsection2, n\textsuperscript{o}~3, cor.\ of prop.~1 and chap.~V, \textsection1, n\textsuperscript{o}~5, prop.~15).
\end{proof}

\begin{example}[17.1.4]
\label{0.17.1.4}
\begin{enumerate}
  \item[{\rm(i)}] A regular local ring of dimension $0$, being integral by \sref{0.17.1.3}, is necessarily a \emph{field}, and conversely.
  \item[{\rm(ii)}] For a Noetherian local ring of dimension $1$ to be regular, it is necessary and sufficient that it be a \emph{discrete valuation ring}:
  indeed, to say that a Noetherian local ring of dimension $1$ is regular means that its maximal ideal is \emph{principal}, and the conclusion follows from Bourbaki, \emph{Alg.~comm.}, chap.~VI, \textsection3, n\textsuperscript{o}~6, prop.~9.
  \item[{\rm(iii)}] Let $k$ be a field; the formal power series ring $A=k[[\mathrm{T}_1,\ldots,\mathrm{T}_n]]$ is a \emph{regular ring of dimension $n$};
  indeed, it is clear that the $\mathrm{T}_i$ ($1\leq i\leq n$) generate the maximal ideal $\mathfrak{m}$ of $A$ and form an $A$-regular sequence, since $A/(\mathrm{T}_1 A+\cdots+\mathrm{T}_i A)$ is isomorphic to $k[[\mathrm{T}_{i+1},\ldots,\mathrm{T}_n]]$.
\end{enumerate}
\end{example}

\begin{proposition}[17.1.5]
\label{0.17.1.5}
For a Noetherian local ring $A$ to be regular, it is necessary and sufficient that its completion $\widehat{A}$ be regular.
\end{proposition}

\begin{proof}
\oldpage[0\textsubscript{IV-1}]{41}
Indeed, the maximal ideal of $\widehat{A}$ is $\mathfrak{m}\widehat{A}$ and we know that $\mathfrak{m}^h/\mathfrak{m}^{h+1}$ and $(\mathfrak{m}\widehat{A})^h/(\mathfrak{m}\widehat{A})^{h+1}$ are isomorphic $k$-vector spaces;
condition a) of \sref{0.17.1.1} is thus the same for $A$ and $\widehat{A}$.
\end{proof}

\begin{definition}[17.1.6]
\label{0.17.1.6}
Given a Noetherian local ring $A$, with maximal ideal $\mathfrak{m}$, we call a \emph{regular system of parameters} of $A$ a system of parameters for $A$ \sref{0.16.3.5} that generates $\mathfrak{m}$.
\end{definition}

The existence of such a system is equivalent to the fact that $A$ is regular \sref{0.17.1.1};
if $(x_i)_{1\leq i\leq n}$ is a regular system of parameters, it is a minimal system of generators of $\mathfrak{m}$, whose classes mod.~$\mathfrak{m}^2$ form a basis of $\mathfrak{m}/\mathfrak{m}^2$ over $k$;
by virtue of \sref{0.17.1.1}[a)] and \sref{0.15.1.9}, $(x_i)$ is a maximal $A$-regular sequence (whence the terminology).

One should note however that, in a regular ring, a system of parameters for $A$ that is also an $A$-regular sequence is \emph{not} necessarily a regular system of parameters, as shown by the example of $k$-th powers of a regular system of parameters \sref{0.15.1.20}.

\begin{proposition}[17.1.7]
\label{0.17.1.7}
Let $A$ be a Noetherian local ring, $\mathfrak{m}$ its maximal ideal, $k=A/\mathfrak{m}$ its residue field, $(x_i)_{1\leq i\leq r}$ a sequence of elements of $\mathfrak{m}$, $\mathfrak{J}$ the ideal $x_1 A+\cdots+x_r A$.
The following conditions are equivalent:
\begin{enumerate}
  \item[{\rm a)}] $A$ is regular and the $x_i$ are part of a regular system of parameters of $A$.
  \item[{\rm a')}] $A$ is regular and the classes $\overline{x}_i$ of the $x_i$ mod.~$\mathfrak{m}^2$ are linearly independent over $k$.
  \item[{\rm b)}] The $x_i$ are part of a system of parameters for $A$ and $A/\mathfrak{J}$ is a regular ring.
\end{enumerate}
Moreover, when these conditions are satisfied, $\mathfrak{J}$ is a prime ideal.
\end{proposition}

\begin{proof}
The last assertion follows trivially from the fact that $A/\mathfrak{J}$, being regular, is integral \sref{0.17.1.3}.
Set $n=\dim(A)$.
The equivalence of a) and a') is immediate;
indeed, the classes mod.~$\mathfrak{m}^2$ of the elements of a regular system of parameters form a basis of $\mathfrak{m}/\mathfrak{m}^2$ over $k$, and conversely, if the $\overline{x}_i$ ($1\leq i\leq r$) are linearly independent over $k$, one can find $n-r$ elements $x_i\in\mathfrak{m}$ ($r+1\leq i\leq n$) such that the $\overline{x}_i$ ($1\leq i\leq n$) form a basis of $\mathfrak{m}/\mathfrak{m}^2$, so $(x_i)_{1\leq i\leq n}$ is a regular system of parameters.
To see that a) implies b), note that since the $x_i$ are part of a system of parameters of $A$, we have $\dim(A/\mathfrak{J})=n-r$ \sref{0.16.3.6};
let $\mathfrak{n}=\mathfrak{m}/\mathfrak{J}$ be the ideal of $A/\mathfrak{J}$.
We have an exact sequence
\[
\label{0.17.1.7.1}
  0\to(\mathfrak{m}^2+\mathfrak{J})/\mathfrak{m}^2\to\mathfrak{m}/\mathfrak{m}^2\to\mathfrak{n}/\mathfrak{n}^2\to 0
  \tag{17.1.7.1}
\]
since $\mathfrak{n}^2=(\mathfrak{m}^2+\mathfrak{J})/\mathfrak{J}$;
hypothesis a) implies that $\operatorname{rg}_k((\mathfrak{m}^2+\mathfrak{J})/\mathfrak{m}^2)=r$, so $\operatorname{rg}_k(\mathfrak{n}/\mathfrak{n}^2)=n-r=\dim(A/\mathfrak{J})$, which proves b) by virtue of \sref{0.17.1.1}[b)].

Conversely, to prove that b) implies a), note that the fact that the $x_i$ are part of a system of parameters implies $\dim(A/\mathfrak{J})=n-r$ \sref{0.16.3.6};
the hypothesis that $A/\mathfrak{J}$ is regular implies on the other hand $\operatorname{rg}_k(\mathfrak{n}/\mathfrak{n}^2)=n-r$ \sref{0.17.1.1};
moreover we evidently have $\operatorname{rg}_k((\mathfrak{m}^2+\mathfrak{J})/\mathfrak{m}^2)\leq r$, so from \sref{0.17.1.7.1} we obtain $\operatorname{rg}_k(\mathfrak{m}/\mathfrak{m}^2)\leq r+(n-r)=n$;
thus $A$ is regular by \sref{0.16.2.6} and \sref{0.17.1.1}.
Furthermore the relation $\operatorname{rg}_k(\mathfrak{m}/\mathfrak{m}^2)=n$ then implies $\operatorname{rg}_k((\mathfrak{m}^2+\mathfrak{J})/\mathfrak{m}^2)=r$, and consequently the $\overline{x}_i$ are linearly independent over $k$.
\end{proof}

\oldpage[0\textsubscript{IV-1}]{42}
\begin{corollary}[17.1.8]
\label{0.17.1.8}
Let $A$ be a Noetherian local ring, $\mathfrak{m}$ its maximal ideal, $t$ an element of $\mathfrak{m}$.
The following conditions are equivalent:
\begin{enumerate}
  \item[{\rm a)}] $A/tA$ is regular and $t$ is a non-zero-divisor in $A$.
  \item[{\rm b)}] $A$ is regular and $t\notin\mathfrak{m}^2$.
\end{enumerate}
\end{corollary}

\begin{proof}
This is the special case $r=1$ of \sref{0.17.1.7} (taking into account that an element of $\mathfrak{m}$ that is a non-zero-divisor is part of a system of parameters \sref{0.16.4.1}).
\end{proof}

\begin{corollary}[17.1.9]
\label{0.17.1.9}
Let $A$ be a regular local ring, $\mathfrak{m}$ its maximal ideal, $\mathfrak{J}$ an ideal of $A$ contained in $\mathfrak{m}$.
The following conditions are equivalent:
\begin{enumerate}
  \item[{\rm a)}] The ring $A/\mathfrak{J}$ is regular.
  \item[{\rm b)}] $\mathfrak{J}$ is generated by a sequence $(x_i)_{1\leq i\leq r}$ forming part of a regular system of parameters of $A$.
\end{enumerate}
\end{corollary}

\begin{proof}
We have already seen \sref{0.17.1.7} that b) implies a).
Conversely, suppose that $A/\mathfrak{J}$ is regular (which implies that $\mathfrak{J}$ is prime) and let $n=\dim(A)$, $n-r=\dim(A/\mathfrak{J})$;
with the notation of \sref{0.17.1.7}, the exact sequence \sref{0.17.1.7.1} gives $\operatorname{rg}_k((\mathfrak{m}^2+\mathfrak{J})/\mathfrak{m}^2)=r$, so there exists a sequence $(x_i)_{1\leq i\leq r}$ of elements of $\mathfrak{J}$ forming part of a regular system of parameters of $A$.
Set $\mathfrak{J}'=x_1 A+\cdots+x_r A$;
it follows from \sref{0.17.1.7} that $\mathfrak{J}'$ is a prime ideal of $A$ and that $A/\mathfrak{J}'$ is regular and of dimension $n-r$;
but since $A/\mathfrak{J}=(A/\mathfrak{J}')/(\mathfrak{J}/\mathfrak{J}')$ and $\mathfrak{J}/\mathfrak{J}'$ is prime in the integral ring $A/\mathfrak{J}'$, the dimensions of $A/\mathfrak{J}$ and $A/\mathfrak{J}'$ can only be equal if $\mathfrak{J}=\mathfrak{J}'$ \sref{0.16.1.2.2}.
\end{proof}

\subsection{Reminders on projective and injective dimension of modules}
\label{subsection:0.17.2}

\begin{env}[17.2.1]
\label{0.17.2.1}
Let $A$ be a ring, $M$ an $A$-module.
Recall (M, VI, 2) that we call \emph{projective dimension} (resp. \emph{injective dimension}) of $M$ and denote $\operatorname{dim.proj}(M)$ or $\operatorname{dim.proj}_A(M)$ (resp. $\operatorname{dim.inj}(M)$ or $\operatorname{dim.inj}_A(M)$) the smallest $n$ (integer or equal to $+\infty$) such that there exists a left projective resolution (resp. a right injective resolution) of $M$ of length $n$.
It is equivalent (loc.\ cit.) to say that $n$ is the smallest number such that for every $A$-module $N$, we have $\operatorname{Ext}^i_A(M,N)=0$ (resp. $\operatorname{Ext}^i_A(N,M)=0$) for all $i>n$, or only for $i=n+1$.
For every direct factor $M'$ of $M$, we thus have $\operatorname{dim.proj}(M')\leq\operatorname{dim.proj}(M)$ since $\operatorname{Ext}^i_A(M',N)$ is a direct factor of $\operatorname{Ext}^i_A(M,N)$;
similarly $\operatorname{dim.inj}(M')\leq\operatorname{dim.inj}(M)$.
A condition equivalent to $\operatorname{dim.proj}(M)=n$ (resp. $\operatorname{dim.inj}(M)=n$) is that $n$ is the smallest number such that, for every exact sequence
\[
  0\to\mathrm{R}\to\mathrm{P}_{n-1}\to\cdots\to\mathrm{P}_0\to M\to 0
\]
(resp. $0\to M\to\mathrm{Q}_0\to\cdots\to\mathrm{Q}_{n-1}\to\mathrm{C}\to 0$), where all the $\mathrm{P}_i$ for $0\leq i<n$ are projective (resp. all the $\mathrm{Q}_i$ for $0\leq i<n$ are injective), $\mathrm{R}$ is projective (resp. $\mathrm{C}$ is injective).
\end{env}

\begin{remark}[17.2.2]
\label{0.17.2.2}
\begin{enumerate}
  \item[{\rm(i)}] To say that $M$ is a projective (resp. injective) $A$-module is therefore equivalent to saying that $\operatorname{dim.proj}(M)=0$ (resp. $\operatorname{dim.inj}(M)=0$).
  \item[{\rm(ii)}] Let $\mathrm{T}$ be a covariant additive functor from the category of $A$-modules to an abelian category.
  It follows immediately from the definitions of \sref{0.17.2.1} and those of derived functors that if $\operatorname{dim.proj}(M)\leq n$ (resp. $\operatorname{dim.inj}(M)\leq n$), we have $\mathrm{L}^i\mathrm{T}(M)=0$ (resp. $\mathrm{R}^i\mathrm{T}(M)=0$) for all $i>n$.
\oldpage[0\textsubscript{IV-1}]{43}
  \item[{\rm(iii)}] If we suppose $A$ Noetherian and $M$ of finite type, the last interpretation of the projective dimension given in \sref{0.17.2.1} shows that if $\operatorname{dim.proj}(M)=n$, $M$ admits a left resolution of length $n$ formed of projective modules \emph{of finite type}.
  If moreover $A$ is a Noetherian \emph{local} ring, $M$ admits a resolution of length $n$ by \emph{free} modules of finite type, a projective module of finite type then being free (Bourbaki, \emph{Alg.~comm.}, chap.~II, \textsection3, n\textsuperscript{o}~2, cor.~2 of prop.~5).
\end{enumerate}
\end{remark}

\begin{lemma}[17.2.3]
\label{0.17.2.3}
Let $A$ be a ring, $M$ an $A$-module.
For $\operatorname{dim.inj}(M)\leq n$, it is necessary and sufficient that for every monogenic $A$-module $N$, we have $\operatorname{Ext}^{n+1}_A(N,M)=0$.
\end{lemma}

\begin{proof}
With the notation of \sref{0.17.2.1}, it suffices to prove that $\mathrm{C}$ is injective;
for every $A$-module $N$, $\operatorname{Ext}^{n+1}_A(N,M)$ is isomorphic to $\operatorname{Ext}^1_A(N,\mathrm{C})$ (M, V, 7), so we have $\operatorname{Ext}^1_A(N,\mathrm{C})=0$ for every module $N$ of the form $A/\mathfrak{J}$, where $\mathfrak{J}$ is any ideal of $A$.
The exact sequence of Ext then shows that the canonical homomorphism $\Hom(A,\mathrm{C})\to\Hom(\mathfrak{J},\mathrm{C})$ is surjective for every ideal $\mathfrak{J}$ of $A$, which implies that $\mathrm{C}$ is an injective $A$-module (M, I, 3.2).
\end{proof}

\begin{lemma}[17.2.4]
\label{0.17.2.4}
Let $A$ be a Noetherian local ring, $M$ an $A$-module of finite type.
For $\operatorname{dim.proj}(M)\leq n$, it is necessary and sufficient that for every monogenic $A$-module $N$, we have $\operatorname{Ext}^{n+1}_A(M,N)=0$.
\end{lemma}

\begin{proof}
We know in fact (M, VI, 2.5) that the condition $\operatorname{dim.proj}(M)\leq n$ is equivalent to $\operatorname{Ext}^{n+1}_A(M,N)=0$ for every $A$-module $N$ of finite type;
to see that the condition of the statement is also sufficient, we argue by induction on the number $m$ of generators of $N$:
there is a submodule $\mathrm{N}_1$ of $N$ generated by $m-1$ elements and such that $\mathrm{N}_2=N/\mathrm{N}_1$ is monogenic;
from the exact sequence $0\to\mathrm{N}_1\to N\to\mathrm{N}_2\to 0$, we then deduce the exact sequence $\operatorname{Ext}^{n+1}_A(M,\mathrm{N}_1)\to\operatorname{Ext}^{n+1}_A(M,N)\to\operatorname{Ext}^{n+1}_A(M,\mathrm{N}_2)$ and the induction hypothesis shows that the condition of the statement indeed implies $\operatorname{Ext}^{n+1}_A(M,N)=0$.
\end{proof}

\begin{corollary}[17.2.5]
\label{0.17.2.5}
Let $A$ be a Noetherian ring, $M$ an $A$-module.
We have
\[
\label{0.17.2.5.1}
  \operatorname{dim.inj}_A(M)=\sup_\mathfrak{m}\big(\operatorname{dim.inj}_{A_\mathfrak{m}}(M_\mathfrak{m})\big)
  \tag{17.2.5.1}
\]
and if $M$ is of finite type
\[
\label{0.17.2.5.2}
  \operatorname{dim.proj}_A(M)=\sup_\mathfrak{m}\big(\operatorname{dim.proj}_{A_\mathfrak{m}}(M_\mathfrak{m})\big)
  \tag{17.2.5.2}
\]
where $\mathfrak{m}$ ranges over the set of prime ideals (or the set of maximal ideals) of $A$.
\end{corollary}

\begin{proof}
\oldpage[0\textsubscript{IV-1}]{44}
Indeed, if $N$ is an $A$-module of finite type (thus of finite presentation), we have, for every multiplicative subset $S$ of $A$ and for all $i\geq 0$
\[
  S^{-1}\operatorname{Ext}^i_A(N,M)=\operatorname{Ext}^i_{S^{-1}\!A}(S^{-1}N,S^{-1}M)
\]
by flatness, considering a free resolution of $N$ and using the fact that the preceding relation is true for $i=0$ (Bourbaki, \emph{Alg.~comm.}, chap.~II, \textsection2, n\textsuperscript{o}~7, prop.~19).
In particular $\operatorname{Ext}^i_{A_\mathfrak{m}}(A_\mathfrak{m}/\mathfrak{J}A_\mathfrak{m},M_\mathfrak{m})=(\operatorname{Ext}^i_A(A/\mathfrak{J},M))_\mathfrak{m}$ for every prime ideal $\mathfrak{m}$ and every ideal $\mathfrak{J}$ of $A$;
taking into account \sref{0.17.2.3} and the fact that every ideal of $A_\mathfrak{m}$ is of the form $\mathfrak{J}A_\mathfrak{m}$ for a suitable ideal $\mathfrak{J}$ of $A$, we deduce formula \sref{0.17.2.5.1} from what precedes and from Bourbaki, \emph{Alg.~comm.}, chap.~II, \textsection3, n\textsuperscript{o}~3, th.~1.
We proceed similarly for \sref{0.17.2.5.2} using this time \sref{0.17.2.4} and interchanging the roles of $M$ and $N$.
\end{proof}

For Noetherian rings and modules of finite type, the study of projective or injective dimension is thus reduced by \sref{0.17.2.5} to the case of local rings.
We then have the

\begin{lemma}[17.2.6]
\label{0.17.2.6}
Let $A$ be a Noetherian local ring, $k$ its residue field, $M$ an $A$-module of finite type.
For $\operatorname{dim.proj}(M)\leq n$, it is necessary and sufficient that $\operatorname{Tor}^A_i(M,k)=0$ for $i>n$, and it suffices that $\operatorname{Tor}^A_{n+1}(M,k)=0$.
\end{lemma}

\begin{proof}
The necessity is a special case of the remark \sref{0.17.2.2}[(ii)], applied to the covariant functor $M\mapsto k\otimes_A M$.
To prove the sufficiency, it suffices, with the notation of \sref{0.17.2.1}, to establish that $\mathrm{R}$ is projective when the $\mathrm{P}_i$ are assumed of finite type;
now $\operatorname{Tor}^A_{n+1}(M,k)$ is isomorphic to $\operatorname{Tor}^A_1(\mathrm{R},k)$ (M, V, 7);
and we know that, since $\mathrm{R}$ is of finite type, the condition $\operatorname{Tor}^A_1(\mathrm{R},k)=0$ implies that $\mathrm{R}$ is free (Bourbaki, \emph{Alg.~comm.}, chap.~II, \textsection3, n\textsuperscript{o}~2, cor.~2 of prop.~5).
\end{proof}

\begin{corollary}[17.2.7]
\label{0.17.2.7}
Under the hypotheses of \sref{0.17.2.6}, let $x$ be an $M$-regular element of $A$ belonging to the maximal ideal of $A$; then we have
\[
\label{0.17.2.7.1}
  \operatorname{dim.proj}(M/xM)=\operatorname{dim.proj}(M)+1.
  \tag{17.2.7.1}
\]
\end{corollary}

\begin{proof}
Indeed, from the exact sequence $0\to M\xrightarrow{x}M\to M/xM\to 0$ defined by multiplication by $x$ in $M$, we deduce by the exact sequence of Tor:
\[
  \operatorname{Tor}^A_i(M,k)\xrightarrow{x}\operatorname{Tor}^A_i(M/xM,k)\to\operatorname{Tor}^A_{i-1}(M,k)\xrightarrow{x}\operatorname{Tor}^A_{i-1}(M,k)
\]
for all $i\geq 1$.
Now, for every $A$-module $N$, the homothety of ratio $x$ in $M\otimes_A N$ arises equally well from the homothety of ratio $x$ in $M$ as from the homothety of ratio $x$ in $N$;
we conclude immediately, since the homothety of ratio $x$ in $k$ is zero by definition, that for every $i$, the homomorphism $\operatorname{Tor}^A_i(M,k)\xrightarrow{x}\operatorname{Tor}^A_i(M,k)$ is zero;
in other words, we have the exact sequence
\[
  0\to\operatorname{Tor}^A_i(M,k)\to\operatorname{Tor}^A_i(M/xM,k)\to\operatorname{Tor}^A_{i-1}(M,k)\to 0.
\]
If $\operatorname{dim.proj}(M)=n$, we have $\operatorname{Tor}^A_n(M,k)\neq 0$ and $\operatorname{Tor}^A_{n+1}(M,k)=\operatorname{Tor}^A_{n+2}(M,k)=0$ by \sref{0.17.2.6}.
It thus follows from what precedes that $\operatorname{Tor}^A_{n+1}(M/xM,k)\neq 0$ and $\operatorname{Tor}^A_{n+2}(M/xM,k)=0$, so $\operatorname{dim.proj}(M/xM)=n+1$ by \sref{0.17.2.6}.
\end{proof}

\begin{proposition}[17.2.8]
\label{0.17.2.8}
\textnormal{(M.~Auslander)}
Let $A$ be a ring.
The following conditions are equivalent:
\begin{enumerate}
  \item[{\rm a)}] Every $A$-module is of projective dimension $\leq n$.
  \item[{\rm a')}] Every $A$-module of finite type is of projective dimension $\leq n$.
  \item[{\rm b)}] Every $A$-module is of injective dimension $\leq n$.
  \item[{\rm c)}] For every pair of $A$-modules $M$, $N$ we have $\operatorname{Ext}^{n+1}_A(M,N)=0$.
  \item[{\rm c')}] For every pair of $A$-modules $M$, $N$ such that $M$ is of finite type (or only monogenic), we have $\operatorname{Ext}^{n+1}_A(M,N)=0$.
\end{enumerate}
\end{proposition}

\begin{proof}
This follows immediately from \sref{0.17.2.1} and \sref{0.17.2.3}.
\end{proof}

\oldpage[0\textsubscript{IV-1}]{45}
The smallest number $n$ (integer or $+\infty$) for which the equivalent conditions of \sref{0.17.2.8} are satisfied is called the \emph{global cohomological dimension} (or simply \emph{cohomological dimension}) of $A$ and denoted $\operatorname{dim.coh}(A)$.

\begin{proposition}[17.2.9]
\label{0.17.2.9}
Let $A$ be a Noetherian ring.
The following conditions are equivalent:
\begin{enumerate}
  \item[{\rm a)}] $\operatorname{dim.coh}(A)\leq n$.
  \item[{\rm b)}] Every $A$-module of finite type is of injective dimension $\leq n$.
  \item[{\rm c)}] For every pair of $A$-modules of finite type, $M$, $N$, we have $\operatorname{Ext}^{n+1}_A(M,N)=0$.
\end{enumerate}
\end{proposition}

\begin{proof}
This follows immediately from the definition and from \sref{0.17.2.4}.
\end{proof}

\begin{corollary}[17.2.10]
\label{0.17.2.10}
If $A$ is a Noetherian ring, we have
\[
\label{0.17.2.10.1}
  \operatorname{dim.coh}(A)=\sup_\mathfrak{m}\big(\operatorname{dim.coh}(A_\mathfrak{m})\big)
  \tag{17.2.10.1}
\]
where $\mathfrak{m}$ ranges over the spectrum of $A$ (or the set of maximal ideals of $A$).
\end{corollary}

\begin{proof}
This follows immediately from \sref{0.17.2.9} and \sref{0.17.2.5}.
\end{proof}

\begin{proposition}[17.2.11]
\label{0.17.2.11}
Let $A$ be a Noetherian local ring, $k$ its residue field.
For $\operatorname{dim.coh}(A)\leq n$, it is necessary that $\operatorname{Tor}^A_i(k,k)=0$ for $i>n$, and it suffices that $\operatorname{Tor}^A_{n+1}(k,k)=0$.
\end{proposition}

\begin{proof}
Taking into account \sref{0.17.2.6}, it suffices to prove that the relations $\operatorname{dim.coh}(A)\leq n$ and $\operatorname{dim.proj}_A(k)\leq n$ are equivalent.
It is clear that the first implies the second by definition.
Conversely, if $\operatorname{dim.proj}_A(k)\leq n$, we have $\operatorname{Tor}^A_i(M,k)=0$ for $i>n$ and for every $A$-module $M$ by \sref{0.17.2.2}[(ii)];
thus $\operatorname{dim.proj}(M)\leq n$, which proves the proposition by \sref{0.17.2.8}.
\end{proof}

\begin{corollary}[17.2.12]
\label{0.17.2.12}
Let $A$ be a Noetherian ring.
For $\operatorname{dim.coh}(A)\leq n$, it is necessary that, for every maximal ideal $\mathfrak{m}$ of $A$, we have $\operatorname{Tor}^{A_\mathfrak{m}}_i(A/\mathfrak{m},A/\mathfrak{m})=0$ for $i>n$, and it suffices that these relations be satisfied for $i=n+1$.
\end{corollary}

\begin{proof}
This follows immediately from \sref{0.17.2.11} and \sref{0.17.2.10}.
\end{proof}

\begin{proposition}[17.2.13]
\label{0.17.2.13}
Let $A$, $B$ be two Noetherian local rings, $\varphi:A\to B$ a local homomorphism making $B$ a flat $A$-module.
Then we have
\[
\label{0.17.2.13.1}
  \operatorname{dim.coh}(A)\leq\operatorname{dim.coh}(B).
  \tag{17.2.13.1}
\]
\end{proposition}

\begin{proof}
Suppose that $\operatorname{dim.coh}(B)=n$ is finite;
it suffices to prove that for every pair $(M,N)$ of $A$-modules of finite type, we have $\operatorname{Tor}^A_i(M,N)=0$ for $i>n$ \sref{0.17.2.11}.
Since $B$ is a faithfully flat $A$-module \sref[0\textsubscript{I}]{0.6.6.2}, it amounts to the same \sref[0\textsubscript{I}]{0.6.4.1} to show that
\[
\label{0.17.2.13.2}
  \operatorname{Tor}^A_i(M,N)\otimes_A B=0.
  \tag{17.2.13.2}
\]
\oldpage[0\textsubscript{IV-1}]{46}
Now, if $\mathrm{L}_\bullet=(\mathrm{L}_j)$ is a right resolution of $M$ by free $A$-modules, it follows from the fact that $B$ is a flat $A$-module that $\mathrm{L}_\bullet\otimes_A B=(\mathrm{L}_j\otimes_A B)$ is a right resolution of the $B$-module $M\otimes_A B$ by free $B$-modules;
moreover, we have $(\mathrm{L}_j\otimes_A B)\otimes_B(N\otimes_A B)=(\mathrm{L}_j\otimes_A N)\otimes_A B$, whence we conclude immediately that the first term of \sref{0.17.2.13.2} is equal to $\operatorname{Tor}^B_i(M\otimes_A B,N\otimes_A B)$;
the hypothesis on $B$ implies that this $B$-module is zero for $i>n$ \sref{0.17.2.2}[(ii)], whence the conclusion.
\end{proof}

\begin{env}[17.2.14]
\label{0.17.2.14}
Let $(X,\sh{O}_X)$ be a ringed space, $\sh{F}$ an $\sh{O}_X$-Module;
we call the \emph{pointwise projective dimension} (resp. \emph{injective dimension}) of $\sh{F}$ and denote $\operatorname{dim.proj}(\sh{F})$ (resp. $\operatorname{dim.inj}(\sh{F})$) the number $\sup_{x\in X}(\operatorname{dim.proj}(\sh{F}_x))$ (resp. $\sup_{x\in X}(\operatorname{dim.inj}(\sh{F}_x))$).
We call the \emph{pointwise cohomological dimension} of $X$ and denote $\operatorname{dim.coh}(X)$ the number $\sup_{x\in X}(\operatorname{dim.coh}(\sh{O}_x))$.
It follows from \sref{0.17.2.5} and \sref{0.17.2.10} that when $A$ is a Noetherian ring, $M$ an $A$-module of finite type and $X=\Spec(A)$, we have $\operatorname{dim.proj}(\widetilde{M})=\operatorname{dim.proj}(M)$, $\operatorname{dim.inj}(\widetilde{M})=\operatorname{dim.inj}(M)$ and $\operatorname{dim.coh}(X)=\operatorname{dim.coh}(A)$.
We call the \emph{projective dimension} (resp. \emph{injective dimension}) of $\sh{F}$ at a point $x\in X$ the projective (resp. injective) dimension of $\sh{F}_x$, and the \emph{cohomological dimension of $X$ at a point $x$} the cohomological dimension of $\sh{O}_x$.
\end{env}

\begin{proposition}[17.2.15]
\label{0.17.2.15}
Let $X$, $Y$ be two ringed spaces in Noetherian local rings, $f:X\to Y$ a flat morphism.
If $\operatorname{dim.coh}(X)\leq n$, then $Y$ is of cohomological dimension $\leq n$ at every point of $f(X)$.
\end{proposition}

\begin{proof}
This follows immediately from \sref{0.17.2.13}.
\end{proof}

\subsection{Cohomological theory of regular rings}
\label{subsection:0.17.3}

\begin{theorem}[17.3.1]
\label{0.17.3.1}
\textnormal{(Hilbert-Serre)}
Let $A$ be a Noetherian local ring.
For $A$ to be of finite cohomological dimension, it is necessary and sufficient that $A$ be regular; we then have
\[
\label{0.17.3.1.1}
  \operatorname{dim.coh}(A)=\dim(A).
  \tag{17.3.1.1}
\]
\end{theorem}

\begin{proof}
Suppose $A$ regular;
let $\mathfrak{m}$ be its maximal ideal, $\mathbf{x}=(x_i)_{1\leq i\leq n}$ a regular system of parameters of $A$ \sref{0.17.1.6};
consider the Koszul complex $\mathrm{K}_\bullet(\mathbf{x})$ (\textbf{III}, 1.1.1), which is formed of free $A$-modules, with $\mathrm{K}_i(\mathbf{x})=0$ for $i>n$;
since $(x_i)$ is a regular sequence, we have $\mathrm{H}_i(\mathrm{K}_\bullet(\mathbf{x}))=0$ for $i>0$ (\textbf{III}, 1.1.4 and 1.1.3.3) and $\mathrm{H}_0(\mathrm{K}_\bullet(\mathbf{x}))=A/(x_1 A+\cdots+x_n A)=A/\mathfrak{m}=k$;
the $\mathrm{K}_i(\mathbf{x})$ thus form a free resolution of $k$ of length $n$.
Now, the fact that the $x_i$ belong to $\mathfrak{m}$ immediately implies that in the complex $\mathrm{K}_\bullet(\mathbf{x},k)=\mathrm{K}_\bullet(\mathbf{x})\otimes_A k$, the boundary operator is zero in all dimensions, so that we have, by definition, $\operatorname{Tor}^A_i(k,k)=\textstyle\bigwedge^i(k^n)$;
the equality \sref{0.17.3.1.1} thus follows immediately from \sref{0.17.2.11} (this result is essentially the ``syzygy theorem'' of Hilbert).

Let us now show that if $A$ is a Noetherian local ring, $\mathfrak{m}$ its maximal ideal, and if $\operatorname{dim.proj}_A(\mathfrak{m})$ is \emph{finite}, then $A$ is regular, which will finish proving \sref{0.17.3.1}.
We proceed by induction on $n=\operatorname{rg}_k(\mathfrak{m}/\mathfrak{m}^2)$.
For $n=0$, we have $\mathfrak{m}=0$ and the assertion is trivial.
\end{proof}

\begin{lemma}[17.3.1.2]
\label{0.17.3.1.2}
\textnormal{(Nagata)}
If every element of $\mathfrak{m}-\mathfrak{m}^2$ is a zero divisor in $A$, there exists $c\neq 0$ in $A$ such that $c\mathfrak{m}=0$ (in other words $\mathfrak{m}\in\operatorname{Ass}(A)$).
\end{lemma}

\begin{proof}
\oldpage[0\textsubscript{IV-1}]{47}
We may restrict to the case $\mathfrak{m}\neq 0$, hence $\mathfrak{m}\neq\mathfrak{m}^2$.
The hypothesis implies that $\mathfrak{m}-\mathfrak{m}^2$ is contained in the union of the ideals $\mathfrak{p}_i$ of $\operatorname{Ass}(A)$ (Bourbaki, \emph{Alg.~comm.}, chap.~IV, \textsection1, n\textsuperscript{o}~1, cor.~3 of prop.~2);
thus $\mathfrak{m}$ is contained in the union of $\mathfrak{m}^2$ and the $\mathfrak{p}_i$, hence in one of the $\mathfrak{p}_i$, since $\mathfrak{m}\not\subset\mathfrak{m}^2$ (Bourbaki, \emph{Alg.~comm.}, chap.~II, \textsection1, n\textsuperscript{o}~1, prop.~2);
$\mathfrak{m}$ being maximal, this proves the lemma.
\end{proof}

\begin{lemma}[17.3.1.3]
\label{0.17.3.1.3}
If $a\in\mathfrak{m}-\mathfrak{m}^2$, $\mathfrak{m}/Aa$ is isomorphic to a direct factor of $\mathfrak{m}/a\mathfrak{m}$.
\end{lemma}

\begin{proof}
There exists a minimal system of generators of $\mathfrak{m}$ containing $a$ (whose images in $\mathfrak{m}/\mathfrak{m}^2$ form a basis of this $k$-vector space);
let $\mathfrak{b}$ be the ideal generated by the elements of this system other than $a$.
Since the relation $xa\in\mathfrak{b}$ implies $x\in\mathfrak{m}$ (considering the image of $xa$ in $\mathfrak{m}/\mathfrak{m}^2$), we have $\mathfrak{b}\cap Aa\subset a\mathfrak{m}$ whence a homomorphism $\mathfrak{b}/(\mathfrak{b}\cap Aa)\to\mathfrak{m}/a\mathfrak{m}$ deduced from the injection $\mathfrak{b}\to\mathfrak{m}$ by passing to quotients, and which is injective.
Moreover, $\mathfrak{b}+Aa=\mathfrak{m}$, and the composite homomorphism
\[
  \mathfrak{m}/Aa=(\mathfrak{b}+Aa)/Aa\xrightarrow{\sim}\mathfrak{b}/(\mathfrak{b}\cap Aa)\to\mathfrak{m}/a\mathfrak{m}\to\mathfrak{m}/Aa
\]
is the identity; whence the lemma.
\end{proof}

\begin{lemma}[17.3.1.4]
\label{0.17.3.1.4}
Let $A$ be a Noetherian local ring, $\mathfrak{m}$ its maximal ideal, $E$ an $A$-module of finite type and of finite projective dimension.
If $a\in\mathfrak{m}$ is $A$-regular and $E$-regular, then $E/aE$ is an $(A/aA)$-module of finite projective dimension, at most equal to $\operatorname{dim.proj}_A(E)$.
\end{lemma}

\begin{proof}
We argue by induction on $h=\operatorname{dim.proj}_A(E)$, the case $h=0$ being trivial since then $E$ is a projective $A$-module, hence $E/aE$ is a projective $(A/aA)$-module.
There exists an exact sequence
\[
  0\to N\to L\to E\to 0
\]
where $L$ is free and $\operatorname{dim.proj}_A(N)=h-1$ \sref{0.17.2.2}[(iii)], $N$ being of finite type.
Moreover, the sequence
\[
  0\to N/aN\to L/aL\to E/aE\to 0
\]
\oldpage[0\textsubscript{IV-1}]{48}
is exact \sref{0.15.1.18}.
Since $a$ is $A$-regular, it is also $N$-regular since $L$ is free;
the induction hypothesis implies that $N/aN$ is an $(A/aA)$-module of projective dimension $\leq h-1$, and since $L/aL$ is a free $(A/aA)$-module, $E/aE$ is an $(A/aA)$-module of projective dimension $\leq h$.
\end{proof}

Let us now examine two cases:

\medskip\noindent\textbf{I.} ---
Suppose first that every element of $\mathfrak{m}-\mathfrak{m}^2$ is a zero divisor in $A$, in which case \sref{0.17.3.1.2} there exists $c\neq 0$ in $A$ such that $c\mathfrak{m}=0$.
Let us show then that $\mathfrak{m}=0$.
If this were not the case, note first that $\mathfrak{m}$ could not be a projective $A$-module, since it would be free (Bourbaki, \emph{Alg.~comm.}, chap.~II, \textsection3, n\textsuperscript{o}~2, cor.~2 of prop.~5), which contradicts the relation $c\mathfrak{m}=0$.
We would thus have $n=\operatorname{dim.coh}(A)\geq 1$.
Since $\mathfrak{m}\in\operatorname{Ass}(A)$, there would exist an exact sequence of $A$-homomorphisms
\[
  0\to k\to A\to E\to 0.
\]
But this is absurd, since, in view of the relation $n\geq 1$, the exact sequence of Tor would give the exact sequence $0\to\operatorname{Tor}^A_{n+1}(E,k)\to\operatorname{Tor}^A_n(k,k)\to 0$;
but we have $\operatorname{Tor}^A_{n+1}(E,k)=0$ \sref{0.17.2.2}[(ii)] and $\operatorname{Tor}^A_n(k,k)\neq 0$ \sref{0.17.2.11} and we have reached a contradiction.

\medskip\noindent\textbf{II.} ---
We can thus restrict to the case where there exists $a\in\mathfrak{m}-\mathfrak{m}^2$ that is $A$-regular, and consequently also $\mathfrak{m}$-regular.
Consider the ring $A'=A/aA$ and its maximal ideal $\mathfrak{m}'=\mathfrak{m}/aA$;
it is clear that $\operatorname{rg}_k(\mathfrak{m}'/\mathfrak{m}'^2)=n-1$.
By \sref{0.17.3.1.4}, $\mathfrak{m}/a\mathfrak{m}$ is an $A'$-module of finite projective dimension, hence so is $\mathfrak{m}'=\mathfrak{m}/aA$, which is a direct factor of it \sref{0.17.3.1.3} and \sref{0.17.2.1}.
The induction hypothesis implies therefore that $A'=A/aA$ is regular, which proves by \sref{0.17.1.8} that $A$ is regular.

\begin{corollary}[17.3.2]
\label{0.17.3.2}
If $A$ is a regular local ring, $A_\mathfrak{p}$ is regular for every prime ideal $\mathfrak{p}$ of $A$.
\end{corollary}

\begin{proof}
Indeed, we have seen \sref{0.17.2.10} that $\operatorname{dim.coh}(A_\mathfrak{p})\leq\operatorname{dim.coh}(A)$, so the conclusion follows from \sref{0.17.3.1}.
\end{proof}

\begin{proposition}[17.3.3]
\label{0.17.3.3}
Let $A$, $B$ be two Noetherian local rings, $\varphi:A\to B$ a local homomorphism, $\mathfrak{m}$ (resp. $\mathfrak{n}$) the maximal ideal of $A$ (resp. $B$), $k=A/\mathfrak{m}$ (resp. $k'=B/\mathfrak{n}$) the residue field of $A$ (resp. $B$).
We thus have a canonical homomorphism of $k'$-vector spaces
\[
\label{0.17.3.3.1}
  \psi:(\mathfrak{m}/\mathfrak{m}^2)\otimes_k k'\to\mathfrak{n}/\mathfrak{n}^2.
  \tag{17.3.3.1}
\]
\begin{enumerate}
  \item[{\rm(i)}] If $B$ is regular and $B$ is a flat $A$-module, $A$ is regular.
  \item[{\rm(ii)}] The following conditions are equivalent:
  \begin{enumerate}
    \item[{\rm a)}] $B$ is regular and the homomorphism $\psi$ \sref{0.17.3.3.1} is injective.
    \item[{\rm b)}] $A$ and $B$ are regular, and for a regular system of parameters $(x_i)_{1\leq i\leq m}$ of $A$, the $\varphi(x_i)$ are part of a regular system of parameters of $B$ (in which case this property holds for every regular system of parameters of $A$).
    \item[{\rm c)}] $B$ and $B\otimes_A k$ are regular, and $B$ is a flat $A$-module.
    \item[{\rm d)}] $A$ and $B\otimes_A k$ are regular, and $B$ is a flat $A$-module.
    \item[{\rm e)}] $A$ and $B\otimes_A k$ are regular, and we have
    \[
    \label{0.17.3.3.2}
      \dim(B)=\dim(A)+\dim(B\otimes_A k).
      \tag{17.3.3.2}
    \]
  \end{enumerate}
\end{enumerate}
\end{proposition}

\begin{proof}
\oldpage[0\textsubscript{IV-1}]{49}
(i) We have $\operatorname{dim.coh}(A)\leq\operatorname{dim.coh}(B)$ by \sref{0.17.2.13}, so it suffices to apply \sref{0.17.3.1}.

(ii) When $A$ is assumed regular and $(x_i)$ is a regular system of parameters of $A$, saying that $B$ is a flat $A$-module is equivalent, by \sref{0.15.1.21}, to saying that the sequence of $y_i=\varphi(x_i)$ is $B$-regular (since $B\otimes_A k$ is a flat $k$-module).
On the other hand, since $\dim(A)=m$ and the sequence $(y_i)$ generates the ideal $\mathfrak{m}B$, the relation \sref{0.17.3.3.2}, which can also be written $\dim(B/\mathfrak{m}B)=\dim(B)-m$, is equivalent to saying that the sequence $(y_i)_{1\leq i\leq m}$ is part of a system of parameters of $B$ \sref{0.16.3.7}.
We thus see that conditions b), d) and e) are respectively equivalent to the following:
\begin{enumerate}
  \item[{\rm b')}] $A$ and $B$ are regular and $(y_i)_{1\leq i\leq m}$ is part of a regular system of parameters of $B$.
  \item[{\rm d')}] $A$ is regular, $B/(\sum y_i B)$ is regular and the sequence $(y_i)$ is $B$-regular.
  \item[{\rm e')}] $A$ is regular, $B/(\sum y_i B)$ is regular, $(y_i)$ is part of a system of parameters of $B$.
\end{enumerate}
Now, b') and e') are equivalent by \sref{0.17.1.7}, and since d') implies e') and is implied by b') \sref{0.17.1.7}, it is equivalent to both.
The conjunction of b) and d) trivially implies c), and by (i), c) implies d);
we have thus proved the equivalence of b), c), d) and e).
Moreover, it is clear that b) implies a), the classes of the $x_i$ in $\mathfrak{m}/\mathfrak{m}^2$ then forming a basis of this $k$-vector space and the classes of the $y_i$ in $\mathfrak{n}/\mathfrak{n}^2$ a free system of this $k'$-vector space.
It thus remains to prove that a) implies e).
Set $V=\mathfrak{m}/\mathfrak{m}^2$, $W=\mathfrak{n}/\mathfrak{n}^2$;
it is immediate that $\psi(V\otimes_k k')=(\mathfrak{n}^2+\mathfrak{m}B)/\mathfrak{n}^2$.
On the one hand, by \sref{0.16.3.9}
\[
\label{0.17.3.3.3}
  \dim(B)\leq\dim(A)+\dim(B\otimes_A k);
  \tag{17.3.3.3}
\]
secondly \sref{0.16.2.6}, we have
\[
\label{0.17.3.3.4}
  \dim(A)\leq\operatorname{rg}_k(V\otimes_k k').
  \tag{17.3.3.4}
\]
Finally, in the local ring $B\otimes_A k=B/\mathfrak{m}B$, $\mathfrak{n}'=\mathfrak{n}/\mathfrak{m}B$ is the maximal ideal, $k'$ the residue field, and $\mathfrak{n}'/\mathfrak{n}'^2$ is isomorphic to $\mathfrak{n}/(\mathfrak{n}^2+\mathfrak{m}B)=W/\psi(V\otimes_k k')$;
we thus have, by \sref{0.16.2.6}
\[
\label{0.17.3.3.5}
  \dim(B\otimes_A k)\leq\operatorname{rg}_{k'}W-\operatorname{rg}_{k'}\psi(V\otimes_k k').
  \tag{17.3.3.5}
\]
Finally, since $B$ is assumed regular, we have $\dim(B)=\operatorname{rg}_{k'}W$ \sref{0.17.1.1};
we thus conclude from \sref{0.17.3.3.3}, \sref{0.17.3.3.4} and \sref{0.17.3.3.5} that
\[
  \operatorname{rg}_{k'}W\leq\dim(A)+\dim(B\otimes_A k)\leq\operatorname{rg}_{k'}W+\operatorname{rg}_{k'}(V\otimes_k k')-\operatorname{rg}_{k'}\psi(V\otimes_k k')
\]
and saying that the two extreme terms of this inequality are equal means that $\psi$ is injective.
Condition a) thus necessarily implies that in each of the relations \sref{0.17.3.3.3}, \sref{0.17.3.3.4} and \sref{0.17.3.3.5}, the two sides are equal;
now, equality in \sref{0.17.3.3.4} (resp. \sref{0.17.3.3.5}) means that $A$ (resp. $B\otimes_A k$) is regular \sref{0.17.1.1};
we have thus proved that a) implies e).
\end{proof}

\begin{proposition}[17.3.4]
\label{0.17.3.4}
Let $A$ be a regular local ring of dimension $n$, $\mathfrak{m}$ its maximal ideal.
For every non-zero $A$-module of finite type $M$, we have
\[
\label{0.17.3.4.1}
  \operatorname{prof}(M)+\operatorname{dim.proj}(M)=n.
  \tag{17.3.4.1}
\]
\end{proposition}

\begin{proof}
We argue by induction on $m=\operatorname{prof}(M)$.
If $m=0$, we know \sref{0.16.4.6}[(i)] that there exists a sub-$A$-module $N$ of $M$ isomorphic to $k=A/\mathfrak{m}$;
applying the Tor exact sequence to the exact sequence $0\to N\to M\to M/N\to 0$, we obtain an exact sequence
\[
  \operatorname{Tor}^A_{n+1}(M/N,k)\to\operatorname{Tor}^A_n(N,k)\to\operatorname{Tor}^A_n(M,k),
\]
and the hypothesis that $A$ is regular implies $\operatorname{Tor}^A_{n+1}(M/N,k)=0$ (\sref{0.17.3.1.1} and \sref{0.17.2.6}), so $\operatorname{Tor}^A_n(k,k)$ is isomorphic to a submodule of $\operatorname{Tor}^A_n(M,k)$;
applying again \sref{0.17.3.1.1} and \sref{0.17.2.6}, we see that $\operatorname{Tor}^A_n(M,k)\neq 0$, so $\operatorname{dim.proj}(M)\geq n$ by \sref{0.17.2.6};
but since $\operatorname{dim.proj}(M)\leq n$ by \sref{0.17.3.1.1}, we indeed have $\operatorname{dim.proj}(M)=n$.
Suppose now $m>0$, and let $x$ be an $M$-regular element belonging to $\mathfrak{m}$;
we then know \sref{0.16.4.6}[(ii)] that $\operatorname{prof}(M/xM)=m-1$, and on the other hand \sref{0.17.2.7} $\operatorname{dim.proj}(M/xM)=\operatorname{dim.proj}(M)+1$;
the induction hypothesis immediately proves the relation \sref{0.17.3.4.1}.
\end{proof}

\oldpage[0\textsubscript{IV-1}]{50}

\begin{corollary}[17.3.5]
\label{0.17.3.5}
\begin{enumerate}
  \item[{\rm(i)}] Let $A$ be a regular local ring of dimension $n$; for an $A$-module $M\neq 0$ of finite type to be free, it is necessary and sufficient that $M$ be a Cohen-Macaulay $A$-module of dimension $n$.
  \item[{\rm(ii)}] Let $A$ be a regular local ring, $B$ a local ring, $\rho:A\to B$ a local homomorphism making $B$ an $A$-module of finite type.
  For $B$ to be a free $A$-module, it is necessary and sufficient that $B$ be a Cohen-Macaulay ring and that $\rho$ be injective (or, equivalently \sref{0.16.1.5}, that $\dim(B)=\dim(A)$).
\end{enumerate}
\end{corollary}

\begin{proof}
(i) This follows from \sref{0.17.3.4.1} and from the fact that for an $A$-module of finite type, it is equivalent to saying that this module is projective or free (Bourbaki, \emph{Alg.~comm.}, chap.~II, \textsection3, n\textsuperscript{o}~2, cor.~2 of prop.~5);
the free $A$-modules of finite type $M$ are thus characterised by the relation $\operatorname{dim.proj}(M)=0$ \sref{0.17.2.2}[(i)].

(ii) To say that $B$ is a Cohen-Macaulay ring is equivalent to saying that $B$ is a Cohen-Macaulay $A$-module \sref{0.16.5.3}, so it suffices to apply (i), since $\dim B=\dim A$ \sref{0.16.1.5}.
\end{proof}

\begin{env}[17.3.6]
\label{0.17.3.6}
We say that a Noetherian ring $A$ is \emph{regular} if for every prime ideal $\mathfrak{p}$ of $A$, the local ring $A_\mathfrak{p}$ is regular;
when $A$ is itself a local ring, it follows from \sref{0.17.3.2} that this definition is equivalent to that of \sref{0.17.1.2}.
For $A$ to be regular, it suffices, by \sref{0.17.3.2}, that $A_\mathfrak{m}$ be regular for every \emph{maximal} ideal $\mathfrak{m}$ of $A$.
Moreover, it follows immediately from this definition that for every multiplicative subset $S$ of $A$, $S^{-1}A$ is regular.
\end{env}

\begin{proposition}[17.3.7]
\label{0.17.3.7}
If $A$ is a regular Noetherian ring, every polynomial ring $A[\mathrm{T}_1,\ldots,\mathrm{T}_n]$ is regular.
\end{proposition}

\begin{proof}
It clearly suffices to prove that the polynomial ring $B=A[\mathrm{T}]$ is regular;
since $B$ is a free $A$-module, for every prime ideal $\mathfrak{q}$ of $B$, $B_\mathfrak{q}$ is a flat $A_\mathfrak{p}$-module, where $\mathfrak{p}=\mathfrak{q}\cap A$ \sref[0\textsubscript{I}]{0.6.3.1};
  it thus suffices, by \erratum[IV]{\sref{0.17.3.3}}, to prove that $B_\mathfrak{q}/\mathfrak{p}B_\mathfrak{q}$ is regular, and since this ring is a local ring at a prime ideal of $A_\mathfrak{p}[\mathrm{T}]/\mathfrak{p}A_\mathfrak{p}[\mathrm{T}]=k[\mathrm{T}]$ (where $k$ is the residue field $A_\mathfrak{p}/\mathfrak{p}A_\mathfrak{p}$), it suffices to prove that $k[\mathrm{T}]$ is regular;
now, $\mathrm{C}=k[\mathrm{T}]$ being principal, the local rings of $\mathrm{C}$ are discrete valuation rings or a field (for the ideal $(0)$), thus regular \sref{0.17.1.4}, which concludes the proof.
\end{proof}

\begin{corollary}[17.3.8]
\label{0.17.3.8}
If $A$ is a regular ring, every formal power series ring $A[[\mathrm{T}_1,\ldots,\mathrm{T}_n]]$ is regular.
\end{corollary}

\begin{proof}
Let $\mathfrak{J}$ be the ideal generated by the $\mathrm{T}_i$ in the polynomial ring $A[\mathrm{T}_1,\ldots,\mathrm{T}_n]$;
since the latter is regular by \sref{0.17.3.7}, and $A[[\mathrm{T}_1,\ldots,\mathrm{T}_n]]$ is the completion of $A[\mathrm{T}_1,\ldots,\mathrm{T}_n]$ for the $\mathfrak{J}$-preadic topology (Bourbaki, \emph{Alg.~comm.}, chap.~III, \textsection2, n\textsuperscript{o}~12, \emph{Example~$1$}), the conclusion follows from the following lemma.
\end{proof}

\begin{lemma}[17.3.8.1]
\label{0.17.3.8.1}
Let $A$ be a regular ring, $\mathfrak{J}$ an ideal of $A$, $\widehat{A}$ the separated completion of $A$ for the $\mathfrak{J}$-preadic topology.
Then $\widehat{A}$ is regular.
\end{lemma}

\begin{proof}
It suffices in fact \sref{0.17.3.6} to see that for every maximal ideal $\mathfrak{n}$ of $\widehat{A}$, the local ring $\widehat{A}_\mathfrak{n}$ is regular;
  we know (Bourbaki, \emph{Alg.~comm.}, chap.~III, \textsection3, n\textsuperscript{o}~4, prop.~8) that $\mathfrak{n}$
\oldpage[0\textsubscript{IV-1}]{51}
is of the form $\mathfrak{m}\widehat{A}$, where $\mathfrak{m}$ is a maximal ideal of $A$ containing $\mathfrak{J}$, and that the canonical homomorphism $A\to\widehat{A}$ gives an injective homomorphism $A_\mathfrak{m}\to\widehat{A}_\mathfrak{n}$, such that the $\mathfrak{n}\widehat{A}_\mathfrak{n}$-preadic topology of $\widehat{A}_\mathfrak{n}$ induces on $A_\mathfrak{m}$ the $\mathfrak{m}A_\mathfrak{m}$-preadic topology, and that $A_\mathfrak{m}$ is dense in $\widehat{A}_\mathfrak{n}$ for the $\mathfrak{n}\widehat{A}_\mathfrak{n}$-preadic topology;
consequently the completions of the Noetherian local rings $A_\mathfrak{m}$ and $\widehat{A}_\mathfrak{n}$ are canonically isomorphic.
Now, by hypothesis $A_\mathfrak{m}$ is regular, hence so is its completion \sref{0.17.1.5} and since the completion of $\widehat{A}_\mathfrak{n}$ is regular, it follows that $\widehat{A}_\mathfrak{n}$ is regular by \sref{0.17.1.5}.
\end{proof}

\begin{corollary}[17.3.9]
\label{0.17.3.9}
Let $A$ be a Noetherian ring, quotient of a regular Noetherian ring $B$.
If $C$ is an $A$-algebra of finite type, every ring of fractions $S^{-1}C$ of $C$ is a quotient of a regular ring.
\end{corollary}

\begin{proof}
Indeed, $C$ is a quotient of a polynomial ring $A[\mathrm{T}_1,\ldots,\mathrm{T}_n]$;
since $A=B/\mathfrak{J}$, where $\mathfrak{J}$ is an ideal of $B$, we have $A[\mathrm{T}_1,\ldots,\mathrm{T}_n]=B[\mathrm{T}_1,\ldots,\mathrm{T}_n]/\mathfrak{J}B[\mathrm{T}_1,\ldots,\mathrm{T}_n]$;
by \sref{0.17.3.7}, we may restrict to the case where $C$ is a quotient of a regular ring $B'$;
but if $S'$ is the inverse image of $S$ in $B'$, $S'$ is a multiplicative subset of $B'$ and $S^{-1}C$ is a quotient ring of $S'^{-1}B'$, so that ultimately everything reduces to seeing that when $C$ is regular, $S^{-1}C$ is also regular, which we saw in \sref{0.17.3.6}.
\end{proof}

The importance of quotients of regular rings comes among other things from the property that they are \emph{catenary} rings \sref{0.16.5.12}.

All the important rings in the applications to algebraic geometry are quotients of regular rings.
% END NATIVE INPUT ega0/ega0-17
% BEGIN NATIVE INPUT ega0/ega0-18
\section{Supplement on extensions of algebras}
\label{section:0.18}

This paragraph assembles a certain number of functorial constructions on rings, which will be used repeatedly in \textsection\textsection~19 and 20; it contains no non-trivial result.

\subsection{Inverse images of augmented rings}
\label{subsection:0.18.1}

\begin{env}[18.1.1]
\label{0.18.1.1}
Given a ring $A$ (not necessarily commutative), the category of \emph{$A$-rings} has as objects the pairs $(B,\rho)$ formed of a ring $B$ and a homomorphism of rings $\rho:A\to B$, for morphisms (also called \emph{$A$-homomorphisms}) the homomorphisms of rings $u:B\to C$ such that, if $\rho:A\to B$ and $\sigma:A\to C$ are the homomorphisms (called \emph{structural}) defining the structure of $A$-ring on $B$ and $C$ respectively, the diagram
\begin{equation}
\tag{18.1.1.1}
  \label{0.18.1.1.1}
  \xymatrix{
    B\ar[r]^u & C\\
    & A\ar[lu]^\rho\ar[u]_\sigma
  }
\end{equation}
is commutative.
The \emph{kernel} $\mathfrak{J}$ of $u$ is a bilateral ideal that is an $A$-\emph{bimodule} (for $\rho$).

If $f:A'\to A$ is a homomorphism of rings, for every $A$-ring $(B,\rho)$, $(B,\rho\circ f)$ is an $A'$-ring, and if $u:B\to C$ is an $A$-homomorphism from the $A$-ring $(B,\rho)$
\oldpage[0\textsubscript{IV-1}]{52}
to the $A$-ring $(C,\sigma)$, it is also an $A'$-homomorphism from the $A'$-ring $(B,\rho\circ f)$ to the $A'$-ring $(C,\sigma\circ f)$; one thus defines a canonical functor from the category of $A$-rings into that of $A'$-rings.
\end{env}

\begin{env}[18.1.2]
\label{0.18.1.2}
Let $B$, $E$, $F$ be three $A$-rings, $f:E\to B$, $g:F\to B$ two $A$-homomorphisms.
Recall that one calls \emph{fibre product} of $E$ and $F$ over $B$ (for the homomorphisms $f$ and $g$) and one denotes $E\times_B F$ the subring of the product $E\times F$ formed of the pairs $(x,y)$ such that $f(x)=g(y)$; the restrictions $p_1:G\to E$, $p_2:G\to F$ of the projections $\operatorname{pr}_1$ and $\operatorname{pr}_2$ to $G$ are still called the \emph{canonical projections};
the structure of $A$-ring of $G$ is defined by the homomorphism $\alpha\mapsto(\rho(\alpha),\sigma(\alpha))$ (where $\rho:A\to E$ and $\sigma:A\to F$ are the structural homomorphisms), which indeed maps $A$ into $G$ by virtue of \sref{0.18.1.1}.
The characteristic property of $E\times_B F$ is that, for every pair of $A$-homomorphisms $u:C\to E$, $v:C\to F$ such that $f\circ u=g\circ v$, there exists one and only one $A$-homomorphism $w:C\to G$ such that $u=p_1\circ w$, $v=p_2\circ w$.
One can still say that $E\times_B F$ is the \emph{projective limit} of the projective system formed by $B$, $E$, $F$ and the $A$-homomorphisms $f$, $g$, in the category of $A$-rings ($\mathbf{0_{III}}$, 8.1.9).
\end{env}

\begin{env}[18.1.3]
\label{0.18.1.3}
Let $\mathfrak{J}$ be the bilateral ideal of $E$, kernel of $f$; it is immediate that the kernel $\mathfrak{J}'$ of $p_2:G\to F$ is the ideal formed of the elements $(x,0)$ where $x\in\mathfrak{J}$; the restriction $i_1:\mathfrak{J}'\to\mathfrak{J}$ of $p_1$ is an isomorphism of rings without identity, and also an isomorphism of $A$-bimodules.
Similarly, if $\mathfrak{R}$ is the kernel of $g$, the kernel $\mathfrak{R}'$ of $p_1:G\to E$ is the ideal of the elements $(0,y)$ where $y\in\mathfrak{R}$, and the restriction $i_2:\mathfrak{R}'\to\mathfrak{R}$ of $p_2$ is an isomorphism (in the same sense).
Finally, it is clear that the kernel of the $A$-homomorphism $q=f\circ p_1=g\circ p_2$ of $E\times_B F$ into $B$ is $\mathfrak{J}\times\mathfrak{R}=\mathfrak{J}'\oplus\mathfrak{R}'$, so that one has a commutative diagram
\begin{equation}
\tag{18.1.3.1}
  \label{0.18.1.3.1}
  \xymatrix@C=2.2pc@R=1.8pc{
    0\ar[dr] & & 0\ar[d] & 0\ar[d]\\
    & \mathfrak{J}'\oplus\mathfrak{R}'\ar[r]\ar[d]\ar[dr]
      & \mathfrak{R}'\ar[r]^{i_2}_{\sim}\ar[d]^j
      & \mathfrak{R}\ar[d]^{p_2}\\
    0\ar[r] & \mathfrak{J}'\ar[r]\ar[d]^{i_1}_{\sim}
      & G\ar[r]\ar[d]^{p_1}\ar[dr]^q & F\ar[d]^g\\
    0\ar[r] & \mathfrak{J}\ar[r] & E\ar[r]_f & B
  }
\end{equation}

The definitions and results of \sref{0.18.1.2} and \sref{0.18.1.3} extend immediately to the fibre product of any \emph{family} $(E_\lambda)_{\lambda\in I}$ of $A$-rings defined by a family of $A$-homomorphisms $f_\lambda:E_\lambda\to B$.
We leave the formulation of these results to the reader.
\end{env}

\begin{env}[18.1.4]
\label{0.18.1.4}
Following the terminology of (M, VIII), we call an $A$-ring $E$ equipped with a \emph{surjective} $A$-homomorphism (called the \emph{augmentation} of $E$), $f:E\to B$, an \emph{$A$-ring augmented over $B$}; the kernel $\mathfrak{J}$ of $f$ is called the \emph{augmentation ideal}.
One says that the $A$-augmented ring $E$ is \emph{trivial} if there exists an $A$-homomorphism $s:B\to E$ that is \emph{right inverse} of the augmentation $f:E\to B$ (in other words $f\circ s=1_B$).
The exact sequence of $A$-bimodules
\[
  0\to\mathfrak{J}\to E\xrightarrow{f} B\to 0
\]
\oldpage[0\textsubscript{IV-1}]{53}
is then \emph{split}; in other words, one can identify the $A$-bimodule $E$ with $B\times\mathfrak{J}$, and the multiplication in $E$ is then given by
\[
  (b,z)(b',z')=(bb',bz'+zb'+zz'),
\]
$\mathfrak{J}$ being considered as a $B$-bimodule by means of $s:B\to E$.
\end{env}

\begin{env}[18.1.5]
\label{0.18.1.5}
With the notations of \sref{0.18.1.2}, suppose that the $A$-homomorphism $f$ makes $E$ into an $A$-ring \emph{augmented over $B$}.
Then it is clear that $p_2:G\to F$ is also surjective, in other words $G$ is given a structure of $A$-ring \emph{augmented over $F$}, that one calls the \emph{inverse image} by $g:F\to B$ of the augmented $A$-ring $E$.
\end{env}

\begin{proposition}[18.1.6]
\label{0.18.1.6}
For the inverse image $G=E\times_B F$ by $g:F\to B$ of the augmented $A$-ring $E$ to be a trivial augmented $A$-ring, it is necessary and sufficient that there exists an $A$-homomorphism $u:F\to E$ making the diagram
  \[
  \xymatrix{
    & F\ar[ld]_u\ar[d]^g\\
    E\ar[r]_f & B
  }
  \]
commutative.
\end{proposition}

\begin{proof}
The condition is evidently necessary, taking $u=p_1\circ s$, where $s:F\to G$ is an $A$-homomorphism right inverse of $p_2$.
Conversely, if there exists such an $A$-homomorphism $u$, the existence of the right inverse $s$ of $p_2$ follows from the universal property of the fibre product \sref{0.18.1.2} applied to the $A$-homomorphisms $u:F\to E$ and $1_F:F\to F$.

This result implies in particular that if $E$ is a trivial augmented $A$-ring, so is every inverse image.
\end{proof}

\begin{env}[18.1.7]
\label{0.18.1.7}
Let us take up again the situation described in \sref{0.18.1.2} and \sref{0.18.1.3}; on $\mathfrak{R}$ one has evidently a structure of $G$-\emph{bimodule}, coming from its structure of $F$-bimodule and the homomorphism of rings $p_2:G\to F$.
As $i_2$ is bijective one has moreover an injection $\theta=j\circ i_2^{-1}:\mathfrak{R}\to G$, which is a homomorphism of $G$-\emph{bimodules}.
We shall now see conversely that, when one supposes moreover that $f$ and $g$ are \emph{surjective}, in other words that $E$ and $F$ are $A$-rings \emph{augmented} over $B$, the data of such a homomorphism $\theta$ allows us to reconstruct the augmented ring $E$ (over $B$) from the augmented ring $G$ (over $F$).

More precisely, suppose we are given an $A$-ring $F$ augmented over $B$ and an $A$-ring $G$ augmented over $F$, the augmentation ideals being denoted $\mathfrak{R}$ and $\mathfrak{J}'$ respectively:
\begin{equation}
\tag{18.1.7.1}
  \label{0.18.1.7.1}
  \xymatrix@C=2pc@R=1.5pc{
    & & 0\ar[d] & & \\
    & & \mathfrak{R}\ar[ld]^\theta\ar[d] & & \\
    0\ar[r] & \mathfrak{J}'\ar[r] & G\ar[r]_h & F\ar[r]\ar[d]^g & 0\\
    & & & B\ar[d] & \\
    & & & 0 &
  }
\end{equation}

\oldpage[0\textsubscript{IV-1}]{54}
The homomorphism $h$ defines on all the ideals of $F$ (and in particular on $\mathfrak{R}$) a structure of $G$-bimodule.
Let $\theta:\mathfrak{R}\to G$ be a homomorphism of $G$-bimodules making the diagram \sref{0.18.1.7.1} commutative; this implies that $\theta$ is \emph{injective} and that $\mathfrak{J}'\cap\theta(\mathfrak{R})=0$; as $\theta(\mathfrak{R})$ is a bilateral ideal of $G$, and $h(\theta(\mathfrak{R}))=\mathfrak{R}$, the composite homomorphism $g\circ h:G\to B$ factors as $G\xrightarrow{q}G/\theta(\mathfrak{R})\xrightarrow{f}B$, where $f$ is surjective; moreover the image $\mathfrak{J}$ of $\mathfrak{J}'$ in $E=G/\theta(\mathfrak{R})$ is the kernel of $f$ (the kernel of $g\circ h$ being $\mathfrak{J}'\oplus\theta(\mathfrak{R})$), and the restriction to $\mathfrak{J}'$ of the canonical homomorphism $q:G\to E$ is injective.
One can therefore form the \emph{fibre product} $G'=E\times_B F$, and as the two $A$-homomorphisms $q:G\to E$ and $h:G\to F$ are such that $f\circ q=g\circ h$, they define a unique $A$-homomorphism $u:G\to G'$ by the universal property of $G'$ \sref{0.18.1.2}; we shall now see that $u$ is \emph{bijective}.
It suffices to prove this when $u$ is considered as a homomorphism of $A$-bimodules; but one then notes that $u$ is compatible with the finite filtrations on $G$ and $G'$ formed respectively by $G$ and $\mathfrak{J}'$, and by $G'$ and $\mathfrak{J}''=\operatorname{Ker}(G'\to F)$; moreover, one has seen \sref{0.18.1.3} that $\operatorname{gr}_0 u:F\to F$ is the identity and that $\operatorname{gr}_1 u:\mathfrak{J}'\to\mathfrak{J}''$ is bijective, hence $u$ itself is bijective (Bourbaki, \emph{Alg.\ comm.}, chap.~III, \S~2, no.~8, cor.~3 of th.~1).
\end{env}

\subsection{Extensions of a ring by a bimodule}
\label{subsection:0.18.2}

\begin{env}[18.2.1]
\label{0.18.2.1}
Let $E$ be an $A$-ring augmented over $B$, $f:E\to B$ the augmentation, $\mathfrak{J}=\operatorname{Ker}(f)$ the augmentation ideal.
If $\mathfrak{J}^2=0$, $\mathfrak{J}$ is not only an $E$-bimodule but also a $B$-\emph{bimodule} since $B$ is isomorphic to $E/\mathfrak{J}$; more precisely, every $b\in B$ is of the form $f(x)$ with $x\in E$, and if $z$, $z'$ are in $\mathfrak{J}$, one has $(x+z)z'=xz'$ (resp.\ $z'(x+z)=z'x$) so that the value of $xz'$ (resp.\ $z'x$) does not depend on the element $x\in f^{-1}(b)$, and one can write $bz'$ (resp.\ $z'b$) which defines the structure of $B$-bimodule considered.
Conversely, if $\mathfrak{J}$ is equipped with a structure of $B$-bimodule such that $xz'=f(x)z'$ and $z'x=z'f(x)$ for $x\in E$ and $z'\in\mathfrak{J}$, it is clear that $\mathfrak{J}^2=0$.
\end{env}

\begin{env}[18.2.2]
\label{0.18.2.2}
One calls \emph{$A$-extension of an $A$-ring $B$ by a $B$-bimodule $L$} an \emph{exact} sequence of homomorphisms of $A$-bimodules
\[
  0\to L\xrightarrow{j} E\xrightarrow{f} B\to 0
\]
where $E$ is an $A$-ring, $f$ an $A$-homomorphism of rings and one has, for $x\in E$ and $z\in L$,
\[
  j(f(x)z)=xj(z),\quad j(zf(x))=j(z)x
\]
hence \sref{0.18.2.1} $j(L)$ is a bilateral ideal of square zero of $E$.
By abuse of language one also says that $E$ is an extension of $B$ by $L$.
One says that two $A$-extensions $E$, $E'$ of $B$ by $L$ are \emph{$A$-equivalent} if there exists an isomorphism of $A$-rings $u:E\xrightarrow{\sim} E'$ making the diagram
\begin{equation}
\tag{18.2.2.1}
  \label{0.18.2.2.1}
  \xymatrix@C=2.1pc@R=1.2pc{
    & & E\ar[dr]\ar[dd]_u\ar@{}[dd]|{\wr} & & \\
    0\ar[r] & L\ar[ur]\ar[dr] & & B\ar[r] & 0\\
    & & E'\ar[ur] & &
  }
\end{equation}
commutative.
\end{env}

\begin{env}[18.2.3]
\label{0.18.2.3}
\oldpage[0\textsubscript{IV-1}]{55}
One says that an $A$-extension $E$ of $B$ by a $B$-bimodule $L$ is \emph{$A$-trivial} if $E$ is a trivial augmented $A$-ring (over $B$) \sref{0.18.1.4}.
One defines on the $A$-bimodule product $B\times L$ a structure of $A$-ring by setting $(x,s)(y,t)=(xy,xt+sy)$, as one easily checks, and the canonical maps $j:L\to B\times L$ and $f:B\times L\to B$ define an $A$-extension of $B$ by $L$ that is $A$-trivial, the canonical map $g:B\to B\times L$ being an $A$-homomorphism right inverse of $f$.
One says this extension is the \emph{trivial type extension} of $B$ by $L$ and one denotes it $D_B(L)$; it is immediate that every $A$-trivial $A$-extension of $B$ by $L$ is $A$-equivalent to $D_B(L)$.

One will note that every extension of the $A$-ring $A$ itself by an $A$-bimodule is necessarily $A$-trivial.
\end{env}

\begin{env}[18.2.4]
\label{0.18.2.4}
Given two $A$-extensions $L\xrightarrow{j} E\xrightarrow{f} B$, $L'\xrightarrow{j'} E'\xrightarrow{f'} B'$, a \emph{morphism} of the first into the second is by definition a triplet of homomorphisms of $A$-bimodules $(u,v,w)$ such that the diagram
\begin{equation}
\tag{18.2.4.1}
  \label{0.18.2.4.1}
  \xymatrix{
    0\ar[r] & L\ar[r]^j\ar[d]_w & E\ar[r]^f\ar[d]_u & B\ar[r]\ar[d]_v & 0\\
    0\ar[r] & L'\ar[r]_{j'} & E'\ar[r]_{f'} & B'\ar[r] & 0
  }
\end{equation}
is commutative, $u$ and $v$ being $A$-homomorphisms of \emph{rings} and $w$ being such that $w(bz)=v(b)w(z)$ and $w(zb)=w(z)v(b)$ for $z\in L$ and $b\in B$ (in other words, the pair $(v,w)$ constitutes a \emph{di-homomorphism} of the $B$-bimodule $L$ into the $B'$-bimodule $L'$);
it is clear that if $(u',v',w')$ is a morphism from $L'\to E'\to B'$ into an $A$-extension $L''\to E''\to B''$, $(u'\circ u,v'\circ v,w'\circ w)$ is still a morphism, which justifies the terminology.

The consideration of the two commutative squares of the diagram \sref{0.18.2.4.1} will lead us to two operations on extensions of $A$-rings.
\end{env}

\begin{env}[18.2.5]
\label{0.18.2.5}
In the first place, consider an $A$-extension $E'$ of $B'$ by $L'$
\[
  0\to L'\xrightarrow{j'} E'\xrightarrow{f'} B'\to 0
\]
and an $A$-homomorphism of rings $v:B\to B'$, and let $F=E'\times_{B'} B$ be the \emph{inverse image} of $E'$ \sref{0.18.1.5}, so that one has a commutative diagram
\[
  \xymatrix{
    0\ar[r] & L_0\ar[r]\ar[d]^{i}\ar@{}[d]|{\wr} & F\ar[r]^{p_2}\ar[d]^{p_1} & B\ar[r]\ar[d]^v & 0\\
    0\ar[r] & L'\ar[r]_{j'} & E'\ar[r]_{f'} & B'\ar[r] & 0
  }
\]
where the two lines are exact, $p_1$ and $p_2$ are canonical homomorphisms and $i$ is bijective; one sees from the definition \sref{0.18.1.2} and \sref{0.18.1.3} that $L_0^2=0$, so that one can consider $F$ as an $A$-extension of $B$ by $L'$, called the \emph{inverse image by $v$} of the extension $E'$ of $B'$ by $L'$ (where $L'$ is naturally considered as a $B$-bimodule by means of the homomorphism of rings $v$).
The functorial character
\oldpage[0\textsubscript{IV-1}]{56}
of the fibre product with respect to each of the factors shows moreover that if one has a morphism of two extensions of $B'$
\[
  \xymatrix{
    0\ar[r] & L'_1\ar[r]\ar[d]^h & E'_1\ar[r]\ar[d]^g & B'\ar[r]\ar[d]^{1_{B'}} & 0\\
    0\ar[r] & L'_2\ar[r] & E'_2\ar[r] & B'\ar[r] & 0
  }
\]
one deduces from it a morphism \emph{inverse image by $v$}
\[
  \xymatrix{
    0\ar[r] & L'_1\ar[r]\ar[d]^h & E'_1\times_{B'}B\ar[r]\ar[d]^{g\times_{B'}1_B} & B\ar[r]\ar[d]^{1_B} & 0\\
    0\ar[r] & L'_2\ar[r] & E'_2\times_{B'}B\ar[r] & B\ar[r] & 0.
  }
\]

In particular, if $E'_1$ and $E'_2$ are $A$-equivalent $A$-extensions of $B'$ by $L'$, their inverse images by $v$ are $A$-equivalent $A$-extensions of $B$ by $L'$.

The definition of the fibre product shows that when one has a morphism \sref{0.18.2.4.1} of $A$-extensions, it \emph{factors through} the inverse image of $E'$ by $v$; more precisely, there exists a unique $A$-homomorphism $u_0:E\to F=E'\times_{B'}B$ making the diagram
\[
  \xymatrix{
    0\ar[r] & L\ar[r]^j\ar[d]^{w_0} & E\ar[r]^f\ar[d]^{u_0} & B\ar[r]\ar[d]^{1_B} & 0\\
    0\ar[r] & L_0\ar[r]\ar[d]^i\ar@{}[d]|{\wr} & F\ar[r]^{p_2}\ar[d]^{p_1} & B\ar[r]\ar[d]^v & 0\\
    0\ar[r] & L'\ar[r]_{j'} & E'\ar[r]_{f'} & B'\ar[r] & 0
  }
\]
commutative, where $w_0$ is the restriction of $u_0$ to $L$ and $p_1\circ u_0=u$, $i\circ w_0=w$.
\end{env}

\begin{env}[18.2.6]
\label{0.18.2.6}
Let us study the inverse image of extensions by surjective $A$-homomorphisms.
Consider a surjective $A$-homomorphism $v:B\to B'$, an $A$-extension $F$ of $B$ by a $B$-bimodule $L$, and the bilateral ideal $\mathfrak{R}$ of $B$, kernel of $v$ (which can be considered as an $F$-bimodule by means of the augmentation $F\to B$); one has seen \sref{0.18.1.7} that every homomorphism of $F$-bimodules $\theta:\mathfrak{R}\to F$ making the diagram
\[
  \xymatrix{
    & \mathfrak{R}\ar@/^/[ld]^\theta\ar[d]\\
    F\ar[r] & B
  }
\]
commutative determines an extension $E'$ of $B'=B/\mathfrak{R}$ by $L$ whose inverse image by $v:B\to B'$ is $A$-equivalent to $F$, and that every $A$-extension $E$ of $B'$ by $L$ having this last property is obtained in this way (up to $A$-equivalence).
Moreover, for two homomorphisms $\theta_1$, $\theta_2$ of $F$-bimodules of $\mathfrak{R}$ into $F$ to give two $A$-equivalent $A$-extensions of $B'$ by $L$, it is necessary and sufficient that there exists an $A$-\emph{equivalence} $u$ of the $A$-extension $F$ onto itself such that
\oldpage[0\textsubscript{IV-1}]{57}
$\theta_2=u\circ\theta_1$; this follows at once from what was seen in \sref{0.18.2.5} and from the definition of the canonical bijection of $F$ onto the fibre product $E'\times_{B'}B$ \sref{0.18.1.7}.
\end{env}

\begin{env}[18.2.7]
\label{0.18.2.7}
Consider now the left square of \sref{0.18.2.4.1}, and recall first the notion of \emph{amalgamated sum} in the category of $A$-bimodules: given three $A$-bimodules $X$, $Y$, $Z$ and two $A$-homomorphisms $f:X\to Y$, $g:X\to Z$, the amalgamated sum $Y\oplus_X Z$ is the \emph{inductive limit} of the inductive system formed by $X$, $Y$, $Z$ and the $A$-homomorphisms $f$, $g$ in the category of $A$-bimodules ($\mathbf{0_{III}}$, 8.1.11).
One defines this $A$-bimodule as quotient of the product $Y\times Z$ by the sub-$A$-bimodule $M$ image of $X$ by the homomorphism $x\mapsto(f(x),-g(x))$.
Its characteristic property is that, for every pair of homomorphisms of $A$-bimodules $u:Y\to T$, $v:Z\to T$ such that $u\circ f=v\circ g$, there exists a unique homomorphism $w:Y\oplus_X Z\to T$ such that $u=w\circ j_1$ and $v=w\circ j_2$, where $j_1:Y\to Y\oplus_X Z$ and $j_2:Z\to Y\oplus_X Z$ are the canonical maps.
\end{env}

\begin{env}[18.2.8]
\label{0.18.2.8}
Consider now an $A$-extension $E$ of $B$ by a $B$-bimodule $L$:
\[
  0\to L\xrightarrow{j} E\xrightarrow{f} B\to 0
\]
and let moreover $w:L\to L'$ be a homomorphism of $B$-\emph{bimodules}.
Let $H$ be the $A$-bimodule amalgamated sum $E\oplus_L L'$; let us show how one can equip it with a structure of $A$-\emph{ring} and define an $A$-\emph{extension}
\[
  0\to L'\to H\to B\to 0.
\]

Let us note for this that $L'$ is equipped with a structure of $E$-\emph{bimodule} by means of the augmentation homomorphism $E\to B$; one can therefore form the $A$-extension of trivial type $G=D_E(L')$ \sref{0.18.2.3}.
Consider then the map $\theta:z\mapsto(j(z),-w(z))$ from $L$ into $G$; it is a homomorphism of $G$-\emph{bimodules} ($L$ being considered as a $G$-bimodule by means of the canonical homomorphism $\rho:G\to E$).
Indeed, for $(x,z')\in G$ and $z\in L$, one has $(j(z),-w(z))(x,z')=(j(z)x,j(z)z'-w(z)f(x))$; but $j(z)z'=f(j(z))z'=0$ by definition of the structure of $E$-bimodule on $L'$, and $w(z)f(x)=w(zf(x))$ by definition of the structure of $B$-bimodule on $L'$; one thus verifies that $\theta$ is also a homomorphism of $G$-modules on the left.
One can then apply the result of \sref{0.18.1.7} to the diagram
\[
  \xymatrix@C=2pc@R=1.4pc{
    & & 0\ar[d] & & \\
    & & L\ar[ld]^\theta\ar[d]^j & & \\
    0\ar[r] & L'\ar[r] & G\ar[r]_p & E\ar[r]\ar[d]^f & 0\\
    & & & B\ar[d]\\
    & & & 0
  }
\]
As $H=G/\theta(L)$ by definition, our assertion is an immediate consequence of \sref{0.18.1.7}.

\oldpage[0\textsubscript{IV-1}]{58}

One says that the $A$-extension $H$ of $B$ by $L'$ is \emph{deduced from $E$ by the homomorphism} $w:L\to L'$.
The functorial character of the amalgamated sum in each of the summands shows moreover that if one has a morphism of extensions
\[
  \xymatrix{
    0\ar[r] & L\ar[r]\ar[d]^{1_L}\ar@{}[d]|{\wr} & E_1\ar[r]\ar[d]^g & B_1\ar[r]\ar[d]^h & 0\\
    0\ar[r] & L\ar[r] & E_2\ar[r] & B_2\ar[r] & 0
  }
\]
one deduces from it canonically a morphism of deduced extensions
\[
  \xymatrix{
    0\ar[r] & L'\ar[r]\ar[d]^{1_{L'}}\ar@{}[d]|{\wr} & E_1\oplus_L L'\ar[r]\ar[d]^{g\oplus_L1_{L'}} & B_1\ar[r]\ar[d]^h & 0\\
    0\ar[r] & L'\ar[r] & E_2\oplus_L L'\ar[r] & B_2\ar[r] & 0
  }
\]

In particular, if $E_1$ and $E_2$ are $A$-equivalent $A$-extensions of $B$ by $L$, the extensions of $B$ by $L'$ obtained by $w$ are $A$-equivalent.

When one has a morphism \sref{0.18.2.4.1} of $A$-extensions, it factors through the $A$-extension $H$ of $B$ by $L'$ deduced from $E$ by the homomorphism $w:L\to L'$ (where $L'$ being considered as $B$-bimodule by means of the homomorphism of rings $v:B\to B'$):
in fact, the definition of the amalgamated sum shows that there exists a unique $A$-homomorphism $u_0:H\to E'$ of $A$-bimodules, making the diagram
\[
  \xymatrix{
    0\ar[r] & L\ar[r]^j\ar[d]^{w_0} & E\ar[r]^f\ar[d]^{j_1} & B\ar[r]\ar[d]^{1_B}\ar@{}[d]|{\wr} & 0\\
    0\ar[r] & L'_0\ar[r]^{j_2}\ar[d]^{j_0}\ar@{}[d]|{\wr} & H\ar[r]^{f_0}\ar[d]^{u_0} & B\ar[r]\ar[d]^v & 0\\
    0\ar[r] & L'\ar[r]_{j'} & E'\ar[r]_{f'} & B'\ar[r] & 0
  }
\]
commutative, with $u=u_0\circ j_1$, $w=j_0\circ w_0$, $j_1$ and $j_2$ being the canonical homomorphisms; one verifies immediately that $u_0$ is also a homomorphism of rings.

Let us note moreover the functorial properties relative to trivial extensions:
\end{env}

\begin{proposition}[18.2.9]
\label{0.18.2.9}
Let $B$, $B'$ be two $A$-rings, $L$ a $B$-bimodule, $L'$ a $B'$-bimodule, $v:B\to B'$ an $A$-homomorphism of rings, $w:L\to L'$ a homomorphism of $A$-bimodules such that $(v,w)$ is a di-homomorphism of bimodules.
Then there exists a unique $A$-homomorphism of rings $u:D_B(L)\to D_{B'}(L')$ making the diagrams commutative
\[
  \xymatrix{
    D_B(L)\ar[r]^u & D_{B'}(L') & & D_B(L)\ar[r]^u & D_{B'}(L')\\
    B\ar[u]\ar[r]_v & B'\ar[u] & & L\ar[u]\ar[r]_w & L'\ar[u]
  }
\]
where the vertical arrows are the canonical injections.
\end{proposition}

\oldpage[0\textsubscript{IV-1}]{59}

\begin{proof}
Indeed, $u$ cannot be other than the map $(x,s)\mapsto(v(x),w(s))$ and it remains to verify that it is an $A$-homomorphism of rings, which follows trivially from the definition \sref{0.18.2.3}.
One notes that $u$ also makes commutative the diagram
\[
  \xymatrix@C=2.2pc@R=1.5pc{
    D_B(L)\ar[r]^u\ar[d] & D_{B'}(L')\ar[d]\\
    B\ar[r]_v & B'
  }
\]
where this time the vertical arrows are the augmentations.
\end{proof}

\begin{proposition}[18.2.10]
\label{0.18.2.10}
Let $B$ be an $A$-ring, $L$ a $B$-bimodule, and $E$ an $A$-extension of $B$ by $L$.
There is a bijection from the set $G$ of homomorphisms of $B$-rings from $B$ into $E$ (in other words, the set of $A$-homomorphisms right inverse to the augmentation $E\to B$) onto the set $G'$ of $A$-equivalences from $D_B(L)$ onto $E$: to every $g\in G$ it assigns the $A$-equivalence $g':(x,s)\mapsto g(x)+s$, and the inverse map assigns to every $g'\in G'$ the $B$-homomorphism $x\mapsto g'(x,0)$.
\end{proposition}

\begin{proof}
This follows immediately from the definitions.
\end{proof}

\subsection{The group of classes of A-extensions}
\label{subsection:0.18.3}

\begin{env}[18.3.1]
\label{0.18.3.1}
Consider a fixed $A$-ring $B$ and a fixed $B$-bimodule $L$; then the relation ``$E$ and $E'$ are $A$-equivalent'' between $A$-extensions $E$, $E'$ of $B$ by $L$ is an equivalence relation and for this relation one can speak of the \emph{set} of classes of $A$-equivalent $A$-extensions of $B$ by $L$.
To see this, it suffices to remark that, for every $x\in B$, if $c_x$ is an element of $E$ whose image in $B$ is $x$, every $z\in E$ can be written in one and only one way in the form $c_x+t$, where $t\in L$, and one can write $c_{x+y}=c_x+c_y+\varphi(x,y)$, $c_{xy}=c_x c_y+\psi(x,y)$, where $\varphi(x,y)$ and $\psi(x,y)$ are elements of $L$, the maps $\varphi$ and $\psi$ from $B\times B$ into $L$ being the expressions of the fact that $E$ is an $A$-ring, which we shall not write out here.
Every $A$-extension of $B$ by $L$ is therefore $A$-equivalent to an $A$-extension whose underlying set is $B\times L$, whence our conclusion follows immediately.
\end{env}

\begin{env}[18.3.2]
\label{0.18.3.2}
The $A$-ring $B$ being fixed, denote provisionally by $T(L)$ the set of classes of $A$-extensions of $B$ by $L$.
For every $B$-homomorphism of $B$-bimodules $w:L\to L'$, one defines canonically a map $T(w):T(L)\to T(L')$ by associating to the class of an $A$-extension $E$ of $B$ by $L$ the class of the $A$-extension $E\oplus_L L'$ deduced from $E$ by $w$, by virtue of \sref{0.18.2.8}.
If $w':L'\to L''$ is a second homomorphism of $B$-bimodules, one has also
\begin{equation}
\tag{18.3.2.1}
  \label{0.18.3.2.1}
  T(w'\circ w)=T(w')\circ T(w).
\end{equation}

Indeed, one knows that there exists a canonical isomorphism of $A$-bimodules
\[
  E\oplus_L L''\xrightarrow{\sim}(E\oplus_L L')\oplus_{L'} L''
\]
by virtue of the general properties of inductive limits (cf.\ for example ($\mathbf{I}$, 3.3.9)), and it is immediate to verify that it is indeed an $A$-equivalence of $A$-extensions, hence

\oldpage[0\textsubscript{IV-1}]{60}

relation (18.3.2.1). If $\mathbf C_B$ denotes the category of $B$-bimodules, one sees that one has thus defined a \emph{covariant functor} $T:\mathbf C_B\to\mathbf{Sets}$.
\end{env}

\begin{env}[18.3.3]
\label{0.18.3.3}
Consider now a family $(L_\alpha)_{\alpha\in I}$ of $B$-bimodules and their product $L=\prod_\alpha L_\alpha$; the projections $\operatorname{pr}_\alpha:L\to L_\alpha$ define applications
\[
  T(\operatorname{pr}_\alpha):T(L)\to T(L_\alpha),
\]
hence a canonical map
\begin{equation}
\tag{18.3.3.1}
  \label{0.18.3.3.1}
  \prod_\alpha T(\operatorname{pr}_\alpha):T(L)\to\prod_{\alpha\in I} T(L_\alpha).
\end{equation}

We shall see that this map is \emph{bijective}.
Indeed, for every $\alpha\in I$, let $E_\alpha$ be an $A$-extension of $B$ by $L_\alpha$, and let $F$ be the \emph{fibre product} of the $E_\alpha$ on $B$ \sref{0.18.1.3}; it is immediate that $F$ is an $A$-extension of $B$ by $L=\prod_\alpha L_\alpha$, and that if one replaces each $E_\alpha$ by an $A$-equivalent $A$-extension $E'_\alpha$, the fibre product $F'$ of the $E'_\alpha$ on $B$ is $A$-equivalent to $F$.
One has thus defined a map $\prod_\alpha T(L_\alpha)\to T(L)$ and it is clear that this map is the inverse of \sref{0.18.3.3.1}, hence our assertion.
One will identify canonically $T(L)$ with $\prod_\alpha T(L_\alpha)$ by the preceding map.
One verifies also immediately that if $(L'_\alpha)_{\alpha\in I}$ is a second family of $B$-bimodules and, for each $\alpha$, $w_\alpha:L_\alpha\to L'_\alpha$ a homomorphism of $B$-bimodules, then, setting $w=\prod_\alpha w_\alpha:L\to L'=\prod_\alpha L'_\alpha$, $T(w)$ identifies with $\prod_\alpha T(w_\alpha)$ when one makes the preceding identification.
\end{env}

\begin{env}[18.3.4]
\label{0.18.3.4}
For a $B$-bimodule $L$, the addition is a homomorphism $s:L\times L\to L$ of $B$-\emph{bimodules}, and so is the symmetry $t:L\to L$ of the additive law of $L$.
One deduces from them a law of composition
\[
  T(s):T(L)\times T(L)\to T(L)
\]
in $T(L)$ by virtue of \sref{0.18.3.3}, and this law is a law of \emph{commutative group} of which $T(t)$ is the symmetry, as it follows from the definition of an object in groups by means of commutative diagrams ($\mathbf 0_{\mathrm{III}}$, 8.2.5 and 8.2.6).
We shall denote by $\operatorname{Exan}_A(B,L)$ the commutative group thus defined and we say it is the \emph{group of classes of $A$-extensions of $B$ by $L$}.
\end{env}

\begin{env}[18.3.5]
\label{0.18.3.5}
Denote by $\boldsymbol{K}$ the category whose objects are the triplets $(A,B,L)$ where $A$ is a ring, $B$ an $A$-ring and $L$ a $B$-bimodule; the \emph{morphisms} of this category are the triplets $(u,v,w)$ where $u:A'\to A$ and $v:B'\to B$ are two homomorphisms of rings making commutative the left square of the diagram
\[
  \xymatrix{
    A\ar[r] & B & L\ar[d]^w\\
    A'\ar[u]^u\ar[r] & B'\ar[u]^v & L'
  }
\]
where the horizontal arrows are the structural homomorphisms; finally $w:L\to L'$ is a homomorphism of commutative groups such that one has $w(v(b')z)=b'w(z)$ and

\oldpage[0\textsubscript{IV-1}]{61}

$w(zv(b'))=w(z)b'$ for all $z\in L$ and $b'\in B'$ (in other words, $w$ is a homomorphism of $B'$-bimodules when $L$ is equipped with its structure of $B'$-bimodule defined by $v$).
The composition of morphisms is defined by $(u',v',w')\circ(u,v,w)=(u\circ u',v\circ v',w'\circ w)$, which is immediately verified.
We propose to show that
\begin{equation}
\tag{18.3.5.1}
  \label{0.18.3.5.1}
  (A,B,L)\mapsto\operatorname{Exan}_A(B,L)
\end{equation}
is a \emph{covariant functor} from the category $\boldsymbol{K}$ into the category $\boldsymbol{Ab}$ of commutative groups.
It is therefore a matter of defining for every triplet $(u,v,w)$ as above, a homomorphism of commutative groups
\[
  (u,v,w)_*:\operatorname{Exan}_A(B,L)\to\operatorname{Exan}_{A'}(B',L').
\]

By virtue of the definition of morphisms in $\boldsymbol{K}$, one can write
\[
  (u,v,w)=(1_{A'},1_{B'},w)\circ(1_{A'},v,1_L)\circ(u,1_B,1_L)
\]
where, in the first factor, $L$ is equipped with its structure of $B'$-bimodule defined by $v$; we shall therefore define $(u,v,w)_*$ when two of the three homomorphisms $u$, $v$, $w$ are reduced to the identity.
\end{env}

\begin{env}[18.3.6]
\label{0.18.3.6}
We take first $(1_A,1_B,w)_*$ the map
\begin{equation}
\tag{18.3.6.1}
  \label{0.18.3.6.1}
  w_*:\operatorname{Exan}_A(B,L)\to\operatorname{Exan}_A(B,L')
\end{equation}
denoted $T(w)$ in \sref{0.18.3.2}; it is immediate to verify that it is a homomorphism of groups, this property expressing the commutativity of diagrams, transformed by $T$ of analogous diagrams for $L$ and $L'$.

The map $(1_A,v,1_L)_*$ is the map
\begin{equation}
\tag{18.3.6.2}
  \label{0.18.3.6.2}
  v^*:\operatorname{Exan}_A(B,L)\to\operatorname{Exan}_A(B',L)
\end{equation}
defined in the following way: if $E$ is an $A$-extension of $B$ by $L$, one has seen that $E\times_B B'$ is an $A$-extension of $B'$ by $L$ \sref{0.18.2.5} and that if one replaces $E$ by an $A$-equivalent $A$-extension $E'$, $E'\times_B B'$ is $A$-equivalent to $E\times_B B'$; the image by $v^*$ of the class of $E$ is the class of $E\times_B B'$.
One verifies immediately that if $w:L\to L'$ is a homomorphism of $B$-bimodules, the diagram
\begin{equation}
\tag{18.3.6.3}
  \label{0.18.3.6.3}
  \xymatrix{
    \operatorname{Exan}_A(B,L)\ar[r]^{v^*}\ar[d]_{w_*} & \operatorname{Exan}_A(B',L)\ar[d]^{w_*}\\
    \operatorname{Exan}_A(B,L')\ar[r]_{v^*} & \operatorname{Exan}_A(B',L')
  }
\end{equation}
is commutative, $L$ and $L'$ being considered as $B'$-bimodules by means of $v$ in the right column.
Replacing $L$ and $L'$ by $L\times L$ and $L$, and $w$ by the addition $s$ in $L$, one concludes as above that $v^*$ is indeed a \emph{homomorphism of groups}.
\end{env}

\oldpage[0\textsubscript{IV-1}]{62}

\begin{env}[18.3.6.4]
\label{0.18.3.6.4}
Finally, the map $(u,1_B,1_L)_*$ is the map
\begin{equation}
  u^*:\operatorname{Exan}_A(B,L)\to\operatorname{Exan}_{A'}(B,L)
\end{equation}
which is obtained by associating to an $A$-extension $E$ of $B$ by $L$ the ring $E$ considered as an $A'$-\emph{ring} by means of $u$ \sref{0.18.1.1}, which is evidently also an $A'$-extension of $B$ by $L$, $B$ being also considered as $A'$-ring by means of $u$; it is clear that an $A$-equivalence is also an $A'$-equivalence, whence the map \sref{0.18.3.6.4}, which, for every homomorphism $w:L\to L'$ of $B$-bimodules, also makes the diagram commutative:
  \begin{equation}
  \tag{18.3.6.5}
  \label{0.18.3.6.5}
  \xymatrix{
    \operatorname{Exan}_A(B,L)\ar[r]^{u^*}\ar[d]_{w_*} & \operatorname{Exan}_{A'}(B,L)\ar[d]^{w_*}\\
    \operatorname{Exan}_A(B,L')\ar[r]_{u^*} & \operatorname{Exan}_{A'}(B,L')
  }
  \end{equation}
One deduces as above that $u^*$ is indeed a homomorphism of groups.

This being so, one sets $(u,v,w)_*=(1_{A'},1_{B'},w)_*\circ(1_{A'},v,1_L)_*\circ(u,1_B,1_L)_*$ and one verifies easily, by reason of the commutativity of the diagrams \sref{0.18.3.6.3} and \sref{0.18.3.6.5}, that one has $(u\circ u',v\circ v',w'\circ w)_*=(u',v',w')_*\circ(u,v,w)_*$, which finishes the proof that \sref{0.18.3.5.1} is a functor.

The existence of the homomorphism of groups \sref{0.18.3.6.1} shows in particular that if $E$ is a trivial $A$-extension of $B$ by $L$, the extension $E\oplus_L L'$ of $B$ by $L'$ defined in \sref{0.18.2.8} is also trivial, which one moreover verifies directly without difficulty.

On the other hand, for every element $z$ of the \emph{centre} $Z$ of $B$, the homothety $h_z:y\mapsto yz$ is an endomorphism of the $B$-bimodule $L$, hence $(h_z)_*$ is an endomorphism of the commutative group $\operatorname{Exan}_A(B,L)$, and by functoriality, these endomorphisms define on $\operatorname{Exan}_A(B,L)$ a structure of $Z$-\emph{module}.
\end{env}

\begin{env}[18.3.7]
\label{0.18.3.7}
Let $A$, $A'$ be two rings, $u:A'\to A$ a homomorphism, $B$ an $A$-ring and $L$ a $B$-bimodule.
The \emph{kernel} of the homomorphism of groups
\[
  u^*:\operatorname{Exan}_A(B,L)\to\operatorname{Exan}_{A'}(B,L)
\]
is formed by definition of classes of $A$-extensions of $B$ by $L$ that are $A'$-trivial when one considers them as $A'$-extensions by means of $u$.
One denotes this kernel by the notation $\operatorname{Exan}_{A/A'}(B,L)$ when this causes no confusion.

If $\Lambda$ is a ring, and if one has a commutative diagram of $\Lambda$-\emph{homomorphisms} of $\Lambda$-rings
\begin{equation}
\tag{18.3.7.1}
  \label{0.18.3.7.1}
  \xymatrix{
    B'\ar[r] & B\\
    A'\ar[u]\ar[r] & A\ar[u]
  }
\end{equation}
\oldpage[0\textsubscript{IV-1}]{63}
one deduces from it canonically homomorphisms
\begin{equation}
\tag{18.3.7.2}
  \label{0.18.3.7.2}
  \operatorname{Exan}_{A/\Lambda}(B,L)\to\operatorname{Exan}_{A'/\Lambda}(B,L)\to\operatorname{Exan}_{A'/\Lambda}(B',L)
\end{equation}
which arise from the commutativity of the diagram
\[
  \xymatrix@C=2.8pc@R=2.3pc{
    \operatorname{Exan}_A(B,L)\ar[r]\ar[d]
      & \operatorname{Exan}_{A'}(B,L)\ar[r]\ar[d]
      & \operatorname{Exan}_{A'}(B',L)\ar[d]\\
    \operatorname{Exan}_{\Lambda}(B,L)\ar[r]
      & \operatorname{Exan}_{\Lambda}(B,L)\ar[r]
      & \operatorname{Exan}_{\Lambda}(B',L)
  }
\]
where the arrows are deduced from those of \sref{0.18.3.7.1} by functoriality.
\end{env}

\begin{proposition}[18.3.8]
\label{0.18.3.8}
Let $B$ be a ring, $\mathfrak{J}$ a bilateral ideal of $B$, and $C=B/\mathfrak{J}$ the quotient ring; $\mathfrak{J}/\mathfrak{J}^2$ is then canonically equipped with a structure of $C$-bimodule.
For every $C$-bimodule $L$, let $\operatorname{Hom}_C(\mathfrak{J}/\mathfrak{J}^2,L)$ denote the additive group of homomorphisms of $C$-bimodules from $\mathfrak{J}/\mathfrak{J}^2$ into $L$.
One defines then an isomorphism of commutative groups
\begin{equation}
\tag{18.3.8.1}
  \label{0.18.3.8.1}
  \eta_L:\operatorname{Hom}_C(\mathfrak{J}/\mathfrak{J}^2,L)\xrightarrow{\sim}\operatorname{Exan}_B(C,L)
\end{equation}
by associating to every $C$-homomorphism $w:\mathfrak{J}/\mathfrak{J}^2\to L$ (which is \emph{a fortiori} a $B$-homomorphism) the class of the extension $(B/\mathfrak{J}^2)\oplus_{(\mathfrak{J}/\mathfrak{J}^2)} L$ deduced from the extension $B/\mathfrak{J}^2$ of $C$ by $\mathfrak{J}/\mathfrak{J}^2$ by means of $w$; the inverse isomorphism associates to a class of $B$-extension $E$ of $C$ by $L$ the homomorphism $w$ such that the composite $\mathfrak{J}\to\mathfrak{J}/\mathfrak{J}^2\xrightarrow{w} L$ is the restriction to $\mathfrak{J}$ of the structural homomorphism $B\to E$.
\end{proposition}

\begin{proof}
Let $F=B/\mathfrak{J}^2$, which is a $B$-extension of $C$ by $\mathfrak{J}/\mathfrak{J}^2$.
For every $B$-extension
$0\to L\xrightarrow{j} E\xrightarrow{p} C\to 0$ of $C$ by $L$, the structural homomorphism $f:B\to E$ is such that the composite $B\xrightarrow{f} E\xrightarrow{p} C$ is the canonical homomorphism $B\to B/\mathfrak{J}=C$.
As the image of $\mathfrak{J}$ by $p\circ f$ is zero, $f(\mathfrak{J})$ is contained in the kernel $j(L)$, and as $j(L)$ is of square zero, $f(\mathfrak{J}^2)=0$; hence $f$ factors as $B\to B/\mathfrak{J}^2=F\xrightarrow{u} E$, and if $w$ is the restriction of $u$ to $\mathfrak{J}/\mathfrak{J}^2$, one has a commutative diagram
\[
  \xymatrix{
    0\ar[r] & \mathfrak{J}/\mathfrak{J}^2\ar[r]\ar[d]^w & F\ar[r]^g\ar[d]^u & C\ar[r]\ar[d]^{1_C}\ar@{}[d]|{\wr} & 0\\
    0\ar[r] & L\ar[r]_j & E\ar[r]_p & C\ar[r] & 0
  }
\]
in other words $(u,1_C,w)$ is a morphism of extensions \sref{0.18.2.4}.
One deduces from it a morphism of extensions
\[
  \xymatrix{
    0\ar[r] & L\ar[r]\ar[d]^{1_L}\ar@{}[d]|{\wr} & E'\ar[r]\ar[d]^{u'} & C\ar[r]\ar[d]^{1_C}\ar@{}[d]|{\wr} & 0\\
    0\ar[r] & L\ar[r]_j & E\ar[r]_p & C\ar[r] & 0
  }
\]

\oldpage[0\textsubscript{IV-1}]{64}

where $E'=F\oplus_{(\mathfrak J/\mathfrak J^2)}L$ \sref{0.18.2.8}; it remains to prove that $u'$ is a $B$-\emph{equivalence}, which will establish that $\eta_L$ is bijective.
By the construction in \sref{0.18.2.8}, it suffices (in view of \sref{0.18.1.7}) to prove that the inverse image $G$ of the extension $E$ of $C$ by the canonical homomorphism $g:F\to C$ is $F$-\emph{trivial}, which is evident since $g$ factors as $F\xrightarrow{u} E\xrightarrow{p} C$ \sref{0.18.1.6}.

It remains to see that $\eta_L$ is a homomorphism of \emph{groups}; now, for every $B$-homomorphism $h:L\to L'$, it is immediate that the diagram
\[
  \xymatrix{
    \operatorname{Hom}_C(\mathfrak{J}/\mathfrak{J}^2,L)\ar[r]^{\eta_L}\ar[d]^{\operatorname{Hom}(1,h)} & \operatorname{Exan}_B(C,L)\ar[d]^{h_*}\\
    \operatorname{Hom}_C(\mathfrak{J}/\mathfrak{J}^2,L')\ar[r]_{\eta_{L'}} & \operatorname{Exan}_B(C,L')
  }
\]
is commutative.
It suffices to apply this remark to the homomorphism $L\times L\to L$ defining the addition to conclude.
\end{proof}

\subsection{Extensions of algebras}
\label{subsection:0.18.4}

\begin{env}[18.4.1]
\label{0.18.4.1}
Let $A$ be a \emph{commutative} ring.
The category of $A$-\emph{algebras} is then a full subcategory of that of $A$-rings, characterized by the fact that the structural homomorphisms $\rho:A\to B$ are \emph{central}.

If $B$ is an $A$-algebra and $L$ a $B$-bimodule whose two $A$-module structures induced through $B$ coincide, it is immediate that the $A$-extension of trivial type $D_B(L)$ is an $A$-algebra.
Similarly, in the construction of \sref{0.18.2.5} (resp.\ of \sref{0.18.2.8}), if $B$, $B'$ and $E'$ (resp.\ $B$ and $E$) are $A$-algebras, it is the same for $E'\times_{B'} B$ (resp.\ for $E\oplus_L L'$).
Finally, if $B$ is an $A$-algebra and $E$ an $A$-extension of $B$ by a $B$-bimodule $L$ that is an $A$-algebra, every $A$-extension $E'$ of $B$ by $L$, $A$-\emph{equivalent} to $E$, is also an $A$-algebra.
One deduces then immediately from the definition \sref{0.18.3.4} that the classes of equivalent $A$-extensions of $B$ by $L$ which are $A$-\emph{algebras} form a \emph{subgroup}, denoted $\operatorname{Exal}_A(B,L)$, of $\operatorname{Exan}_A(B,L)$.
Let $\boldsymbol{K}'$ be the full subcategory of $\boldsymbol{K}$ defined in \sref{0.18.3.5}, whose objects $(A,B,L)$ are such that $A$ is commutative, $B$ is an $A$-algebra, and the two $A$-module structures on $L$ induced through $B$ coincide.
Then what precedes shows that
\[
  (A,B,L)\mapsto\operatorname{Exal}_A(B,L)
\]
is a \emph{covariant functor} from $\boldsymbol{K}'$ to $\boldsymbol{Ab}$.
The results of \sref{0.18.3.7} are unchanged when one replaces $\operatorname{Exan}$ everywhere by $\operatorname{Exal}$.
\end{env}

\begin{env}[18.4.2]
\label{0.18.4.2}
Suppose always $A$ commutative.
The remarks of \sref{0.18.4.1} remain valid when one replaces ``algebra'' by ``commutative algebra'' and ``bimodule'' by ``module''.
If $B$ is a commutative $A$-algebra and $L$ a $B$-module,

\oldpage[0\textsubscript{IV-1}]{65}

the classes of equivalent $A$-extensions of $B$ by $L$ which are commutative $A$-\emph{algebras} (or, what amounts to the same thing, commutative $A$-\emph{rings}) form a \emph{subgroup} of $\operatorname{Exal}_A(B,L)$, denoted $\operatorname{Exalcom}_A(B,L)$.
If $\boldsymbol{K}''$ is the full subcategory of $\boldsymbol{K}'$ where $A$ and $B$ are commutative and $L$ a $B$-module,
\[
  (A,B,L)\mapsto\operatorname{Exalcom}_A(B,L)
\]
is still a \emph{covariant functor} from $\boldsymbol{K}''$ to $\boldsymbol{Ab}$.
One can also replace $\operatorname{Exan}$ by $\operatorname{Exalcom}$ in \sref{0.18.3.7}.
Finally, if in \sref{0.18.3.8}, one supposes $B$ to be a commutative ring and $L$ a $C$-\emph{module}, the same reasoning gives a canonical isomorphism
\begin{equation}
\tag{18.4.2.1}
  \label{0.18.4.2.1}
  \operatorname{Hom}_C(\mathfrak{J}/\mathfrak{J}^2,L)\xrightarrow{\sim}\operatorname{Exalcom}_B(C,L)
\end{equation}
where the first member is the group of homomorphisms of $C$-\emph{modules}.
\end{env}

\begin{env}[18.4.3]
\label{0.18.4.3}
Let $A$ be a commutative ring.
An important case of extensions of $A$-algebras is formed of $A$-algebras $E$, extensions of an $A$-algebra $B$ by a $B$-bimodule $L$ such that $E$, as $A$-\emph{module}, is the \emph{trivial} extension of the $A$-module $B$ by the $A$-module $L$; in other words, the exact sequence of $A$-\emph{modules} $0\to L\to E\to B\to 0$ is \emph{split}.
One says then that $E$ is a \emph{Hochschild extension} of $B$ by $L$.
This will always be the case when $B$ is a \emph{projective} $A$-module, and in particular when $A$ is a commutative \emph{field}.
As an $A$-module, one can identify $E$ with $B\times L$, the multiplication in $E$ being given by $(x,s)(y,t)=(xy,xt+sy+f(x,y))$ with $f(x,y)\in L$.
If one writes that this multiplication defines on $B\times L$ a structure of $A$-algebra, one finds (M, XIV, 2) that $f$ must be an $A$-\emph{bilinear} map of $B\times B$ into $L$, such that
\begin{equation}
\tag{18.4.3.1}
  \label{0.18.4.3.1}
  f(xy,z)+f(x,y)z=xf(y,z)+f(x,yz)
\end{equation}
in other words, $f$ is a $2$-\emph{cocycle} on $B$ with values in $L$, in the sense of Hochschild; for the extension $E$ to be $A$-trivial, it is necessary and sufficient that
\begin{equation}
\tag{18.4.3.2}
  \label{0.18.4.3.2}
  f(x,y)=xg(y)-g(xy)+g(x)y
\end{equation}
where $g$ is an $A$-\emph{linear} map of $B$ into $L$, in other words $f$ must be a $2$-\emph{coboundary} in the sense of Hochschild.
One deduces from it that the classes of Hochschild extensions of $B$ by $L$ form a subgroup of $\operatorname{Exal}_A(B,L)$, isomorphic to the \emph{Hochschild cohomology group} $\mathrm{H}^2_A(B,L)$.

If $B$ is a commutative $A$-algebra, and $L$ a $B$-module, the commutative Hochschild extensions of $B$ by $L$ correspond to the \emph{symmetric} $2$-cocycles, i.e.\ those such that $f(y,x)=f(x,y)$.
The classes of commutative Hochschild extensions of $B$ by $L$ form therefore a subgroup of $\operatorname{Exalcom}_A(B,L)$, isomorphic to the subgroup of $\mathrm{H}^2_A(B,L)$ image of the group of symmetric $2$-cocycles, which we shall denote $\mathrm{H}^2_A(B,L)^s$.
\end{env}

\begin{env}[18.4.4]
\label{0.18.4.4}
Let us restrict to the case where $A$ and $B$ are \emph{commutative} and $L$ a $B$-module, and recall in this case the equivalent definition of the Hochschild cohomology groups $\mathrm{H}^i_A(B,L)$ (M, IX, 6).
One considers a complex $\mathrm{P}_\bullet=(\mathrm{P}_n)_{n\geqslant 0}$ of $B$-modules, where $\mathrm{P}_n=B^{\otimes(n+1)}=B\otimes_A B\otimes\cdots\otimes_A B$ ($n+1$ times), the structure of $B$-module being defined

\oldpage[0\textsubscript{IV-1}]{66}

by $x(y_1\otimes y_2\otimes\cdots\otimes y_{n+1})=(xy_1)\otimes y_2\otimes\cdots\otimes y_{n+1}$; the differential, of degree $-1$, $d_n:\mathrm{P}_n\to\mathrm{P}_{n-1}$, is defined by
\begin{align*}
  d_n(x_1\otimes x_2\otimes\cdots\otimes x_{n+1})&=(x_1 x_2)\otimes x_3\otimes\cdots\otimes x_{n+1}-x_1\otimes(x_2 x_3)\otimes\cdots\otimes x_{n+1}+\cdots\\
  &\quad +(-1)^{n-1} x_1\otimes x_2\otimes\cdots\otimes(x_n x_{n+1})+(-1)^n(x_{n+1} x_1)\otimes x_2\otimes\cdots\otimes x_n
\end{align*}
which is $B$-linear since $B$ is commutative.

A $2$-cocycle of this complex with values in $L$ is a $B$-linear map $h$ from $B\otimes_A B\otimes_A B$ to $L$ such that $h(d_3(x\otimes y\otimes z\otimes t))=0$; but since $h(x\otimes y\otimes z)=xh(1\otimes y\otimes z)$, the cocycle $h$ is determined by the $A$-bilinear map $(y,z)\mapsto f(y,z)=h(1\otimes y\otimes z)$ from $B\times B$ to $L$, and writing the preceding condition for $h$ gives back condition \sref{0.18.4.3.1} for $f$.
Similarly, a $2$-coboundary is a map of the form $x\otimes y\otimes z\mapsto h'(d_2(x\otimes y\otimes z))$, where $h':B\otimes_A B\to L$ is linear; here again $h'$ is determined by the $A$-linear map $g:y\mapsto h'(1\otimes y)$ from $B$ to $L$.
One then obtains $h'(d_2(1\otimes x\otimes y))=xg(y)-g(xy)+yg(x)$, which gives back \sref{0.18.4.3.2}.
One proceeds similarly for the $1$-cocycles and one sees that one has
\begin{equation}
\tag{18.4.4.1}
  \label{0.18.4.4.1}
  \mathrm{H}^i_A(B,L)=\mathrm{H}^i(\operatorname{Hom}_B(\mathrm{P}_\bullet,L)).
\end{equation}
\end{env}

\begin{env}[18.4.5]
\label{0.18.4.5}
Under the conditions of \sref{0.18.4.4}, one can reinterpret in the same way the group $\mathrm{H}^2_A(B,L)^s$ \sref{0.18.4.3}.
For this, modify the complex $\mathrm{P}_\bullet$ in degree $3$ considering a new complex
\[
  \mathrm{P}'_\bullet:\mathrm{P}'_3\xrightarrow{d'_3}\mathrm{P}_2\xrightarrow{d_2}\mathrm{P}_1;
\]
we take $\mathrm{P}'_3=\mathrm{P}_3\oplus(B\otimes_A B\otimes_A B)$, $d'_3$ coinciding with $d_3$ on $\mathrm{P}_3$, and being given by
\[
  d'_3(x\otimes y\otimes z)=x\otimes y\otimes z-x\otimes z\otimes y.
\]

The relation $d_2 d'_3=0$ follows from the above with the commutativity of $B$.
With the notations introduced above, a $2$-cocycle of $\mathrm{P}'_\bullet$ corresponds now to an $A$-bilinear map $f$ of $B\times B$ into $L$ that is \emph{symmetric} and satisfies \sref{0.18.4.3.1}; consequently, one has $\mathrm{H}^2_A(B,L)^s=\mathrm{H}^2(\operatorname{Hom}_B(\mathrm{P}'_\bullet,L))$.
\end{env}

\begin{env}[18.4.6]
\label{0.18.4.6}
In the particular case where one considers a commutative \emph{field} $k$, an extension $K$ of $k$, one has $\operatorname{Ext}^1_K(M,L)=0$ whatever the $K$-vector spaces $L$, $M$ and consequently (M, VI, 3.3 $a$), one has a canonical isomorphism
\begin{equation}
\tag{18.4.6.1}
  \label{0.18.4.6.1}
  \mathrm{H}^i_k(K,L)\simeq\operatorname{Hom}_K(\mathrm{H}_i(\mathrm{P}_\bullet),L)
\end{equation}
and similarly
\begin{equation}
\tag{18.4.6.2}
  \label{0.18.4.6.2}
  \mathrm{H}^2_k(K,L)^s\simeq\operatorname{Hom}_K(\mathrm{H}_2(\mathrm{P}'_\bullet),L).
\end{equation}
\end{env}

\subsection{Case of topological rings}
\label{subsection:0.18.5}

\begin{env}[18.5.1]
\label{0.18.5.1}
Let $A$ and $B$ be two \emph{topological} rings whose topologies are \emph{linear}, $\rho:A\to B$ a continuous homomorphism, $L$ a topological $B$-bimodule, and suppose that there exists an \emph{open} bilateral ideal $\mathfrak{R}_0$ of $B$ such that $\mathfrak{R}_0.L=L.\mathfrak{R}_0=0$, so that $L$ can be
\oldpage[0\textsubscript{IV-1}]{67}
considered as a $(B/\mathfrak{R})$-bimodule for every open bilateral ideal $\mathfrak{R}\subset\mathfrak{R}_0$.
Let then $\mathfrak{R}\subset\mathfrak{R}_0$ be an open bilateral ideal of $B$ and $\mathfrak{J}$ an open bilateral ideal of $A$ such that $\rho(\mathfrak{J})\subset\mathfrak{R}$, so that $B/\mathfrak{R}$ can be considered as a discrete $(A/\mathfrak{J})$-ring; the preceding remark proves that the group $\operatorname{Exan}_{A/\mathfrak{J}}(B/\mathfrak{R},L)$ is defined.
Moreover, if $\mathfrak{R}'\subset\mathfrak{R}$ and $\mathfrak{J}'\subset\mathfrak{J}$ are open bilateral ideals of $B$ and $A$, respectively, such that $\rho(\mathfrak{J}')\subset\mathfrak{R}'$, one has by \sref{0.18.3.5.1} a canonical homomorphism
\begin{equation}
\tag{18.5.1.1}
  \label{0.18.5.1.1}
  \operatorname{Exan}_{A/\mathfrak{J}}(B/\mathfrak{R},L)\to\operatorname{Exan}_{A/\mathfrak{J}'}(B/\mathfrak{R}',L).
\end{equation}

The set of pairs of open bilateral ideals $(\mathfrak{J},\mathfrak{R})$ such that $\rho(\mathfrak{J})\subset\mathfrak{R}\subset\mathfrak{R}_0$ is evidently directed for the relation ``$\mathfrak{J}\supset\mathfrak{J}'$ and $\mathfrak{R}\supset\mathfrak{R}'$'' and the maps \sref{0.18.5.1.1} define an \emph{inductive system} of additive groups with this set as index set.
One sets, by abuse of notation (since it is no longer a matter of a group of extensions in bijective natural correspondence with a set of extensions)
\begin{equation}
\tag{18.5.1.2}
  \label{0.18.5.1.2}
  \operatorname{Exantop}_A(B,L)=\varinjlim\operatorname{Exan}_{A/\mathfrak{J}}(B/\mathfrak{R},L).
\end{equation}

To say that the second member of \sref{0.18.5.1.2} is \emph{zero} means therefore that, for every pair of open bilateral ideals $\mathfrak{J}\subset A$, $\mathfrak{R}\subset B$ such that $\rho(\mathfrak{J})\subset\mathfrak{R}\subset\mathfrak{R}_0$ and every $(A/\mathfrak{J})$-extension $E$ of $B/\mathfrak{R}$ by $L$, there exist two open ideals $\mathfrak{J}'\subset\mathfrak{J}$, $\mathfrak{R}'\subset\mathfrak{R}$ such that $\rho(\mathfrak{J}')\subset\mathfrak{R}'$ and the inverse image of $E$ by the homomorphism $B/\mathfrak{R}'\to B/\mathfrak{R}$ is \emph{trivial}.

We leave to the reader the analogous definition of the inductive limits $\operatorname{Exaltop}_A(B,L)$ and $\operatorname{Exalcotop}_A(B,L)$ from $\operatorname{Exal}_{A/\mathfrak{J}}(B/\mathfrak{R},L)$ and $\operatorname{Exalcom}_{A/\mathfrak{J}}(B/\mathfrak{R},L)$ when $A$ is commutative and $B$ is a topological $A$-algebra (resp.\ a commutative topological $A$-algebra).
\end{env}

\begin{env}[18.5.2]
\label{0.18.5.2}
If one has a commutative diagram of continuous homomorphisms of rings
\[
  \xymatrix{
    B'\ar[r] & B\\
    A'\ar[u]\ar[r] & A\ar[u]
  }
\]
one deduces from it canonically two homomorphisms of additive groups
\[
  \operatorname{Exantop}_A(B,L)\to\operatorname{Exantop}_{A'}(B,L)\to\operatorname{Exantop}_{A'}(B',L)
\]
by passage to the inductive limit from \sref{0.18.3.5.1}.

By virtue of the exactness of the $\varinjlim$ functor in the category of commutative groups, the kernel of the homomorphism
\[
  \operatorname{Exantop}_A(B,L)\to\operatorname{Exantop}_{A'}(B,L)
\]
is the inductive limit of the kernels of the homomorphisms
\[
  \operatorname{Exan}_{A/\mathfrak{J}}(B/\mathfrak{R},L)\to\operatorname{Exan}_{A'/\mathfrak{J}'}(B/\mathfrak{R},L)
\]
where one took for $\mathfrak{J}'$ the inverse image of $\mathfrak{J}$; one denotes this kernel $\operatorname{Exantop}_{A/A'}(B,L)$.
One defines similarly $\operatorname{Exaltop}_{A/A'}(B,L)$ and $\operatorname{Exalcotop}_{A/A'}(B,L)$.
Finally, if one has a homomorphism

\oldpage[0\textsubscript{IV-1}]{68}

of $B$-bimodules $L\to L'$, one obtains canonically a homomorphism of additive groups
\[
  \operatorname{Exantop}_A(B,L)\to\operatorname{Exantop}_A(B,L')
\]
by passage to the inductive limit from \sref{0.18.3.6.1}.
\end{env}

\begin{env}[18.5.3]
\label{0.18.5.3}
Given a topological ring $C$ and two topological $C$-bimodules $M$, $N$, one denotes $\operatorname{Hom.\ cont}_C(M,N)$ the additive group of $C$-homomorphisms \emph{continuous} from $M$ into $N$.
\end{env}

\begin{lemma}[18.5.3.1]
\label{0.18.5.3.1}
Let $C$ be a topological ring, $E$, $L$ two topological $C$-bimodules; suppose that the topologies are linear, that $L$ is discrete and annihilated by an open bilateral ideal of $C$.
Then one has a canonical isomorphism
\begin{equation}
\tag{18.5.3.2}
  \label{0.18.5.3.2}
  \varinjlim\operatorname{Hom}_{C/\mathfrak{R}}(E/V,L)\xrightarrow{\sim}\operatorname{Hom.\ cont}_C(E,L)
\end{equation}
where in the first member the inductive limit is taken over the directed set of pairs $(\mathfrak{R},V)$ such that $\mathfrak{R}$ is an open bilateral ideal of $C$, $V$ an open sub-$C$-bimodule of $E$, such that $\mathfrak{R}.L=L.\mathfrak{R}=0$, $\mathfrak{R}.E\subset V$, $E.\mathfrak{R}\subset V$.
\end{lemma}

\begin{proof}
As $C/\mathfrak{R}$ and $E/V$ are discrete, one has canonical homomorphisms
\[
  w_{\mathfrak{R},V}:\operatorname{Hom}_{C/\mathfrak{R}}(E/V,L)
  \longrightarrow\operatorname{Hom.\ cont}_C(E,L)
\]
forming an inductive system, hence a homomorphism from the first member of \sref{0.18.5.3.2} to $\operatorname{Hom.\ cont}_C(E,L)$.
As the homomorphism $E/V'\to E/V$ is surjective for $V'\subset V$, it follows immediately from the definition that
\[
  \operatorname{Hom}_{C/\mathfrak{R}}(E/V,L)\to
  \operatorname{Hom}_{C/\mathfrak{R}}(E/V',L)
\]
is injective whenever $\mathfrak{R}.L=L.\mathfrak{R}=0$, $\mathfrak{R}.E\subset V'$, and $E.\mathfrak{R}\subset V'$.
The same is clearly true of
\[
  \operatorname{Hom}_{C/\mathfrak{R}}(E/V,L)\to
  \operatorname{Hom}_{C/\mathfrak{R}'}(E/V,L)
\]
for $\mathfrak{R}'\subset\mathfrak{R}$; hence the homomorphism \sref{0.18.5.3.2} is injective.
On the other hand, if $u$ is a continuous $C$-homomorphism from $E$ to $L$, its kernel is an open sub-bimodule $V_0$ of $E$.
Choose an open bilateral ideal $\mathfrak{R}_0$ of $C$ such that $\mathfrak{R}_0.L=L.\mathfrak{R}_0=0$, $\mathfrak{R}_0.E\subset V_0$, and $E.\mathfrak{R}_0\subset V_0$.
Then $u$ is the canonical image of a $(C/\mathfrak{R}_0)$-homomorphism from $E/V_0$ to $L$, so \sref{0.18.5.3.2} is surjective.
\end{proof}

This being so, the proposition \sref{0.18.3.8} generalizes as follows to topological rings:

\begin{proposition}[18.5.4]
\label{0.18.5.4}
Let $B$ be a topological ring linearly topologized, $\mathfrak{J}$ a bilateral ideal of $B$, $C=B/\mathfrak{J}$ the quotient topological ring; $\mathfrak{J}/\mathfrak{J}^2$ \textup{(}where $\mathfrak{J}$ is equipped with the topology induced by that of $B$ and $\mathfrak{J}/\mathfrak{J}^2$ with the quotient topology of that of $\mathfrak{J}$\textup{)} is then canonically equipped with a structure of topological $C$-bimodule.
For every discrete $C$-bimodule $L$ annihilated by an open bilateral ideal of $C$, there exists then a canonical isomorphism
\begin{equation}
\tag{18.5.4.1}
  \label{0.18.5.4.1}
  \operatorname{Hom.\ cont}_C(\mathfrak{J}/\mathfrak{J}^2,L)\xrightarrow{\sim}\operatorname{Exantop}_B(C,L).
\end{equation}
\end{proposition}

\begin{proof}
Indeed, for every open $\mathfrak{R}$ of $B$ such that $(\mathfrak{J}+\mathfrak{R})/\mathfrak{J}$ annihilates $L$, one has, by \sref{0.18.3.8}, a canonical isomorphism\footnote{The French source prints ``(18.3.9)'', but \S~18.3 ends at (18.3.8), which is precisely the isomorphism used here.}
\[
  \operatorname{Hom}_{B/(\mathfrak{J}+\mathfrak{R})}((\mathfrak{J}+\mathfrak{R})/(\mathfrak{J}^2+\mathfrak{R}),L)\xrightarrow{\sim}\operatorname{Exan}_{B/\mathfrak{R}}(B/(\mathfrak{J}+\mathfrak{R}),L)
\]
and it suffices to pass to the inductive limit, taking into account \sref{0.18.5.3.1}.
\end{proof}
% END NATIVE INPUT ega0/ega0-18
% BEGIN NATIVE INPUT ega0/ega0-19
\section{Formally smooth algebras and Cohen rings}
\label{section:0.19}
\label{0.19}

\setcounter{subsection}{-1}
\subsection{Introduction}
\label{subsection:0.19.0}
\label{0.19.0}

\begin{env}[19.0.1]
\label{0.19.0.1}
\oldpage[0\textsubscript{IV-1}]{69}
In Chapter~IV, we will introduce and study, among other things, an important class of morphisms of preschemes, the \emph{smooth} morphisms\footnote{Also called \emph{simple} morphisms in certain recent works (inspired by the classical terminology of ``simple point''); this terminology lends itself however to confusion, notably in the theory of algebraic groups.}.
One of their fundamental properties (and which, together with a finiteness condition, characterizes them) is a property of \emph{lifting of morphisms}:
if $f:X\to Y$ is a smooth morphism, $g:Y'\to Y$ a morphism, then for every morphism $h:Y''\to Y'$ making $Y''$ a prescheme ``slightly different'' from $Y'$, every $Y$-morphism $Y''\to X$ factorizes as $Y''\xrightarrow{h}Y'\to X$.
More precisely, let us restrict to the case where $Y=\Spec(A)$, $X=\Spec(B)$ are affine;
$B$ is then called a \emph{formally smooth} $A$-algebra if, for every $A$-algebra $C$ and every nilpotent ideal $\mathfrak{J}$ of $C$, every $A$-homomorphism $B\to C/\mathfrak{J}$ \emph{lifts to} $B\to C\to C/\mathfrak{J}$.
In other words, the map
\[
  \Hom_A(B,C) \to \Hom_A(B,C/\mathfrak{J})
\]
is \emph{surjective}.
In many applications, $X$ will appear as a contravariant representable functor $Y'\mapsto F(Y')$ from the category of $Y$-preschemes to that of sets, so that ($\mathbf{0_{III}}$, 8.1.8) one will have $F(Y')=\Hom_Y(Y',X)$.
In the affine case, setting $F^0(C)=F(\Spec(C))$, the verification of the fact that, under the conditions above, $F^0(C)\to F^0(C/\mathfrak{J})$ is \emph{surjective} (which can be done even without knowing that $F$ is representable) will establish that $f$ is smooth, provided the additional finiteness condition is satisfied.
\end{env}

\begin{env}[19.0.2]
\label{0.19.0.2}
For \emph{topological algebras} over a topological ring $A$, there is an analogous notion of \emph{formally smooth algebra} that we will not specify here (cf.~\sref{0.19.3.1}).
The study of these notions is first done from an elementary point of view in \S~19, then, with the aid of the properties of differentials that will be developed in \S\S~20 and 21, more delicate theorems will be proved in \S~22.
Let us summarize here the main results on formally smooth algebras:

\medskip\noindent\textbf{I.} ---
Let $A$, $B$ be two Noetherian local rings, $\varphi:A\to B$ a local homomorphism, $k$ the residue field of $A$, and let $B_0=B\otimes_A k$;
we equip $A$, $B$, and $B_0$ with their preadic topologies and $k$ with the discrete topology.
Then, for $B$ to be a formally smooth $A$-algebra, it is necessary and sufficient that $B$ be a \emph{flat} $A$-module and that $B_0$ be a formally smooth $k$-algebra \sref{0.19.7.1}.
This theorem thus reduces formal smoothness, for Noetherian local rings, to the same question for Noetherian local rings that are algebras over a field.

\medskip\noindent\textbf{II.} ---
Let $k$ be a field, $A$ a Noetherian local ring that is a $k$-algebra.
For $A$ to be formally smooth, it is necessary and sufficient that $A$ be \emph{geometrically regular} over $k$, that is to say that for every finite extension $k'$ of $k$, the semi-local ring $A\otimes_k k'$ be

\oldpage[0\textsubscript{IV-1}]{70}

regular (22.5.8); the completion $\widehat{A}$ of $A$ is then isomorphic to a power series ring $K[[T_1,\ldots,T_n]]$ \sref{0.19.6.5}.
Moreover, the $k$-algebra structure of $A$, when $A$ is assumed complete and formally smooth over $k$, is entirely determined by the residue field $K$ of $A$ and the dimension of $A$;
the latter can furthermore be arbitrary subject to the inequality $\dim(A)\geq\rg_k(\Gamma_{K/k})$, where $\Gamma_{K/k}$ is the ``module of imperfection'' of $K$ (22.2.6).

In particular, for an extension $K$ of $k$ to be formally smooth, it is necessary and sufficient that $K$ be a \emph{separable} extension of $k$ \sref{0.19.6.1}.

\medskip\noindent\textbf{III.} ---
Let $A$ be a Noetherian local ring, $\mathfrak{J}$ an ideal of $A$ distinct from $A$, $A_0=A/\mathfrak{J}$, $B_0$ a complete Noetherian local ring, $A_0\to B_0$ a local homomorphism making $B_0$ a formally smooth $A_0$-algebra.
Then there exists a complete Noetherian local ring $B$, a local homomorphism $A\to B$ making $B$ a flat $A$-module, and an $A_0$-isomorphism $B\otimes_A A_0\xrightarrow{\sim}B_0$ (so that, by~I), $B$ is a formally smooth $A$-algebra);
moreover $B$ is determined by these properties up to isomorphism \sref{0.19.7.2}.
This theorem contains in particular the \emph{theorems of Cohen} on the structure of complete Noetherian local rings \sref{0.19.8}, which will play an important role in \S\S~6 and 7 of Chapter~IV.

\medskip\noindent\textbf{IV.} ---
The interest in the study of formally smooth Noetherian local rings over another comes from the following ``pointwise'' characterization of smooth morphisms:
if $X$ and $Y$ are locally Noetherian preschemes, $f:X\to Y$ a morphism locally of finite type, then, for $f$ to be smooth, it is necessary and sufficient that for every $x\in X$, the ring $\sh{O}_x$ be formally smooth over $\sh{O}_{f(x)}$.
In particular, if $Y=\Spec(k)$, where $k$ is a perfect field, to say that $f$ is smooth is equivalent to saying that $X$ is a \emph{regular} prescheme.

\medskip\noindent\textbf{V.} ---
Finally, we will see in \S\S~20, 21, and 22 that the notion of formally smooth algebra is introduced naturally in the theory of \emph{K\"ahler differentials}, the two theories shedding light on each other.
\end{env}

\begin{env}[19.0.3]
\label{0.19.0.3}
Throughout this section and the following ones, topological rings and modules will be assumed to be \emph{linearly topologized} ($\mathbf{0_I}$, 7.1.1);
the topological rings considered will be assumed to be \emph{commutative}, unless explicit mention to the contrary.
Recall that if $A$ and $B$ are two topological rings, $\rho:A\to B$ a ring homomorphism defining on $B$ a structure of $A$-algebra, we say that $B$ is a \emph{topological $A$-algebra} if $\rho$ is continuous for the topologies considered.

For brevity, in a topological ring $A$ (resp. a topological $A$-module $M$), we will say ``fundamental system of open ideals (resp. submodules)'' instead of ``fundamental system of neighbourhoods of $0$ consisting of ideals (resp. submodules)''.

Given a topological ring $A$ and a topological $A$-module $M$, the sets $\mathfrak{J}M$, where $\mathfrak{J}$ ranges over a fundamental system of open ideals in $A$, form a fundamental system of open submodules of a topology on $M$ making $M$ a topological $A$-module, called the topology \emph{derived} from that of $A$.

Let $M$ be a topological $A$-module whose topology is \emph{less fine} than the topology derived from that of $A$;
then, if $N$ is an open submodule of $M$, the discrete $A$-module $M/N$ is annihilated by an open ideal $\mathfrak{R}$ of $A$, since by hypothesis there exists such an ideal $\mathfrak{R}$ with $\mathfrak{R}\cdot M\subset N$.

\oldpage[0\textsubscript{IV-1}]{71}

If $M$ and $N$ are two topological $A$-modules whose topologies are both derived from that of $A$, then every $A$-homomorphism $u:M\to N$ is \emph{continuous}, since for every neighbourhood $V$ of $0$ in $N$, there exists by hypothesis an open ideal $\mathfrak{J}$ of $A$ such that $\mathfrak{J}N\subset V$, and one thus has $u(\mathfrak{J}M)\subset\mathfrak{J}N\subset V$.

When $B$ is a topological $A$-algebra, the topology on $B$ \emph{derived} from that of $A$ is \emph{finer} than the topology given on $B$, since for every open ideal $\mathfrak{R}$ of $B$, there is by hypothesis an open ideal $\mathfrak{J}$ of $A$ such that $\mathfrak{J}B\subset\mathfrak{R}$.
\end{env}

\subsection{Formal epimorphisms and monomorphisms}
\label{subsection:0.19.1}

\begin{proposition}[19.1.1]
\label{0.19.1.1}
Let $A$ be a topological ring, $M$, $N$ two topological $A$-modules, $(W_\lambda)$ a fundamental system of open submodules in $N$, $u:M\to N$ a continuous $A$-homomorphism, $\widehat{M}$, $\widehat{N}$ the separated completions of $M$ and $N$, $\hat{u}:\widehat{M}\to\widehat{N}$ the continuous extension of $u$ to the separated completions.
\begin{enumerate}
  \item[(i)] The following conditions are equivalent:
  \begin{enumerate}
    \item[a)] $u(M)$ is dense in $N$.
    \item[a$'$)] $\hat{u}(\widehat{M})$ is dense in $\widehat{N}$.
    \item[a$''$)] For all $\lambda$, the composite homomorphism $M\xrightarrow{u}N\to N/W_\lambda$ is surjective.
  \end{enumerate}
  \item[(ii)] The following conditions are equivalent:
  \begin{enumerate}
    \item[b)] The inverse image by $u$ of the topology of $N$ is equal to the topology of $M$.
    \item[b$'$)] $\hat{u}$ is an isomorphism of the topological $\widehat{A}$-module $\widehat{M}$ onto the sub-$\widehat{A}$-module $\hat{u}(\widehat{M})$ of $\widehat{N}$ (which is necessarily closed).
    \item[b$''$)] The $u^{-1}(W_\lambda)$ form a fundamental system of neighbourhoods of $0$ in $M$.
  \end{enumerate}
\end{enumerate}
This follows immediately from the definition of separated completions, and from the fact that $N/W_\lambda$ is discrete.
\end{proposition}

\begin{env}[19.1.2]
\label{0.19.1.2}
When the equivalent conditions of (i) (resp. (ii)) in \sref{0.19.1.1} are satisfied, we say that $u$ is a \emph{formal epimorphism} (resp. a \emph{formal monomorphism}).
We say that $u$ is a \emph{formal bimorphism} if it is both a formal epimorphism and a formal monomorphism;
this amounts to the same, by virtue of \sref{0.19.1.1}, as saying that $\hat{u}$ is an \emph{isomorphism} of the topological $\widehat{A}$-module $\widehat{M}$ onto the topological $\widehat{A}$-module $\widehat{N}$.
\end{env}

\begin{proposition}[19.1.3]
\label{0.19.1.3}
Let $A$ be a topological ring, $M$, $N$ two topological $A$-modules, $u:M\to N$ a continuous $A$-homomorphism, $(V_\lambda)$, $(W_\lambda)$ two fundamental systems of open submodules in $M$ and $N$ respectively, having the same index set, and such that $u(V_\lambda)\subset W_\lambda$ for all $\lambda$;
let $u_\lambda:M/V_\lambda\to N/W_\lambda$ be the homomorphism induced by $u$ by passing to quotients.
Then:
\begin{enumerate}
  \item[(i)] For $u$ to be a formal epimorphism, it is necessary and sufficient that, for all $\lambda$, $u_\lambda$ be surjective.
  \item[(ii)] For $u$ to be a formal monomorphism, it is necessary and sufficient that, for all $\lambda$, there exist $\mu$ such that $V_\mu\subset V_\lambda$ and that $\Ker(u_\mu)$ be contained in the kernel $V_\lambda/V_\mu$ of the canonical map $M/V_\mu\to M/V_\lambda$.
  \item[(iii)] For $u$ to be a formal bimorphism, it is necessary and sufficient that, for all $\lambda$, $u_\lambda$ be surjective and that there exist $\mu$ such that $V_\mu\subset V_\lambda$ and that the canonical homomorphism $M/V_\mu\to M/V_\lambda$

  \oldpage[0\textsubscript{IV-1}]{72}

  factorize as $M/V_\mu\xrightarrow{u_\mu}N/W_\mu\xrightarrow{h}M/V_\lambda$ (where the homomorphism $N/W_\mu\to M/V_\lambda$ is necessarily unique).
\end{enumerate}
\end{proposition}

\begin{proof}
Assertion (i) is immediate;
on the other hand, $\Ker(u_\mu)=u^{-1}(W_\mu)/V_\mu$, and the kernel of the canonical map $M/V_\mu\to M/V_\lambda$ is $V_\lambda/V_\mu$;
the condition of (ii) thus amounts to saying that $u^{-1}(W_\mu)\subset V_\lambda$, and assertion (ii) follows immediately;
finally, when $u_\mu$ is surjective, to say that $\Ker(u_\mu)$ is contained in $V_\lambda/V_\mu$, or to say that $M/V_\mu\to M/V_\lambda$ factorizes as $v\circ u_\mu$, where $v$ is a homomorphism $N/W_\mu\to M/V_\lambda$, is the same thing, since $N/W_\mu$ then identifies with $(M/V_\mu)/\Ker(u_\mu)$.
\end{proof}

\begin{corollary}[19.1.4]
\label{0.19.1.4}
Let $A$ be a topological ring, $M$, $N$ two topological $A$-modules whose topologies are derived from that of $A$, $u:M\to N$ a formal epimorphism.
Let $(\mathfrak{J}_\lambda)$ be a fundamental system of open ideals in $A$.
For $u$ to be a formal bimorphism, it is necessary and sufficient that for all $\lambda$, the homomorphism $u_\lambda:M/\mathfrak{J}_\lambda M\to N/\mathfrak{J}_\lambda N$ induced by $u$ by passing to quotients is bijective.
\end{corollary}

\begin{proof}
Indeed, $u(\mathfrak{J}_\lambda M)\subset\mathfrak{J}_\lambda N$ and one can apply the criterion \sref{0.19.1.3}, (iii);
but if one has a factorization $M/\mathfrak{J}_\mu M\xrightarrow{u_\mu}N/\mathfrak{J}_\mu N\xrightarrow{v}M/\mathfrak{J}_\lambda M$, the image of $\mathfrak{J}_\lambda M/\mathfrak{J}_\mu M$ by $u_\mu$ is $\mathfrak{J}_\lambda N/\mathfrak{J}_\mu N$, hence the image of $\mathfrak{J}_\lambda N/\mathfrak{J}_\mu N$ by $v$ is $0$ and $v$ factorizes as
\[
  N/\mathfrak{J}_\mu N\to N/\mathfrak{J}_\lambda N\xrightarrow{v'}M/\mathfrak{J}_\lambda M;
\]
but then the automorphism identical to that of $M/\mathfrak{J}_\lambda M$ factorizes as
\[
  M/\mathfrak{J}_\lambda M\xrightarrow{u_\lambda}N/\mathfrak{J}_\lambda N\xrightarrow{v'}M/\mathfrak{J}_\lambda M,
\]
which shows that $u_\lambda$ is injective, and as we already know it is surjective, this proves the corollary.
\end{proof}

\begin{proposition}[19.1.5]
\label{0.19.1.5}
Let $A$ be a topological ring, $M$, $N$ two topological $A$-modules, and $u:M\to N$ a continuous $A$-homomorphism.
The following conditions are equivalent:
\begin{enumerate}
  \item[a)] For every discrete topological $A$-module $E$ and every continuous $A$-homomorphism $v:M\to E$, there exists a continuous $A$-homomorphism $w:N\to E$ such that $v=w\circ u$.
  \item[b)] For every open submodule $V'$ of $M$, there exist an open submodule $W''$ of $N$, an open submodule $V''\subset V'\cap u^{-1}(W'')$, and a homomorphism $h:N/W''\to M/V'$ such that the canonical homomorphism $M/V''\to M/V'$ factorizes as
  \[
    M/V''\xrightarrow{u''}N/W''\xrightarrow{h}M/V',
  \]
  where $u''$ is induced by $u$ by passing to quotients.
\end{enumerate}
\end{proposition}

\begin{proof}
To see that a) implies b), apply a) to $E=M/V'$ and to the canonical homomorphism $v:M\to M/V'$.
Then $W''=w^{-1}(0)$ is an open submodule of $N$ and $V''=u^{-1}(W'')$ is an open submodule of $M$ contained in $V'$;
these submodules and the homomorphism $h:N/W''\to M/V'$ induced by $w$ answer the question.
Conversely, to show that b) implies a), one may restrict to the case where $v$ is surjective, replacing $E$ by $v(M)$;
then $V'=\operatorname{Ker}(v)$ is an open submodule of $M$, so $E$ is isomorphic to the discrete $A$-module $M/V'$, and applying b) gives

\oldpage[0\textsubscript{IV-1}]{73}

the desired continuous $A$-homomorphism $w$ by taking the composite $N\to N/W''\xrightarrow{h}M/V'$, the diagram
\[
  \xymatrix{
    M\ar[r]^u\ar[d] & N\ar[d]\\
    M/V''\ar[r]_{u''} & N/W''
  }
\]
being commutative.
\end{proof}

When the equivalent conditions of \sref{0.19.1.5} are satisfied, we say that $u$ is \emph{formally left invertible};
since condition~b) of \sref{0.19.1.5} implies $u^{-1}(W'')\subset V'$, $u$ is then a formal monomorphism.
The terminology is further justified by the following corollaries.

\begin{corollary}[19.1.6]
\label{0.19.1.6}
If there exists a continuous $\widehat A$-homomorphism $s:\widehat N\to\widehat M$ such that $s\circ\widehat u=1_{\widehat M}$, then $u$ is formally left invertible.
\end{corollary}

Note that $\widehat u$ is then a topological isomorphism of $\widehat M$ onto a \emph{topological direct summand} of $\widehat N$.
To prove \sref{0.19.1.6}, it suffices to note that, with the notation of \sref{0.19.1.5}, a), $v:M\to E$ extends to a continuous $\widehat A$-homomorphism $\widehat v:\widehat M\to E$ since $E$ is discrete;
let $j:M\to\widehat M$ and $j':N\to\widehat N$ be the canonical homomorphisms.
Then, setting $w=\widehat v\circ s\circ j'$, one has $w\circ u=\widehat v\circ s\circ\widehat u\circ j=\widehat v\circ j=v$.

\begin{corollary}[19.1.7]
\label{0.19.1.7}
Suppose that the topologies of $M$ and $N$ are derived from that of $A$, and let $(\mathfrak J_\lambda)$ be a fundamental system of open ideals in $A$.
For $u$ to be formally left invertible, it is necessary and sufficient that, for every $\lambda$, the homomorphism
\[
  u_\lambda:M/\mathfrak J_\lambda M\longrightarrow N/\mathfrak J_\lambda N
\]
induced by $u$ by passing to quotients be left invertible (in other words, be an isomorphism of $M/\mathfrak J_\lambda M$ onto a direct summand of $N/\mathfrak J_\lambda N$).
\end{corollary}

Indeed, the condition is sufficient, since for $V'=\mathfrak J_\lambda M$ in condition~b) of \sref{0.19.1.5}, one answers the question by taking $W''=\mathfrak J_\lambda N$, $V''=V'$, and $h$ such that $h\circ u_\lambda$ is the identity on $M/\mathfrak J_\lambda M$.
Conversely, if $u$ is formally left invertible, then for every $\lambda$, by \sref{0.19.1.5}, b), there are $\mu$ with $\mathfrak J_\mu\subset\mathfrak J_\lambda$ and a homomorphism $h:N/\mathfrak J_\mu N\to M/\mathfrak J_\lambda M$ such that the canonical homomorphism $M/\mathfrak J_\mu M\to M/\mathfrak J_\lambda M$ factorizes as
\[
  M/\mathfrak J_\mu M\xrightarrow{u_\mu}N/\mathfrak J_\mu N
  \xrightarrow{h}M/\mathfrak J_\lambda M.
\]
But since
$h(\mathfrak J_\lambda(N/\mathfrak J_\mu N))
 \subset\mathfrak J_\lambda(M/\mathfrak J_\lambda M)=0$,
$h$ factorizes canonically as
\[
  N/\mathfrak J_\mu N\longrightarrow N/\mathfrak J_\lambda N
  \xrightarrow{s}M/\mathfrak J_\lambda M,
\]
and it is immediate that $s$ is a left inverse of $u_\lambda$.

\begin{proposition}[19.1.8]
\label{0.19.1.8}
Let $A$ be a topological ring admitting a countable decreasing fundamental system $(\mathfrak J_n)$ of open ideals.
Let $M$, $N$ be two topological $A$-modules whose topologies are derived from that of $A$, and let $u:M\to N$ be an $A$-homomorphism.
Set $M_n=M/\mathfrak J_nM$, $N_n=N/\mathfrak J_nN$, and let $u_n:M_n\to N_n$ be the $(A/\mathfrak J_n)$-homomorphism induced by $u$ by passing to quotients.
Suppose that, for every $n$, $P_n=\operatorname{Coker}(u_n)$ is a projective $(A/\mathfrak J_n)$-module and that $M$ is separated and complete.
Then $u$ is formally left invertible if and only if $u$ is left invertible (and $u$ is then a topological isomorphism of $M$ onto a topological direct summand of $N$).
\end{proposition}

\oldpage[0\textsubscript{IV-1}]{74}

\begin{proof}
By \sref{0.19.1.7}, one has commutative diagrams
\[
  \xymatrix{
    0\ar[r] & M_n\ar[r]^{u_n} & N_n\ar[r]^{p_n} & P_n\ar[r] & 0\\
    0\ar[r] & M_{n+1}\ar[u]^f\ar[r]_{u_{n+1}} &
      N_{n+1}\ar[u]^g\ar[r]_{p_{n+1}} & P_{n+1}\ar[u]^h\ar[r] & 0
  }
\]
whose rows are split exact sequences;
in other words, for every $n$ there exists a homomorphism $s_n:P_n\to N_n$ such that $p_n\circ s_n=1_{P_n}$.
We show by induction on $n$ that homomorphisms $s'_n:P_n\to N_n$ can be defined so that $p_n\circ s'_n=1_{P_n}$ and the diagrams
\[
  \xymatrix{
    P_n\ar[r]^{s'_n} & N_n\\
    P_{n+1}\ar[u]^h\ar[r]_{s'_{n+1}} & N_{n+1}\ar[u]^g
  }
\]
are commutative.
Indeed, $g\circ s_{n+1}-s'_n\circ h$ is a homomorphism from $P_{n+1}$ into
\[
  u_n(M_n)=u_{n+1}(M_{n+1})/
  \mathfrak J_nu_{n+1}(M_{n+1});
\]
since $P_{n+1}$ is a projective $(A/\mathfrak J_{n+1})$-module, this homomorphism factorizes as
\[
  P_{n+1}\xrightarrow{t_{n+1}}u_{n+1}(M_{n+1})
  \longrightarrow
  u_{n+1}(M_{n+1})/\mathfrak J_nu_{n+1}(M_{n+1}),
\]
whence $s'_{n+1}=s_{n+1}-t_{n+1}$ answers the question.
The decomposition of $N_n$ as the direct sum of $u_n(M_n)$ and $s'_n(P_n)$ then immediately gives a homomorphism $w_n:N_n\to M_n$ which is a left inverse of $u_n$ and makes the diagrams
\[
  \xymatrix{
    N_n\ar[r]^{w_n} & M_n\\
    N_{n+1}\ar[u]^g\ar[r]_{w_{n+1}} & M_{n+1}\ar[u]^f
  }
\]
commutative.
The inverse system $(w_n)$ therefore has an inverse limit $\widehat w:\widehat N\to\widehat M=M$.
Composing with the canonical homomorphism $N\to\widehat N$ gives a homomorphism $w:N\to M$ such that, for every $n$, the endomorphism $(w\circ u)_n$ of $M_n=M/\mathfrak J_nM$ induced by $w\circ u$ is the identity.
Thus $w\circ u=1_M$ because $M$ is separated and complete.
\end{proof}

\begin{proposition}[19.1.9]
\label{0.19.1.9}
Let $A$ be a preadmissible topological ring ($\mathbf{0_I}$, 7.1.2), $\mathfrak Q$ an ideal of definition of $A$, and $(\mathfrak J_\lambda)$ a fundamental system of open ideals in $A$.
Let $M$, $N$ be two topological $A$-modules whose topologies are derived from that of $A$, such that for every $\lambda$, $N/\mathfrak J_\lambda N$ is a projective $(A/\mathfrak J_\lambda)$-module (cf.~\sref{0.19.2.3}).
Let $u:M\to N$ be an $A$-homomorphism.
Then the following conditions are equivalent:
\begin{enumerate}
  \item[a)] $u$ is formally left invertible.

  \oldpage[0\textsubscript{IV-1}]{75}

  \item[b)] The homomorphism $u_0:M/\mathfrak QM\to N/\mathfrak QN$ induced by $u$ by passing to quotients is left invertible.
\end{enumerate}
\end{proposition}

\begin{proof}
By \sref{0.19.1.7}, condition~a) is equivalent to $u_\lambda$ being left invertible for every $\lambda$.
Since $\mathfrak Q$ is open, and therefore contains some $\mathfrak J_\lambda$, it follows at once that $u_0$ is left invertible.
Conversely, to show that b) implies a), note that for every $\lambda$, $\mathfrak Q/\mathfrak J_\lambda$ is by hypothesis a nilpotent ideal of $A/\mathfrak J_\lambda$.
The assertion follows from the next proposition.
\end{proof}

\begin{proposition}[19.1.10]
\label{0.19.1.10}
Let $A$ be a ring, $M$, $N$ two $A$-modules, with $N$ projective, and let $u:M\to N$ be an $A$-homomorphism.
Let $\mathfrak J$ be an ideal of $A$;
suppose that one of the following conditions holds:
\begin{enumerate}
  \item[(i)] $\mathfrak J$ is nilpotent.
  \item[(ii)] $\mathfrak J$ is contained in the radical of $A$ and $M$ is of finite type.
\end{enumerate}
Then $u$ is left invertible if and only if the homomorphism
$u_0:M/\mathfrak JM\to N/\mathfrak JN$ of $(A/\mathfrak J)$-modules induced by $u$ by passing to quotients is left invertible.
\end{proposition}

\begin{proof}
The condition is plainly necessary; let us prove that it is sufficient.
Let $v_0$ be a left inverse of $u_0$.
The composite $N\to N/\mathfrak JN\xrightarrow{v_0}M/\mathfrak JM$ factorizes as $N\xrightarrow{v}M\to M/\mathfrak JM$ because $N$ is projective.
Then $w=v\circ u$ is an endomorphism of $M$ whose induced endomorphism $w_0$ of $M/\mathfrak JM$ is the identity, and it is enough to prove that $w$ is bijective (for then $w^{-1}\circ v$ is a left inverse of $u$).
We distinguish the two cases.

(i) For every $n$ one has the commutative diagram
\[
  \xymatrix{
    (\mathfrak J^n/\mathfrak J^{n+1})\otimes_{A/\mathfrak J}
      (M/\mathfrak JM)\ar[r]\ar[d]_{1\otimes\operatorname{gr}^0(w)} &
      \mathfrak J^nM/\mathfrak J^{n+1}M\ar[d]^{\operatorname{gr}^n(w)}\\
    (\mathfrak J^n/\mathfrak J^{n+1})\otimes_{A/\mathfrak J}
      (M/\mathfrak JM)\ar[r] & \mathfrak J^nM/\mathfrak J^{n+1}M
  }
\]
in which the horizontal arrows are surjective.
Since $\operatorname{gr}^0(w)=w_0$ is the identity, so is $\operatorname{gr}^n(w)$, which is therefore, a fortiori, bijective.
The $\mathfrak J$-preadic filtration on $M$ is finite because $\mathfrak J$ is nilpotent, so $w$ is bijective (Bourbaki, \emph{Alg. comm.}, Chap.~III, \S~2, no.~8, Cor.~3 of Th.~1).

(ii) It suffices to show that, for every maximal ideal $\mathfrak m$ of $A$, the endomorphism $w_{\mathfrak m}$ of $M_{\mathfrak m}$ is bijective (Bourbaki, \emph{Alg. comm.}, Chap.~II, \S~3, no.~3, Th.~1).
Since $\mathfrak JA_{\mathfrak m}\subset\mathfrak mA_{\mathfrak m}$ and
$A_{\mathfrak m}/\mathfrak JA_{\mathfrak m}=(A/\mathfrak J)_{\mathfrak m}$, we are reduced to proving the proposition when $A$ is local.
Moreover, one may suppose that $\mathfrak J$ is the maximal ideal $\mathfrak J_0$ of $A$, because if $u_0$ is left invertible, so is

\oldpage[0\textsubscript{IV-1}]{76}

$u_{00}:M/\mathfrak J_0M\to N/\mathfrak J_0N$, obtained by tensoring $u_0$ with $1_{A/\mathfrak J_0}$, since
\[
  M/\mathfrak J_0M=(M/\mathfrak JM)\otimes_{A/\mathfrak J}(A/\mathfrak J_0).
\]
\footnote{The French source prints $N/\mathfrak J_0M$ as the codomain of $u_{00}$; the module type requires $N/\mathfrak J_0N$.}
Suppose then that $\mathfrak J$ is maximal, so that $A/\mathfrak J$ is a field.
It plainly suffices to prove that $M$ is a free $A$-module under the hypotheses of the statement:
indeed, $\det(w_0)$ is then the canonical image of $\det(w)$ in $A/\mathfrak J$, hence $\det(w)\notin\mathfrak J$ and is therefore invertible.
Now $M/\mathfrak JM$ is a finite-dimensional free $(A/\mathfrak J)$-vector space, so there are a finite free $A$-module $L$ and an $A$-homomorphism $f:L\to M$ such that the induced homomorphism $f_0:L/\mathfrak JL\to M/\mathfrak JM$ is bijective.
Since $M$ is of finite type, $f$ is surjective by Nakayama's lemma (Bourbaki, \emph{Alg. comm.}, Chap.~II, \S~3, no.~2, Cor.~1 of Prop.~4).
Moreover, if $g=u\circ f$, then the induced homomorphism $g_0:L/\mathfrak JL\to N/\mathfrak JN$ is left invertible, and since $L$ is free, the preceding observation shows that $g$ itself is left invertible.
This plainly implies that $f$ is injective, completing the proof.
\end{proof}

Let us mention in passing the following corollaries of \sref{0.19.1.10}.

\begin{corollary}[19.1.11]
\label{0.19.1.11}
Let $A$ be a local ring, $k$ its residue field, $M$ a finite-type $A$-module, $N$ a projective $A$-module, and $u:M\to N$ a homomorphism.
Then $u$ is left invertible if and only if there exists a system of generators $(x_i)_{1\leq i\leq m}$ of $M$ such that the images of the $x_i\otimes1$ under $u\otimes1:M\otimes_Ak\to N\otimes_Ak$ are linearly independent in $N\otimes_Ak$;
the $x_i$ then form a basis of $M$.
\end{corollary}

The condition is plainly necessary, since if $u$ is left invertible, $M$ is a finite-type projective $A$-module and hence free.
Conversely, if the condition holds, the $x_i\otimes1$ plainly form a basis of $M\otimes_Ak$ and $u\otimes1$ is left invertible;
apply \sref{0.19.1.10} to the maximal ideal $\mathfrak J$ of $A$.

\begin{corollary}[19.1.12]
\label{0.19.1.12}
Let $A$ be a ring, $M$ a finite-type $A$-module, $N$ a projective $A$-module, and $u:M\to N$ a homomorphism.
For every prime ideal $\mathfrak p\in\Spec(A)$, the following conditions are equivalent:
\begin{enumerate}
  \item[a)] $u_{\mathfrak p}:M_{\mathfrak p}\to N_{\mathfrak p}$ is left invertible.
  \item[b)] $u\otimes1:M\otimes_Ak(\mathfrak p)\to N\otimes_Ak(\mathfrak p)$ is injective (where $k(\mathfrak p)$ denotes the residue field of $A$ at $\mathfrak p$).
  \item[c)] There is a finite system of elements $x_i\in M$ $(1\leq i\leq m)$ whose images generate $M_{\mathfrak p}$, and a system of $m$ linear forms $y_i^*$ on $N$ $(1\leq i\leq m)$ such that $\det(\langle y_i^*,u(x_j)\rangle)\notin\mathfrak p$.
  \item[d)] There exists $f\in A-\mathfrak p$ such that $u_f:M_f\to N_f$ is left invertible.
\end{enumerate}
Moreover, the set of $\mathfrak p\in\Spec(A)$ satisfying these conditions is open in $\Spec(A)$.
\end{corollary}

The last assertion is an immediate consequence of d).
Since $N$ is projective, it is a direct summand of a free $A$-module $A^{(I)}$;
and since $M$ is of finite type, $u(M)$ is contained in a submodule of $A^{(I)}$ of the form $A^n$ for finite $n$.
Since each of a), b), c), and d) is equivalent to the corresponding assertion with $N$ replaced by $N\oplus P$ (with $u$ still mapping $M$ into $N$), we are reduced to the case where $N$ is free of finite type.
Clearly d) implies a), and a) and b) are equivalent by \sref{0.19.1.11}.
Moreover, a) implies that $M_{\mathfrak p}$ is free by \sref{0.19.1.11}, so a) implies c) by choosing the $x_i$ so that their images form a basis of this

\oldpage[0\textsubscript{IV-1}]{77}

$A_{\mathfrak p}$-module and noting that, because $N$ is free of finite type, every linear form on $N_{\mathfrak p}$ can be written $v/f$, where $f\in A-\mathfrak p$ and $v$ is a linear form on $N$.
Clearly c) implies b), so it remains to see that a) implies d).
Since $N$ is finitely presented, $(\operatorname{Hom}_A(N,M))_{\mathfrak p}$ identifies canonically with $\operatorname{Hom}_{A_{\mathfrak p}}(N_{\mathfrak p},M_{\mathfrak p})$ (Bourbaki, \emph{Alg. comm.}, Chap.~II, \S~2, no.~7, Prop.~19).
If $w_{\mathfrak p}$ is a left inverse of $u_{\mathfrak p}$, there are a homomorphism $w:N\to M$ and an element $f\in A-\mathfrak p$ such that $w_{\mathfrak p}=w\otimes(1/f)1_{A_{\mathfrak p}}$.
The relation $w_{\mathfrak p}\circ u_{\mathfrak p}=1_{M_{\mathfrak p}}$ can therefore also be written $(w\circ u)\otimes1_{A_{\mathfrak p}}=f.1_{M_{\mathfrak p}}$.
Since $M$ is of finite type, there is $g\in A-\mathfrak p$ such that the endomorphism $g((w\circ u)-f.1_M)$ vanishes on all generators of $M$, and hence is zero.
Put $h=fg$, $u_h=u\otimes1_{A_h}$, and $w_h=w\otimes1_{A_h}$.
Then $w_h(u_h(z))=(f/1)z$ for every $z\in M_h$;
since $h/1$ is invertible in $A_h$, so is $f/1$, and $(f/1)^{-1}w_h$ is a left inverse of $u_h$.
This completes the proof.

\begin{proposition}[19.1.13]
\label{0.19.1.13}
Suppose the hypotheses of \sref{0.19.1.9} hold, and suppose in addition that $(\mathfrak J_\lambda)$ is a decreasing sequence $(\mathfrak J_n)$ and that $M$ is separated and complete.
Then conditions~a) and~b) of \sref{0.19.1.9} are also equivalent to:
\begin{enumerate}
  \item[c)] $u$ is left invertible.
\end{enumerate}
\end{proposition}

We already know that c) implies a) by \sref{0.19.1.6}.
Conversely, if a) holds, then with the notation of \sref{0.19.1.8}, $M_n$ is a direct summand of the projective $(A/\mathfrak J_n)$-module $N_n$;
hence $P_n$ is also isomorphic to a direct summand of $N_n$ and is therefore projective.
It remains only to apply \sref{0.19.1.8}.

Finally, note the following proposition.

\begin{proposition}[19.1.14]
\label{0.19.1.14}
Let $A$ be a ring, $M$ a finite-type $A$-module, $N$ a projective $A$-module, and $u:M\to N$ an $A$-homomorphism.
\begin{enumerate}
  \item[(i)] The homomorphism $u$ is left invertible if and only if, for every maximal ideal $\mathfrak m$ of $A$, $u_{\mathfrak m}:M_{\mathfrak m}\to N_{\mathfrak m}$ is left invertible.
  \item[(ii)] Let $A'$ be an $A$-algebra which is a faithfully flat $A$-module.
  Then $u$ is left invertible if and only if $u\otimes1:M\otimes_AA'\to N\otimes_AA'$ is left invertible.
\end{enumerate}
\end{proposition}

As in \sref{0.19.1.12}, it suffices to consider the case where $N$ is free of finite type.
Then saying that $u$ is left invertible means that $u$ is injective and that the quotient module $P=N/u(M)$ is projective, since $u(M)$ is then a direct summand of $N$.
Also, since $M$ is of finite type, $P$ is finitely presented.

(i) Necessity is clear.
Conversely, if the condition holds, $u$ is injective (Bourbaki, \emph{Alg. comm.}, Chap.~II, \S~3, no.~3, Th.~1), and since
$P_{\mathfrak m}=N_{\mathfrak m}/u_{\mathfrak m}(M_{\mathfrak m})$ is projective for every $\mathfrak m$, it follows that $P$ is projective (\emph{loc. cit.}, \S~5, no.~2, Th.~1).

(ii) Again, necessity is immediate.
Conversely, if the condition holds, $u$ is injective ($\mathbf{0_I}$, 6.4.1), and since
$P\otimes_AA'=\operatorname{Coker}(u\otimes1)$ is projective, hence flat, it follows that $P$ is a flat $A$-module ($\mathbf{0_I}$, 6.6.3), and therefore projective because it is finitely presented (Bourbaki, \emph{Alg. comm.}, Chap.~II, \S~5, no.~2, Cor.~2 of Th.~1).

\oldpage[0\textsubscript{IV-1}]{78}

\begin{remark}[19.1.15]
\label{0.19.1.15}
The notion dual to that of a formally left invertible homomorphism is that of a \emph{formally right invertible} homomorphism.
Such a continuous $A$-homomorphism $u:M\to N$ satisfies, by definition, the following condition:
for every open submodule $W'$ of $N$, there exist an open submodule $V'\subset u^{-1}(W')$ of $M$, an open submodule $W''\subset W'$ of $N$, and a homomorphism $h:N/W''\to M/V'$ such that the canonical homomorphism $N/W''\to N/W'$ factorizes as
\[
  N/W''\xrightarrow{h}M/V'\xrightarrow{u'}N/W',
\]
where $u'$ is induced by $u$ by passing to quotients.
This implies that $u$ is a formal epimorphism.
If there exists a continuous $\widehat A$-homomorphism $r:\widehat N\to\widehat M$ such that $\widehat u\circ r=1_{\widehat N}$, then $u$ is immediately seen to be formally right invertible.
Under the hypotheses of \sref{0.19.1.7}, $u$ is formally right invertible if and only if every $u_\lambda$ is right invertible (that is, if $\operatorname{Ker}(u_\lambda)$ is a direct summand of $M/\mathfrak J_\lambda M$ and $u_\lambda$ is an isomorphism from a complement of $\operatorname{Ker}(u_\lambda)$ onto $N/\mathfrak J_\lambda N$).
We leave to the reader the statements and proofs of the propositions corresponding to \sref{0.19.1.8} and \sref{0.19.1.9}:
in the analogue of \sref{0.19.1.8}, one must suppose $M$ separated and complete and $M_n$ a projective $(A/\mathfrak J_n)$-module;
in the analogue of \sref{0.19.1.9}, one must omit the hypothesis on the $N/\mathfrak J_\lambda N$, but instead suppose that, for every $\lambda$, $M/\mathfrak J_\lambda M$ is a projective $(A/\mathfrak J_\lambda)$-module.
\end{remark}

\subsection{Formally projective modules}
\label{subsection:0.19.2}

\begin{definition}[19.2.1]
\label{0.19.2.1}
Let $A$ be a topological ring.
A topological $A$-module $M$ is said to be \emph{formally projective} if it satisfies the following condition:
for every open ideal $\mathfrak J$ of $A$, every pair of discrete $(A/\mathfrak J)$-modules $P$, $Q$, every surjective $A$-homomorphism $u:P\to Q$, and every continuous $A$-homomorphism $v:M\to Q$, there exists a continuous $A$-homomorphism $w:M\to P$ such that $v=u\circ w$.
\end{definition}

\begin{env}[19.2.2]
\label{0.19.2.2}
To verify the condition of \sref{0.19.2.1}, one may plainly restrict, replacing $Q$ by $v(M)$ and $P$ by $u^{-1}(v(M))$, to the case where $v$ itself is surjective.
Then $Q$ is isomorphic to $M/V$, where $V$ is an open submodule of $M$ such that $\mathfrak JM\subset V$.
The condition of \sref{0.19.2.1} is therefore equivalent to requiring that, for every discrete $(A/\mathfrak J)$-module $P$ and every surjective $A$-homomorphism $u:P\to M/V$, there exist an open submodule $V'\subset V$ of $M$ and an $A$-homomorphism $w:M/V'\to P$ such that the canonical homomorphism $M/V'\to M/V$ factorizes as
\[
  M/V'\xrightarrow{w}P\xrightarrow{u}M/V.
\]
It suffices to verify this condition when $\mathfrak J$ ranges over a fundamental system of neighbourhoods of $0$ in $A$ formed of ideals.
\end{env}

\begin{env}[19.2.3]
\label{0.19.2.3}
Suppose that there exist a fundamental system $(\mathfrak J_\lambda)$ of open ideals in $A$ and a fundamental system $(V_\lambda)$ of open submodules of $M$, with the same index set, such that for every $\lambda$, $M/V_\lambda$ is a projective $(A/\mathfrak J_\lambda)$-module.
Then $M$ is formally projective.
Indeed, with the notation of \sref{0.19.2.2}, choose $\lambda$ such that $\mathfrak J_\lambda\subset\mathfrak J$ and $V_\lambda\subset V$.
Since $P$ is also an $(A/\mathfrak J_\lambda)$-module, the factorization of the canonical homomorphism $M/V_\lambda\to M/V$ as $M/V_\lambda\to P\to M/V$ follows from its being an $(A/\mathfrak J_\lambda)$-homomorphism and from the projectivity of $M/V_\lambda$.

When the stronger condition in this paragraph holds, $M$ is said to be \emph{strictly formally projective}.
\end{env}

\begin{proposition}[19.2.4]
\label{0.19.2.4}
Let $A$ be a topological ring, and let $M$ be a topological $A$-module whose topology is derived from that of $A$.
Then $M$ is formally projective if and only if it is strictly formally projective.
\end{proposition}

\oldpage[0\textsubscript{IV-1}]{79}

\begin{proof}
We have just seen that the condition is sufficient.
Conversely, suppose that $M$ is formally projective, and let $(\mathfrak J_\lambda)$ be a fundamental system of open ideals in $A$.
For each $\lambda$, let $P$ be a free $(A/\mathfrak J_\lambda)$-module and let $u:P\to M/\mathfrak J_\lambda M$ be a surjective $(A/\mathfrak J_\lambda)$-homomorphism.
There is then a $\mathfrak J_\mu\subset\mathfrak J_\lambda$ such that the canonical homomorphism $M/\mathfrak J_\mu M\to M/\mathfrak J_\lambda M$ factorizes as
\[
  M/\mathfrak J_\mu M\xrightarrow{w}P\xrightarrow{u}M/\mathfrak J_\lambda M.
\]
But since $w(\mathfrak J_\lambda(M/\mathfrak J_\mu M))\subset\mathfrak J_\lambda P=0$, $w$ factorizes through a map $v:M/\mathfrak J_\lambda M\to P$, and it is clear that $u\circ v$ is the identity on $M/\mathfrak J_\lambda M$.
Thus this $(A/\mathfrak J_\lambda)$-module is projective.
\end{proof}

\begin{proposition}[19.2.5]
\label{0.19.2.5}
Let $A$ be a topological ring, $M$ a topological $A$-module.
\begin{enumerate}
  \item[(i)] The module $M$ is formally projective (resp. strictly formally projective) if and only if the topological $\widehat A$-module $\widehat M$ is so.
  \item[(ii)] Let $A'$ be a topological $A$-algebra.
  If $M$ is formally projective (resp. strictly formally projective), then $M\otimes_A A'$, equipped with the tensor-product topology, is a formally projective (resp. strictly formally projective) topological $A'$-module.
\end{enumerate}
\end{proposition}

\begin{proof}
\begin{enumerate}
  \item[(i)] When $\mathfrak J$ (resp. $V$) ranges over the open ideals of $A$ (resp. the open submodules of $M$), the separated completion $\widehat{\mathfrak J}$ (resp. $\widehat V$) ranges over the open ideals of $\widehat A$ (resp. the open submodules of $\widehat M$), and $\widehat A/\widehat{\mathfrak J}=A/\mathfrak J$ (resp. $\widehat M/\widehat V=M/V$), up to canonical isomorphism.
  Since formal projectivity (resp. strict formal projectivity) involves only the $A/\mathfrak J$ and the $M/V$, assertion~(i) follows immediately.

  \item[(ii)] Suppose first that $M$ is formally projective and put $M'=M\otimes_A A'$.
  Let $\mathfrak J'$ be an open ideal of $A'$, let $P'$ and $Q'$ be two discrete $(A'/\mathfrak J')$-modules, let $u':P'\to Q'$ be a surjective $A'$-homomorphism, and let $v':M'\to Q'$ be a continuous $A'$-homomorphism.
  There is an open ideal $\mathfrak J$ of $A$ such that $\mathfrak JA'\subset\mathfrak J'$, so $P'$ and $Q'$ are also discrete $(A/\mathfrak J)$-modules.
  Applying formal projectivity to the continuous composite $v:M\xrightarrow{j}M'\xrightarrow{v'}Q'$, one obtains a continuous $A$-homomorphism $w:M\to P'$ such that $v=u'\circ w$.
  Since $P'$ is a topological $A'$-module, $w$ factorizes as $M\xrightarrow{j}M'\xrightarrow{w'}P'$, where $w'$ is a continuous $A'$-homomorphism; from $v'\circ j=(u'\circ w')\circ j$ one concludes that $v'=u'\circ w'$.

  Now suppose that $M$ is strictly formally projective.
  Let $(\mathfrak J_\lambda)$ (resp. $(V_\lambda)$) be a fundamental system of open ideals (resp. open submodules) in $A$ (resp. $M$), with $M/V_\lambda$ projective over $A/\mathfrak J_\lambda$, and let $(\mathfrak J'_\mu)$ be a fundamental system of open ideals in $A'$.
  A fundamental system of neighbourhoods of $0$ in $M'$ is formed by
  \[
    W_{\lambda\mu}=\operatorname{Im}(M\otimes_A\mathfrak J'_\mu)+\operatorname{Im}(V_\lambda\otimes_A A'),
  \]
  where $\lambda$ and $\mu$ satisfy $\mathfrak J_\lambda A'\subset\mathfrak J'_\mu$.
  Since
  \[
    (M\otimes_A A')/W_{\lambda\mu}=(M/V_\lambda)\otimes_{A/\mathfrak J_\lambda}(A'/\mathfrak J'_\mu)
  \]
  and $M/V_\lambda$ is projective over $A/\mathfrak J_\lambda$, the module $M'/W_{\lambda\mu}$ is projective over $A'/\mathfrak J'_\mu$.
\end{enumerate}
\end{proof}

\subsection{Formally smooth algebras}
\label{subsection:0.19.3}

\begin{definition}[19.3.1]
\label{0.19.3.1}
Let $A$ be a topological ring, $B$ a topological $A$-algebra.
We say that $B$ is a \emph{formally smooth} $A$-algebra if, for every discrete topological $A$-algebra $C$, and every nilpotent ideal $\mathfrak{J}$ of $C$, every continuous $A$-homomorphism $u:B\to C/\mathfrak{J}$ \emph{factorizes} as $B\xrightarrow{v}C\xrightarrow{\varphi}C/\mathfrak{J}$, where $v$ is a continuous $A$-homomorphism and $\varphi$ the canonical homomorphism.
\end{definition}

\oldpage[0\textsubscript{IV-1}]{80}
The definition \sref{0.19.3.1} amounts to saying that the following property holds:
\begin{enumerate}
  \item[(P)] For every open ideal $\mathfrak{R}$ of the $A$-algebra $B$ and every $A$-homomorphism $u':B/\mathfrak{R}\to C/\mathfrak{J}$, there exists an open ideal $\mathfrak{R}'\subset\mathfrak{R}$ of $B$ such that the homomorphism
  $B\to B/\mathfrak{R}\xrightarrow{u'}C/\mathfrak{J}$ factorizes as
  \[
    B\to B/\mathfrak{R}'\xrightarrow{v'}C\to C/\mathfrak{J},
  \]
  where $v'$ is an $A$-homomorphism.
\end{enumerate}

Indeed, if $u:B\to C/\mathfrak{J}$ is a continuous $A$-homomorphism, its kernel $\mathfrak{R}$ is an open ideal of $B$, and if (P) is satisfied, it suffices to take $u'$ and to take for $v$ the composition $B\to B/\mathfrak{R}\xrightarrow{v'}C$ to satisfy the conditions of \sref{0.19.3.1}.
Conversely, suppose $B$ is a formally smooth $A$-algebra;
let there be given an open ideal $\mathfrak{R}$ of $B$ and a continuous $A$-homomorphism $u':B/\mathfrak{R}\to C/\mathfrak{J}$;
if $v:B\to C$ is a continuous $A$-homomorphism such that $u$ factorizes as $B\xrightarrow{v}C\to C/\mathfrak{J}$, the ideal $\mathfrak{R}'=\Ker(v)\cap\mathfrak{R}$ of $B$ is open and contained in $\mathfrak{R}$;
by taking $v'$ to factorize as $B\to B/\mathfrak{R}'\xrightarrow{v'}C$, and $v'$ clearly satisfies the condition (P).

\begin{proposition}[19.3.2]
\label{0.19.3.2}
Let $A$ be a discrete ring, $V$ a projective $A$-module;
the symmetric algebra $B=\mathrm{S}_A^\bullet(V)$, equipped with the discrete topology, is a formally smooth $A$-algebra.
\end{proposition}

\begin{proof}
Indeed, the notations $C$ and $\mathfrak{J}$ having the same meaning as in \sref{0.19.3.1}, let $u:B\to C/\mathfrak{J}$;
as $V$ is projective, $u_1$, the restriction of $u$ to $V$, lifts to a homomorphism of $A$-modules $v_1:V\to C$ such that $u_1$ factorizes as $V\xrightarrow{v_1}C\xrightarrow{\varphi}C/\mathfrak{J}$;
and $v_1$ extends to a homomorphism of $A$-algebras $v:\mathrm{S}_A^\bullet(V)\to C$, and $v$ coincides with $u$ in $V$, hence $\varphi\circ v=u$, which proves the proposition.
\end{proof}

\begin{corollary}[19.3.3]
\label{0.19.3.3}
If $A$ is a discrete ring, every polynomial algebra $B=A[T_\alpha]_{\alpha\in I}$, equipped with the discrete topology, is a formally smooth $A$-algebra.
\end{corollary}

\begin{proposition}[19.3.4]
\label{0.19.3.4}
Let $A$ be a topological ring, and let $B=A[[T_\alpha]]_{\alpha\in I}$ be a ring of formal power series $\sum c_\lambda T^\lambda$ (algebra large over $A$ of the monoid $\bb{N}^{(I)}$, identified as a topological $A$-module with the product $A^{\bb{N}^{(I)}}$).
If one equips $B$ with the product topology, $B$ is a formally smooth topological $A$-algebra.
\end{proposition}

Let $(\mathfrak{L}_\mu)_{\mu\in M}$ be a fundamental system of open ideals of $A$.
For every finite subset $H$ of $I$, every $\mu\in M$, and every integer $n$, let $\mathfrak{K}_{H,\mu,n}$ be the neighbourhood of $0$ in $B$ consisting of the families $(c_\lambda)$ such that $c_\lambda\in\mathfrak{L}_\mu$ whenever $\lambda=(\lambda_\alpha)_{\alpha\in I}$ satisfies $\lambda_\alpha=0$ for $\alpha\notin H$ and $\sum_{\alpha\in H}\lambda_\alpha\leq n$.
The $\mathfrak{K}_{H,\mu,n}$ plainly form a fundamental system of neighbourhoods of $0$ in $B$ and are ideals of $B$; hence the product topology is compatible with the $A$-algebra structure of $B$.

We first record the following lemma.

\begin{lemma}[19.3.4.1]
\label{0.19.3.4.1}
Let $E$ be a discrete $A$-algebra.
\begin{enumerate}
  \item[\rm(i)] If $f:B\to E$ is a continuous $A$-homomorphism, there exists a finite subset $H$ of $I$ such that $f(T_\alpha)=0$ for $\alpha\notin H$, and $f(T_\alpha)$ is nilpotent in $E$ for every $\alpha\in H$.
  \item[\rm(ii)] Conversely, let $H$ be a finite subset of $I$ and $(z_\alpha)_{\alpha\in H}$ a family of nilpotent elements of $E$.
  There exists a unique continuous $A$-homomorphism $g:B\to E$ such that $g(T_\alpha)=z_\alpha$ for $\alpha\in H$ and $g(T_\alpha)=0$ for $\alpha\notin H$.
\end{enumerate}
\end{lemma}

\oldpage[0\textsubscript{IV-1}]{81}

\begin{proof}
Assertion (i) follows immediately from the fact that $f^{-1}(0)$ is a neighbourhood of $0$ in the product $A$-module $A^{\bb{N}^{(I)}}$: thus $f(T^\lambda)=0$ except for finitely many $\lambda\in\bb{N}^{(I)}$.

For (ii), the polynomial ring $B'=A[T_\alpha]_{\alpha\in I}$ is dense in $B$.
Existence and uniqueness on $B'$ are immediate.
Continuity follows from nilpotence: if $z_\alpha^n=0$ for all $\alpha\in H$, then $g(T^\lambda)=0$ unless $\lambda_\alpha=0$ for $\alpha\notin H$ and $\lambda_\alpha\leq n$ for $\alpha\in H$, and there are only finitely many such $\lambda$.
\end{proof}

\begin{proof}[Proof of Proposition~{\rm\sref{0.19.3.4}}]
Let $C$ and $\mathfrak{J}$ have the meaning of \sref{0.19.3.1}, and let $u:B\to C/\mathfrak{J}$ be a continuous $A$-homomorphism.
By \sref{0.19.3.4.1}, $u(T_\alpha)=0$ except for the indices in a finite subset $H$ of $I$, and the elements $z_\alpha=u(T_\alpha)$, $\alpha\in H$, are nilpotent in $C/\mathfrak{J}$.
Since $\mathfrak{J}$ is nilpotent, choose nilpotent elements $x_\alpha\in C$ lifting the $z_\alpha$.
The lemma gives a continuous $A$-homomorphism $v:B\to C$ such that $v(T_\alpha)=x_\alpha$ for $\alpha\in H$ and $v(T_\alpha)=0$ otherwise.
Then $u$ is the composite $B\xrightarrow{v}C\to C/\mathfrak{J}$.
\end{proof}

\begin{proposition}[19.3.5]
\label{0.19.3.5}
Let $A$ be a topological ring.
\begin{enumerate}
  \item[(i)] $A$ is a formally smooth $A$-algebra.
  \item[(ii)] If $B$ is a formally smooth $A$-algebra and $C$ a formally smooth $B$-algebra, then $C$ is a formally smooth $A$-algebra.
  \item[(iii)] Let $B$ be a formally smooth $A$-algebra, $A'$ a topological $A$-algebra;
  then the topological $A'$-algebra $B\otimes_A A'$ ($\mathbf{0_I}$, 7.7.5 and 7.7.6) is formally smooth.
  \item[(iv)] Let $B$ be a topological $A$-algebra, $S$ (resp. $T$) a multiplicative subset of $A$ (resp. $B$) such that the canonical image of $S$ in $B$ is contained in $T$.
  If $B$ is a formally smooth $A$-algebra, then $T^{-1}B$ is an $S^{-1}A$-algebra formally smooth.
  \item[(v)] Let $B_i$ ($1\leq i\leq n$) be topological $A$-algebras.
  For $\prod_{i=1}^n B_i$ to be a formally smooth $A$-algebra, it is necessary and sufficient that each of the $B_i$ be so.
\end{enumerate}
\end{proposition}

\oldpage[0\textsubscript{IV-1}]{82}

\begin{proof}
\begin{enumerate}
  \item[(i)] If $C$ is a discrete $A$-algebra and
  $\varphi:C\to C/\mathfrak{J}$ is the canonical homomorphism onto any
  quotient $A$-algebra of $C$, the \emph{only} $A$-homomorphism from $A$ to
  $C/\mathfrak{J}$ is the composite
  $A\xrightarrow{\psi}C\xrightarrow{\varphi}C/\mathfrak{J}$, where $\psi$
  defines the $A$-algebra structure of $C$; since $\psi$ is continuous, the
  condition of \sref{0.19.3.1} is trivially satisfied.

  \item[(ii)] Let $\alpha:A\to B$ and $\beta:B\to C$ be the continuous
  homomorphisms defining respectively the $A$-algebra structure on $B$ and the
  $B$-algebra structure on $C$, so that $\beta\circ\alpha$ defines the
  $A$-algebra structure on $C$. Let $E$ be a discrete $A$-algebra,
  $\mathfrak{L}$ a nilpotent ideal of $E$, and
  $u:C\to E/\mathfrak{L}$ a continuous $A$-homomorphism. Since $B$ is a
  formally smooth $A$-algebra, $u\circ\beta:B\to E/\mathfrak{L}$ factorizes as
  $B\xrightarrow{v}E\to E/\mathfrak{L}$, where $v$ is a continuous
  $A$-homomorphism. The homomorphisms $v$ and $u\circ\beta$ then define on $E$
  and $E/\mathfrak{L}$ respectively topological $B$-algebra structures for
  which $E/\mathfrak{L}$ is still the discrete quotient algebra of $E$.
  Moreover, $u$ is a continuous $B$-homomorphism, and hence factorizes as
  $C\xrightarrow{w}E\to E/\mathfrak{L}$, where $w$ is a continuous
  $B$-homomorphism; since $w\circ\beta\circ\alpha$ is the homomorphism defining
  the $A$-algebra structure on $E$, $w$ is indeed a continuous
  $A$-homomorphism.

  \item[(iii)] Let $C$ be a discrete topological $A'$-algebra,
  $\mathfrak{J}$ a nilpotent ideal of $C$, and
  $u:B\otimes_A A'\to C/\mathfrak{J}$ a continuous $A'$-homomorphism. Composing
  $u$ with the canonical homomorphism
  $\rho:B\to B\otimes_A A'$ gives a continuous $A$-homomorphism, which by
  hypothesis factorizes as $B\xrightarrow{v}C\to C/\mathfrak{J}$. If
  $\sigma:A'\to B\otimes_A A'$ and $w:A'\to C$ are the canonical
  homomorphisms, equality of the composites $A\to B\xrightarrow{v}C$ and
  $A\to A'\xrightarrow{w}C$ gives a continuous ring homomorphism
  $f:B\otimes_A A'\to C$ such that $v=f\circ\rho$ and $w=f\circ\sigma$
  ($\mathbf{0_I}$, 7.7.6). Comparing the composites from $B$ and $A'$ shows
  that $u$ factorizes as
  $B\otimes_A A'\xrightarrow{f}C\to C/\mathfrak{J}$, proving~(iii).

  \item[(iv)] The topology on $S^{-1}A$ (resp. $T^{-1}B$) has as a fundamental
  system of neighbourhoods of $0$ the ideals $S^{-1}\mathfrak{J}$ (resp.
  $T^{-1}\mathfrak{R}$), where $\mathfrak{J}$ (resp. $\mathfrak{R}$) ranges
  over a fundamental system of open ideals of $A$ (resp. $B$)
  ($\mathbf{0_I}$, 7.6.1). Let $C$ be a discrete topological
  $S^{-1}A$-algebra, $\mathfrak{J}$ a nilpotent ideal of $C$, and
  $u:T^{-1}B\to C/\mathfrak{J}$ a continuous $S^{-1}A$-homomorphism. The
  composite $B\to T^{-1}B\xrightarrow{u}C/\mathfrak{J}$ factorizes as
  $B\xrightarrow{v}C\to C/\mathfrak{J}$. For every $t\in T$, $u(t/1)$ is
  invertible in $C/\mathfrak{J}$; since $\mathfrak{J}$ is nilpotent, $v(t)$ is
  invertible in $C$ ($\mathbf{0_I}$, 7.1.12). Thus $v$ factorizes through a
  continuous homomorphism $w:T^{-1}B\to C$ ($\mathbf{0_I}$, 7.6.6). The
  canonical diagrams show that $w$ is an $S^{-1}A$-homomorphism and that
  $u$ is the composite $T^{-1}B\xrightarrow{w}C\to C/\mathfrak{J}$.

  \item[(v)] Necessity follows from (iv), since each $B_i$ is the localization of
  $B=\prod_{i=1}^n B_i$ at its canonical idempotent. Conversely, let $C$ be a
  discrete $A$-algebra, $\mathfrak{J}$ a nilpotent ideal of $C$, and
  $u:B\to C/\mathfrak{J}$ a continuous $A$-homomorphism. The images under $u$
  of the canonical idempotents of $B$ form a finite family of pairwise
  orthogonal idempotents with sum $1$; lift this family to pairwise orthogonal
  idempotents $e_i$ of $C$ with sum $1$. Writing $C_i=e_iC$ and
  $\mathfrak{J}_i=e_i\mathfrak{J}$, the restriction of $u$ to $B_i$ is a
  continuous $A$-homomorphism $B_i\to C_i/\mathfrak{J}_i$. Formal smoothness of
  $B_i$ lifts it to a continuous $A$-homomorphism $B_i\to C_i$. The product of
  these lifts is a continuous $A$-homomorphism $B\to C$ lifting $u$.
\end{enumerate}
\end{proof}

\oldpage[0\textsubscript{IV-1}]{83}

\begin{proposition}[19.3.6]
\label{0.19.3.6}
Let $A$ be a topological ring, $B$ a topological $A$-algebra, $\widehat{A}$ and $\widehat{B}$ the separated completions of $A$ and $B$, so that $\widehat{B}$ is a topological $\widehat{A}$-algebra.
The following conditions are equivalent:
\begin{enumerate}
  \item[a)] $B$ is a formally smooth $A$-algebra.
  \item[b)] $\widehat{B}$ is a formally smooth $A$-algebra.
  \item[c)] $\widehat{B}$ is a formally smooth $\widehat{A}$-algebra.
\end{enumerate}
\end{proposition}

\begin{proof}
Of course, the structure of $\widehat{A}$-algebra on $\widehat{B}$ is defined by the homomorphism $\hat{\varphi}$, where $\varphi:A\to B$ is the homomorphism defining the structure of $A$-algebra on $B$.
As every discrete $A$-algebra $C$ is separated and complete, it is also a $\widehat{A}$-algebra (by canonical prolongation of the homomorphism from $A$ to $C$), and a continuous $A$-homomorphism from $B$ to $C$ gives by extension a canonical $\widehat{A}$-homomorphism (and \emph{a fortiori} an $A$-homomorphism) from $\widehat{B}$ to $C$ (in other words, every continuous $A$-homomorphism $B\to C$ factorizes as $B\to\widehat{B}\to C$ uniquely).
These remarks and the definition \sref{0.19.3.1} immediately imply the proposition.
\end{proof}

\begin{corollary}[19.3.7]
\label{0.19.3.7}
Under the hypotheses of \sref{0.19.3.5}, (iv), the topological $A\{S^{-1}\}$-algebra $B\{T^{-1}\}$ is formally smooth.
\end{corollary}

This follows from the definitions ($\mathbf{0_I}$, 7.6.1 and 7.6.7),
\sref{0.19.3.5},~(iv), and \sref{0.19.3.6}.

\begin{proposition}[19.3.8]
\label{0.19.3.8}
Let $A$ be a topological ring, $B$ a topological $A$-algebra, and suppose that for every open ideal $\mathfrak{R}$ of $B$, $\mathfrak{R}^2$ is also open.
Let $A'$, $B'$ be the rings $A$ and $B$ equipped with topologies finer than the initial topologies, and suppose that the canonical homomorphism $\varphi:A\to B$ is still a continuous homomorphism from $A'$ to $B'$.
Then, if $B'$ is a formally smooth $A'$-algebra, $B$ is a formally smooth $A$-algebra.
\end{proposition}

\begin{proof}
It suffices to apply the following lemma:
\end{proof}

\begin{lemma}[19.3.8.1]
\label{0.19.3.8.1}
Let $C$ be a discrete $A$-algebra, $\mathfrak{J}$ a nilpotent ideal of $C$.
Suppose that the square of every open ideal of $B$ is also open.
Then, if a composite homomorphism $v:B\xrightarrow{u}C\to C/\mathfrak{J}$ is continuous, the homomorphism $u$ is continuous.
\end{lemma}

\begin{proof}
Indeed, $\mathfrak{R}=u^{-1}(\mathfrak{J})$ is the kernel of $v$, and it is open by hypothesis;
as $\mathfrak{J}^n=0$, $\mathfrak{R}^n$ is contained in $\Ker(u)$;
but by induction on $h$, it follows from the hypothesis that $\mathfrak{R}^{2^h}$ is also open, hence $\mathfrak{R}^n$ for all $n$, and consequently $\Ker(u)$ is also open, which proves our assertion.
\end{proof}

Note that the hypothesis on $B$ means that the topology of $B$ is the
supremum of the $\mathfrak{R}$-preadic topologies, where $\mathfrak{R}$ ranges
over the open ideals of $B$. If $B$ is a preadic ring ($\mathbf{0_I}$, 7.1.9),
this condition is always satisfied.

\begin{env}[19.3.9]
\label{0.19.3.9}
We will now see that the property of being formally smooth implies properties of ``lifting'' of homomorphisms under more general conditions than those of the definition \sref{0.19.3.1}.
\end{env}

\oldpage[0\textsubscript{IV-1}]{84}

\begin{proposition}[19.3.10]
\label{0.19.3.10}
Let $A$ be a topological ring, $B$ a formally smooth topological $A$-algebra.
Let $C$ be a topological $A$-algebra, $\mathfrak{J}$ an ideal of $C$, satisfying the following conditions:
\begin{enumerate}
  \item[$1^\circ$] $C$ is metrizable and complete.
  \item[$2^\circ$] $\mathfrak{J}$ is closed and the sequence $(\mathfrak{J}^n)$ tends to $0$.
\end{enumerate}
Then, every continuous $A$-homomorphism $u:B\to C/\mathfrak{J}$ factorizes as $B\xrightarrow{v}C\to C/\mathfrak{J}$, where $v$ is a continuous $A$-homomorphism.
\end{proposition}

\begin{proof}
Let $(\mathfrak{Q}_n)$ be a decreasing sequence of ideals of $C$ forming a fundamental system of neighbourhoods of $0$ in this algebra;
as there exists $k$ such that $\mathfrak{J}^k\subset\mathfrak{Q}_n$, $(\mathfrak{J}+\mathfrak{Q}_n)/\mathfrak{Q}_n$ is nilpotent.
For all $n$, let $u_n$ be the continuous $A$-homomorphism $B\xrightarrow{u}C/\mathfrak{J}\to C/(\mathfrak{J}+\mathfrak{Q}_n)=(C/\mathfrak{Q}_n)/((\mathfrak{J}+\mathfrak{Q}_n)/\mathfrak{Q}_n)$;
by hypothesis $u_n$ factorizes as $B\xrightarrow{v_n}C/\mathfrak{Q}_n\xrightarrow{\varphi_n}C/(\mathfrak{J}+\mathfrak{Q}_n)$, where $v_n$ is a continuous $A$-homomorphism.
Let us show that one can choose step by step the $v_n$ so that $v_n$ factorizes as
\[
  B\xrightarrow{v_{n+1}}C/\mathfrak{Q}_{n+1}\to C/\mathfrak{Q}_n
\]
for all $n$ (in other words, the $v_n$ form a \emph{projective system} of homomorphisms).
This will follow from the following lemma:
\end{proof}

\begin{lemma}[19.3.10.1]
\label{0.19.3.10.1}
Let $B$ be a formally smooth $A$-algebra, $E$, $E'$ two discrete topological $A$-algebras, $\mathfrak{R}$ (resp. $\mathfrak{R}'$) a nilpotent ideal of $E$ (resp. $E'$), $f:E\to E'$ a surjective $A$-homomorphism such that $f(\mathfrak{R})\subset\mathfrak{R}'$, $f':E/\mathfrak{R}\to E'/\mathfrak{R}'$ the homomorphism induced by $f$ by passing to quotients.
Let $u:B\to E/\mathfrak{R}$ be a continuous $A$-homomorphism, $u'$ the composite homomorphism $B\xrightarrow{u}E/\mathfrak{R}\xrightarrow{f'}E'/\mathfrak{R}'$, and let $v':B\to E'$ be a continuous $A$-homomorphism such that $u'$ factorizes as $B\xrightarrow{v'}E'\to E'/\mathfrak{R}'$.
Then there exists a continuous $A$-homomorphism $v:B\to E$ such that $v'$ factorizes as $B\xrightarrow{v}E\xrightarrow{f}E'$.
\end{lemma}

\begin{proof}
In the commutative diagram
\[
\xymatrix{
  E \ar[r]^f \ar[d] & E' \ar[d] \\
  E/\mathfrak{R} \ar[r]_{f'} & E'/\mathfrak{R}'
}
\]
all the homomorphisms are surjective.
If $\mathfrak{L}=\Ker(f)$, then $E'$ is identified with $E/\mathfrak{L}$ and $\mathfrak{R}'$ with $(\mathfrak{R}+\mathfrak{L})/\mathfrak{L}$.
We shall use the following elementary lemma.
\end{proof}

\begin{lemma}[19.3.10.2]
\label{0.19.3.10.2}
Let $F$ be a ring, not necessarily commutative, and let $\mathfrak{a}$ and $\mathfrak{b}$ be two-sided ideals of $F$.
Then the fibre product
\[
  (F/\mathfrak{a})\times_{F/(\mathfrak{a}+\mathfrak{b})}(F/\mathfrak{b})
\]
is canonically identified with $F/(\mathfrak{a}\cap\mathfrak{b})$.
\end{lemma}

\begin{proof}
This is the special case of \sref{0.18.1.7} obtained by replacing, in diagram \sref{0.18.1.7.1}, $G$ by $F/(\mathfrak{a}\cap\mathfrak{b})$, $\mathfrak{J}'$ by $\mathfrak{b}/(\mathfrak{a}\cap\mathfrak{b})$, $F$ by $F/\mathfrak{b}$, $\mathfrak{R}$ by $(\mathfrak{a}+\mathfrak{b})/\mathfrak{b}$, $B$ by $F/(\mathfrak{a}+\mathfrak{b})$, and $\theta$ by the canonical isomorphism
\[
  (\mathfrak{a}+\mathfrak{b})/\mathfrak{b}\xrightarrow{\sim}\mathfrak{a}/(\mathfrak{a}\cap\mathfrak{b}).
\]
\end{proof}

\oldpage[0\textsubscript{IV-1}]{85}

\begin{proof}[Continuation of the proof of Lemma~{\rm\sref{0.19.3.10.1}}]
By \sref{0.19.3.10.2}, the commutativity of
\[
\xymatrix{
  B \ar[r]^{v'} \ar[d]_u & E'=E/\mathfrak{L} \ar[d] \\
  E/\mathfrak{R} \ar[r]_{f'} & E'/\mathfrak{R}'
}
\]
gives a homomorphism $w:B\to E/(\mathfrak{R}\cap\mathfrak{L})$ through which both $v'$ and $u$ factor.
Its kernel contains $\Ker(u)\cap\Ker(v')$, and is therefore open; thus $w$ is continuous.
The ideal $\mathfrak{R}\cap\mathfrak{L}$ is nilpotent, so formal smoothness of $B$ gives a continuous $A$-homomorphism $v:B\to E$ lifting $w$.
It follows at once that $f\circ v=v'$, as required.
\end{proof}

\begin{proof}[Continuation of the proof of Proposition~{\rm\sref{0.19.3.10}}]
Apply \sref{0.19.3.10.1} with
\[
  E=C/\mathfrak{Q}_{n+1},\quad E'=C/\mathfrak{Q}_n,
\]
\[
  \mathfrak{R}=(\mathfrak{J}+\mathfrak{Q}_{n+1})/\mathfrak{Q}_{n+1},\quad
  \mathfrak{R}'=(\mathfrak{J}+\mathfrak{Q}_n)/\mathfrak{Q}_n,
\]
and with $u,u',v'$ replaced respectively by $u_{n+1},u_n,v_n$.
It produces $v_{n+1}$ compatible with $v_n$.

Since $C$ is separated and complete, $C=\varprojlim C/\mathfrak{Q}_n$, and $v=\varprojlim v_n$ is a continuous $A$-homomorphism $B\to C$.
Because $\mathfrak{J}$ is closed, $\mathfrak{J}=\bigcap_n(\mathfrak{J}+\mathfrak{Q}_n)$.
Moreover, $C/\mathfrak{J}$ is metrizable and complete and the ideals $(\mathfrak{J}+\mathfrak{Q}_n)/\mathfrak{J}$ form a fundamental system of neighbourhoods of $0$; hence
\[
  C/\mathfrak{J}=\varprojlim C/(\mathfrak{J}+\mathfrak{Q}_n).
\]
Thus $\varprojlim\varphi_n$ is the canonical map $\varphi:C\to C/\mathfrak{J}$, while $\varprojlim u_n=u$.
Consequently $u=\varphi\circ v$.
\end{proof}

\begin{corollary}[19.3.11]
\label{0.19.3.11}
Let $A$ be a topological ring, $B$ a formally smooth topological $A$-algebra, $C$ a Noetherian complete local ring, $\mathfrak{J}$ an ideal distinct from $C$, $\varphi:A\to C$ a continuous homomorphism making $C$ a topological $A$-algebra.
Then every continuous $A$-homomorphism $u_0:B\to C_0=C/\mathfrak{J}$ factorizes as $B\xrightarrow{u}C\to C/\mathfrak{J}$, where $u$ is a continuous $A$-homomorphism.
\end{corollary}

\begin{proof}
All the conditions of \sref{0.19.3.10} are indeed satisfied ($\mathbf{0_I}$, 7.3.5).
\end{proof}

\begin{proposition}[19.3.12]
\label{0.19.3.12}
Let $A$ be a topological ring, $B$ a formally smooth $A$-algebra, $C$ a topological $A$-algebra, and $\mathfrak{J}$ an ideal of $C$, satisfying the following conditions:
\begin{enumerate}
  \item[$1^\circ$] There exists a fundamental system of open ideals $\mathfrak{Q}_\lambda$ of $C$, such that the $C_\lambda=C/\mathfrak{Q}_\lambda$ are Artinian rings and that $C\xrightarrow{\sim}\varprojlim C_\lambda$.
  \item[$2^\circ$] $\mathfrak{J}$ is closed in $C$ and topologically nilpotent.
  \item[$3^\circ$] The square of every open ideal of $B$ is open.
\end{enumerate}
Under these conditions, every continuous $A$-homomorphism $u:B\to C/\mathfrak{J}$ factorizes as $B\xrightarrow{v}C\to C/\mathfrak{J}$, where $v$ is a continuous $A$-homomorphism.
\end{proposition}

\begin{proof}
Let $\mathfrak{J}_\lambda=(\mathfrak{J}+\mathfrak{Q}_\lambda)/\mathfrak{Q}_\lambda$
be the canonical image of $\mathfrak{J}$ in $C_\lambda$, which is a nilpotent
ideal, and let $u_\lambda$ be the composite
$B\xrightarrow{u}C/\mathfrak{J}\to C_\lambda/\mathfrak{J}_\lambda
=C/(\mathfrak{J}+\mathfrak{Q}_\lambda)$. We first show that every
$A$-homomorphism $w:B\to C$ through which $u$ factorizes is necessarily
continuous. For every $\lambda$, there is an open ideal $\mathfrak{R}$ of $B$
such that
$u(\mathfrak{R})\subset(\mathfrak{J}+\mathfrak{Q}_\lambda)/\mathfrak{J}$,
and hence $w(\mathfrak{R})\subset\mathfrak{J}+\mathfrak{Q}_\lambda$. There is
therefore a power $\mathfrak{R}^{2^m}$ such that
$w(\mathfrak{R}^{2^m})\subset\mathfrak{Q}_\lambda$; this proves the assertion,
since $\mathfrak{R}^{2^m}$ is open in $B$. It is thus enough to find an
$A$-homomorphism $w$ with the preceding property without separately requiring
continuity.

Now
\[
  \operatorname{Hom}_A(B,C)
  =\varprojlim_\lambda\operatorname{Hom}_A(B,C_\lambda).
\]
For each $\lambda$, $\operatorname{Hom}_A(B,C_\lambda)$ is closed in the
$C_\lambda$-module $C_\lambda^B$ with its product topology, for which it is
linearly compact because $C_\lambda$ is Artinian. For each $\lambda$, let
$W_\lambda$
\oldpage[0\textsubscript{IV-1}]{86}
be the set of $w_\lambda\in\operatorname{Hom}_A(B,C_\lambda)$ such that
$u_\lambda$ factorizes as
$B\xrightarrow{w_\lambda}C_\lambda\to C_\lambda/\mathfrak{J}_\lambda$.
The set $W_\lambda$ is a closed affine subspace of the $C_\lambda$-module
$\operatorname{Hom}_A(B,C_\lambda)$ and is nonempty because $B$ is formally
smooth. In
$E_\lambda=\prod_{\mu\leq\lambda}\operatorname{Hom}_A(B,C_\mu)$, let
$U_\lambda$ be the closed affine subspace of the families
$(w_\mu)_{\mu\leq\lambda}$ such that
$w_\mu=\varphi_{\mu\lambda}\circ w_\lambda$ for $\mu\leq\lambda$ and
$w_\lambda\in W_\lambda$, where
$\varphi_{\mu\lambda}:C_\lambda\to C_\mu$ is the canonical homomorphism.
Finally let $V_\lambda$ be the closed affine subspace in
$\prod_\mu\operatorname{Hom}_A(B,C_\mu)$ equal to the product of $U_\lambda$
and the $\operatorname{Hom}_A(B,C_\mu)$ for $\mu\not\leq\lambda$. It remains to
observe that the intersection of the $V_\lambda$ is nonempty, for an element
$w$ of this intersection belongs by definition to
$\varprojlim_\lambda\operatorname{Hom}_A(B,C_\lambda)$.
\footnote{The printed proof has a gap here: sets of algebra homomorphisms are
not modules under pointwise addition, and for a general nilpotent ideal the
lift fibre need not be an affine subspace. Thus the stated linear-compactness
argument does not by itself prove that this intersection is nonempty; a
separate compatible-lifting argument is required.}

Moreover, since $\mathfrak{J}$ is closed and
$C=\varprojlim_\lambda C_\lambda$ is linearly compact,
$C/\mathfrak{J}$ identifies with
$\varprojlim_\lambda C_\lambda/\mathfrak{J}_\lambda$, and the canonical map
$\psi:C\to C/\mathfrak{J}$ with $\varprojlim_\lambda\psi_\lambda$, where
$\psi_\lambda:C_\lambda\to C_\lambda/\mathfrak{J}_\lambda$ is canonical.
As in \sref{0.19.3.10}, we conclude that $\psi\circ w=u$.
\end{proof}

\subsection{First criteria for formal smoothness}
\label{subsection:0.19.4}

\begin{proposition}[19.4.1]
\label{0.19.4.1}
Let $A$ be a topological ring, $B$ a topological $A$-algebra;
suppose that there exist two decreasing filtering families $(\mathfrak{J}_\alpha)_{\alpha\in I}$, $(\mathfrak{R}_\alpha)_{\alpha\in I}$ of ideals of $A$ and $B$ respectively, such that:
\begin{enumerate}
  \item[$1^\circ$] $(\mathfrak{J}_\alpha)$ tends to $0$ in $A$ and $(\mathfrak{R}_\alpha)$ tends to $0$ in $B$;
  \item[$2^\circ$] for all $\alpha\in I$, we have $\mathfrak{J}_\alpha B\subset\mathfrak{R}_\alpha$ (so that $B/\mathfrak{R}_\alpha$ is an $(A/\mathfrak{J}_\alpha)$-topological algebra);
  \item[$3^\circ$] for all $\alpha\in I$, $B/\mathfrak{R}_\alpha$ is a formally smooth $(A/\mathfrak{J}_\alpha)$-algebra.
\end{enumerate}
Then $B$ is a formally smooth $A$-algebra.
\end{proposition}

\begin{proof}
Let $C$ be a discrete $A$-algebra, $\mathfrak{L}$ a nilpotent ideal of $C$, and $\mathfrak{R}$ an open ideal of $B$;
by hypothesis there exists $\alpha\in I$ such that $\mathfrak{R}_\alpha\subset\mathfrak{R}$, so $B/\mathfrak{R}$ is a quotient of $B/\mathfrak{R}_\alpha$.
Every $A$-homomorphism $u':B/\mathfrak{R}\to C/\mathfrak{L}$ is also an
$(A/\mathfrak{J}_\alpha)$-homomorphism, so there exists an open ideal
$\mathfrak{R}'$ of $B$ such that
$\mathfrak{R}_\alpha\subset\mathfrak{R}'\subset\mathfrak{R}$ and the composite
$B/\mathfrak{R}_\alpha\to B/\mathfrak{R}\xrightarrow{u'}C/\mathfrak{L}$
factorizes as
$B/\mathfrak{R}_\alpha\to B/\mathfrak{R}'\to C\to C/\mathfrak{L}$, where the
second homomorphism is an $(A/\mathfrak{J}_\alpha)$-homomorphism. The conclusion
follows from the remark after \sref{0.19.3.1}.
\end{proof}

\begin{corollary}[19.4.2]
\label{0.19.4.2}
Let $A$ be a topological ring, $B$ a topological $A$-algebra, $(\mathfrak{J}_\alpha)_{\alpha\in I}$ a decreasing filtering family of ideals of $A$ tending to~$0$.
For $B$ to be a formally smooth $A$-algebra, it is necessary and sufficient that for all $\alpha\in I$, setting $A_\alpha=A/\mathfrak{J}_\alpha$, $B_\alpha=B/\mathfrak{J}_\alpha B$, $B_\alpha$ be a formally smooth $A_\alpha$-algebra.
\end{corollary}

\begin{proof}
The condition is sufficient by \sref{0.19.4.1}, and it is also necessary by \sref{0.19.3.5},~(iii).
\end{proof}

\begin{proposition}[19.4.3]
\label{0.19.4.3}
Let $A$ be a topological ring, $B$ a topological $A$-algebra.
Suppose that for every discrete $A$-algebra $C$ and every ideal $\mathfrak{J}$ of $C$ such that $\mathfrak{J}^2=0$, every continuous $A$-homomorphism $u:B\to C/\mathfrak{J}$ factorizes as $B\xrightarrow{v}C\to C/\mathfrak{J}$, where $v$ is a continuous $A$-homomorphism.
Then $B$ is a formally smooth $A$-algebra.
\end{proposition}

\begin{proof}
With the same notation, let $\mathfrak{R}$ be any nilpotent ideal of $C$,
and consider a continuous $A$-homomorphism $u':B\to C/\mathfrak{R}$.
Suppose $\mathfrak{R}^m=0$, and set $C_i=C/\mathfrak{R}^i$ for $1\leq i\leq m$, so that $C_m=C$, that $\mathfrak{R}^{i-1}/\mathfrak{R}^i$ is an ideal of square zero in $C_i$, and $C_i=C_{i+1}/\mathfrak{J}_{i+1}$;
the hypothesis then gives by induction on $i$ the existence of continuous $A$-homomorphisms $u_i:B\to C_i$ such that $u_1=u'$ and that $u_i$ factorizes as $B\xrightarrow{u_{i+1}}C_{i+1}\to C_{i+1}/\mathfrak{J}_{i+1}=C_i$;
whence the conclusion.
\end{proof}

\begin{proposition}[19.4.4]
\label{0.19.4.4}
Let $A$ be a topological ring, $B$ a topological $A$-algebra (commutative).
For $B$ to be a formally smooth $A$-algebra, it is necessary and sufficient that for
\oldpage[0\textsubscript{IV-1}]{87}
every discrete topological $B$-module $L$, annihilated by an open ideal of $B$, we have (cf.~\sref{0.18.5.1})
\[
  \operatorname{Exalcotop}_A(B,L)=0.
\]
\end{proposition}

\begin{proof}
Let $(\mathfrak{R}_\lambda)$ be a fundamental decreasing system of open ideals of $B$ and set $B_\lambda=B/\mathfrak{R}_\lambda$.
Consider a discrete topological $A$-algebra $C$ and an ideal $\mathfrak{J}$ of $C$ such that $\mathfrak{J}^2=0$, so that $C$ is an $A$-extension of $C/\mathfrak{J}$ by $\mathfrak{J}$;
suppose given an $A$-homomorphism $u:B_\lambda\to C/\mathfrak{J}$ and form the $A$-extension $E_\lambda$ of $B_\lambda$ by $\mathfrak{J}$, inverse image of $C$ by the homomorphism $u$ \sref{0.18.2.5};
it is a topological $A$-algebra for the discrete topology.
If $\operatorname{Exalcotop}_A(B,\mathfrak{J})=0$, there exists by definition \sref{0.18.5.1} a $\mu$ such that $\mathfrak{R}_\mu\subset\mathfrak{R}_\lambda$ and the inverse image extension $E_\mu$ is $A$-trivial;
but this means \sref{0.18.1.6} that there exists a continuous $A$-homomorphism $v:B_\mu\to E_\lambda$ such that the canonical homomorphism $B_\mu\to B_\lambda$ factorizes as $B_\mu\xrightarrow{v}E_\lambda\to B_\lambda$;
we conclude immediately that $B_\mu\to B_\lambda\xrightarrow{u}C/\mathfrak{J}$ factorizes as $B_\mu\xrightarrow{v}E_\lambda\to C\to C/\mathfrak{J}$,
and this proves, by virtue of \sref{0.19.3.1} and \sref{0.19.4.3}, that $B$ is a formally smooth $A$-algebra.
The converse is immediate, applying \sref{0.19.3.1} to the case where $C$ is a topological $A$-algebra which is an $A$-extension of $B_\lambda$ by $L$, and to the identical homomorphism $B_\lambda\to B_\lambda=C/L$.
\end{proof}

When $A$ and $B$ are discrete rings, the criterion of formal smoothness \sref{0.19.4.4} reduces to
\begin{equation}
\tag{19.4.4.1}
\label{0.19.4.4.1}
\operatorname{Exalcom}_A(B,L)=0 \qquad\text{for all $B$-modules $L$};
\end{equation}
in other words, \emph{every commutative $A$-extension of $B$ by a $B$-module must be $A$-trivial}.

\begin{corollary}[19.4.5]
\label{0.19.4.5}
Let $A$ be a topological ring (commutative), $B$ a topological $A$-algebra (commutative).
\begin{enumerate}
  \item[\rm(i)] Suppose that $B$ is a formally smooth $A$-algebra. Then, for every open ideal $\mathfrak{R}$ of $B$, every $(B/\mathfrak{R})$-module $L$, and every Hochschild $A$-bilinear symmetric $2$-cocycle $f$ of $(B/\mathfrak{R})\times(B/\mathfrak{R})$ in $L$ \sref{0.18.4.3}, there exists an open ideal $\mathfrak{R}'\subset\mathfrak{R}$ such that, if $\varphi:B/\mathfrak{R}'\to B/\mathfrak{R}$ is the canonical homomorphism, $f\circ(\varphi\times\varphi)$ is a Hochschild $A$-bilinear $2$-coboundary of $(B/\mathfrak{R}')\times(B/\mathfrak{R}')$ in $L$.
  \item[\rm(ii)] If $B$ is a formally projective $A$-module \sref{0.19.2.1} and if the condition of~{\rm(i)} is satisfied, $B$ is a formally smooth $A$-algebra.
\end{enumerate}
\end{corollary}

\begin{proof}
\begin{enumerate}
  \item[\rm(i)] The $2$-cocycle $f$ defines an $A$-extension of Hochschild $C$ of $B/\mathfrak{R}$ by $L$ \sref{0.18.4.3}.
  Applying \sref{0.19.3.1} to $C$, to the ideal of square zero $L$ of $C$, and to the identical homomorphism $B/\mathfrak{R}\to B/\mathfrak{R}=C/L$, we deduce the condition~(i) by virtue of \sref{0.18.4.3}.

  \item[\rm(ii)] Let us apply the criterion \sref{0.19.4.3}, considering an open ideal $\mathfrak{J}$ of $A$, an open ideal $\mathfrak{R}$ of $B$ such that $\mathfrak{J}B\subset\mathfrak{R}$, and an $(A/\mathfrak{J})$-extension $E$ of $B/\mathfrak{R}$ by $L$.
  As $B$ is a formally projective $A$-module, the canonical continuous $A$-linear map $\rho:B\to B/\mathfrak{R}$ factorizes as $B\xrightarrow{w}E\to B/\mathfrak{R}$, where $w$ is a continuous $A$-linear map \sref{0.19.2.1},
  which itself factorizes as $B\to B/\mathfrak{R}'\xrightarrow{w'}E$ where $\mathfrak{R}'$ is an open ideal of $B$ contained in $\mathfrak{R}$;
  replacing $\mathfrak{J}$ by a smaller ideal, one can suppose that $\mathfrak{J}B\subset\mathfrak{R}'$.
  Then the inverse image of the extension $E$ of $B/\mathfrak{R}$ by $L$, via the canonical homomorphism $B/\mathfrak{R}'\to B/\mathfrak{R}$, is an $(A/\mathfrak{J})$-extension \emph{of Hochschild} $E'$ of $B/\mathfrak{R}'$ by $L$;
  \oldpage[0\textsubscript{IV-1}]{88}
  applying~(i) to a cocycle defining this extension \sref{0.18.4.3}, we conclude that there exists an open ideal $\mathfrak{R}''\subset\mathfrak{R}'$ of $B$ such that the inverse image $E''$ of $E$ by $B/\mathfrak{R}''\to B/\mathfrak{R}$ is $A$-trivial, Q.E.D.
\end{enumerate}
\end{proof}

\begin{corollary}[19.4.6]
\label{0.19.4.6}
Let $A$ be a topological ring, $B$ a topological $A$-algebra which is a formally projective $A$-module.
Let $A'$ be an $A$-algebra equipped with the topology derived from that of $A$.
Suppose furthermore that $A'$ is a faithfully flat $A$-module, and that one of the following conditions is satisfied:
\begin{enumerate}
  \item[$1^\circ$] There exists a fundamental system $(\mathfrak{J}_\lambda)$ of open ideals of $A$ and a fundamental system $(\mathfrak{M}_\lambda)$ of open ideals of $B$, with the same index set, such that for all $\lambda$ we have $\mathfrak{J}_\lambda B\subset\mathfrak{M}_\lambda$ and $B/\mathfrak{M}_\lambda$ is a projective $(A/\mathfrak{J}_\lambda)$-module of finite type.
  \item[$2^\circ$] $A'$ is a projective $A$-module of finite type.
\end{enumerate}
Then, for $B'=B\otimes_A A'$ (equipped with the tensor product topology) to be a formally smooth $A'$-algebra, it is necessary and sufficient that $B$ be a formally smooth $A$-algebra.
\end{corollary}

\begin{proof}
The sufficiency of the condition is contained in \sref{0.19.3.5},~(iii), without any additional hypothesis on $B$ or $A'$.
To prove the converse, we apply the criterion \sref{0.19.4.5};
under hypothesis~$2^\circ$, we will again denote by $(\mathfrak{M}_\lambda)$ a fundamental system of open ideals of $B$, and for each $\lambda$, by $\mathfrak{J}_\lambda$ an open ideal of $A$ such that $\mathfrak{J}_\lambda B\subset\mathfrak{M}_\lambda$;
in both cases, we set $A_\lambda=A/\mathfrak{J}_\lambda$, $B_\lambda=B/\mathfrak{M}_\lambda$, $A'_\lambda=A_\lambda\otimes_A A'$, $B'_\lambda=B_\lambda\otimes_A A'$.
Let $f$ be a Hochschild $A_\lambda$-bilinear symmetric $2$-cocycle of $B_\lambda\times B_\lambda$ in a $B_\lambda$-module $L$;
by extension of scalars, we deduce a Hochschild $A'_\lambda$-bilinear symmetric $2$-cocycle $f'=f\otimes1$ of $B'_\lambda\times B'_\lambda$ in $L'=L\otimes_A A'$.
As $B'$ is by hypothesis a formally smooth $A'$-algebra, there exists by \sref{0.19.4.5},~(i), a $\mu$ such that $\mathfrak{M}_\mu\subset\mathfrak{M}_\lambda$ and such that, if $\varphi':B'_\mu\to B'_\lambda$ is the canonical homomorphism, the image $c'$ of $f'\circ(\varphi'\times\varphi')$ in the Hochschild group $\mathrm{H}^2_{A'_\mu}(B'_\mu,L')^s$ is zero;
it is clear that if $\varphi:B_\mu\to B_\lambda$ is the canonical homomorphism, and $c$ the class of $f\circ(\varphi\times\varphi)$ in the group $\mathrm{H}^2_{A_\mu}(B_\mu,L)^s$, then $c'$ is the canonical image of $c$.
Now, if $\mathrm{P}_\bullet$ is the complex relative to the rings $A_\mu$ and $B_\mu$ defined in \sref{0.18.4.5}, the analogous complex relative to the rings $A'_\mu$ and $B'_\mu$ is evidently $\mathrm{P}_\bullet\otimes_A A'$;
under hypothesis~$1^\circ$, the construction of $\mathrm{P}_\bullet$ shows that it is a projective $A_\mu$-module of finite type.
One concludes from Bourbaki, \emph{Alg.}, chap.~II, $3^{\text{e}}$~\'{e}d., \S\,5, n$^\circ$\,3, prop.~7 that, under both hypotheses, one has, up to canonical isomorphism,
\[
  \operatorname{Hom}_{A'_\mu}(\mathrm{P}_\bullet\otimes_A A',L\otimes_A A')
  =\operatorname{Hom}_{A_\mu}(\mathrm{P}_\bullet,L)\otimes_A A'.
\]
As $A'$ is a flat $A$-module, one has by \sref{0.18.4.5}
\[
  \mathrm{H}^2_{A'_\mu}(B'_\mu,L')^s = (\mathrm{H}^2_{A_\mu}(B_\mu,L)^s)\otimes_A A'
\]
and one can therefore write $c'=c\otimes 1$;
but as $A'$ is a faithfully flat $A$-module, the hypothesis $c'=0$ implies $c=0$, which completes the proof.
\end{proof}

\begin{proposition}[19.4.7]
\label{0.19.4.7}
Let $A$ be a preadmissible topological ring, $\mathfrak{J}$ an ideal of definition of $A$ ($\mathbf{0_I}$, 7.1.2), $B$ a topological $A$-algebra which is a formally projective $A$-module.
Consider the topological quotient rings $A_0=A/\mathfrak{J}$, $B_0=B/\mathfrak{J}B$;
then, for $B$ to be a formally smooth $A$-algebra, it is necessary and sufficient that $B_0$ be a formally smooth $A_0$-algebra.
\end{proposition}

\begin{proof}
\oldpage[0\textsubscript{IV-1}]{89}
The necessity of the condition follows from \sref{0.19.4.2}.
To see that it is sufficient, we first note that, by considering a fundamental system of open neighbourhoods of $0$ in $A$ formed by ideals $\mathfrak{J}_\lambda$ contained in $\mathfrak{J}$, one can by virtue of \sref{0.19.4.2} reduce to the case where $A$ is discrete, since $B/\mathfrak{J}_\lambda B$ is a formally projective $(A/\mathfrak{J}_\lambda)$-module \sref{0.19.2.5};
by definition of a preadmissible ring, $\mathfrak{J}$ is then nilpotent.
It furthermore suffices to prove the proposition when $\mathfrak{J}^2=0$, for if $\mathfrak{J}^m=0$, one applies it successively to the rings $A_k=A/\mathfrak{J}^k$ and $B_k=B/\mathfrak{J}^k B$ and to the ideal $\mathfrak{J}^{k-1}/\mathfrak{J}^k$ of $A_k$ for $k=2,3,\ldots,m$, noting \sref{0.19.2.5} that $B_k=B\otimes_A A_k$ is a formally projective $A_k$-module.
Let us apply the criterion \sref{0.19.4.5},~(ii), considering an open ideal $\mathfrak{R}$ of $B$ and a Hochschild $A$-bilinear symmetric $2$-cocycle $f$ of $(B/\mathfrak{R})\times(B/\mathfrak{R})$ in a $(B/\mathfrak{R})$-module $L$.
Consider first the special case where $\mathfrak{J}L=0$, so that $L$ can also be viewed as a $(B_0/\mathfrak{R}B_0)$-module, and $f$ factorizes as
\[
  (B/\mathfrak{R})\times(B/\mathfrak{R}) \to (B_0/\mathfrak{R}B_0)\times(B_0/\mathfrak{R}B_0) \xrightarrow{f_0} L
\]
where $f_0$ is a Hochschild bilinear symmetric $2$-cocycle.
By hypothesis, there exists an open ideal $\mathfrak{R}'\subset\mathfrak{R}$ in $B$ and an $A_0$-linear map $g_0:B_0/\mathfrak{R}'B_0\to L$ such that $f_0(\varphi_0(x),\varphi_0(y))=xg_0(y)-g_0(xy)+g_0(x)y$ for $x,y$ in $B_0/\mathfrak{R}'B_0$, where $\varphi_0:B_0/\mathfrak{R}'B_0\to B_0/\mathfrak{R}B_0$ is the canonical homomorphism.
One concludes immediately that the composite $A$-linear map $g:B/\mathfrak{R}'\to B_0/\mathfrak{R}'B_0\xrightarrow{g_0}L$ is such that $dg=f\circ(\varphi\times\varphi)$, where $\varphi:B/\mathfrak{R}'\to B/\mathfrak{R}$ is the canonical homomorphism.

Let us now pass to the general case, and consider first the $(B/\mathfrak{R})$-module $L/\mathfrak{J}L=L'$, for which $\mathfrak{J}L'=0$;
if $f'$ is the composite $A$-bilinear map $(B/\mathfrak{R})\times(B/\mathfrak{R})\xrightarrow{f}L\to L'$, $f'$ is again a symmetric Hochschild $2$-cocycle, and by the preceding, there exists an open ideal $\mathfrak{R}'\subset\mathfrak{R}$ in $B$ and an $A$-linear map $g':B/\mathfrak{R}'\to L'$ satisfying $dg'=f'\circ(\varphi'\times\varphi')$ for the canonical map $\varphi':B/\mathfrak{R}'\to B/\mathfrak{R}$.
As $B$ is a formally projective $A$-module, there exists an open ideal $\mathfrak{R}'_1\subset\mathfrak{R}'$ and an $A$-linear map $g_1:B/\mathfrak{R}'_1\to L$ lifting $B/\mathfrak{R}'_1\to B/\mathfrak{R}'\xrightarrow{g'}L'$.
Let $\varphi_1:B/\mathfrak{R}'_1\to B/\mathfrak{R}$ be the canonical homomorphism. Consider then the Hochschild $2$-cocycle
\[
  f_1(x,y) = f(\varphi_1(x),\varphi_1(y)) - xg_1(y) + g_1(xy) - g_1(x)y,
\]
a symmetric $A$-bilinear map of $(B/\mathfrak{R}'_1)\times(B/\mathfrak{R}'_1)$ in $L$.
The choice of $g_1$ implies that $f_1$ takes its values in $\mathfrak{J}L$.
As $\mathfrak{J}(\mathfrak{J}L)=0$, one can again apply the first case, so there exists an open ideal $\mathfrak{R}''\subset\mathfrak{R}'_1$ and an $A$-linear map $g_2:B/\mathfrak{R}''\to\mathfrak{J}L$ such that
\[
  f_1(\varphi_2(x),\varphi_2(y)) = xg_2(y) - g_2(xy) + g_2(x)y
\]
in $(B/\mathfrak{R}'')\times(B/\mathfrak{R}'')$, where $\varphi_2:B/\mathfrak{R}''\to B/\mathfrak{R}'_1$ is the canonical map;
the $A$-linear map $g=g_2+g_1\circ\varphi_2:B/\mathfrak{R}''\to L$ therefore satisfies $dg=f\circ(\varphi\times\varphi)$ for the canonical map $\varphi:B/\mathfrak{R}''\to B/\mathfrak{R}$.
\end{proof}

\subsection{Formal smoothness and associated graded rings}
\label{subsection:0.19.5}

\begin{env}[19.5.1]
\label{0.19.5.1}
\oldpage[0\textsubscript{IV-1}]{90}
Let $C$ be a topological ring (commutative), $V$ a topological $C$-module, and consider the symmetric algebra $\mathbf{S}_C^*(V)=\bigoplus_{n\geq 0}\mathbf{S}_C^n(V)$, which we will canonically equip with a \emph{linear topology}, compatible with its $C$-algebra structure.
For this, consider an open submodule $U$ of $V$, and the \emph{graded ideal} $U.\mathbf{S}_C^*(V)$ it generates in $\mathbf{S}_C^*(V)$;
we take as fundamental system of neighbourhoods of $0$ in $\mathbf{S}_C^*(V)$ the set of sums $\mathfrak{a}.\mathbf{S}_C^*(V)+U.\mathbf{S}_C^*(V)$, where $\mathfrak{a}$ (resp.\ $U$) runs through a fundamental system of open ideals (resp.\ of open submodules) of $C$ (resp.\ $V$).
We note that if the topology of $V$ is less fine than the topology derived from that of $C$, one can restrict to pairs $(\mathfrak{a},U)$ such that $\mathfrak{a}V\subset U$, so $\mathfrak{a}.\mathbf{S}_C^*(V)+U.\mathbf{S}_C^*(V)=\mathfrak{a}.1_C+U.\mathbf{S}_C^*(V)$;
in this case the topology induced on each $\mathbf{S}_C^n(V)$ for $n\geq 1$ has as a fundamental system of neighbourhoods of $0$ the $U.\mathbf{S}_C^{n-1}(V)$, where $U$ runs through the open submodules of $V$;
in particular, on $V=\mathbf{S}_C^1(V)$, it coincides with the given topology (in general it is less fine than the latter).
Furthermore, in all cases, the \emph{topological algebra} $\mathbf{S}_C^*(V)$ thus defined is, for the categories of topological $C$-modules and of topological $C$-algebras, the solution of the same universal problem as for the categories of $C$-modules and $C$-algebras.
Indeed, let $E$ be a topological $C$-algebra, $u:V\to E$ a homomorphism of $C$-modules, $\mathbf{S}(u)=\widetilde{u}$ its canonical extension to $\mathbf{S}_C^*(V)$.
Suppose $u$ continuous; then, if $\mathfrak{b}$ is an open ideal of $E$, its inverse image $U=u^{-1}(\mathfrak{b})$ is an open $C$-submodule of $V$, and the image by $\widetilde{u}$ of $U.\mathbf{S}_C^*(V)$ is therefore contained in $\mathfrak{b}$;
as moreover there exists an open ideal $\mathfrak{a}$ of $C$ such that $\mathfrak{a}.E\subset\mathfrak{b}$, hence $\widetilde{u}(\mathfrak{a}.\mathbf{S}_C^*(V))\subset\mathfrak{b}$, this proves that $\widetilde{u}$ is continuous.
Conversely, if $\widetilde{u}$ is continuous, and if $\mathfrak{b}$ is an open ideal of $E$, there exists an open submodule $U$ of $V$ such that $\widetilde{u}(U.\mathbf{S}_C^*(V))\subset\mathfrak{b}$, and in particular $\widetilde{u}(U.1_C)\subset\mathfrak{b}$, that is $u(U)\subset\mathfrak{b}$, so $u$ is continuous.
We also recall that one has a canonical isomorphism of topological $C$-modules (discrete)

\[
\mathbf{S}_C^n(V)/U.\mathbf{S}_C^{n-1}(V) \xrightarrow{\sim} \mathbf{S}_C^n(V/U). \tag{19.5.1.1}
\]

\label{0.19.5.1.1}
\end{env}

\begin{env}[19.5.2]
\label{0.19.5.2}
Let $A$ be a topological ring, $B$ a topological $A$-algebra (commutative), $\mathfrak{J}$ an ideal of $B$ (not necessarily open or closed);
in all that follows, we equip $\mathfrak{J}$ and the $\mathfrak{J}^n$ ($n\geq 2$) with the topology induced from that of $B$, the quotients $C=B/\mathfrak{J}$, $\mathfrak{J}^n/\mathfrak{J}^{n+1}=\gr_\mathfrak{J}^n(B)$ with the quotient topology, so that the $\mathfrak{J}^n/\mathfrak{J}^{n+1}$ are topological $C$-modules;
the canonical injection $\mathfrak{J}/\mathfrak{J}^2\to\gr_\mathfrak{J}^*(B)$ extends to a homomorphism (not topological at first) of $C$-algebras

\[
\varphi:\mathbf{S}_C^*(\mathfrak{J}/\mathfrak{J}^2) \to \gr_\mathfrak{J}^*(B)
\]

which for each $n$ gives a surjective homomorphism of $C$-modules

\[
\varphi_n:\mathbf{S}_C^n(\mathfrak{J}/\mathfrak{J}^2) \to \mathfrak{J}^n/\mathfrak{J}^{n+1}=\gr_\mathfrak{J}^n(B). \tag{19.5.2.1}
\]

\label{0.19.5.2.1}

\oldpage[0\textsubscript{IV-1}]{91}
When $\mathbf{S}_C^n(\mathfrak{J}/\mathfrak{J}^2)$ is equipped with the topology defined in \sref{0.19.5.1}, the homomorphisms $\varphi_n$ are continuous, by virtue of the universal property \sref{0.19.5.1} of $\mathbf{S}_C^*(\mathfrak{J}/\mathfrak{J}^2)$ applied to the topological algebra $E=\gr_{\mathfrak{J}/\mathfrak{J}^{n+1}}(B/\mathfrak{J}^{n+1})$ equipped with the product topology of those of the $\mathfrak{J}^k/\mathfrak{J}^{k+1}$, and to the canonical injection $\mathfrak{J}/\mathfrak{J}^2\to E$.
Note that here the topology of $\mathfrak{J}/\mathfrak{J}^2$ is less fine than that of $C$ (this latter topology on $\mathfrak{J}/\mathfrak{J}^2$ being also the topology derived from that of $B$).
\end{env}

\begin{theorem}[19.5.3]
\label{0.19.5.3}
Let $A$ be a topological ring, $B$ a topological $A$-algebra, $\mathfrak{J}$ an ideal of $B$, $C=B/\mathfrak{J}$ the topological quotient $A$-algebra.
Suppose that the discrete $A$-algebra $C$ is formally smooth.
Then:
\begin{enumerate}
  \item[\rm(i)] If $B$ is a formally smooth $A$-algebra, $\mathfrak{J}/\mathfrak{J}^2$ is a topologically formally projective $C$-module \sref{0.19.2.1}.
  \item[\rm(ii)] If $B$ is a formally smooth $A$-algebra and a preadmissible ring ($\mathbf{0_I}$, 7.1.2), the homomorphisms $\varphi_n$ \sref{0.19.5.2} are formal bimorphisms \sref{0.19.1.2}.
  \item[\rm(iii)] Conversely, suppose that $B$ is preadmissible, that the sequence $(\mathfrak{J}^n)$ tends to $0$ in $B$, that $\mathfrak{J}/\mathfrak{J}^2$ is a formally projective $C$-module, and that the $\varphi_n$ are formal bimorphisms. Then $B$ is a formally smooth $A$-algebra.
\end{enumerate}
\end{theorem}

The proof of this theorem is long and encumbered with technical details;
we will therefore begin by proving a simpler corollary, where the guiding ideas appear more clearly;
this corollary is moreover the particular case of the theorem \sref{0.19.5.3} that will be most frequently used in the sequel.

\begin{corollary}[19.5.4]
\label{0.19.5.4}
Let $A$ be a topological ring, $B$ a topological $A$-algebra, $\mathfrak{J}$ an ideal of $B$ such that the topology of $B$ is the $\mathfrak{J}$-preadic topology.
Suppose that the discrete $A$-algebra $C=B/\mathfrak{J}$ is formally smooth.
Then the three following conditions are equivalent:
\begin{enumerate}
  \item[a)] $B$ is a formally smooth $A$-algebra.
  \item[b)] $\mathfrak{J}/\mathfrak{J}^2$ is a projective $C$-module and the canonical homomorphism
  
\[
\varphi:\mathbf{S}_C^*(\mathfrak{J}/\mathfrak{J}^2) \to \gr_\mathfrak{J}(B)
\]

  is bijective.
  \item[c)] The separated completion $\widehat{B}$ of $B$ is a topological $\widehat{A}$-algebra isomorphic to an $\widehat{A}$-algebra of the form $\widehat{B'}$, where $B'=\mathbf{S}_C^*(V)$, $V$ being a projective $C$-module and $B'$ being equipped with the ${B'}^+$-preadic topology, where ${B'}^+$ is the ideal of augmentation.
\end{enumerate}
\end{corollary}

\begin{env}[19.5.4.1]
\label{0.19.5.4.1}
Let us first prove that $a)$ implies that $\mathfrak{J}/\mathfrak{J}^2$ is a projective $C$-module.
Let $P$ and $Q$ be two $C$-modules, $u:P\to Q$ a surjective $C$-homomorphism, and $v:\mathfrak{J}/\mathfrak{J}^2\to Q$ a $C$-homomorphism.
The ring $B/\mathfrak{J}^2$ being viewed as a $B$-extension of $C$ by $\mathfrak{J}/\mathfrak{J}^2$, we deduce from $v$ a $B$-extension $E=(B/\mathfrak{J}^2)\oplus_{(\mathfrak{J}/\mathfrak{J}^2)}Q$ of $C$ by $Q$ \sref{0.18.2.8}.
As $C$ is a discrete formally smooth $A$-algebra, the extension $E$ is $A$-trivial \sref{0.19.4.4} and can therefore be identified with $\mathrm{D}_C(Q)$ \sref{0.18.2.3}.
The surjective homomorphism $u:P\to Q$ then canonically defines a surjective $A$-homomorphism $\psi:\mathrm{D}_C(P)\to\mathrm{D}_C(Q)$ \sref{0.18.2.9} whose kernel is an ideal $\mathfrak{L}$ of the extension $E'=\mathrm{D}_C(P)$, contained in $P$, and \emph{a fortiori} of square zero.
Let $f:B\to E=E'/\mathfrak{L}$ be the homomorphism defining the $B$-algebra structure of $E$: as $\mathfrak{L}$ is of square zero and $B$ is a formally smooth $A$-algebra,
\oldpage[0\textsubscript{IV-1}]{92}
$f$ factorizes as $B\xrightarrow{g}E'\to E'/\mathfrak{L}$, where $g$ is a continuous $A$-homomorphism.
The diagram

\[
\xymatrix{
    E' \ar[r]^{\psi} & E \\
    B \ar[u]^{g} \ar[r]_{\theta} & B/\mathfrak{J}^2 \ar[u]_{v'}
  }
\]

where $v'$ is deduced from $v$ \sref{0.18.2.8}, is commutative.
Moreover, the image of $\mathfrak{J}$ by $\psi\circ g$ is contained in $Q$, so the image of $\mathfrak{J}$ by $g$ is contained in $P$, and the image of $\mathfrak{J}^2$ by $g$ is zero.
Restricting $g$ and $\theta$ to $\mathfrak{J}$, one obtains a commutative diagram

\[
\xymatrix{
    P \ar[r]^{u} & Q \\
    \mathfrak{J}/\mathfrak{J}^2 \ar[u]^{w} \ar[ur]_{v}
  }
\]

which proves that the $C$-module $\mathfrak{J}/\mathfrak{J}^2$ is projective.
\end{env}

\begin{env}[19.5.4.2]
\label{0.19.5.4.2}
Let us prove secondly that $a)$ implies that $\varphi$ is bijective, which will complete the proof that $a)$ implies $b)$.
Set $E_n=B/\mathfrak{J}^{n+1}$, and denote by $F_n$ the quotient of $\mathbf{S}_C^*(\mathfrak{J}/\mathfrak{J}^2)$ by the $(n+1)$-th power of its ideal of augmentation.
As $\mathfrak{J}/\mathfrak{J}^{n+1}$ is nilpotent in $E_n$ and $C$ is a discrete formally smooth $A$-algebra, the identity automorphism of $C$ factorizes as $C\xrightarrow{f}E_n\to C=E_n/(\mathfrak{J}/\mathfrak{J}^{n+1})$ where $f$ is an $A$-homomorphism;
 moreover, $\mathfrak{J}/\mathfrak{J}^2$ being a projective $C$-module, the identity of $\mathfrak{J}/\mathfrak{J}^2$ factorizes as $\mathfrak{J}/\mathfrak{J}^2\xrightarrow{g}\mathfrak{J}/\mathfrak{J}^{n+1}\to\mathfrak{J}/\mathfrak{J}^2$, and from $f$ and $g$ one obtains canonically a homomorphism of $C$-algebras $\mathbf{S}_C^*(\mathfrak{J}/\mathfrak{J}^2)\to E_n$ which annihilates on the $(n+1)$-th power of the ideal of augmentation, whence, passing to the quotient, a surjective homomorphism of algebras $v:F_n\to E_n$ such that $\gr^0(v)$ and $\gr^1(v)$ are the identity automorphisms of $C$ and of $\mathfrak{J}/\mathfrak{J}^2$.
By definition of the canonical homomorphism $\varphi$, we see that $\gr^j(v)=\varphi_j$ for all $j\leq n$.
Note now that the kernel $\mathfrak{R}$ is a nilpotent ideal of $F_n$, so that $E_n$ is identified with $F_n/\mathfrak{R}$.
As $B$ is a formally smooth $A$-algebra, the canonical $A$-homomorphism $p_n:B\to E_n=B/\mathfrak{J}^{n+1}$, which is continuous, factorizes as $B\xrightarrow{w}F_n\xrightarrow{v}E_n$, where $w$ is a continuous $A$-homomorphism;
as $\gr^0(v)$ is the identity, $w(\mathfrak{J})$ is contained in the ideal of augmentation of $F_n$, hence $w(\mathfrak{J}^{n+1})=0$,
so $w$ factorizes as $B\xrightarrow{p_n}E_n\xrightarrow{w'}F_n$, and the composite $E_n\xrightarrow{w'}F_n\xrightarrow{v}E_n$ is the identity.
Moreover, since $\gr^0(v)$ and $\gr^1(v)$ are the identity automorphisms of $C$ and $\mathfrak{J}/\mathfrak{J}^2$, the same is true of $\gr^0(w')$ and $\gr^1(w')$.
Since $\gr^*(E_n)$ is generated by $\gr^1(E_n)$, the composite
\[
  \gr^j(F_n)\xrightarrow{\varphi_j}\gr^j(E_n)
  \xrightarrow{\gr^j(w')}\gr^j(F_n)
\]
is the identity for every $j\leq n$, because this is true for $j=0$ and $j=1$.
Taking $j=n$ proves that $\varphi_n$ is injective, which completes the proof that $a)$ implies $b)$.
\end{env}

\begin{env}[19.5.4.3]
\label{0.19.5.4.3}
\oldpage[0\textsubscript{IV-1}]{93}
Let us next prove that $b)$ implies $a)$.
The same reasoning as at the beginning of \sref{0.19.5.4.2} proves the existence of a surjective $A$-homomorphism of algebras $v:F_n\to E_n$ such that $\gr^j(v)=\varphi_j$ for all $j$; since $\varphi_j$ is bijective for every $j$ and the filtrations of $F_n$ and $E_n$ are finite, we conclude that $v$ is \emph{bijective} (Bourbaki, \emph{Alg.\ comm.}, chap.~III, \S\,2, n$^\circ$\,8, cor.~3 of th.~1).
Let now $G$ be a discrete topological $A$-algebra, $\mathfrak{R}$ an ideal of square zero in $G$, $f:B\to G/\mathfrak{R}$ a continuous $A$-homomorphism of algebras.
As $G$ is discrete, there exists an integer $m$ such that $f$ annihilates on $\mathfrak{J}^m$, so $f$ factorizes as $B\to E_n\xrightarrow{f_n}G/\mathfrak{R}$, where we take $n=2m$.
One has therefore by composition a continuous $A$-homomorphism $r:C\to F_n\xrightarrow{v}E_n\xrightarrow{f_n}G/\mathfrak{R}$, and as $C$ is a discrete formally smooth $A$-algebra, $r$ factorizes as $C\xrightarrow{r'}G\to G/\mathfrak{R}$, so $G$ is thus equipped with a structure of $C$-algebra.
On the other hand, the restriction to $\mathfrak{J}/\mathfrak{J}^2$ of the $C$-homomorphism $f'_n:F_n\xrightarrow{v}E_n\xrightarrow{f_n}G/\mathfrak{R}$ is a $C$-linear map $g:\mathfrak{J}/\mathfrak{J}^2\to G/\mathfrak{R}$.
As $\mathfrak{J}/\mathfrak{J}^2$ is a projective $C$-module, $g$ factorizes as $\mathfrak{J}/\mathfrak{J}^2\xrightarrow{h}G\to G/\mathfrak{R}$, where $h$ is a $C$-linear map.
By extension, one deduces a homomorphism of $C$-algebras $w:\mathbf{S}_C^*(\mathfrak{J}/\mathfrak{J}^2)\to G$, and by construction every element of degree $m$ has as its image under $w$ an element of $\mathfrak{R}$, so every element of degree $n=2m$ has zero image since $\mathfrak{R}^2=0$;
in other words, $w$ factorizes as $\mathbf{S}_C^*(\mathfrak{J}/\mathfrak{J}^2)\to F_n\xrightarrow{w'}G$.
By construction, the composite $F_n\xrightarrow{w'}G\xrightarrow{\theta}G/\mathfrak{R}$ coincides with $f'_n$ in $C$ and in $\mathfrak{J}/\mathfrak{J}^2$, hence is equal to $f'_n$.
Finally, the composite $B\to E_n\xrightarrow{w'\circ v^{-1}}G$ being equal to $f$, one sees that $B$ is a formally smooth $A$-algebra.
\end{env}

\begin{env}[19.5.4.4]
\label{0.19.5.4.4}
It remains to prove the equivalence of $a)$ and $c)$.
In the first place, $c)$ implies that $B'$ is a formally smooth $C$-algebra for the discrete topology \sref{0.19.3.2}, hence also for the ${B'}^+$-preadic topology \sref{0.19.3.8}.
As $C$ is a formally smooth $A$-algebra, $B'$ is also a formally smooth $A$-algebra \sref{0.19.3.5},~{\rm(ii)}, and finally $\widehat{B'}$ is a formally smooth $\widehat{A}$-algebra \sref{0.19.3.6}; hence $B$ is a formally smooth $A$-algebra \sref{0.19.3.6}.
It remains to see that $a)$ implies $c)$.
Note first that $C$ is a formally smooth $A$-algebra, the identical homomorphism $C\to B/\mathfrak{J}$ factorizes as $C\to\widehat{B}\to B/\mathfrak{J}$, where $C\to\widehat{B}$ is an $A$-homomorphism \sref{0.19.3.10};
hence $\widehat{B}$, and therefore all the $B/\mathfrak{J}^{n+1}$, are equipped with structures of $C$-algebras.
On the other hand, as $\mathfrak{J}/\mathfrak{J}^2$ is a projective $C$-module by $b)$, the canonical injection $\mathfrak{J}/\mathfrak{J}^2\to B/\mathfrak{J}^2$ allows one to form a projective system of $C$-homomorphisms $u_n:\mathfrak{J}/\mathfrak{J}^2\to B/\mathfrak{J}^{n+1}$ for $n\geq 1$, hence also (by the universal property of the symmetric algebra) a projective system of homomorphisms of $C$-algebras $v_n:\mathbf{S}_C^*(\mathfrak{J}/\mathfrak{J}^2)=B'\to B/\mathfrak{J}^{n+1}$;
moreover it is clear that $v_n$ annihilates on $({B'}^+)^{n+1}$ and as $\gr^0(v_n)=\varphi_0$ and $\gr^1(v_n)=\varphi_1$, one has $\gr^j(v_n)=\varphi_j$ for $0\leq j\leq n$.
As the $\varphi_j$ are isomorphisms by $b)$, and the filtrations of $B'/({B'}^+)^{n+1}$ and of $B/\mathfrak{J}^{n+1}$ are finite, one concludes that $v_n$ is \emph{bijective} for
\oldpage[0\textsubscript{IV-1}]{94}
all $n$ (Bourbaki, \emph{Alg.\ comm.}, chap.~III, \S\,2, n$^\circ$\,8, cor.~3 of th.~1); whence $c)$ by passage to the projective limit.
\end{env}

\begin{remark}[19.5.5]
\label{0.19.5.5}
Under the general hypotheses of \sref{0.19.5.4}, suppose furthermore that $\mathfrak{J}$ is a \emph{maximal} ideal of $B$, so that $C=B/\mathfrak{J}$ is a field, and that $\mathfrak{J}/\mathfrak{J}^2$ is a vector space of \emph{finite} dimension.
Then the conditions $a)$, $b)$ and $c)$ of \sref{0.19.5.4} are also equivalent to the following:
\begin{enumerate}
  \item[d)] \emph{Given} a polynomial algebra $E=C[T_1,\ldots,T_n]$ over the field $C$, and denoting by $\mathfrak{n}$ the maximal ideal of $E$ generated by the $T_i$ ($1\leq i\leq n$), then for every ideal $\mathfrak{b}$ of $E$ containing a power $\mathfrak{n}^k$, every homomorphism $v:B\to E/\mathfrak{b}$ of $C$-augmented $A$-algebras factorizes as
  \[
    B\xrightarrow{w}E/\mathfrak{n}^k\to E/\mathfrak{b},
  \]
  where $w$ is an $A$-homomorphism.
\end{enumerate}
Indeed, $\mathfrak{J}/\mathfrak{J}^2$ being here a free $C$-module, it suffices to verify that condition $d)$ implies that the canonical homomorphism
$\varphi:\mathbf{S}_C^*(\mathfrak{J}/\mathfrak{J}^2)\to\gr_\mathfrak{J}^*(B)$
is bijective.
Here $\mathbf{S}_C^*(\mathfrak{J}/\mathfrak{J}^2)$ identifies with the polynomial algebra $E=C[T_1,\ldots,T_n]$, where $n=\rg_C(\mathfrak{J}/\mathfrak{J}^2)$, and the ideal of augmentation of $E$ is the ideal $\mathfrak{n}$ generated by the $T_i$.
For every integer $k\geq 0$, formal smoothness of the discrete $A$-algebra $C$ gives, as in \sref{0.19.5.4.2}, a surjective $A$-homomorphism
\[
  f:E/\mathfrak{n}^{k+1}\to B/\mathfrak{J}^{k+1}
\]
such that $\gr^j(f)=\varphi_j$ for every $j\leq k$.
If $\mathfrak{b}/\mathfrak{n}^{k+1}=\Ker(f)$, then $B/\mathfrak{J}^{k+1}$ identifies with $E/\mathfrak{b}$, and condition $d)$ factorizes the canonical homomorphism $B\to B/\mathfrak{J}^{k+1}$ as
\[
  B\xrightarrow{w}E/\mathfrak{n}^{k+1}
  \xrightarrow{f}B/\mathfrak{J}^{k+1}.
\]
Since $\gr^0(f)$ is the identity, $w(\mathfrak{J})\subset\mathfrak{n}/\mathfrak{n}^{k+1}$, hence $w(\mathfrak{J}^{k+1})=0$.
As in \sref{0.19.5.4.2}, it follows that $\varphi_k$ is injective for every $k$, which proves the assertion.
\end{remark}

\begin{env}[19.5.6]
\label{0.19.5.6}
\emph{Proof of the theorem} \sref{0.19.5.3}. ---
Let $(\mathfrak{b}_\alpha)$ be a fundamental decreasing system of open ideals of $B$.
Set
\[
  B_\alpha=B/\mathfrak{b}_\alpha,\qquad
  C_\alpha=B/(\mathfrak{b}_\alpha+\mathfrak{J}),\qquad
  \mathfrak{J}_\alpha=(\mathfrak{b}_\alpha+\mathfrak{J})/\mathfrak{b}_\alpha,
\]
so that
\[
  C_\alpha=B_\alpha/\mathfrak{J}_\alpha
  \qquad\text{and}\qquad
  \mathfrak{J}_\alpha^n/\mathfrak{J}_\alpha^{n+1}
  =(\mathfrak{b}_\alpha+\mathfrak{J}^n)/
   (\mathfrak{b}_\alpha+\mathfrak{J}^{n+1}).
\]
The $C$-modules
\[
  ((\mathfrak{b}_\alpha\cap\mathfrak{J}^n)+\mathfrak{J}^{n+1})/
  \mathfrak{J}^{n+1}
\]
of $\mathfrak{J}^n/\mathfrak{J}^{n+1}$ form a fundamental system of open submodules of the topological $C$-module $\mathfrak{J}^n/\mathfrak{J}^{n+1}$.
Since
\[
  (\mathfrak{b}_\alpha\cap\mathfrak{J}^n)+\mathfrak{J}^{n+1}
  =\mathfrak{J}^n\cap(\mathfrak{b}_\alpha+\mathfrak{J}^{n+1}),
\]
we have
\[
\begin{aligned}
  \mathfrak{J}^n/((\mathfrak{b}_\alpha\cap\mathfrak{J}^n)+\mathfrak{J}^{n+1})
  &=
  \mathfrak{J}^n/(\mathfrak{J}^n\cap(\mathfrak{b}_\alpha+\mathfrak{J}^{n+1}))\\
  &=
  (\mathfrak{b}_\alpha+\mathfrak{J}^n)/
  (\mathfrak{b}_\alpha+\mathfrak{J}^{n+1})
  =\mathfrak{J}_\alpha^n/\mathfrak{J}_\alpha^{n+1}.
\end{aligned}
\]
\end{env}

\begin{env}[19.5.6.1]
\label{0.19.5.6.1}
Let $P,Q$ be two discrete $C_\alpha$-modules and let
$u:P\to Q$ be a surjective $C_\alpha$-homomorphism.
The discrete ring $B/(\mathfrak{b}_\alpha+\mathfrak{J}^2)$ is a $B$-extension of $C_\alpha$ by the square-zero ideal
$\mathfrak{J}_\alpha/\mathfrak{J}_\alpha^2$.
Let $v:\mathfrak{J}/\mathfrak{J}^2\to Q$ be a continuous $C$-homomorphism.
After replacing $\mathfrak{b}_\alpha$, if necessary, by a smaller open ideal of $B$, we may suppose that the kernel of $v$ contains the open $C$-submodule
$((\mathfrak{b}_\alpha\cap\mathfrak{J})+\mathfrak{J}^2)/\mathfrak{J}^2$.
Passing to the quotient gives a $C_\alpha$-homomorphism of discrete modules
$v':\mathfrak{J}_\alpha/\mathfrak{J}_\alpha^2\to Q$.
Let $E_\alpha$ be the $B_\alpha$-extension of $C_\alpha$ by $Q$ deduced from
$B/(\mathfrak{b}_\alpha+\mathfrak{J}^2)$ by means of $v'$, and let
$v'_*:B/(\mathfrak{b}_\alpha+\mathfrak{J}^2)\to E_\alpha$
be the corresponding $B_\alpha$-homomorphism \sref{0.18.2.8}.
The algebra $E_\alpha$ is a discrete topological $A$-algebra, and the canonical isomorphism
$C_\alpha\xrightarrow{\sim}E_\alpha/Q$ gives by composition a continuous
$A$-homomorphism $f:C\to C_\alpha\to E_\alpha/Q$.
Since $Q$ is square zero in $E_\alpha$ and $C$ is a formally smooth $A$-algebra,
$f$ factorizes as
$C\xrightarrow{g}E_\alpha\to E_\alpha/Q$, where $g$ is a continuous
$A$-homomorphism.
As $g$ is continuous and $E_\alpha$ discrete, there exists $\beta\geq\alpha$
such that $g$ factorizes as
$C\to C_\beta\xrightarrow{g'}E_\alpha$.
Let $E_\beta$ be the $A$-extension of $C_\beta$ by $Q$ which is the inverse image
of $E_\alpha$ under the canonical homomorphism $C_\beta\to C_\alpha$.
The existence of $g'$ (with composite $C_\beta\to E_\alpha\to C_\alpha$ equal
to the canonical homomorphism) is equivalent to $E_\beta$ being an
$A$-\emph{trivial} extension, so it may be identified with
$\mathrm{D}_{C_\beta}(Q)$.
The surjective homomorphism $u:P\to Q$ then canonically defines a

\oldpage[0\textsubscript{IV-1}]{95}

surjective homomorphism
$\mathrm{D}_{C_\beta}(P)\to\mathrm{D}_{C_\beta}(Q)$ \sref{0.18.2.9}, whose
kernel is an ideal $\mathfrak{L}$ of the extension
$E'_\beta=\mathrm{D}_{C_\beta}(P)$, contained in $P$, and \emph{a fortiori}
square zero.
Let $f':B\to E_\beta=E'_\beta/\mathfrak{L}$ be the continuous $A$-homomorphism
defining the topological $B$-algebra structure of $E_\beta$.
Since $\mathfrak{L}$ is square zero and $B$ is a formally smooth $A$-algebra,
$f'$ factorizes as
$B\xrightarrow{h'}E'_\beta\to E'_\beta/\mathfrak{L}$, where $h'$ is a
continuous $A$-homomorphism.
The construction of $E'_\beta$ gives by composition an $A$-homomorphism
$\psi:E'_\beta\to E_\beta\to E_\alpha$, and the diagram
\[
  \xymatrix{
    E'_\beta \ar[r]^\psi & E_\alpha \\
    B \ar[u]^{h'} \ar[r]_\theta &
    B/(\mathfrak{b}_\alpha+\mathfrak{J}^2) \ar[u]_{v'_*}
  }
\]
is commutative.
Moreover, the image of $\mathfrak{J}$ under $\psi\circ h'$ is contained in $Q$,
so its image under $h'$ is contained in $P$, and the image of
$\mathfrak{J}^2$ under $h'$ is zero.
Restricting $h'$ and $\theta$ to $\mathfrak{J}$ gives a commutative diagram
\[
  \xymatrix{
    P \ar[r]^u & Q \\
    \mathfrak{J}/\mathfrak{J}^2 \ar[u]^w \ar[ur]_v
  }
\]
where $w$ is continuous.
Thus $\mathfrak{J}/\mathfrak{J}^2$ is a formally projective $C$-module.
\end{env}

\begin{env}[19.5.6.2]
\label{0.19.5.6.2}
For all $n\geq 0$, we set $E_n=B/\mathfrak{J}^{n+1}$, so that $E_0=C$;
and we set

\[
E_{\alpha,n}=B/(\mathfrak{b}_\alpha+\mathfrak{J}^{n+1})=B_\alpha/\mathfrak{J}_\alpha^{n+1},
\]

a discrete quotient ring.
Likewise, in $\gr_\mathfrak{J}^n(B)=\mathfrak{J}^n/\mathfrak{J}^{n+1}$, we have seen that the $((\mathfrak{b}_\alpha\cap\mathfrak{J}^n)+\mathfrak{J}^{n+1})/\mathfrak{J}^{n+1}$ form a fundamental system of open submodules.
Consider the symmetric algebra $\mathbf{S}_{C_\alpha}^*(\mathfrak{J}_\alpha/\mathfrak{J}_\alpha^2)$;
we denote by $F_{\alpha,n}$ the quotient of $\mathbf{S}_{C_\alpha}^*(\mathfrak{J}_\alpha/\mathfrak{J}_\alpha^2)$ by the $(n+1)$-th power of its ideal of augmentation.
For a fixed $n\geq 1$, it follows from \sref{0.19.5.1.1} that the $\mathbf{S}_{C_\alpha}^n(\mathfrak{J}_\alpha/\mathfrak{J}_\alpha^2)$ are the quotients of $\mathbf{S}_C^n(\mathfrak{J}/\mathfrak{J}^2)$ by a fundamental system of open submodules within this topological $C$-module.

To abbreviate the language, we will say that for $\alpha\leq\beta$, the canonical homomorphisms $B_\beta\to B_\alpha$, $C_\beta\to C_\alpha$, $\mathfrak{J}_\beta/\mathfrak{J}_\beta^2\to\mathfrak{J}_\alpha/\mathfrak{J}_\alpha^2$, $E_{\beta,n}\to E_{\alpha,n}$, $F_{\beta,n}\to F_{\alpha,n}$, etc., are \emph{homomorphisms of transition}.
\end{env}

\begin{lemma}[19.5.6.3]
\label{0.19.5.6.3}
Suppose that $C$ is a formally smooth $A$-algebra, that the ring $B$ is preadmissible, and that the $C$-module $\mathfrak{J}/\mathfrak{J}^2$ is formally projective.
Then:
\begin{enumerate}
  \item[\rm(i)] For every $\alpha$, there exists $\beta\geq\alpha$ and a surjective $A$-algebra homomorphism
  
\[
v_{\alpha\beta}:F_{\beta,n}\to E_{\alpha,n}
\]

  such that $\gr^0(v_{\alpha\beta}):C_\beta\to C_\alpha$ and $\gr^1(v_{\alpha\beta}):\mathfrak{J}_\beta/\mathfrak{J}_\beta^2\to\mathfrak{J}_\alpha/\mathfrak{J}_\alpha^2$ are the homomorphisms of transition;
  \item[\rm(ii)] If $\beta$ and $v_{\alpha\beta}$ satisfy the conditions of~{\rm(i)}, then for every $\gamma\geq\beta$, there exists $\delta\geq\gamma$ and a surjective $A$-algebra homomorphism $v_{\gamma\delta}:F_{\delta,n}\to E_{\gamma,n}$ satisfying the conditions of~{\rm(i)} (for $\gamma$ and $\delta$) and making commutative the diagram
  
\[
\xymatrix@C=3em{
      F_{\beta,n} \ar[r]^{v_{\alpha\beta}} & E_{\alpha,n} \\
      F_{\delta,n} \ar[u] \ar[r]_{v_{\gamma\delta}} & E_{\gamma,n} \ar[u]
    }
\]

  where the vertical arrows are the homomorphisms of transition.
\end{enumerate}
\end{lemma}

\begin{env}[19.5.6.4]
\label{0.19.5.6.4}
\begin{enumerate}
  \item[\rm(i)] In the discrete topological $A$-algebra $E_{\alpha,n}$, the ideal $\mathfrak{J}_\alpha/\mathfrak{J}_\alpha^{n+1}$ is nilpotent, and the isomorphism $C_\alpha\xrightarrow{\sim}E_{\alpha,n}/(\mathfrak{J}_\alpha/\mathfrak{J}_\alpha^{n+1})$ gives by composition a continuous $A$-homomorphism $C\to C_\alpha\to E_{\alpha,n}/(\mathfrak{J}_\alpha/\mathfrak{J}_\alpha^{n+1})$;

\oldpage[0\textsubscript{IV-1}]{96}

  as $C$ is a formally smooth $A$-algebra, this homomorphism factorizes as $C\xrightarrow{f_\alpha}E_{\alpha,n}\to E_{\alpha,n}/(\mathfrak{J}_\alpha/\mathfrak{J}_\alpha^{n+1})$,
  where $f_\alpha$ is a continuous $A$-homomorphism;
  hence $\mathfrak{J}_\alpha/\mathfrak{J}_\alpha^{n+1}$ becomes, by means of $f_\alpha$, a topological $C$-module annihilated by an open ideal of $C$.
  The hypothesis that $\mathfrak{J}/\mathfrak{J}^2$ is a formally projective $C$-module then implies the existence of a continuous $C$-linear map $g_\alpha:\mathfrak{J}/\mathfrak{J}^2\to\mathfrak{J}_\alpha/\mathfrak{J}_\alpha^{n+1}$ making commutative the diagram
  
\[
\xymatrix{
      \mathfrak{J}_\alpha/\mathfrak{J}_\alpha^{n+1} \ar[r] & \mathfrak{J}_\alpha/\mathfrak{J}_\alpha^2 \\
      & \mathfrak{J}/\mathfrak{J}^2 \ar[ul]^{g_\alpha} \ar[u]
    }
\]


  As $f_\alpha$ and $g_\alpha$ are continuous, there exists $\beta\geq\alpha$ such that these homomorphisms factorize through
  
\[
C\to C_\beta\xrightarrow{f'_{\alpha\beta}}E_{\alpha,n}
\]

  
\[
\mathfrak{J}/\mathfrak{J}^2\to\mathfrak{J}_\beta/\mathfrak{J}_\beta^2\xrightarrow{g'_{\alpha\beta}}\mathfrak{J}_\alpha/\mathfrak{J}_\alpha^{n+1}
\]

  and from $f'_{\alpha\beta}$ and $g'_{\alpha\beta}$, one obtains a homomorphism of $C_\beta$-algebras
  
\[
\mathbf{S}_{C_\beta}^*(\mathfrak{J}_\beta/\mathfrak{J}_\beta^2) \to E_{\alpha,n}
\]

  which, by the definition of $g'_{\alpha\beta}$, vanishes on the $(n+1)$-th power of the augmentation ideal of $\mathbf{S}_{C_\beta}^*(\mathfrak{J}_\beta/\mathfrak{J}_\beta^2)$; passing to the quotient by this power gives the desired homomorphism $v_{\alpha\beta}$, taking into account the definitions of $f_\alpha$ and $g_\alpha$;
  the surjectivity of $v_{\alpha\beta}$ follows from that of the two homomorphisms $\gr^0(v_{\alpha\beta})$ and $\gr^1(v_{\alpha\beta})$,
  since this implies that $\gr(v_{\alpha\beta})$ is surjective (the graded algebra $\gr^*(E_{\alpha,n})$ being generated by $\gr^0(E_{\alpha,n})$ and $\gr^1(E_{\alpha,n})$), and
  as the filtrations considered are finite, one can apply Bourbaki, \emph{Alg.\ comm.}, chap.~III, \S\,2, n$^\circ$\,8, cor.~3 of th.~1.

  \item[\rm(ii)] The hypothesis that $B$ is preadmissible means that one can suppose all the $\mathfrak{b}_\alpha$ contained in the same $\mathfrak{b}_{\alpha_0}$, whose powers tend to $0$.
  This implies in particular that the kernel of the transition homomorphism $E_{\gamma,n}\to E_{\alpha,n}$, equal to $(\mathfrak{b}_\alpha+\mathfrak{J}^{n+1})/(\mathfrak{b}_\gamma+\mathfrak{J}^{n+1})$, is \emph{nilpotent};
  applying the lemma \sref{0.19.3.10.1}, one sees that one can choose $f_\gamma$ so that the diagram
  
\[
\xymatrix{
      & E_{\alpha,n} \\
      C \ar[ur]^{f_\alpha} \ar[r]_{f_\gamma} & E_{\gamma,n} \ar[u]
    }
\]

  (where the vertical arrow is the transition homomorphism) is commutative.
  The hypothesis that $\mathfrak{J}/\mathfrak{J}^2$ is a formally projective $C$-module furthermore allows one to choose $g_\gamma:\mathfrak{J}/\mathfrak{J}^2\to\mathfrak{J}_\gamma/\mathfrak{J}_\gamma^{n+1}$ so that the diagram
  \[
    \xymatrix@C=5em@R=3em{
      \mathfrak{J}_\alpha/\mathfrak{J}_\alpha^{n+1} \ar[r] &
      \mathfrak{J}_\alpha/\mathfrak{J}_\alpha^2 \\
      \mathfrak{J}_\gamma/\mathfrak{J}_\gamma^{n+1} \ar[u] \ar[r] &
      \mathfrak{J}_\gamma/\mathfrak{J}_\gamma^2 \ar[u] \\
      & \mathfrak{J}/\mathfrak{J}^2
        \ar[ul]_{g_\gamma} \ar[u]_{p_\gamma} \ar[uul]^{g_\alpha}
    }
  \]
  is commutative.
  Indeed, the image of $\mathfrak{J}/\mathfrak{J}^2$ under $(g_\alpha,p_\gamma)$ in
  $(\mathfrak{J}_\alpha/\mathfrak{J}_\alpha^{n+1})
  \times(\mathfrak{J}_\gamma/\mathfrak{J}_\gamma^2)$ is contained, by the definition of $g_\alpha$ and the relation $\beta\leq\gamma$, in the canonical image $Q$ of
  $P=\mathfrak{J}_\gamma/\mathfrak{J}_\gamma^{n+1}$ in this product module.
  Apply the definition \sref{0.19.2.1} to the surjective homomorphism $P\to Q$ and to $(g_\alpha,p_\gamma):\mathfrak{J}/\mathfrak{J}^2\to Q$.
  This choice of $f_\gamma$ and $g_\gamma$ allows $v_{\gamma\delta}$ to be constructed as in~{\rm(i)} while also making the diagram \sref{0.19.5.6.4} commutative.
\end{enumerate}
\end{env}

\oldpage[0\textsubscript{IV-1}]{97}

\begin{lemma}[19.5.6.5]
\label{0.19.5.6.5}
Suppose that the $A$-algebras $B$ and $C$ are formally smooth and that $B$ is preadmissible (so that, by \sref{0.19.5.6.1}, the conditions of \sref{0.19.5.6.3} are satisfied).
For every pair of indices $\alpha\leq\beta$ and every homomorphism
$v_{\alpha\beta}:F_{\beta,n}\to E_{\alpha,n}$ satisfying the conditions of
\sref{0.19.5.6.3},~{\rm(i)}, there exists an index $\lambda\geq\beta$ and an
$A$-algebra homomorphism
\[
  w_{\beta\lambda}:E_{\lambda,n}\to F_{\beta,n}
\]
such that:
$1^\circ$ $\gr^0(w_{\beta\lambda}):C_\lambda\to C_\beta$ and
$\gr^1(w_{\beta\lambda}):\mathfrak{J}_\lambda/\mathfrak{J}_\lambda^2
\to\mathfrak{J}_\beta/\mathfrak{J}_\beta^2$ are the transition homomorphisms;
$2^\circ$ the composite
$E_{\lambda,n}\xrightarrow{w_{\beta\lambda}}F_{\beta,n}
\xrightarrow{v_{\alpha\beta}}E_{\alpha,n}$ is the transition homomorphism.
\end{lemma}

\begin{proof}
Apply \sref{0.19.5.6.3},~{\rm(ii)}, with $\gamma=\beta$.
This gives $\delta\geq\beta$ and
$v_{\beta\delta}:F_{\delta,n}\to E_{\beta,n}$.
Recall that $v_{\beta\delta}$ is surjective.
Its kernel $\mathfrak{R}$ is nilpotent: the augmentation ideal
$F_{\delta,n}^+$ is nilpotent by definition, while the kernel of
$\gr^0(v_{\beta\delta}):C_\delta\to C_\beta$ is nilpotent because $B$ is
preadmissible.
Thus $E_{\beta,n}$ identifies with $F_{\delta,n}/\mathfrak{R}$.
Since $B$ is a formally smooth $A$-algebra, the continuous canonical
$A$-homomorphism $p_{\beta,n}:B\to E_{\beta,n}$ factorizes as
\[
  B\xrightarrow{w}F_{\delta,n}
  \xrightarrow{v_{\beta\delta}}E_{\beta,n},
\]
where $w$ is continuous.
As $F_{\delta,n}$ is discrete, there exists $\lambda\geq\delta$ such that $w$
vanishes on $\mathfrak{b}_\lambda$, and hence $w$ factorizes as
\[
  B\to B/\mathfrak{b}_\lambda
  \xrightarrow{w'_\lambda}F_{\delta,n}.
\]
Consider the composite
\[
  B/\mathfrak{b}_\lambda\xrightarrow{w'_\lambda}F_{\delta,n}
  \xrightarrow{q_{\beta\delta}}F_{\beta,n},
\]
where $q_{\beta\delta}$ is the transition homomorphism.
We have
$\gr^0(q_{\beta\delta})=\gr^0(v_{\beta\delta}):C_\delta\to C_\beta$,
the transition homomorphism.
Since
$v_{\alpha\beta}\circ q_{\beta\delta}\circ w'_\lambda$
is the canonical homomorphism
$B/\mathfrak{b}_\lambda\to
B/(\mathfrak{b}_\alpha+\mathfrak{J}^{n+1})$,
the image of
$\mathfrak{J}_\lambda=(\mathfrak{b}_\lambda+\mathfrak{J})/
\mathfrak{b}_\lambda$
under $q_{\beta\delta}\circ w'_\lambda$ is contained in the augmentation ideal
$F_{\beta,n}^+$.
Consequently the image of
$\mathfrak{J}_\lambda^{n+1}
=(\mathfrak{b}_\lambda+\mathfrak{J}^{n+1})/\mathfrak{b}_\lambda$
is zero.
Thus $q_{\beta\delta}\circ w'_\lambda$ factorizes as
\[
  B/\mathfrak{b}_\lambda\to
  B/(\mathfrak{b}_\lambda+\mathfrak{J}^{n+1})=E_{\lambda,n}
  \xrightarrow{w_{\beta\lambda}}F_{\beta,n},
\]
and
$E_{\lambda,n}\xrightarrow{w_{\beta\lambda}}F_{\beta,n}
\xrightarrow{v_{\alpha\beta}}E_{\alpha,n}$
is the transition homomorphism.
The preceding argument also shows that $\gr^0(w_{\beta\lambda})$ is the
transition homomorphism $C_\lambda\to C_\beta$, because it is the composite
\[
  B/(\mathfrak{b}_\lambda+\mathfrak{J})
  \to\gr^0(F_{\delta,n})
  \xrightarrow{\gr^0(q_{\beta\delta})}\gr^0(F_{\beta,n})
\]
and $v_{\beta\delta}\circ w$ is canonical.
Moreover,
$\gr^1(q_{\beta\delta})=\gr^1(v_{\beta\delta})$ by
\sref{0.19.5.6.3}, and the same argument shows that
$\gr^1(w_{\beta\lambda})$ is the transition homomorphism.
\end{proof}

\begin{env}[19.5.6.6]
\label{0.19.5.6.6}
\emph{Proof of} \sref{0.19.5.3},~(ii). ---
To show that $\varphi_n$ is a formal bimorphism, we apply the criterion of \sref{0.19.1.3}; the conditions of \sref{0.19.5.6.5} being satisfied by hypothesis, choose, for every $\alpha$, homomorphisms $v_{\alpha\beta}$ and $w_{\beta\lambda}$ satisfying the conclusions of that lemma.
The homomorphism
\[
  \gr^n(v_{\alpha\beta}):\gr^n(F_{\beta,n})\to\gr^n(E_{\alpha,n})
\]
is the homomorphism
\[
  \varphi_{\alpha\beta,n}:\mathbf{S}_{C_\beta}^n(\mathfrak{J}_\beta/\mathfrak{J}_\beta^2)\to\mathfrak{J}_\alpha^n/\mathfrak{J}_\alpha^{n+1}
\]
deduced from \sref{0.19.5.2.1} by passage to quotients.
Indeed, \sref{0.19.5.6.3} and the definition of $\varphi$ show inductively that $\gr^j(v_{\alpha\beta})=\varphi_{\alpha\beta,j}$ for all $j\leq n$.
Since $\varphi_n$ is surjective, it is \emph{a fortiori} a formal epimorphism.
Moreover, the composite
\[
  \gr^j(F_{\lambda,n})\xrightarrow{\varphi_{\lambda\lambda,j}}\gr^j(E_{\lambda,n})\xrightarrow{\gr^j(w_{\beta\lambda})}\gr^j(F_{\beta,n})
\]
is the transition homomorphism: this holds for $j=0,1$ by \sref{0.19.5.6.5}, and then for every $j\leq n$ because $\gr^*(E_{\lambda,n})$ is generated by $\gr^1(E_{\lambda,n})$.
After composing with the transition map to $\gr^n(F_{\alpha,n})$, the transition homomorphism $\gr^n(F_{\lambda,n})\to\gr^n(F_{\alpha,n})$ therefore factors as
\[
  \gr^n(F_{\lambda,n})\xrightarrow{\varphi_{\lambda\lambda,n}}\gr^n(E_{\lambda,n})\to\gr^n(F_{\alpha,n}),
\]
which is precisely the criterion of \sref{0.19.1.3}.
\end{env}

\begin{lemma}[19.5.6.7]
\label{0.19.5.6.7}
\oldpage[0\textsubscript{IV-1}]{98}
Suppose that $B$ is preadmissible, that $\mathfrak{J}/\mathfrak{J}^2$ is a formally projective $C$-module, that $C$ is a formally smooth $A$-algebra, and that the $\varphi_n$ are formal bimorphisms.
Then, for every pair $\alpha\leq\beta$ and every homomorphism $v_{\alpha\beta}:F_{\beta,n}\to E_{\alpha,n}$ satisfying \sref{0.19.5.6.3},~{\rm(i)}, there exists $\gamma\geq\beta$ such that, for every $\mu\geq\gamma$, there is a commutative diagram
\[
\xymatrix@C=4em@R=3em{
  F_{\beta,n} \ar[r]^{v_{\alpha\beta}} & E_{\alpha,n} \\
  & E_{\mu,n} \ar[ul]_{u_{\beta\mu}} \ar[u]_{p_{\alpha\mu}}
}
\]
of $A$-homomorphisms, where $p_{\alpha\mu}$ is the transition homomorphism.
\end{lemma}

\begin{proof}
Apply the criterion of \sref{0.19.1.3} to each $\varphi_j$, $0\leq j\leq n$.
There exist $\gamma\geq\beta$ and uniquely determined surjective $C_\gamma$-linear homomorphisms
\[
  w_{\beta\gamma}^j:\gr^j(E_{\gamma,n})\to\gr^j(F_{\beta,n})
\]
such that
\[
  \gr^j(F_{\gamma,n})\xrightarrow{\varphi_{\gamma\gamma,j}}\gr^j(E_{\gamma,n})\xrightarrow{w_{\beta\gamma}^j}\gr^j(F_{\beta,n})
\]
is the transition homomorphism.
The same $\gamma$ may be chosen for all $j$, since there are only finitely many of them.
By uniqueness and because $\varphi$ is a homomorphism of graded algebras, the family
$w_{\beta\gamma}=(w_{\beta\gamma}^j)_{0\leq j\leq n}$ is a homomorphism of graded $C_\gamma$-algebras
\[
  \gr^\bullet(E_{\gamma,n})\to\gr^\bullet(F_{\beta,n})=F_{\beta,n}.
\]
Since $\varphi_{\gamma\gamma,0}$ and $\varphi_{\gamma\gamma,1}$ are the identity homomorphisms,
$w_{\beta\gamma}^0:C_\gamma\to C_\beta$ and
$w_{\beta\gamma}^1:\mathfrak{J}_\gamma/\mathfrak{J}_\gamma^2
\to\mathfrak{J}_\beta/\mathfrak{J}_\beta^2$ are the transition homomorphisms.
The same therefore holds for
$\gr^0(v_{\alpha\beta}\circ w_{\beta\gamma})$ and
$\gr^1(v_{\alpha\beta}\circ w_{\beta\gamma})$; since
$\gr^\bullet(E_{\gamma,n})$ is generated by $\gr^1(E_{\gamma,n})$, it follows that
$\gr^j(v_{\alpha\beta}\circ w_{\beta\gamma})$ is also the transition homomorphism for $0\leq j\leq n$.

Apply now \sref{0.19.5.6.3},~(ii), to $\alpha$, $\beta$, and $\gamma$;
this gives diagram \sref{0.19.5.6.4} with $\delta\geq\gamma$.
Repeat the argument at the beginning after replacing $\alpha$ and $\beta$ by $\gamma$ and $\delta$.
One obtains $\lambda\geq\delta$ and a commutative diagram
\[
\xymatrix@C=4em@R=4em{
  \gr^\bullet(E_{\gamma,n})\ar[r]^{w_{\beta\gamma}} &
  F_{\beta,n}\ar[r]^{v_{\alpha\beta}} & E_{\alpha,n} \\
  \gr^\bullet(E_{\lambda,n})\ar[u]_{\gr^\bullet(p_{\gamma\lambda})}
    \ar[r]_{w_{\delta\lambda}} &
  F_{\delta,n}\ar[u]_{q_{\beta\delta}}\ar[r]_{v_{\gamma\delta}} &
  E_{\gamma,n}\ar[u]_{p_{\alpha\gamma}}
    \ar@{-->}[ul]_{u_{\beta\gamma}}
}
\]
whose vertical arrows are the transition homomorphisms.
It remains to prove the existence of $u_{\beta\gamma}$ making the diagram commute, and for this it plainly suffices to show
\[
  \Ker(v_{\gamma\delta})\subset\Ker(q_{\beta\delta}),
\]
since $v_{\gamma\delta}$ and $q_{\beta\delta}$ are surjective.
As $w_{\delta\lambda}$ is surjective, this relation is equivalent to
\[
\begin{split}
  \Ker(v_{\gamma\delta}\circ w_{\delta\lambda})
  &\subset\Ker(q_{\beta\delta}\circ w_{\delta\lambda})\\
  &=\Ker\bigl(w_{\beta\gamma}\circ\gr(p_{\gamma\lambda})\bigr).
\end{split}
\]
But we saw above that
$\gr(p_{\gamma\lambda})=\gr(v_{\gamma\delta}\circ w_{\delta\lambda})$, and clearly
\[
\begin{split}
  \Ker(v_{\gamma\delta}\circ w_{\delta\lambda})
  &\subset\Ker\bigl(\gr(v_{\gamma\delta}\circ w_{\delta\lambda})\bigr)\\
  &=\Ker(\gr(p_{\gamma\lambda}))
  \subset\Ker\bigl(w_{\beta\gamma}\circ\gr(p_{\gamma\lambda})\bigr).
\end{split}
\]
This proves the lemma: for every $\mu\geq\gamma$, define $u_{\beta\mu}$ as the composite
\[
  E_{\mu,n}\xrightarrow{p_{\gamma\mu}}E_{\gamma,n}
  \xrightarrow{u_{\beta\gamma}}F_{\beta,n}.
\]
\end{proof}

\begin{env}[19.5.6.8]
\label{0.19.5.6.8}
\emph{Proof of} \sref{0.19.5.3},~(iii). ---
Let $G$ be a discrete topological $A$-algebra, $\mathfrak{R}$ an ideal of square zero in $G$, and $f:B\to G/\mathfrak{R}$ a continuous $A$-homomorphism of algebras.
As $G/\mathfrak{R}$ is discrete, there exists $\alpha$ such that $f$ annihilates $\mathfrak{b}_\alpha$; moreover, as $\mathfrak{J}^m\subset\mathfrak{b}_\alpha$ for a sufficiently large $m$, $f$ factors as
\[
  B\xrightarrow{p_\alpha}E_{\alpha,n}\xrightarrow{f_{\alpha,n}}G/\mathfrak{R},
\]
where $n=2m$.
By \sref{0.19.5.6.3}, there exist $\beta\geq\alpha$ and $v_{\alpha\beta}:F_{\beta,n}\to E_{\alpha,n}$; thus
\begin{equation}
\label{0.19.5.6.9}
  F_{\beta,n}\xrightarrow{v_{\alpha\beta}}E_{\alpha,n}\xrightarrow{f_{\alpha,n}}G/\mathfrak{R}
  \tag{19.5.6.9}
\end{equation}
is a continuous $A$-homomorphism.
Since $F_{\beta,n}$ is a $C_\beta$-algebra, hence \emph{a fortiori} a discrete topological $C$-algebra, \sref{0.19.5.6.9} gives by composition a continuous $A$-homomorphism
$r:C\to F_{\beta,n}\to G/\mathfrak{R}$.
Since, on the other hand, $C$ is a formally smooth $A$-algebra,

\oldpage[0\textsubscript{IV-1}]{99}

this homomorphism factorizes as
$r:C\xrightarrow{r'}G\to G/\mathfrak{R}$, where $r'$ is a continuous $A$-homomorphism, thereby giving $G$ the structure of a discrete topological $C$-algebra.
Composing \sref{0.19.5.6.9} with
\[
  \mathfrak{J}/\mathfrak{J}^2\to\mathfrak{J}_\beta/\mathfrak{J}_\beta^2\to F_{\beta,n}
\]
gives a continuous $C$-homomorphism $g:\mathfrak{J}/\mathfrak{J}^2\to G/\mathfrak{R}$.
Since $\mathfrak{J}/\mathfrak{J}^2$ is formally projective, $g$ lifts to a continuous $C$-homomorphism $h:\mathfrak{J}/\mathfrak{J}^2\to G$.
As $G$ is discrete, there exists $\gamma\geq\beta$ such that $h$ annihilates the image of $((\mathfrak{b}_\gamma\cap\mathfrak{J})+\mathfrak{J}^2)/\mathfrak{J}^2$ and $r'$ annihilates the image of $(\mathfrak{b}_\gamma+\mathfrak{J})/\mathfrak{J}$.
Thus $h$ factors through a $C_\gamma$-homomorphism $\mathfrak{J}_\gamma/\mathfrak{J}_\gamma^2\to G$, and hence extends to
\[
  w_\gamma:\mathbf{S}_{C_\gamma}^*(\mathfrak{J}_\gamma/\mathfrak{J}_\gamma^2)\to G.
\]
Every element of degree $m$ has image in $\mathfrak{R}$, so every element of degree $n=2m$ has image zero; consequently $w_\gamma$ factors as $F_{\gamma,n}\xrightarrow{w'_\gamma}G$.

Applying \sref{0.19.5.6.7} to
$v_{\alpha\gamma}=v_{\alpha\beta}\circ q_{\beta\gamma}:F_{\gamma,n}\to E_{\alpha,n}$
gives $\delta\geq\gamma$ and $u_{\gamma\delta}:E_{\delta,n}\to F_{\gamma,n}$ such that
$v_{\alpha\gamma}\circ u_{\gamma\delta}=p_{\alpha\delta}$.
One obtains the commutative diagram
\[
\xymatrix@C=3em@R=2.5em{
  && E_{\alpha,n} \ar[r]^{f_{\alpha,n}} & G/\mathfrak{R} \\
  && F_{\beta,n} \ar[u]_{v_{\alpha\beta}} & \\
  B \ar[r]_{p_\delta} & E_{\delta,n} \ar[uur]^{p_{\alpha\delta}} \ar[r]_{u_{\gamma\delta}} & F_{\gamma,n} \ar[u]_{q_{\beta\gamma}} \ar[r]^{w'_\gamma} & G \ar[uu]
}
\]
of continuous $A$-homomorphisms.
The composite $f':B\to G$ along the bottom row satisfies $f=(G\to G/\mathfrak{R})\circ f'$, proving that $B$ is a formally smooth $A$-algebra.
The proof of \sref{0.19.5.3} is complete.
\end{env}

\begin{corollary}[19.5.7]
\label{0.19.5.7}
Let $A$ be a topological ring, $B$ a topological $A$-algebra, $(\mathfrak{b}_\lambda)$ a fundamental system of open ideals of $B$, $\mathfrak{J}$ an ideal of $B$, and $C=B/\mathfrak{J}$ the topological quotient $A$-algebra.
Set $C_\lambda=B/(\mathfrak{b}_\lambda+\mathfrak{J})$.
Suppose that:
\begin{enumerate}
  \item[$1^\circ$] for all $n$, the topology induced on $\mathfrak{J}^n$ by that of $B$ is also the topology of the $B$-module $\mathfrak{J}^n$ deduced from the topology of $B$ \sref{0.19.0.2}; this condition is satisfied in particular if $B$ is Noetherian and its topology is preadic (\textbf{0\textsubscript{I}}, 7.3.2);
  \item[$2^\circ$] $C$ is a formally smooth $A$-algebra.
\end{enumerate}
Under these conditions:
\begin{enumerate}
  \item[\rm(i)] If $B$ is a formally smooth $A$-algebra, then, for all $\lambda$, $(\mathfrak{J}/\mathfrak{J}^2)\otimes_C C_\lambda$ is a projective $C_\lambda$-module.
  \item[\rm(ii)] If $B$ is a formally smooth $A$-algebra and a preadmissible ring, then, for all $\lambda$, the canonical homomorphism
  \[
    \label{0.19.5.7.1}
    \varphi_\lambda=\varphi\otimes 1_{C_\lambda}:
    \mathbf{S}_{C_\lambda}^*((\mathfrak{J}/\mathfrak{J}^2)\otimes_C C_\lambda)
    \to \gr_\mathfrak{J}(B)\otimes_C C_\lambda
    \tag{19.5.7.1}
  \]
  is bijective.
  \oldpage[0\textsubscript{IV-1}]{100}
  \item[\rm(iii)] Conversely, suppose that $B$ is preadmissible, that the sequence $(\mathfrak{J}^n)$ tends to $0$, and that, for all $\lambda$, $(\mathfrak{J}/\mathfrak{J}^2)\otimes_C C_\lambda$ is a projective $C_\lambda$-module and the homomorphism \sref{0.19.5.7.1} is bijective.
  Then $B$ is a formally smooth $A$-algebra.
\end{enumerate}
\end{corollary}

\begin{proof}
Indeed, the topology of $\mathfrak{J}/\mathfrak{J}^2$ and those of the $\mathbf{S}_C^n(\mathfrak{J}/\mathfrak{J}^2)$ are then also deduced from that of $B$ \sref{0.19.5.1};
the conclusions are immediate consequences of \sref{0.19.5.3} and of the criteria of \sref{0.19.1.4} and \sref{0.19.2.4} specialized to this type of topology.
\end{proof}

\begin{remark}[19.5.8]
\label{0.19.5.8}
Suppose that $\mathfrak{J}/\mathfrak{J}^2$ is a $C$-module of finite type and that, for the quotient topology inherited from $B$, $C$ is a Zariski ring.
Let $\mathfrak{r}$ be an ideal of definition of $C$, so that the $\mathfrak{r}^n$ form a fundamental system of neighbourhoods of $0$ in $C$.
Then the conclusions of {\rm(i)} and {\rm(ii)} in \sref{0.19.5.7} may be replaced by the following:
\begin{enumerate}
  \item[\rm(i')] $\mathfrak{J}/\mathfrak{J}^2$ is a projective $C$-module.
  \item[\rm(ii')] The canonical homomorphism
  \[
    \varphi:\mathbf{S}_C^*(\mathfrak{J}/\mathfrak{J}^2)\to\gr_\mathfrak{J}(B)
  \]
  is bijective.
\end{enumerate}
Indeed, it is clear that {\rm(i')} implies the conclusion of {\rm(i)} in \sref{0.19.5.7}.
Conversely, if $(\mathfrak{J}/\mathfrak{J}^2)\otimes_C C_\lambda$ is a projective $C_\lambda$-module for all $\lambda$, then $(\mathfrak{J}/\mathfrak{J}^2)\otimes_C(C/\mathfrak{r}^n)$ is a projective, hence flat, $(C/\mathfrak{r}^n)$-module for every $n$;
one concludes that $\mathfrak{J}/\mathfrak{J}^2$ is a flat $C$-module (\textbf{0\textsubscript{III}}, 10.2.2), hence projective since it is finitely presented (Bourbaki, \emph{Alg. comm.}, chap.~II, \textsection5, n\textsuperscript{o}~2, cor.~2 of th.~1).
On the other hand, the $C$-modules $\mathbf{S}_C^n(\mathfrak{J}/\mathfrak{J}^2)$ and $\gr_\mathfrak{J}^n(B)$ are of finite type, and when $C$ is a Zariski ring it is equivalent to saying that $\varphi_n$ is bijective or that $\varphi_n\otimes 1_{C_\lambda}$ is bijective for all $\lambda$ (Bourbaki, \emph{Alg. comm.}, chap.~III, \textsection3, n\textsuperscript{o}~5, prop.~9);
therefore {\rm(ii)} is equivalent to {\rm(ii')}.
\end{remark}

\subsection{The case of algebras over a field}
\label{subsection:0.19.6}

\begin{theorem}[19.6.1 (Cohen)]
\label{0.19.6.1}
Let $k$ be a field, $K$ an extension of $k$, $k$ and $K$ being equipped with the discrete topologies.
For $K$ to be a formally smooth $k$-algebra, it is necessary and sufficient that $K$ be a separable extension of $k$.
\end{theorem}

\begin{proof}
The necessity of the condition will be established in \sref{0.19.6.5.1} (and of course is not used until then);
we will content ourselves here with proving that the condition is \emph{sufficient}.
We distinguish two cases:

\emph{I. --- $K$ is a separable extension of finite type of $k$.}
We know then (Bourbaki, \emph{Alg.}, chap.~V, \S\,9, n$^\circ$\,3, th.~2) that there exists a pure sub-extension $K'=k(T_1,\ldots,T_n)$ of $K$ such that $K$ is a finite separable algebraic extension of $K'$.
  In view of \sref{0.19.3.5},~(iii), one can therefore reduce to the case where $K=K'$, or to the case where $K$ is algebraically finite over $k$.
In the first case, $A=k[T_1,\ldots,T_n]$ is a formally smooth $k$-algebra \sref{0.19.3.3}, hence $k(T_1,\ldots,T_n)=S^{-1}A$ where $S=A-\{0\}$ is also formally smooth \sref{0.19.3.5},~(iv).
In the second case, \emph{all} Hochschild cohomology groups $\mathrm{H}^n_k(K,L)$ are zero for every $(K,K)$-bimodule $L$:
indeed, if one considers the $k$-algebra tensor product $C=K\otimes_k K$, $L$ is a $C$-module on the left and the cohomology group $\mathrm{H}^n_k(K,L)$ is equal to $\operatorname{Ext}^n_C(K,L)$, where $K$ is also viewed as $(K,K)$-bimodule,
\oldpage[0\textsubscript{IV-1}]{101}
hence as a left $C$-module (M, IX, 4).
Now, as $K$ is a finite separable extension of $k$, we know that $K\otimes_k K$ is a direct product of extensions of $k$, one of which is $K$ itself (Bourbaki, \emph{Alg.}, chap.~VIII, \S\,8, prop.~3);
this implies that $K$ is a \emph{projective} $C$-module, whence our assertion.
Every $k$-extension of $K$ of zero kernel $L$ is therefore $k$-trivial, and \emph{a fortiori} all commutative $k$-extensions are trivial, whence the theorem in this case \sref{0.19.4.4}.

\emph{II. --- General case.} ---
With the notation of \sref{0.18.4.5}, it suffices to prove that $\Hom_K(\mathrm{H}_2(\mathrm{P}'_\bullet),L)=0$ for \emph{every} $K$-vector space $L$, which means that $\mathrm{H}_2(\mathrm{P}'_\bullet)=0$.
If $K$ is a union of a filtering family of sub-extensions $K^{(\alpha)}$ of $k$, $\mathrm{P}'_\bullet$ is the \emph{inductive limit} of the corresponding complexes $\mathrm{P}'^{(\alpha)}_\bullet$, since the $\varinjlim$ functor commutes with the tensor product in the category of $k$-modules;
by the exactness of the $\varinjlim$ functor in this category, one has $\mathrm{H}_2(\mathrm{P}'_\bullet)=\varinjlim\mathrm{H}_2(\mathrm{P}'^{(\alpha)}_\bullet)$.
As $K$ is separable, it follows that the same holds for every sub-extension of $K$ and that $K$ is a union of sub-extensions of finite type, so the first part of the proof implies that $K$ is a formally smooth $k$-algebra, Q.E.D.
\end{proof}

\begin{corollary}[19.6.2]
\label{0.19.6.2}
Let $A$ be a separated complete local ring, $K$ its residue field, $k$ a subfield of $K$ such that $K$ is a separable extension of $k$.
Then there exists a subfield $K'$ of $A$ containing $k$, such that the restriction to $K'$ of the canonical homomorphism $A\to K$ is an isomorphism of $K'$ onto $K$.
\end{corollary}

\begin{proof}
Let $\mathfrak{m}$ be the maximal ideal of $A$.
By the hypothesis and \sref{0.19.6.1}, $K$ is a formally smooth $k$-algebra;
applying \sref{0.19.3.10} (replacing $C$ by $A$ and $\mathfrak{J}$ by $\mathfrak{m}$), one sees that the identity automorphism of $K=A/\mathfrak{m}$ factorizes as $K\xrightarrow{u}A\to A/\mathfrak{m}$, where $u$ is a $k$-homomorphism, hence necessarily a $k$-isomorphism of $K$ onto a subfield $K'$ of $A$ containing $k$.
\end{proof}

\begin{corollary}[19.6.3]
\label{0.19.6.3}
Let $A$ be a Noetherian complete local ring, $\mathfrak{m}$ its maximal ideal, $k$ its residue field.
The following conditions are equivalent:
\begin{enumerate}
  \item[a)] $A$ contains a subfield.
  \item[b)] If $p$ is the characteristic of $k$, one has $pA=0$.
  \item[c)] There exists a field $K$ such that $A$ is isomorphic to a quotient ring of a formal power series ring $B=K[[T_1,\ldots,T_n]]$.
\end{enumerate}
When this is the case, one can suppose that $A$ is isomorphic to $B/\mathfrak{b}$, where $\mathfrak{b}$ is contained in the square of the maximal ideal of $B$.
\end{corollary}

\begin{proof}
It is immediate that $c)$ implies $a)$, since $A$ is $\neq 0$ because it is a local ring;
as $K$ is a field and the composite homomorphism $K\xrightarrow{\varphi}B\to A$ (where $\varphi$ is the canonical injection) is not zero, it is injective.
To see that $a)$ implies $b)$, it suffices to remark that if $K$ is a subfield of $A$, $K$ is isomorphic to a subfield of $k$ and has therefore the same characteristic $p$;
as $p.1=0$ in $K$, hence in $A$, one has $pA=0$.
Conversely, $b)$ implies $a)$, for the composite of the canonical homomorphisms $\mathbf{Z}\xrightarrow{j}A\to k$ has kernel $p\mathbf{Z}$, and one has $j(p\mathbf{Z})=0$, hence $\Ker(j)=p\mathbf{Z}$ and $A$ contains a ring $A'$ isomorphic to $\mathbf{Z}/p\mathbf{Z}$ and not meeting $\mathfrak{m}$;
if $p>0$, $A'$ is a field; otherwise,

\oldpage[0\textsubscript{IV-1}]{102}

every nonzero element of $A'$ being invertible in $A$, the field of fractions of $A'$ (isomorphic to $\mathbf{Q}$) is also contained in $A$.

Let us prove finally that $a)$ implies $c)$;
the condition $a)$ implies that $A$ contains a prime subfield $k_0$, isomorphic to the prime subfield of $k$, and whose $k$ is therefore a separable extension.
Applying \sref{0.19.6.2}, one sees that there exists an isomorphism $f$ of $k$ onto a subfield $K$ of $A$.
On the other hand, let $(x_i)_{1\leq i\leq n}$ be a system of elements of $\mathfrak{m}$ such that the classes mod.\ $\mathfrak{m}^2$ of the $x_i$ form a basis (over $k$) of $\mathfrak{m}/\mathfrak{m}^2$.
Since $A$ is complete, there exists a continuous ring homomorphism $u:k[[T_1,\ldots,T_n]]\to A$ such that $u$ equals $f$ on $k$ and $u(T_i)=x_i$ for all $i$, and this homomorphism is surjective by virtue of the choice of the $x_i$ (Bourbaki, \emph{Alg.\ comm.}, chap.~III, \S\,2, n$^\circ$\,9, prop.~11).
\end{proof}

\begin{theorem}[19.6.4]
\label{0.19.6.4}
Let $k$ be a field, $A$ a local Noetherian $k$-algebra, $\mathfrak{m}$ its maximal ideal, $K=A/\mathfrak{m}$ its residue field, $A$ being equipped with the $\mathfrak{m}$-preadic topology.
Suppose that $K$ is a separable extension of $k$.
Then the following conditions are equivalent:
\begin{enumerate}
  \item[a)] $A$ is a formally smooth $k$-algebra.
  \item[b)] The completion $\widehat{A}$ of $A$ is $k$-isomorphic to a formal power series ring $K[[T_1,\ldots,T_n]]$ (where the $k$-algebra structure is defined by the composite homomorphism $k\xrightarrow{\varphi}K\xrightarrow{\psi}K[[T_1,\ldots,T_n]]$, where $\psi$ is the canonical injection and $\varphi$ the homomorphism defining the extension structure of $K$).
  \item[b')] $\widehat{A}$ is a formal power series ring isomorphic to $K[[T_1,\ldots,T_n]]$.
  \item[c)] $A$ is a regular local ring.
\end{enumerate}
\end{theorem}

\begin{proof}
The fact that $a)$ implies $b)$ is the particular case of \sref{0.19.5.4} (equivalence of $a)$ and $c)$) obtained by replacing $A$ by $k$, $B$ by $A$, and $\mathfrak{J}$ by $\mathfrak{m}$, taking into account \sref{0.19.6.1}.
It is clear that $b)$ implies $b'$) and $b')$ implies $c)$ by \sref{0.17.1.1}.
Finally, if $c)$ holds, it follows from \sref{0.17.1.1},~$a$), and from \sref{0.19.5.4} (equivalence of $b)$ and $a)$), applied with $A=k$, $B=A$, $\mathfrak{J}=\mathfrak{m}$, and $C=K=A/\mathfrak{m}$, that $A$ is a formally smooth $k$-algebra.
\end{proof}

\begin{corollary}[19.6.5]
\label{0.19.6.5}
Let $k$ be a field, $A$ a local Noetherian $k$-algebra, $\mathfrak{m}$ its maximal ideal, $A$ being equipped with the $\mathfrak{m}$-preadic topology.
Suppose that $A$ is a formally smooth $k$-algebra;
then $A$ is \emph{geometrically regular} over $k$, in other words (cf.~$\mathbf{IV}$, 6.7.6), for every finite extension $k'$ of $k$, the semi-local ring $A'=A\otimes_k k'$ is regular \sref{0.17.3.6}.
Furthermore, if $K$ is the residue field of $A$, $\widehat{A}$ is isomorphic to a formal power series ring $K[[T_1,\ldots,T_n]]$.
\end{corollary}

\begin{proof}
As $\widehat{A}'=\widehat{A}\otimes_k k'$, it follows from \sref{0.19.3.6} and \sref{0.17.1.5} (applied to the local rings of $A'$ at its maximal ideals) that one can reduce to the case where $A$ is complete;
then $A'$ is a formally smooth $k'$-algebra \sref{0.19.3.5},~(iii) and is moreover a direct product of $k'$-local algebras, which are also formally smooth \sref{0.19.3.5},~(v).
One is therefore reduced to proving that $A$ is \emph{regular}.
Let $k_0$ be the prime subfield of $k$;
as $k$ is a separable extension of $k_0$, it is a formally smooth $k_0$-algebra \sref{0.19.6.1}, and by the hypothesis, so is $A$ \sref{0.19.3.5},~(ii).
As the residue field $K$ of $A$ is a separable extension of $k_0$, one can then apply \sref{0.19.6.4} to $A$ and $k_0$, whence the conclusion.
\end{proof}

\oldpage[0\textsubscript{IV-1}]{103}

\begin{corollary}[19.6.5.1]
\label{0.19.6.5.1}
Let $A$ be a Noetherian local ring, $k$ a subfield of $A$ such that $A$ is a formally smooth $k$-algebra (for its preadic topology).
Then every field $K$ such that $k\subset K\subset A$ is a separable extension of $k$.
\end{corollary}

\begin{proof}
Indeed, for every finite extension $k'$ of $k$, the ring $K\otimes_k k'$ is identified with a subring of $A'=A\otimes_k k'$;
as $A'$ is a regular ring, hence reduced, so is $K\otimes_k k'$, which proves that $K$ is a separable extension of $k$ (Bourbaki, \emph{Alg.}, chap.~VIII, \S\,7, n$^\circ$\,3, th.~1).
\end{proof}

We note that this proves the necessity of the condition in the statement of \sref{0.19.6.1}.

\begin{remark}[19.6.5.2]
\label{0.19.6.5.2}
If $K$ is a \emph{non-separable} extension of $k$, the formal power series ring $K[[T_1,\ldots,T_n]]$ is therefore \emph{not} a formally smooth $k$-algebra (for the usual $k$-algebra structure recalled in \sref{0.19.6.4}).
On the other hand, there exist formally smooth $k$-algebras $A$ which are complete Noetherian local rings whose residue field $K$ is a non-separable extension of $k$ (for example the completions of the localizations of $B=k[T_1,\ldots,T_n]$ at maximal ideals $\mathfrak{n}$ such that $B/\mathfrak{n}$ is a finite non-separable algebraic extension of $k$);
such a ring cannot therefore be $k$-isomorphic to $K[[T_1,\ldots,T_n]]$ even though it is \emph{isomorphic} to it by virtue of \sref{0.19.6.5}.

One can in certain cases sharpen the corollary \sref{0.19.6.5}:
\end{remark}

\begin{proposition}[19.6.6]
\label{0.19.6.6}
Let $k$ be a field, $A$ a local Noetherian $k$-algebra, $\mathfrak{m}$ its maximal ideal, $A$ being equipped with the $\mathfrak{m}$-preadic topology.
Suppose that the residue field $K=A/\mathfrak{m}$ of $A$ is such that there exists a finite radicial extension $k_0$ of $k$ for which the residue field of the local ring $K\otimes_k k_0$ is a separable extension of $k_0$ (we will express this property later by saying that $K$ is of \emph{finite radicial multiplicity} over $k$ ($\mathbf{IV}$, 4.7.4) and we will see that this condition is always satisfied when $K$ is an extension of finite type of $k$).
Then the following conditions are equivalent:
\begin{enumerate}
  \item[a)] $A$ is a formally smooth $k$-algebra.
  \item[b)] There exists a finite radicial extension $k'$ of $k$ such that $\widehat{A}\otimes_k k'$ is $k'$-isomorphic to a formal power series algebra $K'[[T_1,\ldots,T_n]]$, where $K'$ is a separable extension of $k'$.
  \item[b')] There exists a finite extension $k'$ of $k$ such that the semi-local ring $\widehat{A}\otimes_k k'$ has at least one local component which is $k'$-isomorphic to a formal power series algebra $K'[[T_1,\ldots,T_n]]$, where $K'$ is a separable extension of $k'$.
  \item[c)] $A$ is geometrically regular over $k$ \sref{0.19.6.5}.
\end{enumerate}
\end{proposition}

\begin{proof}
Note first that if $k'$ is a finite radicial extension of $k$, there is only a single maximal ideal of $A'=A\otimes_k k'$ above $\mathfrak{m}$, formed of elements whose $p^h$-th power ($p$ being the characteristic of $k$) is in $\mathfrak{m}$ for a suitable $h$ (Bourbaki, \emph{Alg.\ comm.}, chap.~V, \S\,2, n$^\circ$\,3, lemma~4);
$A'$ is therefore a local ring, and so is
$K\otimes_k k'=(A\otimes_k k')/(\mathfrak{m}\otimes_k k')$; moreover, these two local rings have the same residue field.
We recall, on the other hand, that if $K$ is a separable extension of $k$, then, for every finite extension $k''$ of $k$, $K\otimes_k k''$ is a direct product of fields (Bourbaki, \emph{Alg.}, chap.~VIII, \S\,7, n$^\circ$\,3, cor.~1 of th.~1);
therefore $\mathfrak{m}\otimes_k k''$ is the radical of $A\otimes_k k''$, and the field components of $K\otimes_k k''$ are the residue fields at the maximal ideals of $A\otimes_k k''$;

\oldpage[0\textsubscript{IV-1}]{104}

furthermore, these fields are separable extensions of $k''$, since for every extension $k_1$ of $k''$,
$K\otimes_k k_1=(K\otimes_k k'')\otimes_{k''}k_1$ has no radical (\emph{loc.\ cit.}, th.~1).
Applying these remarks to the field $K$ of finite radicial multiplicity over $k$, one sees that for \emph{every} finite extension $k''$ of $k_0$, $(K\otimes_k k'')_{\mathrm{red}}$ is separable over $k''$.

Let us pass to the proof of \sref{0.19.6.6}.
It is trivial that $b)$ implies $b')$.
Let us show that $b')$ implies $a)$.
Set $B'=K'[[T_1,\ldots,T_n]]$.
The hypothesis that $K'$ is separable over $k'$ entails, by the preceding remarks, that for every finite extension $k''$ of $k'$, the components of the complete semi-local ring
$B'\otimes_{k'}k''$ (equal to the formal power series ring
$(K'\otimes_{k'}k'')[[T_1,\ldots,T_n]]$) are the rings
$K_i''[[T_1,\ldots,T_n]]$, where the $K_i''$ are the field components of
$K'\otimes_{k'}k''$.
Since the local components of $B'\otimes_{k'}k''$ are also local components of
\[
  \widehat{A}\otimes_k k''
  =(\widehat{A}\otimes_k k')\otimes_{k'}k'',
\]
the hypothesis $b')$ remains true when $k'$ is replaced by any one of its finite extensions.
In particular, one can suppose $k'$ \emph{quasi-Galois}, hence a Galois extension of a radicial extension $k_0'$ of $k$;
$A'=\widehat{A}\otimes_k k'$ can then be written
$A_0'\otimes_{k_0'}k'$, where $A_0'=\widehat{A}\otimes_k k_0'$ is a complete local ring.
The Galois group of $k'/k_0'$ acts transitively on the local components of $A'$
(Bourbaki, \emph{Alg.}, chap.~VIII, \S\,8, prop.~4), so these components are
mutually Galois-conjugate. By hypothesis one of them is $k'$-isomorphic to a
formal power series ring $K'[[T_1,\ldots,T_n]]$, where $K'$ is a separable
extension of $k'$, and is therefore a formally smooth $k'$-algebra by
\sref{0.19.6.4}. Its conjugates are likewise formally smooth over $k'$, hence
their product $A'$ is formally smooth over $k'$ by \sref{0.19.3.5}, (v). Since
$k'$ is faithfully flat, projective and of finite type over $k_0'$, it follows
from \sref{0.19.4.6} that $A_0'$ is formally smooth over $k_0'$. Finally, as
$k_0'$ has the same properties over $k$, a second application of
\sref{0.19.4.6} shows that $\widehat{A}$ is a formally smooth $k$-algebra.

We have already seen \sref{0.19.6.5} that $a)$ implies $c)$; it remains to verify that $c)$ implies $b)$.
If $c)$ holds, then, for every finite radicial extension $k'$ of $k$ containing $k_0$, $A\otimes_k k'$ is a regular local ring whose residue field $K'$ is a separable extension of $k'$.
It follows from \sref{0.19.6.4} that its completion
$\widehat{A}\otimes_k k'$ is $k'$-isomorphic to a formal power series ring
$K'[[T_1,\ldots,T_n]]$.
\end{proof}

\begin{remarks}[19.6.7]
\label{0.19.6.7}
\begin{enumerate}
  \item[\rm(i)] We note that the hypothesis that $K$ is of finite radicial multiplicity over $k$ was only used in the proof of the implication $b')\Rightarrow a)$.
  \item[\rm(ii)] We will prove later (22.5.8) that $a)$ and $c)$ are equivalent, \emph{without} any hypothesis on the extension $K$ of $k$.
\end{enumerate}
\end{remarks}

\subsection{The case of local homomorphisms; existence and uniqueness theorems}
\label{subsection:0.19.7}

In this section, whenever a semi-local ring is considered as a topological ring, it is always understood that it is endowed with its $\mathfrak{r}$-preadic topology, where $\mathfrak{r}$ is its radical. Every local homomorphism of local rings is thus automatically continuous.

\begin{theorem}[19.7.1]
\label{0.19.7.1}
Let $A$, $B$ be two local Noetherian rings, $\mathfrak{m}$, $\mathfrak{n}$ their respective maximal ideals, $k=A/\mathfrak{m}$ the residue field of $A$; suppose that $A$ and $B$ are equipped respectively with the $\mathfrak{m}$-preadic and $\mathfrak{n}$-preadic topologies.
Let $\varphi:A\to B$ be a local homomorphism, and set $B_0=B\otimes_A k$.
The following properties are equivalent:
\begin{enumerate}
  \item[\rm a)] $B$ is a formally smooth $A$-algebra.
  \oldpage[0\textsubscript{IV-1}]{105}
  \item[\rm b)] $B$ is a flat $A$-module and $B_0=B/\mathfrak{m}B$ (equipped with the quotient topology) is a formally smooth $k$-algebra.
\end{enumerate}
\end{theorem}

\begin{proof}
The proof is carried out in several steps.

\medskip
\noindent\phantomsection\label{0.19.7.1.1}\textbf{(19.7.1.1)}\enspace
Let us first show that $b)$ implies $a)$; we will apply the criterion \sref{0.19.4.7}, with $\mathfrak{J}=\mathfrak{m}$; by virtue of the second hypothesis in $b)$, it suffices to show that $B$ is a \emph{formally projective} $A$-module.
The hypothesis implies that for all $h>0$, $B/\mathfrak{m}^h B$ is a flat $(A/\mathfrak{m}^h)$-module ($\mathbf{0_{III}}$, 10.2.1); as the $\mathfrak{m}^h$ form a fundamental system of neighbourhoods of $0$ in $A$ and $(B/\mathfrak{m}^h B)\otimes_{A/\mathfrak{m}^h} k = B_0$, we can replace $A$ and $B$ by $A/\mathfrak{m}^h$ and $B/\mathfrak{m}^h B$ respectively, and thus suppose that $A$ is \emph{Artinian} (hence discrete).
As $B_0$ is a $k$-algebra formally smooth, it is a regular ring \sref{0.19.6.5}; let $(x_i^0)_{1\leq i\leq n}$ be a regular system of parameters for $B_0$ \sref{0.17.1.6},
and for each $i$, let $x_i\in B$ be such that $x_i^0$ is its image in $B_0=B/\mathfrak{m}B$; as the $x_i^0$ generate the maximal ideal $\mathfrak{n}_0=\mathfrak{n}/\mathfrak{m}B$ of $B_0$, the ideals $\mathfrak{Q}_m=\sum_{i=1}^n (x_i^0)^m B_0$ (for $m>0$) form a fundamental system of neighbourhoods of $0$ in $B_0$, since $\mathfrak{Q}_m$ is evidently contained in $\mathfrak{n}_0^m$, and on the other hand contains $\mathfrak{n}_0^{mn}$.
Set $\mathfrak{J}_m=\sum_{i=1}^n x_i^m B$ for all $m>0$;
it is clear that $\mathfrak{n}=\mathfrak{J}_1+\mathfrak{m}B$; as there exists $h>0$ such that $\mathfrak{m}^h=0$, one has $\mathfrak{n}^m\subset\mathfrak{J}_1^{m-h}\subset\mathfrak{n}^{m-h}$ for $m>h$, and as we have seen that $\mathfrak{J}_m\supset\mathfrak{J}_1^{mn}$, the $\mathfrak{J}_m$ form a fundamental system of neighbourhoods of $0$ in $B$.
The whole argument then amounts to proving that the $B/\mathfrak{J}_m$ are \emph{free} $A$-modules, and it amounts to the same to verify that these are \emph{flat} $A$-modules ($\mathbf{0_{III}}$, 10.1.3).
Now, the hypothesis that $(x_i^0)$ is a $B_0$-regular sequence of elements of the maximal ideal of $B_0$ implies the same property for the sequence $(x_i^0)^m$ ($1\leq i\leq n$) for all $m>0$ \sref{0.15.1.20}; the conclusion thus follows from \sref{0.15.1.16}, $b)$ and $c)$.
\end{proof}

\begin{lemma}[19.7.1.2]
\label{0.19.7.1.2}
Let $A$ be a topological ring, $B$, $C$ two topological $A$-algebras that are local Noetherian rings, $\mathfrak{m}$, $\mathfrak{n}'$ their respective maximal ideals.
Suppose further that $C$ is complete and that the residue field $B/\mathfrak{m}$ of $B$ is an $A$-module of finite type.
Let $E$ be the completed tensor product $B\widehat{\otimes}_A C$.
Then:
\begin{enumerate}
  \item[\rm(i)] $E$ is a complete semi-local Noetherian ring.
  \item[\rm(ii)] The ideal $\mathfrak{m}E$ is contained in the radical of $E$, and for all $h>0$, $E/\mathfrak{m}^h E$ is isomorphic to $(B/\mathfrak{m}^h)\otimes_A C$.
  \item[\rm(iii)] If $C$ is a flat $A$-module, $E$ is a flat $B$-module.
\end{enumerate}
\end{lemma}

\begin{proof}
By definition, $E$ is the separated completion of the tensor product $B\otimes_A C$ for the topology defined by the ideals $\operatorname{Im}(\mathfrak{m}^i\otimes_A C)+\operatorname{Im}(B\otimes_A \mathfrak{n}'{}^j)$ ($\mathbf{0_I}$, 7.7.5).
If one sets $\mathfrak{r}=\operatorname{Im}(\mathfrak{m}\otimes_A C)+\operatorname{Im}(B\otimes_A \mathfrak{n}')$, one has $\mathfrak{r}^{2i}\subset\operatorname{Im}(\mathfrak{m}^i\otimes_A C)+\operatorname{Im}(B\otimes_A \mathfrak{n}'{}^i)\subset\mathfrak{r}^i$, so $E$ is also the separated completion of $B\otimes_A C$ for the $\mathfrak{r}$-preadic topology.
By hypothesis, $(B\otimes_A C)/\mathfrak{r}=(B/\mathfrak{m})\otimes_A (C/\mathfrak{n}')$ is a $(C/\mathfrak{n}')$-module of finite type, hence an Artinian ring;
moreover, $\mathfrak{r}/\mathfrak{r}^2$, being a quotient of $((\mathfrak{m}/\mathfrak{m}^2)\otimes_A C)\oplus(B\otimes_A (\mathfrak{n}'/\mathfrak{n}'{}^2))$, is a $(B\otimes_A C)$-module of finite type; applying ($\mathbf{0_I}$, 7.2.11), one sees that $E$ is a \emph{Noetherian} ring; moreover, $E/\mathfrak{r}E$, isomorphic to $(B\otimes_A C)/\mathfrak{r}$, being Artinian, $E$ is \emph{semi-local}.
Note now that $E$, which is isomorphic to $\varprojlim_{i,j}((B/\mathfrak{m}^i)\otimes_A(C/\mathfrak{n}'{}^j))$, is also isomorphic, by the theorem of the double projective limit, to $\varprojlim_i(\varprojlim_j((B/\mathfrak{m}^i)\otimes_A (C/\mathfrak{n}'{}^j)))$; but $\varprojlim_j((B/\mathfrak{m}^i)\otimes_A (C/\mathfrak{n}'{}^j))$
\oldpage[0\textsubscript{IV-1}]{106}
is the separated completion of $(B/\mathfrak{m}^i)\otimes_A C$, and as $C$ is complete, this is none other than $(B/\mathfrak{m}^i)\otimes_A C$ itself, since $B/\mathfrak{m}^i$ is an $A$-module of finite type (as $\mathfrak{m}/\mathfrak{m}^i$ is a $(B/\mathfrak{m})$-module of finite type ($\mathbf{0_I}$, 7.3.6)).
One therefore has $E=\varprojlim_i((B/\mathfrak{m}^i)\otimes_A C)$.
For every integer $h>0$, $\mathfrak{m}^h E$, being an ideal of $E$, is closed in $E$ ($\mathbf{0_I}$, 7.3.5), hence complete, and on the other hand it is evidently dense in $\varprojlim_i(\operatorname{Im}(\mathfrak{m}^h/\mathfrak{m}^{h+i})\otimes_A C)$, hence equal to this projective limit.
Moreover, all the projective systems considered are defined by \emph{surjective} homomorphisms; it follows therefore from ($\mathbf{0_{III}}$, 13.2.2) that $E/\mathfrak{m}^h E$ is isomorphic to $(B/\mathfrak{m}^h)\otimes_A C = (B\otimes_A C)/\operatorname{Im}(\mathfrak{m}^h\otimes_A C)$.
In particular $\operatorname{Im}(\mathfrak{m}\otimes_A C)\subset\mathfrak{r}$, which shows that $\mathfrak{m}E$ is contained in $\mathfrak{r}E$, hence in the radical of $E$.
Finally, the hypothesis that $C$ is a flat $A$-module implies that $(B/\mathfrak{m}^h)\otimes_A C = E/\mathfrak{m}^h E$ is a flat $(B/\mathfrak{m}^h)$-module for all $h>0$: since $B$ and $E$ are Noetherian and $\mathfrak{m}E$ is contained in the radical of $E$, it follows from ($\mathbf{0_{III}}$, 10.2.2) that $E$ is a flat $B$-module.
\end{proof}

\begin{lemma}[19.7.1.3]
\label{0.19.7.1.3}
Let $A$ be a local Noetherian ring, $\mathfrak{m}$ its maximal ideal, $k$ its residue field, $B_0$ a local Noetherian $k$-algebra; suppose that $B_0$ is complete and regular.
Then there exists a topological $A$-algebra $B$ that is a complete local Noetherian ring, a flat $A$-module, and such that $B_0$ is $k$-isomorphic to $B\otimes_A k = B/\mathfrak{m}B$.
\end{lemma}

\begin{proof}
As $\widehat{A}$ is flat over $A$ and has the same residue field, we can restrict to the case where $A$ is \emph{complete}.

Let $K$ be the residue field of $B_0$, and let us distinguish two cases:

\medskip
\noindent\emph{Case~I}: $K$ is a \emph{separable} extension of $k$.
By virtue of \sref{0.19.6.4}, $B_0$ is $k$-isomorphic to a formal power series ring $K[[T_1,\ldots,T_n]]$.
When $B_0=K$, the lemma has already been proved ($\mathbf{0_{III}}$, 10.3.1); let $C$ be a complete local Noetherian ring that is a flat $A$-module and such that $C\otimes_A k$ is isomorphic to $K$.
For $n\geq 1$, it suffices to take (using the preceding notation) $B=C[[T_1,\ldots,T_n]]$; indeed it is known (Bourbaki, \emph{Alg.\ comm.}, chap.~III, \S\,3, n$^\circ$\,4, cor.~3 of th.~1) that $B$ is a flat $C$-module, hence also a flat $A$-module, and on the other hand, it is immediate that $C[[T_1,\ldots,T_n]]\otimes_A k$ is isomorphic to $(C/\mathfrak{m}C)[[T_1,\ldots,T_n]]=B_0$.

\medskip
\noindent\emph{Case~II}: $K$ is of \emph{characteristic} $p>0$, and is therefore of the same characteristic as $k$.
Denote by $P$ the prime field $\mathbf{F}_p$, and by $W(P)$ the complete local ring of $p$-adic numbers $\mathbf{Z}_p$, which is a complete discrete valuation ring (hence \emph{regular}), with $P$ as residue field.
Let us first show that there exists a continuous ring homomorphism $W(P)\to A$, thus making $A$ into a $W(P)$-\emph{topological algebra}.
Indeed, if $j:\mathbf{Z}\to A$ is the canonical homomorphism, one has $j(p\mathbf{Z})\subset\mathfrak{m}$ by hypothesis, whence $j^{-1}(\mathfrak{m})=p\mathbf{Z}$, and $j$ therefore factors as $\mathbf{Z}\to\mathbf{Z}_{p\mathbf{Z}}\xrightarrow{f}A$, where $f$ is a local homomorphism, hence continuous, which (since $A$ is complete) extends by continuity to the desired homomorphism $W(P)\to A$.

As $k$ is a separable extension of $P$, case~I) shows that there is a local homomorphism $W(P)\to W(k)$, where $W(k)$ is a complete local Noetherian ring and a flat $W(P)$-module, such that $W(k)\otimes_{W(P)}P$ is isomorphic to $k$.
Moreover, as the uniformiser $p$ of $W(P)$ is a $W(k)$-regular element by flatness ($\mathbf{0_I}$, 6.3.4), and as
\oldpage[0\textsubscript{IV-1}]{107}
$W(k)/pW(k)=k$, the maximal ideal of $W(k)$ is $pW(k)$, which implies that this ring is a \emph{complete discrete valuation ring} (Bourbaki, \emph{Alg.\ comm.}, chap.~VI, \S\,3, n$^\circ$\,5, prop.~9).
By \sref{0.19.7.1.1} (as $k$ is separable over $P$, hence a formally smooth $P$-algebra \sref{0.19.6.1}), one sees that $W(k)$ is a $W(P)$-algebra \emph{formally smooth}.
The continuous $W(P)$-homomorphism $W(k)\to k$ therefore factors as $W(k)\to A\to k$ \sref{0.19.3.11},
which allows us to consider $A$ as a $W(k)$-topological algebra.
Applying now case~I) to $B_0$ considered as a $P$-algebra and to $W(P)$, one sees that there exists a $W(P)$-algebra $B_P$ which is a complete local Noetherian ring, a flat $W(P)$-module, and such that $B_P\otimes_{W(P)}P$ is $P$-isomorphic to $B_0$.
Using again the fact that $W(k)$ is a $W(P)$-algebra formally smooth, one sees by \sref{0.19.3.11} that the homomorphism $W(k)\to k\to B_0$ factors as $W(k)\to B_P\to B_0$; moreover, as $k=W(k)/pW(k)$, one has $B_P\otimes_{W(k)}k = B_P/pB_P = B_P\otimes_{W(P)}P = B_0$.
Let us show that $B_P$ is a flat $W(k)$-module; as $W(k)$ is a discrete valuation ring with $p$ as uniformiser, it suffices to verify that $B_P$ is a $W(k)$-module \emph{without torsion} ($\mathbf{0_I}$, 6.3.4), or equivalently that $p$ is a $B_P$-regular element, which follows from the fact that $B_P$ is a flat $W(P)$-module ($\mathbf{0_I}$, 6.3.4).

Set now
\[
  B = B_P\widehat{\otimes}_{W(k)} A
\]
and note that the residue field of $A$ being equal to that of $W(k)$, is \emph{a fortiori} a $W(k)$-module of finite type.
It follows first from \sref{0.19.7.1.2} that $B$ is a complete semi-local Noetherian ring, $\mathfrak{m}B$ is contained in the radical of $B$; moreover $B/\mathfrak{m}B$ is $k$-isomorphic to $B_P\otimes_{W(k)}k = B_0$, hence $B$ is in fact a \emph{local} ring.
As $B_P$ is a flat $W(k)$-module, \sref{0.19.7.1.2} shows finally that $B$ is a flat $A$-module. \textsc{q.e.d.}
\end{proof}

\begin{lemma}[19.7.1.4]
\label{0.19.7.1.4}
Let $A$ be a ring, $\mathfrak{J}$ an ideal of $A$, $M$, $N$ two $A$-modules separated for the $\mathfrak{J}$-preadic topology.
Suppose further that $N$ is complete for the $\mathfrak{J}$-preadic topology and that $M$ is a flat $A$-module.
Let $u:N\to M$ be an $A$-homomorphism; if $u\otimes 1:N\otimes_A(A/\mathfrak{J})\to M\otimes_A(A/\mathfrak{J})$ is bijective, then $u$ is bijective.
\end{lemma}

\begin{proof}
The graded modules being taken relative to the $\mathfrak{J}$-preadic filtrations, it follows from the hypotheses on $M$ and $N$ relative to the $\mathfrak{J}$-preadic topologies that it suffices to prove that $\gr(u):\gr_\bullet(N)\to\gr_\bullet(M)$ is \emph{bijective} (Bourbaki, \emph{Alg.\ comm.}, chap.~III, \S\,2, n$^\circ$\,8, cor.~3 of th.~1).
Now, there is a commutative diagram
\[
\xymatrix{
  \gr_0(N)\otimes_{A_0}\gr_\bullet(A) \ar[r]^-{\gr_0(u)\otimes 1} \ar[d]_{\varphi_N} & \gr_0(M)\otimes_{A_0}\gr_\bullet(A) \ar[d]^{\varphi_M} \\
  \gr_\bullet(N) \ar[r]_-{\gr(u)} & \gr_\bullet(M)
}
\]
\oldpage[0\textsubscript{IV-1}]{108}
where $A_0=A/\mathfrak{J}=\gr_0(A)$, and $\varphi_M$ and $\varphi_N$ are the canonical maps ($\mathbf{0_{III}}$, 10.1.1.2).
By hypothesis, $\gr_0(u)$ is bijective as well as $\varphi_M$ ($\mathbf{0_{III}}$, 10.2.1), and $\varphi_N$ is surjective; one deduces first that $\gr_0(u)\otimes 1$ is bijective, then that $\varphi_N$ is injective, hence bijective, and finally that $\gr(u)$ is bijective.
\end{proof}

\begin{lemma}[19.7.1.5]
\label{0.19.7.1.5}
Let $A$ be a Noetherian ring, $\mathfrak{J}$ an ideal of $A$, $B$, $B'$ two $A$-algebras that are local Noetherian rings, the homomorphisms $A\to B$, $A\to B'$ being continuous for the $\mathfrak{J}$-preadic topology on $A$.
Suppose that:
\begin{enumerate}
  \item[$1^\circ$] $B$ and $B'$ are complete for the $\mathfrak{J}$-preadic topologies;
  \item[$2^\circ$] $B$ is a formally smooth $A$-algebra;
  \item[$3^\circ$] $B'$ is a flat $A$-module.
\end{enumerate}
Set $A_0=A/\mathfrak{J}$, and let $u_0:B\otimes_A A_0\to B'\otimes_A A_0$ be an $A_0$-isomorphism; then there exists an $A$-isomorphism $u:B\to B'$ such that $u_0=u\otimes 1$ (which implies that $B'$ is also a formally smooth $A$-algebra and $B$ is a flat $A$-module).
\end{lemma}

\begin{proof}
Set $B_0=B\otimes_A A_0$, $B_0'=B'\otimes_A A_0$.
Note that if $\mathfrak{m}$ and $\mathfrak{m}'$ are the maximal ideals of $B$ and $B'$, the $\mathfrak{J}$-preadic topologies on $B$ and $B'$ are separated since $\mathfrak{J}B\subset\mathfrak{m}$ and $\mathfrak{J}B'\subset\mathfrak{m}'$; moreover, as $\mathfrak{J}B'$ is closed in $B'$ for the $\mathfrak{J}$-preadic topology, the composite homomorphism $B\to B_0\xrightarrow{u_0}B_0'$, which is continuous for the $\mathfrak{J}$-preadic topologies, factors as $B\xrightarrow{u}B'\to B'_0$, where $u$ is a continuous $A$-homomorphism \sref{0.19.3.10}.
One has $u_0=u\otimes 1$, and the hypothesis that $u_0$ is bijective implies that $u$ is bijective, by virtue of \sref{0.19.7.1.4}.
\end{proof}

\medskip
\noindent\textbf{(19.7.1.6)}\enspace
\emph{End of the proof of} \sref{0.19.7.1}. ---
To finish proving \sref{0.19.7.1}, it remains to show that $a)$ implies $b)$; it is already known that $a)$ implies that $B_0$ is a $k$-algebra formally smooth \sref{0.19.3.5}, (iii), hence it suffices to prove that $B$ is a flat $A$-module.
It amounts to the same to establish that $\widehat{B}$ is a flat $\widehat{A}$-module (Bourbaki, \emph{Alg.\ comm.}, chap.~III, \S\,5, n$^\circ$\,4, prop.~4), and one knows that $\widehat{B}$ is an $\widehat{A}$-algebra formally smooth \sref{0.19.3.6}; one can therefore restrict to the case where $A$ and $B$ are \emph{complete}.
As $B_0$ is a $k$-algebra formally smooth, it is a regular ring \sref{0.19.6.5} and complete ($\mathbf{0_I}$, 6.3.5); applying \sref{0.19.7.1.3}, one sees that there exists an $A$-algebra $B'$ which is a complete local Noetherian ring and a flat $A$-module, a local homomorphism $A\to B'$, and a $k$-isomorphism $B_0\simeq B_0'=B'\otimes_A k$.
It then suffices to apply \sref{0.19.7.1.5} (taking $\mathfrak{J}$ to be the maximal ideal of $A$) to obtain that $B$ is $A$-isomorphic to $B'$, hence is a flat $A$-module. \textsc{q.e.d.}

\begin{theorem}[19.7.2]
\label{0.19.7.2}
Let $A$ be a local Noetherian ring, $\mathfrak{J}$ an ideal contained in the maximal ideal of $A$, $A_0=A/\mathfrak{J}$, $B_0$ a complete local Noetherian ring, $A_0\to B_0$ a local homomorphism making $B_0$ into a formally smooth $A_0$-algebra.
Then there exists a complete local Noetherian ring $B$, a local homomorphism $A\to B$ making $B$ into a flat $A$-module, and an $A_0$-isomorphism $u:B\otimes_A A_0\xrightarrow{\sim} B_0$.
If $(B',u')$ is a pair satisfying the same conditions as $(B,u)$, there exists an $A$-isomorphism $v:B\xrightarrow{\sim} B'$ making the diagram
\[
\xymatrix{
  B\otimes_A A_0 \ar[rr]^{v\otimes 1}_{\sim} \ar[dr]^{\sim}_u && B'\otimes_A A_0 \ar[dl]^{u'}_{\sim} \\
  & B_0
}
\]
commutative.
\end{theorem}

\oldpage[0\textsubscript{IV-1}]{109}
\begin{proof}
Let $\mathfrak{m}$ be the maximal ideal of $A$, so that $\mathfrak{m}_0=\mathfrak{m}/\mathfrak{J}$ is the maximal ideal of $A_0$, $A$ and $A_0$ having the same residue field $k$.
Set $B_{00}=B_0\otimes_{A_0} k$; as $B_{00}$ is a formally smooth $k$-algebra \sref{0.19.3.5}, (iii), it is a regular local ring \sref{0.19.6.5}; applying \sref{0.19.7.1.3}, one sees that there exists a topological $A$-algebra $B$ which is a complete local Noetherian ring, a flat $A$-module, and for which one has a $k$-isomorphism $u_{00}:B\otimes_A k\xrightarrow{\sim} B_{00}$.
Note that by virtue of \sref{0.19.7.1}, $B$ is a formally smooth $A$-algebra, hence $B\otimes_A A_0=B/\mathfrak{J}B$ is a formally smooth $A_0$-algebra \sref{0.19.3.5}, (iii) and a complete local Noetherian ring; moreover, $B_0$ is a flat $A_0$-module by virtue of the hypothesis and of \sref{0.19.7.1}; as one has a $k$-isomorphism $u_{00}:B\otimes_A k=(B\otimes_A A_0)\otimes_{A_0} k\xrightarrow{\sim} B_0\otimes_{A_0} k = B_{00}$, one deduces from \sref{0.19.7.1.5}, applied by replacing $A$ by $A_0$ and $\mathfrak{J}$ by $\mathfrak{m}_0$, that there exists an $A_0$-isomorphism $u:B\otimes_A A_0\xrightarrow{\sim} B_0$ such that $u_{00}=u\otimes 1$.
As for the uniqueness assertion, note that the ideals $\mathfrak{J}^h B$ (resp.\ $\mathfrak{J}^h B'$) are closed in $B$ (resp.\ $B'$) ($\mathbf{0_I}$, 7.3.5), hence $B$ and $B'$ are separated and complete for the $\mathfrak{J}$-preadic topologies (Bourbaki, \emph{Top.\ g\'{e}n.}, chap.~III, 3$^\text{e}$ \'{e}d., \S\,3, n$^\circ$\,5, cor.~2 of prop.~9); one has by hypothesis an $A_0$-isomorphism $v_0:B\otimes_A A_0\xrightarrow{\sim} B'\otimes_A A_0$ such that $u'v_0=u$; as $B$ is a formally smooth $A$-algebra and $B'$ a flat $A$-module, one can apply \sref{0.19.7.1.5}, whence the existence of the $A$-isomorphism $v$ answering the question.
\end{proof}

\begin{remarks}[19.7.3]
\label{0.19.7.3}
\begin{enumerate}
  \item[\rm(i)] We note that the uniqueness assertion in \sref{0.19.7.2} is still valid if one only supposes that $B$ and $B'$ are complete \emph{for the $\mathfrak{J}$-preadic topologies}.
We do not know whether one can similarly improve the existence assertion, that is, whether one can dispense with assuming the local ring $B_0$ to be complete (for its $\mathfrak{n}_0$-preadic topology, where $\mathfrak{n}_0$ denotes its maximal ideal) by only requiring that $B$ be complete for the $\mathfrak{J}$-preadic topology.
When $A_0$ is complete for the $\mathfrak{m}$-preadic topology, one can see that this problem reduces to the following: if $B_0$ is a regular local Noetherian ring (not necessarily complete) containing the prime field $\mathbf{F}_p=\mathbf{Z}/p\mathbf{Z}$, does there exist for every $n>1$ a flat $(\mathbf{Z}/p^n\mathbf{Z})$-algebra $B$ such that $B/pB$ is isomorphic to $B_0$?
  \item[\rm(ii)] We note that in general, the isomorphism $v$ whose existence is asserted in \sref{0.19.7.2} is \emph{not unique} (cf.\ \sref{0.19.8.7}).
\end{enumerate}
\end{remarks}

\subsection{Cohen algebras and Cohen rings; application to the structure of complete local rings}
\label{subsection:0.19.8}
\label{0.19.8}

The results of this section are immediate applications of the theorems of \sref{subsection:0.19.7}, but deserve to be stated explicitly because of their practical importance.

\begin{definition}[19.8.1]
\label{0.19.8.1}
Let $A$, $B$ be two local Noetherian rings, $\mathfrak{m}$ the maximal ideal of $A$, $k=A/\mathfrak{m}$ its residue field, $\varphi:A\to B$ a local homomorphism, making $B$ into an $A$-algebra.
We say that $B$ is a \emph{Cohen $A$-algebra} if it satisfies the following conditions:
\begin{enumerate}
  \item[\rm(i)] $B$ is a complete ring.
  \item[\rm(ii)] $B$ is a flat $A$-module.
  \item[\rm(iii)] $B\otimes_A k$ is a field (in other words, $\mathfrak{m}B$ is the maximal ideal of $B$) which is a separable extension of $k$.
\end{enumerate}
\end{definition}

\oldpage[0\textsubscript{IV-1}]{110}
\begin{theorem}[19.8.2]
\label{0.19.8.2}
Let $A$ be a local Noetherian ring, $k$ its residue field.
\begin{enumerate}
  \item[\rm(i)] If $B$ is a Cohen $A$-algebra, $B$ is a formally smooth $A$-algebra.
For every complete local Noetherian ring $C$, every local homomorphism $A\to C$ and every ideal $\mathfrak{J}\neq C$ in $C$, every $A$-homomorphism $B\to C/\mathfrak{J}$ therefore factors as $B\xrightarrow{v}C\to C/\mathfrak{J}$, where $v$ is an $A$-homomorphism (necessarily local).
  \item[\rm(ii)] For every field $K$, separable extension of $k$, there exists a Cohen $A$-algebra $B$ such that $B\otimes_A k$ is $k$-isomorphic to $K$, and such an $A$-algebra is unique up to isomorphism.
\end{enumerate}
\end{theorem}

\begin{proof}
As $B\otimes_A k$ is a separable field extension of $k$, it is a formally smooth $k$-algebra \sref{0.19.6.1}, so assertion (i) follows from \sref{0.19.7.1}.
To prove (ii), one can restrict to the case where $A$ is \emph{complete}, for it amounts to the same to say that $B$ is a flat $A$-module or a flat $\widehat{A}$-module ($\mathbf{0_{III}}$, 10.2.3), one has $\mathfrak{m}B=\mathfrak{m}\widehat{A}B$ and $k$ is the residue field of $\widehat{A}$.
It then suffices to apply \sref{0.19.7.2} taking $\mathfrak{J}=\mathfrak{m}$ and $B_0=K$ (and using \sref{0.19.6.1}).
\end{proof}

\begin{definition}[19.8.3]
\label{0.19.8.3}
A \emph{prime local ring} is a local ring of the form $\mathbf{Z}_{p\mathbf{Z}}$, where $p\mathbf{Z}$ is a prime ideal of $\mathbf{Z}$.
A \emph{complete prime local ring} is the completion of a prime local ring.
\end{definition}

The prime local rings are therefore of two kinds:
\begin{enumerate}
  \item[$1^\circ$] Those corresponding to maximal ideals $p\mathbf{Z}$ where $p\neq 0$ is a prime number; $\mathbf{Z}_{p\mathbf{Z}}$ is a discrete valuation ring, whose completion is the ring of $p$-adic integers usually denoted $\mathbf{Z}_p$.\footnote{This notation, now universally used, is in conflict with the notation $A_f$ adopted in ($\mathbf{0_I}$, 1.2.3): with $A=\mathbf{Z}$ and $f=p$, $A_f$ denotes the ring of rational numbers of the form $k/p^n$ ($k\in\mathbf{Z}$, $n$ an integer $\geq 0$); we will always avoid using the notation $\mathbf{Z}_p$ to denote this latter ring.}
  \item[$2^\circ$] For the prime ideal $p\mathbf{Z}=(0)$, $\mathbf{Z}_{p\mathbf{Z}}$ is the field of rational numbers $\mathbf{Q}$, which is identical to its completion (the topology being naturally the topology of a local Noetherian ring, hence discrete here).
\end{enumerate}

The terminology of \sref{0.19.8.3}, analogous to that of ``prime fields'', is justified in the same way: for every local ring $A$, consider the canonical homomorphism $\varphi:\mathbf{Z}\to A$, and let $p\mathbf{Z}$ be the inverse image under this homomorphism of the maximal ideal $\mathfrak{m}$ of $A$; $p\mathbf{Z}$ is a prime ideal of $\mathbf{Z}$ and the preceding homomorphism therefore factors as $\mathbf{Z}\to\mathbf{Z}_{p\mathbf{Z}}\xrightarrow{\psi}A$; moreover, as $\varphi$ is the \emph{unique} homomorphism of $\mathbf{Z}$ into $A$, $p$ and $\psi$ are determined in a unique way.
In other words, for every local ring $A$, there is a unique homomorphism $\psi:P\to A$, where $P$ is a prime local ring; if moreover $A$ is separated and complete, one can extend this homomorphism by completion, and there is thus a unique homomorphism $\hat{\psi}:\widehat{P}\to A$, where $\widehat{P}$ is a complete prime local ring.
Moreover, $\psi$ gives by passage to quotients a homomorphism of the residue field $\mathbf{F}_p$ if $p>0$ (resp.\ $\mathbf{Q}$ if $p=0$) into the residue field $k$ of $A$, and $p$ is therefore the \emph{characteristic} of $k$.

If one takes in particular $A$ to be a prime local ring (resp.\ complete prime local ring), one sees that there exists in such a ring only \emph{a single endomorphism}, namely the identity.

\begin{definition}[19.8.4]
\label{0.19.8.4}
\oldpage[0\textsubscript{IV-1}]{111}
Let $A$ be a local ring, $P\to A$ the unique homomorphism from a prime local ring $P$ into $A$, $p$ the characteristic of the residue fields of $P$ and $A$.
We say that $A$ is a \emph{Cohen ring} if it is a Cohen $P$-algebra, that is \sref{0.19.8.1}, if:
\begin{enumerate}
  \item[$1^\circ$] $A$ is Noetherian and complete.
  \item[$2^\circ$] $A$ is a flat $P$-module (which also amounts to saying that $A$ is a flat $\widehat{P}$-module (Bourbaki, \emph{Alg.\ comm.}, chap.~III, \S\,5, n$^\circ$\,4, prop.~4)).
  \item[$3^\circ$] $A/pA$ is a field (necessarily separable over the residue field of $P$, this field being prime).
\end{enumerate}
\end{definition}

If $p=0$, these conditions are equivalent to saying that $A$ is a field of characteristic $0$.
If $p>0$, one necessarily has $pA\neq 0$; condition $3^\circ$ means that $pA$ is the maximal ideal $\mathfrak{m}$ of $A$; condition $2^\circ$ means that $p$ is $A$-regular, since $P$ is a discrete valuation ring ($\mathbf{0_I}$, 6.3.4).
Hence $A$ is a regular ring \sref{0.17.1.1}, $d)$) of dimension $1$, and therefore a discrete valuation ring, complete by virtue of $1^\circ$; in summary:

\begin{proposition}[19.8.5]
\label{0.19.8.5}
The Cohen rings are the fields of characteristic $0$ and the complete discrete valuation rings whose residue field has characteristic $p>0$, and whose maximal ideal is generated by $p\cdot 1$ ($1$ being the unity of the ring).
\end{proposition}

We note that in the second case, $p\cdot 1\neq 0$ since $p$ is $A$-regular, hence one can identify $p\cdot 1$ with the integer $p$; the canonical homomorphism $\mathbf{Z}_p\to A$ is injective, and one identifies $p\cdot 1$ with the element $p$ of $\mathbf{Z}_p$; one says in this case that $A$ is a \emph{$p$-ring} (or \emph{Cohen $p$-ring}).

\begin{theorem}[19.8.6]
\label{0.19.8.6}
\textup{(Cohen).} ---
\begin{enumerate}
  \item[\rm(i)] Let $W$ be a Cohen ring, $C$ a complete local Noetherian ring, $\mathfrak{J}$ an ideal of $C$ distinct from $C$.
Then every local homomorphism $u:W\to C/\mathfrak{J}$ factors as $W\xrightarrow{v}C\to C/\mathfrak{J}$ where $v$ is a local homomorphism.
  \item[\rm(ii)] Let $K$ be a field.
There exists a Cohen ring $W$ whose residue field is isomorphic to $K$.
If $W'$ is a second Cohen ring, $K'$ its residue field, every isomorphism $u:K\xrightarrow{\sim}K'$ arises by passage to quotients from an isomorphism $v:W\xrightarrow{\sim}W'$.
\end{enumerate}
\end{theorem}

\begin{proof}
This is nothing other than \sref{0.19.8.2} applied to the case where $A$ is a prime local ring.
\end{proof}

\begin{remarks}[19.8.7]
\label{0.19.8.7}
\begin{enumerate}
  \item[\rm(i)] When $K$ is of characteristic $0$, part (ii) of \sref{0.19.8.6} becomes trivial.
  \item[\rm(ii)] The homomorphism $v$ of \sref{0.19.8.6}, (i) is \emph{not} necessarily determined uniquely by $u$, as is shown already by the case where $W$ is a field of characteristic $0$, $\mathfrak{J}^2=0$ and $u$ is an isomorphism (cf.\ (21.5.5)).
Similarly, in \sref{0.19.8.6}, (ii), the isomorphism $v$ is not necessarily determined uniquely by $u$ (cf.\ (21.5.5)).
However, when $K$ is \emph{perfect} and of characteristic $p>0$, one will see (21.5.5) that in \sref{0.19.8.6}, (ii) the isomorphism $v$ is unique.
One will also see later that in this case $W$ identifies with the ring $W_\infty(K)$ of \emph{Witt vectors of infinite length} over $K$.
  \item[\rm(iii)] In \sref{0.19.8.6}, (i), one can weaken the hypotheses on $C$ by using \sref{0.19.3.10} and \sref{0.19.3.12}.
\end{enumerate}
\end{remarks}

\begin{theorem}[19.8.8]
\label{0.19.8.8}
\oldpage[0\textsubscript{IV-1}]{112}
\textup{(Cohen).} ---
Let $A$ be a complete local Noetherian ring, $k$ its residue field.
\begin{enumerate}
  \item[\rm(i)] There exists a Cohen ring $W$ such that $A$ is isomorphic to a quotient ring of a formal power series ring $W[[T_1,\ldots,T_n]]$ (and in particular $A$ is isomorphic to a quotient of a complete regular local ring \sref{0.17.3.8}).
If $A$ contains a field, it is isomorphic to a quotient ring of $k[[T_1,\ldots,T_n]]$.
  \item[\rm(ii)] Suppose moreover that $A$ is an integral domain.
Then there exists a subring $B$ of $A$ such that: $1^\circ$ $B$ is isomorphic to a formal power series ring over a ring $C$ which is a field or a Cohen ring (which implies that $B$ is a complete regular local ring \sref{0.17.3.8}); $2^\circ$ $B$ has the same residue field as $A$ and the injection $B\to A$ is a local homomorphism; $3^\circ$ $A$ is a finite $B$-algebra.
\end{enumerate}
\end{theorem}

\begin{proof}
Let $\mathfrak{m}$ be the maximal ideal of $A$.
There exists a Cohen ring $W$ whose residue field is isomorphic to $k$ \sref{0.19.8.6}, (ii); one therefore has a homomorphism $W\to A/\mathfrak{m}$, which then factors as $W\xrightarrow{u}A\to A/\mathfrak{m}$, where $u$ is a local homomorphism \sref{0.19.8.6}, (i).
For every finite family $(x_i)_{1\leq i\leq n}$ of elements of $\mathfrak{m}$, there then exists a local homomorphism $v:W[[T_1,\ldots,T_n]]\to A$ extending $u$ and such that $v(T_i)=x_i$ for all $i$ (Bourbaki, \emph{Alg.\ comm.}, chap.~III, \S\,4, n$^\circ$\,5, prop.~6).
When $A$ contains a field, it contains a prime field $P$, of which $k$ is an extension (necessarily separable), and therefore $A$ contains a field isomorphic to $k$ \sref{0.19.6.2}; one can then replace $W$ by $k$ in the preceding definition of $v$.

\medskip
\noindent(i)\enspace
Let us first take the $x_i$ to be a system of generators of $\mathfrak{m}$.
As $W$ has the same residue field as $A$, and the classes of the $x_i$ generate the graded ring $\gr_\bullet(A)$ as a $k$-algebra, $\gr(v):\gr_\bullet(W[[T_1,\ldots,T_n]])\to\gr_\bullet(A)$ is surjective; one deduces that $v$ is itself surjective (Bourbaki, \emph{Alg.\ comm.}, chap.~III, \S\,2, n$^\circ$\,8, cor.~2 of th.~1).
We recall that the case where $A$ contains a field has already been seen and is included here only for convenience \sref{0.19.6.3}.

\medskip
\noindent(ii)\enspace
If $A$ contains a field, it contains a field $k'$ isomorphic to $k$ as we have seen; one then considers a system of parameters $(y_j)_{1\leq j\leq m}$ of $A$ \sref{0.16.3.6}, one takes $B=k[[T_1,\ldots,T_m]]$ and one considers the local homomorphism $w:B\to A$ which coincides in $k$ with an isomorphism $k\to k'$ and is such that $w(T_j)=y_j$ for $1\leq j\leq m$.
If $A$ does not contain a field, the unique homomorphism $\mathbf{Z}_{p\mathbf{Z}}\to A$ \sref{0.19.8.3} is necessarily injective (for otherwise, as $A$ is an integral domain, its kernel would be the maximal ideal of $\mathbf{Z}_{p\mathbf{Z}}$ and its image would be isomorphic to a field; moreover one has $p>0$ by hypothesis, $\mathbf{Z}_{p\mathbf{Z}}$ being a field if $p=0$); the element $p\cdot 1$ of $A$ (identified with $p$) is a non-zero-divisor in $A$, and is contained in $\mathfrak{m}$, hence \sref{0.16.3.4} and \sref{0.16.3.7} there exists a family $(z_i)_{1\leq i\leq m-1}$ which, with $p$, forms a system of parameters of $A$.
The Cohen ring $W$ considered at the beginning of the proof is then a discrete valuation ring whose residue field is isomorphic to $k$, in which $p$ generates the maximal ideal \sref{0.19.8.5}, and the homomorphism $u:W\to A$ defined at the beginning maps $p$ to itself.
One then takes $B=W[[T_1,\ldots,T_{m-1}]]$ and considers the local homomorphism $w:B\to A$ which coincides with $u$ in $W$ and is such that $w(T_i)=z_i$ for $1\leq i\leq m-1$.
\oldpage[0\textsubscript{IV-1}]{113}
In both cases, if $\mathfrak{n}$ is the maximal ideal of $B$, it is clear that $\mathfrak{n}A$ is an ideal of definition of $A$; as moreover $B/\mathfrak{n}$ and $A/\mathfrak{m}$ are isomorphic, $A$ is a quasi-finite $B$-module ($\mathbf{0_I}$, 7.4.4), hence a $B$-module of finite type since $B$ is complete and $A$ is separated for the $\mathfrak{n}$-preadic topologies ($\mathbf{0_I}$, 7.4.1).
On the other hand, in both cases one has $\dim(B)=\dim(A)=m$; in the first case, this follows from \sref{0.17.1.4}, (iii); in the second, one sees directly that $p$ and the $T_i$ ($1\leq i\leq m-1$) form a $B$-regular sequence generating $\mathfrak{n}$, or one can also use the fact that these elements generate $\mathfrak{n}$ and that $\dim(B)\geq\dim(A)$ by \sref{0.16.3.10}.
As $A$ and $B$ are integral domains, one finally deduces from \sref{0.16.3.10} that $w$ is injective, which completes the proof.
\end{proof}

\begin{corollary}[19.8.9]
\label{0.19.8.9}
Let $A$ be a complete local Noetherian integral domain containing a field $k_0$; let $k$ be the residue field of $A$, and suppose that $k$ is finite over $k_0$.
Then, in the conclusion of \sref{0.19.8.8}, (ii), one can replace $1^\circ$ and $2^\circ$ by the condition that $B$ is of the form $k_0[[T_1,\ldots,T_n]]$, the canonical injection $B\to A$ being a local $k_0$-homomorphism (for the usual $k_0$-algebra structure of $B$).
\end{corollary}

\begin{proof}
Indeed, taking up the proof of \sref{0.19.8.8}, (ii), one defines this time $w:k_0[[T_1,\ldots,T_n]]\to A$ as coinciding in $k_0$ with the identity and mapping $T_j$ to $y_j$ for $1\leq j\leq n$.
The hypothesis that $k$ is of finite degree over $k_0$ again implies that $A$ is a quasi-finite $B$-module, hence of finite type by ($\mathbf{0_I}$, 7.4.1), and one concludes as in \sref{0.19.8.8}.
\end{proof}

\begin{corollary}[19.8.10]
\label{0.19.8.10}
Let $A$ be a local Artinian ring whose maximal ideal $\mathfrak{m}$ is of square zero; then there exists a regular local Noetherian ring $B$, with maximal ideal $\mathfrak{n}$, such that $A$ is isomorphic to $B/\mathfrak{n}^2$.
\end{corollary}

\begin{proof}
Let $K$ be the residue field $A/\mathfrak{m}$ of $A$, and $n$ the rank of $\mathfrak{m}/\mathfrak{m}^2=\mathfrak{m}$ over $K$.
If $A$ contains a field, it follows from \sref{0.19.6.3} that $A$ is isomorphic to $B/\mathfrak{b}$, where $B=K[[T_1,\ldots,T_n]]$ and $\mathfrak{b}$ is contained in the square of the maximal ideal $\mathfrak{n}$ of $B$; but as $\operatorname{long}(B/\mathfrak{n}^2)=n+1=\operatorname{long}(A)$, one necessarily has $\mathfrak{b}=\mathfrak{n}^2$.

Suppose next that $A$ does not contain a field; this implies that $K$ is of characteristic $p>0$ and that $p\cdot 1\neq 0$ in $A$ \sref{0.19.6.3}; hence $p\cdot 1$ is an element of $\mathfrak{m}$, and there are therefore $n-1$ other elements $x_i$ ($2\leq i\leq n$) of $\mathfrak{m}$ forming with $p\cdot 1$ a basis of $\mathfrak{m}$ over $K$.
Let $W$ be a Cohen ring whose residue field is isomorphic to $K$; $W$ is a discrete valuation ring in which $p$ generates the maximal ideal; one has seen in the proof of \sref{0.19.8.8} that there is a homomorphism $u:W\to A$ mapping $p$ to itself and which by passage to quotients gives the identity on $K$.
One takes $B=W[[T_2,\ldots,T_n]]$ and considers the local homomorphism $w:B\to A$ which coincides with $u$ in $W$ and is such that $w(T_i)=x_i$ for $2\leq i\leq n$.
It is clear that $w$ is surjective and that its kernel $\mathfrak{b}$ is contained in the square of the maximal ideal $\mathfrak{n}=pB+BT_2+\ldots+BT_n$ of $B$; as $\operatorname{long}(B/\mathfrak{n}^2)=n+1=\operatorname{long}(A)$, one again has $\mathfrak{b}=\mathfrak{n}^2$.
\end{proof}

\begin{proposition}[19.8.11]
\label{0.19.8.11}
Let $A$ be a local Artinian ring, $\mathfrak{m}$ its maximal ideal, $k$ its residue field.
For $A$ to be isomorphic to a quotient ring of a Cohen ring, it is necessary and sufficient that $\mathfrak{m}$ be generated by $p\cdot 1$, where $p$ is the characteristic of $k$.
\end{proposition}

\begin{proof}
\oldpage[0\textsubscript{IV-1}]{114}
The condition is evidently necessary \sref{0.19.8.5}.
To see that it is sufficient, one observes, as at the beginning of the proof of \sref{0.19.8.8}, that there exists a Cohen ring $W$ whose residue field is isomorphic to $k$ and a local homomorphism $u:W\to A$.
Moreover, if one considers the composite homomorphism $\mathbf{Z}_{p\mathbf{Z}}\to W\xrightarrow{u}A$ (which is necessarily the unique homomorphism of $\mathbf{Z}_{p\mathbf{Z}}$ into $A$), one sees that the image under $u$ of the element $p\cdot 1$ of $W$ is the element $p\cdot 1$ of $A$; as the element $p\cdot 1$ of $W$ generates the maximal ideal of this ring, one deduces immediately from the hypothesis that $\gr(u):\gr_\bullet(W)\to\gr_\bullet(A)$ is surjective, and therefore so is $u$ (Bourbaki, \emph{Alg.\ comm.}, chap.~III, \S\,2, n$^\circ$\,8, cor.~2 of th.~1).
\end{proof}

\subsection{Relatively formally smooth algebras}
\label{subsection:0.19.9}

\begin{definition}[19.9.1]
\label{0.19.9.1}
Let $\Lambda$ be a topological ring, $A$ a topological $\Lambda$-algebra, $B$ a topological $A$-algebra.
We say that $B$ is a \emph{formally smooth $A$-algebra relative to $\Lambda$} if, for every discrete topological $A$-algebra $C$, and every nilpotent ideal $\mathfrak{J}$ of $C$, every continuous $A$-homomorphism $u_0:B\to C/\mathfrak{J}$ which factors as $B\xrightarrow{u}C\to C/\mathfrak{J}$, where $u$ is a continuous $\Lambda$-homomorphism, also factors as $B\xrightarrow{v}C\to C/\mathfrak{J}$, where $v$ is a continuous $A$-homomorphism.
\end{definition}

It follows from this definition that if $B$ is a \emph{formally smooth} $A$-algebra, then $B$ is also formally smooth relative to $\Lambda$, for every structure of topological $\Lambda$-algebra defined on $A$ (in other words, every continuous ring homomorphism $\Lambda\to A$).

\begin{proposition}[19.9.2]
\label{0.19.9.2}
Let $\Lambda$ be a topological ring, $A$ a topological $\Lambda$-algebra.
\begin{enumerate}
  \item[\rm(i)] $A$ is a $\Lambda$-algebra formally smooth relative to $\Lambda$.
  \item[\rm(ii)] If $B$ is an $A$-algebra formally smooth relative to $\Lambda$ and $C$ is a $B$-algebra formally smooth relative to $\Lambda$, then $C$ is an $A$-algebra formally smooth relative to $\Lambda$.
  \item[\rm(iii)] Let $B$ be an $A$-algebra formally smooth relative to $\Lambda$, $A'$ a topological $A$-algebra, then the topological $A'$-algebra $B\otimes_A A'$ is formally smooth relative to $\Lambda$.
  \item[\rm(iv)] Let $B$ be a topological $A$-algebra, $S$ (resp.\ $T$) a multiplicative subset of $A$ (resp.\ $B$) such that the canonical image of $S$ in $B$ is contained in $T$.
If $B$ is an $A$-algebra formally smooth relative to $\Lambda$, then $T^{-1}B$ is an $S^{-1}A$-algebra formally smooth relative to $\Lambda$.
  \item[\rm(v)] Let $B_i$ ($1\leq i\leq n$) be topological $A$-algebras.
For $\prod_{i=1}^n B_i$ to be an $A$-algebra formally smooth relative to $\Lambda$, it is necessary and sufficient that each of the $B_i$ be so.
\end{enumerate}
\end{proposition}

\begin{proof}
Assertion (i) is trivial, and the proof of the others closely follows the proofs of \sref{0.19.3.5}; it is therefore left to the reader.
\end{proof}

\begin{corollary}[19.9.3]
\label{0.19.9.3}
Let $A$ be a topological ring, $A'$ and $B$ two topological $A$-algebras.
Then the topological $A$-algebra $B'=B\otimes_A A'$ is a $B$-algebra formally smooth relative to $A$.
\end{corollary}

\begin{proof}
This follows from \sref{0.19.9.2}, (i) and (iii).
\end{proof}

\begin{proposition}[19.9.4]
\label{0.19.9.4}
Let $\Lambda$ be a topological ring, $A$ a topological $\Lambda$-algebra, $B$ a topological $A$-algebra.
The following conditions are equivalent:
\begin{enumerate}
  \item[\rm a)] $B$ is an $A$-algebra formally smooth relative to $\Lambda$.
  \item[\rm b)] $\widehat{B}$ is an $A$-algebra formally smooth relative to $\Lambda$.
  \item[\rm c)] $\widehat{B}$ is an $\widehat{A}$-algebra formally smooth relative to $\Lambda$.
  \item[\rm d)] $\widehat{B}$ is an $\widehat{A}$-algebra formally smooth relative to $\widehat{\Lambda}$.
\end{enumerate}
\end{proposition}

\begin{proof}
The proof is again left to the reader, modelled on that of \sref{0.19.3.6}.
\end{proof}

\oldpage[0\textsubscript{IV-1}]{115}
\begin{env}[19.9.5]
\label{0.19.9.5}
Similarly, the statement \sref{0.19.3.8} is still valid (with the same proof) when one replaces ``formally smooth'' by ``formally smooth relative to $\Lambda$''.
If in the statement of \sref{0.19.3.10} one replaces ``formally smooth'' by ``formally smooth relative to $\Lambda$'', the conclusion is replaced by the following (the proof remaining essentially unchanged): \emph{every continuous $A$-homomorphism $u_0:B\to C/\mathfrak{J}$ which factors as $B\xrightarrow{u}C\to C/\mathfrak{J}$, where $u$ is a continuous $\Lambda$-homomorphism, also factors as $B\xrightarrow{v}C\to C/\mathfrak{J}$, where $v$ is a continuous $A$-homomorphism}.
\end{env}

\begin{env}[19.9.6]
\label{0.19.9.6}
The criteria for formal smoothness \sref{0.19.4.1} and \sref{0.19.4.2} are also valid when one replaces ``formally smooth'' by ``formally smooth relative to $\Lambda$'', the proofs remaining practically unchanged.
\end{env}

\begin{proposition}[19.9.7]
\label{0.19.9.7}
Let $\Lambda$ be a topological ring, $A$ a topological $\Lambda$-algebra, $B$ a topological $A$-algebra.
Suppose that for every discrete $A$-algebra $C$ and every ideal $\mathfrak{J}$ of $C$ such that $\mathfrak{J}^2=0$, every continuous $A$-homomorphism $u_0:B\to C/\mathfrak{J}$ which factors as $B\xrightarrow{u}C\to C/\mathfrak{J}$, where $u$ is a continuous $\Lambda$-homomorphism, also factors as $B\xrightarrow{v}C\to C/\mathfrak{J}$, where $v$ is a continuous $A$-homomorphism.
Then $B$ is an $A$-algebra formally smooth relative to $\Lambda$.
\end{proposition}

\begin{proof}
The proof of \sref{0.19.4.3} transcribes immediately.
\end{proof}

\begin{proposition}[19.9.8]
\label{0.19.9.8}
Let $\Lambda$ be a topological ring, $A$ a topological $\Lambda$-algebra, $B$ a topological $A$-algebra.
For $B$ to be an $A$-algebra formally smooth relative to $\Lambda$, it is necessary and sufficient that for every discrete topological $B$-module $L$, annihilated by an open ideal of $B$, one has (cf.\ 18.4.2)
\[
  \operatorname{Exalcotop}_{A/\Lambda}(B,L) = 0.
\]
\end{proposition}

\begin{proof}
Using the notation of the proof of \sref{0.19.4.4}, it suffices to note here that one can suppose that the extension $E_\Lambda$ of $B_\Lambda$ is $\Lambda$-trivial; the rest of the proof is then unchanged.
\end{proof}

When $\Lambda$, $A$ and $B$ are discrete rings, the criterion \sref{0.19.9.8} reduces to
\begin{equation}
\tag{19.9.8.1}
  \label{0.19.9.8.1}
  \operatorname{Exalcom}_{A/\Lambda}(B,L) = 0 \quad\text{for every $B$-module $L$};
\end{equation}
in other words, \emph{every commutative $A$-extension of $B$ by a $B$-module, which is also $\Lambda$-trivial, is also $A$-trivial}.

\subsection{Formally unramified algebras and formally \'{e}tale algebras}
\label{subsection:0.19.10}
\label{0.19.10}

\begin{env}[19.10.1]
\label{0.19.10.1}
We will also need two notions closely related to that of a formally smooth algebra.
\end{env}

\begin{definition}[19.10.2]
\label{0.19.10.2}
Let $A$ be a topological ring, $B$ a topological $A$-algebra.
We say that $B$ is a \emph{formally unramified} $A$-algebra if, for every $A$-homomorphism $p:E\to C$ of discrete $A$-algebras, whose kernel is nilpotent, and every continuous $A$-homomorphism $u:B\to C$, there exists \emph{at most one} continuous $A$-homomorphism $v:B\to E$ such that $u$ factors as $B\xrightarrow{v}E\xrightarrow{p}C$.
We say that $B$ is a \emph{formally \'{e}tale} $A$-algebra if it is formally smooth and formally unramified, in other words, if, for every surjective $A$-homomorphism $p:E\to C$ of discrete $A$-algebras, whose kernel is nilpotent, and for every continuous $A$-homomorphism $u:B\to C$, there exists \emph{one and only one} continuous $A$-homomorphism $v:B\to E$ such that $u$ factors as $B\xrightarrow{v}E\xrightarrow{p}C$.
\end{definition}

\oldpage[0\textsubscript{IV-1}]{116}
\begin{examples}[19.10.3]
\label{0.19.10.3}
We restrict here to \emph{discrete} rings.
\begin{enumerate}
  \item[\rm(i)] If the structural homomorphism $\rho:A\to B$ is \emph{surjective}, $B$ is a \emph{formally unramified} $A$-algebra, for with the notation of \sref{0.19.10.2}, the composite homomorphism $A\xrightarrow{\rho}B\xrightarrow{v}E$ must be the structural homomorphism, hence $v$ is unique (if it exists).
On the other hand, $B$ is formally smooth if and only if $\ker(\rho)=\ker(\rho)^2$.
  \item[\rm(ii)] Let $A$ be a ring, $S$ a multiplicative subset of $A$; then $B=S^{-1}A$ is a \emph{formally \'{e}tale} $A$-algebra.
Indeed, let $j:A\to B$ be the canonical homomorphism, $\rho:A\to E$ a ring homomorphism, $\mathfrak{J}$ a nilpotent ideal of $E$, $p:E\to C=E/\mathfrak{J}$ the canonical homomorphism.
If $u:B\to C$ is an $A$-homomorphism, $u(s/1)$ is invertible in $C$ for $s\in S$ since $s/1$ is invertible in $B$; but $u(s/1)=p(\rho(s))$, and as $\mathfrak{J}$ is nilpotent, the fact that the image of $\rho(s)$ in $E/\mathfrak{J}$ is invertible implies that $\rho(s)$ is itself invertible in $E$ ($\mathbf{0_I}$, 7.1.12); one concludes that $\rho$ factors in a unique way as $A\xrightarrow{j}B\xrightarrow{v}E$, and one evidently has $u\circ j=(p\circ v)\circ j$, whence $u=p\circ v$, which proves our assertion.
  \item[\rm(iii)] A finite separable extension $k'$ of a field $k$ is a formally smooth $k$-algebra \sref{0.19.6.1}; we will see later that for a \emph{finite type} extension $K$ of a field $k$ to be a formally \'{e}tale $k$-algebra, it is necessary and sufficient that it be formally unramified, and this is also equivalent to $K$ being a finite and separable extension of $k$ (21.7.4).
  \item[\rm(iv)] A polynomial algebra $A[X_\alpha]_{\alpha\in I}$ over a ring $A$ is \emph{formally smooth} \sref{0.19.3.3} but \emph{not formally \'{e}tale} if $I$ is non-empty, as follows immediately from the definitions.
\end{enumerate}
\end{examples}

\begin{remark}[19.10.4]
\label{0.19.10.4}
The same reasoning as in \sref{0.19.4.3} shows that in order to verify that a topological $A$-algebra $B$ is formally unramified, one can, in the definition \sref{0.19.10.2}, restrict to the case where the kernel of $p:E\to C$ is of \emph{square zero}; one can moreover restrict to the case where $u(B)\subset p(E)$ (for otherwise there is no $v$ factoring $u$) and, replacing $C$ by $u(B)$ and $E$ by $p^{-1}(u(B))$, suppose $u$ and $p$ \emph{surjective}.
\end{remark}
% END NATIVE INPUT ega0/ega0-19
% BEGIN NATIVE INPUT ega0/ega0-20
\section{Derivations and differentials}
\label{section:0.20}

\oldpage[0\textsubscript{IV-1}]{116}
The notions introduced in this section will be taken up again in geometric form in chapter~IV, \S\,16, and will play an important role in the study of preschemes.
Their importance in this chapter stems first from their relations with the notion of a formally smooth algebra, and in particular from the theorems \sref{0.20.4.9}, \sref{0.20.5.7} and \sref{0.20.5.12}, which will be translated into geometric language in the section of chapter~IV devoted to smooth morphisms, and are in constant use in the applications.
On the other hand, the differential notions will serve to prove in \S\,22 important criteria for regularity which will play an essential role in the detailed study of Noetherian local rings made in \S\,7 of chapter~IV.

\subsection{Derivations and algebra extensions}
\label{subsection:0.20.1}

\oldpage[0\textsubscript{IV-1}]{117}
\begin{proposition}[20.1.1]
\label{0.20.1.1}
Let $A$ be a ring (not necessarily commutative), $B$, $E$, $C$ $A$-rings, $p:E\to C$ an $A$-homomorphism whose kernel $\mathfrak{Q}$ is of square zero, $u:B\to C$ an $A$-homomorphism.
Suppose that there exists an $A$-homomorphism $v_0:B\to E$ such that $u$ factors as $B\xrightarrow{v_0}E\xrightarrow{p}C$.
Then the set of $A$-homomorphisms $v:B\to E$ such that $u=p\circ v$ is identical to the set of maps $v_0+D$, where $D:B\to\mathfrak{Q}$ is a map satisfying the two following conditions:
\begin{enumerate}
  \item[\rm(i)] $D$ is a homomorphism of $A$-bimodules.
  \item[\rm(ii)] For all $f$, $g$ in $B$, one has
\end{enumerate}
\begin{equation}
\tag{20.1.1.1}
  \label{0.20.1.1.1}
  D(fg) = f.D(g) + D(f).g.
\end{equation}
\end{proposition}

\begin{proof}
To say that $p(v(f))=p(v_0(f))$ for $f\in B$ means that $D(f)=v(f)-v_0(f)$ belongs to $\mathfrak{Q}$; if one writes $v(fg)=v(f)v(g)$, one obtains the relation \sref{0.20.1.1.1}, $\mathfrak{Q}$ being of square zero, and condition (i) follows from the fact that $v$ and $v_0$ are $A$-homomorphisms.

If $\rho:A\to B$ is the structural homomorphism, one deduces from \sref{0.20.1.1.1} that one has $D(\rho(a)f)=\rho(a)D(f)+D(\rho(a))f$ for all $a\in A$; but one must also have $D(\rho(a)f)=\rho(a)D(f)$ by (i), hence, setting $f=1$, one obtains
\[
  D(\rho(a)) = 0 \qquad\text{for}\quad a\in A;
\]
conversely, if $D$ is null on $\rho(A)$ and satisfies \sref{0.20.1.1.1}, it also satisfies (i).
\end{proof}

\begin{definition}[20.1.2]
\label{0.20.1.2}
Given an $A$-ring $B$ and a $B$-bimodule $L$, we call an \emph{$A$-derivation} of $B$ into $L$ a map $D:B\to L$ satisfying conditions (i) and (ii) of \sref{0.20.1.1}.
\end{definition}

It follows from \sref{0.20.1.1} that the kernel of an $A$-derivation $D:B\to L$ is a sub-$A$-ring of $B$.

One sometimes calls a \emph{derivation} of $B$ into $L$ a $\mathbf{Z}$-derivation, that is, an additive map of $B$ into $L$ satisfying \sref{0.20.1.1.1}; one then has $D(1)=0$ for such a map.
If $B$ is an algebra over a \emph{prime field} $P$, every $\mathbf{Z}$-derivation of $B$ is a $P$-derivation: this is evident from the above if $P$ is of characteristic $>0$, and otherwise, the relation $n\cdot n^{-1}=1$ for $n\in\mathbf{Z}-\{0\}$ gives
\[
  nD(n^{-1}) = 0
\]
hence $D(n^{-1})=0$ since $P$ is of characteristic $0$.

It follows immediately from the definition \sref{0.20.1.2} that if $D$, $D'$ are two $A$-derivations of $B$ into $L$, then so is $D-D'$.
In other words, the set of $A$-derivations of $B$ into $L$ is endowed with a structure of \emph{additive group}; we denote this group $\Der_A(B,L)$.
If $A$ is \emph{commutative}, $B$ a commutative $A$-\emph{algebra}, and $L$ a $B$-module, then, for every $A$-derivation $D$ of $B$ into $L$ and every $b\in B$, $bD$ is again an $A$-derivation of $B$ into $L$, in other words $\Der_A(B,L)$ is then endowed with a structure of $B$-\emph{module}.

\oldpage[0\textsubscript{IV-1}]{118}
Proposition \sref{0.20.1.1} can be interpreted in the following way:

\begin{corollary}[20.1.3]
\label{0.20.1.3}
Given two $A$-homomorphisms of $A$-rings $p:E\to C$, $u:B\to C$ where the first has kernel $\mathfrak{Q}$ of square zero, the set of $A$-homomorphisms $v:B\to E$ such that $u=p\circ v$ is either empty or a principal homogeneous space for the group $\Der_A(B,\mathfrak{Q})$.
\end{corollary}

In particular:

\begin{corollary}[20.1.4]
\label{0.20.1.4}
Let $A$ be a ring, $C$ an $A$-ring, $L$ a $C$-bimodule, $E$ an $A$-extension of $C$ by $L$, $p:E\to C$ the augmentation.
The map which, to every derivation $D\in\Der_A(C,L)$, associates the map $x\mapsto x+D(p(x))$ is an isomorphism of the group $\Der_A(C,L)$ onto the group of $A$-equivalences of $E$ with itself.
\end{corollary}

\begin{proof}
Apply \sref{0.20.1.1} for $B=E$, $v_0=\mathrm{I}_E$: the set of $A$-homomorphisms $v:E\to E$ such that $p\circ v=p$ is identical to the set of maps $v=\mathrm{I}_E+D'$, where $D'\in\Der_A(E,L)$.
To say that such an $A$-homomorphism $v$ is an $A$-equivalence amounts to saying that $v$ restricts to the canonical injection on $L$, or equivalently that $D'(x)=0$ in $L$; but this also means that $D'$ factors as $E\xrightarrow{p}C\xrightarrow{D}L$, where $D$ is an $A$-derivation.
Whence the corollary.
\end{proof}

For trivial extensions, \sref{0.20.1.3} gives:

\begin{corollary}[20.1.5]
\label{0.20.1.5}
Let $A$ be a ring, $B$, $C$ two $A$-rings, $L$ a $C$-bimodule, $u:B\to C$ an $A$-homomorphism; the map which, to every derivation $D\in\Der_A(B,L)$, associates the map $\alpha:x\mapsto(u(x),D(x))$ is a bijection onto the set of $A$-homomorphisms $v:B\to D_C(L)$ \textup{(}cf.\ 18.2.3\textup{)} such that $u$ factors as $B\xrightarrow{v}D_C(L)\to C$.
\end{corollary}

More particularly:

\begin{corollary}[20.1.6]
\label{0.20.1.6}
Let $A$ be a ring, $B$ an $A$-ring, $L$ a $B$-bimodule.
If, to every derivation $D\in\Der_A(B,L)$, one associates: $1^\circ$ the $A$-equivalence $(x,y)\mapsto(x,y+D(x))$ of the extension $D_B(L)$ with itself; $2^\circ$ the $A$-homomorphism $x\mapsto(x,D(x))$ of $B$ into $D_B(L)$, inverse on the right of the augmentation homomorphism $D_B(L)\to B$; one defines canonical bijections between:
\begin{enumerate}
  \item[\rm(i)] the set $\Der_A(B,L)$;
  \item[\rm(ii)] the set of $A$-equivalences of $D_B(L)$ with itself;
  \item[\rm(iii)] the set of $A$-homomorphisms inverse on the right of the augmentation homomorphism $D_B(L)\to B$.
\end{enumerate}
\end{corollary}

We note that the bijection thus established between (i) and (ii) respects the group structures, and that the bijection between (ii) and (iii) is none other than the bijection already defined in (18.3.8).

\begin{corollary}[20.1.7]
\label{0.20.1.7}
Let $A$ be a ring, $C$ an $A$-ring, $L$ a $C$-bimodule, $E$ a trivial $A$-extension of $C$ by $L$, $p:E\to C$ the augmentation.
The set of $A$-derivations $D$ of $E$ into $L$ such that $D(x)=x$ in $L$ is identical to the set of maps $x\mapsto x-w(p(x))$, where $w$ runs through the set of $A$-homomorphisms inverse on the right of $p$.
\end{corollary}

\begin{proof}
It suffices to apply \sref{0.20.1.1} for $B=E$, $v_0=\mathrm{I}_E$; if $v=\mathrm{I}_E-D$, the condition $D(x)=x$ for $x\in L$ is equivalent to $v(x)=0$ for $x\in L$, that is, $v=w\circ p$, where $w:C\to E$ is an $A$-homomorphism; the condition $p\circ v=p$ then gives $p\circ w=\mathrm{I}_C$, i.e., $w$ is inverse on the right of $p$.
\end{proof}

\subsection{Functorial properties of derivations}
\label{subsection:0.20.2}

\oldpage[0\textsubscript{IV-1}]{119}
\begin{env}[20.2.1]
\label{0.20.2.1}
Let $A$ be a ring, $B$ an $A$-ring, $L$ a $B$-bimodule; if $L'$ is a second $B$-bimodule and $w:L\to L'$ a homomorphism of $B$-bimodules, it is clear that the map $D\mapsto w\circ D$ is a homomorphism of additive groups
\begin{equation}
\tag{20.2.1.1}
  \label{0.20.2.1.1}
  w_0:\Der_A(B,L)\to\Der_A(B,L')
\end{equation}
and if $w':L'\to L''$ is a second homomorphism of $B$-bimodules, one has $(w'\circ w)_0=w_0'\circ w_0$.
When $A$ is commutative, $B$ a commutative $A$-algebra and $L$ a $B$-module, \sref{0.20.2.1.1} is a homomorphism of $B$-\emph{modules}.

In the second place, let $B'$ be an $A$-ring, $v:B'\to B$ an $A$-homomorphism which makes $L$ into a $B'$-bimodule; then the map $D\mapsto D\circ v$ is a homomorphism of additive groups,
\begin{equation}
\tag{20.2.1.2}
  \label{0.20.2.1.2}
  v^0:\Der_A(B,L)\to\Der_A(B',L)
\end{equation}
as follows from \sref{0.20.1.1.1}; if $v':B''\to B'$ is a second $A$-homomorphism, one has $(v\circ v')^0=v'^0\circ v^0$.
When $A$, $B$ and $B'$ are commutative and $L$ a $B$-module, \sref{0.20.2.1.2} is a homomorphism of $B$-modules.

Finally, let $u:A'\to A$ be a ring homomorphism making $B$ into an $A'$-ring; every $A$-derivation $D\in\Der_A(B,L)$ is also an $A'$-derivation, whence a canonical injection of groups
\begin{equation}
\tag{20.2.1.3}
  \label{0.20.2.1.3}
  u^0:\Der_A(B,L)\to\Der_{A'}(B,L)
\end{equation}
and if $u':A''\to A'$ is a second ring homomorphism, one has $(u\circ u')^0=u'^0\circ u^0$;
when $A$, $A'$ and $B$ are commutative and $L$ a $B$-module, \sref{0.20.2.1.3} is a homomorphism of $B$-modules.

One can also say that
\[
  (A,B,L)\leadsto\Der_A(B,L)
\]
is a \emph{covariant functor} from the category $\boldsymbol{K}$ defined in (18.3.5) into the category $\boldsymbol{Ab}$ of commutative groups, by associating to every triplet $(u,v,w)$ constituting a morphism of $\boldsymbol{K}$ the homomorphism $w_0\circ v^0\circ u^0$; the verification of the functoriality follows from the commutativity of the diagrams
\[
\xymatrix{
  \Der_A(B,L) \ar[r]^{v^0} \ar[d]_{w_0} & \Der_A(B',L) \ar[d]^{w_0} \\
  \Der_A(B,L') \ar[r]^{v^0} & \Der_A(B',L')
}
\qquad
\xymatrix{
  \Der_A(B,L) \ar[r]^{u^0} \ar[d]_{w_0} & \Der_{A'}(B,L) \ar[d]^{w_0} \\
  \Der_A(B,L') \ar[r]^{u^0} & \Der_{A'}(B,L')
}
\]
for every homomorphism $w:L\to L'$ of $B$-bimodules.
\end{env}

\oldpage[0\textsubscript{IV-1}]{120}
\begin{theorem}[20.2.2]
\label{0.20.2.2}
Let $u:A\to B$, $v:B\to C$ be two ring homomorphisms, $L$ a $C$-bimodule.
There is a canonical exact sequence of commutative groups
\begin{align}
\tag{20.2.2.1}
  \label{0.20.2.2.1}
  0\to\Der_B(C,L)\xrightarrow{u^0}\Der_A(C,L)&\xrightarrow{v^0}\Der_A(B,L)\xrightarrow{\partial} \\
  &\to\operatorname{Exan}_B(C,L)\xrightarrow{u^1}\operatorname{Exan}_A(C,L)\xrightarrow{v^1}\operatorname{Exan}_A(B,L) \notag
\end{align}
where $u^0$, $v^0$ are the homomorphisms \textup{(}20.2.1.3\textup{)} and \textup{(}20.2.1.2\textup{)} respectively, $u^1$, $v^1$ are the homomorphisms defined in (18.3.4.1) and (18.3.3.1) respectively, and where $\partial$ is defined as follows: for every $A$-derivation $D$ of $B$ into $L$, $\partial(D)$ is the class of the $B$-extension of $C$ by $L$ defined by the $A$-homomorphism $\alpha:x\mapsto(v(x),D(x))$ \textup{(}cf.\ \sref{0.20.1.5}\textup{)}.
Moreover, the exact sequence \textup{(}20.2.2.1\textup{)} is functorial in $L$ \textup{(}for the homomorphisms defined in \textup{(}20.2.1.1\textup{)} and \textup{(}18.3.1.1\textup{)} respectively\textup{)}.
\end{theorem}

\begin{proof}
As $D$ is an $A$-derivation (and \emph{a fortiori} a $\mathbf{Z}$-derivation) of $B$ into $L$, the $A$-homomorphism $\alpha:x\mapsto(v(x),D(x))$ indeed defines on $D_C(L)$ a structure of $B$-extension, hence $\partial(D)$ is well defined \sref{0.20.1.5}.
We must verify the exactness at the $5$ places:

\medskip
\noindent 1)\enspace
The exactness at $\Der_B(C,L)$ is trivial (cf.\ \sref{0.20.2.1}).

\medskip
\noindent 2)\enspace
By definition \sref{0.20.2.1} the kernel of $v^0$ is the set of $A$-derivations of $C$ which vanish on $v(B)$, that is, those which are also $B$-derivations; whence the exactness at $\Der_A(C,L)$.

\medskip
\noindent 3)\enspace
The kernel of $\partial$ is formed of the derivations $D\in\Der_A(B,L)$ for which the $B$-extension defined by $\alpha:x\mapsto(v(x),D(x))$ is $B$-trivial; this means \sref{0.18.2.3} that there exists a $B$-homomorphism $z\mapsto(z,w(z))$ of $C$ into $D_C(L)$ (the structure of $B$-ring of $D_C(L)$ being defined by $\alpha$); but such a homomorphism, being \emph{a fortiori} an $A$-homomorphism, is of the form $z\mapsto(z,D'(z))$ where $D'\in\Der_A(C,L)$ \sref{0.20.1.6}; and writing that this is a $B$-homomorphism, we get
\[
  D(x)z+v(x)D'(z) = D'(v(x)z) = v(x)D'(z)+D'(v(x))z
\]
for $x\in B$, $z\in C$, which gives $D'(v(x))=D(x)$; the kernel of $\partial$ is therefore the image of $v^0$.

\medskip
\noindent 4)\enspace
The kernel of $u^1$ is the set of classes of $B$-extensions of $C$ by $L$ which are $A$-trivial, hence (up to equivalence) of the form $D_C(L)$, where the structure of $A$-ring of $D_C(L)$ is defined by the homomorphism $t\mapsto(v(u(t)),0)$.
Now, every structure of $B$-extension on $D_C(L)$ is defined by a homomorphism $\alpha:x\mapsto(v(x),D(x))$, where $D$ is a $\mathbf{Z}$-derivation of $B$ into $L$ \sref{0.20.1.5}; to say that the structure of $A$-ring of this $B$-extension is deduced from its structure of $B$-ring by means of $u:A\to B$ means that $D(u(t))=0$ for $t\in A$, hence that $D$ is an $A$-derivation, or equivalently that the class of the $B$-extension considered is of the form $\partial(D)$; whence the exactness at $\operatorname{Exan}_B(C,L)$.

\medskip
\noindent 5)\enspace
The kernel of $v^1$ is the set of classes of $A$-extensions $E$ of $C$ by $L$ which trivialise on $v(B)$, that is, for which there is an $A$-homomorphism $w:B\to E$ such that $v$ factors as $B\xrightarrow{w}E\to C$; now such an $A$-homomorphism defines on $E$ a structure of $B$-extension whose class has as its image under $u^1$ the class of the given $A$-extension; the converse being trivially true, this proves the exactness at $\operatorname{Exan}_A(C,L)$.

Finally, the functoriality in $L$ follows trivially from the definitions.
\end{proof}

\oldpage[0\textsubscript{IV-1}]{121}
\begin{corollary}[20.2.3]
\label{0.20.2.3}
Let $A$, $B$, $C$ be three commutative rings, $u:A\to B$, $v:B\to C$ two ring homomorphisms, $L$ a $C$-module.
There is a canonical exact sequence of $A$-modules
\begin{align}
\tag{20.2.3.1}
  \label{0.20.2.3.1}
  0\to\Der_B(C,L)\xrightarrow{u^0}\Der_A(C,L)&\xrightarrow{v^0}\Der_A(B,L)\xrightarrow{\partial} \\
  &\to\operatorname{Exalcom}_B(C,L)\xrightarrow{u^1}\operatorname{Exalcom}_A(C,L)\xrightarrow{v^1}\operatorname{Exalcom}_A(B,L) \notag
\end{align}
functorial in $L$.
\end{corollary}

\begin{proof}
The reasoning is the same as in \sref{0.20.2.2} once one has verified that for every derivation $D\in\Der_A(B,L)$, $\partial(D)$ is indeed the class of a $B$-extension of $C$ by $L$ which is a commutative $B$-algebra; but this follows immediately from the commutativity of $C$ and from the fact that $L$ is a $C$-module.
\end{proof}

\begin{corollary}[20.2.4]
\label{0.20.2.4}
Under the hypotheses of \sref{0.20.2.2} \textup{(}resp.\ \sref{0.20.2.3}\textup{)}, there is a canonical exact sequence, functorial in $L$,
\begin{equation}
\tag{20.2.4.1}
  \label{0.20.2.4.1}
  0\to\Der_B(C,L)\xrightarrow{u^0}\Der_A(C,L)\xrightarrow{v^0}\Der_A(B,L)\xrightarrow{\partial}\operatorname{Exan}_{B/A}(C,L)\to 0
\end{equation}
\textup{(}resp.
\begin{equation}
\tag{20.2.4.2}
  \label{0.20.2.4.2}
  0\to\Der_B(C,L)\xrightarrow{u^0}\Der_A(C,L)\xrightarrow{v^0}\Der_A(B,L)\xrightarrow{\partial}\operatorname{Exalcom}_{B/A}(C,L)\to 0\text{).}
\end{equation}
This follows from the definition of $\operatorname{Exan}_{B/A}(C,L)$ \textup{(}resp.\ $\operatorname{Exalcom}_{B/A}(C,L)$\textup{)} \textup{(}(18.3.7) and (18.4.2)\textup{)}.
\end{corollary}

\begin{remark}[20.2.5]
\label{0.20.2.5}
Suppose that one has a commutative diagram of ring homomorphisms
\[
\xymatrix{
  A \ar[r] & B \ar[r] & C \\
  A' \ar[u] \ar[r] & B' \ar[u] \ar[r] & C' \ar[u]
}
\]
Then there is a commutative diagram
\[
\xymatrix@C=0.5em{
  0 \ar[r] & \Der_B(C,L) \ar[r] \ar[d] & \Der_A(C,L) \ar[r] \ar[d] & \Der_A(B,L) \ar[r]^-\partial \ar[d] & \operatorname{Exan}_B(C,L) \ar[r] \ar[d] & \operatorname{Exan}_A(C,L) \ar[r] \ar[d] & \operatorname{Exan}_A(B,L) \ar[d] \\
  0 \ar[r] & \Der_{B'}(C',L) \ar[r] & \Der_{A'}(C',L) \ar[r] & \Der_{A'}(B',L) \ar[r]^-\partial & \operatorname{Exan}_{B'}(C',L) \ar[r] & \operatorname{Exan}_{A'}(C',L) \ar[r] & \operatorname{Exan}_{A'}(B',L)
}
\]
and similarly for the exact sequences \sref{0.20.2.3.1}, \sref{0.20.2.4.1} and \sref{0.20.2.4.2}.
\end{remark}

\subsection{Continuous derivations in topological rings}
\label{subsection:0.20.3}

\begin{env}[20.3.1]
\label{0.20.3.1}
Given two topological rings $A$, $B$ (linearly topologised, as always), we will denote by $\operatorname{Hom.\ cont}(A,B)$ the set of continuous homomorphisms from $A$ to $B$.
Given a topological ring $A$, the category of topological $A$-rings is defined as that of $A$-rings \sref{0.18.1.4} by replacing throughout ``ring'' by ``topological ring'' and ``homomorphism'' by ``continuous homomorphism''; if $B$

\oldpage[0\textsubscript{IV-1}]{122}
and $C$ are two topological $A$-rings, we denote $\operatorname{Hom.\ cont}_A(B,C)$ the set of continuous $A$-homomorphisms from $B$ to $C$.

Let $A$ be a topological ring, $B$ a topological $A$-ring, $L$ a topological $B$-bimodule; we denote $\operatorname{D\acute{e}r.\ cont}_A(B,L)$ the set of \emph{continuous} $A$-derivations of $B$ into $L$; it is clear that this is a \emph{subgroup} of $\Der_A(B,L)$ (and a sub-$A$-module when $A$ and $B$ are commutative and $L$ a $B$-module).
\end{env}

\begin{env}[20.3.2]
\label{0.20.3.2}
It is immediate that the proposition \sref{0.20.1.1} remains valid when one replaces ``ring'' by ``topological ring'' and ``homomorphism'' by ``continuous homomorphism''.
Similarly, \sref{0.20.1.3} and \sref{0.20.1.4} remain valid by replacing $\Der$ by $\operatorname{D\acute{e}r.\ cont}$, ``ring'' (resp.\ ``bimodule'') by ``topological ring'' (resp.\ ``topological bimodule''), ``homomorphism'' by ``continuous homomorphism''; one must naturally suppose in \sref{0.20.1.4} that $p$ is continuous, and replace ``$A$-equivalences'' by ``continuous $A$-equivalences''.
One has the same results for \sref{0.20.1.5} and \sref{0.20.1.6}, provided one takes as topology on $D_C(L)$ (resp.\ $D_B(L)$) the product topology (on $C\times L$, resp.\ $B\times L$).
\end{env}

\begin{proposition}[20.3.3]
\label{0.20.3.3}
Let $A$ be a topological ring, $B$ a topological $A$-ring, $L$ a discrete topological $B$-bimodule annihilated by an open bilateral ideal of $B$.
If in $B$ the square of every open bilateral ideal is open, one has $\operatorname{D\acute{e}r.\ cont}_A(B,L)=\Der_A(B,L)$.
\end{proposition}

\begin{proof}
Indeed, if $\mathfrak{R}$ is an open bilateral ideal of $B$ annihilating $L$, and $D$ an $A$-derivation of $B$ into $L$, one has $D(\mathfrak{R}^2)\subset\mathfrak{R}.D(\mathfrak{R})+D(\mathfrak{R}).\mathfrak{R}=0$ \sref{0.20.1.1.1}, hence $D$ is continuous.
\end{proof}

\begin{env}[20.3.4]
\label{0.20.3.4}
All the results of \sref{0.20.2.1} remain valid when one replaces ``ring'' by ``topological ring'', ``bimodule'' by ``topological bimodule'', ``homomorphism'' by ``continuous homomorphism'' and $\Der$ by $\operatorname{D\acute{e}r.\ cont}$.
\end{env}

\begin{proposition}[20.3.5]
\label{0.20.3.5}
Let $A$ be a topological ring, $B$ a topological $A$-ring, $L$ a discrete topological $B$-bimodule annihilated by an open bilateral ideal of $B$.
There is then a canonical isomorphism
\begin{equation}
\tag{20.3.5.1}
  \label{0.20.3.5.1}
  \varinjlim\Der_{A/\mathfrak{J}}(B/\mathfrak{R},L)\xrightarrow{\sim}\operatorname{D\acute{e}r.\ cont}_A(B,L)
\end{equation}
where in the first member the inductive limit is taken over the ordered filtered set of pairs $(\mathfrak{J},\mathfrak{R})$ of bilateral ideals such that $\mathfrak{R}.L=L.\mathfrak{R}=0$, $\mathfrak{J}.B\subset\mathfrak{R}$, $B.\mathfrak{J}\subset\mathfrak{R}$.
\end{proposition}

\begin{proof}
As $A/\mathfrak{J}$ and $B/\mathfrak{R}$ are discrete, one has canonical homomorphisms $w_{\mathfrak{J},\mathfrak{R}}:\Der_{A/\mathfrak{J}}(B/\mathfrak{R},L)\to\operatorname{D\acute{e}r.\ cont}_A(B,L)$ forming an inductive system \sref{0.20.3.4}, whence the homomorphism \sref{0.20.3.5.1} by passage to the inductive limit.
As the homomorphism $B/\mathfrak{R}'\to B/\mathfrak{R}$ is surjective for $\mathfrak{R}\supset\mathfrak{R}'$, it follows immediately from the definition that the homomorphism $\Der_{A/\mathfrak{J}}(B/\mathfrak{R},L)\to\Der_{A/\mathfrak{J}}(B/\mathfrak{R}',L)$ (with $\mathfrak{R}.L=L.\mathfrak{R}=0$, $\mathfrak{J}.B\subset\mathfrak{R}'$, $B.\mathfrak{J}\subset\mathfrak{R}'$) is injective, and it follows similarly that the homomorphism $\Der_{A/\mathfrak{J}}(B/\mathfrak{R},L)\to\Der_{A/\mathfrak{J}'}(B/\mathfrak{R},L)$ for $\mathfrak{J}'\subset\mathfrak{J}$ \sref{0.20.2.1} is injective; one concludes that the homomorphism \sref{0.20.3.5.1} is injective.
On the other hand, if $D$ is a continuous $A$-derivation of $B$ into $L$, its kernel contains an open bilateral ideal $\mathfrak{R}_0$ of $B$, and if $\mathfrak{J}_0$ is an open bilateral ideal of $A$ such that $\mathfrak{J}_0 B\subset\mathfrak{R}_0$ and $B\mathfrak{J}_0\subset\mathfrak{R}_0$, it is clear that $D$ is the canonical image of an $(A/\mathfrak{J}_0)$-derivation of $B/\mathfrak{R}_0$ into $L$, hence \sref{0.20.3.5.1} is surjective.
\end{proof}

\oldpage[0\textsubscript{IV-1}]{123}
\begin{proposition}[20.3.6]
\label{0.20.3.6}
Let $u:A\to B$, $v:B\to C$ be two continuous homomorphisms of topological rings, $L$ a discrete $C$-bimodule annihilated by an open bilateral ideal of $C$.
There is a canonical exact sequence
\begin{align}
\tag{20.3.6.1}
  \label{0.20.3.6.1}
  0\to\operatorname{D\acute{e}r.\ cont}_B(C,L)&\xrightarrow{u^0}\operatorname{D\acute{e}r.\ cont}_A(C,L)\xrightarrow{v^0}\operatorname{D\acute{e}r.\ cont}_A(B,L)\xrightarrow{\partial} \\
  &\to\operatorname{Exantop}_B(C,L)\xrightarrow{u^1}\operatorname{Exantop}_A(C,L)\xrightarrow{v^1}\operatorname{Exantop}_A(B,L) \notag
\end{align}
where $\partial$ is defined by passage to the inductive limit from the homomorphism $\partial$ of \sref{0.20.2.2.1}; this exact sequence is functorial in $L$ (in the category of discrete $C$-bimodules annihilated by open bilateral ideals).
\end{proposition}

\begin{proof}
This follows from the exactness of the functor $\varinjlim$, from \sref{0.20.2.2}.
\end{proof}

\begin{corollary}[20.3.7]
\label{0.20.3.7}
Let $A$, $B$, $C$ be three commutative topological rings, $u:A\to B$, $v:B\to C$ two continuous homomorphisms, $L$ a discrete $C$-module annihilated by an open ideal of $C$.
There is a canonical exact sequence of $A$-modules, functorial in $L$,
\begin{align}
\tag{20.3.7.1}
  \label{0.20.3.7.1}
  0\to\operatorname{D\acute{e}r.\ cont}_B(C,L)&\xrightarrow{u^0}\operatorname{D\acute{e}r.\ cont}_A(C,L)\xrightarrow{v^0}\operatorname{D\acute{e}r.\ cont}_A(B,L)\xrightarrow{\partial} \\
  &\to\operatorname{Exalcotop}_B(C,L)\xrightarrow{u^1}\operatorname{Exalcotop}_A(C,L)\xrightarrow{v^1}\operatorname{Exalcotop}_A(B,L). \notag
\end{align}
\end{corollary}

\begin{corollary}[20.3.8]
\label{0.20.3.8}
Under the hypotheses of \sref{0.20.3.6} \textup{(}resp.\ \sref{0.20.3.7}\textup{)} there is a canonical exact sequence, functorial in $L$
\begin{equation}
\tag{20.3.8.1}
  \label{0.20.3.8.1}
  0\to\operatorname{D\acute{e}r.\ cont}_B(C,L)\xrightarrow{u^0}\operatorname{D\acute{e}r.\ cont}_A(C,L)\xrightarrow{v^0}\operatorname{D\acute{e}r.\ cont}_A(B,L)\to\operatorname{Exantop}_{B/A}(C,L)\to 0
\end{equation}
\textup{(}resp.
\begin{equation}
\tag{20.3.8.2}
  \label{0.20.3.8.2}
  0\to\operatorname{D\acute{e}r.\ cont}_B(C,L)\xrightarrow{u^0}\operatorname{D\acute{e}r.\ cont}_A(C,L)\xrightarrow{v^0}\operatorname{D\acute{e}r.\ cont}_A(B,L)\to\operatorname{Exalcotop}_{B/A}(C,L)\to 0\text{).}
\end{equation}
\end{corollary}

We leave to the reader the task of writing the diagrams analogous to those of \sref{0.20.2.5}.

\subsection{Principal parts and differentials}
\label{subsection:0.20.4}

In the remainder of this section and in the three following sections, all rings are assumed \emph{commutative}.

\begin{env}[20.4.1]
\label{0.20.4.1}
Let $A$ be a topological ring, $B$ a topological $A$-algebra; the $A$-algebra $B\otimes_A B$ will be equipped with the \emph{tensor product} topology, which in fact makes it into a topological $A$-algebra; we denote by $p$ (or $p_{B/A}$) the surjective continuous $A$-homomorphism
\begin{equation}
\tag{20.4.1.1}
  \label{0.20.4.1.1}
  p:B\otimes_A B\to B
\end{equation}
such that $p(b\otimes b')=bb'$; it is immediate that $p$ is continuous.
The \emph{kernel} of $p$ will be denoted $\mathfrak{J}_{B/A}$ (or simply $\mathfrak{J}$ if there is no risk of confusion).
We denote by $j_1:B\to B\otimes_A B$ and $j_2:B\to B\otimes_A B$ the two canonical $A$-homomorphisms, such that
\[
  j_1(b) = b\otimes 1, \qquad j_2(b) = 1\otimes b
\]
which are continuous.
\end{env}

\oldpage[0\textsubscript{IV-1}]{124}
\begin{definition}[20.4.2]
\label{0.20.4.2}
The \emph{augmented $B$-algebra of principal parts of order $1$} of $B$ over $A$ is the topological $A$-algebra quotient
\begin{equation}
\tag{20.4.2.1}
  \label{0.20.4.2.1}
  \mathrm{P}^1_{B/A} = (B\otimes_A B)/\mathfrak{J}^2
\end{equation}
equipped with the structure of $B$-algebra defined by the homomorphism $\bar{j}_1:B\to\mathrm{P}^1_{B/A}$ (deduced from $j_1$ by composition with the canonical map $B\otimes_A B\to\mathrm{P}^1_{B/A}$), and with the augmentation of $B$-algebra $\varepsilon:\mathrm{P}^1_{B/A}\to B$ (also denoted $\varepsilon_{B/A}$) deduced from $p$ by passage to the quotient.
As $p(b\otimes 1)=b$ by definition, it is clear that $\varepsilon$ is indeed an augmentation of $B$-algebra.
\end{definition}

\begin{definition}[20.4.3]
\label{0.20.4.3}
The kernel of the augmentation $\varepsilon:\mathrm{P}^1_{B/A}\to B$,
\begin{equation}
\tag{20.4.3.1}
  \label{0.20.4.3.1}
  \Omega^1_{B/A} = \mathfrak{J}_{B/A}/(\mathfrak{J}_{B/A})^2
\end{equation}
equipped with the topology induced by that of $\mathrm{P}^1_{B/A}$, which in fact makes it into a topological $B$-module, is called the $B$-module of \emph{$1$-differentials} (or simply \emph{differentials}) of $B$ over $A$.
\end{definition}

We note that the topology of $\Omega^1_{B/A}$ is also the quotient topology of the topology induced on $\mathfrak{J}_{B/A}$ by that of $B\otimes_A B$ (Bourbaki, \emph{Top.\ g\'{e}n.}, chap.~III, 3$^\text{e}$ \'{e}d., \S\,2, n$^\circ$\,7, prop.~20).
If $B$ is discrete it follows that so is $\Omega^1_{B/A}$.
We denote $\widehat{\Omega}^1_{B/A}$ the \emph{separated completion} of the topological $B$-module $\Omega^1_{B/A}$.

Every topological ring $B$ can be considered as a topological $\mathbf{Z}$-algebra (where $\mathbf{Z}$ is equipped with the discrete topology), so one can define the topological $B$-module $\Omega^1_{B/\mathbf{Z}}$, which is sometimes also called the $B$-module of \emph{absolute differentials} over $B$ and which one denotes $\Omega^1_B$.
If $B$ is a topological algebra over a prime field $P$ (discrete), one has $B\otimes_P B=B\otimes_{\mathbf{Z}}B$, hence $\Omega^1_{B/P}=\Omega^1_B$.

\begin{lemma}[20.4.4]
\label{0.20.4.4}
Let $A$ be a ring, $B$ an $A$-algebra.
The ideal $\mathfrak{J}_{B/A}$ of $B\otimes_A B$ is generated by the elements $1\otimes s-s\otimes 1$, where $s$ runs through a set of generators of the $A$-algebra $B$.
\end{lemma}

\begin{proof}
It is clear that for every $x\in B$, one has $x\otimes 1-1\otimes x\in\mathfrak{J}$; on the other hand, for all $x$, $y$ in $B$, one has $x\otimes y=xy\otimes 1+(x\otimes 1)(1\otimes y-y\otimes 1)$.
If $\sum_i (x_i\otimes y_i)\in\mathfrak{J}$, one has by definition $\sum_i x_i y_i=0$, hence
\begin{equation}
\tag{20.4.4.1}
  \label{0.20.4.4.1}
  \sum_i (x_i\otimes y_i) = \sum_i (x_i\otimes 1)(1\otimes y_i-y_i\otimes 1)
\end{equation}
which proves that $\mathfrak{J}$ is generated by the elements $1\otimes x-x\otimes 1$.
Moreover, if $x=st$, one has
\begin{equation}
\tag{20.4.4.2}
  \label{0.20.4.4.2}
  x\otimes 1-1\otimes x = (s\otimes 1)(t\otimes 1-1\otimes t)+(s\otimes 1-1\otimes s)(1\otimes t)
\end{equation}
which immediately completes the proof by induction.
\end{proof}

\begin{proposition}[20.4.5]
\label{0.20.4.5}
Let $A$ be a topological ring, $B$ a topological $A$-algebra.
The topology of $\Omega^1_{B/A}$ is less fine than the topology deduced from that of $B$ \sref{0.19.0.2}; if in $B$ the square of every open ideal is open, these two topologies are identical.
\end{proposition}

\begin{proof}
The first assertion is trivial, the tensor product topology on $B\otimes_A B$ being less fine than the topology deduced from that of $B$; \emph{a fortiori} the topology induced on $\mathfrak{J}=\mathfrak{J}_{B/A}$

\oldpage[0\textsubscript{IV-1}]{125}
by that of $B\otimes_A B$ is less fine than the topology on $\mathfrak{J}$ deduced from that of $B$.
To prove the second assertion, write, by abuse of notation, $M\otimes N$ for the submodule $\operatorname{Im}(M\otimes_A N)$ for two sub-$A$-modules $M$, $N$ of $B$.
Using the relation
\[
  \begin{aligned}
  (xy)\otimes z-x\otimes(yz)
  &=(x\otimes 1)(1\otimes z)(y\otimes 1-1\otimes y)\\
  &=xz\cdot(y\otimes 1-1\otimes y)\\
  &\quad-x\cdot(z\otimes 1-1\otimes z)(y\otimes 1-1\otimes y)
  \end{aligned}
\]
in the $B$-module $B\otimes_A B$ (defined by $j_1$), and taking into account \sref{0.20.4.4}, one sees immediately that, if $\mathfrak{R}$ is an ideal of $B$, one has
\[
  ((\mathfrak{R}^2\otimes B)+(B\otimes\mathfrak{R}^2))\cap\mathfrak{J}
  \subset(\mathfrak{R}\otimes\mathfrak{R})\cap\mathfrak{J}
  +\mathfrak{R}\mathfrak{J}+\mathfrak{J}^2.
\]
On the other hand, if $x_i$, $y_i$ are elements of $\mathfrak{R}$ such that
$\sum_i(x_i\otimes y_i)\in\mathfrak{J}$, it follows from
\sref{0.20.4.4.1} that $\sum_i(x_i\otimes y_i)\in\mathfrak{R}\mathfrak{J}$,
and hence finally
\begin{equation}
\tag{20.4.5.1}
  \label{0.20.4.5.1}
  (\mathfrak{R}^2\otimes B+B\otimes\mathfrak{R}^2)\cap\mathfrak{J}\subset\mathfrak{R}\mathfrak{J}+\mathfrak{J}^2.
\end{equation}
Now, there is a fundamental system of neighbourhoods of $0$ in $\mathfrak{J}$ (for the topology induced by that of $B\otimes_A B$) obtained by taking for neighbourhoods of $0$ the sets $(\mathfrak{R}\otimes B+B\otimes\mathfrak{R})\cap\mathfrak{J}$, where $\mathfrak{R}$ runs through the set of open ideals.
As the topology of $\Omega^1_{B/A}$ deduced from that of $B$ is also the quotient by $\mathfrak{J}^2$ of the topology of $\mathfrak{J}$ deduced from that of $B$, the hypothesis on the open ideals of $B$ and the relation \sref{0.20.4.5.1} complete the proof of the second assertion.
\end{proof}

\begin{definition}[20.4.6]
\label{0.20.4.6}
Let $p_1$ and $p_2$ be the composite maps $B\xrightarrow{j_1}B\otimes_A B\to\mathrm{P}^1_{B/A}$, $B\xrightarrow{j_2}B\otimes_A B\to\mathrm{P}^1_{B/A}$, which are continuous $A$-homomorphisms such that $\varepsilon\circ p_1=\varepsilon\circ p_2=\mathrm{I}_B$.
The continuous $A$-module homomorphism
\begin{equation}
\tag{20.4.6.1}
  \label{0.20.4.6.1}
  d_{B/A} = p_1-p_2:B\to\Omega^1_{B/A}
\end{equation}
is called the \emph{exterior differential} of $B$ relative to $A$; for every $x\in B$, $d_{B/A}(x)$ (also denoted $d(x)$ or $dx$) is called the \emph{differential} of $x$ (relative to $A$).
\end{definition}

When $A=\mathbf{Z}$, one writes $d_B$ instead of $d_{B/\mathbf{Z}}$; if $B$ is an algebra over a prime field $P$, one has $d_B=d_{B/P}$.

\begin{proposition}[20.4.7]
\label{0.20.4.7}
The $B$-module $\Omega^1_{B/A}$ is generated by the elements $d_{B/A}(x)$, where $x$ runs through a system of generators of the $A$-algebra $B$.
\end{proposition}

\begin{proof}
As $d_{B/A}(x)$ is the canonical image of $x\otimes 1-1\otimes x$ in $\Omega^1_{B/A}$, the proposition is an immediate consequence of \sref{0.20.4.4}.
\end{proof}

\begin{theorem}[20.4.8]
\label{0.20.4.8}
Let $A$ be a topological ring, $B$ a topological $A$-algebra.
\begin{enumerate}
  \item[\rm(i)] There exists a unique isomorphism of augmented topological $B$-algebras
\begin{equation}
\tag{20.4.8.1}
  \label{0.20.4.8.1}
  \varphi:\mathrm{P}^1_{B/A}\xrightarrow{\sim}D_B(\Omega^1_{B/A})
\end{equation}
which reduces to the identity on $\Omega^1_{B/A}$.
  \item[\rm(ii)] The homomorphism $d_{B/A}$ is an $A$-derivation of $B$ into $\Omega^1_{B/A}$, having the following universal property: for every topological $B$-module $L$, the map $u\mapsto u\circ d_{B/A}$ is an isomorphism of $A$-modules
\begin{equation}
\tag{20.4.8.2}
  \label{0.20.4.8.2}
  \operatorname{Hom.\ cont}_B(\Omega^1_{B/A},L)\xrightarrow{\sim}\operatorname{D\acute{e}r.\ cont}_A(B,L).
\end{equation}
\end{enumerate}
\end{theorem}

\begin{proof}
\oldpage[0\textsubscript{IV-1}]{126}
(i) It is immediate that $\varphi$ (using the notation of \sref{0.20.4.6}) is necessarily the map $z\mapsto(\varepsilon(z),z-p_1(\varepsilon(z)))$, and the inverse isomorphism the map $(b,x)\mapsto p_1(b)+x$; these two maps being continuous, this proves the first assertion.
We note that the topology of $B$ identifies (by $p_1$) with the \emph{quotient} of the topology of $B\otimes_A B$ by $\mathfrak{J}_{B/A}$.

(ii) The fact that $d_{B/A}$ is an $A$-derivation of $B$ follows from the definition \sref{0.20.4.6} and from \sref{0.20.1.1}.

  To prove the universal property, recall that $\operatorname{D\acute{e}r.\ cont}_A(B,L)$ identifies canonically with the set $G$ of continuous $A$-algebra homomorphisms $u:B\to D_B(L)$ such that the composite $B\xrightarrow{u}D_B(L)\xrightarrow{q}B$ is the identity (where $q(b,x)=b$ is the projection; cf.\ \sref{0.20.1.6} and \sref{0.20.3.2}).
On the other hand, by virtue of the isomorphism $\varphi$, $\operatorname{Hom.\ cont}_B(\Omega^1_{B/A},L)$ identifies canonically with the set of continuous $B$-algebra homomorphisms $v:\mathrm{P}^1_{B/A}\to D_B(L)$ such that the composite $\mathrm{P}^1_{B/A}\xrightarrow{v}D_B(L)\xrightarrow{q}B$ is the augmentation $\varepsilon$.
As $p_2=p_1-d_{B/A}$ by definition, it suffices to prove that every $u\in G$ factors as
\[
\xymatrix{
  B \ar[r]^u \ar[d]_{p_2} & D_B(L) \\
  \mathrm{P}^1_{B/A} \ar[ur]_v &
}
\]
where $v$ is a continuous $B$-homomorphism.
Now, one already has a continuous $A$-algebra homomorphism $j:b\mapsto(b,0)$ from $B$ into $D_B(L)$, which belongs to $G$; by definition of the topological tensor product of topological algebras ($\mathbf{0_I}$, 7.7.6), there therefore exists a continuous $A$-algebra homomorphism $w:B\otimes_A B\to D_B(L)$ making the diagram
\[
\xymatrix{
  B\otimes_A B \ar[dr]^w & B \ar[l]_{j_1} \ar[d]^j \\
  B \ar[u]^{j_2} \ar[r]_u & D_B(L)
}
\]
commutative.
\mbox{One} therefore has by definition $w(b\otimes 1-1\otimes b)=j(b)-u(b)\in L$, and by virtue of \sref{0.20.4.4}, this implies $w(\mathfrak{J})\subset L$, hence $w(\mathfrak{J}^2)=0$; $w$ therefore factors as $B\otimes_A B\to\mathrm{P}^1_{B/A}\xrightarrow{v}D_B(L)$, where $v$ is a continuous $A$-algebra homomorphism; moreover, as $v\circ p_1=j$ is a homomorphism of $B$-algebras, $v$ is a $B$-algebra homomorphism by the definition of the $B$-algebra structure of $\mathrm{P}^1_{B/A}$; as one has by definition $u=v\circ p_2$, this completes the proof.
\end{proof}

\begin{theorem}[20.4.9]
\label{0.20.4.9}
Suppose that $B$ is a topological $A$-algebra formally smooth.
Then the topological $B$-module $\Omega^1_{B/A}$ is \emph{formally projective}.
\end{theorem}

\begin{proof}
Indeed, the hypothesis implies that $B\otimes_A B$ (equipped with the $B$-algebra structure defined by $j_1$) is a topological $B$-algebra formally smooth \sref{0.19.3.5}, (iii), and therefore also a topological $A$-algebra formally smooth \sref{0.19.3.5}, (ii); as $B$ is topologically isomorphic to the $A$-algebra quotient $(B\otimes_A B)/\mathfrak{J}$ and is a formally smooth $A$-algebra, the conclusion follows from \sref{0.19.5.3}.
\end{proof}

\oldpage[0\textsubscript{IV-1}]{127}
\begin{corollary}[20.4.10]
\label{0.20.4.10}
Suppose moreover that in $B$ the square of every open ideal is open; then, for every open ideal $\mathfrak{R}$ of $B$, $\Omega^1_{B/A}\otimes_B(B/\mathfrak{R})=\Omega^1_{B/A}/\mathfrak{R}.\Omega^1_{B/A}$ is a projective $(B/\mathfrak{R})$-module.
\end{corollary}

\begin{proof}
Indeed, the topology of $\Omega^1_{B/A}$ is then deduced from that of $B$ \sref{0.20.4.5}, and it suffices to apply \sref{0.19.2.4}.
\end{proof}

\begin{corollary}[20.4.11]
\label{0.20.4.11}
Let $A$, $B$ be two local Noetherian rings, $\rho:A\to B$ a local homomorphism, making $B$ a topological $A$-algebra formally smooth \textup{(}for the preadic topologies\textup{)}.
Then, for every ideal of definition $\mathfrak{b}$ of $B$, $\Omega^1_{B/A}\otimes_B(B/\mathfrak{b})$ is a \emph{free} $(B/\mathfrak{b})$-module.
\end{corollary}

\begin{proof}
Indeed, it follows from \sref{0.20.4.10} that this module is projective, and as $B/\mathfrak{b}$ is an Artinian ring, every projective $(B/\mathfrak{b})$-module is free ($\mathbf{0_{III}}$, 10.1.3).
\end{proof}

\begin{proposition}[20.4.12]
\label{0.20.4.12}
Let $A$ be a topological ring, $B$ a topological $A$-algebra.
If the structural homomorphism $A\to B$ is surjective, one has $\Omega^1_{B/A}=0$.
\end{proposition}

\begin{proof}
Indeed, one has $\Der_A(B,L)=0$ for every $B$-module $L$ by virtue of \sref{0.20.1.1}, and the proposition follows immediately from \sref{0.20.4.8}.
\end{proof}

\begin{examples}[20.4.13]
\label{0.20.4.13}
\begin{enumerate}
  \item[\rm(i)] Let $A$ be a ring, $B=A[X_\alpha]_{\alpha\in I}$ a polynomial algebra over $A$.
Then $\Omega^1_{B/A}$ is a free $B$-module, of which the $dX_\alpha$ form a \emph{basis}.
Indeed, the $dX_\alpha$ generate this $B$-module \sref{0.20.4.7}.
On the other hand, if $L$ is a free $B$-module with basis $(z_\alpha)_{\alpha\in I}$ whose index set $I$ is the same, there exists an $A$-homomorphism $u$ of $B$ into $D_B(L)$ such that $u(X_\alpha)=(X_\alpha,z_\alpha)$ for every $\alpha\in I$, hence \sref{0.20.1.5} a derivation $D$ of $B$ into $L$ such that $D(X_\alpha)=z_\alpha$ for every $\alpha$; by virtue of \sref{0.20.4.8.1}, this proves that the $dX_\alpha$ are linearly independent.
  \item[\rm(ii)] Let $A$ be a ring, $L$ an $A$-module, and $B$ the square-zero
  $A$-algebra $D_A(L)$. For every $B$-module $M$, restriction to $L$ identifies
  $\Der_A(B,M)$ with the $A$-linear maps $u:L\to M$ such that
  \[
    (0,x)u(y)+(0,y)u(x)=0\quad\text{for all }x,y\in L.
  \]
  Consequently, $\Omega^1_{B/A}$ has the presentation
  \[
    B\otimes_A(L\otimes_A L)\longrightarrow B\otimes_A L
    \longrightarrow\Omega^1_{B/A}\longrightarrow0,
  \]
  where the first arrow sends $b\otimes(x\otimes y)$ to
  $b((0,x)\otimes y+(0,y)\otimes x)$ and the second sends $b\otimes x$ to
  $b\,d_{B/A}(0,x)$.
  \item[\rm(iii)] Suppose that $B$ is the product of two topological $A$-algebras $B_1$, $B_2$ (identified with ideals of $B$).
Then the images of $B_1\otimes_A B_2$ and $B_2\otimes_A B_1$ by the homomorphism $p$ \sref{0.20.4.1.1} are zero, whence it follows immediately that $\mathrm{P}^1_{B/A}$ identifies with the \emph{product} (topological) $\mathrm{P}^1_{B_1/A}\times\mathrm{P}^1_{B_2/A}$, and that the $B$-module $\Omega^1_{B/A}$ is the \emph{direct sum} (topological) of the $B$-modules $\Omega^1_{B_1/A}$ and $\Omega^1_{B_2/A}$ (annihilated by $B_2$ and $B_1$ respectively).
  \item[\rm(iv)] Let $A$, $B$ be two integral domains such that $A\subset B$, $B$ integral over $A$, and the field of fractions of $B$ a \emph{separable} extension of that of $A$.
Then the $B$-module $\Omega^1_{B/A}$ is a \emph{torsion} module.
Indeed, for every $x\in B$, let $f$ be the minimal polynomial of $x$ over the field of fractions $K$ of $A$; there exists a nonzero $a\in A$ such that $af[T]\in A[T]$; as $x$ is separable over $K$, one has $f'(x)\neq 0$, and on the other hand from the relation $af(x)=0$ one deduces that $af'(x)d_{B/A}x=0$, whence our assertion by virtue of \sref{0.20.4.7}.
\end{enumerate}
\end{examples}

\begin{remarks}[20.4.14]
\label{0.20.4.14}
\begin{enumerate}
  \item[\rm(i)] We note that the $B$-module $\Omega^1_{B/A}$, stripped of its topology, is independent of the topologies of $A$ and of $B$.
  \item[\rm(ii)] We will introduce later the ``algebra of principal parts of order $n$'' of $B$ over $A$, $\mathrm{P}^n_{B/A}=(B\otimes_A B)/(\mathfrak{J}_{B/A})^{n+1}$, which is the basis of the ``differential calculus of order $n$''.
\end{enumerate}
\end{remarks}

\subsection{Fundamental functorial properties of \texorpdfstring{$\Omega^1_{B/A}$}{Ω¹(B/A)}}
\label{subsection:0.20.5}

\oldpage[0\textsubscript{IV-1}]{128}
\begin{env}[20.5.1]
\label{0.20.5.1}
In this section and the following, unless expressly stated otherwise, all rings and modules considered are assumed equipped with the \emph{discrete} topology.
\end{env}

\begin{env}[20.5.2]
\label{0.20.5.2}
Let $A$ be a ring, $B$, $C$ two $A$-algebras, $u:B\to C$ an $A$-homomorphism; there is a commutative diagram
\[
\xymatrix{
  B\otimes_A B \ar[r]^{u\otimes u} \ar[d]_{p_{B/A}} & C\otimes_A C \ar[d]^{p_{C/A}} \\
  B \ar[r]^u & C
}
\tag{20.5.2.1}
  \label{0.20.5.2.1}
\]
whence by passage to quotients, an $A$-homomorphism of algebras
\begin{equation}
\tag{20.5.2.2}
  \label{0.20.5.2.2}
  u':\mathrm{P}^1_{B/A}\to\mathrm{P}^1_{C/A}
\end{equation}
such that the diagram
\[
\xymatrix{
  \mathrm{P}^1_{B/A} \ar[r]^{u'} & \mathrm{P}^1_{C/A} \\
  B \ar[u]^{p_1} \ar[r]^u & C \ar[u]_{p_1}
}
\]
is commutative; as $u\otimes u$ maps $\mathfrak{J}_{B/A}$ into $\mathfrak{J}_{C/A}$, one obtains, by restricting $u'$ to $\Omega^1_{B/A}$, a map
\begin{equation}
\tag{20.5.2.3}
  \label{0.20.5.2.3}
  u'':\Omega^1_{B/A}\to\Omega^1_{C/A}
\end{equation}
such that the pair $(u'',u)$ is a \emph{di-homomorphism} for the $B$-module structure of $\Omega^1_{B/A}$ and the $C$-module structure of $\Omega^1_{C/A}$; this last fact allows one to deduce canonically a \emph{homomorphism of $C$-modules}
\begin{equation}
\tag{20.5.2.4}
  \label{0.20.5.2.4}
  u_{C/B/A}:\Omega^1_{B/A}\otimes_B C\to\Omega^1_{C/A}.
\end{equation}

\oldpage[0\textsubscript{IV-1}]{129}
Moreover, as the diagram
\[
\xymatrix{
  \mathrm{P}^1_{B/A} \ar[r]^{u'} & \mathrm{P}^1_{C/A} \\
  B \ar[u]^{p_2} \ar[r]^u & C \ar[u]_{p_2}
}
\tag{20.5.2.5}
  \label{0.20.5.2.5}
\]
is also commutative, one deduces that the diagram
\[
\xymatrix{
  \Omega^1_{B/A} \ar[r]^{u''} & \Omega^1_{C/A} \\
  B \ar[u]^{d_{B/A}} \ar[r]^u & C \ar[u]_{d_{C/A}}
}
\tag{20.5.2.6}
  \label{0.20.5.2.6}
\]
is commutative.
Finally, if $w:C\to D$ is a second homomorphism of $A$-algebras, one has the transitivity property
\begin{equation}
\tag{20.5.2.7}
  \label{0.20.5.2.7}
  (w\circ u)_{D/B/A} = w_{D/C/A}\circ(u_{C/B/A}\otimes\mathrm{I}_D)
\end{equation}
as follows from the definition.
\end{env}

\begin{env}[20.5.3]
\label{0.20.5.3}
Let now $A$, $B$ be two rings, $v:A\to B$ a homomorphism of rings, $C$ a $B$-algebra which becomes an $A$-algebra by means of $v$; then the canonical map $v_0:C\otimes_A C\to C\otimes_B C$ is a surjective \emph{di-homomorphism} of algebras (relative to $v:A\to B$) such that the diagram
\[
\xymatrix{
  C\otimes_A C \ar[r]^{v_0} \ar[d]_{p_{C/A}} & C\otimes_B C \ar[d]^{p_{C/B}} \\
  C \ar[r]^{\mathrm{I}_C} & C
}
\tag{20.5.3.1}
  \label{0.20.5.3.1}
\]
is commutative; by passage to quotients, one deduces a di-homomorphism of algebras
\begin{equation}
\tag{20.5.3.2}
  \label{0.20.5.3.2}
  v':\mathrm{P}^1_{C/A}\to\mathrm{P}^1_{C/B}
\end{equation}
such that the diagram
\[
\xymatrix{
  \mathrm{P}^1_{C/A} \ar[r]^{v'} & \mathrm{P}^1_{C/B} \\
  C \ar[u]^{p_1} \ar[r]^{\mathrm{I}_C} & C \ar[u]_{p_1}
}
\]
is commutative.
As $v_0$ maps $\mathfrak{J}_{C/A}$ into $\mathfrak{J}_{C/B}$, one obtains, by restricting $v'$ to $\Omega^1_{C/A}$, a map
\begin{equation}
\tag{20.5.3.3}
  \label{0.20.5.3.3}
  v_{C/B/A}:\Omega^1_{C/A}\to\Omega^1_{C/B}
\end{equation}
which is a \emph{homomorphism of $C$-modules}.
\oldpage[0\textsubscript{IV-1}]{130}
Moreover, since the diagram
\begin{equation}
\label{0.20.5.3.4}
\xymatrix{
  \mathrm{P}^1_{C/A} \ar[r]^{v'} & \mathrm{P}^1_{C/B} \\
  C \ar[u]^{p_2} \ar[r]^{\mathrm{I}_C} & C \ar[u]_{p_2}
}
\tag{20.5.3.4}
\end{equation}
is commutative, it follows that the diagram
\[
\xymatrix{
  \Omega^1_{C/A} \ar[r]^{v_{C/B/A}} & \Omega^1_{C/B} \\
  C \ar[u]^{d_{C/A}} \ar[r]^{\mathrm{I}_C} & C \ar[u]_{d_{C/B}}
}
\tag{20.5.3.5}
  \label{0.20.5.3.5}
\]
is commutative.
Finally, if $s:A'\to A$ is a second ring homomorphism, one has the transitivity property
\begin{equation}
\tag{20.5.3.6}
  \label{0.20.5.3.6}
  (v\circ s)_{C/B/A'} = v_{C/B/A}\circ s_{C/A/A'}.
\end{equation}
\end{env}

\begin{env}[20.5.4]
\label{0.20.5.4}
If now one has a commutative diagram of ring homomorphisms
\[
\xymatrix{
  B \ar[r]^v & B' \\
  A \ar[u]^f \ar[r]^u & A' \ar[u]_g
}
\]
one deduces from \sref{0.20.5.2.4} and \sref{0.20.5.3.3}, by composition, a homomorphism of $B'$-modules
\begin{equation}
\tag{20.5.4.1}
  \label{0.20.5.4.1}
  \Omega^1_{B/A}\otimes_B B'\xrightarrow{v_{B'/B/A}}\Omega^1_{B'/A}\xrightarrow{u_{B'/A'/A}}\Omega^1_{B'/A'}
\end{equation}
such that the diagram of $A$-homomorphisms
\[
\xymatrix{
  \Omega^1_{B/A}\otimes_B B' \ar[r] & \Omega^1_{B'/A'} \\
  B \ar[u]^{x\mapsto d_{B/A}x\otimes 1} \ar[r]^v & B' \ar[u]_{d_{B'/A'}}
}
\tag{20.5.4.2}
  \label{0.20.5.4.2}
\]
is commutative.
The homomorphism \sref{0.20.5.4.1} corresponds moreover to a di-homomorphism of $B$-modules
\begin{equation}
\tag{20.5.4.3}
  \label{0.20.5.4.3}
  \Omega^1_{B/A}\to\Omega^1_{B'/A'}.
\end{equation}
\end{env}

\begin{proposition}[20.5.5]
\label{0.20.5.5}
If $A'$, $B$ are two $A$-algebras and $B'=B\otimes_A A'$, the canonical homomorphism \sref{0.20.5.4.1}
\begin{equation}
\tag{20.5.5.1}
  \label{0.20.5.5.1}
  \Omega^1_{B/A}\otimes_B B'\to\Omega^1_{B'/A'}
\end{equation}
is bijective.
\end{proposition}

\begin{proof}
\oldpage[0\textsubscript{IV-1}]{131}
The first member of \sref{0.20.5.5.1} is none other than $\Omega^1_{B/A}\otimes_A A'$.
One can write $B'\otimes_{A'}B'=(B\otimes_A B)\otimes_A A'$ up to canonical isomorphism, and $p_{B'/A'}$ then identifies with $p_{B/A}\otimes\mathrm{I}_{A'}$; hence (as $p_{B/A}$ is surjective) $\mathfrak{J}_{B'/A'}=\operatorname{Im}(\mathfrak{J}_{B/A}\otimes_A A')$ and $\mathfrak{J}^2_{B'/A'}=\operatorname{Im}(\mathfrak{J}^2_{B/A}\otimes_A A')$, whence $\mathrm{P}^1_{B'/A'}=\mathrm{P}^1_{B/A}\otimes_A A'$ up to canonical isomorphism, which transforms the augmentation ideals into each other; as the $A$-module $\mathrm{P}^1_{B/A}$ identifies canonically with the direct sum of $B$ and $\Omega^1_{B/A}$, one has
\begin{equation}
\tag{20.5.5.2}
  \label{0.20.5.5.2}
  \Omega^1_{B/A}\otimes_A A' = \Omega^1_{B'/A'}
\end{equation}
by the same isomorphism, and one verifies immediately that the composite of this isomorphism and the canonical isomorphism $\Omega^1_{B/A}\otimes_B B'\xrightarrow{\sim}\Omega^1_{B/A}\otimes_A A'$ is none other than \sref{0.20.5.5.1}.
\end{proof}

\begin{env}[20.5.6]
\label{0.20.5.6}
The canonical homomorphisms \sref{0.20.5.2.4} and \sref{0.20.5.3.3} give, by functoriality, for every $C$-module $L$, canonical homomorphisms
\begin{align}
\tag{20.5.6.1}
  \label{0.20.5.6.1}
  &\Hom_C(\Omega^1_{C/A},L)\to\Hom_C(\Omega^1_{B/A}\otimes_B C,L)=\Hom_B(\Omega^1_{B/A},L) \\
\tag{20.5.6.2}
  \label{0.20.5.6.2}
  &\Hom_C(\Omega^1_{C/B},L)\to\Hom_C(\Omega^1_{C/A},L).
\end{align}
In view of \sref{0.20.4.8.2} and of the commutative diagrams \sref{0.20.5.2.6} and \sref{0.20.5.3.5}, these homomorphisms are none other (up to canonical identification) than the homomorphisms \sref{0.20.2.1.2} and \sref{0.20.2.1.3} respectively.
\end{env}

\begin{theorem}[20.5.7]
\label{0.20.5.7}
Let $u:A\to B$, $v:B\to C$ be two ring homomorphisms.
\begin{enumerate}
  \item[\rm(i)] The sequence of $C$-modules
\begin{equation}
\tag{20.5.7.1}
  \label{0.20.5.7.1}
  \Omega^1_{B/A}\otimes_B C\xrightarrow{v_{C/B/A}}\Omega^1_{C/A}\xrightarrow{u_{C/B/A}}\Omega^1_{C/B}\to 0
\end{equation}
is exact.
  \item[\rm(ii)] For $v_{C/B/A}$ to be left-invertible, it is necessary and sufficient that $C$ be a $B$-algebra formally smooth relative to $A$ \textup{(}for the discrete topologies \textup{(}cf.\ \sref{0.19.9.1}\textup{))}; in particular, it suffices that $C$ be a $B$-algebra formally smooth \textup{(}for the discrete topologies\textup{)}.
\end{enumerate}
\end{theorem}

\begin{proof}
(i) The exactness of the sequence \sref{0.20.2.4.2} shows, taking into account \sref{0.20.5.6}, that the sequence
\[
  0\to\Hom_C(\Omega^1_{C/B},L)\to\Hom_C(\Omega^1_{C/A},L)\to\Hom_C(\Omega^1_{B/A}\otimes_B C,L)
\]
is exact for every $C$-module $L$.
It is known that this implies the exactness of the sequence \sref{0.20.5.7.1} (Bourbaki, \emph{Alg.}, chap.~II, 3$^\text{e}$ \'{e}d., \S\,2, n$^\circ$\,1, th.~1).

\medskip
(ii) By virtue of the exactness of \sref{0.20.5.7.1}, to say that $v_{C/B/A}$ is left-invertible means that the sequence
\begin{equation}
\tag{20.5.7.2}
  \label{0.20.5.7.2}
  0\to\Omega^1_{B/A}\otimes_B C\xrightarrow{v_{C/B/A}}\Omega^1_{C/A}\xrightarrow{u_{C/B/A}}\Omega^1_{C/B}\to 0
\end{equation}
is \emph{exact and split}; it is known (Bourbaki, \emph{loc.\ cit.}, n$^\circ$\,1, prop.~1) that this is equivalent to saying that for every $C$-module $L$, the sequence
\[
  0\to\Hom_C(\Omega^1_{C/B},L)\to\Hom_C(\Omega^1_{C/A},L)\to\Hom_C(\Omega^1_{B/A}\otimes_B C,L)\to 0
\]
is exact; taking into account \sref{0.20.5.6} and \sref{0.20.2.4.2}, this condition is equivalent to
\[
  \operatorname{Exalcom}_{B/A}(C,L)=0
\]
for every $C$-module $L$, and the conclusion follows from \sref{0.19.9.8.1}.
\end{proof}

\oldpage[0\textsubscript{IV-1}]{132}
We note moreover that if one has a commutative diagram of ring homomorphisms
\[
\xymatrix{
  A' \ar[r] & B' \ar[r] & C' \\
  A \ar[u] \ar[r] & B \ar[u] \ar[r] & C \ar[u]
}
\]
one deduces a commutative diagram
\[
\xymatrix{
  \Omega^1_{B/A}\otimes_B C \ar[r] \ar[d] & \Omega^1_{C/A} \ar[r] \ar[d] & \Omega^1_{C/B} \ar[r] \ar[d] & 0 \\
  \Omega^1_{B'/A'}\otimes_{B'} C' \ar[r] & \Omega^1_{C'/A'} \ar[r] & \Omega^1_{C'/B'} \ar[r] & 0
}
\tag{20.5.7.3}
  \label{0.20.5.7.3}
\]
where the vertical arrows come from the di-homomorphisms \sref{0.20.5.4.3}.

\begin{corollary}[20.5.8]
\label{0.20.5.8}
Suppose that the homomorphism $v:B\to C$ makes $C$ into a $B$-algebra formally \'{e}tale \textup{(}for the discrete topologies \sref{0.19.10.2}\textup{)}; then the homomorphism \sref{0.20.5.2.4}
\[
  v_{C/B/A}:\Omega^1_{B/A}\otimes_B C\to\Omega^1_{C/A}
\]
is bijective.
\end{corollary}

\begin{proof}
Indeed, if $C$ is a $B$-algebra formally unramified for the discrete topologies, it follows from \sref{0.19.10.4}, \sref{0.20.4.8} and \sref{0.20.1.1} that one has $\Hom_C(\Omega^1_{C/B},L)=0$ for \emph{every} $C$-module $L$, hence $\Omega^1_{C/B}=0$ (cf.\ \sref{0.20.7.4}); on the other hand, if $C$ is a $B$-algebra formally smooth for the discrete topologies, the sequence \sref{0.20.5.7.2} is exact; whence the corollary.
\end{proof}

\begin{corollary}[20.5.9]
\label{0.20.5.9}
Let $A$ be a ring, $B$ an $A$-algebra, $S$ a multiplicative subset of $B$; then the canonical homomorphism
\begin{equation}
\tag{20.5.9.1}
  \label{0.20.5.9.1}
  S^{-1}\Omega^1_{B/A}\to\Omega^1_{S^{-1}B/A}
\end{equation}
is bijective.
\end{corollary}

\begin{proof}
It suffices to apply \sref{0.20.5.8} to $C=S^{-1}B$, which is a $B$-algebra formally \'{e}tale for the discrete topologies \sref{0.19.10.3}, (ii).
\end{proof}

Taking into account \sref{0.20.5.5}, one can therefore write
\begin{equation}
\tag{20.5.9.2}
  \label{0.20.5.9.2}
  \Omega^1_{S^{-1}B/S^{-1}A} = S^{-1}\Omega^1_{B/A} = \Omega^1_{S^{-1}B/A},
\end{equation}
up to canonical isomorphisms.

\begin{corollary}[20.5.10]
\label{0.20.5.10}
If $k$ is a field and $K=k(X_\alpha)_{\alpha\in I}$ a pure transcendental extension of $k$, the $dX_\alpha$ form a basis of the $K$-vector space $\Omega^1_{K/k}$.
\end{corollary}

\begin{proof}
As $K$ is the field of fractions of the polynomial ring $k[X_\alpha]_{\alpha\in I}$, this follows from \sref{0.20.4.13}, (i) and \sref{0.20.5.9}.
\end{proof}

\oldpage[0\textsubscript{IV-1}]{133}
\begin{env}[20.5.11]
\label{0.20.5.11}
Let $A$ be a ring, $B$ an $A$-algebra, $\mathfrak{R}$ an ideal of $B$, $C$ the $A$-algebra quotient $B/\mathfrak{R}$, and consider the composite homomorphism of $A$-modules
\begin{equation}
\tag{20.5.11.1}
  \label{0.20.5.11.1}
  \delta:\mathfrak{R}\to B\xrightarrow{d_{B/A}}\Omega^1_{B/A}
\end{equation}
where the first arrow is the canonical injection; as $d(xy)=xdy+ydx$, we see that $\delta(\mathfrak{R}^2)\subset\mathfrak{R}.\Omega^1_{B/A}$, whence, by passing to quotients, a homomorphism of $A$-modules
\begin{equation}
\tag{20.5.11.2}
  \label{0.20.5.11.2}
  \delta_{C/B/A}:\mathfrak{R}/\mathfrak{R}^2\to\Omega^1_{B/A}\otimes_B C=\Omega^1_{B/A}/(\mathfrak{R}.\Omega^1_{B/A}).
\end{equation}
But in fact, $\delta_{C/B/A}$ is a homomorphism of \emph{$C$-modules}, since for $x\in B$ and $y\in\mathfrak{R}$, we have $ydx\in\mathfrak{R}.\Omega^1_{B/A}$, hence $d(xy)\equiv xdy$ $(\operatorname{mod}\mathfrak{R}.\Omega^1_{B/A})$, which proves first that \sref{0.20.5.11.2} is a homomorphism of $B$-modules, and as $\mathfrak{R}$ annihilates both members, this establishes our assertion.

If $B'$ is a second $A$-algebra, $u:B\to B'$ an $A$-homomorphism, $\mathfrak{R}'$ an ideal of $B'$ such that $u(\mathfrak{R})\subset\mathfrak{R}'$ and $C'=B'/\mathfrak{R}'$ the quotient algebra, one has a commutative diagram
\begin{equation}
\tag{20.5.11.3}
  \label{0.20.5.11.3}
\xymatrix{
  \mathfrak{R}/\mathfrak{R}^2 \ar[r]^-{\delta_{C/B/A}} \ar[d] & \Omega^1_{B/A}\otimes_B C \ar[d] \\
  \mathfrak{R}'/\mathfrak{R}'^2 \ar[r]^-{\delta_{C'/B'/A}} & \Omega^1_{B'/A}\otimes_{B'} C'
}
\end{equation}
where the vertical arrows come from $u$ \sref{0.20.5.2.4}.
\end{env}

\begin{theorem}[20.5.12]
\label{0.20.5.12}
Let $A$ be a ring, $B$ an $A$-algebra, $C$ the $A$-algebra quotient $B/\mathfrak{R}$, $u:B\to C$ the canonical homomorphism.
\begin{enumerate}
  \item[\rm(i)] One has an exact sequence of $C$-modules
\begin{equation}
\tag{20.5.12.1}
  \label{0.20.5.12.1}
  \mathfrak{R}/\mathfrak{R}^2\xrightarrow{\delta_{C/B/A}}\Omega^1_{B/A}\otimes_B C\xrightarrow{u_{C/B/A}}\Omega^1_{C/A}\to 0
\end{equation}
where $u_{C/B/A}$ and $\delta_{C/B/A}$ are defined by \sref{0.20.5.2.4} and \sref{0.20.5.11.2} respectively.
  \item[\rm(ii)] If one sets $E=B/\mathfrak{R}^2$, the canonical homomorphism \sref{0.20.5.2.4}
\[
  \Omega^1_{B/A}\otimes_B C\to\Omega^1_{E/A}\otimes_E C
\]
is bijective.
  \item[\rm(iii)] The three following conditions are equivalent:
  \begin{enumerate}
    \item[\rm a)] $\delta_{C/B/A}$ is left-invertible.
    \item[\rm b)] Every $A$-extension of $C$ by a $C$-module $L$, whose inverse image by $u:B\to C$ is $A$-trivial, is itself $A$-trivial.
    \item[\rm c)] The $A$-algebra $E=B/\mathfrak{R}^2$ is an $A$-trivial extension of $C$ by $\mathfrak{R}/\mathfrak{R}^2$.
  \end{enumerate}
  \item[\rm(iv)] There is a canonical bijective correspondence between the left inverses of $\delta_{C/B/A}$ and the right inverses of the canonical homomorphism $E\to C$.
\end{enumerate}
\end{theorem}

\begin{proof}
\oldpage[0\textsubscript{IV-1}]{134}
(i) As $u$ is surjective, one has $\Der_B(C,L)=0$ for every $C$-module $L$ by virtue of \sref{0.20.1.1}. The exact sequence \sref{0.20.2.3.1} therefore becomes
\begin{equation}
\tag{20.5.12.2}
  \label{0.20.5.12.2}
  0\to\Der_A(C,L)\xrightarrow{u^0}\Der_A(B,L)\xrightarrow{\partial}\operatorname{Exalcom}_B(C,L)
  \xrightarrow{v^1}\operatorname{Exalcom}_A(C,L)\xrightarrow{u^1}\operatorname{Exalcom}_A(B,L)
\end{equation}
where $v$ denotes the homomorphism $A\to B$.
Recall on the other hand that $\operatorname{Exalcom}_B(C,L)$ is canonically identified with $\Hom_C(\mathfrak{R}/\mathfrak{R}^2,L)$ \sref{0.18.3.8}; one deduces therefore from \sref{0.20.5.12.2} and \sref{0.20.4.8} the exact sequence
\begin{equation}
\tag{20.5.12.3}
  \label{0.20.5.12.3}
  0\to\Hom_C(\Omega^1_{C/A},L)\to\Hom_C(\Omega^1_{B/A}\otimes_B C,L)\xrightarrow{\varphi}\Hom_C(\mathfrak{R}/\mathfrak{R}^2,L)\xrightarrow{\psi}\Ker(u^1)\to 0
\end{equation}
with $\varphi=\partial\circ\eta^{-1}$ and $\psi=v^1\circ\eta$.
Returning to the definitions of $\partial$ \sref{0.20.2.2} and of $\eta$ \sref{0.18.3.8}, one verifies immediately that $\varphi$ is precisely the homomorphism $\Hom(\delta_{C/B/A},1_L)$.
The existence of the exact sequence formed by the first 4 terms of \sref{0.20.5.12.3} shows therefore that the sequence \sref{0.20.5.12.1} is exact (Bourbaki, \emph{Alg.}, chap.~II, 3$^\text{e}$ \'{e}d., \S\,2, n$^\circ$\,1, th.~1).

\medskip
(ii) Apply to $B$ and the ideal $\mathfrak{R}^2$ the exact sequence \sref{0.20.5.12.1}, which gives
\begin{equation}
\tag{20.5.12.4}
  \label{0.20.5.12.4}
  \mathfrak{R}^2/\mathfrak{R}^4\to\Omega^1_{B/A}\otimes_B E\to\Omega^1_{E/A}\to 0
\end{equation}
whence, by tensoring with $C$ (considered as an $E$-algebra), the exact sequence
\[
  \mathfrak{R}^2/\mathfrak{R}^4\otimes_E C\to\Omega^1_{B/A}\otimes_B C\to\Omega^1_{E/A}\otimes_E C\to 0.
\]
But if $x,y$ are two elements of $\mathfrak{R}$ and $\zeta$ the class of $xy$ modulo $\mathfrak{R}^4$, the image $\delta'(\zeta\otimes 1)$ is $(xd_{B/A}(y)+yd_{B/A}(x))\otimes 1$; but as the images of $x$ and of $y$ in $C$ are zero, one has also $d_{B/A}(xy)\otimes 1=0$ in $\Omega^1_{B/A}\otimes_B C$, which proves our assertion.

\medskip
(iii) To say that $\delta_{C/B/A}$ is left-invertible amounts, taking into account the exactness of \sref{0.20.5.12.1}, to saying that the sequence
\begin{equation}
\tag{20.5.12.5}
  \label{0.20.5.12.5}
  0\to\mathfrak{R}/\mathfrak{R}^2\xrightarrow{\delta_{C/B/A}}\Omega^1_{B/A}\otimes_B C\xrightarrow{u_{C/B/A}}\Omega^1_{C/A}\to 0
\end{equation}
is \emph{exact and split}, and it is known (Bourbaki, \emph{loc.\ cit.}) that to say that $\Ker(u^1)=0$ in the exact sequence \sref{0.20.5.12.3} for every $L$, which shows the equivalence of the conditions $a)$ and $b)$ (cf.\ \sref{0.18.3.6.2}).

The fact that $b)$ implies $c)$ comes from the fact that the inverse image by $u:B\to C$ of the $A$-extension $E$ of $C$ by $\mathfrak{R}/\mathfrak{R}^2$ is $A$-trivial, $u$ being composed of the canonical homomorphisms $B\to E\to C$ \sref{0.18.1.6}.
Conversely, $c)$ implies $b)$, since every $B$-extension of $C$ by $L$ is $B$-equivalent to $E\otimes_{\mathfrak{R}/\mathfrak{R}^2}L$ \sref{0.18.3.8}, in other words its class is the image of the class of $E$ (considered as $B$-extension) by the homomorphism
\oldpage[0\textsubscript{IV-1}]{135}
$w_*:\operatorname{Exan}_B(C,\mathfrak{R}/\mathfrak{R}^2)\to\operatorname{Exan}_B(C,L)$ corresponding to a $B$-homomorphism $w:\mathfrak{R}/\mathfrak{R}^2\to L$.
The fact that $c)$ implies $b)$ then follows from the commutativity of the diagram \sref{0.18.3.6.5}
\[
\xymatrix{
  \operatorname{Exan}_B(C,\mathfrak{R}/\mathfrak{R}^2) \ar[r]^{w_*} \ar[d]_{u^*} & \operatorname{Exan}_B(C,L) \ar[d]^{u^*} \\
  \operatorname{Exan}_A(C,\mathfrak{R}/\mathfrak{R}^2) \ar[r]_{w_*} & \operatorname{Exan}_A(C,L)
}
\]

\medskip
(iv) We have seen \sref{0.20.1.7} that the right inverses of $E\to C$ correspond bijectively, in a canonical fashion, to the set of elements $D\in\Der_A(E,\mathfrak{R}/\mathfrak{R}^2)$ such that $D(x)=x$ in $\mathfrak{R}/\mathfrak{R}^2$, hence also, by \sref{0.20.4.8}, to the set of $E$-homomorphisms $h:\Omega^1_{E/A}\to\mathfrak{R}/\mathfrak{R}^2$ such that the composite $\mathfrak{R}/\mathfrak{R}^2\xrightarrow{d_{E/A}}\Omega^1_{E/A}\xrightarrow{h}\mathfrak{R}/\mathfrak{R}^2$ is the identity.
By tensorization with $C$, one deduces that the composite
\[
  \mathfrak{R}/\mathfrak{R}^2\xrightarrow{\delta_{C/B/A}}\Omega^1_{B/A}\otimes_B C\xrightarrow{h\otimes 1}\mathfrak{R}/\mathfrak{R}^2
\]
is the identity; or, since $\mathfrak{R}/\mathfrak{R}^2$ is a $C$-module, $h\mapsto h\otimes 1$ is an isomorphism of the set $\Hom_E(\Omega^1_{E/A},\mathfrak{R}/\mathfrak{R}^2)$ onto $\Hom_C(\Omega^1_{E/A}\otimes_E C,\mathfrak{R}/\mathfrak{R}^2)$; and by (ii) elsewhere one can canonically identify $\Omega^1_{E/A}\otimes_E C$ and $\Omega^1_{B/A}\otimes_B C$, $\delta_{C/E/A}$ being identified with $\delta_{C/B/A}$.
\end{proof}

\begin{example}[20.5.13]
\label{0.20.5.13}
Let $B=A[X_\alpha]_{\alpha\in I}$ be a polynomial algebra over $A$, $\mathfrak{R}$ an ideal of $B$ and $C=B/\mathfrak{R}$; we know that $\Omega^1_{B/A}$ is a free $B$-module whose $dX_\alpha$ form a basis \sref{0.20.4.13}, (i), hence the $dX_\alpha$ also form a basis of the free $C$-module $\Omega^1_{B/A}\otimes_B C$.
On the other hand, from the definition it follows immediately that the image of $\mathfrak{R}/\mathfrak{R}^2$ by $\delta_{C/B/A}$ is the sub-$C$-module generated by the
\[
  dP_\lambda=\sum_\alpha\frac{\partial P_\lambda}{\partial X_\alpha}\,dX_\alpha.
\]
One concludes that $\Omega^1_{C/A}$ is isomorphic to the quotient of the free $C$-module with basis the $dX_\alpha$ by the sub-$C$-module generated by the $dP_\lambda$, which gives a description of the module of differentials of an \emph{arbitrary} $A$-algebra $C$ that can always be obtained in the above manner.
\end{example}

\begin{corollary}[20.5.14]
\label{0.20.5.14}
If $C$ is an $A$-algebra formally smooth \textup{(}for the discrete topologies\textup{)}, the sequence
\begin{equation}
\tag{20.5.14.1}
  \label{0.20.5.14.1}
  0\to\mathfrak{R}/\mathfrak{R}^2\xrightarrow{\delta_{C/B/A}}\Omega^1_{B/A}\otimes_B C\xrightarrow{u_{C/B/A}}\Omega^1_{C/A}\to 0
\end{equation}
is exact and split.
\end{corollary}

\begin{proof}
Indeed, every $A$-extension of $C$ by a $C$-module is then trivial \sref{0.19.4.1}.
\end{proof}

\oldpage[0\textsubscript{IV-1}]{136}
\begin{remark}[20.5.15]
\label{0.20.5.15}
Let $u:A\to B$ be a \emph{surjective} homomorphism of rings; then, for every homomorphism of rings $v:B\to C$, the canonical homomorphism
\begin{equation}
\tag{20.5.15.1}
  \label{0.20.5.15.1}
  u_{C/B/A}:\Omega^1_{C/A}\to\Omega^1_{C/B}
\end{equation}
is \emph{bijective}; this indeed follows from the exact sequence \sref{0.20.5.7.1}, since $\Omega^1_{B/A}=0$ \sref{0.20.4.12}.
\end{remark}

\subsection{Modules of imperfection and characteristic homomorphisms}
\label{subsection:0.20.6}
\label{0.20.6}

\begin{definition}[20.6.1]
\label{0.20.6.1}
Given two ring homomorphisms $u:A\to B$, $v:B\to C$, one calls the \emph{module of imperfection} of the $B$-algebra $C$ relative to $A$, and one denotes by $\Upsilon_{C/B/A}$, the kernel of the $C$-module homomorphism $v_{C/B/A}:\Omega^1_{B/A}\otimes_B C\to\Omega^1_{C/A}$.
One therefore has by definition (cf.\ \sref{0.20.5.7}) the exact sequence
\begin{equation}
\tag{20.6.1.1}
  \label{0.20.6.1.1}
  0\to\Upsilon_{C/B/A}\to\Omega^1_{B/A}\otimes_B C\xrightarrow{v_{C/B/A}}\Omega^1_{C/A}\xrightarrow{u_{C/B/A}}\Omega^1_{C/B}\to 0.
\end{equation}
\end{definition}

When $A=\bb{Z}$ (so that $\Omega^1_{B/A}$ and $\Omega^1_{C/A}$ are the modules of ``absolute'' differentials), one writes $\Upsilon_{C/B}$ instead of $\Upsilon_{C/B/\bb{Z}}$; when $B$ and $C$ are algebras over a prime field $P$, one has $\Upsilon_{C/B/P}=\Upsilon_{C/B}$.

Let $R$, $S$ be multiplicative subsets of $B$ and $C$ respectively, such that the image of $R$ is contained in $S$.
It then follows from the exact sequence \sref{0.20.6.1.1} and \sref{0.20.5.9} that one has
\begin{equation}
\tag{20.6.1.2}
  \label{0.20.6.1.2}
  \Upsilon_{S^{-1}C/R^{-1}B/A}=S^{-1}\Upsilon_{C/B/A}.
\end{equation}

\begin{proposition}[20.6.2]
\label{0.20.6.2}
If $C$ is a $B$-algebra formally smooth relative to $A$ \textup{(}and in particular if $C$ is a $B$-algebra formally smooth\textup{)}, one has $\Upsilon_{C/B/A}=0$.
\end{proposition}

\begin{proof}
This follows from \sref{0.20.5.7}, (ii).
\end{proof}

\begin{proposition}[20.6.3]
\label{0.20.6.3}
Let $k$ be a field, $K$ an extension of $k$.
For $K$ to be a separable extension of $k$, it is necessary and sufficient that $\Upsilon_{K/k}=0$ \textup{(}or equivalently, the canonical homomorphism $\Omega^1_k\otimes_k K\to\Omega^1_K$ is injective, or again \sref{0.20.4.8}, every derivation of $k$ in $K$ extends to a derivation of $K$ in $K$\textup{)}.
\end{proposition}

\begin{proof}
Indeed, it is equivalent to saying that $K$ is separable over $k$ or that $K$ is a $k$-algebra formally smooth \sref{0.19.6.1}; on the other hand, if $P$ is the prime field of $k$, $K$ is separable over $P$, hence also a $P$-algebra formally smooth, and it follows that $K$ is a $k$-algebra formally smooth relative to $P$.
Finally, to say that $K$ is a $k$-algebra formally smooth relative to $P$ amounts, by \sref{0.20.5.7}, (ii), to saying that the homomorphism $v_{K/k/P}:\Omega^1_k\otimes_k K\to\Omega^1_K$ is left-invertible; but as $K$ is a field, this last condition is equivalent to saying that the kernel of $v_{K/k/P}$, that is to say $\Upsilon_{K/k/P}$, is zero.
\end{proof}

\begin{env}[20.6.4]
\label{0.20.6.4}
Consider a commutative diagram of ring homomorphisms
\begin{equation}
\tag{20.6.4.1}
  \label{0.20.6.4.1}
\xymatrix{
  A' \ar[r]^{u'} & B' \ar[r]^{v'} & C' \\
  A \ar[u]^{f} \ar[r]^{u} & B \ar[u]^{g} \ar[r]^{v} & C \ar[u]^{h}
}
\end{equation}
The commutativity of the diagram corresponding to \sref{0.20.5.7.3} implies the existence of a unique $C$-homomorphism
\begin{equation}
\tag{20.6.4.2}
  \label{0.20.6.4.2}
  \Upsilon_{C/B/A}\to\Upsilon_{C'/B'/A'}
\end{equation}
deduced canonically from \sref{0.20.6.4.1} and rendering commutative the diagram
\begin{equation}
\tag{20.6.4.3}
  \label{0.20.6.4.3}
\xymatrix{
  0 \ar[r] & \Upsilon_{C/B/A} \ar[r] \ar[d] & \Omega^1_{B/A}\otimes_B C \ar[r]^{v_{C/B/A}} \ar[d] & \Omega^1_{C/A} \ar[r]^{u_{C/B/A}} \ar[d] & \Omega^1_{C/B} \ar[r] \ar[d] & 0 \\
  0 \ar[r] & \Upsilon_{C'/B'/A'} \ar[r] & \Omega^1_{B'/A'}\otimes_{B'} C' \ar[r]_{v_{C'/B'/A'}} & \Omega^1_{C'/A'} \ar[r]_{u_{C'/B'/A'}} & \Omega^1_{C'/B'} \ar[r] & 0
}
\end{equation}
where the vertical arrows come from the di-homomorphisms \sref{0.20.5.4.3}.

The definition of \sref{0.20.6.4.2} is also equivalent to that of a $C'$-homomorphism
\begin{equation}
\tag{20.6.4.4}
  \label{0.20.6.4.4}
  \Upsilon_{C/B/A}\otimes_C C'\to\Upsilon_{C'/B'/A'}
\end{equation}
which, composed with the canonical homomorphism $\Upsilon_{C/B/A}\to\Upsilon_{C/B/A}\otimes_C C'$, gives back \sref{0.20.6.4.2}.
It is clear that \sref{0.20.6.4.2} enjoys an evident transitivity property, allowing us to say that $\Upsilon_{C/B/A}$ is a \emph{functor} in the triple $(A,B,C)$.
\end{env}

\begin{env}[20.6.5]
\label{0.20.6.5}
It will be convenient for us in what follows, under the conditions of \sref{0.20.6.1}, to introduce a \emph{complex} (of chains) \emph{of $C$-modules} $\mathrm{K}_\bullet(C/B/A)$ whose terms are zero except in degrees $0$ and $1$, where one takes
\begin{equation}
\tag{20.6.5.1}
  \label{0.20.6.5.1}
  \begin{cases}
  \mathrm{K}_0(C/B/A)=\Omega^1_{C/A} \\
  \mathrm{K}_1(C/B/A)=\Omega^1_{B/A}\otimes_B C
  \end{cases}
\end{equation}
the differential operator being $v_{C/B/A}:\mathrm{K}_1\to\mathrm{K}_0$.
This allows writing (up to canonical isomorphisms) $\Omega^1_{C/B}$ and $\Upsilon_{C/B/A}$ as the \emph{homology modules} of this complex
\begin{equation}
\tag{20.6.5.2}
  \label{0.20.6.5.2}
  \mathrm{H}_0(\mathrm{K}_\bullet(C/B/A))=\Omega^1_{C/B},\qquad \mathrm{H}_1(\mathrm{K}_\bullet(C/B/A))=\Upsilon_{C/B/A}.
\end{equation}
\end{env}

Similarly:
\begin{proposition}[20.6.6]
\label{0.20.6.6}
Under the hypotheses of \sref{0.20.6.1}, for every $C$-module $L$, one has canonical $C$-isomorphisms
\begin{align}
\tag{20.6.6.1}
  \label{0.20.6.6.1}
  \mathrm{H}^0(\mathrm{K}_\bullet(C/B/A),L) &\simeq\Der_B(C,L) \\
\tag{20.6.6.2}
  \label{0.20.6.6.2}
  \mathrm{H}^1(\mathrm{K}_\bullet(C/B/A),L) &\simeq\operatorname{Exalcom}_{B/A}(C,L).
\end{align}
\end{proposition}

\begin{proof}
Indeed, the complex of cochains $\Hom_C(\mathrm{K}_\bullet(C/B/A),L)$ is none other, by virtue of \sref{0.20.4.8} and \sref{0.20.5.6}, than the complex
\[
  \cdots\to 0\to\Der_A(C,L)\to\Der_A(B,L)\to 0\to\cdots
\]
\oldpage[0\textsubscript{IV-1}]{138}
where the differential operator is $v^0$ (with the notations of \sref{0.20.2.1}).
The proposition follows from the exact sequence \sref{0.20.2.4.2} and from the definition of the cohomology modules
\begin{equation}
\tag{20.6.6.3}
  \label{0.20.6.6.3}
  \mathrm{H}^i(\mathrm{K}_\bullet,L)=\mathrm{H}^i(\Hom_C(\mathrm{K}_\bullet,L)).
\end{equation}
\end{proof}

If one has a commutative diagram of ring homomorphisms
\[
\xymatrix{
  A' \ar[r] & B' \ar[r] & C' \\
  A \ar[u] \ar[r] & B \ar[u] \ar[r] & C \ar[u]
}
\]
the di-homomorphisms \sref{0.20.5.4.3} define a di-homomorphism of complexes of modules
\begin{equation}
\tag{20.6.6.4}
  \label{0.20.6.6.4}
  \mathrm{K}_\bullet(C/B/A)\to\mathrm{K}_\bullet(C'/B'/A')
\end{equation}
and the di-homomorphisms that one deduces for the homology or the cohomology are identified, by the formulas \sref{0.20.6.5.2}, \sref{0.20.6.6.1} and \sref{0.20.6.6.2}, with the di-homomorphisms already defined in \sref{0.20.5.4.3}, \sref{0.20.6.4.2}, \sref{0.20.2.1} and \sref{0.18.3.7.2} (taking into account, for the last, the definition of the operator $\partial$ in the exact sequence \sref{0.20.2.4.2}).

\begin{env}[20.6.7]
\label{0.20.6.7}
One knows that for a complex $\mathrm{K}_\bullet$ of $C$-modules and a $C$-module $L$, one has canonical homomorphisms
\[
  \alpha_i:\mathrm{H}^i(\mathrm{K}_\bullet,L)\to\Hom_C(\mathrm{H}_i(\mathrm{K}_\bullet),L)
\]
(M, IV, 6).
Here, the canonical homomorphism
\[
  \alpha_1:\mathrm{H}^1(\mathrm{K}_\bullet(C/B/A),L)\to\Hom_C(\mathrm{H}_1(\mathrm{K}_\bullet(C/B/A)),L)
\]
is defined immediately as obtained by passing to the quotient by the image of $\Hom_C(\mathrm{K}_0,L)$ of the restriction homomorphism
\[
  \Hom_C(\mathrm{K}_1(C/B/A),L)\to\Hom_C(\mathrm{H}_1(\mathrm{K}_\bullet(C/B/A)),L)
\]
since $\mathrm{H}_1(\mathrm{K}_\bullet(C/B/A))$ is none other than the kernel of $\mathrm{K}_1\to\mathrm{K}_0$; taking into account \sref{0.20.6.6.2} and \sref{0.20.6.5.2}, one obtains therefore a canonical $C$-homomorphism
\begin{equation}
\tag{20.6.7.1}
  \label{0.20.6.7.1}
  \operatorname{Exalcom}_{B/A}(C,L)\to\Hom_C(\Upsilon_{C/B/A},L)
\end{equation}
which can be described explicitly as follows: by virtue of \sref{0.20.2.4.2}, every $B$-extension of $C$ by $L$ which is $A$-trivial comes from the datum of an $A$-derivation $D$ of $B$ in $L$, hence \sref{0.20.4.8} of a $C$-homomorphism $f$ of $\Omega^1_{B/A}\otimes_B C$ in $L$; one associates to this extension the \emph{restriction} of $f$ to $\Upsilon_{C/B/A}$, which in fact depends only on the class of this extension and not on the choice of $D$.
\end{env}

\begin{definition}[20.6.8]
\label{0.20.6.8}
Let $u:A\to B$, $v:B\to C$ be two ring homomorphisms.
For every $B$-extension $E$ of $C$ by a $C$-module $L$, which is $A$-trivial \sref{0.18.3.7}, one calls the \emph{characteristic homomorphism} of $E$, and one denotes by $\chi_E$, the $C$-homomorphism $\Upsilon_{C/B/A}\to L$, image of the class of $E$ by the canonical homomorphism \sref{0.20.6.7.1}.
\end{definition}

\oldpage[0\textsubscript{IV-1}]{139}
One can define the homomorphism $\chi_E$ in another way:
\begin{proposition}[20.6.9]
\label{0.20.6.9}
Let $E$ be a $B$-extension of $C$ by a $C$-module $L$, which is $A$-trivial \sref{0.18.3.7}; then the diagram
\begin{equation}
\tag{20.6.9.1}
  \label{0.20.6.9.1}
\xymatrix{
  0 \ar[r] & \Upsilon_{C/B/A} \ar[r] \ar[d]^{\chi_E} & \Omega^1_{B/A}\otimes_B C \ar[r]^{v_{C/B/A}} \ar[d]^{q_{E/B/A}\otimes 1_C} & \Omega^1_{C/A} \ar[d]^{1} \\
  0 \ar[r] & L \ar[r]_{\delta_{C/E/A}} & \Omega^1_{E/A}\otimes_E C \ar[r]_{p_{C/E/A}} & \Omega^1_{C/A} \ar[r] & 0
}
\end{equation}
where $q:B\to E$ defines the structure of $B$-extension of $E$ and $p:E\to C$ is the augmentation, is commutative and its rows are exact.
\end{proposition}

\begin{proof}
The bottom row of the diagram is the sequence \sref{0.20.5.12.5} relative to the two homomorphisms $A\xrightarrow{q\circ u}E$ and $p:E\to C=E/L$; as $p$ is surjective and $E$ is an $A$-trivial extension, this sequence is exact and split by virtue of \sref{0.20.5.12}, (iii).
The commutativity of the right-hand square of \sref{0.20.6.9.1} follows from the relation $v=p\circ q$ \sref{0.20.5.2.7}; the image by $q_{E/B/A}\otimes 1_C$ of the kernel $\Upsilon_{C/B/A}$ of $v_{C/B/A}$ is therefore contained in the kernel $L$ of $p_{C/E/A}$.
On the other hand, let $h:C\to E$ be an $A$-homomorphism that is a right inverse of $p$, and let $j:L\to E$ be the canonical injection, so that one has $q(b)=h(v(b))+j(D(b))$ for $b\in B$, where $D$ is the $A$-derivation of $B$ in $L$ defining the $B$-extension $E$; one can write $D=f\circ d_{B/A}$, where $f:\Omega^1_{B/A}\to L$ is a $B$-homomorphism.
By virtue of \sref{0.20.5.2.6}, one has $q_{E/B/A}(d_{B/A}(b))=d_{E/A}(q(b))=d_{E/A}(h(v(b)))+d_{E/A}(j(f(d_{B/A}(b))))$ for $b\in B$.
Let then $z=\sum_i d_{B/A}(b_i)\otimes c_i$, where $b_i\in B$ and $c_i\in C$, be an element of $\Omega^1_{B/A}\otimes_B C$; one has
\[
  (q_{E/B/A}\otimes 1_C)(z)=\sum_i d_{E/A}(h(v(b_i)))\otimes c_i+\sum_i d_{E/A}(j(f(d_{B/A}(b_i))))\otimes c_i.
\]
In the first sum, one has $d_{E/A}(h(v(b_i)))=h_{E/C/A}(d_{C/A}(v(b_i)))$, hence this sum equals $(h_{E/C/A}\otimes 1_C)(v_{C/B/A}(z))$.
So if one takes $z\in\Upsilon_{C/B/A}$, this sum is zero, and it remains, by definition of $\delta_{C/E/A}$,
\[
  (q_{E/B/A}\otimes 1_C)(z)=\sum_i\delta_{C/E/A}(c_i f(d_{B/A}(b_i)))=\delta_{C/E/A}(f(z))
\]
which proves the commutativity of the left-hand square in \sref{0.20.6.9.1}.
\end{proof}

One notes that when one only assumes that $\delta_{C/E/A}$ is \emph{injective} (and not necessarily left-invertible), this interpretation would still allow one to \emph{define} $\chi_E$ as the restriction of $q_{E/B/A}\otimes 1_C$ to $\Upsilon_{C/B/A}$.

\begin{corollary}[20.6.10]
\label{0.20.6.10}
If $\mathfrak{R}$ is an ideal of $B$, $C=B/\mathfrak{R}$, $E=B/\mathfrak{R}^2$, and if $E$ is a $B$-extension $A$-trivial \sref{0.18.3.7} of $C$ by $\mathfrak{R}/\mathfrak{R}^2$, then the characteristic homomorphism $\chi_E:\Upsilon_{C/B/A}\to\mathfrak{R}/\mathfrak{R}^2$ is bijective.
\end{corollary}

\begin{proof}
\oldpage[0\textsubscript{IV-1}]{140}
Indeed, in the diagram \sref{0.20.6.9.1}, the two right-hand vertical arrows are bijective \sref{0.20.5.12}, (ii).
\end{proof}

\begin{theorem}[20.6.11]
\label{0.20.6.11}
Let $u:A\to B$, $v:B\to C$ be two ring homomorphisms, $L$ a $C$-module.
Suppose one of the following conditions satisfied:
\begin{enumerate}
  \item[\rm(i)] $L$ is an injective $C$-module.
  \item[\rm(ii)] $\Upsilon_{C/B/A}$ is a direct factor of $\Omega^1_{B/A}\otimes_B C$ and $u_{C/B/A}:\Omega^1_{C/A}\to\Omega^1_{C/B}$ is right-invertible.
\end{enumerate}
Then the canonical homomorphism \sref{0.20.6.7.1}
\[
  \operatorname{Exalcom}_{B/A}(C,L)\to\Hom_C(\Upsilon_{C/B/A},L)
\]
is bijective.

In particular, if $\Omega^1_{C/B}$ and $\Omega^1_{C/A}$ are projective $C$-modules, the canonical homomorphism \sref{0.20.6.7.1} is bijective.
\end{theorem}

\begin{proof}
The fact that each of the conditions (i), (ii) implies that \sref{0.20.6.7.1} is bijective follows in both cases from the definition of $\alpha_1$.
We note moreover that condition (ii) is \emph{necessary and sufficient} for the homomorphism \sref{0.20.6.7.1} to be bijective for \emph{every} $C$-module $L$ (Bourbaki, \emph{Alg.}, chap.~II, 3$^\text{e}$ \'{e}d., \S\,2, n$^\circ$\,1, prop.~1).
If one assumes that $\Omega^1_{C/B}$ and $\Omega^1_{C/A}$ are projective $C$-modules, then, in the exact sequence \sref{0.20.6.1.1}, $\Ker(u_{C/B/A})$ is a projective $C$-module, since the sequence
\[
  0\to\Ker(u_{C/B/A})\to\Omega^1_{C/A}\to\Omega^1_{C/B}\to 0
\]
is split, $\Omega^1_{C/B}$ being projective; as $\Ker(u_{C/B/A})=\operatorname{Im}(v_{C/B/A})$, the sequence
\[
  0\to\Upsilon_{C/B/A}\to\Omega^1_{B/A}\otimes_B C\to\operatorname{Im}(v_{C/B/A})\to 0
\]
is split.
\end{proof}

\begin{corollary}[20.6.12]
\label{0.20.6.12}
Suppose that $C$ is an $A$-algebra formally smooth.
Then there exists a canonical homomorphism
\begin{equation}
\tag{20.6.12.1}
  \label{0.20.6.12.1}
  \operatorname{Exalcom}_B(C,L)\to\Hom_C(\Upsilon_{C/B/A},L).
\end{equation}
Moreover, this homomorphism is bijective if one of the conditions \textup{(i)}, \textup{(ii)} of \sref{0.20.6.11} is verified.
\end{corollary}

\begin{proof}
Indeed, it follows from the hypothesis on $C$ that $\operatorname{Exalcom}_B(C,L)=\operatorname{Exalcom}_{B/A}(C,L)$, and the homomorphism \sref{0.20.6.12.1} is none other than \sref{0.20.6.7.1}.
\end{proof}

\begin{corollary}[20.6.13]
\label{0.20.6.13}
If $C$ is an $A$-algebra formally smooth and if $\Omega^1_{C/B}$ is a projective $C$-module, the homomorphism \sref{0.20.6.12.1} is bijective.
\end{corollary}

\begin{proof}
Indeed, one knows then that $\Omega^1_{C/A}$ is a projective $C$-module \sref{0.20.4.9}, the discrete topologies being considered.
\end{proof}

\begin{env}[20.6.14]
\label{0.20.6.14}
The notations remaining the same, suppose now in addition that $A$, $B$, $C$ are $\Lambda$-algebras and $u$, $v$ are $\Lambda$-homomorphisms, which amounts to giving three ring homomorphisms $\Lambda\xrightarrow{s}A\xrightarrow{u}B\xrightarrow{v}C$.

\oldpage[0\textsubscript{IV-1}]{141}
One has therefore, apart from the module of imperfection $\Upsilon_{C/B/A}$, the modules of imperfection $\Upsilon_{B/A/\Lambda}$, $\Upsilon_{C/A/\Lambda}$ and $\Upsilon_{C/B/\Lambda}$, and one has already defined canonical $C$-module homomorphisms \sref{0.20.6.4.2}
\begin{align}
\tag{20.6.14.1}
  \label{0.20.6.14.1}
  u':\Upsilon_{C/A/\Lambda}&\to\Upsilon_{C/B/\Lambda} \\
\tag{20.6.14.2}
  \label{0.20.6.14.2}
  s':\Upsilon_{C/B/\Lambda}&\to\Upsilon_{C/B/A}.
\end{align}
As in the commutative diagram \sref{0.20.5.7.3}
\begin{equation}
\tag{20.6.14.3}
  \label{0.20.6.14.3}
\xymatrix{
  \Omega^1_{A/\Lambda}\otimes_A B \ar[r] \ar[d] & \Omega^1_{B/\Lambda} \ar[r] \ar[d] & \Omega^1_{B/A} \ar[r] \ar[d] & 0 \\
  \Omega^1_{A/\Lambda}\otimes_A C \ar[r] & \Omega^1_{C/\Lambda} \ar[r] & \Omega^1_{C/A} \ar[r] & 0
}
\end{equation}
the bottom row is formed of $C$-modules, one deduces by tensorization a commutative diagram
\begin{equation}
\tag{20.6.14.4}
  \label{0.20.6.14.4}
\xymatrix{
  \Omega^1_{A/\Lambda}\otimes_A C \ar[r] \ar@{=}[d] & \Omega^1_{B/\Lambda}\otimes_B C \ar[r] \ar[d] & \Omega^1_{B/A}\otimes_B C \ar[r] \ar[d] & 0 \\
  \Omega^1_{A/\Lambda}\otimes_A C \ar[r] & \Omega^1_{C/\Lambda} \ar[r] & \Omega^1_{C/A} \ar[r] & 0
}
\end{equation}
where the first row is still exact and the left vertical arrow is the identity; if one puts
\begin{equation}
\tag{20.6.14.5}
  \label{0.20.6.14.5}
  \Upsilon^C_{B/A/\Lambda}=\Ker(u_{B/A/\Lambda}\otimes 1_C)=\Ker(\Omega^1_{A/\Lambda}\otimes_A C\to\Omega^1_{B/\Lambda}\otimes_B C)
\end{equation}
one sees, taking into account the definition of $\Upsilon_{C/A/\Lambda}$, that one has a unique $C$-homomorphism
\begin{equation}
\tag{20.6.14.6}
  \label{0.20.6.14.6}
  v':\Upsilon^C_{B/A/\Lambda}\to\Upsilon_{C/A/\Lambda}
\end{equation}
rendering commutative the diagram
\begin{equation}
\tag{20.6.14.7}
  \label{0.20.6.14.7}
\xymatrix{
  0 \ar[r] & \Upsilon^C_{B/A/\Lambda} \ar[r] \ar[d]^{v'} & \Omega^1_{A/\Lambda}\otimes_A C \ar[r] \ar@{=}[d] & \Omega^1_{B/\Lambda}\otimes_B C \ar[r] \ar[d] & \Omega^1_{B/A}\otimes_B C \ar[r] \ar[d] & 0 \\
  0 \ar[r] & \Upsilon_{C/A/\Lambda} \ar[r] & \Omega^1_{A/\Lambda}\otimes_A C \ar[r] & \Omega^1_{C/\Lambda} \ar[r] & \Omega^1_{C/A} \ar[r] & 0
}
\end{equation}
whose rows are exact.

When $\Lambda=\bb{Z}$, one writes $\Upsilon^C_{B/A}$ instead of $\Upsilon^C_{B/A/\bb{Z}}$.
If $\Lambda$ is a prime field, one has $\Upsilon^C_{B/A}=\Upsilon^C_{B/A/\Lambda}$.
\end{env}

\begin{env}[20.6.15]
\label{0.20.6.15}
To study the relations between the preceding modules, we introduce on the one hand the complex of $B$-modules $\mathrm{K}_\bullet(B/A/\Lambda)$ \sref{0.20.6.5}, and on the other hand the complexes
\oldpage[0\textsubscript{IV-1}]{142}
of $C$-modules $\mathrm{K}_\bullet(C/A/\Lambda)$ and $\mathrm{K}_\bullet(C/B/\Lambda)$.
We shall also set
\begin{equation}
\tag{20.6.15.1}
  \label{0.20.6.15.1}
  \mathrm{K}^C_\bullet(B/A/\Lambda)=\mathrm{K}_\bullet(B/A/\Lambda)\otimes_B C.
\end{equation}
On the other hand, we denote by $\mathrm{T}_\bullet(C/B/\Lambda)$ the complex of $C$-modules whose terms are zero except in degrees $0$ and $1$, and where
\begin{equation}
\tag{20.6.15.2}
  \label{0.20.6.15.2}
  \mathrm{T}_0(C/B/\Lambda)=\mathrm{T}_1(C/B/\Lambda)=\Omega^1_{B/\Lambda}\otimes_B C
\end{equation}
the differential operator being the identity, so that this complex is homotopic to $0$; finally let us set
\begin{equation}
\tag{20.6.15.3}
  \label{0.20.6.15.3}
  \mathrm{K}'_\bullet(C/A/\Lambda)=\mathrm{K}_\bullet(C/A/\Lambda)\oplus\mathrm{T}_\bullet(C/B/\Lambda).
\end{equation}
By virtue of the trivial character of $\mathrm{T}_\bullet$, it is clear that one has
\begin{equation}
\tag{20.6.15.4}
  \label{0.20.6.15.4}
  \mathrm{H}^i(\mathrm{K}'_\bullet,L)=\mathrm{H}^i(\mathrm{K}_\bullet,L)\quad\text{and}\quad\mathrm{H}_i(\mathrm{K}'_\bullet,L)=\mathrm{H}_i(\mathrm{K}_\bullet,L)
\end{equation}
for every $C$-module $L$ and every $i$.
\end{env}

\begin{env}[20.6.16]
\label{0.20.6.16}
Let us now define an exact sequence of complexes, \emph{split} in each degree
\begin{equation}
\tag{20.6.16.1}
  \label{0.20.6.16.1}
  0\to\mathrm{K}^C_\bullet(B/A/\Lambda)\xrightarrow{j}\mathrm{K}'_\bullet(C/A/\Lambda)\xrightarrow{p}\mathrm{K}_\bullet(C/B/\Lambda)\to 0
\end{equation}
in the following way: let us denote temporarily by
\[
  f:\Omega^1_{A/\Lambda}\otimes_A C\to\Omega^1_{B/\Lambda}\otimes_B C, \quad
  g:\Omega^1_{B/\Lambda}\otimes_B C\to\Omega^1_{C/\Lambda}
\]
the canonical homomorphisms $u_{B/A/\Lambda}\otimes 1_C$ and $v_{C/B/\Lambda}$ respectively, whose composite is $g\circ f=(v\circ u)_{C/A/\Lambda}$ (cf.\ \sref{0.20.6.14.4}).
One will take $j_1(x)=(x,f(x))$, $p_1(y,z)=z-f(y)$, $j_0(x)=(g(x),x)$, $p_0(y,z)=g(z)-y$ so that $\operatorname{Im}(j_1)=\Ker(p_1)$ is the graph of $f$, and $\operatorname{Im}(j_0)=\Ker(p_0)$ is the graph of $g$, complementary to $\{0\}\oplus\mathrm{T}_1$ and $\mathrm{K}_1(C/A/\Lambda)\oplus\{0\}$ respectively; the verification of the commutativity of the diagram
\[
\xymatrix{
  0 \ar[r] & \mathrm{K}^C_1(B/A/\Lambda) \ar[r]^{j_1} \ar[d] & \mathrm{K}'_1(C/A/\Lambda) \ar[r]^{p_1} \ar[d] & \mathrm{K}_1(C/B/\Lambda) \ar[r] \ar[d] & 0 \\
  0 \ar[r] & \mathrm{K}^C_0(B/A/\Lambda) \ar[r]_{j_0} & \mathrm{K}'_0(C/A/\Lambda) \ar[r]_{p_0} & \mathrm{K}_0(C/B/\Lambda) \ar[r] & 0
}
\]
where the vertical arrows are the differential operators, is immediate.
\end{env}

\begin{theorem}[20.6.17]
\label{0.20.6.17}
One has an exact sequence of $C$-modules
\begin{equation}
\tag{20.6.17.1}
  \label{0.20.6.17.1}
  0\to\Upsilon^C_{B/A/\Lambda}\xrightarrow{v'}\Upsilon_{C/A/\Lambda}\xrightarrow{u'}\Upsilon_{C/B/\Lambda}\xrightarrow{\partial}\Omega^1_{B/A}\otimes_B C\xrightarrow{v_{C/B/A}}\Omega^1_{C/A}\xrightarrow{u_{C/B/A}}\Omega^1_{C/B}\to 0
\end{equation}
\oldpage[0\textsubscript{IV-1}]{143}
where the boundary operator $\partial$ is the composite
\begin{equation}
\tag{20.6.17.2}
  \label{0.20.6.17.2}
  \Upsilon_{C/B/\Lambda}\xrightarrow{s'}\Upsilon_{C/B/A}\to\Omega^1_{B/A}\otimes_B C
\end{equation}
the second arrow being the canonical injection.
\end{theorem}

\begin{proof}
Writing the long exact homology sequence for the exact sequence of complexes \sref{0.20.6.16.1}, one obtains \sref{0.20.6.17.1}, the homology being zero in degrees distinct from $0$ and $1$; the fact that the homomorphisms of this exact sequence come from the functoriality of $j$ and $p$ is immediate.
It remains to verify that $\partial$ is equal to \sref{0.20.6.17.2}; but an element $z\in\Upsilon_{C/B/\Lambda}$ is the image by $p_1$ of $(0,z)$, whence one deduces immediately that $\partial(z)$ is the image of $z$ by the canonical homomorphism $s_{B/A/\Lambda}\otimes 1_C:\Omega^1_{B/\Lambda}\otimes_B C\to\Omega^1_{B/A}\otimes_B C$.
Our assertion follows from the commutativity of the diagram
\begin{equation}
\tag{20.6.17.3}
  \label{0.20.6.17.3}
\xymatrix{
  \Upsilon_{C/B/\Lambda} \ar[r]^{s'} \ar[dd] & \Upsilon_{C/B/A} \ar[dd] \\
  & \\
  \Omega^1_{B/\Lambda}\otimes_B C \ar[r] & \Omega^1_{B/A}\otimes_B C
}
\end{equation}
(cf.\ \sref{0.20.6.4.3}).
\end{proof}

\begin{corollary}[20.6.18]
\label{0.20.6.18}
\begin{enumerate}
  \item[\rm(i)] The sequence of $C$-modules
\begin{equation}
\tag{20.6.18.1}
  \label{0.20.6.18.1}
  0\to\Upsilon^C_{B/A/\Lambda}\xrightarrow{v'}\Upsilon_{C/A/\Lambda}\xrightarrow{u'}\Upsilon_{C/B/\Lambda}\xrightarrow{s'}\Upsilon_{C/B/A}\to 0
\end{equation}
is exact.
  \item[\rm(ii)] If $B$ is an $A$-algebra formally smooth relative to $\Lambda$, one has $\Upsilon^C_{B/A/\Lambda}=0$.
\end{enumerate}
\end{corollary}

\begin{proof}
(i) In \sref{0.20.6.17.1}, the image of $\Upsilon_{C/B/\Lambda}$ by $\partial$ is the kernel of $v_{C/B/A}$, hence $\Upsilon_{C/B/A}$ by definition.

(ii) The hypothesis implies that the sequence
\[
  0\to\Omega^1_{A/\Lambda}\otimes_A B\to\Omega^1_{B/\Lambda}\to\Omega^1_{B/A}\to 0
\]
is exact and split \sref{0.20.5.7}; by tensorization with $C$, the sequence
\[
  0\to\Omega^1_{A/\Lambda}\otimes_A C\to\Omega^1_{B/\Lambda}\otimes_B C\to\Omega^1_{B/A}\otimes_B C\to 0
\]
remains therefore exact, whence our assertion.
\end{proof}

\begin{corollary}[20.6.19]
\label{0.20.6.19}
Let $K$ be a field, $E$, $F$ two extensions of $K$ such that $K\subset E\subset F$.
\begin{enumerate}
  \item[\rm(i)] If $F$ is a separable extension of $E$, one has $\Upsilon_{F/E/K}=0$ \textup{(}or equivalently, the canonical homomorphism $\Omega^1_{E/K}\otimes_E F\to\Omega^1_{F/K}$ is injective\textup{)}.
  \item[\rm(ii)] Conversely, if $F$ is a separable extension of $K$ and if $\Upsilon_{F/E/K}=0$, then $F$ is a separable extension of $E$.
\end{enumerate}
\end{corollary}

\begin{proof}
If $P$ is the prime field of $K$, one has the exact sequence \sref{0.20.6.18}
\[
  \Upsilon_{F/K/P}\to\Upsilon_{F/E/P}\to\Upsilon_{F/E/K}\to 0.
\]
If $\Upsilon_{F/E/P}=\Upsilon_{F/E}=0$, one therefore has $\Upsilon_{F/E/K}=0$, whence (i) by virtue of \sref{0.20.6.3}; conversely, if $\Upsilon_{F/E/K}=0$ and $\Upsilon_{F/K/P}=\Upsilon_{F/K}=0$, one has $\Upsilon_{F/E/P}=\Upsilon_{F/E}=0$, whence (ii) by virtue of \sref{0.20.6.3}.
\end{proof}

\oldpage[0\textsubscript{IV-1}]{144}
\begin{corollary}[20.6.20]
\label{0.20.6.20}
\begin{enumerate}
  \item[\rm(i)] If $K$ is an algebraic separable extension of $k$, one has $\Omega^1_{K/k}=0$.
  \item[\rm(ii)] Let $k$ be a field of characteristic $0$, $K$ an extension of $k$.
  For $\Omega^1_{K/k}=0$, it is necessary and sufficient that $K$ be an algebraic extension of $k$.
  In particular, for a field $K$ of characteristic $0$ to satisfy $\Omega^1_K=0$, it is necessary and sufficient that $K$ be an algebraic extension of $\bb{Q}$ (cf.\ \sref{0.21.4.4} and \sref{0.21.7.5}).
\end{enumerate}
\end{corollary}

\begin{proof}
(i) For every $x\in K$, let $f$ be the minimal polynomial of $x$ over $k$; as $f'(x)\neq 0$ and $d_{K/k}(f(x))=f'(x)d_{K/k}(x)=0$, one has $d_{K/k}(x)=0$ and our assertion follows from \sref{0.20.4.7}.

(ii) There exists a pure extension $L$ of $k$ such that $k\subset L\subset K$ and $K$ is an algebraic extension of $L$.
As $K$ is separable over $L$, it follows from \sref{0.20.6.19}, (i) that the sequence
\[
  0\to\Omega^1_{L/k}\otimes_L K\to\Omega^1_{K/k}\to\Omega^1_{K/L}\to 0
\]
is exact, and from (i) that $\Omega^1_{K/L}=0$.
The relation $\Omega^1_{K/k}=0$ is therefore equivalent to $\Omega^1_{L/k}=0$, and as $L$ is a pure extension of $k$, it follows from \sref{0.20.5.10} that the relation $\Omega^1_{L/k}=0$ is equivalent to $L=k$.
\end{proof}

\begin{remarks}[20.6.21]
\label{0.20.6.21}
\begin{enumerate}
  \item[\rm(i)] As $\Upsilon_{B/A/\Lambda}$ is the kernel of $u_{B/A/\Lambda}:\Omega^1_{A/\Lambda}\otimes_A B\to\Omega^1_{B/\Lambda}$ and $\Upsilon^C_{B/A/\Lambda}$ the kernel of $u_{B/A/\Lambda}\otimes 1_C$, one has a canonical homomorphism
\begin{equation}
\tag{20.6.21.1}
  \label{0.20.6.21.1}
  \Upsilon_{B/A/\Lambda}\otimes_B C\to\Upsilon^C_{B/A/\Lambda}.
\end{equation}
This homomorphism is bijective when the sequence
\begin{equation}
\tag{20.6.21.2}
  \label{0.20.6.21.2}
  0\to\Upsilon_{B/A/\Lambda}\otimes_B C\to\Omega^1_{A/\Lambda}\otimes_A C\to\Omega^1_{B/\Lambda}\otimes_B C\to\Omega^1_{B/A}\otimes_B C\to 0
\end{equation}
is exact, which occurs in the following cases:
\begin{enumerate}
  \item[$1^\circ$] $C$ is a flat $B$-module.
  \item[$2^\circ$] The $B$-modules $\Omega^1_{B/\Lambda}$ and $\Omega^1_{B/A}$ are \emph{flat}: indeed, the kernel of $\Omega^1_{B/\Lambda}\to\Omega^1_{B/A}$ is then also flat ($\mathbf{0_I}$, 6.1.2), and the sequence \sref{0.20.6.21.2} is then exact by virtue of ($\mathbf{0_I}$, 6.1.2).
\end{enumerate}
  \item[\rm(ii)] Consider a commutative diagram of ring homomorphisms
\[
\xymatrix{
  \Lambda' \ar[r] & A' \ar[r] & B' \ar[r] & C' \\
  \Lambda \ar[u] \ar[r] & A \ar[u] \ar[r] & B \ar[u] \ar[r] & C \ar[u]
}
\]
Then the definitions of \sref{0.20.6.16} show that one has a commutative diagram (where the vertical arrows come from \sref{0.20.6.6.4})
\[
\xymatrix{
  0 \ar[r] & \mathrm{K}^C_\bullet(B/A/\Lambda) \ar[r] \ar[d] & \mathrm{K}'_\bullet(C/A/\Lambda) \ar[r] \ar[d] & \mathrm{K}_\bullet(C/B/\Lambda) \ar[r] \ar[d] & 0 \\
  0 \ar[r] & \mathrm{K}^{C'}_\bullet(B'/A'/\Lambda') \ar[r] & \mathrm{K}'_\bullet(C'/A'/\Lambda') \ar[r] & \mathrm{K}_\bullet(C'/B'/\Lambda') \ar[r] & 0
}
\]
\oldpage[0\textsubscript{IV-1}]{145}
whence, by passage to homology, a \emph{commutative diagram}
\begin{equation}
\tag{20.6.21.3}
  \label{0.20.6.21.3}
\xymatrix@C=0.7em{
  0 \ar[r] & \Upsilon^C_{B/A/\Lambda} \ar[r] \ar[d] & \Upsilon_{C/A/\Lambda} \ar[r] \ar[d] & \Upsilon_{C/B/\Lambda} \ar[r] \ar[d] & \Omega^1_{B/A}\otimes_B C \ar[r] \ar[d] & \Omega^1_{C/A} \ar[r] \ar[d] & \Omega^1_{C/B} \ar[r] \ar[d] & 0 \\
  0 \ar[r] & \Upsilon^{C'}_{B'/A'/\Lambda'} \ar[r] & \Upsilon_{C'/A'/\Lambda'} \ar[r] & \Upsilon_{C'/B'/\Lambda'} \ar[r] & \Omega^1_{B'/A'}\otimes_{B'} C' \ar[r] & \Omega^1_{C'/A'} \ar[r] & \Omega^1_{C'/B'} \ar[r] & 0
}
\end{equation}
One has an analogous commutative diagram for \sref{0.20.6.18.1}.
\end{enumerate}
\end{remarks}

\begin{proposition}[20.6.22]
\label{0.20.6.22}
Let $s:\Lambda\to A$, $u:A\to B$ be two ring homomorphisms, $\mathfrak{R}$ an ideal of $B$, $C$ the quotient ring $B/\mathfrak{R}$.
Suppose that $E=B/\mathfrak{R}^2$ is a $B$-extension $\Lambda$-trivial of $C$ by $\mathfrak{R}/\mathfrak{R}^2$.
One then has an exact sequence
\begin{equation}
\tag{20.6.22.1}
  \label{0.20.6.22.1}
  0\to\Upsilon^C_{B/A/\Lambda}\xrightarrow{v'}\Upsilon_{C/A/\Lambda}\xrightarrow{\chi_E}\mathfrak{R}/\mathfrak{R}^2\xrightarrow{\delta_{C/B/A}}\Omega^1_{B/A}\otimes_B C\xrightarrow{v_{C/B/A}}\Omega^1_{C/A}\to 0.
\end{equation}
\end{proposition}

\begin{proof}
Indeed, as $v:B\to C$ is surjective, one has $\Omega^1_{C/B}=0$ \sref{0.20.4.12}.
Moreover, it follows from \sref{0.20.6.10} that $\Upsilon_{C/B/\Lambda}$ is canonically identified with $\mathfrak{R}/\mathfrak{R}^2$.
It suffices then to apply the exact sequences \sref{0.20.6.17.1} and \sref{0.20.6.18.1}, noting that one has a commutative diagram
\[
\xymatrix{
  \Upsilon_{C/A/\Lambda} \ar[r] \ar[d]_{\chi_E} & \Omega^1_{A/\Lambda}\otimes_A C \ar[d] \\
  \Upsilon_{C/E/\Lambda}=\Upsilon_{C/B/\Lambda}=\mathfrak{R}/\mathfrak{R}^2 \ar[r]^-{\delta_{C/B/A}} \ar[d]_{\ell} & \Omega^1_{B/A}\otimes_B C=\Omega^1_{E/A}\otimes_E C \ar[d] \\
  \mathfrak{R}/\mathfrak{R}^2 \ar[r]^{\delta_{C/B/A}} & \Omega^1_{B/A}\otimes_B C
}
\]
and using the commutativity of the diagram \sref{0.20.6.17.3}.
\end{proof}

One has thus determined the \emph{kernel} of $\delta_{C/B/A}$ in the case where there exists a ring $\Lambda$ such that $A$ is a $\Lambda$-algebra and $B/\mathfrak{R}^2$ is $\Lambda$-trivial; one observes that this is the case in particular when $C$ is a \emph{$\Lambda$-algebra formally smooth} (for the discrete topology).

Suppose in addition that one has a commutative diagram of ring homomorphisms
\[
\xymatrix{
  \Lambda' \ar[r] & A' \ar[r] & B' \\
  \Lambda \ar[u] \ar[r] & A \ar[u] \ar[r] & B \ar[u]^f
}
\]
\oldpage[0\textsubscript{IV-1}]{146}
that $\mathfrak{R}'$ is an ideal of $B'$ such that $f(\mathfrak{R})\subset\mathfrak{R}'$, and that $E'=B'/\mathfrak{R}'^2$ is a $B'$-extension $\Lambda'$-trivial of $C'=B'/\mathfrak{R}'$ by $\mathfrak{R}'/\mathfrak{R}'^2$.
One then has a commutative diagram
\begin{equation}
\tag{20.6.22.2}
\label{0.20.6.22.2}
\xymatrix@C=0.7em{
  0 \ar[r] & \Upsilon^{C'}_{B'/A'/\Lambda'} \ar[r] & \Upsilon_{C'/A'/\Lambda'} \ar[r] & \mathfrak{R}'/\mathfrak{R}'^2 \ar[r] & \Omega^1_{B'/A'}\otimes_{B'} C' \ar[r] & \Omega^1_{C'/A'} \ar[r] & 0 \\
  0 \ar[r] & \Upsilon^C_{B/A/\Lambda} \ar[r] \ar[u] & \Upsilon_{C/A/\Lambda} \ar[r] \ar[u] & \mathfrak{R}/\mathfrak{R}^2 \ar[r] \ar[u] & \Omega^1_{B/A}\otimes_B C \ar[r] \ar[u] & \Omega^1_{C/A} \ar[r] \ar[u] & 0
}
\end{equation}
as follows from \sref{0.20.6.21.3} and \sref{0.20.5.11.3}.

\begin{corollary}[20.6.23]
\label{0.20.6.23}
Under the hypotheses of \sref{0.20.6.22}, suppose in addition that $B$ is an $A$-algebra formally smooth.
Then one has an exact sequence
\begin{equation}
\tag{20.6.23.1}
  \label{0.20.6.23.1}
  0\to\Upsilon_{C/A/\Lambda}\xrightarrow{\chi_E}\mathfrak{R}/\mathfrak{R}^2\xrightarrow{\delta_{C/B/A}}\Omega^1_{B/A}\otimes_B C\to\Omega^1_{C/A}\to 0.
\end{equation}
\end{corollary}

\begin{proof}
This follows from \sref{0.20.6.18}, (ii).
\end{proof}

\begin{env}[20.6.24]
\label{0.20.6.24}
When the hypotheses of \sref{0.20.6.22} are satisfied, we also say that the characteristic homomorphism $\chi_E$ is the \emph{characteristic homomorphism of the $A$-algebra $B$, relative to $\Lambda$ and to the ideal $\mathfrak{R}$} (suppressing these last details when there is no danger of confusion); it is sometimes denoted by $\chi_B$ or $\chi_{B/A}$.
\end{env}

\begin{proposition}[20.6.25]
\label{0.20.6.25}
Let $s:\Lambda\to A$, $u:A\to B$, $v:B\to C$ be three ring homomorphisms, $L$ a $C$-module.
One then has an exact sequence
\begin{equation}
\tag{20.6.25.1}
  \label{0.20.6.25.1}
\begin{split}
  0\to\Der_B(C,L)\xrightarrow{u^1}\Der_A(C,L)\xrightarrow{v^1}\Der_A(B,L)\xrightarrow{\partial} \\
  \to\operatorname{Exalcom}_{B/\Lambda}(C,L)\xrightarrow{u^1}\operatorname{Exalcom}_{A/\Lambda}(C,L)\xrightarrow{v^1}\operatorname{Exalcom}_{A/\Lambda}(B,L)\to 0
\end{split}
\end{equation}
where $u^1$, $v^1$ are homomorphisms defined in \sref{0.18.3.6.4} and \sref{0.18.3.6.2}, and $\partial$ is defined as in \sref{0.20.2.2}.
\end{proposition}

\begin{proof}
Indeed, as the exact sequence \sref{0.20.6.16.1} is split, one deduces from it an exact sequence
\[
  0\to\Hom_C(\mathrm{K}_\bullet(C/B/\Lambda),L)\to\Hom_C(\mathrm{K}'_\bullet(C/A/\Lambda),L)\to\Hom_C(\mathrm{K}^C_\bullet(B/A/\Lambda),L)\to 0.
\]
If one applies to this complex the long exact cohomology sequence, taking into account \sref{0.20.6.15.4} and \sref{0.20.6.6}, one obtains \sref{0.20.6.25.1} since one has
\[
  \Hom_C(\mathrm{K}^C_\bullet(B/A/\Lambda),L)=\Hom_B(\mathrm{K}_\bullet(B/A/\Lambda),L)
\]
by definition; the identification of $u^1$ and $v^1$ with the homomorphisms \sref{0.18.3.6.4} and \sref{0.18.3.6.2}, respectively, follows from \sref{0.20.6.6.4}.\footnote{The French source prints the nonexistent locator ``(18.3.4.2)''; the two homomorphisms are defined in (18.3.6.4) and (18.3.6.2).}
\end{proof}

\begin{remark}[20.6.26]
\label{0.20.6.26}
In this section, the complexes $\mathrm{K}_\bullet(C/B/A)$ have appeared as a technical device intended to simplify the exposition of certain functorial behaviours.
In reality, these complexes, considered as objects of the category of
\oldpage[0\textsubscript{IV-1}]{147}
complexes of $C$-modules ``up to homotopy'' (that is to say where the morphisms are the classes of homotopic homomorphisms) are remarkable invariants, finer than the pair formed by $\Omega^1_{C/B}$ and $\Upsilon_{C/B/A}$.
When $k$ is a prime field, and $C$ a $k$-algebra formally smooth (for example a regular ring of finite type over an extension of $k$) (cf.\ ($\mathbf{IV}$, 6.8.6)), one can show that the complex $\mathrm{K}_\bullet(C/A/k)$ can be described solely by making intervene $C$ and $A$ (to the exclusion of $k$): one expresses $C$ as a quotient of a polynomial algebra $B$ over $A$ by an ideal $\mathfrak{Q}$, and one considers the complex $\mathrm{F}_\bullet(C/A)$ with 2 nonzero terms
\[
  \cdots\to 0\to\mathfrak{Q}/\mathfrak{Q}^2\to\Omega^1_{B/A}\otimes_B C\to 0\to\cdots
\]
(whose homology coincides with that of $\mathrm{K}_\bullet(C/A/k)$ by virtue of \sref{0.20.6.23.1}).
These complexes $\mathrm{F}_\bullet(C/A)$, from the point of view of homological algebra, play the role of a conormal bundle for $\operatorname{Spec}(C)$ over $\operatorname{Spec}(A)$, and will be of great importance in the chapters of this work devoted to the duality of coherent sheaves and to the theorem of Riemann-Roch.
\end{remark}

\subsection{Generalizations to topological rings}
\label{subsection:0.20.7}

\begin{env}[20.7.1]
\label{0.20.7.1}
It follows immediately from the definitions that if, in \sref{0.20.5.2} and \sref{0.20.5.3}, the rings $A$, $B$, $C$ are assumed to be topological rings and the ring homomorphisms $u$, $v$ continuous, then the homomorphisms $u_{C/B/A}$ and $v_{C/B/A}$ are \emph{continuous} (one must naturally take on $\Omega^1_{B/A}\otimes_B C$ the tensor product topology).

Moreover, $u_{C/B/A}:\Omega^1_{C/A}\to\Omega^1_{C/B}$ is a \emph{strict surjective morphism} of topological $C$-modules; indeed, it is also so of the canonical homomorphism $C\otimes_A C\to C\otimes_B C$, taking into account the definition of the tensor product topology; by passing to quotients one deduces that the corresponding homomorphism $\mathrm{P}^1_{C/A}\to\mathrm{P}^1_{C/B}$ is a strict surjective morphism, and $u_{C/B/A}$ is the restriction of this last to $\Omega^1_{C/A}$.

In \sref{0.20.5.5}, if $A'$ and $B$ are topological $A$-algebras, and if $B'$ and $\Omega^1_{B/A}\otimes_B B'$ are equipped with the tensor product topologies, the homomorphism \sref{0.20.5.5.1} is a topological isomorphism, taking into account that $\mathrm{P}^1_{B/A}$ is the same as the direct topological sum of $B$ and $\Omega^1_{B/A}$.
\end{env}

\begin{proposition}[20.7.2]
\label{0.20.7.2}
Let $u:A\to B$, $v:B\to C$ be two continuous homomorphisms of topological rings.
For the continuous homomorphism $v_{C/B/A}:\Omega^1_{B/A}\otimes_B C\to\Omega^1_{C/A}$ to be formally left-invertible \textup{(}cf.\ \sref{0.19.1.5}\textup{)}, it is necessary and sufficient that $C$ be a $B$-algebra formally smooth relative to $A$ \sref{0.19.9.1} \textup{(}and \emph{a fortiori} it suffices that $C$ be a $B$-algebra formally smooth\textup{)}.
\end{proposition}

\begin{proof}
To say that $v_{C/B/A}$ is formally left-invertible means, given that the topologies of $\Omega^1_{C/A}$ and $\Omega^1_{B/A}\otimes_B C$ are less fine than those deduced from the topology of $C$ \sref{0.20.4.5}, that for every discrete $C$-module $L$, annihilated by an open ideal of $C$, the canonical homomorphism
\[
  \Hom.\operatorname{cont}_C(\Omega^1_{C/A},L)\to\Hom.\operatorname{cont}_C(\Omega^1_{B/A}\otimes_B C,L)
\]
\oldpage[0\textsubscript{IV-1}]{148}
is \emph{surjective} \sref{0.19.1.5}; as $\Hom.\operatorname{cont}_C(\Omega^1_{B/A}\otimes_B C,L)=\Hom.\operatorname{cont}_C(\Omega^1_{B/A},L)$ by definition of the tensor product topology, it amounts to the same, by virtue of \sref{0.20.4.8}, to say that the canonical homomorphism
\[
  \operatorname{D\acute{e}r.\ cont}_A(C,L)\to\operatorname{D\acute{e}r.\ cont}_A(B,L)
\]
is surjective.
But the exact sequence \sref{0.20.3.8.2} shows that this condition is equivalent to $\operatorname{Exalcotop}_{B/A}(C,L)=0$, that is to say precisely that $C$ is formally smooth relative to $A$ \sref{0.19.9.8}.
\end{proof}

\begin{corollary}[20.7.3]
\label{0.20.7.3}
Suppose that in $B$ and in $C$, the square of every open ideal is also open.
For $C$ to be a $B$-algebra formally smooth relative to $A$, it is necessary and sufficient that, if one denotes by $(\mathfrak{R}_\lambda)$ a fundamental system of neighbourhoods of $0$ formed of ideals of $C$, then for every $\lambda$, the homomorphism
\begin{equation}
\tag{20.7.3.1}
  \label{0.20.7.3.1}
  v_{C/B/A}\otimes 1_{C/\mathfrak{R}_\lambda}:\Omega^1_{B/A}\otimes_B(C/\mathfrak{R}_\lambda)\to\Omega^1_{C/A}\otimes_C(C/\mathfrak{R}_\lambda)
\end{equation}
is left-invertible.
\end{corollary}

\begin{proof}
Indeed, one then knows that the topology of $\Omega^1_{B/A}$ (respectively $\Omega^1_{C/A}$) is deduced from that of $B$ (respectively $C$) \sref{0.20.4.5}; it follows immediately that the topology of $\Omega^1_{B/A}\otimes_B C$ is also deduced from that of $C$. The corollary then follows from \sref{0.20.7.2} and \sref{0.19.1.7}.
\end{proof}

\begin{proposition}[20.7.4]
\label{0.20.7.4}
Let $A$ be a topological ring, $B$ a topological $A$-algebra.
For $B$ to be formally non-ramified \sref{0.19.10.2}, it is necessary and sufficient that the separated completion $\widehat{\Omega}^1_{B/A}$ is zero.
\end{proposition}

\begin{proof}
Indeed, it follows immediately from \sref{0.19.10.4} and \sref{0.20.1.1} that, for $B$ to be formally non-ramified, it is necessary and sufficient that for every open ideal $\mathfrak{R}$ of $B$, every open ideal $\mathfrak{H}$ of $A$ such that $\mathfrak{H}.B\subset\mathfrak{R}$ and every $(B/\mathfrak{R})$-module $L$, one has $\Der_{A/\mathfrak{H}}(B/\mathfrak{R},L)=0$, that is to say $\operatorname{D\acute{e}r.\ cont}_A(B,L)=0$ for every discrete $B$-module $L$ annihilated by an open ideal; by virtue of \sref{0.20.4.8}, this is equivalent to $\Hom.\operatorname{cont}_B(\Omega^1_{B/A},L)=0$ for every such $B$-module, whence the proposition.

When $B$ is discrete, the condition of \sref{0.20.7.4} is equivalent to $\Omega^1_{B/A}=0$.
\end{proof}

\begin{corollary}[20.7.5]
\label{0.20.7.5}
Let $A$ be a ring, $\mathfrak{m}$ an ideal of $A$, and $B$ an $A$-algebra; equip $A$ with the $\mathfrak{m}$-preadic topology and $B$ with the topology deduced from that of $A$, and set $A_0=A/\mathfrak{m}$, $B_0=B/\mathfrak{m}B=B\otimes_A A_0$.
Then, for $B$ to be an $A$-algebra formally non-ramified, it is necessary and sufficient that $\Omega^1_{B_0/A_0}=0$ \textup{(}or equivalently that $B_0$ is an $A_0$-algebra formally non-ramified\textup{)}.
\end{corollary}

\begin{proof}
Indeed, one has $\Omega^1_{B/A}\otimes_B B_0=\Omega^1_{B/A}/\mathfrak{m}.\Omega^1_{B/A}=\Omega^1_{B_0/A_0}$ by virtue of \sref{0.20.5.5}; to say that every open submodule of $\Omega^1_{B/A}$ is equal to $\Omega^1_{B/A}$ is equivalent by \sref{0.20.4.5} to writing that $\mathfrak{m}.\Omega^1_{B/A}=\Omega^1_{B/A}$, whence the conclusion.
\end{proof}

\begin{proposition}[20.7.6]
\label{0.20.7.6}
Let $u:A\to B$, $v:B\to C$ be two continuous homomorphisms of topological rings, and suppose that $v$ makes $C$ into a $B$-algebra formally \'{e}tale; then $v_{C/B/A}:\Omega^1_{B/A}\otimes_B C\to\Omega^1_{C/A}$ is a formal bimorphism \sref{0.19.1.2}.
\end{proposition}

\begin{proof}
Indeed, for every discrete $C$-module $L$ annihilated by an open ideal of $C$, one has then \sref{0.20.7.4} $\operatorname{D\acute{e}r.\ cont}_B(C,L)=0$ and by virtue of \sref{0.20.3.6.1} the canonical homomorphism
\oldpage[0\textsubscript{IV-1}]{149}
$\operatorname{D\acute{e}r.\ cont}_A(C,L)\to\operatorname{D\acute{e}r.\ cont}_A(B,L)$ is injective; it is moreover surjective by virtue of \sref{0.20.7.2}, hence it is bijective. It follows that the image of $v_{C/B/A}$ is necessarily dense in $\Omega^1_{C/A}$, otherwise the quotient of $\Omega^1_{C/A}$ by the closure of $\operatorname{Im}(v_{C/B/A})$ would be separated and nonzero and would therefore have a nonzero discrete quotient $L$, contrary to what has just been seen (taking into account \sref{0.20.4.8}). Consequently $v_{C/B/A}$, which is a formal monomorphism by virtue of \sref{0.20.7.2}, is also a formal epimorphism \sref{0.19.1.2}, hence a formal bimorphism.
\end{proof}

\begin{corollary}[20.7.7]
\label{0.20.7.7}
Suppose that in $B$ and in $C$, the square of every open ideal is also open.
If $C$ is a $B$-algebra formally \'{e}tale, then, for every open ideal $\mathfrak{R}_\lambda$ of $C$, the homomorphism \sref{0.20.7.3.1} is bijective.
\end{corollary}

\begin{proof}
Indeed, by virtue of \sref{0.19.1.1}, this homomorphism is surjective, and it is injective by \sref{0.20.7.3}.
\end{proof}

\begin{proposition}[20.7.8]
\label{0.20.7.8}
Let $u:A\to B$ be a continuous homomorphism of topological rings, $\mathfrak{R}$ an ideal of $B$, $C$ the topological quotient $B/\mathfrak{R}$, $v:B\to C$ the canonical homomorphism.
\begin{enumerate}
  \item[\rm(i)] In the exact sequence \sref{0.20.5.12.1}
\[
  \mathfrak{R}/\mathfrak{R}^2\xrightarrow{\delta_{C/B/A}}\Omega^1_{B/A}\otimes_B C\xrightarrow{v_{C/B/A}}\Omega^1_{C/A}\to 0
\]
the homomorphism $\delta_{C/B/A}$ is continuous and the homomorphism $v_{C/B/A}$ is a strict morphism of topological $C$-modules.
  \item[\rm(ii)] For $\delta_{C/B/A}$ to be formally left-invertible \sref{0.19.1.5}, it is necessary and sufficient that for every discrete $C$-module $L$ annihilated by an open ideal of $C$, the canonical homomorphism
\begin{equation}
\tag{20.7.8.1}
  \label{0.20.7.8.1}
  \operatorname{Exalcotop}_A(C,L)\to\operatorname{Exalcotop}_A(B,L)
\end{equation}
is injective.
\end{enumerate}
\end{proposition}

\begin{proof}
(i) The first assertion is evident.
On the other hand, the canonical homomorphism $v\otimes v:B\otimes_A B\to C\otimes_A C$ is a strict morphism by definition of the tensor product topology, and one deduces immediately that it is the same for $v_{C/B/A}$.

(ii) To say that $\delta_{C/B/A}$ is formally left-invertible means that for every discrete $C$-module $L$ annihilated by an open ideal of $C$, the canonical homomorphism
\[
  \Hom.\operatorname{cont}_C(\Omega^1_{B/A}\otimes_B C,L)\to\Hom.\operatorname{cont}_C(\mathfrak{R}/\mathfrak{R}^2,L)
\]
is surjective.
But, taking into account \sref{0.18.4.3} and \sref{0.20.4.8}, this amounts to saying that the canonical homomorphism
\[
  \operatorname{D\acute{e}r.\ cont}_A(B,L)\to\operatorname{Exalcotop}_B(C,L)
\]
is surjective, and the conclusion follows from the exact sequence \sref{0.20.3.6.1}.

The injectivity of \sref{0.20.7.8.1} can also be expressed as follows, taking into account the definition of its two members \sref{0.18.4.1}: for every open ideal $\mathfrak{M}$ of $A$, every open ideal $\mathfrak{N}$ of $B$ such that $\mathfrak{M}B\subset\mathfrak{N}$, and every $(A/\mathfrak{M})$-extension $E$ of $B/(\mathfrak{R}+\mathfrak{N})$ by a $(B/(\mathfrak{R}+\mathfrak{N}))$-module $L$, such that the inverse image of $E$ under the canonical homomorphism $B/\mathfrak{N}\to B/(\mathfrak{R}+\mathfrak{N})$ is $(A/\mathfrak{M})$-trivial, there exists an open ideal $\mathfrak{M}'\subset\mathfrak{M}$ of $A$,
\oldpage[0\textsubscript{IV-1}]{150}
an open ideal $\mathfrak{N}'\subset\mathfrak{N}$ of $B$ such that $\mathfrak{M}'B\subset\mathfrak{N}'$ and such that the inverse image of $E$ under the canonical homomorphism $B/(\mathfrak{R}+\mathfrak{N}')\to B/(\mathfrak{R}+\mathfrak{N})$ is $(A/\mathfrak{M}')$-trivial.
\end{proof}

\begin{corollary}[20.7.9]
\label{0.20.7.9}
If the topological $A$-algebra $C=B/\mathfrak{R}$ is formally smooth, the canonical homomorphism $\delta_{C/B/A}$ is formally left-invertible.
\end{corollary}

\begin{env}[20.7.10]
\label{0.20.7.10}
In \sref{0.20.6.1}, when $A$, $B$, $C$ are topological rings and $u$ and $v$ continuous homomorphisms, one equips $\Upsilon_{C/B/A}$ with the topology induced by that of $\Omega^1_{B/A}\otimes_B C$; the homomorphisms \sref{0.20.6.4.2} and \sref{0.20.6.4.4} are then continuous provided the same is true of those in diagram \sref{0.20.6.4.1}. Moreover, if in \sref{0.20.6.7} one assumes that $L$ is a discrete $C$-module annihilated by an open ideal of $C$, one deduces, by passage to the inductive limit, the canonical $C$-homomorphism
\begin{equation}
\tag{20.7.10.1}
  \label{0.20.7.10.1}
  \operatorname{Exalcotop}_{B/A}(C,L)\to\Hom.\operatorname{cont}_C(\Upsilon_{C/B/A},L)
\end{equation}
which can be described explicitly as follows (taking into account \sref{0.18.5.3.1}): for every open ideal $\mathfrak{M}$ of $A$, every open ideal $\mathfrak{R}$ of $B$ such that $\mathfrak{M}B\subset\mathfrak{R}$, every open ideal $\mathfrak{P}$ of $C$ such that $\mathfrak{R}C\subset\mathfrak{P}$ and $\mathfrak{P}.L=0$, every $(B/\mathfrak{R})$-extension $E$ of $C/\mathfrak{P}$ by $L$ which is $(A/\mathfrak{M})$-trivial comes from the datum of a continuous $C$-homomorphism $f$ of $\Omega^1_{B/A}\otimes_B C$ in $L$, and the homomorphism \sref{0.20.7.10.1} associates to the class of $E$ the restriction of $f$ to $\Upsilon_{C/B/A}$, called the \emph{characteristic homomorphism of $E$} and again denoted $\chi_E$.
\end{env}

\begin{proposition}[20.7.11]
\label{0.20.7.11}
Suppose that the topology of $C$ is such that the square of every open ideal is also open.
If $C$ is a topological $A$-algebra formally smooth and if $\Omega^1_{C/B}$ is a formally projective $C$-module, one has a canonical isomorphism
\begin{equation}
\tag{20.7.11.1}
  \label{0.20.7.11.1}
  \operatorname{Exalcotop}_B(C,L)\xrightarrow{\sim}\Hom.\operatorname{cont}_C(\Upsilon_{C/B/A},L)
\end{equation}
for every discrete $C$-module $L$ annihilated by an open ideal of $C$.
\end{proposition}

\begin{proof}
One has indeed, in the exact sequence \sref{0.20.3.7.1}, $\operatorname{Exalcotop}_A(C,L)=0$ \sref{0.19.4.4}, hence $\operatorname{Exalcotop}_B(C,L)=\operatorname{Exalcotop}_{B/A}(C,L)$ and the homomorphism \sref{0.20.7.10.1} is none other than \sref{0.20.7.8.1}; the fact that it is bijective is deduced from \sref{0.20.6.13} by passage to the inductive limit by \sref{0.20.4.5} and \sref{0.19.2.4}.
\end{proof}

\begin{env}[20.7.12]
\label{0.20.7.12}
In \sref{0.20.6.14} one also equips $\Upsilon^C_{B/A/\Lambda}$ with the topology induced by that of $\Omega^1_{A/\Lambda}\otimes_A C$ and the homomorphism \sref{0.20.6.14.6} is continuous, when the considered rings are topological and the homomorphisms of rings are continuous.
\end{env}

\begin{env}[20.7.13]
\label{0.20.7.13}
If, in \sref{0.20.6.23}, one assumes that the rings $\Lambda$, $A$, $B$, $C$ are topological, the homomorphisms $s$, $u$, $v$ continuous and $L$ a discrete $C$-module annihilated by an open ideal of $C$, one can pass to the inductive limit as in \sref{0.20.3.6}, and one obtains an exact sequence
\begin{equation}
\tag{20.7.13.1}
  \label{0.20.7.13.1}
\begin{split}
  0\to\operatorname{D\acute{e}r.\ cont}_B(C,L)\to\operatorname{D\acute{e}r.\ cont}_A(C,L)\to\operatorname{D\acute{e}r.\ cont}_A(B,L)\to \\
  \to\operatorname{Exalcotop}_{B/\Lambda}(C,L)\to\operatorname{Exalcotop}_{A/\Lambda}(C,L)\to\operatorname{Exalcotop}_{A/\Lambda}(B,L)\to 0.
\end{split}
\end{equation}
\end{env}

\begin{env}[20.7.14]
\label{0.20.7.14}
Let $A$ be a topological ring, $B$ a topological $A$-algebra, $\mathfrak{M}'$ an open ideal of $A$, $\mathfrak{R}'$ an open ideal of $B$ such that $\mathfrak{M}'B\subset\mathfrak{R}'$; set $A'=A/\mathfrak{M}'$, $B'=B/\mathfrak{R}'$;
\oldpage[0\textsubscript{IV-1}]{151}
the kernel of the homomorphism $B\otimes_A B\to B'\otimes_{A'} B'$ is $\mathfrak{U}=\operatorname{Im}(\mathfrak{R}'\otimes B+B\otimes\mathfrak{R}')$, whence it follows immediately that the kernel of the homomorphism
\begin{equation}
\tag{20.7.14.1}
  \label{0.20.7.14.1}
  \varphi_{(\mathfrak{M}',\mathfrak{R}')}:\Omega^1_{B/A}\to\Omega^1_{B'/A'}
\end{equation}
is $((\mathfrak{J}\cap\mathfrak{U})+\mathfrak{J}^2)/\mathfrak{J}^2$; moreover, the homomorphism \sref{0.20.7.14.1} is \emph{surjective}, as it follows from \sref{0.20.4.7}.
If $\mathfrak{M}''$ (resp.\ $\mathfrak{R}''$) is a second open ideal of $A$ (resp.\ $B$) such that $\mathfrak{M}''\subset\mathfrak{M}'$, $\mathfrak{R}''\subset\mathfrak{R}'$ and $\mathfrak{M}''B\subset\mathfrak{R}''$, and if one sets $A''=A/\mathfrak{M}''$, $B''=B/\mathfrak{R}''$, one has a surjective homomorphism
\[
  \varphi_{(\mathfrak{M}'',\mathfrak{R}''),({\mathfrak{M}',\mathfrak{R}'})}:\Omega^1_{B''/A''}\to\Omega^1_{B'/A'}
\]
and these homomorphisms form a \emph{projective system}.
If one notices that the $((\mathfrak{J}\cap\mathfrak{U})+\mathfrak{J}^2)/\mathfrak{J}^2$ form a fundamental system of neighbourhoods of $0$ in $\Omega^1_{B/A}$, one sees therefore that the \emph{separated completion} $\widehat{\Omega}^1_{B/A}$ of the topological $B$-module $\Omega^1_{B/A}$ is given, up to a canonical isomorphism, by
\begin{equation}
\tag{20.7.14.2}
  \label{0.20.7.14.2}
  \widehat{\Omega}^1_{B/A}=\varprojlim(\Omega^1_{B'/A'}).
\end{equation}
Moreover, the canonical homomorphism $j:\Omega^1_{B/A}\to\widehat{\Omega}^1_{B/A}$ is the projective limit of the $\varphi_{(\mathfrak{M}',\mathfrak{R}')}$, hence $\varphi_{(\mathfrak{M}',\mathfrak{R}')}$ factors as $\Omega^1_{B/A}\xrightarrow{j}\widehat{\Omega}^1_{B/A}\to\Omega^1_{B'/A'}$, and as it is surjective one concludes that the canonical homomorphism
\begin{equation}
\tag{20.7.14.3}
  \label{0.20.7.14.3}
  \widehat{\Omega}^1_{B/A}\to\Omega^1_{B'/A'}
\end{equation}
is \emph{surjective} for every pair $(\mathfrak{M}',\mathfrak{R}')$.
One can moreover replace $A'$ by $A$ everywhere in the above.

Finally, if $L$ is a separated and \emph{complete} topological $B$-module, every continuous $B$-homomorphism from $\Omega^1_{B/A}$ to $L$ extends uniquely to a continuous $\widehat{B}$-homomorphism from $\widehat{\Omega}^1_{B/A}$ to $L$ and conversely such a homomorphism gives by composition with $\Omega^1_{B/A}\to\widehat{\Omega}^1_{B/A}$ a continuous $B$-homomorphism, so that one has a canonical isomorphism
\[
  \Hom.\operatorname{cont}_B(\Omega^1_{B/A},L)\xrightarrow{\sim}\Hom.\operatorname{cont}_{\widehat{B}}(\widehat{\Omega}^1_{B/A},L).
\]
More particularly, if $L$ is a discrete $B$-module annihilated by an open ideal of $B$, one sees that the canonical isomorphism \sref{0.20.4.8.2} can also be written
\begin{equation}
\tag{20.7.14.4}
  \label{0.20.7.14.4}
  \Hom.\operatorname{cont}_{\widehat{B}}(\widehat{\Omega}^1_{B/A},L)\xrightarrow{\sim}\operatorname{D\acute{e}r.\ cont}_A(B,L).
\end{equation}
\end{env}

\begin{proposition}[20.7.15]
\label{0.20.7.15}
Let $A$ be a discrete ring, $B$ a topological adic $A$-algebra ($\mathbf{0_I}$, 7.1.9), $\mathfrak{n}$ an ideal of definition of $B$, $B_0=B/\mathfrak{n}$.
Suppose that $\Omega^1_{B_0/A}$ and $\mathfrak{n}/\mathfrak{n}^2$ are $B_0$-modules of finite type.
Then $\widehat{\Omega}^1_{B/A}$ is a $\widehat{B}$-module of finite type.
\end{proposition}

\begin{proof}
As the square of every open ideal of $B$ is open, the topology of $\Omega^1_{B/A}$ is the $\mathfrak{n}$-preadic topology \sref{0.20.4.5}.
Taking into account the hypothesis that $B$ is an adic ring, it therefore suffices, by virtue of ($\mathbf{0_I}$, 7.2.7 and 7.2.9), to prove
\oldpage[0\textsubscript{IV-1}]{152}
that $\Omega^1_{B/A}/\mathfrak{n}\Omega^1_{B/A}=\Omega^1_{B/A}\otimes_B B_0$ is a $B_0$-module of finite type. But this follows from the hypothesis and from the exact sequence \sref{0.20.5.12.1}
\[
  \mathfrak{n}/\mathfrak{n}^2\longrightarrow \Omega^1_{B/A}\otimes_B B_0\longrightarrow \Omega^1_{B_0/A}\longrightarrow 0.
\]
\end{proof}

\begin{env}[20.7.16]
\label{0.20.7.16}
Proposition \sref{0.20.7.15} applies, for example, when $A$ is a field $k$, $B=k'[[T_1,\ldots,T_n]]$ is a formal power series algebra equipped with its usual topology, and $k'$ is a finite extension of $k$ (cf. \sref{0.21.9.2}). One will note that the field of fractions $K$ of $B$ has infinite transcendence degree over $k$; when $k$ has characteristic $0$, it follows immediately (with the aid, in particular, of \sref{0.20.4.13}, (i), and \sref{0.20.5.9}, also using the fact that every derivation of a field of characteristic $0$ extends to every extension) that $\Omega^1_{B/k}$ is \emph{not} a $B$-module of finite type.
\end{env}

\begin{env}[20.7.17]
\label{0.20.7.17}
Let $A$, $B$, $C$ be three topological rings, and let $u:A\to B$, $v:B\to C$ be two continuous homomorphisms. Replacing $A$, $B$, $C$ by quotients by open ideals $A'=A/\mathfrak{M}'$, $B'=B/\mathfrak{N}'$, $C'=C/\mathfrak{P}'$, with $u(\mathfrak{M}')\subset\mathfrak{N}'$ and $v(\mathfrak{N}')\subset\mathfrak{P}'$, so that one has homomorphisms $u':A'\to B'$, $v':B'\to C'$, one obtains canonical homomorphisms $u_{C'/B'/A'}$ and $v_{C'/B'/A'}$ which, by virtue of \sref{0.20.5.4}, form projective systems and consequently give, by passage to the limit, canonical homomorphisms extending to the separated completions the homomorphisms of the exact sequence \sref{0.20.5.7.1}
\[
  \Omega^1_{B/A}\otimes_B C\xrightarrow{v_{C/B/A}}\Omega^1_{C/A}\xrightarrow{u_{C/B/A}}\Omega^1_{C/B}\longrightarrow 0:
\]
\begin{align}
  \widehat{v}_{C/B/A}:\widehat{\Omega}^1_{B/A}\widehat{\otimes}_{\widehat{B}}\widehat{C}&\longrightarrow\widehat{\Omega}^1_{C/A}, \tag{20.7.17.1}\label{0.20.7.17.1}\\
  \widehat{u}_{C/B/A}:\widehat{\Omega}^1_{C/A}&\longrightarrow\widehat{\Omega}^1_{C/B}. \tag{20.7.17.2}\label{0.20.7.17.2}
\end{align}
In the sequence
\begin{equation}
\tag{20.7.17.3}
\label{0.20.7.17.3}
  \widehat{\Omega}^1_{B/A}\widehat{\otimes}_{\widehat{B}}\widehat{C}
  \xrightarrow{\widehat{v}_{C/B/A}}\widehat{\Omega}^1_{C/A}
  \xrightarrow{\widehat{u}_{C/B/A}}\widehat{\Omega}^1_{C/B}\longrightarrow 0,
\end{equation}
the composite of two consecutive homomorphisms is $0$, but the sequence is \emph{not necessarily exact}. However, if $B$ and $C$ are metrizable, the homomorphism $\widehat{u}_{C/B/A}$ is \emph{surjective}, and $\operatorname{Im}(\widehat{v}_{C/B/A})$ is \emph{dense} in $\ker(\widehat{u}_{C/B/A})$: this follows immediately (cf. $\mathbf{0_I}$, 7.3.1) from the fact that, if $(\mathfrak{M}_k)$ (resp. $(\mathfrak{N}_k)$) is a decreasing sequence of ideals of $B$ (resp. $C$) forming a fundamental system of neighbourhoods of $0$, and if one sets $B_k=B/\mathfrak{M}_k$, $C_k=C/\mathfrak{N}_k$, the transition homomorphisms $\Omega^1_{C_{k+1}/B_{k+1}}\to\Omega^1_{C_k/B_k}$, $\Omega^1_{C_{k+1}/A}\to\Omega^1_{C_k/A}$ and $\Omega^1_{B_{k+1}/A}\to\Omega^1_{B_k/A}$ are surjective \sref{0.20.7.14}.
\end{env}

\begin{proposition}[20.7.18]
\label{0.20.7.18}
Let $A$, $B$, $C$ be three topological rings, and let $u:A\to B$, $v:B\to C$ be two continuous homomorphisms. Suppose that $B$ and $C$ are admissible ($\mathbf{0_I}$, 7.1.2) and metrizable. For the canonical homomorphism
\[
  \widehat{v}_{C/B/A}:\widehat{\Omega}^1_{B/A}\widehat{\otimes}_{\widehat{B}}\widehat{C}\longrightarrow\widehat{\Omega}^1_{C/A}
\]
to admit a left inverse which is a continuous $C$-homomorphism, it is necessary and sufficient that $C$ be a $B$-algebra formally smooth relative to $A$.
\end{proposition}

\begin{proof}
The condition is necessary by virtue of \sref{0.20.7.2} and \sref{0.19.1.6}. To see that it is sufficient, note that the topological $C$-module
\[
  L=\widehat{\Omega}^1_{B/A}\widehat{\otimes}_{\widehat{B}}\widehat{C}
\]
is metrizable and complete by virtue of the hypotheses, hence $E=D_C(L)$, equipped with the product topology, is
\oldpage[0\textsubscript{IV-1}]{153}
metrizable and complete; it is moreover admissible, for if $\mathfrak{R}$ is an ideal of definition in $C$, the sequence $(\mathfrak{R}\oplus L)^n=\mathfrak{R}^n\oplus\mathfrak{R}^{n-1}L$ tends to $0$. Since the composite map
\[
  D:B\xrightarrow{d_{B/A}}\Omega^1_{B/A}\longrightarrow
  \widehat{\Omega}^1_{B/A}\widehat{\otimes}_{\widehat{B}}\widehat{C}=L
\]
is a continuous $A$-derivation of $B$ into $L$, the continuous $A$-homomorphism
\[
  f:x\longmapsto(v(x),D(x))
\]
from $B$ into $E$ defines on $E$ a structure of $B$-extension. Since $L$ is a closed ideal in $E$, it follows from \sref{0.19.9.5} and from the hypothesis that the identity map $C\to E/L$ (which is a $B$-homomorphism) factors as $C\xrightarrow{g}E\to E/L$, where $g$ is a continuous homomorphism such that $g\circ v=f$; consequently \sref{0.20.1.3}, $g$ is of the form $y\mapsto(y,D'(y))$, where $D'$ is a continuous $B$-derivation of $C$ into $L$, in other words $D'\circ v=D$. Taking into account \sref{0.20.7.14.4}, one has $D'=h\circ\widehat{d}_{C/A}$, where $h:\widehat{\Omega}^1_{C/A}\to L$ is a continuous $C$-homomorphism; but one has $\widehat{d}_{C/A}\circ v=\widehat{v}_{C/B/A}\circ D$ by definition, and since the image of $B$ under $D$ generates topologically the $C$-module $L$ \sref{0.20.4.7}, the relation $D'\circ v=D$ gives $h\circ\widehat{v}_{C/B/A}=1_L$.
\end{proof}

\begin{corollary}[20.7.19]
\label{0.20.7.19}
Under the hypotheses of \sref{0.20.7.18}, if one supposes in addition that in $B$ and $C$ the square of every open ideal is open, then, for $C$ to be a $B$-algebra formally smooth relative to $A$, it is necessary and sufficient that $\widehat{v}_{C/B/A}$ be left-invertible.
\end{corollary}

\begin{proof}
Indeed, the topologies of $\widehat{\Omega}^1_{B/A}\widehat{\otimes}_{\widehat{B}}\widehat{C}$ and $\widehat{\Omega}^1_{C/A}$ are then deduced from that of $C$ \sref{0.20.4.5}, and every $C$-homomorphism from one to the other is necessarily continuous.
\end{proof}

\begin{env}[20.7.20]
\label{0.20.7.20}
Let $A$ be a topological ring, $B$ a topological $A$-algebra, metrizable and complete, $\mathfrak{R}$ a closed ideal of $B$, and $C=B/\mathfrak{R}$ the quotient topological ring, which is metrizable and complete. Let $(\mathfrak{M}_k)$ be a decreasing fundamental system of neighbourhoods of $0$ in $B$ formed of ideals, and set $B_k=B/\mathfrak{M}_k$, $\mathfrak{R}_k=(\mathfrak{R}+\mathfrak{M}_k)/\mathfrak{M}_k$, $C_k=B_k/\mathfrak{R}_k$. One has a projective system of homomorphisms
\[
  \delta_{C_k/B_k/A}:\mathfrak{R}_k/\mathfrak{R}_k^2\longrightarrow\Omega^1_{B_k/A}\otimes_{B_k}C_k
\]
\sref{0.20.5.11.3}, from which one obtains by passage to the limit a canonical homomorphism
\[
  \widehat{\delta}_{C/B/A}:\mathfrak{R}/\mathfrak{R}^2\longrightarrow
  \widehat{\Omega}^1_{B/A}\widehat{\otimes}_{\widehat{B}}\widehat{C};
\]
and, reasoning as in \sref{0.20.7.17}, one sees that the canonical homomorphism
\[
  \widehat{u}_{C/B/A}:\widehat{\Omega}^1_{B/A}\widehat{\otimes}_{\widehat{B}}\widehat{C}\longrightarrow\widehat{\Omega}^1_{C/A}
\]
is \emph{surjective} and that $\operatorname{Im}(\widehat{\delta}_{C/B/A})$ is \emph{dense} in $\ker(\widehat{u}_{C/B/A})$.
\end{env}
% END NATIVE INPUT ega0/ega0-20
% BEGIN NATIVE INPUT ega0/ega0-21
\section{Differentials in rings of characteristic \texorpdfstring{$p$}{p}}
\label{section:0.21}
\label{0.21}

The results of the present section, of a more specialized and technical nature than those of \S\S\,19, 20 and 22, will only serve occasionally in chapter~IV.
Their principal role here is in the proofs of three theorems of \S\,22 (22.3.3, 22.5.8 and 22.7.3), the first and last of which intervene in an essential way in the ``fine'' theory of Noetherian local rings of chapter~IV, \S\,7.

\oldpage[0\textsubscript{IV-1}]{154}
\subsection{Systems of \texorpdfstring{$p$}{p}-generators and \texorpdfstring{$p$}{p}-bases}
\label{subsection:0.21.1}

\begin{env}[21.1.1]
\label{0.21.1.1}
Given a number $p$ which is either $0$, or a prime number, we say that a ring $A$ is \emph{of characteristic $p$} if there exists a ring homomorphism $P\to A$, where $P$ is the \emph{prime field} of characteristic $p$; we note that this homomorphism is then unique, the composite $\bb{Z}\to P\to A$ being the unique homomorphism of $\bb{Z}$ into $A$.
If $A\neq 0$, this amounts to saying that $A$ \emph{contains a field of characteristic $p$}, the image of $P$ being necessarily the only subfield of $A$ isomorphic to $P$ (and moreover the only subfield of $A$ isomorphic to $P$).
\end{env}

\begin{env}[21.1.2]
\label{0.21.1.2}
If $p>0$, to say that $A$ is of characteristic $p$ is equivalent to saying that, in $A$, $p.1=0$, or equivalently $pA=0$.
If $p=0$, to say that $A$ is of characteristic $p$ is equivalent to saying that for every integer $n\neq 0$, $n.1$ is \emph{invertible} in $A$.
If $A\neq 0$, there can exist only one $p$ (prime or $0$) such that $A$ is of characteristic $p$; this follows from the identity of Bezout $ap+bq=1$ for two distinct primes $p$, $q$.
On the other hand, the reduced ring $0$ is of characteristic $p$ for \emph{every} $p$.
\end{env}

\begin{env}[21.1.3]
\label{0.21.1.3}
If $A$ is of characteristic $p$, then so is every \emph{algebra} over $A$.
In particular, for every prime ideal $\mathfrak{p}$ of $A$, the \emph{residue field} of $A$ at $\mathfrak{p}$ is of characteristic $p$.
Conversely, if $p=0$ and if for every \emph{maximal} ideal $\mathfrak{m}$ of $A$, the residue field of $A$ at $\mathfrak{m}$ is of characteristic $0$, then so is $A$, since for every integer $n\neq 0$, $n.1$ is then invertible in all the $A_\mathfrak{m}$, hence also in $A$.
On the other hand, if $p>0$, a local ring can have its residue field of characteristic $p$ without being itself of characteristic $p$, as shown by the example of the \emph{prime integer ring} $\bb{Z}_{p\bb{Z}}$ \sref{0.19.8.3} or of the local Artinian ring $\bb{Z}/p^n\bb{Z}$ for $n\geq 2$, which do not contain a field.

We note finally that for a ring (even reduced), the residue fields at its prime ideals can have different characteristics, as shown by the example of $\bb{Z}$.
\end{env}

\begin{env}[21.1.4]
\label{0.21.1.4}
In all the rest of this section, we fix a \emph{prime number $p$} and all rings will be assumed \emph{of characteristic $p$}, unless express mention to the contrary.
For such a ring $A$, the map $x\mapsto x^p$ is an \emph{endomorphism} of $A$, which we denote $\mathrm{F}_A$; if $A$ is \emph{reduced}, $\mathrm{F}_A$ is \emph{injective}.
We set $A^p=\mathrm{F}_A(A)$ (subring of $A$ formed by the $x^p$ for $x\in A$); one can naturally consider $A$ as an $A^p$-algebra.

One can also consider $A$ as an $A$-algebra \emph{by means of the homomorphism} $\mathrm{F}_A:x\mapsto x^p$ of $A$ into $A$; in other words, this is the $A$-algebra $A_{[\mathrm{F}_A]}$, for which the product $\lambda.x$ of $x\in A$ by a scalar $\lambda\in A$ is the product $\lambda^p x$ in the ring $A$; we denote this $A$-algebra $A^{(p)}$.
It is clear that for every ring homomorphism $u:A\to B$ the pair $(u,u)$ is a \emph{di-homomorphism} of $A$-algebras $A^{(p)}\to B^{(p)}$.
For every $A$-module $E$, we set $E^{(p)}=E\otimes_A A^{(p)}$ where $A^{(p)}$ is considered as an $(A,A)$-bimodule, the structure of $A$-module on the left being that which one has just defined, and the structure of $A$-module on the right coming from the structure of $A$-module of $A^{(p)}$ provenant of the ring structure, so that for $x\in E$, $\alpha$, $\beta$ in $A$, one has
\[
  \alpha(x\otimes\beta)=x\otimes(\alpha\beta) \quad\text{and}\quad (\alpha x)\otimes\beta=x\otimes\alpha^p\beta.
\]
\end{env}

\oldpage[0\textsubscript{IV-1}]{155}
Setting $x^{(p)}=x\otimes 1$, one therefore has $(\alpha x)^{(p)}=\alpha^p x^{(p)}$.
When $\mathrm{F}_A$ is \emph{injective}, one can identify $A^{(p)}$ as an algebra over its \emph{subring} $A^p$.

\begin{proposition}[21.1.5]
\label{0.21.1.5}
Let $A$, $B$ be two rings, $u:A\to B$ a homomorphism.
\begin{enumerate}
  \item[\rm(i)] Every $A$-derivation of $B$ in a $B$-module $L$ is zero in $A[B^p]$ \textup{(}hence is an $A[B^p]$-derivation\textup{)}.
  \item[\rm(ii)] For every sub-$A$-algebra $A'$ of $A[B^p]$, if $j:A'\to B$ is the canonical injection, the canonical homomorphism
\[
  j_{B/A'/A}:\Omega^1_{B/A}\to\Omega^1_{B/A'}
\]
is bijective.
  \item[\rm(iii)] Suppose that there exists an integer $s\geq 0$ such that $u(A)\supset B^{p^s}$.
  Then, in the ring $B\otimes_A B$, $\mathfrak{J}_{B/A}$ is a nilideal.
\end{enumerate}
\end{proposition}

\begin{proof}
(i) By induction from \sref{0.20.1.1}, one deduces, for every $A$-derivation $D:B\to L$, that $D(x^k)=kx^{k-1}D(x)$, whence in particular $D(x^p)=0$.

(ii) In view of \sref{0.20.4.8}, the assertion (ii) is only a translation of (i), this last being written $\Der_{A[B^p]}(B,L)=\Der_A(B,L)$ for every $B$-module $L$.

(iii) For every $x\in B$, one has $(x\otimes 1-1\otimes x)^{p^s}=x^{p^s}\otimes 1-1\otimes x^{p^s}=0$, since $x^{p^s}\in u(A)$.
The conclusion follows from the fact that the elements $1\otimes x-x\otimes 1$ generate $\mathfrak{J}_{B/A}$ \sref{0.20.4.4}.
\end{proof}

It follows immediately from \sref{0.21.1.5} that one has, for every pair of ring homomorphisms $A\to B\to C$,
\begin{equation}
\tag{21.1.5.1}
  \label{0.21.1.5.1}
  \Upsilon_{C/B/A}=\Upsilon_{C/B/A'}
\end{equation}
for every sub-$A$-algebra $A'$ of $A[B^p]$.

On the other hand, \sref{0.21.1.5} also shows in particular that for the modules of ``absolute'' differentials
\begin{equation}
\tag{21.1.5.2}
  \label{0.21.1.5.2}
  \Omega^1_A=\Omega^1_{A/A^p}
\end{equation}

\begin{corollary}[21.1.6]
\label{0.21.1.6}
Suppose that $B$ is an $A$-algebra of finite type and that there exists an integer $s\geq 0$ such that $u(A)\supset B^{p^s}$.
If $\Omega^1_{B/A}=0$, $u$ is surjective.
\end{corollary}

\begin{proof}
\oldpage[0\textsubscript{IV-1}]{156}
By \sref{0.21.1.5}, (iii), $\mathfrak{J}_{B/A}$ is a nilideal; moreover, by virtue of \sref{0.20.4.4} and the hypothesis $\Omega^1_{B/A}=0$, the relation $\Omega^1_{B/A}=0$ signifies that $\mathfrak{J}_{B/A}=\mathfrak{J}^2_{B/A}$, whence one concludes that $\mathfrak{J}_{B/A}=0$, or equivalently that $\pi:B\otimes_A B\to B$ is bijective.
But every element of $B$ has its $p^s$-th power in $u(A)$, $B$ is \emph{integral} over $A$, and being of finite type over $A$, $B$ is thus reduced to proving the following lemma (where one \emph{does not assume} that the rings considered are of characteristic $p$):
\end{proof}

\begin{lemma}[21.1.6.1]
\label{0.21.1.6.1}
Let $R$ be a ring, $S$ a finite $R$-algebra; if the canonical homomorphism $\pi:S\otimes_R S\to S$ is bijective, then the structural homomorphism $u:R\to S$ is surjective.
\end{lemma}

\begin{proof}
It suffices to show that for every maximal ideal $\mathfrak{m}$ of $R$, if one sets $T=R-\mathfrak{m}$, the homomorphism $u_\mathfrak{m}:R_\mathfrak{m}\to T^{-1}S$ is surjective (Bourbaki, \emph{Alg.\ comm.}, chap.~II, \S\,3, n$^\circ$\,3, th.~1); or the hypothesis implies that the homomorphism $(T^{-1}S)\otimes_{R_\mathfrak{m}}(S/\mathfrak{m}S)\to T^{-1}S$ is bijective, and as $T^{-1}S$ is a finite $R_\mathfrak{m}$-algebra, one sees that one can reduce to the case where $R$ is a local ring.
Denoting again by $\mathfrak{m}$ its maximal ideal, it suffices to prove that $u\otimes 1:R/\mathfrak{m}\to S/\mathfrak{m}S$ is surjective, by virtue of the lemma of Nakayama; as the canonical homomorphism $(S/\mathfrak{m}S)\otimes_{R/\mathfrak{m}}(S/\mathfrak{m}S)\to S/\mathfrak{m}S$ is bijective and $S/\mathfrak{m}S$ is a finite $(R/\mathfrak{m})$-algebra, one is finally reduced to the case where $R$ is a \emph{field}; but then the ranks of $S\otimes_R S$ and of $S$ over $R$ must be equal, which is possible only if $S$ is of rank $0$ or $1$, hence $u:R\to S$ is surjective.
\end{proof}

\begin{proposition}[21.1.7]
\label{0.21.1.7}
Let $A$ be a ring, $B$ an $A$-algebra, $(x_\alpha)_{\alpha\in I}$ a family of elements of $B$.
Consider the following properties:
\begin{enumerate}
  \item[\rm a)] $B=A[B^p,(x_\alpha)]$, or equivalently the $A[B^p]$-algebra $B$ is generated by the family $(x_\alpha)_{\alpha\in I}$.
  \item[\rm b)] The $A[B^p]$-module $B$ is generated by the monomials $\prod_\alpha x_\alpha^{n(\alpha)}$, where $(n(\alpha))_{\alpha\in I}$ is a family of integers with finite support such that $0\leq n(\alpha)<p$ for all $\alpha\in I$.
  \item[\rm c)] The $B$-module $\Omega^1_{B/A}$ is generated by the $d_{B/A}(x_\alpha)$ $(\alpha\in I)$.
\end{enumerate}
Then the properties \textup{a)} and \textup{b)} are equivalent and imply \textup{c)}; if in addition $B$ is an $A[B^p]$-algebra of finite type, \textup{c)} is equivalent to \textup{a)} and \textup{b)}.
\end{proposition}

\begin{proof}
It is clear that b) implies a), and conversely a) implies b), since every monomial $\prod_\alpha x_\alpha^{m(\alpha)}$, where the $m(\alpha)$ are integers $\geq 0$ (forming a family with finite support), can be written $(\prod_\alpha x_\alpha^{q(\alpha)})^p\prod_\alpha x_\alpha^{r(\alpha)}$, by dividing each $m(\alpha)$ by $p$, which gives $m(\alpha)=p.q(\alpha)+r(\alpha)$ with $0\leq r(\alpha)<p$.
The fact that a) implies c) follows from \sref{0.21.1.5}, (ii) and \sref{0.20.4.7}.

Suppose conversely c) satisfied and that $B$ is an $A[B^p]$-algebra of finite type; let $B'$ be the sub-$A[B^p]$-algebra of $B$ generated by the $x_\alpha$; in the exact sequence \sref{0.20.5.7.1}
\[
  \Omega^1_{B'/A[B^p]}\otimes_{B'}B\to\Omega^1_{B/A[B^p]}\to\Omega^1_{B/B'}\to 0
\]
the hypothesis implies that the left arrow is surjective (taking into account \sref{0.21.1.5}, (ii)); one therefore has $\Omega^1_{B/B'}=0$, and as $B^p\subset B'\subset B$ and $B$ is a $B'$-algebra of finite type, one necessarily has $B'=B$ by \sref{0.21.1.6}.
\end{proof}

\begin{remark}[21.1.8]
\label{0.21.1.8}
When $B$ is a \emph{field}, we will demonstrate the equivalence of properties a), b) and c) without hypothesis of finiteness \sref{0.21.4.5}.
\end{remark}

\begin{definition}[21.1.9]
\label{0.21.1.9}
Let $A$ be a ring (of characteristic $p$), $B$ an $A$-algebra.
One says that a family $(x_\alpha)_{\alpha\in I}$ of elements of $B$ is \emph{$p$-free over $A$} (resp.\ a \emph{system of $p$-generators} of $B$ over $A$, resp.\ a \emph{$p$-base} of $B$ over $A$) if the family of monomials $\prod_\alpha x_\alpha^{n(\alpha)}$ $(0\leq n(\alpha)<p$, $(n(\alpha))_{\alpha\in I}$ with finite support) is a free family (resp.\ a system of generators, resp.\ a base) in the $A[B^p]$-module $B$.
\end{definition}

One has corresponding definitions for a \emph{subset} $M$ of elements of $B$, by considering the family defined by the canonical injection $M\to B$.
When one takes for $A$ the \emph{prime field} $\bb{F}_p$ (in which case $A[B^p]=B^p$), one omits the mention of $A$ in the preceding definition (or also says that a family is ``absolutely'' $p$-free, resp.\ an absolute system of $p$-generators, resp.\ an absolute $p$-base).

\oldpage[0\textsubscript{IV-1}]{157}
\begin{lemma}[21.1.10]
\label{0.21.1.10}
Let $C$ be a sub-$A$-algebra of $B$ such that $B^p\subset C$, $M$ a subset of $C$, $N$ a subset of $B$.
\begin{enumerate}
  \item[\rm(i)] If $M$ is a system of $p$-generators of $C$ over $A$ and $N$ a system of $p$-generators of $B$ over $C$, then $M\cup N$ is a system of $p$-generators of $B$ over $A$.
  \item[\rm(ii)] Suppose that $M$ is a $p$-base of $C$ over $A$; then, for $N$ to be $p$-free over $C$, it is necessary and sufficient that $M\cup N$ is $p$-free over $A$.
\end{enumerate}
\end{lemma}

\begin{proof}
One can reduce to the case where $B^p\subset A\subset C\subset B$.
Then (i) is a particular case from the fact that if $P$ (resp.\ $Q$) is a system of generators of the $A$-module $C$ (resp.\ the $C$-module $B$), the set of products $xy$, $x\in P$ and $y\in Q$, is a system of generators of the $A$-module $B$.
Preserving the same notations, if $P$ is a base of the $A$-module $C$, to say that $Q$ is a free family on $C$ signifies that the relation $\sum_{\lambda,\mu}a_{\lambda\mu}x_\lambda y_\mu=0$, where $a_{\lambda\mu}\in A$, $x_\lambda\in P$, $y_\mu\in Q$ (with $x_\lambda$ (resp.\ $y_\mu$) pairwise distinct) is equivalent to $a_{\lambda\mu}=0$ for every pair $(\lambda,\mu)$, or equivalently $\sum_\lambda a_{\lambda\mu}x_\lambda=0$ for every $\mu$, whence the assertion (ii).
\end{proof}

\subsection{\texorpdfstring{$p$}{p}-bases and formal smoothness}
\label{subsection:0.21.2}

\begin{theorem}[21.2.1]
\label{0.21.2.1}
Let $B$ be a ring, $A$ a subring of $B$ such that $B^p\subset A\subset B$, $(x_\alpha)_{\alpha\in I}$ a $p$-base of $B$ over $A$.
Let $E$ be an $A$-algebra, $u:A\to E$ the structural homomorphism, $q:E\to B$ a surjective $A$-homomorphism, $\mathfrak{R}$ its kernel, and suppose that $\mathfrak{R}^p=0$.
Then:
\begin{enumerate}
  \item[\rm(i)] For there to exist an $A$-homomorphism right inverse $v:B\to E$ of the homomorphism $q:E\to B$, it is necessary and sufficient that $u(q(z)^p)=z^p$ for every $z\in E$.
  \item[\rm(ii)] When the condition of \textup{(i)} is satisfied, for every family $(z_\alpha)_{\alpha\in I}$ of elements of $E$ such that $q(z_\alpha)=x_\alpha$ for all $\alpha\in I$, there exists an $A$-homomorphism $v:B\to E$ and \emph{a unique one} such that $v(x_\alpha)=z_\alpha$ for all $\alpha\in I$, and $v$ is right inverse to $q$.
\end{enumerate}
\end{theorem}

\begin{proof}
If there exists an $A$-homomorphism $v:B\to E$ right inverse to $q$, then $v(a)=u(a)$ for every $a\in A\subset B$, and hence $v(q(z)^p)=u(q(z)^p)$ for every $z\in E$, since $B^p\subset A$.
Moreover $v(q(z))\equiv z\pmod{\mathfrak{R}}$, so $(v(q(z)))^p=z^p$ because $\mathfrak{R}^p=0$; as $v(q(z)^p)=(v(q(z)))^p$, this proves the necessity of (i).
The sufficiency of condition (i) follows from (ii).

Under the hypotheses of (ii), uniqueness is clear because the monomials $\prod_\alpha x_\alpha^{n(\alpha)}$ generate the $A$-module $B$.
As these monomials, for finitely supported families with $0\leq n(\alpha)<p$, form an $A$-basis of $B$, there is an $A$-linear map $v:B\to E$ such that
\[
  v\left(\prod_\alpha x_\alpha^{n(\alpha)}\right)=\prod_\alpha z_\alpha^{n(\alpha)}.
\]
It remains to prove that $v$ is multiplicative. One can write
\[
  \left(\prod_\alpha x_\alpha^{m(\alpha)}\right)\left(\prod_\alpha x_\alpha^{n(\alpha)}\right)=a\prod_\alpha x_\alpha^{r(\alpha)},
\]
where $r(\alpha)=m(\alpha)+n(\alpha)$ if this sum is less than $p$, and $r(\alpha)=m(\alpha)+n(\alpha)-p$ otherwise, while $a\in A$
\oldpage[0\textsubscript{IV-1}]{158}
is the product of the $x_\alpha^p$ over the indices for which $m(\alpha)+n(\alpha)\geq p$.
Thus it suffices to prove $u(a)=\prod_\alpha z_\alpha^p$, and this follows from $q(z_\alpha)=x_\alpha$ and condition (i).
\end{proof}

\begin{corollary}[21.2.2]
\label{0.21.2.2}
If there exists a $p$-base of $B$ over $A$, $B$ is an $A$-algebra formally smooth relative to $B^p$ \textup{(}for the discrete topologies\textup{)}.
\end{corollary}

\begin{proof}
Indeed, let $E$ be an $A$-extension of $B$ by a $B$-module $L$, $q:E\to B$ the augmentation, $u:A\to E$ the structural homomorphism.
To say that $E$ is $B^p$-trivial means that there exists an $A$-homomorphism $v:B\to E$ such that $q(v(b))=b$ and $v(b^p)=u(b^p)$ for all $b\in B$.
But one has $u(q(z)^p)=v(q(z)^p)=(v(q(z)))^p$; but by the relation $L^2=0$, one also has $L^p=0$, and as $v(q(z))-z\in L$, one has $(v(q(z)))^p=z^p$; the condition of \sref{0.21.2.1}, (i) is therefore satisfied, and $E$ is $B$-trivial.
\end{proof}

\begin{corollary}[21.2.3]
\label{0.21.2.3}
Let $B$ be a ring, $A$ a subring of $B$ such that $B^p\subset A\subset B$, $L$ a $B$-module.
Then:
\begin{enumerate}
  \item[\rm(i)] For a derivation $D$ of $A$ in $L$ to extend to a derivation of $B$ in $L$, it is necessary and sufficient that $D$ annihilate $B^p$.
  \item[\rm(ii)] If $D$ annihilates $B^p$, then, for every family $(y_\alpha)_{\alpha\in I}$ of elements of $L$, there exists a unique derivation $D'$ of $B$ in $L$ extending $D$ and such that $D'(x_\alpha)=y_\alpha$ for all $\alpha\in I$.
\end{enumerate}
\end{corollary}

\begin{proof}
Given a derivation $D$ of $A$ in $L$, consider the ring $E=\mathrm{D}_B(L)$ and the homomorphism $u:A\to E$ defined by $u(a)=(a,D(a))$; an $A$-homomorphism right inverse $v:B\to E$ is then of the form $x\mapsto(x,D'(x))$, where $D'$ is a derivation of $B$ in $L$ extending $D$ \sref{0.20.1.5}; as $u(q(z)^p)=(q(z)^p,D(q(z)^p))$ for $z\in E$, the corollary follows immediately from \sref{0.21.2.1} applied to $E$.
\end{proof}

It follows from \sref{0.21.2.3} that the sequence
\begin{equation}
\tag{21.2.3.1}
  \label{0.21.2.3.1}
  0\to\Der_A(B,L)\to\Der_{B^p}(B,L)\to\Der_{B^p}(A,L)\to 0
\end{equation}
(cf.\ \sref{0.20.2.2.1}) is \emph{exact}, and that the map
\begin{equation}
\tag{21.2.3.2}
  \label{0.21.2.3.2}
  D'\mapsto(D'(x_\alpha))_{\alpha\in I}
\end{equation}
is an isomorphism of the $B$-module $\Der_A(B,L)$ onto the $B$-module product $L^I$ (by taking $D=0$ in \sref{0.21.2.3}).

\begin{corollary}[21.2.4]
\label{0.21.2.4}
Under the hypotheses of \sref{0.21.2.3}, the sequence
\begin{equation}
\tag{21.2.4.1}
  \label{0.21.2.4.1}
  0\to\Omega^1_{A/B^p}\otimes_A B\to\Omega^1_{B/B^p}\to\Omega^1_{B/A}\to 0
\end{equation}
is exact and split, and the family $(d_{B/A}(x_\alpha))_{\alpha\in I}$ is a base of the $B$-module $\Omega^1_{B/A}$.
\end{corollary}

\begin{proof}
This follows immediately from \sref{0.21.2.3} and from the formula $\Der_A(B,L)=\Hom_B(\Omega^1_{B/A},L)$.
\end{proof}

\begin{corollary}[21.2.5]
\label{0.21.2.5}
Let $A$ be a ring, $B$ an $A$-algebra, $(x_\alpha)_{\alpha\in I}$ a $p$-base of $B$ over $A$.
Then $(d_{B/A}(x_\alpha))_{\alpha\in I}$ is a base of the $B$-module $\Omega^1_{B/A}$.
\end{corollary}

\begin{proof}
By using \sref{0.21.1.5}, (ii), one reduces to the case where $A=A[B^p]$ and it suffices then to apply \sref{0.21.2.4}.
\end{proof}

\oldpage[0\textsubscript{IV-1}]{159}
\begin{env}[21.2.6]
\label{0.21.2.6}
Let $A$, $B$ be two rings, $u:A\to B$ a ring homomorphism; one has (with the notations of \sref{0.21.1.4}) a commutative diagram
\begin{equation}
\tag{21.2.6.1}
  \label{0.21.2.6.1}
\xymatrix{
  A^{(p)} \ar[r]^{u} & B^{(p)} \\
  A \ar[u]^{\mathrm{F}_A} \ar[r]^{u} & B \ar[u]_{\mathrm{F}_B}
}
\end{equation}
One deduces therefore canonically a homomorphism of $B$-algebras
\begin{equation}
\tag{21.2.6.2}
  \label{0.21.2.6.2}
  u\otimes 1:A^{(p)}\otimes_A B\to B^{(p)}
\end{equation}
whose image is the ring $A[B^p]$ (which is a $B$-algebra for the homomorphism $\mathrm{F}_B:B\to B^p$).
\end{env}

\begin{theorem}[21.2.7]
\label{0.21.2.7}
Let $A$ be a ring, $B$ an $A$-algebra, $u:A\to B$ the structural homomorphism.
One assumes the following conditions satisfied:
\begin{enumerate}
  \item[\rm(i)] The homomorphism \sref{0.21.2.6.2} deduced from $u$ is injective.
  \item[\rm(ii)] $B$ admits a $p$-base $(x_\alpha)_{\alpha\in I}$ relative to $A$.
\end{enumerate}
Then $B$ is an $A$-algebra formally smooth \textup{(}for the discrete topologies\textup{)}.
More precisely, let $E$ be a topological $A$-algebra admissible ($\mathbf{0_I}$, 7.1.2), $\mathfrak{R}$ an ideal of definition of $E$, $C=E/\mathfrak{R}$, $v:B\to C$ an $A$-homomorphism, $q:E\to C$ the augmentation.
Then, for every family $(z_\alpha)_{\alpha\in I}$ of elements of $E$ such that $q(z_\alpha)=v(x_\alpha)$ for all $\alpha\in I$, there exists an $A$-homomorphism $w:B\to E$ and \emph{a unique one} such that $v=q\circ w$ and $w(x_\alpha)=z_\alpha$ for all $\alpha\in I$.
\end{theorem}

\begin{proof}
Consider first the case where $E$ is discrete, hence $\mathfrak{R}$ is nilpotent.
One can restrict to the case where $v$ is surjective; moreover, the reasoning of
\sref{0.19.4.3} allows one to suppose $\mathfrak{R}^2=0$ and consequently
$\mathfrak{R}^p=0$. Finally, by considering the inverse image under $v$ of the
extension $E$ of $C$ by $\mathfrak{R}$, one can restrict to the case where $C=B$
and $v$ is the identity \sref{0.19.4.4}. As the ring homomorphism
$\mathrm{F}_E:z\rightsquigarrow z^p$ vanishes on $\mathfrak{R}$, it factors as
$E\xrightarrow q B\xrightarrow{\mathrm{F}'}E$, and $\mathrm{F}'$, considered as a
homomorphism from $B$ to $E^{(p)}$, is an $A$-\emph{homomorphism} by definition
of the $A$-algebra structure on $E^{(p)}$ using the composite homomorphism
$A\xrightarrow r E\xrightarrow{\mathrm{F}_E}E^{(p)}$, where $r:A\to E$ is the
structural homomorphism of the $A$-algebra $E$; moreover,
$r:A^{(p)}\to E^{(p)}$ is also an $A$-homomorphism. There is therefore a unique
$A$-homomorphism $f:A^{(p)}\otimes_A B\to E^{(p)}$ such that the composites of $f$
with the canonical homomorphisms $A^{(p)}\to A^{(p)}\otimes_A B$ and
$B\to A^{(p)}\otimes_A B$ are respectively $r$ and $\mathrm{F}'$. Now, by hypothesis,
one can identify $A^{(p)}\otimes_A B$ with $A[B^p]$, the canonical homomorphisms
$A^{(p)}\to A^{(p)}\otimes_A B$ and $B\to A^{(p)}\otimes_A B$ being identified
respectively with $u$ and $\mathrm{F}_B$. One can therefore now consider $E$
as an $A[B^p]$-algebra by means of the homomorphism $f$, and, by construction,
one has $f(q(z)^p)=z^p$ for every $z\in E$; the conditions for applying
\sref{0.21.2.1} are therefore satisfied, whence the theorem in this case.

To pass to the general case, consider a fundamental system
$(\mathfrak{J}_\lambda)$ of open ideals of $E$, and put
$E_\lambda=E/\mathfrak{J}_\lambda$,
$\mathfrak{R}_\lambda=(\mathfrak{R}+\mathfrak{J}_\lambda)/\mathfrak{J}_\lambda$,
$C_\lambda=E_\lambda/\mathfrak{R}_\lambda=E/(\mathfrak{R}+\mathfrak{J}_\lambda)$;
as $E$ is admissible, one has $E=\varprojlim E_\lambda$ ($\mathbf{0_I}$, 7.2.2).
For every pair $(\alpha,\lambda)$, denote by $z_{\alpha\lambda}$ the canonical
image of $z_\alpha$ in $E_\lambda$, by $p_\lambda:C\to C_\lambda$ the canonical
homomorphism,
\oldpage[0\textsubscript{IV-1}]{160}
and by $q_\lambda$ the canonical homomorphism $E_\lambda\to C_\lambda$.
Since by hypothesis $\mathfrak{R}_\lambda$ is nilpotent in $E_\lambda$, the first
part of the proof establishes the existence and uniqueness of an
$A$-homomorphism $w_\lambda:B\to E_\lambda$ such that
$p_\lambda\circ v=q_\lambda\circ w_\lambda$ and
$w_\lambda(x_\alpha)=z_{\alpha\lambda}$ for every $\alpha$.
The uniqueness of the $w_\lambda$ then immediately shows that $(w_\lambda)$ is a
projective system of homomorphisms, and $w=\varprojlim w_\lambda$ answers the
question.
\end{proof}

\begin{remark}[21.2.8]
\label{0.21.2.8}
The verification of the existence and uniqueness of the homomorphism $w:B\to E$ such that $w(x_\alpha)=z_\alpha$ for all $\alpha$ can also be deduced directly from the fact that $B$ is an $A$-algebra formally smooth and from the fact that the $d_{B/A}(x_\alpha)$ form a base of the $B$-module $\Omega^1_{B/A}$ \sref{0.21.2.4}, without invoking the existence of a $p$-base (so that the result is valid without supposing the rings of characteristic $p$).
Indeed, one can reduce to the case where $\mathfrak{R}^2=0$; as $\Der_A(B,\mathfrak{R})=\Hom_B(\Omega^1_{B/A},\mathfrak{R})$ by \sref{0.20.4.8}, there exists an $A$-derivation $D$ of $B$ in $\mathfrak{R}$ and a unique one such that $D(x_\alpha)=z'_\alpha$ for every family $(z'_\alpha)$ of elements of $\mathfrak{R}$; the conclusion follows from \sref{0.20.1.1}.
\end{remark}

\subsection{\texorpdfstring{$p$}{p}-bases and modules of imperfection}
\label{subsection:0.21.3}

\begin{env}[21.3.1]
\label{0.21.3.1}
Let $A$, $B$ be two rings (of characteristic $p$), $i:A\to B$, $j:B\to A$ two ring homomorphisms such that one has
\begin{align}
\tag{21.3.1.1}
  \label{0.21.3.1.1}
  j(i(a))&=a^p \quad\text{for all } a\in A, \\
\tag{21.3.1.2}
  \label{0.21.3.1.2}
  i(j(b))&=b^p \quad\text{for all } b\in B.
\end{align}
Most often, $i$ will be \emph{injective}, so that $A$ is identified with a subring of $B$; once this identification is made, the existence of $j$ implies that $B^p\subset A$, and $j$ is then identified with $\mathrm{F}_B$.
\end{env}

\begin{env}[21.3.2]
\label{0.21.3.2}
If $h:A^p\to A$ is the canonical injection, one has, by \sref{0.21.3.1.2} and \sref{0.21.3.1.1}, a commutative diagram
\begin{equation}
\tag{21.3.2.1}
  \label{0.21.3.2.1}
\xymatrix{
  B \ar[r]^{j} & A \\
  A \ar[u]^{i} \ar[r]^{\mathrm{F}_A} & A^p \ar[u]_{h}
}
\end{equation}
so that the pair $(j,\mathrm{F}_A)$ can be considered as a \emph{di-homomorphism} of the $A$-algebra $B$ (for $i$) in the $A^p$-algebra $A$ (for $h$).
One deduces a \emph{canonical homomorphism of $A$-modules}
\begin{equation}
\tag{21.3.2.2}
  \label{0.21.3.2.2}
  \pi_{B/A}:\Omega^1_{B/A}\otimes_B A_{[j]}\to\Omega^1_{A/A^p}=\Omega^1_A
\end{equation}
(cf.\ \sref{0.20.5.4}; the identification of $\Omega^1_{A/A^p}$ and of the module of absolute differentials $\Omega^1_A$ comes from \sref{0.21.1.5}, (ii)).

One will set
\begin{align}
\tag{21.3.2.3}
  \label{0.21.3.2.3}
  \Theta_{B/A}&=\Omega^1_{B/A}\otimes_B A_{[j]} \\
\tag{21.3.2.4}
  \label{0.21.3.2.4}
  \Xi_{B/A}&=\Ker(\pi_{B/A})
\end{align}
\oldpage[0\textsubscript{IV-1}]{161}
so that one has the exact sequence
\begin{equation}
\tag{21.3.2.5}
  \label{0.21.3.2.5}
  0\to\Xi_{B/A}\to\Theta_{B/A}\xrightarrow{\pi_{B/A}}\Omega^1_A
\end{equation}
(one notes that these notations can be confusing since $\Theta_{B/A}$ and $\pi_{B/A}$ depend not only on $A$ and $B$, but also on $i$ and $j$).
\end{env}

\begin{env}[21.3.3]
\label{0.21.3.3}
As by \sref{0.21.3.1.2}, one has $\mathrm{F}_B=i\circ j$, one can write for every $B$-module $M$ (cf.\ \sref{0.21.1.4})
\begin{equation}
\tag{21.3.3.1}
  \label{0.21.3.3.1}
  M^{(p)}=M\otimes_B B^{(p)}=(M\otimes_B A_{[j]})\otimes_A B_{[i]}
\end{equation}
whence in particular
\begin{equation}
\tag{21.3.3.2}
  \label{0.21.3.3.2}
  (\Omega^1_{B/A})^{(p)}=\Theta_{B/A}\otimes_A B_{[i]}
\end{equation}
and one deduces therefore from \sref{0.21.3.2.2} a \emph{canonical homomorphism of $B$-modules}
\begin{equation}
\tag{21.3.3.3}
  \label{0.21.3.3.3}
  \pi_{B/A}\otimes 1_B:(\Omega^1_{B/A})^{(p)}\to\Omega^1_A\otimes_A B_{[i]}.
\end{equation}

Taking into account \sref{0.20.5.4.2}, the image of this homomorphism is the $B$-module generated by the elements $d_A(j(b))\otimes 1$ for $b\in B$.
Their image by the canonical homomorphism $i_{B/A}:\Omega^1_A\otimes_A B\to\Omega^1_B$ deduced from $i$ \sref{0.20.5.2} is therefore contained in the sub-$B$-module of $\Omega^1_B$ generated by the $d_B(i(j(b)))$ for $b\in B$; by virtue of \sref{0.21.3.1.2}, this image is \emph{zero}.
In other words, one has a sequence of homomorphisms
\begin{equation}
\tag{21.3.3.4}
  \label{0.21.3.3.4}
  0\to\Xi_{B/A}\otimes_A B_{[i]}\to(\Omega^1_{B/A})^{(p)}\to\Upsilon_{B/A}
\end{equation}
which is not necessarily exact, but where the composite of two consecutive homomorphisms is zero.
\end{env}

\begin{proposition}[21.3.4]
\label{0.21.3.4}
If $B$ is a flat $A$-module \textup{(}for $i$\textup{)}, the sequence \sref{0.21.3.3.4} is exact.
\end{proposition}

\begin{proof}
This follows from the definition of flatness.
The proposition \sref{0.21.3.4} applies in particular when $B^p\subset A\subset B$, $i$ and $j$ being respectively the canonical injection and $\mathrm{F}_B$ (so that $B$ is then a free $A$-module).
But even in this case the kernel $\Xi_{B/A}$ is not necessarily zero.
However:
\end{proof}

\begin{proposition}[21.3.5]
\label{0.21.3.5}
Let $B$ be a \emph{reduced} ring, $A$ a subring of $B$ such that $B^p\subset A\subset B$.
Suppose that there exists a $p$-base $(x_\lambda)_{\lambda\in L}$ of $B$ over $A$ and a $p$-base $(y_\mu)_{\mu\in M}$ of $A$ over $B^p$; then the canonical homomorphism \sref{0.21.3.3.4}
\begin{equation}
\tag{21.3.5.1}
  \label{0.21.3.5.1}
  (\Omega^1_{B/A})^{(p)}\to\Upsilon_{B/A}
\end{equation}
is bijective, and the elements $d_A(x^p_\lambda)\otimes 1$ form a base of the $B$-module $\Upsilon_{B/A}$.
\end{proposition}

\begin{proof}
\oldpage[0\textsubscript{IV-1}]{162}
As $B$ is reduced, $x\mapsto x^p$ is an isomorphism of structure of $B$ onto $B^p$, and by transport of structure via this isomorphism, one sees that $(x^p_\lambda)$ and the $y_\mu$ form a $p$-base of $B^p$ over $A^p$ \sref{0.21.1.10}; by consequence \sref{0.21.1.5}, (ii) and \sref{0.21.2.5} the $d_A(x^p_\lambda)$ and the $d_A(y_\mu)$ form a \emph{base} of the $A$-module $\Omega^1_A$.
But, the image of $d_A(x_\lambda)\otimes 1$ by $i_{B/A}$ is $d_B(x^p_\lambda)=0$; on the other hand, the image of $d_A(y_\mu)$ by $i_{B/A}$ is $d_B(i(y_\mu))$, and as the $y_\mu$ and the $x_\lambda$ form a $p$-base of $B$ over $B^p$ \sref{0.21.1.10}, the $d_B(y_\mu)$ are part of a base of $\Omega^1_B$ \sref{0.21.2.5}, and are thus linearly independent over $B$.
Hence the kernel of $i_{B/A}$ has a base formed by the $d_A(x^p_\lambda)\otimes 1$, and as these are the images by \sref{0.21.3.5.1}, these last form a base of the $B$-module $\Omega^1_{B/A}\otimes_B B^{(p)}$ \sref{0.21.2.5}, the homomorphism \sref{0.21.3.5.1} of $B$-modules is bijective.
\end{proof}

\begin{remarks}[21.3.6]
\label{0.21.3.6}
\begin{enumerate}
  \item[\rm(i)] Let $A'$, $B'$ be two rings of characteristic $p$, and suppose one has two homomorphisms $i':A'\to B'$, $j':B'\to A'$ verifying \sref{0.21.3.1.1} and \sref{0.21.3.1.2}, and ring homomorphisms $f:A\to A'$, $g:B\to B'$ rendering commutative the diagram
\[
\xymatrix{
  A' \ar[r]^{i'} & B' \ar[r]^{j'} & A' \\
  A \ar[u]^f \ar[r]^{i} & B \ar[u]^g \ar[r]^{j} & A \ar[u]^f
}
\]
Then the canonical homomorphism $\Omega^1_{B/A}\to\Omega^1_{B'/A'}$ \sref{0.20.5.4.3} gives a di-homomorphism $\Theta_{B/A}\to\Theta_{B'/A'}$ rendering commutative the diagram
\[
\xymatrix{
  0 \ar[r] & \Xi_{B/A} \ar[r] \ar[d] & \Theta_{B/A} \ar[r]^{\pi_{B/A}} \ar[d] & \Omega^1_A \ar[d] \\
  0 \ar[r] & \Xi_{B'/A'} \ar[r] & \Theta_{B'/A'} \ar[r]^{\pi_{B'/A'}} & \Omega^1_{A'}
}
\]

  \item[\rm(ii)] Let $A'$ be any $A$-algebra, and set $B'=B\otimes_A A'$; then $i'=i\otimes 1:A'\to B'$ and $j'=j\otimes 1:B'\to A'$ verify \sref{0.21.3.1.1} and \sref{0.21.3.1.2}, and one has
\[
  \Omega^1_{B'/A'}=\Omega^1_{B/A}\otimes_A A'=\Omega^1_{B/A}\otimes_B B'
\]
\sref{0.20.5.5}; one deduces a canonical $A'$-isomorphism
\begin{equation}
\tag{21.3.6.1}
  \label{0.21.3.6.1}
  \Theta_{B/A}\otimes_A A'\xrightarrow{\sim}\Theta_{B'/A'}
\end{equation}
and also a canonical $B'$-homomorphism
\begin{equation}
\tag{21.3.6.2}
  \label{0.21.3.6.2}
  (\Omega^1_{B/A})^{(p)}\otimes_B B'\xrightarrow{\sim}(\Omega^1_{B'/A'})^{(p)}.
\end{equation}
\end{enumerate}
\end{remarks}

\subsection{The case of field extensions}
\label{subsection:0.21.4}

\begin{env}[21.4.0]
\label{0.21.4.0}
Let $K$ be a field of characteristic $p>0$, $k$ a subfield of $K$; then the ring $k[K^p]$ is equal to the field $k(K^p)$ since $k$ is algebraic over $K^p$.
One can therefore apply the results of the preceding sections everywhere replacing $A$, $B$ and $A[B^p]$ by $k$, $K$ and $k(K^p)$.
\end{env}

\begin{lemma}[21.4.1]
\label{0.21.4.1}
Let $k$ be a field of characteristic $p>0$, $K$ an extension of $k$.
For an element $x\in K$ to be $p$-free over $k$, it is necessary and sufficient that $x\notin k(K^p)$.
\end{lemma}

\begin{proof}
Indeed, $x$ is a root of the polynomial $X^p-x^p$ of $k(K^p)[X]$, and one knows (Bourbaki, \emph{Alg.}, chap.~V, \S\,8, n$^\circ$\,1, prop.~1) that if $x\notin k(K^p)$ this polynomial is irreducible, so that the elements $1,x,\ldots,x^{p-1}$ form a base of the $k(K^p)$-module $k(K^p)(x)$.
\end{proof}

\begin{theorem}[21.4.2]
\label{0.21.4.2}
Let $k$ be a field of characteristic $p>0$, $K$ an extension of $k$, $S$ a system of $p$-generators of $K$ over $k$, $L\subset S$ a $p$-free subset over $k$.
There exists then a $p$-base $B$ of $K$ over $k$ such that $L\subset B\subset S$.
In particular, every extension of $k$ admits a $p$-base over $k$.
\end{theorem}

\begin{proof}
One can reduce to the case where $K^p\subset k$.
By virtue of Zorn's theorem, there exists in $K$ a subset $B$ such that $L\subset B\subset S$, $p$-free over $k$ and \emph{maximal} amongst the subsets of $S$ having these properties.
It suffices to see that the subfield $K'$ of $K$ generated by $k$ and $B$ is equal to $K$.
In the opposite case, there would exist $x\in S$ not in $K'$; as $K^p\subset k$, one has $K'=K'(K^p)$; hence $x$ would be $p$-free over $K'$ \sref{0.21.4.1}, and hence $B\cup\{x\}$ would be $p$-free over $k$ \sref{0.21.1.10}, contradicting the maximality of $B$.
The last assertion of \sref{0.21.4.2} is obtained by taking $L=\emptyset$, $S=K$.
\qed
\end{proof}

\begin{corollary}[21.4.3]
\label{0.21.4.3}
Let $k$ be a field of characteristic $p>0$, $K$ an extension of $k$.
For a family $(x_\alpha)$ of elements of $K$ to be $p$-free over $k$, it is necessary and sufficient that, for all $\alpha$, $x_\alpha$ does not belong to the field $K_\alpha$ generated by $k(K^p)$ and the $x_\beta$ of index $\beta\neq\alpha$.
\end{corollary}

\begin{proof}
The condition is necessary by virtue of \sref{0.21.1.10}.
Conversely, suppose it is satisfied; one can reduce to the case where $(x_\alpha)$ is a system of $p$-generators of $K$.
There exists then a sub-family of $(x_\alpha)$ which is a $p$-base of $K$ \sref{0.21.4.2}, but this sub-family cannot be distinct from $(x_\alpha)$, since by hypothesis it would not be a family of $p$-generators of $K$.
\end{proof}

\oldpage[0\textsubscript{IV-1}]{163}
\begin{corollary}[21.4.4]
\label{0.21.4.4}
Let $k$ be a field of characteristic $p>0$, $K$ an extension of $k$.
For $\Omega^1_{K/k}=0$, it is necessary and sufficient that $K=k(K^p)$.
In particular, for $\Omega^1_K=0$, it is necessary and sufficient that $K$ be a \emph{perfect field}.
\end{corollary}

\begin{proof}
Indeed, if $(x_\alpha)$ is a $p$-base of $K$ over $k$, the $d_{K/k}(x_\alpha)$ form a base of the $K$-vector space $\Omega^1_{K/k}$ \sref{0.21.2.5}.
\end{proof}

\begin{theorem}[21.4.5]
\label{0.21.4.5}
Let $k$ be a field of characteristic $p>0$, $K$ an extension of $k$, $(x_\alpha)$ a family of elements of $K$.
For $(x_\alpha)$ to be $p$-free over $k$ \textup{(}resp.\ a family of $p$-generators of $K$ over $k$, resp.\ a $p$-base of $K$ over $k$\textup{)}, it is necessary and sufficient that the family $(d_{K/k}(x_\alpha))$ be a free family \textup{(}resp.\ a system of generators, resp.\ a base\textup{)} in the $K$-vector space $\Omega^1_{K/k}$.
\end{theorem}

\begin{proof}
\oldpage[0\textsubscript{IV-1}]{164}
Let $K'$ be the subfield (equal to the subring) of $K$ generated by $k(K^p)$ and the $x_\alpha$; taking into account \sref{0.21.4.3}, it suffices to see that, for each $\alpha$, $x_\alpha$ does not belong to $K_\alpha$.
But, in the exact sequence
\[
  \Omega^1_{K_\alpha/k(K^p)}\otimes_{K_\alpha} K\to\Omega^1_{K/k(K^p)}\to\Omega^1_{K/K_\alpha}\to 0
\]
the images by the left arrow of the $d_{K_\alpha/k}(x_\beta)$ for $\beta\neq\alpha$ are the $d_{K/k}(x_\beta)$, hence generate a sub-vector space of $\Omega^1_{K/k(K^p)}$ not containing $d_{K/k}(x_\alpha)$, and as this sub-vector space is the kernel of the right arrow, we see that $d_{K/K_\alpha}(x_\alpha)\neq 0$, hence $x_\alpha\notin K_\alpha$.

If now $(x_\alpha)$ is a $p$-free family, it is a sub-family of a $p$-base of $K$ over $k$ \sref{0.21.4.2}, hence the $d_{K/k}(x_\alpha)$ are part of a base of $\Omega^1_{K/k}$ \sref{0.21.2.5}, and are thus linearly independent over $K$.
Let us prove conversely that if the $d_{K/k}(x_\alpha)$ are linearly independent over $K$, the family $(x_\alpha)$ is $p$-free over $k$.
\end{proof}

\begin{corollary}[21.4.6]
\label{0.21.4.6}
Let $k$ be a field of characteristic $p>0$, $K$ an extension of $k$, $x$ an element of $K$.
The three following conditions are equivalent:
\begin{enumerate}
  \item[\rm a)] $x\notin k(K^p)$.
  \item[\rm b)] $d_{K/k}(x)\neq 0$.
  \item[\rm c)] The element $x$ is $p$-free over $k$.
\end{enumerate}
\end{corollary}

\begin{proposition}[21.4.7]
\label{0.21.4.7}
Let $L$ be a field of characteristic $p>0$, $K$ a subfield of $L$ such that $L^p\subset K\subset L$.
Then the sequence of $L$-vector spaces
\begin{equation}
\tag{21.4.7.1}
  \label{0.21.4.7.1}
  0\to\Omega^1_{K/L^p}\otimes_K L\to\Omega^1_{L/L^p}\to\Omega^1_{L/K}\to 0
\end{equation}
is exact; in other words, one has $\Upsilon_{L/K/L^p}=0$.
\end{proposition}

\begin{proof}
This is a particular case of \sref{0.21.2.4}, taking into account \sref{0.21.4.2}.
\end{proof}

\begin{proposition}[21.4.8]
\label{0.21.4.8}
Under the hypotheses of \sref{0.21.4.7}, the canonical homomorphism
\begin{equation}
\tag{21.4.8.1}
  \label{0.21.4.8.1}
  (\Omega^1_{L/K})^{(p)}\to\Upsilon_{L/K}
\end{equation}
is bijective; if $(x_\lambda)$ is a $p$-base of $L$ over $K$, the elements $d_K(x^p_\lambda)\otimes 1$ form a base of the $L$-vector space $\Upsilon_{L/K}$.
\end{proposition}

\begin{proof}
This is a particular case of \sref{0.21.3.5}.
\end{proof}

\subsection{Application: criteria for separability}
\label{subsection:0.21.5}

In this and the following numbers, \emph{we no longer assume} that the rings considered are of characteristic $p>0$.

\begin{env}[21.5.1]
\label{0.21.5.1}
Let us first note that the criterion \sref{0.21.2.7} allows us to prove one part of the Cohen theorem on separable extensions \sref{0.19.6.1}, namely that if $k$ is a field of characteristic $p>0$ and if $K$ is a separable extension of $k$, then $K$ is a formally smooth $k$-algebra.
Indeed, $K$ admits a $p$-base over $k$ \sref{0.21.4.2}, and on the other hand, it follows from the criterion of MacLane (Bourbaki, \emph{Alg.}, chap.~V, \S~8, no.~2, prop.~3) that in an algebraic closure of $K$, $k^{p^{-1}}$ and $K$ are linearly disjoint over $k$, and consequently the canonical homomorphism $k^{p^{-1}}\otimes_k K\to k^{p^{-1}}(K)$ is bijective, which is precisely the condition (i) of \sref{0.21.2.7}, after transport of structure by the isomorphism $x\mapsto x^p$.
\end{env}

\begin{proposition}[21.5.2]
\label{0.21.5.2}
\textnormal{(MacLane).}
Let $A$, $B$ be two complete discrete valuation rings, $\mathfrak{m}$, $\mathfrak{n}$ their respective maximal ideals, $K=A/\mathfrak{m}$ and $L=B/\mathfrak{n}$ their residue fields, $u:A\to B$ a homomorphism such that $Bu(\mathfrak{m})=\mathfrak{n}$, $u_0=u\otimes 1:K\to L$ the corresponding homomorphism for the residue fields.
Consider the following conditions:
\begin{enumerate}
  \item[\rm a)] $L$ is a separable extension of $K$ \textup{(}for $u_0$\textup{)}.
  \item[\rm b)] For every complete discrete valuation ring $B'$ with maximal ideal $\mathfrak{n}'$ and residue field $L'$,
    every homomorphism $u':A\to B'$ such that $B'u'(\mathfrak{m})=\mathfrak{n}'$ and every $K$-isomorphism $\sigma:L\to L'$ \textup{(}relative to $u_0$ and $u'_0=u'\otimes 1:K\to L'$\textup{)}, there exists an isomorphism $w:B\to B'$ such that $u'=w\circ u$ and that $w_0=w\otimes 1:L\to L'$ is equal to $\sigma$.
  \item[\rm b$'$)] For every homomorphism $u':A\to B$ such that $Bu'(\mathfrak{m})=\mathfrak{n}$ and such that $u'_0=u'\otimes 1:K\to L$ is equal to $u_0$, there exists an automorphism $w$ of $B$ such that $u'=w\circ u$ and that $w_0=w\otimes 1:L\to L$ is the identity.
  \item[\rm c)] Denoting by $u_1:A/\mathfrak{m}^2\to B/\mathfrak{n}^2$ the homomorphism deduced from $u$ by passage to quotients, then for every local homomorphism $u'_1:A/\mathfrak{m}^2\to B/\mathfrak{n}^2$ such that $\operatorname{gr}_0(u'_1)=\operatorname{gr}_0(u_1)(=u_0)$, there exists an automorphism $w_1$ of $B/\mathfrak{n}^2$ such that $\operatorname{gr}_0(w_1)$ is the identity and that $u'_1=w_1\circ u_1$.
  \item[\rm c$'$)] For every local homomorphism $u'_1:A/\mathfrak{m}^2\to B/\mathfrak{n}^2$ such that $\operatorname{gr}_0(u'_1)=\operatorname{gr}_0(u_1)$ and $\operatorname{gr}_1(u'_1)=\operatorname{gr}_1(u_1)$ \textup{(}homomorphism of $\mathfrak{m}/\mathfrak{m}^2$ into $\mathfrak{n}/\mathfrak{n}^2$\textup{)}, there exists an automorphism $w_1$ of $B/\mathfrak{n}^2$ such that $\operatorname{gr}_0(w_1)$ is the identity and that $u'_1=w_1\circ u_1$.
\end{enumerate}
Then one has the implications $c)\Leftrightarrow c')\Leftrightarrow a)\Rightarrow b)\Rightarrow b')$.

If moreover $A$ is a Cohen ring, the five conditions are equivalent.
\end{proposition}

\begin{proof}
The implications $b)\Rightarrow b')$ and $c)\Rightarrow c')$ are trivial.
Let us show that $a)$ implies $b)$.
The homomorphism $u$ makes $B$ into a torsion-free $A$-module, since it transforms a uniformizer of $A$ into a uniformizer of $B$; it follows that $B$ is a \emph{flat} $A$-module ($\mathbf{0_I}$, 6.3.4), hence a \emph{Cohen $A$-algebra} \sref{0.19.8.1}; the fact that $a)$ implies $b)$ is then a consequence of \sref{0.19.8.2}, (i) applied to $C=B'$, $\mathfrak{J}=\mathfrak{n}'$, $B'$ being considered as an $A$-algebra for $u'$ and the homomorphism $B\to B'/\mathfrak{n}'$ being the composite $B\to B/\mathfrak{n}\xrightarrow{\sigma} B'/\mathfrak{n}'$, which is an $A$-homomorphism by virtue of the hypothesis $u'_0=\sigma\circ u_0$.
The same reasoning proves that $a)$ implies $c)$, taking $C=B/\mathfrak{n}^2$, $\mathfrak{J}=\mathfrak{n}/\mathfrak{n}^2$, $C$ being considered as an $A$-algebra for $u'_1$.

Let us prove next that $c')$ implies $a)$.
The two homomorphisms $u_1$ and $u'_1$ are such that $u'_1=u_1+D$, where $D$ is a derivation of $A/\mathfrak{m}^2$ into $\mathfrak{n}/\mathfrak{n}^2$ \sref{0.20.1.1}; moreover, the hypothesis $\operatorname{gr}_1(u'_1)=\operatorname{gr}_1(u_1)$ means that $D$ is zero in $\mathfrak{m}/\mathfrak{m}^2$, and can thus be considered as a derivation of $K=A/\mathfrak{m}$ into $\mathfrak{n}/\mathfrak{n}^2$, and $\mathfrak{n}/\mathfrak{n}^2$ identifies (by choice of a uniformizer of $B$) with the $K$-module $L$ (for $u_0$).
On the other hand, the conditions imposed on $w_1$ imply that $\operatorname{gr}_1(w_1)$ is also the identity (since $u_1(\mathfrak{m}/\mathfrak{m}^2)$ generates $\mathfrak{n}/\mathfrak{n}^2$);
$w_1$ is thus of the form $z\mapsto z+D'(z)$, where $D'$ is this time a derivation of $B/\mathfrak{n}^2$ in the $(B/\mathfrak{n}^2)$-module $\mathfrak{n}/\mathfrak{n}^2$; as $w_1$ is the identity on $\mathfrak{n}/\mathfrak{n}^2$, $D'$ can again (by the preceding identification of $\mathfrak{n}/\mathfrak{n}^2$ and of $L$) be considered as a derivation \emph{of $L$ into $L$};
finally, the relation $u'_1=w_1\circ u_1$ means that $D'$ \emph{extends} $D$.
Let us note, on the other hand, that every derivation $D$ of $K$ into $L$ corresponds to a homomorphism $u'_1$ satisfying the conditions of $c')$ \sref{0.20.1.1}; condition $c')$ therefore means that every derivation of $K$ into $L$ extends to a derivation of $L$ into itself, that is to say \sref{0.20.6.5}, that $L$ is separable over $K$.

Let us prove finally that, when $A$ is a Cohen ring, $b')$ implies $c')$.
Let us show first that under the hypotheses of $c')$, there exists a homomorphism $u':A\to B$ satisfying the hypotheses of $b'$) and that, by passage to quotients, gives $u'_1:A/\mathfrak{m}^2\to B/\mathfrak{n}^2$.
Indeed, this follows from \sref{0.19.8.6}, (i) applied to the composite homomorphism $v':A\to A/\mathfrak{m}^2\xrightarrow{u'_1} B/\mathfrak{n}^2$;
\oldpage[0\textsubscript{IV-1}]{166}
this last thus factors as $A\xrightarrow{u'} B\to B/\mathfrak{n}^2$ and the hypotheses on $\operatorname{gr}_0(u'_1)$ and $\operatorname{gr}_1(u'_1)$ imply $\operatorname{gr}_0(u')=\operatorname{gr}_0(u)=u_0$ and $\operatorname{gr}_1(u')=\operatorname{gr}_1(u)$, hence the image by $u'$ of a uniformizer of $A$ is a uniformizer of $B$, and therefore $Bu'(\mathfrak{m})=\mathfrak{n}$.
It then suffices, to obtain an automorphism $w_1$ satisfying the question, to take the automorphism deduced by passage to quotients from the automorphism $w$ of $B$ provided by the application of $b')$ to the homomorphism~$u'$.
\end{proof}

\begin{remarks}[21.5.3]
\label{0.21.5.3}
\begin{enumerate}
  \item[(i)] The differential properties of fields allow us to solve the question of the \emph{uniqueness} of the representative field in a complete local Noetherian ring $C$ \sref{0.19.8.7}, (ii).
    Let $\mathfrak{J}$ be the maximal ideal of $C$, $K=C/\mathfrak{J}$ the residue field of $C$; we can restrict to the case $\mathfrak{J}\neq 0$.
    Suppose there exists a homomorphism $u:K\to C$ that, composed with the augmentation $C\to K$, gives the identity; then, \emph{for this homomorphism to be unique, it is necessary and sufficient that} $\Omega^1_K=0$.
    Indeed, the condition $\Omega^1_K=0$ implies that $K$ is formally unramified on its prime field \sref{0.20.7.4}; if there existed a second homomorphism $v\neq u$ satisfying the question, there would exist a largest integer $n$ such that $\mathfrak{J}^n$ contains the set of the $u(x)-v(x)$ for $x\in K$; by passage to the quotient, $u$ and $v$ would give two distinct homomorphisms $u'$, $v'$ of $K$ into $C/\mathfrak{J}^{n+1}$, whose composites with $C/\mathfrak{J}^{n+1}\to C/\mathfrak{J}^n$ would be equal, which contradicts the definition \sref{0.19.10.2}.
    Conversely, suppose $\Omega^1_K\neq 0$; then there exists a derivation $D\neq 0$ of $K$ into $\mathfrak{J}/\mathfrak{J}^2$ \sref{0.20.4.8}, hence a homomorphism $v_1:K\to C/\mathfrak{J}^2$ such that, if $u_1:K\to C/\mathfrak{J}^2$ is obtained by passage to quotients from $u$, one has $v_1=u_1+D$ \sref{0.20.1.1}.
    If $k$ is the prime field of $K$, $C$ is a $k$-algebra and $K$ a formally smooth $k$-algebra \sref{0.19.6.1}, and the restrictions of $u_1$ and $v_1$ to $k$ coincide, hence $v_1$ factors as $K\xrightarrow{v} C\to C/\mathfrak{J}^2$ and one has $v\neq u$.

    Let us recall \sref{0.20.6.20} and \sref{0.21.4.4} that the condition $\Omega^1_K=0$ means that $K$ is \emph{perfect} if it is of characteristic $\neq 0$, and an algebraic extension of $\mathbf{Q}$ if it is of characteristic~$0$.
  \item[(ii)] In the same way, let $W$ be a Cohen ring, $C$ a complete local Noetherian ring, $\mathfrak{J}$ a nonzero ideal of $C$ contained in the maximal ideal; for the factorization $W\xrightarrow{v} C\to C/\mathfrak{J}$ in \sref{0.19.8.6}, (i) to be \emph{unique}, it is necessary and sufficient that the residue field $K$ of the Cohen ring $W$ is such that $\Omega^1_K=0$.
    Indeed, if $\Omega^1_K\neq 0$ and if $\mathfrak{m}$ denotes the maximal ideal of $W$, it suffices to compose a derivation $D\neq 0$ of $K$ into $\mathfrak{J}/\mathfrak{m}\mathfrak{J}$ to obtain a derivation $D'\neq 0$ of $W$ into $\mathfrak{J}/\mathfrak{m}\mathfrak{J}$ and one finishes the reasoning as in (i) by forming the homomorphism $v_1:W\to C/\mathfrak{m}\mathfrak{J}$ distinct from the homomorphism $u_1:W\to C/\mathfrak{m}\mathfrak{J}$ deduced from $u$ by passage to quotients; one finishes by invoking \sref{0.19.8.6}, (i).
    If on the contrary $\Omega^1_K=0$, the uniqueness of $v$ already follows from (i) when $W$ is a field of characteristic $0$.
    In the opposite case, one has $\widehat{\Omega}^1_W=0$; in fact, one then has $\mathfrak{m}=pW$, $p$ being the characteristic of $K$ \sref{0.19.8.5}, and the canonical homomorphism $\mathfrak{m}/\mathfrak{m}^2\to\Omega^1_W/\mathfrak{m}\Omega^1_W$ \sref{0.20.5.11.2} is consequently zero. The exact sequence \sref{0.20.5.12.1} applied to $W$ and to $K=W/\mathfrak{m}$ then implies that $\Omega^1_W=\mathfrak{m}\Omega^1_W$, whence our assertion.
    But then \sref{0.20.7.4} $W$ is formally unramified (for its adic topology) over $\mathbf{Z}$, and the uniqueness of $v$ is proved as in (i).
\end{enumerate}
\end{remarks}

\subsection{Admissible fields for an extension}
\label{subsection:0.21.6}

\begin{env}[21.6.1]
\label{0.21.6.1}
Given four fields $k_0\subset k\subset K\subset L$, it follows from \sref{0.20.6.16} and \sref{0.20.6.17} that one has an exact sequence
\begin{equation}
\tag{21.6.1.1}
  \label{0.21.6.1.1}
  0\to\Upsilon_{K/k/k_0}\otimes_K L\xrightarrow{v}\Upsilon_{L/k/k_0}\xrightarrow{u}\Upsilon_{L/K/k_0}\xrightarrow{s}\Upsilon_{L/K/k}\to 0.
\end{equation}

When one fixes $k_0$, $K$, and $L$ and one lets the intermediate field $k$ between $k_0$ and $K$ ``vary'', one evidently has $\Upsilon_{L/K/k_0}=\Upsilon_{L/K/k}$ when $k=k_0$.
When the canonical homomorphism $s$ of \sref{0.21.6.1.1} is still \emph{bijective}, one says that $k$ is a \emph{$k_0$-admissible field for the extension $L$ of $K$}.
The interest of the existence, under certain conditions, of such fields $k$, which are however ``sufficiently close'' to $K$ (for example such that $[K:k]$ is finite) is that they allow us to replace in certain questions the modules of differentials $\Omega^1_{K/k_0}$ and $\Omega^1_{L/k_0}$ (which can be ``too large'', for example when $k_0$ is the prime field) by $\Omega^1_{K/k}$ and $\Omega^1_{L/k}$, which are more easily handled.

When $k_0$ is the \emph{prime field}, one says ``admissible field'' instead of ``$k_0$-admissible field''.

We set
\end{env}

\begin{env}[21.6.1.2]
\label{0.21.6.1.2}
\[
  \Delta(L/K,k/k_0)=\operatorname{Coker}(\Upsilon_{K/k/k_0}\otimes_K L\to\Upsilon_{L/k/k_0})\simeq\operatorname{Ker}(\Upsilon_{L/K/k_0}\to\Upsilon_{L/K/k})
\]
(vector space over $L$); its rank will be denoted $d(L/K,k/k_0)$ and called the \emph{defect of $k_0$-admissibility of $k$ for the extension $L$ of $K$} (it is evidently zero if and only if $k$ is $k_0$-admissible for this extension).
When $k_0$ is the prime field, one writes $\Delta(L/K,k)$ and $d(L/K,k)$ instead of $\Delta(L/K,k/k_0)$ and $d(L/K,k/k_0)$.
\end{env}

\begin{proposition}[21.6.2]
\label{0.21.6.2}
Let $k_0\subset k\subset K\subset L$ be four fields.
\begin{enumerate}
  \item[\rm (i)] The following conditions are equivalent:
  \begin{enumerate}
    \item[\rm a)] The field $k$ is $k_0$-admissible for the extension $L$ of $K$ \textup{(}in other words, the homomorphism $\Upsilon_{L/K/k_0}\to\Upsilon_{L/K/k}$ is injective, hence bijective\textup{)}.
    \item[\rm b)] The canonical homomorphism $u:\Upsilon_{K/k/k_0}\to\Upsilon_{L/k/k_0}$ is zero.
    \item[\rm c)] The canonical homomorphism $v:\Upsilon_{K/k/k_0}\otimes_K L\to\Upsilon_{L/k/k_0}$ is surjective \textup{(}hence bijective\textup{)}.
    \item[\rm d)] One has $d(L/K,k/k_0)=0$ \textup{(}or $\Delta(L/K,k/k_0)=0$\textup{)}.
  \end{enumerate}
  \item[\rm (ii)] The equivalent conditions of \textnormal{(i)} are satisfied when one is in one of the following cases: $\alpha)$ $L$ is separable over $k$; $\beta)$ $L$ is separable over $K$; $\gamma)$ one has $k\subset k_0(K^p)$, where $p$ denotes the characteristic exponent of $k_0$.
\end{enumerate}
\end{proposition}

\begin{proof}
\begin{enumerate}
  \item[(i)] The assertions follow trivially from the exactness of the sequence \sref{0.21.6.1.1}.
  \item[(ii)] If $L$ is separable over $k$, one has $\Upsilon_{L/k/k_0}=0$ \sref{0.20.6.19}, hence condition $b)$ of (i) is satisfied; if $L$ is separable over $K$, one has $\Upsilon_{L/K/k_0}=0$ \sref{0.20.6.19}, hence condition $a)$ of (i) is satisfied; finally, if $k\subset k_0(K^p)$, it follows that $\Upsilon_{L/K/k_0}=\Upsilon_{L/K/k}$ by virtue of \sref{0.21.1.5.1}, and the condition $a)$ of (i) is satisfied.
\end{enumerate}
\end{proof}

\begin{env}[21.6.3]
\label{0.21.6.3}
Suppose that one has a commutative diagram of monomorphisms of fields
\[
  \xymatrix{
    k'_0\ar[r] & k'\ar[r] & K'\ar[r] & L'\\
    k_0\ar[u]\ar[r] & k\ar[u]\ar[r] & K\ar[u]\ar[r] & L\ar[u]
  }
\]
It then follows from \sref{0.20.6.17}, (ii) that one has a canonical homomorphism
\begin{equation}
\tag{21.6.3.1}
  \label{0.21.6.3.1}
  \Delta(L/K,k/k_0)\to\Delta(L'/K',k'/k'_0)
\end{equation}
with an evident transitivity property, so that one can say that $\Delta(L/K,k/k_0)$ is a \emph{functor} in the quadruplet $(k_0,k,K,L)$.
\end{env}

\begin{proposition}[21.6.4]
\label{0.21.6.4}
\begin{enumerate}
  \item[\rm (i)] Let $k\subset k'\subset k''\subset K\subset L$ be five fields.
  One has an exact sequence of canonical homomorphisms
  \begin{equation}
  \tag{21.6.4.1}
  \label{0.21.6.4.1}
    0\to\Delta(L/K,k'/k)\to\Delta(L/K,k''/k)\to\Delta(L/K,k''/k')\to 0
  \end{equation}
  and consequently the equality
  \begin{equation}
  \tag{21.6.4.2}
  \label{0.21.6.4.2}
    d(L/K,k''/k)=d(L/K,k''/k')+d(L/K,k'/k).
  \end{equation}
  \item[\rm (ii)] Let $k_0\subset k\subset K\subset L\subset M$ be five fields.
  One has an exact sequence of canonical homomorphisms
  \begin{equation}
  \tag{21.6.4.3}
  \label{0.21.6.4.3}
    0\to\Delta(L/K,k/k_0)\otimes_L M\to\Delta(M/K,k/k_0)\to\Delta(M/L,k/k_0)\to 0
  \end{equation}
  and consequently the equality
  \begin{equation}
  \tag{21.6.4.4}
  \label{0.21.6.4.4}
    d(M/K,k/k_0)=d(M/L,k/k_0)+d(L/K,k/k_0).
  \end{equation}
\end{enumerate}
\end{proposition}

\begin{proof}
\begin{enumerate}
  \item[(i)] Consider the commutative diagram
  \[
    \xymatrix{
      0\ar[r] & \Upsilon_{K/k'/k}\otimes_K L\ar[r]\ar[d] & \Upsilon_{K/k''/k}\otimes_K L\ar[r]\ar[d] & \Upsilon_{K/k''/k'}\otimes_K L\ar[r]\ar[d] & 0\\
      0\ar[r] & \Upsilon_{L/k'/k}\ar[r]\ar[d] & \Upsilon_{L/k''/k}\ar[r]\ar[d] & \Upsilon_{L/k''/k'}\ar[r]\ar[d] & 0\\
      0\ar[r] & \Delta(L/K,k'/k)\ar[r] & \Delta(L/K,k''/k)\ar[r] & \Delta(L/K,k''/k')\ar[r] & 0
    }
  \]
  where one can consider the three rows as \emph{complexes} $T^\bullet_1$, $T^\bullet_2$, $T^\bullet_3$ respectively; the exact sequence \sref{0.21.6.1.1} and the definition of $\Delta$ \sref{0.21.6.1.2} show that one has an exact sequence of complexes $0\to T^\bullet_1\to T^\bullet_2\to T^\bullet_3\to 0$; applying the exact cohomology sequence to it, and noting that, by virtue of the exactness of \sref{0.21.6.1.1}, the cohomology of $T^\bullet_1$ and that of $T^\bullet_2$ are zero except in a single degree, for which the cohomology modules are both equal to $\Upsilon_{k''/k'/k}\otimes_{k''}L$; as $T^\bullet_1$ and $T^\bullet_2$ have the same cohomology, that of $T^\bullet_3$ is necessarily zero, which proves (i).
  \item[(ii)] Consider similarly the commutative diagram
  \[
    \xymatrix{
      0\ar[r] & \Delta(L/K,k/k_0)\otimes_L M\ar[r]\ar[d] & \Delta(M/K,k/k_0)\ar[r]\ar[d] & \Delta(M/L,k/k_0)\ar[r] & 0\\
      0\ar[r] & \Upsilon_{L/K/k_0}\otimes_L M\ar[r]\ar[d] & \Upsilon_{M/K/k_0}\ar[r]\ar[d] & \Upsilon_{M/L/k_0}\ar[r] & 0\\
      0\ar[r] & \Upsilon_{L/K/k}\otimes_L M\ar[r] & \Upsilon_{M/K/k}\ar[r] & \Upsilon_{M/L/k}\ar[r] & 0
    }
  \]
  where again one considers the three rows as \emph{complexes} $T'^\bullet_1$, $T'^\bullet_2$, and $T'^\bullet_3$; the exact sequence \sref{0.21.6.1.1} and the definition of $\Delta$ \sref{0.21.6.1.2} give here a short exact sequence of complexes $0\to T'^\bullet_1\to T'^\bullet_2\to T'^\bullet_3\to 0$; applying the exact cohomology sequence to it again, this time, by virtue of the exactness of \sref{0.21.6.1.1}, the cohomology of $T'^\bullet_2$ and that of $T'^\bullet_3$ are zero except in a single degree, for which the cohomology modules are both equal to $\Upsilon_{M/L/k_0}$; one concludes therefore that the cohomology of $T'^\bullet_1$ is zero, which establishes (ii).
\end{enumerate}
\end{proof}

\begin{corollary}[21.6.5]
\label{0.21.6.5}
\begin{enumerate}
  \item[\rm (i)] Given five fields $k\subset k'\subset k''\subset K\subset L$, for $k''$ to be $k$-admissible for the extension $L$ of $K$, it is necessary and sufficient that $k'$ be $k$-admissible and that $k''$ be $k'$-admissible for the extension $L$ of $K$.
  \item[\rm (ii)] Given five fields $k_0\subset k\subset K\subset L\subset M$, for $k$ to be $k_0$-admissible for the extension $M$ of $K$, it is necessary and sufficient that it be so for the extension $L$ of $K$ and for the extension $M$ of $L$.
\end{enumerate}

This follows immediately from the relations \sref{0.21.6.4.2} and \sref{0.21.6.4.4}, the values of $d$ being $\geqslant 0$.
\end{corollary}

\begin{corollary}[21.6.6]
\label{0.21.6.6}
Let $k_0\subset k\subset K\subset L$ be four fields, and suppose that $k$ is $k_0$-admissible for the extension $L$ of $K$.
Then, if $k'_0$, $k'$, $K'$, $L'$ are four fields such that $k_0\subset k'_0\subset k'\subset k\subset K\subset K'\subset L'\subset L$, $k'$ is $k'_0$-admissible for the extension $L'$ of $K'$.
\end{corollary}

\subsection{Cartier's equality}
\label{subsection:0.21.7}

The following result translates in terms of differentials a theorem of MacLane on derivations:

\begin{theorem}[21.7.1]
\label{0.21.7.1}
\textnormal{(Cartier).}
Let $K$ be a field, $L$ a finite type extension of $K$.
Then $\Omega^1_{L/K}$ and $\Upsilon_{L/K}$ are $L$-vector spaces of finite rank, and one has
\begin{equation}
\tag{21.7.1.1}
  \label{0.21.7.1.1}
  \operatorname{rg}\Omega^1_{L/K}-\operatorname{rg}\Upsilon_{L/K}=\deg.\operatorname{tr}_K L.
\end{equation}
\end{theorem}

\begin{proof}
If $L'$ is a field such that $K\subset L'\subset L$, one has the exact sequence \sref{0.20.6.15.1} (applied to $\Lambda=P$, $A=K$, $B=L'$, $C=L$)
\[
  0\to\Upsilon_{L'/K}\otimes_{L'} L\to\Upsilon_{L/K}\to\Omega^1_{L'/K}\otimes_{L'} L\to\Omega^1_{L/K}\to\Omega^1_{L/L'}\to 0
\]
hence ($\mathbf{0_{III}}$, 11.10.2)
\[
  \operatorname{rg}_L\Omega^1_{L/K}-\operatorname{rg}_L\Upsilon_{L/K}=(\operatorname{rg}_{L'}\Omega^1_{L/L'}-\operatorname{rg}_{L'}\Upsilon_{L/L'})+(\operatorname{rg}_{L'}\Omega^1_{L'/K}-\operatorname{rg}_{L'}\Upsilon_{L'/K}).
\]
Since moreover $\deg.\operatorname{tr}_K L=\deg.\operatorname{tr}_K L'+\deg.\operatorname{tr}_{L'} L$, one sees, by induction on the number of generators of the extension $L$, that one is reduced to proving \sref{0.21.7.1.1} when $L=K(x)$.
We then distinguish three cases:

\emph{a)} $x$ is transcendent over $K$; as $L$ is separable over $K$, one has $\Upsilon_{L/K}=0$ \sref{0.20.6.3}; on the other hand, \sref{0.20.5.9} and \sref{0.20.4.13} show that $\Omega^1_{L/K}$ is of rank $1$, hence \sref{0.21.7.1.1} in this case.

\emph{b)} $L$ is a finite algebraic separable extension of $K$, so that one has again $\Upsilon_{L/K}=0$ \sref{0.20.6.3}.
On the other hand, one has $\Omega^1_{L/K}=0$ by \sref{0.20.6.20}, hence again \sref{0.21.7.1.1}.

\emph{c)} $L$ is a finite algebraic inseparable extension of $K$; the reasoning at the beginning shows that one can restrict to the case where $L^p\subset K$; it then follows from \sref{0.21.4.8} that one has $\operatorname{rg}_L\Upsilon_{L/K}=\operatorname{rg}_L\Omega^1_{L/K}$, hence again \sref{0.21.7.1.1}.
\end{proof}

\begin{corollary}[21.7.2]
\label{0.21.7.2}
Let $K$ be a field, $L$ a finite type extension of $K$, $k$ a subfield of $K$.
Then $\Omega^1_{L/K}$ and $\Upsilon_{L/K/k}$ are vector spaces of finite rank over $L$, and one has
\begin{equation}
\tag{21.7.2.1}
  \label{0.21.7.2.1}
  \operatorname{rg}\Omega^1_{L/K}-\operatorname{rg}\Upsilon_{L/K/k}=\deg.\operatorname{tr}_K L+d(L/K,k).
\end{equation}

Moreover, one has the inequality
\begin{equation}
\tag{21.7.2.2}
  \label{0.21.7.2.2}
  \operatorname{rg}\Omega^1_{L/K}-\operatorname{rg}\Upsilon_{L/K/k}\geqslant\deg.\operatorname{tr}_K L.
\end{equation}

Furthermore, for the two members of \sref{0.21.7.2.2} to be equal, it is necessary and sufficient that $k$ be an admissible field for the extension $L$ of $K$.
\end{corollary}

\begin{proof}
Indeed, as the homomorphism $\Upsilon_{L/K}\to\Upsilon_{L/K/k}$ is surjective, one has by definition \sref{0.21.6.1.2} $\operatorname{rg}\Upsilon_{L/K}-\operatorname{rg}\Upsilon_{L/K/k}=d(L/K,k)$, and the corollary follows immediately from \sref{0.21.7.1} and from \sref{0.21.6.2}.
\end{proof}

\begin{corollary}[21.7.3]
\label{0.21.7.3}
Let $k\subset K\subset L$ be three fields such that $L$ is a finite type extension of $k$.
Then one has
\begin{equation}
\tag{21.7.3.1}
  \label{0.21.7.3.1}
  \operatorname{rg}_L\Omega^1_{L/k}-\operatorname{rg}_K\Omega^1_{K/k}=\deg.\operatorname{tr}_K L+d(L/K,k).
\end{equation}

\oldpage[0\textsubscript{IV-1}]{171}
Moreover, one has the inequality
\begin{equation}
\tag{21.7.3.2}
  \label{0.21.7.3.2}
  \operatorname{rg}_L\Omega^1_{L/k}-\operatorname{rg}_K\Omega^1_{K/k}\geqslant\deg.\operatorname{tr}_K L
\end{equation}
the equality being attained if and only if $k$ is an admissible field for the extension $L$ of $K$.
\end{corollary}

\begin{proof}
One has the exact sequence \sref{0.20.6.1.1}
\[
  0\to\Upsilon_{L/K/k}\to\Omega^1_{K/k}\otimes_K L\to\Omega^1_{L/k}\to\Omega^1_{L/K}\to 0
\]
hence the first member of \sref{0.21.7.3.1} is equal to
\[
  \operatorname{rg}_L\Omega^1_{L/k}-\operatorname{rg}_L\Upsilon_{L/K/k}
\]
and it suffices to apply \sref{0.21.7.2}.
\end{proof}

The interest of this last corollary is that it only involves modules of differentials, to the exclusion of modules of imperfection.
We shall see moreover below \sref{0.21.8.6} that for every finite type extension $L$ of $K$, there exists a subfield $k$ of $K$ such that $[K:k]$ is finite, and which is admissible for the extension $L$ of $K$, so that the two members of \sref{0.21.7.3.2} are then equal.

\begin{corollary}[21.7.4]
\label{0.21.7.4}
Let $K$ be a finite type extension of a field $k$.
The following conditions are equivalent:
\begin{enumerate}
  \item[\rm a)] $K$ is a finite separable extension of $k$.
  \item[\rm b)] $K$ is a formally unramified $k$-algebra \sref{0.19.10.2}.
  \item[\rm c)] $K$ is a formally \'etale $k$-algebra \sref{0.19.10.2}.
  \item[\rm d)] One has $\Omega^1_{K/k}=0$.
\end{enumerate}
\end{corollary}

\begin{proof}
Indeed, one knows (the topologies being discrete) that the conditions $b)$ and $d)$ are equivalent \sref{0.20.7.4} and that $a)$ implies that $K$ is a formally smooth $k$-algebra \sref{0.19.6.1}, hence the conjunction of $a)$ and $b)$ is equivalent to $c)$.
The whole thing reduces to seeing that $d)$ implies $a)$.
Now, by virtue of \sref{0.21.7.1.1}, the relation $\Omega^1_{K/k}=0$ implies that $K$ is algebraic (hence finite) over $k$ and that $\Upsilon_{K/k}=0$, that is to say \sref{0.20.6.3} that $K$ is separable over $k$.
\end{proof}

\begin{remark}[21.7.5]
\label{0.21.7.5}
By virtue of \sref{0.20.6.5.2} and ($\mathbf{0_{III}}$, 11.10.2) the first member of \sref{0.21.7.1.1} is none other than the \emph{Euler-Poincar\'e characteristic} of the complex $\mathrm{K}_\bullet(C/B/A)$ introduced in \sref{0.20.6.5}, for $C=L$, $B=K$ and $A=P$ (prime field of $K$).\footnote{The French source prints ``(20.6.3.2)'', which does not exist; the homology identities used here are (20.6.5.2).}
In the chapter devoted to intersection theory and to the theorem of Riemann-Roch, an important role will be played equally by generalized Euler-Poincar\'e characteristics (with values in groups of classes of $\sh{O}_X$-Modules) of complexes generalizing the complexes $\mathrm{F}_\bullet(C/A)$ considered in \sref{0.20.6.22}.
\end{remark}

\subsection{Criteria for admissibility}
\label{subsection:0.21.8}

We return to our previous conventions and suppose therefore that all fields considered in this number are of characteristic $p>0$.

\begin{lemma}[21.8.1]
\label{0.21.8.1}
Let $K$ be a field, $k$ a subfield of $K$, $(k_\alpha)_{\alpha\in I}$ a decreasing filtering family of subfields of $K$ such that $k=\bigcap_\alpha k_\alpha$.
Let $V$ be a vector space over $K$, $(a_i)_{1\leqslant i\leqslant n}$ a finite family of vectors of $V$; if the family $(a_i)$ is free over $k$, there exists an index $\gamma$ such that it is also free over $k_\gamma$.
\end{lemma}

\begin{proof}
\oldpage[0\textsubscript{IV-1}]{172}
Let $r$ be the rank of the family $(a_i)_{1\leqslant i\leqslant n}$ over $K$, and let us reason by induction on $n-r$; the proposition is evident for $n=r$, for the family $(a_i)_{1\leqslant i\leqslant r}$ is then free over $K$, hence over every subfield of $K$.
Suppose for example that $(a_i)_{1\leqslant i\leqslant r}$ is free over $K$, and write $a_{r+1}=\sum_{i=1}^r\lambda_i a_i$ with $\lambda_i\in K$; the family $(a_i)_{1\leqslant i\leqslant r+1}$ being free over $k$, the $\lambda_i$ cannot all belong to $k$; suppose for example that $\lambda_{i_0}\notin k$.
Then there is an index $\beta$ such that $\lambda_{i_0}\notin k_\beta$; one concludes that the family $(a_i)_{1\leqslant i\leqslant r+1}$ is \emph{free over $k_\beta$}; indeed, if the family $(a_i)_{1\leqslant i\leqslant r+1}$ were not free over $k_\beta$, $a_{r+1}$ would be equal to a linear combination of the $a_i$ with coefficients in $k_\beta$, and as these coefficients are necessarily the $\lambda_i$, one would arrive at a contradiction.
It now suffices to apply the induction hypothesis by replacing $K$ by $k_\beta$ and the family $(k_\alpha)_{\alpha\in I}$ by the sub-family of the $k_\alpha$ contained in $k_\beta$.
\end{proof}

\begin{lemma}[21.8.2]
\label{0.21.8.2}
Let $K$ be a field, $p$ its exponential characteristic, $k_0$ a subfield of $K$, $(k_\alpha)_{\alpha\in I}$ a decreasing filtering family of subfields of $K$ such that $\bigcap_\alpha k_\alpha(K^p)=k_0(K^p)$.
If $(x_i)_{1\leqslant i\leqslant n}$ is a finite family of elements of $K$ that is $p$-free over $k_0$, there exists an index $\alpha$ such that $(x_i)$ is also $p$-free over $k_\alpha$.
\end{lemma}

\begin{proof}
Indeed, to say that the family $(x_i)$ is $p$-free over a subfield $k$ of $K$ means that the finite family of monomials $\prod_{i=1}^n x_i^{m(i)}$ with $0\leqslant m(i)<p$ is free over $k(K^p)$; it suffices therefore to apply the lemma \sref{0.21.8.1} to this family of monomials in the vector space $V=K$, and to the subfields $k_\alpha(K^p)$ and $k_0(K^p)$ of $K$.
\end{proof}

\begin{theorem}[21.8.3]
\label{0.21.8.3}
Let $K$ be a field of characteristic $p>0$, $k_0$ a subfield of $K$, $(k_\alpha)_{\alpha\in I}$ a decreasing filtering family of subfields of $K$ containing $k_0$.
The following conditions are equivalent:
\begin{enumerate}
  \item[\rm a)] $\bigcap_\alpha k_\alpha(K^p)=k_0(K^p)$.
  \item[\rm b)] For every extension $L$ of $K$ such that $\Upsilon_{L/K/k_\alpha}$ is an $L$-vector space of finite rank \textup{(}which holds in particular if $L$ is a finite type extension by virtue of \sref{0.21.7.2}\textup{)}, there exists $\alpha\in I$ such that $k_\alpha$ is $k_0$-admissible for the extension $L$ of $K$.
  \item[\rm b$'$)] For every extension $L=K(x)$ of $K$, with $x^p\in K$, there exists $\alpha\in I$ such that $k_\alpha$ is $k_0$-admissible for the extension $L$ of $K$.
  \item[\rm c)] The canonical map
  \begin{equation}
  \tag{21.8.3.1}
  \label{0.21.8.3.1}
    \Omega^1_{K/k_0}\to\varprojlim_\alpha\Omega^1_{K/k_\alpha}
  \end{equation}
  is injective.
\end{enumerate}
\end{theorem}

\begin{proof}
The canonical map \sref{0.21.8.3.1} is understood as passage to the projective limit in the projective system of homomorphisms $\Omega^1_{K/k_0}\to\Omega^1_{K/k_\alpha}$ \sref{0.20.5.3.3}.
We shall prove the logical scheme $c)\Rightarrow b)\Rightarrow b')\Rightarrow a)\Rightarrow c)$.

To say that $k_\alpha$ is $k_0$-admissible for the extension $L$ of $K$ means that the canonical homomorphism $\Upsilon_{L/K/k_0}\to\Upsilon_{L/K/k_\alpha}$ is injective; or, if $N_\alpha$ is the kernel of this homomorphism, the $N_\alpha$ form a decreasing filtering system of sub-vector spaces of $\Upsilon_{L/K/k_0}$, and as $\Upsilon_{L/K/k_0}$ is by hypothesis of \emph{finite} rank, it amounts to the same to say that one of the $N_\alpha$ is $0$ or that their \emph{intersection} is $0$.
But this intersection is none other than the kernel of the projective limit homomorphism $\Upsilon_{L/K/k_0}\to\varprojlim_\alpha\Upsilon_{L/K/k_\alpha}$.
Now, one has the commutative diagram
\[
  \xymatrix{
    \Upsilon_{L/K/k_0}\ar[r]\ar[d] & \varprojlim\Upsilon_{L/K/k_\alpha}\ar[d]\\
    \Omega^1_{K/k_0}\otimes_K L\ar[r] & \varprojlim(\Omega^1_{K/k_\alpha}\otimes_K L)
  }
\]
where the left vertical arrow is injective by definition.
To prove that $c)$ implies $b)$, it suffices therefore to show that the canonical map
\begin{equation}
\tag{21.8.3.2}
  \label{0.21.8.3.2}
  \Omega^1_{K/k_0}\otimes_K V\to\varprojlim_\alpha(\Omega^1_{K/k_\alpha}\otimes_K V)
\end{equation}
is injective for every vector space $V$ over $K$ (and in particular for $V=L$).
Now this is evident if $V=K^n$ since then $W\otimes_K V=W^n$ for every vector space $W$ over $K$ and projective limits and products commute.
On the other hand, for every element $z$ of the first member of \sref{0.21.8.3.2} there exists a subspace $V'$ of $V$ of finite rank such that $z\in\Omega^1_{K/k_0}\otimes_K V'$ (identified canonically with a subspace of $\Omega^1_{K/k_0}\otimes_K V$); if $z\neq 0$, its image in $\varprojlim(\Omega^1_{K/k_\alpha}\otimes_K V')$ is also non-zero; as the $\varprojlim$ functor is exact on the left, the image of $z$ in the second member of \sref{0.21.8.3.2} is therefore also $\neq 0$, which finishes the proof that $c)$ implies $b)$.

It is trivial that $b)$ implies $b')$; let us show that $b')$ implies $a)$.
Using the notations of $b')$, one can suppose $L\neq K$.
It then follows from \sref{0.21.4.7} that one has $\operatorname{rg}\Upsilon_{L/K}\leqslant 1$.
If $\Upsilon_{L/K/k_\alpha}=0$ there is nothing to prove by virtue of \sref{0.21.6.1.1}.
Otherwise, $\Upsilon_{L/K/k_\alpha}$ is identified canonically with $\Upsilon_{L/K}$, and if one sets $a=x^p$, $\Upsilon_{L/K}$ is a one-dimensional space generated by the single element $d_K(a)\otimes 1$ \sref{0.21.4.7}; $\Upsilon_{L/K/k_\alpha}$ is thus a one-dimensional space generated by the single element $d_{K/k_\alpha}(a)\otimes 1$, i.e., to say that a subfield $k_\alpha$ is $k_0$-admissible for the extension $L$ of $K$ means that one has $d_{K/k_\alpha}(a)\neq 0$, or equivalently \sref{0.21.4.6} that $a\notin k_\alpha(K^p)$.
Now, for every $a\notin k_0(K^p)$, one has $d_{K/k_0}(a)\neq 0$ and $K(a^{1/p})\neq K$, hence one can apply $b')$, which proves the existence of an $\alpha$ such that $a\notin k_\alpha(K^p)$; as this holds for every $a\notin k_0(K^p)$, we have shown $b')\Rightarrow a)$.

It remains to show that $a)$ implies $c)$.
Let $(x_\mu)_{\mu\in M}$ be a $p$-base of $K$ over $k_0$; then the $d_{K/k_0}(x_\mu)$ form a base of the $K$-vector space $\Omega^1_{K/k_0}$ \sref{0.21.4.5}, and the condition $c)$ means that for every finite subset $J$ of the set of indices $M$, there exists $\alpha\in I$ such that the $d_{K/k_\alpha}(x_\mu)$ for $\mu\in J$ are linearly independent in $\Omega^1_{K/k_\alpha}$; but this means also \sref{0.21.4.5} that the $x_\mu$ for $\mu\in J$ form a $p$-free family \emph{over $k_\alpha$}, and the existence of an $\alpha$ having this property follows from \sref{0.21.8.2} and from the hypothesis $a)$.
\end{proof}

\oldpage[0\textsubscript{IV-1}]{174}

\begin{corollary}[21.8.4]
\label{0.21.8.4}
Let $K$ be a field of characteristic $p>0$, $k_0$ a subfield of $K$, $(k_\alpha)_{\alpha\in I}$ a decreasing filtering family of subfields of $K$, containing $k_0$, such that
\[
  \bigcap_\alpha k_\alpha(K^p)=k_0(K^p).
\]

Let $L$ be an extension of $K$ such that $\Upsilon_{L/K/k_0}$ is an $L$-vector space of finite rank; then, for every field $K'$ such that $K\subset K'\subset L$, there exists $\alpha\in I$ such that $k_\alpha$ is $k_0$-admissible for the extension $L$ of $K'$.
\end{corollary}

\begin{proof}
It suffices to apply \sref{0.21.8.3} and \sref{0.21.6.6}.
\end{proof}

\begin{corollary}[21.8.5]
\label{0.21.8.5}
Let $K$ be a field of characteristic $p>0$, $k_0$ a subfield of $K$, $(k_\alpha)_{\alpha\in I}$ a decreasing filtering family of subfields of $K$ containing $k_0$, such that
\[
  \bigcap_\alpha k_\alpha(K^p)=k_0(K^p).
\]

Then, for every extension $L$ of $K$ such that $\Upsilon_{L/K/k_0}$ is an $L$-vector space of finite rank, one has $\bigcap_\alpha k_\alpha(L^p)=k_0(L^p)$.
\end{corollary}

\begin{proof}
Suppose indeed that $M$ is an extension of $L$ such that $\Upsilon_{M/L/k_0}$ is an $M$-vector space of finite rank.
The exact sequence \sref{0.21.6.1.1}
\[
  0\to\Upsilon_{L/K/k_0}\otimes_L M\to\Upsilon_{M/K/k_0}\to\Upsilon_{M/L/k_0}
\]
and the hypothesis show that $\Upsilon_{M/K/k_0}$ is also an $M$-vector space of finite rank.
There exists therefore an index $\alpha$ such that $k_\alpha$ is $k_0$-admissible for the extension $M$ of $K$ \sref{0.21.8.3}; this holding also for the extension $M$ of $L$ \sref{0.21.6.6}; the equivalence of $a)$ and $b)$ in \sref{0.21.8.3} gives the corollary.
\end{proof}

\begin{corollary}[21.8.6]
\label{0.21.8.6}
Let $K$ be a field, $p$ its characteristic exponent, and $k_0$ a subfield of $K$.
If $L$ is an extension of $K$ such that $\Upsilon_{L/K/k_0}$ is an $L$-vector space of finite rank, then there exists a subfield $k$ of $K$, containing $k_0(K^p)$, such that $[K:k]$ is finite, and which is $k_0$-admissible for the extension $L$ of $K$.
\end{corollary}

\begin{proof}
It suffices indeed, by virtue of \sref{0.21.8.3}, to construct a decreasing filtering family $(k_\alpha)_{\alpha\in I}$ of subfields of $K$, containing $K^p$ and $k_0$, for which $[K:k_\alpha]<+\infty$ and $\bigcap_\alpha k_\alpha=k_0(K^p)$.
For this one considers a $p$-base $(x_\lambda)_{\lambda\in J}$ of $K$ over $k_0$ and, for every finite subset $H$ of $J$, one considers the subfield $k_H$ of $K$ generated by $k_0(K^p)$ and the $x_\lambda$ for $\lambda\in J-H$; it follows from this definition that $(x_\lambda)_{\lambda\in H}$ is a $p$-base of $K$ over $k_H$, and one concludes immediately that the $k_H$ satisfy the required conditions.
\end{proof}

\begin{remarks}[21.8.7]
\label{0.21.8.7}
\begin{enumerate}
  \item[(i)] One has already seen \sref{0.21.7.2} that if $L$ is a \emph{finite type} extension of $K$, $\Upsilon_{L/K/k_0}$ is of finite rank for every subfield $k_0$ of $K$.
    It is the same if $L$ is a \emph{separable} extension of $K$, for by virtue of \sref{0.20.6.19}, one has $\Upsilon_{L/K/k_0}=0$.
    Finally, if $L$ is a \emph{finite type extension of a separable extension $L_0$ of $K$}, the same reasoning as in \sref{0.21.8.5} shows that $\Upsilon_{L/K/k_0}$ is still of finite rank (and in fact is isomorphic to a subspace of $\Upsilon_{L/L_0/k_0}$).
  \oldpage[0\textsubscript{IV-1}]{175}
  \item[(ii)] In the statement of \sref{0.21.8.5}, and consequently in that of \sref{0.21.8.3}, $b)$), one cannot omit the hypothesis that $\Upsilon_{L/K/k_0}$ is of finite rank over $L$.
    Let us take for $k_0$ the prime field $\mathbf{F}_p$, for $K$ a field such that $[K:K^p]$ is countably infinite (for example the field of rational functions $\mathbf{F}_p(X_1,\ldots,X_n,\ldots)$ with a countable infinity of indeterminates);
    the construction procedure of \sref{0.21.8.6} then shows that there exists an infinite strictly decreasing sequence $(K_n)$ of subfields of $K$ such that $K_0=K$ and $\bigcap_n K_n=K^p$.
    We shall construct increasing finite extensions $M_n$ of $K^p$, such that if $M$ is the union of the $M_n$, the extension $L=M^{1/p}$ of $K$ gives \sref{0.21.8.5} a counterexample.
    Namely, take an element $x$ of $K-K^p$, set $M_0=K^p$, and $M_n=K^p(a_1 x,a_2 x,\ldots,a_n x)$ for $n\geqslant 1$, where the $a_n$ are constructed by induction so that $a_n\in K_n$ and $a_{n+1}\notin M_n(x)$ for all $n$: this is possible, since $M_n(x)$ is of finite degree over $K^p$, while $K_n$ is not of finite degree over $K^p$.
    One concludes also that $x\notin M=L^p$, but as $x=(a_n x)a_n^{-1}$, one has $x\in K_n(M)=K_n(L^p)$ for all $n$, hence $x\in\bigcap_n K_n(L^p)$, which proves our assertion.
\end{enumerate}
\end{remarks}

\begin{proposition}[21.8.8]
\label{0.21.8.8}
Let $k$ be a field of characteristic $p>0$, and let $A=k[[T_1,\ldots,T_r]]$ be the ring of formal power series in $r$ indeterminates over $k$, $K=k((T_1,\ldots,T_r))$ its field of fractions.
Then there exists a decreasing filtering family $(A_\alpha)$ of Noetherian subrings of $A$ such that $A$ is an $A_\alpha$-module \emph{free of finite type} for all $\alpha$ and that, if $K_\alpha$ is the field of fractions of $A_\alpha$, one has $\bigcap_\alpha K_\alpha=K^p$.
\end{proposition}

\begin{proof}
One can write $K^p=k^p((T_1^p,\ldots,T_r^p))$; as one has seen in the proof of \sref{0.21.8.6}, there exists a decreasing family $(k_\alpha)$ of subfields of $k$ such that $[k:k_\alpha]$ is finite for all $\alpha$ and $\bigcap_\alpha k_\alpha=k^p$; it is clear that if one sets $A_\alpha=k_\alpha[[T_1^p,\ldots,T_r^p]]$, $A$ is an $A_\alpha$-module \emph{free of finite type}; the whole thing therefore reduces to proving the relation
\begin{equation}
\tag{21.8.8.1}
  \label{0.21.8.8.1}
  \bigcap_\alpha k_\alpha((T_1^p,\ldots,T_r^p))=k^p((T_1^p,\ldots,T_r^p)).
\end{equation}

Since $\bigcap_\alpha k_\alpha=k^p$, this will follow from the following two lemmas:
\end{proof}

\begin{lemma}[21.8.8.2]
\label{0.21.8.8.2}
If $k'$ is an extension of a field $k$, one has
\[
  k'[[T_1,\ldots,T_r]]\cap k((T_1,\ldots,T_r))=k[[T_1,\ldots,T_r]].
\]
\end{lemma}

\begin{proof}
Indeed, set $C=k[[T_1,\ldots,T_r]]$, $D=k'[[T_1,\ldots,T_r]]$; as $k((T_1,\ldots,T_r))$ is the field of fractions of $C$, it will suffice to prove that $D$ is a $C$-module \emph{faithfully flat} (Bourbaki, \emph{Alg.\ comm.}, chap.~I, \S~3, no.~5, prop.~10).
Now, $C$ and $D$ are Noetherian local rings, and if $\mathfrak{m}$ is the maximal ideal of $C$, one has $D/\mathfrak{m}^j D=(C/\mathfrak{m}^j)\otimes_k k'$, hence $D/\mathfrak{m}^j D$ is a flat $(C/\mathfrak{m}^j)$-module; it suffices to apply ($\mathbf{0_{III}}$, 10.2.1) and ($\mathbf{0_I}$, 6.6.2).
\end{proof}

\begin{lemma}[21.8.8.3]
\label{0.21.8.8.3}
Let $k$ be a field, $(k_\alpha)$ a decreasing filtering family of subfields of $k$, and set $k_0=\bigcap_\alpha k_\alpha$.
Suppose there exists a power $q$ of the characteristic exponent of $k$ such that $k^q\subset k_0$.
Then one has
\begin{equation}
\tag{21.8.8.4}
  \label{0.21.8.8.4}
  \bigcap_\alpha k_\alpha((T_1,\ldots,T_r))=k_0((T_1,\ldots,T_r)).
\end{equation}
\end{lemma}

\begin{proof}
\oldpage[0\textsubscript{IV-1}]{176}
It suffices to prove that a nonzero element $f$ of the first member of \sref{0.21.8.8.4} belongs to the second member.
Let $\gamma$ be an index, $g\in k_\gamma[[T_1,\ldots,T_r]]$ such that $gf\in k_\gamma[[T_1,\ldots,T_r]]$; replacing $g$ by $g^q$, one can suppose that $g\in k_0[[T_1,\ldots,T_r]]$.
Then, for all $\alpha\geqslant\gamma$, one has
\[
  gf\in k_\gamma[[T_1,\ldots,T_r]]\cap k_\alpha((T_1,\ldots,T_r))=k_\alpha[[T_1,\ldots,T_r]]
\]
by virtue of the lemma \sref{0.21.8.8.2}.
But it is clear that the intersection of the rings $k_\alpha[[T_1,\ldots,T_r]]$ is $k_0[[T_1,\ldots,T_r]]$, and one has thus $f\in k_0((T_1,\ldots,T_r))$.
\end{proof}

\subsection{Completed modules of differentials in formal power series rings}
\label{subsection:0.21.9}

In this number, the fields are no longer necessarily supposed to be of characteristic $>0$.

\begin{lemma}[21.9.1]
\label{0.21.9.1}
Let $k$ be a field, $A$ a complete Noetherian local ring that is a $k$-algebra, $K$ the residue field of $A$.
If $K$ is a finite type extension of $k$, the $A$-module $\widehat{\Omega}^1_{A/k}$ is of finite type.
\end{lemma}

\begin{proof}
By virtue of \sref{0.20.7.15}, it suffices to prove that $\Omega^1_{K/k}$ is a $K$-vector space of finite rank.
Now, by hypothesis, $K$ is the field of fractions of a $k$-algebra of finite type $B$; as $\Omega^1_{B/k}$ is a $B$-module of finite type by virtue of \sref{0.20.4.7}, the conclusion follows from \sref{0.20.5.9}.
\end{proof}

\begin{proposition}[21.9.2]
\label{0.21.9.2}
Let $k_0$ be a field, $A$ a complete Noetherian local ring that is a formally smooth $k_0$-algebra, $\mathfrak{m}$ its maximal ideal, $K=A/\mathfrak{m}$ its residue field \textup{(}so that in particular $A$ is isomorphic to a formal power series ring $K[[T_1,\ldots,T_r]]$ \sref{0.19.6.5}\textup{)}.
If $K$ is a finite type extension of $k_0$, then $\widehat{\Omega}^1_{A/k_0}$ is an $A$-module \emph{free} of rank equal to $\dim(A)+\deg.\operatorname{tr}_{k_0} K$.
Furthermore, for every subfield $k$ of $k_0$ such that $\Omega^1_{k_0/k}$ is of finite rank, $\widehat{\Omega}^1_{A/k}$ is an $A$-free module of rank equal to $\dim(A)+\deg.\operatorname{tr}_{k_0} K+\operatorname{rg}_{k_0}(\Omega^1_{k_0/k})$.
\end{proposition}

\begin{proof}
It is clear that, if $P$ is the prime subfield of $k_0$, $K$ is a formally smooth $P$-algebra \sref{0.19.6.1}.
Let us note, on the other hand, that $\Upsilon^K_{A/k_0/P}=0$: indeed, this amounts to seeing that the homomorphism
\[
  u_{A/k_0/P}\otimes 1_K:\Omega^1_{k_0/P}\otimes_{k_0}K
  \longrightarrow\Omega^1_{A/P}\otimes_AK
\]
is injective \sref{0.20.6.14.5}, where $u$ denotes the structural homomorphism $k_0\to A$.
But since $A$ is a formally smooth $k_0$-algebra by hypothesis, it follows from \sref{0.20.7.2} that $u_{A/k_0/P}$ is formally left invertible; hence \sref{0.19.1.7}, $u_{A/k_0/P}\otimes1_K$ is left invertible and \emph{a fortiori} injective, proving our assertion.
Applying the exact sequence \sref{0.20.6.22.1}, with $\Lambda,A,B,C$ replaced by $P,k_0,A,K$, gives the exact sequence
\[
  0\longrightarrow\Upsilon^K_{A/k_0}
  \longrightarrow\mathfrak{m}/\mathfrak{m}^2
  \longrightarrow\Omega^1_{A/k_0}\otimes_AK
  \longrightarrow\Omega^1_{K/k_0}\longrightarrow0,
\]
so that
\[
\begin{split}
  \operatorname{rg}_K(\Omega^1_{A/k_0}\otimes_AK)
  &={\operatorname{rg}}_K(\mathfrak{m}/\mathfrak{m}^2)
    +\bigl(\operatorname{rg}_K(\Omega^1_{K/k_0})
    -\operatorname{rg}_K(\Upsilon_{K/k_0})\bigr)\\
  &={\operatorname{rg}}_K(\mathfrak{m}/\mathfrak{m}^2)
    +\deg.\operatorname{tr}_{k_0}K
\end{split}
\]
by virtue of \sref{0.21.7.1}.
On the other hand, $\widehat{\Omega}^1_{A/k_0}$ is a formally projective $A$-module \sref{0.20.4.9}; since its topology is the $\mathfrak{m}$-preadic topology \sref{0.20.4.5},
$\widehat{\Omega}^1_{A/k_0}/\mathfrak{m}^{j+1}\widehat{\Omega}^1_{A/k_0}$ is a projective $(A/\mathfrak{m}^{j+1})$-module for every $j$ \sref{0.19.2.4}, hence free ($\mathbf{0_{III}}$, 10.1.2) and of rank
\oldpage[0\textsubscript{IV-1}]{177}
$n=\operatorname{rg}_K(\Omega^1_{A/k_0}\otimes_AK)$.
By virtue of \sref{0.21.9.1}, $\widehat{\Omega}^1_{A/k_0}$ is an $A$-module of finite type; moreover this $A$-module is flat by virtue of ($\mathbf{0_{III}}$, 10.2.1), and is therefore free ($\mathbf{0_{III}}$, 10.1.3) of rank $n$, proving the first assertion.

Since, by hypothesis, $\Omega^1_{k_0/k}$ has finite rank over $k_0$, the completed tensor product $\Omega^1_{k_0/k}\widehat{\otimes}_{k_0}A$ identifies with $\Omega^1_{k_0/k}\otimes_{k_0}A$.
As $A$ is a formally smooth $k_0$-algebra, it follows from \sref{0.20.7.18} that the homomorphism
\[
  \widehat{u}_{A/k_0/k}:\Omega^1_{k_0/k}\otimes_{k_0}A
  \longrightarrow\widehat{\Omega}^1_{A/k}
\]
has a left inverse which is a continuous $A$-homomorphism; in particular, its image is closed in $\widehat{\Omega}^1_{A/k}$.
Consequently the exact sequence
\[
  0\longrightarrow\Omega^1_{k_0/k}\otimes_{k_0}A
  \xrightarrow{\widehat{u}_{A/k_0/k}}\widehat{\Omega}^1_{A/k}
  \longrightarrow\widehat{\Omega}^1_{A/k_0}\longrightarrow0
\]
is exact and split by virtue of \sref{0.20.7.17} and what precedes; this completes the proof.
\end{proof}

\begin{corollary}[21.9.3]
\label{0.21.9.3}
Let $k$ be a field, $A$ the ring of formal power series $k[[T_1,\ldots,T_r]]$, equipped with its usual structure as a $k$-algebra.
Then $\widehat{\Omega}^1_{A/k}$ is an $A$-free module of rank $r=\dim(A)$.
\end{corollary}

\begin{proof}
It suffices to note that $A$ is a $k$-algebra formally smooth \sref{0.19.3.4}, and to apply \sref{0.21.9.2} with $k_0=k=K$.
\end{proof}

\begin{lemma}[21.9.4]
\label{0.21.9.4}
Let $k$ be a field of characteristic $p>0$, $A$ a $k$-algebra that is a local Noetherian complete ring, $B$ a sub-$k$-algebra of $A$ isomorphic to a topological algebra of the form $k[[T_1,\ldots,T_r]]$ and such that $A$ is a $B$-algebra of finite type.
If $B_1$ is the subalgebra $k[[T_1^p,\ldots,T_r^p]]$ of $B$, $\widehat{\Omega}^1_{A/k}$ identifies canonically with $\Omega^1_{A/B_1}$.
\end{lemma}

\begin{proof}
Every continuous derivation of $B_1$ into a topological $A$-module $P$, which is the restriction of a continuous $k$-derivation of $A$ into $P$, is \emph{zero}, for it is zero in the subring $k[T_1^p,\ldots,T_r^p]$ of polynomials, and this last is dense in $B_1$.
The exact sequence \sref{0.20.3.7.1} then shows that the canonical homomorphism
\[
  \operatorname{D\acute{e}r.\ cont}_{B_1}(A,P)\to\operatorname{D\acute{e}r.\ cont}_k(A,P)
\]
is bijective; the conclusion follows from \sref{0.20.7.14.4}, taking into account that $\Omega^1_{A/B_1}$ is an $A$-module of finite type \sref{0.20.4.7}, hence separated and complete since $A$ is a local Noetherian complete ring.
\end{proof}

\begin{proposition}[21.9.5]
\label{0.21.9.5}
Let $A$ be a local Noetherian complete integral ring, $A_0$ a subring of $A$ isomorphic to a formal power series ring $k[[T_1,\ldots,T_r]]$ over a field $k$, such that $A$ is a finite $A_0$-algebra and that the field of fractions $E$ of $A$ is a separable extension of the field of fractions $L_0$ of $A_0$.
Then one has
\begin{equation}
\tag{21.9.5.1}
  \label{0.21.9.5.1}
  \operatorname{rg}_A\widehat{\Omega}^1_{A/k}=\operatorname{rg}_E(\widehat{\Omega}^1_{A/k}\otimes_A E)=\dim(A).
\end{equation}
\end{proposition}

\begin{proof}
Let us note that if $\mathfrak{m}_0$ is the maximal ideal of $A_0$, the topology of $A$ is the $\mathfrak{m}_0$-adic topology since $A$ is a finite $A_0$-algebra (Bourbaki, \emph{Alg.\ comm.}, chap.~IV, \S~2, no.~5, cor.~3 of prop.~9) and it induces on $A_0$ the $\mathfrak{m}_0$-adic topology (Bourbaki, \emph{Alg.\ comm.}, chap.~III, \S~3, no.~4, th.~3).
One knows \sref{0.21.9.1} that $\widehat{\Omega}^1_{A/k}$ is an $A$-module of finite type.
\oldpage[0\textsubscript{IV-1}]{178}
By virtue of \sref{0.21.9.3}, $\widehat{\Omega}^1_{A_0/k}$ is a free $A_0$-module of rank $r=\dim(A_0)=\dim(A)$ \sref{0.16.1.5}; the tensor product $\widehat{\Omega}^1_{A_0/k}\otimes_{A_0}A$ is therefore a free $A$-module of rank $r$, identical with its separated completion.
Finally, $\widehat{\Omega}^1_{A/A_0}$ is an $A$-module of finite type \sref{0.20.4.7}, hence identical with its separated completion \sref{0.7.3.6}.
This being so, in the sequence of homomorphisms
\begin{equation}
\tag{21.9.5.2}
  \label{0.21.9.5.2}
  \widehat{\Omega}^1_{A_0/k}\otimes_{A_0} A\xrightarrow{u}\widehat{\Omega}^1_{A/k}\xrightarrow{v}\widehat{\Omega}^1_{A/A_0}\to 0,
\end{equation}
one knows that $v$ is surjective and that $\operatorname{Im}(u)$ is dense in $\ker(v)$ \sref{0.20.7.17}; but every sub-$A_0$-module of $\widehat{\Omega}^1_{A/k}$ is closed for the $\mathfrak{m}_0$-adic topology \sref{0.7.3.5}, so the sequence \sref{0.21.9.5.2} is exact.
On the other hand, since $A$ is integral over $A_0$, $A_0$ is integrally closed, and $E$ is separable over $L_0$, $\widehat{\Omega}^1_{A/A_0}$ is a torsion $A$-module \sref{0.20.4.13}, (iv).
Tensoring the exact sequence \sref{0.21.9.5.2} by $E$, one obtains an exact sequence
\[
  \widehat{\Omega}^1_{A_0/k}\otimes_{A_0}E\longrightarrow
  \widehat{\Omega}^1_{A/k}\otimes_A E\longrightarrow 0,
\]
whence $\operatorname{rg}_E(\widehat{\Omega}^1_{A/k}\otimes_AE)\leqslant r$.

It remains to show that there exist in $\widehat{\Omega}^1_{A/k}\otimes_AE$ $r$ elements linearly independent over $E$.
Let $D_i=\partial/\partial T_i$ be the canonical derivations of $A_0$ into itself $(1\leqslant i\leqslant r)$.
They extend uniquely to derivations (still denoted $D_i$) of $E$ into itself, since $E$ is a finite separable extension of $L_0$.
By restriction to $A$, these derivations give derivations of $A$ into $E$, and it is immediate that they take their values in one and the same finite type sub-$A$-module of $E$.
Since they are continuous on $A_0$ and the topology of $A$ is the $\mathfrak{m}_0$-adic topology, one has thus obtained $r$ continuous derivations of $A$ into a topological $A$-module; they are evidently linearly independent, since $D_i(T_j)=\delta_{ij}$.
This completes the proof.
\end{proof}

\begin{proposition}[21.9.6]
\label{0.21.9.6}
Let $k_0\subset k\subset K_0$ be three fields of characteristic $p>0$, such that $[K_0:k]<+\infty$; set
\[
  L_0=K_0((T_1,\ldots,T_r)),\quad M_0=K_0((T_1^p,\ldots,T_r^p))
\]
\[
  L=k((T_1,\ldots,T_r)),\quad M=k((T_1^p,\ldots,T_r^p)),\quad N=k((T_1^{p^2},\ldots,T_r^{p^2})).
\]
Let $E$ be a finite type extension of $L_0$.
\begin{enumerate}
  \item[\rm (i)] If $\Upsilon_{K_0/k_0}$ is of finite rank, one has
  \begin{equation}
  \tag{21.9.6.1}
  \label{0.21.9.6.1}
    \operatorname{rg}_E\Omega^1_{E/M}-\operatorname{rg}_{K_0}\Omega^1_{K_0/k}=(r+\deg.\operatorname{tr}_{L_0} E)+d(E/K_0,k_0)+d(E/M_0,N/k_0)
  \end{equation}
  \textup{(}notations of \sref{0.21.6.1}\textup{)}, and in particular one has
  \begin{equation}
  \tag{21.9.6.2}
  \label{0.21.9.6.2}
    \operatorname{rg}_E\Omega^1_{E/M}-\operatorname{rg}_{K_0}\Omega^1_{K_0/k}\geqslant r+\deg.\operatorname{tr}_{L_0} E+d(E/K_0,k_0).
  \end{equation}
  Furthermore, if the two members of \sref{0.21.9.6.2} are equal, they remain so when one replaces $k$ by a subfield $k'$ such that $k_0\subset k'\subset k$.
  \item[\rm (ii)] Suppose moreover that $[k_0:k_0^p]<+\infty$.
  Let $(k_\alpha)$ be a decreasing filtering family of subfields of $K_0$ containing $k_0$, such that $[K_0:k_\alpha]<+\infty$ for all $\alpha$ and that $\bigcap_\alpha k_\alpha(K_0^p)=k_0(K_0^p)$; then there exists $\alpha$ such that, for $k=k_\alpha$, the two members of \sref{0.21.9.6.2} are equal.
\end{enumerate}
\end{proposition}

\begin{proof}
\begin{enumerate}
  \oldpage[0\textsubscript{IV-1}]{179}
  \item[(i)] One knows \sref{0.21.1.5}, (ii) that $\operatorname{rg}_E\Omega^1_{E/M}$ does not change when one replaces $M$ by $M(L_0^p)$; as $L_0^p=K_0^p((T_1^p,\ldots,T_r^p))$ and that $[k(K_0^p):k]<+\infty$, one has $M(L_0^p)=k(K_0^p)((T_1^p,\ldots,T_r^p))$;
  similarly $\Omega^1_{K_0/k}$ does not change when one replaces $k$ by $k(K_0^p)$, so that one can \emph{suppose that $K_0^p\subset k$}, which we shall do in what follows.
  We introduce also the fields
  \[
    k'_0=k_0(K_0^p),\quad N'=k'_0((T_1^{p^2},\ldots,T_r^{p^2}))
  \]
  so that one has the diagram of fields
  \[
    \xymatrix{
      K_0\ar[r] & M_0\ar[r] & L_0\ar[r] & E\\
      k\ar[u]\ar[r] & N\ar[u]\ar[r] & M\ar[u]\ar[r] & L\ar[u] \\
      k_0\ar[u]\ar[r] & k'_0\ar[u]\ar[r] & N'\ar[u] &
    }
  \]

  We propose to evaluate the difference $\delta$ of the first and second members of \sref{0.21.9.6.2}.

  It follows from \sref{0.21.7.3} that one has
  \[
    \operatorname{rg}_E\Omega^1_{E/M}-\operatorname{rg}_{M_0}\Omega^1_{M_0/M}=\deg.\operatorname{tr}_{M_0}E+d(E/M_0,M)
  \]
  and as $L_0$ is a finite extension of $M_0$, one has $\deg.\operatorname{tr}_{L_0} E=\deg.\operatorname{tr}_{M_0} E$.
  On the other hand, a $p$-base of $K_0$ over $k=k(K_0^p)$ is also a $p$-base of $M_0$ over $M$, hence \sref{0.21.4.5}, one has
  \[
    \operatorname{rg}_{M_0}\Omega^1_{M_0/M}=\operatorname{rg}_{K_0}\Omega^1_{K_0/k},
  \]
  so that
  \begin{equation}
  \tag{21.9.6.3}
  \label{0.21.9.6.3}
    \delta=d(E/M_0,M)-d(E/K_0,k_0)-r.
  \end{equation}

  Let us now note the following classical lemma.

  \begin{lemma}[21.9.6.4]
  \label{0.21.9.6.4}
  For every field $k$, the formal power series field $K=k((T_1,\ldots,T_r))$ is a separable extension of $k$.
  \end{lemma}

  \begin{proof}
  Let us briefly recall the proof for completeness.
  It suffices (Bourbaki, \emph{Alg.}, chap.~VIII, \S~7, no.~3, proof of th.~1) to prove that for every finite extension $k'$ of $k$, $K\otimes_k k'$ has no nilpotent element.
  Set $A=k[[T_1,\ldots,T_r]]$ and $S=A-\{0\}$.
  Then $K\otimes_k k'=S^{-1}(A\otimes_k k')$, while $A\otimes_k k'$ identifies canonically with the integral ring $A'=k'[[T_1,\ldots,T_r]]$ and $A$ with a subring of $A'$; the conclusion follows.
  \end{proof}

  Using this lemma and \sref{0.21.6.2}, (ii), one has $d(L_0/K_0,k_0)=0$, and formula \sref{0.21.6.4.4} therefore gives
  \[
    d(E/K_0,k_0)=d(E/L_0,k_0).
  \]

  \oldpage[0\textsubscript{IV-1}]{180}
  Formulas \sref{0.21.6.4.4} and \sref{0.21.6.4.2} give, on the other hand,
  \[
    d(E/M_0,M)=d(E/L_0,M)+d(L_0/M_0,M),
  \]
  \[
    d(L_0/M_0,M)=d(L_0/M_0,M/N)+d(L_0/M_0,N),
  \]
  whence
  \[
    \delta=d(E/L_0,M)+d(L_0/M_0,N)-d(E/L_0,k_0)+d(L_0/M_0,M/N)-r.
  \]

  We shall show that
  \begin{equation}
  \tag{21.9.6.5}
  \label{0.21.9.6.5}
    d(L_0/M_0,M/N)=r.
  \end{equation}
  Indeed, by definition,
  \[
    d(L_0/M_0,M/N)=\operatorname{rg}_{L_0}\Upsilon_{L_0/M/N}-\operatorname{rg}_{M_0}\Upsilon_{M_0/M/N}.
  \]
  Taking account of the two exact sequences
  \[
  \begin{aligned}
    0&\longrightarrow\Upsilon_{L_0/M/N}\longrightarrow\Omega^1_{M/N}\otimes_ML_0
      \longrightarrow\Omega^1_{L_0/N}\longrightarrow\Omega^1_{L_0/M}\longrightarrow0,\\
    0&\longrightarrow\Upsilon_{M_0/M/N}\longrightarrow\Omega^1_{M/N}\otimes_MM_0
      \longrightarrow\Omega^1_{M_0/N}\longrightarrow\Omega^1_{M_0/M}\longrightarrow0,
  \end{aligned}
  \]
  one obtains
  \[
  \begin{split}
    d(L_0/M_0,M/N)={}&\operatorname{rg}_{L_0}\Omega^1_{L_0/M}
      -\operatorname{rg}_{L_0}\Omega^1_{L_0/N}\\
      &-\operatorname{rg}_{M_0}\Omega^1_{M_0/M}
      +\operatorname{rg}_{M_0}\Omega^1_{M_0/N}.
  \end{split}
  \]
  Now $N(L_0^p)\supset M$, hence $\Omega^1_{L_0/N}=\Omega^1_{L_0/M}$ \sref{0.21.1.5}, (ii).
  On the other hand, the $T_i^p$ $(1\leqslant i\leqslant r)$ form a $p$-basis of $M$ over $N$, and a $p$-basis of $K_0$ over $k$ is also a $p$-basis of $M_0$ over $M$.
  Since $M_0^p\subset N$, it follows from \sref{0.21.4.5} and \sref{0.21.1.10} that
  \[
    \operatorname{rg}_{M_0}\Omega^1_{M_0/N}
      =r+\operatorname{rg}_{M_0}\Omega^1_{M_0/M}
      =r+\operatorname{rg}_{K_0}\Omega^1_{K_0/k},
  \]
  proving \sref{0.21.9.6.5}.

  One has again, by \sref{0.21.6.4.2},
  \[
    d(E/L_0,M)=d(E/L_0,k_0)+d(E/L_0,M/k_0)
  \]
  and
  \[
    d(L_0/M_0,N)=d(L_0/M_0,N/k)+d(L_0/M_0,k).
  \]
  Since $L_0$ is separable over $K_0$ \sref{0.21.9.6.4}, one has $d(L_0/K_0,k)=0$, and consequently also $d(L_0/M_0,k)=0$ by \sref{0.21.6.2} and \sref{0.21.6.4}.
  Thus
  \[
    \delta=d(E/L_0,M/k_0)+d(L_0/M_0,N/k).
  \]

  Applying \sref{0.21.6.4.2} twice more to the first term on the right gives
  \[
    d(E/L_0,M/k_0)=d(E/L_0,M/N)+d(E/L_0,N/k)+d(E/L_0,k/k_0).
  \]
  Since $M\subset N(L_0^p)$, one has $d(E/L_0,M/N)=0$ \sref{0.21.6.2}; hence, by \sref{0.21.6.4.4},
  \[
    \delta=d(E/L_0,k/k_0)+d(E/M_0,N/k).
  \]
  But $M_0$ and $L_0$ are separable extensions of $K_0$ \sref{0.21.9.6.4}, so
  \[
    d(L_0/K_0,k/k_0)=d(M_0/K_0,k/k_0)=0
  \]
  by \sref{0.21.6.2}.
  Applying \sref{0.21.6.4.4} twice more,
  \oldpage[0\textsubscript{IV-1}]{181}
  one has
  \[
    d(E/L_0,k/k_0)=d(E/K_0,k/k_0)=d(E/M_0,k/k_0),
  \]
  and a final application of \sref{0.21.6.4.2} gives
  \[
    \delta=d(E/M_0,N/k_0),
  \]
  which proves \sref{0.21.9.6.1}.

  Furthermore, when $k$ is replaced by an intermediate field $k'$ with $k_0\subset k'\subset k$, $N$ is replaced by a subextension $N'$, and therefore $d(E/M_0,N'/k_0)\leqslant d(E/M_0,N/k_0)$ \sref{0.21.6.4.2}.
  This proves the last assertion of (i).

  \item[(ii)] For all $\alpha$, set $N_\alpha=k_\alpha((T_1^{p^2},\ldots,T_r^{p^2}))$.
  By virtue of \sref{0.21.9.6.1} and of \sref{0.21.8.3}, it suffices to show (taking into account the fact that $E$ is a finite type extension of $M_0$) that one has
  \[
    \bigcap_\alpha N_\alpha=k_0(M_0^p).
  \]
  Or, one has
  \begin{equation}
  \tag{21.9.6.6}
  \label{0.21.9.6.6}
    \bigcap_\alpha N_\alpha=N'=(k_0(K_0^p))((T_1^{p^2},\ldots,T_r^{p^2}))
  \end{equation}
  by virtue of \sref{0.21.8.8.3} and of the relation $\bigcap_\alpha k_\alpha=k_0(K_0^p)$ \sref{0.21.8.6}.
  On the other hand, one has $k_0(M_0^p)=k_0(K_0^p((T_1^{p^2},\ldots,T_r^{p^2})))$.
  To finish the proof of \sref{0.21.9.6}, (ii), it therefore remains to prove the following lemma.

  \begin{lemma}[21.9.6.7]
  \label{0.21.9.6.7}
  Let $k_0$ be a field of characteristic $p>0$ such that $[k_0:k_0^p]<+\infty$, and let $K_0$ be an extension of $k_0$.
  Then
  \begin{equation}
  \tag{21.9.6.8}
    \label{0.21.9.6.8}
    (k_0(K_0^p))((T_1,\ldots,T_r))=k_0\bigl(K_0^p((T_1,\ldots,T_r))\bigr).
  \end{equation}
  \end{lemma}

  \begin{proof}
  Since $k_0^p\subset K_0^p$, one has
  \[
    k_0^p\bigl(K_0^p((T_1,\ldots,T_r))\bigr)=K_0^p((T_1,\ldots,T_r)).
  \]
  If $(c_i)_{1\leqslant i\leqslant n}$ is a basis of $k_0$ over $k_0^p$, it is immediate that $(c_i)$ is also a system of generators of each side of \sref{0.21.9.6.8} over $K_0^p((T_1,\ldots,T_r))$.
  \end{proof}
\end{enumerate}
\end{proof}

\begin{remark}[21.9.7]
\label{0.21.9.7}
The assertion of \sref{0.21.9.6}, (ii) \emph{does not extend} to the case where $[k_0:k_0^p]$ is infinite.
Take for $k_0$ the field $\mathbf{F}_p(X_n)_{n\geqslant0}$ of rational functions in infinitely many indeterminates, so that the $X_n$ form a $p$-basis of $k_0$ over $\mathbf{F}_p$, and consequently $[k_0:k_0^p]=+\infty$.
In the statement of \sref{0.21.9.6}, take $K_0=k_0$ (hence necessarily $k=k_0$), $L_0=k_0((T))$, and let $E=L_0(y)$, where $y$ is a root of
\[
  Y^p-\sum_{n=0}^{\infty}X_nT^{np}\in L_0[Y].
\]
Then the difference $\delta$ between the two sides of \sref{0.21.9.6.2} is nonzero.
Indeed, $E$ is separable over $k_0$; otherwise $y^p$ would belong to $k_0(L_0^p)$, which is plainly impossible because the $X_n$ are infinitely many algebraically independent elements.
Thus $d(E/K_0,k_0)=0$, formula \sref{0.21.9.6.3} reduces to $\delta=d(E/M_0,M_0)-1$, and it is clear that $d(E/M_0,M_0)=\operatorname{rg}_E\Upsilon_{E/M_0}$.
It remains to check that this rank is $2$.
Now $E^p\subset M_0$; a $p$-basis of $E$ over $L_0$ consists of the single element $y$, and $L_0$ has the single-element $p$-basis $T$ over $M_0$.
Since $E$ is algebraic over $M_0$, the assertion follows from \sref{0.21.7.1} and \sref{0.21.4.5}.
\end{remark}

\begin{proposition}[21.9.8]
\label{0.21.9.8}
Let $k_0$ be a field \textup{(}of arbitrary characteristic\textup{)}, $K_0$ a separable extension of $k_0$, $A$ a local Noetherian complete ring, which is a $k_0$-algebra and whose residue field is a finite extension of $K_0$.
\oldpage[0\textsubscript{IV-1}]{182}
For every prime ideal $\mathfrak{p}$ of $A$, let $k(\mathfrak{p})$ be the residue field of $A$ at $\mathfrak{p}$.
Then, for every field $k$ such that $k_0\subset k\subset K_0$ and such that $[K_0:k]<+\infty$, one has
\begin{equation}
\tag{21.9.8.1}
  \label{0.21.9.8.1}
  \operatorname{rg}_{k(\mathfrak{p})}(\widehat{\Omega}^1_{(A/\mathfrak{p})/k}\otimes_{A/\mathfrak{p}} k(\mathfrak{p}))-\operatorname{rg}_{K_0}\Omega^1_{K_0/k}\geqslant\dim(A/\mathfrak{p})+\operatorname{rg}_{k(\mathfrak{p})}\Upsilon_{k(\mathfrak{p})/k_0}.
\end{equation}

Suppose moreover that $[k_0:k_0^p]<+\infty$ \textup{(}where $p$ is the exponential characteristic of $k_0$\textup{)}.
Then one can find a field $k$ such that $k_0(K_0^p)\subset k\subset K_0$ and $[K_0:k]<+\infty$, for which the two members of \sref{0.21.9.8.1} are equal.
\end{proposition}

\begin{proof}
As $A/\mathfrak{p}$ has the same residue field $K$ as $A$ and is complete, one can restrict to the case where $A$ is integral and $\mathfrak{p}=0$; one will then set $E=k(\mathfrak{p})$, field of fractions of $A$.
As $K_0$ is formally smooth over $k_0$ \sref{0.19.6.1}, there exists a $k_0$-monomorphism $K_0\to A$ making $A$ into a $K_0$-algebra; one knows in particular that there exists a sub-$K_0$-algebra $A_0$ of $A$, $K_0$-isomorphic to a formal power series algebra $K_0[[T_1,\ldots,T_r]]$ and such that $A$ is an $A_0$-algebra \emph{finite} \sref{0.19.8.9}, which implies that $E$ is a finite extension of the field of fractions $L_0=K_0((T_1,\ldots,T_r))$ of $A_0$.
One has $\Upsilon_{K_0/k_0}=0$ since $K_0$ is separable over $k_0$, and $\deg.\operatorname{tr}_{L_0}E=0$; moreover, as $A$ is a finite $A_0$-algebra, $\dim(A)=\dim(A_0)$ \sref{0.16.1.5}.

If $p=1$, one has $\operatorname{rg}_E(\widehat{\Omega}^1_{A/k}\otimes_A E)=r$ by \sref{0.21.9.5} and \sref{0.17.1.4}, (iii), and $\Omega^1_{K_0/k}=0$ since $K_0$ being a separable finite extension of $k$ \sref{0.21.7.4}; on the other hand $\Upsilon_{E/k_0}=0$, $E$ being a separable extension of $k_0$; the two members of \sref{0.21.9.8.1} are therefore equal in this case.

Suppose now $p>1$; if $B=k[[T_1,\ldots,T_r]]$, $A$ is a finite $B$-algebra. Setting $B_1=k[[T_1^p,\ldots,T_r^p]]$, $\widehat{\Omega}^1_{A/k}$ identifies with $\Omega^1_{A/B_1}$ \sref{0.21.9.4}; and if $M$ denotes the field of fractions $k((T_1^p,\ldots,T_r^p))$, it follows from \sref{0.20.5.9} that $\widehat{\Omega}^1_{A/k}\otimes_A E$ identifies with $\Omega^1_{E/M}$. The inequality \sref{0.21.9.8.1} is then none other than \sref{0.21.9.6.2}, and the last assertion follows from \sref{0.21.9.6}, (ii).
\end{proof}
% END NATIVE INPUT ega0/ega0-21
% BEGIN NATIVE INPUT ega0/ega0-22
\section{Differential criteria for smoothness and regularity}
\label{section:0.22}

\oldpage[0\textsubscript{IV-1}]{182}
The main results of this section are:
\begin{enumerate}
  \item[a)] The criteria \sref{0.22.2.2} and \sref{0.22.5.8}, which give necessary and sufficient conditions for a complete Noetherian local ring $A$ containing a field $k$ to be formally smooth over $k$; the criterion \sref{0.22.5.8} is stated without involving differential notions, and completes \sref{0.19.6.6}; the statement \sref{0.22.2.2}, of a slightly more technical nature, permits giving a \emph{classification}, for fixed $k$, of the local rings $A$ in question \sref{0.22.2.5}: they are determined by their dimension $n$ and their residue field $\mathrm{K}$, under the sole condition $n\geq\rg_\mathrm{K}\Upsilon_{\mathrm{K}/k}$.
  \item[b)] The theorem \sref{0.22.3.3}, and the corollary \sref{0.22.7.6} of the Jacobian criterion of Nagata \sref{0.22.7.3}, which give conditions allowing one to assert that a \emph{localized} ring of a complete Noetherian local ring $A$ containing a field $k$ is geometrically regular over $k$; they will play an important role in the study of the ``fine'' properties of local rings carried out in chap.~IV, \textsection~7. We do not possess for the moment any proof of \sref{0.22.7.6} independent of the Jacobian criterion of Nagata.
  \item[c)] The Jacobian criterion of Zariski \sref{0.22.6.7} and its variants, which are easy consequences of \sref{0.20.7.8}.
\end{enumerate}

\subsection{Lifting of formal smoothness}
\label{subsection:0.22.1}

\oldpage[0\textsubscript{IV-1}]{183}
\begin{theorem}[22.1.1]
\label{0.22.1.1}
Let $s:\Lambda\to\mathrm{A}$, $u:\mathrm{A}\to\mathrm{B}$ be two ring homomorphisms, $\mathfrak{J}$ an ideal of $\mathrm{B}$, $\mathrm{C}$ the quotient ring $\mathrm{B}/\mathfrak{J}$.
Suppose that $\mathrm{C}$ is a $\Lambda$-algebra formally smooth for the discrete topology.
\begin{enumerate}
  \item[\rm{(i)}] \emph{Suppose the following conditions are satisfied:}
  \begin{enumerate}
    \item[$\alpha$)] $\mathfrak{J}/\mathfrak{J}^2$ \emph{is a projective $\mathrm{C}$-module.}
    \item[$\beta$)] \emph{The canonical homomorphism $\mathbf{S}^*_\mathrm{C}(\mathfrak{J}/\mathfrak{J}^2)\to\gr^*_\mathfrak{J}(\mathrm{B})$ \sref{0.19.5.2} is bijective.}
    \item[$\gamma$)] \emph{The characteristic homomorphism $\chi:\Upsilon_{\mathrm{C}/\mathrm{B}/\Lambda}\to\mathfrak{J}/\mathfrak{J}^2$ \sref{0.20.6.20} is injective.}
    \item[$\delta$)] \emph{The $\mathrm{C}$-module $\Omega^1_{\mathrm{B}/\mathrm{A}}\otimes_\mathrm{B}\mathrm{C}$ is projective.}
  \end{enumerate}
  \emph{Then $\mathrm{B}$, equipped with the $\mathfrak{J}$-preadic topology, is a formally smooth $\Lambda$-algebra.}
  \item[\rm{(ii)}] \emph{Conversely, suppose that $\mathrm{B}$ is a $\Lambda$-algebra formally smooth for the $\mathfrak{J}$-preadic topology, and that $\mathrm{A}$ is a $\Lambda$-algebra formally smooth for the discrete topology.
  Then the conditions $\alpha$), $\beta$), $\gamma$), $\delta$) of \rm{(i)} are satisfied.}
\end{enumerate}
\end{theorem}

\begin{proof}
By virtue of the exact sequence \sref{0.20.6.22.1} (which is applicable, since $\mathrm{B}/\mathfrak{J}^2$ is a $\Lambda$-trivial extension of $\mathrm{C}$ by $\mathfrak{J}/\mathfrak{J}^2$, $\mathrm{C}$ being assumed to be a formally smooth $\Lambda$-algebra \sref{0.19.4.4}), the condition $\gamma$) is equivalent to $\Upsilon^{\mathrm{C}}_{\mathrm{B}/\mathrm{A}/\Lambda}=0$, or equivalently to the following:
\begin{enumerate}
  \item[$\gamma'$)] \emph{The homomorphism $u_{\mathrm{B}/\mathrm{A}/\Lambda}\otimes 1_\mathrm{C}:\Omega^1_{\mathrm{A}/\Lambda}\otimes_\mathrm{A}\mathrm{C}\to\Omega^1_{\mathrm{B}/\Lambda}\otimes_\mathrm{B}\mathrm{C}$ is injective.}
\end{enumerate}

(i) The conditions $\alpha$) and $\beta$) imply that $\mathrm{B}$ is a $\Lambda$-algebra formally smooth for the $\mathfrak{J}$-preadic topology \sref{0.19.5.4}.
It amounts to the same thing to say that $\mathrm{B}$ is an $\mathrm{A}$-algebra formally smooth, or an $\mathrm{A}$-algebra formally smooth \emph{relative to $\Lambda$}, for the $\mathfrak{J}$-preadic topology.
To see that $\mathrm{B}$ is an $\mathrm{A}$-algebra formally smooth, it suffices therefore, by virtue of \sref{0.20.7.2}, to prove that the continuous homomorphism of topological $\mathrm{B}$-modules $u_{\mathrm{B}/\mathrm{A}/\Lambda}:\Omega^1_{\mathrm{A}/\Lambda}\otimes_\mathrm{A}\mathrm{B}\to\Omega^1_{\mathrm{B}/\Lambda}$ is \emph{formally invertible on the left}.
But since $\mathrm{B}$ is a $\Lambda$-algebra formally smooth, the topological $\mathrm{B}$-module $\Omega^1_{\mathrm{B}/\Lambda}$ is formally projective \sref{0.20.4.9}, and its topology is deduced from that of $\mathrm{B}$ \sref{0.20.4.5}.
By virtue of \sref{0.19.1.9}, it suffices therefore, for $u_{\mathrm{B}/\mathrm{A}/\Lambda}$ to be formally invertible on the left, that $u_{\mathrm{B}/\mathrm{A}/\Lambda}\otimes 1_\mathrm{C}$ be invertible on the left.
But by virtue of $\gamma'$), this last map is injective, and one therefore has the exact sequence \sref{0.20.6.14.7}
\[
  0\to\Omega^1_{\mathrm{A}/\Lambda}\otimes_\mathrm{A}\mathrm{C}\to\Omega^1_{\mathrm{B}/\Lambda}\otimes_\mathrm{B}\mathrm{C}\to\Omega^1_{\mathrm{B}/\mathrm{A}}\otimes_\mathrm{B}\mathrm{C}\to 0.
\]

Finally, by virtue of $\delta$), the $\mathrm{C}$-module $\Omega^1_{\mathrm{B}/\mathrm{A}}\otimes_\mathrm{B}\mathrm{C}$ is projective, so the preceding exact sequence is split, which finishes proving (i).

\oldpage[0\textsubscript{IV-1}]{184}
(ii) If $\mathrm{B}$ is a $\Lambda$-algebra formally smooth for the $\mathfrak{J}$-preadic topology, and $\mathrm{A}$ is a $\Lambda$-algebra formally smooth for the discrete topology, then $\mathrm{B}$ is an $\mathrm{A}$-algebra formally smooth for the $\mathfrak{J}$-preadic topology \sref{0.19.3.5}; the conditions $\alpha$) and $\beta$) therefore follow from \sref{0.19.5.6}. Moreover, \sref{0.20.7.2} shows that $u_{\mathrm{B}/\mathrm{A}/\Lambda}$ is formally invertible on the left, so $u_{\mathrm{B}/\mathrm{A}/\Lambda}\otimes 1_\mathrm{C}$ is invertible on the left \sref{0.19.1.7}, and \emph{a fortiori} injective. Finally, $\Omega^1_{\mathrm{B}/\mathrm{A}}$ is a formally projective $\mathrm{B}$-module \sref{0.20.4.9}, and consequently $\Omega^1_{\mathrm{B}/\mathrm{A}}\otimes_\mathrm{B}\mathrm{C}$ is a projective $\mathrm{C}$-module \sref{0.19.2.4}.
\end{proof}

\begin{corollary}[22.1.2]
\label{0.22.1.2}
Let $s:\Lambda\to\mathrm{A}$, $u:\mathrm{A}\to\mathrm{B}$ be two ring homomorphisms, $\mathfrak{m}$ a maximal ideal of $\mathrm{B}$, $\mathrm{K}=\mathrm{B}/\mathfrak{m}$ the quotient field.
Suppose that the $\Lambda$-algebras $\mathrm{A}$ and $\mathrm{K}$ are formally smooth for the discrete topology.
Then the following conditions are equivalent:
\begin{enumerate}
  \item[a)] $\mathrm{B}$ is an $\mathrm{A}$-algebra formally smooth for the $\mathfrak{m}$-preadic topology.
  \item[b)] \emph{The canonical homomorphism}
\end{enumerate}
\[
\label{0.22.1.2.1}
  \mathbf{S}^*_\mathrm{K}(\mathfrak{m}/\mathfrak{m}^2)\to\gr^*_\mathfrak{m}(\mathrm{B})
  \tag{22.1.2.1}
\]
\emph{is bijective, and the characteristic homomorphism}
\[
\label{0.22.1.2.2}
  \chi:\Upsilon_{\mathrm{K}/\mathrm{B}/\Lambda}\to\mathfrak{m}/\mathfrak{m}^2
  \tag{22.1.2.2}
\]
\emph{is injective.}
\end{corollary}

\begin{proof}
This follows immediately from \sref{0.22.1.1}, the conditions $\alpha$) and $\delta$) being trivially satisfied here, since $\mathrm{K}$ is a field.
\end{proof}

\begin{remark}[22.1.3]
\label{0.22.1.3}
Suppose the hypotheses of \sref{0.22.1.2} are satisfied, and suppose in addition that $\mathfrak{m}/\mathfrak{m}^2$ is a $\mathrm{K}$-vector space of \emph{finite} rank.
Then, for $\mathrm{B}$ to be an $\mathrm{A}$-algebra formally smooth for the $\mathfrak{m}$-preadic topology, it suffices that the following conditions be satisfied:

\begin{env}[22.1.3.1]
\label{0.22.1.3.1}
Given a polynomial algebra $\mathrm{E}=\mathrm{K}[\mathrm{T}_1,\ldots,\mathrm{T}_n]$, the maximal ideal $\mathfrak{n}$ of $\mathrm{E}$ generated by the $\mathrm{T}_i$ $(1\leq i\leq n)$, and an ideal $\mathfrak{b}\supset\mathfrak{n}^k$ in $\mathrm{E}$, every $\Lambda$-homomorphism $v:\mathrm{B}\to\mathrm{E}/\mathfrak{b}$ which, by passing to quotients, gives the identity $\mathrm{K}\to\mathrm{K}$, factors as $\mathrm{B}\xrightarrow{w}\mathrm{E}/\mathfrak{n}^k\to\mathrm{E}/\mathfrak{b}$, where $w$ is a $\Lambda$-homomorphism.
\end{env}

\begin{env}[22.1.3.2]
\label{0.22.1.3.2}
If $\mathrm{F}$ is the trivial extension $\mathrm{D}_\mathrm{K}(\mathrm{K})$ (``algebra of dual numbers over $\mathrm{K}$''), $\rho:\mathrm{A}\to\mathrm{F}$ a ring homomorphism such that the composition $\mathrm{A}\xrightarrow{\rho}\mathrm{F}\to\mathrm{K}$ is the canonical homomorphism, then every $\Lambda$-homomorphism $v:\mathrm{B}\to\mathrm{K}$ which, by passing to quotients, gives the identity $\mathrm{K}\to\mathrm{K}$, factors as $\mathrm{B}\xrightarrow{w}\mathrm{F}\to\mathrm{K}$, where $w$ is a $\Lambda$-homomorphism (for the $\mathrm{A}$-algebra structure on $\mathrm{F}$ defined by $\rho$).
\end{env}

Indeed, it was seen \sref{0.19.5.5} that the condition \sref{0.22.1.3.1} implies that the canonical homomorphism \sref{0.22.1.2.1} is bijective.
By virtue of \sref{0.22.1.2}, it suffices then to see that $u_{\mathrm{B}/\mathrm{A}/\Lambda}\otimes 1_\mathrm{K}:\Omega^1_{\mathrm{A}/\Lambda}\otimes_\mathrm{A}\mathrm{K}\to\Omega^1_{\mathrm{B}/\Lambda}\otimes_\mathrm{B}\mathrm{K}$ is injective, or equivalently (since these are vector spaces over $\mathrm{K}$) that for every $\mathrm{K}$-vector space $\mathrm{L}$, the canonical homomorphism $\Der_\Lambda(\mathrm{B},\mathrm{L})\to\Der_\Lambda(\mathrm{A},\mathrm{L})$ \emph{is surjective} (by virtue of \sref{0.20.4.8}); it is clear that for this, it is necessary and sufficient that the canonical homomorphism $\Der_\Lambda(\mathrm{B},\mathrm{K})\to\Der_\Lambda(\mathrm{A},\mathrm{K})$ be surjective, and our assertion follows, by virtue of \sref{0.20.1.5}.

Note that the two conditions \sref{0.22.1.3.1} and \sref{0.22.1.3.2} are implied by the following:

\begin{env}[22.1.3.3]
\label{0.22.1.3.3}
For every local Artinian $\Lambda$-algebra $\mathrm{E}$ with residue field equal to $\mathrm{K}$, and every nilpotent ideal $\mathfrak{r}$ of $\mathrm{E}$, every $\Lambda$-homomorphism $\mathrm{B}\to\mathrm{E}/\mathfrak{r}$ which, by passing to quotients, gives the identity $\mathrm{K}\to\mathrm{K}$, factors as $\mathrm{B}\xrightarrow{u}\mathrm{E}\to\mathrm{E}/\mathfrak{r}$, where $u$ is a $\Lambda$-homomorphism.
\end{env}
\end{remark}

\oldpage[0\textsubscript{IV-1}]{185}
\begin{proposition}[22.1.4]
\label{0.22.1.4}
Let $\varphi:\mathrm{A}\to\mathrm{B}$ be a local homomorphism of local Noetherian rings, $k$ (resp.~$\mathrm{K}$) the residue field of $\mathrm{A}$ (resp.~$\mathrm{B}$).
For $\mathrm{B}$ to be a formally smooth $\mathrm{A}$-algebra (for the preadic topologies), it is necessary and sufficient that for every local Artinian ring $\mathrm{C}$, with residue field isomorphic to $\mathrm{K}$, every local homomorphism $\mathrm{A}\to\mathrm{C}$ (making $\mathrm{C}$ a topological $\mathrm{A}$-algebra), and every nilpotent ideal $\mathfrak{J}$ of $\mathrm{C}$, every local $\mathrm{A}$-homomorphism $\mathrm{B}\to\mathrm{C}/\mathfrak{J}$ such that the homomorphism $\mathrm{K}\to\mathrm{K}$ obtained by passing to quotients is the identity, factors as $\mathrm{B}\xrightarrow{f}\mathrm{C}\to\mathrm{C}/\mathfrak{J}$, where $f$ is a local $\mathrm{A}$-homomorphism.
\end{proposition}

\begin{proof}
The condition being necessary by definition \sref{0.19.3.1}, let us show that it is sufficient.
Since the local $\widehat{\mathrm{A}}$-homomorphisms from $\widehat{\mathrm{B}}$ into a discrete local $\widehat{\mathrm{A}}$-algebra arise by extension from the local $\mathrm{A}$-homomorphisms from $\mathrm{B}$ into that algebra, we may restrict ourselves to the case where $\mathrm{A}$ and $\mathrm{B}$ are \emph{complete}, by virtue of \sref{0.19.3.6}.
When $\mathrm{A}$ is a \emph{field} $k$, the proposition follows from \sref{0.22.1.3.3}, where one takes $\Lambda$ to be the prime subfield of $k$, and \sref{0.19.6.1}.
In the general case, set $\mathrm{B}_0=\mathrm{B}\otimes_\mathrm{A}k=\mathrm{B}/\mathfrak{m}\mathrm{B}$ ($\mathfrak{m}$ the maximal ideal of $\mathrm{A}$), which is complete.
If $\mathrm{C}_0$ is a local Artinian $k$-algebra with residue field isomorphic to $\mathrm{K}$, $\mathfrak{J}_0$ a nilpotent ideal of $\mathrm{C}_0$, every $k$-local homomorphism $g_0:\mathrm{B}_0\to\mathrm{C}_0/\mathfrak{J}_0$ giving the identity on $\mathrm{K}$ by passing to quotients, gives by composition an $\mathrm{A}$-local homomorphism
\[
  g:\mathrm{B}\to\mathrm{B}_0\xrightarrow{g_0}\mathrm{C}_0/\mathfrak{J}_0
\]
which by hypothesis factors as $\mathrm{B}\xrightarrow{f}\mathrm{C}_0\to\mathrm{C}_0/\mathfrak{J}_0$, where $f$ is a local $\mathrm{A}$-homomorphism; but as $\mathrm{C}_0$ is a $k$-algebra, $f$ factors as $\mathrm{B}\to\mathrm{B}_0\xrightarrow{f_0}\mathrm{C}_0$, where $f_0$ is a local $k$-homomorphism.
We thus see that the hypotheses of the statement are still fulfilled when one substitutes $\mathrm{B}_0$ and $k$ for $\mathrm{B}$ and $\mathrm{A}$ respectively; therefore $\mathrm{B}_0$ is a $k$-algebra \emph{formally smooth} from what we have just seen.
Applying \sref{0.19.7.2} and \sref{0.19.7.1}, we see that there exists a local Noetherian complete ring $\mathrm{B}'$, a local homomorphism $\mathrm{A}\to\mathrm{B}'$ making $\mathrm{B}'$ a flat $\mathrm{A}$-module and an $\mathrm{A}$-algebra \emph{formally smooth}, and a $k$-isomorphism $u_0':\mathrm{B}_0'=\mathrm{B}'\otimes_\mathrm{A}k\xrightarrow{\sim}\mathrm{B}_0$.
Let $v_0'$ be the inverse isomorphism of $u_0'$; we now propose to show that there exists a local $\mathrm{A}$-homomorphism $v:\mathrm{B}\to\mathrm{B}'$ such that $v_0'$ comes from $v$ by passing to quotients.
For this, let us denote by $\mathfrak{n}$ and $\mathfrak{n}'$ the maximal ideals of $\mathrm{B}$ and $\mathrm{B}'$ respectively, and for every $j$, let $w_j:\mathrm{B}/(\mathfrak{n}^{j+1}+\mathfrak{m}\mathrm{B})\to\mathrm{B}'/(\mathfrak{n}'^{j+1}+\mathfrak{m}\mathrm{B}')$ be the $\mathrm{A}$-local homomorphism deduced from $v_0'$ by passing to quotients; it will suffice to form for each $j$ an $\mathrm{A}$-local homomorphism $v_j:\mathrm{B}/\mathfrak{n}^{j+1}\to\mathrm{B}'/\mathfrak{n}'^{j+1}$ such that the diagrams
\[
  \begin{tikzcd}
    \mathrm{B}/\mathfrak{n}^{j+1} \ar[r,"v_j"] \ar[d] & \mathrm{B}'/\mathfrak{n}'^{j+1} \ar[d] & \mathrm{B}/\mathfrak{n}^{j+1} \ar[r,"v_j"] \ar[d] & \mathrm{B}'/\mathfrak{n}'^{j+1} \ar[d]\\
    \mathrm{B}/\mathfrak{n}^j \ar[r,"v_{j-1}"'] & \mathrm{B}'/\mathfrak{n}'^j & \mathrm{B}/(\mathfrak{n}^{j+1}+\mathfrak{m}\mathrm{B}) \ar[r,"w_j"'] & \mathrm{B}'/(\mathfrak{n}'^{j+1}+\mathfrak{m}\mathrm{B}')
  \end{tikzcd}
\]
are commutative; $\mathrm{B}$ and $\mathrm{B}'$ being complete, the homomorphism $v=\varprojlim v_j$ will give the answer. The inductive construction of $v_j$ follows from the hypothesis on $\mathrm{B}$ and from

\oldpage[0\textsubscript{IV-1}]{186}

the same reasoning as in \sref{0.19.3.10.1}, using the lemma \sref{0.19.3.10.2}, and the fact that $\mathfrak{n}'^j+(\mathfrak{n}'^{j+1}+\mathfrak{m}\mathrm{B}')=\mathfrak{n}'^j+\mathfrak{m}\mathrm{B}'$.
It follows then from \sref{0.19.7.1.4} that $v$ is an $\mathrm{A}$-\emph{isomorphism}, for $\mathrm{B}'$ is a flat $\mathrm{A}$-module, and $\mathrm{B}$ is complete for the $\mathfrak{m}$-preadic topology, the ideals $\mathfrak{m}^i\mathrm{B}$ being closed in $\mathrm{B}$ for the $\mathfrak{n}$-adic topology $(\mathbf{0}_\mathrm{I},~7.3.5)$.
\hfill Q.E.D.
\end{proof}

\subsection{Differential characterization of local algebras formally smooth over a field}
\label{subsection:0.22.2}

\begin{env}[22.2.1]
\label{0.22.2.1}
Let $k$ be a field, $\mathrm{P}$ its prime subfield, $\mathrm{A}$ a $k$-algebra that is a \emph{local} ring, $\mathfrak{m}$ its maximal ideal, $\mathrm{K}=\mathrm{A}/\mathfrak{m}$ its residue field.
As $\mathrm{K}$ is separable over $\mathrm{P}$, $\mathrm{K}$ is a $\mathrm{P}$-algebra formally smooth for the discrete topology \sref{0.19.6.1}, and consequently the $\mathrm{P}$-extension $\mathrm{A}/\mathfrak{m}^2$ of $\mathrm{K}$ by $\mathfrak{m}/\mathfrak{m}^2$ is $\mathrm{P}$-\emph{trivial}, and the characteristic homomorphism
\[
\label{0.22.2.1.1}
  \chi_{\mathrm{A}/k}:\Upsilon_{\mathrm{K}/k}\to\mathfrak{m}/\mathfrak{m}^2
  \tag{22.2.1.1}
\]
is therefore defined \sref{0.20.6.23}.
\end{env}

\begin{theorem}[22.2.2]
\label{0.22.2.2}
Under the conditions of \sref{0.22.2.1}, for $\mathrm{A}$ to be a $k$-algebra formally smooth for the $\mathfrak{m}$-preadic topology, it is necessary and sufficient that the following conditions be satisfied:
\begin{enumerate}
  \item[\rm{(i)}] \emph{The canonical homomorphism $\mathbf{S}^*_\mathrm{K}(\mathfrak{m}/\mathfrak{m}^2)\to\gr^*_\mathfrak{m}(\mathrm{A})$ is bijective.}
  \item[\rm{(ii)}] \emph{The characteristic homomorphism $\chi_{\mathrm{A}/k}:\Upsilon_{\mathrm{K}/k}\to\mathfrak{m}/\mathfrak{m}^2$ is injective.}
\end{enumerate}
This is a particular case of \sref{0.22.1.2}, where one replaces $\Lambda$ by $\mathrm{P}$, $\mathrm{A}$ by $k$, and $\mathrm{B}$ by $\mathrm{A}$; $\mathrm{K}$ and $k$, being separable over $\mathrm{P}$, are indeed formally smooth $\mathrm{P}$-algebras \sref{0.19.6.1}.
\end{theorem}

\begin{remark}[22.2.3]
\label{0.22.2.3}
The condition (ii) of \sref{0.22.2.2} can be replaced by any one of the following:
\begin{enumerate}
  \item[a)] The canonical homomorphism $\Omega^1_k\otimes_k\mathrm{K}\to\Omega^1_\mathrm{A}\otimes_\mathrm{A}\mathrm{K}$ is injective.
  \item[b)] For every $n$, the canonical homomorphism $\Omega^1_k\otimes_k(\mathrm{A}/\mathfrak{m}^{n+1})\to\Omega^1_\mathrm{A}\otimes_\mathrm{A}(\mathrm{A}/\mathfrak{m}^{n+1})$ is invertible on the left.
  \item[c)] The canonical homomorphism $\Omega^1_k\widehat{\otimes}_k\mathrm{A}\to\widehat{\Omega}^1_\mathrm{A}$ is invertible on the left.
  \item[d)] $\mathrm{A}$ is a $k$-algebra formally smooth (for the $\mathfrak{m}$-preadic topology) \emph{relative to the prime field}.
\end{enumerate}
Indeed, it was already noted in \sref{0.22.1.1} that the injectivity of $\chi_{\mathrm{A}/k}$ is equivalent to $a$), and that the conjunction of $a$) and condition (i) of \sref{0.22.2.2} is equivalent to that of (i) and $d$).
Finally, we have also seen in \sref{0.22.1.1} that $b$) and $d$) are equivalent, and the equivalence of $b$) and $c$) follows from \sref{0.19.1.8}.
\end{remark}

\begin{env}[22.2.4]
\label{0.22.2.4}
The main application of \sref{0.22.2.2} concerns Noetherian local $k$-algebras that are \emph{complete} (the case where $\mathfrak{m}/\mathfrak{m}^2$ is a $\mathrm{K}$-vector space of \emph{finite} rank).
More precisely, let $k$ be a field, $\mathrm{K}$ an extension of $k$, $\mathrm{V}$ a $\mathrm{K}$-vector space of
\oldpage[0\textsubscript{IV-1}]{187}
finite rank; we consider the triplets $(\mathrm{A},\alpha,\beta)$, where $\mathrm{A}$ is a complete Noetherian local $k$-algebra \emph{formally smooth} with maximal ideal $\mathfrak{m}$, $\alpha:\mathrm{A}/\mathfrak{m}\xrightarrow{\sim}\mathrm{K}$ a $k$-algebra isomorphism, $\beta:\mathfrak{m}/\mathfrak{m}^2\xrightarrow{\sim}\mathrm{V}$ a di-isomorphism of vector spaces (for the isomorphism $\alpha$).
Given two such triplets $(\mathrm{A},\alpha,\beta)$, $(\mathrm{A}',\alpha',\beta')$, an \emph{equivalence} from the first to the second is a $k$-isomorphism $u:\mathrm{A}\xrightarrow{\sim}\mathrm{A}'$ such that (if $\mathfrak{m}'$ is the maximal ideal of $\mathrm{A}'$), the diagrams
\[
  \begin{tikzcd}
    \mathrm{A}/\mathfrak{m} \ar[r,"\gr^0(u)"] \ar[d,"\alpha"] & \mathrm{A}'/\mathfrak{m}' \ar[d,"\alpha'"] & \mathfrak{m}/\mathfrak{m}^2 \ar[r,"\gr^1(u)"] \ar[d,"\beta"] & \mathfrak{m}'/\mathfrak{m}'^2 \ar[d,"\beta'"]\\
    \mathrm{K} \ar[r,"1_\mathrm{K}"] & \mathrm{K} & \mathrm{V} \ar[r,"1_\mathrm{V}"] & \mathrm{V}
  \end{tikzcd}
\]
are commutative.
It is immediate that the \emph{equivalence classes} of triplets $(\mathrm{A},\alpha,\beta)$ form a \emph{set} that we shall denote $\operatorname{FL}(\mathrm{K}/k,\mathrm{V})$.
We now note that to every triplet $(\mathrm{A},\alpha,\beta)$ is associated the $\mathrm{K}$-homomorphism composition of vector spaces $\Upsilon_{\mathrm{K}/k}\xrightarrow{\chi_{\mathrm{A}/k}}\mathfrak{m}/\mathfrak{m}^2\xrightarrow{\beta}\mathrm{V}$ (where $\chi_{\mathrm{A}/k}$ is considered as a di-homomorphism relative to $\alpha^{-1}$); the preceding definition shows (by \sref{0.20.6.22.2}) that this $\mathrm{K}$-homomorphism \emph{depends only on the equivalence class} of $(\mathrm{A},\alpha,\beta)$ in $\operatorname{FL}(\mathrm{K}/k,\mathrm{V})$; in other words we have thus defined a canonical map
\[
\label{0.22.2.4.1}
  \operatorname{FL}(\mathrm{K}/k,\mathrm{V})\to\Hom_\mathrm{K}(\Upsilon_{\mathrm{K}/k},\mathrm{V}).
  \tag{22.2.4.1}
\]
\end{env}

\begin{proposition}[22.2.5]
\label{0.22.2.5}
Let $k$ be a field, $\mathrm{K}$ an extension of $k$, $\mathrm{V}$ a $\mathrm{K}$-vector space of finite rank.
The canonical map \sref{0.22.2.4.1} is a bijection from $\operatorname{FL}(\mathrm{K}/k,\mathrm{V})$ to the set of injective $\mathrm{K}$-homomorphisms from $\Upsilon_{\mathrm{K}/k}$ to $\mathrm{V}$.
\end{proposition}

\begin{proof}
(i) To show that \sref{0.22.2.4.1} is injective, it must be proved that the following is true: let $\mathrm{A}$, $\mathrm{A}'$ be two complete Noetherian local $k$-algebras formally smooth, $\mathfrak{m}$, $\mathfrak{m}'$ their respective maximal ideals, $\alpha:\mathrm{A}/\mathfrak{m}\xrightarrow{\sim}\mathrm{K}$, $\alpha':\mathrm{A}'/\mathfrak{m}'\xrightarrow{\sim}\mathrm{K}$ two $k$-isomorphisms; let us be given two $k$-isomorphisms $v^0:\mathrm{A}/\mathfrak{m}\xrightarrow{\sim}\mathrm{A}'/\mathfrak{m}'$, $v^1:\mathfrak{m}/\mathfrak{m}^2\xrightarrow{\sim}\mathfrak{m}'/\mathfrak{m}'^2$ such that the diagrams
\[
\label{0.22.2.5.1}
  \begin{tikzcd}
    \mathrm{A}/\mathfrak{m} \ar[r,"v^0"] \ar[d,"\alpha"] & \mathrm{A}'/\mathfrak{m}' \ar[dl,"\alpha'"] & \mathfrak{m}/\mathfrak{m}^2 \ar[r,"v^1"] & \mathfrak{m}'/\mathfrak{m}'^2\\
    \mathrm{K} & & \Upsilon_{\mathrm{K}/k} \ar[u,"\chi_\mathrm{A}"] \ar[ur,"\chi_{\mathrm{A}'}"'] &
  \end{tikzcd}
  \tag{22.2.5.1}
\]
are commutative.
One must prove that there exists a $k$-isomorphism $u:\mathrm{A}\xrightarrow{\sim}\mathrm{A}'$ such that $\gr^0(u)=v^0$ and $\gr^1(u)=v^1$.
\oldpage[0\textsubscript{IV-1}]{188}
We note first that since $\mathrm{K}$ is a field, the canonical homomorphism \sref{0.20.6.7.1} is \emph{bijective} \sref{0.20.6.11}, so there already exists a $k$-isomorphism $v:\mathrm{A}/\mathfrak{m}^2\xrightarrow{\sim}\mathrm{A}'/\mathfrak{m}'^2$ such that $\gr^0(v)=v^0$ and $\gr^1(v)=v^1$.
We now note that since $\mathrm{A}'$ is metrizable and complete, and $\mathrm{A}$ a $k$-algebra formally smooth, the composite homomorphism $\mathrm{A}\to\mathrm{A}/\mathfrak{m}^2\xrightarrow{v}\mathrm{A}'/\mathfrak{m}'^2$ factors as $\mathrm{A}\xrightarrow{u}\mathrm{A}'\to\mathrm{A}'/\mathfrak{m}'^2$, where $u$ is a $k$-homomorphism \sref{0.19.3.11}, and one has thus already $\gr^0(u)=v^0$ and $\gr^1(u)=v^1$.
One has the commutative diagram
\[
  \begin{tikzcd}
    \mathbf{S}^*_\mathrm{K}(\mathfrak{m}/\mathfrak{m}^2) \ar[r] \ar[d,"w"'] & \gr^*_\mathfrak{m}(\mathrm{A}) \ar[d,"\gr^*(u)"]\\
    \mathbf{S}^*_\mathrm{K}(\mathfrak{m}'/\mathfrak{m}'^2) \ar[r] & \gr^*_{\mathfrak{m}'}(\mathrm{A}')
  \end{tikzcd}
\]
where the horizontal arrows are the canonical homomorphisms and $w$ is induced by $v^0$ and $v^1$. Since $v^0$ and $v^1$ are bijective, so is $w$; and since $\mathrm{A}$ and $\mathrm{A}'$ are $k$-algebras formally smooth, one deduces from \sref{0.22.2.2}, (i) that $\gr(u)$ is \emph{bijective}; but since $\mathrm{A}$ and $\mathrm{A}'$ are separated and complete, this implies that $u$ itself is \emph{bijective} (Bourbaki, \emph{Alg.~comm.}, chap.~III, \textsection~2, n\textsuperscript{o}~8, cor.~3 of th.~1).

(ii) It remains to show that the image of \sref{0.22.2.4.1} is the set of \emph{injective} homomorphisms of $\Upsilon_{\mathrm{K}/k}$ in $\mathrm{V}$; one already knows by \sref{0.22.2.2}, (ii) that it is contained in this set, so it remains to prove the
\end{proof}

\begin{lemma}[22.2.5.2]
\label{0.22.2.5.2}
For every $\mathrm{K}$-homomorphism $h:\Upsilon_{\mathrm{K}/k}\to\mathrm{V}$, there exists a $k$-algebra $\mathrm{A}$ that is a local Noetherian ring, regular and complete, with maximal ideal $\mathfrak{m}$, with residue field $\mathrm{K}$, and a $\mathrm{K}$-isomorphism $\beta:\mathfrak{m}/\mathfrak{m}^2\xrightarrow{\sim}\mathrm{V}$ such that $h=\beta\circ\chi_{\mathrm{A}/k}$.
\end{lemma}

\begin{proof}
Indeed, if this lemma is proved, the fact that $\mathrm{A}$ is regular implies that condition (i) of \sref{0.22.2.2} is satisfied \sref{0.17.1.1}; if moreover $h$ is injective, \sref{0.22.2.2} will show that $\mathrm{A}$ is a $k$-algebra formally smooth whose equivalence class maps to $h$ under \sref{0.22.2.4.1} (while if $h$ is not injective, the $k$-algebra $\mathrm{A}$ of which \sref{0.22.2.5.2} proves the existence is not formally smooth).

To prove \sref{0.22.2.5.2}, let $\mathrm{B}$ be the $\mathrm{K}$-algebra $\mathbf{S}^*_\mathrm{K}(\mathrm{V})$, let $\mathfrak{n}$ be the augmentation ideal of $\mathrm{B}$, and let $\mathrm{A}=\widehat{\mathrm{B}}$ be the completion of $\mathrm{B}$ for the $\mathfrak{n}$-preadic topology (so that $\mathrm{A}$ is isomorphic to a formal power series algebra $\mathrm{K}[[\mathrm{T}_1,\ldots,\mathrm{T}_n]]$ if $n=\rg_\mathrm{K}\mathrm{V}$).
By virtue of \sref{0.20.6.11}, there exists on $\mathrm{A}/\mathfrak{m}^2$ (where $\mathfrak{m}=\mathfrak{n}\mathrm{A}$) a structure of $k$-extension of $\mathrm{K}$ by $\mathrm{V}$, in general distinct from the $k$-trivial structure defined by the canonical injection $k\to\mathrm{K}\to\mathrm{A}$, such that $h$ is the characteristic homomorphism of this extension.
If $f:k\to\mathrm{A}/\mathfrak{m}^2$ is the structural homomorphism, the fact that $k$ is separable over the prime field permits \sref{0.19.6.1} factoring $f$ as $k\xrightarrow{g}\mathrm{A}\to\mathrm{A}/\mathfrak{m}^2$, and the $k$-algebra $\mathrm{A}$ thus defined answers the question \sref{0.20.6.22.2}.
\end{proof}

\begin{theorem}[22.2.6]
\label{0.22.2.6}
Let $k$ be a field, $\mathrm{K}$ an extension of $k$.
For there to exist a $k$-algebra $\mathrm{A}$, which is a local Noetherian ring, with maximal ideal $\mathfrak{m}$, such that $\mathrm{A}/\mathfrak{m}=\mathrm{K}$ and such that $\mathrm{A}$ is formally smooth (for the $\mathfrak{m}$-preadic topology), it is necessary and sufficient that $\Upsilon_{\mathrm{K}/k}$ is a $\mathrm{K}$-vector space of finite rank.
The $k$-algebra $\mathrm{A}$ is then determined (up to $k$-isomorphism inducing the identity $\mathrm{K}\to\mathrm{K}$ on residue fields) by its dimension alone, which is subject to the sole condition of being $\geq\rg_\mathrm{K}\Upsilon_{\mathrm{K}/k}$.
\end{theorem}

\begin{proof}
Since $\mathfrak{m}/\mathfrak{m}^2$ is of finite rank over $\mathrm{K}$, the necessity of the condition follows from \sref{0.22.2.2}, (ii); conversely, \sref{0.22.2.5} shows that the condition is sufficient and that one can take for $\mathfrak{m}/\mathfrak{m}^2$ a $\mathrm{K}$-vector space of rank any $\geq\rg_\mathrm{K}(\Upsilon_{\mathrm{K}/k})$.
\end{proof}

\oldpage[0\textsubscript{IV-1}]{189}
In particular, the identity homomorphism of $\Upsilon_{\mathrm{K}/k}$ associates \emph{canonically} to $\mathrm{K}$ a complete Noetherian $k$-algebra, formally smooth, for which $\Upsilon_{\mathrm{K}/k}$ is isomorphic to $\mathfrak{m}/\mathfrak{m}^2$.

\begin{remark}[22.2.7]
\label{0.22.2.7}
\begin{enumerate}
  \item[\rm{(i)}] Let $\Lambda$ be a local Noetherian ring, with residue field $k$; it follows from \sref{0.19.7.1} and \sref{0.19.7.2} that the determination of the $\Lambda$-algebras $\mathrm{A}$ that are complete Noetherian local rings and are formally smooth is \emph{equivalent} to the same problem where $\Lambda$ is replaced by $k$, and is thus in principle entirely solved by \sref{0.22.2.5}, which reduces the question to the structure of the \emph{residue fields} of the $k$-algebras sought, as extensions of $k$.
  \item[\rm{(ii)}] The fact that $\Upsilon_{\mathrm{K}/k}$ is of finite rank over $\mathrm{K}$ does \emph{not imply} that $\mathrm{K}$ is of ``finite radicial multiplicity'' over $k$ (cf.~\sref{0.19.6.6}).
  For example, let $k_0$ be a perfect field of characteristic $p>0$, $k=k_0(\mathrm{T})$ the field of rational functions in one indeterminate over $k_0$, so that $\mathrm{K}=k^{p^{-\infty}}$ is the smallest perfect extension of $k$.
  Then $\Omega^1_\mathrm{K}=\Omega^1_{\mathrm{K}/k}=0$ \sref{0.21.4.4}, therefore $\Upsilon_{\mathrm{K}/k}=\Omega^1_k\otimes_k\mathrm{K}$, and since $k^p=k_0(\mathrm{T}^p)$, $\mathrm{T}$ is a $p$-base of $k$ over $k^p$, therefore \sref{0.21.4.5}, $\Omega^1_k$ is of rank $1$ over $k$, and hence $\Upsilon_{\mathrm{K}/k}$ is of rank $1$ over $\mathrm{K}$.
  But for every $n$ the residue field of the local ring $\mathrm{K}\otimes_kk^{p^{-n}}$ is isomorphic to $\mathrm{K}$, which is therefore not separable over $k^{p^{-n}}$.
\end{enumerate}
\end{remark}

\begin{proposition}[22.2.8]
\label{0.22.2.8}
Under the conditions of \sref{0.22.2.1}, suppose that the characteristic $p$ of $k$ is $>0$; let $(k_\alpha)$ be a decreasing filtering family of subfields of $k$, such that $\bigcap k_\alpha(k^p)=k^p$.
Then the condition (ii) of \sref{0.22.2.2} can be replaced by any one of the following:
\begin{enumerate}
  \item[a)] There exists $\alpha_0$ such that the canonical homomorphism
  \[
    \Omega^1_{k/k_\alpha}\otimes_k\mathrm{K}\to\Omega^1_{\mathrm{A}/k_\alpha}\otimes_\mathrm{A}\mathrm{K}
  \]
  is injective for every $\alpha\geq\alpha_0$.
  \item[b)] There exists $\alpha_0$ such that $\mathrm{A}$ is a $k$-algebra formally smooth (for the $\mathfrak{m}$-preadic topology) relative to $k_\alpha$, for every $\alpha\geq\alpha_0$.
  \item[c)] For every $\alpha$ and every integer $n\geq 0$, the canonical homomorphism
  \[
    \Omega^1_{k/k_\alpha}\otimes_k(\mathrm{A}/\mathfrak{m}^{n+1})\to\Omega^1_{\mathrm{A}/k_\alpha}\otimes_\mathrm{A}(\mathrm{A}/\mathfrak{m}^{n+1})
  \]
  is invertible on the left.
\end{enumerate}
\end{proposition}

\begin{proof}
Indeed, if $\mathrm{A}$ is a formally smooth $k$-algebra, it is so \emph{a fortiori} relative to every subfield of $k$, which implies $c$) by virtue of \sref{0.20.7.2} and \sref{0.19.1.7}; it is clear that $c$) implies $b$), and that $b$) implies $a$).
Finally, taking into account \sref{0.22.2.3}, it remains to prove that $a$) implies that the canonical homomorphism $u:\Omega^1_k\otimes_k\mathrm{K}\to\Omega^1_\mathrm{A}\otimes_\mathrm{A}\mathrm{K}$ is \emph{injective}.
But for every $\alpha$, one has a commutative diagram
\[
  \begin{tikzcd}
    \Omega^1_k\otimes_k\mathrm{K} \ar[r,"u"] \ar[d,"v_\alpha"'] & \Omega^1_\mathrm{A}\otimes_\mathrm{A}\mathrm{K} \ar[d]\\
    \Omega^1_{k/k_\alpha}\otimes_k\mathrm{K} \ar[r,"u_\alpha"'] & \Omega^1_{\mathrm{A}/k_\alpha}\otimes_\mathrm{A}\mathrm{K}
  \end{tikzcd}
\]

\oldpage[0\textsubscript{IV-1}]{190}

and the hypothesis that $u_\alpha$ is injective for $\alpha\geq\alpha_0$ implies that $\Ker(u)$ is contained in the intersection of the $\Ker(v_\alpha)$.
But one has seen \sref{0.21.8.3} that this intersection is zero by virtue of the hypothesis $\bigcap k_\alpha(k^p)=k^p$, which finishes the proof.
\end{proof}

\begin{proposition}[22.2.9]
\label{0.22.2.9}
Under the hypotheses of \sref{0.22.2.1}, the following conditions are equivalent:
\begin{enumerate}
  \item[a)] $\mathrm{A}/\mathfrak{m}^2$ is a $k$-extension $k$-trivial of $\mathrm{K}$ by $\mathfrak{m}/\mathfrak{m}^2$.
  \item[b)] $\chi_{\mathrm{A}/k}=0$.
  \item[c)] The canonical homomorphism \sref{0.20.5.11.2}
  \[
    \delta_{\mathrm{K}/\mathrm{A}/k}:\mathfrak{m}/\mathfrak{m}^2\to\Omega^1_{\mathrm{A}/k}\otimes_\mathrm{A}\mathrm{K}
  \]
  is injective (in other words, the sequence
\end{enumerate}
\[
\label{0.22.2.9.1}
  0\to\mathfrak{m}/\mathfrak{m}^2\to\Omega^1_{\mathrm{A}/k}\otimes_\mathrm{A}\mathrm{K}\to\Omega^1_{\mathrm{K}/k}\to 0
  \tag{22.2.9.1}
\]
\emph{is exact}).
\end{proposition}

\begin{proof}
The equivalence of $b$) and $c$) follows from the exact sequence \sref{0.20.6.22.1}; the equivalence of $c$) and $a$) follows from \sref{0.20.5.12}, this sequence being split if it is exact, since these are vector spaces.

Note that if $\mathrm{K}$ is \emph{separable} over $k$, the equivalent conditions of \sref{0.22.2.9} are satisfied, by virtue of the definition of $\chi_{\mathrm{A}/k}$ \sref{0.22.2.1} and of \sref{0.20.6.3}.
\end{proof}

\begin{corollary}[22.2.10]
\label{0.22.2.10}
Under the hypotheses of \sref{0.22.2.1}, consider a subfield $k_0$ of $k$.
The following conditions are equivalent:
\begin{enumerate}
  \item[a)] $\mathrm{A}/\mathfrak{m}^2$ is a $k_0$-extension $k_0$-trivial of $\mathrm{K}$ by $\mathfrak{m}/\mathfrak{m}^2$.
  \item[b)] The composite homomorphism $\Upsilon_{\mathrm{K}/k_0}\to\Upsilon_{\mathrm{K}/k}\xrightarrow{\chi_{\mathrm{A}/k}}\mathfrak{m}/\mathfrak{m}^2$ is zero.
  \item[c)] The characteristic homomorphism $\chi_{\mathrm{A}/k}:\Upsilon_{\mathrm{K}/k}\to\mathfrak{m}/\mathfrak{m}^2$ factors through $\Upsilon_{\mathrm{K}/k}\to\Upsilon_{\mathrm{K}/k/k_0}\xrightarrow{\chi'}\mathfrak{m}/\mathfrak{m}^2$.
\end{enumerate}
Moreover, when these conditions are satisfied, the homomorphism $\chi'$ is determined in a unique and canonical way, and coincides with the characteristic homomorphism of the $k_0$-extension $k_0$-trivial $\mathrm{A}/\mathfrak{m}^2$ of $\mathrm{K}$ by $\mathfrak{m}/\mathfrak{m}^2$.
\end{corollary}

\begin{proof}
Since the composition in $b$) is nothing other than $\chi_{\mathrm{A}/k_0}$, the equivalence of $a$) and $b$) follows from \sref{0.22.2.9}; on the other hand, $b$) and $c$) are equivalent by virtue of the exact sequence \sref{0.20.6.18.1}
\[
  \Upsilon_{\mathrm{K}/k_0}\to\Upsilon_{\mathrm{K}/k}\to\Upsilon_{\mathrm{K}/k/k_0}\to0,
\]
and since the canonical homomorphism $\Upsilon_{\mathrm{K}/k}\to\Upsilon_{\mathrm{K}/k/k_0}$ is surjective, this implies the uniqueness of $\chi'$; the fact that $\chi'$ is the characteristic homomorphism of the $k_0$-extension $\mathrm{A}/\mathfrak{m}^2$ follows from \sref{0.20.6.22.2}.
\end{proof}

\begin{corollary}[22.2.11]
\label{0.22.2.11}
Under the hypotheses of \sref{0.22.2.1}, let $p$ be the characteristic exponent of $k$, and let $(k_\alpha)$ be a decreasing filtering family of subfields of $k$ such that $\bigcap k_\alpha(k^p)=k^p$.
If $\Upsilon_{\mathrm{K}/k}$ is of finite rank over $\mathrm{K}$, there exists an index $\alpha$ such that $\mathrm{A}/\mathfrak{m}^2$ is a $k_\alpha$-extension $k_\alpha$-trivial of $\mathrm{K}$ by $\mathfrak{m}/\mathfrak{m}^2$.
\end{corollary}

\begin{proof}
It is known \sref{0.21.8.3} that there exists then $\alpha$ such that the canonical homomorphism $\Upsilon_{\mathrm{K}/k}\to\Upsilon_{\mathrm{K}/k/k_\alpha}$ is bijective, so the condition $b$) of \sref{0.22.2.10} (where $k_\alpha$ replaces $k_0$) is satisfied.
\end{proof}

\subsection{Application to relations between certain local rings and their completions}
\label{subsection:0.22.3}

\oldpage[0\textsubscript{IV-1}]{191}
\begin{lemma}[22.3.1]
\label{0.22.3.1}
Let $\mathrm{A}_0$ be a semi-local Noetherian ring, $\mathrm{A}$ a finite $\mathrm{A}_0$-algebra (which is thus a semi-local Noetherian ring, cf.~Bourbaki, \emph{Alg.~comm.}, chap.~IV, \textsection~2, n\textsuperscript{o}~5, cor.~3 of prop.~9).
If $\widehat{\mathrm{A}}_0$ and $\widehat{\mathrm{A}}$ are the completions of $\mathrm{A}_0$ and $\mathrm{A}$ for their preadic topologies respectively, then, when one equips $\mathrm{A}_0$, $\mathrm{A}$ and $\widehat{\mathrm{A}}$ with discrete topologies, $\widehat{\mathrm{A}}$ is an $\mathrm{A}$-algebra formally smooth relative to $\mathrm{A}_0$ \sref{0.19.9.1}.
\end{lemma}

\begin{proof}
It is known that $\widehat{\mathrm{A}}=\mathrm{A}\otimes_{\mathrm{A}_0}\widehat{\mathrm{A}}_0$ (Bourbaki, \emph{Alg.~comm.}, chap.~IV, \textsection~2, n\textsuperscript{o}~5, cor.~3 of prop.~9 and chap.~III, \textsection~3, n\textsuperscript{o}~4, th.~1), so the proposition follows from \sref{0.19.9.3}.
\end{proof}

\begin{proposition}[22.3.2]
\label{0.22.3.2}
Let $\mathrm{A}$ be a local Noetherian regular ring, $\mathrm{K}$ its field of fractions, $p$ the characteristic of $\mathrm{K}$.
Suppose that one of the following hypotheses is satisfied:
\begin{enumerate}
  \item[\rm{(i)}] $p=0$.
  \item[\rm{(ii)}] $p>0$ and there exists a filtering decreasing family $(\mathrm{A}_\alpha)$ of Noetherian subrings of $\mathrm{A}$, such that $\mathrm{A}$ is a finite $\mathrm{A}_\alpha$-algebra for every $\alpha$, such that $\mathrm{A}^{p^{h(\alpha)}}\subset\mathrm{A}_\alpha$ for a suitable integer $h(\alpha)$, and such that, if $\mathrm{K}_\alpha$ denotes the field of fractions of $\mathrm{A}_\alpha$, we have $\bigcap\mathrm{K}_\alpha(\mathrm{K}^p)=\mathrm{K}^p$.
\end{enumerate}

Let then $\mathrm{B}$ be a finite $\mathrm{A}$-algebra that is an integral domain and contains $\mathrm{A}$, $\mathfrak{n}$ a prime ideal of $\mathrm{B}$, $\mathrm{C}$ the local ring $\mathrm{B}_\mathfrak{n}$, $\mathfrak{q}$ a prime ideal of the completion $\widehat{\mathrm{C}}$, such that $\mathfrak{q}\cap\mathrm{C}=0$, so that the local ring $\widehat{\mathrm{C}}_\mathfrak{q}$ is an algebra over the field of fractions $\mathrm{L}$ of $\mathrm{B}$.
Then $\widehat{\mathrm{C}}_\mathfrak{q}$ is an $\mathrm{L}$-algebra formally smooth for its $\mathfrak{q}$-adic topology, and is therefore a geometrically regular ring over $\mathrm{L}$ \sref{0.19.6.5}.
\end{proposition}

\begin{proof}
We distinguish two cases.

\emph{I)} Suppose that $\mathfrak{n}$ contains the maximal ideal $\mathfrak{m}$ of $\mathrm{A}$, so that $\mathfrak{n}\cap\mathrm{A}=\mathfrak{m}$.
Then $\mathfrak{n}$ is maximal (Bourbaki, \emph{Alg.~comm.}, chap.~V, \textsection~2, n\textsuperscript{o}~1, prop.~1).
If $\mathfrak{r}$ is the radical of the semi-local ring $\mathrm{B}$, $\widehat{\mathrm{C}}$ is the separated completion of $\mathrm{B}$ for the $\mathfrak{n}$-preadic topology, and a component of the semi-local ring $\widehat{\mathrm{B}}$, the completion of $\mathrm{B}$ for the $\mathfrak{r}$-preadic topology (Bourbaki, \emph{Alg.~comm.}, chap.~III, \textsection~3, n\textsuperscript{o}~4, prop.~8);
one can thus write $\widehat{\mathrm{C}}_\mathfrak{q}=\widehat{\mathrm{B}}_{\mathfrak{q}'}$, where $\mathfrak{q}'$ is the prime ideal of $\widehat{\mathrm{B}}$ that is the inverse image of $\mathfrak{q}$.
Let us now note that $\mathrm{B}$ is a subring of $\widehat{\mathrm{B}}$, and the hypothesis $\mathfrak{q}\cap\mathrm{C}=0$ implies $\mathfrak{q}'\cap\mathrm{B}=0$, so $\widehat{\mathrm{B}}_{\mathfrak{q}'}$ is also a localization of $\widehat{\mathrm{B}}\otimes_\mathrm{B}\mathrm{L}$ at one of its prime ideals $\mathfrak{q}''$, of which $\mathfrak{q}'$ is the inverse image in $\widehat{\mathrm{B}}$;
moreover we have $\widehat{\mathrm{B}}=\widehat{\mathrm{A}}\otimes_\mathrm{A}\mathrm{B}$ (Bourbaki, \emph{Alg.~comm.}, chap.~IV, \textsection~2, n\textsuperscript{o}~5, cor. of prop.~9 and chap.~III, \textsection~3, n\textsuperscript{o}~4, th.~1), so $\widehat{\mathrm{B}}\otimes_\mathrm{B}\mathrm{L}=\widehat{\mathrm{A}}\otimes_\mathrm{A}\mathrm{L}$; finally, if $\mathfrak{p}$ is the prime ideal of $\widehat{\mathrm{A}}$ that is the inverse image of $\mathfrak{q}''$, the fact that $\mathrm{A}\to\mathrm{B}$ is injective and the commutativity of the diagram
\[
  \begin{tikzcd}
    \mathrm{B} \ar[r] & \widehat{\mathrm{B}}\\
    \mathrm{A} \ar[u] \ar[r] & \widehat{\mathrm{A}} \ar[u]
  \end{tikzcd}
\]
imply that $\mathfrak{p}\cap\mathrm{A}=0$, so $(\widehat{\mathrm{A}}\otimes_\mathrm{A}\mathrm{L})_{\mathfrak{q}''}$ is also a localization of the ring $\widehat{\mathrm{A}}_\mathfrak{p}\otimes_\mathrm{A}\mathrm{L}=(\widehat{\mathrm{A}}_\mathfrak{p}\otimes_\mathrm{A}\mathrm{K})\otimes_\mathrm{K}\mathrm{L}$.

\oldpage[0\textsubscript{IV-1}]{192}

To show that this local ring is an $\mathrm{L}$-algebra formally smooth, it suffices therefore to prove that $\widehat{\mathrm{A}}_\mathfrak{p}\otimes_\mathrm{A}\mathrm{K}$ is a $\mathrm{K}$-algebra formally smooth \sref{0.19.3.5}, (iii) and (iv); moreover $\widehat{\mathrm{A}}_\mathfrak{p}\otimes_\mathrm{A}\mathrm{K}=\widehat{\mathrm{A}}_\mathfrak{p}$ since $\mathfrak{p}\cap\mathrm{A}=0$.
We shall apply the criteria \sref{0.22.2.2} and \sref{0.22.2.8}.
First of all, as $\mathrm{A}$ is regular, so is $\widehat{\mathrm{A}}$ \sref{0.17.1.5}, and hence so is $\widehat{\mathrm{A}}_\mathfrak{p}$ \sref{0.17.3.2}; condition (i) of \sref{0.22.2.2} is therefore satisfied \sref{0.17.1.1}.
It therefore suffices (when $p>0$) to verify condition $b$) of \sref{0.22.2.8}.
But \sref{0.22.3.1} shows that, for every $\alpha$, $\widehat{\mathrm{A}}$ is an $\mathrm{A}$-algebra formally smooth relative to $\mathrm{A}_\alpha$ (for the discrete topologies); hence $\widehat{\mathrm{A}}\otimes_{\mathrm{A}_\alpha}\mathrm{K}_\alpha=\widehat{\mathrm{A}}\otimes_\mathrm{A}(\mathrm{A}\otimes_{\mathrm{A}_\alpha}\mathrm{K}_\alpha)$ is an $(\mathrm{A}\otimes_{\mathrm{A}_\alpha}\mathrm{K}_\alpha)$-algebra formally smooth relative to $\mathrm{K}_\alpha$, for the discrete topologies \sref{0.19.9.2}, (iii).
Since $\mathrm{A}$ is an $\mathrm{A}_\alpha$-algebra \emph{finite}, we have $\mathrm{A}\otimes_{\mathrm{A}_\alpha}\mathrm{K}_\alpha=\mathrm{K}$; as $\widehat{\mathrm{A}}_\mathfrak{p}$ is a ring of fractions of $\widehat{\mathrm{A}}\otimes_\mathrm{A}\mathrm{K}$, we deduce \sref{0.19.9.2}, (iv) that $\widehat{\mathrm{A}}_\mathfrak{p}$ is a $\mathrm{K}$-algebra formally smooth \emph{relative} to $\mathrm{K}_\alpha$, for the discrete topology, and \emph{a fortiori} for its $\mathfrak{p}$-adic topology \sref{0.19.3.8} and \sref{0.19.9.5}.
But this is precisely condition $b$) of \sref{0.22.2.8}, so the proposition is proved in this case for $p>0$.
When $p=0$, the residue field of $\widehat{\mathrm{A}}_\mathfrak{p}$ is separable over $\mathrm{K}$, and since $\widehat{\mathrm{A}}_\mathfrak{p}$ is a regular ring, it is a $\mathrm{K}$-algebra formally smooth by virtue of \sref{0.19.6.4}.

\emph{II)} General case.
Let $\mathfrak{s}$ be the prime ideal $\mathfrak{n}\cap\mathrm{A}$ of $\mathrm{A}$, and set $\mathrm{S}=\mathrm{A}-\mathfrak{s}$, so that $\mathrm{S}^{-1}\mathrm{A}=\mathrm{A}_\mathfrak{s}$, and $\mathrm{C}=(\mathrm{S}^{-1}\mathrm{B})_{\mathrm{S}^{-1}\mathfrak{n}}$, where this time $\mathrm{S}^{-1}\mathfrak{n}$ contains the maximal ideal $\mathfrak{s}\mathrm{A}_\mathfrak{s}$ of $\mathrm{A}_\mathfrak{s}$.
As $\mathrm{A}_\mathfrak{s}$ is regular \sref{0.17.3.2} and $\mathrm{S}^{-1}\mathrm{B}$ is a finite $\mathrm{A}_\mathfrak{s}$-algebra that is an integral domain and contains $\mathrm{A}_\mathfrak{s}$, it remains only to verify condition (ii) for $\mathrm{A}_\mathfrak{s}$ when $p>0$: set $\mathfrak{s}_\alpha=\mathfrak{s}\cap\mathrm{A}_\alpha$ and consider the local ring $(\mathrm{A}_\alpha)_{\mathfrak{s}_\alpha}$.
For every $t\notin\mathfrak{s}$, one has by hypothesis, setting $q=p^{h(\alpha)}$, $t^q\in\mathrm{A}_\alpha$ and $t^q\notin\mathfrak{s}$, hence $t^q\notin\mathfrak{s}_\alpha$; if one sets $\mathrm{S}_\alpha=\mathrm{A}_\alpha-\mathfrak{s}_\alpha$, one sees therefore that $\mathrm{A}_\mathfrak{s}=\mathrm{S}_\alpha^{-1}\mathrm{A}$, and the hypothesis implies that $\mathrm{A}_\mathfrak{s}$ is a finite $(\mathrm{A}_\alpha)_{\mathfrak{s}_\alpha}$-algebra; moreover, $\mathrm{K}_\alpha$ is also the field of fractions of $(\mathrm{A}_\alpha)_{\mathfrak{s}_\alpha}$, which finishes the proof.
\end{proof}

\begin{theorem}[22.3.3]
\label{0.22.3.3}
Let $\mathrm{A}$ be a local Noetherian complete ring, $\mathfrak{p}$ a prime ideal of $\mathrm{A}$, $\mathrm{B}$ the localized ring $\mathrm{A}_\mathfrak{p}$, $\mathfrak{q}$ a prime ideal of the completion $\widehat{\mathrm{B}}$, $\mathfrak{r}=\mathfrak{q}\cap\mathrm{B}$.
Then the localized ring $\widehat{\mathrm{B}}_\mathfrak{q}$ is a $\mathrm{B}_\mathfrak{r}$-algebra formally smooth for the preadic topologies.
\end{theorem}

\begin{proof}
The prime ideal $\mathfrak{r}$ of $\mathrm{B}$ is of the form $\mathfrak{n}_\mathfrak{p}$, where $\mathfrak{n}\subset\mathfrak{p}$ is a prime ideal of $\mathrm{A}$, and $\mathrm{B}_\mathfrak{r}=\mathrm{A}_\mathfrak{n}$; let $k$ be the residue field of $\mathrm{B}_\mathfrak{r}$, which is also the field of fractions of the complete integral local ring $\mathrm{A}'=\mathrm{A}/\mathfrak{n}$.
As $\widehat{\mathrm{B}}$ is a flat $\mathrm{B}$-module $(\mathbf{0}_\mathrm{I},~7.3.5)$, $\widehat{\mathrm{B}}_\mathfrak{q}$ is a flat $\mathrm{B}_\mathfrak{r}$-module $(\mathbf{0}_\mathrm{I},~6.3.2)$, and to prove that $\widehat{\mathrm{B}}_\mathfrak{q}$ is a $\mathrm{B}_\mathfrak{r}$-algebra formally smooth, it suffices to prove that
\[
  \widehat{\mathrm{B}}_\mathfrak{q}\otimes_{\mathrm{B}_\mathfrak{r}}k
  =\widehat{\mathrm{B}}_\mathfrak{q}/\mathfrak{r}\widehat{\mathrm{B}}_\mathfrak{q}
\]
is a $k$-algebra formally smooth \sref{0.19.7.1}.
But $\widehat{\mathrm{B}}_\mathfrak{q}/\mathfrak{r}\widehat{\mathrm{B}}_\mathfrak{q}=(\widehat{\mathrm{B}}/\mathfrak{r}\widehat{\mathrm{B}})_{\mathfrak{q}'}$, where $\mathfrak{q}'$ is the canonical image of $\mathfrak{q}$ in $\widehat{\mathrm{B}}/\mathfrak{r}\widehat{\mathrm{B}}$; one has by definition $\mathfrak{q}'\cap(\mathrm{B}/\mathfrak{r})=0$ and $\mathrm{B}/\mathfrak{r}=\mathrm{A}_{\mathfrak{p}'}$, where $\mathfrak{p}'=\mathfrak{p}/\mathfrak{n}$; moreover $\widehat{\mathrm{B}}/\mathfrak{r}\widehat{\mathrm{B}}$ is the completion of $\mathrm{B}/\mathfrak{r}$.

We are thus reduced to proving the following corollary:
\end{proof}

\begin{corollary}[22.3.4]
\label{0.22.3.4}
Let $\mathrm{A}$ be a complete Noetherian local integral domain, $\mathrm{K}$ its field of fractions, $\mathfrak{p}$ a prime ideal of $\mathrm{A}$, $\mathrm{B}$ the localized ring $\mathrm{A}_\mathfrak{p}$.
Then, for every prime ideal $\mathfrak{q}$

\oldpage[0\textsubscript{IV-1}]{193}

of the completion $\widehat{\mathrm{B}}$, such that $\mathfrak{q}\cap\mathrm{B}=0$, the local ring $\widehat{\mathrm{B}}_\mathfrak{q}$ is a $\mathrm{K}$-algebra formally smooth for its $\mathfrak{q}$-adic topology, and is therefore a geometrically regular ring over $\mathrm{K}$.
\end{corollary}

\begin{proof}
Indeed, it follows from \sref{0.19.8.9} that there exists a subring $\mathrm{A}_0$ of $\mathrm{A}$ that is a power series ring over a field of characteristic $p>0$ or over a Cohen ring (recall that Cohen rings have a field of fractions of characteristic $0$), and such that $\mathrm{A}$ is a finite $\mathrm{A}_0$-algebra.
Now, we know that the ring $\mathrm{A}_0$ is \emph{regular} \sref{0.17.3.8}, and it satisfies one of the hypotheses (i), (ii) of \sref{0.22.3.2}, by virtue of \sref{0.21.8.8}; it suffices therefore to apply \sref{0.22.3.2}, replacing $\mathrm{B}$ by $\mathrm{A}$ and $\mathrm{A}$ by $\mathrm{A}_0$.
\end{proof}

\subsection{Preliminary results on finite extensions of local rings whose maximal ideal has square zero}
\label{subsection:0.22.4}
\label{0.22.4}

\begin{proposition}[22.4.1]
\label{0.22.4.1}
Let $\mathrm{A}$ be a local ring whose maximal ideal $\mathfrak{m}$ has square zero, $\mathrm{K}=\mathrm{A}/\mathfrak{m}$ its residue field; let us denote by $\mathrm{V}$ the ideal $\mathfrak{m}$ considered as a $\mathrm{K}$-vector space.
Let $\mathrm{T}_i$ $(1\leq i\leq r)$ be indeterminates, and for each $i$, let $\mathrm{F}_i$ be a monic polynomial of $\mathrm{A}[\mathrm{T}_i]$, of degree $d_i$; let $\mathrm{A}'$ be the quotient of the polynomial ring $\mathrm{A}[\mathrm{T}_1,\ldots,\mathrm{T}_r]$ by the ideal $\mathfrak{b}$ generated by the $r$ polynomials $\mathrm{F}_i$.
Suppose that $\mathrm{A}'$ is a local ring, and let $\mathfrak{m}'$ and $\mathrm{K}'=\mathrm{A}'/\mathfrak{m}'$ be its maximal ideal and residue field.
Then:
\begin{enumerate}
  \item[\rm{(i)}] If one sets $d=\prod_{i=1}^r d_i$, $\mathrm{A}'$ is a free $\mathrm{A}$-module of rank $d$, and we have $[\mathrm{K}':\mathrm{K}]\leq d$.
  \item[\rm{(ii)}] For $[\mathrm{K}':\mathrm{K}]=d$, it is necessary and sufficient that $\mathrm{A}'\otimes_\mathrm{A}\mathrm{K}=\mathrm{A}'/\mathfrak{m}\mathrm{A}'$ be a field (necessarily isomorphic to $\mathrm{K}'$), in other words, that $\mathfrak{m}'=\mathfrak{m}\mathrm{A}'$.
  In this case, $\mathfrak{m}'^2=0$, and if $\mathrm{V}'$ denotes the ideal $\mathfrak{m}'$ considered as a $\mathrm{K}'$-vector space, the canonical homomorphism $\mathrm{V}\otimes_\mathrm{K}\mathrm{K}'\to\mathrm{V}'$ is bijective (and consequently $\rg_{\mathrm{K}'}\mathrm{V}'=\rg_\mathrm{K}\mathrm{V}$).
\end{enumerate}
\end{proposition}

\begin{proof}
Without assuming $\mathrm{A}'$ local, it is evident by Euclidean division and induction on $r$ that, if $t_i$ is the class of $\mathrm{T}_i$ modulo $\mathfrak{b}$, the monomials $\mathrm{M}_\nu=t_1^{\nu_1}t_2^{\nu_2}\cdots t_r^{\nu_r}$, where $0\leq\nu_i\leq d_i-1$, form a basis of $\mathrm{A}'$ as an $\mathrm{A}$-module.
If $\mathrm{A}'$ is assumed local, $\mathrm{K}'$ is a quotient of $\mathrm{A}'/\mathfrak{m}\mathrm{A}'=\mathrm{A}'\otimes_\mathrm{A}\mathrm{K}$, so $[\mathrm{K}':\mathrm{K}]\leq[\mathrm{A}'\otimes_\mathrm{A}\mathrm{K}:\mathrm{K}]=d$, and there cannot be equality unless $\mathfrak{m}'=\mathfrak{m}\mathrm{A}'$, which proves (i) and the first assertion of (ii).
On the other hand, it is clear that when $\mathfrak{m}'=\mathfrak{m}\mathrm{A}'$, we have $\mathfrak{m}'^2=0$, and $\mathrm{V}\otimes_\mathrm{K}\mathrm{K}'=\mathfrak{m}\otimes_\mathrm{A}\mathrm{A}'=\mathfrak{m}\mathrm{A}'=\mathrm{V}'$ since $\mathrm{A}'$ is a free $\mathrm{A}$-module.
\end{proof}

\begin{remark}[22.4.2]
\label{0.22.4.2}
When $\mathrm{A}$ is Artinian (that is to say $\rg_\mathrm{K}(\mathrm{V})$ is finite), the hypothesis that $\mathrm{A}'$ is local always implies $\rg_{\mathrm{K}'}(\mathfrak{m}'/\mathfrak{m}'^2)\geq\rg_\mathrm{K}(\mathrm{V})$, as we shall see further on \sref{0.22.5.2}.
Note that these two ranks can be equal without having $\mathfrak{m}'=\mathfrak{m}\mathrm{A}'$.
\end{remark}

\begin{proposition}[22.4.3]
\label{0.22.4.3}
With the hypotheses and notations of \sref{0.22.4.1}, suppose that the $\mathrm{F}_i$ are of the form
\[
\label{0.22.4.3.1}
  \mathrm{F}_i(\mathrm{T}_i)=\mathrm{T}_i^{d_i}+\sum_{1\leq k\leq d_i}c_{ik}\mathrm{T}_i^{d_i-k}
  \tag{22.4.3.1}
\]
where $d_i\geq 2$ for every $i$ and $c_{ik}\in\mathfrak{m}$ (``Eisenstein polynomials'').
Then:
\begin{enumerate}
  \item[\rm{(i)}] $\mathrm{A}'$ is a local ring, and its residue field $\mathrm{K}'=\mathrm{A}'/\mathfrak{m}'$ is isomorphic to $\mathrm{K}$.

  \oldpage[0\textsubscript{IV-1}]{194}

  \item[\rm{(ii)}] The kernel of the canonical homomorphism $\mathrm{V}\to\mathrm{V}'=\mathfrak{m}'/\mathfrak{m}'^2$ is the ideal $\mathfrak{n}$ of $\mathrm{A}$ generated by the elements $\xi_i=c_{i,d_i}$, and its cokernel has as a basis over $\mathrm{K}'$ the images of the $\mathrm{T}_i$.
  In particular,
  \begin{equation}
  \label{0.22.4.3.2}
    \rg_{\mathrm{K}'}(\mathrm{V}')=\rg_\mathrm{K}(\mathrm{V})+r-r',
    \tag{22.4.3.2}
  \end{equation}
  where $r'=\rg_\mathrm{K}(\mathfrak{n})$.
  When $\rg_\mathrm{K}(\mathrm{V})$ is finite, equivalently when $\mathrm{A}$ is Artinian, equality $\rg_{\mathrm{K}'}(\mathrm{V}')=\rg_\mathrm{K}(\mathrm{V})$ holds if and only if the $\xi_i$ are linearly independent over $\mathrm{K}$.
\end{enumerate}
\end{proposition}

\begin{proof}
(i) It is clear that $\mathrm{A}'\otimes_\mathrm{A}\mathrm{K}=\mathrm{A}'/\mathfrak{m}\mathrm{A}'$ is isomorphic to $\mathrm{K}[\mathrm{T}_1,\ldots,\mathrm{T}_r]/\mathfrak{b}'$, where $\mathfrak{b}'$ is the ideal generated by the $r$ polynomials $\mathrm{T}_i^{d_i}$; this is therefore a local ring of residue field $\mathrm{K}$, and it is the same for $\mathrm{A}'$, from which (i) follows.

(ii) Let $t_i$ $(1\leq i\leq r)$ be the class of $\mathrm{T}_i$ mod.~$\mathfrak{b}$; it follows from (i) that the maximal ideal $\mathfrak{m}'$ of $\mathrm{A}'$ is given by
\[
  \mathfrak{m}'=\mathfrak{m}\mathrm{A}'+\sum_i t_i\mathrm{A}'
\]
from which, taking into account that $\mathfrak{m}^2=0$,
\[
  \mathfrak{m}'^2=\sum_{i,j}t_it_j\mathrm{A}'+\sum_i t_i\mathfrak{m}\mathrm{A}'.
\]
One deduces that $\mathrm{A}'/\mathfrak{m}'^2$ is isomorphic to the ring $\mathrm{A}[\mathrm{T}_1,\ldots,\mathrm{T}_r]/\mathfrak{c}$, where $\mathfrak{c}$ is the ideal generated by the elements
\[
  \mathrm{F}_i\ (1\leq i\leq r),\quad \mathrm{T}_i\mathrm{T}_j\ (1\leq i\leq r,\ 1\leq j\leq r),\quad x\mathrm{T}_i\ (1\leq i\leq r,\ x\in\mathfrak{m}).
\]

Since $d_i\geq 2$ and $\mathrm{T}_i^2\in\mathfrak{c}$, the hypotheses show that $\mathfrak{c}$ is also generated by
\[
  \xi_i\ (1\leq i\leq r),\quad \mathrm{T}_i\mathrm{T}_j\ (1\leq i\leq r,\ 1\leq j\leq r),\quad \mathrm{T}_i x\ (1\leq i\leq r,\ x\in\mathfrak{m}).
\]

Let $\mathfrak{c}_1=\sum_{i,j}\mathrm{T}_i\mathrm{T}_j\mathrm{A}[\mathrm{T}_1,\ldots,\mathrm{T}_r]$, $\mathfrak{c}_2=\mathfrak{c}_1+\sum_i\mathfrak{m}\mathrm{T}_i\mathrm{A}[\mathrm{T}_1,\ldots,\mathrm{T}_r]$, then $\mathfrak{c}=\mathfrak{c}_2+\sum_i\xi_i\mathrm{A}[\mathrm{T}_1,\ldots,\mathrm{T}_r]$.
It is clear that $\mathrm{A}[\mathrm{T}_1,\ldots,\mathrm{T}_r]/\mathfrak{c}_1$ is the $\mathrm{A}$-module direct sum of $\mathrm{A}$ and of $\mathrm{A}t_i^{(1)}$, where $t_i^{(1)}$ is the class of $\mathrm{T}_i$ mod.~$\mathfrak{c}_1$.
One deduces that $\mathrm{A}[\mathrm{T}_1,\ldots,\mathrm{T}_r]/\mathfrak{c}_2$ is the direct sum of $\mathrm{A}$ and of the $\mathrm{K}t_i^{(2)}$, where $t_i^{(2)}$ is the class of $\mathrm{T}_i$ mod.~$\mathfrak{c}_2$.
Finally $\mathrm{A}[\mathrm{T}_1,\ldots,\mathrm{T}_r]/\mathfrak{c}$ is the direct sum of $\mathrm{A}/\mathfrak{n}$ and of the $\mathrm{K}t_i'$, where $t_i'$ is the class of $\mathrm{T}_i$ mod.~$\mathfrak{c}$, and $\mathfrak{n}$ is the ideal of $\mathrm{A}$ generated by the $\xi_i$; in the identification of $\mathrm{A}'/\mathfrak{m}'^2$ with $\mathrm{A}[\mathrm{T}_1,\ldots,\mathrm{T}_r]/\mathfrak{c}$, the image of $\mathrm{V}$ identifies with $\mathfrak{m}/\mathfrak{n}$, and this proves (ii).
\end{proof}

\begin{proposition}[22.4.4]
\label{0.22.4.4}
With the hypotheses and notations of \sref{0.22.4.1}, suppose that $\mathrm{K}$ is of characteristic $p>0$ (hence $p.1\in\mathfrak{m}$ in $\mathrm{A}$) and that the $\mathrm{F}_i$ are of the form
\[
\label{0.22.4.4.1}
  \mathrm{F}_i(\mathrm{T}_i)=\mathrm{T}_i^p+p\sum_{k=1}^{p-1}c_{ik}\mathrm{T}_i^{p-k}-f_i
  \tag{22.4.4.1}
\]
with $c_{ik}\in\mathrm{A}$ and $f_i\in\mathrm{A}$ for every $i$ and every $k$ such that $1\leq k\leq p-1$; furthermore, if $a_i$ is the class of $f_i$ modulo $\mathfrak{m}$, suppose that the family $(a_i)_{1\leq i\leq r}$ is $p$-free over $\mathrm{K}^p$ (``absolutely $p$-free'').
Then:
\begin{enumerate}
  \item[\rm{(i)}] $\mathrm{A}'$ is a local ring, with maximal ideal $\mathfrak{m}\mathrm{A}'$, whose residue field $\mathrm{K}'$ is isomorphic to the extension of $\mathrm{K}$ obtained by adjunction of the $a_i^{1/p}$ $(1\leq i\leq r)$.

  \oldpage[0\textsubscript{IV-1}]{195}

  \item[\rm{(ii)}] The elements $d_\mathrm{A}f_i\otimes 1_{\mathrm{K}'}$ form a basis of the kernel $\mathrm{N}=\Upsilon^{\mathrm{K}'}_{\mathrm{A}'/\mathrm{A}}$ of the canonical homomorphism
  \[
    \Omega^1_\mathrm{A}\otimes_\mathrm{A}\mathrm{K}'
    \longrightarrow\Omega^1_{\mathrm{A}'}\otimes_{\mathrm{A}'}\mathrm{K}'.
  \]
\end{enumerate}
\end{proposition}

\begin{proof}
(i) This time $\mathrm{A}'\otimes_\mathrm{A}\mathrm{K}=\mathrm{A}'/\mathfrak{m}\mathrm{A}'$ is isomorphic to $\mathrm{K}[\mathrm{T}_1,\ldots,\mathrm{T}_r]/\mathfrak{b}'$, where $\mathfrak{b}'$ is the ideal generated by the polynomials $\mathrm{T}_i^p-a_i$ $(1\leq i\leq r)$; the hypothesis made on the $a_i$ implies that this quotient ring is a field \sref{0.21.1.8}, from which the assertions of (i) follow.

Before proving (ii), we shall establish the

\begin{lemma}[22.4.4.2]
\label{0.22.4.4.2}
Let $\Lambda$ be a ring, $\mathrm{A}$ a $\Lambda$-algebra, $\mathfrak{J}$ an ideal of $\mathrm{A}$ such that the ring $\mathrm{C}=\mathrm{A}/\mathfrak{J}$ is of characteristic $p>0$ (which is equivalent to saying that $p.1\in\mathfrak{J}$ in $\mathrm{A}$).
Let $\pi$ be the canonical image of $p.1$ in $\mathfrak{J}/\mathfrak{J}^2$.
If $\mathrm{C}$ is a $(\Lambda/p\Lambda)$-algebra formally smooth for the discrete topologies, one has a canonical exact sequence
\begin{equation}
\label{0.22.4.4.3}
  0\longrightarrow(\mathfrak{J}/\mathfrak{J}^2)/\mathrm{C}\cdot\pi
  \longrightarrow\Omega^1_{\mathrm{A}/\Lambda}\otimes_\mathrm{A}\mathrm{C}
  \longrightarrow\Omega^1_{\mathrm{C}/\Lambda}\longrightarrow0.
  \tag{22.4.4.3}
\end{equation}
\end{lemma}

\begin{proof}[Proof of the lemma]
Indeed, set $\mathrm{A}'=\mathrm{A}/p\mathrm{A}$, $\mathfrak{J}'=\mathfrak{J}/p\mathrm{A}$, so that $\mathrm{C}=\mathrm{A}'/\mathfrak{J}'$.
As $\mathrm{A}'=\mathrm{A}\otimes_\Lambda(\Lambda/p\Lambda)$, we have $\Omega^1_{\mathrm{A}/\Lambda}\otimes_\mathrm{A}\mathrm{A}'=\Omega^1_{\mathrm{A}'/\Lambda'}$ up to canonical isomorphism, where $\Lambda'$ denotes $\Lambda/p\Lambda$ \sref{0.20.5.5}.
Applying the exact sequence \sref{0.20.5.14.1} to the formally smooth $\Lambda'$-algebra $\mathrm{C}=\mathrm{A}'/\mathfrak{J}'$, one obtains the exact sequence
\[
  0\longrightarrow\mathfrak{J}'/\mathfrak{J}'^2
  \longrightarrow\Omega^1_{\mathrm{A}'/\Lambda'}\otimes_{\mathrm{A}'}\mathrm{C}
  \longrightarrow\Omega^1_{\mathrm{C}/\Lambda'}\longrightarrow0.
\]
Since the homomorphism $\Lambda\to\Lambda'$ is surjective, one has $\Omega^1_{\mathrm{C}/\Lambda'}=\Omega^1_{\mathrm{C}/\Lambda}$ up to canonical isomorphism \sref{0.20.5.15}; hence the exact sequence \sref{0.22.4.4.3} follows from the preceding observations.
\end{proof}

(ii) As $\mathfrak{m}'^2=0$, we shall also denote by $\mathrm{V}'$ the ideal $\mathfrak{m}'$ considered as a $\mathrm{K}'$-vector space, and we shall set
\begin{equation}
\label{0.22.4.4.4}
  \mathrm{V}_0=\mathrm{V}/\mathrm{K}\cdot p=\mathfrak{m}/(\mathfrak{m}^2+p\mathrm{A}),\qquad
  \mathrm{V}_0'=\mathrm{V}'/\mathrm{K}'\cdot p=\mathfrak{m}'/(\mathfrak{m}'^2+p\mathrm{A}').
  \tag{22.4.4.4}
\end{equation}

Applying the exact sequence \sref{0.22.4.4.3} in the case where $\Lambda=\mathbf{Z}$, to the $\mathbf{Z}$-algebra $\mathrm{A}$ (resp.~$\mathrm{A}'$) and to its ideal $\mathfrak{m}$ (resp.~$\mathfrak{m}'$), it follows, since $\mathrm{K}$ (resp.~$\mathrm{K}'$) is separable over the prime field $\mathbf{Z}/p\mathbf{Z}$, hence a $(\mathbf{Z}/p\mathbf{Z})$-algebra formally smooth \sref{0.19.6.1}, that the following sequences are exact:
\begin{equation}
\label{0.22.4.4.5}
  0\longrightarrow\mathrm{V}_0\longrightarrow\Omega^1_\mathrm{A}\otimes_\mathrm{A}\mathrm{K}
  \longrightarrow\Omega^1_\mathrm{K}\longrightarrow0,\qquad
  0\longrightarrow\mathrm{V}_0'\longrightarrow\Omega^1_{\mathrm{A}'}\otimes_{\mathrm{A}'}\mathrm{K}'
  \longrightarrow\Omega^1_{\mathrm{K}'}\longrightarrow0.
  \tag{22.4.4.5}
\end{equation}

Consider then the diagram \sref{0.22.4.4.6} in which the two middle columns are the second exact sequence \sref{0.22.4.4.5} and the first exact sequence \sref{0.22.4.4.5} tensorized by $\mathrm{K}'$; as the first exact sequence \sref{0.22.4.4.5} is formed of $\mathrm{K}$-vector spaces, the two middle columns of \sref{0.22.4.4.6} are exact; moreover, the diagram formed by these columns and the horizontal arrows connecting them is commutative, by virtue of the proof of \sref{0.22.4.4.2} and the commutativity of the diagram \sref{0.20.5.11.3}.
The last line of \sref{0.22.4.4.6} is the exact sequence obtained from \sref{0.20.6.1.1} by replacing $\mathrm{A}$, $\mathrm{B}$, $\mathrm{C}$ by $\mathbf{Z}$, $\mathrm{K}$ and $\mathrm{K}'$; the middle line is obtained by tensorization with $\mathrm{K}'$ from the $\mathrm{K}$-module exact sequence deduced from \sref{0.20.6.1.1} by replacing $\mathrm{A}$, $\mathrm{B}$, $\mathrm{C}$ by $\mathbf{Z}$, $\mathrm{A}$, $\mathrm{K}$ respectively.
The diagram formed by these two lines and the vertical arrows connecting them is commutative by virtue of the commutativity of \sref{0.20.6.4.3}.
We note moreover that we are in the conditions of application of \sref{0.22.4.1}, (ii), so the upper line of the diagram \sref{0.22.4.4.6} is \emph{exact}; the extreme columns of the diagram

\oldpage[0\textsubscript{IV-1}]{196}

\begin{equation}
\label{0.22.4.4.6}
  \xymatrix{
    && 0 \ar[d] & 0 \ar[d] && \\
    & 0 \ar[r] & \mathrm{V}_0\otimes_\mathrm{K}\mathrm{K}' \ar[r] \ar[d] & \mathrm{V}_0' \ar[r] \ar[d] & 0 & \\
    0 \ar[r] & \mathrm{N} \ar[r] \ar[d] & \Omega^1_\mathrm{A}\otimes_\mathrm{A}\mathrm{K}' \ar[r] \ar[d] &
      \Omega^1_{\mathrm{A}'}\otimes_{\mathrm{A}'}\mathrm{K}' \ar[r] \ar[d] &
      \Omega^1_{\mathrm{A}'/\mathrm{A}}\otimes_{\mathrm{A}'}\mathrm{K}' \ar[r] \ar[d] & 0 \\
    0 \ar[r] & \Upsilon_{\mathrm{K}'/\mathrm{K}} \ar[r] \ar[d] &
      \Omega^1_\mathrm{K}\otimes_\mathrm{K}\mathrm{K}' \ar[r] \ar[d] & \Omega^1_{\mathrm{K}'} \ar[r] \ar[d] &
      \Omega^1_{\mathrm{K}'/\mathrm{K}} \ar[r] \ar[d] & 0 \\
    & 0 & 0 & 0 & 0 &
  }
  \tag{22.4.4.6}
\end{equation}
are therefore formed respectively of the \emph{kernels} and the \emph{cokernels} of the middle horizontal arrows, which finishes proving that the diagram is \emph{commutative}.
Furthermore:

\begin{lemma}[22.4.4.7]
\label{0.22.4.4.7}
The homomorphisms
\[
  \mathrm{N}\longrightarrow\Upsilon_{\mathrm{K}'/\mathrm{K}}\qquad\text{and}\qquad
  \Omega^1_{\mathrm{A}'/\mathrm{A}}\otimes_{\mathrm{A}'}\mathrm{K}'
  \longrightarrow\Omega^1_{\mathrm{K}'/\mathrm{K}}
\]
of diagram \sref{0.22.4.4.6} are bijective.
\end{lemma}

\begin{proof}
Indeed, the snake diagram (Bourbaki, \emph{Alg.~comm.}, chap.~I, \textsection~1, n\textsuperscript{o}~4, prop.~2) applied to the two middle columns of \sref{0.22.4.4.6} gives an exact sequence
\[
  0\longrightarrow\mathrm{N}\longrightarrow\Upsilon_{\mathrm{K}'/\mathrm{K}}\longrightarrow0
  \longrightarrow\Omega^1_{\mathrm{A}'/\mathrm{A}}\otimes_{\mathrm{A}'}\mathrm{K}'
  \longrightarrow\Omega^1_{\mathrm{K}'/\mathrm{K}}\longrightarrow0.
\]
But, if $t_i$ is the canonical image of $\mathrm{T}_i$ in $\mathrm{A}'$, the relation $\mathrm{F}_i(t_i)=0$, together with the fact that $p.1\in\mathfrak{m}\mathrm{A}'\subset\mathfrak{m}'$, shows immediately that $d_{\mathrm{A}'}f_i\in\mathfrak{m}'\Omega^1_{\mathrm{A}'}$, hence $d_{\mathrm{A}'}f_i\otimes 1_{\mathrm{K}'}=0$;
this proves that the $d_\mathrm{A}f_i\otimes 1_{\mathrm{K}'}$ belong to $\mathrm{N}$.
Moreover, their images in $\Upsilon_{\mathrm{K}'/\mathrm{K}}$ are the $d_\mathrm{K}a_i\otimes 1_{\mathrm{K}'}$, which form a basis of $\Upsilon_{\mathrm{K}'/\mathrm{K}}$ over $\mathrm{K}'$ \sref{0.21.4.7}; taking into account \sref{0.22.4.4.7}, this finishes proving \sref{0.22.4.4}.
\end{proof}
\end{proof}

\begin{env}[22.4.5]
\label{0.22.4.5}
Now consider two rings $\mathrm{A}$, $\mathrm{B}$ of \emph{characteristic} $p>0$, and two homomorphisms
\[
  \mathrm{A}\xrightarrow{i}\mathrm{B}\xrightarrow{j}\mathrm{A}
\]
\oldpage[0\textsubscript{IV-1}]{197}
satisfying the conditions of \sref{0.21.3.1}.
Let furthermore $\mathfrak{m}$ be an ideal \emph{of square zero} in $\mathrm{A}$; set
\[
  \mathrm{K}=\mathrm{A}/\mathfrak{m},\qquad \mathrm{B}_{(\mathrm{K})}=\mathrm{B}\otimes_\mathrm{A}\mathrm{K}=\mathrm{B}/\mathrm{B}i(\mathfrak{m})
\]
and let
\[
  \varphi:\mathrm{A}\to\mathrm{K},\qquad \psi:\mathrm{B}\to\mathrm{B}_{(\mathrm{K})}
\]
be the canonical homomorphisms.
Since $\varphi$ is surjective, $\Omega^1_{\mathrm{B}_{(\mathrm{K})}/\mathrm{A}}$ identifies canonically with $\Omega^1_{\mathrm{B}_{(\mathrm{K})}/\mathrm{K}}$ \sref{0.20.5.15}.
We note that, since $p\geq 2$, for every $x\in\mathfrak{m}$, one has $j(i(x))=x^p=0$, in other words $j(i(\mathfrak{m}))=0$; by passing to quotients, one deduces from $j$ a homomorphism
\[
  j':\mathrm{B}_{(\mathrm{K})}\to\mathrm{A}
\]
such that, if $i'=\psi\circ i:\mathrm{A}\to\mathrm{B}_{(\mathrm{K})}$, we have
\[
  j'(i'(x))=x^p \quad\text{and}\quad i'(j'(y))=y^p
  \quad\text{for}\quad x\in\mathrm{A}\quad\text{and}\quad y\in\mathrm{B}_{(\mathrm{K})}.
\]
Setting, following the notations of \sref{0.21.3.2}
\[
  \Theta_{\mathrm{B}_{(\mathrm{K})}/\mathrm{A}}=\Omega^1_{\mathrm{B}_{(\mathrm{K})}/\mathrm{A}}\otimes_{\mathrm{B}_{(\mathrm{K})}}\mathrm{A}_{[j']}=\Omega^1_{\mathrm{B}_{(\mathrm{K})}/\mathrm{K}}\otimes_{\mathrm{B}_{(\mathrm{K})}}\mathrm{A}_{[j']}
\]
one has therefore \sref{0.21.3.2.2} a canonical homomorphism
\[
\label{0.22.4.5.1}
  \pi_{\mathrm{B}_{(\mathrm{K})}/\mathrm{A}}:\Theta_{\mathrm{B}_{(\mathrm{K})}/\mathrm{A}}\to\Omega^1_\mathrm{A}.
  \tag{22.4.5.1}
\]

On the other hand, one deduces from $i$ and $j$, by passing to quotients (taking into account that $j(i(\mathfrak{m}))=0$) homomorphisms
\[
  \mathrm{K}\xrightarrow{i_0}\mathrm{B}_{(\mathrm{K})}\xrightarrow{j_0}\mathrm{K}
\]
so that one has the commutative diagram
\[
  \begin{tikzcd}
    \mathrm{A} \ar[r,"i"] \ar[d,"\varphi"'] & \mathrm{B} \ar[r,"j"] \ar[d,"\psi"] & \mathrm{A} \ar[d,"\varphi"]\\
    \mathrm{K} \ar[r,"i_0"'] & \mathrm{B}_{(\mathrm{K})} \ar[r,"j_0"'] \ar[ur,"j'"] & \mathrm{K}
  \end{tikzcd}
\]
It is clear furthermore that
\[
  j_0(i_0(s))=s^p \quad\text{and}\quad i_0(j_0(t))=t^p
  \quad\text{for}\quad s\in\mathrm{K}\quad\text{and}\quad t\in\mathrm{B}_{(\mathrm{K})}.
\]

One therefore defines in the same way
\[
  \Theta_{\mathrm{B}_{(\mathrm{K})}/\mathrm{K}}=\Omega^1_{\mathrm{B}_{(\mathrm{K})}/\mathrm{K}}\otimes_{\mathrm{B}_{(\mathrm{K})}}\mathrm{K}_{[j_0]}=\Theta_{\mathrm{B}_{(\mathrm{K})}/\mathrm{A}}\otimes_\mathrm{A}\mathrm{K}
\]
hence, tensorizing \sref{0.22.4.5.1} by $\mathrm{K}$, one obtains a \emph{canonical} $\mathrm{K}$-homomorphism
\[
\label{0.22.4.5.2}
  \pi'_{\mathrm{B}_{(\mathrm{K})}/\mathrm{K}}=\pi_{\mathrm{B}_{(\mathrm{K})}/\mathrm{A}}\otimes 1_\mathrm{K}:\Theta_{\mathrm{B}_{(\mathrm{K})}/\mathrm{K}}\to\Omega^1_\mathrm{A}\otimes_\mathrm{A}\mathrm{K}.
  \tag{22.4.5.2}
\]

\oldpage[0\textsubscript{IV-1}]{198}
It follows immediately from the definitions that the diagram
\[
  \begin{tikzcd}
    \Theta_{\mathrm{B}_{(\mathrm{K})}/\mathrm{K}} \ar[r,"\pi'_{\mathrm{B}_{(\mathrm{K})}/\mathrm{K}}"] \ar[dr,"\pi_{\mathrm{B}_{(\mathrm{K})}/\mathrm{K}}",swap] & \Omega^1_\mathrm{A}\otimes_\mathrm{A}\mathrm{K} \ar[d,"\varphi_{\mathrm{K}/\mathrm{A}}"]\\
    & \Omega^1_\mathrm{K}
  \end{tikzcd}
\]
is commutative, $\pi_{\mathrm{B}_{(\mathrm{K})}/\mathrm{K}}$ being the canonical homomorphism corresponding to $i_0$ and $j_0$ \sref{0.21.3.2.2}.
By \emph{restriction} to $\Xi_{\mathrm{B}_{(\mathrm{K})}/\mathrm{K}}=\Ker(\pi_{\mathrm{B}_{(\mathrm{K})}/\mathrm{K}})$, the homomorphism $\pi'_{\mathrm{B}_{(\mathrm{K})}/\mathrm{K}}$ therefore defines a canonical homomorphism
\[
\label{0.22.4.5.3}
  \Xi_{\mathrm{B}_{(\mathrm{K})}/\mathrm{K}}\to\Ker(\Omega^1_\mathrm{A}\otimes_\mathrm{A}\mathrm{K}\to\Omega^1_\mathrm{K})=\Upsilon_{\mathrm{K}/\mathrm{A}}.
  \tag{22.4.5.3}
\]
\end{env}

\begin{env}[22.4.6]
\label{0.22.4.6}
With the hypotheses and notations of \sref{0.22.4.5}, suppose in addition that $\mathfrak{m}$ be a \emph{maximal} ideal of square zero, hence $\mathrm{K}$ a \emph{field}, and let $\mathrm{V}$ denote the ideal $\mathfrak{m}$ considered as a $\mathrm{K}$-vector space.
As $\mathrm{K}$ is a formally smooth algebra over its prime field $\mathbf{F}_p$ \sref{0.21.5.2} (so that $\mathrm{A}$ is an algebra over $\mathbf{F}_p$, in the sense that $\Omega^1_\mathrm{A}=\Omega^1_{\mathrm{A}/\mathbf{F}_p}$), one deduces from \sref{0.20.5.14} that one has an \emph{exact sequence}
\[
\label{0.22.4.6.1}
  0\to\mathrm{V}\to\Omega^1_\mathrm{A}\otimes_\mathrm{A}\mathrm{K}\to\Omega^1_\mathrm{K}\to 0.
  \tag{22.4.6.1}
\]

One has therefore by \sref{0.22.4.5.3} a homomorphism denoted
\[
\label{0.22.4.6.2}
  \chi'_{\mathrm{B}/\mathrm{A}}:\Xi_{\mathrm{B}_{(\mathrm{K})}/\mathrm{K}}\to\mathrm{V}
  \tag{22.4.6.2}
\]
that we shall call the \emph{characteristic homomorphism} of the $\mathrm{A}$-algebra $\mathrm{B}$ (relative to $i$, $j$ and $\mathfrak{m}$).
It follows from the definitions that for $\chi'_{\mathrm{B}/\mathrm{A}}$ to be \emph{injective}, it is necessary and sufficient that $\pi'_{\mathrm{B}_{(\mathrm{K})}/\mathrm{K}}$ is injective, the kernel of the latter being obviously contained in $\Xi_{\mathrm{B}_{(\mathrm{K})}/\mathrm{K}}$.
\end{env}

\begin{proposition}[22.4.7]
\label{0.22.4.7}
Let $\mathrm{A}$ be a local Artinian ring whose maximal ideal $\mathfrak{m}$ is of square zero, $\mathrm{K}=\mathrm{A}/\mathfrak{m}$ its residue field, $\mathrm{V}$ the ideal $\mathfrak{m}$ considered as a $\mathrm{K}$-vector space; we suppose that $\mathrm{A}$ is of characteristic $p>0$; let
\[
\label{0.22.4.7.1}
  \mathrm{F}_i(\mathrm{T}_i)=\mathrm{T}_i^p-f_i\qquad\text{with}\quad f_i\in\mathrm{A}\quad(1\leq i\leq r)
  \tag{22.4.7.1}
\]
and let $\mathrm{B}$ be the quotient ring of $\mathrm{A}[\mathrm{T}_1,\ldots,\mathrm{T}_r]$ by the ideal generated by the $r$ polynomials $\mathrm{F}_i$.
Then:
\begin{enumerate}
  \item[\rm{(i)}] $\mathrm{B}$ is a local ring.
  \item[\rm{(ii)}] If $\mathfrak{m}'$ is the maximal ideal of $\mathrm{B}$, $\mathrm{K}'=\mathrm{B}/\mathfrak{m}'$ its residue field, $\mathrm{V}'$ the $\mathrm{B}$-module $\mathfrak{m}'/\mathfrak{m}'^2$ considered as a $\mathrm{K}'$-vector space, then, for $\rg_\mathrm{K}(\mathrm{V})=\rg_{\mathrm{K}'}(\mathrm{V}')$, it is necessary and sufficient that the characteristic homomorphism of $\mathrm{K}$-vector spaces
\[
\label{0.22.4.7.2}
  \chi'_{\mathrm{B}/\mathrm{A}}:\Xi_{\mathrm{B}_{(\mathrm{K})}/\mathrm{K}}\to\mathrm{V}
  \tag{22.4.7.2}
\]
  (cf.~\sref{0.22.4.6.2}) is injective.
\end{enumerate}
\end{proposition}

\begin{proof}
Let us again denote by $a_i$ $(1\leq i\leq r)$ the class of $f_i$ mod.~$\mathfrak{m}$; the $a_i$ are not necessarily $p$-independent over $\mathrm{K}^p$, but one can always suppose that the subfamily
\oldpage[0\textsubscript{IV-1}]{199}
$(a_i)_{s+1\leq i\leq r}$ is $p$-free over $\mathrm{K}^p$ and forms a $p$-base over $\mathrm{K}^p$ of the field $\mathrm{K}^p(a_1,\ldots,a_r)$ so that
\[
  a_i\in\mathrm{K}^p(a_{s+1},\ldots,a_r)\qquad\text{for}\quad 1\leq i\leq s.
\]

Let us denote by $\mathrm{A}''$ the quotient ring of $\mathrm{A}[\mathrm{T}_{s+1},\ldots,\mathrm{T}_r]$ by the ideal generated by the $\mathrm{F}_i$ of index $i\geq s+1$; then, by \sref{0.22.4.4}, $\mathrm{A}''$ is a local ring whose maximal ideal is $\mathfrak{m}''=\mathfrak{m}\mathrm{A}''$ and the residue field is $\mathrm{K}''=\mathrm{K}(a_{s+1}^{1/p},\ldots,a_r^{1/p})$.
One has thus $a_i\in\mathrm{K}''^p$ for $1\leq i\leq s$ and hence, in $\mathrm{A}''$, the element $f_i$, for $1\leq i\leq s$, can be written
\[
\label{0.22.4.7.3}
  f_i=g_i^p+h_i
  \tag{22.4.7.3}
\]
where $g_i\in\mathrm{A}''$ and $h_i\in\mathfrak{m}''$ (since $\mathfrak{m}''^2=0$ and $p\mathfrak{m}''=0$, it is immediate that $h_i$ is uniquely determined by these conditions).
Replacing $\mathrm{T}_i$ by $\mathrm{T}_i+g_i$, one sees thus that if one sets
\[
  \mathrm{G}_i(\mathrm{T}_i)=\mathrm{T}_i^p+\sum_{k=1}^{p-1}\binom{p}{k}g_i^{p-k}\mathrm{T}_i^k-h_i\qquad(1\leq i\leq s)
\]
$\mathrm{B}$ is the quotient ring of $\mathrm{A}''[\mathrm{T}_1,\ldots,\mathrm{T}_s]$ by the ideal generated by the $\mathrm{G}_i$; but, all the coefficients of $\mathrm{G}_i$ except the leading coefficient belong to $\mathfrak{m}''$, we are in the situation of \sref{0.22.4.3}, which proves (i) (which indeed follows directly from the fact that $\mathrm{B}^p\subset\mathrm{A}$, and by consequence $\mathrm{B}$ has only a single ideal above $\mathfrak{m}$), and furthermore the residue field $\mathrm{K}'$ is \emph{isomorphic} to $\mathrm{K}''$.

Now let us note that for \sref{0.22.4.7.2} to be injective, it is necessary and sufficient that
\[
  \pi'_{\mathrm{B}_{(\mathrm{K})}/\mathrm{K}}:\Omega^1_{\mathrm{B}_{(\mathrm{K})}/\mathrm{K}}\otimes_{\mathrm{B}_{(\mathrm{K})}}\mathrm{K}\to\Omega^1_\mathrm{A}\otimes_\mathrm{A}\mathrm{K}
\]
is injective (by \sref{0.22.4.6}).
But $\mathrm{B}_{(\mathrm{K})}=\mathrm{B}/\mathfrak{m}\mathrm{B}$ is the quotient algebra $\mathrm{C}=\mathrm{K}[\mathrm{T}_1,\ldots,\mathrm{T}_r]$ by the ideal generated by the polynomials $\overline{\mathrm{F}}_i=\mathrm{T}_i^p-a_i$, and since $d_{\mathrm{C}/\mathrm{K}}\overline{\mathrm{F}}_i=0$ it follows from \sref{0.20.5.13} that $\Omega^1_{\mathrm{B}_{(\mathrm{K})}/\mathrm{K}}$ is a free $\mathrm{B}_{(\mathrm{K})}$-module with base $d_{\mathrm{B}_{(\mathrm{K})}/\mathrm{K}}t_i$, where $t_i$ is the canonical image of $\mathrm{T}_i$.
But by definition \sref{0.22.4.5}, the canonical image under $j':\mathrm{B}_{(\mathrm{K})}\to\mathrm{A}$ of the class modulo $\mathfrak{m}\mathrm{B}$ of an element $x\in\mathrm{B}$ is the element $x^p\in\mathrm{A}$; one deduces immediately, since $t_i^p=f_i$, that the image of $d_{\mathrm{B}_{(\mathrm{K})}/\mathrm{K}}t_i\otimes 1$ under $\pi'_{\mathrm{B}_{(\mathrm{K})}/\mathrm{K}}$ is $d_\mathrm{A}f_i\otimes 1$ in $\Omega^1_\mathrm{A}\otimes_\mathrm{A}\mathrm{K}$; the condition that $\pi'_{\mathrm{B}_{(\mathrm{K})}/\mathrm{K}}$ is injective is thus equivalent to the fact that the $d_\mathrm{A}f_i\otimes 1_\mathrm{K}$ are \emph{linearly independent over $\mathrm{K}$}, or equivalently that their images $d_\mathrm{A}f_i\otimes 1_{\mathrm{K}'}$ are \emph{linearly independent over $\mathrm{K}'$} in $\Omega^1_\mathrm{A}\otimes_\mathrm{A}\mathrm{K}'$.

Now, if one applies \sref{0.22.4.4} to the $\mathrm{A}$-algebra $\mathrm{A}''$, one sees that the kernel of the canonical homomorphism $\Omega^1_\mathrm{A}\otimes_\mathrm{A}\mathrm{K}'\to\Omega^1_{\mathrm{A}''}\otimes_{\mathrm{A}''}\mathrm{K}'$ has as its \emph{basis} the $d_\mathrm{A}f_i\otimes 1_{\mathrm{K}'}$ for $s+1\leq i\leq r$.
The preceding condition thus also amounts to saying that the $d_{\mathrm{A}''}f_i\otimes 1_{\mathrm{K}'}$ for $1\leq i\leq s$ are linearly independent in $\Omega^1_{\mathrm{A}''}\otimes_{\mathrm{A}''}\mathrm{K}'$ (taking into account that $d_{\mathrm{A}''}f_i=d_{\mathrm{A}''}h_i$ by \sref{0.22.4.7.3});
on the other hand, if $\mathrm{V}''$ is the $\mathrm{K}'$-vector space $\mathfrak{m}''$, one has $\rg_{\mathrm{K}'}(\mathrm{V}'')=\rg_\mathrm{K}(\mathrm{V})$ \sref{0.22.4.1}, (ii), and the condition $\rg_{\mathrm{K}'}(\mathrm{V}')=\rg_\mathrm{K}(\mathrm{V})$ is therefore equivalent to $\rg_{\mathrm{K}'}(\mathrm{V}'')=\rg_{\mathrm{K}'}(\mathrm{V}')$.
But it follows from \sref{0.22.4.3} that this last relation is equivalent to the classes $h_i$
\oldpage[0\textsubscript{IV-1}]{200}
being linearly independent over $\mathrm{K}'$ in $\mathrm{V}''$.
Now, in $\Omega^1_{\mathrm{A}''}\otimes_{\mathrm{A}''}\mathrm{K}'$, the $d_{\mathrm{A}''}h_i\otimes 1_{\mathrm{K}'}$ are the canonical images \sref{0.20.5.11.2} of the $h_i$ under the injection $\mathrm{V}''\to\Omega^1_{\mathrm{A}''}\otimes_{\mathrm{A}''}\mathrm{K}'$ (cf.~\sref{0.22.4.6.1}), and this finishes the proof of \sref{0.22.4.7}.
\end{proof}

\begin{corollary}[22.4.7.4]
\label{0.22.4.7.4}
For $\rg_\mathrm{K}(\mathrm{V})=\rg_{\mathrm{K}'}(\mathrm{V}')$, it is necessary and sufficient that the elements $d_\mathrm{A}f_i\otimes 1_\mathrm{K}$ $(1\leq i\leq r)$ be linearly independent in the $\mathrm{K}$-vector space $\Omega^1_\mathrm{A}\otimes_\mathrm{A}\mathrm{K}$.
\end{corollary}

\begin{remark}[22.4.8]
\label{0.22.4.8}
One can define the homomorphisms \sref{0.22.4.5.2} and \sref{0.22.4.6.2} in slightly different conditions.
Let $\mathrm{A}$ be a ring, $\mathfrak{m}$ an ideal \emph{of square zero} in $\mathrm{A}$, $\mathrm{K}=\mathrm{A}/\mathfrak{m}$, $p$ a prime number such that $p.1\in\mathfrak{m}$ (in other words, $\mathrm{K}$ is of characteristic $p$, but not necessarily $\mathrm{A}$).
Let furthermore $\mathrm{B}$ be an $\mathrm{A}$-algebra that is an $\mathrm{A}$-module \emph{faithfully flat} (so that $\mathrm{A}\subset\mathrm{B}$), and suppose that we have
\[
\label{0.22.4.8.1}
  y^p\in\mathrm{A}+p\mathrm{B}\subset\mathrm{A}+\mathfrak{m}\mathrm{B}
  \tag{22.4.8.1}
\]
for every $y\in\mathrm{B}$.
If $y^p=x+pz$ with $x\in\mathrm{A}$, $z\in\mathrm{B}$, the \emph{class modulo} $(\mathrm{A}\cap p\mathrm{B})$ of $x$ is well determined, and since $\mathrm{A}\cap p\mathrm{B}=p\mathrm{A}$ as $\mathrm{B}$ is a faithfully flat $\mathrm{A}$-module, one has thus defined a canonical map $j:\mathrm{B}\to\mathrm{A}/p\mathrm{A}$, such that, for $x\in\mathrm{A}$, $j(x)$ is the class of $x^p$ mod.~$p\mathrm{A}$.
Furthermore, if $y'\in\mathrm{B}$ and $y'^p=x'+pz'$ with $x'\in\mathrm{A}$, $z'\in\mathrm{B}$, it is immediate that the elements $(y+y')^p-x-x'$ and $(yy')^p-xx'$ belong to $p\mathrm{B}$ by virtue of the fact that the binomial coefficients $\binom{p}{i}$ are multiples of $p$ for $1\leq i\leq p-1$.
The map $j$ is therefore a \emph{ring homomorphism}.
By composition, one deduces from $j$ a homomorphism $j':\mathrm{B}\xrightarrow{j}\mathrm{A}/p\mathrm{A}\to\mathrm{A}/\mathfrak{m}=\mathrm{K}$ and hence a homomorphism $j_0:\mathrm{B}_{(\mathrm{K})}\to\mathrm{K}$ such that $j'$ is the composite $\mathrm{B}\to\mathrm{B}_{(\mathrm{K})}=\mathrm{B}\otimes_\mathrm{A}\mathrm{K}\xrightarrow{j_0}\mathrm{K}$; moreover, if $i_0:\mathrm{K}\to\mathrm{B}_{(\mathrm{K})}$ is the canonical map, $i_0$ and $j_0$ satisfy the conditions of \sref{0.21.3.1} for the rings \emph{of characteristic} $p$, $\mathrm{K}$ and $\mathrm{B}_{(\mathrm{K})}$.
This being so, in order to define a $\mathrm{K}$-homomorphism $\pi'_{\mathrm{B}_{(\mathrm{K})}/\mathrm{K}}:\Omega^1_{\mathrm{B}_{(\mathrm{K})}/\mathrm{K}}\otimes_{\mathrm{B}_{(\mathrm{K})}}\mathrm{K}_{[j_0]}\to\Omega^1_\mathrm{A}\otimes_\mathrm{A}\mathrm{K}$, it suffices \sref{0.20.4.8} to define a $\mathrm{K}$-\emph{derivation}
\[
\label{0.22.4.8.2}
  \mathrm{D}_0:\mathrm{B}_{(\mathrm{K})}\to\Omega^1_\mathrm{A}\otimes_\mathrm{A}\mathrm{K}
  \tag{22.4.8.2}
\]
where the second member is considered as a $\mathrm{B}_{(\mathrm{K})}$-module by means of $j_0$.
For this, it suffices to define an $\mathrm{A}$-\emph{derivation}
\[
\label{0.22.4.8.3}
  \mathrm{D}:\mathrm{B}\to\Omega^1_\mathrm{A}\otimes_\mathrm{A}\mathrm{K}
  \tag{22.4.8.3}
\]
that \emph{vanishes on} $\mathfrak{m}\mathrm{B}$ (the second member being considered as a $\mathrm{B}$-module by means of $j'$).
Indeed, if $y^p=x+pz$ with $y\in\mathrm{B}$, $x\in\mathrm{A}$, $z\in\mathrm{B}$, the element $d_\mathrm{A}x\otimes 1_\mathrm{K}$ of $\Omega^1_\mathrm{A}\otimes_\mathrm{A}\mathrm{K}$ depends only on the \emph{class} $j(y)$ mod.~$p\mathrm{A}$ of $x$ (for since $d_\mathrm{A}(pt)=pd_\mathrm{A}(t)$, hence $d_\mathrm{A}(pt)\otimes 1=0$ in $\Omega^1_\mathrm{A}\otimes_\mathrm{A}\mathrm{K}=\Omega^1_\mathrm{A}/\mathfrak{m}\Omega^1_\mathrm{A}$, since $p.1\in\mathfrak{m}$).
One can therefore take
\[
  \mathrm{D}(y)=d_\mathrm{A}x\otimes 1_\mathrm{K}.
\]

The fact that $\mathrm{D}$ is a derivation follows easily from the fact that $j$ is a ring homomorphism; moreover, for $t\in\mathrm{A}$, one has $d_\mathrm{A}(t^px)=t^pd_\mathrm{A}(x)+pt^{p-1}xd_\mathrm{A}(t)$ and one deduces that $\mathrm{D}(ty)=t\mathrm{D}(y)$, hence $\mathrm{D}$ is an $\mathrm{A}$-derivation; finally, if in addition $t\in\mathfrak{m}$, one has
\oldpage[0\textsubscript{IV-1}]{201}
$\mathrm{D}(ty)=t\mathrm{D}(y)=0$ since $\Omega^1_\mathrm{A}\otimes_\mathrm{A}\mathrm{K}$ is annihilated by $\mathfrak{m}$.
We have thus defined the desired homomorphism \sref{0.22.4.5.2}, and it is immediate to verify that the composition
\[
  \Theta_{\mathrm{B}_{(\mathrm{K})}/\mathrm{K}}\xrightarrow{\pi'_{\mathrm{B}_{(\mathrm{K})}/\mathrm{K}}}\Omega^1_\mathrm{A}\otimes_\mathrm{A}\mathrm{K}\to\Omega^1_\mathrm{K}
\]
is again the canonical homomorphism $\pi_{\mathrm{B}_{(\mathrm{K})}/\mathrm{K}}$ (relative to $i_0$ and $j_0$).
One therefore also has a canonical homomorphism analogous to \sref{0.22.4.5.3}.

Now if $\mathrm{K}$ is a \emph{field}, $\mathrm{A}/p\mathrm{A}$ is an $\mathbf{F}_p$-algebra, and one has the exact sequence \sref{0.20.5.14}
\[
  0\to\mathrm{V}/p\mathrm{A}\to\Omega^1_{\mathrm{A}/p\mathrm{A}}\otimes_{\mathrm{A}/p\mathrm{A}}\mathrm{K}\to\Omega^1_\mathrm{K}\to 0
\]
and, on the other hand, one has the exact sequence \sref{0.20.5.12.1}
\[
  p\mathrm{A}\to\Omega^1_\mathrm{A}\otimes_\mathrm{A}(\mathrm{A}/p\mathrm{A})\to\Omega^1_{\mathrm{A}/p\mathrm{A}}\to 0
\]
which shows, tensorizing with $\mathrm{K}$, that $\Omega^1_\mathrm{A}\otimes_\mathrm{A}\mathrm{K}$ and $\Omega^1_{\mathrm{A}/p\mathrm{A}}\otimes_{\mathrm{A}/p\mathrm{A}}\mathrm{K}$ are isomorphic.
The analog of \sref{0.22.4.6.2} is therefore a homomorphism
\[
\label{0.22.4.8.4}
  \Xi_{\mathrm{B}_{(\mathrm{K})}/\mathrm{K}}\to\mathrm{V}/p\mathrm{A}=\mathfrak{m}/(\mathfrak{m}^2+p\mathrm{A}).
  \tag{22.4.8.4}
\]

That said, the \emph{criterion} \sref{0.22.4.7}, in which one replaces $\mathrm{V}$ and $\mathrm{V}'$ by $\mathrm{V}/p\mathrm{A}$ and $\mathrm{V}'/p\mathrm{B}$ respectively, \emph{remains valid assuming only that $\mathrm{K}$ is of characteristic $p>0$} (in other words, that $p.1\in\mathfrak{m}$ in $\mathrm{A}$), \emph{and that the $\mathrm{F}_i(\mathrm{T}_i)$ are of the more general form} \sref{0.22.4.4.1}: indeed, $\mathrm{B}$ is a free $\mathrm{A}$-module (hence faithfully flat), and it suffices to redo the argument of \sref{0.22.4.7}, replacing \sref{0.22.4.6.2} by \sref{0.22.4.8.4} and $\mathrm{V}''$ by $\mathrm{V}''/p\mathrm{A}''$.
\end{remark}

\subsection{Geometrically regular algebras and formally smooth algebras}
\label{subsection:0.22.5}

\begin{proposition}[22.5.1]
\label{0.22.5.1}
Let $\mathrm{A}$ be a regular local ring, $\mathrm{A}'$ a local ring containing $\mathrm{A}$ and being a finite $\mathrm{A}$-algebra, $\mathfrak{m}$ (resp.~$\mathfrak{m}'$) the maximal ideal of $\mathrm{A}$ (resp.~$\mathrm{A}'$), $\mathrm{K}=\mathrm{A}/\mathfrak{m}$ (resp.~$\mathrm{K}'=\mathrm{A}'/\mathfrak{m}'$) its residue field.
For $\mathrm{A}'$ to be a regular ring, it is necessary and sufficient that
\[
\label{0.22.5.1.1}
  \rg_\mathrm{K}(\mathfrak{m}/\mathfrak{m}^2)=\rg_{\mathrm{K}'}(\mathfrak{m}'/\mathfrak{m}'^2).
  \tag{22.5.1.1}
\]
\end{proposition}

\begin{proof}
Indeed, the first member of \sref{0.22.5.1.1} is $\dim(\mathrm{A})$ \sref{0.17.1.1} and since $\mathrm{A}'\supset\mathrm{A}$ and $\mathrm{A}'$ is a finite $\mathrm{A}$-algebra, $\dim(\mathrm{A}')=\dim(\mathrm{A})$ \sref{0.16.1.5}; the equality \sref{0.22.5.1.1} is therefore necessary and sufficient for $\mathrm{A}'$ to be regular \sref{0.17.1.1}.
\end{proof}

\begin{remark}[22.5.2]
\label{0.22.5.2}
\begin{enumerate}
  \item[\rm{(i)}] Let $\mathrm{A}_1=\mathrm{A}/\mathfrak{m}^2$, $\mathrm{A}_1'=\mathrm{A}'\otimes_\mathrm{A}\mathrm{A}_1=\mathrm{A}'/\mathfrak{m}^2\mathrm{A}'$; as $\mathfrak{m}^2\mathrm{A}'\subset\mathfrak{m}'^2$, $\mathfrak{m}_1'=\mathfrak{m}'/\mathfrak{m}^2\mathrm{A}'$ is the maximal ideal of $\mathrm{A}_1'$ and $\mathfrak{m}_1'/\mathfrak{m}_1'^2$ is a $\mathrm{K}'$-vector space isomorphic to $\mathfrak{m}'/\mathfrak{m}'^2$; the condition of regularity of $\mathrm{A}'$ depends therefore only on the structure of the $\mathrm{A}_1$-\emph{algebra} $\mathrm{A}_1'$, which will allow us below to apply the preliminary results of \sref{0.22.4} on algebras finite over Artinian rings.
  \item[\rm{(ii)}] It follows from the proof of \sref{0.22.5.1} and from \sref{0.16.2.6.1} that one has in all cases
\[
\label{0.22.5.2.1}
  \rg_\mathrm{K}(\mathfrak{m}/\mathfrak{m}^2)\leq\rg_{\mathrm{K}'}(\mathfrak{m}'/\mathfrak{m}'^2).
  \tag{22.5.2.1}
\]
  \oldpage[0\textsubscript{IV-1}]{202}
  \item[\rm{(iii)}] Suppose that $\mathrm{A}$ and $\mathrm{A}'$ verify the general hypotheses of \sref{0.22.4.1}.
  It is known \sref{0.19.8.10} that there exists a regular local ring $\mathrm{B}$, with maximal ideal $\mathfrak{n}$, such that $\mathrm{A}$ is isomorphic to $\mathrm{B}/\mathfrak{n}^2$; let us denote by $\mathrm{G}_i\in\mathrm{B}[\mathrm{T}_i]$ a monic polynomial whose canonical image in $\mathrm{A}[\mathrm{T}_i]$ is $\mathrm{F}_i$ $(1\leq i\leq r)$; then, if $\mathrm{B}'$ is the quotient ring of $\mathrm{B}[\mathrm{T}_1,\ldots,\mathrm{T}_r]$ by the ideal generated by the $\mathrm{G}_i$, it is clear that $\mathrm{B}'$ is a free $\mathrm{B}$-module of rank $d$ and $\mathrm{A}'=\mathrm{B}'\otimes_\mathrm{B}\mathrm{A}=\mathrm{B}'/\mathfrak{n}^2\mathrm{B}'$; as $\mathrm{A}'$ is assumed to be a local ring, so is $\mathrm{B}'$, and if $\mathfrak{n}'$ is the maximal ideal of $\mathrm{B}'$, $\mathrm{B}'/\mathfrak{n}'$ is isomorphic to $\mathrm{K}'$ and $\mathfrak{n}'/\mathfrak{n}'^2$ is isomorphic to $\mathfrak{m}'/\mathfrak{m}'^2$ as a $\mathrm{K}'$-vector space;
  the inequality \sref{0.22.5.2.1} shows thus that we have, under the hypotheses of \sref{0.22.4.1}
\[
\label{0.22.5.2.2}
  \rg_\mathrm{K}(\mathrm{V})\leq\rg_{\mathrm{K}'}(\mathfrak{m}'/\mathfrak{m}'^2).
  \tag{22.5.2.2}
\]
\end{enumerate}
\end{remark}

\begin{corollary}[22.5.3]
\label{0.22.5.3}
Let $\mathrm{A}$ be a regular local ring, $\mathfrak{m}$ its maximal ideal, $\mathrm{T}_i$ $(1\leq i\leq r)$ indeterminates, and for each $i$, let
\[
\label{0.22.5.3.1}
  \mathrm{F}_i(\mathrm{T}_i)=\mathrm{T}_i^{d_i}+\sum_{1\leq k\leq d_i}c_{ik}\mathrm{T}_i^{d_i-k}
  \tag{22.5.3.1}
\]
be a unitary polynomial of $\mathrm{A}[\mathrm{T}_i]$, with $d_i\geq 2$ for every $i$ and $c_{ik}\in\mathfrak{m}$.
Let $\mathrm{A}'$ be the quotient of the polynomial ring $\mathrm{A}[\mathrm{T}_1,\ldots,\mathrm{T}_r]$ by the ideal $\mathfrak{b}$ generated by the $r$ polynomials $\mathrm{F}_i$.
Then:
\begin{enumerate}
  \item[\rm{(i)}] $\mathrm{A}'$ is a local ring, and its residue field $\mathrm{K}'$ is isomorphic to the residue field $\mathrm{K}$ of $\mathrm{A}$.
  \item[\rm{(ii)}] For $\mathrm{A}'$ to be regular, it is necessary and sufficient that the classes $\xi_i$ mod.~$\mathfrak{m}^2$ of the elements $c_{i,d_i}$ $(1\leq i\leq r)$ be linearly independent over $\mathrm{K}$.
\end{enumerate}
\end{corollary}

\begin{proof}
Indeed, with the notations of \sref{0.22.5.2}, (i), if $\overline{\mathrm{F}}_i$ denotes the polynomial of $\mathrm{A}_1[\mathrm{T}_i]$ obtained by applying to the coefficients of $\mathrm{F}_i$ the homomorphism $\mathrm{A}\to\mathrm{A}_1$, $\mathrm{A}_1'$ is defined from $\mathrm{A}_1$ and the polynomials $\overline{\mathrm{F}}_i$ as in \sref{0.22.4.1}, and the conclusion follows from \sref{0.22.5.1} and \sref{0.22.4.3}.
\end{proof}

\begin{theorem}[22.5.4]
\label{0.22.5.4}
Let $\mathrm{A}$ be a regular local ring, $\mathrm{K}$ its residue field; suppose that $\mathrm{K}$ is of characteristic $p>0$; let
\[
\label{0.22.5.4.1}
  \mathrm{F}_i(\mathrm{T}_i)=\mathrm{T}_i^p+p\sum_{k=1}^{p-1}c_{ik}\mathrm{T}_i^{p-k}-f_i
  \tag{22.5.4.1}
\]
with $c_{ik}\in\mathrm{A}$ and $f_i\in\mathrm{A}$ $(1\leq i\leq r)$; let $\mathrm{B}$ be the quotient ring of $\mathrm{A}[\mathrm{T}_1,\ldots,\mathrm{T}_r]$ by the ideal generated by the $r$ polynomials $\mathrm{F}_i$.
Then $\mathrm{B}$ is a local ring, and the following conditions are equivalent:
\begin{enumerate}
  \item[a)] $\mathrm{B}$ is regular;
  \item[b)] the elements $d_\mathrm{A}f_i\otimes 1_\mathrm{K}$ $(1\leq i\leq r)$ are linearly independent in the $\mathrm{K}$-vector space $\Omega^1_\mathrm{A}\otimes_\mathrm{A}\mathrm{K}$;
  \item[c)] if $\mathrm{B}'=\mathrm{B}/\mathfrak{m}^2\mathrm{B}$ and if one sets $\Xi_{\mathrm{B}_{(\mathrm{K})}/\mathrm{K}}=\Xi_{\mathrm{B}'_{(\mathrm{K})}/\mathrm{K}}$ \sref{0.22.4.5}, the characteristic homomorphism $\Xi_{\mathrm{B}_{(\mathrm{K})}/\mathrm{K}}\to\mathfrak{m}/\mathfrak{m}^2$ \sref{0.22.4.6.2} is injective.
\end{enumerate}
\end{theorem}

\begin{proof}
Indeed, $\mathrm{B}'$ is defined again from $\mathrm{A}_1=\mathrm{A}/\mathfrak{m}^2$ and the polynomials $\overline{\mathrm{F}}_i$ obtained by applying the homomorphism $\mathrm{A}\to\mathrm{A}_1$ to the coefficients.
Taking into account that $\Omega^1_\mathrm{A}\otimes_\mathrm{A}\mathrm{K}$ is canonically isomorphic to $\Omega^1_{\mathrm{A}_1}\otimes_{\mathrm{A}_1}\mathrm{K}$ \sref{0.20.5.12}, (ii), the theorem follows from \sref{0.22.5.1} and \sref{0.22.4.7.4}, using remark \sref{0.22.4.8} when $\mathrm{A}_1$ is not of characteristic $p$.
\end{proof}

\oldpage[0\textsubscript{IV-1}]{203}
\begin{env}[22.5.5]
\label{0.22.5.5}
Let $k$ be a field of characteristic $p>0$, $\mathrm{A}$ a local ring, $k'$ a \emph{finite radicial} extension of $k$ such that $k'^p\subset k$, $\mathfrak{m}$ the maximal ideal of $\mathrm{A}$, $\mathrm{K}=\mathrm{A}/\mathfrak{m}$ its residue field; recall that $\mathrm{A}'=\mathrm{A}\otimes_kk'$ is a local ring whose residue field is the same as that of the local ring $\mathrm{K}\otimes_kk'$ (Bourbaki, \emph{Alg.~comm.}, chap.~V, \textsection~2, No.~3, lemma~4).
We note that one has $\mathrm{A}'_{(\mathrm{K})}=\mathrm{A}'\otimes_\mathrm{A}\mathrm{K}=\mathrm{K}\otimes_kk'$, and hence \sref{0.21.3.6}
\[
\label{0.22.5.5.1}
  \Theta_{\mathrm{A}'_{(\mathrm{K})}/\mathrm{K}}=\Theta_{k'/k}\otimes_k\mathrm{K}
  \tag{22.5.5.1}
\]
to a canonical isomorphism.
Furthermore, it is known \sref{0.21.4.7} that one has $\Xi_{k'/k}=0$, in other words the canonical homomorphism \sref{0.21.3.2.2}
\[
  \pi_{k'/k}:\Theta_{k'/k}\to\Omega^1_k
\]
is injective, so that one has a canonical \emph{injection}
\[
\label{0.22.5.5.2}
  \Theta_{\mathrm{A}'_{(\mathrm{K})}/\mathrm{K}}\to\Omega^1_k\otimes_k\mathrm{K}
  \tag{22.5.5.2}
\]
which can be described explicitly as follows (taking into account \sref{0.22.5.5.1}): for every $x'\in k'$, to the element $d_{\mathrm{A}'_{(\mathrm{K})}/\mathrm{K}}(x'\otimes 1_\mathrm{K})\otimes 1_\mathrm{K}$ of $\Omega^1_{\mathrm{A}'_{(\mathrm{K})}/\mathrm{K}}\otimes_{\mathrm{A}'_{(\mathrm{K})}}\mathrm{K}$, one associates the element $d_k(x'^p)\otimes 1_\mathrm{K}$ of $\Omega^1_k\otimes_k\mathrm{K}$.
\end{env}

\begin{lemma}[22.5.6]
\label{0.22.5.6}
With the notations of \sref{0.22.5.5}:
\begin{enumerate}
  \item[\rm{(i)}] The kernel of the canonical homomorphism $\pi_{\mathrm{A}'_{(\mathrm{K})}/\mathrm{K}}$ \sref{0.21.3.2.2} identifies by \sref{0.22.5.5.2} with $(\Theta_{k'/k}\otimes_k\mathrm{K})\cap\Upsilon_{\mathrm{K}/k}$.
  \item[\rm{(ii)}] Let $\mathrm{A}_1=\mathrm{A}/\mathfrak{m}^2$, $\mathrm{A}_1'=\mathrm{A}_1\otimes_kk'=\mathrm{A}'/\mathfrak{m}^2\mathrm{A}'$; then the characteristic homomorphism $\chi'_{\mathrm{A}_1'/\mathrm{A}_1}$ \sref{0.22.4.6.2} identifies with the restriction of the characteristic homomorphism $\chi_{\mathrm{A}/k}:\Upsilon_{\mathrm{K}/k}\to\mathrm{V}=\mathfrak{m}/\mathfrak{m}^2$ to $(\Theta_{k'/k}\otimes_k\mathrm{K})\cap\Upsilon_{\mathrm{K}/k}$ \sref{0.20.6.24}.
\end{enumerate}
\end{lemma}

\begin{proof}
(i) By virtue of \sref{0.20.6.9}, applied by replacing $\mathrm{A}$ with the prime field, $\mathrm{B}$ with $k$, $\mathrm{C}$ with $\mathrm{K}$ and $\mathrm{E}$ with $\mathrm{A}_1=\mathrm{A}/\mathfrak{m}^2$ (extension of $\mathrm{K}$ by $\mathrm{V}=\mathfrak{m}/\mathfrak{m}^2$), it suffices (cf.~\sref{0.22.5.5}) to verify that if $x'\in k'$ and if $z\in\mathrm{A}_1$ is an element whose class mod.~$\mathfrak{m}/\mathfrak{m}^2$ is $x'^p\in k$, then the image of $d_k(x'^p)\otimes 1_\mathrm{K}$ under the canonical homomorphism $\Omega^1_k\otimes_k\mathrm{K}\to\Omega^1_{\mathrm{A}_1}\otimes_{\mathrm{A}_1}\mathrm{K}$ is $d_{\mathrm{A}_1}z\otimes 1_\mathrm{K}$, which follows from the definitions.

(ii) As $\mathrm{A}'_{(\mathrm{K})}=\mathrm{K}\otimes_kk'=(\mathrm{A}_1')_{(\mathrm{K})}$ (by virtue of the relation $\mathrm{K}=\mathrm{A}_1/(\mathfrak{m}/\mathfrak{m}^2)$), $\Theta_{\mathrm{A}'_{(\mathrm{K})}/\mathrm{K}}$ identifies with $\Theta_{(\mathrm{A}_1')_{(\mathrm{K})}/\mathrm{K}}$, and the verification follows from the same calculation as in (i) and from the definition of $\pi'_{(\mathrm{A}_1')_{(\mathrm{K})}/\mathrm{K}}$.
\end{proof}

\begin{proposition}[22.5.7]
\label{0.22.5.7}
Let $k$ be a field of characteristic $p>0$, $\mathrm{A}$ a $k$-algebra that is a regular local ring, $\mathfrak{m}$ its maximal ideal, $\mathrm{K}=\mathrm{A}/\mathfrak{m}$ its residue field; let $\mathrm{V}=\mathfrak{m}/\mathfrak{m}^2$, considered as a $\mathrm{K}$-vector space.
Let $k'$ be a finite extension of $k$ such that $k'^p\subset k$, and $\mathrm{A}'=\mathrm{A}\otimes_kk'$; for $\mathrm{A}'$ to be a regular local ring, it is necessary and sufficient that the restriction of the characteristic homomorphism $\chi_{\mathrm{A}/k}:\Upsilon_{\mathrm{K}/k}\to\mathrm{V}$ \sref{0.20.6.24} to $(\Theta_{k'/k}\otimes_k\mathrm{K})\cap\Upsilon_{\mathrm{K}/k}$ is injective.
\end{proposition}

\begin{proof}
We must prove that the condition in the statement is equivalent to \sref{0.22.5.1.1}, denoting by $\mathfrak{m}'$ the maximal ideal of $\mathrm{A}'$ and by $\mathrm{K}'=\mathrm{A}'/\mathfrak{m}'$ its residue field.
With the notations of remark \sref{0.22.5.2}, (i), one has $\mathrm{A}_1'=\mathrm{A}_1\otimes_kk'$.
Now, let $(a_i)_{1\leq i\leq r}$ be a $p$-basis of $k'$ over $k$, so that $a_i^p=b_i\in k$, and $k'$ is isomorphic to the quotient of the polynomial ring
\oldpage[0\textsubscript{IV-1}]{204}
$k[\mathrm{T}_1,\ldots,\mathrm{T}_r]$ by the ideal generated by the polynomials $\mathrm{F}_i=\mathrm{T}_i^p-b_i$ $(1\leq i\leq r)$; it follows that $\mathrm{A}_1'$ is isomorphic to the quotient of the polynomial ring $\mathrm{A}_1[\mathrm{T}_1,\ldots,\mathrm{T}_r]$ by the ideal generated by the $\mathrm{F}_i$.
One can then apply the criterion \sref{0.22.4.7}, and the condition in the statement is therefore equivalent to $\chi'_{\mathrm{A}_1'/\mathrm{A}_1}$ being injective; the conclusion follows from \sref{0.22.5.6}.
\end{proof}

\begin{theorem}[22.5.8]
\label{0.22.5.8}
Let $k$ be a field, $p$ its characteristic exponent, $\mathrm{A}$ a local Noetherian $k$-algebra.
The following conditions are equivalent:
\begin{enumerate}
  \item[a)] $\mathrm{A}$ is a $k$-algebra formally smooth (for its preadic topology).
  \item[b)] $\mathrm{A}$ is geometrically regular over $k$.
  \item[b')] For every finite extension $k'$ of $k$ such that $k'^p\subset k$, $\mathrm{A}'=\mathrm{A}\otimes_kk'$ is a regular ring.
\end{enumerate}
\end{theorem}

\begin{proof}
Taking into account \sref{0.19.6.6}, one can restrict oneself to the case $p>1$.

The fact that $a$) implies $b$) is none other than \sref{0.19.6.5}, and $b$) trivially implies $b'$).
Let us show that $b'$) implies the conditions of \sref{0.22.2.2} are satisfied; it is evident for the first \sref{0.17.1.1} since $b'$) implies first of all that $\mathrm{A}$ is regular.
Next, let $(x_\alpha)_{\alpha\in\mathrm{I}}$ be a $p$-base of $k^{1/p}$ over $k$; we know that the elements $d_k(x_\alpha^p)$ form a base of the $k$-vector space $\Omega^1_k$ \sref{0.21.4.5}, hence the elements $d_k(x_\alpha^p)\otimes 1_\mathrm{K}$ form a base of the $\mathrm{K}$-vector space $\Omega^1_k\otimes_k\mathrm{K}$.
One concludes (cf.~\sref{0.22.5.5}) that when $k'$ ranges over the finite subextensions of $k^{1/p}$, the family of subspaces $\Theta_{k'/k}\otimes_k\mathrm{K}$ is filtering increasing and has as union $\Omega^1_k\otimes_k\mathrm{K}$.
It follows then from \sref{0.22.5.7} that the condition $b'$) implies that $\chi_{\mathrm{A}/k}$ is injective, which is none other than condition (ii) of \sref{0.22.2.2}.
\end{proof}

\begin{corollary}[22.5.9]
\label{0.22.5.9}
Let $k$ be a field, $\mathrm{A}$ a local Noetherian $k$-algebra that is formally smooth.
Then, for every prime ideal $\mathfrak{p}$ of $\mathrm{A}$, the local ring $\mathrm{A}_\mathfrak{p}$ is a $k$-algebra formally smooth (for the $\mathfrak{p}$-preadic topology).
\end{corollary}

\begin{proof}
Indeed, with the notations of \sref{0.22.5.8}, $b'$), it suffices to show that the local ring $\mathrm{A}_\mathfrak{p}\otimes_kk'$ is regular; but since this is a ring of fractions of the local ring $\mathrm{A}\otimes_kk'$ and the latter is regular by hypothesis, it is the same by \sref{0.17.3.6}.
\end{proof}

\begin{remark}[22.5.10]
\label{0.22.5.10}
\begin{enumerate}
  \item[\rm{(i)}] The conclusion of \sref{0.22.5.8} fails when, in condition $b'$), one restricts to \emph{monogenic} extensions $k'=k(x)$ of $k$ (with $x^p\in k$).
  It can indeed happen that for all these extensions $k'$, one has $\Upsilon_{\mathrm{K}/k}\cap(\Theta_{k'/k}\otimes_k\mathrm{K})=0$, although $\Upsilon_{\mathrm{K}/k}\neq 0$; this means that for every $x\in k^{1/p}$ such that $x\notin k$, one must have $d_k(x^p)\otimes 1_\mathrm{K}\notin\Upsilon_{\mathrm{K}/k}$, or equivalently $d_\mathrm{K}(x^p)\neq 0$, that is, $x^p\notin\mathrm{K}^p$; in other words, one must have $\mathrm{K}^p\cap k=k^p$, even though $\mathrm{K}$ is an \emph{inseparable} extension of $k$.
  Now, one easily constructs examples of such extensions: start from a perfect field $k_0$ of characteristic $p>0$, let $\mathrm{X}$, $\mathrm{Y}$, $\mathrm{Z}$ be three indeterminates and set $k=k_0(\mathrm{X},\mathrm{Y})$ and $\mathrm{K}=k(\mathrm{X}^{1/p}+\mathrm{Z}\mathrm{Y}^{1/p})$; one verifies easily that $\mathrm{K}$ satisfies the preceding conditions, hence $\Upsilon_{\mathrm{K}/k}\neq 0$.
  Applying now the lemma \sref{0.22.2.5.2} taking $\mathrm{V}=\mathrm{K}$ and for $h$ the zero homomorphism; $\mathrm{A}$ is then a valuation ring of discrete valuation having $\mathrm{K}$ as residue field, with $\chi_{\mathrm{A}/k}=0$, and is not a $k$-algebra formally smooth by \sref{0.22.2.2}; however $\mathrm{A}\otimes_kk'$ is a regular ring for every \emph{monogenic} subextension $k'\subset k^{1/p}$ of $k$.
  \oldpage[0\textsubscript{IV-1}]{205}
  \item[\rm{(ii)}] The corollary \sref{0.22.5.9} leads to the following question: if $\mathrm{A}$ is a Noetherian $k$-algebra, under what conditions is the set of prime ideals $\mathfrak{p}\in\Spec(\mathrm{A})$ such that $\mathrm{A}_\mathfrak{p}$ is a $k$-algebra formally smooth an \emph{open} set?
  We will address this later in certain particular cases.
\end{enumerate}
\end{remark}

\subsection{Zariski's Jacobian criterion}
\label{subsection:0.22.6}

\begin{theorem}[22.6.1]
\label{0.22.6.1}
\textnormal{(Jacobian criterion for formal smoothness.)}
Let $\mathrm{A}$, $\mathrm{B}$ be two topological rings, $u:\mathrm{A}\to\mathrm{B}$ a continuous homomorphism making $\mathrm{B}$ a formally smooth $\mathrm{A}$-algebra, $\mathfrak{J}$ an ideal of $\mathrm{B}$ (not necessarily closed), $\mathrm{C}=\mathrm{B}/\mathfrak{J}$ the quotient topological algebra.
For $\mathrm{C}$ to be a formally smooth $\mathrm{A}$-algebra, it is necessary and sufficient that the canonical homomorphism (cf.~\sref{0.20.5.11.2})
\[
\label{0.22.6.1.1}
  \delta_{\mathrm{C}/\mathrm{B}/\mathrm{A}}:\mathfrak{J}/\mathfrak{J}^2\to\Omega^1_{\mathrm{B}/\mathrm{A}}\otimes_\mathrm{B}\mathrm{C}
  \tag{22.6.1.1}
\]
be \emph{formally invertible on the left} (cf.~\sref{0.19.1.5}).
\end{theorem}

\begin{proof}
Indeed, saying that $\mathrm{B}$ (resp.~$\mathrm{C}$) is a formally smooth $\mathrm{A}$-algebra signifies \sref{0.19.4.4} that $\operatorname{Exalcotop}_\mathrm{A}(\mathrm{B},\mathrm{L})=0$ (resp.~$\operatorname{Exalcotop}_\mathrm{A}(\mathrm{C},\mathrm{L})=0$) for every $\mathrm{B}$-module (resp.~$\mathrm{C}$-module) discrete $\mathrm{L}$ annihilated by an open ideal of $\mathrm{B}$ (resp.~$\mathrm{C}$).
By hypothesis $\operatorname{Exalcotop}_\mathrm{A}(\mathrm{B},\mathrm{L})=0$ for every such discrete $\mathrm{C}$-module $\mathrm{L}$ annihilated by an open ideal of $\mathrm{C}$, saying that $\mathrm{C}$ is a formally smooth $\mathrm{A}$-algebra thus signifies that the canonical homomorphism $\operatorname{Exalcotop}_\mathrm{A}(\mathrm{C},\mathrm{L})\to\operatorname{Exalcotop}_\mathrm{A}(\mathrm{B},\mathrm{L})$ is \emph{injective}, and the theorem follows from \sref{0.20.7.8}.
\end{proof}

Recall \sref{0.19.5.3} that when $\mathrm{B}$ and $\mathrm{C}$ are formally smooth $\mathrm{A}$-algebras, $\mathfrak{J}/\mathfrak{J}^2$ is a $\mathrm{C}$-module topologically \emph{formally projective}; moreover, the canonical homomorphisms $\varphi_n:\mathbf{S}^n_\mathrm{C}(\mathfrak{J}/\mathfrak{J}^2)\to\mathfrak{J}^n/\mathfrak{J}^{n+1}$ are \emph{formal bimorphisms} when $\mathrm{B}$ is assumed to be a preadmissible ring.

\begin{corollary}[22.6.2]
\label{0.22.6.2}
With the hypotheses and notations of \sref{0.22.6.1}, suppose furthermore
that the square of every open ideal of $\mathrm{B}$ is open, and that on
$\mathfrak{J}$ the topology induced by that of $\mathrm{B}$ is identical to
the topology deduced from that of $\mathrm{B}$ \sref{0.19.0} (note that these
two conditions are satisfied when $\mathrm{B}$ is Noetherian and its topology
is preadic $(\mathbf{0}_\mathrm{I},~7.3.2)$, or if the topology of $\mathrm{B}$
is the $\mathfrak{J}$-preadic topology). Let $(\mathfrak{b}_\lambda)$ be a
fundamental system of open ideals of $\mathrm{B}$ and set
$\mathrm{B}_\lambda=\mathrm{B}/\mathfrak{b}_\lambda$ for every $\lambda$.
Then:
\begin{enumerate}
  \item[\rm{(i)}] The following conditions are equivalent:
  \begin{enumerate}
    \item[a)] $\mathrm{C}$ is a formally smooth $\mathrm{A}$-algebra.
    \item[b)] For every $\lambda$, the homomorphism of
    $\mathrm{B}_\lambda$-modules
    \[
    \label{0.22.6.2.1}
      (\mathfrak{J}/\mathfrak{J}^2)\otimes_\mathrm{B}\mathrm{B}_\lambda
      \longrightarrow
      \Omega^1_{\mathrm{B}/\mathrm{A}}\otimes_\mathrm{B}\mathrm{B}_\lambda
      \tag{22.6.2.1}
    \]
    deduced from $\delta_{\mathrm{C}/\mathrm{B}/\mathrm{A}}$ by tensorization,
    is invertible on the left (in other words, an isomorphism onto a direct
    factor of
    $\Omega^1_{\mathrm{B}/\mathrm{A}}\otimes_\mathrm{B}\mathrm{B}_\lambda$).
  \end{enumerate}
  \item[\rm{(ii)}] Suppose furthermore that the topological ring $\mathrm{C}$
  is preadmissible $(\mathbf{0}_\mathrm{I},~7.1.2)$ and

  \oldpage[0\textsubscript{IV-1}]{206}

  let $\mathfrak{L}$ be an ideal of definition of $\mathrm{C}$; then
  conditions $a$) and $b$) of (i) are also equivalent to:
  \begin{enumerate}
    \item[c)] The homomorphism of $(\mathrm{C}/\mathfrak{L})$-modules
    \[
      (\mathfrak{J}/\mathfrak{J}^2)\otimes_\mathrm{C}
      (\mathrm{C}/\mathfrak{L})\longrightarrow
      \Omega^1_{\mathrm{B}/\mathrm{A}}\otimes_\mathrm{B}
      (\mathrm{C}/\mathfrak{L})
    \]
    is invertible on the left.
  \end{enumerate}
\end{enumerate}
\end{corollary}

\begin{proof}
The hypotheses imply that on $\Omega^1_{\mathrm{B}/\mathrm{A}}$ the topology
is deduced from that of $\mathrm{B}$ \sref{0.20.4.5}, hence (i) follows from
\sref{0.22.6.1} and \sref{0.19.1.7}. On the other hand, recall
\sref{0.20.4.9} that since $\mathrm{B}$ is a formally smooth
$\mathrm{A}$-algebra, $\Omega^1_{\mathrm{B}/\mathrm{A}}$ is a formally
projective $\mathrm{B}$-module, hence
$\Omega^1_{\mathrm{B}/\mathrm{A}}\otimes_\mathrm{B}\mathrm{B}_\lambda$ is a
projective $\mathrm{B}_\lambda$-module for every $\lambda$
\sref{0.19.2.4}; the equivalence of $c$) with $a$) and $b$), when $\mathrm{C}$
is preadmissible, then follows from \sref{0.19.1.9}.
\end{proof}

In particular:
\begin{corollary}[22.6.3]
\label{0.22.6.3}
Let $\mathrm{A}$ be a ring, $\mathrm{B}$ an $\mathrm{A}$-algebra, $\mathfrak{J}$ an ideal of $\mathrm{B}$, $\mathrm{C}=\mathrm{B}/\mathfrak{J}$; we suppose $\mathrm{A}$, $\mathrm{B}$, $\mathrm{C}$ equipped with the discrete topologies and that $\mathrm{B}$ is a formally smooth $\mathrm{A}$-algebra.
For $\mathrm{C}$ to be a formally smooth $\mathrm{A}$-algebra, it is necessary and sufficient that the canonical homomorphism \sref{0.22.6.1.1} be invertible on the left.
\end{corollary}

Note that since \emph{every} $\mathrm{A}$-algebra $\mathrm{C}$ is a quotient of a polynomial algebra $\mathrm{A}[\mathrm{T}_\alpha]_{\alpha\in\mathrm{I}}$, and the latter is formally smooth \sref{0.19.3.3}, the criterion \sref{0.22.6.3} allows in principle to determine whether $\mathrm{C}$ is a formally smooth $\mathrm{A}$-algebra.

\begin{proposition}[22.6.4]
\label{0.22.6.4}
Let $\mathrm{A}$ be a ring, $\mathrm{B}$ a formally smooth $\mathrm{A}$-algebra (for the discrete topologies), $\mathfrak{J}$ an ideal of $\mathrm{B}$, $\mathrm{C}=\mathrm{B}/\mathfrak{J}$; suppose that $\mathfrak{J}/\mathfrak{J}^2$ is a $\mathrm{C}$-module of finite type.
Let $\mathfrak{p}$ be a prime ideal of $\mathrm{C}$,
$k(\mathfrak{p})$ the residue field of $\mathrm{C}_\mathfrak{p}$,
$\mathfrak{p}'$ the inverse image of $\mathfrak{p}$ in $\mathrm{B}$, and
$\mathfrak{q}$ the prime ideal of $\mathrm{A}$ which is the inverse image of
$\mathfrak{p}'$.
The following conditions are equivalent:
\begin{enumerate}
  \item[a)] $\mathrm{C}_\mathfrak{p}$ is a formally smooth $\mathrm{A}_{\mathfrak{q}}$-algebra (for the discrete topologies).
  \item[a')] $\mathrm{C}_\mathfrak{p}$ is a formally smooth $\mathrm{A}$-algebra (for the discrete topologies).
  \item[b)] The canonical homomorphism
\[
\label{0.22.6.4.1}
  (\mathfrak{J}/\mathfrak{J}^2)\otimes_\mathrm{C}k(\mathfrak{p})\to\Omega^1_{\mathrm{B}/\mathrm{A}}\otimes_\mathrm{B}k(\mathfrak{p})
  \tag{22.6.4.1}
\]
  is injective.
  \item[c)] There exists $f\in\mathrm{C}-\mathfrak{p}$ such that
  $\mathrm{C}_f$ is a formally smooth $\mathrm{A}$-algebra.
\end{enumerate}

If furthermore $\mathrm{B}$ is \emph{Noetherian}, the preceding conditions are also equivalent to:
\begin{enumerate}
  \item[d)] $\mathrm{C}_\mathfrak{p}$ is a formally smooth
  $\mathrm{A}_{\mathfrak{q}}$-algebra for the $\mathfrak{p}$-preadic topology
  on $\mathrm{C}_\mathfrak{p}$ and the discrete topology (or the
  $\mathfrak{q}$-preadic topology) on $\mathrm{A}_{\mathfrak{q}}$.
\end{enumerate}
\end{proposition}

\begin{proof}
Since $\mathrm{A}_{\mathfrak{q}}$ is a formally smooth
$\mathrm{A}$-algebra \sref{0.19.3.5}, (iv) and (ii), conditions $a$) and
$a'$) are already known to be equivalent \sref{0.19.3.5}, (ii). We have
$\mathrm{C}_\mathfrak{p}=\mathrm{B}_{\mathfrak{p}'}/
\mathfrak{J}\mathrm{B}_{\mathfrak{p}'}$, and
$\mathrm{B}_{\mathfrak{p}'}$ is a formally smooth $\mathrm{A}$-algebra for
the discrete topology \sref{0.19.3.5}, (iv); moreover,
$\Omega^1_{\mathrm{B}_{\mathfrak{p}'}/\mathrm{A}}
=(\Omega^1_{\mathrm{B}/\mathrm{A}})_{\mathfrak{p}'}$ \sref{0.20.5.9}.
The equivalence of $a'$) and $b$) therefore follows by applying
\sref{0.22.6.3} to $\mathrm{B}_{\mathfrak{p}'}$ and
$\mathrm{C}_\mathfrak{p}$, taking into account that
$\Omega^1_{\mathrm{B}/\mathrm{A}}$ is a projective $\mathrm{B}$-module
\sref{0.20.4.9}, \sref{0.19.2.11}, that
$\mathfrak{J}/\mathfrak{J}^2$ is a $\mathrm{B}$-module of finite type, and
using \sref{0.19.1.12}. Applying \sref{0.19.1.12} also proves the equivalence
of $b$) and $c$). Finally, since $\mathrm{B}_{\mathfrak{p}'}$ is a formally
smooth $\mathrm{A}_{\mathfrak{q}}$-algebra for the

\oldpage[0\textsubscript{IV-1}]{207}

discrete topologies \sref{0.19.3.5}, (iv), it remains formally smooth when
$\mathrm{B}_{\mathfrak{p}'}$ is given its $\mathfrak{p}'$-preadic topology
and $\mathrm{A}_{\mathfrak{q}}$ is given either the discrete or the
$\mathfrak{q}$-preadic topology \sref{0.19.3.8}. To prove the equivalence of
$b$) and $d$) when $\mathrm{B}$ is Noetherian, apply \sref{0.22.6.2} to
$\mathrm{A}_{\mathfrak{q}}$ with the discrete (or $\mathfrak{q}$-preadic)
topology and to $\mathrm{B}_{\mathfrak{p}'}$ with the
$\mathfrak{p}'$-preadic topology. Since $\mathrm{C}_\mathfrak{p}$ is
preadmissible and $\mathfrak{p}\mathrm{C}_\mathfrak{p}$ is an ideal of
definition for its topology, the equivalence of $c$) and $a$) in
\sref{0.22.6.2} applies.
\end{proof}

\begin{corollary}[22.6.5]
\label{0.22.6.5}
Under the general hypotheses of \sref{0.22.6.4}, the set of
$\mathfrak{p}\in\Spec(\mathrm{C})$ such that $\mathrm{C}_\mathfrak{p}$ is a
formally smooth $\mathrm{A}$-algebra (or a formally smooth
$\mathrm{A}_{\mathfrak{q}}$-algebra, where $\mathfrak{q}$ denotes the inverse
image of $\mathfrak{p}$ in $\mathrm{A}$), for the discrete topologies, is open
in $\Spec(\mathrm{C})$.
\end{corollary}

\begin{proof}
This follows from form $c$) of \sref{0.22.6.4}.
Note that when $\mathrm{B}$ is Noetherian, one can replace the discrete topologies by the preadic topologies.
\end{proof}

\begin{corollary}[22.6.6]
\label{0.22.6.6}
Under the general hypotheses of \sref{0.22.6.4}, the following conditions are equivalent:
\begin{enumerate}
  \item[a)] $\mathrm{C}$ is a formally smooth $\mathrm{A}$-algebra (for the discrete topologies).
  \item[b)] For every $\mathfrak{p}\in\Spec(\mathrm{C})$ (or only for every \emph{maximal} ideal $\mathfrak{p}$ of $\mathrm{C}$), $\mathrm{C}_\mathfrak{p}$ is a formally smooth $\mathrm{A}$-algebra (or a formally smooth $\mathrm{A}_{\mathfrak{q}}$-algebra, where $\mathfrak{q}$ denotes the inverse image of $\mathfrak{p}$ in $\mathrm{A}$) \emph{for the discrete topologies}.
\end{enumerate}
Furthermore, when $\mathrm{B}$ is Noetherian, one can replace in $b$) the discrete topologies by the preadic topologies on $\mathrm{A}_{\mathfrak{q}}$ and $\mathrm{C}_\mathfrak{p}$.
\end{corollary}

\begin{proof}
The last assertion follows from \sref{0.22.6.4}.
By virtue of \sref{0.22.6.3}, condition $a$) is equivalent to the fact that
the homomorphism \sref{0.22.6.1.1} is invertible on the left; the equivalence
of $a$) and $b$) therefore follows from \sref{0.19.1.14}, taking into account
that $\Omega^1_{\mathrm{B}/\mathrm{A}}$ is a projective $\mathrm{B}$-module
and $\mathfrak{J}/\mathfrak{J}^2$ a $\mathrm{C}$-module of finite type.
\end{proof}

\begin{proposition}[22.6.7]
\label{0.22.6.7}
Let $k_0$ be a field, $k$ a separable extension of $k_0$, $\mathrm{C}$ a $k$-algebra of finite type.
\begin{enumerate}
  \item[\rm{(i)}] Let $\mathfrak{p}$ be a prime ideal of $\mathrm{C}$.
  The following conditions are equivalent:
  \begin{enumerate}
    \item[a)] $\mathrm{C}_\mathfrak{p}$ is a $k_0$-algebra formally smooth for the discrete topology.
    \item[b)] $\mathrm{C}_\mathfrak{p}$ is a $k_0$-algebra formally smooth for the $\mathfrak{p}$-preadic topology.
    \item[c)] There exists $f\in\mathrm{C}-\mathfrak{p}$ such that $\mathrm{C}_f$ is a $k_0$-algebra formally smooth for the discrete topology.
    \item[d)] $\mathrm{C}_\mathfrak{p}$ is a geometrically regular ring over $k_0$ \sref{0.19.6.5}.
  \end{enumerate}
  Furthermore the set of $\mathfrak{p}\in\Spec(\mathrm{C})$ having any of these properties is open in $\Spec(\mathrm{C})$.
  \item[\rm{(ii)}] The following conditions are equivalent:
  \begin{enumerate}
    \item[a)] $\mathrm{C}$ is a $k_0$-algebra formally smooth for the discrete topology.
    \item[b)] Every $\mathfrak{p}\in\Spec(\mathrm{C})$ satisfies the equivalent conditions of (i).
    \item[c)] Every maximal ideal $\mathfrak{m}$ of $\mathrm{C}$ satisfies the equivalent conditions of (i).
    \item[d)] For every extension $k'$ of $k_0$, $\mathrm{C}\otimes_{k_0}k'$ is a regular ring \sref{0.17.3.6}; we also say that $\mathrm{C}$ is a geometrically regular ring over $k_0$.
  \end{enumerate}

  \oldpage[0\textsubscript{IV-1}]{208}

  \item[\rm{(iii)}] Let
  $\mathrm{B}=k[\mathrm{T}_1,\ldots,\mathrm{T}_r]$ be a polynomial algebra and
  $\mathfrak{J}$ an ideal of $\mathrm{B}$ such that $\mathrm{C}$ is
  isomorphic to $\mathrm{B}/\mathfrak{J}$. Let $\mathfrak{q}$ be a prime
  ideal of $\mathrm{B}$ containing $\mathfrak{J}$ and set
  $\mathfrak{p}=\mathfrak{q}/\mathfrak{J}$. The conditions of (i) are then
  also equivalent to the following (the ``Jacobian criterion of Zariski''):
  \begin{enumerate}
    \item[e)] There exists a system $(f_i)_{1\leq i\leq m}$ of elements of
    $\mathfrak{J}$ generating the ideal $\mathfrak{J}_\mathfrak{q}$ in
    $\mathrm{B}_\mathfrak{q}$, and $k_0$-derivations $\mathrm{D}_j$ of
    $\mathrm{B}$ into itself $(1\leq j\leq m)$, such that
    $\det(\mathrm{D}_j(f_i))\notin\mathfrak{q}$.
  \end{enumerate}
\end{enumerate}
\end{proposition}

\begin{proof}
One knows that $\mathrm{C}$ is always isomorphic to a $k$-algebra of the form
$\mathrm{B}/\mathfrak{J}$. Since $\mathrm{B}$ is a formally smooth
$k$-algebra \sref{0.19.3.3}, and $k$, being separable over $k_0$, is a
formally smooth $k_0$-algebra \sref{0.19.6.1}, $\mathrm{B}$ is a formally
smooth $k_0$-algebra \sref{0.19.3.5}, (ii). The equivalence of conditions
$a$), $b$), $c$) of (i) therefore follows from \sref{0.22.6.4}, and that of
$a$), $b$), $c$) in (ii) follows from \sref{0.22.6.6}; the equivalence of
$a$) and $d$) in (i) follows from \sref{0.22.5.8}. Since every localization
of $\mathrm{C}\otimes_{k_0}k'$ is also a localization of
$\mathrm{C}_\mathfrak{p}\otimes_{k_0}k'$ for a suitable prime ideal
$\mathfrak{p}\in\Spec(\mathrm{C})$, the equivalence of $d$) and $b$) in (ii)
follows from the equivalence of $d$) and $a$) in (i). Finally, since
$\Omega^1_{\mathrm{B}/k_0}$ is a projective $\mathrm{B}$-module, the
equivalence of the Zariski criterion $e$) and the other conditions in (i)
follows from \sref{0.19.1.12} and \sref{0.22.6.3}, taking
\sref{0.20.4.8} into account.
\end{proof}

Zariski was in fact interested in a differential criterion of regularity for
the local rings $\mathrm{C}_\mathfrak{p}$. Since saying that a Noetherian
local ring containing a field is regular is equivalent to saying that it is
formally smooth as an algebra over its prime subfield \sref{0.19.6.4}, such a
criterion follows immediately by replacing $k_0$ by the prime subfield of
$k$ in \sref{0.22.6.7}. In particular, one obtains the following result,
which will be recovered later in (IV,~6.12.5) as a special case of more
general results of Nagata:

\begin{corollary}[22.6.8]
\label{0.22.6.8}
\textnormal{(Zariski.)}
Let $\mathrm{C}$ be an algebra of finite type over a field.
Then the set of $\mathfrak{p}\in\Spec(\mathrm{C})$ such that $\mathrm{C}_\mathfrak{p}$ is a regular local ring is open in $\Spec(\mathrm{C})$.
\end{corollary}

\begin{remark}[22.6.9]
\label{0.22.6.9}
\begin{enumerate}
  \item[\rm{(i)}] The statements of \sref{0.22.6.7} are still valid if instead of supposing $k$ separable over $k_0$, one assumes only that it is of \emph{finite radiciel multiplicity}, that is to say \sref{0.19.6.6} that there exists a finite radiciel extension $k_0'$ of $k_0$ such that the residue field of the local Artinian ring $k\otimes_{k_0}k_0'$ is a \emph{separable} extension $k'$ of $k_0'$.
  There exists then a $k_0'$-monomorphism
  $k'\to k\otimes_{k_0}k_0'$ which, composed with the canonical homomorphism
  $k\otimes_{k_0}k_0'\to k'$, gives the identity, since $k'$ is a formally
  smooth $k_0'$-algebra \sref{0.19.6.1}. One concludes that
  $\mathrm{C}\otimes_{k_0}k_0'$ is equipped with the structure of a
  $k'$-algebra of finite type, and since $k'$ is separable over $k_0'$, one
  can apply \sref{0.22.6.7} to this $k'$-algebra; the assertion follows by
  applying \sref{0.19.4.7}. Since we shall not need this generalization, we
  leave the details to the reader.
  We do not know whether the results of \sref{0.22.6.7} are valid without \emph{any} hypothesis on the extension $k$ of $k_0$.
  \item[\rm{(ii)}] Under the general hypotheses of \sref{0.22.6.7}, let $(k_\alpha)$ be a decreasing filtering family of subfields of $k$ containing $k_0$ and such that $\bigcap k_\alpha(k^p)=k_0(k^p)$ (where $p$ is the characteristic exponent of $k_0$).
  Using a dimension counting method due to Nagata, one can show that, in the
  Jacobian criterion \sref{0.22.6.7}, (iii), one may restrict the derivations
  $\mathrm{D}_j$ of $\mathrm{B}$ to be $k_\alpha$-derivations for a suitable
  $\alpha$.

  \oldpage[0\textsubscript{IV-1}]{209}

  The interest of this result is that there always exist such families
  $(k_\alpha)$ for which $[k:k_\alpha]$ is finite \sref{0.21.8.6}. We shall
  not prove this refinement of the Zariski criterion, which we shall not need.
  In \sref{0.22.7}, we give, for complete local rings, a variant (also due to
  Nagata) of the Zariski criterion, proved essentially by the same method
  (with somewhat greater technical difficulties).
\end{enumerate}
\end{remark}

\subsection{Nagata's Jacobian criterion}
\label{subsection:0.22.7}

\begin{env}[22.7.1]
\label{0.22.7.1}
The Jacobian criterion of Nagata is the analog of the Jacobian criterion of Zariski, but for quotient rings of \emph{formal power series} rings over a field.
We shall give, as Nagata \cite{IV-31}, two versions, presented here as criteria of formal smoothness.
\end{env}

\begin{proposition}[22.7.2]
\label{0.22.7.2}
Let $k$ be a field, $\mathrm{A}=k[[\mathrm{T}_1,\ldots,\mathrm{T}_r]]$ a formal power series ring, equipped with its usual $k$-algebra structure, $\mathfrak{q}$ an ideal of $\mathrm{A}$, $\mathrm{B}=\mathrm{A}/\mathfrak{q}$.
Let $\mathfrak{p}$ be a prime ideal of $\mathrm{A}$ containing $\mathfrak{q}$.
Suppose that there exists a sub-$k$-algebra $\mathrm{C}'$ of $\mathrm{C}=\mathrm{A}/\mathfrak{p}$, isomorphic to a formal power series ring $k[[\mathrm{X}_1,\ldots,\mathrm{X}_s]]$, such that $\mathrm{C}$ is finite over $\mathrm{C}'$ and that the field of fractions $\mathrm{L}$ of $\mathrm{C}$ is a separable extension of $\mathrm{L}'=k((\mathrm{X}_1,\ldots,\mathrm{X}_s))$ (this hypothesis is always satisfied if $k$ is of characteristic $0$).
Then the following conditions are equivalent:
\begin{enumerate}
  \item[a)] $\mathrm{B}_\mathfrak{p}=\mathrm{A}_\mathfrak{p}/\mathfrak{q}\mathrm{A}_\mathfrak{p}$ is a $k$-algebra formally smooth for the $\mathfrak{p}$-preadic topology.
  \item[b)] There exist $k$-derivations $\mathrm{D}_i$ of $\mathrm{A}$ into itself $(1\leq i\leq m)$, and elements $f_i$ of $\mathfrak{q}$ $(1\leq i\leq m)$, such that the images of the $f_i$ in $\mathfrak{q}\mathrm{A}_\mathfrak{p}$ generate this ideal of $\mathrm{A}_\mathfrak{p}$, and that $\det(\mathrm{D}_if_j)\notin\mathfrak{p}$.
  \item[c)] $\mathrm{B}_\mathfrak{p}$ is a regular ring.
\end{enumerate}
\end{proposition}

\begin{proof}
The residue field of $\mathrm{B}_\mathfrak{p}$ is $k(\mathfrak{p})$.
Since $k((\mathrm{X}_1,\ldots,\mathrm{X}_s))$ is separable over $k$ \sref{0.21.9.6.4}, the field $k(\mathfrak{p})=\mathrm{L}$ is separable over $k$ by hypothesis; hence conditions $a$) and $c$) are equivalent by \sref{0.19.6.4}.
Furthermore, $\mathrm{A}_\mathfrak{p}$ is regular \sref{0.17.1.4}, \sref{0.17.3.2}, and therefore formally smooth over $k$ for its $\mathfrak{p}$-preadic topology \sref{0.19.6.4}.
By \sref{0.22.6.2},~(ii), condition $a$) is consequently equivalent to injectivity of
\begin{equation}
\label{0.22.7.2.1}
  j_0:(\mathfrak{q}/\mathfrak{q}^2)\otimes_\mathrm{A}\mathrm{L}
  \longrightarrow \Omega^1_{\mathrm{A}_\mathfrak{p}/k}\otimes_{\mathrm{A}_\mathfrak{p}}\mathrm{L}
  =\widehat{\Omega}^1_{\mathrm{A}/k}\otimes_\mathrm{A}\mathrm{L}.
  \tag{22.7.2.1}
\end{equation}

Consider also
\begin{equation}
\label{0.22.7.2.2}
  j:(\mathfrak{q}/\mathfrak{q}^2)\otimes_\mathrm{A}\mathrm{L}
  \xrightarrow{j_0}\Omega^1_{\mathrm{A}/k}\otimes_\mathrm{A}\mathrm{L}
  \longrightarrow\widehat{\Omega}^1_{\mathrm{A}/k}\otimes_\mathrm{A}\mathrm{L}.
  \tag{22.7.2.2}
\end{equation}
We show first that $b$) is equivalent to injectivity of $j$.
Indeed,
\[
  (\mathfrak{q}/\mathfrak{q}^2)\otimes_\mathrm{A}\mathrm{L}
  =(\mathfrak{q}\mathrm{A}_\mathfrak{p}/\mathfrak{q}^2\mathrm{A}_\mathfrak{p})\otimes_{\mathrm{A}_\mathfrak{p}}\mathrm{L}.
\]
If $F_i=j(f_i'\otimes 1)$, condition $b$) says that the matrix $(\langle F_i,\mathrm{D}_j\otimes 1\rangle)$ is invertible.
Thus the $F_i$, and hence the generators $f_i'\otimes 1$, are linearly independent, so $j$ is injective.
Conversely, if $j$ is injective, choose $f_i$ so that the
$f_i'\otimes 1$ form

\oldpage[0\textsubscript{IV-1}]{210}

a basis of the source.
Then the $F_i$ are linearly independent.
By \sref{0.21.9.3}, $\widehat{\Omega}^1_{\mathrm{A}/k}$ is free of rank $r$, and the $k$-derivations of $\mathrm{A}$ into itself generate its dual; hence there exist derivations $\mathrm{D}_j$ for which $(\langle F_i,\mathrm{D}_j\otimes 1\rangle)$ is invertible.
This is condition $b$).

These observations show already that $b$) implies $a$).
Conversely, assume $a$), so that $j_0$ is injective, and consider
\begin{equation}
\label{0.22.7.2.3}
\begin{tikzcd}[column sep=large]
  (\mathfrak{q}/\mathfrak{q}^2)\otimes_\mathrm{A}\mathrm{L} \ar[r,"j"] \ar[d,"\alpha"'] & \widehat{\Omega}^1_{\mathrm{A}/k}\otimes_\mathrm{A}\mathrm{L} \ar[d,equal] \\
  (\mathfrak{p}/\mathfrak{p}^2)\otimes_\mathrm{A}\mathrm{L} \ar[r,"i"'] & \widehat{\Omega}^1_{\mathrm{A}/k}\otimes_\mathrm{A}\mathrm{L}.
\end{tikzcd}
\tag{22.7.2.3}
\end{equation}
We shall use the following lemma.
\end{proof}

\begin{lemma}[22.7.2.4]
\label{0.22.7.2.4}
Let $\mathrm{R}$ be a regular local ring, $\mathfrak{m}$ its maximal ideal, $\mathrm{K}=\mathrm{R}/\mathfrak{m}$ its residue field.
If $\mathfrak{n}$ is an ideal of $\mathrm{R}$ such that $\mathrm{R}/\mathfrak{n}$ is regular, then the canonical homomorphism
\[
\label{0.22.7.2.5}
  \alpha:(\mathfrak{n}/\mathfrak{n}^2)\otimes_\mathrm{R}\mathrm{K}\to(\mathfrak{m}/\mathfrak{m}^2)\otimes_\mathrm{R}\mathrm{K}
  \tag{22.7.2.5}
\]
is injective.
\end{lemma}

\begin{proof}
Indeed, one knows \sref{0.17.1.6} that $\mathfrak{n}$ is generated by a subset $(x_i)_{1\leq i\leq l}$ forming part of a regular system of parameters $(x_i)_{1\leq i\leq n}$ of $\mathrm{R}$; the images of the $x_i$ for $1\leq i\leq l$ form a base of $(\mathfrak{n}/\mathfrak{n}^2)\otimes_\mathrm{R}\mathrm{K}$, and their images by $\alpha$ form part of the base of $(\mathfrak{m}/\mathfrak{m}^2)\otimes_\mathrm{R}\mathrm{K}$ formed by the images of the $x_i$ for $1\leq i\leq n$, from which the lemma.
\end{proof}

\begin{proof}[Continuation of the proof of Proposition~{\rm\sref{0.22.7.2}}]
By the lemma and the diagram \sref{0.22.7.2.3}, it remains only to prove that $i$ is injective.
Consider the sequence
\begin{equation}
\label{0.22.7.2.6}
  \mathfrak{p}/\mathfrak{p}^2\xrightarrow{h}
  \widehat{\Omega}^1_{\mathrm{A}/k}\widehat{\otimes}_\mathrm{A}(\mathrm{A}/\mathfrak{p})
  \xrightarrow{g}\widehat{\Omega}^1_{(\mathrm{A}/\mathfrak{p})/k}\to 0.
  \tag{22.7.2.6}
\end{equation}
By \sref{0.20.7.20}, $g$ is surjective and the image of $h$ is dense in $\Ker(g)$ for the $\mathfrak{m}$-adic topology.
Since $\widehat{\Omega}^1_{\mathrm{A}/k}$ is an $\mathrm{A}$-module of finite type, every submodule of
\[
  \widehat{\Omega}^1_{\mathrm{A}/k}\otimes_\mathrm{A}(\mathrm{A}/\mathfrak{p})
  =\widehat{\Omega}^1_{\mathrm{A}/k}\widehat{\otimes}_\mathrm{A}(\mathrm{A}/\mathfrak{p})
\]
is closed for the $\mathfrak{m}$-adic topology \sref[I]{I.7.3.5}; therefore \sref{0.22.7.2.6} is exact.
After tensoring with $\mathrm{L}$, one obtains the exact sequence
\[
  (\mathfrak{p}/\mathfrak{p}^2)\otimes_\mathrm{A}\mathrm{L}\xrightarrow{i}
  \widehat{\Omega}^1_{\mathrm{A}/k}\otimes_\mathrm{A}\mathrm{L}\to
  \widehat{\Omega}^1_{(\mathrm{A}/\mathfrak{p})/k}\otimes_{\mathrm{A}/\mathfrak{p}}\mathrm{L}\to 0.
\]
The separability hypothesis and \sref{0.21.9.5} show that the last module has rank $s$ over $\mathrm{L}$; the middle module has rank $r$ by \sref{0.21.9.3}.
Thus $\rg_\mathrm{L}(\operatorname{Im}i)=r-s$.
Since $\mathrm{A}/\mathfrak{p}$ is finite over a subalgebra isomorphic to $k[[\mathrm{T}_1,\ldots,\mathrm{T}_s]]$, one has $\dim(\mathrm{A}/\mathfrak{p})=s$ by \sref{0.17.1.4} and \sref{0.16.1.5}, and hence
\[
  r-s=\dim(\mathrm{A})-\dim(\mathrm{A}/\mathfrak{p})=\dim(\mathrm{A}_\mathfrak{p})
\]

\oldpage[0\textsubscript{IV-1}]{211}

by \sref{0.16.5.11}.
Finally, $\mathrm{A}_\mathfrak{p}$ is regular, so $\dim(\mathrm{A}_\mathfrak{p})$ is the rank over $\mathrm{L}$ of
$(\mathfrak{p}/\mathfrak{p}^2)\otimes_\mathrm{A}\mathrm{L}$ \sref{0.17.1.1}.
The source and image of $i$ therefore have the same rank, and $i$ is injective.
\end{proof}

\begin{theorem}[22.7.3]
\label{0.22.7.3}
Let $k$ be a field, $\mathrm{A}$ a local Noetherian complete ring with residue field $\mathrm{K}$.
Suppose that:
\begin{enumerate}
  \item[$1^\circ$] $[k:k^p]<+\infty$ (where $p$ is the characteristic exponent of $k$);
  \item[$2^\circ$] $\mathrm{K}$ is a finite extension of a separable extension $\mathrm{K}_0$ of $k$;
  \item[$3^\circ$] $\mathrm{A}$ is equipped with a structure of $\mathrm{K}_0$-algebra formally smooth (for the preadic topology).
\end{enumerate}
Let $\mathfrak{q}$ be an ideal of $\mathrm{A}$, $\mathrm{B}=\mathrm{A}/\mathfrak{q}$, $\mathfrak{p}$ a prime ideal of $\mathrm{A}$ containing $\mathfrak{q}$.
The following conditions are equivalent:
\begin{enumerate}
  \item[a)] $\mathrm{B}_\mathfrak{p}$ is a $k$-algebra formally smooth for the $\mathfrak{p}$-preadic topology.
  \item[b)] There exist $k$-derivations $\mathrm{D}_i$ of $\mathrm{A}$ into itself $(1\leq i\leq m)$, and elements $f_i$ $(1\leq i\leq m)$ of $\mathfrak{q}$, such that the images of the $f_i$ in $\mathfrak{q}\mathrm{A}_\mathfrak{p}$ generate this ideal of $\mathrm{A}_\mathfrak{p}$, and that $\det(\mathrm{D}_if_j)\notin\mathfrak{p}$.
  \item[b')] There exists a sub-extension $k'$ of $\mathrm{K}_0$, containing $\mathrm{K}_0^p$, such that $[\mathrm{K}_0:k']<+\infty$, $k'$-derivations $\mathrm{D}_i$ of $\mathrm{A}$ into itself, and elements $f_i$ of $\mathfrak{q}$ $(1\leq i\leq m)$, such that the images of the $f_i$ in $\mathfrak{q}\mathrm{A}_\mathfrak{p}$ generate this ideal of $\mathrm{A}_\mathfrak{p}$, and that $\det(\mathrm{D}_if_j)\notin\mathfrak{p}$.
\end{enumerate}
\end{theorem}

\begin{proof}
We distinguish two cases, according to whether $p=1$ or $p>1$.

\emph{A)} $k$ \emph{is of characteristic} $0$.
Since $\mathrm{K}$ is then a separable extension of $\mathrm{K}_0$,
\sref{0.19.6.4} shows that $\mathrm{A}$ is $\mathrm{K}_0$-isomorphic to a
formal power series ring $\mathrm{K}[[\mathrm{T}_1,\ldots,\mathrm{T}_r]]$
equipped with its usual $\mathrm{K}$-algebra structure. Taking
\sref{0.19.8.8}, (ii), applied to $\mathrm{A}/\mathfrak{p}$ into account, the
general hypotheses of \sref{0.22.7.2} are satisfied after replacing $k$ by
$\mathrm{K}$. Moreover, by \sref{0.19.6.4}, saying that
$\mathrm{B}_\mathfrak{p}$ is formally smooth as a $k$-algebra or as a
$\mathrm{K}$-algebra amounts to the same thing, both conditions being
equivalent to regularity of $\mathrm{B}_\mathfrak{p}$. We may therefore apply
\sref{0.22.7.2}, which immediately proves the equivalence of $a$), $b$), and
$b'$) (with $k'=\mathrm{K}_0$ in $b'$)).

\emph{B)} $k$ \emph{is of characteristic} $p>0$.
Since $\mathrm{A}$ is a formally smooth $\mathrm{K}_0$-algebra and
$\mathrm{K}_0$ is separable over $k$, it follows from \sref{0.19.3.5}, (ii),
and \sref{0.19.6.1} that $\mathrm{A}$ is a formally smooth $k$-algebra; by
\sref{0.22.5.9}, $\mathrm{A}_\mathfrak{p}$ is also formally smooth over $k$
for the $\mathfrak{p}$-preadic topology.
It follows therefore from \sref{0.22.6.2}, (ii) that the condition $a$) is equivalent to the fact that the homomorphism $j_0$ defined in \sref{0.22.7.2.1} (where $\mathrm{L}=k(\mathfrak{p})$) is injective.
Therefore \sref{0.19.1.12}, $c$) shows already that $b$) implies $a$); and since moreover $b'$) trivially implies $b$), it all amounts to showing that the hypothesis that $j_0$ is injective implies $b'$).

Let us first denote by $k'$ a sub-extension of $\mathrm{K}_0$ such that $[\mathrm{K}_0:k']<+\infty$.
Note that $\widehat{\Omega}^1_{\mathrm{A}/k'}$ is a free $\mathrm{A}$-module of finite type.
By the argument of \sref{0.22.7.2}, it is enough to prove that (for a suitable
choice of $k'$) the composite homomorphism
\[
\label{0.22.7.3.1}
  j':(\mathfrak{q}/\mathfrak{q}^2)\otimes_\mathrm{A}\mathrm{L}\xrightarrow{j_0}\Omega^1_{\mathrm{A}/k}\otimes_\mathrm{A}\mathrm{L}\to\widehat{\Omega}^1_{\mathrm{A}/k'}\otimes_\mathrm{A}\mathrm{L}
  \tag{22.7.3.1}
\]

\oldpage[0\textsubscript{IV-1}]{212}

is injective.
Consider again the commutative diagram
\[
\label{0.22.7.3.2}
  \begin{tikzcd}
    (\mathfrak{q}/\mathfrak{q}^2)\otimes_\mathrm{A}\mathrm{L}
      \ar[r,"j_0"] \ar[d,"\alpha"'] &
    \Omega^1_{\mathrm{A}/k}\otimes_\mathrm{A}\mathrm{L}
      \ar[r] \ar[d,equal] &
    \widehat{\Omega}^1_{\mathrm{A}/k'}\otimes_\mathrm{A}\mathrm{L}
      \ar[d,equal]\\
    (\mathfrak{p}/\mathfrak{p}^2)\otimes_\mathrm{A}\mathrm{L}
      \ar[r,"i_0"'] &
    \Omega^1_{\mathrm{A}/k}\otimes_\mathrm{A}\mathrm{L} \ar[r] &
    \widehat{\Omega}^1_{\mathrm{A}/k'}\otimes_\mathrm{A}\mathrm{L}
  \end{tikzcd}
  \tag{22.7.3.2}
\]

Since $j_0$ is injective, we have $\Ker(i_0\circ\alpha)=0$.
But as $\mathrm{A}_\mathfrak{p}$ is regular, and the hypothesis $a$) implies that $\mathrm{B}_\mathfrak{p}=\mathrm{A}_\mathfrak{p}/\mathfrak{q}\mathrm{A}_\mathfrak{p}$ is also regular \sref{0.19.6.5}, the lemma \sref{0.22.7.2.4} shows that $\alpha$ is injective, from which $\alpha^{-1}(\Ker(i_0))=0$.
If $i'$ is the composite homomorphism of the second line of
\sref{0.22.7.3.2}, one likewise has
$\Ker(j')=\alpha^{-1}(\Ker(i'))$; it therefore remains to prove that, for a
suitable choice of $k'$,
\begin{equation}
\label{0.22.7.3.3}
  \Ker(i')=\Ker(i).
  \tag{22.7.3.3}
\end{equation}

But the same argument as in \sref{0.22.7.2} gives an exact sequence
\[
  (\mathfrak{p}/\mathfrak{p}^2)\otimes_\mathrm{A}\mathrm{L}
  \xrightarrow{i'}\widehat{\Omega}^1_{\mathrm{A}/k'}
  \otimes_\mathrm{A}\mathrm{L}\longrightarrow
  \widehat{\Omega}^1_{(\mathrm{A}/\mathfrak{p})/k'}
  \otimes_{\mathrm{A}/\mathfrak{p}}\mathrm{L}\longrightarrow 0.
\]

The hypotheses allow us to apply \sref{0.21.9.8}; hence there exists a
sub-extension $k'$ of $\mathrm{K}_0$ containing $\mathrm{K}_0^p$, with
$[\mathrm{K}_0:k']<+\infty$, for which
\[
  \rg_\mathrm{L}\bigl(
  \widehat{\Omega}^1_{(\mathrm{A}/\mathfrak{p})/k'}
  \otimes_{\mathrm{A}/\mathfrak{p}}\mathrm{L}\bigr)
  =\dim(\mathrm{A}/\mathfrak{p})
  +\rg_\mathrm{L}\Upsilon_{\mathrm{L}/k}
  +\rg_{\mathrm{K}_0}\Omega^1_{\mathrm{K}_0/k'}.
\]
But
\[
  \dim(\mathrm{A}/\mathfrak{p})
  =\dim(\mathrm{A})-\dim(\mathrm{A}_\mathfrak{p})
\]
by \sref{0.16.5.11},
\[
  \dim(\mathrm{A}_\mathfrak{p})
  =\rg_\mathrm{L}\bigl((\mathfrak{p}/\mathfrak{p}^2)
  \otimes_\mathrm{A}\mathrm{L}\bigr)
\]
because $\mathrm{A}_\mathfrak{p}$ is regular \sref{0.17.1.1}, and
\[
  \dim(\mathrm{A})+\rg_{\mathrm{K}_0}\Omega^1_{\mathrm{K}_0/k'}
  =\rg_\mathrm{L}\bigl(
  \widehat{\Omega}^1_{\mathrm{A}/k'}\otimes_\mathrm{A}\mathrm{L}\bigr)
\]
by \sref{0.21.9.2}. Consequently
\begin{equation}
\label{0.22.7.3.4}
  \rg_\mathrm{L}(\Ker(i'))=\rg_\mathrm{L}\Upsilon_{\mathrm{L}/k}.
  \tag{22.7.3.4}
\end{equation}

On the other hand, since the residue field of $\mathrm{A}_\mathfrak{p}$ is
formally smooth over the prime field of $k$, there is an exact sequence
\sref{0.20.6.22.1}
\[
  \Upsilon_{\mathrm{L}/k}
  \xrightarrow{\chi_{\mathrm{A}_\mathfrak{p}/k}}
  (\mathfrak{p}\mathrm{A}_\mathfrak{p})/
  (\mathfrak{p}\mathrm{A}_\mathfrak{p})^2
  \xrightarrow{i}\Omega^1_{\mathrm{A}_\mathfrak{p}/k}
  \otimes_{\mathrm{A}_\mathfrak{p}}\mathrm{L}.
\]
Since $\mathrm{A}_\mathfrak{p}$ is a formally smooth $k$-algebra for the
$\mathfrak{p}$-preadic topology, \sref{0.22.2.2} shows that
$\chi_{\mathrm{A}_\mathfrak{p}/k}$ is injective. Thus $\Ker(i)$ is isomorphic
to $\Upsilon_{\mathrm{L}/k}$; taking \sref{0.22.7.3.4} and
$\Ker(i)\subset\Ker(i')$ into account proves \sref{0.22.7.3.3}, and hence the
theorem.
\end{proof}

\oldpage[0\textsubscript{IV-1}]{213}
\begin{remark}[22.7.4]
\label{0.22.7.4}
It was shown in fact (using \sref{0.21.9.6}) that condition $b'$) is even equivalent to the other conditions of \sref{0.22.7.3} when one subjects $k'$ to be one of the fields of a decreasing filtering family $(k_\alpha)$ of subfields of $\mathrm{K}_0$, containing $\mathrm{K}_0^p$, such that $[\mathrm{K}_0:k_\alpha]<+\infty$ for every $\alpha$ and that $\bigcap k_\alpha=\mathrm{K}_0^p$.
\end{remark}

\begin{corollary}[22.7.5]
\label{0.22.7.5}
Under the general hypotheses of \sref{0.22.7.3}, the set of prime ideals $\mathfrak{n}\in\Spec(\mathrm{B})$ such that $\mathrm{B}_\mathfrak{n}$ is a $k$-algebra formally smooth (for the $\mathfrak{n}$-preadic topology) is open in $\Spec(\mathrm{B})$.
\end{corollary}

\begin{proof}
With the notations of \sref{0.22.7.3}, it suffices to see that if $\mathrm{B}_\mathfrak{p}$ is formally smooth over $k$, then the same is true of $\mathrm{B}_{\mathfrak{p}'}$ for every $\mathfrak{p}'\in\mathrm{V}(\mathfrak{q})$ sufficiently near to $\mathfrak{p}$ in $\Spec(\mathrm{A})$.
But, if the $k$-derivations $\mathrm{D}_i$ and the elements $f_i$ of $\mathfrak{q}$ satisfy criterion $b$) of \sref{0.22.7.3}, one has also $f=\det(\mathrm{D}_if_j)\notin\mathfrak{p}'$ for every $\mathfrak{p}'\in\mathrm{D}(f)$; on the other hand, there exists $g\notin\mathfrak{p}$ such that the images of $f_i$ in $\mathfrak{q}\mathrm{A}_{\mathfrak{p}'}$ generate this ideal of $\mathrm{A}_{\mathfrak{p}'}$ for every $\mathfrak{p}'\in\mathrm{D}(g)$, which finishes the proof of the corollary.
\end{proof}

\begin{corollary}[22.7.6]
\label{0.22.7.6}
\textnormal{(Nagata.)}
Let $\mathrm{B}$ be a local Noetherian complete ring containing a field.
Then the set of $\mathfrak{n}\in\Spec(\mathrm{B})$ such that $\mathrm{B}_\mathfrak{n}$ is regular is open in $\Spec(\mathrm{B})$.
\end{corollary}

\begin{proof}
Indeed, $\mathrm{B}$ is an algebra over a \emph{prime} field $k$, hence by \sref{0.19.8.8}[ii] it is of the form $\mathrm{A}/\mathfrak{q}$, where $\mathrm{A}=\mathrm{K}_0[[\mathrm{T}_1,\ldots,\mathrm{T}_n]]$ is a formal power series ring and $\mathrm{K}_0$ an extension of $k$; on the other hand, saying that $\mathrm{B}_\mathfrak{n}$ is regular is equivalent to saying that it is a $k$-algebra formally smooth for the $\mathfrak{n}$-preadic topology \sref{0.19.6.4}.
All the conditions of application of \sref{0.22.7.5} are thus realized.
\end{proof}

\begin{remark}[22.7.7]
\label{0.22.7.7}
\begin{enumerate}
  \item[\rm{(i)}] We shall see further on, using \sref{0.22.7.6}, that the conclusion of \sref{0.22.7.6} can be extended to the case of any complete Noetherian local ring (\textbf{IV},~6.12.7).
  \item[\rm{(ii)}] The conclusion of the theorem \sref{0.22.7.3} is no longer necessarily exact when one no longer assumes that $[k:k^p]$ is finite.
  Take for example $k=\mathbf{F}_p(\mathrm{X}_n)_{n\geq 0}$, $\mathrm{A}=k[[\mathrm{T},\mathrm{U}]]$, $\mathfrak{q}$ the principal ideal $\mathrm{A}f$, where $f=\mathrm{U}^p-\sum_{n=0}^\infty\mathrm{X}_n\mathrm{T}^{np}$; $\mathrm{A}$ is a factorial ring in which $f$ is an extremal element, for it is extremal in the polynomial ring $k((\mathrm{T}))[\mathrm{U}]$, hence also in the formal power series ring $k((\mathrm{T}))[[\mathrm{U}]]$ (Bourbaki, \emph{Alg.~comm.}, chap.~VII, \textsection~3), and \emph{a fortiori} in $\mathrm{A}$.
  Take $\mathfrak{p}=\mathfrak{q}$, so that $\mathrm{B}_\mathfrak{p}$ is the field of fractions denoted by $\mathrm{E}$ in \sref{0.21.9.7}, which is separable over $k$ (\emph{loc.~cit.}), hence a $k$-algebra formally smooth.
  But one has, with the notations of \sref{0.22.7.3}, $k=\mathrm{K}_0=\mathrm{K}$, and the $k$-derivations of $\mathrm{A}$ into itself are combinations of $\partial/\partial\mathrm{T}$ and $\partial/\partial\mathrm{U}$; as one has $\partial f/\partial\mathrm{T}=\partial f/\partial\mathrm{U}=0$, the criteria $b$) and $b'$) of \sref{0.22.7.3} are not satisfied.
\end{enumerate}
\end{remark}
% END NATIVE INPUT ega0/ega0-22
% BEGIN NATIVE INPUT ega0/ega0-23
\section{Japanese rings}
\label{section:0.23}

\oldpage[0\textsubscript{IV-1}]{213}
The results of this section will be completed in \textbf{IV}, \sref[IV]{subsection:IV.7.6} and \sref[IV]{subsection:IV.7.7}.

\subsection{Japanese rings}
\label{subsection:0.23.1}

\begin{definition}[23.1.1]
\label{0.23.1.1}
We say that an integral ring $A$ is a \emph{Japanese ring} if, for every finite extension $K'$ of its field of fractions $K$, the integral closure $A'$ of $A$ in $K'$ is
\oldpage[0\textsubscript{IV-1}]{214}
an $A$-module of finite type (in other words, a finite $A$-algebra).
We say that a ring $A$ is \emph{universally Japanese} if every integral $A$-algebra of finite type is a Japanese ring.
\end{definition}

It is clear that if $A$ is a Japanese ring, then every ring of fractions $S^{-1}A$ is also a Japanese ring, since (with the preceding notation) $S^{-1}A'$ is the integral closure of $S^{-1}A$ in $K'$, and is evidently an $S^{-1}A$-module of finite type.

A Noetherian integral ring, even a \emph{discrete valuation ring}, is not necessarily a Japanese ring \cite{IV-30}.

\begin{proposition}[23.1.2]
\label{0.23.1.2}
Let $A$ be a Noetherian integral ring, $K$ its field of fractions.
If, for every finite radicial extension $K'$ of $K$, the integral closure of $A$ in $K'$ is an $A$-module of finite type, then $A$ is a Japanese ring.
\end{proposition}

\begin{proof}
Since $A$ is Noetherian, it suffices, to verify that $A$ is a Japanese ring, to see that for every finite \emph{quasi-Galois} extension $L$ of $K$, the integral closure $B$ of $A$ in $L$ is an $A$-module of finite type.
Now $L$ is a \emph{Galois} extension of the largest \emph{radicial} extension $K'$ of $K$ contained in $L$; and if $A'$ is the integral closure of $A$ in $K'$, $B$ is the integral closure of $A'$ in $L$.
But we know that in a \emph{separable} extension, the integral closure of $A'$ is an $A'$-module of finite type (Bourbaki, \emph{Alg.~comm.}, chap.~V, \textsection1, n\textsuperscript{o}~6, cor.~1 of prop.~20), whence the proposition.
\end{proof}

It follows from \sref{0.23.1.2} that when $K$ is of \emph{characteristic}~$0$, saying that $A$ is a Japanese ring means that its \emph{integral closure} $A'$ is an $A$-module of finite type.

\begin{theorem}[23.1.3]
\label{0.23.1.3}
\textnormal{(Tate)}
Let $A$ be a Noetherian integral ring, $x\neq 0$ an element of $A$.
Suppose the following conditions are satisfied:
\begin{enumerate}
  \item[{\rm(i)}] $A$ is integrally closed.
  \item[{\rm(ii)}] $\mathfrak{p}=xA$ is prime and $A$ is separated and complete for the $\mathfrak{p}$-preadic topology.
  \item[{\rm(iii)}] $A/xA$ is a Japanese ring.
\end{enumerate}
Then $A$ is a Japanese ring.
\end{theorem}

We will use the following lemma:

\begin{lemma}[23.1.3.1]
\label{0.23.1.3.1}
Let $A$ be a ring, $x$ an element of $A$ that is a non-zero-divisor in $A$ and such that $xA=\mathfrak{p}$ is prime; then, for every integer $n>0$, the inverse image of $x^nA_\mathfrak{p}$ in $A$ is $x^nA$.
\end{lemma}

\begin{proof}
Indeed, suppose that $b\in A$ is an element such that $b/1=x^na/s$ in $A_\mathfrak{p}$, where $a\in A$ and $s\notin\mathfrak{p}$; then there exists $s'\notin\mathfrak{p}$ such that $s'sb=x^nas'$, whence $s'sb\in\mathfrak{p}$, and since $s's\notin\mathfrak{p}$, this implies $b\in\mathfrak{p}$, that is to say $b=xb'$ with $b'\in A$; since $x$ is a non-zero-divisor, we conclude that $s'sb'=x^{n-1}as'$\footnote{The printed source has $as$ here; cancelling $x$ in $s'sxb'=x^nas'$ retains the factor $s'$.}, and it suffices to argue by induction on $n$.
\end{proof}

\begin{proof}[Proof of Theorem~{\rm\sref{0.23.1.3}}]
To prove the theorem, we can restrict to the case where the field of fractions $K$ of $A$ is of characteristic $p>0$, since $A$ is integrally closed.
Let $K'$ be a finite radicial extension of $K$, so that there exists a power $q=p^e$ such that $K'^q\subset K$; by replacing $K'$ by a larger radicial extension, we may even suppose that there exists $y\in K'$ such that $y^q=x$.
Let $A'$ be the integral closure of $A$ in $K'$; since $A$ is integrally closed, $A'$ is the set of $x'\in K'$ such that $x'^q\in A'\cap K=A$.
Set $V=A_\mathfrak{p}$; the maximal ideal $\mathfrak{m}=xV$ of $V$ being principal, we know that the local Noetherian ring $V$ is a \emph{discrete valuation ring} \sref{0.17.1.4}; since $V$ is integrally closed, the same
\oldpage[0\textsubscript{IV-1}]{215}
argument as above shows that the integral closure $V'$ of $V$ in $K'$ is the set of $x'\in K'$ such that $x'^q\in V$; we know (Bourbaki, \emph{Alg.~comm.}, chap.~VI, \textsection8, n\textsuperscript{o}~6, prop.~6, and chap.~V, \textsection2, n\textsuperscript{o}~3, lemma~4) that $V'$ is a \emph{discrete valuation ring}, whose maximal ideal $\mathfrak{m}'$ is the set of $x'\in K'$ such that $x'^q\in\mathfrak{m}$, and moreover the residue field $V'/\mathfrak{m}'$ is a \emph{finite} extension of $V/\mathfrak{m}$ (\emph{loc.~cit.}, chap.~VI, \textsection8, n\textsuperscript{o}~1, lemma~2).
Let us show that, for every integer $n>0$, we have
\[
\label{0.23.1.3.2}
  \mathfrak{m}'^n\cap A'=y^nA'.
  \tag{23.1.3.2}
\]

Indeed, since $y^q=x$, we have $y\in\mathfrak{m}'$, hence $y^nA'\subset\mathfrak{m}'^n\cap A'$.
Conversely, let $x'\in\mathfrak{m}'^n\cap A'$, and set $x'=z'y^n$ with $z'\in K'$.
We have $x'^q\in A'$; on the other hand, $x'$ is a sum of products $t'_1\cdots t'_n$ with $t'_i\in\mathfrak{m}'$, hence $t_i'^q\in\mathfrak{m}$, and we conclude that $x'^q\in\mathfrak{m}^n\cap A$.
Now the lemma \sref{0.23.1.3.1} proves that $\mathfrak{m}^n\cap A=x^nA$, hence we have $z'^qx^n\in x^nA$, or equivalently $z'^q\in A$ since $x\neq 0$, which establishes \sref{0.23.1.3.2}.

Let us show secondly that on $A'$ the $xA'$-preadic topology is \emph{separated}; this topology is also the $yA'$-preadic topology of $A'$, and it follows from \sref{0.23.1.3.2} that this topology is induced by the $\mathfrak{m}'$-\emph{preadic} topology of $V'$, which is separated since $V'$ is a discrete valuation ring.

Let us then prove that $A'/xA'$ is an $A$-module of \emph{finite type}.
Since $x=y^q$ and $y^kA'/y^{k+1}A'$ is isomorphic to $A'/yA'$, it suffices to show that $A'/yA'$ is an $A$-module of finite type; but the formula \sref{0.23.1.3.2} applied for $n=1$ shows that $A'/yA'$ is a subring of $V'/\mathfrak{m}'$, that is to say of the \emph{residue field} of $V'$, which is a \emph{finite} extension of the residue field $V/\mathfrak{m}$.
Now $V/\mathfrak{m}$ is the field of fractions of $A/\mathfrak{p}$, and since $A'$ is integral over $A$, $A'/yA'$ is integral over $A/\mathfrak{p}$, hence contained in the integral closure of $A/\mathfrak{p}$ in $V'/\mathfrak{m}'$; since by hypothesis $A/\mathfrak{p}$ is a Japanese Noetherian ring, $A'/yA'$ is an $(A/\mathfrak{p})$-module of finite type, and \emph{a fortiori} an $A$-module of finite type.

This being the case, the $xA'$-preadic topology of $A'$ being separated, $A'$ is a subring of its completion $\hat{A}'$ for this topology; but since $A'/xA'$ is an $A$-module of finite type and $A$ is separated and complete for the $xA$-preadic topology, $\hat{A}'$ is an $A$-module of finite type by virtue of \sref[0\textsubscript{I}]{0.7.2.9}, hence so is $A'$.
\end{proof}

\begin{corollary}[23.1.4]
\label{0.23.1.4}
Let $A$ be a Japanese Noetherian and integrally closed ring.
Then every formal power series ring $A[[T_1,\ldots,T_r]]$ is a Japanese ring.
\end{corollary}

\begin{proof}
We know that for every integrally closed Noetherian ring $B$, the formal power series ring $B[[T]]$ is Noetherian and integrally closed (Bourbaki, \emph{Alg.~comm.}, chap.~V, \textsection1, n\textsuperscript{o}~4, prop.~14); we can thus, by induction on $r$\footnote{The printed source has $n$; the induction is on the number $r$ of variables in the statement.}, restrict to proving that $A[[T]]$ is a Japanese ring.
Now the element $x=T$ satisfies all the conditions of \sref{0.23.1.3}, whence the conclusion.
\end{proof}

\begin{theorem}[23.1.5]
\label{0.23.1.5}
\textnormal{(Nagata)}
Every complete local Noetherian integral ring $A$ is a Japanese ring.
\end{theorem}

\begin{proof}
We know \sref{0.19.8.8}[ii] that there exists a subring $B$ of $A$ that is a complete local
\oldpage[0\textsubscript{IV-1}]{216}
ring and is in addition \emph{regular}, such that $A$ is a finite $B$-algebra; it clearly suffices to prove that $B$ is a Japanese ring, that is to say, we can restrict to the case where $A$ is in addition regular.
Let us argue by induction on $n=\dim(A)$, the theorem being evident for $n=0$.
Suppose then $n>0$, and, denoting by $\mathfrak{m}$ the maximal ideal of $A$, let $x$ be an element of $\mathfrak{m}\setmin\mathfrak{m}^2$; we know \sref{0.17.1.8} that $A/xA$ is a regular ring, and in addition (\sref{0.17.1.7} and \sref{0.16.3.4}) that $\dim(A/xA)=n-1$; moreover, the ideal $xA$ is closed in $A$ \sref[0\textsubscript{I}]{0.7.3.5}, hence $A/xA$ is complete.
By the induction hypothesis, $A/xA$ is thus a Japanese ring, and it then follows from \sref{0.23.1.3} that the same is true of $A$.
\end{proof}

\begin{corollary}[23.1.6]
\label{0.23.1.6}
Let $A$ be a complete local Noetherian integral ring, $K$ its field of fractions, $K'$ a finite extension of $K$; then the integral closure $A'$ of $A$ in $K'$ is a complete local ring.
\end{corollary}

\begin{proof}
We already know from \sref{0.23.1.5} that $A'$ is a finite $A$-algebra, hence a complete semi-local Noetherian ring, and consequently a \emph{direct product} of local rings; but since $A'$ is integral, it is necessarily a local ring.
\end{proof}

\begin{proposition}[23.1.7]
\label{0.23.1.7}
Let $A$ be a local Noetherian integral ring, $K$ its field of fractions, $K'$ a finite extension of $K$, $A'$ the integral closure of $A$ in $K'$.
Suppose that the completion $\hat{A}$ is reduced, and let $R$ denote its total ring of fractions.
\begin{enumerate}
  \item[{\rm(i)}] If the ring $R\otimes_K K'$ is \emph{reduced} (which will be the case in particular when $K'$ is a \emph{separable} extension of $K$, $R$ being a direct product of fields, hence extensions of $K$ \textnormal{(Bourbaki, \emph{Alg.}, chap.~VIII, \textsection7, n\textsuperscript{o}~3, cor.~1 of th.~1)}), \emph{then $A'$ is an $A$-module of finite type}.
  \item[{\rm(ii)}] \emph{If in particular $R$ is a separable $K$-algebra, then $A$ is a Japanese ring.}
\end{enumerate}
\end{proposition}

\begin{proof}
\medskip\noindent
(i).
Set $A_1=\hat{A}$, $A'_1=A'\otimes_A A_1$, $K'_1=K'\otimes_A A_1$.
Since $A_1$ is a flat $A$-module \sref[0\textsubscript{I}]{0.7.3.5}, $A'_1$ is identified with a subring of $K'_1$ and is evidently integral over $A_1$.
On the other hand, if we set $K_1=K\otimes_A A_1$, we can write $K'_1=K'\otimes_K K_1$; since $A$ is integral and $A_1$ is a flat $A$-module, every element $\neq 0$ of $A$ is $A_1$-regular \sref[0\textsubscript{I}]{0.6.3.4}; since $K=S^{-1}A$, where $S=A\setmin\{0\}$, we have $K_1=S^{-1}A_1$, and the preceding remark proves that $K_1$ is identified with a \emph{subring} of the total ring of fractions $R$ of $A_1$.
Since every $K$-module is flat, $K'_1$ is identified with a \emph{subring} of $K'\otimes_K R$, which by hypothesis is \emph{reduced} and is a \emph{finite} $R$-algebra.
Denoting by $\mathfrak{q}_i$ ($1\leq i\leq n$) the minimal prime ideals of $A_1$, by $B_i$ the integral ring $A_1/\mathfrak{q}_i$, by $L_i$ the field of fractions of $B_i$, we see that $R$ is the direct product of the $L_i$, and $K'\otimes_K R$ the direct product of the $K'\otimes_K L_i$; these latter $K$-algebras, being reduced, are direct products of fields, \emph{finite} extensions of the $L_i$.
Since $B_i$ is a complete local Noetherian integral ring, the theorem of Nagata \sref{0.23.1.5} proves that the integral closure of $B_i$ in $K'\otimes_K L_i$ is a $B_i$-module of finite type, and \emph{a fortiori} an $A_1$-module of finite type; since every element of $K'\otimes_K L_i$ that is integral over $A_1$ is also integral over $B_i$ (being annihilated by $\mathfrak{q}_i$), we conclude finally that the integral closure of $A_1$ in $K'\otimes_K R$ is an $A_1$-module of finite type.
But since $A'_1$ is contained in this integral closure and $A_1$ is Noetherian, $A'_1$ is also an $A_1$-module of finite type.
Finally, since $A_1$ is a faithfully flat $A$-module \sref[0\textsubscript{I}]{0.7.3.5}, we conclude that $A'$ is an $A$-module of finite type (Bourbaki, \emph{Alg.~comm.}, chap.~I, \textsection3, n\textsuperscript{o}~6, prop.~11).

\medskip\noindent
\oldpage[0\textsubscript{IV-1}]{217}
(ii).
The hypothesis implies that the $L_i$ are separable extensions of $K$, hence $K'\otimes_K L_i$ is reduced (Bourbaki, \emph{Alg.}, chap.~VIII, \textsection7, n\textsuperscript{o}~3, th.~1) and we can apply~(i) to \emph{every} finite extension $K'$ of $K$, which proves our assertion.
\end{proof}

\subsection{Integral closure of a local Noetherian integral ring}
\label{subsection:0.23.2}

\begin{env}[23.2.1]
\label{0.23.2.1}
In what follows, for every ring $A$, we will write $A_\red$ for short for the quotient of $A$ by its nilradical (so that if $X=\Spec(A)$, we have $X_\red=\Spec(A_\red)$).
We will say that a \emph{local} ring $A$ is \emph{unibranch} if $A_\red$ is integral and if the integral closure of $A_\red$ is a \emph{local} ring, which generalizes the definition given in \sref[III]{III.4.3.6}.
We will say that $A$ is \emph{geometrically unibranch} if it is unibranch and if the residue field of the local ring that is the integral closure of $A_\red$ is a \emph{radicial} extension of that of $A$.
It is clear that a local ring that is integrally closed is geometrically unibranch; it follows from \sref{0.23.1.6} that a complete local Noetherian integral ring is unibranch.
\end{env}

\begin{lemma}[23.2.2]
\label{0.23.2.2}
Let $A$ be an integral ring, $A'$ its integral closure; for the canonical morphism $f:\Spec(A')\to\Spec(A)$ to be radiciel, it is necessary and sufficient that for every $\mathfrak{p}\in\Spec(A)$, $A_\mathfrak{p}$ be geometrically unibranch; when this is so, $f$ is a homeomorphism.\footnote{In the remainder of this item, we use notions and results set out in chap.~IV, \textsection\textsection2 and~5; since the results of this item are not used before chap.~IV, \textsection6, this does not lead to circular reasoning.}
\end{lemma}

\begin{proof}
Indeed, for every $\mathfrak{p}\in\Spec(A)$, the integral closure of $A_\mathfrak{p}$ is $A'_\mathfrak{p}$ (Bourbaki, \emph{Alg.~comm.}, chap.~V, \textsection1, n\textsuperscript{o}~5, prop.~16), and all the prime ideals of $A'_\mathfrak{p}$ above $\mathfrak{p}A_\mathfrak{p}$ are maximal (\emph{loc.~cit.}, \textsection2, n\textsuperscript{o}~1, prop.~1).
Saying that $f$ is \emph{injective} thus means that for every $\mathfrak{p}\in\Spec(A)$, $A'_\mathfrak{p}$ is a \emph{local} ring, that is to say that the $A_\mathfrak{p}$ are unibranch.
Saying that every $A_\mathfrak{p}$ is geometrically unibranch then means that $f$ is radiciel by virtue of \sref[I]{I.3.5.8}.
When this is so, $f$ is surjective and closed \sref[II]{II.6.1.10}, hence a homeomorphism.
\end{proof}

\begin{lemma}[23.2.3]
\label{0.23.2.3}
Let $A$ be an integral ring, $K$ its field of fractions, $B$ a Noetherian ring, $\vphi:A\to B$ a ring homomorphism making $B$ an $A$-module that is flat.
Let $K'$ be an extension of $K$; set $B_1=B_\red$, and let $R$ be the total ring of fractions of $B_1$.
Then, for every sub-$A$-algebra $A'$ of $K'$, the canonical homomorphism
\[
\label{0.23.2.3.1}
  (A'\otimes_A B)_\red\to(K'\otimes_A R)_\red
  \tag{23.2.3.1}
\]
is injective.\footnote{The French source prints ``(23.1.8.1)''; the formula occurs in Lemma~23.2.3 and is therefore numbered (23.2.3.1) here.}
\end{lemma}

\begin{proof}
We can in fact consider the canonical homomorphism $h:A'\otimes_A B\to K'\otimes_A R$ as the composition of the following homomorphisms
\[
  A'\otimes_A B\xrightarrow{u}K'\otimes_A B\xrightarrow{v}K'\otimes_A B_1\xrightarrow{w}K'\otimes_A R.
\]
Since $B$ is a flat $A$-module, $u$ is injective; similarly, $K'$ being a flat $K$-module and $K$ a flat $A$-module, $K'$ is a flat $A$-module, and $w$ is therefore injective since $B_1$ is reduced, hence the homomorphism $B_1\to R$ is injective.
Finally, for the same reason, if $\mathfrak{R}$ is
\oldpage[0\textsubscript{IV-1}]{218}
the nilradical of $B$, the kernel of $v$ is $K'\otimes_A\mathfrak{R}$, hence it is contained in the nilradical of $K'\otimes_A B$; we immediately conclude that if the image by $h=w\circ v\circ u$ of an element $x\in A'\otimes_A B$ is nilpotent, then $x$ is nilpotent, which proves the lemma.
\end{proof}

\begin{proposition}[23.2.4]
\label{0.23.2.4}
Let $A$ be an integral ring, $K$ its field of fractions, $B$ a Noetherian ring, $\mathfrak{q}_i$ ($1\leq i\leq n$) its minimal prime ideals, $\vphi:A\to B$ a ring homomorphism.
Suppose that the $B/\mathfrak{q}_i$ are Japanese rings, and that $B$ is a flat $A$-module.
Let $K'$ be a finite extension of $K$, and let $(C_\alpha)$ be the increasing filtering family of subrings of $K'$ that are finite $A$-algebras and admit $K'$ as their field of fractions; the union $A'$ of the $C_\alpha$ is thus the integral closure of $A$ in $K'$.
Then:
\begin{enumerate}
  \item[{\rm(i)}] There exists an index $\alpha$ such that for $\lambda\geq\alpha$, the canonical homomorphism
    \[
    \label{0.23.2.4.1}
      (C_\alpha\otimes_A B)_\red\to(C_\lambda\otimes_A B)_\red
      \tag{23.2.4.1}
    \]
    is bijective.
  \item[{\rm(ii)}] If in addition $B$ is a faithfully flat $A$-module, the morphism $\Spec(A')\to\Spec(C_\alpha)$ is \emph{radiciel}.
\end{enumerate}
\end{proposition}

\begin{proof}
\medskip\noindent
(i).
Let $B_1=B_\red$, and let $R$ be the total ring of fractions of $B_1$; since $A$ is integral and $B$ is a flat $A$-module, every element $\neq 0$ of $A$ is $B$-regular \sref[0\textsubscript{I}]{0.6.3.4}, hence the composite homomorphism $A\to B\to B_1$ extends to a homomorphism $K\to R$. We can therefore write $K'\otimes_A R=(K'\otimes_A K)\otimes_K R$, and since $K'\otimes_A K$ is identified with $K'$, we have $K'\otimes_A R=K'\otimes_K R$. Since $R$ is the direct product of the fields of fractions $L_i$ of the $B/\mathfrak{q}_i$, $K'\otimes_A R$ is the direct product of the $K'\otimes_K L_i$, and hence $(K'\otimes_A R)_\red$ is the direct product of a finite number of finite extensions of the $L_i$. By hypothesis, the integral closure of $B/\mathfrak{q}_i$ in a finite extension of $L_i$ is a finite $(B/\mathfrak{q}_i)$-module, hence a finite $B$-module; it follows that the integral closure $B'$ of $B$ in $(K'\otimes_A R)_\red$ is a finite $B$-module.
By virtue of \sref{0.23.2.3}, the $(C_\lambda\otimes_A B)_\red$ are canonically identified with subrings of $(K'\otimes_A R)_\red$ which are \emph{finite} $B$-algebras, hence contained in $B'$.
Since $B$ is Noetherian and $B'$ a $B$-module of finite type, the filtering family of the $(C_\alpha\otimes_A B)_\red$ admits a \emph{greatest element} $(C_\alpha\otimes_A B)_\red$, which proves~(i).

\medskip\noindent
(ii).
Since $B$ is a faithfully flat $A$-module, it suffices \sref[IV]{IV.2.6.1}[v] to show that the morphism $\Spec(A'\otimes_A B)\to\Spec(C_\alpha\otimes_A B)$ is radiciel, or, what amounts to the same \sref[I]{I.5.1.6}, that the morphism $\Spec(A'\otimes_A B)_\red\to\Spec(C_\alpha\otimes_A B)_\red$ is radiciel; but we have $A'\otimes_A B=\varinjlim(C_\lambda\otimes_A B)$, hence \sref[IV]{IV.5.13.2} $(A'\otimes_A B)_\red=\varinjlim(C_\lambda\otimes_A B)_\red$, and the conclusion follows from~(i).
\end{proof}

\begin{corollary}[23.2.5]
\label{0.23.2.5}
Let $A$ be a local Noetherian integral ring, $K$ its field of fractions, $K'$ a finite extension of $K$.
Let $(C_\alpha)$ be the increasing filtering family of subrings of $K'$ that are finite $A$-algebras (hence Noetherian semi-local rings \textnormal{(Bourbaki, \emph{Alg.~comm.}, chap.~IV, \textsection2, n\textsuperscript{o}~5, cor.~3 of prop.~9)}) and admit $K'$ as their field of fractions.
Then there exists $\alpha$ such that the homomorphism $(\hat{C}_\alpha)_\red\to(\hat{C}_\lambda)_\red$ is an isomorphism for $\lambda\geq\alpha$, and if $A'$ is the integral closure of $A$ in $K'$, the morphism $\Spec(A')\to\Spec(C_\alpha)$ is radiciel \textnormal{(cf.~\sref{0.23.2.2})}.
\end{corollary}

\oldpage[0\textsubscript{IV-1}]{219}
\begin{proof}
We apply \sref{0.23.2.4} taking $B=\hat{A}$; since the $B/\mathfrak{q}_i$ are complete local Noetherian rings, they are Japanese rings \sref{0.23.1.5} and $B$ is a faithfully flat $A$-module \sref[0\textsubscript{I}]{0.7.3.5}; moreover we have $C_\lambda\otimes_A B=\hat{C}_\lambda$ (Bourbaki, \emph{Alg.~comm.}, chap.~III, \textsection3, n\textsuperscript{o}~4, th.~3 and chap.~IV, \textsection2, n\textsuperscript{o}~5, cor.~3 of prop.~9).
\end{proof}

\begin{corollary}[23.2.6]
\label{0.23.2.6}
Under the hypotheses of \sref{0.23.2.5}, the integral closure $A'$ of $A$ in $K'$ is a semi-local ring; if $\mathfrak{m}$ is the maximal ideal of $A$, then, for every maximal ideal $\mathfrak{m}'$ of $A'$, $A'/\mathfrak{m}'$ is a finite extension of $A/\mathfrak{m}$.
\end{corollary}

\begin{proof}
The fact that the homomorphism $(\hat{C}_\alpha)_\red\to(\hat{C}_\lambda)_\red$ is bijective for $\lambda\geq\alpha$ implies that the number of maximal ideals of $C_\lambda$ is \emph{constant} for $\lambda\geq\alpha$, and that if $\mathfrak{m}_\alpha$ is a maximal ideal of $C_\alpha$ and $\mathfrak{m}_\lambda$ is the unique maximal ideal of $C_\lambda$ above $\mathfrak{m}_\alpha$, the fields $C_\alpha/\mathfrak{m}_\alpha$ and $C_\lambda/\mathfrak{m}_\lambda$ are canonically isomorphic.
The conclusion follows from the fact that $\mathfrak{m}'=\varinjlim\mathfrak{m}_\lambda$ and $A'/\mathfrak{m}'=\varinjlim C_\lambda/\mathfrak{m}_\lambda$ \sref[0\textsubscript{III}]{0.10.3.1.3}.
\end{proof}

\begin{proposition}[23.2.7]
\label{0.23.2.7}
\textnormal{(Y.~Mori)}
Let $A$ be a local Noetherian integral ring, $K$ its field of fractions, and $A'$ the integral closure of $A$.
Then $A'$ is a Krull ring \textnormal{(Bourbaki, \emph{Alg.~comm.}, chap.~VII, \textsection1)} that is semi-local; in other words, there exists a family $(V_\lambda)$ of discrete valuation rings having $K$ as their field of fractions, such that $A'=\bigcap V_\lambda$ and that every $x\in K$ belongs to all but a finite number of the $V_\lambda$.
Moreover, there exists a subalgebra $C$ of $K$ which is a finite $A$-algebra, and such that the $V_\lambda$ are the integral closures of the rings $C_{\mathfrak{p}_\lambda}$, where $(\mathfrak{p}_\lambda)$ is the family of prime ideals of height $1$ of the local ring $C$.
\end{proposition}

We will first prove the following lemma:

\begin{lemma}[23.2.7.1]
\label{0.23.2.7.1}
Let $A$, $B$ be two rings, $\vphi:A\to B$ a homomorphism making $B$ a faithfully flat $A$-module.
Suppose that $A$ is integral; then, if $C=B_\red$, the composite homomorphism $\psi:A\to B\to B_\red=C$ is injective and extends to an injective homomorphism of the field of fractions $K$ of $A$ into the total ring of fractions $R$ of $C$; moreover, if $C'$ is the integral closure of $C$ in $R$, $C'\cap K$ is the integral closure of $A$.
\end{lemma}

\begin{proof}
Since $A$ is integral and $\vphi$ is injective, the intersection of $\vphi(A)$ and the nilradical $\mathfrak{R}$ of $B$ is $0$, and since $C=B/\mathfrak{R}$, $\psi$ is injective.
Every $a\neq 0$ in $A$ being a non-zero-divisor does not become a zero-divisor in $B$ by flatness \sref[0\textsubscript{I}]{0.6.3.4}; we deduce that $a$ is also a non-zero-divisor in $C$, since if $ax\in\mathfrak{R}$ for an $x\notin\mathfrak{R}$ in $B$, we would have $a^nx^n=0$ for some integer $n$, which contradicts the above since $x^n\neq 0$.
We can thus extend $\psi$ to an injective homomorphism of $K$ into $R$.
To prove the last assertion, note that the integral closure $A'$ of $A$ is contained in $C'\cap K$.
Conversely, let $x\in C'\cap K$; $x$ is then integral over $C$, and \emph{a fortiori} over $B$, that is to say $B[x]$ is a finite $B$-algebra; moreover, $K$ is identified with a subring of $B\otimes_A K$, and the subring $B[x]$ of $B\otimes_A K$ is identified with $B\otimes_A A[x]$ by flatness; we conclude that $A[x]$ is an $A$-module of finite type (Bourbaki, \emph{Alg.~comm.}, chap.~I\textsuperscript{er}, \textsection3, n\textsuperscript{o}~6, prop.~11), hence $x\in A'$.
\end{proof}

\begin{proof}[Proof of Proposition~{\rm\sref{0.23.2.7}}]
This lemma being established, let us apply it to the canonical injection of $A$ into its completion $\hat{A}$, which is a faithfully flat $A$-module; $B=(\hat{A})_\red=\hat{A}/\mathfrak{R}$ (where $\mathfrak{R}$ is the nilradical of $\hat{A}$) is a complete local Noetherian reduced ring, whose total ring of
\oldpage[0\textsubscript{IV-1}]{220}
fractions $R$ is a direct product of a finite number of fields $L_i$; the canonical image $B_i$ of $B$ in $L_i$ is a complete local Noetherian \emph{integral} ring, whose field of fractions is $L_i$, and the integral closure $B'$ of $B$ in $R$ is the \emph{direct product} of the $B'_i$, where $B'_i$ is the integral closure of $B_i$.
But by virtue of \sref{0.23.1.5}, $B'_i$ is a finite $B_i$-algebra, hence an integrally closed local Noetherian ring.
We know (Bourbaki, \emph{Alg.~comm.}, chap.~VII, \textsection1, n\textsuperscript{o}~3, cor.~of th.~2) that $B'_i$ is a Krull ring.
Now, for every $i$, there is a homomorphism $\vphi_i:K\to L_i$, which is injective, and so $\vphi_i^{-1}(B'_i)$ is a Krull ring in $K$.
But since $A'=B'\cap K$ by virtue of the lemma and $A'=\bigcap\vphi_i^{-1}(B'_i)$, $A'$ is an intersection of a finite number of Krull rings and is therefore a Krull ring.\footnote{The French source prints $B'=\bigcap_i\vphi_i^{-1}(B'_i)$ here; the required intersection inside $K$ is $A'=\bigcap_i\vphi_i^{-1}(B'_i)$.}
To prove the last assertion of the proposition, note that there exists a finite $A$-algebra $C\subset K$ such that the morphism $\Spec(A')\to\Spec(C)$ is radiciel and a homeomorphism \sref{0.23.2.4}; for every prime ideal $\mathfrak{p}'\in\Spec(A')$, $\mathfrak{p}'$ is thus the only prime ideal of $A'$ above $\mathfrak{p}=\mathfrak{p}'\cap C$ and $A'_{\mathfrak{p}'}$ the integral closure of $C_\mathfrak{p}$; on the other hand the fact that $\Spec(A')\to\Spec(C)$ is a homeomorphism implies that the map $\mathfrak{p}'\mapsto\mathfrak{p}'\cap C$ is a \emph{bijection} of the set of prime ideals of height~$1$ of $A'$ onto the set of prime ideals of height~$1$ of $C$; whence the conclusion, by virtue of Bourbaki, \emph{Alg.~comm.}, chap.~VII, \textsection1, n\textsuperscript{o}~6, th.~4.
\end{proof}

\begin{remarks}[23.2.8]
\label{0.23.2.8}
\medskip\noindent
\begin{enumerate}
  \item[{\rm(i)}] We should take care to note, in the application of \sref{0.23.2.6} and \sref{0.23.2.7}, that the ring $A'$ is \emph{not necessarily Noetherian}, as an example of Nagata shows with $\dim(A)=3$ \cite{IV-30}.
  \item[{\rm(ii)}] The conclusions of \sref{0.23.2.7} remain valid if $A'$ is the integral closure of $A$ in a finite extension $K'$ of $K$; it suffices in fact to consider a finite $A$-algebra $B$ whose field of fractions is $K'$; $B$ is a semi-local Noetherian ring, hence an intersection of a finite number of local Noetherian rings $B_{\mathfrak{m}_i}$ ($\mathfrak{m}_i$ the maximal ideals of $B$), and its integral closure (which is equal to $A'$) is the intersection of the $A'_{\mathfrak{m}_i}$ (Bourbaki, \emph{Alg.~comm.}, chap.~II, \textsection3, n\textsuperscript{o}~3, cor.~4 of th.~1) which are the integral closures of the $B_{\mathfrak{m}_i}$, hence Krull rings by \sref{0.23.2.7}; consequently $A'$ is a Krull ring.
  \item[{\rm(iii)}] One can show (\cite{IV-30}, 3.3.10) that for every Noetherian integral ring $A$ (not necessarily local), the integral closure of $A$ is a Krull ring.
\end{enumerate}
\end{remarks}

\begin{corollary}[23.2.9]
\label{0.23.2.9}
Let $A$ be a Noetherian integral ring, $A'$ its integral closure.
Suppose that for every ring $C$ such that $A\subset C\subset A'$ and such that $C$ is a finite $A$-algebra, and for every prime ideal $\mathfrak{q}$ of height~$1$ in $C$, $\mathfrak{q}\cap A$ is a prime ideal of height~$1$ in $A$.
Let $P$ be the set of prime ideals of height~$1$ in $A$; we then have
\[
  A'=\bigcap_{\mathfrak{p}\in P}A'_\mathfrak{p}.
\]
\end{corollary}

\begin{proof}
Since $A'$ is an $A$-module without torsion, we have $A'=\bigcap A'_\mathfrak{m}$, where $\mathfrak{m}$ ranges over the set of maximal ideals of $A$ (Bourbaki, \emph{Alg.~comm.}, chap.~II, \textsection3, n\textsuperscript{o}~3, cor.~4 of th.~1);
\oldpage[0\textsubscript{IV-1}]{221}
it will thus suffice to prove that $A'_\mathfrak{m}$ is the intersection of the $A'_\mathfrak{p}$ for the prime ideals $\mathfrak{p}\subset\mathfrak{m}$ of height~$1$, and since $A'_\mathfrak{m}$ is the integral closure of $A_\mathfrak{m}$, we see that it suffices to prove the corollary for $A_\mathfrak{m}$.
Now there is, by virtue of \sref{0.23.2.7}, a finite $A_\mathfrak{m}$-algebra $B$ contained in $A'_\mathfrak{m}$ such that $A'_\mathfrak{m}$ is the intersection of the integral closures of the rings $B_\mathfrak{q}$, where $\mathfrak{q}$ ranges over the set of prime ideals of height~$1$ of $B$.
But $B$ is of the form $C_\mathfrak{m}$, where $C$ is a finite $A$-algebra; hence, for every prime ideal $\mathfrak{q}$ of height~$1$ in $B$, $\mathfrak{q}\cap A$ is by hypothesis a prime ideal of height~$1$ in $A$; since the integral closure of $B_\mathfrak{q}$ contains that of $A_{\mathfrak{q}\cap A}$ and the latter contains $A'_\mathfrak{m}$, it follows that $A'_\mathfrak{m}$ is indeed the intersection of the $A'_\mathfrak{p}$ for $\mathfrak{p}\in P$ and $\mathfrak{p}\subset\mathfrak{m}$, which finishes the proof.
\end{proof}

\begin{remarks}[23.2.10]
\label{0.23.2.10}
\medskip\noindent
\begin{enumerate}
  \item[{\rm(i)}] We will see \sref[IV]{IV.5.6.3} and \sref[IV]{IV.5.6.10} that the hypothesis of \sref{0.23.2.9} is satisfied when $A$ is \emph{universally catenary}, in particular \sref[IV]{IV.5.10.17} and \sref[IV]{IV.5.11.2} when $A$ is a \emph{quotient of a regular ring}; finally, it is evidently satisfied if $\Spec(A')\to\Spec(A)$ is a homeomorphism, in particular if this morphism is radiciel \sref{0.23.2.2}.
  \item[{\rm(ii)}] The conclusion of \sref{0.23.2.9} is no longer necessarily exact when one omits the hypothesis on the finite sub-$A$-algebras of $A'$.
    Indeed, one can define a local Noetherian integral ring $A$, of dimension~$2$, whose integral closure $A'$ is a finite $A$-algebra, but such that there is a prime ideal $\mathfrak{p}'$ of $A'$ of height~$1$ above the maximal ideal $\mathfrak{m}$ (of height~$2$) of $A$, and such moreover that for every prime ideal $\mathfrak{p}$ of height~$1$ in $A$, $A_\mathfrak{p}$ is integrally closed \sref[IV]{IV.5.6.11}; the intersection of these rings $A_\mathfrak{p}$ can then not be equal to $A'$ since $A'$ is a Krull ring (Bourbaki, \emph{Alg.~comm.}, chap.~VII, \textsection1, n\textsuperscript{o}~5, prop.~9).
\end{enumerate}
\end{remarks}
% END NATIVE INPUT ega0/ega0-23

% Chapter IV begins with section 1 in this part.
\setcounter{section}{0}
% BEGIN NATIVE INPUT ega4-1.source_aligned_r1
\cleardoublepage
\oldpage[IV-1]{222}
\thispagestyle{plain}

\begin{center}
{\Large CHAPTER IV}

\vspace{1.5cm}
{\LARGE\bfseries LOCAL STUDY OF SCHEMES\\[0.3em]
AND MORPHISMS OF SCHEMES}

\vspace{1.5cm}
{\large\bfseries Summary}
\end{center}

\medskip
\begin{tabular}{@{}r p{0.84\textwidth}}
§ 1. & Relative finiteness conditions. Constructible sets in preschemes.\\
§ 2. & Base change and flatness.\\
§ 3. & Associated prime cycles and primary decompositions.\\
§ 4. & Change of the base field in preschemes.\\
§ 5. & Dimension and depth in preschemes.\\
§ 6. & Flat morphisms of locally Noetherian preschemes.\\
§ 7. & Application to relations between a Noetherian local ring and its completion.
Excellent rings.\\
§ 8. & Projective limits of preschemes.\\
§ 9. & Constructible properties.\\
§ 10. & Jacobson preschemes.\\
§ 11.\footnotemark[1] & Topological properties of flat morphisms of finite presentation.
Local criteria for flatness.\\
§ 12. & Study of the fibres of flat morphisms of finite presentation.\\
§ 13. & Equidimensional morphisms.\\
§ 14. & Universally open morphisms.\\
§ 15. & Study of the fibres of a universally open morphism.\\
§ 16. & Differential invariants. Differentially smooth morphisms.\\
§ 17. & Smooth, unramified, and étale morphisms.\\
§ 18. & Complements on étale morphisms. Henselian local rings.\\
§ 19. & Regular and transversally regular immersions.\\
§ 20. & Hyperplane sections; generic projections.\\
§ 21. & Infinitesimal extensions.\\
\end{tabular}
\footnotetext[1]{The order and content of §§ 11 to 21 are given only as an
indication and may be modified before their publication.}

\clearpage
\leavevmode\oldpage[IV-1]{223}

The subjects treated in this Chapter call for the following remarks:

\medskip
\noindent
\emph{a)} Their common feature is that they concern \emph{local} properties
of preschemes or morphisms, i.e. properties considered at a point, at the
points of a fibre, or in an unspecified neighbourhood of a point or a fibre.
These properties are generally \emph{topological}, \emph{differential}, or
\emph{dimensional} in nature (i.e. involving the notions of \emph{dimension}
and \emph{depth}), and are linked to properties of the \emph{local rings} at
the points under consideration. One typical problem is, for a given morphism
$f:X\to Y$ and a point $x\in X$, to relate the properties of $X$ at $x$ to
those of $Y$ at $y=f(x)$ and of the fibre $X_y=f^{-1}(y)$ at $x$. Another is
to determine the topological nature (for example, constructibility, or being
open or closed) of the set of points $x\in X$ at which $X$ has a given
property, or for which the fibre $X_{f(x)}$ passing through $x$ has a given
property at $x$. Similarly, we shall be interested in the topological nature
of the set of points $y\in Y$ such that $X$ has a given property at all
points of the fibre $X_y$, or such that this fibre has a given property.

\medskip
\noindent
\emph{b)} The most important notions for the following chapters are that of
a \emph{flat morphism of finite presentation} and those of a \emph{smooth
morphism} and an \emph{étale morphism}, which are special cases. Their
detailed study, together with related questions, properly begins in §~11.

\medskip
\noindent
\emph{c)} Sections 1 to 10 may be regarded as preliminary in nature and
develop three types of techniques, used not only in the other sections of
this chapter but, of course, also in the following chapters:

\medskip
\noindent
\emph{c 1)} Sections 1 to 4 consider various aspects of the notion of
\emph{base change}, especially in relation to finiteness or flatness
conditions; the technique of \emph{descent} is initiated there in its most
elementary aspects (questions of “effectivity” connected with this
technique will be studied in Chapter~V).

\medskip
\noindent
\emph{c 2)} Sections 5 to 7 centre on what may be called \emph{Noetherian
techniques}: the preschemes considered are always assumed locally Noetherian,
whereas, on the contrary, no finiteness condition is generally imposed on
the \emph{morphisms}. This is essentially because the notions of dimension
and depth are manageable only for Noetherian local rings. Recall that §~7
constitutes a “fine” theory of Noetherian local rings, used relatively
little in the remainder of the chapter.

\medskip
\noindent
\emph{c 3)} Sections 8 to 10 provide, among other things, a means of
\emph{eliminating Noetherian hypotheses} on the preschemes studied, by
replacing them with suitable \emph{finiteness} hypotheses (“finite
presentation”) on the \emph{morphisms} considered. The advantage of this
replacement is that the latter hypotheses are \emph{stable under base
change}, whereas this is not the case for Noetherian hypotheses on
preschemes. The technique permitting this passage rests, on the one hand, on
the use of the notion of a \emph{projective limit} of preschemes, by which a
question may be reduced to the same question under \emph{Noetherian}
hypotheses, and, on the other hand, on the systematic use of
\emph{constructible sets}, which have the double advantage of being
preserved by inverse image (under an arbitrary morphism)

\oldpage[IV-1]{224}

and by direct image (under a morphism of finite presentation), and of having
manageable topological properties in locally Noetherian preschemes. The same
techniques often even permit reduction to more particular Noetherian rings,
for example finitely generated $\mathbf Z$-algebras, and it is here that the
properties of “excellent” rings, studied in §~7, intervene decisively.
Quite apart from the elimination of Noetherian hypotheses, the very
elementary techniques of §§~8 to 10 are in constant use in almost all
applications.

\setcounter{section}{0}

\section{Relative finiteness conditions. Constructible sets in preschemes}
\label{section:IV.1}

In this section, we will resume and complete the expos\'e of ``finiteness conditions'' for a morphism of preschemes $f:X\to Y$ given in~(\textbf{I},~\hyperref[subsection:I.6.3]{6.3}~and~\hyperref[subsection:I.6.6]{6.6}).
There are essentially two notions of ``finiteness'' of a \emph{global} nature on $X$, that of \emph{quasi-compact} morphism (defined in~\sref[I]{I.6.6.1}) and that of a \emph{quasi-separated} morphism; on the other hand, there are two notions of ``finiteness'' of a \emph{local} nature on $X$, that of a morphism \emph{locally of finite type} (defined in~\sref[I]{I.6.6.2}) and that of a morphism \emph{locally of finite presentation}.
By combining these local notions with the preceding global notions, we obtain the notion of a morphism \emph{of finite type} (defined in~\sref[I]{I.6.3.1}) and of a morphism \emph{of finite presentation}.
For the convenience of the reader, we will give again in this section the principal properties stated in~(\textbf{I},~\hyperref[subsection:I.6.3]{6.3}~and~\hyperref[subsection:I.6.6]{6.6}), referring, of course, to those numbered passages in Chapter~I for their proofs.

In~n\textsuperscript{os}{\hyperlink{ega4.c04.subsec.1-8}{1.8}} and~{\hyperlink{ega4.c04.subsec.1-9}{1.9}}, we complete, in the context of preschemes, and making use of the previous notions of finiteness, the results on constructible sets given in~(\textbf{0}\textsubscript{III},~\textsection\hyperref[section:0.9]{9}).

\subsection{Quasi-compact morphisms}
\label{subsection:IV.1.1}

\begin{definition}[1.1.1]
\label{IV.1.1.1}
We say that a morphism of preschemes $f:X\to Y$ is \emph{quasi-compact} if the continuous map $f$ from the topological space $X$ to the topological space $Y$ is quasi-compact~\sref[0]{0.9.1.1}, in other words, if the inverse image $f^{-1}(U)$ of every quasi-compact open subset $U$ of $Y$ is quasi-compact~(cf.~\sref[I]{I.6.6.1}).
\end{definition}

If $\mathfrak{B}$ is a basis for the topology of $Y$ consisting of affine open sets, then for $f$ to be quasi-compact, it is necessary and sufficient that for all $V\in\mathfrak{B}$, $f^{-1}(V)$ is a \emph{finite union of affine open sets}.
For example, if $Y$ is affine and $X$ is quasi-compact, \emph{every} morphism $f:X\to Y$ is quasi-compact~\sref[I]{I.6.6.1}.

If $f:X\to Y$ is a quasi-compact morphism, then it is clear that for every open subset $V$ of $Y$, the restriction of $f$ to $f^{-1}(V)$ is a quasi-compact morphism $f^{-1}(V)\to V$.
Conversely, if $(U_\alpha)$ is an open cover of $Y$ and $f:X\to Y$ is a morphism such that the restrictions $f^{-1}(U_\alpha)\to U_\alpha$ are quasi-compact, then $f$ is quasi-compact.
As
\oldpage[IV-1]{225}
a result, if $f:X\to Y$ is an $S$-morphism of $S$-preschemes, and if there exists an open cover $(S_\lambda)$ of $S$ such that the restrictions $g^{-1}(S_\lambda)\to h^{-1}(S_\lambda)$ of $f$ (where $g$ and $h$ are the structure morphisms) are quasi-compact, then $f$ is quasi-compact.

\begin{proposition}[1.1.2]
\label{IV.1.1.2}
\medskip\noindent
\begin{enumerate}[(i)]
  \item An immersion $X\to Y$ is quasi-compact if it is closed,
    or if the underlying space of $Y$ is locally Noetherian or if the underlying space of $X$ is Noetherian.
  \item The composition of two quasi-compact morphisms is quasi-compact.
  \item If $f:X\to Y$ is a quasi-compact $S$-morphism, the same is true of $f_{(S')}:X_{(S')}\to Y_{(S')}$
    for every extension $g:S'\to S$ of the base prescheme.
  \item If $f:X\to X'$ and $g:Y\to Y'$ are two quasi-compact $S$-morphisms,
    \[
      f\times_S g: X\times_S Y \to X'\times_S Y'
    \]
    is quasi-compact.
  \item If the composition $g\circ f$ of two morphisms $f:X\to Y$, $g:Y\to Z$ is quasi-compact, and
    if either $g$ is separated or the underlying space of $X$ is locally Noetherian, then $f$ is quasi-compact.
  \item For a morphism $f$ to be quasi-compact, it is necessary and sufficient that $f_{\mathrm{red}}$ be so.
\end{enumerate}
\end{proposition}

For the proof, see \sref[I]{I.6.6.4}.
We note that the assertion~(vi) is also a
consequence of the following more general proposition:

\begin{proposition}[1.1.3]
\label{IV.1.1.3}
Let $f:X\to Y$, $g:Y\to Z$ be two morphisms. If $g\circ f$ is quasi-compact and $f$ is surjective, then $g$ is quasi-compact.
\end{proposition}

\begin{proof}
Indeed, if $V$ is a quasi-compact open subset of $Z$, $f^{-1}(g^{-1}(V))$ is quasi-compact
by hypothesis, and we have $g^{-1}(V)=f(f^{-1}(g^{-1}(V)))$, since $f$ is surjective; so $g^{-1}(V)$
is quasi-compact.
\end{proof}

\begin{corollary}[1.1.4]
\label{IV.1.1.4}
Let $f:X\to Y$, $g:Y'\to Y$ be two morphisms, $X'=X\times_Y Y'$,
$f'=f_{(Y')}:X'\to Y'$; suppose $g$ is quasi-compact and surjective \textnormal{(}resp.\ surjective\textnormal{)}. Then, for $f$ to be
quasi-compact \textnormal{(}resp.\ dominant\textnormal{)}, it suffices that $f'$ be so.
\end{corollary}

\begin{proof}
If $g':X'\to X$ is the canonical projection, $g'$ is a surjective morphism~\sref[I]{I.3.5.2},
and we have $f\circ g'=g\circ f'$; if $f'$ is quasi-compact (resp.\ dominant), the same is true of $g\circ f'$
since $g$ is quasi-compact~\sref{IV.1.1.2} (resp.\ surjective); as $f\circ g'$ is quasi-compact
(resp.\ dominant) and $g'$ is surjective, we deduce that $f$ is quasi-compact~\sref{IV.1.1.3}
(resp.\ dominant).
\end{proof}

\oldpage[IV-1]{226}

To abbreviate, we will call \emph{maximal point} of a prescheme $X$ every generic point
of an irreducible component of $X$; if $X=\Spec(A)$ is affine, the maximal points
of $X$ are the \emph{minimal prime ideals} of $A$.

\begin{proposition}[1.1.5]
\label{IV.1.1.5}
Let $f:X\to Y$ be a quasi-compact morphism. The following
conditions are equivalent:
\begin{enumerate}[a)]
  \item $f$ is dominant;
  \item for every maximal point $y$ of $Y$, we have $f^{-1}(y)\neq\varnothing$;
  \item for every maximal point $y$ of $Y$, $f^{-1}(y)$ contains a maximal point of $X$.
\end{enumerate}
\end{proposition}

The equivalence of \emph{a)} and \emph{c)} was proved in~\sref[I]{I.6.6.5}. It is clear that \emph{c)}
implies \emph{b)}; on the other hand, if $X'$ is an irreducible component of $X$ meeting $f^{-1}(y)$,
$f(X')$ is irreducible and contains $y$, so it is contained in the irreducible component
$Y'=\overline{\{y\}}$ of $Y$~\sref[0]{0.2.1.6}; if $x$ is the generic point of $X'$, we therefore necessarily have
$f(x)=y$~\sref[0]{0.2.1.5}; so \emph{b)} implies \emph{c)}.

\begin{proposition}[1.1.6]
\label{IV.1.1.6}
Let $f':X'\to Y$, $f'':X''\to Y$ be two morphisms, $X$ the coproduct prescheme
$X'\coprod X''$, $f:X\to Y$ the morphism equal to $f'$ on $X'$ and to $f''$ on $X''$. For $f$ to be
quasi-compact, it is necessary and sufficient that $f'$ and $f''$ be so.
\end{proposition}

This follows immediately from the definition.

\subsection{Quasi-separated morphisms}
\label{subsection:IV.1.2}

\begin{definition}[1.2.1]
\label{IV.1.2.1}
Let $X$, $Y$ be two preschemes. We say that a morphism $f:X\to Y$ is
\emph{quasi-separated} (or that $X$ is \emph{quasi-separated over} $Y$) if the diagonal morphism $\Delta_f=(1_X,1_X)_Y:X\to X\times_Y X$
is quasi-compact. We say that a prescheme $X$ is \emph{quasi-separated} if it is quasi-separated over $\Spec(\mathbf{Z})$.
\end{definition}

By definition, every separated morphism is quasi-separated, $\Delta_f$ being a closed immersion~\sref{IV.1.1.2}[(i)]; a \emph{scheme} $X$ is a quasi-separated prescheme, being separated over
$\Spec(\mathbf{Z})$~\sref[I]{I.5.4.1}.

\begin{proposition}[1.2.2]
\label{IV.1.2.2}
\medskip\noindent
\begin{enumerate}[(i)]
  \item Every monomorphism of preschemes $f:X\to Y$ \textnormal{(}in particular
    every immersion\textnormal{)} is quasi-separated.
  \item If $f:X\to Y$ and $g:Y\to Z$ are two quasi-separated morphisms, $g\circ f:X\to Z$ is
    quasi-separated.
  \item Let $X$, $Y$ be two $S$-preschemes, $f:X\to Y$ a quasi-separated $S$-morphism. Then,
    for every base change $g:S'\to S$, the morphism $f_{(S')}:X_{(S')}\to Y_{(S')}$ is quasi-separated.
  \item If $f:X\to Y$, $f':X'\to Y'$ are two quasi-separated $S$-morphisms, then
    \[
      f\times_S f': X\times_S X'\to Y\times_S Y'
    \]
    is quasi-separated.
  \item If $f:X\to Y$, $g:Y\to Z$ are two morphisms such that $g\circ f$ is quasi-separated, then $f$
    is quasi-separated.
  \item If $f:X\to Y$ is a quasi-separated morphism, $f_{\mathrm{red}}:X_{\mathrm{red}}\to Y_{\mathrm{red}}$ is quasi-separated.
\end{enumerate}
\end{proposition}

\begin{proof}
By virtue of~\sref[I]{I.5.5.12}, it suffices to prove~(i), (ii) and~(iii).

The assertion~(i) is trivial, every monomorphism being separated~\sref[I]{I.5.5.1}. To
prove~(iii), we reduce immediately to the case where $Y=S$~\sref[I]{I.3.3.11} and the assertion follows
from the fact that $\Delta_{f_{(S')}}=(\Delta_f)_{(S')}$~\sref[I]{I.5.3.4} and from~\sref{IV.1.1.2}[(iii)]. To prove~(ii), consider
the projections $p$ and $q$ of $X\times_Y X$ over $X$; if $i=(p,q)_Z$, we know that the diagram
\[
\xymatrix{
  X\times_Y X \ar[r]^-{i} \ar[d]_-{\pi} & X\times_Z X \ar[d]^-{f\times_Z f} \\
  Y \ar[r]_-{\Delta_g} & Y\times_Z Y
}
\]
\oldpage[IV-1]{227}
(where $\pi$ is the structure morphism) is commutative and identifies $X\times_Y X$ with the product,
as $(Y\times_Z Y)$-preschemes, of $Y$ and $X\times_Z X$~\sref[I]{I.5.3.5}. If $g$ is
quasi-separated, $\Delta_g$ is quasi-compact, and hence so is $i$~\sref{IV.1.1.2}[(iii)];
if in addition $f$ is quasi-separated, $\Delta_f$ is quasi-compact, and hence so is
$i\circ\Delta_f$~\sref{IV.1.1.2}[(ii)], which is equal to $\Delta_{g\circ f}$.
\end{proof}

\begin{corollary}[1.2.3]
\label{IV.1.2.3}
\medskip\noindent
\begin{enumerate}[(i)]
  \item Let $X$ be a quasi-separated prescheme. Then every morphism $f:X\to Y$
    is quasi-separated.
  \item Let $Y$ be a quasi-separated prescheme. For a morphism $f:X\to Y$ to be quasi-separated,
    it is necessary and sufficient that the prescheme $X$ be quasi-separated.
\end{enumerate}
\end{corollary}

This follows immediately from~\sref{IV.1.2.2}[(ii) and~(v)] applied to $X$, $Y$ and $\Spec(\mathbf{Z})$.

\begin{proposition}[1.2.4]
\label{IV.1.2.4}
Let $f:X\to Y$ be a morphism, $g:Y\to Z$ a quasi-separated morphism.
If $g\circ f$ is quasi-compact, then so is $f$.
\end{proposition}

\begin{proof}
We know that $f$ is the composite morphism $X\xrightarrow{\Gamma_f} X\times_Z Y\xrightarrow{\mathrm{pr}_2} Y$~\sref[I]{I.5.3.13}; on the other hand,
$\mathrm{pr}_2$ is identified with $(g\circ f)\times_Z 1_Y$~\sref[I]{I.3.3.4} and since $g\circ f$ is quasi-compact, so is
$\mathrm{pr}_2$~\sref{IV.1.1.2}[(iv)]. Finally, we have the commutative diagram
\[
\xymatrix{
  X \ar[r]^-{\Gamma_f} \ar[d]_-{f} & X\times_Z Y \ar[d]^-{f\times 1_Y} \\
  Y \ar[r]_-{\Delta_g} & Y\times_Z Y
}
\label{IV.1.2.4.1}
\tag{1.2.4.1}
\]
which identifies $X$ with the product, as $(Y\times_Z Y)$-preschemes, of $Y$ and
$X\times_Z Y$~\sref[I]{I.5.3.7}. As
by hypothesis $\Delta_g$ is a quasi-compact morphism, so is $\Gamma_f$~\sref{IV.1.1.2}[(iii)];
the conclusion then follows from~\sref{IV.1.1.2}[(ii)].
\end{proof}

\begin{proposition}[1.2.5]
\label{IV.1.2.5}
Let $f:X\to Y$, $g:Y'\to Y$ be two morphisms, $X'=X\times_Y Y'$,
$f'=f_{(Y')}:X'\to Y'$. If $g$ is quasi-compact and surjective and $f'$ quasi-separated, then $f$ is
quasi-separated.
\end{proposition}

\begin{proof}
We have indeed $(X'\times_{Y'} X')=(X\times_Y X)_{(Y')}$ and $\Delta_{f'}=(\Delta_f)_{(Y')}$~\sref[I]{I.5.3.4}; as the
projection $X'\times_{Y'} X'\to X\times_Y X$ is quasi-compact~\sref{IV.1.1.2}[(iii)] and surjective~\sref[I]{I.3.5.2},
we can, by virtue of~\sref[I]{I.3.3.11}, apply~\sref{IV.1.1.4}, which shows that since $\Delta_{f'}$ is
a quasi-compact morphism, the same is true of $\Delta_f$.
\end{proof}

\begin{proposition}[1.2.6]
\label{IV.1.2.6}
Let $f:X\to Y$ be a morphism, $(U_\alpha)$ a covering of $Y$
by open subsets such that the induced preschemes $U_\alpha$ are quasi-separated. For $f$ to be quasi-separated, it is necessary and sufficient that each of the induced preschemes on the open subsets $f^{-1}(U_\alpha)$ be
quasi-separated.
\end{proposition}

\begin{proof}
The inverse image in $X\times_Y X$ of $U_\alpha$ is $X_\alpha\times_{U_\alpha} X_\alpha$, writing $X_\alpha=f^{-1}(U_\alpha)$,
and the restriction $X_\alpha\to X_\alpha\times_{U_\alpha} X_\alpha$ of $\Delta_f$ is none other than $\Delta_{f_\alpha}$, denoting by $f_\alpha$ the restriction
$X_\alpha\to U_\alpha$ of $f$; by virtue of the definition~\sref{IV.1.2.1} and the local character of the notion of
quasi-compact morphism~\sref{IV.1.1.1}, for $f$ to be quasi-separated, it is necessary and sufficient that
each of the morphisms $f_\alpha$ be so. But since by hypothesis the morphism $U_\alpha\to\Spec(\mathbf{Z})$
is quasi-separated, saying that $f_\alpha$ is quasi-separated is equivalent to saying that the composite
$X_\alpha\xrightarrow{f_\alpha} U_\alpha\to\Spec(\mathbf{Z})$ is quasi-separated~\sref{IV.1.2.2}[(ii) and~(v)];
hence the proposition.
\end{proof}

\oldpage[IV-1]{228}

To verify that a morphism is quasi-separated, we are therefore reduced to giving
criteria for a prescheme to be quasi-separated:

\begin{proposition}[1.2.7]
\label{IV.1.2.7}
Let $X$ be a prescheme, $(U_\alpha)$ a covering of $X$ formed by quasi-compact open subsets. The following properties are equivalent:
\begin{enumerate}[a)]
  \item $X$ is quasi-separated.
  \item For every quasi-compact open subset $U$ of $X$, the canonical immersion $U\to X$ is quasi-compact
    (in other words, $U$ is retrocompact in $X$).
  \item[b')] The intersection of two quasi-compact open subsets of $X$ is quasi-compact.
  \item For every pair of indices $(\alpha,\beta)$, the intersection $U_\alpha\cap U_\beta$ is quasi-compact.
\end{enumerate}
\end{proposition}

\begin{proof}
The properties \emph{b)} and \emph{b')} are trivially equivalent by definition of a quasi-compact
morphism. As a quasi-compact open subset is a finite union of affine open subsets,
for two quasi-compact open subsets $U$, $V$ of $X$, $U\times_{\mathbf{Z}} V$ is a quasi-compact open subset
of $X\times_{\mathbf{Z}} X$~\sref[I]{I.3.2.7}, whose inverse image under $\Delta_X$ is $U\cap V$, so \emph{a)} implies \emph{b'}).
It is trivial that \emph{b')} implies \emph{c)}; finally, if \emph{c)} is satisfied, the $U_\alpha\times_{\mathbf{Z}} U_\beta$ form a
covering of $X\times_{\mathbf{Z}} X$ by quasi-compact open subsets, and we know that for the
morphism $\Delta_X:X\to X\times_{\mathbf{Z}} X$ to be quasi-compact, it suffices that the inverse images
$U_\alpha\cap U_\beta$ under $\Delta_X$ be quasi-compact~\sref{IV.1.1.1}; so \emph{c)} implies \emph{a)}.
\end{proof}

\begin{corollary}[1.2.8]
\label{IV.1.2.8}
Every prescheme $X$ whose underlying space is locally Noetherian
\textnormal{(}for example a locally Noetherian prescheme\textnormal{)} is quasi-separated; every morphism $f:X\to Y$ is
then quasi-separated.
\end{corollary}

\begin{proposition}[1.2.9]
\label{IV.1.2.9}
Let $X'$, $X''$ be two preschemes, $X=X'\coprod X''$ their coproduct prescheme,
$f':X'\to Y$, $f'':X''\to Y$ two morphisms, $f:X\to Y$ the morphism which coincides with $f'$ on $X'$
and with $f''$ on $X''$. For $f$ to be quasi-separated, it is necessary and sufficient that $f'$ and $f''$ be so.
\end{proposition}

\begin{proof}
Indeed, $X\times_Y X$ is the coproduct of the four preschemes $X'\times_Y X'$, $X'\times_Y X''$, $X''\times_Y X'$
and $X''\times_Y X''$, and $\Delta_f$ is the morphism which coincides with $\Delta_{f'}$ on $X'$ and with $\Delta_{f''}$ on $X''$;
the proposition then follows immediately from the definitions.
\end{proof}

\subsection{Morphisms locally of finite type}
\label{subsection:IV.1.3}

\begin{env}[1.3.1]
\label{IV.1.3.1}
Recall that if $A$ is a ring, a (commutative) $A$-algebra $B$ is
said to be \emph{of finite type} if it is generated by a finite number of elements of $B$, or, what amounts to
the same thing, if it is isomorphic to a quotient of a polynomial algebra $A[T_1,\ldots,T_n]$
by an ideal of this algebra. It is clear that for every (commutative) $A$-algebra $A'$,
$B\otimes_A A'$ is then an $A'$-algebra of finite type. If $B$ is an $A$-algebra of finite type and $C$
a $B$-algebra of finite type, $C$ is an $A$-algebra of finite type, for if $C$ is a quotient
of $B[T_1,\ldots,T_n]$ and $B$ a quotient of $A[S_1,\ldots,S_m]$, $B[T_1,\ldots,T_n]$ is isomorphic to $A[S_1,\ldots,S_m,T_1,\ldots,T_n]/\mathfrak{a}$,
where $\mathfrak{a}$ is the kernel of the surjection $A[S_1,\ldots,S_m,T_1,\ldots,T_n]\to B[T_1,\ldots,T_n]$; hence $C$ is also a quotient of
$A[S_1,\ldots,S_m,T_1,\ldots,T_n]$, and so it is indeed an $A$-algebra of finite type.
\end{env}

\begin{definition}[1.3.2]
\label{IV.1.3.2}
Let $f:X\to Y$ be a morphism of preschemes, $x$ a point of $X$,
$y=f(x)$. We say that $f$ is \emph{of finite type at $x$} if there exists an affine open neighbourhood $V$ of $y$ and an affine open
neighbourhood $U$ of $x$ such that $f(U)\subset V$ and the ring $A(U)$ is an $A(V)$-algebra of finite type.
We say that $f$ is \emph{locally of finite type} if it is of finite type at every point of $X$.
\end{definition}

\oldpage[IV-1]{229}

\begin{proposition}[1.3.3]
\label{IV.1.3.3}
If $Y$ is a locally Noetherian prescheme and $f:X\to Y$ a
morphism locally of finite type, then $X$ is locally Noetherian.
\end{proposition}

For the proof, see~\sref[I]{I.6.3.7}.

\begin{proposition}[1.3.4]
\label{IV.1.3.4}
\medskip\noindent
\begin{enumerate}[(i)]
  \item Every local immersion is locally of finite type.
  \item If two morphisms $f:X\to Y$, $g:Y\to Z$ are locally of finite type, then so is $g\circ f$.
  \item If $f:X\to Y$ is an $S$-morphism locally of finite type, $f_{(S')}:X_{(S')}\to Y_{(S')}$ is
    locally of finite type for every extension $S'\to S$ of the base prescheme.
  \item If $f:X\to X'$ and $g:Y\to Y'$ are two $S$-morphisms locally of finite type, $f\times_S g$
    is locally of finite type.
  \item If the composition $g\circ f$ of two morphisms is locally of finite type, $f$ is locally of
    finite type.
  \item If a morphism $f$ is locally of finite type, the same is true of $f_{\mathrm{red}}$.
\end{enumerate}
\end{proposition}

For the proof, see~\sref[I]{I.6.6.6}.

\begin{corollary}[1.3.5]
\label{IV.1.3.5}
Let $f:X\to Y$ be a morphism locally of finite type. For every
morphism $Y'\to Y$ such that $Y'$ is locally Noetherian, $X\times_Y Y'$ is locally Noetherian.
\end{corollary}

\begin{proof}
Indeed, $f_{(Y')}:X\times_Y Y'\to Y'$ is locally of finite type by~\sref{IV.1.3.4}[(iii)], and it
suffices to apply~\sref{IV.1.3.3}.
\end{proof}

\begin{proposition}[1.3.6]
\label{IV.1.3.6}
Let $\varphi:A\to B$ be a homomorphism of rings. For the
corresponding morphism $f:\Spec(B)\to\Spec(A)$ to be locally of finite type, it is necessary and sufficient
that $B$ be an $A$-algebra of finite type.
\end{proposition}

For the proof, see~\sref[I]{I.6.3.3}, taking into account that $f$ is quasi-compact and using~\sref[I]{I.6.6.3}.

\begin{proposition}[1.3.7]
\label{IV.1.3.7}
For a morphism locally of finite type $f:X\to Y$ to be
surjective, it is necessary and sufficient that, for every algebraically closed field $\Omega$, the map $X(\Omega)\to Y(\Omega)$
corresponding to $f$~\sref[I]{I.3.4.1} be surjective.
\end{proposition}

For the proof, see~\sref[I]{I.6.3.10}.

\begin{env}[1.3.8]
\label{IV.1.3.8}
Let $A$ be a ring. We say that an $A$-algebra $B$ is \emph{essentially of finite type}
if $B$ is $A$-isomorphic to an $A$-algebra of the form $S^{-1}C$, where $C$ is an $A$-algebra \emph{of
finite type} and $S$ a multiplicative subset of $C$.
\end{env}

\begin{proposition}[1.3.9]
\label{IV.1.3.9}
\medskip\noindent
\begin{enumerate}[(i)]
  \item If $B$ is an $A$-algebra essentially of finite type and $C$
    a $B$-algebra essentially of finite type, then $C$ is an $A$-algebra essentially of finite type.
  \item Let $B$ be an $A$-algebra essentially of finite type, $A'$ an $A$-algebra; then $B'=B\otimes_A A'$
    is an $A'$-algebra essentially of finite type.
\end{enumerate}
\end{proposition}

\begin{proof}
\begin{enumerate}[(i)]
  \item Let $B=S'^{-1}B'$ and $C=T'^{-1}C'$, where $B'$ (resp.\ $C'$) is an $A$-algebra (resp.\
    $B$-algebra) of finite type and $S'$ (resp.\ $T'$) a multiplicative subset of $B'$ (resp.\ $C'$);
    then $C'$ is of the form $S'^{-1}C''$, where $C''$ is also a $B'$-algebra of finite type, so $C$ is also
    a ring of fractions of $C''$; as $C''$ is an $A$-algebra of finite type, this proves
    our assertion.
  \item If $B$ is a ring of fractions of an $A$-algebra of finite type $C$, $B'$ is a ring
    of fractions of the $A'$-algebra of finite type $C\otimes_A A'$, which proves~(ii).
\end{enumerate}
\end{proof}

\begin{env}[1.3.10]
\label{IV.1.3.10}
If $B$ is a \emph{local} $A$-algebra essentially of finite type, $B$ is of the
form $C_{\mathfrak{q}}$, where $C$ is an $A$-algebra of finite type and $\mathfrak{q}$ a prime ideal of $C$ (Bourbaki,
\emph{Alg.\ comm.}, chap.~II, \S\,5, prop.~11). Let $\mathfrak{p}$ be the inverse image in $A$ of the ideal $\mathfrak{q}$; if we set
$S=A-\mathfrak{p}$, $C_{\mathfrak{q}}$ is also the localization of $S^{-1}C$ at a prime ideal; as $S^{-1}C$
is an algebra of finite type over $A_{\mathfrak{p}}=S^{-1}A$, we see that $B$ is also an $A_{\mathfrak{p}}$-algebra \emph{essentially
of finite type}, and furthermore the homomorphism $A_{\mathfrak{p}}\to B$ is \emph{local}.
\end{env}

\begin{proposition}[1.3.11]
\label{IV.1.3.11}
If $B$ is a local $A$-algebra essentially of finite type, $B$ is
$A$-isomorphic to a quotient of an $A$-algebra of the form $C_{\mathfrak{q}}$, where
$C=A[T_1,\ldots,T_n]$ is a polynomial algebra over $A$ and $\mathfrak{q}$ a prime ideal of $C$.
\end{proposition}

\begin{proof}
Indeed, by definition, $B$ is isomorphic to $C'_{\mathfrak{q}'}$, where $C'$ is an $A$-algebra of finite type
and $\mathfrak{q}'$ a prime ideal of $C'$; but $C'=C/\mathfrak{b}$, where $C=A[T_1,\ldots,T_n]$ and $\mathfrak{b}$ is an
ideal of $C$; so $\mathfrak{q}'=\mathfrak{q}/\mathfrak{b}$, where $\mathfrak{q}$ is a prime ideal of $C$, and $C'_{\mathfrak{q}'}$ is isomorphic to $C_{\mathfrak{q}}/\mathfrak{b} C_{\mathfrak{q}}$.
\end{proof}

\subsection{Morphisms locally of finite presentation}
\label{subsection:IV.1.4}

\begin{env}[1.4.1]
\label{IV.1.4.1}
\oldpage[IV-1]{230}
Given a ring $A$, we say that an $A$-algebra $B$ (commutative) is \emph{of finite presentation} if it is isomorphic to a quotient $A[T_1,\ldots,T_n]/\mathfrak{a}$ of a polynomial algebra over $A$ by an ideal \emph{of finite type} $\mathfrak{a}$ of $A[T_1,\ldots,T_n]$.
For every (commutative) $A$-algebra $A'$, $B\otimes_A A'$ is then an $A'$-algebra of finite presentation, being isomorphic to the quotient of $A'[T_1,\ldots,T_n]$ by the image of the ideal $\mathfrak{a}\otimes_A A'$ in this ring, which is evidently an $A'[T_1,\ldots,T_n]$-module of finite type.
If $B$ is an $A$-algebra of finite presentation and $C$ a $B$-algebra of finite presentation, then $C$ is an $A$-algebra of finite presentation.
Indeed, let $B = A[S_1,\ldots,S_m]/\mathfrak{a}$ and $C = B[T_1,\ldots,T_n]/\mathfrak{b}$, where $\mathfrak{a}$ (resp.~$\mathfrak{b}$) is an ideal of finite type of $A[S_1,\ldots,S_m]$ (resp.~$B[T_1,\ldots,T_n]$); the ring $B[T_1,\ldots,T_n]$ is isomorphic to $A[S_1,\ldots,S_m,T_1,\ldots,T_n]/\mathfrak{a}'$, where $\mathfrak{a}'$ is the canonical image of $\mathfrak{a}\otimes_A A[T_1,\ldots,T_n]$, and is therefore an ideal of finite type of $A[S_1,\ldots,S_m,T_1,\ldots,T_n]$; the ideal $\mathfrak{b}$ of $B[T_1,\ldots,T_n]$ is of the form $\mathfrak{b}'/\mathfrak{a}'$, where $\mathfrak{b}'$ is an ideal of $A[S_1,\ldots,S_m,T_1,\ldots,T_n]$; as $\mathfrak{a}'$ and $\mathfrak{b}'/\mathfrak{a}'$ are modules of finite type over $A[S_1,\ldots,S_m,T_1,\ldots,T_n]$, so is $\mathfrak{b}'$, and $C$, being isomorphic to $A[S_1,\ldots,S_m,T_1,\ldots,T_n]/\mathfrak{b}'$, is therefore an $A$-algebra of finite presentation.
We note finally that if $A$ is \emph{Noetherian}, then so is $A[T_1,\ldots,T_n]$; hence every $A$-algebra of finite type is an algebra of finite presentation.
\end{env}

\begin{definition}[1.4.2]
\label{IV.1.4.2}
Let $f:X\to Y$ be a morphism of preschemes, $x$ a point of $X$, $y=f(x)$.
We say that $f$ is \emph{of finite presentation at $x$} if there exists an affine open neighbourhood $V$ of $y$ and an affine open neighbourhood $U$ of $x$ such that $f(U)\subset V$ and such that the ring $\sh{O}(U)$ is an $\sh{O}(V)$-algebra of finite presentation.
We say that $f$ is \emph{locally of finite presentation} if it is of finite presentation at every point of $X$.
\end{definition}

If $Y$ is locally Noetherian, it amounts to the same to say that $f$ is locally of finite type or locally of finite presentation.

\begin{proposition}[1.4.3]
\label{IV.1.4.3}
\begin{enumerate}
  \item[\rm(i)] Every local isomorphism is locally of finite presentation.
  \item[\rm(ii)] If two morphisms $f:X\to Y$, $g:Y\to Z$ are locally of finite presentation, so is $g\circ f$.
\end{enumerate}
\end{proposition}

\oldpage[IV-1]{231}
\begin{enumerate}
  \item[\rm(iii)] If $f:X\to Y$ is an $S$-morphism locally of finite presentation, $f_{(S')}:X_{(S')}\to Y_{(S')}$ is locally of finite presentation for every base extension $S'\to S$ of the base prescheme.
  \item[\rm(iv)] If $f:X\to X'$ and $g:Y\to Y'$ are two $S$-morphisms locally of finite presentation, $f\times_S g$ is locally of finite presentation.
  \item[\rm(v)] If the composition $g\circ f$ of two morphisms is locally of finite presentation and if $g$ is locally of finite type, then $f$ is locally of finite presentation.
\end{enumerate}

\begin{proof}
Assertion~{\hyperlink{ega4.c04.prop.1-4-3}{(i)}} is trivial.
To prove~{\hyperlink{ega4.c04.prop.1-4-3}{(iii)}}, we may limit to the case where $S=Y$ \sref[I]{I.3.3.11}; let $x'$ be a point of $X_{(S')}$, $s'$ and $x$ its projections on $S'$ and $X$ respectively, $s$ the common projection of $s'$ and $x$ on $S$.
By hypothesis, there exists an affine open neighbourhood $V$ (resp.~$U$) of $s$ (resp.~$x$) such that the image of $U$ is contained in $V$ and such that $\sh{O}(U)$ is an $\sh{O}(V)$-algebra of finite presentation; let $V'$ be an affine open neighbourhood of $s'$ whose image in $S$ is contained in $V$; then $U\times_V V'$ is an affine open neighbourhood of $x'$ whose image in $S'$ is contained in $V'$ and whose ring $\sh{O}(U)\otimes_{\sh{O}(V)}\sh{O}(V')$ is an $\sh{O}(V')$-algebra of finite presentation~\sref{IV.1.4.1}.

To prove~{\hyperlink{ega4.c04.prop.1-4-3}{(ii)}}, consider a point $x\in X$, and let $y=f(x)$, $z=g(f(x))$.
There is by hypothesis an affine open neighbourhood $W=\Spec(C)$ (resp.~$V=\Spec(B)$) of $z$ (resp.~$y$) such that $g(V)\subset W$ and such that $B$ is a $C$-algebra of finite presentation.
On the other hand, there is an affine open neighbourhood $V'=\Spec(B')$ (resp.~$U=\Spec(A)$) of $y$ (resp.~$x$) such that $f(U)\subset V'$ and such that $A$ is a $B'$-algebra of finite presentation.
Let $s\in B$ be such that the open set $D(s)=\Spec(B_s)$ in $V$ is a neighbourhood of $y$ contained in $V'$, and let $U'=f^{-1}(D(s))\subset U$; we have $U'=\Spec(A\otimes_{B'}B_s)$, and $A\otimes_{B'}B_s$ is a $B_s$-algebra of finite presentation~\sref{IV.1.4.1}; on the other hand $B_s=B[T]/(1-sT)$ is a $B$-algebra of finite presentation by definition; hence by~\sref{IV.1.4.1}, $A\otimes_{B'}B_s$ is a $C$-algebra of finite presentation, which proves~{\hyperlink{ega4.c04.prop.1-4-3}{(ii)}}.

We know that~{\hyperlink{ega4.c04.prop.1-4-3}{(iv)}} follows from~{\hyperlink{ega4.c04.prop.1-4-3}{(i)}}, {\hyperlink{ega4.c04.prop.1-4-3}{(ii)}}, and~{\hyperlink{ega4.c04.prop.1-4-3}{(iii)}} \sref[I]{I.3.5.1}.
To prove~{\hyperlink{ega4.c04.prop.1-4-3}{(v)}}, consider the commutative diagram
\[
\xymatrix{
  X \ar[r]^{\Gamma_f} & X\times_Z Y \ar[d]^{f\times 1} \\
  Y \ar[r]_{\Delta_g} & Y\times_Z Y
}
\]
which identifies $X$ with the product of the $(Y\times_Z Y)$-preschemes $Y$ and $X\times_Z Y$ \sref[I]{I.5.3.7}, $\Delta_g$ being the diagonal morphism.
If we know that $\Delta_g$ is locally of finite presentation, it will follow by~{\hyperlink{ega4.c04.prop.1-4-3}{(iii)}} that so is $\Gamma_f$.
On the other hand, we have the factorisation of $f$ as $X\xrightarrow{\Gamma_f} X\times_Z Y\xrightarrow{p_2} Y$ \sref[I]{I.5.3.13}, where the projection $p_2$ is equal to $(g\circ f)\times_Z 1_Y$; as $g\circ f$ is assumed locally of finite presentation, so is $p_2$ by~{\hyperlink{ega4.c04.prop.1-4-3}{(iv)}}, and it will follow finally from~{\hyperlink{ega4.c04.prop.1-4-3}{(ii)}} that $f$ is also locally of finite presentation.
We thus see that we are reduced to proving the following.
\end{proof}

\begin{corollary}[1.4.3.1]
\label{IV.1.4.3.1}
Let $g:Y\to Z$ be a morphism locally of finite type; then the diagonal morphism $\Delta_g:Y\to Y\times_Z Y$ is locally of finite presentation.
\end{corollary}

\oldpage[IV-1]{232}
\begin{proof}
We may reduce to the case where $Z=\Spec(A)$, $Y=\Spec(B)$, $B$ being an $A$-algebra of finite type; the morphism $\Delta_g$ then corresponds to the \emph{augmentation homomorphism} $\delta:B\otimes_A B\to B$ such that $\delta(x\otimes y)=xy$.
In view of the definition~\sref{IV.1.4.2}, it suffices to show that the kernel $\mathfrak{J}$ of $\delta$ is an ideal \emph{of finite type}, and this follows from \sref[0]{0.20.4.4}.
\end{proof}

\begin{proposition}[1.4.4]
\label{IV.1.4.4}
Let $A$ be a ring, $B$ an $A$-algebra, $B'$ a polynomial algebra $A[T_1,\ldots,T_n]$, $u:B'\to B$ a surjective homomorphism of $A$-algebras.
For $B$ to be an $A$-algebra of finite presentation, it is necessary and sufficient that the kernel $\mathfrak{J}$ of $u$ be an ideal of finite type of $B'$.
\end{proposition}

\begin{proof}
The condition is sufficient by definition~\sref{IV.1.4.1}.
To see that it is necessary, we remark that the morphism $g:\Spec(B')\to\Spec(A)$ is locally of finite type; if $B$ is an $A$-algebra of finite presentation, the morphism $f:\Spec(B)\to\Spec(B')$ corresponding to $u$ is such that $g\circ f:\Spec(B)\to\Spec(A)$ is locally of finite presentation, and so it follows from~\sref{IV.1.4.3}[(v)] that $f$ is also locally of finite presentation.
Now, to show that the ideal $\mathfrak{J}$ is of finite type, it suffices to prove that $\widetilde{\mathfrak{J}}$ is a $\widetilde{B'}$-module of finite type \sref[II]{II.6.1.4.1}.
Returning to the definition~\sref{IV.1.4.2}, we are finally reduced to proving the following.
\end{proof}

\begin{lemma}[1.4.4.1]
\label{IV.1.4.4.1}
Let $\mathfrak{a}$ be an ideal of a ring $C$, $s$ an element of $C$, such that $C_s/\mathfrak{a}_s$ is a $C$-algebra of finite presentation; then $\mathfrak{a}_s$ is an ideal of finite type in the ring $C_s$.
\end{lemma}

\begin{proof}
By hypothesis, there exists a polynomial $C$-algebra $C'=C[T_1,\ldots,T_n]$ and a surjective $C$-homomorphism $v:C'\to C_s/\mathfrak{a}_s$ whose kernel $\mathfrak{b}$ is of finite type.
Let $p:C_s\to C_s/\mathfrak{a}_s$ be the canonical homomorphism; for each $i$, there is a $t_i\in C$ such that $p(t_i/s^k)=v(T_i)$ (we may assume that the exponent $k$ is the same for all indices $i$).
Consider then the $C_s$-algebra $D=C_s[T_1,\ldots,T_n]$, and the $C_s$-algebra homomorphism $w:D\to C_s$ such that $w(T_i)=t_i/s^k$; $w$ is evidently surjective, and so is $v'=p\circ w:D\to C_s/\mathfrak{a}_s$.
On the other hand, every polynomial $P$ of $D$ can be written as $P/s^m$, where $P\in C'=C[T_1,\ldots,T_n]$, and we have $v'(P/s^m)=(1/s^m)v(P)$; as $1/s$ is invertible in $C_s$, the relation $v'(P/s^m)=0$ in $C_s/\mathfrak{a}_s$ is equivalent to $v(P)=0$, and hence the kernel $\mathfrak{b}'$ of $v'$ is generated by the image of $\mathfrak{b}$ in the ring $D$, and \emph{a fortiori} is \emph{of finite type} in $D$.
But $\mathfrak{b}'=v'^{-1}(0)=w^{-1}(p^{-1}(0))=w^{-1}(\mathfrak{a}_s)$, and as $w$ is surjective, $\mathfrak{a}_s=w(\mathfrak{b}')$, which proves that $\mathfrak{a}_s$ is an ideal of finite type of $C_s$.
\end{proof}

\begin{corollary}[1.4.5]
\label{IV.1.4.5}
Let $X$, $Y$ be two preschemes, $j:X\to Y$ an immersion, $U$ an open subset of $Y$ such that $j(X)$ is closed in $U$, $\sh{J}$ the quasi-coherent ideal of $\sh{O}_U$ that defines the closed subprescheme of $Y$ associated to $j$.
For $j$ to be locally of finite presentation, it is necessary and sufficient that $\sh{J}$ be an $\sh{O}_U$-module of finite type.
\end{corollary}

\begin{proposition}[1.4.6]
\label{IV.1.4.6}
Let $\varphi:A\to B$ be a homomorphism of rings.
For the corresponding morphism $f:\Spec(B)\to\Spec(A)$ to be locally of finite presentation, it is necessary and sufficient that $B$ be an $A$-algebra of finite presentation.
\end{proposition}

\begin{proof}
The condition being trivially sufficient, it remains to prove that it is necessary.
If $f$ is locally of finite presentation, it already follows from~\sref{IV.1.3.6} that $B$ is an $A$-algebra
\oldpage[IV-1]{233}
of finite type, so there exists a surjective homomorphism $u:B'=A[T_1,\ldots,T_n]\to B$ of $A$-algebras; it follows from~\sref{IV.1.4.3}[(v)] that the immersion $j:\Spec(B)\to\Spec(B')$ is locally of finite presentation; hence by~\sref{IV.1.4.5}, if $\mathfrak{J}$ is the kernel of $u$, the $\widetilde{B'}$-module $\widetilde{\mathfrak{J}}$ is of finite type, and hence by \sref[II]{II.6.1.4.1} $\mathfrak{J}$ is an ideal of finite type.
\end{proof}

\begin{proposition}[1.4.7]
\label{IV.1.4.7}
Let $\rho:A\to B$ be a homomorphism of rings making $B$ a finite $A$-algebra.
For $B$ to be an $A$-algebra of finite presentation, it is necessary and sufficient that $B$ be an $A$-module of finite presentation.
\end{proposition}

We will use the following.

\begin{lemma}[1.4.7.1]
\label{IV.1.4.7.1}
If $B$ is a finite $A$-algebra, there exists a surjective homomorphism of $A$-algebras $u:B'\to B$, where $B'$ is a finite $A$-algebra of finite presentation and a free $A$-module.
\end{lemma}

\begin{proof}
Indeed, there is a finite system $(a_i)_{1\leq i\leq m}$ of generators of the $A$-algebra $B$ such that each $a_i$ satisfies a relation $F_i(a_i)=0$, where $F_i\in A[X]$ is a monic polynomial of degree $>0$; the $A$-algebra quotient $B'_i=A[X]/(F_i)$ is therefore free of finite rank over $A$; let $c_i$ be the class of $X$ in $B'_i$.
There exists an $A$-homomorphism of algebras $u_i:B'_i\to B$ such that $u_i(c_i)=a_i$; it then suffices to take for $B'$ the tensor product $B'_1\otimes_A B'_2\otimes\ldots\otimes_A B'_m$ and to consider the homomorphism $u_1\otimes u_2\otimes\ldots\otimes u_m:B'\to B$, which is surjective by construction.
Furthermore, if $B''=A[T_1,\ldots,T_m]$ and if $b'_i$ denotes the canonical image of $c_i$ in $B'$, the $A$-homomorphism of algebras $v:B''\to B'$ such that $v(T_i)=b'_i$ ($1\leq i\leq m$) is surjective, and its kernel $\mathfrak{b}'$ is the ideal of $B''$ generated by the $F_i(T_i)$, and therefore of finite type; this proves that $B'$ is an $A$-algebra of finite presentation.
\end{proof}

\begin{proof}[Proof of {\rm\sref{IV.1.4.7}}]
This lemma being proved, let us keep the same notation, and let $w=u\circ v:B''\to B$.\footnote{The French source prints $v\circ u$; with the displayed domains and codomains, the typed composite is $u\circ v$.}
Let $\mathfrak{b}$ (resp.~$\mathfrak{a}$) be the kernel of $w$ (resp.~$u$); we have $\mathfrak{a}=v(\mathfrak{b})$ since $v$ is surjective.
Now by~\sref{IV.1.4.4}, if $B$ is an $A$-algebra of finite presentation, $\mathfrak{b}$ is an ideal of finite type of $B''$, so $\mathfrak{a}$ is an ideal of finite type of $B'$, hence an $A$-module of finite type since $B'$ is a finite $A$-algebra.
As $B'$ is a free $A$-module, $B$ is an $A$-module of finite presentation.
Conversely, if $B$ is an $A$-module of finite presentation,\footnote{The French source prints $B'$ here; the proposition's hypothesis and the exact sequence $0\to\mathfrak a\to B'\to B\to0$ require $B$.} $\mathfrak{a}$ is an $A$-module of finite type (Bourbaki, \emph{Alg.~comm.}, chap.~I, \S\,2, n\textsuperscript{o}\,8, lemme~9) and \emph{a fortiori} an ideal of finite type of $B'$; hence, $B$ is by definition a $B'$-algebra of finite presentation, and as $B'$ is an $A$-algebra of finite presentation, $B$ is an $A$-algebra of finite presentation.
\end{proof}

\subsection{Morphisms of finite type}
\label{subsection:IV.1.5}

\begin{proposition}[1.5.1]
\label{IV.1.5.1}
Let $f:X\to Y$ be a morphism, $(U_\alpha)$ a covering of $Y$ by open affines.
The following conditions are equivalent:
\begin{enumerate}
  \item[\rm a)] $f$ is locally of finite type and quasi-compact.
  \item[\rm b)] For all $\alpha$, $f^{-1}(U_\alpha)$ is a finite union of affine open sets $V_{\alpha i}$ such that each ring $\sh{O}(V_{\alpha i})$ is an $\sh{O}(U_\alpha)$-algebra of finite type.
  \item[\rm c)] For every affine open subset $U$ of $Y$, $f^{-1}(U)$ is a finite union of affine open sets $V_j$ such that the rings $\sh{O}(V_j)$ are $\sh{O}(U)$-algebras of finite type.
\end{enumerate}
\end{proposition}

\begin{proof}
For the proof, see \sref[I]{I.6.3.2} and \sref[I]{I.6.6.3}.
\end{proof}

\begin{definition}[1.5.2]
\label{IV.1.5.2}
We say that a morphism $f$ is \emph{of finite type} if it satisfies the equivalent conditions of the proposition~\sref{IV.1.5.1}.
\end{definition}

\oldpage[IV-1]{234}
\begin{proposition}[1.5.3]
\label{IV.1.5.3}
Let $f:X\to Y$ be a morphism of finite type; if $Y$ is Noetherian, so is $X$.
\end{proposition}

\begin{proof}
For the proof, see \sref[I]{I.6.3.7}.
\end{proof}

\begin{proposition}[1.5.4]
\label{IV.1.5.4}
\begin{enumerate}
  \item[\rm(i)] Let $j:X\to Y$ be an immersion.
  If $j$ is quasi-compact (which is the case if $j$ is a closed immersion, or if the underlying space of $X$ is Noetherian, or if the space underlying $Y$ is locally Noetherian), then $j$ is of finite type.
  \item[\rm(ii)] The composition of two morphisms of finite type is of finite type.
  \item[\rm(iii)] If $f:X\to Y$ is an $S$-morphism of finite type, $f_{(S')}:X_{(S')}\to Y_{(S')}$ is of finite type for every base extension $S'\to S$ of the base prescheme.
  \item[\rm(iv)] If $f:X\to X'$ and $g:Y\to Y'$ are two $S$-morphisms of finite type, $f\times_S g$ is of finite type.
  \item[\rm(v)] If the composition $g\circ f$ of two morphisms is of finite type and if $g$ is quasi-separated or if $X$ is Noetherian, $f$ is of finite type.
  \item[\rm(vi)] If a morphism $f$ is of finite type, so is $f_{\mathrm{red}}$.
\end{enumerate}
\end{proposition}

\begin{proof}
This follows immediately (except for the case~{\hyperlink{ega4.c04.prop.1-5-4}{(v)}} where $X$ is Noetherian, proved in \sref[I]{I.6.3.6}) from the definitions and from \sref{IV.1.1.2}, \sref{IV.1.2.4}, and \sref{IV.1.3.4}.
\end{proof}

\begin{corollary}[1.5.5]
\label{IV.1.5.5}
Let $f:X\to Y$ be a morphism of finite type.
For every morphism $Y'\to Y$ such that $Y'$ is Noetherian, $X\times_Y Y'$ is Noetherian.
\end{corollary}

\begin{proof}
This follows from \sref{IV.1.5.4}[(iii)] and from \sref{IV.1.5.3}.
\end{proof}

\begin{corollary}[1.5.6]
\label{IV.1.5.6}
Let $X$ be a prescheme of finite type over a locally Noetherian prescheme $S$.
Then every $S$-morphism $f:X\to Y$ is of finite type.
\end{corollary}

\begin{proof}
For the proof, see \sref[I]{I.6.3.9}.
\end{proof}

\begin{proposition}[1.5.7]
\label{IV.1.5.7}
Let $\varphi:A\to B$ be a homomorphism of rings.
For the corresponding morphism $f:\Spec(B)\to\Spec(A)$ to be of finite type, it is necessary and sufficient that $\varphi$ make $B$ an $A$-algebra of finite type.
\end{proposition}

\begin{proof}
For the proof, see \sref[I]{I.6.3.3}.
\end{proof}

\subsection{Morphisms of finite presentation}
\label{subsection:IV.1.6}

\begin{definition}[1.6.1]
\label{IV.1.6.1}
Let $X$, $Y$ be two preschemes.
We say that a morphism $f:X\to Y$ is \emph{of finite presentation} if it satisfies the following three conditions:
\begin{enumerate}
  \item[\rm(i)] $f$ is locally of finite presentation;
  \item[\rm(ii)] $f$ is quasi-compact (which, when~{\hyperlink{ega4.c04.def.1-6-1}{(i)}} is satisfied, is equivalent to saying that $f$ is of finite type~\sref{IV.1.5.2});
  \item[\rm(iii)] $f$ is quasi-separated.
\end{enumerate}
We then also say that $X$ is \emph{of finite presentation over $Y$}, or a \emph{$Y$-prescheme of finite presentation}.
\end{definition}

Condition~{\hyperlink{ega4.c04.def.1-6-1}{(iii)}} is automatically satisfied if $f$ is separated, or if $X$ is locally Noetherian~\sref{IV.1.2.8}; when $Y$ is locally Noetherian, it amounts to the same to say that $f$ is \emph{of finite presentation} or that $f$ is \emph{of finite type} (the latter condition implying that $X$ is locally Noetherian~\sref{IV.1.3.3}).

\oldpage[IV-1]{235}
\begin{proposition}[1.6.2]
\label{IV.1.6.2}
\begin{enumerate}
  \item[\rm(i)] Every quasi-compact immersion that is locally of finite presentation (in particular every quasi-compact open immersion) is of finite presentation.
  \item[\rm(ii)] The composition of two morphisms of finite presentation is of finite presentation.
  \item[\rm(iii)] Let $X$, $Y$ be two $S$-preschemes, $f:X\to Y$ an $S$-morphism of finite presentation.
  Then, for every morphism $S'\to S$, the morphism $f_{(S')}:X_{(S')}\to Y_{(S')}$ is of finite presentation.
  \item[\rm(iv)] Let $f:X\to Y$, $f':X'\to Y'$ be two $S$-morphisms of finite presentation; then $f\times_S f':X\times_S X'\to Y\times_S Y'$ is of finite presentation.
  \item[\rm(v)] If the composition $g\circ f$ of two morphisms is of finite presentation and if $g$ is quasi-separated and locally of finite presentation, then $f$ is of finite presentation.
\end{enumerate}
\end{proposition}

\begin{proof}
This follows immediately from \sref{IV.1.1.2}, \sref{IV.1.2.2}, \sref{IV.1.2.4}, and \sref{IV.1.4.3}.
\end{proof}

It follows in particular from \sref{IV.1.6.2}[(iii)] that if $f$ is a morphism of finite presentation and $U$ an open subset of $Y$, the restriction $f^{-1}(U)\to U$ of $f$ is a morphism of finite presentation.
Conversely, let $(U_\alpha)$ be a covering of $Y$ by open affines, and suppose that the restrictions $f^{-1}(U_\alpha)\to U_\alpha$ of $f$ are morphisms of finite presentation; then $f$ is evidently locally of finite presentation, quasi-compact by virtue of~\sref{IV.1.1.1}, and quasi-separated by virtue of~\sref{IV.1.2.6}.
We can therefore say that the property for a morphism $f:X\to Y$ of being of finite presentation is \emph{local on $Y$}.

If $X$ is a quasi-separated prescheme, every morphism $f:X\to Y$ is quasi-separated~\sref{IV.1.2.3}.
Hence, if $f$ is quasi-compact and locally of finite presentation, $f$ is of finite presentation.

In particular:

\begin{corollary}[1.6.3]
\label{IV.1.6.3}
Let $\varphi:A\to B$ be a homomorphism of rings.
For the corresponding morphism $f:\Spec(B)\to\Spec(A)$ to be of finite presentation, it is necessary and sufficient that $\varphi$ make $B$ an $A$-algebra of finite presentation.
\end{corollary}

\begin{proof}
The necessity of the condition follows from~\sref{IV.1.4.6}.
\end{proof}

\begin{remark}[1.6.4]
\label{IV.1.6.4}
In the definition~\sref{IV.1.6.1}, condition~{\hyperlink{ega4.c04.def.1-6-1}{(iii)}} is not a consequence of the two others.
Let for example $Y$ be an affine scheme whose underlying space is not Noetherian, and let $U$ be an open subset that is not quasi-compact in $Y$.
Let $X$ be the prescheme obtained by gluing two preschemes $Y_1$, $Y_2$ isomorphic to $Y$ along the open subsets $U_1$, $U_2$ corresponding to $U$ \sref[0\textsubscript{I}]{0.4.1.7}, so that $X$ is the union of two open subsets isomorphic respectively to $Y_1$ and $Y_2$ (and that we identify with them), with $Y_1\cap Y_2=U$.
Furthermore, there is a morphism $f:X\to Y$ that coincides in $Y_i$ with the isomorphism $Y_i\to Y$ ($i=1,2$).
It is clear that $f$ satisfies condition~{\hyperlink{ega4.c04.def.1-6-1}{(i)}} of~\sref{IV.1.6.1}, being a local isomorphism; it also satisfies~{\hyperlink{ega4.c04.def.1-6-1}{(ii)}}, the inverse image of an open quasi-compact subset of $Y$ being the union of its images in $Y_1$ and $Y_2$; but since $Y_1\cap Y_2$ is not quasi-compact, $f$ is not quasi-separated~\sref{IV.1.2.7}.
\end{remark}

\begin{proposition}[1.6.5]
\label{IV.1.6.5}
Let $X'$, $X''$ be two preschemes, $X=X'\amalg X''$ their sum, $f':X'\to Y$, $f'':X''\to Y$ two morphisms, $f:X\to Y$ the morphism that coincides with $f'$ in $X'$ and with $f''$ in $X''$.
For $f$ to be of finite presentation, it is necessary and sufficient that $f'$ and $f''$ be so.
\end{proposition}

\begin{proof}
It suffices to show that for $f$ to possess one of the three properties of the definition~\sref{IV.1.6.1},
\oldpage[IV-1]{236}
it is necessary and sufficient that $f'$ and $f''$ possess it; this is evident for the property of being locally of finite presentation, which is local on $X$; for the property of being quasi-compact, this was seen in~\sref{IV.1.1.6}, and for the property of being quasi-separated, in~\sref{IV.1.2.9}.
\end{proof}

\subsection{Improvements of earlier results}
\label{subsection:IV.1.7}

We will give in this number a list of propositions proved in the preceding chapters, whose statements can be improved with the aid of the new finiteness conditions introduced above.

\begin{env}[1.7.1]
\label{IV.1.7.1}
In the statements of \sref[I]{I.6.4.2}, \sref[I]{I.6.4.3}, and \sref[I]{I.6.4.9}, we can, instead of assuming that $X$ and $Y$ are of finite type over the field $K$, assume only that they are \emph{locally of finite type} over $K$; for \sref[I]{I.6.4.2}, it suffices to observe that every point of a prescheme $X$ that is closed in every affine open subset of $X$ is closed in $X$, and an affine scheme locally of finite type over $K$ is \emph{ipso facto} of finite type over $K$ \sref{IV.1.5.1}.
For \sref[I]{I.6.4.9}, we use \sref{IV.1.3.7} and \sref{IV.1.3.4}[(v)].
\end{env}

\begin{env}[1.7.2]
\label{IV.1.7.2}
In \sref[I]{I.6.5.1}[(ii)], we can, instead of assuming that $S$ is locally Noetherian and $Y$ of finite type over $S$, assume only that $Y$ is \emph{locally of finite presentation} over $S$, as one shows immediately by applying the definition~\sref{IV.1.4.2}.
Similarly in \sref[I]{I.6.5.4}[(ii)] and \sref[I]{I.6.5.5}[(ii)] it suffices to assume that $f$ is an $S$-morphism, $Y$ being \emph{locally of finite presentation} over $S$.
\end{env}

\begin{env}[1.7.3]
\label{IV.1.7.3}
In \sref[I]{I.7.1.11}, for the second assertion, instead of assuming $S$ locally Noetherian and $Y$ of finite type over $S$, we can assume only $Y$ locally of finite presentation over $S$ (the proof depending on \sref[I]{I.6.5.1}[(ii)]); similarly in \sref[I]{I.7.1.12} to \sref[I]{I.7.1.15}.
\end{env}

\begin{env}[1.7.4]
\label{IV.1.7.4}
In \sref[I]{I.9.2.2}, we can replace the three hypotheses by the single hypothesis (implied by each of the others) that $f$ is \emph{quasi-compact and quasi-separated}, by virtue of~\sref{IV.1.2.6} and~\sref{IV.1.2.7}.
We note that in \sref[I]{I.9.2.1}, the hypothesis means exactly that $f$ is quasi-compact and quasi-separated.
\end{env}

\begin{env}[1.7.5]
\label{IV.1.7.5}
In \sref[I]{I.9.3.1}, \sref[I]{I.9.3.2}, and \sref[I]{I.9.3.3}, we can replace the two hypotheses on $X$ by the (weaker) hypothesis that $X$ is a prescheme \emph{quasi-compact and quasi-separated}, the proof using these two properties (by virtue of~\sref{IV.1.2.7}).
\end{env}

\begin{env}[1.7.6]
\label{IV.1.7.6}
In \sref[I]{I.9.3.4} and \sref[I]{I.9.3.5}, it suffices to assume $X$ quasi-compact, $\sh{F}$ and $\sh{J}$ quasi-coherent and of finite type.
\end{env}

\begin{env}[1.7.7]
\label{IV.1.7.7}
In \sref[I]{I.9.4.7}, we can weaken hypothesis~b) by assuming only $X$ \emph{quasi-separated}, since it suffices that the canonical immersion $U\to X$ be quasi-compact~\sref{IV.1.2.7}.
We can make the same modification in \sref[I]{I.9.4.8}, \sref[I]{I.9.4.9}, and \sref[I]{I.9.4.10}.
\end{env}

\begin{env}[1.7.8]
\label{IV.1.7.8}
In \sref[I]{I.9.5.1}, the conditions indicated in parentheses in the statement can be replaced by the condition that $f$ be quasi-compact and quasi-separated.
\end{env}

\begin{env}[1.7.9]
\label{IV.1.7.9}
\oldpage[IV-1]{237}
In \sref[I]{I.9.6.5}, we can replace the hypotheses on $X$ by the weaker hypothesis that $X$ is \emph{quasi-compact and quasi-separated}, as the proof shows, which uses only \sref[I]{I.9.4.7}.
In \sref[I]{I.9.6.6}, the hypothesis ``$X$ is a quasi-compact scheme'' may be replaced by ``$X$ is a prescheme quasi-compact and quasi-separated''.
\end{env}

\begin{env}[1.7.10]
\label{IV.1.7.10}
In \sref[II]{II.1.7.15}, we can suppose only that $Y$ is quasi-compact and quasi-separated, the proof using only \sref[I]{I.9.6.5}.
\end{env}

\begin{env}[1.7.11]
\label{IV.1.7.11}
In the second assertion of \sref[II]{II.3.4.5}, we can substitute for the two hypotheses on $Y$ the weaker hypothesis that $Y$ is quasi-compact and quasi-separated.
Similarly in \sref[II]{II.3.4.8}.
\end{env}

\begin{env}[1.7.12]
\label{IV.1.7.12}
In \sref[II]{II.3.8.5}, we can similarly suppose only that $Y$ is quasi-compact and quasi-separated, this hypothesis entering through \sref[I]{I.9.4.7}.
\end{env}

\begin{env}[1.7.13]
\label{IV.1.7.13}
In \sref[II]{II.4.4.1}, \sref[II]{II.4.4.6}, and \sref[II]{II.4.4.7}, it suffices to assume that $Y$ is \emph{quasi-compact and quasi-separated}, in view of~\sref{IV.1.2.3} in the proof of \sref[II]{II.4.4.7}.
Similarly, in \sref[II]{II.4.4.10.1}, it suffices to assume that $Z$ is quasi-compact and quasi-separated.
\end{env}

\begin{env}[1.7.14]
\label{IV.1.7.14}
In \sref[II]{II.4.5.2} and \sref[II]{II.4.5.5}, it suffices to assume $X$ \emph{quasi-compact and quasi-separated}, this hypothesis entering through \sref[I]{I.9.3.1} and \sref[I]{I.9.3.2}.
\end{env}

\begin{env}[1.7.15]
\label{IV.1.7.15}
In \sref[II]{II.4.6.8}, it suffices still to assume $X$ \emph{quasi-compact and quasi-separated}, this hypothesis entering through \sref[I]{I.9.4.7}.
\end{env}

\begin{env}[1.7.16]
\label{IV.1.7.16}
In \sref[II]{II.5.1.2}, it suffices to assume that $X$ is quasi-compact and quasi-separated (the hypothesis entering through \sref[II]{II.4.5.2} and \sref[II]{II.4.5.5}).
Similarly in \sref[II]{II.5.1.9}, where we use \sref[I]{I.9.6.5}.
\end{env}

\begin{env}[1.7.17]
\label{IV.1.7.17}
In \sref[II]{II.5.2.1}, it suffices always to assume $X$ quasi-compact and quasi-separated (the hypothesis entering through \sref[II]{II.4.5.2}).
\end{env}

\begin{env}[1.7.18]
\label{IV.1.7.18}
In \sref[II]{II.5.2.2}, one must assume $f$ \emph{quasi-compact} and $X$ \emph{quasi-separated}; the current proof is in fact insufficient (with the hypotheses made), for from the fact that for every quasi-coherent $\sh{O}_X$-module $\sh{F}$ we have $\mathrm{R}^1 f_*(\sh{F})=0$, it does not follow, necessarily, for an affine open subset $U$ of $Y$, that the same property holds for quasi-coherent $(\sh{O}_X|f^{-1}(U))$-modules, if one does not know that such a module is the restriction of a quasi-coherent $\sh{O}_X$-module (the same remark applying to conditions~b) and~c')); when $f$ is quasi-compact and $X$ quasi-separated, this extension is possible by virtue of \sref[I]{I.9.4.7}, modified above~\sref{IV.1.7.7}.
\end{env}

\begin{env}[1.7.19]
\label{IV.1.7.19}
In \sref[II]{II.5.3.2}, it suffices to assume $Y$ \emph{quasi-compact and quasi-separated} (this hypothesis entering through \sref[II]{II.4.4.6} and \sref[II]{II.4.4.7}).
Similarly in \sref[II]{II.5.5.3}[(ii)].
\end{env}

\begin{env}[1.7.20]
\label{IV.1.7.20}
In \sref[III]{III.1.2.2}, \sref[III]{III.1.2.3}, and \sref[III]{III.1.2.4}, we can replace the two hypotheses on $X$ by the weaker hypothesis that $X$ is a prescheme \emph{quasi-compact and quasi-separated}, the proof using only \sref[I]{I.9.3.3}.
\end{env}

\begin{env}[1.7.21]
\label{IV.1.7.21}
In \sref[III]{III.1.4.10} to \sref[III]{III.1.4.14}, we can replace ``separated'' by ``quasi-separated'', this hypothesis serving only to be able to apply \sref[I]{I.9.2.2} (see above~\sref{IV.1.7.4}).
Similarly, in \sref[III]{III.1.4.15}, it suffices to assume that the morphism $u$ is \emph{quasi-separated and quasi-compact}.
\end{env}

\begin{env}[1.7.22]
\label{IV.1.7.22}
In \sref[III]{III.2.1.3}, we can again replace the hypotheses by the weaker hypothesis that $X$ is a prescheme quasi-compact and quasi-separated, the reasoning being the same as in \sref[III]{III.1.2.2}.
\end{env}

\subsection{Morphisms of finite presentation and constructible sets}
\label{subsection:IV.1.8}

\oldpage[IV-1]{238}

\begin{env}[1.8.1]
\label{IV.1.8.1}
Let $X$ be a prescheme; we know that every quasi-compact open subset of $X$ is a finite union of affine open subsets, and conversely.
For an open subset $U$ of $X$ to be \emph{retrocompact} in $X$ \sref[0\textsubscript{III}]{0.9.1.1}, it is necessary and sufficient that for every \emph{open affine} subset $V$ of $X$, $U\cap V$ be quasi-compact.
When $X$ is \emph{quasi-separated} \sref{IV.1.2.1}, every quasi-compact open subset of $X$ is retrocompact in $X$ \sref{IV.1.2.7}, so every locally constructible part of $X$ is retrocompact in $X$ \sref[0\textsubscript{III}]{0.9.1.13}; if \emph{in addition} $X$ is \emph{quasi-compact}, constructible subsets and locally constructible subsets of $X$ coincide \sref[0\textsubscript{III}]{0.9.1.12}, as do constructible open subsets and quasi-compact open subsets \sref[0\textsubscript{III}]{0.9.1.5}.
\end{env}

\begin{proposition}[1.8.2]
\label{IV.1.8.2}
Let $X$, $Y$ be two preschemes, $f:X\to Y$ a morphism.
For every constructible (resp. locally constructible) part $Z$ of $Y$, $f^{-1}(Z)$ is a constructible (resp. locally constructible) part of $X$.
\end{proposition}

\begin{proof}
Suppose $Z$ locally constructible; for every $x\in X$, there will be an affine open neighbourhood $V$ of $f(x)$ such that $Z\cap V$ is constructible in $V$; if the proposition is proved for constructible parts, $f^{-1}(Z)\cap f^{-1}(V)$ will be constructible in $f^{-1}(V)$, so $f^{-1}(Z)$ will be locally constructible.
It therefore suffices to consider the case where $Z$ is constructible, and taking into account \sref[0\textsubscript{III}]{0.9.1.3}, we are reduced to the case where $Z$ is open and \emph{retrocompact} in $Y$, in other words such that the canonical injection $Z\to Y$ is a quasi-compact morphism; it then suffices to see that $f^{-1}(Z)$ is \emph{retrocompact in} $X$, in other words such that the canonical injection $f^{-1}(Z)\to X$ is a quasi-compact morphism; but this follows from \sref{IV.1.1.2}[(iii)], since $f^{-1}(Z)=Z\times_Y X$ \sref[I]{I.4.4.1}.
\end{proof}

\begin{lemma}[1.8.3]
\label{IV.1.8.3}
Let $X$ be a prescheme quasi-compact and quasi-separated, $Z$ a constructible part of $X$.
There exists then an affine scheme $X'$ and a morphism of finite presentation $f:X'\to X$ such that $f(X')=Z$.
\end{lemma}

\begin{proof}
Since $X$ is quasi-compact, it is a finite union of affine open subsets $X_i$ ($i\in\mathrm{I}$), and since $X$ is quasi-separated, it follows from \sref{IV.1.2.7} that each of the canonical immersions $g_i:X_i\to X$ is of finite presentation; one concludes that if $f_i:X_i'\to X_i$ is a morphism of finite presentation, the same is true of $g_i\circ f_i:X_i'\to X$ \sref{IV.1.6.2}, and if $X'$ is the prescheme sum of the $X_i'$, the morphism $h:X'\to X$ which coincides with $g_i\circ f_i$ in each $X_i'$ is also of finite presentation \sref{IV.1.6.5}.
Since $Z\cap X_i$ is constructible in $X_i$ \sref[0\textsubscript{III}]{0.9.1.8}, we see that we are reduced to proving the lemma when $X_i$ is \emph{affine} of ring $A$.
Since $Z$ is then a finite union of sets of the form $U\cap\complement V$, where $U$ and $V$ are quasi-compact open subsets \sref[0\textsubscript{III}]{0.9.1.3}, one sees, by considering yet another suitable sum of preschemes, that one can reduce to the case where $Z=U\cap\complement V$; since moreover $U$ is a finite union of affine open subsets, one can even restrict ourselves to the case where $U$ is \emph{affine}.
Since $V$ is quasi-compact, it is a finite union of open subsets of the form $X_{f_i}$, where $f_i\in A$; let $\mathfrak{J}$ be the ideal of $A$ generated by the $f_i$, and let $Z'$ be the closed affine subscheme of $X$ that it defines (and which is by construction of finite presentation over $X$); one has by definition $V=\complement Z'$, so $U\cap\complement V=U\cap Z'$.
Consider the affine scheme $X'=Z'\times_X U$

\oldpage[IV-1]{239}

and let $f:X'\to X$ be the structural morphism; we just saw that the canonical immersion $Z'\to X$ is of finite presentation, and it is the same for the open immersion $U\to X$, which is quasi-compact since $U$ and $X$ are affine; one concludes therefore from \sref{IV.1.6.2}[(iv)] that $f$ is of finite presentation, and one has $f(X')=Z'\cap U$ \sref[I]{I.3.4.8}, which completes the proof.
\end{proof}

\begin{theorem}[1.8.4]
\label{IV.1.8.4}
\textup{(Chevalley).} ---
Let $f:X\to Y$ be a morphism of finite presentation.
For every locally constructible part $Z$ of $X$, $f(Z)$ is locally constructible in $Y$.
\end{theorem}

\begin{proof}
Let $y\in Y$ and $V$ an open affine neighbourhood of $y$; since $f$ is quasi-compact and quasi-separated, the same is true of its restriction $f^{-1}(V)\to V$, so $f^{-1}(V)$ is a quasi-compact and quasi-separated prescheme \sref{IV.1.2.3}; since $Z\cap f^{-1}(V)$ is constructible in $f^{-1}(V)$ \sref{IV.1.8.1}, we see that one can restrict ourselves to the case where $Y$ is affine and $Z$ constructible; $X$ is then quasi-compact and quasi-separated, so by \sref{IV.1.8.3} there exists a morphism of finite presentation $g:X'\to X$ such that $Z=g(X')$; since $f\circ g$ is of finite presentation, we see that we are reduced to the case where $Z=X$; in other words, it will suffice to prove the
\end{proof}

\begin{lemma}[1.8.4.1]
\label{IV.1.8.4.1}
Let $Y$ be an affine scheme, $f:X\to Y$ a morphism quasi-compact and locally of finite presentation; then $f(X)$ is a constructible part of $Y$.
\end{lemma}

\begin{proof}
Since $X$ is quasi-compact, it is a finite union of affine open subsets; one can therefore restrict ourselves to the case where $Y=\Spec(A)$, $X=\Spec(B)$, $B$ being an $A$-algebra of finite presentation \sref{IV.1.4.6}.
Now we have the
\end{proof}

\begin{lemma}[1.8.4.2]
\label{IV.1.8.4.2}
Let $A_0$ be a ring, $(A_\lambda)_{\lambda\in\mathrm{L}}$ an inductive system of $A_0$-algebras, $A=\varinjlim A_\lambda$; if $B$ is an $A$-algebra of finite presentation, there exist a $\lambda$ and an $A_\lambda$-algebra of finite presentation $B_\lambda$ such that $B$ is isomorphic to $B_\lambda\otimes_{A_\lambda} A$.
\end{lemma}

\begin{proof}
By hypothesis, $B$ is isomorphic to an algebra of the form $A[\mathrm{T}_1,\ldots,\mathrm{T}_n]/\mathfrak{J}$, where the $\mathrm{T}_i$ are indeterminates and $\mathfrak{J}$ an ideal of finite type; let $(\mathrm{F}_j)$ ($1\leq j\leq m$) be a system of generators of $\mathfrak{J}$.
Let $\varphi_\lambda:A_\lambda\to A$ be the canonical homomorphism.
Since $\mathrm{L}$ is filtering, there exists $\lambda$ such that each of the coefficients of each of the polynomials $\mathrm{F}_j$ is the image by $\varphi_\lambda$ of an element of $A_\lambda$.
In other words, there exists a system of polynomials $(\mathrm{F}_{j\lambda})_{1\leq j\leq m}$ of $A_\lambda[\mathrm{T}_1,\ldots,\mathrm{T}_n]$ such that $\mathrm{F}_j=\varphi_\lambda(\mathrm{F}_{j\lambda})$ for $1\leq j\leq m$.
If $\mathfrak{J}_\lambda$ is the ideal of $A_\lambda[\mathrm{T}_1,\ldots,\mathrm{T}_n]$ generated by the $\mathrm{F}_{j\lambda}$, $\mathfrak{J}$ is the image of $\mathfrak{J}_\lambda\otimes_{A_\lambda} A$ in $A[\mathrm{T}_1,\ldots,\mathrm{T}_n]=A_\lambda[\mathrm{T}_1,\ldots,\mathrm{T}_n]\otimes_{A_\lambda} A$; the ring $B_\lambda=A_\lambda[\mathrm{T}_1,\ldots,\mathrm{T}_n]/\mathfrak{J}_\lambda$ has the required property.
\end{proof}

We apply this lemma here by considering $A$ as the inductive limit of its sub-$\bb{Z}$-algebras \emph{of finite type}.
One sees that there exist such a sub-$\bb{Z}$-algebra $A_0$ of $A$, and an $A_0$-algebra of finite presentation $B_0$ such that $B$ is isomorphic to $B_0\otimes_{A_0} A$; let us set $X_0=\Spec(B_0)$, $Y_0=\Spec(A_0)$, so that $X=X_0\times_{Y_0} Y$, $f$ being the projection $X\to Y$; let $f_0:X_0\to Y_0$, $g_0:Y\to Y_0$ be the structural morphisms; it follows from \sref[I]{I.3.4.8} that $f(X)=g_0^{-1}(f_0(X_0))$; taking into account \sref{IV.1.8.2}, it therefore suffices to show that $f_0(X_0)$ is constructible.
In other words, one is finally reduced to proving the

\begin{corollary}[1.8.5]
\label{IV.1.8.5}
Let $Y$ be a Noetherian prescheme, $f:X\to Y$ a morphism of finite type.
For every constructible part $Z$ of $X$, $f(Z)$ is a constructible part of $Y$.
\end{corollary}

\oldpage[IV-1]{240}

\begin{proof}
One reduces as above to proving that $f(X)$ is constructible.
Let us apply the criterion \sref[0\textsubscript{III}]{0.9.2.3}: one must prove that for every irreducible closed part $T$ of $Y$, $T\cap f(X)=f(f^{-1}(T))$ is rare in $T$ or contains a non-empty open subset of $T$; taking a subscheme of $Y$ whose underlying space is $T$ \sref[I]{I.5.2.1} and considering its inverse image by $f$, one sees finally (taking into account \sref{IV.1.5.4}[(iii)]) that one is reduced to proving that if one supposes $Y$ \emph{irreducible and} $f$ \emph{dominant}, $f(X)$ contains a non-empty open subset of $Y$.
One can moreover suppose $Y$ affine; then $X$ is a finite union of affine open subsets $V_i$, and since $f(X)$ is dense in $Y$, the same is true of at least one of the $f(V_i)$.
One can therefore also suppose $X$ affine; if $(X_j)$ is the (finite) family of irreducible components of $X$, one sees as above that at least one of the $f(X_j)$ is dense in $Y$; one can therefore suppose $X$ irreducible.
Finally, by replacing $f$ by $f_{\mathrm{red}}$ \sref{IV.1.5.4}[(vi)] one can suppose $X$ and $Y$ \emph{integral}.
Then $X=\Spec(B)$, $Y=\Spec(A)$, where $A$ is an integral Noetherian ring, $B$ an integral $A$-algebra of finite type containing $A$ \sref[I]{I.1.2.7}.
It suffices to show that there exists $g\in A$ such that, for every $y\in\mathrm{D}(g)$ (that is to say such that $g\notin\mathfrak{p}_y$) the ideal $\mathfrak{p}_y$ is the intersection of $A$ and a prime ideal of $B$, for this will show that $\mathrm{D}(g)\subset f(X)$.
Finally it suffices to prove the
\end{proof}

\begin{lemma}[1.8.5.1]
\label{IV.1.8.5.1}
Let $A$ be an integral ring, $B$ an integral $A$-algebra of finite type.
There exists $g\in A$ such that every homomorphism of $A$ into an algebraically closed field $\Omega$, non-zero on $g$, extends to a homomorphism of $B$ into $\Omega$.
\end{lemma}

\begin{proof}
This is a classical result of commutative algebra (Bourbaki, \emph{Alg. comm.}, chap.~V, \S\,3, n\textsuperscript{o}\,1, cor.~3 of th.~1).
\end{proof}

\begin{corollary}[1.8.6]
\label{IV.1.8.6}
Let $Y$ be an irreducible prescheme, $\eta$ its generic point, $f:X\to Y$ a morphism locally of finite type.
If $f^{-1}(\eta)$ is not empty, there exists a neighbourhood $U$ open of $\eta$ in $Y$ such that $f^{-1}(y)$ is non-empty for every $y\in U$.
If $\sh{F}$ is a quasi-coherent $\sh{O}_X$-module of finite type, and if $\sh{F}_\eta=\sh{F}\otimes_{\sh{O}_Y} k(\eta)$ is not zero, then there exists a neighbourhood $U'$ of $\eta$ in $Y$ such that $\sh{F}_y=\sh{F}\otimes_{\sh{O}_Y} k(y)$ is not zero for every $y\in U'$.
\end{corollary}

\begin{proof}
If $p_y$ is the canonical projection of the fibre $f^{-1}(y)=X\times_Y\Spec(k(y))$ to $X$, one has $\sh{F}_y=p_y^*(\sh{F})$, so $\mathrm{Supp}(\sh{F}_y)=p_y^{-1}(\mathrm{Supp}(\sh{F}))=\mathrm{Supp}(\sh{F})\cap f^{-1}(y)$ by virtue of \sref[I]{I.9.1.13.1} and \sref[I]{I.3.6.1}, since $\sh{F}$ is of finite type; moreover $\mathrm{Supp}(\sh{F})$ is closed in $X$ \sref[0\textsubscript{I}]{0.5.2.2}, and if $Z$ is a closed subscheme of $X$ having as underlying space $\mathrm{Supp}(\sh{F})$ \sref[I]{I.5.2.1}, and $j$ the canonical immersion $Z\to X$, $f\circ j$ is locally of finite type \sref{IV.1.3.4}; this proves that the first assertion implies the second.
To prove the first assertion, let us first note that one can suppose $X$ and $Y$ reduced \sref{IV.1.3.4}[(vi)], in other words $Y$ \emph{integral}.
Let $\xi$ be a point of $f^{-1}(\eta)$; one can replace $X$ and $Y$ by affine open neighbourhoods of $\xi$ and $\eta$ respectively, and thus suppose that $X=\Spec(B)$, $Y=\Spec(A)$, $B$ being an $A$-algebra of finite type.
Moreover, if $Z$ is the reduced closed subscheme of $X$ having $\overline{\{\xi\}}$ as underlying set \sref[I]{I.5.2.1}, one can, as above, replace $X$ by $Z$, in other words suppose $B$ \emph{integral} \sref[I]{I.5.1.4}.
Since the morphism $f$ is then dominant, the corresponding homomorphism $\varphi:A\to B$ is injective \sref[I]{I.1.2.7}; the corollary therefore follows finally from the lemma \sref{IV.1.8.5.1}.
\end{proof}

\oldpage[IV-1]{241}

\begin{proposition}[1.8.7]
\label{IV.1.8.7}
Let $S$ be a prescheme, $g:X\to S$, $h:Y\to S$ two morphisms of finite presentation, $f:X\to Y$ an $S$-morphism.
For every $s\in S$, let $X_s=g^{-1}(s)=X\times_S\Spec(k(s))$, $Y_s=h^{-1}(s)=Y\times_S\Spec(k(s))$, $f_s=f\times 1:X_s\to Y_s$.
Then the set of $s\in S$ such that $f_s$ is surjective (resp. radicial) is locally constructible.
\end{proposition}

\begin{proof}
Let $E$ be the set of $s\in S$ such that $f_s$ is surjective; the set $S-E$ is equal to $h(Y-f(X))$; now $f$ is of finite presentation \sref{IV.1.6.2}[(v)] so $f(X)$ is locally constructible in $Y$ \sref{IV.1.8.4}; since $h$ is of finite presentation, $h(Y-f(X))$ is locally constructible in $S$, and therefore also $E$.

To show that the set $E'$ of $s\in S$ such that $f_s$ is radicial is locally constructible, we will use the
\end{proof}

\begin{lemma}[1.8.7.1]
\label{IV.1.8.7.1}
Let $U$, $V$ be two preschemes; for a morphism $f:U\to V$ to be radicial, it is necessary and sufficient that the diagonal morphism $\Delta_f:U\to U\times_V U$ is surjective.
Consequently, every radicial morphism is separated.
\end{lemma}

\begin{proof}
Indeed, $\Delta_f$ being an immersion \sref[I]{I.5.3.9}, is a morphism locally of finite type \sref{IV.1.3.4}; to say that $\Delta_f$ is surjective therefore means that for every algebraically closed field $K$, the corresponding map $U(K)\to(U\times_V U)(K)=U(K)\times_{V(K)} U(K)$ is surjective \sref{IV.1.3.7} and \sref[I]{I.3.4.2.1}; but by definition of a fibre product of sets, this means that the map $U(K)\to V(K)$ corresponding to $f$ is injective, and this is equivalent to saying that $f$ is radicial \sref[I]{I.3.5.5}.

This being so, $\Delta_{f_s}$ is deduced from $\Delta_f:X\to X\times_Y X$ by the base change $\Spec(k(s))\to S$ \sref[I]{I.5.3.4}.
Since $f$ is of finite presentation, it is the same for the structural morphism $X\times_Y X\to Y$ \sref{IV.1.6.2}[(iv)], so also for $\Delta_f$ \sref{IV.1.6.2}[(v)]; it therefore suffices to apply the first part of the proposition to $\Delta_f$, using the lemma \sref{IV.1.8.7.1}.
\end{proof}

\subsection{Pro-constructible sets and ind-constructible sets}
\label{subsection:IV.1.9}

\begin{lemma}[1.9.1]
\label{IV.1.9.1}
Let $(A_\alpha)_{\alpha\in\mathrm{I}}$ be an inductive system of rings with index set $\mathrm{I}$ filtering to the right, and let $A=\varinjlim A_\alpha$.
For $A=0$, it is necessary and sufficient that there exists an index $\gamma$ such that $A_\gamma=0$ (and one then has $A_\beta=0$ for every $\beta\geq\gamma$).
\end{lemma}

\begin{proof}
Indeed, for every $\alpha\in\mathrm{I}$, the canonical homomorphism $\varphi_\alpha:A_\alpha\to A$ sends the unit element to the unit element; to say that $A=0$ means that $\varphi_\alpha(1)=0$, or again $\varphi_\alpha(1)=\varphi_\alpha(0)$; one knows that this is equivalent to saying that there exists an index $\gamma\geq\alpha$ such that the homomorphism $\varphi_{\gamma\alpha}:A_\alpha\to A_\gamma$ is such that $\varphi_{\gamma\alpha}(1)=\varphi_{\gamma\alpha}(0)$, and this last relation is equivalent to $A_\gamma=0$, whence the lemma.
\end{proof}

\begin{proposition}[1.9.2]
\label{IV.1.9.2}
Let $B$ be a ring, $(A_\alpha)_{\alpha\in\mathrm{I}}$ an inductive system of $B$-algebras with index set $\mathrm{I}$ filtering to the right, and let $A=\varinjlim A_\alpha$; set $Y=\Spec(B)$, $X=\Spec(A)$, $X_\alpha=\Spec(A_\alpha)$, and let $u:X\to Y$, $u_\alpha:X_\alpha\to Y$ be the morphisms corresponding to the structural homomorphisms $B\to A$, $B\to A_\alpha$.
Then:
\end{proposition}

\begin{env}[1.9.2.i]
\label{IV.1.9.2.i}
\emph{For $X=\varnothing$, it is necessary and sufficient that there exists $\gamma\in\mathrm{I}$ such that $X_\gamma=\varnothing$, and one then has $X_\alpha=\varnothing$ for $\alpha\geq\gamma$.}
\end{env}

\oldpage[IV-1]{242}

\begin{env}[1.9.2.ii]
\label{IV.1.9.2.ii}
\emph{We have}
\[
  u(X)=\bigcap_{\alpha\in\mathrm{I}} u_\alpha(X_\alpha). \tag{1.9.2.1}
\label{IV.1.9.2.1}
\]
\end{env}

The assertion (i) is nothing other than the translation of \sref{IV.1.9.1}.
To prove (ii), note that, since $u$ factors as $X\longrightarrow X_\alpha\xrightarrow{u_\alpha}Y$, the left-hand side is contained in the right-hand side.
Conversely, let $y\in Y-u(X)$; one has $u^{-1}(y)=\varnothing$, and $u^{-1}(y)$ is the space underlying the $k(y)$-prescheme $\Spec(A\otimes_B k(y))$; since in the category of $B$-modules the functor $\varinjlim$ commutes with the tensor product, $A\otimes_B k(y)$ is the inductive limit of the system of rings $A_\alpha\otimes_B k(y)$; the lemma \sref{IV.1.9.1} shows that there exists $\alpha$ such that $A_\alpha\otimes_B k(y)=0$, that is to say $u_\alpha^{-1}(y)=\varnothing$, so $y\notin u_\alpha(X_\alpha)$.

\begin{proposition}[1.9.3]
\label{IV.1.9.3}
Let $X$ be a quasi-compact and quasi-separated prescheme, let $E$ be a part of $X$, let $(X_\alpha)_{\alpha\in\mathrm{I}}$ be an affine open cover of $X$, and, for every $\alpha\in\mathrm{I}$, let $E_\alpha=E\cap X_\alpha$.
The following conditions are equivalent:
\begin{enumerate}
  \item[a)] There exists a quasi-compact morphism $f:X'\to X$ such that $E=f(X')$.
  \item[a$'$)] For every $\alpha$, there exists a quasi-compact morphism $f_\alpha:X_\alpha'\to X_\alpha$ such that $E_\alpha=f_\alpha(X_\alpha')$.
  \item[a$''$)] For every $\alpha$, there exist an affine scheme $X_\alpha'$ and a morphism $f_\alpha:X_\alpha'\to X_\alpha$ such that $f_\alpha(X_\alpha')=E_\alpha$.
  \item[b)] $E$ is an intersection of constructible parts of $X$.
  \item[b$'$)] $F=X-E$ is a union of constructible parts of $X$.
\end{enumerate}
\end{proposition}

\begin{proof}
It is clear that $b$) and $b'$) are equivalent.
The condition $a$) implies $a'$) by taking for $X_\alpha'$ the prescheme induced on the open subset $f^{-1}(X_\alpha)$ of $X'$ and for $f_\alpha$ the restriction of $f$.
To see that $a'$) implies $a''$), it suffices to remark that $X_\alpha'$ is a union of affine open subsets $X_{\alpha\lambda}''$ ($\lambda\in\mathrm{L}$); let $X_\alpha''$ be the prescheme \emph{sum} of the $X_{\alpha\lambda}''$ ($\lambda\in\mathrm{L}$), which is an affine scheme; if $g_\alpha:X_\alpha''\to X_\alpha'$ is the morphism coinciding with the identity in each $X_{\alpha\lambda}''$, it is clear that if one sets $f_\alpha'=f_\alpha\circ g_\alpha$, one has $E_\alpha=f_\alpha'(X_\alpha'')$ since $g_\alpha$ is surjective.
To show that $a'')$ implies $a$), note that there is a finite part $\mathrm{J}$ of $\mathrm{I}$ such that the $X_\alpha$ for $\alpha\in\mathrm{J}$ cover $X$; let $X'$ be the prescheme sum of the $X_\alpha'$ for $\alpha\in\mathrm{J}$, and let $f:X'\to X$ be the morphism coinciding with $j_\alpha\circ f_\alpha$ for every $\alpha\in\mathrm{J}$, where $j_\alpha:X_\alpha\to X$ is the canonical injection.
Since $E$ is the union of the $E\cap X_\alpha$ for $\alpha\in\mathrm{J}$, one has $E=f(X')$ and it remains to see that $f$ is a quasi-compact morphism; but the hypothesis that $X$ is quasi-separated implies that $j_\alpha$ is quasi-compact \sref{IV.1.2.7}, so $j_\alpha\circ f_\alpha$ is quasi-compact \sref{IV.1.1.1} and \sref{IV.1.1.2}, and hence so is $f$ \sref{IV.1.1.6}.

It remains therefore to prove the equivalence of $a'')$ and $b$).
Let us first prove that $a'')$ implies $b$): let us consider a finite part $\mathrm{J}$ of $\mathrm{I}$ such that the $X_\alpha$ for $\alpha\in\mathrm{J}$ cover $X$; it will suffice to show that, for every $\alpha\in\mathrm{J}$, $E_\alpha$ is an intersection of constructible parts $\mathrm{F}_{\alpha\lambda}$ of $X_\alpha$ ($\lambda\in\mathrm{L}_\alpha$); indeed, every $\mathrm{F}_{\alpha\lambda}$ is also a constructible part of $X$ by virtue of \sref[0\textsubscript{III}]{0.9.1.8}[(ii)], since the hypothesis that $X$ is quasi-separated implies that every $X_\alpha$ is retrocompact in $X$ \sref{IV.1.2.7}; since $E$ is the union of the $E_\alpha$ for $\alpha\in\mathrm{J}$, it is the intersection of the finite unions $\bigcup_{\alpha\in\mathrm{J}} \mathrm{F}_{\alpha,\lambda(\alpha)}$, over all choices of $\lambda(\alpha)\in\mathrm{L}_\alpha$; these finite unions are constructible parts of $X$, which proves our assertion.
We may therefore reduce to the case
\oldpage[IV-1]{243}
where $X=\Spec(A)$ and $E=f(X')$, with $X'=\Spec(A')$ affine and $f:X'\to X$ corresponding to a homomorphism of rings $A\to A'$.
Let us note now that $A'$ is the inductive limit of the family $(A_\lambda')$ of its sub-$A$-algebras \emph{of finite type}, ordered by inclusion; if $X_\lambda'=\Spec(A_\lambda')$, $f$ factors as $X'\to X_\lambda'\xrightarrow{f_\lambda} X$ and it follows from \sref{IV.1.9.2}[(ii)] that $f(X')=\bigcap_\lambda f_\lambda(X_\lambda')$; if one shows that $f_\lambda(X_\lambda')$ is an intersection of constructible parts of $X$, it will be the same for $f(X')$.
We may therefore reduce to the case where $A'$ is an $A$-algebra \emph{of finite type}.
Now, we have the following lemma:
\end{proof}

\begin{lemma}[1.9.3.1]
\label{IV.1.9.3.1}
Let $A$ be a ring, $A'$ an $A$-algebra of finite type.
Then $A'$ is $A$-isomorphic to a filtering inductive limit $\varinjlim A_\mu''$ where the $A_\mu''$ are $A$-algebras of finite presentation \sref{IV.1.4.1}.
\end{lemma}

\begin{proof}
Indeed, one can write $A'=B/\mathfrak{J}$, with $B=A[\mathrm{T}_1,\ldots,\mathrm{T}_n]$ and $\mathfrak{J}$ an ideal of $B$.
But $\mathfrak{J}$ is inductive limit of the ideals \emph{of finite type} $\mathfrak{J}_\mu'$ of $B$, contained in $\mathfrak{J}$, ordered by inclusion; since in the category of $B$-modules the functor $\varinjlim$ is exact, one concludes that $A'=\varinjlim B/\mathfrak{J}_\mu'$ up to $A$-isomorphism.
\end{proof}

The same reasoning as above then permits us to reduce to the case where $A'$ is an $A$-algebra of finite presentation; but then $f(X')$ is constructible in $X$ by virtue of the theorem of Chevalley \sref{IV.1.8.4},\footnote{The French source cites (1.8.5), but Chevalley's theorem is (1.8.4).} which proves $b$).

Let us prove finally that $b$) implies $a''$).
If $E$ is an intersection of a family $(G_\lambda)_{\lambda\in\mathrm{L}}$ of constructible parts of $X$, each $E_\alpha$ is an intersection of the $G_\lambda\cap X_\alpha$, which are constructible in $X_\alpha$ \sref[0\textsubscript{III}]{0.9.1.8}[(i)], and one is reduced to the case where $X=\Spec(A)$ is affine.

One knows then \sref{IV.1.8.3} that for every $\lambda\in\mathrm{L}$ there exists a morphism $f_\lambda:X_\lambda'\to X$, where $X_\lambda'=\Spec(A_\lambda')$ is affine, such that $f_\lambda(X_\lambda')=G_\lambda$.
It therefore suffices to prove the following lemma:

\begin{lemma}[1.9.3.2]
\label{IV.1.9.3.2}
Let $(A_\lambda')_{\lambda\in\mathrm{L}}$ be a family of $A$-algebras, $X=\Spec(A)$, $X_\lambda'=\Spec(A_\lambda')$, and let $f_\lambda:X_\lambda'\to X$ be the structural morphism.
For every finite part $\mathrm{J}$ of $\mathrm{L}$, set $A_\mathrm{J}'=\bigotimes_{\lambda\in\mathrm{J}} A_\lambda'$ (tensor product of $A$-algebras), $X_\mathrm{J}'=\Spec(A_\mathrm{J}')$, and let $f_\mathrm{J}:X_\mathrm{J}'\to X$ be the structural morphism.
One then has $f_\mathrm{J}(X_\mathrm{J}')=\bigcap_{\lambda\in\mathrm{J}} f_\lambda(X_\lambda')$.
If $A'=\varinjlim A_\mathrm{J}'$, $\mathrm{J}$ running through the set of finite parts of $\mathrm{L}$, $X'=\Spec(A')$ and $f:X'\to X$ is the structural morphism, one has $f(X')=\bigcap_{\lambda\in\mathrm{L}} f_\lambda(X_\lambda')$.
\end{lemma}

\begin{proof}
The first assertion follows from \sref[I]{I.3.4.7}; the second follows from \sref{IV.1.9.2}[(ii)], which gives the relation $f(X')=\bigcap_\mathrm{J} f_\mathrm{J}(X_\mathrm{J}')$.
\end{proof}

\begin{definition}[1.9.4]
\label{IV.1.9.4}
Let $X$ be a topological space.
We say that a part $E$ of $X$ is \emph{pro-constructible} (resp. \emph{ind-constructible}) in $X$ if, for every $x\in X$, there exists an open neighbourhood $U$ of $x$ in $X$ such that $E\cap U$ is an intersection (resp. a union) of locally constructible parts of $U$.
\end{definition}

Taking into account \sref{IV.1.2.7}, \sref[0\textsubscript{III}]{0.9.1.8}, and \sref[0\textsubscript{III}]{0.9.1.12}, the equivalent conditions of the proposition \sref{IV.1.9.3} are therefore expressed by saying that $E$ is \emph{pro-constructible} in $X$, the locally constructible parts of $X$ being identical to the constructible parts of $X$ under the hypotheses of \sref{IV.1.9.3}.

\oldpage[IV-1]{244}

\begin{proposition}[1.9.5]
\label{IV.1.9.5}
Let $X$ be a prescheme.
\begin{enumerate}
  \item[\rm(i)] For a part $E$ of $X$ to be pro-constructible, it is necessary and sufficient that $X-E$ be ind-constructible.
  \item[\rm(ii)] Every finite union and every intersection of pro-constructible parts of $X$ are pro-constructible.
  Every finite intersection and every union of ind-constructible parts of $X$ are ind-constructible.
  \item[\rm(iii)] Every intersection (resp. union) of locally constructible parts of $X$ is pro-constructible (resp. ind-constructible).
  Conversely, if $X$ is quasi-compact and quasi-separated, every pro-constructible (resp. ind-constructible) part of $X$ is intersection (resp. union) of constructible parts of $X$.
  \item[\rm(iv)] Let $(U_\alpha)$ be an open cover of $X$.
  For a part $E$ of $X$ to be pro-constructible (resp. ind-constructible) in $X$, it is necessary and sufficient that for every $\alpha$, $E\cap U_\alpha$ be pro-constructible (resp. ind-constructible) in $U_\alpha$.\footnote{The French source prints \emph{rep.} in place of \emph{resp.} in this item.}
  \item[\rm(v)] Every pro-constructible part $E$ of $X$ is retrocompact in $X$.
  For a locally closed part $E$ of $X$ to be retrocompact, it is necessary and sufficient that it be pro-constructible in $X$.
  \item[\rm(vi)] Let $f:X'\to X$ be a morphism of preschemes; for every pro-constructible (resp. ind-constructible) part $E$ of $X$, $f^{-1}(E)$ is pro-constructible (resp. ind-constructible) in $X'$.
  \item[\rm(vii)] Let $f:X'\to X$ be a quasi-compact morphism; for every pro-constructible part $E'$ of $X'$, $f(E')$ is pro-constructible in $X$; in particular $f(X')$ is pro-constructible in $X$.
  \item[\rm(viii)] Let $f:X'\to X$ be a morphism locally of finite presentation; for every ind-constructible part $E'$ of $X'$, $f(E')$ is ind-constructible in $X$; in particular $f(X')$ is ind-constructible in $X$.
  \item[\rm(ix)] Suppose that $X$ is quasi-compact and quasi-separated; then, for a part $E$ of $X$ to be pro-constructible, it is necessary and sufficient that there exist an affine scheme $X'$ and a (necessarily quasi-compact) morphism $f:X'\to X$ such that $E=f(X')$.
\end{enumerate}
\end{proposition}

\begin{proof}
Properties (i), (ii), (iv) and the first assertion of (iii) follow from the definition \sref{IV.1.9.4} and are valid for any topological space, using the mutual distributivity of intersection and union and the fact that if $T$ is locally constructible in $X$, $T\cap U$ is locally constructible in $U$ for every open subset $U$ of $X$.
If $X$ is quasi-compact and quasi-separated and $E$ is pro-constructible in $X$, there is a finite open cover $(U_i)$ of $X$ by affine open subsets such that $E\cap U_i$ is an intersection of constructible parts $\mathrm{M}_{i\lambda}$ of $U_i$ ($\lambda\in\mathrm{L}_i$); the $\mathrm{M}_{i\lambda}$ are also constructible parts of $X$ by virtue of \sref{IV.1.2.7} and \sref[0\textsubscript{III}]{0.9.1.8}[(ii)], and, since $E$ is the union of the $E\cap U_i$, it is the intersection of the finite unions $\bigcup_i \mathrm{M}_{i,\lambda(i)}$ (where $\lambda(i)\in\mathrm{L}_i$ for every $i$), which are constructible parts of $X$; this proves the second assertion of (iii).
The assertion (ix) follows from (iii) and \sref{IV.1.9.3}.
To prove the first assertion of (v), we may restrict ourselves to the case where $X$ is affine, and then prove that $E$ is quasi-compact \sref[0\textsubscript{III}]{0.9.1.1}; but since $X$ is then quasi-separated, there exists by virtue of (ix) a quasi-compact morphism $f:X'\to X$ such that $E=f(X')$; since $X'$ is quasi-compact, the same is true of $E$ under a continuous map.

To prove (vi), one can restrict ourselves, by virtue of (iv), to the case where $X$ is affine; the conclusion follows then from (iii) and \sref{IV.1.8.2}.

\oldpage[IV-1]{245}

To prove (vii) as well, one can restrict ourselves to the case where $X$ is affine; then $X'$ is a finite union of affine open subsets $X_i'$ and $f(E')$ is the union of the $f(E'\cap X_i')$, so one can also suppose that $X'$ is affine, by virtue of (ii); but then $E'=g(X'')$, where $g:X''\to X'$ is a quasi-compact morphism, by virtue of (ix), and one has $f(E')=f(g(X''))$, so $f\circ g$\footnote{The French source prints $g\circ f$ here; the displayed domains require $f\circ g:X''\to X$.} is quasi-compact, and therefore $f(E')$ is pro-constructible.

One deduces from (vii) the proof of the second assertion of (v): indeed, let $E$ be a locally closed and retrocompact part of $X$, and consider the sub-prescheme of $X$ having $E$ as underlying space \sref[I]{I.5.2.1}; the canonical injection $j:E\to X$ is a quasi-compact morphism, and it follows from (vii) that $E=j(E)$ is pro-constructible in $X$.

Finally, to prove (viii), one can still suppose $X$ affine; moreover, if $(X_\alpha')$ is an open affine cover of $X'$ such that $E'\cap X_\alpha'$ is union of constructible parts of $X_\alpha'$ ((iii) and (iv)), $f(E')$ is union of the $f(E'\cap X_\alpha')$, and by virtue of the theorem of Chevalley \sref{IV.1.8.4}, each $f(E'\cap X_\alpha')$ is union of constructible parts, whence the conclusion.
\end{proof}

\begin{env}[1.9.6]
\label{IV.1.9.6}
\textbf{Examples} (1.9.6). ---
Every finite part of a prescheme $X$ is pro-constructible: indeed, it suffices to consider the parts $\{x\}$ reduced to a single point \sref{IV.1.9.5}[(ii)], and $\{x\}$ is the image of $\Spec(k(x))$ by the canonical morphism $\Spec(k(x))\to X$, which is quasi-compact; whence the conclusion by \sref{IV.1.9.5}[(vii)].
On the other hand, a part $\{x\}$ reduced to a point is not necessarily ind-constructible; for example, let $A$ be a Noetherian integral ring having an infinity of maximal ideals, and let $x$ be the generic point of $X=\Spec(A)$; if $\{x\}$ were ind-constructible in $X$, it would be constructible by virtue of \sref{IV.1.9.5}[(iii)] and therefore would contain a non-empty open subset of $X$ \sref[0\textsubscript{III}]{0.9.2.2}; but this contradicts the hypothesis that $A$ has an infinity of maximal ideals, by virtue of the Artin--Tate theorem \sref[0]{0.16.3.3}.

Every closed part of a prescheme $X$ is pro-constructible, by \sref{IV.1.9.5}[(v)].
Every open part of $X$ is therefore ind-constructible, but an open part $U$ of $X$ is pro-constructible only if it is \emph{retrocompact}, by virtue again of \sref{IV.1.9.5}[(v)].

Finally, if $X$ is a \emph{Noetherian} prescheme, therefore quasi-separated \sref{IV.1.2.8}, it follows from \sref{IV.1.9.5}[(iii)] and \sref[0\textsubscript{III}]{0.9.1.7} that the ind-constructible parts of $X$ are the \emph{unions of locally closed parts} of $X$.
\end{env}

\begin{theorem}[1.9.7]
\label{IV.1.9.7}
Let $X$ be a quasi-compact prescheme, $E$ an ind-constructible part of $X$, $(\mathrm{F}_\lambda)_{\lambda\in\mathrm{L}}$ a family of pro-constructible parts of $X$ such that $\bigcap_{\lambda\in\mathrm{L}} \mathrm{F}_\lambda\subset E$; then there exists a finite part $\mathrm{J}$ of $\mathrm{L}$ such that $\bigcap_{\lambda\in\mathrm{J}} \mathrm{F}_\lambda\subset E$.
\end{theorem}

\begin{proof}
Let us write $F=X-E$; then the sets $F$ and $F\cap\mathrm{F}_\lambda$ are pro-constructible, so one is reduced to the following corollary.
\end{proof}

\begin{corollary}[1.9.8]
\label{IV.1.9.8}
Let $X$ be a quasi-compact prescheme, $(\mathrm{F}_\lambda)_{\lambda\in\mathrm{L}}$ a family of pro-constructible parts of $X$ whose intersection is empty; then there exists a finite subfamily $(\mathrm{F}_\lambda)_{\lambda\in\mathrm{J}}$ whose intersection is empty.
\end{corollary}

\begin{proof}
Covering $X$ by a finite number of affine open subsets, and using \sref{IV.1.9.5}[(iv)], one is reduced to the case where $X$ is affine; by virtue of \sref{IV.1.9.5}[(ix)], one has $\mathrm{F}_\lambda=f_\lambda(X_\lambda')$,
\oldpage[IV-1]{246}
where $X_\lambda'=\Spec(A_\lambda')$ is affine, and $f_\lambda$ is a morphism $X_\lambda'\to X$; then, with the notations of \sref{IV.1.9.3.2}, on the one hand, by hypothesis $f(X')=\varnothing$, hence $X'=\varnothing$; but this implies by \sref{IV.1.9.2}[(i)] that there exists a finite part $\mathrm{J}$ of $\mathrm{L}$ such that $X_\mathrm{J}'=\varnothing$, hence $f_\mathrm{J}(X_\mathrm{J}')=\varnothing$; but $f_\mathrm{J}(X_\mathrm{J}')=\bigcap_{\lambda\in\mathrm{J}} \mathrm{F}_\lambda=\varnothing$.
\end{proof}

\begin{corollary}[1.9.9]
\label{IV.1.9.9}
Let $X$ be a quasi-compact prescheme, $\mathrm{F}$ a pro-constructible part of $X$, $(\mathrm{E}_\lambda)_{\lambda\in\mathrm{L}}$ a family of ind-constructible parts of $X$ such that $\mathrm{F}\subset\bigcup_{\lambda\in\mathrm{L}} \mathrm{E}_\lambda$; then there exists a finite part $\mathrm{J}$ of $\mathrm{L}$ such that $\mathrm{F}\subset\bigcup_{\lambda\in\mathrm{J}} \mathrm{E}_\lambda$.
In particular, from every covering of $X$ by ind-constructible parts, one can extract a finite covering of $X$.
\end{corollary}

\begin{proof}
This follows from \sref{IV.1.9.7} by passing to complements.
\end{proof}

\begin{proposition}[1.9.10]
\label{IV.1.9.10}
Let $X$ be a prescheme.
For a part $E$ of $X$ to be ind-constructible, it is necessary that, for every $x\in X$, the intersection $E\cap\overline{\{x\}}$ be a neighbourhood of $x$ in $\overline{\{x\}}$.
If $X$ is locally Noetherian, this condition is sufficient.
\end{proposition}

\begin{proof}
Suppose $E$ ind-constructible, and let $Y$ be the intersection of $\overline{\{x\}}$ and an open affine subset of $X$ containing $x$; then there is therefore a sub-prescheme of $X$ having $Y$ for underlying space \sref[I]{I.5.2.1}, and if $j:Y\to X$ is the canonical injection, $E\cap Y=j^{-1}(E)$ is ind-constructible in $Y$ \sref{IV.1.9.5}[(vi)].
As $Y$ is quasi-compact and separated, $E\cap Y$ is therefore a union of constructible parts of $Y$ \sref{IV.1.9.5}[(iii)]; consequently, there are two open subsets $U$, $V$ in $Y$ such that $x\in U\cap\complement V\subset E\cap Y$ \sref[0\textsubscript{III}]{0.9.1.3}.
But as $x$ is the generic point of the irreducible space $Y$, $V$ is necessarily empty and we have $U\subset E\cap Y$.

Suppose now $X$ locally Noetherian, and let $E$ be a part of $X$ satisfying the condition of the statement.
By virtue of the definition \sref{IV.1.9.4}, one can restrict to the case where $X$ is \emph{Noetherian}.
Then, for every $x\in E$, $E\cap\overline{\{x\}}$ contains a non-empty open part of the form $U\cap\overline{\{x\}}$, where $U$ is open in $X$; this shows that $E$ is a union of locally closed parts of $X$, and hence is ind-constructible \sref{IV.1.9.6}.
\end{proof}

\begin{proposition}[1.9.11]
\label{IV.1.9.11}
Let $X$ be a prescheme, $E$ a part of $X$.
The following properties are equivalent:
\begin{enumerate}
  \item[\rm a)] $E$ is locally constructible.
  \item[\rm b)] $E$ is ind-constructible and pro-constructible.
  \item[\rm c)] $E$ and $X-E$ are pro-constructible.
  \item[\rm d)] $E$ and $X-E$ are ind-constructible.
\end{enumerate}
\end{proposition}

\begin{proof}
It clearly suffices to prove that $d$) implies $a$), and one can restrict to the case where $X$ is affine.
Then, by \sref{IV.1.9.5}[(iii)], $E$ (resp. $X-E$) is a union of a family $(\mathrm{E}_\lambda)$ (resp. $(\mathrm{E}_\mu')$) of constructible parts of $X$; as the $\mathrm{E}_\lambda$ and the $\mathrm{E}_\mu'$ form a covering of $X$, it follows from \sref{IV.1.9.9} that there are finitely many indices $\lambda_i$, $\mu_j$ such that the $\mathrm{E}_{\lambda_i}$ and the $\mathrm{E}_{\mu_j}'$ form a covering of $X$; this implies that $E$ is a union of the $\mathrm{E}_{\lambda_i}$, and hence is constructible.
\end{proof}

\begin{corollary}[1.9.12]
\label{IV.1.9.12}
Let $f:X\to Y$ be a surjective morphism, either quasi-compact, or locally of finite presentation.
For a part $E$ of $Y$ to be locally constructible (resp. pro-constructible, resp. ind-constructible), it is necessary and sufficient that $f^{-1}(E)$ be so in $X$.
\end{corollary}

\oldpage[IV-1]{247}

\begin{proof}
We know that the condition is necessary by \sref{IV.1.8.2} and \sref{IV.1.9.5}[(vi)]; to prove sufficiency, one is reduced to the case where $X$ is affine by virtue of \sref{IV.1.9.5}[(iv)]; moreover, by virtue of \sref{IV.1.9.11}, one can restrict to the case where $f^{-1}(E)$ is pro-constructible, \emph{or} the case where $f^{-1}(E)$ is ind-constructible.
If $f$ is surjective and quasi-compact, and $f^{-1}(E)$ is pro-constructible, it follows from \sref{IV.1.9.5}[(vii)] that $E=f(f^{-1}(E))$ is pro-constructible.
If $f$ is surjective and locally of finite presentation, and $f^{-1}(E)$ is ind-constructible, $E=f(f^{-1}(E))$ is ind-constructible by \sref{IV.1.9.5}[(viii)], which completes the proof.
\end{proof}

\begin{env}[1.9.13]
\label{IV.1.9.13}
Let $X$ be a prescheme, $\mathscr{T}$ its topology; it follows from \sref{IV.1.9.5}[(i) and (ii)] that the ind-constructible (resp. pro-constructible) parts of $X$ are the open parts (resp. closed parts) for a topology on $X$, called the \emph{constructible topology}, which we will denote by $\mathscr{T}^{\mathrm{cons}}$; we will denote by $X^{\mathrm{cons}}$ the set $X$ equipped with the topology $\mathscr{T}^{\mathrm{cons}}$.
\end{env}

\begin{proposition}[1.9.14]
\label{IV.1.9.14}
Let $X$ be a prescheme, $\mathscr{T}$ its topology, $\mathscr{T}^{\mathrm{cons}}$ the constructible topology on $X$.
\begin{enumerate}
  \item[\rm(i)] The topology $\mathscr{T}^{\mathrm{cons}}$ is finer than $\mathscr{T}$.
  \item[\rm(ii)] The locally constructible parts of $X$ are identical to the parts that are both open and closed in the space $X^{\mathrm{cons}}$.
  \item[\rm(iii)] For every morphism $f:X\to Y$, the underlying map of $X^{\mathrm{cons}}$ to $Y^{\mathrm{cons}}$ is continuous; we denote it $f^{\mathrm{cons}}$.
  \item[\rm(iv)] If the morphism $f:X\to Y$ is quasi-compact, $f^{\mathrm{cons}}$ is a closed map; in particular, if $f$ is quasi-compact and bijective, $f^{\mathrm{cons}}$ is a homeomorphism.
  \item[\rm(v)] If the morphism $f:X\to Y$ is locally of finite presentation, $f^{\mathrm{cons}}$ is an open map; in particular, if $f$ is bijective and locally of finite presentation, $f^{\mathrm{cons}}$ is a homeomorphism.
  \item[\rm(vi)] For every open subset $U$ of $X$, the topology induced by $\mathscr{T}^{\mathrm{cons}}$ on $U$ is identical to the topology of $U^{\mathrm{cons}}$.
\end{enumerate}
\end{proposition}

\begin{proof}
Indeed, (i) follows from the fact that every open set for $\mathscr{T}$ is open for $\mathscr{T}^{\mathrm{cons}}$ \sref{IV.1.9.6}, and (ii) is a restatement of \sref{IV.1.9.11}; (iii), (iv), and (v) are restatements respectively of \sref{IV.1.9.5}[(vi)], \sref{IV.1.9.5}[(vii)], and \sref{IV.1.9.5}[(viii)].
Finally, to prove (vi), it suffices to remark that the canonical injection $j:U\to X$ is locally of finite presentation \sref{IV.1.4.3}, hence $j^{\mathrm{cons}}:U^{\mathrm{cons}}\to X^{\mathrm{cons}}$ is a continuous open map.
\end{proof}

\begin{proposition}[1.9.15]
\label{IV.1.9.15}
Let $X$ be a prescheme.
\begin{enumerate}
  \item[\rm(i)] For $X^{\mathrm{cons}}$ to be quasi-compact, it is necessary and sufficient that $X$ be quasi-compact.
  \item[\rm(ii)] For $X^{\mathrm{cons}}$ to be separated, it is necessary and sufficient that $X$ be quasi-separated, and then $X^{\mathrm{cons}}$ is locally compact and totally disconnected.
  \item[\rm(iii)] For $X^{\mathrm{cons}}$ to be compact, it is necessary and sufficient that $X$ be quasi-compact and quasi-separated.
  \item[\rm(iv)] Every point of $X$ admits, for the topology $\mathscr{T}^{\mathrm{cons}}$, an open and compact neighbourhood.
  \item[\rm(v)] For a morphism $f:X\to Y$ to be quasi-compact, it is necessary and sufficient that the map $f^{\mathrm{cons}}:X^{\mathrm{cons}}\to Y^{\mathrm{cons}}$ be proper.
\end{enumerate}
\end{proposition}

\begin{proof}
(i) As $\mathscr{T}^{\mathrm{cons}}$ is finer than $\mathscr{T}$, it is clear that if $X^{\mathrm{cons}}$ is quasi-compact, so is $X$; the converse follows from \sref{IV.1.9.9}.

(ii) Suppose $X$ quasi-separated, and let us show that $X^{\mathrm{cons}}$ is totally disconnected: indeed, if $x$, $y$ are two distinct points of $X$ and if, for example, $x\notin\overline{\{y\}}$, there exists an
\oldpage[IV-1]{248}
affine open neighbourhood $U$ of $x$ which does not contain $y$; as $X$ is quasi-separated, $U$ is retrocompact in $X$ \sref{IV.1.2.7}, hence $U$ and $X-U$ are locally constructible in $X$, and therefore open in $X^{\mathrm{cons}}$ by virtue of \sref{IV.1.9.11}, whence our assertion.
For every open affine subset $U$ of $X$, the topology induced by $\mathscr{T}^{\mathrm{cons}}$ is that of $U^{\mathrm{cons}}$ by virtue of \sref{IV.1.9.14}[(vi)]; it follows from the above that $X^{\mathrm{cons}}$ is locally compact, since $U$ is open in $X^{\mathrm{cons}}$ and $U^{\mathrm{cons}}$ is compact.
It remains to show that if $X^{\mathrm{cons}}$ is separated, then $X$ is quasi-separated; consider indeed an open affine subset $U$ of $X$; the topology induced on $U$ by $\mathscr{T}^{\mathrm{cons}}$ is that of $U^{\mathrm{cons}}$ \sref{IV.1.9.14}[(vi)], hence $U$ is compact for this induced topology, since $U^{\mathrm{cons}}$ is compact by the first part of the reasoning.
If $V$ is a second open affine subset in $X$, $U\cap V$ is therefore a compact subset of the separated space $X^{\mathrm{cons}}$, being the intersection of two compact subsets of this space; as the topology induced by $\mathscr{T}^{\mathrm{cons}}$ on $U\cap V$ is that of $(U\cap V)^{\mathrm{cons}}$ \sref{IV.1.9.14}[(vi)] and the identity $(U\cap V)^{\mathrm{cons}}\to U\cap V$ is continuous, one concludes that $U\cap V$ is quasi-compact for the topology induced by $\mathscr{T}$; $X$ is therefore quasi-separated by virtue of \sref{IV.1.2.7}.

(iii) is an immediate corollary of (i) and (ii).

(iv) For every $x\in X$, an open affine neighbourhood $U$ of $x$ for $\mathscr{T}$ is also a neighbourhood of $x$ for $\mathscr{T}^{\mathrm{cons}}$ and it is compact by virtue of (iii) and \sref{IV.1.9.14}[(vi)], the topology induced on $U$ by $\mathscr{T}^{\mathrm{cons}}$ being identical to that of $U^{\mathrm{cons}}$.

(v) Suppose $f$ quasi-compact; then we already know \sref{IV.1.9.14}[(iv)] that $f^{\mathrm{cons}}$ is a closed map.
Let on the other hand $y$ be a point of $Y$, and let us set
\[
  Z=f^{-1}(y)=X\times_Y\Spec(k(y));
\]
$Z$ is quasi-compact \sref{IV.1.1.2}[(iii)] and, as the canonical morphism $p:Z\to X$ is injective, it follows from (i) and from the fact that the map $p^{\mathrm{cons}}$ is continuous that the topology induced by that of $X^{\mathrm{cons}}$ on $f^{-1}(y)$ makes $f^{-1}(y)$ a quasi-compact space.
This proves that $f^{\mathrm{cons}}$ is a proper map (Bourbaki, \emph{Top.\ g\'{e}n.}, chap.~I\textsuperscript{er}, 3\textsuperscript{e}~\'{e}d., \S\,10, n\textsuperscript{o}\,2, th.~1).
Conversely, suppose the continuous map $f^{\mathrm{cons}}$ proper, and let $V$ be an open quasi-compact subset of $Y$; if $h:V\to Y$ is the canonical injection, $h^{\mathrm{cons}}:V^{\mathrm{cons}}\to Y^{\mathrm{cons}}$ is continuous and injective and $V^{\mathrm{cons}}$ is quasi-compact by (i), hence the topology induced on $V$ by that of $Y^{\mathrm{cons}}$ makes $V$ a quasi-compact space.
The hypothesis that $f^{\mathrm{cons}}$ is proper implies that the topology induced on $f^{-1}(V)$ by that of $X^{\mathrm{cons}}$ makes $f^{-1}(V)$ a quasi-compact space (\emph{loc.\ cit.}, prop.~6), hence $f^{-1}(V)$ is also a quasi-compact subspace of $X$, which shows that the morphism $f$ is quasi-compact.
\end{proof}

\begin{env}[1.9.16]
\label{IV.1.9.16}
We will show later how one can, for every prescheme $X$, equip the space $X^{\mathrm{cons}}$ with a sheaf of rings which makes it a prescheme whose local rings are all fields, identical to the residue fields at the points of $X$.
Such preschemes arise naturally in the translation, into the language of schemes, of the constructions of N\'{e}ron~\cite{IV-32} in his theory of the reduction of abelian varieties.
\end{env}

\oldpage[IV-1]{249}

\subsection{Application to open morphisms}
\label{subsection:IV.1.10}

\begin{theorem}[1.10.1]
\phantomsection
\label{IV.1.10.1}
Let $X$ be a prescheme, $E$ an ind-constructible part of $X$ \sref{IV.1.9.4}, $x$ a point of $X$.
For $x$ to be interior to $E$, it is necessary and sufficient that every generisation \sref[0\textsubscript{I}]{0.2.1.2} $x'$ of $x$ belong to $E$, in other words \sref[I]{I.2.4.2} that $\Spec(\mathcal{O}_{X,x})\subset E$.
\end{theorem}

\begin{proof}
The condition is evidently necessary, since every neighbourhood of $x$ contains the generisations of $x$.
To show that it is sufficient, we may evidently restrict to the case where $X=\Spec(A)$ is affine, $x$ being a prime ideal $\mathfrak{p}$ of $A$.
We know then \sref{IV.1.9.5}[(ix)] that there exists an affine scheme $X'=\Spec(A')$ and a morphism $f:X'\to X$ such that $f(X')=X-E$.
Let $Y=\Spec(\mathcal{O}_{X,x})$ and $Y'=X'\times_X Y$; the hypothesis means (by virtue of \sref[I]{I.3.6.5}) that $Y'=\varnothing$, in other words $A'\otimes_A A_\mathfrak{p}=0$.
As $A_\mathfrak{p}=\varinjlim A_t$, where $t$ runs through $A-\mathfrak{p}$ \sref[0\textsubscript{I}]{0.1.4.5}, we have $\varinjlim A_t'=0$, since the functor $\varinjlim$ commutes with the tensor product.
Consequently, by \sref{IV.1.9.1}, there exists $t\in A-\mathfrak{p}$ such that $A_t'=0$, and $\mathrm{D}(t)$ is then an open neighbourhood of $x$ in $X$ contained in $E$.
\end{proof}

\begin{corollary}[1.10.2]
\phantomsection
\label{IV.1.10.2}
Let $f:X\to Y$ be a morphism locally of finite presentation, and let $y\in Y$.
For $y$ to be interior to $f(X)$, it is necessary and sufficient that every generisation $y'$ of $y$ belong to $f(X)$.
\end{corollary}

\begin{proof}
We know indeed \sref{IV.1.9.5}[(viii)] that $f(X)$ is ind-constructible in $Y$ and it suffices to apply \sref{IV.1.10.1}.
\end{proof}

We will say that a continuous map $f:X\to Y$ is \emph{open at a point} $x\in X$ if the image under $f$ of \emph{every} neighbourhood of $x$ in $X$ is a neighbourhood of $f(x)$ in $Y$.

\begin{proposition}[1.10.3]
\phantomsection
\label{IV.1.10.3}
Let $f:X\to Y$ be a morphism, $x$ a point of $X$, $y=f(x)$.
Consider the following conditions:
\begin{enumerate}
  \item[\rm a)] $f$ is open at the point $x$.
  \item[\rm b)] For every generisation $y'$ of $y$, there exists a generisation $x'$ of $x$ such that $f(x')=y'$.
  \item[\rm b$'$)] The image under $f$ of $\Spec(\mathcal{O}_{X,x})$ is $\Spec(\mathcal{O}_{Y,y})$.
  \item[\rm c)] For every irreducible closed part $Y'$ of $Y$ containing $y$, there exists an irreducible component of $X'=f^{-1}(Y')$ containing $x$ and dominating $Y'$.
\end{enumerate}
Then we have the implications $\mathrm{a)}\Rightarrow\mathrm{b)}\Leftrightarrow\mathrm{b')}\Leftrightarrow\mathrm{c)}$.
If in addition $f$ is locally of finite presentation, the four conditions are equivalent.
\end{proposition}

\begin{proof}
By definition of a generisation, the image under $f$ of $\Spec(\mathcal{O}_{X,x})$ is always contained in $\Spec(\mathcal{O}_{Y,y})$, so b) and b$'$) are equivalent.
It is immediate that c) implies b), since if $y'$ is a generisation of $y$, $Y'=\overline{\{y'\}}$ is an irreducible closed part of $Y$ containing $y$; if $x'$ is the generic point of an irreducible component of $f^{-1}(Y')$ containing $x$ and dominating $Y'$, we have $f(x')=y'$ \sref[0\textsubscript{I}]{0.2.1.5} and $x'$ is a generisation of $x$.
Conversely, b) implies c): indeed, let $y'$ be the generic point of $Y'$ and let $x'$ be a generisation of $x$ such that $f(x')=y'$; let $Z'$ be an irreducible component of $f^{-1}(Y')$ containing $x'$ (and hence $x$); as $f(x')$ is already dense in $Y'$, so \emph{a fortiori} is $f(Z')$.

To show that a) implies b$'$), one can evidently restrict to the case where $X=\Spec(B)$ and $Y=\Spec(A)$ are affine, $x$ being a prime ideal $\mathfrak{q}$ of $B$, so that $\mathcal{O}_{X,x}=B_\mathfrak{q}=\varinjlim B_t$,
\oldpage[IV-1]{250}
where $t$ runs through $B-\mathfrak{q}$.
One has, by \sref{IV.1.9.2}, $f(\Spec(\mathcal{O}_{X,x}))=\bigcap_t f(\Spec(B_t))=\bigcap_t f(\mathrm{D}(t))$, and by hypothesis, for every $t\in B-\mathfrak{q}$, $y$ is interior to $f(\mathrm{D}(t))$, hence $f(\mathrm{D}(t))$ contains $\Spec(\mathcal{O}_{Y,y})$; whence $\Spec(\mathcal{O}_{Y,y})\subset f(\Spec(\mathcal{O}_{X,x}))$, which establishes our assertion.
Finally, when $f$ is locally of finite presentation, b) implies a) by virtue of \sref{IV.1.10.2} applied to the restriction $U\to Y$ of $f$ to an arbitrary open neighbourhood $U$ of $x$.
\end{proof}

\begin{corollary}[1.10.4]
\phantomsection
\label{IV.1.10.4}
Let $f:X\to Y$ be a morphism.
Consider the following conditions:
\begin{enumerate}
  \item[\rm a)] $f$ is open.
  \item[\rm b)] For every $x\in X$ and every generisation $y'$ of $y=f(x)$, there exists a generisation $x'$ of $x$ such that $f(x')=y'$.
  \item[\rm b$'$)] For every $x\in X$, the image under $f$ of $\Spec(\mathcal{O}_{X,x})$ is $\Spec(\mathcal{O}_{Y,f(x)})$.
  \item[\rm c)] For every irreducible closed part $Y'$ of $Y$, every irreducible component of $f^{-1}(Y')$ dominates $Y'$.
\end{enumerate}
Then we have the implications $\mathrm{a)}\Rightarrow\mathrm{b)}\Leftrightarrow\mathrm{b')}\Leftrightarrow\mathrm{c)}$.
If in addition $f$ is locally of finite presentation, the four conditions are equivalent.
\end{corollary}

\begin{proof}
To say that $f$ is open means that $f$ is open at every point $x\in X$, hence the implications $\mathrm{a)}\Rightarrow\mathrm{b)}\Leftrightarrow\mathrm{b')}$ follow from the analogous implications in \sref{IV.1.10.3}, as well as the implication $\mathrm{b)}\Rightarrow\mathrm{a)}$ when $f$ is locally of finite presentation.
The implication $\mathrm{c)}\Rightarrow\mathrm{b)}$ also follows from the analogous implication in \sref{IV.1.10.3}; let us finally prove that b) implies c).
Indeed, let $x'$ be the generic point of an irreducible component of $f^{-1}(Y')$, and let us show that $y'=f(x')$ is the generic point of $Y'$.
Let $y''$ be a generisation of $y'$; by hypothesis there exists a generisation $x''$ of $x'$ such that $f(x'')=y''$, and as $x''\in f^{-1}(Y')$, we have necessarily $x''=x'$, hence $y''=y'$, which completes the proof.
\end{proof}

\begin{center}
\emph{(To be continued.)}
\end{center}

\clearpage
\oldpage[IV-1]{251}
\phantomsection
\addcontentsline{toc}{section}{Bibliography (continued)}
\section*{BIBLIOGRAPHY (\emph{continued})}

\begingroup
\small
\setlength{\parindent}{0pt}
\setlength{\parskip}{3pt}
[30] M.~Nagata, \emph{Local rings}, \emph{Interscience Tracts}, 13 (1962),
Interscience, New York.

[31] M.~Nagata, \emph{A Jacobian criterion of simple points},
\emph{Ill. Journ. of Math.}, 1 (1957), pp.~427--432.

[32] A.~Néron, \emph{Modèles minimaux des variétés abéliennes sur les corps
locaux et globaux}, \emph{Publ. Inst. Hautes Études Sci.}, no.~21 (1964).

\clearpage
\oldpage[IV-1]{252}
\phantomsection
\addcontentsline{toc}{section}{Index of notation}
\section*{INDEX OF NOTATION}

$\dim_c(X)$, $\dim(X)$ ($X$ a topological space) : 0, 14.1.2.\par
$\dim_x(X)$ ($X$ a topological space, $x$ a point of $X$) : 0, 14.1.2.\par
$\operatorname{codim}(Y,X)$ ($X$ a topological space, $Y$ a closed subset of $X$) : 0, 14.2.1.\par
$\operatorname{codim}_x(Y,X)$ ($X$ a topological space, $Y$ a closed subset of $X$, $x$ a point of $X$) : 0, 14.2.4.\par
$N[T_1,\ldots,T_n]$ ($N$ an $A$-module) : 0, 15.1.1.\par
$\dim(A)$ ($A$ a ring) : 0, 16.1.1.\par
$\operatorname{ht}(\mathfrak a)$ ($\mathfrak a$ an ideal) : 0, 16.1.3.\par
$\dim_A(M)$, $\dim(M)$ ($M$ an $A$-module) : 0, 16.1.7.\par
$P_{\mathfrak q}(M,n)$ ($M$ a finite type module over a noetherian semilocal ring $A$, $\mathfrak q$ an ideal of definition of $A$) : 0, 16.2.1.\par
$\operatorname{prof}_A(M)$, $\operatorname{prof}(M)$ ($M$ an $A$-module) : 0, 16.4.5.\par
$\operatorname{coprof}_A(M)$, $\operatorname{coprof}(M)$ ($M$ an $A$-module) : 0, 16.4.9.\par
$\operatorname{dim.proj}_A(M)$, $\operatorname{dim.proj}(M)$,
$\operatorname{dim.inj}_A(M)$, $\operatorname{dim.inj}(M)$ ($M$ an $A$-module) : 0, 17.2.1.\par
$\operatorname{dim.coh}(A)$ ($A$ a ring) : 0, 17.2.8.\par
$\operatorname{dim.proj}(\mathscr F)$, $\operatorname{dim.inj}(\mathscr F)$
($\mathscr F$ a module on a ringed space) : 0, 17.2.14.\par
$\operatorname{dim.coh}(X)$ ($X$ a ringed space) : 0, 17.2.14.\par
$E\times_B F$ ($B$, $E$, $F$ $A$-rings) : 0, 18.1.2.\par
$D_B(L)$ ($B$ a ring, $L$ a $B$-bimodule) : 0, 18.2.3.\par
$Y\oplus_X Z$ ($X$, $Y$, $Z$ $A$-bimodules) : 0, 18.2.7.\par
$\operatorname{Exan}_A(B,L)$ ($B$ an $A$-ring, $L$ a $B$-bimodule) : 0, 18.3.4.\par
$\operatorname{Exan}_{A/A'}(B,L)$ : 0, 18.3.7.\par
$\operatorname{Exal}_A(B,L)$, $\operatorname{Exal}_{A/A'}(B,L)$ ($A$ a commutative ring, $B$ an $A$-algebra, $L$ a $B$-bimodule) : 0, 18.4.1.\par
$\operatorname{Exalcom}_A(B,L)$, $\operatorname{Exalcom}_{A/A'}(B,L)$ ($A$ a commutative ring, $B$ a commutative $A$-algebra, $L$ a $B$-module) : 0, 18.4.2.\par
$\mathrm H_A^2(B,L)$, $\mathrm H_A^2(B,L)^s$ : 0, 18.4.3.\par
$\operatorname{Exantop}_A(B,L)$, $\operatorname{Exaltop}_A(B,L)$,
$\operatorname{Exalcotop}_A(B,L)$ : 0, 18.5.1.\par
$\operatorname{Exantop}_{A/A'}(B,L)$, $\operatorname{Exaltop}_{A/A'}(B,L)$,
$\operatorname{Exalcotop}_{A/A'}(B,L)$ : 0, 18.5.2.\par
$\operatorname{Hom.\ cont}_C(M,N)$ ($M$, $N$ topological $C$-bimodules) : 0, 18.5.3.\par
$\operatorname{Der}_A(B,L)$ ($B$ an $A$-ring, $L$ a $B$-bimodule) : 0, 20.1.2.\par
$\operatorname{Der.\ cont}_A(B,L)$ : 0, 20.3.1.\par
$p_{B/A}$, $\mathfrak J_{B/A}$ ($B$ a commutative topological $A$-algebra) : 0, 20.4.1.\par
$P^1_{B/A}$, $\varepsilon_{B/A}$ ($B$ a commutative topological $A$-algebra) : 0, 20.4.2.\par
$\Omega^1_{B/A}$, $\widehat{\Omega}^1_{B/A}$ ($B$ a commutative topological $A$-algebra) : 0, 20.4.3.\par
$\Omega^1_B$ ($B$ a topological ring) : 0, 20.4.3.\par
$d_{B/A}$, $d(x)$, $dx$, $d_B$ : 0, 20.4.6.\par
$\Upsilon_{C/B/A}$, $\Upsilon_{C/B}$ ($A$, $B$, $C$ rings) : 0, 20.6.1.\par
$\chi_E$ ($E$ a $B$-extension of $C$ by $L$ that is $A$-trivial) : 0, 20.6.8.\par

\clearpage
\oldpage[IV-1]{253}
$\Upsilon^C_{B/A/\Lambda}$, $\Upsilon^C_{B/A}$ ($\Lambda$, $A$, $B$, $C$ rings) : 0, 20.6.14.\par
$\chi_{B/A}$, $\chi_B$ ($B$ an $A$-algebra) : 0, 20.6.24.\par
$\chi_E$ ($E$ a topological $B$-algebra) : 0, 20.7.10.\par
$F_A$, $A^p$, $A^{(p)}$, $E^{(p)}$ ($A$ a ring of characteristic $p>0$, $E$ an $A$-module) : 0, 21.1.4.\par
$\pi_{B/A}$, $\Theta_{B/A}$, $\Xi_{B/A}$ ($A$, $B$ rings of characteristic $p>0$) : 0, 21.3.2.\par
$\Delta(L/K,k/k_0)$, $d(L/K,k/k_0)$, $\Delta(L/K,k)$, $d(L/K,k)$ ($k_0$, $k$, $K$, $L$ fields) : 0, 21.6.1.\par
$\chi'_{B/A}$ ($A$, $B$ rings of characteristic $p>0$) : 0, 22.4.6.\par
$\mathscr T^{\mathrm{cons}}$, $X^{\mathrm{cons}}$ ($X$ a prescheme, $\mathscr T$ the topology of $X$) : IV, 1.9.12.\par
$f^{\mathrm{cons}}$ ($f$ a morphism of preschemes) : IV, 1.9.13.\par

\clearpage
\oldpage[IV-1]{254}
\phantomsection
\addcontentsline{toc}{section}{Terminological index}
\section*{TERMINOLOGICAL INDEX}

Augmented algebra of principal parts of order 1 : 0, 20.4.2.\par
Cohen $A$-algebra : 0, 19.8.1.\par
Algebra of finite presentation : IV, 1.4.1.\par
Essentially finite type $A$-algebra : IV, 1.3.8.\par
$A$-ring : 0, 18.1.1.\par
$A$-ring augmented over $B$ : 0, 18.1.4.\par
Biequidimensional ring : 0, 16.1.4.\par
Catenary ring : 0, 16.1.4.\par
Ring of characteristic $p$ : 0, 21.1.1.\par
Cohen ring : 0, 19.8.4.\par
Cohen-Macaulay ring : 0, 16.5.1 and 16.5.13.\par
Equicodimensional ring, equidimensional ring : 0, 16.1.4.\par
Geometrically unibranch local ring : 0, 23.1.7.\par
Japanese ring : 0, 23.1.1.\par
Primary local ring, complete primary local ring : 0, 19.8.3.\par
Regular ring : 0, 17.1.1 and 17.3.6.\par
Unibranch local ring : 0, 23.1.7.\par
Universally Japanese ring : 0, 23.1.1.\par
Open map at a point : IV, 1.10.2.\par
Augmentation : 0, 18.1.4.\par

Formal bimorphism : 0, 19.1.2.\par

Chain in an ordered set, chain of irreducible closed subsets of a topological space : 0, 14.1.1.\par
Saturated chain : 0, 14.3.1.\par
Combinatorial codimension, codimension of a closed subset of a topological space : 0, 14.2.1.\par
Codimension at a point of a closed subset of a topological space : 0, 14.2.4.\par
Chain condition : 0, 14.3.2.\par
Codepth of a module over a noetherian local ring : 0, 16.4.9.\par
$k_0$-admissible field, field admissible for an extension : 0, 21.6.1.\par
Field of finite radical multiplicity over a subfield : 0, 19.6.6.\par

Defect of $k_0$-admissibility of an extension : 0, 21.6.1.\par
Derivation, $A$-derivation : 0, 20.1.2.\par
Differentials, 1-differentials, absolute differentials : 0, 20.4.3.\par
Exterior differential : 0, 20.4.6.\par
Combinatorial dimension, dimension of a topological space : 0, 14.1.2.\par
Dimension of a topological space at a point : 0, 14.1.2.\par
Krull dimension, dimension of a ring : 0, 16.1.1.\par
Dimension of a module : 0, 16.1.7.\par
Cohomological dimension, global cohomological dimension of a ring : 0, 17.2.8.\par
Cohomological dimension of a ringed space at a point, cohomological dimension of a ringed space : 0, 17.2.14.\par

\clearpage
\oldpage[IV-1]{255}
Injective dimension, projective dimension of a module : 0, 17.2.1.\par
Injective dimension at a point, injective dimension of a sheaf : 0, 17.2.14.\par
Projective dimension at a point, projective dimension of a sheaf : 0, 17.2.14.\par

Element of $A$ that is not a zero divisor in an $A$-module : 0, 15.1.1.\par
$M$-regular element, $M$-quasi-regular element : 0, 15.1.4.\par
Ind-constructible set, pro-constructible set : IV, 1.9.4.\par
Formal epimorphism : 0, 19.1.2.\par
$A$-equivalence of $A$-extensions : 0, 18.2.2.\par
Biequidimensional space : 0, 14.3.3.\par
Catenary space : 0, 14.3.2.\par
Equicodimensional space : 0, 14.2.1.\par
Equidimensional space : 0, 14.1.3.\par
$A$-extension of an $A$-ring $B$ by a $B$-bimodule : 0, 18.2.2.\par
Extension deduced from another by a bimodule homomorphism : 0, 18.2.8.\par
Hochschild extension : 0, 18.4.3.\par
$B$-extension that is $A$-trivial : 0, 18.3.7.\par
Standard trivial extension : 0, 18.2.3.\par
$A$-equivalent $A$-extensions : 0, 18.2.2.\par

$p$-free family, absolutely $p$-free family : 0, 21.1.9.\par
Formally étale (algebra) : 0, 19.10.2.\par
Formally right-invertible (homomorphism) : 0, 19.1.15.\par
Formally left-invertible (homomorphism) : 0, 19.1.5.\par
Formally unramified (algebra) : 0, 19.10.2.\par
Formally projective (module) : 0, 19.2.1.\par
Formally smooth (algebra) : 0, 19.3.1.\par
Formally smooth relative to a ring (algebra) : 0, 19.9.1.\par

Geometrically regular over a field (local ring) : 0, 19.6.5.\par
Group of classes of $A$-extensions : 0, 18.3.4.\par

Height of an ideal : 0, 16.1.3.\par
$A$-homomorphism of $A$-rings : 0, 18.1.1.\par
Characteristic homomorphism of a $B$-extension that is $A$-trivial : 0, 20.6.8.\par
Characteristic homomorphism of an $A$-algebra relative to a ring and an ideal : 0, 20.6.24.\par
Characteristic homomorphism of an algebra of characteristic $p>0$ : 0, 22.4.6.\par

Augmentation ideal : 0, 18.1.4.\par
Inverse image of an augmented ring : 0, 18.1.5.\par
Inverse image of an $A$-extension : 0, 18.2.5.\par

Length of a chain : 0, 14.1.1.\par

Biequidimensional, catenary, equicodimensional, equidimensional module : 0, 16.1.7.\par
Cohen-Macaulay module : 0, 16.5.1 and 16.5.13.\par
Imperfection module : 0, 20.6.1.\par
Formal monomorphism : 0, 19.1.2.\par
Morphism of $A$-extensions : 0, 18.2.4.\par
Morphism of finite presentation at a point, morphism locally of finite presentation : IV, 1.4.2.\par
Morphism of finite presentation : IV, 1.6.1.\par
Quasi-separated morphism : IV, 1.2.1.\par

Cohen $p$-ring : 0, 19.8.4.\par
$p$-basis, absolute $p$-basis : 0, 21.1.9.\par
Maximal point of a prescheme : IV, 1.1.4.\par

\clearpage
\oldpage[IV-1]{256}
Hilbert-Samuel polynomial : 0, 16.2.1.\par
Prescheme of finite presentation over another : IV, 1.6.1.\par
Quasi-separated prescheme : IV, 1.2.1.\par
Fiber product of $A$-rings : 0, 18.1.2.\par
Depth of a module : 0, 16.4.5.\par

Amalgamated sum of bimodules : 0, 18.2.7.\par
Strictly formally projective (module) : 0, 19.2.3.\par
$M$-regular sequence, $M$-quasi-regular sequence : 0, 15.1.7.\par
$\mathscr F$-regular sequence, $\mathscr F$-quasi-regular sequence : 0, 15.2.2.\par
System of parameters for a module over a noetherian semilocal ring : 0, 16.3.6.\par
System of $p$-generators, absolute system of $p$-generators : 0, 21.1.9.\par
Fundamental system of open ideals, of open submodules : 0, 19.0.2.\par
Regular system of parameters of a regular local ring : 0, 17.1.6.\par

Constructible topology on a prescheme : IV, 1.9.12.\par
Topology on an $A$-module deduced from the topology of $A$ : 0, 19.0.2.\par
Trivial (augmented $A$-ring) : 0, 18.1.4.\par
$A$-trivial (extension) : 0, 18.2.3 and 18.3.7.\par

\clearpage
\oldpage[IV-1]{257}
\phantomsection
\addcontentsline{toc}{section}{Table of contents}
\section*{TABLE OF CONTENTS}

\newcommand{\ivonecontentsline}[3]{\noindent\makebox[17mm][l]{#1}\parbox[t]{0.77\linewidth}{#2\dotfill}\makebox[12mm][r]{#3}\par}
\ivonecontentsline{Chapter 0.}{\textbf{Preliminaries} (\emph{continued})}{5}
\ivonecontentsline{\S~14.}{Combinatorial dimension of a topological space}{6}
\ivonecontentsline{14.1.}{Combinatorial dimension of a topological space}{6}
\ivonecontentsline{14.2.}{Codimension of a closed subset}{8}
\ivonecontentsline{14.3.}{The chain condition}{10}
\ivonecontentsline{\S~15.}{$M$-regular sequences and $\mathscr F$-regular sequences}{12}
\ivonecontentsline{15.1.}{$M$-regular sequences and $M$-quasi-regular sequences}{12}
\ivonecontentsline{15.2.}{$\mathscr F$-regular sequences}{20}
\ivonecontentsline{\S~16.}{Dimension and depth in noetherian local rings}{22}
\ivonecontentsline{16.1.}{Dimension of a ring}{22}
\ivonecontentsline{16.2.}{Dimension of a noetherian semilocal ring}{25}
\ivonecontentsline{16.3.}{Systems of parameters in a noetherian local ring}{28}
\ivonecontentsline{16.4.}{Depth and codepth}{32}
\ivonecontentsline{16.5.}{Cohen-Macaulay modules}{36}
\ivonecontentsline{\S~17.}{Regular rings}{39}
\ivonecontentsline{17.1.}{Definition of regular rings}{39}
\ivonecontentsline{17.2.}{Review of the projective and injective dimensions of modules}{42}
\ivonecontentsline{17.3.}{Cohomological theory of regular rings}{46}
\ivonecontentsline{\S~18.}{Complements on extensions of algebras}{51}
\ivonecontentsline{18.1.}{Inverse images of augmented rings}{51}
\ivonecontentsline{18.2.}{Extensions of a ring by a bimodule}{54}
\ivonecontentsline{18.3.}{The group of classes of $A$-extensions}{59}
\ivonecontentsline{18.4.}{Extensions of algebras}{64}
\ivonecontentsline{18.5.}{The case of topological rings}{66}
\ivonecontentsline{\S~19.}{Formally smooth algebras and Cohen rings}{69}
\ivonecontentsline{19.0.}{Introduction}{69}
\ivonecontentsline{19.1.}{Formal epimorphisms and monomorphisms}{71}
\ivonecontentsline{19.2.}{Formally projective modules}{78}

\clearpage
\oldpage[IV-1]{258}
\ivonecontentsline{19.3.}{Formally smooth algebras}{79}
\ivonecontentsline{19.4.}{First criteria of formal smoothness}{86}
\ivonecontentsline{19.5.}{Formal smoothness and associated graded rings}{90}
\ivonecontentsline{19.6.}{The case of algebras over a field}{100}
\ivonecontentsline{19.7.}{The case of local homomorphisms: existence and uniqueness theorems}{104}
\ivonecontentsline{19.8.}{Cohen algebras and Cohen $p$-rings; application to the structure of complete local rings}{109}
\ivonecontentsline{19.9.}{Relatively formally smooth algebras}{114}
\ivonecontentsline{19.10.}{Formally unramified algebras and formally étale algebras}{115}
\ivonecontentsline{\S~20.}{Derivations and differentials}{116}
\ivonecontentsline{20.1.}{Derivations and extensions of algebras}{117}
\ivonecontentsline{20.2.}{Functorial properties of derivations}{119}
\ivonecontentsline{20.3.}{Continuous derivations in topological rings}{121}
\ivonecontentsline{20.4.}{Principal parts and differentials}{123}
\ivonecontentsline{20.5.}{Fundamental functorial properties of $\Omega^1_{B/A}$}{128}
\ivonecontentsline{20.6.}{Imperfection modules and characteristic homomorphisms}{136}
\ivonecontentsline{20.7.}{Generalizations to topological rings}{147}
\ivonecontentsline{\S~21.}{Differentials in rings of characteristic $p$}{153}
\ivonecontentsline{21.1.}{Systems of $p$-generators and $p$-bases}{154}
\ivonecontentsline{21.2.}{$p$-bases and formal smoothness}{157}
\ivonecontentsline{21.3.}{$p$-bases and imperfection modules}{160}
\ivonecontentsline{21.4.}{The case of field extensions}{162}
\ivonecontentsline{21.5.}{Application: separability criteria}{164}
\ivonecontentsline{21.6.}{Fields admissible for an extension}{167}
\ivonecontentsline{21.7.}{Cartier's equality}{169}
\ivonecontentsline{21.8.}{Admissibility criteria}{171}
\ivonecontentsline{21.9.}{Completed modules of differentials in formal power series rings}{176}
\ivonecontentsline{\S~22.}{Differential criteria of formal smoothness and regularity}{182}
\ivonecontentsline{22.1.}{Lifting formal smoothness}{183}
\ivonecontentsline{22.2.}{Differential characterization of local algebras formally smooth over a field}{186}
\ivonecontentsline{22.3.}{Application to relations between certain local rings and their completions}{191}
\ivonecontentsline{22.4.}{Preliminary results on finite extensions of local rings whose maximal ideal has square zero}{193}

\clearpage
\oldpage[IV-1]{259}
\ivonecontentsline{22.5.}{Geometrically regular algebras and formally smooth algebras}{201}
\ivonecontentsline{22.6.}{Zariski's Jacobian criterion}{205}
\ivonecontentsline{22.7.}{Nagata's Jacobian criterion}{209}
\ivonecontentsline{\S~23.}{Japanese rings}{213}
\ivonecontentsline{23.1.}{Japanese rings}{213}
\ivonecontentsline{23.2.}{Integral closure of a noetherian integral local ring}{217}
\ivonecontentsline{}{\textsc{Chapter IV.} --- \textbf{Local study of schemes and morphisms of schemes}}{222}
\ivonecontentsline{\S~1.}{Relative finiteness conditions. Constructible sets in preschemes}{224}
\ivonecontentsline{1.1.}{Quasi-compact morphisms}{224}
\ivonecontentsline{1.2.}{Quasi-separated morphisms}{226}
\ivonecontentsline{1.3.}{Morphisms locally of finite type}{228}
\ivonecontentsline{1.4.}{Morphisms locally of finite presentation}{230}
\ivonecontentsline{1.5.}{Morphisms of finite type}{233}
\ivonecontentsline{1.6.}{Morphisms of finite presentation}{234}
\ivonecontentsline{1.7.}{Improvement of earlier results}{236}
\ivonecontentsline{1.8.}{Morphisms of finite presentation and constructible sets}{238}
\ivonecontentsline{1.9.}{Pro-constructible and ind-constructible sets}{241}
\ivonecontentsline{1.10.}{Applications to open morphisms}{249}
\ivonecontentsline{}{Bibliography (\emph{continued})}{251}
\ivonecontentsline{}{Index of notation}{252}
\ivonecontentsline{}{Terminological index}{254}

\vspace{2em}
\begin{flushright}
\emph{Received March 15, 1963.}
\end{flushright}
\let\ivonecontentsline\relax
\endgroup
% END NATIVE INPUT ega4-1.source_aligned_r1

\clearpage
\phantomsection
\addcontentsline{toc}{section}{Consolidated bibliography}
\section*{Consolidated bibliography}
\nocite{*}
\bibliography{the}
\bibliographystyle{amsalpha}

\end{document}
